% ============================================
%  Abstract Algebra — Translation Template
%  Dummit & Foote
%  基于大笔记模板，仅保留基础设置
% ============================================
\documentclass{book}

% ── 基础配置 ──
% ============================================
%  packages.tex — 基础宏包
%  基于大笔记模板，精简至最小必要集合
% ============================================

% ── 中文支持 ──
\usepackage[scheme=plain,heading=true,fontset=none]{ctex}

% ── 数学 ──
\usepackage{amsmath, amsthm, amssymb, amsfonts, mathtools}
\usepackage{mathrsfs}

% ── 页边距 ──
\usepackage{geometry}

% ── 颜色（必须在 hyperref 之前，明确 xetex 驱动）──
\usepackage[xetex,svgnames]{xcolor}
\usepackage{tikz}
\usetikzlibrary{calc}

% ── 文档结构 ──
\usepackage{setspace}
\usepackage{enumitem}
\usepackage{graphicx}
\usepackage{bookmark}
\usepackage{hyperref}
\usepackage{cleveref}
\usepackage{etoolbox}

% ============================================
%  geometry.tex — 页面布局与字体
%  直接引自大笔记模板
% ============================================

% ── 页面尺寸 ──
\geometry{
  textheight=9in,
  textwidth=5.5in,
  top=1in,
  headheight=12pt,
  headsep=25pt,
  footskip=30pt
}

% ── 字体 ──
\usepackage{fontspec}
\setmainfont{Times New Roman}
\setCJKmainfont{SimSun}[AutoFakeBold={3}]

% 定理环境专用仿宋
\setCJKfamilyfont{theoremCJK}{FangSong}[AutoFakeBold={2.5}]
\newcommand{\theoremFont}{\CJKfamily{theoremCJK}}

% 证明环境专用楷体
\setCJKfamilyfont{proofCJK}{KaiTi}[AutoFakeBold={2.5}]
\newcommand{\proofFont}{\CJKfamily{proofCJK}}

% 封面专用黑体
\setCJKfamilyfont{heiti}{SimHei}[AutoFakeBold={3}]
\newcommand{\heiti}{\CJKfamily{heiti}}

% ── 行距 ──
\setstretch{1.2}

% ── 容忍小幅溢出 ──
\hfuzz=50pt

% ── 页面底部允许自然高度，避免 Underfull vbox ──
\raggedbottom
\vbadness=9999

% ============================================
%  theorems.tex — 定理环境（标准 amsthm，无色框）
%  基于大笔记模板，去掉 tcolorbox，改用纯 amsthm
% ============================================

% ── 定理样式 ──
\theoremstyle{definition}
\newtheorem{theorem}{定理}[chapter]
\newtheorem{proposition}[theorem]{命题}
\newtheorem{lemma}[theorem]{引理}
\newtheorem{corollary}[theorem]{推论}
\newtheorem{definition}[theorem]{定义}
\newtheorem{example}[theorem]{例}
\newtheorem{remark}[theorem]{注记}
\newtheorem{exercise}[theorem]{练习}
\newtheorem{axiom}[theorem]{公理}
\newtheorem{problem}[theorem]{问题}

% ── example 环境楷体 ──
\AddToHook{env/example/begin}{\proofFont}

% ── 证明环境 ──
\renewenvironment{proof}[1][证明]{%
  \par\proofFont
  \noindent\textbf{#1.}\quad
}{%
  \hfill $\square$\par
}

% ── 解答环境 ──
\newcounter{solution}
\newenvironment{solution}{%
  \par\proofFont
  \refstepcounter{solution}%
  \noindent\textbf{解.}\quad
}{%
  \hfill $\square$\par
}

% ── 方程编号按章 ──
\numberwithin{equation}{chapter}

% ============================================
%  commands.tex — 数学快捷命令 + 翻译辅助
% ============================================

% ── 基础符号 ──
\newcommand{\eps}{\varepsilon}
\newcommand{\vp}{\varphi}
\renewcommand{\leq}{\leqslant}
\renewcommand{\geq}{\geqslant}

% ── 分隔符 ──
\DeclarePairedDelimiter{\abs}{\lvert}{\rvert}
\DeclarePairedDelimiter{\norm}{\lVert}{\rVert}
\DeclarePairedDelimiter{\inner}{\langle}{\rangle}
\DeclarePairedDelimiter{\gen}{\langle}{\rangle}

% ── 数集（黑板体）──
\newcommand{\RR}{\mathbb{R}}
\newcommand{\NN}{\mathbb{N}}
\newcommand{\ZZ}{\mathbb{Z}}
\newcommand{\QQ}{\mathbb{Q}}
\newcommand{\CC}{\mathbb{C}}
\newcommand{\FF}{\mathbb{F}}
\newcommand{\HH}{\mathbb{H}}
\newcommand{\PP}{\mathbb{P}}
\newcommand{\EE}{\mathbb{E}}

% ── 花体 ──
\newcommand{\mcA}{\mathcal{A}}
\newcommand{\mcB}{\mathcal{B}}
\newcommand{\mcC}{\mathcal{C}}
\newcommand{\mcD}{\mathcal{D}}
\newcommand{\mcF}{\mathcal{F}}
\newcommand{\mcG}{\mathcal{G}}
\newcommand{\mcI}{\mathcal{I}}
\newcommand{\mcL}{\mathcal{L}}
\newcommand{\mcM}{\mathcal{M}}
\newcommand{\mcN}{\mathcal{N}}
\newcommand{\mcO}{\mathcal{O}}
\newcommand{\mcP}{\mathcal{P}}
\newcommand{\mcS}{\mathcal{S}}
\newcommand{\mcT}{\mathcal{T}}
\newcommand{\mcU}{\mathcal{U}}

% ── 哥特体 ──
\newcommand{\mfa}{\mathfrak{a}}
\newcommand{\mfb}{\mathfrak{b}}
\newcommand{\mfg}{\mathfrak{g}}
\newcommand{\mfh}{\mathfrak{h}}
\newcommand{\mfm}{\mathfrak{m}}
\newcommand{\mfn}{\mathfrak{n}}
\newcommand{\mfp}{\mathfrak{p}}
\newcommand{\mfq}{\mathfrak{q}}
\newcommand{\mfr}{\mathfrak{r}}
\newcommand{\mfs}{\mathfrak{s}}

% ── 群论 ──
\DeclareMathOperator{\Aut}{Aut}
\DeclareMathOperator{\Hol}{Hol}
\DeclareMathOperator{\Inn}{Inn}
\DeclareMathOperator{\End}{End}
\DeclareMathOperator{\Hom}{Hom}
\DeclareMathOperator{\Syl}{Syl}
\DeclareMathOperator{\Stab}{Stab}
\DeclareMathOperator{\sgn}{sgn}
\DeclareMathOperator{\GL}{GL}
\DeclareMathOperator{\SL}{SL}
\DeclareMathOperator{\PSL}{PSL}
\DeclareMathOperator{\SO}{SO}
\DeclareMathOperator{\charac}{char}

% ── 环论与模论 ──
\DeclareMathOperator{\Spec}{Spec}
\DeclareMathOperator{\nilrad}{nil}
\DeclareMathOperator{\rad}{rad}
\DeclareMathOperator{\Jac}{Jac}
\DeclareMathOperator{\Ann}{Ann}
\DeclareMathOperator{\Frac}{Frac}
\DeclareMathOperator{\coker}{coker}
\DeclareMathOperator{\Ext}{Ext}
\DeclareMathOperator{\Tor}{Tor}

% ── 域论 ──
\DeclareMathOperator{\Gal}{Gal}
\DeclareMathOperator{\Tr}{Tr}
\DeclareMathOperator{\Nm}{Nm}

% ── 常用 ──
\DeclareMathOperator{\id}{id}
\DeclareMathOperator{\im}{im}
\newcommand{\divides}{\mid}
\newcommand{\ndivides}{\nmid}
\newcommand{\normal}{\trianglelefteq}
\newcommand{\semidirect}{\rtimes}
\newcommand{\dsum}{\bigoplus}
\newcommand{\tensor}{\otimes}

% ── 偏导数 ──
\newcommand{\del}[2]{\frac{\partial #1}{\partial #2}}

% ============================================
%  翻译辅助命令
% ============================================

% 原文边注
\usepackage{marginnote}
\renewcommand*{\marginfont}{\footnotesize\color{gray}\sffamily}
\newcommand{\orig}[1]{\marginnote{\textsf{#1}}[0pt]}

% 首次术语中英对照
\newcommand{\tr}[2]{#2\marginnote{\scriptsize\textsf{#1}}[-2pt]}

% 译者注
\newcommand{\transnote}[1]{%
  \par\smallskip
  \noindent{\footnotesize\sffamily\color{gray}【译者注】#1}%
  \par\smallskip
}

% 待确认
\newcommand{\tocheck}[1]{\textcolor{red}{【待确认：#1】}}

% 人名
\newcommand{\person}[3]{#2（\textsf{#1}, #3）}

% preamble/part.title.tex
% ============================================
% Part 扉页设计 —— 与封面统一的橙红 / 金黄 / 深绿 / 蓝线配色
%
% 接口：
%   \setpartimage{<path>}
%   \setpartintro{<text>}
% ============================================

% ---- 色彩定义：参考封面 ----
\definecolor{pCoverRedDark}{RGB}{108, 31, 18}
\definecolor{pCoverRed}{RGB}{150, 46, 20}
\definecolor{pCoverOrange}{RGB}{209, 86, 24}
\definecolor{pCoverAmber}{RGB}{247, 183, 43}
\definecolor{pCoverAmberLight}{RGB}{255, 221, 98}
\definecolor{pCoverGreen}{RGB}{0, 78, 43}
\definecolor{pCoverBlue}{RGB}{0, 74, 178}
\definecolor{pCoverInk}{RGB}{64, 24, 14}
\definecolor{pCoverCream}{RGB}{255, 239, 196}

% 保留旧名称，避免其他文件引用 BlueMain/BlueDark/GoldMain 时失效
\colorlet{BlueMain}{pCoverBlue}
\colorlet{BlueDark}{pCoverGreen}
\colorlet{GoldMain}{pCoverAmber}

\makeatletter

% -------- 每 Part 图片接口 --------
\gdef\partimage{}
\newcommand{\setpartimage}[1]{\gdef\partimage{#1}}

% -------- Part 简介接口 --------
\gdef\partintro{}
\newcommand{\setpartintro}[1]{\gdef\partintro{#1}}

% -------- 版式参数 --------
\newlength{\PartTopBandHeight}
\setlength{\PartTopBandHeight}{2.15cm}

\newlength{\PartTitleTop}
\setlength{\PartTitleTop}{2.50cm}

\newlength{\PartTitlePanelHeight}
\setlength{\PartTitlePanelHeight}{5.35cm}

\newlength{\PartIntroTop}
\setlength{\PartIntroTop}{8.00cm}

\newlength{\PartIntroHeight}
\setlength{\PartIntroHeight}{19.55cm}
\newlength{\PartIntroBHeight}
\setlength{\PartIntroBHeight}{19.55cm}

\newlength{\PartImageTop}
\setlength{\PartImageTop}{29.35cm}

\newlength{\PartSideMargin}
\setlength{\PartSideMargin}{1.55cm}

\newlength{\PartTextMargin}
\setlength{\PartTextMargin}{2.05cm}

% -------- 蓝色横线：呼应封面 --------
\newcommand{\PartBlueRule}[2]{%
    \draw[pCoverBlue,line width=#2]
        ([xshift=\PartSideMargin,yshift=-#1]current page.north west) --
        ([xshift=-\PartSideMargin,yshift=-#1]current page.north east);%
}

% -------- 左上 Part 编号 + 标题 --------
\newcommand{\PartTitleBlock}[2]{%
    \ifdefempty{#1}{%
        \node[anchor=north west,text=pCoverGreen,font=\fontsize{26}{32}\bfseries,
              text width=\dimexpr\paperwidth-4.1cm\relax,align=left]
            at ([xshift=\PartTextMargin,yshift=-\PartTitleTop-0.65cm]current page.north west)
            {#2};%
    }{%
        \node[anchor=north west,text=pCoverGreen,font=\fontsize{54}{60}\bfseries]
            at ([xshift=\PartTextMargin,yshift=-\PartTitleTop-0.15cm]current page.north west)
            {#1};%
        \node[anchor=north west,text=pCoverInk!72,font=\Large\bfseries]
            at ([xshift=\PartTextMargin,yshift=-\PartTitleTop-3.25cm]current page.north west)
            {PART};%
        \node[anchor=north west,text=pCoverGreen,font=\fontsize{26}{32}\bfseries,
              text width=\dimexpr\paperwidth-4.1cm\relax,align=left]
            at ([xshift=\PartTextMargin+2.95cm,yshift=-\PartTitleTop-0.55cm]current page.north west)
            {#2};%
    }%
}

% -------- Part 简介 block：改为横向大文本区，适合长简介 --------
\newcommand{\PartIntroBlock}{%
    \ifdefempty{\partintro}{}{%
        \fill[pCoverCream!88]
            ([xshift=\PartTextMargin,yshift=-\PartIntroTop]current page.north west)
            rectangle
            ([xshift=-\PartTextMargin,yshift=-\dimexpr\PartIntroTop+\PartIntroHeight\relax]current page.north east);%
        \draw[pCoverBlue,line width=1.05pt]
            ([xshift=\PartTextMargin,yshift=-\PartIntroTop]current page.north west) --
            ([xshift=-\PartTextMargin,yshift=-\PartIntroTop]current page.north east);%
        \draw[pCoverBlue,line width=1.05pt]
            ([xshift=\PartTextMargin,yshift=-\dimexpr\PartIntroTop+\PartIntroHeight\relax]current page.north west) --
            ([xshift=-\PartTextMargin,yshift=-\dimexpr\PartIntroTop+\PartIntroHeight\relax]current page.north east);%
        \node[anchor=north west,text=pCoverGreen,font=\large\bfseries]
            at ([xshift=\PartTextMargin+0.32cm,yshift=-\PartIntroTop-0.42cm]current page.north west)
            {简介};%
        \node[anchor=north west,inner sep=0pt,text=pCoverInk]
            at ([xshift=\PartTextMargin+0.32cm,yshift=-\PartIntroTop-1.18cm]current page.north west)
            {\parbox[t][\dimexpr\PartIntroHeight-1.55cm\relax][t]{\dimexpr\paperwidth-\PartTextMargin-\PartTextMargin-0.64cm\relax}{%
                \setstretch{1.12}\fontsize{8.2}{10.8}\selectfont\partintro
            }};%
    }%
}

% -------- Part 扉页 --------
\newcommand{\PartTitlePage}[2]{%
    \thispagestyle{empty}%
    \vspace*{-\baselineskip}%
    \begin{tikzpicture}[remember picture,overlay]
        % 全页底色：暖色，与封面统一但降低饱和度，保证正文可读
        \shade[top color=pCoverAmberLight!72,bottom color=pCoverRed!28]
            (current page.north west) rectangle (current page.south east);

        % 顶部橙红横条
        \shade[left color=pCoverRedDark,right color=pCoverOrange]
            (current page.north west) rectangle
            ([yshift=-\PartTopBandHeight]current page.north east);

        % 左侧深红竖带
        \shade[top color=pCoverRedDark,bottom color=pCoverRed]
            (current page.north west) rectangle
            ([xshift=0.72cm]current page.south west);

        % 标题金黄面板
        \shade[left color=pCoverAmberLight,right color=pCoverAmber]
            ([xshift=\PartSideMargin,yshift=-\PartTitleTop]current page.north west)
            rectangle
            ([xshift=-\PartSideMargin,yshift=-\dimexpr\PartTitleTop+\PartTitlePanelHeight\relax]current page.north east);
        \draw[pCoverBlue,line width=1.2pt]
            ([xshift=\PartSideMargin,yshift=-\PartTitleTop]current page.north west) --
            ([xshift=-\PartSideMargin,yshift=-\PartTitleTop]current page.north east);
        \draw[pCoverBlue,line width=1.2pt]
            ([xshift=\PartSideMargin,yshift=-\dimexpr\PartTitleTop+\PartTitlePanelHeight\relax]current page.north west) --
            ([xshift=-\PartSideMargin,yshift=-\dimexpr\PartTitleTop+\PartTitlePanelHeight\relax]current page.north east);

      
        % 左上 Part 编号 + 标题
        \PartTitleBlock{#1}{#2}

        % 简介：横向大文本区，给长简介更多行宽与高度
        \PartIntroBlock

        
    \end{tikzpicture}%
    \clearpage%
}

% ============================================================
%  重定义 \part
% ============================================================
\def\part{%
  \@ifstar{\part@TP@star}{\part@TP}%
}%
\def\part@TP@star#1{%
  \addcontentsline{toc}{part}{#1}%
  \PartTitlePage{}{#1}%
  \@endpart@TP
}%
\def\part@TP{%
  \@ifnextchar[{\part@TP@opt}{\part@TP@noopt}%
}%
\def\part@TP@opt[#1]#2{%
  \refstepcounter{part}%
  \addcontentsline{toc}{part}{\thepart\hspace{1em}#1}%
  \PartTitlePage{\thepart}{#2}%
  \@endpart@TP
}%
\def\part@TP@noopt#1{%
  \refstepcounter{part}%
  \addcontentsline{toc}{part}{\thepart\hspace{1em}#1}%
  \PartTitlePage{\thepart}{#1}%
  \@endpart@TP
}%
\def\@endpart@TP{%
  \if@twoside
    \if@openright
      \null
      \thispagestyle{empty}%
      \newpage
    \fi
  \fi
}%

\makeatother


% ── 章节/Part 介绍接口 ──
\newcommand{\setchapterintro}[1]{\setpartintro{#1}}
\newcommand{\setchapterimage}[1]{\setpartimage{#1}}

% ── 按需编译 ──
%\includeonly{chapters/chapter01}

% ── 封面信息（在 frontmatter/cover.tex 中定义）──
% ============================================
% cover.tex — Abstract Algebra 翻译稿封面
% 参考原书第三版封面，用 TikZ 绘制。
%
% 用法：在 main.tex 的 \begin{document} 后调用
%   % ============================================
% cover.tex — Abstract Algebra 翻译稿封面
% 参考原书第三版封面，用 TikZ 绘制。
%
% 用法：在 main.tex 的 \begin{document} 后调用
%   % ============================================
% cover.tex — Abstract Algebra 翻译稿封面
% 参考原书第三版封面，用 TikZ 绘制。
%
% 用法：在 main.tex 的 \begin{document} 后调用
%   \input{frontmatter/cover}
%   \makecover
%
% 依赖：preamble 中需已加载 tikz 与 xcolor，并建议：
%   \usepackage{tikz}
%   \usetikzlibrary{calc}
% ============================================

% ── 封面信息：可在导言区或本文件中修改 ──
\providecommand{\CoverTitleCN}{抽象代数}
\providecommand{\CoverTitleEN}{Abstract Algebra}
\providecommand{\CoverTitleENLineA}{A B S T R A C T}
\providecommand{\CoverTitleENLineB}{A L G E B R A}
\providecommand{\CoverEdition}{THIRD EDITION}
\providecommand{\CoverAuthorA}{D A V I D\quad S.\quad D U M M I T}
\providecommand{\CoverAuthorB}{R I C H A R D\quad M.\quad F O O T E}
\providecommand{\CoverSubtitle}{Dummit \& Foote 中译本}
\providecommand{\CoverTranslator}{译者：万一之}
\providecommand{\CoverDate}{\today}
\providecommand{\CoverNote}{仅供学习与交流使用}

% 若未使用 ctex，给 \heiti 一个安全后备定义；在 ctex 中不会覆盖原定义。
\providecommand{\heiti}{\bfseries}

% ── 封面配色，均使用独立名称，避免与模板其他颜色冲突 ──
\definecolor{CoverOuterTop}{HTML}{B94A22}
\definecolor{CoverOuterBottom}{HTML}{7F2417}
\definecolor{CoverPanelTop}{HTML}{FFC835}
\definecolor{CoverPanelMid}{HTML}{E88920}
\definecolor{CoverPanelBottom}{HTML}{A73518}
\definecolor{CoverRuleBlue}{HTML}{0046A8}
\definecolor{CoverRuleEdge}{HTML}{7B2B24}
\definecolor{CoverTitleGreen}{HTML}{006B34}
\definecolor{CoverDarkText}{HTML}{5B211D}
\definecolor{CoverAuthorWhite}{HTML}{FFF5E9}

\providecommand{\makecover}{}
\renewcommand{\makecover}{%
\begin{titlepage}
  \thispagestyle{empty}%
  \begingroup
  \begin{tikzpicture}[remember picture, overlay]
    % 整页深红棕背景
    \shade[top color=CoverOuterTop, bottom color=CoverOuterBottom]
      (current page.south west) rectangle (current page.north east);

    % 中央橙色主面板
    \coordinate (PanelNW) at ([xshift=1.45cm,yshift=-1.42cm]current page.north west);
    \coordinate (PanelNE) at ([xshift=-1.45cm,yshift=-1.42cm]current page.north east);
    \coordinate (PanelSW) at ([xshift=1.45cm,yshift=1.55cm]current page.south west);
    \coordinate (PanelSE) at ([xshift=-1.45cm,yshift=1.55cm]current page.south east);

    \shade[top color=CoverPanelTop, bottom color=CoverPanelBottom]
      (PanelSW) rectangle (PanelNE);
    \shade[top color=white, bottom color=CoverPanelMid, opacity=.16]
      (PanelNW) rectangle ([yshift=-7.15cm]PanelNE);

    % 轻微暗边，模拟原封面的印刷边界
    \draw[CoverOuterBottom!55!black, line width=1.1pt, opacity=.35]
      (PanelSW) rectangle (PanelNE);
    \draw[CoverPanelTop!70!white, line width=.45pt, opacity=.42]
      ([xshift=.08cm,yshift=-.08cm]PanelNW) rectangle ([xshift=-.08cm,yshift=.08cm]PanelSE);

    % 左右与下方的暗红过渡，增加书封质感
    \shade[left color=CoverOuterBottom, right color=CoverOuterBottom!0, opacity=.20]
      (PanelSW) rectangle ([xshift=1.10cm]PanelNW);
    \shade[right color=CoverOuterBottom, left color=CoverOuterBottom!0, opacity=.13]
      ([xshift=-1.10cm]PanelSE) rectangle (PanelNE);
    \shade[bottom color=CoverOuterBottom, top color=CoverOuterBottom!0, opacity=.18]
      (PanelSW) rectangle ([yshift=2.60cm]PanelSE);

    % 蓝色横线：仿原书封面的分隔线
    \foreach \Y in {5.18,8.50,11.82,20.82,24.08}{
      \draw[CoverRuleEdge, line width=4.6pt, opacity=.70]
        ([xshift=1.62cm,yshift=-\Y cm]current page.north west) --
        ([xshift=-1.62cm,yshift=-\Y cm]current page.north east);
      \draw[CoverRuleBlue, line width=2.7pt]
        ([xshift=1.62cm,yshift=-\Y cm]current page.north west) --
        ([xshift=-1.62cm,yshift=-\Y cm]current page.north east);
      \draw[CoverPanelTop!75!white, line width=.35pt, opacity=.55]
        ([xshift=1.62cm,yshift=-\Y cm]current page.north west) --
        ([xshift=-1.62cm,yshift=-\Y cm]current page.north east);
    }

    % 主标题阴影与文字
    \node[anchor=west, inner sep=0pt, text=black, opacity=.12,
      font=\sffamily\bfseries\fontsize{26}{31}\selectfont]
      at ([xshift=2.22cm,yshift=-4.50cm]current page.north west)
      {\CoverTitleENLineA};
    \node[anchor=west, inner sep=0pt, text=CoverTitleGreen,
      font=\sffamily\bfseries\fontsize{26}{31}\selectfont]
      at ([xshift=2.15cm,yshift=-4.42cm]current page.north west)
      {\CoverTitleENLineA};

    \node[anchor=west, inner sep=0pt, text=black, opacity=.12,
      font=\sffamily\bfseries\fontsize{26}{31}\selectfont]
      at ([xshift=2.22cm,yshift=-7.80cm]current page.north west)
      {\CoverTitleENLineB};
    \node[anchor=west, inner sep=0pt, text=CoverTitleGreen,
      font=\sffamily\bfseries\fontsize{26}{31}\selectfont]
      at ([xshift=2.15cm,yshift=-7.72cm]current page.north west)
      {\CoverTitleENLineB};

    % 版本信息
    \node[anchor=east, inner sep=0pt, text=CoverDarkText,
      font=\sffamily\bfseries\fontsize{10.8}{13}\selectfont]
      at ([xshift=-2.45cm,yshift=-12.36cm]current page.north east)
      {\CoverEdition};

    % 作者名
    \node[anchor=east, inner sep=0pt, text=black, opacity=.18,
      font=\sffamily\bfseries\fontsize{16.8}{21}\selectfont]
      at ([xshift=-2.02cm,yshift=-21.36cm]current page.north east)
      {\CoverAuthorA};
    \node[anchor=east, inner sep=0pt, text=CoverAuthorWhite,
      font=\sffamily\bfseries\fontsize{16.8}{21}\selectfont]
      at ([xshift=-2.08cm,yshift=-21.28cm]current page.north east)
      {\CoverAuthorA};

    \node[anchor=east, inner sep=0pt, text=black, opacity=.18,
      font=\sffamily\bfseries\fontsize{16.8}{21}\selectfont]
      at ([xshift=-2.02cm,yshift=-24.62cm]current page.north east)
      {\CoverAuthorB};
    \node[anchor=east, inner sep=0pt, text=CoverAuthorWhite,
      font=\sffamily\bfseries\fontsize{16.8}{21}\selectfont]
      at ([xshift=-2.08cm,yshift=-24.54cm]current page.north east)
      {\CoverAuthorB};

    % 中文译本信息：保持低调，不破坏原封面构图
    \node[anchor=west, inner sep=0pt, text=CoverAuthorWhite!88!CoverPanelTop,
      font=\heiti\fontsize{13.5}{18}\selectfont]
      at ([xshift=2.05cm,yshift=2.18cm]current page.south west)
      {\CoverTitleCN\quad \CoverSubtitle};

    \node[anchor=west, inner sep=0pt, text=CoverAuthorWhite!78!CoverPanelTop,
      font=\heiti\fontsize{9.8}{13}\selectfont]
      at ([xshift=2.05cm,yshift=1.55cm]current page.south west)
      {\CoverTranslator};

    \node[anchor=east, inner sep=0pt, text=CoverAuthorWhite!62!CoverPanelTop,
      font=\sffamily\fontsize{8.8}{12}\selectfont]
      at ([xshift=-2.05cm,yshift=1.55cm]current page.south east)
      {\CoverDate};

    % 底部注释
    \node[anchor=south, inner sep=0pt, text=CoverAuthorWhite!45!CoverPanelTop,
      font=\sffamily\fontsize{7.5}{10}\selectfont]
      at ([yshift=0.42cm]current page.south)
      {\CoverNote};
  \end{tikzpicture}%
  \null
  \endgroup
\end{titlepage}%
\clearpage
}

%   \makecover
%
% 依赖：preamble 中需已加载 tikz 与 xcolor，并建议：
%   \usepackage{tikz}
%   \usetikzlibrary{calc}
% ============================================

% ── 封面信息：可在导言区或本文件中修改 ──
\providecommand{\CoverTitleCN}{抽象代数}
\providecommand{\CoverTitleEN}{Abstract Algebra}
\providecommand{\CoverTitleENLineA}{A B S T R A C T}
\providecommand{\CoverTitleENLineB}{A L G E B R A}
\providecommand{\CoverEdition}{THIRD EDITION}
\providecommand{\CoverAuthorA}{D A V I D\quad S.\quad D U M M I T}
\providecommand{\CoverAuthorB}{R I C H A R D\quad M.\quad F O O T E}
\providecommand{\CoverSubtitle}{Dummit \& Foote 中译本}
\providecommand{\CoverTranslator}{译者：万一之}
\providecommand{\CoverDate}{\today}
\providecommand{\CoverNote}{仅供学习与交流使用}

% 若未使用 ctex，给 \heiti 一个安全后备定义；在 ctex 中不会覆盖原定义。
\providecommand{\heiti}{\bfseries}

% ── 封面配色，均使用独立名称，避免与模板其他颜色冲突 ──
\definecolor{CoverOuterTop}{HTML}{B94A22}
\definecolor{CoverOuterBottom}{HTML}{7F2417}
\definecolor{CoverPanelTop}{HTML}{FFC835}
\definecolor{CoverPanelMid}{HTML}{E88920}
\definecolor{CoverPanelBottom}{HTML}{A73518}
\definecolor{CoverRuleBlue}{HTML}{0046A8}
\definecolor{CoverRuleEdge}{HTML}{7B2B24}
\definecolor{CoverTitleGreen}{HTML}{006B34}
\definecolor{CoverDarkText}{HTML}{5B211D}
\definecolor{CoverAuthorWhite}{HTML}{FFF5E9}

\providecommand{\makecover}{}
\renewcommand{\makecover}{%
\begin{titlepage}
  \thispagestyle{empty}%
  \begingroup
  \begin{tikzpicture}[remember picture, overlay]
    % 整页深红棕背景
    \shade[top color=CoverOuterTop, bottom color=CoverOuterBottom]
      (current page.south west) rectangle (current page.north east);

    % 中央橙色主面板
    \coordinate (PanelNW) at ([xshift=1.45cm,yshift=-1.42cm]current page.north west);
    \coordinate (PanelNE) at ([xshift=-1.45cm,yshift=-1.42cm]current page.north east);
    \coordinate (PanelSW) at ([xshift=1.45cm,yshift=1.55cm]current page.south west);
    \coordinate (PanelSE) at ([xshift=-1.45cm,yshift=1.55cm]current page.south east);

    \shade[top color=CoverPanelTop, bottom color=CoverPanelBottom]
      (PanelSW) rectangle (PanelNE);
    \shade[top color=white, bottom color=CoverPanelMid, opacity=.16]
      (PanelNW) rectangle ([yshift=-7.15cm]PanelNE);

    % 轻微暗边，模拟原封面的印刷边界
    \draw[CoverOuterBottom!55!black, line width=1.1pt, opacity=.35]
      (PanelSW) rectangle (PanelNE);
    \draw[CoverPanelTop!70!white, line width=.45pt, opacity=.42]
      ([xshift=.08cm,yshift=-.08cm]PanelNW) rectangle ([xshift=-.08cm,yshift=.08cm]PanelSE);

    % 左右与下方的暗红过渡，增加书封质感
    \shade[left color=CoverOuterBottom, right color=CoverOuterBottom!0, opacity=.20]
      (PanelSW) rectangle ([xshift=1.10cm]PanelNW);
    \shade[right color=CoverOuterBottom, left color=CoverOuterBottom!0, opacity=.13]
      ([xshift=-1.10cm]PanelSE) rectangle (PanelNE);
    \shade[bottom color=CoverOuterBottom, top color=CoverOuterBottom!0, opacity=.18]
      (PanelSW) rectangle ([yshift=2.60cm]PanelSE);

    % 蓝色横线：仿原书封面的分隔线
    \foreach \Y in {5.18,8.50,11.82,20.82,24.08}{
      \draw[CoverRuleEdge, line width=4.6pt, opacity=.70]
        ([xshift=1.62cm,yshift=-\Y cm]current page.north west) --
        ([xshift=-1.62cm,yshift=-\Y cm]current page.north east);
      \draw[CoverRuleBlue, line width=2.7pt]
        ([xshift=1.62cm,yshift=-\Y cm]current page.north west) --
        ([xshift=-1.62cm,yshift=-\Y cm]current page.north east);
      \draw[CoverPanelTop!75!white, line width=.35pt, opacity=.55]
        ([xshift=1.62cm,yshift=-\Y cm]current page.north west) --
        ([xshift=-1.62cm,yshift=-\Y cm]current page.north east);
    }

    % 主标题阴影与文字
    \node[anchor=west, inner sep=0pt, text=black, opacity=.12,
      font=\sffamily\bfseries\fontsize{26}{31}\selectfont]
      at ([xshift=2.22cm,yshift=-4.50cm]current page.north west)
      {\CoverTitleENLineA};
    \node[anchor=west, inner sep=0pt, text=CoverTitleGreen,
      font=\sffamily\bfseries\fontsize{26}{31}\selectfont]
      at ([xshift=2.15cm,yshift=-4.42cm]current page.north west)
      {\CoverTitleENLineA};

    \node[anchor=west, inner sep=0pt, text=black, opacity=.12,
      font=\sffamily\bfseries\fontsize{26}{31}\selectfont]
      at ([xshift=2.22cm,yshift=-7.80cm]current page.north west)
      {\CoverTitleENLineB};
    \node[anchor=west, inner sep=0pt, text=CoverTitleGreen,
      font=\sffamily\bfseries\fontsize{26}{31}\selectfont]
      at ([xshift=2.15cm,yshift=-7.72cm]current page.north west)
      {\CoverTitleENLineB};

    % 版本信息
    \node[anchor=east, inner sep=0pt, text=CoverDarkText,
      font=\sffamily\bfseries\fontsize{10.8}{13}\selectfont]
      at ([xshift=-2.45cm,yshift=-12.36cm]current page.north east)
      {\CoverEdition};

    % 作者名
    \node[anchor=east, inner sep=0pt, text=black, opacity=.18,
      font=\sffamily\bfseries\fontsize{16.8}{21}\selectfont]
      at ([xshift=-2.02cm,yshift=-21.36cm]current page.north east)
      {\CoverAuthorA};
    \node[anchor=east, inner sep=0pt, text=CoverAuthorWhite,
      font=\sffamily\bfseries\fontsize{16.8}{21}\selectfont]
      at ([xshift=-2.08cm,yshift=-21.28cm]current page.north east)
      {\CoverAuthorA};

    \node[anchor=east, inner sep=0pt, text=black, opacity=.18,
      font=\sffamily\bfseries\fontsize{16.8}{21}\selectfont]
      at ([xshift=-2.02cm,yshift=-24.62cm]current page.north east)
      {\CoverAuthorB};
    \node[anchor=east, inner sep=0pt, text=CoverAuthorWhite,
      font=\sffamily\bfseries\fontsize{16.8}{21}\selectfont]
      at ([xshift=-2.08cm,yshift=-24.54cm]current page.north east)
      {\CoverAuthorB};

    % 中文译本信息：保持低调，不破坏原封面构图
    \node[anchor=west, inner sep=0pt, text=CoverAuthorWhite!88!CoverPanelTop,
      font=\heiti\fontsize{13.5}{18}\selectfont]
      at ([xshift=2.05cm,yshift=2.18cm]current page.south west)
      {\CoverTitleCN\quad \CoverSubtitle};

    \node[anchor=west, inner sep=0pt, text=CoverAuthorWhite!78!CoverPanelTop,
      font=\heiti\fontsize{9.8}{13}\selectfont]
      at ([xshift=2.05cm,yshift=1.55cm]current page.south west)
      {\CoverTranslator};

    \node[anchor=east, inner sep=0pt, text=CoverAuthorWhite!62!CoverPanelTop,
      font=\sffamily\fontsize{8.8}{12}\selectfont]
      at ([xshift=-2.05cm,yshift=1.55cm]current page.south east)
      {\CoverDate};

    % 底部注释
    \node[anchor=south, inner sep=0pt, text=CoverAuthorWhite!45!CoverPanelTop,
      font=\sffamily\fontsize{7.5}{10}\selectfont]
      at ([yshift=0.42cm]current page.south)
      {\CoverNote};
  \end{tikzpicture}%
  \null
  \endgroup
\end{titlepage}%
\clearpage
}

%   \makecover
%
% 依赖：preamble 中需已加载 tikz 与 xcolor，并建议：
%   \usepackage{tikz}
%   \usetikzlibrary{calc}
% ============================================

% ── 封面信息：可在导言区或本文件中修改 ──
\providecommand{\CoverTitleCN}{抽象代数}
\providecommand{\CoverTitleEN}{Abstract Algebra}
\providecommand{\CoverTitleENLineA}{A B S T R A C T}
\providecommand{\CoverTitleENLineB}{A L G E B R A}
\providecommand{\CoverEdition}{THIRD EDITION}
\providecommand{\CoverAuthorA}{D A V I D\quad S.\quad D U M M I T}
\providecommand{\CoverAuthorB}{R I C H A R D\quad M.\quad F O O T E}
\providecommand{\CoverSubtitle}{Dummit \& Foote 中译本}
\providecommand{\CoverTranslator}{译者：万一之}
\providecommand{\CoverDate}{\today}
\providecommand{\CoverNote}{仅供学习与交流使用}

% 若未使用 ctex，给 \heiti 一个安全后备定义；在 ctex 中不会覆盖原定义。
\providecommand{\heiti}{\bfseries}

% ── 封面配色，均使用独立名称，避免与模板其他颜色冲突 ──
\definecolor{CoverOuterTop}{HTML}{B94A22}
\definecolor{CoverOuterBottom}{HTML}{7F2417}
\definecolor{CoverPanelTop}{HTML}{FFC835}
\definecolor{CoverPanelMid}{HTML}{E88920}
\definecolor{CoverPanelBottom}{HTML}{A73518}
\definecolor{CoverRuleBlue}{HTML}{0046A8}
\definecolor{CoverRuleEdge}{HTML}{7B2B24}
\definecolor{CoverTitleGreen}{HTML}{006B34}
\definecolor{CoverDarkText}{HTML}{5B211D}
\definecolor{CoverAuthorWhite}{HTML}{FFF5E9}

\providecommand{\makecover}{}
\renewcommand{\makecover}{%
\begin{titlepage}
  \thispagestyle{empty}%
  \begingroup
  \begin{tikzpicture}[remember picture, overlay]
    % 整页深红棕背景
    \shade[top color=CoverOuterTop, bottom color=CoverOuterBottom]
      (current page.south west) rectangle (current page.north east);

    % 中央橙色主面板
    \coordinate (PanelNW) at ([xshift=1.45cm,yshift=-1.42cm]current page.north west);
    \coordinate (PanelNE) at ([xshift=-1.45cm,yshift=-1.42cm]current page.north east);
    \coordinate (PanelSW) at ([xshift=1.45cm,yshift=1.55cm]current page.south west);
    \coordinate (PanelSE) at ([xshift=-1.45cm,yshift=1.55cm]current page.south east);

    \shade[top color=CoverPanelTop, bottom color=CoverPanelBottom]
      (PanelSW) rectangle (PanelNE);
    \shade[top color=white, bottom color=CoverPanelMid, opacity=.16]
      (PanelNW) rectangle ([yshift=-7.15cm]PanelNE);

    % 轻微暗边，模拟原封面的印刷边界
    \draw[CoverOuterBottom!55!black, line width=1.1pt, opacity=.35]
      (PanelSW) rectangle (PanelNE);
    \draw[CoverPanelTop!70!white, line width=.45pt, opacity=.42]
      ([xshift=.08cm,yshift=-.08cm]PanelNW) rectangle ([xshift=-.08cm,yshift=.08cm]PanelSE);

    % 左右与下方的暗红过渡，增加书封质感
    \shade[left color=CoverOuterBottom, right color=CoverOuterBottom!0, opacity=.20]
      (PanelSW) rectangle ([xshift=1.10cm]PanelNW);
    \shade[right color=CoverOuterBottom, left color=CoverOuterBottom!0, opacity=.13]
      ([xshift=-1.10cm]PanelSE) rectangle (PanelNE);
    \shade[bottom color=CoverOuterBottom, top color=CoverOuterBottom!0, opacity=.18]
      (PanelSW) rectangle ([yshift=2.60cm]PanelSE);

    % 蓝色横线：仿原书封面的分隔线
    \foreach \Y in {5.18,8.50,11.82,20.82,24.08}{
      \draw[CoverRuleEdge, line width=4.6pt, opacity=.70]
        ([xshift=1.62cm,yshift=-\Y cm]current page.north west) --
        ([xshift=-1.62cm,yshift=-\Y cm]current page.north east);
      \draw[CoverRuleBlue, line width=2.7pt]
        ([xshift=1.62cm,yshift=-\Y cm]current page.north west) --
        ([xshift=-1.62cm,yshift=-\Y cm]current page.north east);
      \draw[CoverPanelTop!75!white, line width=.35pt, opacity=.55]
        ([xshift=1.62cm,yshift=-\Y cm]current page.north west) --
        ([xshift=-1.62cm,yshift=-\Y cm]current page.north east);
    }

    % 主标题阴影与文字
    \node[anchor=west, inner sep=0pt, text=black, opacity=.12,
      font=\sffamily\bfseries\fontsize{26}{31}\selectfont]
      at ([xshift=2.22cm,yshift=-4.50cm]current page.north west)
      {\CoverTitleENLineA};
    \node[anchor=west, inner sep=0pt, text=CoverTitleGreen,
      font=\sffamily\bfseries\fontsize{26}{31}\selectfont]
      at ([xshift=2.15cm,yshift=-4.42cm]current page.north west)
      {\CoverTitleENLineA};

    \node[anchor=west, inner sep=0pt, text=black, opacity=.12,
      font=\sffamily\bfseries\fontsize{26}{31}\selectfont]
      at ([xshift=2.22cm,yshift=-7.80cm]current page.north west)
      {\CoverTitleENLineB};
    \node[anchor=west, inner sep=0pt, text=CoverTitleGreen,
      font=\sffamily\bfseries\fontsize{26}{31}\selectfont]
      at ([xshift=2.15cm,yshift=-7.72cm]current page.north west)
      {\CoverTitleENLineB};

    % 版本信息
    \node[anchor=east, inner sep=0pt, text=CoverDarkText,
      font=\sffamily\bfseries\fontsize{10.8}{13}\selectfont]
      at ([xshift=-2.45cm,yshift=-12.36cm]current page.north east)
      {\CoverEdition};

    % 作者名
    \node[anchor=east, inner sep=0pt, text=black, opacity=.18,
      font=\sffamily\bfseries\fontsize{16.8}{21}\selectfont]
      at ([xshift=-2.02cm,yshift=-21.36cm]current page.north east)
      {\CoverAuthorA};
    \node[anchor=east, inner sep=0pt, text=CoverAuthorWhite,
      font=\sffamily\bfseries\fontsize{16.8}{21}\selectfont]
      at ([xshift=-2.08cm,yshift=-21.28cm]current page.north east)
      {\CoverAuthorA};

    \node[anchor=east, inner sep=0pt, text=black, opacity=.18,
      font=\sffamily\bfseries\fontsize{16.8}{21}\selectfont]
      at ([xshift=-2.02cm,yshift=-24.62cm]current page.north east)
      {\CoverAuthorB};
    \node[anchor=east, inner sep=0pt, text=CoverAuthorWhite,
      font=\sffamily\bfseries\fontsize{16.8}{21}\selectfont]
      at ([xshift=-2.08cm,yshift=-24.54cm]current page.north east)
      {\CoverAuthorB};

    % 中文译本信息：保持低调，不破坏原封面构图
    \node[anchor=west, inner sep=0pt, text=CoverAuthorWhite!88!CoverPanelTop,
      font=\heiti\fontsize{13.5}{18}\selectfont]
      at ([xshift=2.05cm,yshift=2.18cm]current page.south west)
      {\CoverTitleCN\quad \CoverSubtitle};

    \node[anchor=west, inner sep=0pt, text=CoverAuthorWhite!78!CoverPanelTop,
      font=\heiti\fontsize{9.8}{13}\selectfont]
      at ([xshift=2.05cm,yshift=1.55cm]current page.south west)
      {\CoverTranslator};

    \node[anchor=east, inner sep=0pt, text=CoverAuthorWhite!62!CoverPanelTop,
      font=\sffamily\fontsize{8.8}{12}\selectfont]
      at ([xshift=-2.05cm,yshift=1.55cm]current page.south east)
      {\CoverDate};

    % 底部注释
    \node[anchor=south, inner sep=0pt, text=CoverAuthorWhite!45!CoverPanelTop,
      font=\sffamily\fontsize{7.5}{10}\selectfont]
      at ([yshift=0.42cm]current page.south)
      {\CoverNote};
  \end{tikzpicture}%
  \null
  \endgroup
\end{titlepage}%
\clearpage
}


\begin{document}

% ── 封面 ──
\pagenumbering{gobble}
\makecover
\cleardoublepage

% ── 前言与目录 ──
\frontmatter
% ============================================
%  preface.tex — 译者序
% ============================================
\chapter*{第三版前言}
\addcontentsline{toc}{chapter}{第三版前言}

第二版以来的主要变化，是本版加入了 Gröbner 基。其基本理论在新的第 9.6 节中介绍。求解多项式方程组的应用（消元理论）出现在该节末尾，从而使这一节成为关于此主题的一个自足基础。随后，更多的应用和例子被融入第 15 章对仿射代数集和 \(k\)-代数同态的处理中。虽然后一章中的理论仍然不依赖于 Gröbner 基，但新增的应用、例子和计算技巧显著加强了该部分的展开；我们建议把第 9.6 节作为进入第 15 章的过渡来阅读，或者与第 15 章并行阅读。第 9.6 节和第 15 章中加入了大量涉及 Gröbner 基的习题，这些习题既有计算性质的，也有理论性质的。关于 Gröbner 基的初步习题可以（而且为了帮助理解算法，应当）手算完成；但较大的计算，特别是第 15 章习题中 Gröbner 基的大部分使用，很可能需要借助计算机完成。

其他变化包括对 168 阶单群分类（第 6.2 节）的处理作了简化，并加入了一个通过 2 阶射影平面给出的唯一性证明。文本中的其他若干证明或部分内容也作了轻微修订。全书各处加入了许多新的习题，主要放在各节末尾，以便尽可能保留早期版本的编号方案。特别地，新增习题涉及非交换环上的自由模（10.3）、Krull 维数（15.3）以及平坦模（10.5 和 17.1）。

与以前各版一样，本书所包含的内容远多于通常一年课程所能讲授的内容。一门基本的入门（一年）课程大概应包括第一部分直到第 5.3 节、第二部分直到第 9.5 节、第 10.1、10.2、10.3、11.1、11.2 节以及第四部分。第 12 章也应讲授，可以放在第四部分之前或之后。第 5、6、9、10 和 11 章中的其他主题可以穿插在这样的课程中，也可以在课程末尾视时间许可讲授。

第 10.4 和 10.5 节的难度比第 10 章开头几节略高；对于按顺序阅读本书的读者，初读时可以推迟阅读。后一节关于正合列性质的内容虽然相当长，但通过在各个小节中对三个基本函子的平行处理而保持了连贯性。

在核心材料之外，第三版为希望学习现代代数若干重要领域的学生和教师提供了相当大的灵活性，无论是作为独立学习还是作为课程内容。例如，对于已有一定代数背景的学生，可以组织结构良好的一学期课程，内容包括：第 9.6 节以及第 15、16 章；或第 10、17 章；或第 5、6 章以及第六部分。这些选择中的任何一种，也都会为后续课程打下坚实基础，使学生能够更深入地学习许多可能领域中的某一个，例如代数数论、代数拓扑、代数几何、表示论、Lie 群等等。

新增材料的选择，以及将这些材料展开并融入正文的方式，与本书的一个基本主题相一致：数学不同领域之间丰富的相互作用所产生的力量与美。全书始终强调尽可能用许多例子来激发重要代数概念的引入和发展。我们并未试图做到百科全书式的包罗万象，而是试图以一种显示这些思想十分自然地发展出来的方式，触及初等代数中的许多核心主题。

许多重要的思想和结果出现在习题中。这并不是因为它们不重要，而是因为它们不容易放入正文的行文脉络中，但又太重要而不能完全略去。关于同一主题的一系列习题前面会有一些说明，并且经过安排，使读者即使不实际做这些习题也可以阅读。在某些情况下，新材料首先在习题中引入——往往比正文中出现早几节——这样学生可以通过做这些习题而较容易地初步接触这些材料（例如，Lagrange 定理出现在第 1.7 节的习题中，而在正文中出现在第 3.2 节）。所有习题都在正文范围之内；在我们认为需要的地方，提示用[方括号]给出。我们认为可能不那么直接的习题，通常会以给出习题答案的方式来表述；同时，许多习题被分解成一串较为常规的习题，以便使其更容易处理。

我们也有意尽量减少正文中的函子语言，以便使叙述尽可能保持初等性。对于进一步阅读的资料，在已有许多优秀选择可供参考的地方，我们没有给出具体文献。并且，虽然我们已经尽力包含尽可能多的基本主题，但若由于篇幅或个人取舍的原因而遗漏了读者特别喜爱的某些内容，我们在此致歉。

我们深切感谢，并愿在这里向世界各地的许多学生和同事表示感谢；十五年多以来，他们提供了宝贵的意见、见解和鼓励——他们持续的关注激励我们写成了这个第三版。

\begin{flushright}
	David Dummit\\
	Richard Foote\\
	June, 2003
\end{flushright}
\chapter*{译者序}
\addcontentsline{toc}{chapter}{译者序}

本文是 David S.\ Dummit 与 Richard M.\ Foote 所著
\textit{Abstract Algebra}（第三版）的中文翻译稿。

该书是抽象代数领域的经典教材，涵盖群论、环论、模论、
域论与 Galois 理论、交换代数、同调代数以及有限群表示论
等核心内容。全书共 19 章，分为六大部分，体系完整，习题丰富。

翻译过程中力求忠实于原著，同时兼顾中文数学表达的流畅性。


\begin{flushright}
	万一之 \\
	\today
\end{flushright}
\tableofcontents

% ── 正文 ──
\mainmatter
\setcounter{page}{1}

% \setcounter{chapter}{-1}
\chapter{预备知识}

为便于使用，本章汇集了贯穿全书将要用到的一些结果和记号。学生起初可以快速浏览本章，随后在正文中遇到相应概念时，再更仔细地阅读各节。

\section{基础}

集合论的基础知识：集合、$\cap$、$\cup$、$\in$ 等，应为读者所熟悉。对于给定集合 $A$ 的子集，我们将使用记号
\[
  B=\{a\in A\mid \cdots\text{（关于 }a\text{ 的条件）}\cdots\}.
\]
集合 $A$ 的\emph{阶}或\emph{基数}记为 $|A|$。若 $A$ 是有限集，则 $A$ 的阶就是 $A$ 中元素的个数。

重要的是要理解如何检验某个特定的 $x\in A$ 是否属于 $A$ 的一个子集 $B$（参见练习 1--4）。两个集合 $A$ 与 $B$ 的\emph{笛卡儿积}是由来自 $A$ 和 $B$ 的有序对组成的集合
\[
  A\times B=\{(a,b)\mid a\in A,\ b\in B\}.
\]

对于一些常见的数集，我们将使用如下记号：
\begin{enumerate}[label=(\arabic*)]
  \item $\ZZ=\{0,\pm1,\pm2,\pm3,\ldots\}$ 表示整数（字母 $Z$ 来自德语中表示“数”的词 ``Zahlen''）。
  \item $\QQ=\{a/b\mid a,b\in\ZZ,\ b\neq0\}$ 表示有理数。
  \item $\RR=\{\text{所有十进制展开 }\pm d_1d_2\cdots d_n.a_1a_2a_3\cdots\}$ 表示实数。
  \item $\CC=\{a+bi\mid a,b\in\RR,\ i^2=-1\}$ 表示复数。
  \item $\ZZ^+$、$\QQ^+$ 和 $\RR^+$ 分别表示 $\ZZ$、$\QQ$ 和 $\RR$ 中的正的（非零）元素。
\end{enumerate}

我们将使用记号 $f:A\to B$ 或 $A\xrightarrow{f}B$ 表示从 $A$ 到 $B$ 的函数 $f$，$f$ 在 $a$ 处的值记为 $f(a)$（即所有函数都写在自变量的左侧）。我们交替使用“函数”和“映射”这两个词。集合 $A$ 称为 $f$ 的\emph{定义域}，$B$ 称为 $f$ 的\emph{陪域}。记号 $f:a\mapsto b$，或者在 $f$ 已经明确时写作 $a\mapsto b$，表示 $f(a)=b$，也就是说，函数是在元素上被指定的。

如果函数 $f$ 不是在元素上给定的，那么通常必须检验 $f$ 是\emph{良定}的，即它是无歧义地确定的。例如，若集合 $A$ 是两个子集 $A_1$ 与 $A_2$ 的并，则可以尝试指定一个从 $A$ 到集合 $\{0,1\}$ 的函数：规定 $A_1$ 中的所有元素都映到 $0$，$A_2$ 中的所有元素都映到 $1$。除非 $A_1$ 与 $A_2$ 有公共元素，否则这就无歧义地定义了 $f$（若它们有公共元素，则不清楚这些元素应映到 $0$ 还是 $1$）。因此，检验这个 $f$ 是否良定，就等同于检验 $A_1$ 与 $A_2$ 没有交集。

集合
\[
  f(A)=\{b\in B\mid b=f(a),\ \text{对某个 }a\in A\}
\]
是 $B$ 的一个子集，称为 $f$ 的\emph{值域}或\emph{像}（也称为 $A$ 在 $f$ 下的像）。对 $B$ 的每个子集 $C$，集合
\[
  f^{-1}(C)=\{a\in A\mid f(a)\in C\}
\]
由 $A$ 中在 $f$ 下映入 $C$ 的元素组成，称为 $C$ 在 $f$ 下的\emph{原像}或\emph{逆像}。对每个 $b\in B$，$\{b\}$ 在 $f$ 下的原像称为 $f$ 在 $b$ 上方的\emph{纤维}。注意，一般而言 $f^{-1}$ 并不是函数，并且 $f$ 的纤维通常含有许多元素，因为可能有许多 $A$ 中的元素映到同一个元素 $b$。

若 $f:A\to B$ 且 $g:B\to C$，则复合映射 $g\circ f:A\to C$ 定义为
\[
  (g\circ f)(a)=g(f(a)).
\]

设 $f:A\to B$。
\begin{enumerate}[label=(\arabic*)]
  \item 若只要 $a_1\neq a_2$ 就有 $f(a_1)\neq f(a_2)$，则称 $f$ 是\emph{单射}的，或称 $f$ 是一个\emph{单射}。
  \item 若对所有 $b\in B$，存在某个 $a\in A$ 使得 $f(a)=b$，即 $f$ 的像是整个 $B$，则称 $f$ 是\emph{满射}的，或称 $f$ 是一个\emph{满射}。注意，因为函数总是映到它的值域上（按定义），所以必须指定陪域 $B$，满射性的问题才有意义。
  \item 若 $f$ 既是单射又是满射，则称 $f$ 是\emph{双射}的，或称 $f$ 是一个\emph{双射}。若存在这样一个从 $A$ 到 $B$ 的双射 $f$，则称 $A$ 与 $B$ 是\emph{双射对应}的。
  \item 若存在函数 $g:B\to A$，使得 $g\circ f:A\to A$ 是 $A$ 上的恒等映射，即对所有 $a\in A$ 都有 $(g\circ f)(a)=a$，则称 $f$ 有一个\emph{左逆}。
  \item 若存在函数 $h:B\to A$，使得 $f\circ h:B\to B$ 是 $B$ 上的恒等映射，则称 $f$ 有一个\emph{右逆}。
\end{enumerate}

\medskip
\noindent\textbf{命题 1.} 设 $f:A\to B$。
\begin{enumerate}[label=(\arabic*)]
  \item 映射 $f$ 是单射，当且仅当 $f$ 有左逆。
  \item 映射 $f$ 是满射，当且仅当 $f$ 有右逆。
  \item 映射 $f$ 是双射，当且仅当存在 $g:B\to A$，使得 $f\circ g$ 是 $B$ 上的恒等映射，并且 $g\circ f$ 是 $A$ 上的恒等映射。
  \item 若 $A$ 与 $B$ 是元素个数相同的有限集（即 $|A|=|B|$），则 $f:A\to B$ 是双射，当且仅当 $f$ 是单射，当且仅当 $f$ 是满射。
\end{enumerate}
\noindent\textit{证明：} 练习。

在上面命题第 (3) 部分的情形中，映射 $g$ 必然唯一；我们称 $g$ 为 $f$ 的\emph{双侧逆}（或简称为\emph{逆}）。

集合 $A$ 的一个\emph{置换}就是从 $A$ 到自身的一个双射。

若 $A\subseteq B$ 且 $f:B\to C$，则 $f$ 到 $A$ 的限制记为 $f|_A$。当所考虑的定义域已明确时，我们偶尔仍把 $f|_A$ 简写为 $f$，尽管形式上它们是不同的函数（因为它们的定义域不同）。

若 $A\subseteq B$，$g:A\to C$，并且存在函数 $f:B\to C$ 使得 $f|_A=g$，则称 $f$ 是 $g$ 到 $B$ 的一个\emph{延拓}（这样的映射 $f$ 不一定存在，也不一定唯一）。

设 $A$ 是非空集合。
\begin{enumerate}[label=(\arabic*)]
  \item 集合 $A$ 上的一个\emph{二元关系}是 $A\times A$ 的一个子集 $R$，若 $(a,b)\in R$，则写作 $a\sim b$。
  \item 称 $A$ 上的关系 $\sim$ 为：
  \begin{enumerate}[label=(\alph*)]
    \item \emph{自反的}，若对所有 $a\in A$，有 $a\sim a$；
    \item \emph{对称的}，若对所有 $a,b\in A$，$a\sim b$ 蕴含 $b\sim a$；
    \item \emph{传递的}，若对所有 $a,b,c\in A$，$a\sim b$ 且 $b\sim c$ 蕴含 $a\sim c$。
  \end{enumerate}
  若一个关系是自反的、对称的并且传递的，则称它是一个\emph{等价关系}。
  \item 若 $\sim$ 在 $A$ 上定义了一个等价关系，则 $a\in A$ 的\emph{等价类}定义为 $\{x\in A\mid x\sim a\}$。$a$ 的等价类中的元素称为与 $a$ 等价。若 $C$ 是一个等价类，则 $C$ 中的任一元素都称为类 $C$ 的一个\emph{代表元}。
  \item $A$ 的一个\emph{划分}是由 $A$ 的非空子集组成的任意集合 $\{A_i\mid i\in I\}$（其中 $I$ 是某个指标集），满足
  \begin{enumerate}[label=(\alph*)]
    \item $A=\bigcup_{i\in I}A_i$，并且
    \item 对所有满足 $i\neq j$ 的 $i,j\in I$，有 $A_i\cap A_j=\varnothing$，
  \end{enumerate}
  也就是说，$A$ 是该划分中各集合的不交并。
\end{enumerate}

$A$ 上的等价关系与 $A$ 的划分是同一回事：

\medskip
\noindent\textbf{命题 2.} 设 $A$ 是非空集合。
\begin{enumerate}[label=(\arabic*)]
  \item 若 $\sim$ 在 $A$ 上定义了一个等价关系，则 $\sim$ 的等价类全体构成 $A$ 的一个划分。
  \item 若 $\{A_i\mid i\in I\}$ 是 $A$ 的一个划分，则 $A$ 上存在一个等价关系，其等价类恰好是这些集合 $A_i$，$i\in I$。
\end{enumerate}
\noindent\textit{证明：} 略。

最后，我们假定读者熟悉归纳法证明。

\subsection*{练习}

在练习 1 到 4 中，令 $\mathcal{A}$ 为实数域上所有 $2\times2$ 矩阵组成的集合。回忆矩阵乘法定义为
\[
  \begin{pmatrix}a&b\\ c&d\end{pmatrix}
  \begin{pmatrix}p&q\\ r&s\end{pmatrix}
  =
  \begin{pmatrix}ap+br&aq+bs\\ cp+dr&cq+ds\end{pmatrix}.
\]
令
\[
  M=\begin{pmatrix}1&1\\0&1\end{pmatrix}
\]
并令
\[
  \mathcal{B}=\{X\in\mathcal{A}\mid MX=XM\}.
\]
\begin{enumerate}[label=\arabic*.]
  \item 判断下列 $\mathcal{A}$ 中的元素哪些属于 $\mathcal{B}$：
  \[
  \begin{pmatrix}1&1\\0&1\end{pmatrix},\quad
  \begin{pmatrix}1&1\\1&1\end{pmatrix},\quad
  \begin{pmatrix}0&0\\0&0\end{pmatrix},\quad
  \begin{pmatrix}1&1\\1&0\end{pmatrix},\quad
  \begin{pmatrix}1&0\\0&1\end{pmatrix},\quad
  \begin{pmatrix}0&1\\1&0\end{pmatrix}.
  \]
  \item 证明：若 $P,Q\in\mathcal{B}$，则 $P+Q\in\mathcal{B}$（其中 $+$ 表示通常的矩阵加法）。
  \item 证明：若 $P,Q\in\mathcal{B}$，则 $P\cdot Q\in\mathcal{B}$（其中 $\cdot$ 表示通常的矩阵乘法）。
  \item 求出关于 $p,q,r,s$ 的条件，它们精确地刻画何时 $\begin{pmatrix}p&q\\ r&s\end{pmatrix}\in\mathcal{B}$。
  \item 判断下列函数 $f$ 是否良定：
  \begin{enumerate}[label=(\alph*)]
    \item $f:\QQ\to\ZZ$，定义为 $f(a/b)=a$。
    \item $f:\QQ\to\QQ$，定义为 $f(a/b)=a^2/b^2$。
  \end{enumerate}
  \item 判断函数 $f:\RR^+\to\ZZ$ 是否良定，其中 $f$ 把实数 $r$ 映到 $r$ 的某个十进制展开中小数点右边的第一位数字。
  \item 设 $f:A\to B$ 是集合之间的满射。证明关系
  \[
    a\sim b\quad\text{当且仅当}\quad f(a)=f(b)
  \]
  是一个等价关系，其等价类就是 $f$ 的纤维。
\end{enumerate}

\section{整数的性质}

整数 $\ZZ$ 的以下性质（其中许多是初等算术中熟悉的性质）将在第 8 章的环论中置于更一般的背景下证明；不过在第一部分中有必要使用它们（当然，环论中对这些性质的证明并不依赖群论）。
\begin{enumerate}[label=(\arabic*)]
  \item （$\ZZ$ 的良序性）若 $A$ 是 $\ZZ^+$ 的任一非空子集，则存在某个元素 $m\in A$，使得对所有 $a\in A$ 都有 $m\leq a$（$m$ 称为 $A$ 的一个\emph{极小元}）。
  \item 若 $a,b\in\ZZ$ 且 $a\neq0$，如果存在元素 $c\in\ZZ$ 使得 $b=ac$，则称 $a$\emph{整除} $b$。此时记为 $a\mid b$；若 $a$ 不整除 $b$，则记为 $a\nmid b$。
  \item 若 $a,b\in\ZZ-\{0\}$，则存在唯一的正整数 $d$，称为 $a$ 与 $b$ 的\emph{最大公因子}（或 g.c.d.），满足：
  \begin{enumerate}[label=(\alph*)]
    \item $d\mid a$ 且 $d\mid b$（因此 $d$ 是 $a$ 与 $b$ 的一个公因子），并且
    \item 若 $e\mid a$ 且 $e\mid b$，则 $e\mid d$（因此 $d$ 是这样的最大因子）。
  \end{enumerate}
  $a$ 与 $b$ 的最大公因子记为 $(a,b)$。若 $(a,b)=1$，则称 $a$ 与 $b$\emph{互素}。
  \item 若 $a,b\in\ZZ-\{0\}$，则存在唯一的正整数 $l$，称为 $a$ 与 $b$ 的\emph{最小公倍数}（或 l.c.m.），满足：
  \begin{enumerate}[label=(\alph*)]
    \item $a\mid l$ 且 $b\mid l$（因此 $l$ 是 $a$ 与 $b$ 的一个公倍数），并且
    \item 若 $a\mid m$ 且 $b\mid m$，则 $l\mid m$（因此 $l$ 是这样的最小倍数）。
  \end{enumerate}
  两个整数 $a$ 与 $b$ 的最大公因子 $d$ 与最小公倍数 $l$ 之间的联系由 $dl=ab$ 给出。
  \item \emph{带余除法}：若 $a,b\in\ZZ-\{0\}$，则存在唯一的 $q,r\in\ZZ$，使得
  \[
    a=qb+r\quad\text{且}\quad 0\leq r<|b|,
  \]
  其中 $q$ 是商，$r$ 是余数。这就是初等算术中熟悉的通常“长除法”。
  \item \emph{欧几里得算法}是一个重要过程，它通过反复使用带余除法来求出两个整数 $a$ 与 $b$ 的最大公因子：若 $a,b\in\ZZ-\{0\}$，则可得到一列商与余数
  \begin{alignat*}{2}
    a     &= q_0b+r_0,       &\qquad &(0)\\
    b     &= q_1r_0+r_1,     &&(1)\\
    r_0   &= q_2r_1+r_2,     &&(2)\\
    r_1   &= q_3r_2+r_3,     &&(3)\\
          &\ \vdots          &&\\
    r_{n-2}&=q_nr_{n-1}+r_n, &&(n)\\
    r_{n-1}&=q_{n+1}r_n,     &&(n+1)
  \end{alignat*}
  其中 $r_n$ 是最后一个非零余数。这样的 $r_n$ 必然存在，因为若余数非零，则 $|b|>|r_0|>|r_1|>\cdots>|r_n|$ 是严格正整数的递减序列，而这样的序列不可能无限继续。于是 $r_n$ 就是 $a$ 与 $b$ 的最大公因子 $(a,b)$。
\end{enumerate}

\paragraph{例}
设 $a=57970$，$b=10353$。应用欧几里得算法得到：
\begin{align*}
  57970&=(5)10353+6205,\\
  10353&=(1)6205+4148,\\
  6205&=(1)4148+2057,\\
  4148&=(2)2057+34,\\
  2057&=(60)34+17,\\
  34&=(2)17,
\end{align*}
这说明 $(57970,10353)=17$。

\begin{enumerate}[label=(\arabic*),resume]
  \item 欧几里得算法的一个后果将被我们经常使用：若 $a,b\in\ZZ-\{0\}$，则存在 $x,y\in\ZZ$，使得
  \[
    (a,b)=ax+by,
  \]
  也就是说，$a$ 与 $b$ 的最大公因子是 $a$ 与 $b$ 的一个 $\ZZ$-线性组合。这可由在欧几里得算法中递归地把元素 $r_n$ 写成前面余数的组合而得出（具体地说，利用上面的方程 $(n)$ 解出 $r_n=r_{n-2}-q_nr_{n-1}$，从而用余数 $r_{n-1}$ 与 $r_{n-2}$ 表示 $r_n$；再利用方程 $(n-1)$ 用 $r_{n-2}$ 与 $r_{n-3}$ 表示 $r_n$，依此类推，最终用 $a$ 与 $b$ 表示 $r_n$）。
\end{enumerate}

\paragraph{例}
设 $a=57970$，$b=10353$，它们的最大公因子上面已经算出为 $17$。对这两个整数应用欧几里得算法时，从第五个方程（倒数第二个方程）解出它们的最大公因子：
\[
  17=2057-(60)34.
\]
第四个方程又给出 $34=4148-(2)2057$，所以把前一个余数 $34$ 的这个表达式代入，得到
\[
  17=2057-(60)[4148-(2)2057],
\]
即
\[
  17=(121)2057-(60)4148.
\]
由第三个方程解出 $2057$ 并代入，得
\[
  17=(121)[6205-(1)4148]-(60)4148=(121)6205-(181)4148.
\]
再用第二个方程解出 $4148$，然后用第一个方程解出 $6205$，最终得到
\[
  17=(302)57970-(1691)10353,
\]
这一点很容易直接检验。因此，在这个例子中，对于最大公因子的方程 $ax+by=(a,b)$，有解 $x=302$ 与 $y=-1691$。注意，单靠猜测找到这个关系的可能性相当小。

上面 (7) 中的整数 $x$ 与 $y$ 并不唯一。例如，在 $a=57970$、$b=10353$ 的例子中，我们求得一个解为 $x=302$、$y=-1691$；同时也比较容易检验 $x=-307$、$y=1719$ 也满足 $57970x+10353y=17$。$x$ 与 $y$ 的一般解是已知的（参见下面的练习以及第 8 章）。

\begin{enumerate}[label=(\arabic*),resume]
  \item $\ZZ^+$ 的一个元素 $p$ 称为\emph{素数}，若 $p>1$ 且 $p$ 的正因子只有 $1$ 和 $p$（起初，“素数”一词只指正整数）。大于 $1$ 但不是素数的整数称为\emph{合数}。例如，$2,3,5,7,11,13,17,19,\ldots$ 是素数，而 $4,6,8,9,10,12,14,15,16,18,\ldots$ 是合数。

  素数的一个重要性质（事实上可以用来定义素数，参见练习 3）如下：若 $p$ 是素数，并且对某些 $a,b\in\ZZ$ 有 $p\mid ab$，则 $p\mid a$ 或 $p\mid b$。

  \item \emph{算术基本定理}说：若 $n\in\ZZ$ 且 $n>1$，则 $n$ 可以唯一地分解为素数的乘积；也就是说，存在互不相同的素数 $p_1,p_2,\ldots,p_s$ 和正整数 $\alpha_1,\alpha_2,\ldots,\alpha_s$，使得
  \[
    n=p_1^{\alpha_1}p_2^{\alpha_2}\cdots p_s^{\alpha_s}.
  \]
  这个分解的唯一性是指：若 $q_1,q_2,\ldots,q_t$ 是任意互不相同的素数，$\beta_1,\beta_2,\ldots,\beta_t$ 是正整数，并且
  \[
    n=q_1^{\beta_1}q_2^{\beta_2}\cdots q_t^{\beta_t},
  \]
  则 $s=t$，并且如果把两组素数都按递增顺序排列，则 $q_i=p_i$ 且 $\alpha_i=\beta_i$，$1\leq i\leq s$。例如
  \[
    n=1852423848=2^3 3^2 11^2 19^3 31,
  \]
  而这个素数乘积分解是唯一的。

  假设正整数 $a$ 与 $b$ 表示为素数幂的乘积：
  \[
    a=p_1^{\alpha_1}p_2^{\alpha_2}\cdots p_s^{\alpha_s},
    \qquad
    b=p_1^{\beta_1}p_2^{\beta_2}\cdots p_s^{\beta_s},
  \]
  其中 $p_1,p_2,\ldots,p_s$ 互不相同，并且指数都 $\geq0$（这里允许指数为 $0$，以便两个乘积取在同一组素数上；若某个素数实际上不是因子，则其指数为 $0$）。于是 $a$ 与 $b$ 的最大公因子为
  \[
    (a,b)=p_1^{\min(\alpha_1,\beta_1)}p_2^{\min(\alpha_2,\beta_2)}\cdots p_s^{\min(\alpha_s,\beta_s)}
  \]
  （最小公倍数则改为取 $\alpha_i$ 与 $\beta_i$ 的最大值，而不是最小值）。
\end{enumerate}

\paragraph{例}
在上面的例子中，$a=57970$ 与 $b=10353$ 可分解为
\[
  a=2\cdot5\cdot11\cdot17\cdot31,
  \qquad
  b=3\cdot7\cdot17\cdot29,
\]
由此可以立即得出它们的最大公因子为 $17$。不过要注意，对大整数而言，确定其素因数分解是极其困难的（事实上，目前使用的若干常见编码正是基于这种困难），因此这通常不是求最大公因子的有效方法。欧几里得算法无需 $a$ 与 $b$ 的素因数分解，就能相当快速地求出最大公因子。

\begin{enumerate}[label=(\arabic*),resume]
  \item \emph{欧拉 $\varphi$ 函数}定义如下：对 $n\in\ZZ^+$，令 $\varphi(n)$ 为满足 $a\leq n$ 且 $a$ 与 $n$ 互素，即 $(a,n)=1$ 的正整数 $a$ 的个数。例如，$\varphi(12)=4$，因为 $1,5,7,11$ 是不超过 $12$ 且与 $12$ 没有公共因子的唯一正整数。同样，$\varphi(1)=1$，$\varphi(2)=1$，$\varphi(3)=2$，$\varphi(4)=2$，$\varphi(5)=4$，$\varphi(6)=2$，等等。对素数 $p$，有 $\varphi(p)=p-1$；更一般地，对所有 $a\geq1$，有公式
  \[
    \varphi(p^a)=p^a-p^{a-1}=p^{a-1}(p-1).
  \]
  函数 $\varphi$ 在如下意义下是\emph{乘性的}：
  \[
    \varphi(ab)=\varphi(a)\varphi(b)\qquad\text{若 }(a,b)=1
  \]
  （注意这里 $a$ 与 $b$ 互素是重要的）。结合上面的公式，可得到 $\varphi$ 的一般取值公式：若 $n=p_1^{\alpha_1}p_2^{\alpha_2}\cdots p_s^{\alpha_s}$，则
  \begin{align*}
    \varphi(n)
      &=\varphi(p_1^{\alpha_1})\varphi(p_2^{\alpha_2})\cdots\varphi(p_s^{\alpha_s})\\
      &=p_1^{\alpha_1-1}(p_1-1)p_2^{\alpha_2-1}(p_2-1)\cdots p_s^{\alpha_s-1}(p_s-1).
  \end{align*}
  例如，$\varphi(12)=\varphi(2^2)\varphi(3)=2^1(2-1)3^0(3-1)=4$。读者应注意，在全书中我们会用字母 $\varphi$ 表示许多不同的函数，因此当我们要用这个字母表示欧拉函数时，会明确说明。
\end{enumerate}

\subsection*{练习}

\begin{enumerate}[label=\arabic*.]
  \item 对下列每一对整数 $a$ 与 $b$，求出它们的最大公因子、最小公倍数，并把它们的最大公因子写成 $ax+by$ 的形式，其中 $x,y$ 是整数。
  \begin{enumerate}[label=(\alph*)]
    \item $a=20,\ b=13$。
    \item $a=69,\ b=372$。
    \item $a=792,\ b=275$。
    \item $a=11391,\ b=5673$。
    \item $a=1761,\ b=1567$。
    \item $a=507885,\ b=60808$。
  \end{enumerate}
  \item 证明：若整数 $k$ 整除整数 $a$ 与 $b$，则对任意一对整数 $s$ 与 $t$，$k$ 整除 $as+bt$。
  \item 证明：若 $n$ 是合数，则存在整数 $a$ 与 $b$，使得 $n$ 整除 $ab$，但 $n$ 不整除 $a$ 也不整除 $b$。
  \item 设 $a,b,N$ 是固定整数，其中 $a$ 与 $b$ 非零；令 $d=(a,b)$ 为 $a$ 与 $b$ 的最大公因子。假设 $x_0$ 与 $y_0$ 是 $ax+by=N$ 的一组特解（即 $ax_0+by_0=N$）。证明：对任意整数 $t$，整数
  \[
    x=x_0+\frac{b}{d}t,\qquad y=y_0-\frac{a}{d}t
  \]
  也是 $ax+by=N$ 的解（这事实上给出了一般解）。
  \item 求出每个整数 $n\leq30$ 的 $\varphi(n)$ 值，其中 $\varphi$ 表示欧拉 $\varphi$ 函数。
  \item 用归纳法证明 $\ZZ$ 的良序性，并证明极小元唯一。
  \item 若 $p$ 是素数，证明不存在非零整数 $a$ 与 $b$ 使得 $a^2=pb^2$（即 $\sqrt{p}$ 不是有理数）。
  \item 设 $p$ 是素数，$n\in\ZZ^+$。求出整除
  \[
    n!=n(n-1)(n-2)\cdots 2\cdot1
  \]
  的最大 $p$ 的幂的公式（它涉及取整函数）。
  \item 编写一个计算机程序，求两个整数 $a$ 与 $b$ 的最大公因子 $(a,b)$，并把 $(a,b)$ 表示为 $ax+by$，其中 $x,y$ 是整数。
  \item 证明：对任意给定的正整数 $N$，只有有限多个整数 $n$ 满足 $\varphi(n)=N$，其中 $\varphi$ 表示欧拉 $\varphi$ 函数。特别地，由此推出当 $n$ 趋于无穷时，$\varphi(n)$ 也趋于无穷。
  \item 证明：若 $d$ 整除 $n$，则 $\varphi(d)$ 整除 $\varphi(n)$，其中 $\varphi$ 表示欧拉 $\varphi$ 函数。
\end{enumerate}

\section{$\ZZ/n\ZZ$：模 $n$ 的整数}

令 $n$ 是一个固定的正整数。在 $\ZZ$ 上定义关系：
\[
  a\sim b\quad\text{当且仅当}\quad n\mid(b-a).
\]
显然，对任意整数 $a$，有 $a\sim a$；并且对任意整数 $a,b$，$a\sim b$ 蕴含 $b\sim a$，所以这个关系显然是自反且对称的。若 $a\sim b$ 且 $b\sim c$，则 $n$ 整除 $a-b$，并且 $n$ 整除 $b-c$，于是 $n$ 也整除这两个整数之和，即 $n$ 整除
\[
  (a-b)+(b-c)=a-c,
\]
所以 $a\sim c$，从而这个关系是传递的。因此这是一个等价关系。若 $a\sim b$，则写作 $a\equiv b\pmod n$（读作：$a$ 模 $n$ 同余于 $b$）。对任意 $k\in\ZZ$，我们把 $a$ 的等价类记为 $\overline{a}$；这称为 $a$ 模 $n$ 的\emph{同余类}或\emph{剩余类}，它由与 $a$ 相差 $n$ 的整数倍的所有整数组成，即
\[
  \overline{a}=\{a+kn\mid k\in\ZZ\}
  =\{a,\ a\pm n,\ a\pm2n,\ a\pm3n,\ldots\}.
\]

模 $n$ 恰有 $n$ 个不同的等价类，即
\[
  \overline{0},\overline{1},\overline{2},\ldots,\overline{n-1},
\]
它们由除以 $n$ 后可能出现的余数决定，并且这些剩余类划分整数集 $\ZZ$。在这个等价关系下的等价类集合记为 $\ZZ/n\ZZ$，称为\emph{模 $n$ 的整数}（或\emph{模 $n$ 整数}）。当我们讨论商群和商环时，这个记号的动机会变得更加清楚。注意，对于不同的 $n$，等价关系和等价类是不同的；因此在使用横线记号以前，我们总是会先固定 $n$。求某个整数 $a$ 模 $n$ 的等价类这一过程常称为把 $a$\emph{化简模 $n$}。这个术语也常指求出与 $a$ 模 $n$ 同余的最小非负整数（$a$ 模 $n$ 的\emph{最小剩余}）。

我们可以为 $\ZZ/n\ZZ$ 的元素定义加法和乘法，从而定义\emph{模算术}如下：对 $\overline{a},\overline{b}\in\ZZ/n\ZZ$，定义它们的和与积为
\[
  \overline{a}+\overline{b}=\overline{a+b},
  \qquad
  \overline{a}\cdot\overline{b}=\overline{ab}.
\]
这意味着如下做法：给定 $\ZZ/n\ZZ$ 中的任意两个元素 $\overline{a}$ 与 $\overline{b}$，要计算它们的和（分别地，积），取类 $\overline{a}$ 中任意一个代表整数 $a$，以及类 $\overline{b}$ 中任意一个代表整数 $b$，在 $\ZZ$ 中照常把整数 $a$ 与 $b$ 相加（分别地，相乘），然后取包含结果的等价类。下面的定理 3 断言这是良定的，即它不依赖于为 $\ZZ/n\ZZ$ 中的元素 $\overline{a}$ 与 $\overline{b}$ 所选择的代表元。

\paragraph{例}
设 $n=12$，考虑 $\ZZ/12\ZZ$，它由十二个剩余类
\[
  \overline{0},\overline{1},\overline{2},\ldots,\overline{11}
\]
组成，这些类由整数除以 $12$ 后的十二种可能余数决定。例如，剩余类 $\overline{5}$ 中的元素，就是除以 $12$ 时余数为 $5$ 的那些整数（即模 $12$ 同余于 $5$ 的整数）。任意一个模 $12$ 同余于 $5$ 的整数（如 $5,17,29,\ldots$ 或 $-7,-19,\ldots$）都可作为剩余类 $\overline{5}$ 的一个代表元。注意，$\ZZ/12\ZZ$ 由上述十二个元素组成（而 $\ZZ/12\ZZ$ 的这些元素中，每一个本身又包含无限多个通常意义下的整数）。

现在设 $\overline{a}=\overline{5}$ 且 $\overline{b}=\overline{8}$。$\overline{a}$ 最显然的代表元是整数 $5$，同样 $\overline{b}$ 最显然的代表元是 $8$。使用这些代表元，得到
\[
  \overline{5}+\overline{8}=\overline{13}=\overline{1},
\]
因为 $13$ 与 $1$ 属于模 $n=12$ 的同一个类。若改用代表元 $17$ 代表 $\overline{a}$（注意 $5$ 与 $17$ 的确属于模 $12$ 的同一个剩余类），并用代表元 $-28$ 代表 $\overline{b}$，则得到
\[
  \overline{5}+\overline{8}=\overline{17-28}=\overline{-11}=\overline{1},
\]
如前所述，结果不依赖于所选择的代表元。这两个类的积为
\[
  \overline{a}\cdot\overline{b}=\overline{5}\cdot\overline{8}=\overline{40}=\overline{4},
\]
同样与所选择的代表元无关。

\medskip
\noindent\textbf{定理 3.} 上面定义的 $\ZZ/n\ZZ$ 上的加法与乘法运算都是良定的，也就是说，它们不依赖于所涉及各类代表元的选择。更准确地说，若 $a_1,a_2\in\ZZ$ 且 $b_1,b_2\in\ZZ$ 满足 $\overline{a_1}=\overline{b_1}$ 与 $\overline{a_2}=\overline{b_2}$，则
\[
  \overline{a_1+a_2}=\overline{b_1+b_2}
  \quad\text{且}\quad
  \overline{a_1a_2}=\overline{b_1b_2};
\]
也就是说，若
\[
  a_1\equiv b_1\pmod n\quad\text{且}\quad a_2\equiv b_2\pmod n,
\]
则
\[
  a_1+a_2\equiv b_1+b_2\pmod n
  \quad\text{且}\quad
  a_1a_2\equiv b_1b_2\pmod n.
\]
\noindent\textit{证明：} 假设 $a_1\equiv b_1\pmod n$，即 $a_1-b_1$ 可被 $n$ 整除。于是 $a_1=b_1+sn$，其中 $s$ 是某个整数。同样，$a_2\equiv b_2\pmod n$ 意味着 $a_2=b_2+tn$，其中 $t$ 是某个整数。因此
\[
  a_1+a_2=(b_1+b_2)+(s+t)n,
\]
所以 $a_1+a_2\equiv b_1+b_2\pmod n$，这说明剩余类的和与所选代表元无关。同样，
\[
  a_1a_2=(b_1+sn)(b_2+tn)=b_1b_2+(b_1t+b_2s+stn)n,
\]
这说明 $a_1a_2\equiv b_1b_2\pmod n$，因此剩余类的积也与所选代表元无关，证明完成。

我们稍后将看到，通过代表元相加来相加等价类的过程，是一个更一般构造（商构造）的特殊情形。等价类相加这个概念，在有理数加法的语境中已经很熟悉：每个有理数 $a/b$ 实际上都是一类表达式：$a/b=2a/2b=(-3a)/(-3b)$ 等等；我们经常改变代表元（例如通分）以便把两个分数相加（例如，计算 $1/2+1/3$ 时，改用与 $1/2$ 等价的代表元 $3/6$ 以及与 $1/3$ 等价的代表元 $2/6$，从而得到 $1/2+1/3=3/6+2/6=5/6$）。模算术的概念也很熟悉：在增加或减少若干小时后求一天中的时刻时，我们模 $12$ 化简并求最小剩余。

重要的是，要能够把某个等价关系的等价类看成可以操作的\emph{元素}（例如我们对分数所做的那样），而不是仅仅看成集合。与这种观点一致，我们常常把 $\ZZ/n\ZZ$ 的元素简单地记为 $\{0,1,\ldots,n-1\}$，其中加法和乘法都模 $n$ 化简。不过必须记住，$\ZZ/n\ZZ$ 的元素不是整数，而是通常整数的集合，并且其中的算术很不相同。例如，在整数 $\ZZ$ 中，$5+8$ 并不是 $1$，而在上面 $\ZZ/12\ZZ$ 的例子中却是如此。

能够在 $\ZZ/n\ZZ$ 中定义算术，在初等数论中有许多重要应用。作为一个简单例子，我们计算 $2^{1000}$ 的最后两位数字。首先注意，最后两位数字给出 $2^{1000}$ 除以 $100$ 的余数，因此我们关心的是包含 $2^{1000}$ 的模 $100$ 剩余类。计算得
\[
  2^{10}=1024\equiv24\pmod{100},
\]
于是
\[
  2^{20}=(2^{10})^2\equiv24^2=576\equiv76\pmod{100}.
\]
然后
\[
  2^{40}=(2^{20})^2\equiv76^2=5776\equiv76\pmod{100}.
\]
同样，$2^{80}\equiv2^{160}\equiv2^{320}\equiv2^{640}\equiv76\pmod{100}$。最后，
\[
  2^{1000}=2^{640}2^{320}2^{40}\equiv76\cdot76\cdot76\equiv76\pmod{100},
\]
所以最后两位数字是 $76$。

$\ZZ/n\ZZ$ 的一个重要子集由在 $\ZZ/n\ZZ$ 中有乘法逆元的剩余类构成：
\[
  (\ZZ/n\ZZ)^\times
  =\{\overline{a}\in\ZZ/n\ZZ\mid \text{存在 }\overline{c}\in\ZZ/n\ZZ\text{ 使得 }\overline{a}\cdot\overline{c}=\overline{1}\}.
\]
后面一些练习给出了一个证明大纲，说明 $(\ZZ/n\ZZ)^\times$ 也正是那些代表元与 $n$ 互素的剩余类全体，这证明了下面的命题。

\medskip
\noindent\textbf{命题 4.}
\[
  (\ZZ/n\ZZ)^\times=\{\overline{a}\in\ZZ/n\ZZ\mid (a,n)=1\}.
\]
容易看出，若 $\overline{a}$ 的任意一个代表元与 $n$ 互素，则所有代表元都与 $n$ 互素，因此命题右端的集合是良定的。

\paragraph{例}
当 $n=9$ 时，由命题可得
\[
  (\ZZ/9\ZZ)^\times=\{\overline{1},\overline{2},\overline{4},\overline{5},\overline{7},\overline{8}\}.
\]
这些元素的乘法逆元分别是
\[
  \{\overline{1},\overline{5},\overline{7},\overline{2},\overline{4},\overline{8}\}.
\]

若整数 $a$ 与 $n$ 互素，则欧几里得算法产生整数 $x$ 与 $y$，使得 $ax+ny=1$；于是 $ax\equiv1\pmod n$，所以 $\overline{x}$ 是 $\ZZ/n\ZZ$ 中 $\overline{a}$ 的乘法逆元。这给出了一个计算 $\ZZ/n\ZZ$ 中乘法逆元的有效方法。

\paragraph{例}
设 $n=60$ 且 $a=17$。应用欧几里得算法得到
\begin{align*}
  60&=(3)17+9,\\
  17&=(1)9+8,\\
  9&=(1)8+1.
\end{align*}
因此 $a$ 与 $n$ 互素，并且 $(-7)17+(2)60=1$。于是 $\overline{-7}=\overline{53}$ 是 $\ZZ/60\ZZ$ 中 $\overline{17}$ 的乘法逆元。

\subsection*{练习}

\begin{enumerate}[label=\arabic*.]
  \item 明确写出 $\ZZ/18\ZZ$ 的各个剩余类中的所有元素。
  \item 证明：$\ZZ/n\ZZ$ 中不同的等价类恰好是 $\overline{0},\overline{1},\overline{2},\ldots,\overline{n-1}$（使用带余除法）。
  \item 证明：若 $a=a_n10^n+a_{n-1}10^{n-1}+\cdots+a_1 10+a_0$ 是任意正整数，则
  \[
    a\equiv a_n+a_{n-1}+\cdots+a_1+a_0\pmod 9
  \]
  （注意这就是通常的算术规则：除以 $9$ 的余数与十进制各位数字之和模 $9$ 的余数相同；特别地，一个整数能被 $9$ 整除，当且仅当它的各位数字之和能被 $9$ 整除）[注意 $10\equiv1\pmod9$]。
  \item 计算 $37^{100}$ 除以 $29$ 的余数。
  \item 计算 $9^{1500}$ 的最后两位数字。
  \item 证明：$\ZZ/4\ZZ$ 中各元素的平方只有 $\overline{0}$ 和 $\overline{1}$。
  \item 证明：对任意整数 $a$ 与 $b$，$a^2+b^2$ 除以 $4$ 的余数永远不是 $3$（使用上一练习）。
  \item 证明：方程 $a^2+b^2=3c^2$ 在非零整数 $a,b,c$ 中没有解。[像前两个练习那样模 $4$ 考虑该方程，并证明 $a,b,c$ 都必须能被 $2$ 整除。于是 $a^2,b^2,c^2$ 各有因子 $4$，除以 $4$ 后说明原方程会有一组更小的解。不断迭代，得到矛盾。]
  \item 证明：任意奇整数的平方除以 $8$ 的余数总是 $1$。
  \item 证明：$(\ZZ/n\ZZ)^\times$ 的元素个数是 $\varphi(n)$，其中 $\varphi$ 表示欧拉 $\varphi$ 函数。
  \item 证明：若 $\overline{a},\overline{b}\in(\ZZ/n\ZZ)^\times$，则 $\overline{a}\cdot\overline{b}\in(\ZZ/n\ZZ)^\times$。
  \item 设 $n\in\ZZ$，$n>1$，且 $a\in\ZZ$ 满足 $1\leq a\leq n$。证明：若 $a$ 与 $n$ 不互素，则存在整数 $b$，满足 $1\leq b<n$ 且 $ab\equiv0\pmod n$；并推出不存在整数 $c$ 使得 $ac\equiv1\pmod n$。
  \item 设 $n\in\ZZ$，$n>1$，且 $a\in\ZZ$ 满足 $1\leq a\leq n$。证明：若 $a$ 与 $n$ 互素，则存在整数 $c$ 使得 $ac\equiv1\pmod n$。[使用两个整数的最大公因子是这两个整数的一个 $\ZZ$-线性组合这一事实。]
  \item 从前两个练习推出：$(\ZZ/n\ZZ)^\times$ 是 $\ZZ/n\ZZ$ 中满足 $(a,n)=1$ 的元素 $\overline{a}$ 的集合，并由此证明命题 4。在 $n=12$ 的情形中直接验证这一点。
  \item 对下列每一对整数 $a$ 与 $n$，证明 $a$ 与 $n$ 互素，并求出 $\ZZ/n\ZZ$ 中 $\overline{a}$ 的乘法逆元。
  \begin{enumerate}[label=(\alph*)]
    \item $a=13,\ n=20$。
    \item $a=69,\ n=89$。
    \item $a=1891,\ n=3797$。
    \item $a=6003722857,\ n=77695236973$。[对这两个整数，欧几里得算法只需要 $3$ 步。]
  \end{enumerate}
  \item 编写一个计算机程序，对任意给定的 $n$ 作模 $n$ 的加法和乘法。这些运算的输出应为两个整数的和与积的最小剩余。同时加入如下功能：若 $(a,n)=1$，则在需要时输出一个介于 $1$ 与 $n-1$ 之间的整数 $c$，使得 $\overline{a}\cdot\overline{c}=\overline{1}$。（当然，你的程序不应直接调用许多系统中已经内置的 ``mod'' 函数。）
\end{enumerate}

 \setchapterintro{抽象代数的现代表述以一个看似极其简单的抽象定义——群的定义——开始。这个简单的定义很快引出涉及这类对象结构的困难问题。群有许多具体例子；当关于所有这些例子的结果都能通过对抽象群证明一个单一结果而获得时，抽象观点的威力便显现出来。
 	
 	然而，群的概念并非突然出现，而是长期数学研究的结晶。我们今天所用形式的抽象群的第一个正式定义出现在 1882 年。\footnote{以下大多数历史评述可参见 B. L. van der Waerden 的优秀著作 \textit{A History of Algebra}，Springer-Verlag，1980，以及其中的参考文献，尤其是 H. Wussing 的 \textit{The Genesis of the Abstract Group Concept: A Contribution to the History of the Origin of Abstract Group Theory}（Abe Shenitzer 由德文译出），MIT Press，1984。另可参见 A. Weil 的 \textit{Number Theory, An Approach Through History from Hammurapai to Legendre}，Birkh\"auser，1984。}抽象群的定义源于代数方程、数论和几何中的一些极为古老的问题；它的产生，是因为人们发现非常相似的技巧可应用于多种情形。正如奥托·赫尔德（Otto H\"older, 1859--1937）所说，数学的一个本质特征是：在应用某种算法或证明方法之后，人们会进一步考察该方法的适用范围与局限。于是，若干有趣对象共同具有的性质常被抽象出来，并引出问题：是否能确定所有具有这些性质的对象？试图回答这样的问题，往往也会极大地增进我们对原来那些对象的理解。正是在这种方式下，抽象群的定义发展成了对我们而言的抽象代数的出发点。
 	
 	下面用几个彼此相当不同的情形，说明后来被形式化为抽象群概念的思想是如何被使用的。
 	
 	\begin{enumerate}[label=(\arabic*)]
 		\item 在数论中，研究对象本身——整数集——就是一个群的例子。例如考虑我们称为“欧拉定理”的结果（参见第 3.2 节练习 22），它的一个极简单例子是：若整数 $a$ 既不能被 $2$ 整除也不能被 $5$ 整除，则 $a^{40}$ 的最后两位数字是 $01$。这由莱昂哈德·欧拉（Leonhard Euler, 1707--1783）于 1761 年利用约瑟夫·路易·拉格朗日（Joseph Louis Lagrange, 1736--1813）的“群论”思想证明，远在群的第一个正式定义出现之前。从我们的观点看，现在先证明“拉格朗日定理”（参见第 3.2 节定理 8），把这些技巧抽象到任意群上，然后把欧拉定理（以及许多其他结果）作为特殊情形恢复出来。
 		
 		\item 对形如 $y^2=x^3-2x$ 的代数方程的有理解问题的研究显示：任意连接两个解的直线与曲线 $y^2=x^3-2x$ 的交点又会给出一个解（这样的解有无穷多个，例如 $(0,0)$、$(-1,1)$、$(2,2)$、$(9/4,-21/8)$、$(-11169,23912197)$）。这样的“丢番图方程”以及其他方程，曾由皮埃尔·德·费马（Pierre de Fermat, 1601--1655）（他在 1644 年解决了这个例子）、欧拉、约 1777 年的拉格朗日等人研究。1730 年，欧拉提出求“双纽线微分” $dx/\sqrt{1-x^4}$ 的不定积分 $\int dx/\sqrt{1-x^4}$ 的问题；这个微分用于确定椭圆上的弧长（戈特弗里德·威廉·莱布尼茨（Gottfried Wilhelm Leibniz, 1646--1716）和约翰·伯努利（Johannes Bernoulli, 1667--1748）也曾考虑过该问题）。1752 年，欧拉利用 G. C. di Fagnano（1682--1766）的思想（欧拉于 1751 年收到）证明了这类椭圆积分的一个“乘法公式”，它说明两个椭圆积分如何产生第三个椭圆积分，从而使分析中的椭圆函数理论诞生。1834 年，卡尔·古斯塔夫·雅可比（Carl Gustav Jacob Jacobi, 1804--1851）观察到，欧拉关于求解某些丢番图方程的工作，实质上相当于为某些椭圆积分写出乘法公式。今天，上述曲线称为一条“椭圆曲线”，而这些问题被看作同一事实的两个不同方面：点上的这种几何运算可使椭圆曲线上的点集具有群结构。研究这些群的“算术”仍是当前研究的活跃领域。\footnote{见 J. Silverman, \textit{The Arithmetic of Elliptic Curves}, Springer-Verlag, 1986。}
 		
 		\item 到 1824 年，人们已经知道二次、三次和四次方程的根可以用公式表示（这推广了熟悉的方程 $ax^2+bx+c=0$ 的二次求根公式）。然而，在 1824 年，尼尔斯·亨利克·阿贝尔（Niels Henrik Abel, 1802--1829）证明五次方程的根不存在这样的公式（参见第 14.7 节推论 40）。证明的思想基于考察根在彼此之间置换时会发生什么（例如交换两个根）。这些置换组成的集合具有群结构（自然地称为“置换群”）。这一思想在埃瓦里斯特·伽罗瓦（\'Evariste Galois, 1811--1832）于 1830--1832 年关于显式“代换”群的优美工作中达到顶点。今天，这项工作称为伽罗瓦理论（也是本书第四部分的主题）。类似的显式群还由阿瑟·凯莱（Arthur Cayley, 1821--1895，约 1850 年）、卡米尔·若尔当（Camille Jordan, 1838--1922，约 1867 年）、费利克斯·克莱因（Felix Klein, 1849--1925，约 1870 年）等人在几何中作为几何变换（平移、反射等）的集合使用；群在几何中的应用，至今仍在三维空间、四维空间等结构的研究中非常活跃。在研究五次方程可解性时出现的同一个群，也出现在几何中正二十面体刚体运动的研究里，以及分析中椭圆函数的研究里。
 	\end{enumerate}
 	
 	今天的抽象群的先驱可以追溯到许多年以前，甚至早于伽罗瓦的“代换”群。作为我们出发点的抽象群的正式定义，于 1882 年出现在沃尔特·迪克（Walter Dyck, 1856--1934）和海因里希·韦伯（Heinrich Weber, 1842--1913）的工作中，前者是费利克斯·克莱因的助手。
 	
 	在数学研究中，常常先发现某个思想的具体应用，然后才把这个思想抽取出来，并作为一个值得单独研究的对象来呈现（例如，伽罗瓦在 1830 年的研究中已经隐含地使用了“商群”的概念，而抽象商群的定义则归功于 1889 年的赫尔德）。不论是否放在历史语境中，重要的是要认识到：之所以给出抽象定义，是因为把某些特定特征分离出来，并研究具有这些特征的对象被强加了怎样的结构，是很有用的。代数对象的结构这一概念（它通过同构的概念变得更精确——同构研究两个表面上不同的对象在何种意义下相同）将成为贯穿全文的一个主题。
 }
\part{群论}
 
\chapter{群导论}

\section{基本公理与例子}

本节引入第一部分将要研究的基本代数结构，并给出一些例子。

\noindent\textbf{定义.}
\begin{enumerate}[label=(\arabic*)]
  \item 集合 $G$ 上的一个\emph{二元运算} $*$ 是一个函数 $*:G\times G\to G$。对任意 $a,b\in G$，我们把 $*(a,b)$ 写作 $a*b$。
  \item 若集合 $G$ 上的二元运算 $*$ 满足对所有 $a,b,c\in G$，都有
  \[
    a*(b*c)=(a*b)*c,
  \]
  则称 $*$ 是\emph{结合的}。
  \item 若 $*$ 是集合 $G$ 上的二元运算，若 $G$ 中元素 $a,b$ 满足 $a*b=b*a$，则称 $a$ 与 $b$ \emph{交换}。若对所有 $a,b\in G$ 都有 $a*b=b*a$，则称 $*$（或称 $G$）是\emph{交换的}。
\end{enumerate}

\paragraph{例}
\begin{enumerate}[label=(\arabic*)]
  \item $+$（通常加法）是 $\ZZ$ 上的交换二元运算（在 $\QQ$、$\RR$ 或 $\CC$ 上亦然）。
  \item $\times$（通常乘法）是 $\ZZ$ 上的交换二元运算（在 $\QQ$、$\RR$ 或 $\CC$ 上亦然）。
  \item $-$（通常减法）是 $\ZZ$ 上的非交换二元运算，其中 $-(a,b)=a-b$。映射 $a\mapsto -a$ 不是二元运算（不是二元的）。
  \item $-$ 不是 $\ZZ^+$ 上的二元运算（也不是 $\QQ^+$ 或 $\RR^+$ 上的二元运算），因为若 $a,b\in\ZZ^+$ 且 $a<b$，则 $a-b\notin\ZZ^+$；也就是说，$-$ 并不把 $\ZZ^+\times\ZZ^+$ 映入 $\ZZ^+$。
  \item 三维空间 $\RR^3$ 中两个向量的叉积是一个既不结合也不交换的二元运算。
\end{enumerate}

设 $*$ 是集合 $G$ 上的二元运算，且 $H$ 是 $G$ 的子集。若 $*$ 到 $H$ 的限制是 $H$ 上的二元运算，即对所有 $a,b\in H$ 都有 $a*b\in H$，则称 $H$ 对 $*$ \emph{封闭}。注意，如果 $*$ 是 $G$ 上的结合（相应地，交换）二元运算，并且 $*$ 限制到 $G$ 的某个子集 $H$ 上是 $H$ 的二元运算，那么 $*$ 在 $H$ 上也自动是结合的（相应地，交换的）。

\noindent\textbf{定义.}
\begin{enumerate}[label=(\arabic*)]
  \item \emph{群}是一个有序对 $(G,*)$，其中 $G$ 是一个集合，$*$ 是 $G$ 上的二元运算，并满足以下公理：
  \begin{enumerate}[label=(\roman*)]
    \item 对所有 $a,b,c\in G$，$(a*b)*c=a*(b*c)$，即 $*$ 是结合的；
    \item 存在 $G$ 中的元素 $e$，称为 $G$ 的\emph{单位元}，使得对所有 $a\in G$，都有 $a*e=e*a=a$；
    \item 对每个 $a\in G$，存在 $G$ 中的元素 $a^{-1}$，称为 $a$ 的\emph{逆元}，使得 $a*a^{-1}=a^{-1}*a=e$。
  \end{enumerate}
  \item 若群 $(G,*)$ 满足对所有 $a,b\in G$ 都有 $a*b=b*a$，则称它是\emph{阿贝尔的}（或\emph{交换的}）。
\end{enumerate}

我们很快会不那么形式化：若 $(G,*)$ 是一个群，就说 $G$ 在 $*$ 下是一个群（当运算 $*$ 从上下文中清楚时，干脆说 $G$ 是一个群）。此外，若 $G$ 还是一个有限集，则称 $G$ 是\emph{有限群}。注意，公理 (ii) 保证群总是非空的。

\paragraph{例}
\begin{enumerate}[label=(\arabic*)]
  \item $\ZZ$、$\QQ$、$\RR$ 和 $\CC$ 在 $+$ 下是群，其中 $e=0$，并且对所有 $a$，$a^{-1}=-a$。
  \item $\QQ-\{0\}$、$\RR-\{0\}$、$\CC-\{0\}$、$\QQ^+$、$\RR^+$ 在 $\times$ 下是群，其中 $e=1$，并且对所有 $a$，$a^{-1}=1/a$。但是要注意，$\ZZ-\{0\}$ 在 $\times$ 下不是群，因为虽然 $\times$ 是 $\ZZ-\{0\}$ 上的结合二元运算，但元素 $2$（例如）在 $\ZZ-\{0\}$ 中没有逆元。
\end{enumerate}

我们在这些熟悉的例子中略过了结合律成立这一事实。对 $\ZZ$ 的加法而言，它是自然数加法结合公理的结果。$\QQ$ 上加法的结合律由 $\ZZ$ 的结合律推出——稍后在从 $\ZZ$ 严格构造 $\QQ$ 时会给出证明概要（参见第 7.5 节）。$\RR$ 以及进而 $\CC$ 上加法的结合律，则在初等分析课程中通过把 $\RR$ 构造为 $\QQ$ 的完备化来证明；最终，结合律仍是 $\ZZ$ 结合律的结果。乘法的结合公理也可通过类似的发展建立，先从 $\ZZ$ 开始。由于 $\RR$ 和 $\CC$ 在本书中主要用于说明性目的，而且我们不会从 $\QQ$ 构造 $\RR$（虽然会从 $\RR$ 构造 $\CC$），所以我们把 $\RR$ 与 $\CC$ 在 $+$ 和 $\times$ 下的结合律作为已知事实。

\paragraph{例（续）}
\begin{enumerate}[label=(\arabic*),start=3]
  \item 向量空间 $V$ 的公理包括规定 $(V,+)$ 是阿贝尔群的那些公理（运算 $+$ 称为向量加法）。因此，任何向量空间，例如 $\RR^n$，特别地都是一个加法群。
  \item 对 $n\in\ZZ^+$，$\ZZ/n\ZZ$ 在第 0 章所述的剩余类加法 $+$ 下是阿贝尔群。我们将在第 3 章（在更一般的语境中）证明这个二元运算 $+$ 是良定且结合的；目前先承认这一点。这个群的单位元是元素 $\overline{0}$，对每个 $\overline{a}\in\ZZ/n\ZZ$，其逆元是 $\overline{-a}$。今后当我们谈到群 $\ZZ/n\ZZ$ 时，默认其群运算是模 $n$ 的类加法。
  \item 对 $n\in\ZZ^+$，集合 $(\ZZ/n\ZZ)^\times$ 由所有在模 $n$ 下有乘法逆元的等价类 $\overline{a}$ 组成；它在第 0 章所述的剩余类乘法下是阿贝尔群。同样，我们暂且承认该运算是良定且结合的。这个群的单位元是元素 $\overline{1}$，并且按 $(\ZZ/n\ZZ)^\times$ 的定义，每个元素都有乘法逆元。今后当我们谈到群 $(\ZZ/n\ZZ)^\times$ 时，默认其群运算是模 $n$ 的类乘法。
  \item 若 $(A,*)$ 和 $(B,\circ)$ 是群，则可以构造一个新群 $A\times B$，称为它们的\emph{直积}，其元素是笛卡儿积
  \[
    A\times B=\{(a,b)\mid a\in A,\ b\in B\},
  \]
  中的元素，运算按分量定义：
  \[
    (a_1,b_1)(a_2,b_2)=(a_1*a_2,\ b_1\circ b_2).
  \]
  例如，若取 $A=B=\RR$（二者的运算都是加法），则 $\RR\times\RR$ 就是熟悉的欧氏平面。两个群的直积仍是群的证明留作一个直接的练习（以后会做）：$A\times B$ 中每条群公理的成立，都是 $A$ 与 $B$ 中相应公理成立以及 $A\times B$ 中运算按分量定义的结果。
\end{enumerate}

不要混淆群 $\ZZ/n\ZZ$（加法下）与 $(\ZZ/n\ZZ)^\times$（乘法下），即使后者是前者的子集；上标 $\times$ 将总表示运算是乘法。

在继续更复杂的例子之前，我们先证明两个基本结果，它们尤其使我们可以谈论单位元和元素的逆元。

\medskip
\noindent\textbf{命题 1.} 若 $G$ 在运算 $*$ 下是群，则
\begin{enumerate}[label=(\arabic*)]
  \item $G$ 的单位元唯一；
  \item 对每个 $a\in G$，$a^{-1}$ 唯一确定；
  \item 对所有 $a\in G$，$(a^{-1})^{-1}=a$；
  \item $(a*b)^{-1}=(b^{-1})*(a^{-1})$；
  \item 对任意 $a_1,a_2,\ldots,a_n\in G$，$a_1*a_2*\cdots*a_n$ 的值与该表达式如何加括号无关（这称为\emph{广义结合律}）。
\end{enumerate}

\begin{proof}
(1) 若 $f$ 和 $g$ 都是单位元，则由群定义的公理 (ii)，$f*g=f$（取 $a=f$、$e=g$）。同一公理也给出 $f*g=g$（取 $a=g$、$e=f$）。故 $f=g$，单位元唯一。

(2) 假设 $b$ 和 $c$ 都是 $a$ 的逆元，令 $e$ 为 $G$ 的单位元。由公理 (iii)，$a*b=e$ 且 $c*a=e$。于是
\begin{align*}
  c &= c*e &&\text{（$e$ 的定义——公理 (ii)）}\\
    &= c*(a*b) &&\text{（因为 $e=a*b$）}\\
    &= (c*a)*b &&\text{（结合律）}\\
    &= e*b &&\text{（因为 $e=c*a$）}\\
    &= b &&\text{（公理 (ii)）。}
\end{align*}

(3) 证明 $(a^{-1})^{-1}=a$，正是要证明 $a$ 是 $a^{-1}$ 的逆元（由于第 (2) 部分，$a^{-1}$ 的逆元唯一）。把 $a$ 与 $a^{-1}$ 的角色互换来看 $a^{-1}$ 的定义，便知 $a$ 满足 $a^{-1}$ 的逆元的定义性质，所以 $a$ 是 $a^{-1}$ 的逆元。

(4) 令 $c=(a*b)^{-1}$，则按 $c$ 的定义，$(a*b)*c=e$。由结合律，
\[
  a*(b*c)=e.
\]
左乘 $a^{-1}$ 得
\[
  a^{-1}*(a*(b*c))=a^{-1}*e.
\]
左边用结合律、右边用 $e$ 的定义，得
\[
  (a^{-1}*a)*(b*c)=a^{-1},
\]
因此
\[
  e*(b*c)=a^{-1},
\]
从而
\[
  b*c=a^{-1}.
\]
现在左乘 $b^{-1}$ 并同样化简：
\begin{align*}
  b^{-1}*(b*c)&=b^{-1}*a^{-1},\
  (b^{-1}*b)*c&=b^{-1}*a^{-1},\
  e*c&=b^{-1}*a^{-1},\
  c&=b^{-1}*a^{-1},
\end{align*}
这就是所要证明的。

(5) 这是一个用 $n$ 作归纳的很好的练习。先证明 $n=1,2,3$ 的情形。然后假设对任意 $k<n$，$k$ 个元素 $b_1*b_2*\cdots*b_k$ 的任意加括号方式都能在不改变乘积值的情况下化为形如
\[
  b_1*(b_2*(b_3*(\cdots*b_k))\cdots)
\]
的表达式。任意给 $a_1*a_2*\cdots*a_n$ 加括号，必然把它分成两个子乘积，比如
\[
  (a_1*a_2*\cdots*a_k)*(a_{k+1}*a_{k+2}*\cdots*a_n),
\]
每个子乘积又以某种方式加括号。对这两个子乘积分别应用归纳假设，最后把结果化为
\[
  a_1*(a_2*(a_3*(\cdots*a_n))\cdots),
\]
即可完成归纳。
\end{proof}

注意，在命题 1 的整个证明中，我们都小心地没有改变任何乘积的顺序（除非公理 (ii) 和 (iii) 允许这样做），因为 $G$ 可能是非阿贝尔群。

\paragraph{记号}
\begin{enumerate}[label=(\arabic*)]
  \item 对抽象群 $G$ 来说，在计算中不断写运算 $*$ 是很麻烦的。今后（除非必要）我们的抽象群 $G,H$ 等都把运算写作 $\cdot$，并把 $a\cdot b$ 写作 $ab$。鉴于广义结合律，三个或更多群元素的乘积不再加括号（虽然运算仍然是二元运算）。最后，对抽象群 $G$（运算为 $\cdot$），我们把 $G$ 的单位元记为 $1$。
  \item 对任意群 $G$（省略运算 $\cdot$）和 $x\in G$，以及 $n\in\ZZ^+$，由于乘积 $xx\cdots x$（$n$ 项）与加括号方式无关，记作 $x^n$。把 $x^{-1}x^{-1}\cdots x^{-1}$（$n$ 项）记作 $x^{-n}$。令 $x^0=1$，即 $G$ 的单位元。
\end{enumerate}

这个新记号非常简洁。当然，处理具体群时，我们仍使用自然的（给定的）运算。例如，当运算是 $+$ 时，单位元记为 $0$；任一元素 $a$ 的逆元 $a^{-1}$ 写作 $-a$；$a+a+\cdots+a$（$n>0$ 项）写作 $na$；$-a-a-\cdots-a$（$n$ 项）写作 $-na$；并且 $0a=0$。

\medskip
\noindent\textbf{命题 2.} 设 $G$ 是群，$a,b\in G$。方程 $ax=b$ 与 $ya=b$ 在 $G$ 中分别有唯一解 $x,y$。特别地，$G$ 中左右消去律成立，即
\begin{enumerate}[label=(\arabic*)]
  \item 若 $au=av$，则 $u=v$；
  \item 若 $ub=vb$，则 $u=v$。
\end{enumerate}

\begin{proof}
通过左乘 $a^{-1}$ 并化简，可解出 $ax=b$ 的解为 $x=a^{-1}b$。$x$ 的唯一性来自 $a^{-1}$ 的唯一性。类似地，若 $ya=b$，则 $y=ba^{-1}$。若 $au=av$，左乘 $a^{-1}$ 并化简得 $u=v$。右消去律同理。
\end{proof}

命题 2 的一个后果是：若 $a$ 是 $G$ 的任一元素，且对某个 $b\in G$ 有 $ab=e$ 或 $ba=e$，则 $b=a^{-1}$；也就是说，我们不必检验两个等式都成立。另外，若对某个 $b\in G$ 有 $ab=a$（或 $ba=a$），则 $b$ 必为 $G$ 的单位元；也就是说，我们不必检验对所有 $x\in G$ 都有 $bx=xb=x$。

\noindent\textbf{定义.} 设 $G$ 是群，$x\in G$。若存在正整数 $n$ 使 $x^n=1$，则把满足此式的最小正整数 $n$ 称为 $x$ 的\emph{阶}，并记为 $|x|$；此时称 $x$ 的阶为 $n$。若 $x$ 的任一正幂都不是单位元，则定义 $x$ 的阶为无穷，并称 $x$ 是\emph{无限阶}的。

元素 $x$ 的阶的符号不应与绝对值符号混淆（当 $G\subseteq\RR$ 时我们会小心区分）。用同一个符号表示一个元素的阶和一个集合的基数（或阶）似乎不太谨慎；不过我们将看到，群中元素的阶等于其所有（不同）幂组成的集合的基数，因此“阶”这个词的两种用法是自然相关的。

\paragraph{例}
\begin{enumerate}[label=(\arabic*)]
  \item 群中的元素阶为 $1$ 当且仅当它是单位元。
  \item 在加法群 $\ZZ$、$\QQ$、$\RR$ 或 $\CC$ 中，每个非零（即非单位）元素都是无限阶的。
  \item 在乘法群 $\RR-\{0\}$ 或 $\QQ-\{0\}$ 中，元素 $-1$ 的阶为 $2$，其他所有非单位元素都是无限阶的。
  \item 在加法群 $\ZZ/9\ZZ$ 中，元素 $\overline{6}$ 的阶为 $3$，因为 $\overline{6}\neq\overline{0}$，$\overline{6}+\overline{6}=\overline{12}=\overline{3}\neq\overline{0}$，但 $\overline{6}+\overline{6}+\overline{6}=\overline{18}=\overline{0}$，即这个群的单位元。回忆，在加法群中，一个元素的幂就是该元素的整数倍。类似地，元素 $\overline{5}$ 的阶为 $9$，因为 $45$ 是能被 $9$ 整除的 $5$ 的最小正倍数。
  \item 在乘法群 $(\ZZ/7\ZZ)^\times$ 中，元素 $\overline{2}$ 的幂为 $\overline{2},\overline{4},\overline{8}=\overline{1}$，其中 $\overline{1}$ 是这个群的单位元，所以 $\overline{2}$ 的阶为 $3$。类似地，元素 $\overline{3}$ 的阶为 $6$，因为 $3^6$ 是同余于 $1\pmod 7$ 的 $3$ 的最小正幂。
\end{enumerate}

\noindent\textbf{定义.} 设 $G=\{g_1,g_2,\ldots,g_n\}$ 是有限群，且 $g_1=1$。$G$ 的\emph{乘法表}或\emph{群表}是一个 $n\times n$ 矩阵，其第 $i,j$ 项为群元素 $g_ig_j$。

对有限群来说，乘法表在某种意义上包含了关于群的全部信息。然而，从计算角度看，它是一个笨重的对象（大小为群阶的平方）；从视觉角度看，它也不是判断群性质的好工具。可以把群表类比为一张记录全国任意两座城市间距离的表。这样的表有用，并且本质上记录了所有距离关系；但一幅地图（更好的是标出所有距离的地图）显然更便于使用。我们对群论（尤其是有限群）最初的展开，部分目的正是寻找一种更概念化的方式来观察群的内部结构。

\subsection*{练习}

设 $G$ 是一个群。
\begin{enumerate}[label=\arabic*.]
  \item 判断下列二元运算哪些是结合的：
  \begin{enumerate}[label=(\alph*)]
    \item $\ZZ$ 上由 $a*b=a-b$ 定义的运算 $*$；
    \item $\RR$ 上由 $a*b=a+b+ab$ 定义的运算 $*$；
    \item $\QQ$ 上由 $a*b=(a+b)/5$ 定义的运算 $*$；
    \item $\ZZ\times\ZZ$ 上由 $(a,b)*(c,d)=(ad+bc,bd)$ 定义的运算 $*$；
    \item $\QQ-\{0\}$ 上由 $a*b=a/b$ 定义的运算 $*$。
  \end{enumerate}
  \item 判断前一题中的哪些二元运算是交换的。
  \item 证明 $\ZZ/n\ZZ$ 中剩余类加法是结合的（可假设它是良定的）。
  \item 证明 $\ZZ/n\ZZ$ 中剩余类乘法是结合的（可假设它是良定的）。
  \item 证明对所有 $n>1$，$\ZZ/n\ZZ$ 在剩余类乘法下不是群。
  \item 判断下列集合在加法下哪些是群：
  \begin{enumerate}[label=(\alph*)]
    \item 最简形式下分母为奇数的有理数集合（包括 $0=0/1$）；
    \item 最简形式下分母为偶数的有理数集合（包括 $0=0/1$）；
    \item 绝对值 $<1$ 的有理数集合；
    \item 绝对值 $\geq 1$ 的有理数集合再加上 $0$；
    \item 分母等于 $1$ 或 $2$ 的有理数集合；
    \item 分母等于 $1$、$2$ 或 $3$ 的有理数集合。
  \end{enumerate}
  \item 令 $G=\{x\in\RR\mid 0\leq x<1\}$，对 $x,y\in G$，令 $x*y$ 为 $x+y$ 的小数部分（即 $x*y=x+y-[x+y]$，其中 $[a]$ 是不超过 $a$ 的最大整数）。证明 $*$ 是 $G$ 上良定的二元运算，并且 $G$ 在 $*$ 下是阿贝尔群（称为模 $1$ 的实数）。
  \item 令 $G=\{z\in\CC\mid z^n=1\text{ 对某个 }n\in\ZZ^+\}$。
  \begin{enumerate}[label=(\alph*)]
    \item 证明 $G$ 在乘法下是群（称为 $\CC$ 中的单位根群）。
    \item 证明 $G$ 在加法下不是群。
  \end{enumerate}
  \item 令 $G=\{a+b\sqrt2\in\RR\mid a,b\in\QQ\}$。
  \begin{enumerate}[label=(\alph*)]
    \item 证明 $G$ 在加法下是群。
    \item 证明 $G$ 的非零元素在乘法下是群。［“有理化分母”以求乘法逆元。］
  \end{enumerate}
  \item 证明有限群是阿贝尔群，当且仅当它的群表是对称矩阵。
  \item 求加法群 $\ZZ/12\ZZ$ 中每个元素的阶。
  \item 求乘法群 $(\ZZ/12\ZZ)^\times$ 中下列元素的阶：$\overline{1},\overline{-1},\overline{5},\overline{7},\overline{-5},\overline{-7}$。
  \item 求加法群 $\ZZ/36\ZZ$ 中下列元素的阶：$\overline{1},\overline{2},\overline{6},\overline{9},\overline{10},\overline{12},\overline{-1},\overline{-10},\overline{-18}$。
  \item 求乘法群 $(\ZZ/36\ZZ)^\times$ 中下列元素的阶：$\overline{1},\overline{-1},\overline{5},\overline{13},\overline{-13},\overline{17}$。
  \item 证明 $(a_1a_2\cdots a_n)^{-1}=a_n^{-1}a_{n-1}^{-1}\cdots a_1^{-1}$，对所有 $a_1,a_2,\ldots,a_n\in G$ 成立。
  \item 设 $x\in G$。证明 $x^2=1$ 当且仅当 $|x|$ 为 $1$ 或 $2$。
  \item 设 $x\in G$。证明若 $|x|=n$，其中 $n$ 为正整数，则 $x^{-1}=x^{n-1}$。
  \item 设 $x,y\in G$。证明 $xy=yx$ 当且仅当 $y^{-1}xy=x$，当且仅当 $x^{-1}y^{-1}xy=1$。
  \item 设 $x\in G$，$a,b\in\ZZ^+$。
  \begin{enumerate}[label=(\alph*)]
    \item 证明 $x^{a+b}=x^ax^b$ 且 $(x^a)^b=x^{ab}$。
    \item 证明 $(x^a)^{-1}=x^{-a}$。
    \item 对任意整数 $a,b$（正、负或零）建立 (a) 的结论。
  \end{enumerate}
  \item 对 $G$ 中元素 $x$，证明 $x$ 与 $x^{-1}$ 有相同的阶。
  \item 设 $G$ 是有限群，$x\in G$ 的阶为 $n$。证明若 $n$ 为奇数，则对某个 $k$ 有 $x=(x^2)^k$。
  \item 若 $x,g\in G$，证明 $|x|=|g^{-1}xg|$。推出对所有 $a,b\in G$，有 $|ab|=|ba|$。
  \item 假设 $x\in G$ 且 $|x|=n<\infty$。若 $n=st$，其中 $s,t$ 为正整数，证明 $|x^s|=t$。
  \item 若 $a$ 与 $b$ 是 $G$ 中可交换的元素，证明对所有 $n\in\ZZ$，$(ab)^n=a^nb^n$。［先对正 $n$ 用归纳法。］
  \item 证明若对所有 $x\in G$ 都有 $x^2=1$，则 $G$ 是阿贝尔群。
  \item 假设 $H$ 是 $(G,*)$ 的非空子集，且对 $G$ 的二元运算封闭并且对取逆封闭，即对所有 $h,k\in H$，有 $hk\in H$ 且 $h^{-1}\in H$。证明 $H$ 在限制到 $H$ 的运算 $*$ 下是群（这样的子集 $H$ 称为 $G$ 的\emph{子群}）。
  \item 证明若 $x$ 是群 $G$ 的一个元素，则 $\{x^n\mid n\in\ZZ\}$ 是 $G$ 的子群（参见前一题），称为由 $x$ 生成的\emph{循环子群}。
  \item 设 $(A,*)$ 与 $(B,\circ)$ 是群，$A\times B$ 是它们的直积（如例 6 定义）。验证 $A\times B$ 的所有群公理：
  \begin{enumerate}[label=(\alph*)]
    \item 证明结合律成立：对所有 $(a_i,b_i)\in A\times B$，$i=1,2,3$，
    \[
      (a_1,b_1)[(a_2,b_2)(a_3,b_3)]=[(a_1,b_1)(a_2,b_2)](a_3,b_3);
    \]
    \item 证明 $(1,1)$ 是 $A\times B$ 的单位元；
    \item 证明 $(a,b)$ 的逆元是 $(a^{-1},b^{-1})$。
  \end{enumerate}
  \item 证明 $A\times B$ 是阿贝尔群，当且仅当 $A$ 与 $B$ 都是阿贝尔群。
  \item 证明 $A\times B$ 中的元素 $(a,1)$ 与 $(1,b)$ 交换，并推出 $(a,b)$ 的阶是 $|a|$ 与 $|b|$ 的最小公倍数。
  \item 证明任一偶数阶有限群 $G$ 含有阶为 $2$ 的元素。［令 $t(G)=\{g\in G\mid g\neq g^{-1}\}$。证明 $t(G)$ 的元素个数为偶数，并且 $G-t(G)$ 中每个非单位元素的阶为 $2$。］
  \item 若 $x$ 是 $G$ 中有限阶 $n$ 的元素，证明 $1,x,x^2,\ldots,x^{n-1}$ 两两不同。推出 $|x|\leq |G|$。
  \item 设 $x$ 是 $G$ 中有限阶 $n$ 的元素。
  \begin{enumerate}[label=(\alph*)]
    \item 证明若 $n$ 为奇数，则对所有 $i=1,2,\ldots,n-1$，有 $x^i\neq x^{-i}$。
    \item 证明若 $n=2k$ 且 $1\leq i<n$，则 $x^i=x^{-i}$ 当且仅当 $i=k$。
  \end{enumerate}
  \item 若 $x$ 是 $G$ 中无限阶元素，证明元素 $x^n$（$n\in\ZZ$）两两不同。
  \item 若 $x$ 是 $G$ 中有限阶 $n$ 的元素，用带余除法证明 $x$ 的任意整数幂都等于集合 $\{1,x,x^2,\ldots,x^{n-1}\}$ 中的某个元素（因此这些就是由 $x$ 生成的循环子群的全部不同元素，参见上面的练习 27）。
  \item 假设 $G=\{1,a,b,c\}$ 是阶为 $4$、单位元为 $1$ 的群。再假设 $G$ 没有阶为 $4$ 的元素（所以由练习 32，每个元素阶都 $\leq 3$）。用消去律证明 $G$ 的群表唯一。推出 $G$ 是阿贝尔群。
\end{enumerate}

\section{二面体群}

群的一类重要例子是几何对象的对称所组成的群。最简单的子类是几何对象为正平面图形的情形。

对每个 $n\in\ZZ^+$，$n\geq3$，令 $D_{2n}$ 为正 $n$ 边形的所有对称组成的集合。这里的对称是指正 $n$ 边形的任一刚体运动，它可以通过取该正 $n$ 边形的一个副本，在三维空间中以任意方式移动它，然后把该副本放回原正 $n$ 边形上，使之恰好覆盖原图形来实现。更精确地说，我们可以先给 $n$ 个顶点编号，例如如下图所示，由此描述这些对称。

\begin{center}
\begin{tikzpicture}[scale=1.0]
  \def\n{7}
  \foreach \i in {1,...,7} {
    \coordinate (P\i) at ({90-360*(\i-1)/7}:1.6);
    \fill (P\i) circle (1.2pt);
    \node at ({90-360*(\i-1)/7}:1.95) {$\i$};
  }
  \draw (P1)--(P2)--(P3)--(P4)--(P5)--(P6)--(P7)--cycle;
\end{tikzpicture}
\end{center}

于是每个对称 $s$ 都可由 $\{1,2,3,\ldots,n\}$ 的相应置换 $\sigma$ 唯一描述：若对称 $s$ 把顶点 $i$ 放到原来顶点 $j$ 所在的位置，则 $\sigma$ 是把 $i$ 送到 $j$ 的置换。例如，若 $s$ 是绕 $n$ 边形中心顺时针旋转 $2\pi/n$ 弧度，则 $\sigma$ 把 $i$ 送到 $i+1$，$1\leq i\leq n-1$，并且 $\sigma(n)=1$。现在把 $D_{2n}$ 作成一个群：对 $s,t\in D_{2n}$，定义 $st$ 为先对 $n$ 边形施行 $t$，再施行 $s$ 所得到的对称（注意我们把对称看作 $n$ 边形上的函数，所以 $st$ 就是函数复合——按通常方式从右向左读）。若 $s,t$ 在顶点上分别实现置换 $\sigma,\tau$，则 $st$ 实现 $\sigma\circ\tau$。$D_{2n}$ 上的二元运算是结合的，因为函数复合是结合的。$D_{2n}$ 的单位元是恒等对称（它固定所有顶点），记作 $1$；$s\in D_{2n}$ 的逆元是把 $s$ 的所有刚体运动反向的对称（所以若 $s$ 在顶点上实现置换 $\sigma$，则 $s^{-1}$ 实现 $\sigma^{-1}$）。下一段我们说明
\[
  |D_{2n}|=2n,
\]
因此 $D_{2n}$ 称为\emph{阶为 $2n$ 的二面体群}。有些教材把这个群记为 $D_n$；不过在群论文献中，$D_{2n}$（下标给出群的阶而不是顶点数）更常见。

为求阶 $|D_{2n}|$，注意给定任一顶点 $i$，存在一个对称把顶点 $1$ 送到位置 $i$。由于顶点 $2$ 与顶点 $1$ 相邻，顶点 $2$ 必须落在位置 $i+1$ 或 $i-1$（其中 $n+1$ 读作 $1$，$1-1$ 读作 $n$，即顶点编号按模 $n$ 读取）。而且，在第一个对称之后再复合关于通过顶点 $i$ 和 $n$ 边形中心的直线的反射，可见顶点 $2$ 可由某个对称送到位置 $i+1$ 或 $i-1$ 中任意一个。因此，应用对称后，有序顶点对 $1,2$ 有 $n\cdot2$ 种可能位置。由于对称是刚体运动，一旦有序顶点对 $1,2$ 的位置被指定，该对称在所有其余顶点上的作用就完全确定。因此正 $n$ 边形恰有 $2n$ 个对称。而且我们可以显式列出这 $2n$ 个对称：绕中心旋转 $2\pi i/n$ 弧度的 $n$ 个旋转，$0\leq i\leq n-1$，以及关于 $n$ 条对称轴的 $n$ 个反射（若 $n$ 为奇数，每条对称轴通过一个顶点和对边的中点；若 $n$ 为偶数，则有 $n/2$ 条对称轴通过两个相对顶点，有 $n/2$ 条对称轴垂直平分两条相对边）。例如，若 $n=4$，并在 $x,y$ 平面中把正方形画在原点处，则对称轴为
\[
  x=0\quad(y\text{-轴}),\qquad y=0\quad(x\text{-轴}),\qquad y=x,
  \qquad y=-x,
\]
（注意，关于原点的“反射”并不是反射，而是旋转 $\pi$ 弧度）。

\begin{center}
\begin{tikzpicture}[scale=1.2]
  \draw[->] (-1.6,0)--(1.7,0) node[right] {$x$};
  \draw[->] (0,-1.6)--(0,1.7) node[above] {$y$};
  \draw (-1,1)--(1,1)--(1,-1)--(-1,-1)--cycle;
  \node at (1.15,1.15) {$1$};
  \node at (1.15,-1.15) {$2$};
  \node at (-1.15,-1.15) {$3$};
  \node at (-1.15,1.15) {$4$};
  \draw[dashed] (-1.4,-1.4)--(1.4,1.4) node[right] {$y=x$};
  \draw[dashed] (-1.4,1.4)--(1.4,-1.4) node[right] {$y=-x$};
\end{tikzpicture}
\end{center}

由于二面体群将在全文中被大量用作例子，我们固定一些记号，并提及一些计算；它们将简化以后的计算，也有助于把 $D_{2n}$ 看作一个抽象群（而不必每次都回到几何背景）。固定一个以原点为中心、位于 $x,y$ 平面中的正 $n$ 边形，并按顺时针方向依次给顶点标号 $1$ 到 $n$。令 $r$ 为绕原点顺时针旋转 $2\pi/n$ 弧度的旋转。令 $s$ 为关于通过顶点 $1$ 和原点的对称轴的反射（对每个 $n$ 都用同样的字母，但上下文总会使 $n$ 明确）。下面计算的细节留作练习（我们大多将处理 $D_6$ 和 $D_8$，读者也许可以先对 $n=3$ 和 $n=4$ 试做这些练习）：
\begin{enumerate}[label=(\arabic*)]
  \item $1,r,r^2,\ldots,r^{n-1}$ 两两不同，并且 $r^n=1$，所以 $|r|=n$。
  \item $|s|=2$。
  \item 对任意 $i$，$s\neq r^i$。
  \item 若 $0\leq i,j\leq n-1$ 且 $i\neq j$，则 $sr^i\neq sr^j$，所以
  \[
    D_{2n}=\{1,r,r^2,\ldots,r^{n-1},s,sr,sr^2,\ldots,sr^{n-1}\},
  \]
  即每个元素可唯一写成 $s^kr^i$ 的形式，其中 $k=0$ 或 $1$，且 $0\leq i\leq n-1$。
  \item $rs=sr^{-1}$。［先求出 $s$ 在 $\{1,2,\ldots,n\}$ 上实现的置换，再分别计算此等式两边对顶点 $1$ 和 $2$ 的作用。］这尤其说明 $r$ 与 $s$ 不交换，因此 $D_{2n}$ 是非阿贝尔的。
  \item 对所有 $0\leq i\leq n$，有 $r^is=sr^{-i}$。［对 $i$ 作归纳，并利用 $r^{i+1}s=r(r^is)$ 以及前面的计算。］这说明如何把 $s$ 与 $r$ 的幂交换位置。
\end{enumerate}

完成这些计算后，我们现在观察到，$D_{2n}$ 的完整乘法表可以仅用 $r$ 和 $s$ 写出；也就是说，$D_{2n}$ 的所有元素都可唯一表示为 $s^kr^i$，其中 $k=0$ 或 $1$，$0\leq i\leq n-1$，而且任意两个这种形式的元素的乘积，都可以只用“关系” (1)、(2)、(6)（把所有指数模 $n$ 化简）化为另一个这种形式的元素。例如，若 $n=12$，则
\[
  (sr^9)(sr^6)=s(r^9s)r^6=s(sr^{-9})r^6=s^2r^{-9+6}=r^{-3}=r^9.
\]

\subsection*{生成元与关系}

二面体群中生成元 $r$ 与 $s$ 的使用，为在 $D_{2n}$ 中计算提供了一种简单而简洁的方式。我们可以类似地为任意群引入生成元与关系的概念。较早拥有这些概念（在其形式化证明之前）是有用的，因为它们提供了描述和计算许多群的简单方式。生成元将在第 2.4 节中更详细地讨论；当第 6.3 节引入自由群概念时，这两个概念都会得到严格处理。

群 $G$ 的元素子集 $S$ 若具有如下性质：$G$ 的每个元素都可以写成 $S$ 中元素及其逆元的一个（有限）乘积，则称 $S$ 是 $G$ 的一个\emph{生成元集}。记号上我们写作 $G=\langle S\rangle$，并说 $G$ \emph{由 $S$ 生成}，或 $S$ \emph{生成} $G$。例如，整数 $1$ 是加法群 $\ZZ$ 的一个生成元，因为每个整数都是有限多个 $+1$ 与 $-1$ 的和，所以 $\ZZ=\langle 1\rangle$。由 $D_{2n}$ 的性质 (4)，集合 $S=\{r,s\}$ 是 $D_{2n}$ 的一个生成元集，所以 $D_{2n}=\langle r,s\rangle$。我们稍后将看到，在有限群 $G$ 中，若 $G$ 的每个元素都是 $S$ 中元素的有限乘积，则 $S$ 生成 $G$（即不必另外把 $S$ 中元素的逆元也包括进来）。

在一般群 $G$ 中，生成元满足的任何方程称为 $G$ 中的\emph{关系}。因此在 $D_{2n}$ 中有关系：$r^n=1$、$s^2=1$ 和 $rs=sr^{-1}$。此外，在 $D_{2n}$ 中这三个关系还具有额外性质：群元素之间的任何其他关系都可由这三个关系推出（这并非立即显然；它来自这样一个事实：只用这三个关系就能确定两个群元素何时相等）。

一般地，若某个群 $G$ 由子集 $S$ 生成，并且有一组关系，例如 $R_1,R_2,\ldots,R_m$（这里每个 $R_i$ 都是由 $S\cup\{1\}$ 中元素写成的一个方程），使得 $S$ 中元素之间的任一关系都能由这些关系推出，则称这些生成元与关系给出了 $G$ 的一个\emph{表示}，并写作
\[
  G=\langle S\mid R_1,R_2,\ldots,R_m\rangle.
\]
因此，二面体群 $D_{2n}$ 的一个表示（使用上述生成元和关系）是
\begin{equation}
  D_{2n}=\langle r,s\mid r^n=s^2=1,\ rs=sr^{-1}\rangle. \tag{1.1}
\end{equation}
我们将看到，用这个表示描述 $D_{2n}$（而不是总回到原始的几何描述）会大大简化对这些群的处理。

表示给出了描述许多群的一种容易方法，但也有若干微妙之处需要考虑。其中之一是：在任意表示中，可能很难（甚至不可能）判断两个用给定生成元表达的群元素是否相等。因此，所表示的群的阶是多少，甚至该群是有限还是无限，可能都不明显。例如，可以证明
\[
  \langle x_1,y_1\mid x_1^2=y_1^2=(x_1y_1)^2=1\rangle
\]
是一个阶为 $4$ 的群的表示，而
\[
  \langle x_2,y_2\mid x_2^3=y_2^3=(x_2y_2)^3=1\rangle
\]
则是一个无限群的表示（参见练习）。

另一个微妙之处是，即使在相当简单的表示中，也可能因为关系以某种不明显的方式交织在一起而发生某些“塌缩”；也就是说，可能存在没有在表示中显式给出、而是指定关系的后果的“隐藏”或隐含关系。一般而言，这种塌缩使得甚至很难确定所表示群大小的下界。例如，假设模仿 $D_{2n}$ 的表示来试图创造另一个群，定义
\begin{equation}
  X_{2n}=\langle x,y\mid x^n=y^2=1,\ xy=yx^2\rangle. \tag{1.2}
\end{equation}
“交换”关系 $xy=yx^2$ 决定了如何交换 $y$ 和 $x$（即如何把 $y$ 从 $x$ 的右侧“移到”左侧），所以与 $D_{2n}$ 中一样，这个群的每个元素都可写成 $y^kx^i$ 的形式，所有 $y$ 的幂在左，所有 $x$ 的幂在右。此外，由前两个关系，$x$ 和 $y$ 的任意幂都可化简，使 $i$ 位于 $0$ 与 $n-1$ 之间，并且 $k$ 为 $0$ 或 $1$。因此人们也许会猜测 $X_{2n}$ 仍是一个阶为 $2n$ 的群。但事实并非如此，因为在这个群中存在一个“隐藏”关系，它由关系 $x=xy^2$（因为 $y^2=1$）通过反复应用交换关系和结合律把 $y$ 移到左侧而得到：
\begin{align*}
  x&=x1=(xy)y=(yx^2)y=(yx)(xy)=(yx)(yx^2)\
   &=y(xy)x^2=y(yx^2)x^2=1x^4=x^4.
\end{align*}
由于 $x^4=x$，由消去律得 $x^3=1$ 于 $X_{2n}$ 中；由上面的讨论可知，对任意 $n$，$X_{2n}$ 的阶至多为 $6$。根据 $n$ 的取值，还可能发生更多塌缩（见练习）。

再例如，考虑表示
\begin{equation}
  Y=\langle u,v\mid u^4=v^3=1,\ uv=v^2u^2\rangle. \tag{1.3}
\end{equation}
在这个情形中，容易猜测 $Y$ 是阶为 $12$ 的群，但仍有额外的隐含关系。事实上，这个群 $Y$ 退化为阶为 $1$ 的平凡群；也就是说，$u$ 和 $v$ 满足额外关系 $u=1$ 与 $v=1$（证明概要见练习）。

这种塌缩不会发生在 $D_{2n}$ 的表示中，因为我们已经用独立的（几何）方法说明存在一个阶为 $2n$、由 $r$ 与 $s$ 生成并满足 (1) 中关系的群。因此，只具有这些关系的群的阶至少为 $2n$。另一方面，容易看出（使用上面对 $X_{2n}$ 的同类论证以及交换关系 $rs=sr^{-1}$），任意由 (1) 中的生成元与关系定义的群，其阶至多为 $2n$。于是，具有表示 (1) 的群阶恰为 $2n$，并且这个群确实就是正 $n$ 边形的对称群。

我们对表示 (1) 所拥有的额外信息，是存在一个已知阶的群满足这些信息。相比之下，对满足 (2) 或 (3) 中关系的任何群，我们没有独立的知识。没有这种独立的“下界”信息，我们甚至可能无法判断给定表示是否只是描述了平凡群，如 (3) 所示。

虽然一般来说，用表示来规定群时必须极其小心，但对已知群使用表示是一种强有力的概念工具和计算工具。关于表示的更多结果，包括更复杂的例子，见第 6.3 节。

\subsection*{练习}

在这些练习中，$D_{2n}$ 取通常的表示
\[
  D_{2n}=\langle r,s\mid r^n=s^2=1,\ rs=sr^{-1}\rangle.
\]
\begin{enumerate}[label=\arabic*.]
  \item 计算下列群中每个元素的阶：
  \begin{enumerate}[label=(\alph*)]
    \item $D_6$；\quad (b) $D_8$；\quad (c) $D_{10}$。
  \end{enumerate}
  \item 用上面的生成元与关系证明：若 $x$ 是 $D_{2n}$ 中任一不是 $r$ 的幂的元素，则 $rx=xr^{-1}$。
  \item 用上面的生成元与关系证明：$D_{2n}$ 中任一不是 $r$ 的幂的元素都有阶 $2$。推出 $D_{2n}$ 由两个元素 $s$ 与 $sr$ 生成，且这两个元素均有阶 $2$。
  \item 若 $n=2k$ 为偶数且 $n\geq4$，证明 $z=r^k$ 是一个阶为 $2$ 且与 $D_{2n}$ 的所有元素交换的元素。还要证明 $z$ 是 $D_{2n}$ 中唯一与所有元素交换的非单位元素。［参见第 1 节练习 33。］
  \item 若 $n$ 为奇数且 $n\geq3$，证明单位元是 $D_{2n}$ 中唯一与所有元素交换的元素。［参见第 1 节练习 33。］
  \item 设 $x$ 与 $y$ 是任意群 $G$ 中两个阶为 $2$ 的元素。证明若 $t=xy$，则 $tx=xt^{-1}$（所以若 $n=|xy|<\infty$，则 $x,t$ 在 $G$ 中满足与 $D_{2n}$ 中 $s,r$ 相同的关系）。
  \item 证明
  \[
    \langle a,b\mid a^2=b^2=(ab)^n=1\rangle
  \]
  给出了 $D_{2n}$ 关于练习 3 中得到的两个阶为 $2$ 的生成元 $a=s$ 与 $b=sr$ 的表示。［证明 $r,s$ 的关系可由 $a,b$ 的关系推出；反过来，$a,b$ 的关系也可由 $r,s$ 的关系推出。］
  \item 求 $D_{2n}$ 中由 $r$ 生成的循环子群的阶（参见第 1 节练习 27）。
\end{enumerate}

在练习 9 到 13 中，你可以仿照求 $D_{2n}$ 阶的证明来求给定柏拉图立体在 $\RR^3$ 中的刚体运动群（也称旋转群）的阶：求一对相邻顶点可被送到的位置数。或者，也可以先求一个给定面可被送到的位置数，再在固定一个面后求该面上一顶点可被送到的位置数。
\begin{enumerate}[label=\arabic*.,resume]
  \item 设 $G$ 是正四面体在 $\RR^3$ 中的刚体运动群。证明 $|G|=12$。
  \item 设 $G$ 是立方体在 $\RR^3$ 中的刚体运动群。证明 $|G|=24$。
  \item 设 $G$ 是正八面体在 $\RR^3$ 中的刚体运动群。证明 $|G|=24$。
  \item 设 $G$ 是正十二面体在 $\RR^3$ 中的刚体运动群。证明 $|G|=60$。
  \item 设 $G$ 是正二十面体在 $\RR^3$ 中的刚体运动群。证明 $|G|=60$。
  \item 找出 $\ZZ$ 的一个生成元集。
  \item 找出 $\ZZ/n\ZZ$ 的一组生成元与关系。
  \item 证明群 $\langle x_1,y_1\mid x_1^2=y_1^2=(x_1y_1)^2=1\rangle$ 是二面体群 $D_4$（其中可把 $x_1$ 改写为字母 $r$，把 $y_1$ 改写为 $s$）。［证明最后一个关系等价于 $x_1y_1=y_1x_1^{-1}$。］
  \item 设 $X_{2n}$ 是由 (1.2) 中表示给出的群。
  \begin{enumerate}[label=(\alph*)]
    \item 证明若 $n=3k$，则 $X_{2n}$ 的阶为 $6$；当把 $x$ 换成 $r$、$y$ 换成 $s$ 时，它与 $D_6$ 有相同的生成元与关系。
    \item 证明若 $(3,n)=1$，则 $x$ 满足额外关系 $x=1$。在此情形下推出 $X_{2n}$ 的阶为 $2$。［用 $x^n=1$ 与 $x^3=1$。］
  \end{enumerate}
  \item 设 $Y$ 是由 (1.3) 中表示给出的群。
  \begin{enumerate}[label=(\alph*)]
    \item 证明 $v^2=v^{-1}$。［用关系 $v^3=1$。］
    \item 证明 $v$ 与 $u^3$ 交换。［把左边写为 $(v^2u^2)(uv)$ 并用关系把它化为右边，从而证明 $v^2u^3v=u^3$；再用 (a)。］
    \item 证明 $v$ 与 $u$ 交换。［证明 $u^9=u$，然后用 (b)。］
    \item 证明 $uv=1$。［用 (c) 与最后一个关系。］
    \item 证明 $u=1$，推出 $v=1$，并得出 $Y=1$。［用 (d) 与等式 $u^4v^3=1$。］
  \end{enumerate}
\end{enumerate}


\section{对称群}

令 $\Omega$ 为任一非空集合，令 $S_\Omega$ 为从 $\Omega$ 到自身的所有双射组成的集合（即 $\Omega$ 的所有置换组成的集合）。集合 $S_\Omega$ 在函数复合 $\circ$ 下是群。注意，$\circ$ 是 $S_\Omega$ 上的二元运算，因为若 $\sigma:\Omega\to\Omega$ 和 $\tau:\Omega\to\Omega$ 都是双射，则 $\sigma\circ\tau$ 也是从 $\Omega$ 到 $\Omega$ 的双射。由于函数复合一般是结合的，$\circ$ 是结合的。$S_\Omega$ 的单位元是置换 $1$，定义为对所有 $a\in\Omega$，$1(a)=a$。每个置换 $\sigma$ 都有一个（双侧）逆函数 $\sigma^{-1}:\Omega\to\Omega$，满足 $\sigma\circ\sigma^{-1}=\sigma^{-1}\circ\sigma=1$。因此，$(S_\Omega,\circ)$ 满足所有群公理。这个群称为集合 $\Omega$ 上的\emph{对称群}。重要的是要认识到，$S_\Omega$ 的元素是 $\Omega$ 的置换，而不是 $\Omega$ 本身的元素。

在特殊情形 $\Omega=\{1,2,3,\ldots,n\}$ 下，$\Omega$ 上的对称群记为 $S_n$，称为\emph{$n$ 次对称群}。\footnote{我们将在第 6 节看到，$S_\Omega$ 的结构只依赖于 $\Omega$ 的基数，而不依赖于 $\Omega$ 的具体元素；因此若 $\Omega$ 是任一含 $n$ 个元素的有限集，则 $S_\Omega$ “看起来像” $S_n$。}群 $S_n$ 将在全文中发挥重要作用：它本身是一个相当重要的研究对象，同时也是说明和激发一般理论的工具。

首先说明 $S_n$ 的阶为 $n!$。集合 $\{1,2,3,\ldots,n\}$ 的置换恰好就是该集合到自身的单射，因为该集合有限（命题 0.1）；于是我们可以数出单射的个数。单射 $\sigma$ 可把数字 $1$ 送到 $\{1,2,3,\ldots,n\}$ 中任意一个元素；$\sigma(2)$ 随后可为该集合中除 $\sigma(1)$ 外的任一元素（所以 $\sigma(2)$ 有 $n-1$ 种选择）；$\sigma(3)$ 可为除 $\sigma(1)$ 和 $\sigma(2)$ 外的任一元素（所以有 $n-2$ 种选择）；依此类推。因此，从 $\{1,2,3,\ldots,n\}$ 到自身的单射恰有
\[
  n(n-1)(n-2)\cdots 2\cdot1=n!
\]
个。因此 $\{1,2,3,\ldots,n\}$ 恰有 $n!$ 个置换，也就是 $S_n$ 中恰有 $n!$ 个元素。

现在描述一种写出 $S_n$ 的元素 $\sigma$ 的有效记号，全文都将使用它，这称为\emph{循环分解}。

一个\emph{循环}是一串整数，用来表示 $S_n$ 中把这些整数循环置换（并固定所有其他整数）的元素。循环 $(a_1\ a_2\ \cdots\ a_m)$ 是这样一个置换：把 $a_i$ 送到 $a_{i+1}$，$1\leq i\leq m-1$，并把 $a_m$ 送到 $a_1$。例如，$(2\ 1\ 3)$ 是把 $2$ 映到 $1$、把 $1$ 映到 $3$、把 $3$ 映到 $2$ 的置换。一般地，对每个 $\sigma\in S_n$，数字 $1$ 到 $n$ 会被重新排列并分组成 $k$ 个循环，形如
\[
  (a_1\ a_2\ \cdots\ a_{m_1})
  (a_{m_1+1}\ \cdots\ a_{m_2})\cdots
  (a_{m_{k-1}+1}\ \cdots\ a_{m_k}),
\]
其中 $m_k=n$；从这个表达式可以很容易读出 $\sigma$ 对 $1$ 到 $n$ 中任一数字的作用。对任意 $x\in\{1,2,3,\ldots,n\}$，先在上述表达式中找到 $x$。若 $x$ 后面不是紧跟一个右括号（即 $x$ 不是某个循环的右端），则 $\sigma(x)$ 是紧挨着 $x$ 右侧的整数。若 $x$ 后面紧跟右括号，则 $\sigma(x)$ 是以 $x$ 结尾的那个循环开头的数（即若 $x=a_{m_i}$，则 $\sigma(x)=a_{m_{i-1}+1}$，其中 $m_0=0$）。所有循环的乘积称为 $\sigma$ 的\emph{循环分解}。

下面给出求 $S_n$ 中元素 $\sigma$ 的循环分解的算法，并通过一个具体置换演示该算法。算法的证明以及循环分解唯一性的完整分析推迟到第 4 章。

令 $n=13$，定义 $\sigma\in S_{13}$ 为
\begin{align*}
&\sigma(1)=12,\quad \sigma(2)=13,\quad \sigma(3)=3,\quad \sigma(4)=1,\quad \sigma(5)=11,\\
&\sigma(6)=9,\quad \sigma(7)=5,\quad \sigma(8)=10,\quad \sigma(9)=6,\quad \sigma(10)=4,\\
&\sigma(11)=7,\quad \sigma(12)=8,\quad \sigma(13)=2.
\end{align*}

\paragraph{循环分解算法}
\begin{center}
\begin{tabular}{p{0.56\linewidth}|p{0.34\linewidth}}
\textbf{方法} & \textbf{例子} \\ \hline
要开始一个新循环，选取 $\{1,2,\ldots,n\}$ 中尚未在前面循环中出现的最小元素，称之为 $a$（若刚开始，则 $a=1$）；开始新的循环：$(a$
& $(1$ \\[0.8em]
从 $\sigma$ 的给定描述中读出 $\sigma(a)$，称之为 $b$。若 $b=a$，则用右括号结束该循环（不写下 $b$）；这完成一个循环，回到第 1 步。若 $b\neq a$，则把 $b$ 写在该循环中 $a$ 的旁边：$(a\ b$
& $\sigma(1)=12=b$，$12\neq1$，所以写成：$(1\ 12$ \\[0.8em]
从 $\sigma$ 的给定描述中读出 $\sigma(b)$，称之为 $c$。若 $c=a$，则用右括号结束该循环，从而完成该循环，并回到第 1 步。若 $c\neq a$，则把 $c$ 写在该循环中 $b$ 的旁边：$(a\ b\ c$。以数字 $c$ 作为新的 $b$，重复这一步，直到循环闭合。
& $\sigma(12)=8$，$8\neq1$，所以继续该循环为：$(1\ 12\ 8$ 
\end{tabular}
\end{center}
自然地，当 $\{1,2,\ldots,n\}$ 中所有数都已出现在某个循环中时，这个过程停止。对该例中的特定 $\sigma$，这给出
\[
  \sigma=(1\ 12\ 8\ 10\ 4)(2\ 13)(3)(5\ 11\ 7)(6\ 9).
\]
一个循环的\emph{长度}是出现在其中的整数个数。长度为 $t$ 的循环称为一个 $t$-循环。若两个循环没有公共数字，则称它们\emph{不交}。因此，上面的元素 $\sigma$ 是 $5$ 个两两不交循环的乘积：一个 $5$-循环、一个 $2$-循环、一个 $1$-循环、一个 $3$-循环和另一个 $2$-循环。

今后我们约定不写出 $1$-循环。因此，若某个整数 $i$ 没有出现在置换 $\tau$ 的循环分解中，则默认 $\tau(i)=i$，即 $\tau$ 固定 $i$。$S_n$ 的恒等置换的循环分解是 $(1)(2)\cdots(n)$，并将简写为 $1$。所以算法的最后一步是：
\[
\boxed{\text{最后一步：删去所有长度为 }1\text{ 的循环}}
\]
所以该例中 $\sigma$ 的循环分解最终写为
\[
  \sigma=(1\ 12\ 8\ 10\ 4)(2\ 13)(5\ 11\ 7)(6\ 9).
\]
这个约定的优点是，$S_n$ 中元素 $\tau$ 的循环分解，同时也是 $S_m$（$m\geq n$）中如下置换的循环分解：它在 $\{1,2,3,\ldots,n\}$ 上按 $\tau$ 作用，并固定 $\{n+1,n+2,\ldots,m\}$ 中每个元素。因此，例如 $(1\ 2)$ 表示交换 $1$ 和 $2$ 并固定所有更大整数的置换，无论把它看作 $S_2$、$S_3$、$S_4$ 等中的元素。

再例如，$S_3$ 的 $6$ 个元素具有如下循环分解：
\[
\begin{array}{c|c|c}
 & \text{置换的取值} & \text{循环分解}\\ \hline
\sigma_1&1\mapsto1,\ 2\mapsto2,\ 3\mapsto3&1\\
\sigma_2&1\mapsto1,\ 2\mapsto3,\ 3\mapsto2&(2\ 3)\\
\sigma_3&1\mapsto3,\ 2\mapsto2,\ 3\mapsto1&(1\ 3)\\
\sigma_4&1\mapsto2,\ 2\mapsto1,\ 3\mapsto3&(1\ 2)\\
\sigma_5&1\mapsto2,\ 2\mapsto3,\ 3\mapsto1&(1\ 2\ 3)\\
\sigma_6&1\mapsto3,\ 2\mapsto1,\ 3\mapsto2&(1\ 3\ 2)
\end{array}
\]
对任意 $\sigma\in S_n$，$\sigma^{-1}$ 的循环分解可通过把 $\sigma$ 的循环分解中每个循环里的数字按相反顺序写出得到。例如，若 $\sigma=(1\ 12\ 8\ 10\ 4)(2\ 13)(5\ 11\ 7)(6\ 9)$，则
\[
  \sigma^{-1}=(4\ 10\ 8\ 12\ 1)(13\ 2)(7\ 11\ 5)(9\ 6).
\]

在 $S_n$ 中计算乘积是直接的，只要记住在计算 $\sigma\circ\tau$ 时从右向左读置换。只需“跟踪”元素在连续置换下的去向。例如，在乘积
\[
  (1\ 2\ 3)\circ(1\ 2)(3\ 4)
\]
中，数字 $1$ 先被第一个置换送到 $2$，然后 $2$ 被第二个置换送到 $3$，所以复合映射把 $1$ 送到 $3$。为求乘积的循环分解，接下来要看 $3$ 的去向：$3$ 先被送到 $4$，而后 $4$ 被固定，所以 $3$ 被复合映射送到 $4$。类似地，$4$ 先被送到 $3$，再被送到 $1$，于是乘积中的这个循环闭合为 $(1\ 3\ 4)$。最后，$2$ 被送到 $1$，再被送到 $2$，所以 $2$ 被这个乘积固定。因此
\[
  (1\ 2\ 3)\circ(1\ 2)(3\ 4)=(1\ 3\ 4).
\]
作为更多例子，
\[
  (1\ 2)\circ(1\ 3)=(1\ 3\ 2),\qquad (1\ 3)\circ(1\ 2)=(1\ 2\ 3).
\]
这特别说明，对所有 $n\geq3$，$S_n$ 都是非阿贝尔群。

循环分解中的每个循环 $(a_1\ a_2\ \cdots\ a_m)$ 可看成把 $a_1,a_2,\ldots,a_m$ 循环置换并固定所有其他整数的置换。由于不交循环置换的数字属于彼此不交的集合，所以
\[
  \text{不交循环彼此交换。}
\]
因此，任意不交循环乘积（特别是循环分解）中循环的排列顺序不会改变该置换。

此外，由于一个给定循环 $(a_1\ a_2\ \cdots\ a_m)$ 循环置换集合 $\{a_1,a_2,\ldots,a_m\}$，所以循环中数字本身可以循环移位而不改变置换，即
\[
  (a_1\ a_2\ \cdots\ a_m)=(a_2\ a_3\ \cdots\ a_m\ a_1)=(a_3\ a_4\ \cdots\ a_m\ a_1\ a_2)=\cdots=(a_m\ a_1\ a_2\ \cdots\ a_{m-1}).
\]
因此，例如 $(1\ 2)=(2\ 1)$，且 $(1\ 2\ 3\ 4)=(3\ 4\ 1\ 2)$。按惯例，通常把循环中出现的最小数字写在最前面。

使用循环时必须稍加小心，因为一个置换可用许多方式写成循环的任意乘积。例如在 $S_3$ 中，
\[
  (1\ 2\ 3)=(1\ 2)(2\ 3)=(1\ 3)(1\ 3\ 2)(1\ 3)
\]
等等。但是（我们将会证明）每个置换的循环分解是把它表示成不交循环之积的唯一方式（只差循环的重新排序和每个循环内部数字的循环移位）。把任意循环乘积化为不交循环之积，使我们能够一眼判断两个置换是否相同。这种记号的另一个优点是：证明一个置换的阶等于其循环分解中各循环长度的最小公倍数，是一个练习（下面给出提示）。

\subsection*{练习}
\begin{enumerate}[label=\arabic*.]
  \item 设 $\sigma$ 是置换
  \[
  1\mapsto3,\\ 2\mapsto4,\\ 3\mapsto5,\\ 4\mapsto2,\\ 5\mapsto1,
  \]
  并设 $\tau$ 是置换
  \[
  1\mapsto5,\\ 2\mapsto3,\\ 3\mapsto2,\\ 4\mapsto4,\\ 5\mapsto1.
  \]
  求下列置换的循环分解：$\sigma,\tau,\sigma^2,\sigma\tau,\tau\sigma,\tau^2\sigma$。
  \item 设 $\sigma$ 是置换
  \[
  \begin{array}{ccccc}
  1\mapsto13&2\mapsto2&3\mapsto15&4\mapsto14&5\mapsto10\\
  6\mapsto6&7\mapsto12&8\mapsto3&9\mapsto4&10\mapsto1\\
  11\mapsto7&12\mapsto9&13\mapsto5&14\mapsto11&15\mapsto8
  \end{array}
  \]
  并设 $\tau$ 是置换
  \[
  \begin{array}{ccccc}
  1\mapsto14&2\mapsto9&3\mapsto10&4\mapsto2&5\mapsto12\\
  6\mapsto6&7\mapsto5&8\mapsto11&9\mapsto15&10\mapsto3\\
  11\mapsto8&12\mapsto7&13\mapsto4&14\mapsto1&15\mapsto13.
  \end{array}
  \]
  求下列置换的循环分解：$\sigma,\tau,\sigma^2,\sigma\tau,\tau\sigma,\tau^2\sigma$。
  \item 对前两题中求得循环分解的每个置换，计算其阶。
  \item 计算下列群中每个元素的阶：\begin{enumerate}[label=(\alph*)]
    \item $S_3$；\quad \item $S_4$。
  \end{enumerate}
  \item 求 $(1\ 12\ 8\ 10\ 4)(2\ 13)(5\ 11\ 7)(6\ 9)$ 的阶。
  \item 写出 $S_4$ 中每个阶为 $4$ 的元素的循环分解。
  \item 写出 $S_4$ 中每个阶为 $2$ 的元素的循环分解。
  \item 证明若 $\Omega=\{1,2,3,\ldots\}$，则 $S_\Omega$ 是无限群（不要说 $\infty!=\infty$）。
  \item
  \begin{enumerate}[label=(\alph*)]
    \item 令 $\sigma$ 为 $12$-循环 $(1\ 2\ 3\ 4\ 5\ 6\ 7\ 8\ 9\ 10\ 11\ 12)$。对哪些正整数 $i$，$\sigma^i$ 仍是 $12$-循环？
    \item 令 $\tau$ 为 $8$-循环 $(1\ 2\ 3\ 4\ 5\ 6\ 7\ 8)$。对哪些正整数 $i$，$\tau^i$ 仍是 $8$-循环？
    \item 令 $\omega$ 为 $14$-循环 $(1\ 2\ 3\ 4\ 5\ 6\ 7\ 8\ 9\ 10\ 11\ 12\ 13\ 14)$。对哪些正整数 $i$，$\omega^i$ 仍是 $14$-循环？
  \end{enumerate}
  \item 证明若 $\sigma$ 是 $m$-循环 $(a_1\ a_2\ \cdots\ a_m)$，则对所有 $i\in\{1,2,\ldots,m\}$，$\sigma^i(a_k)=a_{k+i}$，其中当 $k+i>m$ 时，把 $k+i$ 换成模 $m$ 的最小剩余。推出 $|\sigma|=m$。
  \item 令 $\sigma$ 为 $m$-循环 $(1\ 2\ \cdots\ m)$。证明 $\sigma^i$ 也是 $m$-循环，当且仅当 $i$ 与 $m$ 互素。
  \item
  \begin{enumerate}[label=(\alph*)]
    \item 若 $\tau=(1\ 2)(3\ 4)(5\ 6)(7\ 8)(9\ 10)$，判断是否存在某个 $n$-循环 $\sigma$（$n\geq10$）以及某个整数 $k$，使得 $\tau=\sigma^k$。
    \item 若 $\tau=(1\ 2)(3\ 4\ 5)$，判断是否存在某个 $n$-循环 $\sigma$（$n\geq5$）以及某个整数 $k$，使得 $\tau=\sigma^k$。
  \end{enumerate}
  \item 证明 $S_n$ 中一个元素阶为 $2$，当且仅当它的循环分解是一些可交换的 $2$-循环的乘积。
  \item 设 $p$ 是素数。证明 $S_n$ 中一个元素阶为 $p$，当且仅当它的循环分解是一些可交换的 $p$-循环的乘积。用一个显式例子说明若 $p$ 不是素数，则这不一定成立。
  \item 证明 $S_n$ 中元素的阶等于其循环分解中各循环长度的最小公倍数。［用练习 10 和第 1 节练习 24。］
  \item 证明若 $n\geq m$，则 $S_n$ 中 $m$-循环的个数为
  \[
    \frac{n(n-1)(n-2)\cdots(n-m+1)}{m}.
  \]
  ［数一数形成 $m$-循环的方法数，再除以一个特定 $m$-循环的表示数。］
  \item 证明若 $n\geq4$，则 $S_n$ 中可写成两个不交 $2$-循环之积的置换个数为
  \[
    \frac{n(n-1)(n-2)(n-3)}{8}.
  \]
  \item 求所有整数 $n$，使得 $S_5$ 含有阶为 $n$ 的元素。［用练习 15。］
  \item 求所有整数 $n$，使得 $S_7$ 含有阶为 $n$ 的元素。［用练习 15。］
  \item 找出 $S_3$ 的一组生成元与关系。
\end{enumerate}

\section{矩阵群}

本节引入系数来自域的矩阵群的概念。这样一族群将在第一部分中作为说明性例子使用，并会在关于向量空间的章节中更详细地研究。

域是能够进行所有算术运算 $+$、$-$、$\times$ 和 $/$（除以非零元素）的“最小”数学结构；特别地，每个非零元素都必须有乘法逆元。我们稍后会更深入地研究域；在本书这一部分中，我们遇到的域 $F$ 只会是 $\QQ$、$\RR$ 和 $\ZZ/p\ZZ$，其中 $p$ 是素数。例子 $\ZZ/p\ZZ$ 是一个有限域；为强调它是域，我们记作 $\FF_p$。为完整起见，这里给出域的精确定义。

\noindent\textbf{定义.}
\begin{enumerate}[label=(\arabic*)]
  \item \emph{域}是一个集合 $F$，连同 $F$ 上两个二元运算 $+$ 与 $\cdot$，使得 $(F,+)$ 是阿贝尔群（其单位元记为 $0$），$(F-\{0\},\cdot)$ 也是阿贝尔群，并且满足分配律
  \[
    a\cdot(b+c)=(a\cdot b)+(a\cdot c),\qquad\text{对所有 }a,b,c\in F.
  \]
  \item 对任意域 $F$，令 $F^\times=F-\{0\}$。
\end{enumerate}

当标量来自 $\RR$ 时，向量空间理论、矩阵与线性变换理论以及行列式理论中的所有结论，在作适当修改后，当标量来自任意域 $F$ 时也都成立。当我们在第一部分中使用这些理论时，会明确说明我们假设了哪些关于域的事实。

对每个 $n\in\ZZ^+$，令 $\GL_n(F)$ 为所有元素来自 $F$ 且行列式非零的 $n\times n$ 矩阵的集合，即
\[
  \GL_n(F)=\{A\mid A\text{ 是元素来自 }F\text{ 的 }n\times n\text{ 矩阵且 }\det(A)\neq0\}.
\]
任意元素来自 $F$ 的矩阵 $A$ 的行列式可用 $F=\RR$ 时的同样公式计算。对任意 $n\times n$ 矩阵 $A$ 与 $B$，令 $AB$ 为按 $F=\RR$ 时同样规则计算的矩阵积。这个乘法是结合的。此外，由 $\det(AB)=\det(A)\det(B)$，若 $\det(A)\neq0$ 且 $\det(B)\neq0$，则 $\det(AB)\neq0$，所以 $\GL_n(F)$ 对矩阵乘法封闭。进一步地，$\det(A)\neq0$ 当且仅当 $A$ 有矩阵逆（并且这个逆可由 $F=\RR$ 时使用的同样伴随矩阵公式计算），所以每个 $A\in\GL_n(F)$ 都有逆 $A^{-1}\in\GL_n(F)$：
\[
  AA^{-1}=A^{-1}A=I,
\]
其中 $I$ 是 $n\times n$ 单位矩阵。因此，$\GL_n(F)$ 在矩阵乘法下是群，称为\emph{$n$ 次一般线性群}。

以下结果将在第三部分中证明，这里为方便先记录：
\begin{enumerate}[label=(\arabic*)]
  \item 若 $F$ 是域且 $|F|<\infty$，则 $|F|=p^m$，其中 $p$ 为某个素数，$m$ 为某个整数；
  \item 若 $|F|=q<\infty$，则
  \[
    |\GL_n(F)|=(q^n-1)(q^n-q)(q^n-q^2)\cdots(q^n-q^{n-1}).
  \]
\end{enumerate}

\subsection*{练习}

设 $F$ 是域，$n\in\ZZ^+$。
\begin{enumerate}[label=\arabic*.]
  \item 证明 $|\GL_2(\FF_2)|=6$。
  \item 写出 $\GL_2(\FF_2)$ 的所有元素，并计算每个元素的阶。
  \item 证明 $\GL_2(\FF_2)$ 是非阿贝尔的。
  \item 证明若 $n$ 不是素数，则 $\ZZ/n\ZZ$ 不是域。
  \item 证明 $\GL_n(F)$ 是有限群，当且仅当 $F$ 含有有限多个元素。
  \item 若 $|F|=q$ 有限，证明 $|\GL_n(F)|<q^{n^2}$。
  \item 设 $p$ 是素数。证明 $\GL_2(\FF_p)$ 的阶为 $p^4-p^3-p^2+p$（不要直接引用本节中的阶公式）。［从 $\FF_p$ 上所有 $2\times2$ 矩阵总数中减去不可逆矩阵的个数。可以使用这样一个事实：$2\times2$ 矩阵不可逆，当且仅当一行是另一行的倍数。］
  \item 证明对任意 $n\geq2$ 和任意 $F$，$\GL_n(F)$ 都是非阿贝尔的。
  \item 证明实数项 $2\times2$ 矩阵的矩阵乘法是结合的。
  \item 令
  \[
    G=\left\{\begin{pmatrix}a&b\\0&c\end{pmatrix}\ \middle|\ a,b,c\in\RR,\ a\neq0,\ c\neq0\right\}.
  \]
  \begin{enumerate}[label=(\alph*)]
    \item 计算
    \[
      \begin{pmatrix}a_1&b_1\\0&c_1\end{pmatrix}
      \begin{pmatrix}a_2&b_2\\0&c_2\end{pmatrix}
    \]
    以说明 $G$ 对矩阵乘法封闭。
    \item 求矩阵 $\begin{pmatrix}a&b\\0&c\end{pmatrix}$ 的逆，并推出 $G$ 对取逆封闭。
    \item 推出 $G$ 是 $\GL_2(\RR)$ 的子群（参见第 1 节练习 26）。
    \item 证明 $G$ 中两个对角元相等（即 $a=c$）的元素组成的集合也是 $\GL_2(\RR)$ 的子群。
  \end{enumerate}
\end{enumerate}

下一个练习引入域 $F$ 上的海森堡群，并发展它的一些基本性质。当 $F=\RR$ 时，这个群在量子力学和信号理论中起重要作用，它给出了海森堡不确定性原理的一个群论解释（归功于 H. Weyl）。还要注意，海森堡群可以更一般地定义——例如，元素可取在 $\ZZ$ 中。
\begin{enumerate}[label=\arabic*.,resume]
  \item 令
  \[
    H(F)=\left\{\begin{pmatrix}1&a&b\\0&1&c\\0&0&1\end{pmatrix}\ \middle|\ a,b,c\in F\right\},
  \]
  称为 $F$ 上的\emph{海森堡群}。令
  \[
    X=\begin{pmatrix}1&a&b\\0&1&c\\0&0&1\end{pmatrix},\qquad
    Y=\begin{pmatrix}1&d&e\\0&1&f\\0&0&1\end{pmatrix}
  \]
  为 $H(F)$ 的元素。
  \begin{enumerate}[label=(\alph*)]
    \item 计算矩阵积 $XY$，并推出 $H(F)$ 对矩阵乘法封闭。给出显式矩阵，使得 $XY\neq YX$（从而 $H(F)$ 总是非阿贝尔的）。
    \item 求矩阵逆 $X^{-1}$ 的显式公式，并推出 $H(F)$ 对取逆封闭。
    \item 证明 $H(F)$ 的结合律，并推出 $H(F)$ 是阶为 $|F|^3$ 的群。（不要假设矩阵乘法结合。）
    \item 求有限群 $H(\ZZ/2\ZZ)$ 中每个元素的阶。
    \item 证明群 $H(\RR)$ 的每个非单位元素都是无限阶的。
  \end{enumerate}
\end{enumerate}

\section{四元数群}

\emph{四元数群} $Q_8$ 定义为
\[
  Q_8=\{1,-1,i,-i,j,-j,k,-k\},
\]
乘法按如下规则计算：
\begin{align*}
  1\cdot a&=a\cdot1=a, &&\text{对所有 }a\in Q_8,\\
  (-1)\cdot(-1)&=1,\qquad (-1)\cdot a=a\cdot(-1)=-a, &&\text{对所有 }a\in Q_8,\\
  i\cdot i&=j\cdot j=k\cdot k=-1,\\
  i\cdot j&=k, & j\cdot i&=-k,\\
  j\cdot k&=i, & k\cdot j&=-i,\\
  k\cdot i&=j, & i\cdot k&=-j.
\end{align*}
照常，今后把 $a\cdot b$ 写作 $ab$。检验结合律相当繁琐（我们稍后会用较少计算的方法证明它），但其他公理容易检验。注意，$Q_8$ 是阶为 $8$ 的非阿贝尔群。

\subsection*{练习}
\begin{enumerate}[label=\arabic*.]
  \item 计算 $Q_8$ 中每个元素的阶。
  \item 写出 $S_3$、$D_8$ 和 $Q_8$ 的群表。
  \item 找出 $Q_8$ 的一组生成元与关系。
\end{enumerate}


\section{同态与同构}

本节精确说明两个群何时“看起来相同”，也就是说，何时具有完全相同的群论结构。这就是两个群之间同构的概念。我们先定义同态的概念，稍后还会对此作大量讨论。

\noindent\textbf{定义.} 设 $(G,*)$ 和 $(H,\circ)$ 是群。若映射 $\varphi:G\to H$ 满足
\[
  \varphi(x*y)=\varphi(x)\circ\varphi(y),\qquad\text{对所有 }x,y\in G,
\]
则称 $\varphi$ 是一个\emph{同态}。

当 $G$ 和 $H$ 的群运算没有显式写出时，同态条件就简写为
\[
  \varphi(xy)=\varphi(x)\varphi(y),
\]
但必须牢记，左边的乘积 $xy$ 是在 $G$ 中计算的，而右边的乘积 $\varphi(x)\varphi(y)$ 是在 $H$ 中计算的。直观地说，若映射 $\varphi$ 保持其定义域与陪域的群结构，则它是同态。

\noindent\textbf{定义.} 映射 $\varphi:G\to H$ 称为一个\emph{同构}，并且称 $G$ 与 $H$ \emph{同构}或具有相同的\emph{同构型}，记作 $G\cong H$，若
\begin{enumerate}[label=(\arabic*)]
  \item $\varphi$ 是同态（即 $\varphi(xy)=\varphi(x)\varphi(y)$），并且
  \item $\varphi$ 是双射。
\end{enumerate}

换言之，群 $G$ 与 $H$ 同构，是指存在一个保持群运算的双射从 $G$ 到 $H$。直观地说，除元素名称和运算符号可能写得不同外，$G$ 与 $H$ 是同一个群。因此，$G$ 所具有的任何只依赖于 $G$ 的群结构的性质（也就是可由群公理推出的性质——例如群的交换性）也在 $H$ 中成立。注意，这从形式上说明了为什么可以把所有群运算都写作 $\cdot$：改变运算符号并不改变同构型。

\paragraph{例}
\begin{enumerate}[label=(\arabic*)]
  \item 对任意群 $G$，$G\cong G$。恒等映射给出了一个明显的同构，但一般并不是从 $G$ 到自身的唯一同构。更一般地，令 $\mathcal{G}$ 为任意非空群族。容易检验关系 $\cong$ 是 $\mathcal{G}$ 上的等价关系，其等价类称为\emph{同构类}。这也解释了“同构”定义中略显对称的措辞。
  \item 指数映射 $\exp:\RR\to\RR^+$，定义为 $\exp(x)=e^x$（其中 $e$ 是自然对数的底），是从 $(\RR,+)$ 到 $(\RR^+,\times)$ 的同构。$\exp$ 是双射，因为它有逆函数（即 $\log_e$）；且 $\exp$ 保持群运算，因为 $e^{x+y}=e^xe^y$。在这个例子中，元素和运算都不同，但两个群同构；也就是说，作为群它们有相同的结构。
  \item 本例说明，对称群的同构型只依赖于被置换的底层集合的基数。令 $\Delta$ 与 $\Omega$ 是非空集合。若 $|\Delta|=|\Omega|$，则对称群 $S_\Delta$ 与 $S_\Omega$ 同构。直观上，既然 $|\Delta|=|\Omega|$，存在一个从 $\Delta$ 到 $\Omega$ 的双射 $\theta$。把 $\Delta$ 与 $\Omega$ 的元素通过 $\theta$ 粘在一起，即每个 $x\in\Delta$ 与 $\theta(x)\in\Omega$ 粘在一起。为了得到映射 $\Phi:S_\Delta\to S_\Omega$，令 $\sigma\in S_\Delta$ 为 $\Delta$ 的一个置换，并令 $\Phi(\sigma)$ 为 $\Omega$ 的一个置换，它以 $\sigma$ 移动相应被粘合的 $\Delta$ 元素的方式来移动 $\Omega$ 的元素；也就是说，若对某些 $x,y\in\Delta$ 有 $\sigma(x)=y$，则在 $\Omega$ 中令 $\Phi(\sigma)(\theta(x))=\theta(y)$。由于集合双射 $\theta$ 有逆，容易检验对称群之间的这个映射也有逆。该映射的精确定义及其作为同构所需性质的直接（虽有些繁琐）验证留给后面的练习。

  反过来，若 $S_\Delta\cong S_\Omega$，则 $|\Delta|=|\Omega|$；我们只在底层集合有限时证明这一点（当 $\Delta$ 与 $\Omega$ 都是无限集时，证明更难，将在第 4 章作为练习给出）。任意两个群 $G$ 与 $H$ 的同构，首先就是二者之间的双射，因此同构的必要条件是 $|S_\Delta|=|S_\Omega|$。当 $\Delta$ 是阶为 $n$ 的有限集时，$|S_\Delta|=n!$；虽然我们只为 $S_n$ 证明过这一点，但同样的论证适用于 $S_\Delta$。类似地，若 $\Omega$ 是阶为 $m$ 的有限集，则 $|S_\Omega|=m!$。因此若 $S_\Delta$ 与 $S_\Omega$ 同构，则 $n!=m!$，所以 $m=n$，即 $|\Delta|=|\Omega|$。
\end{enumerate}

全文还会出现更多同构的例子。当我们研究不同结构（环、域、向量空间等）时，也会为相应结构之间制定对应的同构概念。数学中的一个核心问题是确定一个结构的哪些性质能够指定它的同构型（即证明：若 $G$ 是具有某种结构的对象（例如群），且 $G$ 具有性质 $P$，则任何其他同类结构对象（群）$X$ 只要具有性质 $P$，就同构于 $G$）。这类定理称为\emph{分类定理}。例如，我们将证明
\[
  \text{任意阶为 }6\text{ 的非阿贝尔群都同构于 }S_3.
\]
（这里 $G$ 是群 $S_3$，$P$ 是性质“非阿贝尔且阶为 $6$”。）由这个分类定理，我们无需寻找这些群之间的显式映射，就得到 $D_6\cong S_3$ 以及 $\GL_2(\FF_2)\cong S_3$。注意，并非任意阶为 $6$ 的群都同构于 $S_3$。事实上，我们将证明在同构意义下，阶为 $6$ 的群恰有两个：$S_3$ 与 $\ZZ/6\ZZ$（也就是说，任意阶为 $6$ 的群都同构于这两个群之一，且 $S_3$ 不同构于 $\ZZ/6\ZZ$）。注意，结论不那么具体（有两种可能类型），但假设更容易检验（只需检验阶是否为 $6$）。后一类结果也称为分类。一般来说，判断两个群（或其他数学对象）是否同构，即使在具体例子中也是微妙而困难的——构造一个保持群运算的显式映射，或证明不存在这样的映射，除极小的情形外通常在计算上不可行；这一点在没有进一步理论而试图证明上述阶 $6$ 群分类时已经有所体现。

有时容易看出两个给定群不同构。例如，下面练习断言若 $\varphi:G\to H$ 是同构，则特别地：
\begin{enumerate}[label=(\alph*)]
  \item $|G|=|H|$；
  \item $G$ 是阿贝尔群，当且仅当 $H$ 是阿贝尔群；
  \item 对所有 $x\in G$，$|x|=|\varphi(x)|$。
\end{enumerate}
因此，$S_3$ 与 $\ZZ/6\ZZ$ 不同构（如上所述），因为一个是阿贝尔的，另一个不是。还有，$(\RR-\{0\},\times)$ 与 $(\RR,+)$ 不可能同构，因为在 $(\RR-\{0\},\times)$ 中元素 $-1$ 的阶为 $2$，而 $(\RR,+)$ 中没有阶为 $2$ 的元素，这与 (c) 矛盾。

最后，记录一个稍后（讨论自由群时）将证明的非常有用的事实，它处理由生成元与关系给出的两个群之间同态与同构的问题：

设 $G$ 是一个有限群，阶为 $n$，我们有它的一个表示，并令 $S=\{s_1,\ldots,s_m\}$ 为生成元。设 $H$ 是另一个群，且 $\{r_1,\ldots,r_m\}$ 是 $H$ 中元素。假设 $G$ 中由 $s_i$ 满足的任一关系，在把每个 $s_i$ 换成 $r_i$ 后也在 $H$ 中成立。则存在唯一同态 $\varphi:G\to H$，把 $s_i$ 映到 $r_i$。若我们有 $G$ 的一个表示，则只需检验该表示所列出的关系（因为按表示的定义，所有关系都能由表示中给出的关系推出）。若 $H$ 由元素 $\{r_1,\ldots,r_m\}$ 生成，则 $\varphi$ 是满射（任意 $r_i$ 的乘积都是相应 $s_i$ 乘积的像）。若此外 $H$ 与 $G$ 有相同的有限阶，则任意满射必为单射，因此 $\varphi$ 是同构：$G\cong H$。直观地说，我们可以把 $G$ 的生成元映到 $H$ 的任意元素上，并在 $G$ 中的关系仍被满足的条件下得到一个同态。

读者可能已经熟悉向量空间中的相应说法。设 $V$ 是一个 $n$ 维有限维向量空间，基为 $S$，$W$ 是另一个向量空间。则可以通过把 $S$ 中元素映到 $W$ 中任意向量来指定从 $V$ 到 $W$ 的一个线性变换（这里没有关系需要满足）。若 $W$ 也维数为 $n$，并且在 $W$ 中选定的这些向量张成 $W$（因而是 $W$ 的一组基），则该线性变换可逆（是向量空间同构）。

\paragraph{例}
\begin{enumerate}[label=(\arabic*)]
  \item 回忆 $D_{2n}=\langle r,s\mid r^n=s^2=1,\ sr=r^{-1}s\rangle$。假设 $H$ 是一个含有元素 $a,b$ 的群，并满足 $a^n=1$、$b^2=1$ 和 $ba=a^{-1}b$。则存在一个从 $D_{2n}$ 到 $H$ 的同态，把 $r$ 映到 $a$，把 $s$ 映到 $b$。例如，令 $k$ 是整除 $n$ 的整数且 $k\geq3$，并令
  \[
    D_{2k}=\langle r_1,s_1\mid r_1^k=s_1^2=1,\ s_1r_1=r_1^{-1}s_1\rangle.
  \]
  定义
  \[
    \varphi:D_{2n}\to D_{2k}\quad\text{由}\quad \varphi(r)=r_1,
    \quad \varphi(s)=s_1.
  \]
  若写 $n=km$，则因为 $r_1^k=1$，也有 $r_1^n=(r_1^k)^m=1$。于是 $D_{2n}$ 中由 $r,s$ 满足的三个关系也由 $D_{2k}$ 中的 $r_1,s_1$ 满足。因此 $\varphi$ 唯一地扩张为从 $D_{2n}$ 到 $D_{2k}$ 的同态。由于 $\{r_1,s_1\}$ 生成 $D_{2k}$，$\varphi$ 是满射。若 $k<n$，这个同态不是同构。
  \item 继续前例，令 $G=D_6$ 如上表示。检验在 $H=S_3$ 中，元素 $a=(1\ 2\ 3)$ 与 $b=(1\ 2)$ 满足关系：$a^3=1$、$b^2=1$ 且 $ba=a^{-1}b$。因此存在从 $D_6$ 到 $S_3$ 的同态，把 $r\mapsto a$，$s\mapsto b$。进一步可检验 $S_3$ 由 $a$ 与 $b$ 生成，所以此同态是满射。由于 $D_6$ 与 $S_3$ 的阶都为 $6$，这个同态是同构：$D_6\cong S_3$。
\end{enumerate}

注意，上述例子中的元素 $a$ 不必有阶 $n$（即 $n$ 不必是使 $a$ 的幂等于 $H$ 中单位元的最小正指数），同样，$b$ 也不必有阶 $2$（例如若 $a=a^{-1}$，则 $b$ 完全可以是单位元）。这使我们能更容易地构造同态，也符合这样的思想：群 $G$ 的生成元与关系构成了 $G$ 的群结构的一套完整数据。

\subsection*{练习}
设 $G$ 与 $H$ 是群。
\begin{enumerate}[label=\arabic*.]
  \item 设 $\varphi:G\to H$ 是同态。
  \begin{enumerate}[label=(\alph*)]
    \item 证明对所有 $n\in\ZZ^+$，$\varphi(x^n)=\varphi(x)^n$。
    \item 对 $n=-1$ 证明 (a)，并推出对所有 $n\in\ZZ$，$\varphi(x^n)=\varphi(x)^n$。
  \end{enumerate}
  \item 若 $\varphi:G\to H$ 是同构，证明对所有 $x\in G$，$|\varphi(x)|=|x|$。推出任意两个同构群对每个 $n\in\ZZ^+$ 都有相同个数的阶为 $n$ 的元素。若 $\varphi$ 仅假设为同态，结论是否仍成立？
  \item 若 $\varphi:G\to H$ 是同构，证明 $G$ 是阿贝尔群当且仅当 $H$ 是阿贝尔群。若 $\varphi:G\to H$ 是同态，还需要对 $\varphi$ 加上什么条件（若需要），才能保证若 $G$ 阿贝尔，则 $H$ 也阿贝尔？
  \item 证明乘法群 $\RR-\{0\}$ 与 $\CC-\{0\}$ 不同构。
  \item 证明加法群 $\RR$ 与 $\QQ$ 不同构。
  \item 证明加法群 $\ZZ$ 与 $\QQ$ 不同构。
  \item 证明 $D_8$ 与 $Q_8$ 不同构。
  \item 证明若 $n\neq m$，则 $S_n$ 与 $S_m$ 不同构。
  \item 证明 $D_{24}$ 与 $S_4$ 不同构。
  \item 按如下方式补全证明：若 $|\Delta|=|\Omega|$，则对称群 $S_\Delta$ 与 $S_\Omega$ 同构。令 $\theta:\Delta\to\Omega$ 为双射，定义
  \[
    \varphi:S_\Delta\to S_\Omega,
    \qquad \varphi(\sigma)=\theta\circ\sigma\circ\theta^{-1},
    \quad \sigma\in S_\Delta.
  \]
  证明：
  \begin{enumerate}[label=(\alph*)]
    \item $\varphi$ 是良定的，即若 $\sigma$ 是 $\Delta$ 的置换，则 $\theta\circ\sigma\circ\theta^{-1}$ 是 $\Omega$ 的置换；
    \item $\varphi$ 是从 $S_\Delta$ 到 $S_\Omega$ 上的双射；［找出 $\varphi$ 的双侧逆。］
    \item $\varphi$ 是同态，即 $\varphi(\sigma\circ\tau)=\varphi(\sigma)\circ\varphi(\tau)$。
  \end{enumerate}
  注意这与矩阵的换基或相似变换的相似性（稍后我们会看到它们之间的联系）。
  \item 设 $A$ 与 $B$ 是群。证明 $A\times B\cong B\times A$。
  \item 设 $A,B,C$ 是群，令 $G=A\times B$，$H=B\times C$。证明 $G\times C\cong A\times H$。
  \item 设 $G,H$ 是群，$\varphi:G\to H$ 是同态。证明 $\varphi$ 的像 $\varphi(G)$ 是 $H$ 的子群（参见第 1 节练习 26）。证明若 $\varphi$ 是单射，则 $G\cong\varphi(G)$。
  \item 设 $G,H$ 是群，$\varphi:G\to H$ 是同态。定义 $\varphi$ 的\emph{核}为 $\{g\in G\mid \varphi(g)=1_H\}$（即 $G$ 中映到 $H$ 的单位元的元素集合，也就是 $H$ 的单位元上方的纤维）。证明 $\varphi$ 的核是 $G$ 的子群（参见第 1 节练习 26）。证明 $\varphi$ 是单射当且仅当 $\varphi$ 的核是 $G$ 的单位子群。
  \item 定义映射 $\pi:\RR^2\to\RR$ 为 $\pi((x,y))=x$。证明 $\pi$ 是同态，并求 $\pi$ 的核（参见练习 14）。
  \item 设 $A$ 与 $B$ 是群，$G=A\times B$ 为它们的直积。证明由 $\pi_1((a,b))=a$ 与 $\pi_2((a,b))=b$ 定义的映射 $\pi_1:G\to A$ 和 $\pi_2:G\to B$ 是同态，并求它们的核（参见练习 14）。
  \item 设 $G$ 是任意群。证明从 $G$ 到自身的映射 $g\mapsto g^{-1}$ 是同态，当且仅当 $G$ 是阿贝尔群。
  \item 设 $G$ 是任意群。证明从 $G$ 到自身的映射 $g\mapsto g^2$ 是同态，当且仅当 $G$ 是阿贝尔群。
  \item 令 $G=\{z\in\CC\mid z^n=1\text{ 对某个 }n\in\ZZ^+\}$。证明对任意固定整数 $k>1$，由 $z\mapsto z^k$ 定义的从 $G$ 到自身的映射是满同态，但不是同构。
  \item 设 $G$ 是群，$\Aut(G)$ 为从 $G$ 到自身的所有同构组成的集合。证明 $\Aut(G)$ 在函数复合下是群（称为 $G$ 的自同构群，$\Aut(G)$ 的元素称为 $G$ 的自同构）。
  \item 证明对每个固定的非零 $k\in\QQ$，由 $q\mapsto kq$ 定义的从 $\QQ$ 到自身的映射是 $\QQ$ 的自同构（参见练习 20）。
  \item 设 $A$ 是阿贝尔群，并固定某个 $k\in\ZZ$。证明映射 $a\mapsto a^k$ 是从 $A$ 到自身的同态。若 $k=-1$，证明这个同态是同构（即 $A$ 的自同构）。
  \item 设 $G$ 是有限群，并且有一个自同构 $\alpha$（参见练习 20）满足 $\alpha(g)=g$ 当且仅当 $g=1$。若 $\alpha^2$ 是从 $G$ 到 $G$ 的恒等映射，证明 $G$ 是阿贝尔群（这样的自同构 $\alpha$ 称为阶为 $2$ 的无不动点自同构）。［证明 $G$ 的每个元素都可写成 $x^{-1}\alpha(x)$ 的形式，并对这样的表达式作用 $\alpha$。］
  \item 设 $G$ 是有限群，且 $x,y$ 是 $G$ 中两个不同的阶为 $2$ 的元素，它们生成 $G$。证明 $G\cong D_{2n}$，其中 $n=|xy|$。［见第 2 节练习 6。］
  \item 设 $n\in\ZZ^+$，$r,s$ 是 $D_{2n}$ 的通常生成元，令 $\theta=2\pi/n$。
  \begin{enumerate}[label=(\alph*)]
    \item 证明矩阵
    \[
      \begin{pmatrix}\cos\theta&-\sin\theta\\ \sin\theta&\cos\theta\end{pmatrix}
    \]
    是一个线性变换的矩阵，该线性变换把 $x,y$ 平面绕原点逆时针旋转 $\theta$ 弧度。
    \item 证明在生成元上由
    \[
      \varphi(r)=\begin{pmatrix}\cos\theta&-\sin\theta\\ \sin\theta&\cos\theta\end{pmatrix},
      \qquad
      \varphi(s)=\begin{pmatrix}0&1\\1&0\end{pmatrix}
    \]
    定义的映射 $\varphi:D_{2n}\to\GL_2(\RR)$ 可扩张为从 $D_{2n}$ 到 $\GL_2(\RR)$ 的同态。
    \item 证明 (b) 中的同态 $\varphi$ 是单射。
  \end{enumerate}
  \item 设 $i$ 和 $j$ 是第 5 节中描述的 $Q_8$ 的生成元。证明从 $Q_8$ 到 $\GL_2(\CC)$ 的映射 $\varphi$ 在生成元上由
  \[
    \varphi(i)=\begin{pmatrix}\sqrt{-1}&0\\0&-\sqrt{-1}\end{pmatrix},
    \qquad
    \varphi(j)=\begin{pmatrix}0&-1\\1&0\end{pmatrix}
  \]
  定义时，可扩张为一个同态。证明 $\varphi$ 是单射。
\end{enumerate}


\section{群作用}

本节给出群作用在集合上的精确定义，并给出一些例子。群作用将是一个强有力的工具，我们将用它来证明抽象群的定理，并揭示具体例子的结构。此外，“作用”这一概念会作为研究代数对象的一种方法反复出现：通过考察该对象如何作用在其他结构上来研究它。

\noindent\textbf{定义.} 群 $G$ 在集合 $A$ 上的一个\emph{群作用}，是一个从 $G\times A$ 到 $A$ 的映射（对所有 $g\in G$ 与 $a\in A$，写作 $g\cdot a$），满足下列性质：
\begin{enumerate}[label=(\arabic*)]
  \item 对所有 $g_1,g_2\in G$、$a\in A$，有 $g_1\cdot(g_2\cdot a)=(g_1g_2)\cdot a$；
  \item 对所有 $a\in A$，有 $1\cdot a=a$。
\end{enumerate}

我们很快会不那么形式化，并说 $G$ 是一个作用在集合 $A$ 上的群。当不致与例如群运算混淆时，表达式 $g\cdot a$ 通常简写为 $ga$（请记住，这里的 $\cdot$ 不是 $A$ 上的二元运算，并且 $ga$ 总是 $A$ 的一个元素）。注意，在性质 (1) 的左边，$g_2\cdot a$ 是 $A$ 的元素，所以用 $g_1$ 作用在它上面是有意义的；在该等式右边，乘积 $g_1g_2$ 是在 $G$ 中计算的，所得群元素再作用在集合元素 $a$ 上。

在给出群作用的例子之前，先作一些观察。设群 $G$ 作用在集合 $A$ 上。对每个固定的 $g\in G$，得到映射 $\sigma_g:A\to A$，定义为
\[
  \sigma_g(a)=g\cdot a.
\]
我们证明两个重要事实：
\begin{enumerate}[label=(\roman*)]
  \item 对每个固定的 $g\in G$，$\sigma_g$ 是 $A$ 的一个置换；
  \item 由 $g\mapsto\sigma_g$ 定义的从 $G$ 到 $S_A$ 的映射是同态。
\end{enumerate}

为说明 $\sigma_g$ 是 $A$ 的置换，我们证明它作为 $A$ 到 $A$ 的集合映射有一个双侧逆，即 $\sigma_{g^{-1}}$（于是由第 0.1 节命题 1，它是置换）。对所有 $a\in A$，
\begin{align*}
  (\sigma_{g^{-1}}\circ\sigma_g)(a)
  &=\sigma_{g^{-1}}(\sigma_g(a)) &&\text{（函数复合的定义）}\\
  &=g^{-1}\cdot(g\cdot a) &&\text{（$\sigma_{g^{-1}}$ 与 $\sigma_g$ 的定义）}\\
  &=(g^{-1}g)\cdot a &&\text{（作用的性质 (1)）}\\
  &=1\cdot a=a &&\text{（作用的性质 (2)）。}
\end{align*}
这证明 $\sigma_{g^{-1}}\circ\sigma_g$ 是 $A$ 到 $A$ 的恒等映射。由于 $g$ 是任意的，交换 $g$ 与 $g^{-1}$ 的角色可知 $\sigma_g\circ\sigma_{g^{-1}}$ 也是 $A$ 上的恒等映射。因此 $\sigma_g$ 有双侧逆，故是 $A$ 的置换。

为检验上面的断言 (ii)，定义 $\varphi:G\to S_A$ 为 $\varphi(g)=\sigma_g$。注意 (i) 说明 $\sigma_g$ 的确是 $S_A$ 的元素。为了证明 $\varphi$ 是同态，必须证明
\[
  \varphi(g_1g_2)=\varphi(g_1)\circ\varphi(g_2)
\]
（回忆 $S_A$ 在函数复合下成群）。置换 $\varphi(g_1g_2)$ 与 $\varphi(g_1)\circ\varphi(g_2)$ 相等，当且仅当它们在每个 $a\in A$ 上的取值相同。对所有 $a\in A$，
\begin{align*}
  \varphi(g_1g_2)(a)
  &=\sigma_{g_1g_2}(a) &&\text{（$\varphi$ 的定义）}\\
  &=(g_1g_2)\cdot a &&\text{（$\sigma_{g_1g_2}$ 的定义）}\\
  &=g_1\cdot(g_2\cdot a) &&\text{（作用的性质 (1)）}\\
  &=\sigma_{g_1}(\sigma_{g_2}(a)) &&\text{（$\sigma_{g_1}$ 与 $\sigma_{g_2}$ 的定义）}\\
  &=(\varphi(g_1)\circ\varphi(g_2))(a) &&\text{（$\varphi$ 的定义）。}
\end{align*}
这证明了断言 (ii)。

直观地说，群 $G$ 在集合 $A$ 上的作用，就是 $G$ 中每个元素 $g$ 都以一种与 $G$ 的群运算相容的方式作为 $A$ 上的置换起作用；上面的断言 (i) 和 (ii) 使这一点精确化。上述从 $G$ 到 $S_A$ 的同态称为给定作用所伴随的\emph{置换表示}。容易看出，这个过程可以反过来：若 $\varphi:G\to S_A$ 是从群 $G$ 到集合 $A$ 上的对称群的任意同态，则由
\[
  g\cdot a=\varphi(g)(a),\qquad g\in G,
  \quad a\in A,
\]
定义的从 $G\times A$ 到 $A$ 的映射满足 $G$ 在 $A$ 上的群作用性质。因此，群 $G$ 在集合 $A$ 上的作用与从 $G$ 到对称群 $S_A$ 的同态一一对应（也就是说，本质上是同一概念，只是用不同术语表述）。还应注意，作用的这个定义更准确地说应称为\emph{左作用}，因为群元素写在集合元素的左边。类似地也可以定义右作用。

\paragraph{例}
设 $G$ 是群，$A$ 是非空集合。下面每个例子中，对作用性质 (1) 与 (2) 的检验都留作练习。
\begin{enumerate}[label=(\arabic*)]
  \item 令 $ga=a$，对所有 $g\in G$ 与 $a\in A$。群作用的性质 (1) 和 (2) 立即成立。这个作用称为\emph{平凡作用}，并称 $G$ 在 $A$ 上\emph{平凡地作用}。注意，$G$ 的不同元素诱导 $A$ 上同一个置换（在此情形中为恒等置换）。相应的置换表示 $G\to S_A$ 是把每个 $G$ 中元素都映到恒等置换的平凡同态。

  若 $G$ 作用在集合 $B$ 上，且 $G$ 的不同元素诱导 $B$ 上不同的置换，则称该作用是\emph{忠实的}。因此，忠实作用正是相应置换表示为单射的作用。

  $G$ 在 $B$ 上作用的\emph{核}定义为
  \[
    \{g\in G\mid gb=b\text{ 对所有 }b\in B\},
  \]
  即固定 $B$ 中所有元素的 $G$ 中元素。对平凡作用，作用的核为整个 $G$；当 $|G|>1$ 时，该作用不是忠实的。

  \item 域 $F$ 上向量空间 $V$ 的公理包括：乘法群 $F^\times$ 作用在集合 $V$ 上。因此，向量空间是域的乘法群作用的熟悉例子，而且这里还有更多结构（特别是 $V$ 必须是阿贝尔群）可加以利用。在特殊情形 $V=\RR^n$ 且 $F=\RR$ 时，该作用由
  \[
    a(r_1,r_2,\ldots,r_n)=(ar_1,ar_2,\ldots,ar_n)
  \]
  给出，其中 $a\in\RR$，$(r_1,r_2,\ldots,r_n)\in\RR^n$，$ar_i$ 只是两个实数的乘法。

  \item 对任意非空集合 $A$，对称群 $S_A$ 通过 $\sigma\cdot a=\sigma(a)$（对所有 $\sigma\in S_A$，$a\in A$）作用在 $A$ 上。相应的置换表示就是从 $S_A$ 到自身的恒等映射。

  \item 若固定正 $n$ 边形顶点的一个标号，则 $D_{2n}$ 的每个元素 $\alpha$ 都根据该对称如何置换对应顶点，给出 $\{1,2,\ldots,n\}$ 的一个置换 $\sigma_\alpha$。由
  \[
    (\alpha,i)\mapsto \sigma_\alpha(i)
  \]
  定义的从 $D_{2n}\times\{1,2,\ldots,n\}$ 到 $\{1,2,\ldots,n\}$ 的映射，给出 $D_{2n}$ 在 $\{1,2,\ldots,n\}$ 上的群作用。按照群作用的记号，我们现在可以不用形式化且笨重的 $\sigma_\alpha(i)$，而写作 $\alpha i$。注意这个作用是忠实的：正 $n$ 边形的不同对称诱导顶点上的不同置换。

  当 $n=3$ 时，$D_6$ 在三角形三个（带标号）顶点上的作用给出了从 $D_6$ 到 $S_3$ 的单同态。由于这两个群阶相同，该映射也必为满射，即为同构：$D_6\cong S_3$。这是前一节通过生成元与关系建立的同一事实的另一个证明。几何上它表示：三角形顶点的任意置换都是一个对称。对任意 $n\geq4$ 的 $n$ 边形，类似说法不成立（仅从阶数考虑，就不可能有 $D_{2n}$ 与 $S_n$ 同构）。

  \item 设 $G$ 是任意群，并令 $A=G$。定义从 $G\times A$ 到 $A$ 的映射为 $g\cdot a=ga$，对每个 $g\in G$ 与 $a\in A$，其中右侧的 $ga$ 是群 $G$ 中 $g$ 与 $a$ 的乘积。这样得到 $G$ 在自身上的一个群作用，每个固定的 $g\in G$ 都通过左乘置换 $G$ 中元素：
  \[
    g:a\mapsto ga\qquad\text{对所有 }a\in G.
  \]
  （若 $G$ 用加法记号，则得到 $a\mapsto g+a$，并称为左平移。）这个作用称为 $G$ 在自身上的\emph{左正则作用}。由消去律可知，这个作用是忠实的（请检验）。
\end{enumerate}

其他作用的例子见练习。

\subsection*{练习}
\begin{enumerate}[label=\arabic*.]
  \item 设 $F$ 是域。证明 $F$ 的非零元素乘法群（记为 $F^\times$）通过 $g\cdot a=ga$ 作用在集合 $F$ 上，其中 $g\in F^\times$、$a\in F$，而 $ga$ 是 $F$ 中两个域元素的通常乘积（请明确指出使用了域定义中的哪些公理）。
  \item 证明加法群 $\ZZ$ 通过 $z\cdot a=z+a$（对所有 $z,a\in\ZZ$）作用在自身上。
  \item 证明加法群 $\RR$ 通过 $r\cdot(x,y)=(x+ry,y)$ 作用在 $x,y$ 平面 $\RR\times\RR$ 上。
  \item 设群 $G$ 作用在集合 $A$ 上，并固定某个 $a\in A$。证明下列集合是 $G$ 的子群（参见第 1 节练习 26）：
  \begin{enumerate}[label=(\alph*)]
    \item 该作用的核；
    \item $\{g\in G\mid ga=a\}$——这个子群称为 $a$ 在 $G$ 中的\emph{稳定子}。
  \end{enumerate}
  \item 证明群 $G$ 在集合 $A$ 上作用的核，与相应置换表示 $G\to S_A$ 的核相同（参见第 6 节练习 14）。
  \item 证明群 $G$ 在集合 $A$ 上忠实作用，当且仅当该作用的核只由单位元组成。
  \item 证明本节例 2 中的作用是忠实的。
  \item 设 $A$ 是非空集合，$k$ 是正整数且 $k\leq |A|$。对称群 $S_A$ 通过
  \[
    \sigma\cdot\{a_1,\ldots,a_k\}=\{\sigma(a_1),\ldots,\sigma(a_k)\}
  \]
  作用在由 $A$ 的所有 $k$ 元子集组成的集合 $B$ 上。
  \begin{enumerate}[label=(\alph*)]
    \item 证明这是一个群作用。
    \item 明确描述元素 $(1\ 2)$ 和 $(1\ 2\ 3)$ 如何作用在 $\{1,2,3,4\}$ 的六个 $2$ 元子集上。
  \end{enumerate}
  \item 把前一题中的“$k$ 元子集”换成“有序 $k$ 元组”，完成同样的两个部分；其中对 $k$ 元组的作用按上式定义，但把集合花括号换成圆括号（注意，例如 $2$ 元组 $(1,2)$ 与 $(2,1)$ 不同，而集合 $\{1,2\}$ 与 $\{2,1\}$ 相同，所以被作用的集合不同）。
  \item 参考前两题，判断：
  \begin{enumerate}[label=(\alph*)]
    \item 对哪些 $k$，$S_n$ 在 $k$ 元子集上的作用是忠实的；
    \item 对哪些 $k$，$S_n$ 在有序 $k$ 元组上的作用是忠实的。
  \end{enumerate}
  \item 写出与 $D_8$ 作用在正方形顶点上相对应的 $S_4$ 中八个置换的循环分解（正方形顶点按第 2 节中的方式标号）。
  \item 假设 $n$ 为偶正整数，证明 $D_{2n}$ 作用在由正 $n$ 边形的相对顶点对组成的集合上。求这个作用的核（照常标号顶点）。
  \item 求左正则作用的核。
  \item 设 $G$ 是群并令 $A=G$。证明若 $G$ 非阿贝尔，则由 $g\cdot a=ag$（对所有 $g,a\in G$）定义的映射不满足 $G$ 在自身上的（左）群作用公理。
  \item 设 $G$ 是任意群并令 $A=G$。证明由 $g\cdot a=ag^{-1}$（对所有 $g,a\in G$）定义的映射满足 $G$ 在自身上的（左）群作用公理。
  \item 设 $G$ 是任意群并令 $A=G$。证明由 $g\cdot a=gag^{-1}$（对所有 $g,a\in G$）定义的映射满足（左）群作用公理（这个 $G$ 在自身上的作用称为\emph{共轭作用}）。
  \item 设 $G$ 是群，并令 $G$ 通过左共轭作用在自身上，所以每个 $g\in G$ 把 $G$ 映到 $G$：
  \[
    x\mapsto gxg^{-1}.
  \]
  对固定的 $g\in G$，证明由 $g$ 的共轭给出的映射是从 $G$ 到自身的同构（即 $G$ 的自同构——参见第 6 节练习 20）。推出对所有 $x\in G$，$x$ 与 $gxg^{-1}$ 有相同的阶；并且对 $G$ 的任意子集 $A$，有 $|A|=|gAg^{-1}|$（其中 $gAg^{-1}=\{gag^{-1}\mid a\in A\}$）。
  \item 设 $H$ 是作用在集合 $A$ 上的群。证明 $A$ 上由
  \[
    a\sim b\quad\text{当且仅当}\quad a=hb\text{ 对某个 }h\in H
  \]
  定义的关系 $\sim$ 是等价关系。（对每个 $x\in A$，$x$ 在 $\sim$ 下的等价类称为 $x$ 在 $H$ 的作用下的\emph{轨道}。$H$ 作用下的轨道划分集合 $A$。）
  \item 设 $H$ 是有限群 $G$ 的子群（参见第 1 节练习 26），并令 $H$ 通过左乘作用在 $G$ 上（这里 $A=G$）。设 $x\in G$，令 $O$ 为 $x$ 在 $H$ 的作用下的轨道。证明映射
  \[
    H\to O,\qquad h\mapsto hx
  \]
  是双射（因此所有轨道的基数都是 $|H|$）。由此和前一题推出拉格朗日定理：
  \[
    \text{若 }G\text{ 是有限群且 }H\text{ 是 }G\text{ 的子群，则 }|H|\text{ 整除 }|G|.
  \]
  \item 说明正四面体的刚体运动群同构于 $S_4$ 的一个子群（参见第 1 节练习 26）。
  \item 说明立方体的刚体运动群同构于 $S_4$。［这个群作用在四对相对顶点组成的集合上。］
  \item 说明正八面体的刚体运动群同构于 $S_4$ 的一个子群（参见第 1 节练习 26）。［这个群作用在四对相对面组成的集合上。］推出立方体与正八面体的刚体运动群同构。（这些群同构是因为这些立体互为“对偶”——见 H. Coxeter, \textit{Introduction to Geometry}, Wiley, 1961。稍后我们会看到正十二面体与正二十面体的刚体运动群也同构——这些立体同样互为对偶。）
  \item 解释为什么立方体的刚体运动群在三对相对面组成的集合上的作用不是忠实的。求该作用的核。
\end{enumerate}

% \chapter{子群}

\section{定义与例子}

研究任何由一组公理定义的数学对象的结构，一个基本方法是研究该对象中同样满足这些公理的子集。我们从讨论群的子群开始这一计划。揭示结构的第二个基本方法，是研究一个对象的商；商群的概念，粗略地说，是把一个群压缩成一个较小群的方法，将在下一章处理。这两个主题将在全书中反复出现：我们会研究群的子群与商群、环的子环与商环、向量空间的子空间与商空间，等等。

\noindent\textbf{定义.} 设 $G$ 是一个群。若 $G$ 的子集 $H$ 非空，并且对乘积和取逆封闭（即 $x,y\in H$ 蕴含 $x^{-1}\in H$ 且 $xy\in H$），则称 $H$ 是 $G$ 的一个\emph{子群}。若 $H$ 是 $G$ 的子群，记作 $H\leq G$。

$G$ 的子群就是 $G$ 的那些子集，它们关于 $G$ 中定义的运算本身也是群；也就是说，$G$ 上的二元运算限制到 $H$ 后给出 $H$ 上的二元运算，该运算结合，在 $H$ 中有单位元，并且 $H$ 的每个元素在 $H$ 中有逆元。

当我们说 $H$ 是 $G$ 的子群时，总是指 $H$ 上的群运算是 $G$ 上的运算限制到 $H$ 得到的运算（一般说来，子集 $H$ 也可能关于某个不同于 $G$ 的限制运算的运算而具有群结构，参见下面例 5(a)）。正如我们对函数限制到子集时所做的那样，我们用同一个符号表示 $G$ 的运算和子群 $H$ 的运算。若 $H\leq G$ 且 $H\neq G$，则写作 $H<G$，以强调包含是真包含。

若 $H$ 是 $G$ 的子群，则由于 $H$ 的运算是 $G$ 的运算在 $H$ 上的限制，$H$ 中的任何等式也可看作 $G$ 中的等式。因此，$G$ 中的消去律说明：$H$ 的单位元与 $G$ 的单位元相同（特别地，每个子群都必须含有 $1$，即 $G$ 的单位元）；并且 $H$ 中元素 $x$ 的逆元，与把 $x$ 看作 $G$ 中元素时的逆元相同（所以记号 $x^{-1}$ 没有歧义）。

\paragraph{例}
\begin{enumerate}[label=(\arabic*)]
  \item 在加法运算下，$\ZZ\leq \QQ$ 且 $\QQ\leq \RR$。
  \item 任意群 $G$ 都有两个子群：$H=G$ 和 $H=\{1\}$；后者称为\emph{平凡子群}，今后记作 $1$。
  \item 若 $G=D_{2n}$ 是阶为 $2n$ 的二面体群，令 $H=\{1,r,r^2,\ldots,r^{n-1}\}$ 为 $G$ 中所有旋转组成的集合。由于两个旋转的乘积仍是旋转，旋转的逆仍是旋转，所以 $H$ 是 $D_{2n}$ 的一个阶为 $n$ 的子群。
  \item 偶整数集合是所有整数在加法下构成的群的一个子群。
  \item 以下是一些不是子群的子集的例子：
  \begin{enumerate}[label=(\alph*)]
    \item $\QQ-\{0\}$ 在乘法下不是 $\RR$ 在加法下的子群，尽管二者都是群且 $\QQ-\{0\}$ 是 $\RR$ 的子集；$\QQ-\{0\}$ 上的乘法不是 $\RR$ 上加法的限制。
    \item $\ZZ^+$（在加法下）不是 $\ZZ$（在加法下）的子群，因为虽然 $\ZZ^+$ 对 $+$ 封闭，但它不含 $\ZZ$ 的单位元 $0$；并且虽然每个 $x\in\ZZ^+$ 在 $\ZZ$ 中有加法逆元 $-x$，但 $-x\notin\ZZ^+$，即 $\ZZ^+$ 对取逆不封闭（特别地，$\ZZ^+$ 在加法下不是群）。基于类似原因，$(\ZZ-\{0\},\times)$ 不是 $(\QQ-\{0\},\times)$ 的子群。
    \item $D_6$ 不是 $D_8$ 的子群，因为前者甚至不是后者的子集。
  \end{enumerate}
  \item “是……的子群”这一关系具有传递性：若 $H$ 是群 $G$ 的子群，$K$ 是 $H$ 的子群，则 $K$ 也是 $G$ 的子群。
\end{enumerate}

正如我们在第 1 章看到的，即使在简单例子中，检查某个给定二元运算是否满足所有群公理（尤其是结合律）也可能相当繁琐。然而，一旦我们知道已经有一个群，要检查其中某个子集是否为子群就容易得多，因为只需检查它是否对乘法和取逆封闭。下一个命题说明这两个条件可合并为一个检验；并且对于有限群，只需检查乘法封闭即可。

\medskip
\noindent\textbf{命题 1（子群判别法）.} 群 $G$ 的子集 $H$ 是子群，当且仅当
\begin{enumerate}[label=(\arabic*)]
  \item $H\neq\varnothing$；
  \item 对所有 $x,y\in H$，$xy^{-1}\in H$。
\end{enumerate}
此外，若 $H$ 是有限集，则只需检查 $H$ 非空且对乘法封闭。

\begin{proof}
若 $H$ 是 $G$ 的子群，则 (1) 和 (2) 显然成立，因为 $H$ 含有 $G$ 的单位元、含有每个元素的逆元，并且对乘法封闭。

反过来，假设 $H$ 满足 (1) 和 (2)，我们证明 $H\leq G$。取 $H$ 中任意元素 $x$（由 (1) 可取）。令 $y=x$ 并应用 (2)，得 $1=xx^{-1}\in H$，所以 $H$ 含有 $G$ 的单位元。再由 (2)，由于 $H$ 含有 $1$ 和 $x$，$H$ 含有元素 $1x^{-1}$，即 $x^{-1}\in H$，所以 $H$ 对取逆封闭。最后，若 $x,y$ 是 $H$ 的任意两个元素，则由刚才证明，$y^{-1}\in H$；再由 (2)，$H$ 也含有 $x(y^{-1})^{-1}=xy$。因此 $H$ 也对乘法封闭，从而 $H$ 是 $G$ 的子群。

现在假设 $H$ 有限且对乘法封闭，取任意 $x\in H$。在 $x,x^2,x^3,\ldots$ 中只有有限多个不同元素，所以存在整数 $a,b$，$b>a$，使得 $x^a=x^b$。若 $n=b-a$，则 $x^n=1$，因此 $H$ 中每个元素都有有限阶。于是 $x^{n-1}=x^{-1}$ 属于 $H$，所以 $H$ 自动对取逆封闭。
\end{proof}

\subsection*{练习}

设 $G$ 是一个群。
\begin{enumerate}[label=\arabic*.]
  \item 在 (a)--(e) 中，证明指定子集是给定群的子群：
  \begin{enumerate}[label=(\alph*)]
    \item 形如 $a+ai$、$a\in\RR$ 的复数集合（在加法下）；
    \item 绝对值为 $1$ 的复数集合，即复平面中的单位圆（在乘法下）；
    \item 对固定 $n\in\ZZ^+$，所有分母整除 $n$ 的有理数集合（在加法下）；
    \item 对固定 $n\in\ZZ^+$，所有分母与 $n$ 互素的有理数集合（在加法下）；
    \item 所有平方为有理数的非零实数集合（在乘法下）。
  \end{enumerate}
  \item 在 (a)--(e) 中，证明指定子集不是给定群的子群：
  \begin{enumerate}[label=(\alph*)]
    \item 当 $n\geq 3$ 时，$S_n$ 中所有 $2$-轮换的集合；
    \item 当 $n\geq 3$ 时，$D_{2n}$ 中所有反射的集合；
    \item 对复合整数 $n>1$ 和一个含有阶为 $n$ 的元素的群 $G$，集合
    \[
      \{x\in G\mid |x|=n\}\cup\{1\};
    \]
    \item $\ZZ$ 中所有正、负奇整数与 $0$ 合在一起的集合；
    \item 平方为有理数的实数集合（在加法下）。
  \end{enumerate}
  \item 证明二面体群 $D_8$ 的下列子集实际上是子群：
  \[
    \text{(a) }\{1,r^2,s,sr^2\},\qquad
    \text{(b) }\{1,r^2,sr,sr^3\}.
  \]
  \item 给出一个群 $G$ 和一个无限子集 $H$ 的具体例子，使得 $H$ 对群运算封闭，但不是 $G$ 的子群。
  \item 证明若 $n=|G|>2$，则 $G$ 不可能有一个满足 $|H|=n-1$ 的子群 $H$。
  \item 设 $G$ 是阿贝尔群。证明
  \[
    \{g\in G\mid |g|<\infty\}
  \]
  是 $G$ 的子群（称为 $G$ 的\emph{挠子群}）。给出一个具体例子说明，当 $G$ 非阿贝尔时，这个集合不一定是子群。
  \item 固定整数 $n>1$。求 $\ZZ\times(\ZZ/n\ZZ)$ 的挠子群（参见上一题）。说明无限阶元素再加上单位元组成的集合不是这个直积的子群。
  \item 设 $H$ 和 $K$ 是 $G$ 的子群。证明 $H\cup K$ 是子群，当且仅当 $H\subseteq K$ 或 $K\subseteq H$。
  \item 设 $G=\GL_n(F)$，其中 $F$ 是任意域。定义
  \[
    \SL_n(F)=\{A\in \GL_n(F)\mid \det(A)=1\}
  \]
  （称为\emph{特殊线性群}）。证明 $\SL_n(F)\leq \GL_n(F)$。
  \item \begin{enumerate}[label=(\alph*)]
    \item 证明若 $H$ 与 $K$ 是 $G$ 的子群，则交 $H\cap K$ 也是 $G$ 的子群。
    \item 证明 $G$ 的任意非空子群族的交仍是 $G$ 的子群（不要假设该族可数）。
  \end{enumerate}
  \item 设 $A$ 和 $B$ 是群。证明下列集合是直积 $A\times B$ 的子群：
  \begin{enumerate}[label=(\alph*)]
    \item $\{(a,1)\mid a\in A\}$；
    \item $\{(1,b)\mid b\in B\}$；
    \item $\{(a,a)\mid a\in A\}$，这里假设 $B=A$（称为\emph{对角子群}）。
  \end{enumerate}
  \item 设 $A$ 是阿贝尔群，固定 $n\in\ZZ$。证明下列集合是 $A$ 的子群：
  \[
    \text{(a) }\{a^n\mid a\in A\},\qquad
    \text{(b) }\{a\in A\mid a^n=1\}.
  \]
  \item 设 $H$ 是有理数加法群的一个子群，并满足对 $H$ 的每个非零元素 $x$ 都有 $1/x\in H$。证明 $H=0$ 或 $H=\QQ$。
  \item 证明 $\{x\in D_{2n}\mid x^2=1\}$ 不是 $D_{2n}$ 的子群（这里 $n\geq 3$）。
  \item 设 $H_1\leq H_2\leq\cdots$ 是 $G$ 的一个升链子群。证明 $\bigcup_{i=1}^{\infty}H_i$ 是 $G$ 的子群。
  \item 设 $n\in\ZZ^+$，$F$ 是一个域。证明集合
  \[
    \{(a_{ij})\in\GL_n(F)\\mid a_{ij}=0\text{ 对所有 }i>j\}
  \]
  是 $\GL_n(F)$ 的一个子群（称为\emph{上三角矩阵群}）。
  \item 设 $n\in\ZZ^+$，$F$ 是一个域。证明集合
  \[
    \{(a_{ij})\in\GL_n(F)\mid a_{ij}=0\text{ 对所有 }i>j,\text{ 且 }a_{ii}=1\text{ 对所有 }i\}
  \]
  是 $\GL_n(F)$ 的一个子群。
\end{enumerate}

\section{中心化子与正规化子，稳定子与核}

现在我们引入任意群 $G$ 中一些重要的子群族；它们尤其提供许多子群的例子。设 $A$ 是 $G$ 的任意非空子集。

\noindent\textbf{定义.} 定义
\[
  C_G(A)=\{g\in G\mid gag^{-1}=a\text{ 对所有 }a\in A\}.
\]
这个 $G$ 的子集称为 $A$ 在 $G$ 中的\emph{中心化子}。由于 $gag^{-1}=a$ 当且仅当 $ga=ag$，$C_G(A)$ 就是 $G$ 中与 $A$ 的每个元素都交换的那些元素组成的集合。

我们证明 $C_G(A)$ 是 $G$ 的子群。首先，$C_G(A)\neq\varnothing$，因为 $1\in C_G(A)$：单位元的定义给出 $1a=a1$ 对所有 $a\in G$ 成立，特别地对所有 $a\in A$ 成立，所以 $1$ 满足属于 $C_G(A)$ 的定义条件。其次，假设 $x,y\in C_G(A)$，即对所有 $a\in A$，有 $xax^{-1}=a$ 且 $yay^{-1}=a$（注意，这并不意味着 $xy=yx$）。先观察到，由 $yay^{-1}=a$，先左乘 $y^{-1}$ 再右乘 $y$ 并化简，得到 $a=y^{-1}ay$；也就是说，$y^{-1}\in C_G(A)$，所以 $C_G(A)$ 对取逆封闭。现在
\begin{align*}
  (xy)a(xy)^{-1}
  &= (xy)a(y^{-1}x^{-1}) &&\text{（由命题 1.1(4) 用于 $(xy)^{-1}$）}\\
  &= x(yay^{-1})x^{-1} &&\text{（结合律）}\\
  &= xax^{-1} &&\text{（因为 $y\in C_G(A)$）}\\
  &= a &&\text{（因为 $x\in C_G(A)$），}
\end{align*}
所以 $xy\in C_G(A)$，$C_G(A)$ 对乘积封闭，故 $C_G(A)\leq G$。

在特殊情形 $A=\{a\}$ 下，我们简写 $C_G(a)$ 而不是 $C_G(\{a\})$。此时对所有 $n\in\ZZ$，都有 $a^n\in C_G(a)$。

例如，在阿贝尔群 $G$ 中，对所有子集 $A$，有 $C_G(A)=G$。通过检查可得 $C_{Q_8}(i)=\{\pm1,\pm i\}$。练习中还会给出若干其他例子。

稍后我们会讨论如何减少计算中心化子（以及其他类似性质的子群）时看起来必不可少的单个群元素之间交换性的检查。

\noindent\textbf{定义.} 定义
\[
  Z(G)=\{g\in G\mid gx=xg\text{ 对所有 }x\in G\},
\]
即与 $G$ 中所有元素都交换的元素集合。这个 $G$ 的子集称为 $G$ 的\emph{中心}。注意 $Z(G)=C_G(G)$，所以上面的论证把 $Z(G)\leq G$ 作为特殊情形证明了。读者也可直接证明 $Z(G)$ 是子群。

\noindent\textbf{定义.} 定义
\[
  gAg^{-1}=\{gag^{-1}\mid a\in A\}.
\]
定义 $A$ 在 $G$ 中的\emph{正规化子}为
\[
  N_G(A)=\{g\in G\mid gAg^{-1}=A\}.
\]

注意若 $g\in C_G(A)$，则对所有 $a\in A$ 有 $gag^{-1}=a\in A$，所以 $C_G(A)\leq N_G(A)$。证明 $N_G(A)$ 是 $G$ 的子群，与证明 $C_G(A)\leq G$ 的步骤相同，只需作适当修改。

\paragraph{例}
\begin{enumerate}[label=(\arabic*)]
  \item 若 $G$ 是阿贝尔群，则 $G$ 中所有元素都交换，所以 $Z(G)=G$。类似地，对 $G$ 的任意子集 $A$，$C_G(A)=N_G(A)=G$，因为对每个 $g\in G$ 和每个 $a\in A$，都有 $gag^{-1}=gg^{-1}a=a$。
  \item 设 $G=D_8$ 是阶为 $8$ 的二面体群，有通常的生成元 $r$ 和 $s$；设 $A=\{1,r,r^2,r^3\}$ 为 $D_8$ 中旋转构成的子群。我们说明 $C_{D_8}(A)=A$。由于 $r$ 的各个幂彼此交换，$A\leq C_{D_8}(A)$。由于 $sr=r^{-1}s\neq rs$，元素 $s$ 不与 $A$ 的所有成员交换，即 $s\notin C_{D_8}(A)$。最后，$D_8$ 中不属于 $A$ 的元素都形如 $sr^i$，其中 $i\in\{0,1,2,3\}$。若元素 $sr^i$ 属于 $C_{D_8}(A)$，则由于 $C_{D_8}(A)$ 是含有 $r$ 的子群，也会有
  \[
    s=(sr^i)(r^{-i})\in C_{D_8}(A),
  \]
  这与上面矛盾。因此 $C_{D_8}(A)=A$。
  \item 仍如前例，令 $G=D_8$ 且 $A=\{1,r,r^2,r^3\}$。我们说明 $N_{D_8}(A)=D_8$。一般地，子集的中心化子包含于其正规化子，所以 $A\leq N_{D_8}(A)$。再计算
  \[
    sAs^{-1}=\{s1s^{-1},srs^{-1},sr^2s^{-1},sr^3s^{-1}\}
    =\{1,r^3,r^2,r\}=A,
  \]
  因此 $s\in N_{D_8}(A)$。（注意，集合 $sAs^{-1}$ 等于集合 $A$，尽管两个集合中的元素出现顺序不同；这是因为 $s$ 在 $A$ 的正规化子中，但不在 $A$ 的中心化子中。）现在 $r$ 和 $s$ 都属于子群 $N_{D_8}(A)$，所以对所有整数 $i,j$，$s^jr^i\in N_{D_8}(A)$；也就是说，$D_8$ 的每个元素都在 $N_{D_8}(A)$ 中（回忆 $r$ 和 $s$ 生成 $D_8$）。既然 $D_8\leq N_{D_8}(A)$，而反向包含由正规化子定义显然成立，故 $N_{D_8}(A)=D_8$。
  \item 我们说明 $D_8$ 的中心是子群 $\{1,r^2\}$。首先，任意群 $G$ 的中心都包含于 $C_G(A)$，其中 $A$ 是 $G$ 的任意子集。于是由例 2，
  \[
    Z(D_8)\leq C_{D_8}(A)=A,
  \]
  其中 $A=\{1,r,r^2,r^3\}$。例 2 中的计算说明 $r$，同样 $r^3$，不在 $Z(D_8)$ 中，所以 $Z(D_8)\leq\{1,r^2\}$。为证明反向包含，注意 $r$ 与 $r^2$ 交换，并且直接计算可得 $s$ 也与 $r^2$ 交换。由于 $r$ 与 $s$ 生成 $D_8$，$D_8$ 中每个元素都与 $r^2$（以及 $1$）交换，因此 $\{1,r^2\}\leq Z(D_8)$，故等号成立。
  \item 设 $G=S_3$，令 $A=\{1,(1\ 2)\}$。我们解释为什么 $C_{S_3}(A)=N_{S_3}(A)=A$。可以直接计算 $C_{S_3}(A)=A$，并用例 2 中的思想减少计算。另一种方法是：由于一个元素与其幂交换，$A\leq C_{S_3}(A)$。由拉格朗日定理（第 1.7 节练习 19），子群 $C_{S_3}(A)$ 的阶整除 $|S_3|=6$。同样把拉格朗日定理应用于 $C_{S_3}(A)$ 的子群 $A$，得 $2$ 整除 $|C_{S_3}(A)|$。唯一的可能是 $|C_{S_3}(A)|=2$ 或 $6$。若为后者，则 $C_{S_3}(A)=S_3$，即 $A\leq Z(S_3)$；这与 $(1\ 2)$ 不与 $(1\ 2\ 3)$ 交换矛盾。因此 $|C_{S_3}(A)|=2$，故 $A=C_{S_3}(A)$。

  接着注意 $N_{S_3}(A)=A$，因为 $\sigma\in N_{S_3}(A)$ 当且仅当
  \[
    \{\sigma1\sigma^{-1},\sigma(1\ 2)\sigma^{-1}\}=\{1,(1\ 2)\}.
  \]
  由于 $\sigma1\sigma^{-1}=1$，这个集合等式成立当且仅当 $\sigma(1\ 2)\sigma^{-1}=(1\ 2)$，也就是当且仅当 $\sigma\in C_{S_3}(A)$。

  $S_3$ 的中心是单位子群，因为 $Z(S_3)\leq C_{S_3}(A)=A$，而 $(1\ 2)\notin Z(S_3)$。
\end{enumerate}

\subsection*{群作用的稳定子与核}

$A$ 在 $G$ 中的正规化子、中心化子以及 $G$ 的中心都是子群这一事实，可作为群作用相关结果的特殊情形推出。这说明 $G$ 的结构会反映在它所作用的集合上。具体地说，若群 $G$ 作用在集合 $S$ 上，且 $s$ 是 $S$ 中的某个固定元素，则 $s$ 在 $G$ 中的\emph{稳定子}是集合
\[
  G_s=\{g\in G\mid g\cdot s=s\}
\]
（参见第 1.7 节练习 4）。我们简要说明 $G_s\leq G$：首先，由作用公理 (2)，$1\in G_s$。此外，若 $y\in G_s$，则
\begin{align*}
  s &= 1\cdot s=(y^{-1}y)\cdot s\\
    &= y^{-1}\cdot (y\cdot s) &&\text{（由作用公理 (1)）}\\
    &= y^{-1}\cdot s &&\text{（因为 $y\in G_s$），}
\end{align*}
所以 $y^{-1}\in G_s$。最后，若 $x,y\in G_s$，则
\begin{align*}
  (xy)\cdot s &= x\cdot(y\cdot s) &&\text{（由作用公理 (1)）}\\
  &=x\cdot s &&\text{（因为 $y\in G_s$）}\\
  &=s &&\text{（因为 $x\in G_s$）。}
\end{align*}
这证明了 $G_s$ 是 $G$ 的子群。类似但更简单的论证证明作用的核也是子群，其中 $G$ 在 $S$ 上作用的\emph{核}定义为
\[
  \{g\in G\mid g\cdot s=s\text{ 对所有 }s\in S\}
\]
（参见第 1.7 节练习 1）。

\paragraph{例}
\begin{enumerate}[label=(\arabic*)]
  \item 群 $G=D_8$ 作用在正方形四个顶点的集合 $A$ 上（参见第 1.7 节例 4）。任意顶点 $a$ 的稳定子是 $D_8$ 的子群 $\{1,t\}$，其中 $t$ 是关于经过顶点 $a$ 与正方形中心的对称轴的反射。这个作用的核是单位子群，因为只有恒等对称固定每个顶点。
  \item 群 $G=D_8$ 也作用在集合 $A$ 上，$A$ 的元素是两对相对顶点的无序对（按第 1.2 节图 2 的标号，$A=\{\{1,3\},\{2,4\}\}$）。$D_8$ 在这个集合 $A$ 上作用的核是子群 $\{1,s,r^2,sr^2\}$；对 $A$ 的任一元素 $a$，$a$ 在 $D_8$ 中的稳定子等于该作用的核。
\end{enumerate}

最后我们观察，中心化子、正规化子以及核是子群这一事实，都是作用的稳定子和核是子群这一事实的特殊情形（这会在第 4 章进一步讨论）。令 $S=\mcP(G)$ 为 $G$ 的所有子集的集合，并令 $G$ 在 $S$ 上按共轭作用：对每个 $g\in G$ 和每个 $B\subseteq G$，令
\[
  g:B\mapsto gBg^{-1},\qquad gBg^{-1}=\{gbg^{-1}\mid b\in B\}
\]
（见第 1.7 节练习 16）。在这个作用下，容易验证 $N_G(A)$ 正是 $A$ 在 $G$ 中的稳定子（即 $N_G(A)=G_s$，其中 $s=A\in\mcP(G)$），所以 $N_G(A)$ 是 $G$ 的子群。

再令群 $N_G(A)$ 在集合 $S=A$ 上按共轭作用，即对所有 $g\in N_G(A)$ 和 $a\in A$，
\[
  g:a\mapsto gag^{-1}.
\]
注意由 $N_G(A)$ 的定义，这确实把 $A$ 映到 $A$，从而给出 $A$ 上的作用。此时容易验证 $C_G(A)$ 正是这个作用的核，因此 $C_G(A)\leq N_G(A)$；由关系“$\leq$”的传递性，$C_G(A)\leq G$。最后，$Z(G)$ 是 $G$ 在 $S=G$ 上按共轭作用的核，所以 $Z(G)\leq G$。

\subsection*{练习}

\begin{enumerate}[label=\arabic*.]
  \item 证明
  \[
    C_G(A)=\{g\in G\mid g^{-1}ag=a\text{ 对所有 }a\in A\}.
  \]
  \item 证明 $C_G(Z(G))=G$，并推出 $N_G(Z(G))=G$。
  \item 证明若 $A$ 与 $B$ 是 $G$ 的子集且 $A\subseteq B$，则 $C_G(B)$ 是 $C_G(A)$ 的子群。
  \item 对 $S_3$、$D_8$ 和 $Q_8$，计算每个元素的中心化子，并求每个群的中心。拉格朗日定理（第 1.7 节练习 19）是否能简化你的工作？
  \item 在 (a)--(c) 中，对指定的群 $G$ 及 $G$ 的子群 $A$，证明 $C_G(A)=A$ 且 $N_G(A)=G$。
  \begin{enumerate}[label=(\alph*)]
    \item $G=S_3$，$A=\{1,(1\ 2\ 3),(1\ 3\ 2)\}$；
    \item $G=D_8$，$A=\{1,s,r^2,sr^2\}$；
    \item $G=D_{10}$，$A=\{1,r,r^2,r^3,r^4\}$。
  \end{enumerate}
  \item 设 $H$ 是群 $G$ 的子群。
  \begin{enumerate}[label=(\alph*)]
    \item 证明 $H\leq N_G(H)$。给出一个例子说明，若 $H$ 不是子群，则该结论不一定成立。
    \item 证明 $H\leq C_G(H)$ 当且仅当 $H$ 是阿贝尔群。
  \end{enumerate}
  \item 设 $n\in\ZZ$ 且 $n\geq 3$。证明：
  \begin{enumerate}[label=(\alph*)]
    \item 若 $n$ 为奇数，则 $Z(D_{2n})=1$；
    \item 若 $n=2k$，则 $Z(D_{2n})=\{1,r^k\}$。
  \end{enumerate}
  \item 令 $G=S_n$，固定 $i\in\{1,2,\ldots,n\}$，并令
  \[
    G_i=\{\sigma\in G\mid \sigma(i)=i\}
  \]
  （即 $i$ 在 $G$ 中的稳定子）。用群作用证明 $G_i$ 是 $G$ 的子群，并求 $|G_i|$。
  \item 对 $G$ 的任意子群 $H$ 和任意非空子集 $A$，定义
  \[
    N_H(A)=\{h\in H\mid hAh^{-1}=A\}.
  \]
  证明 $N_H(A)=N_G(A)\cap H$，并推出 $N_H(A)$ 是 $H$ 的子群（注意 $A$ 不必是 $H$ 的子集）。
  \item 设 $H$ 是 $G$ 中阶为 $2$ 的子群。证明 $N_G(H)=C_G(H)$。推出若 $N_G(H)=G$，则 $H\leq Z(G)$。
  \item 证明对 $G$ 的任意子集 $A$，都有 $Z(G)\leq N_G(A)$。
  \item 设 $R$ 是所有以独立变量 $x_1,x_2,x_3,x_4$ 为变量、整数为系数的多项式集合；也就是说，$R$ 的元素是形如 $a x_1^{r_1}x_2^{r_2}x_3^{r_3}x_4^{r_4}$ 的项的有限和，其中 $a$ 是整数，$r_1,\ldots,r_4$ 是非负整数。例如
  \[
    12x_1^5x_2^7x_4-18x_2^3x_3+11x_1^6x_2^3x_3^2x_4^3 \tag{*}
  \]
  是 $R$ 的一个典型元素。每个 $\sigma\in S_4$ 通过定义 $\sigma\cdot x_i=x_{\sigma(i)}$ 给出 $\{x_1,\ldots,x_4\}$ 的一个置换。将其扩展为从 $R$ 到 $R$ 的映射：
  \[
    \sigma\cdot p(x_1,x_2,x_3,x_4)=p(x_{\sigma(1)},x_{\sigma(2)},x_{\sigma(3)},x_{\sigma(4)})
  \]
  对所有 $p(x_1,x_2,x_3,x_4)\in R$ 成立（即 $\sigma$ 只是置换变量下标）。例如，若 $\sigma=(1\ 2)(3\ 4)$，且 $p$ 是 (*) 中的多项式，则
  \[
    \sigma\cdot p=12x_2^5x_1^7x_3-18x_1^3x_4+11x_2^6x_1^3x_4^2x_3^3
    =12x_1^7x_2^5x_3-18x_1^3x_4+11x_1^3x_2^6x_3^3x_4^2.
  \]
  \begin{enumerate}[label=(\alph*)]
    \item 令 $p=p(x_1,\ldots,x_4)$ 为 (*) 中的多项式，令 $\sigma=(1\ 2\ 3\ 4)$，$\tau=(1\ 2\ 3)$。计算 $\sigma\cdot p$、$\tau\cdot(\sigma\cdot p)$、$(\tau\circ\sigma)\cdot p$ 和 $(\sigma\circ\tau)\cdot p$。
    \item 证明上述定义给出 $S_4$ 在 $R$ 上的一个（左）群作用。
    \item 列出 $S_4$ 中所有稳定 $x_4$ 的置换，并证明它们构成一个同构于 $S_3$ 的子群。
    \item 列出 $S_4$ 中所有稳定 $x_1+x_2$ 的置换，并证明它们构成一个阶为 $4$ 的阿贝尔子群。
    \item 列出 $S_4$ 中所有稳定 $x_1x_2+x_3x_4$ 的置换，并证明它们构成一个同构于阶为 $8$ 的二面体群的子群。
    \item 说明稳定 $(x_1+x_2)(x_3+x_4)$ 的 $S_4$ 中的置换，恰好与 (e) 中得到的置换相同。（(e) 和 (f) 中出现的两个多项式及稳定它们的子群，将在第 14.6 节研究四次方程的根时起重要作用。）
  \end{enumerate}
  \item 设 $n$ 是正整数，$R$ 是所有以独立变量 $x_1,x_2,
  \ldots,x_n$ 为变量、整数为系数的多项式集合。对每个 $\sigma\in S_n$ 定义映射
  \[
    \sigma:R\to R,
    \qquad
    \sigma\cdot p(x_1,x_2,\ldots,x_n)=p(x_{\sigma(1)},x_{\sigma(2)},\ldots,x_{\sigma(n)}).
  \]
  证明这定义了 $S_n$ 在 $R$ 上的一个（左）群作用。
  \item 设 $H(F)$ 为第 1.4 节练习 11 中引入的域 $F$ 上的海森堡群。确定哪些矩阵属于 $H(F)$ 的中心，并证明 $Z(H(F))$ 同构于 $F$ 的加法群。
\end{enumerate}

\section{循环群与循环子群}

设 $G$ 是任意群，$x$ 是 $G$ 的任意元素。构造 $G$ 的子群 $H$ 的一种方式，是令 $H$ 为 $x$ 的所有整数次幂（正次幂、负次幂和零次幂）组成的集合；至少就 $x$ 而言，这保证了对取逆和乘积封闭。本节研究由一个元素生成的群。

\noindent\textbf{定义.} 群 $H$ 称为\emph{循环的}，如果 $H$ 可由单个元素生成；也就是说，存在某个 $x\in H$，使得
\[
  H=\{x^n\mid n\in\ZZ\}
\]
（像往常一样，这里运算写成乘法）。

若用加法记号，则 $H$ 循环表示 $H=\{nx\mid n\in\ZZ\}$。两种情形中我们都写 $H=\langle x\rangle$，并称 $H$ 由 $x$ 生成（$x$ 是 $H$ 的生成元）。一个循环群可以有不止一个生成元。例如，若 $H=\langle x\rangle$，则也有 $H=\langle x^{-1}\rangle$，因为 $(x^{-1})^n=x^{-n}$，而当 $n$ 遍历所有整数时，$-n$ 也遍历所有整数，所以
\[
  \{x^n\mid n\in\ZZ\}=\{(x^{-1})^n\mid n\in\ZZ\}.
\]
我们很快会说明如何确定给定循环群 $H$ 的全部生成元。应注意，$\langle x\rangle$ 的元素是 $x$ 的幂（或在加法群中是 $x$ 的倍数），而不是整数。$x$ 的所有幂不一定彼此不同。此外，由指数律（第 1.1 节练习 19），循环群都是阿贝尔群。

\paragraph{例}
\begin{enumerate}[label=(\arabic*)]
  \item 设 $G=D_{2n}=\langle r,s\mid r^n=s^2=1,
  rs=sr^{-1}\rangle$，$n\geq 3$，并令 $H$ 为正 $n$ 边形所有旋转构成的子群。因此 $H=\langle r\rangle$，且 $H$ 的不同元素为
  \[
    1,r,r^2,\ldots,r^{n-1}
  \]
  （这些是 $r$ 的所有不同幂）。特别地，$|H|=n$，$H$ 的生成元 $r$ 的阶为 $n$。$r$ 的幂以周期 $n$“循环”（向前或向后），即
  \[
    r^n=1,
    \quad r^{n+1}=r,
    \quad r^{n+2}=r^2,\ldots,
  \]
  且
  \[
    r^{-1}=r^{n-1},\quad r^{-2}=r^{n-2},\ldots
  \]
  一般地，要把任意幂 $r^t$ 写成 $r^k$，其中 $0\leq k<n$，用除法算法写
  \[
    t=nq+k,
    \qquad 0\leq k<n,
  \]
  则
  \[
    r^t=r^{nq+k}=(r^n)^q r^k=1^q r^k=r^k.
  \]
  例如，在 $D_8$ 中，$r^4=1$，所以 $r^{105}=r^{4(26)+1}=r$，而 $r^{-42}=r^{4(-11)+2}=r^2$。注意 $D_{2n}$ 本身不是循环群，因为它非阿贝尔。
  \item 令 $H=\ZZ$，运算为 $+$。于是 $H=\langle 1\rangle$（这里 $1$ 是整数 $1$，而 $H$ 的单位元是 $0$），并且 $H$ 的每个元素都能唯一写成 $n\cdot 1$ 的形式，其中 $n\in\ZZ$。与前例相反，生成元的各个倍数都不同，并且要用该生成元的正倍数、负倍数和零倍数才能得到 $H$ 的所有元素。在这个例子中，$|H|$ 以及生成元 $1$ 的阶都是 $\infty$。还要注意 $H=\langle -1\rangle$，因为每个整数 $x$ 都可唯一写成 $(-x)(-1)$。
\end{enumerate}

在进一步讨论循环群之前，我们先证明前两个例子中观察到的有限循环群和无限循环群的各种性质是一般性的。这个命题也说明（在第 1 章中）对元素的“阶”使用术语和符号 $|\ |$ 与集合的阶的概念是一致的。

\medskip
\noindent\textbf{命题 2.} 若 $H=\langle x\rangle$，则 $|H|=|x|$（若等式一边为无穷，则另一边也为无穷）。更具体地：
\begin{enumerate}[label=(\arabic*)]
  \item 若 $|H|=n<\infty$，则 $x^n=1$，且 $1,x,x^2,\ldots,x^{n-1}$ 是 $H$ 的全部不同元素；
  \item 若 $|H|=\infty$，则对所有 $n\neq 0$ 有 $x^n\neq 1$，且对 $\ZZ$ 中不同的 $a,b$ 有 $x^a\neq x^b$。
\end{enumerate}

\begin{proof}
设 $|x|=n$。先考虑 $n<\infty$ 的情形。元素 $1,x,x^2,\ldots,x^{n-1}$ 彼此不同，因为若 $x^a=x^b$，不妨设 $0\leq a<b<n$，则 $x^{b-a}=x^0=1$，这与 $n$ 是使 $x$ 的正幂等于单位元的最小正整数矛盾。因此 $H$ 至少有 $n$ 个元素，还需说明这些就是全部元素。与例 1 中相同，若 $x^t$ 是 $x$ 的任意幂，用除法算法写 $t=nq+k$，其中 $0\leq k<n$，于是
\[
  x^t=x^{nq+k}=(x^n)^qx^k=1^qx^k=x^k\in\{1,x,x^2,\ldots,x^{n-1}\},
\]
如所需。

接着假设 $|x|=\infty$，即 $x$ 没有正幂等于单位元。若对某些 $a,b$ 且不妨设 $a<b$，有 $x^a=x^b$，则 $x^{b-a}=1$，矛盾。故 $x$ 的不同幂给出 $H$ 的不同元素，因此 $|H|=\infty$。命题得证。
\end{proof}

注意，命题证明也给出了在有限循环群中把生成元的任意幂化为“最小剩余”幂的方法。阶为 $n$ 的循环群中生成元不同幂的计算要通过 $\ZZ/n\ZZ$ 中的算术来完成，这并非巧合。下面的定理 4 证明这两个群同构。

首先需要一个简单命题。

\medskip
\noindent\textbf{命题 3.} 设 $G$ 是任意群，$x\in G$，$m,n\in\ZZ$。若 $x^n=1$ 且 $x^m=1$，则 $x^d=1$，其中 $d=(m,n)$。特别地，若对某个 $m\in\ZZ$ 有 $x^m=1$，则 $|x|$ 整除 $m$。

\begin{proof}
由欧几里得算法（见第 0.2 节 (6)），存在整数 $r,s$，使得
\[
  d=mr+ns,
\]
其中 $d$ 是 $m$ 和 $n$ 的最大公因数。因此
\[
  x^d=x^{mr+ns}=(x^m)^r(x^n)^s=1^r1^s=1.
\]
这证明了第一断言。

若 $x^m=1$，令 $n=|x|$。若 $m=0$，则当然 $n\mid m$，所以可假设 $m\neq0$。由于 $x$ 的某个非零幂等于单位元，$n<\infty$。令 $d=(m,n)$，则由前一结论 $x^d=1$。由于 $0<d\leq n$，而 $n$ 是使 $x$ 的正幂等于单位元的最小正整数，必有 $d=n$，即 $n\mid m$，如所断言。
\end{proof}

\medskip
\noindent\textbf{定理 4.} 任意两个同阶循环群同构。更具体地：
\begin{enumerate}[label=(\arabic*)]
  \item 若 $n\in\ZZ^+$，且 $\langle x\rangle$ 与 $\langle y\rangle$ 都是阶为 $n$ 的循环群，则映射
  \[
    \varphi:\langle x\rangle\to \langle y\rangle,
    \qquad x^k\mapsto y^k
  \]
  是良定的，并且是同构；
  \item 若 $\langle x\rangle$ 是无限循环群，则映射
  \[
    \varphi:\ZZ\to\langle x\rangle,
    \qquad k\mapsto x^k
  \]
  是良定的，并且是同构。
\end{enumerate}

\begin{proof}
假设 $\langle x\rangle$ 与 $\langle y\rangle$ 都是阶为 $n$ 的循环群。定义 $\varphi:\langle x\rangle\to\langle y\rangle$ 为 $\varphi(x^k)=y^k$。首先证明 $\varphi$ 良定，即
\[
  \text{若 }x^r=x^s,\text{ 则 }\varphi(x^r)=\varphi(x^s).
\]
由于 $x^{r-s}=1$，命题 3 蕴含 $n\mid r-s$。写 $r=tn+s$，则
\[
  \varphi(x^r)=\varphi(x^{tn+s})=y^{tn+s}=(y^n)^t y^s=y^s=\varphi(x^s).
\]
这证明了 $\varphi$ 良定。由指数律立即可得 $\varphi(x^a x^b)=\varphi(x^a)\varphi(x^b)$（请检查），即 $\varphi$ 是同态。$\langle y\rangle$ 的元素 $y^k$ 是 $x^k$ 在 $\varphi$ 下的像，所以该映射满射。由于两个群同为有限且阶相同，从一个到另一个的满射必为双射，因此 $\varphi$ 是同构（也可说 $\varphi$ 显然有双侧逆）。

若 $\langle x\rangle$ 是无限循环群，定义 $\varphi:\ZZ\to\langle x\rangle$ 为 $\varphi(k)=x^k$。注意该映射已然良定，因为定义域中元素的表示没有歧义。由命题 2，不同的 $a,b\in\ZZ$ 满足 $x^a\neq x^b$，故 $\varphi$ 单射。由循环群定义，$\varphi$ 满射。如上，指数律保证 $\varphi$ 是同态，因此 $\varphi$ 是同构。定理得证。
\end{proof}

我们选择旋转群 $\langle r\rangle$ 作为阶为 $n$ 的有限循环群的原型例子（而不是同构的群 $\ZZ/n\ZZ$），因为我们通常把循环群写成乘法记号。

\paragraph{记号} 对每个 $n\in\ZZ^+$，令 $Z_n$ 表示阶为 $n$ 的循环群（以乘法记号书写）。

在同构意义下，$Z_n$ 是唯一的阶为 $n$ 的循环群，且 $Z_n\cong \ZZ/n\ZZ$。有时当加法记号更方便时，我们会用后者作为阶为 $n$ 的循环群同构类的代表。我们偶尔会说“令 $\langle x\rangle$ 是无限循环群”（写成乘法记号），但我们总用 $\ZZ$（加法记号）来表示无限循环群。

如前所述，一个给定循环群可以有多个生成元。下面两个命题精确确定 $x$ 的哪些幂生成群 $\langle x\rangle$。

\medskip
\noindent\textbf{命题 5.} 设 $G$ 是群，$x\in G$，且 $a\in\ZZ-\{0\}$。
\begin{enumerate}[label=(\arabic*)]
  \item 若 $|x|=\infty$，则 $|x^a|=\infty$。
  \item 若 $|x|=n<\infty$，则
  \[
    |x^a|=\frac{n}{(n,a)}.
  \]
  \item 特别地，若 $|x|=n<\infty$ 且正整数 $a$ 整除 $n$，则
  \[
    |x^a|=\frac{|x|}{a}.
  \]
\end{enumerate}

\begin{proof}
(1) 反证。假设 $|x|=\infty$ 但 $|x^a|=m<\infty$。由阶的定义，
\[
  1=(x^a)^m=x^{am}.
\]
又有
\[
  x^{-am}=(x^{am})^{-1}=1^{-1}=1.
\]
由于 $am$ 与 $-am$ 中有一个为正（$a,m$ 都非零），$x$ 的某个正幂为单位元。这与 $|x|=\infty$ 矛盾，所以 $|x^a|<\infty$ 的假设为假，即 (1) 成立。

(2) 在 (2) 的记号下令
\[
  y=x^a,
  \qquad (n,a)=d,
  \qquad n=db,
  \qquad a=dc,
\]
其中 $b,c\in\ZZ$ 且 $b>0$。由于 $d$ 是 $n$ 和 $a$ 的最大公因数，整数 $b$ 和 $c$ 互素：
\[
  (b,c)=1.
\]
要证明 (2)，需证明 $|y|=b$。首先
\[
  y^b=x^{ab}=x^{dcb}=(x^{db})^c=(x^n)^c=1^c=1,
\]
所以由命题 3 应用于 $\langle y\rangle$，$|y|$ 整除 $b$。令 $k=|y|$。则
\[
  x^{ak}=y^k=1,
\]
故由命题 3 应用于 $\langle x\rangle$，$n\mid ak$，即 $db\mid dck$。于是 $b\mid ck$。由于 $b$ 与 $c$ 互素，必有 $b\mid k$。既然正整数 $b$ 与 $k$ 互相整除，所以 $b=k$，这证明了 (2)。

(3) 是 (2) 的一个特殊情形，列出以备后用。
\end{proof}

\medskip
\noindent\textbf{命题 6.} 设 $H=\langle x\rangle$。
\begin{enumerate}[label=(\arabic*)]
  \item 若 $|x|=\infty$，则 $H=\langle x^a\rangle$ 当且仅当 $a=\pm1$。
  \item 若 $|x|=n<\infty$，则 $H=\langle x^a\rangle$ 当且仅当 $(a,n)=1$。特别地，$H$ 的生成元个数是 $\varphi(n)$（其中 $\varphi$ 是欧拉 $\varphi$ 函数）。
\end{enumerate}

\begin{proof}
(1) 留作练习。在 (2) 中，若 $|x|=n<\infty$，命题 2 说明 $x^a$ 生成的 $H$ 的子群阶为 $|x^a|$。这个子群等于整个 $H$ 当且仅当 $|x^a|=|x|$。由命题 5，
\[
  |x^a|=|x| \quad\text{当且仅当}\quad \frac{n}{(a,n)}=n,
\]
也就是当且仅当 $(a,n)=1$。由于 $\varphi(n)$ 的定义就是集合 $\{1,2,\ldots,n\}$ 中与 $n$ 互素的 $a$ 的个数，这也就是 $H$ 的生成元个数。
\end{proof}

\paragraph{例}
命题 6 精确说明哪些模 $n$ 的剩余类生成 $\ZZ/n\ZZ$：$\overline{a}$ 生成 $\ZZ/n\ZZ$ 当且仅当 $(a,n)=1$。例如，$\overline{1},\overline{5},\overline{7}$ 和 $\overline{11}$ 是 $\ZZ/12\ZZ$ 的生成元，并且 $\varphi(12)=4$。

本节最后一个定理给出循环群的完整子群结构。

\medskip
\noindent\textbf{定理 7.} 设 $H=\langle x\rangle$ 是循环群。
\begin{enumerate}[label=(\arabic*)]
  \item $H$ 的每个子群都是循环的。更精确地，若 $K\leq H$，则要么 $K=\{1\}$，要么 $K=\langle x^d\rangle$，其中 $d$ 是满足 $x^d\in K$ 的最小正整数。
  \item 若 $|H|=\infty$，则对任意不同的非负整数 $a,b$，有 $\langle x^a\rangle\neq \langle x^b\rangle$。此外，对任意整数 $m$，有 $\langle x^m\rangle=\langle x^{|m|}\rangle$，其中 $|m|$ 表示 $m$ 的绝对值；因此 $H$ 的非平凡子群与整数 $1,2,3,\ldots$ 一一对应。
  \item 若 $|H|=n<\infty$，则对每个整除 $n$ 的正整数 $a$，$H$ 有唯一一个阶为 $a$ 的子群。这个子群是循环群 $\langle x^d\rangle$，其中 $d=n/a$。此外，对任意整数 $m$，
  \[
    \langle x^m\rangle=\langle x^{(n,m)}\rangle,
  \]
  因而 $H$ 的子群与 $n$ 的正因子一一对应。
\end{enumerate}

\begin{proof}
(1) 设 $K\leq H$。若 $K=\{1\}$，则命题对该子群成立，所以假设 $K\neq\{1\}$。于是存在某个 $a\neq0$ 使得 $x^a\in K$。若 $a<0$，则由于 $K$ 是群，也有 $x^{-a}=(x^a)^{-1}\in K$。因此 $K$ 总含有 $x$ 的某个正幂。令
\[
  \mcP=\{b\mid b\in\ZZ^+\text{ 且 }x^b\in K\}.
\]
由上可知，$\mcP$ 是非空正整数集合。由良序原理（第 0.2 节），$\mcP$ 有最小元，记为 $d$。由于 $K$ 是子群且 $x^d\in K$，$\langle x^d\rangle\leq K$。又因 $K$ 是 $H$ 的子群，$K$ 的任一元素都形如 $x^a$，其中 $a$ 为整数。由除法算法写
\[
  a=qd+r,
  \qquad 0\leq r<d.
\]
于是
\[
  x^r=x^{a-qd}=x^a(x^d)^{-q}
\]
属于 $K$，因为 $x^a$ 与 $x^d$ 都属于 $K$。由 $d$ 的最小性，必有 $r=0$；即 $a=qd$，所以 $x^a=(x^d)^q\in\langle x^d\rangle$。这给出反向包含 $K\leq\langle x^d\rangle$，从而证明 (1)。

(2) 的证明留作练习（其推理与下面 (3) 的证明类似，但更容易）。

(3) 假设 $|H|=n<\infty$ 且 $a\mid n$。令 $d=n/a$，由命题 5(3) 得 $\langle x^d\rangle$ 是阶为 $a$ 的子群，故存在性成立。为证明唯一性，设 $K$ 是 $H$ 中任意一个阶为 $a$ 的子群。由 (1)，有
\[
  K=\langle x^b\rangle,
\]
其中 $b$ 是满足 $x^b\in K$ 的最小正整数。由命题 5，
\[
  \frac{n}{d}=a=|K|=|x^b|=\frac{n}{(n,b)},
\]
所以 $d=(n,b)$。特别地，$d\mid b$。由于 $b$ 是 $d$ 的倍数，$x^b\in\langle x^d\rangle$，因此
\[
  K=\langle x^b\rangle\leq \langle x^d\rangle.
\]
由于 $|\langle x^d\rangle|=a=|K|$，故 $K=\langle x^d\rangle$。

(3) 的最后断言来自以下观察：$\langle x^m\rangle$ 是 $\langle x^{(n,m)}\rangle$ 的子群（请检查），并且由命题 5(2) 和命题 2 可知二者阶相同。因为 $(n,m)$ 当然是 $n$ 的因子，这说明 $H$ 的每个子群都来自 $n$ 的某个因子。定理得证。
\end{proof}

\paragraph{例}
\begin{enumerate}[label=(\arabic*)]
  \item 可以用命题 6 和定理 7 列出任意给定 $n$ 的 $\ZZ/n\ZZ$ 的全部子群。例如，$\ZZ/12\ZZ$ 的子群为
  \begin{enumerate}[label=(\alph*)]
    \item $\ZZ/12\ZZ=\langle\overline{1}\rangle=\langle\overline{5}\rangle=\langle\overline{7}\rangle=\langle\overline{11}\rangle$（阶 $12$）；
    \item $\langle\overline{2}\rangle=\langle\overline{10}\rangle$（阶 $6$）；
    \item $\langle\overline{3}\rangle=\langle\overline{9}\rangle$（阶 $4$）；
    \item $\langle\overline{4}\rangle=\langle\overline{8}\rangle$（阶 $3$）；
    \item $\langle\overline{6}\rangle$（阶 $2$）；
    \item $\langle\overline{0}\rangle$（阶 $1$）。
  \end{enumerate}
  它们之间的包含关系由
  \[
    \langle\overline{a}\rangle\leq \langle\overline{b}\rangle
    \quad\Longleftrightarrow\quad
    (b,12)\mid (a,12),
    \qquad 1\leq a,b\leq 12
  \]
  给出。
  \item 我们还可以把本节结果与前一节的结果结合起来。例如，对群 $G$ 的每个 $x$，可构造子群 $C_G(\langle x\rangle)$ 和 $N_G(\langle x\rangle)$。容易检查，$G$ 中元素 $g$ 与 $x$ 交换当且仅当 $g$ 与 $x$ 的所有幂交换，因此
  \[
    C_G(\langle x\rangle)=C_G(x).
  \]
  如第 2 节练习 6 所述，$\langle x\rangle\leq N_G(\langle x\rangle)$，但等号不一定成立。例如，若 $G=Q_8$ 且 $x=i$，则
  \[
    C_G(\langle i\rangle)=\{\pm1,\pm i\}=\langle i\rangle,
    \qquad
    N_G(\langle i\rangle)=Q_8.
  \]
  注意，我们已经观察到了上面两个等式中的第一个，而第二个最容易用后面练习 24 的结果计算。
\end{enumerate}

\subsection*{练习}

\begin{enumerate}[label=\arabic*.]
  \item 求 $Z_{45}=\langle x\rangle$ 的全部子群，并给出每个子群的一个生成元。描述这些子群之间的包含关系。
  \item 若 $x$ 是有限群 $G$ 的元素且 $|x|=|G|$，证明 $G=\langle x\rangle$。给出一个具体例子说明若 $G$ 是无限群，则该结论不一定成立。
  \item 求 $\ZZ/48\ZZ$ 的所有生成元。
  \item 求 $\ZZ/202\ZZ$ 的所有生成元。
  \item 求 $\ZZ/49000\ZZ$ 的生成元个数。
  \item 在 $\ZZ/48\ZZ$ 中，对每个 $\overline{a}$ 写出 $\langle\overline{a}\rangle$ 的全部元素。求 $\ZZ/48\ZZ$ 中子群之间的所有包含关系。
  \item 令 $Z_{48}=\langle x\rangle$，并利用由 $\overline{1}\mapsto x$ 给出的同构 $\ZZ/48\ZZ\cong Z_{48}$，列出前一题计算得到的 $Z_{48}$ 的所有子群。
  \item 令 $Z_{48}=\langle x\rangle$。对哪些整数 $a$，由 $\varphi_a:\overline{1}\mapsto x^a$ 定义的映射 $\varphi_a$ 可扩张为从 $\ZZ/48\ZZ$ 到 $Z_{48}$ 的同构？
  \item 令 $Z_{36}=\langle x\rangle$。对哪些整数 $a$，由 $\psi_a:\overline{1}\mapsto x^a$ 定义的映射 $\psi_a$ 可扩张为从 $\ZZ/48\ZZ$ 到 $Z_{36}$ 的良定同态？$\psi_a$ 有可能是满同态吗？
  \item $\overline{30}$ 在 $\ZZ/54\ZZ$ 中的阶是多少？写出 $\langle\overline{30}\rangle$ 的全部元素及其阶。
  \item 求 $D_8$ 的所有循环子群。求 $D_8$ 的一个非循环的真子群。
  \item 证明下列群不是循环群：
  \[
    \text{(a) }Z_2\times Z_2,
    \qquad
    \text{(b) }Z_2\times\ZZ,
    \qquad
    \text{(c) }\ZZ\times\ZZ.
  \]
  \item 证明下列两对群不同构：
  \[
    \text{(a) }\ZZ\times Z_2\text{ 与 }\ZZ,
    \qquad
    \text{(b) }\QQ\times Z_2\text{ 与 }\QQ.
  \]
  \item 令 $\sigma=(1\ 2\ 3\ 4\ 5\ 6\ 7\ 8\ 9\ 10\ 11\ 12)$。对下列每个整数 $a$，计算 $\sigma^a$：
  \[
    a=13,65,626,1195,-6,-81,-570,-1211.
  \]
  \item 证明 $\QQ\times\QQ$ 不是循环群。
  \item 假设 $|x|=n$ 且 $|y|=m$。若 $x$ 与 $y$ 交换，即 $xy=yx$，证明 $|xy|$ 整除 $m$ 与 $n$ 的最小公倍数。若 $x$ 与 $y$ 不交换，该结论是否仍必然成立？给出一个交换元素 $x,y$ 的例子，使得 $xy$ 的阶不等于 $|x|$ 与 $|y|$ 的最小公倍数。
  \item 给出 $Z_n$ 的一个单生成元表示。
  \item 设 $H$ 是任意群，$h\in H$ 且 $h^n=1$。证明存在唯一同态从 $Z_n=\langle x\rangle$ 到 $H$，使得 $x\mapsto h$。
  \item 设 $H$ 是任意群，$h\in H$。证明存在唯一同态从 $\ZZ$ 到 $H$，使得 $1\mapsto h$。
  \item 设 $p$ 是素数，$n$ 是正整数。证明若 $x$ 是群 $G$ 中满足 $x^{p^n}=1$ 的元素，则 $|x|=p^m$，其中某个 $m\leq n$。
  \item 设 $p$ 是奇素数，$n$ 是正整数。用二项式定理证明
  \[
    (1+p)^{p^{n-1}}\equiv 1\pmod{p^n},
    \qquad
    (1+p)^{p^{n-2}}\not\equiv 1\pmod{p^n}.
  \]
  推出 $1+p$ 是乘法群 $(\ZZ/p^n\ZZ)^\times$ 中阶为 $p^{n-1}$ 的元素。
  \item 设 $n\geq3$。用二项式定理证明
  \[
    (1+2^2)^{2^{n-2}}\equiv1\pmod{2^n},
    \qquad
    (1+2^2)^{2^{n-3}}\not\equiv1\pmod{2^n}.
  \]
  推出 $5$ 是乘法群 $(\ZZ/2^n\ZZ)^\times$ 中阶为 $2^{n-2}$ 的元素。
  \item 证明对任意 $n\geq3$，$(\ZZ/2^n\ZZ)^\times$ 都不是循环群。[找出两个不同的阶为 $2$ 的子群。]
  \item 设 $G$ 是有限群，$x\in G$。
  \begin{enumerate}[label=(\alph*)]
    \item 证明若 $g\in N_G(\langle x\rangle)$，则对某个 $a\in\ZZ$ 有 $gxg^{-1}=x^a$。
    \item 反过来，证明若对某个 $a\in\ZZ$ 有 $gxg^{-1}=x^a$，则 $g\in N_G(\langle x\rangle)$。[先证明对任意整数 $k$，$gx^kg^{-1}=(gxg^{-1})^k=x^{ak}$，所以 $g\langle x\rangle g^{-1}\leq\langle x\rangle$。若 $x$ 的阶为 $n$，说明元素 $gx^ig^{-1}$，$i=0,1,\ldots,n-1$，彼此不同；于是 $|g\langle x\rangle g^{-1}|=|\langle x\rangle|=n$，并推出 $g\langle x\rangle g^{-1}=\langle x\rangle$。]
  \end{enumerate}
  注意，这可以减少计算循环子群正规化子时的一些工作，因为不必对每个 $h\in\langle x\rangle$ 都检查 $ghg^{-1}\in\langle x\rangle$。
  \item 设 $G$ 是阶为 $n$ 的循环群，$k$ 是与 $n$ 互素的整数。证明映射 $x\mapsto x^k$ 是满射。用拉格朗日定理（第 1.7 节练习 19）证明同样结论对任意阶为 $n$ 的有限群都成立。（对这样的 $k$，$G$ 的每个元素都有一个 $k$ 次根。由第 3.2 节的柯西定理可知，若 $k$ 与 $G$ 的阶不互素，则映射 $x\mapsto x^k$ 不是满射。）
  \item 设 $Z_n$ 是阶为 $n$ 的循环群。对每个整数 $a$，定义
  \[
    \sigma_a:Z_n\to Z_n,
    \qquad
    \sigma_a(x)=x^a\quad (x\in Z_n).
  \]
  \begin{enumerate}[label=(\alph*)]
    \item 证明 $\sigma_a$ 是 $Z_n$ 的自同构，当且仅当 $a$ 与 $n$ 互素（自同构见第 1.6 节练习 20）。
    \item 证明 $\sigma_a=\sigma_b$ 当且仅当 $a\equiv b\pmod n$。
    \item 证明 $Z_n$ 的每个自同构都等于某个 $\sigma_a$。
    \item 证明 $\sigma_a\circ\sigma_b=\sigma_{ab}$。推出映射 $\overline{a}\mapsto \sigma_a$ 是从 $(\ZZ/n\ZZ)^\times$ 到 $Z_n$ 的自同构群的同构（因此 $\Aut(Z_n)$ 是阶为 $\varphi(n)$ 的阿贝尔群）。
  \end{enumerate}
\end{enumerate}

\section{由群的子集生成的子群}

构造给定群的循环子群的方法，是一般技术的一个特殊情形：由群的任意子集生成子群。在循环子群的情形，我们取群 $G$ 的单元素子集 $\{x\}$，并形成 $x$ 的所有整数幂；这相当于把集合 $\{x\}$ 对群运算以及取逆过程封闭起来。所得子群是 $G$ 中包含 $\{x\}$ 的最小子群（“最小”的意思是：若 $H$ 是任意包含 $\{x\}$ 的子群，则 $H$ 包含 $\langle x\rangle$）。换句话说，$\langle x\rangle$ 是所有包含 $x$ 的子群所组成的集合（按包含关系排序）中的唯一极小元。本节研究当 $\{x\}$ 换成 $G$ 的任意子集时的类似问题。

在整个数学中，下面这个主题反复出现：给定一个对象 $G$（例如群、域、向量空间等）和 $G$ 的一个子集 $A$，是否存在唯一最小的包含 $A$ 的 $G$ 的子对象（子群、子域、子空间等）；若存在，该子对象的元素如何计算？学生可能已经在向量空间的学习中遇到过这个问题。当 $G$ 是向量空间（比如标量为实数）且 $A=\{v_1,v_2,\ldots,v_n\}$ 时，存在唯一最小的包含 $A$ 的子空间，即 $v_1,v_2,\ldots,v_n$ 的（线性）张成空间；该张成空间中的每个向量都能写成
\[
  k_1v_1+k_2v_2+\cdots+k_nv_n
\]
的形式，其中 $k_1,\ldots,k_n\in\RR$。当 $A$ 是单个非零向量 $v$ 时，$\{v\}$ 的张成空间就是包含 $v$ 的一维子空间或直线，该子空间的每个元素都形如 $kv$，其中 $k\in\RR$。这是向量空间理论中循环子群的类比。注意，一维子空间包含 $kv$，其中 $k\in\RR$，而不仅仅包含 $k\in\ZZ$ 时的 $kv$；原因是子空间必须对所有向量空间运算（例如标量乘法）封闭，而不只是对向量加法这一群运算封闭。

设 $G$ 是任意群，$A$ 是 $G$ 的任意子集。我们现在精确给出由 $A$ 生成的 $G$ 的子群的概念。我们证明：因为 $G$ 的任意子群族的交仍是 $G$ 的子群，由 $A$ 生成的子群就是 $G$ 中包含 $A$ 的唯一最小子群；它在“最小”的意义下，是所有包含 $A$ 的子群组成的集合中的极小元。我们还会说明，该子群的元素是通过将给定子集对群运算（及取逆）封闭得到的。以后在本书后面部分发展其他代数对象理论时，我们将把本节作为范例，来证明某个给定子集被包含于唯一最小的子对象中，并且该子对象的元素由将该子集对定义对象的运算封闭而得到。由于在后续章节中细节会被省略，学生应在此处扎实理解这一过程。

为继续进行，我们只需要下面这个命题。

\medskip
\noindent\textbf{命题 8.} 若 $\mcA$ 是 $G$ 的任意非空子群族，则 $\mcA$ 中所有成员的交仍是 $G$ 的子群。

\begin{proof}
这是子群判别法的一个简单应用（另见第 1 节练习 10）。令
\[
  K=\bigcap_{H\in\mcA}H.
\]
由于每个 $H\in\mcA$ 都是子群，$1\in H$，所以 $1\in K$，即 $K\neq\varnothing$。若 $a,b\in K$，则对所有 $H\in\mcA$，有 $a,b\in H$。由于每个 $H$ 是群，$ab^{-1}\in H$ 对所有 $H$ 成立，所以 $ab^{-1}\in K$。由命题 1，$K\leq G$。
\end{proof}

\noindent\textbf{定义.} 若 $A$ 是群 $G$ 的任意子集，定义
\[
  \langle A\rangle=\bigcap_{\substack{A\subseteq H\\ H\leq G}}H.
\]
这称为由 $A$ \emph{生成}的 $G$ 的子群。

因此，$\langle A\rangle$ 是所有包含 $A$ 的 $G$ 的子群的交。由命题 8 应用于集合
\[
  \mcA=\{H\leq G\mid A\subseteq H\}
\]
可知它是 $G$ 的子群（$\mcA$ 非空，因为 $G\in\mcA$）。由于 $A$ 属于每个 $H\in\mcA$，$A$ 是它们交 $\langle A\rangle$ 的子集。注意 $\langle A\rangle$ 是 $\mcA$ 的唯一极小元：$\langle A\rangle$ 是含有 $A$ 的 $G$ 的子群，因此 $\langle A\rangle\in\mcA$；而 $\mcA$ 的任意元素都包含 $\mcA$ 中所有元素的交，即包含 $\langle A\rangle$。

当 $A$ 是有限集 $\{a_1,a_2,\ldots,a_n\}$ 时，我们写 $\langle a_1,a_2,\ldots,a_n\rangle$ 表示由 $a_1,a_2,\ldots,a_n$ 生成的群，而不写 $\langle\{a_1,a_2,\ldots,a_n\}\rangle$。若 $A$ 与 $B$ 是 $G$ 的两个子集，我们写 $\langle A,B\rangle$ 代替 $\langle A\cup B\rangle$。

这种“自上而下”的方式定义 $\langle A\rangle$，证明了包含 $A$ 的最小子群的存在性和唯一性，但并不太说明如何构造其中的元素。正如“生成”一词所暗示的，我们现在定义一个把 $A$ 对群运算（及取逆过程）封闭起来的集合，并证明该集合等于 $\langle A\rangle$。令
\[
  \overline{A}=\{a_1^{\epsilon_1}a_2^{\epsilon_2}\cdots a_n^{\epsilon_n}\mid
  n\in\ZZ,
  n\geq0,
  a_i\in A,
  \epsilon_i=\pm1\text{ 对每个 }i\},
\]
其中若 $A=\varnothing$ 则令 $\overline{A}=\{1\}$。于是 $\overline{A}$ 是 $A$ 的元素及其逆的所有有限乘积（称为\emph{字}）的集合。注意 $a_i$ 不必彼此不同，所以在定义 $\overline{A}$ 的记号中，$a^2$ 写成 $aa$。还要注意，不假定 $A$ 是有限集（甚至不假定可数）。

\medskip
\noindent\textbf{命题 9.} $\overline{A}=\langle A\rangle$。

\begin{proof}
先证明 $\overline{A}$ 是子群。注意 $\overline{A}\neq\varnothing$（即使 $A=\varnothing$ 也是如此）。若 $a,b\in\overline{A}$，且
\[
  a=a_1^{\epsilon_1}a_2^{\epsilon_2}\cdots a_n^{\epsilon_n},
  \qquad
  b=b_1^{\delta_1}b_2^{\delta_2}\cdots b_m^{\delta_m},
\]
则
\[
  ab^{-1}=a_1^{\epsilon_1}a_2^{\epsilon_2}\cdots a_n^{\epsilon_n}
  b_m^{-\delta_m}b_{m-1}^{-\delta_{m-1}}\cdots b_1^{-\delta_1}
\]
（这里用第 1.1 节练习 15 计算 $b^{-1}$）。因此 $ab^{-1}$ 是 $A$ 的元素的 $\pm1$ 次幂的乘积，故 $ab^{-1}\in\overline{A}$。由命题 1，$\overline{A}$ 是 $G$ 的子群。

由于每个 $a\in A$ 都可写成 $a^1$，有 $A\subseteq\overline{A}$，从而 $\langle A\rangle\leq\overline{A}$。另一方面，$\langle A\rangle$ 是含有 $A$ 的群，并且对群运算和取逆封闭，所以 $\langle A\rangle$ 含有每个形如 $a_1^{\epsilon_1}a_2^{\epsilon_2}\cdots a_n^{\epsilon_n}$ 的元素，即 $\overline{A}\subseteq\langle A\rangle$。命题得证。
\end{proof}

现在我们用 $\langle A\rangle$ 代替 $\overline{A}$，并可把 $\overline{A}$ 的定义作为 $\langle A\rangle$ 的等价定义。如上所述，在 $\langle A\rangle$ 的这个等价定义中，形如 $a\cdot a$、$a\cdot a\cdot a$、$a\cdot a^{-1}$ 等乘积可分别化简为 $a^2$、$a^3$、$1$ 等。因此另一种写法是
\[
  \langle A\rangle=
  \{a_1^{\alpha_1}a_2^{\alpha_2}\cdots a_n^{\alpha_n}
  \mid a_i\in A,
  \alpha_i\in\ZZ,
  a_i\neq a_{i+1},
  n\in\ZZ^+\}.
\]
事实上，当 $A=\{x\}$ 时，这就是我们对 $\langle A\rangle$ 的定义。

若 $G$ 是阿贝尔群，我们可以交换各个 $a_i$，于是把同一生成元的所有幂收集在一起。例如，若 $A$ 是阿贝尔群 $G$ 的有限子集 $\{a_1,a_2,\ldots,a_k\}$，则容易检查
\[
  \langle A\rangle=\{a_1^{\alpha_1}a_2^{\alpha_2}\cdots a_k^{\alpha_k}
  \mid \alpha_i\in\ZZ\text{ 对每个 }i\}.
\]
若在这种情形下进一步假设每个 $a_i$ 的阶为有限数 $d_i$，则由于 $a_i$ 恰有 $d_i$ 个不同幂，形如 $a_1^{\alpha_1}a_2^{\alpha_2}\cdots a_k^{\alpha_k}$ 的不同乘积总数至多为 $d_1d_2\cdots d_k$，也就是说
\[
  |\langle A\rangle|\leq d_1d_2\cdots d_k.
\]
可能发生 $a^\alpha b^\beta=a^\gamma b^\delta$，尽管 $a^\alpha\neq a^\gamma$ 且 $b^\beta\neq b^\delta$。在第 5 章研究直积时，我们会精确考察这种现象何时发生。

当 $G$ 非阿贝尔时，情况复杂得多。例如，令 $G=D_8$，并令 $r,s$ 是 $D_8$ 的通常生成元（注意记号 $D_8=\langle r,s\rangle$ 与第 1.2 节引入的记号一致）。令 $a=s$，$b=rs$，并令 $A=\{a,b\}$。由于 $s$ 和 $r=rs\cdot s$ 都属于 $\langle a,b\rangle$，有 $G=\langle a,b\rangle$，即 $G$ 也由 $a$ 与 $b$ 生成。$a$ 与 $b$ 都有阶 $2$，但 $D_8$ 有阶 $8$。这意味着不可能把 $D_8$ 的每个元素都写成 $a^\alpha b^\beta$ 的形式，其中 $\alpha,\beta\in\ZZ$。更具体地，乘积 $aba$ 不能化简为形如 $a^\alpha b^\beta$ 的乘积。事实上，若对任意 $n>2$ 有 $G=D_{2n}$，并按上面同样方式定义 $r,s,a,b$，仍有
\[
  |a|=|b|=2,
  \qquad
  D_{2n}=\langle a,b\rangle,
  \qquad
  |D_{2n}|=2n.
\]
这意味着当 $n$ 很大时，形如 $abab\cdots ab$ 的长乘积不能进一步化简。特别地，这说明与阿贝尔群（更不用说循环群）情形不同，即使知道某个生成集里元素的阶，也无法给出一个有限群阶的上界。

这种现象的另一个例子是 $S_n$：
\[
  S_n=\langle (1\ 2),(1\ 2\ 3\ \cdots\ n)\rangle.
\]
因此 $S_n$ 由一个阶为 $2$ 的元素和一个阶为 $n$ 的元素生成，但 $|S_n|=n!$（在发展更多技术之后，我们将在后面证明这些陈述）。

最后一个例子强调，若 $G$ 非阿贝尔，则由 $G$ 中不止一个元素生成的子群可能相当复杂。令
\[
  G=\GL_2(\RR),
  \qquad
  a=\begin{pmatrix}0&1\\1&0\end{pmatrix},
  \qquad
  b=\begin{pmatrix}0&2\\1/2&0\end{pmatrix}.
\]
则 $a^2=b^2=1$，但
\[
  ab=\begin{pmatrix}1/2&0\\0&2\end{pmatrix}.
\]
容易看出 $ab$ 有无限阶，所以 $\langle a,b\rangle$ 是 $\GL_2(\RR)$ 的一个无限子群，而它由两个阶为 $2$ 的元素生成。

这些例子说明，当 $|A|\geq2$ 时，一般来说，甚至计算由 $A$ 生成的子群的阶都很困难，更不用说其他结构性质。因此，对非阿贝尔群中任意取的子集 $A$，试图写出 $\langle A\rangle$ 的元素（或关于它的其他信息）通常并不实际。对非阿贝尔群 $G$ 的某些“选得好”的子集 $A$，我们仍然能够在理论和计算中利用由 $A$ 生成的子群。一个例子是，当我们想找一个真正包含 $\langle x\rangle$ 的 $G$ 的子群时，可以寻找某个与 $x$ 交换的元素 $y$（即 $y\in C_G(x)$），并构成 $\langle x,y\rangle$。容易检查后者是阿贝尔群，所以其阶由 $|x||y|$ 约束。或者，我们可以取 $y\in N_G(\langle x\rangle)$；在这种情形下，同样的阶上界成立，并且 $\langle x,y\rangle$ 的结构也不太复杂（正如下一章将看到的）。

非阿贝尔群中出现的复杂性，在研究其他基本代数系统时通常不会这么严重，因为额外的代数结构会施加更多限制。

\subsection*{练习}

\begin{enumerate}[label=\arabic*.]
  \item 证明若 $H$ 是 $G$ 的子群，则 $\langle H\rangle=H$。
  \item 证明若 $A$ 是 $B$ 的子集，则 $\langle A\rangle\leq\langle B\rangle$。给出一个例子，其中 $A\subseteq B$ 且 $A\neq B$，但 $\langle A\rangle=\langle B\rangle$。
  \item 证明若 $H$ 是群 $G$ 的阿贝尔子群，则 $\langle H,Z(G)\rangle$ 是阿贝尔群。给出一个具体例子，说明存在群 $G$ 的阿贝尔子群 $H$，使得 $\langle H,C_G(H)\rangle$ 不是阿贝尔群。
  \item 证明若 $H$ 是 $G$ 的子群，则 $H$ 由集合 $H-\{1\}$ 生成。
  \item 证明 $S_3$ 中任意两个不同的阶为 $2$ 的元素生成的子群都是整个 $S_3$。
  \item 证明 $S_4$ 中由 $(1\ 2)$ 和 $(1\ 2)(3\ 4)$ 生成的子群是一个阶为 $4$ 的非循环群。
  \item 证明 $S_4$ 中由 $(1\ 2)$ 和 $(1\ 3)(2\ 4)$ 生成的子群同构于阶为 $8$ 的二面体群。
  \item 证明 $S_4=\langle(1\ 2\ 3\ 4),(1\ 2\ 4\ 3)\rangle$。
  \item 证明 $\SL_2(\FF_3)$ 是 $\GL_2(\FF_3)$ 中由
  \[
    \begin{pmatrix}1&1\\0&1\end{pmatrix}
    \quad\text{和}\quad
    \begin{pmatrix}1&0\\1&1\end{pmatrix}
  \]
  生成的子群。[回忆第 1 节练习 9，$\SL_2(\FF_3)$ 是行列式为 $1$ 的矩阵子群。可假设该子群阶为 $24$——这将是第 3.2 节的一个练习。]
  \item 证明 $\SL_2(\FF_3)$ 中由
  \[
    \begin{pmatrix}0&1\\-1&0\end{pmatrix}
    \quad\text{和}\quad
    \begin{pmatrix}0&1\\1&0\end{pmatrix}
  \]
  生成的子群同构于阶为 $8$ 的四元数群。[使用 $Q_8$ 的一个表示。]
  \item 说明 $\SL_2(\FF_3)$ 与 $S_4$ 是两个阶为 $24$ 的不同构群。
  \item 证明 $\GL_3(\FF_2)$ 中上三角矩阵组成的子群同构于阶为 $8$ 的二面体群（参见第 1 节练习 16）。[先求出该子群的阶。]
  \item 证明正有理数的乘法群由集合
  \[
    \{1/p\mid p\text{ 是素数}\}
  \]
  生成。
  \item 若存在有限集 $A$ 使 $H=\langle A\rangle$，则称群 $H$ \emph{有限生成}。
  \begin{enumerate}[label=(\alph*)]
    \item 证明每个有限群都是有限生成的。
    \item 证明 $\ZZ$ 是有限生成的。
    \item 证明加法群 $\QQ$ 的每个有限生成子群都是循环的。[若 $H$ 是 $\QQ$ 的有限生成子群，说明 $H\leq\langle 1/k\rangle$，其中 $k$ 是 $H$ 的一组生成元中出现的所有分母的乘积。]
    \item 证明 $\QQ$ 不是有限生成的。
  \end{enumerate}
  \item 给出 $\QQ$ 的一个非循环真子群。
  \item 群 $G$ 的子群 $M$ 称为\emph{极大子群}，如果 $M\neq G$，且 $G$ 中包含 $M$ 的子群只有 $M$ 和 $G$。
  \begin{enumerate}[label=(\alph*)]
    \item 证明若 $H$ 是有限群 $G$ 的真子群，则存在 $G$ 的一个极大子群包含 $H$。
    \item 说明二面体群中所有旋转组成的子群是极大子群。
    \item 说明若 $G=\langle x\rangle$ 是阶 $n\geq1$ 的循环群，则子群 $H$ 是极大子群，当且仅当 $H=\langle x^p\rangle$，其中 $p$ 是整除 $n$ 的某个素数。
  \end{enumerate}
  \item 这是一个关于佐恩引理（见附录 I）的练习，用于证明每个非平凡有限生成群都有极大子群。设 $G$ 是有限生成群，比如 $G=\langle g_1,g_2,\ldots,g_n\rangle$，令 $\mcS$ 为 $G$ 的所有真子群组成的集合。则 $\mcS$ 按包含关系偏序。令 $\mcC$ 是 $\mcS$ 中的一条链。
  \begin{enumerate}[label=(\alph*)]
    \item 证明 $\mcC$ 中所有子群的并 $H$ 是 $G$ 的子群。
    \item 证明 $H$ 是真子群。[若不是，则每个 $g_i$ 都在 $H$ 中，因此都在链 $\mcC$ 的某个元素中。利用链的定义得到矛盾。]
    \item 用佐恩引理说明 $\mcS$ 有极大元（按定义即极大子群）。
  \end{enumerate}
  \item 设 $p$ 是素数，令
  \[
    Z=\{z\in\CC\mid z^{p^n}=1\text{ 对某个 }n\in\ZZ^+\}
  \]
  （所以 $Z$ 是 $\CC$ 中所有 $p$ 次幂阶单位根组成的乘法群）。对每个 $k\in\ZZ^+$，令
  \[
    H_k=\{z\in Z\mid z^{p^k}=1\}
  \]
  （即 $p^k$ 次单位根群）。证明：
  \begin{enumerate}[label=(\alph*)]
    \item $H_k\leq H_m$ 当且仅当 $k\leq m$；
    \item 每个 $H_k$ 都是循环群（可假设对任意 $n\in\ZZ^+$，$\{e^{2\pi it/n}\mid t=0,1,\ldots,n-1\}$ 是 $\CC$ 中所有 $n$ 次单位根的集合）；
    \item $Z$ 的每个真子群都等于某个 $H_k$（特别地，$Z$ 的每个真子群都是有限循环群）；
    \item $Z$ 不是有限生成的。
  \end{enumerate}
  \item 非平凡阿贝尔群 $A$（写作乘法记号）称为\emph{可除的}，如果对每个 $a\in A$ 和每个非零整数 $k$，都存在 $x\in A$ 使得 $x^k=a$；也就是说，$A$ 的每个元素在 $A$ 中都有 $k$ 次根（在加法记号中，每个元素都是某个元素的 $k$ 倍）。
  \begin{enumerate}[label=(\alph*)]
    \item 证明有理数加法群 $\QQ$ 是可除的。
    \item 证明没有有限阿贝尔群是可除的。
  \end{enumerate}
  \item 证明若 $A$ 与 $B$ 是非平凡阿贝尔群，则 $A\times B$ 可除当且仅当 $A$ 与 $B$ 都是可除群。
\end{enumerate}


\section{群的子群格}

本节描述与一个群相关的图，它描绘该群各子群之间的关系。这个图称为该群的\emph{子群格}\footnote{“格”一词在偏序集理论中有精确的数学含义。}，是一种“可视化”群的好方法——它显然比群表更能揭示群的结构。在群论各章中，我们将用格图，或其中的部分，描述具体群以及一般群的某些性质。此外，群的子群格将在伽罗瓦理论中起重要作用。

给定有限群 $G$ 的子群格构造如下：把 $G$ 的所有子群画出来，底部是 $1$，顶部是 $G$，粗略地说，阶较大的子群在页面上放得比阶较小的子群高。按如下规则在子群之间向上画路径：若 $A\leq B$，并且没有子群真地位于 $A$ 与 $B$ 之间，则从 $A$ 向上画一条线到 $B$。因此，若 $A\leq B$，则从 $A$ 到 $B$ 有一条（可能有多条）向上的路径，经过一串中间子群；若 $B\geq A$，则也有从 $B$ 到 $A$ 的向下路径。子群在页面上的初始位置原则上有些任意，但通过练习通常可以选择得使图形简单。

注意，对于 $G$ 的任意两个子群 $H$ 和 $K$，包含二者的唯一最小子群，即 $\langle H,K\rangle$（称为 $H$ 与 $K$ 的\emph{上确界}），可从格中读出如下：从 $H$ 和 $K$ 沿路径向上追踪，直到到达一个同时包含 $H$ 和 $K$ 的公共子群 $A$（注意 $G$ 本身总包含所有子群，所以至少存在这样的 $A$）。要保证 $A=\langle H,K\rangle$，需确认没有 $A_1\leq A$（在图中表现为从 $A$ 到 $A_1$ 的向下路径）同时包含 $H$ 和 $K$；否则把 $A$ 换成 $A_1$，并重复这一过程，查看是否 $A_1=\langle H,K\rangle$。对称地，也可读出 $G$ 中同时包含于 $H$ 与 $K$ 的最大子群，即它们的交（由命题 8，交是子群）。

这一过程有一些限制，特别是它本身不能直接用于无限群。即使对阶相对较小的有限群，格也可能相当复杂（可参见 M. Hall 和 J. Senior 的 \textit{Groups of Order $2^n$, $n\leq6$}，Macmillan，1964，其中有一些令人头痛的例子）。本节末尾将说明，即使对无限群，也可画出并使用格的一部分。

注意，同构群有相同的格（即相同的有向图）。不同构群也可能有相同的格（这会发生在两个阶为 $16$ 的群中——见后面的练习）。由于子群格只是我们描述群时携带的信息的一部分，这并不是严重缺陷（事实上，在观察两个不同构群何时具有某些共同性质时，它甚至可能有用）。

\paragraph{例}
除循环群（例 1）外，我们尚未证明下面的格都是正确的（例如，包含了给定群的全部子群，或具有正确的上确界与交）。目前我们先把这些事实作为已知；随着本书理论的建立，我们将把证明这些格确实正确作为练习。

\begin{enumerate}[label=(\arabic*)]
  \item 对 $G=Z_n\cong \ZZ/n\ZZ$，由定理 7，$G$ 的子群格就是 $n$ 的因子格（即把 $n$ 的因子写在页面上，底部为 $n$，顶部为 $1$，若 $b\mid a$ 则从 $a$ 向上连到 $b$）。下面是若干 $n$ 的具体例子。

  \begin{center}
  \begin{minipage}{0.32\linewidth}
  \centering
  \begin{tikzpicture}[baseline=(current bounding box.center), every node/.style={inner sep=1pt}]
    \node (a) at (0,0.8) {$\ZZ/2\ZZ=\langle \overline{1}\rangle$};
    \node (b) at (0,-0.2) {$\langle\overline{2}\rangle=\{0\}$};
    \draw (a)--(b);
  \end{tikzpicture}
  \end{minipage}
  \begin{minipage}{0.42\linewidth}
  \centering
  \begin{tikzpicture}[baseline=(current bounding box.center), every node/.style={inner sep=1pt}]
    \node (a) at (0,1.2) {$\ZZ/4\ZZ=\langle \overline{1}\rangle$};
    \node (note) at (2.5,1.2) {\small（注意：$\langle\overline{1}\rangle=\langle\overline{3}\rangle$）};
    \node (b) at (0,0.1) {$\langle\overline{2}\rangle$};
    \node (c) at (0,-1.0) {$\langle\overline{4}\rangle=\{0\}$};
    \draw (a)--(b)--(c);
  \end{tikzpicture}
  \end{minipage}
  \end{center}

  \begin{center}
  \begin{tikzpicture}[every node/.style={inner sep=1pt}]
    \node (a) at (0,1.8) {$\ZZ/8\ZZ=\langle\overline{1}\rangle$};
    \node (note) at (3.8,1.8) {\small（注意：$\langle\overline{1}\rangle=\langle\overline{3}\rangle=\langle\overline{5}\rangle=\langle\overline{7}\rangle$）};
    \node (b) at (0,0.7) {$\langle\overline{2}\rangle$};
    \node (c) at (0,-0.4) {$\langle\overline{4}\rangle$};
    \node (d) at (0,-1.5) {$\langle\overline{8}\rangle=\{0\}$};
    \draw (a)--(b)--(c)--(d);
  \end{tikzpicture}
  \end{center}

  一般地，若 $p$ 是素数，则 $\ZZ/p^n\ZZ$ 的格为
  \begin{center}
  \begin{tikzpicture}[every node/.style={inner sep=1pt}]
    \node (a) at (0,2.8) {$\ZZ/p^n\ZZ=\langle\overline{1}\rangle$};
    \node (b) at (0,1.95) {$\langle\overline{p}\rangle$};
    \node (c) at (0,1.1) {$\langle\overline{p^2}\rangle$};
    \node (d) at (0,0.25) {$\langle\overline{p^3}\rangle$};
    \node (e) at (0,-0.55) {$\vdots$};
    \node (f) at (0,-1.35) {$\langle\overline{p^{n-1}}\rangle$};
    \node (g) at (0,-2.2) {$\langle\overline{p^n}\rangle=\{0\}$};
    \draw (a)--(b)--(c)--(d)--(e)--(f)--(g);
  \end{tikzpicture}
  \end{center}

  \begin{center}
  \begin{minipage}{0.45\linewidth}
  \centering
  \begin{tikzpicture}[scale=.95, every node/.style={inner sep=1pt}]
    \node (top) at (0,2) {$\ZZ/6\ZZ$};
    \node (a) at (-1.3,1) {$\langle\overline{2}\rangle$};
    \node (b) at (1.3,0.3) {$\langle\overline{3}\rangle$};
    \node (bot) at (0,-1.2) {$\langle\overline{6}\rangle=\{0\}$};
    \draw (top)--(a)--(bot)--(b)--(top);
  \end{tikzpicture}
  \end{minipage}
  \begin{minipage}{0.45\linewidth}
  \centering
  \begin{tikzpicture}[scale=.95, every node/.style={inner sep=1pt}]
    \node (top) at (0,2.2) {$\ZZ/12\ZZ$};
    \node (a) at (-1.4,1.2) {$\langle\overline{2}\rangle$};
    \node (b) at (1.5,.4) {$\langle\overline{3}\rangle$};
    \node (c) at (-2.4,0) {$\langle\overline{4}\rangle$};
    \node (d) at (0,-.8) {$\langle\overline{6}\rangle$};
    \node (bot) at (-1.35,-2.0) {$\langle\overline{12}\rangle=\{0\}$};
    \draw (top)--(a)--(c)--(bot)--(d)--(b)--(top);
    \draw (a)--(d);
  \end{tikzpicture}
  \end{minipage}
  \end{center}

  \item \emph{克莱因四元群}（Viergruppe）$V_4$ 是阶为 $4$ 的群，其乘法表与格为
  \begin{center}
  \begin{tabular}{c@{\qquad}c}
  $\begin{array}{c|cccc}
    \cdot&1&a&b&c\\ \hline
    1&1&a&b&c\\
    a&a&1&c&b\\
    b&b&c&1&a\\
    c&c&b&a&1
  \end{array}$
  &
  \begin{tikzpicture}[baseline=(current bounding box.center), scale=.9, every node/.style={inner sep=1pt}]
    \node (top) at (0,2) {$V_4$};
    \node (a) at (-1.4,.8) {$\langle a\rangle$};
    \node (b) at (0,.8) {$\langle b\rangle$};
    \node (c) at (1.4,.8) {$\langle c\rangle$};
    \node (bot) at (0,-.6) {$1$};
    \draw (top)--(a)--(bot)--(c)--(top);
    \draw (top)--(b)--(bot);
  \end{tikzpicture}
  \end{tabular}
  \end{center}
  注意 $V_4$ 是阿贝尔群，且不同构于 $Z_4$（为什么？）。我们会看到 $D_8$ 含有一个同构于 $V_4$ 的子群，因此不必检查上述定义的二元运算满足结合律。

  \item $S_3$ 的格为
  \begin{center}
  \begin{tikzpicture}[scale=1.05, every node/.style={inner sep=1pt}]
    \node (top) at (0,2.2) {$S_3$};
    \node (a) at (-2.7,.6) {$\langle(1\ 2)\rangle$};
    \node (b) at (-1.2,.6) {$\langle(1\ 3)\rangle$};
    \node (c) at (.3,.6) {$\langle(2\ 3)\rangle$};
    \node (d) at (2.3,1.4) {$\langle(1\ 2\ 3)\rangle$};
    \node (bot) at (1.2,-.5) {$1$};
    \foreach \x in {a,b,c,d} {\draw (top)--(\x); \draw (\x)--(bot);}
  \end{tikzpicture}
  \end{center}

  \item 用通常记号 $D_8=\langle r,s\rangle$，$D_8$ 的格为
  \begin{center}
  \begin{tikzpicture}[scale=1.05, every node/.style={inner sep=1pt}]
    \node (top) at (0,3) {$D_8$};
    \node (a) at (-1.7,1.8) {$\langle s,r^2\rangle$};
    \node (b) at (0,1.8) {$\langle r\rangle$};
    \node (c) at (1.7,1.8) {$\langle rs,r^2\rangle$};
    \node (d) at (-3,.5) {$\langle s\rangle$};
    \node (e) at (-1.5,.5) {$\langle r^2s\rangle$};
    \node (f) at (0,.5) {$\langle r^2\rangle$};
    \node (g) at (1.5,.5) {$\langle rs\rangle$};
    \node (h) at (3,.5) {$\langle r^3s\rangle$};
    \node (bot) at (0,-.8) {$1$};
    \draw (top)--(a)--(d)--(bot);
    \draw (top)--(b)--(f)--(bot);
    \draw (top)--(c)--(h)--(bot);
    \draw (a)--(e)--(bot);
    \draw (a)--(f);
    \draw (c)--(f);
    \draw (c)--(g)--(bot);
  \end{tikzpicture}
  \end{center}

  \item $Q_8$ 的子群格为
  \begin{center}
  \begin{tikzpicture}[scale=1.05, every node/.style={inner sep=1pt}]
    \node (top) at (0,2.8) {$Q_8$};
    \node (a) at (-1.4,1.6) {$\langle i\rangle$};
    \node (b) at (0,1.6) {$\langle j\rangle$};
    \node (c) at (1.4,1.6) {$\langle k\rangle$};
    \node (d) at (0,.3) {$\langle -1\rangle$};
    \node (bot) at (0,-1) {$1$};
    \draw (top)--(a)--(d)--(bot);
    \draw (top)--(b)--(d);
    \draw (top)--(c)--(d);
  \end{tikzpicture}
  \end{center}

  \item $D_{16}$ 的子群格不是平面图（不能在平面上画出而不使线相交）。它的一种画法是
  \begin{center}
  \begin{tikzpicture}[scale=.9, every node/.style={inner sep=1pt}]
    \node (top) at (0,5) {$D_{16}$};
    \node (sr2) at (-2.6,3.8) {$\langle s,r^2\rangle$};
    \node (r) at (0,3.8) {$\langle r\rangle$};
    \node (srr2) at (2.6,3.8) {$\langle sr,r^2\rangle$};
    \node (sr2r4) at (-4.2,2.3) {$\langle sr^2,r^4\rangle$};
    \node (sr4) at (-2,2.3) {$\langle s,r^4\rangle$};
    \node (r2) at (0,2.3) {$\langle r^2\rangle$};
    \node (sr3r4) at (2,2.3) {$\langle sr^3,r^4\rangle$};
    \node (sr5r4) at (4.2,2.3) {$\langle sr^5,r^4\rangle$};
    \node (sr6) at (-5.1,.8) {$\langle sr^6\rangle$};
    \node (sr2s) at (-3.7,.8) {$\langle sr^2\rangle$};
    \node (sr4s) at (-2.3,.8) {$\langle sr^4\rangle$};
    \node (ss) at (-1.1,.8) {$\langle s\rangle$};
    \node (r4) at (0,.8) {$\langle r^4\rangle$};
    \node (sr3) at (1.3,.8) {$\langle sr^3\rangle$};
    \node (sr7) at (2.6,.8) {$\langle sr^7\rangle$};
    \node (sr5) at (3.8,.8) {$\langle sr^5\rangle$};
    \node (sr) at (5.1,.8) {$\langle sr\rangle$};
    \node (bot) at (0,-.7) {$1$};
    \draw (top)--(sr2) (top)--(r) (top)--(srr2);
    \draw (sr2)--(sr2r4) (sr2)--(sr4) (sr2)--(r2);
    \draw (r)--(r2);
    \draw (srr2)--(r2) (srr2)--(sr3r4) (srr2)--(sr5r4);
    \draw (sr2r4)--(sr6) (sr2r4)--(sr2s) (sr2r4)--(r4);
    \draw (sr4)--(sr4s) (sr4)--(ss) (sr4)--(r4);
    \draw (r2)--(r4);
    \draw (sr3r4)--(r4) (sr3r4)--(sr3) (sr3r4)--(sr7);
    \draw (sr5r4)--(r4) (sr5r4)--(sr5) (sr5r4)--(sr);
    \foreach \x in {sr6,sr2s,sr4s,ss,r4,sr3,sr7,sr5,sr} {\draw (\x)--(bot);}
  \end{tikzpicture}
  \end{center}
\end{enumerate}

在理论证明和具体例子中，我们常常只关心给定群的两个（或少数几个）子群及其相互关系。为了把这些关系画出来，我们将画出整个群格的一个\emph{子格}，其中包含相关的上确界与交。这样的子格中的一条不断开的线一般并不表示该线两端之间没有其他子群。这些部分格也会在处理无限群时使用。例如，若只想讨论 $D_{16}$ 的子群 $\langle sr^2,r^4\rangle$ 与 $\langle r^2\rangle$ 的关系，可画出如下子格：
\begin{center}
\begin{tikzpicture}[scale=1.0, every node/.style={inner sep=1pt}]
  \node (top) at (0,4) {$D_{16}$};
  \node (a) at (0,3) {$\langle s,r^2\rangle$};
  \node (b) at (-1.5,1.8) {$\langle sr^2,r^4\rangle$};
  \node (c) at (1.5,1.8) {$\langle r^2\rangle$};
  \node (d) at (0,.6) {$\langle r^4\rangle$};
  \node (bot) at (0,-.5) {$1$};
  \draw (top)--(a)--(b)--(d)--(bot);
  \draw (a)--(c)--(d);
\end{tikzpicture}
\end{center}
注意 $\langle s,r^2\rangle$ 和 $\langle r^4\rangle$ 分别正是这两个子群在 $D_{16}$ 中的上确界与交。

最后，给定一个群的子群格后，计算正规化子和中心化子相对容易。例如，在 $D_8$ 中，可以看出 $C_{D_8}(s)=\langle s,r^2\rangle$，因为先计算得 $r^2\in C_{D_8}(s)$（见第 2 节）。这证明 $\langle s,r^2\rangle\leq C_{D_8}(s)$（注意一个元素总属于它自身的中心化子）。包含 $\langle s,r^2\rangle$ 的子群只有它自身和整个 $D_8$。不可能有 $C_{D_8}(s)=D_8$，因为 $r$ 不与 $s$ 交换（即 $r\notin C_{D_8}(s)$）。于是只剩所声称的一种可能。

\subsection*{练习}

\begin{enumerate}[label=\arabic*.]
  \item 设 $H$ 和 $K$ 是 $G$ 的子群。列出所有可能的只显示 $G$、$1$、$H$、$K$ 及其上确界与交的子格。不同图形的区别是什么？
  \item 在 (a)--(d) 中，列出 $D_{16}$ 中满足给定条件的所有子群：
  \begin{enumerate}[label=(\alph*)]
    \item 包含于 $\langle sr^2,r^4\rangle$ 的子群；
    \item 包含于 $\langle sr^7,r^4\rangle$ 的子群；
    \item 含有 $\langle r^4\rangle$ 的子群；
    \item 含有 $\langle s\rangle$ 的子群。
  \end{enumerate}
  \item 证明 $D_8$ 的子群 $\langle s,r^2\rangle$ 同构于 $V_4$。
  \item 用给定的格找出所有生成 $D_8$ 的元素对（共有 $12$ 对）。
  \item 用给定的格找出所有 $x\in D_{16}$，使得 $D_{16}=\langle x,s\rangle$（共有 $16$ 个这样的元素 $x$）。
  \item 用给定的格帮助求下列群中每个元素的中心化子：
  \[
    \text{(a) }D_8,\qquad
    \text{(b) }Q_8,\qquad
    \text{(c) }S_3,\qquad
    \text{(d) }D_{16}.
  \]
  \item 求 $D_{16}$ 的中心。
  \item 在下列每个群中，求每个子群的正规化子：
  \[
    \text{(a) }S_3,
    \qquad
    \text{(b) }Q_8.
  \]
  \item 画出下列群的子群格：
  \[
    \text{(a) }\ZZ/16\ZZ,
    \qquad
    \text{(b) }\ZZ/24\ZZ,
    \qquad
    \text{(c) }\ZZ/48\ZZ.
  \]
  [参见第 3 节练习 6。]
  \item 通过证明若 $|G|=4$，则 $G\cong Z_4$ 或 $G\cong V_4$，来分类阶为 $4$ 的群。[见第 1.1 节练习 36。]
  \item 考虑具有如下表示的阶为 $16$ 的群：
  \[
    QD_{16}=\langle \sigma,\tau\mid \sigma^8=\tau^2=1,
    \sigma\tau=\tau\sigma^3\rangle
  \]
  （称为阶为 $16$ 的\emph{拟二面体群}或\emph{半二面体群}）。这个群有三个阶为 $8$ 的子群：
  \[
    \langle \tau,\sigma^2\rangle\cong D_8,
    \qquad
    \langle\sigma\rangle\cong Z_8,
    \qquad
    \langle\sigma^2,\sigma\tau\rangle\cong Q_8,
  \]
  并且每个真子群都包含于这三个子群之一。在下列拟二面体群全部子群的格中填入缺失的子群，并用至多两个生成元表示每个子群。（这是另一个非平面格的例子。）

  \begin{center}
  \begin{tikzpicture}[scale=.9, every node/.style={inner sep=1pt}]
    \node (top) at (0,5) {$G$};
    \node (left8) at (-2.7,3.8) {$\langle\sigma^2,\tau\rangle$};
    \node (mid8) at (0,3.8) {$\langle\sigma\rangle$};
    \node (right8) at (2.7,3.8) {$\langle\sigma^2,\tau\sigma\rangle$};
    \node (l1) at (-4.0,2.45) {$\langle **\rangle$};
    \node (l2) at (-1.6,2.45) {$\langle\sigma^4,\tau\rangle$};
    \node (m1) at (0,2.45) {$\langle **\rangle$};
    \node (r1) at (1.8,2.45) {$\langle **\rangle$};
    \node (r2) at (4.0,2.45) {$\langle **\rangle$};
    \node (a) at (-4.8,1.1) {$\langle\tau\sigma^2\rangle$};
    \node (b) at (-3.2,1.1) {$\langle **\rangle$};
    \node (c) at (-1.6,1.1) {$\langle **\rangle$};
    \node (d) at (-.35,1.1) {$\langle\tau\rangle$};
    \node (e) at (0,.1) {$\langle\sigma^4\rangle$};
    \node (bot) at (0,-1.2) {$1$};
    \draw (top)--(left8) (top)--(mid8) (top)--(right8);
    \draw (left8)--(l1) (left8)--(l2) (left8)--(m1);
    \draw (mid8)--(m1);
    \draw (right8)--(m1) (right8)--(r1) (right8)--(r2);
    \draw (l1)--(a) (l1)--(b) (l1)--(e);
    \draw (l2)--(c) (l2)--(d) (l2)--(e);
    \draw (m1)--(e) (r1)--(e) (r2)--(e);
    \draw (a)--(bot) (b)--(bot) (c)--(bot) (d)--(bot) (e)--(bot);
  \end{tikzpicture}
  \end{center}

  下面三个例子导向两个不同构但具有相同子群格的群。

  \item 群
  \[
    A=Z_2\times Z_4=\langle a,b\mid a^2=b^4=1,
    ab=ba\rangle
  \]
  有阶 $8$，并有三个阶为 $4$ 的子群：
  \[
    \langle a,b^2\rangle\cong V_4,
    \qquad
    \langle b\rangle\cong Z_4,
    \qquad
    \langle ab\rangle\cong Z_4,
  \]
  且每个真子群都包含于这三个子群之一。画出 $A$ 的全部子群格，并用至多两个生成元表示每个子群。
  \item 群
  \[
    G=Z_2\times Z_8=\langle x,y\mid x^2=y^8=1,
    xy=yx\rangle
  \]
  有阶 $16$，并有三个阶为 $8$ 的子群：
  \[
    \langle x,y^2\rangle\cong Z_2\times Z_4,
    \qquad
    \langle y\rangle\cong Z_8,
    \qquad
    \langle xy\rangle\cong Z_8,
  \]
  且每个真子群都包含于这三个子群之一。画出 $G$ 的全部子群格，并用至多两个生成元表示每个子群（参见练习 12）。
  \item 设 $M$ 是具有如下表示的阶为 $16$ 的群：
  \[
    \langle u,v\mid u^2=v^8=1,
    vu=uv^5\rangle,
  \]
  （有时称为阶为 $16$ 的\emph{模群}）。它有三个阶为 $8$ 的子群：$\langle u,v^2\rangle$、$\langle v\rangle$ 和 $\langle uv\rangle$，并且每个真子群都包含于这三个子群之一。证明
  \[
    \langle u,v^2\rangle\cong Z_2\times Z_4,
    \qquad
    \langle v\rangle\cong Z_8,
    \qquad
    \langle uv\rangle\cong Z_8.
  \]
  说明 $M$ 的子群格与 $Z_2\times Z_8$ 的子群格相同（参见练习 13），但这两个群不同构。
  \item 描述 $D_{16}$ 中三个阶为 $8$ 的子群各自的同构类型。
  \item 用阶为 $16$ 的拟二面体群的子群格说明，每个阶为 $2$ 的元素都包含在真子群 $\langle\tau,\sigma^2\rangle$ 中（参见练习 11）。
  \item 用阶为 $16$ 的模群 $M$ 的子群格说明集合
  \[
    \{x\in M\mid x^2=1\}
  \]
  是 $M$ 的一个同构于克莱因四元群的子群（参见练习 14）。
  \item 用格帮助求 $QD_{16}$ 中每个元素的中心化子（参见练习 11）。
  \item 用格帮助求 $N_{D_{16}}(\langle s,r^4\rangle)$。
  \item 用 $QD_{16}$ 的子群格（参见练习 11）帮助求正规化子：
  \[
    \text{(a) }N_{QD_{16}}(\langle\tau\sigma\rangle),
    \qquad
    \text{(b) }N_{QD_{16}}(\langle\tau,\sigma^4\rangle).
  \]
\end{enumerate}

 \chapter{商群与同态}

\section{定义与例子}

本章引入群 $G$ 的\emph{商群}概念，这是从群 $G$ 得到一个“较小”群的另一种方法；正如我们对待子群那样，我们将用商群来研究 $G$ 的结构。群 $G$ 的结构反映在它的商群和子群的结构之中。例如，我们将看到，$G$ 的一个商群的子群格会以精确的方式反映在 $G$ 的子群格的“顶部”，而 $G$ 的一个子群的子群格则自然出现在“底部”。因此，人们可以把这些信息结合起来以获得关于 $G$ 的信息；我们也将说明某些分类定理如何由此产生。

研究 $G$ 的商群本质上等价于研究 $G$ 的同态，即把群 $G$ 映到另一个群并保持群结构的映射。回忆：若 $\varphi$ 是从 $G$ 到群 $H$ 的同态，则 $\varphi$ 的纤维是 $G$ 中映到 $H$ 的单个元素上的元素集合。可以把这种情形图示为图 1，其中点 $a$ 上方方框中的竖线表示 $\varphi$ 在 $a$ 上方的纤维。

\begin{center}
\begin{tikzpicture}[scale=.95, every node/.style={inner sep=1pt}]
  \draw (-3.8,1) rectangle (3.8,2.4);
  \foreach \x in {-3.1,-2.8,-.55,.55,2.8,3.1}{\draw (\x,1.15)--(\x,2.25);} 
  \node at (-1.65,1.7) {$\cdots$};
  \node at (1.65,1.7) {$\cdots$};
  \node[right] at (4.0,1.7) {$G$};
  \draw (-3.8,0)--(3.8,0);
  \foreach \x in {-3.1,-2.8,0,2.8,3.1}{\fill (\x,0) circle (1.7pt);}
  \foreach \x in {-3.1,-2.8,0,2.8,3.1}{\draw[->] (\x,1.0)--(\x,.13);} 
  \node[below] at (0,0) {$a$};
  \node[right] at (4.0,0) {$H$};
  \node at (3.95,-.55) {图 1};
\end{tikzpicture}
\end{center}

$H$ 中的群运算给出一种方法来相乘 $\varphi$ 的像中的两个元素。这提示我们用一种自然方法相乘位于这两个点之上的纤维，使纤维集合本身成为一个群：若 $X_a$ 是 $a$ 上方的纤维，$X_b$ 是 $b$ 上方的纤维，则定义 $X_a$ 与 $X_b$ 的乘积为 $ab$ 上方的纤维 $X_{ab}$，即
\[
  X_aX_b=X_{ab}.
\]
这种乘法是结合的，因为 $H$ 中的乘法结合；单位元是 $H$ 的单位元上方的纤维；$a$ 上方纤维的逆是 $a^{-1}$ 上方的纤维。例如，结合律证明如下：
\[
  (X_aX_b)X_c=(X_{ab})X_c=X_{(ab)c},\qquad
  X_a(X_bX_c)=X_a(X_{bc})=X_{a(bc)}.
\]
由于 $(ab)c=a(bc)$ 在 $H$ 中成立，故 $(X_aX_b)X_c=X_a(X_bX_c)$。粗略地说，群 $G$ 被分划成若干片（纤维），而这些片本身具有群结构，称为 $G$ 的商群（正式定义见下面的例子之后）。

由于纤维的乘法由 $H$ 中的乘法定义，按构造，带有这种乘法的商群自然同构于 $G$ 在同态 $\varphi$ 下的像（纤维 $X_a$ 与它在 $H$ 中的像 $a$ 对应）。

\paragraph{例}
令 $G=\ZZ$，令 $H=Z_n=\langle x\rangle$ 是阶为 $n$ 的循环群，定义
\[
  \varphi:\ZZ\longrightarrow Z_n,
  \qquad \varphi(a)=x^a.
\]
由于
\[
  \varphi(a+b)=x^{a+b}=x^a x^b=\varphi(a)\varphi(b),
\]
$\varphi$ 是同态（注意 $\ZZ$ 中的运算是加法，而 $Z_n$ 中的运算是乘法）。显然 $\varphi$ 是满射。$\varphi$ 在 $x^a$ 上方的纤维是
\begin{align*}
  \varphi^{-1}(x^a)
  &=\{m\in\ZZ\mid x^m=x^a\}  \\
  &=\{m\in\ZZ\mid x^{m-a}=1\} \\
  &=\{m\in\ZZ\mid n\text{ 整除 }m-a\}  \\
  &=\{m\in\ZZ\mid m\equiv a\pmod n\}=\overline a.
\end{align*}
也就是说，$\varphi$ 的纤维正是模 $n$ 的剩余类。图 1 在这里成为图 2。

\begin{center}
\begin{tikzpicture}[scale=.82, every node/.style={inner sep=1pt}]
  \draw (-4,0) rectangle (4,2.7);
  \draw (-2.8,0)--(-2.8,2.7); \draw (-1.6,0)--(-1.6,2.7); \draw (1.4,0)--(1.4,2.7); \draw (2.8,0)--(2.8,2.7);
  \node at (-3.4,2.35) {$0$}; \node at (-3.4,1.95) {$\pm n$}; \node at (-3.4,1.55) {$\pm2n$}; \node at (-3.4,1.15) {$\pm3n$}; \node at (-3.4,.55) {$\vdots$};
  \node at (-2.2,2.35) {$1$}; \node at (-2.2,1.95) {$1\pm n$}; \node at (-2.2,1.55) {$1\pm2n$}; \node at (-2.2,1.15) {$1\pm3n$}; \node at (-2.2,.55) {$\vdots$};
  \node at (-.8,1.35) {$\cdots$};
  \node at (.0,2.35) {$a$}; \node at (.0,1.95) {$a\pm n$}; \node at (.0,1.55) {$a\pm2n$}; \node at (.0,1.15) {$a\pm3n$}; \node at (.0,.55) {$\vdots$};
  \node at (2.1,1.35) {$\cdots$};
  \node at (3.4,2.35) {$n-1$}; \node at (3.4,1.95) {$(n-1)\pm n$}; \node at (3.4,1.55) {$(n-1)\pm2n$}; \node at (3.4,1.15) {$(n-1)\pm3n$}; \node at (3.4,.55) {$\vdots$};
  \node[right] at (4.2,1.35) {$\ZZ$};
  \foreach \x/\lab in {-3.4/x^0,-2.2/x^1,0/x^a,3.4/x^{n-1}}{\draw[->] (\x,0)--(\x,-1.0); \node[below] at (\x,-1.0) {$\lab$};}
  \node[right] at (4.2,-1.0) {$Z_n$};
  \node at (3.8,-1.55) {图 2};
\end{tikzpicture}
\end{center}

$Z_n$ 中的乘法就是 $x^a x^b=x^{a+b}$。对应的纤维为 $\overline a$、$\overline b$ 和 $\overline{a+b}$，所以纤维上的对应群运算为
\[
  \overline a\cdot\overline b=\overline{a+b}.
\]
这正是加法下的群 $\ZZ/n\ZZ$，它同构于 $\varphi$ 的像（即整个 $Z_n$）。

这个群的单位元（$Z_n$ 中单位元上方的纤维）由 $\ZZ$ 中所有 $n$ 的倍数组成，即 $n\ZZ$，这是 $\ZZ$ 的一个子群；其余纤维只是这个子群的平移 $a+n\ZZ$。群运算也可直接定义为：从这些纤维中取代表元，在 $\ZZ$ 中相加，再取含有此和的纤维（这正是群 $\ZZ/n\ZZ$ 的原始定义）。从计算角度看，直接把代表元 $a$ 与 $b$ 相加来计算 $\overline a$ 与 $\overline b$ 的乘积，比先求这些纤维在 $\varphi$ 下的像（$x^a$ 与 $x^b$）、在 $H$ 中相乘（得到 $x^{a+b}$）、再取这个乘积上方的纤维要容易得多。

我们先考虑同态及其纤维的一些基本性质。给同态 $\varphi:G\to H$ 位于 $H$ 的单位元上方的纤维一个名称：

\noindent\textbf{定义.} 若 $\varphi:G\to H$ 是同态，则 $\varphi$ 的\emph{核}是集合
\[
  \{g\in G\mid \varphi(g)=1\},
\]
记为 $\ker\varphi$（这里 $1$ 是 $H$ 的单位元）。

\medskip
\noindent\textbf{命题 1.} 设 $G$ 和 $H$ 是群，$\varphi:G\to H$ 是同态。
\begin{enumerate}[label=(\arabic*)]
  \item $\varphi(1_G)=1_H$，其中 $1_G$ 与 $1_H$ 分别为 $G$ 与 $H$ 的单位元。
  \item 对所有 $g\in G$，$\varphi(g^{-1})=\varphi(g)^{-1}$。
  \item 对所有 $n\in\ZZ$，$\varphi(g^n)=\varphi(g)^n$。
  \item $\ker\varphi$ 是 $G$ 的子群。
  \item $\im(\varphi)$，即 $G$ 在 $\varphi$ 下的像，是 $H$ 的子群。
\end{enumerate}

\begin{proof}
(1) 由于
\[
  \varphi(1_G)=\varphi(1_G1_G)=\varphi(1_G)\varphi(1_G),
\]
由消去律得 $\varphi(1_G)=1_H$。

(2) 由 (1)，
\[
  1_H=\varphi(1_G)=\varphi(gg^{-1})=\varphi(g)\varphi(g^{-1}).
\]
左乘 $\varphi(g)^{-1}$ 并化简，得结论。

(3) 当 $n\in\ZZ^+$ 时由归纳法即可；负整数情形由 (2) 推出。

(4) 由 (1)，$1_G\in\ker\varphi$，所以核非空。若 $x,y\in\ker\varphi$，则 $\varphi(x)=\varphi(y)=1_H$，于是
\[
  \varphi(xy^{-1})=\varphi(x)\varphi(y^{-1})
  =\varphi(x)\varphi(y)^{-1}=1_H1_H^{-1}=1_H.
\]
所以 $xy^{-1}\in\ker\varphi$。由子群判别法，$\ker\varphi\leq G$。

(5) 由 $\varphi(1_G)=1_H$，$H$ 的单位元属于 $\im(\varphi)$。若 $x,y\in\im(\varphi)$，设 $x=\varphi(a)$，$y=\varphi(b)$，则由 (2)，$y^{-1}=\varphi(b^{-1})$，所以
\[
  xy^{-1}=\varphi(a)\varphi(b^{-1})=\varphi(ab^{-1})
\]
仍在像中。由子群判别法，$\im(\varphi)\leq H$。
\end{proof}

现在可以定义与商群相关的术语。

\noindent\textbf{定义.} 设 $\varphi:G\to H$ 是同态，其核为 $K$。\emph{商群}或\emph{因子群} $G/K$（读作 $G$ 模 $K$，或简称 $G$ mod $K$）是一个群，其元素是 $\varphi$ 的纤维，群运算定义如上：若 $X$ 是 $a$ 上方的纤维，$Y$ 是 $b$ 上方的纤维，则 $X$ 与 $Y$ 的乘积定义为 $ab$ 上方的纤维。

记号强调：核 $K$ 是群 $G/K$ 中的一个单元素；下面的命题 2 将说明，正如上面 $\ZZ/n\ZZ$ 的例子一样，$G/K$ 的其他元素只是核 $K$ 的“平移”。因此我们可把 $G/K$ 看作由把 $K$ 压缩掉或“除掉”而得到（更精确地说，是模 $K$ 等价）。这解释了为何 $G/K$ 称为“商”群。

上面关于商群 $G/K$ 的定义显式要求给出映射 $\varphi$，因为纤维的乘法先通过 $\varphi$ 投影到 $H$，在 $H$ 中相乘，再确定这个乘积上方的纤维。与 $\ZZ/n\ZZ$ 一样，也可直接用纤维中的代表元定义纤维乘法；这在计算上更简单，且不显式出现 $\varphi$。我们先证明同态的纤维可用同态核来表示，就像上面的例子中核为 $n\ZZ$，纤维为 $a+n\ZZ$ 的平移一样。

\medskip
\noindent\textbf{命题 2.} 设 $\varphi:G\to H$ 是群同态，核为 $K$。设 $X\in G/K$ 是 $a$ 上方的纤维，即 $X=\varphi^{-1}(a)$。则
\begin{enumerate}[label=(\arabic*)]
  \item 对任意 $u\in X$，有 $X=\{uk\mid k\in K\}$；
  \item 对任意 $u\in X$，有 $X=\{ku\mid k\in K\}$。
\end{enumerate}

\begin{proof}
证明 (1)，(2) 留作练习。取 $u\in X$，于是 $\varphi(u)=a$。令
\[
  uK=\{uk\mid k\in K\}.
\]
先证 $uK\subseteq X$。若 $k\in K$，则
\[
  \varphi(uk)=\varphi(u)\varphi(k)=\varphi(u)1=a,
\]
所以 $uk\in X$。反过来，设 $g\in X$，令 $k=u^{-1}g$。则
\[
  \varphi(k)=\varphi(u^{-1})\varphi(g)=\varphi(u)^{-1}\varphi(g)=a^{-1}a=1.
\]
因此 $k\in K$。又 $k=u^{-1}g$，故 $g=uk\in uK$。于是 $X=uK$。
\end{proof}

命题 2 中描述同态纤维的集合，可对 $G$ 的任意子群 $K$ 定义，不必假设它是某个同态的核（稍后我们会确定一个子群成为某个同态核的充要条件）。这些集合有专门名称。

\noindent\textbf{定义.} 对任意 $N\leq G$ 和任意 $g\in G$，令
\[
  gN=\{gn\mid n\in N\},\qquad Ng=\{ng\mid n\in N\},
\]
分别称为 $N$ 在 $G$ 中的\emph{左陪集}和\emph{右陪集}。陪集中的任意元素称为该陪集的一个\emph{代表元}。

我们已在命题 2 中看到，若 $N$ 是同态的核，且 $g_1$ 是陪集 $gN$ 的任意代表元，则 $g_1N=gN$（若 $g_1\in Ng$，则 $Ng_1=Ng$）。命题 4 会说明这一事实对任意子群 $N$ 都成立，这解释了“代表元”这个术语。

若 $G$ 是加法群，我们写 $g+N$ 和 $N+g$ 分别表示代表元为 $g$ 的左、右陪集。一般地，可把左陪集 $gN$ 看作 $N$ 被 $g$ 左平移所得的集合。

用这个定义来说，命题 2 表明同态的纤维是核的左陪集（也是核的右陪集），即商 $G/K$ 的元素是左陪集 $gK$，$g\in G$。在 $\ZZ/n\ZZ$ 的例子中，商群中的乘法也可用代表元来定义。下面的结果说明一般的 $G/K$ 中也如此（只要知道 $K$ 是某个同态的核）：两个左陪集 $X$ 和 $Y$ 在 $G/K$ 中的乘积，通过任取 $X$ 的代表元 $u$、$Y$ 的代表元 $v$，在 $G$ 中相乘，并取陪集 $(uv)K$ 来计算。

\medskip
\noindent\textbf{定理 3.} 设 $G$ 是群，$K$ 是从 $G$ 到另一群的某个同态的核。则 $K$ 在 $G$ 中的左陪集组成的集合，在运算
\[
  uK\circ vK=(uv)K
\]
下构成群 $G/K$。特别地，此运算是良定义的：若 $u_1$ 是 $uK$ 中任意元素，$v_1$ 是 $vK$ 中任意元素，则 $u_1v_1\in uvK$，即 $u_1v_1K=uvK$，所以乘法不依赖于陪集代表元的选择。把“左陪集”换成“右陪集”也有相同结论。

\begin{proof}
设 $X,Y\in G/K$，令 $Z=XY$。由命题 2(1)，$X,Y,Z$ 都是 $K$ 的左陪集。按假设，$K$ 是某个同态 $\varphi:G\to H$ 的核，所以 $X=\varphi^{-1}(a)$，$Y=\varphi^{-1}(b)$，其中 $a,b\in H$。根据 $G/K$ 中的运算定义，$Z=\varphi^{-1}(ab)$。任取 $X$、$Y$ 的代表元 $u$、$v$，则 $\varphi(u)=a$，$\varphi(v)=b$，且 $X=uK$、$Y=vK$。要证 $uv\in Z$。确实，
\[
  uv\in Z \Longleftrightarrow uv\in\varphi^{-1}(ab)
  \Longleftrightarrow \varphi(uv)=ab
  \Longleftrightarrow \varphi(u)\varphi(v)=ab,
\]
而最后一个等式成立。故 $uv\in Z$，因此 $Z=uvK$。这证明了对任意代表元 $u\in X$、$v\in Y$，$X$ 与 $Y$ 的乘积都是陪集 $uvK$。最后关于右陪集的断言由命题 2 中 $uK=Ku$、$vK=Kv$ 立即推出。
\end{proof}

用代表元作 $G/K$ 中的乘法，可如图 3 所示。

\begin{center}
\begin{tikzpicture}[scale=.9, every node/.style={inner sep=1pt}]
  \draw (-3.8,1) rectangle (3.8,2.55);
  \foreach \x in {-3.1,-.6,2.6}{\draw (\x,1.1)--(\x,2.45);} 
  \node at (.85,1.78) {$\cdots$};
  \node at (-3.1,2.62) {$u$}; \node at (-.6,2.62) {$v$}; \node at (2.6,2.62) {$uv$};
  \draw[dashed,->] (-3.0,2.35) .. controls (-2.0,2.95) and (-1.2,2.95) .. (-.65,2.35);
  \draw[dashed,->] (-.55,2.35) .. controls (.4,2.95) and (1.7,2.95) .. (2.55,2.35);
  \node[right] at (4.0,1.78) {$G$};
  \draw (-3.8,0)--(3.8,0);
  \foreach \x/\lab in {-3.1/a,-.6/b,2.6/ab}{\fill (\x,0) circle (1.7pt); \node[below] at (\x,0) {$\lab$}; \draw[->] (\x,1.0)--(\x,.13);} 
  \draw[dashed,->] (-3.0,.3) .. controls (-2.0,.75) and (1.2,.75) .. (2.55,.3);
  \draw[dashed,->] (-.55,.25) .. controls (.6,.55) and (1.55,.55) .. (2.55,.25);
  \node[right] at (4.0,0) {$H$};
  \node at (3.9,-.55) {图 3};
\end{tikzpicture}
\end{center}

我们强调：乘法与所选代表元无关。也就是说，两个陪集 $X$ 和 $Y$ 的乘积（若群写作加法，则为和）是含有乘积 $uv$ 的陪集 $uvK$，其中 $u$ 和 $v$ 可分别任取为 $X$ 和 $Y$ 的代表元。只考虑某元素所在陪集，或“模 $K$ 化简”的过程，正是我们在 $\ZZ/n\ZZ$ 中一直在做的事情。表示包含代表元 $u$ 的陪集 $uK$ 的一个有用记号是 $\overline u$。用这个记号，商群 $G/K$ 可记为 $\overline G$，元素 $\overline u$ 与 $\overline v$ 的乘积就是包含 $uv$ 的陪集，即 $\overline{uv}$。此记号也强调：$G/K$ 中的陪集 $uK$ 是 $G/K$ 中的元素 $\overline u$。

\paragraph{例}
\begin{enumerate}[label=(\arabic*)]
  \item 本章开头从 $\ZZ$ 到 $Z_n$ 的同态 $\varphi$ 的纤维是核 $n\ZZ$ 的左陪集（也是右陪集）$a+n\ZZ$。定理 3 证明这些陪集在代表元相加下构成一个群，即 $\ZZ/n\ZZ$，这解释了该群的记号。此群自然同构于它在 $\varphi$ 下的像，所以我们重新得到第 2 章的同构 $\ZZ/n\ZZ\cong Z_n$。
  \item 若 $\varphi:G\to H$ 是同构，则 $K=1$，$\varphi$ 的纤维是 $G$ 的单元素子集，因此 $G/1\cong G$。
  \item 设 $G$ 是任意群，令 $H=1$ 为阶 $1$ 的群，定义 $\varphi:G\to H$ 为 $\varphi(g)=1$。显然 $\varphi$ 是同态，称为\emph{平凡同态}。此时 $\ker\varphi=G$，而 $G/G$ 是只有一个元素 $G$ 的群，即 $G/G\cong Z_1=\{1\}$。
  \item 设 $G=\RR^2$（向量加法），$H=\RR$（加法），定义
  \[
    \varphi:\RR^2\to\RR,
    \qquad \varphi((x,y))=x.
  \]
  这是到 $x$ 轴的投影。它是同态，因为
  \[
    \varphi((x_1,y_1)+(x_2,y_2))=x_1+x_2
    =\varphi((x_1,y_1))+\varphi((x_2,y_2)).
  \]
  其核为
  \[
    \ker\varphi=\{(x,y)\mid x=0\},
  \]
  即 $y$ 轴。$a\in\RR$ 上方的纤维是直线 $x=a$，也就是以 $(a,0)$（或任何投影到 $a$ 的点）为代表元的核的陪集：
  \[
    (a,0)+(y\text{轴}).
  \]
  因此图 1 在此例中成为图 4。

\begin{center}
\begin{tikzpicture}[scale=.75, every node/.style={inner sep=1pt}]
  \draw (-4,0) rectangle (4,3.1);
  \draw[->] (-3.6,1.55)--(3.7,1.55) node[right] {$x$};
  \draw[->] (0,.25)--(0,2.9) node[above] {$y$};
  \foreach \x/\lab in {-2/-1,0/0,1.3/1,3/a}{\draw (\x,.4)--(\x,2.7); \node[below] at (\x,1.45) {$\lab$};}
  \node[right] at (4.2,1.55) {$\RR^2$};
  \foreach \x/\lab in {-2/-1,0/0,1.3/1,3/a}{\draw[->] (\x,0)--(\x,-1.1);}
  \draw[->] (-3.6,-1.1)--(3.7,-1.1) node[right] {$\RR$};
  \foreach \x/\lab in {-2/-1,0/0,1.3/1,3/a}{\fill (\x,-1.1) circle (1.7pt); \node[below] at (\x,-1.1) {$\lab$};}
  \node at (3.6,-1.75) {图 4};
\end{tikzpicture}
\end{center}

  群运算（此处为加法）既可通过 $\varphi$ 描述：直线 $x=a$ 与直线 $x=b$ 的和是直线 $x=a+b$；也可用陪集代表元直接描述：含 $(a,y_1)$ 的竖直线与含 $(b,y_2)$ 的竖直线之和，是含 $(a+b,y_1+y_2)$ 的竖直线。注意代表元的选择不重要（即 $y$ 坐标不重要）。
  \item （一个 $G$ 非阿贝尔的例子。）令 $G=Q_8$，令 $H=V_4$ 为克莱因四元群。定义
  \[
    \varphi:Q_8\to V_4,
    \qquad \varphi(\pm1)=1,
    \quad \varphi(\pm i)=a,
    \quad \varphi(\pm j)=b,
    \quad \varphi(\pm k)=c.
  \]
  验证 $\varphi$ 是同态留作练习；利用对称性可减少检查 $\varphi(xy)=\varphi(x)\varphi(y)$ 的工作。显然 $\varphi$ 是满射，且 $\ker\varphi=\{\pm1\}$。可把 $\varphi$ 看作 $Q_8$ 上的一个“绝对值”函数；其纤维是
  \[
    E=\{\pm1\},\quad A=\{\pm i\},\quad B=\{\pm j\},\quad C=\{\pm k\},
  \]
  它们在 $Q_8/\langle -1\rangle$ 中分别被压缩为 $1,a,b,c$，并且这些纤维正是 $\ker\varphi$ 的左陪集（也是右陪集）。
\end{enumerate}

由定理 3，若给定群 $G$ 的子群 $K$，并知道 $K$ 是某个同态的核，则可不借助该同态而通过乘法 $uK\cdot vK=uvK$ 定义商 $G/K$。这引出一个问题：是否可对 $G$ 的任意子群 $N$ 都类似地定义商群 $G/N$？答案一般是否定的，因为这种乘法一般不是良定义的（见下面命题 5）。事实上，我们将看到，$N$ 的陪集集合能成为群，当且仅当 $N$ 是某个同态的核（命题 7）。我们也会给出判断一个子群 $N$ 是否为这种核的准则，这就是正规子群的概念；后续章节将研究非正规子群。

我们先说明任意子群的陪集都会分划 $G$（即并为全 $G$，且不同陪集不交）。

\medskip
\noindent\textbf{命题 4.} 设 $N$ 是群 $G$ 的任意子群。$N$ 在 $G$ 中的左陪集组成 $G$ 的一个分划。此外，对所有 $u,v\in G$，
\[
  uN=vN\quad\Longleftrightarrow\quad v^{-1}u\in N,
\]
特别地，$uN=vN$ 当且仅当 $u$ 和 $v$ 是同一陪集的代表元。

\begin{proof}
首先，由于 $N$ 是 $G$ 的子群，$1\in N$。因此对所有 $g\in G$，$g=g\cdot1\in gN$，所以
\[
  G=\bigcup_{g\in G} gN.
\]
为证明不同左陪集不交，设 $uN\cap vN\neq\varnothing$。取 $x\in uN\cap vN$，写作
\[
  x=un=vm,
\]
其中 $n,m\in N$。在等式右端两边右乘 $n^{-1}$，得
\[
  u=vmn^{-1}=vm_1,
\]
其中 $m_1=mn^{-1}\in N$。于是对 $uN$ 的任意元素 $ut$（$t\in N$），有
\[
  ut=(vm_1)t=v(m_1t)\in vN.
\]
故 $uN\subseteq vN$。交换 $u$ 与 $v$ 可得 $vN\subseteq uN$。因此有非空交的两个陪集相等。

由前面的结论，$uN=vN$ 当且仅当 $u\in vN$，当且仅当 $u=vn$ 对某个 $n\in N$ 成立，当且仅当 $v^{-1}u\in N$。最后，$v\in uN$ 等价于说 $v$ 是 $uN$ 的代表元，所以 $uN=vN$ 当且仅当 $u$ 与 $v$ 是同一个陪集的代表元。
\end{proof}

\medskip
\noindent\textbf{命题 5.} 设 $G$ 是群，$N$ 是 $G$ 的子群。
\begin{enumerate}[label=(\arabic*)]
  \item 在 $N$ 的左陪集集合上由
  \[
    uN\cdot vN=(uv)N
  \]
  描述的运算是良定义的，当且仅当对所有 $g\in G$ 和所有 $n\in N$，都有 $gng^{-1}\in N$。
  \item 若上述运算是良定义的，则它使 $N$ 的左陪集集合成为一个群。特别地，该群的单位元是陪集 $1N$，$gN$ 的逆是 $g^{-1}N$，即 $(gN)^{-1}=g^{-1}N$。
\end{enumerate}

\begin{proof}
(1) 先假设该运算良定义。取任意 $g\in G$ 与 $n\in N$。令 $u=1$，$u_1=n$，且 $v=v_1=g^{-1}$。由良定义性得
\[
  1g^{-1}N=ng^{-1}N,
\]
即 $g^{-1}N=ng^{-1}N$。由于 $1\in N$，$ng^{-1}\in ng^{-1}N$，所以 $ng^{-1}\in g^{-1}N$，即 $ng^{-1}=g^{-1}n_1$，其中 $n_1\in N$。左乘 $g$，得 $gng^{-1}=n_1\in N$。

反过来，假设 $gng^{-1}\in N$ 对所有 $g\in G$ 和 $n\in N$ 成立。设 $u_1\in uN$，$v_1\in vN$。写作 $u_1=un$，$v_1=vm$，其中 $n,m\in N$。要证 $u_1v_1\in uvN$：
\[
  u_1v_1=(un)(vm)=u(vv^{-1})nvm=(uv)(v^{-1}nv)m=(uv)(n_1m),
\]
其中 $n_1=v^{-1}nv\in N$。由于 $N$ 对乘积封闭，$n_1m\in N$。故 $u_1v_1=(uv)n_2$，其中 $n_2\in N$。于是左陪集 $uvN$ 与 $u_1v_1N$ 有公共元素 $u_1v_1$，由命题 4 它们相等。运算良定义。

(2) 若陪集上的运算良定义，则群公理由 $G$ 中对应公理诱导。例如，对所有 $u,v,w\in G$，
\begin{align*}
  (uN)(vNwN)&=uN(vwN)=u(vw)N\\
  &=(uv)wN=(uNvN)(wN),
\end{align*}
因为 $u(vw)=(uv)w$ 在 $G$ 中成立。由乘法定义可立即看出单位元是 $1N$，$gN$ 的逆是 $g^{-1}N$。
\end{proof}

如前所述，满足命题 5 条件、使商 $G/N$ 自然具有群结构的子群 $N$ 有专门名称。

\noindent\textbf{定义.} 元素 $gng^{-1}$ 称为 $n\in N$ 关于 $g$ 的\emph{共轭}。集合
\[
  gNg^{-1}=\{gng^{-1}\mid n\in N\}
\]
称为 $N$ 关于 $g$ 的\emph{共轭}。若 $gNg^{-1}=N$，则称 $g$ \emph{正规化} $N$。群 $G$ 的子群 $N$ 称为\emph{正规子群}，如果 $G$ 的每个元素都正规化 $N$，即对所有 $g\in G$ 都有 $gNg^{-1}=N$。若 $N$ 是 $G$ 的正规子群，记为 $N\normal G$。

注意，当 $N$ 是正规子群时，$G$ 的结构反映在商 $G/N$ 的结构中（例如，$G/N$ 中乘法的结合律由 $G$ 中的结合律诱导，$G/N$ 中的逆由 $G$ 中的逆诱导）。第 3 节讨论同构定理时，我们会看到更多 $G$ 与其商 $G/N$ 的关系。

把上面的结果概括为：

\medskip
\noindent\textbf{定理 6.} 设 $N$ 是群 $G$ 的子群。以下条件等价：
\begin{enumerate}[label=(\arabic*)]
  \item $N\normal G$；
  \item $N_G(N)=G$（回忆 $N_G(N)$ 是 $N$ 在 $G$ 中的正规化子）；
  \item 对所有 $g\in G$，$gN=Ng$；
  \item 命题 5 中描述的左陪集运算使左陪集集合成为群；
  \item 对所有 $g\in G$，$gNg^{-1}\subseteq N$。
\end{enumerate}

\begin{proof}
困难的等价性已在上面完成；其余留作练习。
\end{proof}

在实际中，人们尽量减少判断某个子群 $N$ 是否在 $G$ 中正规所需的计算。特别地，应尽量避免对所有 $n\in N$ 和 $g\in G$ 计算全部共轭 $gng^{-1}$。例如，$N$ 中元素本身正规化 $N$，因为 $N$ 是子群。此外，若已知 $N$ 的一组生成元，则只需检查这些生成元的所有共轭都在 $N$ 中，即可证明 $N$ 是正规子群（这是因为乘积的共轭是共轭的乘积，逆的共轭是共轭的逆）——见后面的练习 26。类似地，若已知 $G$ 的生成元，只需检查 $G$ 的这些生成元正规化 $N$。特别地，若 $N$ 和 $G$ 的生成元都已知，则只需检查少量共轭。如果 $N$ 是有限群，则只需检查 $N$ 的一组生成元关于 $G$ 的一组生成元的共轭仍在 $N$ 中（练习 29）。最后，常常可以直接证明 $N_G(N)=G$ 而不作过多计算，从而证明 $N$ 是 $G$ 的正规子群。

现在证明正规子群恰好就是前面所说的同态核。

\medskip
\noindent\textbf{命题 7.} 群 $G$ 的子群 $N$ 是正规子群，当且仅当它是某个同态的核。

\begin{proof}
若 $N$ 是同态 $\varphi$ 的核，则命题 2 表明 $N$ 的左陪集与右陪集相同（并且都是 $\varphi$ 的纤维）。由定理 6(3)，$N$ 是正规子群。

反过来，若 $N\normal G$，令 $H=G/N$，并定义
\[
  \pi:G\to G/N,
  \qquad \pi(g)=gN.
\]
由 $G/N$ 中运算的定义，
\[
  \pi(g_1g_2)=(g_1g_2)N=g_1N\,g_2N=\pi(g_1)\pi(g_2),
\]
故 $\pi$ 是同态。并且
\[
  \ker\pi=\{g\in G\mid \pi(g)=1N\}
  =\{g\in G\mid gN=1N\}=N.
\]
所以 $N$ 是同态 $\pi$ 的核。
\end{proof}

上面构造的、把正规子群 $N$ 表示为某个同态的核的同态有一个名称：

\noindent\textbf{定义.} 设 $N\normal G$。同态
\[
  \pi:G\to G/N,
  \qquad \pi(g)=gN
\]
称为 $G$ 到 $G/N$ 上的\emph{自然投影（同态）}\footnote{“自然”一词在范畴论中有精确含义；这里我们用它表示这个同态的定义是“坐标无关”的，即只由元素本身描述，而不是由 $G$ 或 $N$ 的生成元描述。}。若 $\overline H\leq G/N$，则 $\overline H$ 在 $G$ 中的\emph{完全原像}是它在自然投影同态下的原像。

商 $G/N$ 的一个子群的完全原像是 $G$ 的子群，并且包含 $N$，因为 $N$ 的元素映到单位元，而单位元属于任意子群。第 3 节的同构定理将说明，包含 $N$ 的 $G$ 的子群与商群 $G/N$ 的子群之间有一种自然对应。

现在我们有了一个“内部”准则，可精确判断 $G$ 的子群 $N$ 何时是某个同态的核，即
\[
  N_G(N)=G.
\]
因此可把子群 $N$ 的正规化子看作 $N$ “离正规子群有多近”的一种度量（这解释了正规化子这个名称）。要记住，正规性是一种\emph{嵌入性质}，也就是说，它取决于 $N$ 与 $G$ 的关系，而不只取决于 $N$ 自身的内部结构（同一个群 $N$ 可以是 $G$ 的正规子群，但在包含 $G$ 的更大群中并不正规）。

我们从存在一个从 $G$ 到 $H$ 的同态 $\varphi$ 开始讨论商群，说明这个同态的核是 $G$ 的正规子群 $N$，并且商 $G/N$（最初用纤维定义）自然同构于 $G$ 在 $H$ 中的像。反过来，若 $N\normal G$，则可以找到一个群 $H$（即 $G/N$）和同态 $\pi:G\to H$，使得 $\ker\pi=N$（即自然投影）。因此，研究 $G$ 的同态像等价于研究 $G$ 的商群；我们将用同态来产生正规子群，也反过来用正规子群来产生同态。

我们借助同态来发展商群理论，而不是简单地先定义正规子群及其商群，是为了强调：商中的元素是原群 $G$ 的子集（即核 $N$ 的纤维或陪集）。前面的图示也强调，$N$（及其陪集）被投影（或压缩）为商 $G/N$ 中的单个元素。商群 $G/N$ 中的计算通过从相关陪集中取代表元完成。

下面给出一些正规子群及其商群的例子。

\paragraph{例}
设 $G$ 是群。
\begin{enumerate}[label=(\arabic*)]
  \item 子群 $1$ 与 $G$ 总是 $G$ 的正规子群；$G/1\cong G$ 且 $G/G\cong1$。
  \item 若 $G$ 是阿贝尔群，则 $G$ 的任意子群 $N$ 都正规，因为对所有 $g\in G$ 与 $n\in N$，
  \[
    gng^{-1}=gg^{-1}n=n\in N.
  \]
  注意重要的是 $G$ 阿贝尔，而不是仅仅 $N$ 阿贝尔。随着子群 $N$ 的不同，$G/N$ 的结构也可能不同。例如，若 $G=\ZZ$，则 $G$ 的每个子群 $N$ 都是循环的：
  \[
    N=\langle n\rangle=\langle -n\rangle=n\ZZ,
    \qquad n\in\ZZ,
  \]
  且 $G/N=\ZZ/n\ZZ$ 是以 $\overline1=1+n\ZZ$ 为生成元的循环群。

  再设 $G=Z_k$ 是阶为 $k$ 的循环群，$x$ 是 $G$ 的生成元，$N\normal G$。由命题 2.6，$N=\langle x^d\rangle$，其中 $d$ 是属于 $N$ 的 $x$ 的最小正幂。于是
  \[
    G/N=\{gN\mid g\in G\}=\{x^aN\mid a\in\ZZ\}.
  \]
  由于 $x^aN=(xN)^a$，得到 $G/N=\langle xN\rangle$，即 $G/N$ 是以 $xN$ 为生成元的循环群。综上：循环群的商群仍是循环群；且 $G$ 的生成元在商中的像是商群的生成元。若 $G$ 还是有限循环群且 $N\normal G$，则 $|G/N|=|G|/|N|$ 给出商群阶的公式。
  \item 若 $N\leq Z(G)$，则 $N\normal G$，因为对所有 $g\in G$ 与 $n\in N$，$gng^{-1}=n\in N$。这推广了上一例（其中 $Z(G)=G$）。因此特别地 $Z(G)\normal G$。

  子群 $\langle -1\rangle$ 在 $Q_8$ 中先前已被看作某个同态的核；现在由于 $\langle -1\rangle=Z(Q_8)$，可由另一种方式得到其正规性。我们已看到 $Q_8/\langle -1\rangle\cong V_4$。

  令 $G=D_8$，令 $Z=\langle r^2\rangle=Z(D_8)$。由于 $Z=\{1,r^2\}$，每个陪集 $gZ$ 是二元素集合 $\{g,gr^2\}$。这些陪集把 $D_8$ 的 $8$ 个元素分成四对，故 $Z$ 在 $D_8$ 中有四个（互不相交的）左陪集：
  \[
    \overline1=1Z,\quad \overline r=rZ,
    \quad \overline s=sZ,
    \quad \overline{rs}=rsZ.
  \]
  由阶 $4$ 群的分类（第 2.5 节练习 10），$D_8/Z(D_8)$ 同构于 $Z_4$ 或 $V_4$。为判断是哪一个，只需注意
  \[
    \overline r^2=r^2Z=Z=\overline1,
    \quad \overline s^2=s^2Z=Z=\overline1,
    \quad \overline{rs}^2=(rs)^2Z=Z=\overline1.
  \]
  所以商中每个非单位元都有阶 $2$。特别地，商中没有阶 $4$ 的元素，因此 $D_8/Z(D_8)$ 不是循环群，故 $D_8/Z(D_8)\cong V_4$。
\end{enumerate}

\subsection*{练习}
设 $G$ 和 $H$ 是群。
\begin{enumerate}[label=\arabic*.]
  \item 设 $\varphi:G\to H$ 是同态，$E$ 是 $H$ 的子群。证明 $\varphi^{-1}(E)\leq G$（即同态下子群的原像或拉回是子群）。若 $E\normal H$，证明 $\varphi^{-1}(E)\normal G$。推出 $\ker\varphi\normal G$。
  \item 设 $\varphi:G\to H$ 是群同态，核为 $K$，且 $a,b\in\varphi(G)$。设 $X\in G/K$ 是 $a$ 上方的纤维，$Y$ 是 $b$ 上方的纤维，即 $X=\varphi^{-1}(a)$、$Y=\varphi^{-1}(b)$。固定 $X$ 的一个元素 $u$。证明若 $XY=Z$ 在商群 $G/K$ 中成立，且 $w\in Z$，则存在某个 $v\in Y$ 使得 $uv=w$。[证明 $u^{-1}w\in Y$。]
  \item 设 $A$ 是阿贝尔群，$B$ 是 $A$ 的子群。证明 $A/B$ 是阿贝尔群。给出一个非阿贝尔群 $G$，它含有一个真正规子群 $N$，使得 $G/N$ 是阿贝尔群。
  \item 证明在商群 $G/N$ 中，对所有 $a\in\ZZ$，$(gN)^a=g^aN$。
  \item 用上一题证明，$G/N$ 中元素 $gN$ 的阶为 $n$，其中 $n$ 是使 $g^n\in N$ 的最小正整数；若不存在这样的正整数，则 $gN$ 有无限阶。举例说明 $gN$ 在 $G/N$ 中的阶可能严格小于 $g$ 在 $G$ 中的阶。
  \item 定义 $\varphi:\RR^\times\to\{\pm1\}$ 为 $\varphi(x)=x/|x|$。描述 $\varphi$ 的纤维并证明 $\varphi$ 是同态。
  \item 定义 $\pi:\RR^2\to\RR$ 为 $\pi((x,y))=x+y$。证明 $\pi$ 是满同态，并几何地描述其核与纤维。
  \item 设 $\varphi:\RR^\times\to\RR^\times$ 为 $x\mapsto |x|$。证明 $\varphi$ 是同态，求其像，并描述其核与纤维。
  \item 定义 $\varphi:\CC^\times\to\RR^\times$ 为 $\varphi(a+bi)=a^2+b^2$。证明 $\varphi$ 是同态，求其像，并几何地描述其核与纤维（作为平面中的子集）。
  \item 令 $\varphi:\ZZ/8\ZZ\to\ZZ/4\ZZ$ 为 $\varphi(\overline a)=\overline a$。说明这是良定义的满同态，并明确描述其纤维与核（说明良定义性时要注意 $\overline a$ 在定义域与值域中的含义不同）。
  \item 设 $F$ 是域，令
  \[
    G=\left\{\begin{pmatrix}a&b\\0&c\end{pmatrix}\mid a,b,c\in F,
    \ ac\neq0\right\}\leq \GL_2(F).
  \]
  \begin{enumerate}[label=(\alph*)]
    \item 证明映射 $\begin{pmatrix}a&b\\0&c\end{pmatrix}\mapsto a$ 是从 $G$ 到 $F^\times$ 的满同态。描述其纤维与核。
    \item 证明映射 $\begin{pmatrix}a&b\\0&c\end{pmatrix}\mapsto (a,c)$ 是从 $G$ 到 $F^\times\times F^\times$ 的满同态。描述其纤维与核。
    \item 令 $H=\left\{\begin{pmatrix}1&b\\0&1\end{pmatrix}\mid b\in F\right\}$。证明 $H$ 同构于 $F$ 的加法群。
  \end{enumerate}
  \item 设 $G$ 为实数加法群，$H$ 为绝对值为 $1$ 的复数乘法群（复平面中的单位圆 $S^1$），令 $\varphi:G\to H$ 为同态 $\varphi:r\mapsto e^{2\pi ir}$。在实线上画出位于 $\ker\varphi$ 中的点。类似地描述 $\varphi$ 在 $H$ 的点 $-1$、$i$ 和 $e^{4\pi i/3}$ 上方纤维中的元素。（这个同态的图 1 通常用下面的图 5 表示。）

\begin{center}
\begin{tikzpicture}[scale=.72, every node/.style={inner sep=1pt}]
  \draw[domain=0:720,smooth,variable=\t] plot ({.009*\t*cos(\t)},{.009*\t*sin(\t)+1.8});
  \node[right] at (2.6,1.8) {$G=\RR$};
  \draw[->] (0,.75)--(0,-.25) node[midway,right] {$\varphi$};
  \draw (0,-1.0) ellipse (1.9 and .55);
  \node[right] at (2.2,-1.0) {$H=S^1$};
  \node at (2.6,-1.7) {图 5};
\end{tikzpicture}
\end{center}
  \item 把前一题中的映射换成 $r\mapsto e^{4\pi ir}$，重复前一题。
  \item 考虑加法商群 $\QQ/\ZZ$。
  \begin{enumerate}[label=(\alph*)]
    \item 说明 $\QQ/\ZZ$ 的每个陪集都含有唯一一个满足 $0\leq q<1$ 的代表元 $q\in\QQ$。
    \item 说明 $\QQ/\ZZ$ 的每个元素都有有限阶，但其中有任意大阶的元素。
    \item 说明 $\QQ/\ZZ$ 是 $\RR/\ZZ$ 的挠子群（参见第 2.1 节练习 6）。
    \item 证明 $\QQ/\ZZ$ 同构于 $\CC^\times$ 中单位根的乘法群。
  \end{enumerate}
  \item 证明可除阿贝尔群关于任意真子群的商仍可除。推出 $\QQ/\ZZ$ 可除（参见第 2.4 节练习 19）。
  \item 设 $G$ 是群，$N\normal G$，令 $\overline G=G/N$。证明若 $G=\langle x,y\rangle$，则 $\overline G=\langle\overline x,\overline y\rangle$。更一般地，若 $G=\langle S\rangle$，其中 $S$ 是 $G$ 的任意子集，则 $\overline G=\langle\overline S\rangle$。
  \item 设 $G$ 是阶为 $16$ 的二面体群
  \[
    G=\langle r,s\mid r^8=s^2=1,rs=sr^{-1}\rangle,
  \]
  令 $\overline G=G/\langle r^4\rangle$。证明并完成以下内容：
  \begin{enumerate}[label=(\alph*)]
    \item $|\overline G|=8$。
    \item 将 $\overline G$ 的每个元素表示成 $\overline s^{\,a}\overline r^{\,b}$ 的形式。
    \item 求 (b) 中各元素的阶。
    \item 把 $rs$、$sr^{-2}s$、$s^{-1}r^{-1}sr$ 写成 (b) 中的形式。
    \item 证明 $\overline H=\langle\overline s,\overline r^{\,2}\rangle$ 是 $\overline G$ 的正规子群，且 $\overline H\cong V_4$。描述 $\overline H$ 在 $G$ 中的完全原像的同构类型。
    \item 求 $\overline G$ 的中心，并描述 $\overline G/Z(\overline G)$ 的同构类型。
  \end{enumerate}
  \item 设 $G$ 是阶为 $16$ 的拟二面体群（其格在第 2.5 节练习 11 中计算过）：
  \[
    G=\langle\sigma,\tau\mid \sigma^8=\tau^2=1,
    \ \sigma\tau=\tau\sigma^3\rangle,
  \]
  并令 $\overline G=G/\langle\sigma^4\rangle$，即 $G$ 关于由 $\sigma^4$ 生成的子群的商（这个子群是 $G$ 的中心，因而正规）。
  \begin{enumerate}[label=(\alph*)]
    \item 说明 $|\overline G|=8$。
    \item 将 $\overline G$ 的每个元素表示成 $\overline\tau^{\,a}\overline\sigma^{\,b}$ 的形式，其中 $a,b$ 为某些整数。
    \item 求 (b) 中列出的 $\overline G$ 的每个元素的阶。
    \item 把下列每个元素写成 (b) 中的形式 $\overline\tau^{\,a}\overline\sigma^{\,b}$：$\sigma\tau$、$\tau\sigma^{-2}\tau$、$\tau^{-1}\sigma^{-1}\tau\sigma$。
    \item 证明 $\overline G\cong D_8$。
  \end{enumerate}
  \item 设 $G$ 是第 2.5 节练习 14 中阶为 $16$ 的模群（其格在该练习中计算过）：
  \[
    G=\langle u,v\mid u^2=v^8=1,
    \ vu=uv^5\rangle,
  \]
  并令 $\overline G=G/\langle v^4\rangle$，即 $G$ 关于由 $v^4$ 生成的子群的商（这个子群包含在 $G$ 的中心中，因而正规）。
  \begin{enumerate}[label=(\alph*)]
    \item 说明 $|\overline G|=8$。
    \item 将 $\overline G$ 的每个元素表示成 $\overline u^{\,a}\overline v^{\,b}$ 的形式，其中 $a,b$ 为某些整数。
    \item 求 (b) 中列出的 $\overline G$ 的每个元素的阶。
    \item 把下列每个元素写成 (b) 中的形式 $\overline u^{\,a}\overline v^{\,b}$：$vu$、$uv^{-2}u$、$u^{-1}v^{-1}uv$。
    \item 证明 $\overline G$ 是阿贝尔群，并且同构于 $Z_2\times Z_4$。
  \end{enumerate}
  \item 设 $G\cong\ZZ/24\ZZ$，令 $\overline G=G/\langle\overline{12}\rangle$，其中为简化记号仍写 $\overline a$。
  \begin{enumerate}[label=(\alph*)]
    \item 说明 $\overline G=\{\overline0,\overline1,\ldots,\overline{11}\}$。
    \item 求 $\overline G$ 中每个元素的阶。
    \item 证明 $\overline G\cong\ZZ/12\ZZ$。
  \end{enumerate}
  \item 令 $G=Z_4\times Z_4$，用以下生成元与关系给出：
  \[
    G=\langle x,y\mid x^4=y^4=1,
    \ xy=yx\rangle.
  \]
  令 $\overline G=G/\langle x^2y^2\rangle$（注意阿贝尔群 $G$ 的每个子群都是正规子群）。
  \begin{enumerate}[label=(\alph*)]
    \item 说明 $|\overline G|=8$。
    \item 将 $\overline G$ 的每个元素表示成 $\overline x^{\,a}\overline y^{\,b}$ 的形式，其中 $a,b$ 为某些整数。
    \item 求 (b) 中列出的 $\overline G$ 的每个元素的阶。
    \item 证明 $\overline G\cong Z_4\times Z_2$。
  \end{enumerate}
  \item \begin{enumerate}[label=(\alph*)]
    \item 证明若 $H$ 和 $K$ 是群 $G$ 的正规子群，则交 $H\cap K$ 也是 $G$ 的正规子群。
    \item 证明任意非空正规子群族的交仍是正规子群（不要假设该族可数）。
  \end{enumerate}
  \item 证明任意非空正规子群族的并（见第 2.5 节的 join，即生成子群）仍是正规子群。
  \item 证明若 $N\normal G$，$H$ 是 $G$ 的任意子群，则 $N\cap H\normal H$。
  \item \begin{enumerate}[label=(\alph*)]
    \item 证明 $G$ 的子群 $N$ 正规，当且仅当对所有 $g\in G$，$gNg^{-1}\subseteq N$。
    \item 令 $G=\GL_2(\QQ)$，令 $N$ 为对角线为 $1$、上三角且整数项的矩阵子群，令 $g=\operatorname{diag}(2,1)$。说明 $gNg^{-1}\subseteq N$，但 $g$ 不正规化 $N$。
  \end{enumerate}
  \item 设 $a,b\in G$。
  \begin{enumerate}[label=(\alph*)]
    \item 证明 $ab$ 的共轭等于 $a$ 的共轭与 $b$ 的共轭的乘积。证明 $a$ 与任意一个共轭元素的阶相同。
    \item 证明 $a^{-1}$ 的共轭等于 $a$ 的共轭的逆。
    \item 若 $N=\langle S\rangle$，证明若对所有 $g\in G$，$gSg^{-1}\subseteq N$，则 $N\normal G$。
    \item 推出若 $N=\langle x\rangle$ 是循环群，则 $N$ 在 $G$ 中正规，当且仅当对每个 $g\in G$，$gxg^{-1}=x^k$ 对某个 $k\in\ZZ$ 成立。
    \item 设 $n$ 是正整数。证明由 $G$ 中所有阶为 $n$ 的元素生成的子群 $N$ 是 $G$ 的正规子群。
  \end{enumerate}
  \item 设 $N$ 是群 $G$ 的有限子群。说明 $gNg^{-1}\subseteq N$ 当且仅当 $gNg^{-1}=N$。推出
  \[
    N_G(N)=\{g\in G\mid gNg^{-1}\subseteq N\}.
  \]
  \item 设 $N$ 是 $G$ 的有限子群，且 $N=\langle S\rangle$。证明 $g\in G$ 正规化 $N$ 当且仅当 $gSg^{-1}\subseteq N$。
  \item 设 $N$ 是 $G$ 的有限子群，且 $G=\langle T\rangle$、$N=\langle S\rangle$。证明 $N\normal G$ 当且仅当对所有 $t\in T$ 和 $s\in S$，$tst^{-1}\in N$。
  \item 设 $N\leq G$，$g\in G$。证明 $gN=Ng$ 当且仅当 $g\in N_G(N)$。
  \item 证明若 $H\leq G$ 且 $N\normal H$，则 $H\leq N_G(N)$。推出 $N_G(N)$ 是 $G$ 中使 $N$ 为正规子群的最大子群。
  \item 证明 $Q_8$ 的每个子群都是正规的。对每个子群，求相应商群的同构类型。
  \item 求 $D_8$ 的所有正规子群，并对每个正规子群求相应商群的同构类型。
  \item 令
  \[
    D_{2n}=\langle r,s\mid r^n=s^2=1,rs=sr^{-1}\rangle
  \]
  是阶 $2n$ 的通常二面体群表示，令 $k$ 是整除 $n$ 的正整数。
  \begin{enumerate}[label=(\alph*)]
    \item 证明 $\langle r^k\rangle$ 是 $D_{2n}$ 的正规子群。
    \item 证明 $D_{2n}/\langle r^k\rangle\cong D_{2k}$。
  \end{enumerate}
  \item 证明 $\SL_n(F)\normal\GL_n(F)$，并描述商群的同构类型（参见第 2.1 节练习 9）。
  \item 证明若 $G/Z(G)$ 是循环群，则 $G$ 是阿贝尔群。[若 $G/Z(G)$ 以 $xZ(G)$ 为生成元，说明 $G$ 的每个元素都可写成 $x^az$，其中 $a\in\ZZ$，$z\in Z(G)$。]
  \item 设 $A$ 与 $B$ 是群。说明 $\{(a,1)\mid a\in A\}$ 是 $A\times B$ 的正规子群，且商群同构于 $B$。
  \item 设 $A$ 是阿贝尔群，$D=\{(a,a)\mid a\in A\}$ 是 $A\times A$ 的对角子群。证明 $D\normal A\times A$ 且 $(A\times A)/D\cong A$。
  \item 设 $A=S_3$ 是非阿贝尔群，$D=\{(a,a)\mid a\in A\}$ 是 $A\times A$ 的对角子群。证明 $D$ 不是 $A\times A$ 的正规子群。
  \item 设 $G$ 是群，$N\normal G$，令 $\overline G=G/N$。证明 $\overline x$ 与 $\overline y$ 在 $\overline G$ 中交换，当且仅当 $x^{-1}y^{-1}xy\in N$。（元素 $x^{-1}y^{-1}xy$ 称为 $x$ 与 $y$ 的\emph{交换子}，记为 $[x,y]$。）
  \item 设 $G$ 是群。证明
  \[
    N=\langle x^{-1}y^{-1}xy\mid x,y\in G\rangle
  \]
  是 $G$ 的正规子群，且 $G/N$ 是阿贝尔群（$N$ 称为 $G$ 的\emph{交换子子群}）。
  \item 假设 $H$ 与 $K$ 都是 $G$ 的正规子群且 $H\cap K=1$。证明对所有 $x\in H$ 与 $y\in K$，$xy=yx$。[证明 $x^{-1}y^{-1}xy\in H\cap K$。]
  \item 设 $\mathcal P=\{A_i\mid i\in I\}$ 是 $G$ 的一个分划，并假设 $\mathcal P$ 在如下“商运算”下成为群：计算 $A_i$ 与 $A_j$ 的乘积时，任取 $a_i\in A_i$ 与 $a_j\in A_j$，令 $A_iA_j$ 为含有 $a_ia_j$ 的 $\mathcal P$ 的元素（假设这个运算良定义）。证明 $\mathcal P$ 中含 $G$ 的单位元的元素是 $G$ 的正规子群，且 $\mathcal P$ 的元素正是这个子群的陪集（所以 $\mathcal P$ 就是通常意义下的商群）。
\end{enumerate}

\section{关于陪集与拉格朗日定理的进一步讨论}

本节继续研究商群。对有限群而言，最重要的不变量之一是群的阶，因此我们首先证明有限群商群的阶很容易计算：
\[
  |G/N|=\frac{|G|}{|N|}.
\]
事实上，我们将把这作为一个更一般结果——拉格朗日定理——的推论得到。这一定理是有限群论中最重要的组合结果之一，后文将反复使用。说明拉格朗日定理的一些简单推论后，我们将研究非正规子群的陪集相关的更细问题。

拉格朗日定理的证明直接且重要。它与上一节例 3 中计算 $|D_8/Z(D_8)|$ 所用的推理完全相同。

\medskip
\noindent\textbf{定理 8（拉格朗日定理）.} 若 $G$ 是有限群，$H$ 是 $G$ 的子群，则 $H$ 的阶整除 $G$ 的阶（即 $|H|\mid |G|$），并且 $H$ 在 $G$ 中的左陪集数等于
\[
  \frac{|G|}{|H|}.
\]

\begin{proof}
设 $|H|=n$，$H$ 在 $G$ 中的左陪集数为 $k$。由命题 4，$H$ 在 $G$ 中的左陪集分划 $G$。按左陪集定义，映射
\[
  H\to gH,
  \qquad h\mapsto gh
\]
是从 $H$ 到左陪集 $gH$ 的满射。左消去律说明该映射也是单射，因为 $gh_1=gh_2$ 蕴含 $h_1=h_2$。所以 $H$ 与 $gH$ 有相同阶：
\[
  |gH|=|H|=n.
\]
由于 $G$ 被分划为 $k$ 个互不相交、每个基数为 $n$ 的子集，$|G|=kn$。故
\[
  k=\frac{|G|}{n}=\frac{|G|}{|H|}.
\]
证明完成。
\end{proof}

\noindent\textbf{定义.} 若 $G$ 是群（可无限），$H\leq G$，则 $H$ 在 $G$ 中左陪集的数目称为 $H$ 在 $G$ 中的\emph{指数}，记为 $|G:H|$。

在有限群情形，$H$ 在 $G$ 中的指数为 $|G|/|H|$。若 $G$ 是无限群，则商 $|G|/|H|$ 没有意义。无限群可有有限指数或无限指数的子群（例如，$\{0\}$ 在 $\ZZ$ 中有无限指数，而对每个 $n>0$，$\langle n\rangle$ 在 $\ZZ$ 中指数为 $n$）。

下面给出拉格朗日定理的一些简单推论。

\medskip
\noindent\textbf{推论 9.} 若 $G$ 是有限群，$x\in G$，则 $x$ 的阶整除 $G$ 的阶。特别地，对所有 $x\in G$，有 $x^{|G|}=1$。

\begin{proof}
由命题 2.2，$|x|=|\langle x\rangle|$。第一部分由拉格朗日定理应用于 $H=\langle x\rangle$ 得出。第二个断言因为 $|G|$ 是 $x$ 的阶的倍数而显然成立。
\end{proof}

\medskip
\noindent\textbf{推论 10.} 若 $G$ 是素数阶 $p$ 的群，则 $G$ 是循环群，从而 $G\cong Z_p$。

\begin{proof}
取 $x\in G$，$x\neq1$。则 $|\langle x\rangle|>1$ 且 $|\langle x\rangle|$ 整除 $|G|$。由于 $|G|$ 是素数，必有 $|\langle x\rangle|=|G|$，故 $G=\langle x\rangle$ 是循环群（任一非单位元都是生成元）。由定理 2.4 得 $G\cong Z_p$。
\end{proof}

有了拉格朗日定理，我们再考察若干正规子群的例子。

\paragraph{例}
\begin{enumerate}[label=(\arabic*)]
  \item 令 $H=\langle(1\ 2\ 3)\rangle\leq S_3$，$G=S_3$。我们说明 $H\normal S_3$。如第 2.2 节所述，
  \[
    H\leq N_G(H)\leq G.
  \]
  由拉格朗日定理，$|H|$ 整除 $|N_G(H)|$，且 $|N_G(H)|$ 整除 $|G|$。由于 $|G|=6$、$|H|=3$，$N_G(H)$ 只能是 $H$ 或 $G$。直接计算
  \[
    (1\ 2)(1\ 2\ 3)(1\ 2)=(1\ 3\ 2)=(1\ 2\ 3)^{-1}.
  \]
  由于 $(1\ 2)=(1\ 2)^{-1}$，这说明 $(1\ 2)$ 把 $H$ 的一个生成元共轭到另一个生成元。由第 2.3 节练习 24，这足以证明 $(1\ 2)\in N_G(H)$。因此 $N_G(H)\neq H$，故 $N_G(H)=G$，即 $H\normal S_3$。这个论证说明，检查子群正规性时往往可减少到很少几个计算。下一个例子给出这一例子的推广。
  \item 设 $G$ 是任意群，$H$ 是指数 $2$ 的子群。证明 $H\normal G$。取 $g\in G-H$。按假设，$H$ 在 $G$ 中有两个左陪集 $1H$ 和 $gH$。由于 $1H=H$ 且陪集分划 $G$，必有 $gH=G-H$。同理，两个右陪集为 $H1=H$ 和 $Hg$，故 $Hg=G-H$。于是 $gH=Hg$，所以 $H$ 的每个左陪集都是右陪集。由定理 6，$H\normal G$。按指数定义，$|G/H|=2$，故 $G/H\cong Z_2$。

  要注意，$H$ 在这种情形下正规，不是因为可以对左右陪集选择同样的代表元 $1$ 与 $g$，而是因为一种抽屉原则在起作用：对任意子群 $H$ 都有 $1H=H=H1$，指数为 $2$ 的假设强迫剩余元素组成另一个陪集（无论是左还是右）。我们将看到，这个结果本身是下一章一个定理的特殊情形。

  这个结果说明 $Q_8$ 中的 $\langle i\rangle$、$\langle j\rangle$、$\langle k\rangle$ 都正规，且 $D_8$ 中的 $\langle s,r^2\rangle$、$\langle r\rangle$、$\langle sr,r^2\rangle$ 也正规。
  \item “是……的正规子群”这一性质不具有传递性。例如
  \[
    \langle s\rangle\normal\langle s,r^2\rangle\normal D_8
  \]
  （每个子群在下一个群中指数为 $2$），但 $\langle s\rangle$ 不是 $D_8$ 的正规子群，因为 $rsr^{-1}=sr^2\notin\langle s\rangle$。
\end{enumerate}

现在考察一些非正规子群的例子。虽然在阿贝尔群中每个子群都正规，但在非阿贝尔群中并非如此（某种意义上 $Q_8$ 是唯一例外）。事实上，有些群 $G$ 的正规子群只有平凡的 $1$ 和 $G$。这种群称为\emph{单群}（simple 并不意味着 easy）。单群在一般群的研究中扮演重要角色，这一角色将在第 4 节讨论。这里先强调：群 $G$ 的子群不一定都是 $G$ 的正规子群；事实上，正规子群在 $G$ 中可能相当稀少。寻找给定群的正规子群通常是一个非常非平凡的问题。

\paragraph{例}
\begin{enumerate}[label=(\arabic*)]
  \item 令 $H=\langle(1\ 2)\rangle\leq S_3$。由于 $H$ 在 $S_3$ 中指数为素数 $3$，由拉格朗日定理，$N_{S_3}(H)$ 只能是 $H$ 或 $S_3$。直接计算
  \[
    (1\ 3)(1\ 2)(1\ 3)^{-1}=(1\ 3)(1\ 2)(1\ 3)=(2\ 3)\notin H,
  \]
  所以 $N_{S_3}(H)\neq S_3$，即 $H$ 不是 $S_3$ 的正规子群。也可以通过左右陪集看出这一点，例如
  \[
    (1\ 3)H=\{(1\ 3),(1\ 2\ 3)\},\qquad
    H(1\ 3)=\{(1\ 3),(1\ 3\ 2)\}.
  \]
  由于左陪集 $(1\ 3)H$ 是包含 $(1\ 3)$ 的唯一左陪集，右陪集 $H(1\ 3)$ 不可能是左陪集。还要注意，用代表元相乘在 $H$ 的左陪集上定义“群运算”甚至不是良定义的。例如考虑两个左陪集 $1H$ 和 $(1\ 3)H$ 的乘积。$1$ 和 $(1\ 2)$ 都是陪集 $1H$ 的代表元，但
  \[
    1\cdot(1\ 3)=(1\ 3),\qquad
    (1\ 2)(1\ 3)=(1\ 3\ 2),
  \]
  它们不在同一个左陪集中；若陪集乘积与代表元选择无关，本应在同一个左陪集中。这是定理 6 的一个例子：子群的陪集形成群，当且仅当该子群正规。
  \item 令 $G=S_n$，固定 $i\in\{1,2,\ldots,n\}$。如第 2.2 节，令
  \[
    G_i=\{\sigma\in G\mid \sigma(i)=i\}
  \]
  为点 $i$ 的稳定子。若 $\tau\in G$ 且 $\tau(i)=j$，由定义可直接看出，对所有 $\sigma\in G_i$，$\tau\sigma(i)=j$。此外，若 $\mu\in G$ 且 $\mu(i)=j$，则 $\tau^{-1}\mu(i)=i$，即 $\tau^{-1}\mu\in G_i$，故 $\mu\in\tau G_i$。因此
  \[
    \tau G_i=\{\mu\in G\mid \mu(i)=j\},
  \]
  即左陪集 $\tau G_i$ 由 $S_n$ 中把 $i$ 送到 $j$ 的所有置换组成。显然不同左陪集不交，且不同左陪集数等于整数 $i$ 在 $G$ 作用下的不同像数，即共有 $n$ 个左陪集。因此 $|G:G_i|=n$。类似地，令 $k=\tau^{-1}(i)$，则 $\tau(k)=i$，有
  \[
    G_i\tau=\{\lambda\in G\mid \lambda(k)=i\}.
  \]
  若 $n>2$，可取非单位元 $\tau$ 使 $\tau G_i\neq G_i\tau$，因为当然存在把 $i$ 送到 $j$ 但不把 $k$ 送到 $i$ 的置换。因此 $G_i$ 不是正规子群。事实上由第 1 节练习 30，$N_G(G_i)=G_i$，所以 $G_i$ 在某种意义上离在 $S_n$ 中正规很远。这个例子推广了前一个例子。
  \item 在 $D_8$ 中，唯一正规的阶 $2$ 子群是中心 $\langle r^2\rangle$。
\end{enumerate}

随着理论的发展，我们会看到更多非正规子群的例子。

拉格朗日定理的完全逆命题不成立：若 $G$ 是有限群且 $n$ 整除 $|G|$，则 $G$ 不一定有阶为 $n$ 的子群。例如，令 $A$ 为正四面体的对称群。由第 1.2 节练习 9，$|A|=12$。若 $A$ 有一个阶为 $6$ 的子群 $H$，则 $H$ 在 $A$ 中指数为 $2$，故 $H\normal A$ 且 $A/H\cong Z_2$。由于商群阶为 $2$，商中每个元素的平方都是单位元，所以对所有 $g\in A$，$(gH)^2=1H$，即 $g^2\in H$。若 $g$ 是 $A$ 中阶为 $3$ 的元素，则 $g=(g^2)^2\in H$。于是 $H$ 必须包含 $A$ 的所有阶为 $3$ 的元素。这矛盾，因为 $|H|=6$，但容易列出正四面体的 $8$ 个阶为 $3$ 的旋转。

拉格朗日定理有一些部分逆命题。对于有限阿贝尔群，拉格朗日定理的完全逆命题成立：对 $|G|$ 的每个因子 $n$，阿贝尔群 $G$ 都有阶为 $n$ 的子群（事实上，这在比“阿贝尔”更弱的假设下也成立；第 6 章会看到）。对任意有限群都成立的一个部分逆命题是：

\medskip
\noindent\textbf{定理 11（柯西定理）.} 若 $G$ 是有限群，$p$ 是整除 $|G|$ 的素数，则 $G$ 有一个阶为 $p$ 的元素。

\begin{proof}
下一章将给出一个证明；练习 9 还勾勒了另一个优美证明。
\end{proof}

适用于任意有限群的最强拉格朗日逆命题是：

\medskip
\noindent\textbf{定理 12（西洛）.} 若有限群 $G$ 的阶为 $p^\alpha m$，其中 $p$ 是素数且 $p\nmid m$，则 $G$ 有一个阶为 $p^\alpha$ 的子群。

我们将在下一章证明该定理，并得到关于阶为 $p^\alpha$ 的子群数目的更多信息。

本节最后给出一些关于陪集的有用结果。

\noindent\textbf{定义.} 设 $H$ 和 $K$ 是一个群的子群，定义
\[
  HK=\{hk\mid h\in H,
  \ k\in K\}.
\]

\medskip
\noindent\textbf{命题 13.} 若 $H$ 和 $K$ 是某群的有限子群，则
\[
  |HK|=\frac{|H|\,|K|}{|H\cap K|}.
\]

\begin{proof}
注意 $HK$ 是 $K$ 的左陪集的并：
\[
  HK=\bigcup_{h\in H}hK.
\]
由于 $K$ 的每个陪集含 $|K|$ 个元素，只需求形如 $hK$（$h\in H$）的不同左陪集数。对 $h_1,h_2\in H$，有
\[
  h_1K=h_2K
  \Longleftrightarrow h_2^{-1}h_1\in K
  \Longleftrightarrow h_2^{-1}h_1\in H\cap K
  \Longleftrightarrow h_1(H\cap K)=h_2(H\cap K).
\]
因此，$hK$（$h\in H$）的不同陪集数等于 $H$ 中 $H\cap K$ 的不同陪集数。由拉格朗日定理，这个数为 $|H|/|H\cap K|$。于是 $HK$ 由 $|H|/|H\cap K|$ 个不同的 $K$ 的陪集组成，每个陪集有 $|K|$ 个元素，得公式。
\end{proof}

注意命题 13 不假设 $HK$ 是子群。例如，若 $G=S_3$，$H=\langle(1\ 2)\rangle$，$K=\langle(2\ 3)\rangle$，则 $|H|=|K|=2$、$|H\cap K|=1$，所以 $|HK|=4$。由拉格朗日定理，$HK$ 不可能是子群。结果我们必须有 $S_3=\langle(1\ 2),(2\ 3)\rangle$。

\medskip
\noindent\textbf{命题 14.} 若 $H$ 和 $K$ 是某群的子群，则 $HK$ 是子群，当且仅当 $HK=KH$。

\begin{proof}
先假设 $HK=KH$。取 $a,b\in HK$，证明 $ab^{-1}\in HK$。设
\[
  a=h_1k_1,
  \qquad b=h_2k_2,
\]
其中 $h_1,h_2\in H$、$k_1,k_2\in K$。则 $b^{-1}=k_2^{-1}h_2^{-1}$，所以
\[
  ab^{-1}=h_1k_1k_2^{-1}h_2^{-1}.
\]
令 $k_3=k_1k_2^{-1}\in K$，$h_3=h_2^{-1}\in H$。由于 $HK=KH$，$k_3h_3=h_4k_4$，其中 $h_4\in H$、$k_4\in K$。于是
\[
  ab^{-1}=h_1h_4k_4\in HK.
\]
由子群判别法，$HK$ 是子群。

反过来，假设 $HK$ 是 $G$ 的子群。由于 $K\leq HK$ 且 $H\leq HK$，由子群封闭性得 $KH\subseteq HK$。为证反向包含，取 $hk\in HK$。因为 $HK$ 是子群，可写 $hk=a^{-1}$，其中 $a\in HK$。若 $a=h_1k_1$，则
\[
  hk=(h_1k_1)^{-1}=k_1^{-1}h_1^{-1}\in KH.
\]
证明完成。
\end{proof}

注意 $HK=KH$ 并不意味着 $H$ 的元素与 $K$ 的元素交换，而是意味着每个乘积 $hk$ 都可写成某个 $k'h'$ 的形式（不必 $h'=h$、$k'=k$），反之亦然。例如，若 $G=D_{2n}$、$H=\langle r\rangle$、$K=\langle s\rangle$，则 $G=HK=KH$，所以 $HK$ 是子群；但 $rs=sr^{-1}$，$H$ 的元素并不与 $K$ 的元素交换。这是下面充分条件的一个例子。

\medskip
\noindent\textbf{推论 15.} 若 $H$ 和 $K$ 是 $G$ 的子群，且 $H\leq N_G(K)$，则 $HK$ 是 $G$ 的子群。特别地，若 $K\normal G$，则对任意 $H\leq G$ 都有 $HK\leq G$。

\begin{proof}
证明 $HK=KH$。若 $h\in H$、$k\in K$，由假设 $hkh^{-1}\in K$，故
\[
  hk=(hkh^{-1})h\in KH.
\]
这说明 $HK\subseteq KH$。类似地，$kh=h(h^{-1}kh)\in HK$，得反向包含。于是由命题 14 得结论。
\end{proof}

\noindent\textbf{定义.} 若 $A$ 是 $N_G(K)$（或 $C_G(K)$）的任意子集，则称 $A$ \emph{正规化} $K$（相应地，\emph{中心化} $K$）。

用此术语，推论 15 说：若 $H$ 正规化 $K$，则 $HK$ 是子群（同样，若 $K$ 正规化 $H$，则 $HK$ 也是子群）。

有时，可仅用命题 13 的阶公式证明一个有限群是它的两个子群的乘积。例如，令 $G=S_4$，令 $H=D_8$，令 $K=\langle(1\ 2\ 3)\rangle$，这里把 $D_8$ 看成 $S_4$ 的子群，即在固定标号下把每个正方形对称变换与其对四个顶点的置换相对应。由拉格朗日定理，$H\cap K=1$。命题 13 给出 $|HK|=24$，所以 $HK=S_4$。既然 $HK$ 是群，$HK=KH$。验证 $H$ 和 $K$ 都不正规化对方留作练习。

最后，本章一直使用子群的左陪集。同样的组合结果也可用右陪集证明。对正规子群这当然显然，因为左右陪集相同；但对非正规子群，某些左陪集不是任何代表元下的右陪集，所以需要一些简单验证。例如，拉格朗日定理给出，在有限群 $G$ 中，子群 $H$ 的右陪集数也为
\[
  \frac{|G|}{|H|}.
\]
因此在有限群中，$H$ 的左陪集数等于右陪集数，尽管一般左陪集不是右陪集。练习 12 说明这对无限群也成立。因此就纯组合目的而言，可以使用左陪集或右陪集（但当需要 $G$ 的一个分划时不能混用）。我们一贯使用左陪集有些任意，不过在下一章讨论陪集上的作用时会有一些好处。读者可能在某些书中见到记号 $H\backslash G$ 表示 $H$ 在 $G$ 中的右陪集集合。某些文章中也可能用 $G/H$ 表示 $H$ 在 $G$ 中的左陪集集合，即使 $H$ 不正规（此时 $G/H$ 称为左陪集的陪集空间）。本书不采用这种记号。

\subsection*{练习}
设 $G$ 是群。
\begin{enumerate}[label=\arabic*.]
  \item 下列哪些数可以作为阶为 $120$ 的群的子群阶：$1,2,5,7,9,15,60,240$？对每个允许的阶，给出相应指数。
  \item 证明第 2.5 节给出的 $S_3$ 子群格是正确的。
  \item 证明第 2.5 节给出的 $Q_8$ 子群格是正确的。
  \item 证明若 $|G|=pq$，其中 $p$、$q$ 是素数（不必相异），则 $G$ 是阿贝尔群或 $Z(G)=1$。[见第 1 节练习 36。]
  \item 设 $H$ 是 $G$ 的子群，固定 $g\in G$。
  \begin{enumerate}[label=(\alph*)]
    \item 证明 $gHg^{-1}$ 是 $G$ 的子群，且与 $H$ 同阶。
    \item 推出若 $n\in\ZZ^+$ 且 $H$ 是 $G$ 的唯一阶为 $n$ 的子群，则 $H\normal G$。
  \end{enumerate}
  \item 设 $H\leq G$，$g\in G$。证明若右陪集 $Hg$ 等于 $G$ 中 $H$ 的某个左陪集，则它等于左陪集 $gH$，且 $g\in N_G(H)$。
  \item 设 $H\leq G$，在 $G$ 上定义关系 $a\sim b$ 当且仅当 $b^{-1}a\in H$。证明 $\sim$ 是等价关系，并描述每个 $a\in G$ 的等价类。用此证明命题 4。
  \item 证明若 $H$ 与 $K$ 是 $G$ 的有限子群且二者阶互素，则 $H\cap K=1$。
  \item 本题勾勒 James McKay 给出的柯西定理证明。令 $G$ 为有限群，$p$ 为整除 $|G|$ 的素数。令 $S$ 表示 $G$ 中 $p$ 元组的集合，且坐标乘积为 $1$：
  \[
    S=\{(x_1,x_2,
    \ldots,x_p)\mid x_i\in G,
    \ x_1x_2\cdots x_p=1\}.
  \]
  \begin{enumerate}[label=(\alph*)]
    \item 说明 $S$ 有 $|G|^{p-1}$ 个元素，因此其阶被 $p$ 整除。
    \item 定义 $S$ 上关系 $\alpha\sim\beta$，若 $\beta$ 是 $\alpha$ 的一个循环置换。证明 $S$ 的元素的循环置换仍在 $S$ 中。
    \item 证明 $\sim$ 是等价关系。
    \item 证明一个等价类含单个元素，当且仅当它形如 $(x,x,
    \ldots,x)$ 且 $x^p=1$。
    \item 证明每个等价类阶为 $1$ 或 $p$（用到 $p$ 是素数）。推出 $|G|^{p-1}=k+pd$，其中 $k$ 是大小为 $1$ 的等价类数，$d$ 是大小为 $p$ 的等价类数。
    \item 由于 $\{(1,1,\ldots,1)\}$ 是大小为 $1$ 的等价类，由 (e) 推出 $G$ 中存在非单位元 $x$ 使 $x^p=1$，即 $G$ 含有阶为 $p$ 的元素。[说明 $p\mid k$，所以 $k>1$。]
  \end{enumerate}
  \item 设 $H$ 与 $K$ 是（可无限）群 $G$ 的有限指数子群，$|G:H|=m$、$|G:K|=n$。证明
  \[
    \operatorname{lcm}(m,n)\leq |G:H\cap K|\leq mn.
  \]
  推出若 $m$ 与 $n$ 互素，则 $|G:H\cap K|=|G:H|\,|G:K|$。
  \item 设 $H\leq K\leq G$。证明 $|G:H|=|G:K|\,|K:H|$（不要假设 $G$ 有限）。
  \item 设 $H\leq G$。证明映射 $x\mapsto x^{-1}$ 把 $H$ 的每个左陪集送到 $H$ 的一个右陪集，并给出左陪集集合与右陪集集合之间的双射。
  \item 固定正方形顶点的任意标号，并用此把 $D_8$ 识别为 $S_4$ 的子群。证明 $D_8$ 的元素与 $\langle(1\ 2\ 3)\rangle$ 的元素在 $S_4$ 中不交换。
  \item 证明 $S_4$ 没有阶为 $8$ 的正规子群，也没有阶为 $3$ 的正规子群。
  \item 令 $G=S_n$，固定 $i\in\{1,2,
  \ldots,n\}$，令 $G_i$ 为 $i$ 的稳定子。证明 $G_i\cong S_{n-1}$。
  \item 在乘法群 $(\ZZ/p\ZZ)^\times$ 中使用拉格朗日定理证明费马小定理：若 $p$ 是素数，则对所有 $a\in\ZZ$，$a^p\equiv a\pmod p$。
  \item 设 $p$ 是素数，$n$ 是正整数。求 $p$ 在 $(\ZZ/(p^n-1)\ZZ)^\times$ 中的阶，并推出 $n\mid\varphi(p^n-1)$（这里 $\varphi$ 是欧拉函数）。
  \item 设 $G$ 是有限群，$H\leq G$，$N\normal G$。证明若 $|H|$ 与 $|G:N|$ 互素，则 $H\leq N$。
  \item 证明若 $N$ 是有限群 $G$ 的正规子群，且 $(|N|,|G:N|)=1$，则 $N$ 是 $G$ 的唯一阶为 $|N|$ 的子群。
  \item 若 $A$ 是阿贝尔群且 $A\normal G$，$B$ 是 $G$ 的任意子群，证明 $A\cap B\normal AB$。
  \item 证明 $\QQ$ 没有有限指数的真子群。推出 $\QQ/\ZZ$ 没有有限指数的真子群。[回忆第 1.6 节练习 21 和第 1 节练习 15。]
  \item 在乘法群 $(\ZZ/n\ZZ)^\times$ 中使用拉格朗日定理证明欧拉定理：若整数 $a$ 与 $n$ 互素，则
  \[
    a^{\varphi(n)}\equiv1\pmod n,
  \]
  其中 $\varphi$ 表示欧拉函数。
  \item 求 $3^{3^{100}}$ 的最后两位。[先求 $3^{100}\bmod\varphi(100)$，再用上一题。]
\end{enumerate}

\section{同构定理}

本节推导第 1 节中商群与同态关系的一些直接后果。特别地，我们考虑商群 $G/N$ 的子群格与群 $G$ 的子群格之间的关系。第一个结果重述第 1 节中关于同态像与核商群之间关系的观察（有时称为同态基本定理）：

\medskip
\noindent\textbf{定理 16（第一同构定理）.} 若 $\varphi:G\to H$ 是群同态，则 $\ker\varphi\normal G$，且
\[
  G/\ker\varphi\cong \varphi(G).
\]

\medskip
\noindent\textbf{推论 17.} 设 $\varphi:G\to H$ 是群同态。
\begin{enumerate}[label=(\arabic*)]
  \item $\varphi$ 是单射，当且仅当 $\ker\varphi=1$。
  \item $|G:\ker\varphi|=|\varphi(G)|$。
\end{enumerate}

\begin{proof}
练习。
\end{proof}

在讨论抽象向量空间时，我们会看到推论 17(2) 给出线性变换理论中可能已熟悉的公式：若 $\varphi:V\to W$ 是向量空间的线性变换，则
\[
  \dim V=\operatorname{rank}\varphi+\operatorname{nullity}\varphi.
\]

\medskip
\noindent\textbf{定理 18（第二同构定理或菱形同构定理）.} 设 $G$ 是群，$A$ 和 $B$ 是 $G$ 的子群，并假设 $A\leq N_G(B)$。则 $AB$ 是 $G$ 的子群，$B\normal AB$，$A\cap B\normal A$，且
\[
  AB/B\cong A/(A\cap B).
\]

\begin{proof}
由推论 15，$AB$ 是 $G$ 的子群。由于 $A\leq N_G(B)$ 且 $B\leq N_G(B)$ 显然成立，所以 $AB\leq N_G(B)$，即 $B$ 是子群 $AB$ 的正规子群。

因为 $B\normal AB$，商群 $AB/B$ 有定义。定义
\[
  \varphi:A\to AB/B,
  \qquad \varphi(a)=aB.
\]
由于 $AB/B$ 中的群运算良定义，容易看出 $\varphi$ 是同态：
\[
  \varphi(a_1a_2)=(a_1a_2)B=a_1B\cdot a_2B=\varphi(a_1)\varphi(a_2).
\]
等价地，$\varphi$ 是自然投影 $\pi:AB\to AB/B$ 限制到子群 $A$ 上，所以也是同态。由 $AB$ 的定义可知 $\varphi$ 是满射。商 $AB/B$ 的单位元是陪集 $1B$，所以 $\ker\varphi$ 由满足 $aB=1B$ 的 $a\in A$ 组成；由命题 4，这正是 $a\in B$，即 $\ker\varphi=A\cap B$。由第一同构定理，$A\cap B\normal A$，且 $A/(A\cap B)\cong AB/B$。
\end{proof}

在 $A\leq N_G(B)$ 的特殊情形下，这给出了命题 13 的阶公式的一个新证明。这个定理称为菱形同构，是因为其中涉及 $G$ 的子群格中的菱形部分（见图 6）。格线上的标记表示哪些商群同构。商 $AB/A$ 未必是群（即 $A$ 未必在 $AB$ 中正规），但仍有

\begin{center}
\begin{tikzpicture}[scale=.9, every node/.style={inner sep=1pt}]
  \node (G) at (0,3.6) {$G$};
  \node (AB) at (0,2.7) {$AB$};
  \node (A) at (-1.05,1.8) {$A$};
  \node (B) at (1.05,1.8) {$B$};
  \node (I) at (0,.9) {$A\cap B$};
  \node (one) at (0,0) {$1$};
  \draw (G)--(AB);
  \draw (AB)--(A) node[midway,above left] {$\sim$};
  \draw (AB)--(B);
  \draw (A)--(I);
  \draw (B)--(I) node[midway,above right] {$\sim$};
  \draw (I)--(one);
  \node at (2.1,.2) {图 6};
\end{tikzpicture}
\end{center}
\[
  |AB:A|=|B:A\cap B|.
\]

第三同构定理讨论商群的商群。

\medskip
\noindent\textbf{定理 19（第三同构定理）.} 设 $G$ 是群，$H$ 和 $K$ 是 $G$ 的正规子群，且 $H\leq K$。则 $K/H\normal G/H$，并且
\[
  (G/H)/(K/H)\cong G/K.
\]
若用横线表示模 $H$ 的商，则可写作
\[
  \overline G/\overline K\cong G/K.
\]

\begin{proof}
$K/H\normal G/H$ 的验证留作简单练习。定义
\[
  \varphi:G/H\to G/K,
  \qquad gH\mapsto gK.
\]
为说明 $\varphi$ 良定义，设 $g_1H=g_2H$。则 $g_1=g_2h$，其中 $h\in H$。由于 $H\leq K$，有 $h\in K$，故 $g_1K=g_2K$，于是 $\varphi(g_1H)=\varphi(g_2H)$。因为 $g$ 可在 $G$ 中任意选取，$\varphi$ 是满同态。最后
\begin{align*}
  \ker\varphi
  &=\{gH\in G/H\mid \varphi(gH)=1K\}\\
  &=\{gH\in G/H\mid gK=1K\}\\
  &=\{gH\in G/H\mid g\in K\}=K/H.
\end{align*}
由第一同构定理，$(G/H)/(K/H)\cong G/K$。
\end{proof}

记忆第三同构定理的一个简单口诀是：“倒过来并消去”（像分式一样）。这个定理说明，对商群再取商群不会带来新的结构信息。

最后一个同构定理描述商群 $G/N$ 的子群格与 $G$ 的子群格之间的关系。$G/N$ 的子群格可以通过把 $G$ 的子群格中的群 $N$ 压缩为单位元直接读出。更精确地说，包含 $N$ 的 $G$ 的子群与 $G/N$ 的子群之间存在一一对应；因此 $G/N$ 的子群格（或一个同构拷贝）出现在 $G$ 的子群格中，即位于 $N$ 与 $G$ 之间的那些子群。特别地，$G/N$ 的子群格出现在 $G$ 的子群格“顶部”，这正是本章开头提到的现象。

\medskip
\noindent\textbf{定理 20（第四同构定理或格同构定理）.} 设 $G$ 是群，$N\normal G$。则从包含 $N$ 的 $G$ 的子群 $A$ 到 $G/N$ 的子群 $\overline A=A/N$ 有一个双射。特别地，$G/N$ 的每个子群都形如 $A/N$，其中 $A$ 是 $G$ 的一个包含 $N$ 的子群（即它在自然投影同态下的原像）。此双射具有以下性质：若 $A,B\leq G$ 且 $N\leq A$、$N\leq B$，则
\begin{enumerate}[label=(\arabic*)]
  \item $A\leq B$ 当且仅当 $\overline A\leq \overline B$；
  \item 若 $A\leq B$，则 $|B:A|=|\overline B:\overline A|$；
  \item $\overline{\langle A,B\rangle}=\langle\overline A,\overline B\rangle$；
  \item $\overline{A\cap B}=\overline A\cap\overline B$；
  \item $A\normal G$ 当且仅当 $\overline A\normal G/N$。
\end{enumerate}

\begin{proof}
$G/N$ 中子群的完全原像由第 1 节练习 1 是 $G$ 的子群。定理中需要检查的许多细节都完全直接，因此证明留作练习。
\end{proof}

\paragraph{例}
\begin{enumerate}[label=(\arabic*)]
  \item 令 $G=Q_8$，令 $N=\langle -1\rangle$。$G/N$ 的子群格的同构拷贝由 $G$ 的子群格中包含 $N$ 的那部分组成。我们此前已证明 $Q_8/\langle -1\rangle\cong V_4$，而两个格确实一致。
  \item 同样的过程可从 $D_8$ 的子群格中得到 $D_8/\langle r^2\rangle$ 的子群格：它由包含 $\langle r^2\rangle$ 的那些子群组成。
\end{enumerate}

注意在第二个例子中，$G$ 的某些子群并不直接对应于商群 $G/N$ 的子群，即那些不包含正规子群 $N$ 的子群。这是因为子群 $N$ 在 $G/N$ 中投影为一个点，因此 $G$ 的若干子群可以投影到商群中的同一个子群。$G$ 的子群 $H$ 在自然投影同态下的像与 $G$ 的子群 $HN$ 的像相同，而 $HN$ 包含 $N$。反过来，$G/N$ 的子群 $\overline H$ 的原像包含 $N$，并且是 $G$ 中唯一一个包含 $N$ 且在 $G/N$ 中像为 $\overline H$ 的子群。显式出现在 $G/N$ 的子群格中的，是 $G$ 中包含 $N$ 的那些子群。

上面两个阶为 $8$ 的群的格强调了这样一个事实：仅知道 $G/N$ 与 $N$ 的同构类型，一般不能确定 $G$ 的同构类型，因为
\[
  Q_8/\langle -1\rangle\cong D_8/\langle r^2\rangle
  \quad\text{且}\quad
  \langle -1\rangle\cong\langle r^2\rangle,
\]
但 $Q_8$ 与 $D_8$ 不同构。下一节还会进一步讨论这个问题。

我们经常在子群格中标出一个子群在另一个子群中的指数，例如
\[
\begin{array}{c}
A\\[-1mm]
|\,n\\[-1mm]
B
\end{array}
\]
其中整数 $n=|A:B|$。例如，$Q_8$ 和 $D_8$ 子群格中的所有实线边都可标为 $2$。于是任一子群 $A$ 的阶等于从单位子群向上到 $A$ 的任意路径上所有标号整数的乘积。又由定理 20(2)，若模去一个包含在 $B$ 中的 $G$ 的正规子群，这些指数在商中保持不变；即对应于商群子群格的那部分 $G$ 的格也具有商群中的正确指数。

最后说明一点关于在商群上定义同态的注意事项。在同构定理的证明过程中，我们遇到过这样的情形：通过给出陪集 $gN$ 的代表元 $g$ 上的值来定义商群 $G/N$ 上的同态 $\varphi$。每次都必须证明 $\varphi$ 良定义，即与 $g$ 的选择无关。实际上，我们是在 $G$ 本身上定义一个同态 $\Phi$，然后用 $\Phi(g)$ 指定 $\varphi(gN)$。与 $g$ 无关等价于要求 $\Phi$ 在 $N$ 上平凡，即
\[
  \varphi \text{ 在 } G/N \text{ 上良定义}
  \quad\Longleftrightarrow\quad N\leq\ker\Phi.
\]
这给出一种在商群上定义同态的简单准则：先在 $G$ 上定义同态，再检查 $N$ 包含于其核。在这种情形下，称同态 $\Phi$ \emph{通过 $N$ 分解}，而 $\varphi$ 是 $G/N$ 上\emph{诱导的同态}。可用交换图表示为图 7，其中图表示 $\Phi=\varphi\circ\pi$，即 $G$ 中元素在 $H$ 中的像不依赖于在图中选择哪条路径。若如此，则称该图\emph{交换}。

\begin{center}
\begin{tikzpicture}[scale=1.0, every node/.style={inner sep=1pt}]
  \node (G) at (0,1.4) {$G$};
  \node (Q) at (3,1.4) {$G/N$};
  \node (H) at (3,0) {$H$};
  \draw[->] (G)--node[above] {$\pi$} (Q);
  \draw[->] (Q)--node[right] {$\varphi$} (H);
  \draw[->] (G)--node[below left] {$\Phi$} (H);
  \node at (4.1,0) {图 7};
\end{tikzpicture}
\end{center}

至此，我们已经发展了足够的背景材料；现在可以阅读第 6.3 节关于自由群与表示的内容。

\subsection*{练习}
设 $G$ 是群。
\begin{enumerate}[label=\arabic*.]
  \item 设 $F$ 是阶为 $q$ 的有限域，$n\in\ZZ^+$。证明 $|\GL_n(F):\SL_n(F)|=q-1$。[见第 1 节练习 35。]
  \item 证明格同构定理的所有部分。
  \item 证明若 $H$ 是 $G$ 的素数指数 $p$ 的正规子群，则对所有 $K\leq G$，要么 $K\leq H$，要么 $G=HK$ 且 $|K:K\cap H|=p$。
  \item 设 $C\normal A$，$D\normal B$。证明 $(C\times D)\normal(A\times B)$，且
  \[
    (A\times B)/(C\times D)\cong (A/C)\times(B/D).
  \]
  \item 设 $QD_{16}=\langle\sigma,\tau\rangle$ 是第 2.5 节练习 11 中的拟二面体群。证明 $\langle\sigma^4\rangle$ 在 $QD_{16}$ 中正规，并用格同构定理画出 $QD_{16}/\langle\sigma^4\rangle$ 的子群格。哪个阶为 $8$ 的群有相同的格？用 $QD_{16}/\langle\sigma^4\rangle$ 的生成元与关系决定该群的同构类型。
  \item 设 $M=\langle v,u\rangle$ 是第 2.5 节练习 14 中阶为 $16$ 的模群。证明 $\langle v^4\rangle$ 在 $M$ 中正规，并用格同构定理画出 $M/\langle v^4\rangle$ 的子群格。哪个阶为 $8$ 的群有相同的格？用生成元与关系决定该商群的同构类型。
  \item 设 $M$ 和 $N$ 是 $G$ 的正规子群，且 $G=MN$。证明
  \[
    G/(M\cap N)\cong (G/M)\times(G/N).
  \]
  [画出格。]
  \item 设 $p$ 是素数，$G$ 是 $\CC$ 中所有 $p$ 次幂阶单位根的群（参见第 2.4 节练习 18）。证明映射 $z\mapsto z^p$ 是满同态。推出 $G$ 同构于它自身的一个真商。
  \item 设 $p$ 是素数，$|G|=p^am$，且 $p\nmid m$。设 $P$ 是 $G$ 的阶为 $p^a$ 的子群，$N\normal G$ 且 $|N|=p^bn$，其中 $p\nmid n$。证明 $|P\cap N|=p^b$，且 $|PN/N|=p^{a-b}$。（子群 $P$ 称为 $G$ 的一个西洛 $p$-子群。此练习说明任一西洛 $p$-子群与正规子群 $N$ 的交是 $N$ 的西洛 $p$-子群。）
  \item 推广上一题如下：有限群 $G$ 的子群 $H$ 称为 Hall 子群，如果它在 $G$ 中的指数与其阶互素：$(|G:H|,|H|)=1$。证明若 $H$ 是 $G$ 的 Hall 子群，且 $N\normal G$，则 $H\cap N$ 是 $N$ 的 Hall 子群，而 $HN/N$ 是 $G/N$ 的 Hall 子群。
\end{enumerate}

\section{合成列与 Hölder 纲领}

上一节关于格的讨论给我们一种直观图像：商群 $G/N$ 的结构（例如格）描述了正规子群 $N$ “上方”的 $G$ 的结构。虽然这种说法有些模糊，但至少说明了有限群论中（甚至某些无限群论分支中）最强有力技术之一的驱动力：归纳法的使用。许多时候，归纳过程的应用类似于下面对柯西定理一个特殊情形的证明。尽管柯西定理对任意群成立（参见第 2 节练习 9），但下面的证明很好地展示了如何利用正规子群 $N$ 与商群 $G/N$ 的信息来确定 $G$ 的信息；第 4 章中我们还会需要这个特殊结果。

\medskip
\noindent\textbf{命题 21.} 若 $G$ 是有限阿贝尔群，$p$ 是整除 $|G|$ 的素数，则 $G$ 含有阶为 $p$ 的元素。

\begin{proof}
证明按 $|G|$ 作归纳：假设结论对每个阶严格小于 $|G|$ 的群成立，然后证明它对 $G$ 成立（这有时称为完全归纳）。由于 $|G|>1$，存在 $x\in G$ 且 $x\neq1$。若 $|G|=p$，则由拉格朗日定理 $x$ 的阶为 $p$，结论成立。因此可假设 $|G|>p$。

若 $p$ 整除 $|x|$，写 $|x|=pn$。由命题 2.5(3)，$|x^n|=p$，再次得到阶为 $p$ 的元素。因此可假设 $p\nmid |x|$。令 $N=\langle x\rangle$。由于 $G$ 阿贝尔，$N\normal G$。由拉格朗日定理，
\[
  |G/N|=\frac{|G|}{|N|},
\]
并且由于 $N\neq1$，有 $|G/N|<|G|$。又因 $p\nmid |N|$，必有 $p\mid |G/N|$。可对更小的群 $G/N$ 应用归纳假设，得到其中有一个阶为 $p$ 的元素 $\overline y=yN$。由于 $\overline y\neq1$ 但 $\overline y^{\,p}=1$，有 $y\notin N$ 且 $y^p\in N$。于是 $\langle y^p\rangle\neq\langle y\rangle$，即 $|y^p|<|y|$。命题 2.5(2) 蕴含 $p\mid |y|$。于是回到前一段情形，再次得到阶为 $p$ 的元素。归纳完成。
\end{proof}

这种证明方法背后的哲学是：若我们对 $G$ 的某个正规子群 $N$ 有足够信息，并且对 $G/N$ 有足够信息，则可以把这些信息拼接起来，迫使 $G$ 本身具有某种所需性质。归纳法能发挥作用，是因为 $N$ 和 $G/N$ 的阶都比 $G$ 小。一般而言，需要多少数据是一个微妙问题，因为我们已经看到，仅从 $N$ 与 $G/N$ 的同构类型不能决定 $G$ 的完整同构类型。

显然，这一方法的基本障碍是必须产生一个正规子群 $N$，且 $N\neq1$、$N\neq G$。在上面的论证中这很容易，因为 $G$ 是阿贝尔群。没有非平凡真正规子群的群是这种证明方法的基本障碍。

\noindent\textbf{定义.} 若 $|G|>1$，且 $G$ 的正规子群只有 $1$ 与 $G$，则称（有限或无限）群 $G$ 为\emph{单群}。

由拉格朗日定理，若 $|G|$ 为素数，则它的子群（更不用说正规子群）只有 $1$ 与 $G$，故 $G$ 是单群。事实上，每个阿贝尔单群都同构于某个 $Z_p$，其中 $p$ 是素数（参见练习 1）。也存在非阿贝尔单群（有限阶和无限阶都有），其中最小的有限非阿贝尔单群阶为 $60$（下一节会把它作为一族无限单群的成员引入）。

按定义，单群不能再被“分解”为 $N$ 与 $G/N$ 这样的部分，因此它们在群论中扮演的角色类似于素数在整数算术中的角色。这个类比由有限群的一个“唯一分解定理”支持，下面描述之。

\noindent\textbf{定义.} 群 $G$ 中的一列子群
\[
  1=N_0\normal N_1\normal N_2\normal\cdots\normal N_{k-1}\normal N_k=G
\]
称为一个\emph{合成列}，如果 $N_i\normal N_{i+1}$ 且 $N_{i+1}/N_i$ 是单群，$0\leq i\leq k-1$。若上列是合成列，则商群 $N_{i+1}/N_i$ 称为 $G$ 的\emph{合成因子}。

要记住，这里并不假设每个 $N_i\normal G$，只要求 $N_i\normal N_{i+1}$。因此
\[
  1\normal\langle s\rangle\normal\langle s,r^2\rangle\normal D_8,
  \qquad
  1\normal\langle r^2\rangle\normal\langle r\rangle\normal D_8
\]
是 $D_8$ 的两个合成列。在每个合成列中都有 $3$ 个合成因子，每个都同构于单群 $Z_2$。

\medskip
\noindent\textbf{定理 22（Jordan--Hölder）.} 设 $G$ 是有限群且 $G\neq1$。则
\begin{enumerate}[label=(\arabic*)]
  \item $G$ 有合成列；
  \item 合成列中的合成因子是唯一的：若
  \[
    1=N_0\normal N_1\normal\cdots\normal N_r=G,
    \qquad
    1=M_0\normal M_1\normal\cdots\normal M_s=G
  \]
  是 $G$ 的两个合成列，则 $r=s$，并存在集合 $\{1,2,\ldots,r\}$ 的一个置换 $\pi$，使得
  \[
    M_{\pi(i)}/M_{\pi(i)-1}\cong N_i/N_{i-1},
    \qquad 1\leq i\leq r.
  \]
\end{enumerate}

\begin{proof}
这个证明相当直接。由于本书不会显式用它来证明其他定理，我们只在本节末尾的一系列练习中勾勒证明。
\end{proof}

因此每个有限群都有一个“分解”（即合成列），而且尽管合成列本身未必唯一（如 $D_8$ 所示），但合成因子的数目及其同构类型是唯一确定的。进一步地，不同构群也可能有相同的合成因子列表（到同构为止；见练习 2）。这激发了按同构分类所有有限群的两部分纲领：

\paragraph{Hölder 纲领}
\begin{enumerate}[label=(\arabic*)]
  \item 分类所有有限单群。
  \item 找出把单群“拼合”成其他群的所有方式。
\end{enumerate}

这两个问题是群论大量发展的潜在动机之一。类似问题也反复出现在整个数学中。下面对这些问题目前的进展状态作一些评论。

有限单群分类（Hölder 纲领的第 (1) 部分）在 1980 年完成，距离 Hölder 纲领提出约 100 年。超过 100 位数学家的努力，跨越约 5000 到 10000 页期刊论文（分布在大约 300 到 500 篇单独论文中），导致了如下结果的证明：

\medskip
\noindent\textbf{定理.} 存在一个列表，由 18 个无限单群族和 26 个不属于这些族的单群（散在单群）组成，使得每个有限单群都同构于这个列表中的某个群。

有限单群的一个族是 $\{Z_p\mid p\text{ 是素数}\}$。另一个无限族是
\[
  \{\SL_n(\FF)/Z(\SL_n(\FF))\mid n\in\ZZ^+,
  \ n\geq2,
  \ \FF\text{ 是有限域}\}.
\]
这些群除 $\SL_2(\FF_2)$ 与 $\SL_2(\FF_3)$ 外都是单群，其中 $\FF_2$ 和 $\FF_3$ 分别是含 $2$、$3$ 个元素的有限域。这是一个二参数族（$n$ 和 $\FF$ 是独立参数）。我们不证明这些群是单群，而是请读者参阅 M. Aschbacher 的 \textit{Finite Group Theory}（Cambridge University Press, 1986）。有限单群的第三个族——交错群——将在下一节讨论；下一章将证明这些群是单群。

为了感受有限单群分类的复杂性，读者可浏览整个分类的基石之一的证明：

\medskip
\noindent\textbf{定理（Feit--Thompson）.} 若 $G$ 是奇阶单群，则 $G\cong Z_p$，其中 $p$ 是某个素数。

这个证明需要 255 页艰深数学。

Hölder 纲领的第 (2) 部分，有时称为扩张问题，表述相当含糊。“把两个群拼合起来”的一个更精确描述是：给定群 $A$ 和 $B$，描述如何得到所有含有正规子群 $N$ 的群 $G$，使得 $N\cong B$ 且 $G/N\cong A$。例如，若 $A=B=Z_2$，则 $G$ 恰有两种可能，即 $Z_4$ 和 $V_4$；Hölder 纲领试图说明如何从两个 $Z_2$ 构造这两个阶为 $4$ 的群，而不预先知道阶为 $4$ 的群的存在。Hölder 纲领的这一部分极其困难，即使涉及的子群阶很小也是如此。例如，一个群 $G$ 的所有合成因子都阶为 $2$，当且仅当 $|G|=2^n$ 对某个 $n$ 成立（一个方向容易，第 6 章将证明两个方向）。然而已知阶为 $2^n$ 的不同构群的数目作为 $2^n$ 的函数指数增长，因此拼合 $2$ 幂阶群的方式数量没有界。尽管如此，这个精细领域有大量有趣且强有力的技术，用于揭示大类群的结构。我们只会讨论从较小群构造较大群的几种方法；但即使如此有限地进入群扩张领域，我们也将构造许多新的群例子，并证明一些分类定理。

在多项式方程理论中，一个重要的群类是可解群：

\noindent\textbf{定义.} 群 $G$ 称为\emph{可解}的，如果存在子群链
\[
  1=G_0\normal G_1\normal G_2\normal\cdots\normal G_s=G
\]
使得对 $i=0,1,
\ldots,s-1$，商群 $G_{i+1}/G_i$ 是阿贝尔群。

术语来自伽罗瓦理论中这些群与可用根式求解的多项式之间的对应（后者本质上意味着根有代数公式）。练习 8 说明，有限可解群正是那些合成因子全为素数阶的群。

有限可解群的一个显著性质是 Philip Hall 对西洛定理的如下推广：

\medskip
\noindent\textbf{定理.} 有限群 $G$ 可解，当且仅当对 $|G|$ 的每个因子 $n$，只要
\[
  (n,|G|/n)=1,
\]
$G$ 就有一个阶为 $n$ 的子群。

作为另一例，说明如何从正规子群 $N$ 与商群 $G/N$ 的综合信息推导 $G$ 的性质，我们证明：若 $N$ 和 $G/N$ 都可解，则 $G$ 可解。

为此令 $\overline G=G/N$。设
\[
  1=N_0\normal N_1\normal\cdots\normal N_n=N
\]
是 $N$ 的一个子群链，使得 $N_{i+1}/N_i$ 阿贝尔；又设
\[
  \overline1=\overline G_0\normal\overline G_1\normal\cdots\normal\overline G_m=\overline G
\]
是 $\overline G$ 的一个子群链，使得 $\overline G_{i+1}/\overline G_i$ 阿贝尔。由格同构定理，存在 $G$ 的子群 $G_i$，且 $N\leq G_i$，使得 $G_i/N=\overline G_i$，并且 $G_i\normal G_{i+1}$。由第三同构定理，$G_{i+1}/G_i$ 阿贝尔。因此
\[
  1=N_0\normal N_1\normal\cdots\normal N_n=N=G_0\normal G_1\normal\cdots\normal G_m=G
\]
是 $G$ 的一个子群链，其所有相邻商群都是阿贝尔群。这证明 $G$ 可解。

说有限群论只关心 Hölder 纲领是不准确的；但说 Hölder 纲领提出了大量问题并推动许多代数技术的发展则是准确的。例如，在扩张问题中，给定群 $A$ 与 $B$，希望寻找 $G$ 与 $N\normal G$ 使得 $N\cong B$、$G/N\cong A$，我们将看到（在某些条件下）这会引向群 $A$ 在集合 $B$ 上的作用。这类作用正是下一章的核心（并将产生关于单群与非单群的信息），而且这一概念在数学中非常有力，并不局限于群论。

本章最后一节引入另一族群；虽然它与我们对单群的兴趣相一致，但它在全书中也有独立重要性，特别是在我们后面研究行列式和多项式方程可解性时。

\subsection*{练习}
\begin{enumerate}[label=\arabic*.]
  \item 证明若 $G$ 是阿贝尔单群，则 $G\cong Z_p$，其中 $p$ 是素数（不要假设 $G$ 有限）。
  \item 列出 $Q_8$ 的全部 $3$ 条合成列和 $D_8$ 的全部 $7$ 条合成列，并列出每条中的合成因子。
  \item 给出阶为 $16$ 的拟二面体群（参见第 2.5 节练习 11）的一个合成列。推出 $QD_{16}$ 可解。
  \item 使用柯西定理和归纳法证明：有限阿贝尔群对其阶的每个正因子 $n$ 都有阶为 $n$ 的子群。
  \item 证明可解群的子群和商群仍可解。
  \item 对 $|G|$ 作归纳，证明 Jordan--Hölder 定理的第 (1) 部分。
  \item 若 $G$ 是有限群且 $H\normal G$，证明 $G$ 有一条合成列，其中某一项是 $H$。
  \item 设 $G$ 是有限群。证明以下条件等价：
  \begin{enumerate}[label=(\roman*)]
    \item $G$ 可解；
    \item $G$ 有子群链 $1=H_0\normal H_1\normal H_2\normal\cdots\normal H_s=G$，使得 $H_{i+1}/H_i$ 循环；
    \item $G$ 的所有合成因子都是素数阶；
    \item $G$ 有子群链 $1=N_0\normal N_1\normal\cdots\normal N_t=G$，其中每个 $N_i$ 都是 $G$ 的正规子群，且 $N_{i+1}/N_i$ 阿贝尔。
  \end{enumerate}
  [对 (iv)，证明 $G$ 的一个极小非平凡正规子群 $M$ 必是阿贝尔的，然后用归纳。]
  \item 证明 Jordan--Hölder 定理第 (2) 部分的如下特殊情形：若有限群 $G$ 有两条合成列
  \[
    1=N_0\normal N_1\normal\cdots\normal N_r=G,
    \qquad
    1=M_0\normal M_1\normal M_2=G,
  \]
  则 $r=2$，且合成因子列表相同。[使用第二同构定理。]
  \item 对 $\min\{r,s\}$ 作归纳，证明 Jordan--Hölder 定理的第 (2) 部分。[令 $H=N_{r-1}\cap M_{s-1}$，并使用前面练习。]
  \item 证明若 $H$ 是可解群 $G$ 的非平凡正规子群，则 $H$ 中存在一个非平凡子群 $A$，使得 $A\normal G$ 且 $A$ 阿贝尔。
  \item 不使用 Feit--Thompson 定理，证明以下条件等价：
  \begin{enumerate}[label=(\roman*)]
    \item 每个奇阶群都是可解群；
    \item 奇阶单群只有素数阶群。
  \end{enumerate}
\end{enumerate}

\section{换位与交错群}

\subsection*{$S_n$ 的换位与生成}

如第 1.3 节所见（并将在下一章证明），$S_n$ 的每个元素都可以以本质唯一的方式写成不交循环的乘积。相反，$S_n$ 的每个元素也可以以许多不同方式写成（不要求不交的）循环乘积。例如，即使在 $S_3$ 中，元素 $\sigma=(1\ 2\ 3)$ 也可写作
\[
  \sigma=(1\ 2\ 3)=(1\ 3)(1\ 2)=(1\ 2)(1\ 3)(1\ 2)(1\ 3)=(1\ 2)(2\ 3),
\]
事实上还有无穷多种不同写法。不要求循环不交会完全破坏置换作为循环乘积表示的唯一性。不过，若把置换（非唯一地）写成 $2$-循环的乘积，我们可以得到一种“奇偶性检查”。

\noindent\textbf{定义.} 一个 $2$-循环称为\emph{换位}。

直观地说，$\{1,2,
\ldots,n\}$ 的每个置换都可以通过一系列换位，即成对元素的简单交换来实现。下面说明如何做到这一点。首先注意对任意 $m$-循环，
\[
  (a_1\ a_2\ \cdots\ a_m)
  =(a_1\ a_m)(a_1\ a_{m-1})(a_1\ a_{m-2})\cdots(a_1\ a_2).
\]
现在 $S_n$ 的任意置换都可写成循环乘积（例如其循环分解）。把每个循环依次按上式写成换位乘积，就得到
\[
  \text{$S_n$ 的每个元素都可写成换位的乘积，}
\]
等价地，
\[
  S_n=\langle T\rangle,
  \qquad T=\{(i\ j)
  \mid 1\leq i<j\leq n\}.
\]
例如第 1.3 节中的置换可写成
\begin{align*}
  \sigma&=(1\ 12\ 8\ 10\ 4)(2\ 13)(5\ 11\ 7)(6\ 9)\\
  &=(1\ 4)(1\ 10)(1\ 8)(1\ 12)(2\ 13)(5\ 7)(5\ 11)(6\ 9).
\end{align*}

\subsection*{交错群}

再次强调，对任意 $\sigma\in S_n$，可有许多方式把 $\sigma$ 写成换位乘积。现在证明，对固定的 $\sigma$，等于 $\sigma$ 的任意换位乘积的项数奇偶性相同。

令 $x_1,
\ldots,x_n$ 是独立变量，令 $\Delta$ 为多项式
\[
  \Delta=\prod_{1\leq i<j\leq n}(x_i-x_j),
\]
即所有 $x_i-x_j$（$i<j$）的乘积。例如，当 $n=4$ 时，
\[
  \Delta=(x_1-x_2)(x_1-x_3)(x_1-x_4)(x_2-x_3)(x_2-x_4)(x_3-x_4).
\]
对每个 $\sigma\in S_n$，令 $\sigma$ 通过按其置换指标的方式置换变量而作用在 $\Delta$ 上：
\[
  \sigma(\Delta)=\prod_{1\leq i<j\leq n}(x_{\sigma(i)}-x_{\sigma(j)}).
\]
例如，若 $n=4$ 且 $\sigma=(1\ 2\ 3\ 4)$，则
\[
  \sigma(\Delta)=(x_2-x_3)(x_2-x_4)(x_2-x_1)(x_3-x_4)(x_3-x_1)(x_4-x_1).
\]
一般地，$\Delta$ 中对每个 $i<j$ 有一个因子 $x_i-x_j$；由于 $\sigma$ 是指标集合的双射，$\sigma(\Delta)$ 对每个 $i<j$ 必含 $x_i-x_j$ 或 $x_j-x_i$ 中恰一个（且当然没有 $x_i-x_i$ 因子）。若 $\sigma(\Delta)$ 中出现因子 $x_j-x_i$ 且 $j>i$，将其写作 $-(x_i-x_j)$。把所有符号变化收集起来，可见 $\Delta$ 与 $\sigma(\Delta)$ 的因子相同，只差若干个 $-1$ 的乘积；即
\[
  \sigma(\Delta)=\pm\Delta,
  \qquad \sigma\in S_n.
\]

对每个 $\sigma\in S_n$，令
\[
  \varepsilon(\sigma)=
  \begin{cases}
    +1,&\text{若 }\sigma(\Delta)=\Delta,\\
    -1,&\text{若 }\sigma(\Delta)=-\Delta.
  \end{cases}
\]
在上面 $n=4$、$\sigma=(1\ 2\ 3\ 4)$ 的例子中，$\sigma(\Delta)$ 中恰有 $3$ 个形如 $x_j-x_i$（$j>i$）的因子，因此每个贡献一个 $-1$。故
\[
  (1\ 2\ 3\ 4)(\Delta)=(-1)^3\Delta=-\Delta,
\]
所以 $\varepsilon((1\ 2\ 3\ 4))=-1$。

\noindent\textbf{定义.}
\begin{enumerate}[label=(\arabic*)]
  \item $\varepsilon(\sigma)$ 称为 $\sigma$ 的\emph{符号}。
  \item 若 $\varepsilon(\sigma)=1$，称 $\sigma$ 为\emph{偶置换}；若 $\varepsilon(\sigma)=-1$，称 $\sigma$ 为\emph{奇置换}。
\end{enumerate}

下一个结果表明，置换符号定义了一个同态。

\medskip
\noindent\textbf{命题 23.} 映射 $\varepsilon:S_n\to\{\pm1\}$ 是同态（其中 $\{\pm1\}$ 是阶为 $2$ 的循环群的乘法形式）。

\begin{proof}
按定义，
\[
  (\tau\sigma)(\Delta)=\prod_{1\leq i<j\leq n}(x_{\tau\sigma(i)}-x_{\tau\sigma(j)}).
\]
假设 $\sigma(\Delta)$ 中恰有 $k$ 个因子形如 $x_j-x_i$ 且 $j>i$，即 $\varepsilon(\sigma)=(-1)^k$。计算 $(\tau\sigma)(\Delta)$ 时，先把 $\sigma$ 作用于指标，可见其中恰有 $k$ 个因子形如 $x_{\tau(j)}-x_{\tau(i)}$，其中 $j>i$。交换这 $k$ 个因子中两项的顺序会产生符号变化 $(-1)^k=\varepsilon(\sigma)$；此时所有因子都形如 $x_{\tau(p)}-x_{\tau(q)}$，其中 $p<q$。因此
\[
  (\tau\sigma)(\Delta)=\varepsilon(\sigma)
  \prod_{1\leq p<q\leq n}(x_{\tau(p)}-x_{\tau(q)}).
\]
由 $\varepsilon$ 的定义，
\[
  \prod_{1\leq p<q\leq n}(x_{\tau(p)}-x_{\tau(q)})=\varepsilon(\tau)\Delta.
\]
所以 $(\tau\sigma)(\Delta)=\varepsilon(\sigma)\varepsilon(\tau)\Delta$。于是
\[
  \varepsilon(\tau\sigma)=\varepsilon(\sigma)\varepsilon(\tau)=\varepsilon(\tau)\varepsilon(\sigma),
\]
如所需。
\end{proof}

为了看这个证明如何运作，令 $n=4$，$\sigma=(1\ 2\ 3\ 4)$，$\tau=(4\ 2\ 3)$，于是 $\tau\sigma=(1\ 3\ 2\ 4)$。由定义（使用此时 $\Delta$ 的显式表达式）有
\begin{align*}
  (\tau\sigma)(\Delta)\\
  &=(1\ 3\ 2\ 4)(\Delta)\\
  &=(x_3-x_4)(x_3-x_2)(x_3-x_1)(x_4-x_2)(x_4-x_1)(x_2-x_1)\\
  &=(-1)^5\Delta,
\end{align*}
其中除了第一个因子外，其余所有因子都要翻转才能恢复 $\Delta$。这说明 $\varepsilon(\tau\sigma)=-1$。另一方面，由于我们已经计算过 $\sigma(\Delta)$，
\begin{align*}
  (\tau\sigma)(\Delta)
  =\tau(\sigma(\Delta))\\
  =\begin{aligned}
  &(x_{\tau(2)}-x_{\tau(3)})(x_{\tau(2)}-x_{\tau(4)})(x_{\tau(2)}-x_{\tau(1)})\\
  &\quad\times(x_{\tau(3)}-x_{\tau(4)})(x_{\tau(3)}-x_{\tau(1)})(x_{\tau(4)}-x_{\tau(1)})
  \end{aligned}\\
  =(-1)^3\!\prod_{1\leq p<q\leq 4}\!(x_{\tau(p)}-x_{\tau(q)})=(-1)^3\tau(\Delta),
\end{align*}
这里第三、第五、第六个因子需要交换两项，才能把所有因子写成 $x_{\tau(p)}-x_{\tau(q)}$ 且 $p<q$ 的形式。我们已经算出 $\varepsilon(\sigma)=(-1)^3=-1$；用同样方法容易看出 $\varepsilon(\tau)=(-1)^2=1$，所以
\[
  \varepsilon(\tau\sigma)=-1=\varepsilon(\tau)\varepsilon(\sigma).
\]

接下来计算任意换位 $(i\ j)$ 的符号。先做 $i=1,j=2$ 的情形，再化归到一般情形。显然，对任意 $n$，$(1\ 2)$ 作用于 $\Delta$ 时恰好翻转一个因子，即 $x_1-x_2$；因此 $\varepsilon((1\ 2))=-1$。对任意换位 $(i\ j)$，令 $\lambda$ 是交换 $1$ 与 $i$、交换 $2$ 与 $j$ 并固定其他数字的置换（若 $i=1$ 或 $j=2$，则相应地固定该数）。容易看出
\[
  (i\ j)=\lambda(1\ 2)\lambda.
\]
由于 $\varepsilon$ 是同态，得
\[
  \varepsilon((i\ j))=\varepsilon(\lambda(1\ 2)\lambda)
  =\varepsilon(\lambda)\varepsilon((1\ 2))\varepsilon(\lambda)
  =-\varepsilon(\lambda)^2=-1.
\]

\medskip
\noindent\textbf{命题 24.} 所有换位都是奇置换，并且 $\varepsilon$ 是满同态。

\noindent\textbf{定义.} \emph{$n$ 次交错群}，记为 $A_n$，是同态 $\varepsilon$ 的核（即所有偶置换的集合）。

由第一同构定理，$S_n/A_n\cong\varepsilon(S_n)=\{\pm1\}$，所以
\[
  |A_n|=\frac12|S_n|=\frac12 n!.
\]
此外，$S_n-A_n$ 是 $A_n$ 的非单位陪集，即所有奇置换的集合。置换符号遵循通常的 $\ZZ/2\ZZ$ 法则：
\[
\begin{array}{c}
(\text{偶})(\text{偶})=(\text{奇})(\text{奇})=\text{偶},\\
(\text{偶})(\text{奇})=(\text{奇})(\text{偶})=\text{奇}.
\end{array}
\]

又因为 $\varepsilon$ 是同态，且每个 $\sigma\in S_n$ 都是换位乘积，设 $\sigma=\tau_1\tau_2\cdots\tau_k$，则
\[
  \varepsilon(\sigma)=\varepsilon(\tau_1)\cdots\varepsilon(\tau_k)=(-1)^k.
\]
因此无论如何把 $\sigma$ 写成换位乘积，项数 $k$ 的模 $2$ 类，即换位个数的奇偶性，都是相同的：
\[
  \varepsilon(\sigma)=
  \begin{cases}
    +1,&\text{若 $\sigma$ 是偶数个换位的乘积},\\
    -1,&\text{若 $\sigma$ 是奇数个换位的乘积}.
  \end{cases}
\]

最后给出一种从循环分解快速计算 $\varepsilon(\sigma)$ 的方法。回忆一个 $m$-循环可写成 $m-1$ 个换位的乘积。因此
\[
  \text{$m$-循环是奇置换，当且仅当 $m$ 是偶数。}
\]
若置换 $\sigma$ 的循环分解为 $\alpha_1\alpha_2\cdots\alpha_k$，则
\[
  \varepsilon(\sigma)=\varepsilon(\alpha_1)\cdots\varepsilon(\alpha_k),
\]
且 $\varepsilon(\alpha_i)=-1$ 当且仅当 $\alpha_i$ 的长度为偶数。于是 $\varepsilon(\sigma)=-1$ 当且仅当这些因子中出现奇数个 $-1$。概括为：

\medskip
\noindent\textbf{命题 25.} 置换 $\sigma$ 是奇置换，当且仅当其循环分解中偶长度循环的个数为奇数。

例如，
\[
  \sigma=(1\ 2\ 3\ 4\ 5\ 6)(7\ 8\ 9)(10\ 11)(12\ 13\ 14\ 15)(16\ 17\ 18)
\]
有 $3$ 个偶长度循环，所以 $\varepsilon(\sigma)=-1$。另一方面，
\[
  \tau=(1\ 12\ 8\ 10\ 4)(2\ 13)(5\ 11\ 7)(6\ 9)
\]
恰有 $2$ 个偶长度循环，故 $\varepsilon(\tau)=1$。

注意不要把置换 $\sigma$ 的“奇”“偶”与其阶的奇偶混淆。事实上，若 $\sigma$ 有奇数阶，则 $\sigma$ 的循环分解中的所有循环长度都是奇数，所以 $\sigma$ 有偶数个（此处为 $0$ 个）偶长度循环，因此是偶置换。若 $|\sigma|$ 是偶数，则 $\sigma$ 可能是偶置换也可能是奇置换；例如 $(1\ 2)$ 是奇置换，$(1\ 2)(3\ 4)$ 是偶置换，但二者阶都为 $2$。

如上一节所述，交错群 $A_n$ 在多项式可解性的研究中十分重要。下一章我们将证明：
\[
  A_n\text{ 对所有 }n\geq5\text{ 都是非阿贝尔单群。}
\]
对较小的 $n$，$A_n$ 已经是我们熟悉的群：$A_1$ 和 $A_2$ 都是平凡群，$|A_3|=3$（所以 $A_3=\langle(1\ 2\ 3)\rangle\cong Z_3$）。群 $A_4$ 的阶为 $12$。练习 7 说明 $A_4$ 同构于正四面体的对称群。$A_4$ 的子群格见图 8（练习 8 断言这是其完整子群格）。这个格的一个好处是（不像“几乎所有群”那样）它是平面图（除了在顶点处，线没有交叉；可参见 $D_{16}$ 的非平面格）。

\begin{center}
\begin{tikzpicture}[scale=.82, every node/.style={inner sep=1pt}]
  \node (A4) at (0,4.2) {$A_4$};
  \node (V) at (-2.8,3.0) {$\langle(1\ 2)(3\ 4),(1\ 3)(2\ 4)\rangle$};
  \node (C1) at (-.6,3.0) {$\langle(1\ 2\ 3)\rangle$};
  \node (C2) at (1.4,3.0) {$\langle(1\ 2\ 4)\rangle$};
  \node (C3) at (3.4,3.0) {$\langle(1\ 3\ 4)\rangle$};
  \node (C4) at (5.4,3.0) {$\langle(2\ 3\ 4)\rangle$};
  \node (D1) at (-4.2,1.5) {$\langle(1\ 2)(3\ 4)\rangle$};
  \node (D2) at (-2.4,1.5) {$\langle(1\ 3)(2\ 4)\rangle$};
  \node (D3) at (-.6,1.5) {$\langle(1\ 4)(2\ 3)\rangle$};
  \node (one) at (.6,0) {$1$};
  \foreach \X in {V,C1,C2,C3,C4}{\draw (A4)--(\X);}
  \foreach \X in {D1,D2,D3}{\draw (V)--(\X); \draw (\X)--(one);}
  \foreach \X in {C1,C2,C3,C4}{\draw (\X)--(one);}
  \node at (5.8,.4) {图 8};
\end{tikzpicture}
\end{center}

\subsection*{练习}
\begin{enumerate}[label=\arabic*.]
  \item 在第 1.3 节练习 1 和 2 中，你被要求求一些置换的循环分解。把每个置换写成换位的乘积，并判断它们是偶置换还是奇置换。
  \item 证明对每个置换 $\sigma$，$\sigma^2$ 都是偶置换。
  \item 证明 $S_n$ 由集合 $\{(i\ i+1)\mid 1\leq i\leq n-1\}$ 生成。[考虑共轭，例如 $(2\ 3)(1\ 2)(2\ 3)^{-1}$。]
  \item 证明对所有 $n\geq2$，
  \[
    S_n=\langle(1\ 2),(1\ 2\ 3\ \cdots\ n)\rangle.
  \]
  \item 证明若 $p$ 是素数，则 $S_p=\langle\sigma,\tau\rangle$，其中 $\sigma$ 是任意换位，$\tau$ 是任意 $p$-循环。
  \item 说明 $\langle(1\ 3),(1\ 2\ 3\ 4)\rangle$ 是 $S_4$ 的真子群。该子群的同构类型是什么？
  \item 证明正四面体的刚体运动群同构于 $A_4$。[回忆第 1.7 节练习 20。]
  \item 证明正文中给出的 $A_4$ 子群格是正确的。[由上一题以及拉格朗日定理之后的讨论，$A_4$ 没有阶为 $6$ 的子群。]
  \item 证明 $A_4$ 中唯一阶为 $4$ 的子群正规，且同构于 $V_4$。
  \item 给出 $A_4$ 的一个合成列，并推出 $A_4$ 可解。
  \item 证明 $S_4$ 没有同构于 $Q_8$ 的子群。
  \item 证明对每个 $n\geq3$，$A_n$ 含有一个同构于 $S_{n-2}$ 的子群。
  \item 证明 $A_n$ 中每个阶为 $2$ 的元素都是 $S_n$ 中某个阶为 $4$ 的元素的平方。[提示：$A_n$ 中阶为 $2$ 的元素是 $2k$ 个两两交换的换位的乘积。]
  \item 证明 $A_4$ 中由任一阶为 $2$ 的元素和任一阶为 $3$ 的元素生成的子群都是整个 $A_4$。
  \item 证明若 $x$ 和 $y$ 是 $S_4$ 中不同的 $3$-循环且 $x\neq y^{-1}$，则由 $x$ 和 $y$ 生成的 $S_4$ 子群是 $A_4$。
  \item 令 $x$ 和 $y$ 是 $S_5$ 中不同的 $3$-循环，且 $x\neq y^{-1}$。
  \begin{enumerate}[label=(\alph*)]
    \item 证明若 $x$ 与 $y$ 固定 $\{1,
    \ldots,5\}$ 的一个共同元素，则 $\langle x,y\rangle\cong A_4$。
    \item 证明若 $x$ 与 $y$ 不固定共同元素，则 $\langle x,y\rangle=A_5$。
  \end{enumerate}
  \item 若 $x$ 与 $y$ 是 $S_n$ 中的 $3$-循环，证明 $\langle x,y\rangle$ 同构于 $Z_3$、$A_4$、$A_5$ 或 $Z_3\times Z_3$。
\end{enumerate}

% \chapter{群作用}

本章我们考察群作用在集合上所产生的一些后果。数学中一个重要而反复出现的思想是：当一个对象作用在另一个对象上时，二者的大量信息都可以由此获得。随着被作用的集合带有更多结构（例如群作用在群上，或群作用在向量空间上，后者将在第 18 章讨论），关于群本身结构的更多信息也会显现出来。这里对群作用的研究将以 Sylow 定理及由它导出的例子和分类为高潮。

作用的概念在我们研究模、向量空间、矩阵的标准形以及 Galois 理论时还会反复出现，它是贯穿本书的基本统一主题之一。

\section{群作用与置换表示}

本节给出群作用的基本理论，然后把这一理论应用于 $S_n$ 的子群在 $\{1,2,\ldots,n\}$ 上的作用，从而证明 $S_n$ 的每个元素都有唯一的循环分解。在第 2 节和第 3 节中，我们会把一般理论应用于另外两个具体的群作用，以推出一些重要结果。

设 $G$ 是一个作用在非空集合 $A$ 上的群。回忆第 1.7 节：对每个 $g\in G$，映射
\[
  \sigma_g:A\to A,
  \qquad \sigma_g(a)=g\cdot a
\]
是 $A$ 的一个置换。我们也看到，与 $G$ 在 $A$ 上的一个作用相伴随，有一个同态
\[
  \varphi:G\to S_A,
  \qquad \varphi(g)=\sigma_g,
\]
称为与给定作用相伴随的\emph{置换表示}。回忆第 1.7 节和第 2.2 节中引入的与群作用有关的一些术语。

\noindent\textbf{定义.}
\begin{enumerate}[label=(\arabic*)]
  \item 作用的\emph{核}是 $G$ 中在 $A$ 的每个元素上都平凡作用的元素集合：
  \[
    \{g\in G\mid g\cdot a=a\text{ 对所有 }a\in A\}.
  \]
  \item 对每个 $a\in A$，$a$ 在 $G$ 中的\emph{稳定子}是 $G$ 中固定 $a$ 的元素集合：
  \[
    \{g\in G\mid g\cdot a=a\},
  \]
  记为 $G_a$。
  \item 若一个作用的核是单位元，则称该作用是\emph{忠实的}。
\end{enumerate}

注意，作用的核恰好就是相伴随置换表示的核；特别地，核是 $G$ 的正规子群。两个群元素在 $A$ 上诱导相同的置换，当且仅当它们属于核的同一个陪集（也即，当且仅当它们属于置换表示 $\varphi$ 的同一条纤维）。特别地，$G$ 在 $A$ 上的一个作用也可看作商群 $G/\ker\varphi$ 在 $A$ 上的一个忠实作用。由第 2.2 节知，$A$ 的一个元素 $a$ 在 $G$ 中的稳定子是 $G$ 的子群。如果 $a$ 是 $A$ 中固定的一个元素，则作用的核包含于稳定子 $G_a$，因为作用的核是稳定每一点的元素集合，即
\[
  \bigcap_{a\in A}G_a.
\]

\paragraph{例}
\begin{enumerate}[label=(\arabic*)]
  \item 设 $n$ 是正整数。群 $G=S_n$ 按
  \[
    \sigma\cdot i=\sigma(i),\qquad i\in\{1,\ldots,n\}
  \]
  作用在集合 $A=\{1,2,\ldots,n\}$ 上。与此作用相伴随的置换表示是恒等映射 $\varphi:S_n\to S_n$。这个作用是忠实的，并且对每个 $i\in\{1,\ldots,n\}$，稳定子 $G_i$（所有固定 $i$ 的置换组成的子群）同构于 $S_{n-1}$（参见第 3.2 节练习 15）。

  \item 令 $G=D_8$ 作用在一个正方形的四个顶点组成的集合 $A$ 上。按第 1.2 节图 2 那样，把这些顶点顺时针标为 $1,2,3,4$。令 $r$ 为正方形顺时针旋转 $\pi/2$ 弧度的旋转，令 $s$ 为关于经过顶点 $1$ 和 $3$ 的直线的反射。则 $r$ 与 $s$ 在顶点上给出的置换为
  \[
    \sigma_r=(1\ 2\ 3\ 4),\qquad \sigma_s=(2\ 4).
  \]
  注意，由于置换表示是同态，与 $sr$ 相对应的四个顶点上的置换为
  \[
    \sigma_{sr}=\sigma_s\sigma_r=(1\ 4)(2\ 3).
  \]
  $D_8$ 在正方形四个顶点上的作用是忠实的，因为只有恒等对称才固定所有四个顶点。任一顶点 $a$ 的稳定子是 $D_8$ 中由经过 $a$ 与正方形中心的反射生成的阶为 $2$ 的子群（例如顶点 $1$ 的稳定子为 $\gen{s}$）。

  \item 仍按前例标记正方形的四个顶点。现在令 $A$ 为由相对顶点的无序对组成的集合：
  \[
    A=\bigl\{\{1,3\},\{2,4\}\bigr\}.
  \]
  则 $D_8$ 也作用在这个集合 $A$ 上，因为正方形的每个对称都把一对相对顶点送到一对相对顶点。旋转 $r$ 交换两对 $\{1,3\}$ 与 $\{2,4\}$；反射 $s$ 固定这两个相对顶点的无序对。因此，若把 $\{1,3\}$ 与 $\{2,4\}$ 分别标为 $1$ 与 $2$，则 $r$ 与 $s$ 在 $A$ 上给出的置换为
  \[
    \sigma_r=(1\ 2),\qquad \sigma_s=\text{恒等置换}.
  \]
  这个 $D_8$ 的作用不是忠实的：它的核是 $\gen{s,r^2}$。而且，对每个 $a\in A$，$a$ 在 $D_8$ 中的稳定子都等于该作用的核。

  \item 仍按例 (2) 标记正方形四个顶点，并令 $A$ 为如下顶点无序对的集合：
  \[
    \bigl\{\{1,2\},\{3,4\}\bigr\}.
  \]
  群 $D_8$ 不作用在这个集合 $A$ 上，因为 $\{1,2\}\in A$，但
  \[
    r\cdot\{1,2\}=\{2,3\}\notin A.
  \]
\end{enumerate}

作用与到对称群中的同态之间的关系可以反过来。也就是说，给定任意非空集合 $A$ 以及群 $G$ 到 $S_A$ 的任意同态 $\varphi$，可通过
\[
  g\cdot a=\varphi(g)(a)
\]
定义 $G$ 在 $A$ 上的一个作用。这个作用的核等于 $\ker\varphi$。与这个作用相伴随的置换表示正是给定的同态 $\varphi$。于是得到下述结果。

\medskip
\noindent\textbf{命题 1.} 对任意群 $G$ 与任意非空集合 $A$，$G$ 在 $A$ 上的作用与 $G$ 到 $S_A$ 中的同态之间存在一一对应。

鉴于命题 1，置换表示的定义可以重新表述如下。

\noindent\textbf{定义.} 若 $G$ 是群，则 $G$ 的一个\emph{置换表示}是 $G$ 到某个非空集合 $A$ 的对称群 $S_A$ 中的任意同态。我们称 $G$ 在 $A$ 上的一个给定作用\emph{给出}或\emph{诱导} $G$ 的相伴随置换表示。

可以把置换表示看作线性变换的矩阵表示的类比。当 $A$ 是一个有 $n$ 个元素的有限集合时，$S_A\cong S_n$（参见第 1.6 节），所以通过固定 $A$ 中元素的一个标号，我们便可把这些置换看作 $S_n$ 的元素（这正是前面例 2 和例 3 中所做的），就像在向量空间中固定一个基便可把线性变换看作矩阵一样。

下面证明一个关于群作用的组合结果；在以后几节把它应用于具体作用时，它会产生重要后果。

\medskip
\noindent\textbf{命题 2.} 设 $G$ 是作用在非空集合 $A$ 上的群。在 $A$ 上定义关系
\[
  a\sim b \quad\Longleftrightarrow\quad a=g\cdot b\text{ 对某个 }g\in G.
\]
则 $\sim$ 是一个等价关系。对每个 $a\in A$，包含 $a$ 的等价类中的元素个数为 $|G:G_a|$，即 $a$ 的稳定子的指数。

\begin{proof}
先证明 $\sim$ 是等价关系。由作用公理 2，对所有 $a\in A$，有 $a=1\cdot a$，即 $a\sim a$，所以该关系自反。若 $a\sim b$，则 $a=g\cdot b$ 对某个 $g\in G$ 成立，于是
\[
  g^{-1}\cdot a=g^{-1}\cdot(g\cdot b)=(g^{-1}g)\cdot b=1\cdot b=b,
\]
即 $b\sim a$，所以该关系对称。最后，若 $a\sim b$ 且 $b\sim c$，则对某些 $g,h\in G$ 有 $a=g\cdot b$ 且 $b=h\cdot c$，于是
\[
  a=g\cdot b=g\cdot(h\cdot c)=(gh)\cdot c,
\]
故 $a\sim c$，所以该关系传递。

为证明命题的最后部分，我们在 $G_a$ 在 $G$ 中的左陪集与 $a$ 所在等价类的元素之间建立一个双射。令 $C_a$ 为 $a$ 的等价类，则
\[
  C_a=\{g\cdot a\mid g\in G\}.
\]
若 $b=g\cdot a\in C_a$，则 $gG_a$ 是 $G_a$ 在 $G$ 中的一个左陪集。映射
\[
  b=g\cdot a\longmapsto gG_a
\]
把 $C_a$ 映到 $G_a$ 的左陪集集合中。它是满射，因为对任意 $g\in G$，元素 $g\cdot a$ 属于 $C_a$。又因为 $g\cdot a=h\cdot a$ 当且仅当 $h^{-1}g\in G_a$，当且仅当 $gG_a=hG_a$，所以该映射也是单射。因此它是双射。证明完毕。
\end{proof}

由命题 2，群 $G$ 作用在集合 $A$ 上会把 $A$ 分割为该作用下的不交等价类。这些等价类有如下名称。

\noindent\textbf{定义.} 设 $G$ 是作用在非空集合 $A$ 上的群。
\begin{enumerate}[label=(\arabic*)]
  \item 等价类 $\{g\cdot a\mid g\in G\}$ 称为包含 $a$ 的 $G$ 的\emph{轨道}。
  \item 若 $G$ 在 $A$ 上的作用只有一个轨道，则称该作用是\emph{传递的}；也就是说，对任意两个元素 $a,b\in A$，存在某个 $g\in G$ 使得 $a=g\cdot b$。
\end{enumerate}

\paragraph{例}
设 $G$ 是作用在集合 $A$ 上的群。
\begin{enumerate}[label=(\arabic*)]
  \item 若 $G$ 平凡作用在 $A$ 上，则对所有 $a\in A$，有 $G_a=G$，而轨道就是 $A$ 的各个单元素集合。这个作用传递当且仅当 $|A|=1$。
  \item 对称群 $G=S_n$ 按通常的置换作用传递地作用在 $A=\{1,2,\ldots,n\}$ 上。注意，$G$ 中任一点 $i$ 的稳定子在 $S_n$ 中的指数为 $n=|A|$。
  \item 当群 $G$ 作用在集合 $A$ 上时，$G$ 的任意子群也作用在 $A$ 上。即使 $G$ 在 $A$ 上传递，$G$ 的子群也未必在 $A$ 上传递。例如，若 $G=\gen{(1\ 2),(3\ 4)}\leq S_4$，则 $G$ 在 $\{1,2,3,4\}$ 上的轨道为 $\{1,2\}$ 与 $\{3,4\}$，且没有 $G$ 中元素把 $2$ 送到 $3$。下面关于循环分解的讨论表明，当 $\gen{\sigma}$ 是 $S_n$ 的任意循环子群时，$\gen{\sigma}$ 的轨道由 $\sigma$ 的循环分解中各个循环里出现的数字集合组成（例如，$\gen{(1\ 2)(3\ 4\ 5)}$ 的轨道是 $\{1,2\}$ 与 $\{3,4,5\}$）。
  \item 群 $D_8$ 在正方形的四个顶点上传递作用，并且任一顶点的稳定子是由经过该点的对称轴反射生成的阶为 $2$（指数为 $4$）的子群。
  \item 群 $D_8$ 也在两对相对顶点组成的集合上传递作用。在此作用下，任一点的稳定子为 $\gen{s,r^2}$（指数为 $2$）。
\end{enumerate}

\paragraph{循环分解}
现在证明第 1.3 节中描述的事实：对称群 $S_n$ 的每个元素都有唯一的循环分解。令 $A=\{1,2,\ldots,n\}$，令 $\sigma\in S_n$，并令 $G=\gen{\sigma}$。则 $\gen{\sigma}$ 作用在 $A$ 上，所以由命题 2，它把 $\{1,2,\ldots,n\}$ 分割为唯一的一组（不交）轨道。令 $O$ 是其中一个轨道，并取 $x\in O$。由命题 2 的证明（应用于 $A=O$）可见，$G_x$ 在 $G$ 中的左陪集与 $O$ 的元素之间有一个双射，显式地由
\[
  \sigma^iG_x\longmapsto \sigma^i(x)
\]
给出。由于 $G$ 是循环群，$G_x\normal G$，且 $G/G_x$ 是阶为 $d$ 的循环群，其中 $d$ 是满足 $\sigma^d\in G_x$ 的最小正整数（参见第 3.1 节命题 7 后的例 2）。并且 $d=|G:G_x|=|O|$。因此 $G_x$ 在 $G$ 中的不同陪集为
\[
  1G_x,\ \sigma G_x,\ \sigma^2G_x,\ldots,\sigma^{d-1}G_x.
\]
这说明 $O$ 的不同元素为
\[
  x,\ \sigma(x),\ \sigma^2(x),\ldots,\sigma^{d-1}(x).
\]
按这种方式排列 $O$ 的元素，说明 $\sigma$ 循环地置换 $O$ 的元素；也就是说，在大小为 $d$ 的轨道上，$\sigma$ 作用为一个 $d$-循环。这证明了每个 $\sigma\in S_n$ 的循环分解存在。

$\gen{\sigma}$ 的轨道由 $\sigma$ 唯一确定。唯一的自由度是列出轨道的顺序。在每个轨道 $O$ 内，我们可以从任意元素开始作为代表元。若选择 $\sigma^i(x)$ 而不是 $x$ 作为初始代表元，就只是把 $O$ 中元素按如下顺序写出：
\[
  \sigma^i(x),\sigma^{i+1}(x),\ldots,\sigma^{d-1}(x),x,\sigma(x),\ldots,\sigma^{i-1}(x),
\]
这只是原列表的一个循环移位（向前移 $i-1$ 项）。于是，上述循环分解唯一到只差循环的重新排序以及每个循环内部数字的循环置换。

对称群的子群称为\emph{置换群}。对 $S_n$ 的任一子群 $G$，$G$ 的轨道指它在 $\{1,2,\ldots,n\}$ 上的轨道。$S_n$ 中元素 $\sigma$ 的轨道指群 $\gen{\sigma}$ 的轨道（也就是其循环分解中各循环所包含的整数集合）。

下面的练习进一步说明如何从置换表示中获得群论信息。

\subsection*{练习}
设 $G$ 是群，$A$ 是非空集合。
\begin{enumerate}[label=\arabic*.]
  \item 设 $G$ 作用在集合 $A$ 上。证明：若 $a,b\in A$ 且 $b=g\cdot a$ 对某个 $g\in G$ 成立，则 $G_b=gG_ag^{-1}$（$G_a$ 是 $a$ 的稳定子）。推出：若 $G$ 在 $A$ 上传递作用，则该作用的核为 $\bigcap_{g\in G}gG_ag^{-1}$。
  \item 设 $G$ 是集合 $A$ 上的一个置换群（即 $G\leq S_A$），令 $\sigma\in G$ 且 $a\in A$。证明 $\sigma G_a\sigma^{-1}=G_{\sigma(a)}$。推出：若 $G$ 在 $A$ 上传递作用，则
  \[
    \bigcap_{\sigma\in G}\sigma G_a\sigma^{-1}=1.
  \]
  \item 假设 $G$ 是 $S_A$ 的阿贝尔传递子群。证明对所有 $\sigma\in G-\{1\}$ 与所有 $a\in A$，都有 $\sigma(a)\neq a$。推出 $|G|=|A|$。［用前一题。］
  \item 令 $S_3$ 作用在有序对集合 $\Omega=\{(i,j)\mid 1\leq i,j\leq 3\}$ 上，作用为 $\sigma((i,j))=(\sigma(i),\sigma(j))$。求 $S_3$ 在 $\Omega$ 上的轨道。对每个 $\sigma\in S_3$，求 $\sigma$ 在这个作用下的循环分解（即先固定这九个有序对的一个标号，再把 $\sigma$ 看作 $S_9$ 的元素来求其循环分解）。对 $S_3$ 作用在这九个点上的每个轨道 $O$，取某个 $a\in O$，求 $a$ 在 $S_3$ 中的稳定子。
  \item 对下列两个集合，重复前一题：
  \begin{enumerate}[label=(\alph*)]
    \item $27$ 个三元组组成的集合 $\{(i,j,k)\mid 1\leq i,j,k\leq 3\}$；
    \item $\mathcal P(\{1,2,3\})-\{\varnothing\}$，即 $\{1,2,3\}$ 的所有 $7$ 个非空子集组成的集合。
  \end{enumerate}
  \item 如第 2.2 节练习 12，令 $R$ 为独立变量 $x_1,x_2,x_3,x_4$ 的所有整系数多项式集合，并令 $S_4$ 通过置换下标作用在 $R$ 上：
  \[
    \sigma\cdot p(x_1,x_2,x_3,x_4)=p(x_{\sigma(1)},x_{\sigma(2)},x_{\sigma(3)},x_{\sigma(4)})
  \]
  对所有 $\sigma\in S_4$ 成立。
  \begin{enumerate}[label=(\alph*)]
    \item 求 $S_4$ 在 $R$ 上的轨道中包含 $x_1+x_2$ 的多项式（回忆第 2.2 节练习 12：该多项式的稳定子阶为 $4$）。
    \item 求 $S_4$ 在 $R$ 上的轨道中包含 $x_1x_2+x_3x_4$ 的多项式（回忆第 2.2 节练习 12：该多项式的稳定子阶为 $8$）。
    \item 求 $S_4$ 在 $R$ 上的轨道中包含 $(x_1+x_2)(x_3+x_4)$ 的多项式。
  \end{enumerate}
  \item 设 $G$ 是有限集合 $A$ 上的传递置换群。一个\emph{块}是 $A$ 的非空子集 $B$，满足对所有 $\sigma\in G$，或者 $\sigma(B)=B$，或者 $\sigma(B)\cap B=\varnothing$（这里 $\sigma(B)=\{\sigma(b)\mid b\in B\}$）。
  \begin{enumerate}[label=(\alph*)]
    \item 证明：若 $B$ 是包含 $A$ 中元素 $a$ 的块，则集合 $G_B=\{\sigma\in G\mid \sigma(B)=B\}$ 是包含 $G_a$ 的 $G$ 的子群。
    \item 证明：若 $B$ 是一个块，并且 $\sigma_1(B),\sigma_2(B),\ldots,\sigma_n(B)$ 是 $B$ 在 $G$ 的元素下的全部不同像，则这些集合构成 $A$ 的一个分割。
    \item 若传递群 $G$ 在集合 $A$ 上的块只有平凡块（大小为 $1$ 的集合与 $A$ 本身），则称 $G$ 在 $A$ 上是\emph{本原的}。证明 $S_4$ 在 $A=\{1,2,3,4\}$ 上是本原的。证明 $D_8$ 作为正方形四个顶点上的置换群不是本原的。
    \item 证明：传递群 $G$ 在 $A$ 上是本原的，当且仅当对每个 $a\in A$，包含 $G_a$ 的 $G$ 的子群只有 $G_a$ 与 $G$（即 $G_a$ 是 $G$ 的极大子群，参见第 2.4 节练习 16）。［用 (a)。］
  \end{enumerate}
  \item 集合 $A$ 上的传递置换群 $G$ 称为\emph{二重传递的}，若对任意（等价地，对所有）$a\in A$，子群 $G_a$ 在集合 $A-\{a\}$ 上传递。
  \begin{enumerate}[label=(\alph*)]
    \item 证明对所有 $n\geq2$，$S_n$ 在 $\{1,2,\ldots,n\}$ 上二重传递。
    \item 证明二重传递群是本原的。推出 $D_8$ 在正方形四个顶点上的作用不是二重传递的。
  \end{enumerate}
  \item 假设 $G$ 传递作用在有限集合 $A$ 上，且 $H$ 是 $G$ 的正规子群。令 $O_1,O_2,\ldots,O_r$ 为 $H$ 在 $A$ 上的不同轨道。
  \begin{enumerate}[label=(\alph*)]
    \item 证明 $G$ 置换集合 $O_1,O_2,\ldots,O_r$，即对每个 $g\in G$ 与每个 $i\in\{1,\ldots,r\}$，存在 $j$ 使得 $gO_i=O_j$，其中 $gO=\{g\cdot a\mid a\in O\}$（也就是说，用练习 7 的术语，集合 $O_1,\ldots,O_r$ 是块）。证明 $G$ 在 $\{O_1,\ldots,O_r\}$ 上传递。推出 $H$ 在 $A$ 上的所有轨道有相同的基数。
    \item 证明：若 $a\in O_i$，则 $|O_i|=|H:H\cap G_a|$；并证明 $r=|G:HG_a|$。［画出子群 $H$ 与 $G_a$ 的第二同构定理所对应的子格。注意 $H\cap G_a=H_a$。］
  \end{enumerate}
  \item 设 $H$ 与 $K$ 是群 $G$ 的子群。对每个 $x\in G$，定义 $x$ 在 $G$ 中的 $HK$\emph{双陪集}为
  \[
    HxK=\{hxk\mid h\in H,\ k\in K\}.
  \]
  \begin{enumerate}[label=(\alph*)]
    \item 证明 $HxK$ 是左陪集 $x_1K,\ldots,x_nK$ 的并，其中 $\{x_1K,\ldots,x_nK\}$ 是 $H$ 通过左乘作用在 $K$ 的左陪集集合上时包含 $xK$ 的轨道。
    \item 证明 $HxK$ 是若干 $H$ 的右陪集的并。
    \item 证明对所有 $x,y\in G$，$HxK$ 与 $HyK$ 要么相同，要么不交。证明 $HK$ 双陪集集合分割 $G$。
    \item 证明 $|HxK|=|K|\,|H:H\cap xKx^{-1}|$。
    \item 证明 $|HxK|=|H|\,|K:K\cap x^{-1}Hx|$。
  \end{enumerate}
\end{enumerate}

\section{群通过左乘作用在自身上——Cayley 定理}

本节中 $G$ 是任意群。我们首先考虑 $G$ 通过\emph{左乘}作用在自身上（即 $A=G$）：
\[
  g\cdot a=ga\qquad (g\in G,\ a\in G),
\]
其中 $ga$ 表示 $G$ 中两个群元素 $g$ 与 $a$ 的乘积（若 $G$ 用加法记号，则该作用写作 $g\cdot a=g+a$，称为左平移）。第 1.7 节中我们已经看到这满足群作用的两个公理。

当 $G$ 是阶为 $n$ 的有限群时，为了描述由此作用给出的置换表示，把 $G$ 的元素用整数 $1,2,\ldots,n$ 标号是方便的。这样把 $G$ 的元素列为 $g_1,g_2,\ldots,g_n$，则对每个 $g\in G$，置换 $\sigma_g$ 可作为下标 $1,2,\ldots,n$ 的置换描述如下：
\[
  \sigma_g(i)=j\quad\Longleftrightarrow\quad gg_i=g_j.
\]
对群元素采用不同标号，会给出 $\sigma_g$ 作为 $\{1,2,\ldots,n\}$ 的置换的不同描述（参见练习）。

\paragraph{例}
令 $G=\{1,a,b,c\}$ 为 Klein 四元群，其群表已在第 2.5 节写出。把群元素 $1,a,b,c$ 分别标为整数 $1,2,3,4$。在这个标号下，我们计算由群元素 $a$ 左乘作用诱导的置换 $\sigma_a$：
\begin{align*}
  a\cdot1&=a1=a,       &\text{所以 }\sigma_a(1)&=2,\\
  a\cdot a&=aa=1,      &\text{所以 }\sigma_a(2)&=1,\\
  a\cdot b&=ab=c,      &\text{所以 }\sigma_a(3)&=4,\\
  a\cdot c&=ac=b,      &\text{所以 }\sigma_a(4)&=3.
\end{align*}
在 $G$ 的这个元素标号下，$\sigma_a=(1\ 2)(3\ 4)$。在 Klein 四元群通过左乘作用在自身上相伴随的置换表示中，类似地可计算出
\[
  a\longmapsto \sigma_a=(1\ 2)(3\ 4),\qquad
  b\longmapsto \sigma_b=(1\ 3)(2\ 4),\qquad
  c\longmapsto \sigma_c=(1\ 4)(2\ 3),
\]
这就在此标号下显式给出了与该作用相伴随的置换表示 $G\to S_4$。

容易看出（稍后我们将在更一般的情形中证明），一个群通过左乘作用在自身上的作用总是传递且忠实的，并且任一点的稳定子都是单位子群（对上面的例子，这些事实可以直接检查）。

现在考虑群通过左乘作用在自身元素集合上的推广。令 $H$ 是 $G$ 的任意子群，令 $A$ 为 $H$ 在 $G$ 中的所有左陪集组成的集合。定义 $G$ 在 $A$ 上的作用为
\[
  g\cdot aH=gaH\qquad (g\in G,
\ aH\in A),
\]
其中 $gaH$ 是以 $ga$ 为代表元的左陪集。容易检验这满足群作用的两个公理，即 $G$ 确实通过左乘作用在 $H$ 的左陪集集合上。当 $H$ 是 $G$ 的单位子群时，陪集 $aH$ 只是 $\{a\}$；若把元素 $a$ 与集合 $\{a\}$ 认同，则左乘作用在单位子群的左陪集上的这个作用，就是 $G$ 通过左乘作用在自身上的作用。

当 $H$ 在 $G$ 中指数为有限数 $m$ 时，为了描述由该作用给出的置换表示，把 $H$ 的左陪集用整数 $1,2,\ldots,m$ 标号是方便的。这样把 $H$ 在 $G$ 中的不同左陪集列为 $a_1H,a_2H,\ldots,a_mH$，则对每个 $g\in G$，置换 $\sigma_g$ 可作为下标 $1,2,\ldots,m$ 的置换描述如下：
\[
  \sigma_g(i)=j\quad\Longleftrightarrow\quad ga_iH=a_jH.
\]
不同的陪集标号会给出 $\sigma_g$ 作为 $\{1,2,\ldots,m\}$ 的置换的不同描述（参见练习）。

\paragraph{例}
令 $G=D_8$，$H=\gen{s}$。把不同左陪集 $1H,rH,r^2H,r^3H$ 分别标为整数 $1,2,3,4$。在这个标号下，我们计算由群元素 $s$ 左乘作用在 $H$ 的左陪集上所诱导的置换 $\sigma_s$：
\begin{align*}
  s\cdot1H&=sH=1H,       &\text{所以 }\sigma_s(1)&=1,\\
  s\cdot rH&=srH=r^3H,  &\text{所以 }\sigma_s(2)&=4,\\
  s\cdot r^2H&=sr^2H=r^2H, &\text{所以 }\sigma_s(3)&=3,\\
  s\cdot r^3H&=sr^3H=rH, &\text{所以 }\sigma_s(4)&=2.
\end{align*}
在这个左陪集标号下，得到 $\sigma_s=(2\ 4)$。在 $D_8$ 通过左乘作用在 $\gen{s}$ 的左陪集上相伴随的置换表示中，类似地可得 $\sigma_r=(1\ 2\ 3\ 4)$。注意，置换表示是同态，所以一旦其在 $D_8$ 的生成元上的值确定，其在任何其他元素上的值也随之确定（例如 $\sigma_{sr^2}=\sigma_s\sigma_r^2$）。

\medskip
\noindent\textbf{定理 3.} 设 $G$ 是群，$H$ 是 $G$ 的子群，并令 $G$ 通过左乘作用在 $H$ 在 $G$ 中的左陪集集合 $A$ 上。令 $\pi_H$ 为由此作用给出的相伴随置换表示。则
\begin{enumerate}[label=(\arabic*)]
  \item $G$ 在 $A$ 上传递作用；
  \item $G$ 中点 $1H\in A$ 的稳定子是子群 $H$；
  \item 作用的核（即 $\pi_H$ 的核）为
  \[
    \bigcap_{x\in G}xHx^{-1},
  \]
  并且 $\ker\pi_H$ 是包含在 $H$ 中的 $G$ 的最大正规子群。
\end{enumerate}

\begin{proof}
为证明 $G$ 在 $A$ 上传递，任取 $A$ 中两个元素 $aH$ 与 $bH$，令 $g=ba^{-1}$。则
\[
  g\cdot aH=(ba^{-1})aH=bH,
\]
所以任意两个元素 $aH$ 与 $bH$ 属于同一轨道，(1) 得证。对 (2)，点 $1H$ 的稳定子按定义是
\[
  \{g\in G\mid g\cdot1H=1H\}=\{g\in G\mid gH=H\}=H.
\]
由 $\pi_H$ 的定义，
\begin{align*}
  \ker\pi_H
  &=\{g\in G\mid gxH=xH\text{ 对所有 }x\in G\}\\
  &=\{g\in G\mid (x^{-1}gx)H=H\text{ 对所有 }x\in G\}\\
  &=\{g\in G\mid x^{-1}gx\in H\text{ 对所有 }x\in G\}\\
  &=\{g\in G\mid g\in xHx^{-1}\text{ 对所有 }x\in G\}
    =\bigcap_{x\in G}xHx^{-1},
\end{align*}
这证明了 (3) 的第一断言。第二断言来自如下观察：首先 $\ker\pi_H\normal G$ 且 $\ker\pi_H\leq H$。现在若 $N$ 是包含在 $H$ 中的 $G$ 的任意正规子群，则对所有 $x\in G$，有
\[
  N=xNx^{-1}\leq xHx^{-1},
\]
所以
\[
  N\leq\bigcap_{x\in G}xHx^{-1}=\ker\pi_H.
\]
这说明 $\ker\pi_H$ 是包含在 $H$ 中的 $G$ 的最大正规子群。
\end{proof}

\medskip
\noindent\textbf{推论 4（Cayley 定理）.} 每个群都同构于某个对称群的子群。若 $G$ 是阶为 $n$ 的群，则 $G$ 同构于 $S_n$ 的一个子群。

\begin{proof}
令 $H=1$，并应用前一定理，得到 $G$ 到 $S_G$ 中的同态（这里我们把单位子群的陪集与 $G$ 的元素认同）。由于这个同态的核包含于 $H=1$，所以 $G$ 同构于它在 $S_G$ 中的像。
\end{proof}

注意，$G$ 同构于一个对称群的\emph{子群}，而不是同构于整个对称群本身。例如，我们已经给出了 Klein 四元群与 $S_4$ 的子群
\[
  \gen{(1\ 2)(3\ 4),(1\ 3)(2\ 4)}
\]
之间的同构。回忆对称群的子群称为置换群，因此 Cayley 定理说明每个群都同构于一个置换群。由 $G$ 在自身元素集合上（即 $H=1$ 的陪集上）左乘给出的置换表示称为 $G$ 的\emph{左正则表示}。也许有人会认为，只需研究对称群的子群（以及对所有 $n$ 研究 $S_n$ 的子群）就能更有效地研究所有群（以及所有有限群）。单靠这种方法在计算上和理论上都不实际，因为研究阶为 $n$ 的群时，我们必须在大得多的群 $S_n$ 中工作（参见练习 7）。

从历史上看，有限群最初并不是在我们所发展的公理化背景下研究的，而是作为 $S_n$ 的子群来研究的。因此，Cayley 定理证明了历史上的群概念与现代的（公理化）群概念是等价的。现代观点的一个优点是：在研究给定群时，我们不必局限于把该群看作某个特定对称群的子群（在某种意义上，我们的群是“无坐标的”）。

下面的结果推广了关于指数为 $2$ 的子群正规的结果。

\medskip
\noindent\textbf{推论 5.} 若 $G$ 是阶为 $n$ 的有限群，且 $p$ 是整除 $|G|$ 的最小素数，则任意指数为 $p$ 的子群都是正规子群。

\noindent\textbf{注记.} 一般来说，阶为 $n$ 的群不一定有指数为 $p$ 的子群（例如 $A_4$ 没有指数为 $2$ 的子群）。

\begin{proof}
设 $H\leq G$ 且 $|G:H|=p$。令 $\pi_H$ 为由 $G$ 在 $H$ 的左陪集集合上左乘给出的置换表示，令 $K=\ker\pi_H$，并令 $|H:K|=k$。则
\[
  |G:K|=|G:H|\,|H:K|=pk.
\]
由于 $H$ 有 $p$ 个左陪集，第一同构定理说明 $G/K$ 同构于 $S_p$ 的一个子群（即 $G$ 在 $\pi_H$ 下的像）。由 Lagrange 定理，
\[
  pk=|G/K|\mid p!.
\]
因此 $k\mid p!/p=(p-1)!$。但 $(p-1)!$ 的所有素因子都小于 $p$，而由 $p$ 的最小性，$k$ 的每个素因子都大于或等于 $p$。这迫使 $k=1$，所以 $H=K\normal G$。证明完毕。
\end{proof}

\subsection*{练习}
设 $G$ 是群，$H$ 是 $G$ 的子群。
\begin{enumerate}[label=\arabic*.]
  \item 设 $G=\{1,a,b,c\}$ 是第 2.5 节写出群表的 Klein 四元群。
  \begin{enumerate}[label=(\alph*)]
    \item 把 $1,a,b,c$ 分别标为整数 $1,2,4,3$，证明在 $G$ 到 $S_4$ 的左正则表示下，非单位元素映为
    \[
      a\mapsto(1\ 2)(3\ 4),\qquad b\mapsto(1\ 4)(2\ 3),\qquad c\mapsto(1\ 3)(2\ 4).
    \]
    \item 把 $1,a,b,c$ 重新标为 $1,4,2,3$，计算 $G$ 到 $S_4$ 的左正则表示下 $G$ 每个元素的像。证明在这个标号下 $G$ 在 $S_4$ 中的像与 (a) 中的像是同一个子群（虽然在两个不同标号下，非单位元素逐个映到不同置换）。
  \end{enumerate}
  \item 把 $S_3$ 的元素列为 $1,(1\ 2),(2\ 3),(1\ 3),(1\ 2\ 3),(1\ 3\ 2)$，并分别标为整数 $1,2,3,4,5,6$。给出 $S_3$ 到 $S_6$ 的左正则表示下每个元素的像。
  \item 令 $r,s$ 是阶为 $8$ 的二面体群的通常生成元。
  \begin{enumerate}[label=(\alph*)]
    \item 把 $D_8$ 的元素列为 $1,r,r^2,r^3,s,sr,sr^2,sr^3$，并分别标为整数 $1,2,\ldots,8$。给出 $D_8$ 到 $S_8$ 的左正则表示下每个元素的像。
    \item 把同一元素列表重新标为整数 $1,3,5,7,2,4,6,8$，并相对于这个新标号重新计算左正则表示下每个元素的像。证明 (a) 与 (b) 得到的 $S_8$ 的两个子群不同。
  \end{enumerate}
  \item 用 $Q_8$ 的左正则表示给出 $S_8$ 中两个元素，它们生成一个同构于四元数群 $Q_8$ 的子群。
  \item 令 $r,s$ 是阶为 $8$ 的二面体群的通常生成元，并令 $H=\gen{s}$。把 $H$ 在 $D_8$ 中的左陪集列为 $1H,rH,r^2H,r^3H$。
  \begin{enumerate}[label=(\alph*)]
    \item 把这些陪集分别标为整数 $1,2,3,4$。给出由 $D_8$ 左乘作用在 $H$ 的四个左陪集上得到的表示 $\pi_H:D_8\to S_4$ 下每个元素的像。推出该表示是忠实的（即所得 $S_4$ 的元素构成一个同构于 $D_8$ 的子群）。
    \item 用整数 $1,3,2,4$ 重新标记这些陪集，重复 (a)。证明由这个标号得到的置换构成 $S_4$ 的一个子群，且不同于 (a) 中得到的子群。
    \item 令 $K=\gen{sr}$，把 $K$ 的陪集列为 $1K,rK,r^2K,r^3K$，并标为整数 $1,2,3,4$。证明在这个标号下，由左乘作用在 $K$ 的陪集上得到的表示 $\pi_K$ 在 $S_4$ 中的像与 (a) 中是同一个子群（尽管子群 $H$ 与 $K$ 不同，并且在两个同态下，$D_8$ 的某些元素映到不同置换）。
  \end{enumerate}
  \item 令 $r,s$ 是阶为 $8$ 的二面体群的通常生成元，并令 $N=\gen{r^2}$。把 $N$ 在 $D_8$ 中的左陪集列为 $1N,rN,sN,srN$，并分别标为整数 $1,2,3,4$。给出由 $D_8$ 左乘作用在 $N$ 的四个左陪集上得到的表示 $\pi_N:D_8\to S_4$ 下每个元素的像。推出该表示不是忠实的，并证明 $\pi_N(D_8)$ 同构于 Klein 四元群。
  \item 设 $Q_8$ 是阶为 $8$ 的四元数群。
  \begin{enumerate}[label=(\alph*)]
    \item 证明 $Q_8$ 同构于 $S_8$ 的一个子群。
    \item 证明对任何 $n\leq7$，$Q_8$ 都不同构于 $S_n$ 的子群。［若 $Q_8$ 作用在任意阶不超过 $7$ 的集合 $A$ 上，证明任一点 $a\in A$ 的稳定子必须包含子群 $\gen{-1}$。］
  \end{enumerate}
  \item 证明若 $H$ 的有限指数为 $n$，则存在 $G$ 的正规子群 $K$，满足 $K\leq H$ 且 $|G:K|\leq n!$。
  \item 证明若 $p$ 是素数且 $G$ 是阶为 $p^a$ 的群（$a\in\ZZ^+$），则 $G$ 的每个指数为 $p$ 的子群在 $G$ 中正规。推出每个阶为 $p^2$ 的群都有一个阶为 $p$ 的正规子群。
  \item 证明每个阶为 $6$ 的非阿贝尔群都有一个非正规的阶为 $2$ 的子群。用此分类阶为 $6$ 的群。［构造一个到 $S_3$ 中的单同态。］
  \item 设 $G$ 是有限群，$\pi:G\to S_G$ 是左正则表示。证明若 $x\in G$ 的阶为 $n$ 且 $|G|=mn$，则 $\pi(x)$ 是 $m$ 个 $n$-循环的乘积。推出 $\pi(x)$ 是奇置换当且仅当 $|x|$ 为偶数且 $|G|/|x|$ 为奇数。
  \item 设 $G$ 与 $\pi$ 如前一题。证明若 $\pi(G)$ 含有奇置换，则 $G$ 有一个指数为 $2$ 的子群。［用第 3.3 节练习 3。］
  \item 证明若 $|G|=2k$ 且 $k$ 为奇数，则 $G$ 有一个指数为 $2$ 的子群。［用 Cauchy 定理产生一个阶为 $2$ 的元素，然后用前两题。］
  \item 设 $G$ 是复合阶 $n$ 的有限群，并且具有如下性质：对每个整除 $n$ 的正整数 $k$，$G$ 都有一个阶为 $k$ 的子群。证明 $G$ 不是单群。
\end{enumerate}

\section{群通过共轭作用在自身上——类方程}

本节中 $G$ 是任意群。我们首先考虑 $G$ 通过\emph{共轭}作用在自身上（即 $A=G$）：
\[
  g\cdot a=gag^{-1}
  \qquad (g\in G,
\ a\in G),
\]
其中 $gag^{-1}$ 按照 $G$ 中的乘法计算。这个定义满足群作用的两个公理，因为
\begin{align*}
  g_1\cdot(g_2\cdot a)
  &=g_1\cdot(g_2ag_2^{-1})
    =g_1(g_2ag_2^{-1})g_1^{-1}\\
  &=(g_1g_2)a(g_1g_2)^{-1}
    =(g_1g_2)\cdot a,
\end{align*}
且
\[
  1\cdot a=1a1^{-1}=a
\]
对所有 $g_1,g_2\in G$ 与 $a\in G$ 成立。

\noindent\textbf{定义.} 若存在某个 $g\in G$ 使得 $b=gag^{-1}$，则称 $G$ 中两个元素 $a$ 与 $b$ 在 $G$ 中\emph{共轭}（等价地，它们属于 $G$ 通过共轭作用在自身上时的同一轨道）。$G$ 通过共轭作用在自身上的轨道称为 $G$ 的\emph{共轭类}。

\paragraph{例}
\begin{enumerate}[label=(\arabic*)]
  \item 若 $G$ 是阿贝尔群，则 $G$ 在自身上的共轭作用是平凡作用：对所有 $g,a\in G$，有 $g\cdot a=a$；并且每个 $a\in G$ 的共轭类为 $\{a\}$。
  \item 若 $|G|>1$，则与左乘作用不同，$G$ 通过共轭作用在自身上不是传递的，因为 $\{1\}$ 总是一个共轭类（即此作用的一个轨道）。更一般地，单元素集合 $\{a\}$ 是一个共轭类，当且仅当对所有 $g\in G$ 有 $gag^{-1}=a$，当且仅当 $a$ 属于 $G$ 的中心。
  \item 在 $S_3$ 中可直接计算出共轭类为
  \[
    \{1\},\qquad \{(1\ 2),(1\ 3),(2\ 3)\},\qquad \{(1\ 2\ 3),(1\ 3\ 2)\}.
  \]
  我们很快会发展出更容易计算共轭类的技术，尤其是在对称群中。
\end{enumerate}

与群通过左乘作用在自身上的情形类似，共轭作用也可推广。若 $S$ 是 $G$ 的任意子集，定义
\[
  gSg^{-1}=\{gsg^{-1}\mid s\in S\}.
\]
群 $G$ 通过 $g\cdot S=gSg^{-1}$ 作用在自身所有子集组成的集合 $\mathcal P(G)$ 上。与上面一样，这定义了 $G$ 在 $\mathcal P(G)$ 上的群作用。注意，若 $S$ 是单元素集合 $\{s\}$，则 $g\cdot S$ 是单元素集合 $\{gsg^{-1}\}$，所以 $G$ 在所有子集上的这个作用可看作 $G$ 通过共轭作用在自身上的扩张。

\noindent\textbf{定义.} 若存在某个 $g\in G$ 使得 $T=gSg^{-1}$，则称 $G$ 的两个子集 $S$ 与 $T$ 在 $G$ 中\emph{共轭}（等价地，它们属于 $G$ 通过共轭作用在其子集上时的同一轨道）。

现在把命题 2 应用于 $G$ 的共轭作用。命题 2 说明，若 $S$ 是 $G$ 的一个子集，则 $S$ 的共轭子集个数等于 $S$ 的稳定子 $G_S$ 的指数 $|G:G_S|$。对共轭作用，
\[
  G_S=\{g\in G\mid gSg^{-1}=S\}=N_G(S),
\]
这正是 $S$ 在 $G$ 中的正规化子。总结如下。

\medskip
\noindent\textbf{命题 6.} 群 $G$ 中一个子集 $S$ 的共轭子集个数为其正规化子的指数 $|G:N_G(S)|$。特别地，$G$ 的一个元素 $s$ 的共轭元素个数为其中心化子的指数 $|G:C_G(s)|$。

\begin{proof}
命题的第二断言来自观察 $N_G(\{s\})=C_G(s)$。
\end{proof}

$G$ 通过共轭作用在自身上把 $G$ 分割为 $G$ 的共轭类，而这些共轭类的阶可由命题 6 计算。由于这些共轭类的阶数之和等于 $G$ 的阶，我们得到这些阶数之间的如下重要关系。

\medskip
\noindent\textbf{定理 7（类方程）.} 设 $G$ 是有限群，且 $g_1,g_2,\ldots,g_r$ 是 $G$ 的不包含于中心 $Z(G)$ 的不同共轭类的代表元。则
\[
  |G|=|Z(G)|+\sum_{i=1}^r |G:C_G(g_i)|.
\]

\begin{proof}
如上面例 2 所指出，元素 $\{x\}$ 是大小为 $1$ 的共轭类，当且仅当 $x\in Z(G)$，因为此时对所有 $g\in G$ 有 $gxg^{-1}=x$。设
\[
  Z(G)=\{1,z_2,\ldots,z_m\},
\]
并设 $K_1,K_2,\ldots,K_r$ 是 $G$ 中不包含于中心的共轭类，且 $g_i$ 是 $K_i$ 的代表元。于是 $G$ 的全部共轭类为
\[
  \{1\},\{z_2\},\ldots,\{z_m\},K_1,K_2,\ldots,K_r.
\]
由于它们分割 $G$，有
\[
  |G|=\sum_{i=1}^m1+\sum_{i=1}^r |K_i|
      =|Z(G)|+\sum_{i=1}^r |G:C_G(g_i)|,
\]
其中 $|K_i|$ 由命题 6 给出。这证明了类方程。
\end{proof}

尤其注意，类方程右边的所有求和项都是群阶的因子，因为它们是 $G$ 的子群指数。这限制了它们可能的取值（参见练习 6）。

\paragraph{例}
\begin{enumerate}[label=(\arabic*)]
  \item 在阿贝尔群中，类方程不给出信息，因为共轭作用是平凡作用，所有共轭类大小都是 $1$。
  \item 在任意群 $G$ 中，都有 $\gen{g}\leq C_G(g)$；这个观察有助于减少计算共轭类时的工作量。例如，在四元数群 $Q_8$ 中，我们看到
  \[
    \gen{i}\leq C_{Q_8}(i)\leq Q_8.
  \]
  由于 $i\notin Z(Q_8)$ 且 $|Q_8:\gen{i}|=2$，必须有 $C_{Q_8}(i)=\gen{i}$。因此 $i$ 在 $Q_8$ 中恰有两个共轭元素，即 $i$ 与 $-i=kik^{-1}$。$Q_8$ 的其他共轭类也类似确定，为
  \[
    \{1\},\quad \{-1\},\quad \{\pm i\},\quad \{\pm j\},\quad \{\pm k\}.
  \]
  前两个类组成 $Z(Q_8)$，该群的类方程为
  \[
    |Q_8|=2+2+2+2.
  \]
  \item 在 $D_8$ 中，也可以利用三个指数为 $2$ 的子群都是阿贝尔的这一事实迅速看出：若 $x\notin Z(D_8)$，则 $|C_{D_8}(x)|=4$。$D_8$ 的共轭类为
  \[
    \{1\},\quad \{r^2\},\quad \{r,r^3\},\quad \{s,sr^2\},\quad \{sr,sr^3\}.
  \]
  前两个类组成 $Z(D_8)$，该群的类方程为
  \[
    |D_8|=2+2+2+2.
  \]
\end{enumerate}

在讨论更多共轭例子之前，我们先给出类方程的两个重要后果。类方程的第一个应用是证明素数幂阶群有非平凡中心，这是研究素数幂阶群（第 6 章还会回到这一主题）的起点。

\medskip
\noindent\textbf{定理 8.} 若 $p$ 是素数，且 $P$ 是阶为 $p^a$ 的群，其中 $a\geq1$，则 $P$ 有非平凡中心：$Z(P)\neq1$。

\begin{proof}
由类方程，
\[
  |P|=|Z(P)|+\sum_{i=1}^r |P:C_P(g_i)|,
\]
其中 $g_1,\ldots,g_r$ 是不同的非中心共轭类的代表元。按定义，对 $i=1,2,\ldots,r$，有 $C_P(g_i)\neq P$，所以 $p$ 整除 $|P:C_P(g_i)|$。由于 $p$ 也整除 $|P|$，可知 $p$ 整除 $|Z(P)|$，从而中心非平凡。
\end{proof}

\medskip
\noindent\textbf{推论 9.} 若 $|P|=p^2$，其中 $p$ 是素数，则 $P$ 是阿贝尔群。更精确地，$P$ 同构于 $\ZZ/p^2\ZZ$ 或 $\ZZ/p\ZZ\times\ZZ/p\ZZ$。

\begin{proof}
由定理，$Z(P)\neq1$，于是 $P/Z(P)$ 为循环群。由第 3.1 节练习 36，$P$ 是阿贝尔群。若 $P$ 有阶为 $p^2$ 的元素，则 $P$ 是循环群。否则，假设 $P$ 的每个非单位元素阶都为 $p$。取任一非单位元素 $x\in P$，并取 $y\in P-\gen{x}$。由于 $|\gen{x,y}|>|\gen{x}|=p$，必有 $P=\gen{x,y}$。$x$ 与 $y$ 的阶都是 $p$，所以 $\gen{x}\times\gen{y}\cong \ZZ/p\ZZ\times\ZZ/p\ZZ$。直接可见映射
\[
  (x^a,y^b)\longmapsto x^ay^b
\]
是从 $\gen{x}\times\gen{y}$ 到 $P$ 上的同构。证明完毕。
\end{proof}

\paragraph{$S_n$ 中的共轭}
下面考察对称群中的共轭。熟悉线性代数的读者会认出，在矩阵群 $\GL_n(F)$ 中，共轭与“换基”相同：$A\mapsto PAP^{-1}$。$S_n$ 中的情形与此类似。

\medskip
\noindent\textbf{命题 10.} 设 $\sigma,\tau$ 是对称群 $S_n$ 的元素，并且 $\sigma$ 的循环分解为
\[
  (a_1\ a_2\ \cdots\ a_{k_1})(b_1\ b_2\ \cdots\ b_{k_2})\cdots.
\]
则 $\tau\sigma\tau^{-1}$ 的循环分解为
\[
  (\tau(a_1)\ \tau(a_2)\ \cdots\ \tau(a_{k_1}))
  (\tau(b_1)\ \tau(b_2)\ \cdots\ \tau(b_{k_2}))\cdots,
\]
也就是说，$\tau\sigma\tau^{-1}$ 由 $\sigma$ 的循环分解中每个条目 $i$ 替换为 $\tau(i)$ 得到。

\begin{proof}
注意若 $\sigma(i)=j$，则
\[
  \tau\sigma\tau^{-1}(\tau(i))=\tau(j).
\]
因此，若有序对 $i,j$ 出现在 $\sigma$ 的循环分解中，则有序对 $\tau(i),\tau(j)$ 出现在 $\tau\sigma\tau^{-1}$ 的循环分解中。证明完毕。
\end{proof}

\paragraph{例}
令 $\sigma=(1\ 2)(3\ 4\ 5)(6\ 7\ 8\ 9)$，并令 $\tau=(1\ 3\ 5\ 7)(2\ 4\ 6\ 8)$。则
\[
  \tau\sigma\tau^{-1}=(3\ 4)(5\ 6\ 7)(8\ 1\ 2\ 9).
\]

\noindent\textbf{定义.}
\begin{enumerate}[label=(\arabic*)]
  \item 若 $\sigma\in S_n$ 是若干不交循环的乘积，这些循环长度为 $n_1,n_2,\ldots,n_r$，且 $n_1\leq n_2\leq\cdots\leq n_r$（包括 $1$-循环），则整数 $n_1,n_2,\ldots,n_r$ 称为 $\sigma$ 的\emph{循环型}。
  \item 若 $n\in\ZZ^+$，则 $n$ 的一个\emph{分拆}是一个和为 $n$ 的正整数非降序列。
\end{enumerate}

注意，由前一节的结果，置换的循环型是唯一的。例如，$S_n$ 中一个 $m$-循环的循环型为
\[
  1,1,\ldots,1,m,
\]
其中 $m$ 前面有 $n-m$ 个 $1$。

\medskip
\noindent\textbf{命题 11.} $S_n$ 的两个元素在 $S_n$ 中共轭，当且仅当它们有相同的循环型。$S_n$ 的共轭类个数等于 $n$ 的分拆数。

\begin{proof}
由命题 10，共轭的置换有相同循环型。反过来，假设置换 $\sigma_1$ 与 $\sigma_2$ 有相同循环型。把它们的循环按长度非降排序，并包括 $1$-循环（若 $\sigma_1$ 与 $\sigma_2$ 有多个相同长度的循环，则这种排序有若干选择）。忽略括号后，每个循环分解都是一个列表，其中整数 $1$ 到 $n$ 各恰好出现一次。定义 $\tau$ 为把 $\sigma_1$ 的列表中第 $i$ 个整数送到 $\sigma_2$ 的列表中第 $i$ 个整数的函数。于是 $\tau$ 是一个置换；并且由于标示循环分解的括号在两个列表中出现在相同位置，命题 10 保证 $\tau\sigma_1\tau^{-1}=\sigma_2$，所以 $\sigma_1$ 与 $\sigma_2$ 共轭。

由于 $S_n$ 的共轭类与允许的循环型之间有一一对应，而 $S_n$ 中置换的每个循环型都是 $n$ 的一个分拆，命题的第二断言随之成立。证明完毕。
\end{proof}

\paragraph{例}
\begin{enumerate}[label=(\arabic*)]
  \item 令
  \[
    \sigma_1=(1)(3\ 5)(8\ 9)(2\ 4\ 7\ 6),
    \quad
    \sigma_2=(3)(4\ 7)(8\ 1)(5\ 2\ 6\ 9).
  \]
  定义 $\tau$ 为 $\tau(1)=3,\tau(3)=4,\tau(5)=7,\tau(8)=8$，等等。则
  \[
    \tau=(1\ 3\ 4\ 2\ 5\ 7\ 6\ 9)(8),
  \]
  且 $\tau\sigma_1\tau^{-1}=\sigma_2$。
  \item 若在前例中，把 $\sigma_2$ 重新排列为
  \[
    \sigma_2=(3)(8\ 1)(4\ 7)(5\ 2\ 6\ 9),
  \]
  即交换两个长度为 $2$ 的循环，则按上述方法得到的相应 $\tau$ 由 $\tau(1)=3,\tau(3)=8,\tau(5)=1,\tau(8)=4$ 等定义，给出
  \[
    \tau=(1\ 3\ 8\ 4\ 2\ 5)(6\ 9\ 7),
  \]
  同样有 $\tau\sigma_1\tau^{-1}=\sigma_2$。这说明把 $\sigma_1$ 共轭到 $\sigma_2$ 的元素有很多。
  \item 若 $n=5$，则 $5$ 的分拆与相应共轭类代表元（不写 $1$-循环）如下：
  \[
  \begin{array}{c|c}
    5\text{ 的分拆} & \text{共轭类代表元}\\ \hline
    1,1,1,1,1 & 1\\
    1,1,1,2 & (1\ 2)\\
    1,1,3 & (1\ 2\ 3)\\
    1,4 & (1\ 2\ 3\ 4)\\
    5 & (1\ 2\ 3\ 4\ 5)\\
    1,2,2 & (1\ 2)(3\ 4)\\
    2,3 & (1\ 2)(3\ 4\ 5)
  \end{array}
  \]
\end{enumerate}

命题 11 与命题 6 可用于给出 $S_n$ 中某些元素的中心化子。例如，若 $\sigma$ 是 $S_n$ 中的 $m$-循环，则 $\sigma$ 的共轭元素个数（即 $m$-循环个数）为
\[
  \frac{n(n-1)\cdots(n-m+1)}{m}.
\]
由命题 6，这是 $\sigma$ 的中心化子的指数：
\[
  \frac{|S_n|}{|C_{S_n}(\sigma)|}.
\]
由于 $|S_n|=n!$，得到
\[
  |C_{S_n}(\sigma)|=m\,(n-m)!.
\]
元素 $\sigma$ 当然与 $1,\sigma,\sigma^2,\ldots,\sigma^{m-1}$ 交换。它也与 $S_n$ 中任何循环与 $\sigma$ 不交的置换交换，而这种置换有 $(n-m)!$ 个（即在不出现在 $\sigma$ 中的数字上的全对称群）。这两类元素的乘积已经给出 $m(n-m)!$ 个与 $\sigma$ 交换的元素。由上面的阶数计算，这就是 $S_n$ 中 $\sigma$ 的整个中心化子。显式地，
\[
  \text{若 }\sigma\text{ 是 }S_n\text{ 中的 }m\text{-循环，则 }
  C_{S_n}(\sigma)=\{\sigma^i\tau\mid 0\leq i\leq m-1,
  \ \tau\in S_{n-m}\},
\]
其中 $S_{n-m}$ 表示 $S_n$ 中固定所有出现在 $m$-循环 $\sigma$ 中的整数的子群（当 $m=n$ 或 $m=n-1$ 时，它是单位子群）。

例如，$S_7$ 中 $\sigma=(1\ 3\ 5)$ 的中心化子为
\[
  \{(1\ 3\ 5)^i\tau\mid i=0,1,2,
  \ \tau\text{ 固定 }1,3,5\}.
\]
注意 $\tau\in S_A$，其中 $A=\{2,4,6,7\}$，所以 $\tau$ 有 $4!$ 种选择，中心化子阶为 $3\cdot4!=72$。

我们将在下面的练习以及第 5 章讨论 $S_n$ 中其他元素的中心化子。

可以用上述关于 $S_n$ 中共轭类的讨论，给出 $A_5$ 单性的一个组合证明。首先观察：群 $G$ 的正规子群是 $G$ 的共轭类的并，即
\[
  \text{若 }H\normal G,\text{ 则对 }G\text{ 的每个共轭类 }K,
  \text{ 有 }K\subseteq H\text{ 或 }K\cap H=\varnothing.
\]
这是因为若 $x\in K\cap H$，则对所有 $g\in G$，有 $gxg^{-1}\in gHg^{-1}$。由于 $H$ 正规，$gHg^{-1}=H$，所以 $H$ 含有 $x$ 的所有共轭元素，即 $K\subseteq H$。

\medskip
\noindent\textbf{定理 12.} $A_5$ 是单群。

\begin{proof}
先求 $A_5$ 的共轭类及其阶。命题 11 不能直接应用，因为相同循环型的两个元素（它们在 $S_5$ 中共轭）未必在 $A_5$ 中共轭。练习 19 到 22 将详细分析 $S_n$ 中的共轭类与 $A_n$ 中的共轭类之间的关系。

我们已经看到，偶置换的循环型可取代表元
\[
  1,\\quad (1\ 2\ 3),\quad (1\ 2\ 3\ 4\ 5),\quad (1\ 2)(3\ 4).
\]
前面已经确定了 $S_5$ 中 $3$-循环与 $5$-循环的中心化子；检查其中哪些元素属于 $A_5$，可见
\[
  C_{A_5}((1\ 2\ 3))=\gen{(1\ 2\ 3)},
  \qquad
  C_{A_5}((1\ 2\ 3\ 4\ 5))=\gen{(1\ 2\ 3\ 4\ 5)}.
\]
这些群的阶分别为 $3$ 与 $5$（指数分别为 $20$ 与 $12$），所以在 $A_5$ 中，$(1\ 2\ 3)$ 有 $20$ 个不同共轭元素，而 $(1\ 2\ 3\ 4\ 5)$ 有 $12$ 个不同共轭元素。由于 $S_5$ 中总共有 $20$ 个 $3$-循环（第 1.3 节练习 16），且它们全在 $A_5$ 中，所以
\[
  \text{所有 }20\text{ 个 }3\text{-循环在 }A_5\text{ 中共轭}.
\]
$A_5$ 中总共有 $24$ 个 $5$-循环，但 $5$-循环 $(1\ 2\ 3\ 4\ 5)$ 在 $A_5$ 中只有 $12$ 个不同共轭元素。因此某个 $5$-循环 $\sigma$ 在 $A_5$ 中不与 $(1\ 2\ 3\ 4\ 5)$ 共轭（事实上，$(1\ 3\ 5\ 2\ 4)$ 在 $A_5$ 中不与 $(1\ 2\ 3\ 4\ 5)$ 共轭，因为命题 11 的证明方法表明，$S_5$ 中任何把 $(1\ 2\ 3\ 4\ 5)$ 共轭到 $(1\ 3\ 5\ 2\ 4)$ 的元素都必须是奇置换）。如上可见，$\sigma$ 在 $A_5$ 中也有 $12$ 个不同共轭元素，因此
\[
  5\text{-循环在 }A_5\text{ 中分成两个共轭类，每类有 }12\text{ 个元素}.
\]
由于 $3$-循环与 $5$-循环已经包含了所有奇数阶的非单位元素，$A_5$ 中剩下的 $15$ 个非单位元素必为阶 $2$，因此循环型为 $(2,2)$。容易看出 $(1\ 2)(3\ 4)$ 与 $(1\ 3)(2\ 4)$ 交换，但不与 $A_5$ 中任何奇数阶元素交换。因此 $|C_{A_5}((1\ 2)(3\ 4))|=4$。于是 $(1\ 2)(3\ 4)$ 在 $A_5$ 中有 $15$ 个不同共轭元素，所以
\[
  A_5\text{ 中所有 }15\text{ 个阶为 }2\text{ 的元素都与 }(1\ 2)(3\ 4)\text{ 共轭}.
\]
综上，$A_5$ 的共轭类阶为 $1,15,20,12,12$。

现在假设 $H$ 是 $A_5$ 的正规子群。由上面的观察，$H$ 必为 $A_5$ 的若干共轭类的并。于是 $|H|$ 既要整除 $60$（$A_5$ 的阶），又要是集合 $\{1,12,12,15,20\}$ 中某些数之和。直接检查可知唯一可能为 $|H|=1$ 或 $|H|=60$。因此 $A_5$ 没有真非平凡正规子群。
\end{proof}

\paragraph{右群作用}
如第 1.7 节所述，在作用的定义中，群元素写在集合元素左侧，因此我们的作用概念更准确地说是\emph{左群作用}。类似地，可以定义群 $G$ 在非空集合 $A$ 上的\emph{右群作用}：这是从 $A\times G$ 到 $A$ 的映射，对 $a\in A$ 与 $g\in G$ 记为 $a\cdot g$，并满足公理
\begin{enumerate}[label=(\arabic*)]
  \item $(a\cdot g_1)\cdot g_2=a\cdot(g_1g_2)$，对所有 $a\in A$ 与 $g_1,g_2\in G$ 成立；
  \item $a\cdot1=a$，对所有 $a\in A$ 成立。
\end{enumerate}

在很多群论文献中，共轭用如下记号写成右群作用：
\[
  a^g=g^{-1}ag\qquad(g,a\in G).
\]
类似地，对 $G$ 的子集 $S$，定义 $S^g=g^{-1}Sg$。在这个记号下，右作用的两个公理验证为
\[
  (a^{g_1})^{g_2}=a^{g_1g_2},
  \qquad a^1=a
\]
对所有 $g_1,g_2,a\in G$ 成立。因此，群在自身上的这个右作用的两个公理呈现为熟悉的“指数法则”。（注意，群元素 $a$ 的整数幂 $a^n$ 很容易与 $a$ 关于 $g$ 的共轭 $a^g$ 区分：前者的指数 $n\in\ZZ$，后者的指数 $g\in G$。）由于共轭在群论中无处不在，这个记号是一种有用而高效的简写（相较于总是写 $gag^{-1}$ 或把共轭看作左作用时写 $g\cdot a$）。

对任意群作用，容易检验：若给定 $G$ 在 $A$ 上的一个左群作用，则由
\[
  a\cdot g=g^{-1}\cdot a
\]
定义的映射 $A\times G\to A$ 是一个右群作用。反过来，给定 $G$ 在 $A$ 上的右群作用，也可由
\[
  g\cdot a=a\cdot g^{-1}
\]
构造一个左群作用。称这两种作用互为\emph{对应的群作用}。换言之，在对应的群作用中，$g$ 在左边的作用与 $g^{-1}$ 在右边的作用相同。对共轭作用来说这一点尤其透明，因为“$a$ 被 $g$ 左共轭”得到的元素 $gag^{-1}$，与“$a$ 被 $g^{-1}$ 右共轭”得到的元素 $a^{g^{-1}}$ 是同一个群元素。因此，群的两个元素或子集“左共轭”当且仅当它们“右共轭”，所以左、右对应作用给出的“共轭”关系相同。更一般地，也容易通过练习 1 看出，对任意对应的左、右作用，其轨道相同。

我们一直使用左作用，因为它与函数在左边作用的记号相容（即与 $\varphi(g)$ 的记号相容）；这样，左乘作用在子群的左陪集上就是左作用。类似地，右乘作用在子群的右陪集上是右作用；若函数 $\varphi:G\to S_A$ 写在右边（如 $(g_1g_2)\varphi$），并且 $S_A$ 中的置换也作为从 $A$ 到自身的函数写在右边，则相伴随的置换表示 $\varphi$ 是同态。有些情形中，一个集合同时承载群 $G$ 的两个作用：一个自然地在左边，另一个自然地在右边，因此熟悉两种作用都是有用的。

\subsection*{练习}
设 $G$ 是群。
\begin{enumerate}[label=\arabic*.]
  \item 假设 $G$ 左作用在集合 $A$ 上，对 $g\in G,a\in A$ 记为 $g\cdot a$。把对应的右作用记为 $a\cdot g$。证明由
  \[
    a\sim b\Longleftrightarrow a=g\cdot b\text{ 对某个 }g\in G,
  \]
  以及
  \[
    a\sim' b\Longleftrightarrow a=b\cdot g\text{ 对某个 }g\in G
  \]
  定义的等价关系相同（即 $a\sim b$ 当且仅当 $a\sim' b$）。
  \item 求下列群的所有共轭类及其大小：
  \begin{enumerate}[label=(\alph*)]
    \item $D_8$；\item $Q_8$；\item $A_4$。
  \end{enumerate}
  \item 求下列群的所有共轭类及其大小：
  \begin{enumerate}[label=(\alph*)]
    \item $\ZZ/2\ZZ\times S_3$；\item $S_3\times S_3$；\item $\ZZ/3\ZZ\times A_4$。
  \end{enumerate}
  \item 证明若 $S\subseteq G$ 且 $g\in G$，则 $gN_G(S)g^{-1}=N_G(gSg^{-1})$，并且 $gC_G(S)g^{-1}=C_G(gSg^{-1})$。
  \item 若 $Z(G)$ 在 $G$ 中指数为 $n$，证明每个共轭类至多有 $n$ 个元素。
  \item 假设 $G$ 是阶为 $15$ 的非阿贝尔群。证明 $Z(G)=1$。利用对所有 $g\in G$ 有 $\gen{g}\leq C_G(g)$ 的事实，说明 $G$ 至多只有一种可能的类方程。［用第 3.1 节练习 36。］
  \item 对 $n=3,4,6,7$，列出 $n$ 的分拆，并给出 $S_n$ 中相应共轭类的代表元。
  \item 证明对所有 $n\geq3$，$Z(S_n)=1$。
  \item 证明对所有 $n\geq4$，$|C_{S_n}((1\ 2)(3\ 4))|=8(n-4)!$。显式确定这个中心化子的元素。
  \item 令 $\sigma=(1\ 2\ 3\ 4\ 5)\in S_5$。在 (a)--(c) 中分别找出满足指定共轭关系的显式元素 $\tau\in S_5$：
  \begin{enumerate}[label=(\alph*)]
    \item $\tau\sigma\tau^{-1}=\sigma^2$；
    \item $\tau\sigma\tau^{-1}=\sigma^{-1}$；
    \item $\tau\sigma\tau^{-1}=\sigma^{-2}$。
  \end{enumerate}
  \item 在 (a)--(d) 中判断 $\sigma_1$ 与 $\sigma_2$ 是否共轭。若共轭，给出显式置换 $\tau$ 使得 $\tau\sigma_1\tau^{-1}=\sigma_2$。
  \begin{enumerate}[label=(\alph*)]
    \item $\sigma_1=(1\ 2)(3\ 4\ 5)$，$\sigma_2=(1\ 2\ 3)(4\ 5)$；
    \item $\sigma_1=(1\ 5)(3\ 7\ 2)(10\ 6\ 8\ 11)$，$\sigma_2=(3\ 7\ 5\ 10)(4\ 9)(1\ 3\ 11\ 2)$；
    \item $\sigma_1=(1\ 5)(3\ 7\ 2)(10\ 6\ 8\ 11)$，$\sigma_2=\sigma_1^2$；
    \item $\sigma_1=(1\ 3)(2\ 4\ 6)$，$\sigma_2=(3\ 5)(2\ 4)(5\ 6)$。
  \end{enumerate}
  \item 分别在 $S_8$ 与 $S_{12}$ 中，为每个阶为 $4$ 的元素共轭类找一个代表元。
  \item 求所有恰有两个共轭类的有限群。
  \item 在第 2 节练习 1 中，对 Klein 四元群 $V=\{1,a,b,c\}$ 的元素选择了两个标号，给出了 $V$ 到 $S_4$ 的两个左正则表示版本。令 $\pi_1$ 为该练习 (a) 中得到的正则表示版本，$\pi_2$ 为该练习 (b) 的标号得到的版本。令 $\tau=(2\ 4)$。证明对每个 $g\in V$，有
  \[
    \tau\circ\pi_1(g)\circ\tau^{-1}=\pi_2(g),
  \]
  即由 $\tau$ 共轭把 $\pi_1$ 的像逐元素送到 $\pi_2$ 的像。
  \item 找出 $S_8$ 的一个元素，它把第 2 节练习 3(a) 得到的 $S_8$ 子群共轭到该练习 3(b) 得到的 $S_8$ 子群（这两个子群都同构于 $D_8$）。
  \item 找出 $S_4$ 的一个元素，它把第 2 节练习 5(a) 得到的 $S_4$ 子群共轭到该练习 5(b) 得到的 $S_4$ 子群（这两个子群都同构于 $D_8$）。
  \item 设 $A$ 是非空集合，$X$ 是 $S_A$ 的任意子集。令
  \[
    F(X)=\{a\in A\mid \sigma(a)=a\text{ 对所有 }\sigma\in X\}
  \]
  为 $X$ 的不动点集。令 $M(X)=A-F(X)$ 为被 $X$ 的某个元素移动的元素集合。令
  \[
    D=\{\sigma\in S_A\mid |M(\sigma)|<\infty\}.
  \]
  证明 $D$ 是 $S_A$ 的正规子群。
  \item 设 $A$ 是集合，$H$ 是 $S_A$ 的子群，并令 $F(H)$ 为前一题所定义的 $H$ 在 $A$ 上的不动点集。证明若 $\tau\in N_{S_A}(H)$，则 $\tau$ 稳定集合 $F(H)$ 及其补集 $A-F(H)$。
  \item 假设 $H\normal G$，$K$ 是 $G$ 中包含于 $H$ 的一个共轭类，且 $x\in K$。证明 $K$ 是 $H$ 中 $k$ 个大小相等的共轭类的并，其中 $k=|G:HC_G(x)|$。推出：$S_n$ 中一个由偶置换组成的共轭类，在 $A_n$ 的作用下要么仍是一个共轭类，要么分裂为两个大小相等的共轭类。［令 $A=C_G(x)$、$B=H$，则 $A\cap B=C_H(x)$。画出第二同构定理所对应的格图，并解释相应指数。亦参见第 1 节练习 9。］
  \item 令 $\sigma\in A_n$。证明 $\sigma$ 在 $S_n$ 中的共轭类（即所有与 $\sigma$ 有相同循环型的元素）在 $A_n$ 中仍为一个共轭类，当且仅当 $\sigma$ 与某个奇置换交换。［用前一题。］
  \item 设 $K$ 是 $S_n$ 的一个共轭类，且 $K\subseteq A_n$。证明 $\sigma\in S_n$ 不与任何奇置换交换，当且仅当 $\sigma$ 的循环型由互不相同的奇整数组成。推出：$K$ 分裂为 $A_n$ 中两个共轭类，当且仅当 $K$ 的元素的循环型由互不相同的奇整数组成。［先假设 $\sigma\in K$ 不与任何奇置换交换。观察 $\sigma$ 与其循环分解中的每个单独循环交换——用这一点说明其所有循环长度都必须为奇数。若两个循环有相同奇长度 $k$，找出 $k$ 个换位的乘积，它交换这两个循环并与 $\sigma$ 交换。反过来，若 $\sigma$ 的循环型由互不相同的整数组成，证明 $\sigma$ 只与由其循环分解中各循环生成的群交换。］
  \item 证明若 $n$ 为奇数，则所有 $n$-循环组成的集合在 $A_n$ 中分成两个大小相等的共轭类。
  \item 回忆（参见第 2.4 节练习 16）：$G$ 的真子群 $M$ 称为极大的，若只要 $M\leq H\leq G$，就有 $H=M$ 或 $H=G$。证明若 $M$ 是 $G$ 的极大子群，则 $N_G(M)=M$ 或 $N_G(M)=G$。推出：若 $M$ 是 $G$ 的不正规的极大子群，则 $G$ 中包含在 $M$ 的共轭子群中的非单位元素个数至多为 $(|M|-1)|G:M|$。
  \item 假设 $H$ 是有限群 $G$ 的真子群。证明
  \[
    G\neq\bigcup_{g\in G}gHg^{-1},
  \]
  即 $G$ 不是任何真子群的所有共轭子群的并。［把 $H$ 放入某个极大子群，并用前一题。］
  \item 令 $G=\GL_2(\CC)$，并令
  \[
    H=\left\{\begin{pmatrix}a&b\\0&c\end{pmatrix}\mid a,b,c\in\CC,\ ac\neq0\right\}.
  \]
  证明 $G$ 的每个元素都与 $H$ 的某个元素共轭，并推出 $G$ 是 $H$ 的共轭子群的并。［证明 $\GL_2(\CC)$ 的每个元素都有特征向量。］
  \item 设 $G$ 是有限集合 $A$ 上的传递置换群且 $|A|>1$。证明存在某个 $\sigma\in G$，使得对所有 $a\in A$ 都有 $\sigma(a)\neq a$（这样的元素称为无不动点元素）。
  \item 令 $g_1,g_2,\ldots,g_r$ 是有限群 $G$ 的共轭类代表元，并假设这些元素两两交换。证明 $G$ 是阿贝尔群。
  \item 设 $p,q$ 是素数且 $p<q$。证明一个阶为 $pq$ 的非阿贝尔群 $G$ 有一个指数为 $q$ 的非正规子群，从而存在一个到 $S_q$ 的单同态。推出 $G$ 同构于 $S_q$ 中由 $q$-循环 $(1\ 2\ \cdots\ q)$ 生成的循环群的正规化子的一个子群。
  \item 设 $p$ 是素数，且 $G$ 是阶为 $p^a$ 的群。证明对每个 $0\leq \beta\leq a$，$G$ 都有阶为 $p^\beta$ 的子群。［用定理 8 并对 $a$ 归纳。］
  \item 若 $G$ 是奇数阶群，证明对任一非单位元素 $x\in G$，$x$ 与 $x^{-1}$ 在 $G$ 中不共轭。
  \item 用二面体群 $D_{2n}$ 的通常生成元与关系（参见第 1.2 节），证明当 $n=2k$ 为偶数时，$D_{2n}$ 中的共轭类为
  \[
    \{1\},\ \{r^k\},\ \{r^{\pm1}\},\ \{r^{\pm2}\},\ldots,\{r^{\pm(k-1)}\},
  \]
  以及
  \[
    \{sr^{2b}\mid b=1,\ldots,k\},\qquad
    \{sr^{2b-1}\mid b=1,\ldots,k\}.
  \]
  写出 $D_{2n}$ 的类方程。
  \item 当 $n=2k+1$ 为奇数时，证明 $D_{2n}$ 中的共轭类为
  \[
    \{1\},\ \{r^{\pm1}\},\ \{r^{\pm2}\},\ldots,\{r^{\pm k}\},
    \ \{sr^b\mid b=1,\ldots,n\}.
  \]
  写出 $D_{2n}$ 的类方程。
  \item 本题给出 $S_n$ 中每个共轭类大小的公式。设 $\sigma$ 是 $S_n$ 中的置换，令 $m_1,m_2,\ldots,m_s$ 为 $\sigma$ 的循环型中出现的不同整数（包括 $1$-循环）。对每个 $i\in\{1,2,\ldots,s\}$，假设 $\sigma$ 有 $k_i$ 个长度为 $m_i$ 的循环（所以 $\sum_{i=1}^s k_im_i=n$）。证明 $\sigma$ 的共轭元素个数为
  \[
    \frac{n!}{(k_1!m_1^{k_1})(k_2!m_2^{k_2})\cdots(k_s!m_s^{k_s})}.
  \]
  ［参见第 1.3 节练习 6 和 7，其中在一些特殊情形下给出了这个公式。］
  \item 证明若 $p$ 是素数，$P$ 是 $S_p$ 中阶为 $p$ 的子群，则 $|N_{S_p}(P)|=p(p-1)$。［论证 $P$ 的每个共轭子群恰含有 $p-1$ 个 $p$-循环，并用 $p$-循环的个数公式计算 $N_{S_p}(P)$ 在 $S_p$ 中的指数。］
  \item 设 $p$ 是素数。求 $S_n$ 中阶为 $p$ 的元素的共轭类个数公式（使用取整函数）。
  \item 令 $\pi:G\to S_G$ 为 $G$ 通过左乘作用在自身上给出的左正则表示。对每个 $g\in G$，记置换 $\pi(g)$ 为 $\sigma_g$，所以对所有 $x\in G$ 有 $\sigma_g(x)=gx$。令 $\lambda:G\to S_G$ 为对应的 $G$ 在自身上的右作用给出的置换表示，对每个 $h\in G$，记置换 $\lambda(h)$ 为 $\rho_h$。于是 $\rho_h(x)=xh^{-1}$ 对所有 $x\in G$ 成立（$\lambda$ 称为 $G$ 的右正则表示）。
  \begin{enumerate}[label=(\alph*)]
    \item 证明对所有 $g,h\in G$，$\sigma_g$ 与 $\rho_h$ 交换。（因此，$S_G$ 中 $\pi(G)$ 的中心化子包含子群 $\lambda(G)$，而 $\lambda(G)\cong G$。）
    \item 证明 $\sigma_g=\rho_g$ 当且仅当 $g$ 是 $G$ 的中心中阶为 $1$ 或 $2$ 的元素。
    \item 证明 $\sigma_g=\rho_h$ 当且仅当 $g$ 与 $h$ 都在 $G$ 的中心中。推出
    \[
      \pi(G)\cap\lambda(G)=\pi(Z(G))=\lambda(Z(G)).
    \]
  \end{enumerate}
\end{enumerate}

\section{自同构}

\noindent\textbf{定义.} 设 $G$ 是群。从 $G$ 到自身上的同构称为 $G$ 的一个\emph{自同构}。$G$ 的所有自同构组成的集合记为 $\Aut(G)$。

把验证 $\Aut(G)$ 在自同构的复合下构成群这一简单事实留作练习；这个群称为 $G$ 的\emph{自同构群}（自同构之间可以复合，因为每个自同构的定义域和值域相同）。注意，群 $G$ 的自同构特别地是集合 $G$ 的置换，所以 $\Aut(G)$ 是 $S_G$ 的子群。

群 $G$ 的自同构的一个最重要例子来自 $G$ 中固定元素的共轭。下一个结果在略一般的背景中讨论这一点。

\medskip
\noindent\textbf{命题 13.} 设 $H$ 是群 $G$ 的正规子群。则 $G$ 通过共轭作用在 $H$ 上，且这些作用是 $H$ 的自同构。更具体地，$G$ 在 $H$ 上的共轭作用对每个 $g\in G$ 定义为
\[
  h\longmapsto ghg^{-1}\qquad(h\in H).
\]
对每个 $g\in G$，由 $g$ 给出的共轭是 $H$ 的一个自同构。此作用给出的置换表示是 $G$ 到 $\Aut(H)$ 中的一个同态，其核为 $C_G(H)$。特别地，$G/C_G(H)$ 同构于 $\Aut(H)$ 的一个子群。

\begin{proof}
（参见第 1.7 节练习 17。）令 $\varphi_g$ 表示由 $g$ 给出的共轭。注意，因为 $g$ 正规化 $H$，所以 $\varphi_g$ 把 $H$ 映到自身。由于我们已经看到共轭定义了一个作用，因而 $\varphi_1=1$（$H$ 上的恒等映射），且对所有 $a,b\in G$，有 $\varphi_a\circ\varphi_b=\varphi_{ab}$。因此每个 $\varphi_g$ 给出从 $H$ 到自身的双射，因为它有双侧逆 $\varphi_{g^{-1}}$。每个 $\varphi_g$ 是从 $H$ 到 $H$ 的同态，因为对所有 $h,k\in H$，
\[
  \varphi_g(hk)=g(hk)g^{-1}=gh(gg^{-1})kg^{-1}
  =(ghg^{-1})(gkg^{-1})=\varphi_g(h)\varphi_g(k).
\]
这证明了由 $G$ 的任一固定元素给出的共轭定义了 $H$ 的一个自同构。

由前面的说明，置换表示 $\psi:G\to S_H$ 由 $\psi(g)=\varphi_g$ 定义（我们已经证明它是同态），其像包含于 $S_H$ 的子群 $\Aut(H)$。最后，
\begin{align*}
  \ker\psi&=\{g\in G\mid \varphi_g=\id\}\\
  &=\{g\in G\mid ghg^{-1}=h\text{ 对所有 }h\in H\}\\
  &=C_G(H).
\end{align*}
第一同构定理给出命题的最后断言。
\end{proof}

命题 13 表明，一个群通过共轭作用在正规子群上时，是以保持结构的置换作用，即以自同构作用。特别地，这个作用必须把子群送到子群，把阶为 $n$ 的元素送到阶为 $n$ 的元素，等等。下面两个推论描述该命题的两个具体应用。

\medskip
\noindent\textbf{推论 14.} 若 $K$ 是群 $G$ 的任意子群，且 $g\in G$，则 $K\cong gKg^{-1}$。共轭元素与共轭子群有相同阶。

\begin{proof}
在命题中令 $H=G$，可知由 $g\in G$ 给出的共轭是 $G$ 的自同构，推论随之成立。
\end{proof}

\medskip
\noindent\textbf{推论 15.} 对群 $G$ 的任意子群 $H$，商群 $N_G(H)/C_G(H)$ 同构于 $\Aut(H)$ 的一个子群。特别地，$G/Z(G)$ 同构于 $\Aut(G)$ 的一个子群。

\begin{proof}
由于 $H$ 是群 $N_G(H)$ 的正规子群，命题 13（令 $N_G(H)$ 扮演 $G$ 的角色）给出第一断言。第二断言是 $H=G$ 的特殊情形，此时 $N_G(G)=G$ 且 $C_G(G)=Z(G)$。
\end{proof}

\noindent\textbf{定义.} 设 $G$ 是群，$g\in G$。由 $g$ 给出的共轭称为 $G$ 的一个\emph{内自同构}；$\Aut(G)$ 中由所有内自同构组成的子群记为 $\Inn(G)$。

注意，$G$ 的内自同构集合确实是 $\Aut(G)$ 的一个子群，并且由推论 15，$\Inn(G)\cong G/Z(G)$。还要注意，若 $H$ 是 $G$ 的正规子群，则 $G$ 的元素给出的共轭限制到 $H$ 上是 $H$ 的自同构，但未必是 $H$ 的内自同构（我们将会看到）。

\paragraph{例}
\begin{enumerate}[label=(\arabic*)]
  \item 群 $G$ 是阿贝尔群，当且仅当每个内自同构都是平凡的。若 $H$ 是 $G$ 的阿贝尔正规子群且 $H\not\subseteq Z(G)$，则存在某个 $g\in G$，使得由 $g$ 给出的共轭限制到 $H$ 上不是 $H$ 的内自同构。一个显式例子是：$G=A_4$，$H$ 是 $G$ 中的 Klein 四元群，而 $g$ 是任一 $3$-循环。
  \item 由于 $Z(Q_8)=\gen{-1}$，有 $\Inn(Q_8)\cong V_4$。
  \item 由于 $Z(D_8)=\gen{r^2}$，有 $\Inn(D_8)\cong V_4$。
  \item 由于对所有 $n\geq3$，$Z(S_n)=1$，有 $\Inn(S_n)\cong S_n$。
\end{enumerate}

推论 15 表明，关于群 $G$ 的子群 $H$ 的自同构群的任何信息，都会转化为关于 $N_G(H)/C_G(H)$ 的信息。例如，若 $H\cong\ZZ/2\ZZ$，则由于 $H$ 有唯一的阶为 $1$ 与阶为 $2$ 的元素，推论 14 迫使 $\Aut(H)=1$。因此若 $H\cong\ZZ/2\ZZ$，则 $N_G(H)=C_G(H)$；若另外 $H\normal G$，则 $H\leq Z(G)$（参见第 2.2 节练习 10）。

虽然前面的例子相当简单，但它说明：$G$ 通过共轭作用在正规子群 $H$ 上的方式会受到 $H$ 的自同构群的限制。反过来，这可用于研究 $G$ 的结构，并在第 5.5 节讨论半直积时引出一些分类定理。

下面引入一个在后续章节中会自然用到的概念。

\noindent\textbf{定义.} 群 $G$ 的子群 $H$ 称为 $G$ 的\emph{特征子群}，记作 $H\charac G$，若 $G$ 的每个自同构都把 $H$ 映到自身，即对所有 $\sigma\in\Aut(G)$，有 $\sigma(H)=H$。

我们以后会用到下列关于特征子群的结果（证明留作练习）：
\begin{enumerate}[label=(\arabic*)]
  \item 特征子群是正规子群；
  \item 若 $H$ 是 $G$ 中给定阶数的唯一子群，则 $H$ 是 $G$ 的特征子群；
  \item 若 $K\charac H$ 且 $H\normal G$，则 $K\normal G$（因此尽管“正规性”不具有传递性，即正规子群的正规子群未必正规，正规子群的特征子群却正规）。
\end{enumerate}
因此，可以把特征子群看作“强正规”子群。例如，性质 (2) 与定理 2.7 表明循环群的每个子群都是特征子群。

本节最后给出一些关于具体群的自同构群的结果。

\medskip
\noindent\textbf{命题 16.} 阶为 $n$ 的循环群的自同构群同构于 $(\ZZ/n\ZZ)^\times$，它是阶为 $\varphi(n)$ 的阿贝尔群（其中 $\varphi$ 为 Euler 函数）。

\begin{proof}
令 $x$ 是循环群 $\ZZ_n$ 的一个生成元。若 $\psi\in\Aut(\ZZ_n)$，则 $\psi(x)=x^a$ 对某个 $a\in\ZZ$ 成立，并且整数 $a$ 唯一决定 $\psi$。把这个自同构记为 $\psi_a$。照常，由于 $|x|=n$，整数 $a$ 只在模 $n$ 意义下确定。由于 $\psi_a$ 是自同构，$x$ 与 $x^a$ 必须有相同阶，所以 $(a,n)=1$。进一步地，对每个与 $n$ 互素的 $a$，映射 $x\mapsto x^a$ 是 $\ZZ_n$ 的一个自同构。因此有满射
\[
  \Psi:\Aut(\ZZ_n)\to(\ZZ/n\ZZ)^\times,
  \qquad \psi_a\mapsto a\pmod n.
\]
映射 $\Psi$ 是同态，因为对所有 $\psi_a,
\psi_b\in\Aut(\ZZ_n)$，有
\[
  \psi_a\circ\psi_b(x)=\psi_a(x^b)=(x^b)^a=x^{ab}=\psi_{ab}(x),
\]
所以
\[
  \Psi(\psi_a\circ\psi_b)=\Psi(\psi_{ab})=ab\pmod n=\Psi(\psi_a)\Psi(\psi_b).
\]
最后，$\Psi$ 显然单射，故为同构。
\end{proof}

关于 $\Aut(\ZZ_n)$ 同构型的完整描述见第 9.5 节末尾。

\paragraph{例}
假设 $G$ 是阶为 $pq$ 的群，其中 $p,q$ 是素数（不一定不同）且 $p\leq q$。若 $p\nmid q-1$，我们证明 $G$ 是阿贝尔群。

若 $Z(G)\neq1$，则 Lagrange 定理迫使 $G/Z(G)$ 为循环群，于是由第 3.1 节练习 36，$G$ 是阿贝尔群。因此可假设 $Z(G)=1$。

若 $G$ 的每个非单位元素阶都为 $p$，则每个非单位元素的中心化子指数为 $q$，所以 $G$ 的类方程形如
\[
  pq=1+kq.
\]
这不可能，因为 $q$ 整除 $pq$ 与 $kq$，但不整除 $1$。因此 $G$ 含有一个阶为 $q$ 的元素 $x$。

令 $H=\gen{x}$。由于 $H$ 的指数为 $p$，且 $p$ 是整除 $|G|$ 的最小素数，推论 5 表明 $H\normal G$。由于 $Z(G)=1$，必有 $C_G(H)=H$。于是由推论 15，
\[
  G/H=N_G(H)/C_G(H)
\]
是一个阶为 $p$ 的群，且同构于 $\Aut(H)$ 的一个子群。但由命题 16，$\Aut(H)$ 的阶为 $\varphi(q)=q-1$，这由 Lagrange 定理会推出 $p\mid q-1$，与假设矛盾。因此 $G$ 必为阿贝尔群。

可以验证，每个阶为 $pq$ 的群（$p,q$ 为不同素数，$p<q$ 且 $p\nmid q-1$）都是循环群（见练习）。这是复合阶群中第一次出现唯一同构类型的情形。例如，每个阶为 $15$ 的群都是循环群。

下一个命题总结一些已知群的自同构群结果；它们将在后面证明。命题的第 (3) 部分说明了向量空间理论如何进入群论。

\medskip
\noindent\textbf{命题 17.}
\begin{enumerate}[label=(\arabic*)]
  \item 若 $p$ 是奇素数且 $n\in\ZZ^+$，则阶为 $p$ 的循环群的自同构群是阶为 $p-1$ 的循环群。更一般地，阶为 $p^n$ 的循环群的自同构群是阶为 $p^{n-1}(p-1)$ 的循环群（参见第 9.5 节推论 20）。
  \item 对所有 $n\geq3$，阶为 $2^n$ 的循环群的自同构群同构于 $\ZZ/2\ZZ\times\ZZ/2^{n-2}\ZZ$；特别地，它不是循环群，但有指数为 $2$ 的循环子群（参见第 9.5 节推论 20）。
  \item 令 $p$ 是素数，令 $V$ 是一个阿贝尔群（用加法记号），并满足对所有 $v\in V$ 有 $pv=0$。若 $|V|=p^n$，则 $V$ 是域 $\FF_p=\ZZ/p\ZZ$ 上的 $n$ 维向量空间。$V$ 的自同构正是从 $V$ 到自身的非奇异线性变换，即
  \[
    \Aut(V)\cong\GL(V)\cong\GL_n(\FF_p).
  \]
  特别地，$\Aut(V)$ 的阶如第 1.4 节所给（参见第 10.2 节与第 11.1 节中的例）。
  \item 对所有 $n\neq6$，有 $\Aut(S_n)=\Inn(S_n)\cong S_n$（参见练习 18）。对 $n=6$，有 $|\Aut(S_6):\Inn(S_6)|=2$（参见下面的练习 19，也参见第 6.3 节练习 10）。
  \item $\Aut(D_8)\cong D_8$ 且 $\Aut(Q_8)\cong S_4$（参见下面练习 4、5，也参见第 6.3 节练习 9）。
\end{enumerate}

命题第 (3) 部分描述的群 $V$ 称为阶为 $p^n$ 的\emph{初等阿贝尔群}（第 5 章中我们将看到，在同构意义下它由 $p$ 与 $n$ 唯一确定）。Klein 四元群 $V_4$ 是阶为 $4$ 的初等阿贝尔群。该命题断言
\[
  \Aut(V_4)\cong\GL_2(\FF_2).
\]
由第 1.4 节练习，后者阶为 $6$。而 $\Aut(V_4)$ 置换 $V_4$ 的三个非单位元素；$\Aut(V_4)$ 在 $V_4-\{1\}$ 上的这个作用给出 $\Aut(V_4)$ 到 $S_3$ 中的单置换表示。由阶数考虑，该同态满射，所以
\[
  \Aut(V_4)\cong\GL_2(\FF_2)\cong S_3.
\]
注意 $V_4$ 是阿贝尔群，所以 $\Inn(V_4)=1$。

对任意素数 $p$，阶为 $p^2$ 的初等阿贝尔群为 $\ZZ/p\ZZ\times\ZZ/p\ZZ$。其自同构群 $\GL_2(\FF_p)$ 阶为 $p(p-1)^2(p+1)$。因此推论 9 表明，对素数 $p$，若 $|P|=p^2$，则
\[
  |\Aut(P)|=p(p-1)\quad\text{或}\quad p(p-1)^2(p+1),
\]
分别对应于 $P$ 为循环群或初等阿贝尔群。

\paragraph{例}
假设 $G$ 是阶为 $45=3^2\cdot5$ 的群，并且有一个阶为 $3^2$ 的正规子群 $P$。我们说明 $G$ 必为阿贝尔群。

由推论 15，商群 $G/C_G(P)$ 同构于 $\Aut(P)$ 的一个子群；而由前一段，$\Aut(P)$ 的阶为 $6$ 或 $48$（分别对应 $P$ 为循环群或初等阿贝尔群）。另一方面，由于 $P$ 的阶是素数的平方，$P$ 是阿贝尔群，所以 $P\leq C_G(P)$。因此 $|C_G(P)|$ 被 $9$ 整除，从而 $|G/C_G(P)|$ 为 $1$ 或 $5$。结合这些事实可得 $|G/C_G(P)|=1$，即 $C_G(P)=G$，且 $P\leq Z(G)$。于是 $G/Z(G)$ 是循环群，所以 $G$ 必为阿贝尔群。

\subsection*{练习}
设 $G$ 是群。
\begin{enumerate}[label=\arabic*.]
  \item 若 $\alpha\in\Aut(G)$ 且 $\varphi_g$ 表示由 $g$ 给出的共轭，证明 $\alpha\varphi_g\alpha^{-1}=\varphi_{\alpha(g)}$。推出 $\Inn(G)\normal\Aut(G)$。（群 $\Aut(G)/\Inn(G)$ 称为 $G$ 的外自同构群。）
  \item 证明若 $G$ 是阶为 $pq$ 的阿贝尔群，其中 $p,q$ 是不同素数，则 $G$ 是循环群。［用 Cauchy 定理产生阶为 $p$ 与 $q$ 的元素，并考虑它们乘积的阶。］
  \item 证明在 $D_8$ 的任一自同构下，$r$ 至多有 $2$ 个可能像，$s$ 至多有 $4$ 个可能像（$r,s$ 是通常生成元，参见第 1.2 节）。推出 $|\Aut(D_8)|\leq8$。
  \item 用与前一题类似的论证证明 $|\Aut(Q_8)|\leq24$。
  \item 利用 $D_8\normal D_{16}$ 的事实证明 $\Aut(D_8)\cong D_8$。
  \item 证明特征子群正规。给出正规但非特征的子群例子。
  \item 若 $H$ 是群 $G$ 中给定阶数的唯一子群，证明 $H$ 是 $G$ 的特征子群。
  \item 设 $G$ 是群，且 $H,K$ 是 $G$ 的子群，满足 $H\leq K$。
  \begin{enumerate}[label=(\alph*)]
    \item 证明若 $H\charac K$ 且 $K\normal G$，则 $H\normal G$。
    \item 证明若 $H\charac K$ 且 $K\charac G$，则 $H\charac G$。用此证明 Klein 四元群 $V_4$ 是 $S_4$ 的特征子群。
    \item 给出例子说明：若 $H\normal K$ 且 $K\charac G$，则 $H$ 未必在 $G$ 中正规。
  \end{enumerate}
  \item 若 $r,s$ 是二面体群 $D_{2n}$ 的通常生成元，利用前两题推出 $\gen{r}$ 的每个子群都在 $D_{2n}$ 中正规。
  \item 设 $G$ 是群，$A$ 是 $G$ 的阿贝尔正规子群，并写 $\overline G=G/A$。证明 $\overline G$ 通过共轭左作用在 $A$ 上：$\overline g\cdot a=gag^{-1}$，其中 $g$ 是陪集 $\overline g$ 的任一代表元（特别要证明该作用良定）。给出一个显式例子说明若 $A$ 非阿贝尔，则此作用未必良定。
  \item 若 $p$ 是素数且 $P$ 是 $S_p$ 中阶为 $p$ 的子群，证明 $N_{S_p}(P)/C_{S_p}(P)$ 同构于 $\Aut(P)$ 的子群。［用第 3 节练习 34。］
  \item 设 $G$ 是阶为 $3825$ 的群。证明若 $H$ 是 $G$ 中阶为 $17$ 的正规子群，则 $H\leq Z(G)$。
  \item 设 $G$ 是阶为 $203$ 的群。证明若 $H$ 是 $G$ 中阶为 $7$ 的正规子群，则 $H\leq Z(G)$。推出此时 $G$ 是阿贝尔群。
  \item 设 $G$ 是阶为 $1575$ 的群。证明若 $H$ 是 $G$ 中阶为 $9$ 的正规子群，则 $H\leq Z(G)$。
  \item 证明下列乘法群都是循环群：$(\ZZ/5\ZZ)^\times$、$(\ZZ/9\ZZ)^\times$ 与 $(\ZZ/18\ZZ)^\times$。
  \item 证明 $(\ZZ/24\ZZ)^\times$ 是阶为 $8$ 的初等阿贝尔群。（稍后我们会看到，$(\ZZ/n\ZZ)^\times$ 是初等阿贝尔群当且仅当 $n\mid24$。）
  \item 设 $G=\gen{x}$ 是阶为 $n$ 的循环群。对 $n=2,3,4,5,6$，显式写出 $\Aut(G)$ 的元素（由上面的命题 16，$\Aut(G)\cong(\ZZ/n\ZZ)^\times$；所以对 $(\ZZ/n\ZZ)^\times$ 的每个元素 $a$，显式写出自同构 $\psi_a$ 如何作用在 $G$ 的元素 $\{1,x,x^2,\ldots,x^{n-1}\}$ 上）。
  \item 本题说明当 $n\neq6$ 时，$S_n$ 的每个自同构都是内自同构。固定整数 $n\geq2$ 且 $n\neq6$。
  \begin{enumerate}[label=(\alph*)]
    \item 证明群 $G$ 的自同构群置换 $G$ 的共轭类；也就是说，对每个 $\alpha\in\Aut(G)$ 与 $G$ 的每个共轭类 $K$，集合 $\alpha(K)$ 也是 $G$ 的一个共轭类。
    \item 令 $K$ 为 $S_n$ 中所有换位组成的共轭类，令 $K'$ 为 $S_n$ 中任何不是换位的阶为 $2$ 的元素的共轭类。证明 $|K|\neq |K'|$。推出 $S_n$ 的任一自同构都把换位送到换位。［参见第 3 节练习 33。］
    \item 证明对每个 $\alpha\in\Aut(S_n)$，有
    \[
      \alpha:(1\ 2)\mapsto(a\ b_2),\quad
      \alpha:(1\ 3)\mapsto(a\ b_3),\quad\ldots,\quad
      \alpha:(1\ n)\mapsto(a\ b_n),
    \]
    其中 $a,b_2,b_3,\ldots,b_n$ 是 $\{1,2,\ldots,n\}$ 中互不相同的整数。
    \item 证明 $(1\ 2),(1\ 3),\ldots,(1\ n)$ 生成 $S_n$，并推出 $S_n$ 的任一自同构都由它在这些元素上的作用唯一确定。用 (c) 说明 $S_n$ 至多有 $n!$ 个自同构，并结论 $\Aut(S_n)=\Inn(S_n)$（$n\neq6$）。
  \end{enumerate}
  \item 本题说明 $|\Aut(S_6):\Inn(S_6)|\leq2$（第 6.3 节练习 10 通过给出一个非内的 $S_6$ 自同构证明等号成立）。
  \begin{enumerate}[label=(\alph*)]
    \item 令 $K$ 为 $S_6$ 中所有换位组成的共轭类，令 $K'$ 为 $S_6$ 中任何不是换位的阶为 $2$ 的元素的共轭类。证明除非 $K'$ 是三个不交换位的乘积组成的共轭类，否则 $|K|\neq|K'|$。推出 $\Aut(S_6)$ 有一个指数至多为 $2$ 的子群，它把换位送到换位。
    \item 证明 $|\Aut(S_6):\Inn(S_6)|\leq2$。［沿着前一题 (c)、(d) 的步骤，说明任何把换位送到换位的自同构都是内自同构。］
  \end{enumerate}
\end{enumerate}

下一个练习引入一个子群 $J(P)$。与 $P$ 的中心一样，它可对任意有限群 $P$ 定义（尽管大多数应用中 $P$ 是一个素数幂阶群）。这个子群由 J. Thompson 于 1964 年定义，现在在有限群研究中，特别是在有限单群分类中，起着关键作用。

\begin{enumerate}[label=\arabic*.,resume]
  \item 对任意有限群 $P$，令 $d(P)$ 为 $P$ 的最少生成元个数（例如，$d(P)=1$ 当且仅当 $P$ 是非平凡循环群，而 $d(Q_8)=2$）。当 $A$ 遍历 $P$ 的所有阿贝尔子群时，令 $m(P)$ 为整数 $d(A)$ 的最大值（例如 $m(Q_8)=1$，$m(D_8)=2$）。定义
  \[
    J(P)=\gen{A\mid A\text{ 是 }P\text{ 的阿贝尔子群且 }d(A)=m(P)}.
  \]
  （$J(P)$ 称为 $P$ 的 Thompson 子群。）
  \begin{enumerate}[label=(\alph*)]
    \item 证明 $J(P)$ 是 $P$ 的特征子群。
    \item 对下列每个群 $P$，列出 $P$ 中所有满足 $d(A)=m(P)$ 的阿贝尔子群 $A$：$Q_8$、$D_8$、$D_{16}$ 和 $QD_{16}$（其中 $QD_{16}$ 是第 2.5 节练习 11 中定义的阶为 $16$ 的拟二面体群）。［用第 2.5 节中这些群的子群格。］
    \item 证明 $J(Q_8)=Q_8$、$J(D_8)=D_8$、$J(D_{16})=D_{16}$，并且 $J(QD_{16})$ 是 $QD_{16}$ 中一个阶为 $8$ 的二面体子群。
    \item 证明若 $Q\leq P$ 且 $J(P)$ 是 $Q$ 的子群，则 $J(P)=J(Q)$。推出：若 $P$ 是有限群 $G$ 的一个子群（不必正规），并且 $J(P)$ 包含在 $P$ 的某个子群 $Q$ 中，且 $Q\normal G$，则 $J(P)\normal G$。
  \end{enumerate}
\end{enumerate}

\section{Sylow 定理}

本节证明 Lagrange 定理的一个部分逆命题，并推出许多后果，其中一些会在下一章导向分类定理。

\noindent\textbf{定义.} 设 $G$ 是群，$p$ 是素数。
\begin{enumerate}[label=(\arabic*)]
  \item 阶为 $p^a$（其中 $a\geq1$）的群称为一个 $p$\emph{-群}。$G$ 的 $p$-群子群称为 $p$\emph{-子群}。
  \item 若 $G$ 是阶为 $p^am$ 的群，其中 $p\nmid m$，则阶为 $p^a$ 的子群称为 $G$ 的一个 \emph{Sylow $p$-子群}。
  \item $G$ 的 Sylow $p$-子群集合记为 $\Syl_p(G)$；Sylow $p$-子群的个数记为 $n_p(G)$（当 $G$ 由上下文清楚时，也简写为 $n_p$）。
\end{enumerate}

\medskip
\noindent\textbf{定理 18（Sylow 定理）.} 设 $G$ 是阶为 $p^am$ 的群，其中 $p$ 是素数且 $p\nmid m$。
\begin{enumerate}[label=(\arabic*)]
  \item $G$ 的 Sylow $p$-子群存在，即 $\Syl_p(G)\neq\varnothing$。
  \item 若 $P$ 是 $G$ 的 Sylow $p$-子群，$Q$ 是 $G$ 的任意 $p$-子群，则存在 $g\in G$ 使得 $Q\leq gPg^{-1}$，即 $Q$ 包含于 $P$ 的某个共轭子群中。特别地，$G$ 的任意两个 Sylow $p$-子群在 $G$ 中共轭。
  \item Sylow $p$-子群的个数形如 $1+kp$，即
  \[
    n_p\equiv1\pmod p.
  \]
  进一步，对任意 Sylow $p$-子群 $P$，$n_p$ 等于其正规化子 $N_G(P)$ 在 $G$ 中的指数，因此 $n_p$ 整除 $m$。
\end{enumerate}

先证明如下引理。

\medskip
\noindent\textbf{引理 19.} 设 $P\in\Syl_p(G)$。若 $Q$ 是 $G$ 的任意 $p$-子群，则
\[
  Q\cap N_G(P)=Q\cap P.
\]

\begin{proof}
令 $H=N_G(P)\cap Q$。由于 $P\leq N_G(P)$，显然 $P\cap Q\leq H$，故只需证明反向包含。又由于按定义 $H\leq Q$，这等价于证明 $H\leq P$。我们通过说明 $PH$ 是 $G$ 的一个同时包含 $P$ 与 $H$ 的 $p$-子群来证明这一点；但 $P$ 是 $G$ 中可能阶数最大的 $p$-子群，所以必有 $PH=P$，即 $H\leq P$。

由于 $H\leq N_G(P)$，由第 3.2 节推论 15，$PH$ 是子群。由同一节命题 13，
\[
  |PH|=\frac{|P|\,|H|}{|P\cap H|}.
\]
上式商中的所有数都是 $p$ 的幂，所以 $PH$ 是一个 $p$-群。而且 $P$ 是 $PH$ 的子群，所以 $|PH|$ 被 $p^a$ 整除，其中 $p^a$ 是整除 $|G|$ 的最大 $p$ 的幂。这两个事实迫使 $|PH|=p^a=|P|$。于是 $P=PH$，且 $H\leq P$。引理得证。
\end{proof}

\noindent\textbf{Sylow 定理的证明.}
(1) 对 $|G|$ 作归纳。若 $|G|=1$，没有什么要证明的。归纳假设所有阶小于 $|G|$ 的群都有 Sylow $p$-子群。

若 $p$ 整除 $|Z(G)|$，则由阿贝尔群的 Cauchy 定理（第 3.4 节命题 21），$Z(G)$ 有一个阶为 $p$ 的子群 $N$。令 $\overline G=G/N$，则 $|\overline G|=p^{a-1}m$。由归纳假设，$\overline G$ 有一个阶为 $p^{a-1}$ 的子群 $\overline P$。若令 $P$ 为 $G$ 中包含 $N$ 且满足 $P/N=\overline P$ 的子群，则
\[
  |P|=|P/N|\,|N|=p^a,
\]
所以 $P$ 是 $G$ 的 Sylow $p$-子群。因此可化归到 $p$ 不整除 $|Z(G)|$ 的情形。

令 $g_1,g_2,\ldots,g_r$ 是 $G$ 的不同非中心共轭类的代表元。$G$ 的类方程为
\[
  |G|=|Z(G)|+\sum_{i=1}^r |G:C_G(g_i)|.
\]
若对所有 $i$，都有 $p\mid |G:C_G(g_i)|$，则由于 $p\mid |G|$，也会有 $p\mid |Z(G)|$，矛盾。因此对某个 $i$，$p$ 不整除 $|G:C_G(g_i)|$。对此 $i$，令 $H=C_G(g_i)$，则
\[
  |H|=p^ak,
  \qquad p\nmid k.
\]
由于 $g_i\notin Z(G)$，有 $|H|<|G|$。由归纳假设，$H$ 有一个 Sylow $p$-子群 $P$；它当然也是 $G$ 的子群。由于 $|P|=p^a$，$P$ 是 $G$ 的 Sylow $p$-子群。归纳完成，(1) 得证。

在证明 (2) 与 (3) 前先作一些计算。由 (1)，$G$ 存在一个 Sylow $p$-子群 $P$。令
\[
  \{P_1,P_2,\ldots,P_r\}=S
\]
为 $P$ 的所有共轭子群组成的集合（即 $S=\{gPg^{-1}\mid g\in G\}$），并令 $Q$ 是 $G$ 的任意 $p$-子群。由 $S$ 的定义，$G$，因而 $Q$，通过共轭作用在 $S$ 上。把 $S$ 写成此 $Q$ 作用下轨道的不交并：
\[
  S=O_1\cup O_2\cup\cdots\cup O_s,
\]
其中 $r=|O_1|+\cdots+|O_s|$。注意 $r$ 不依赖于 $Q$，但 $Q$-轨道的个数 $s$ 依赖于 $Q$（按定义，$G$ 在 $S$ 上只有一个轨道，但 $G$ 的子群 $Q$ 可能有多个轨道）。必要时重新编号 $S$ 的元素，使得 $S$ 的前 $s$ 个元素为 $Q$-轨道的代表元：$P_i\in O_i$，$1\leq i\leq s$。由命题 2，
\[
  |O_i|=|Q:N_Q(P_i)|.
\]
按定义，$N_Q(P_i)=N_G(P_i)\cap Q$；由引理 19，$N_G(P_i)\cap Q=P_i\cap Q$。合并这两个事实得
\[
  |O_i|=|Q:P_i\cap Q|,
  \qquad 1\leq i\leq s. \tag{4.1}
\]

现在证明 $r\equiv1\pmod p$。由于 $Q$ 是任意的，可在上面取 $Q=P_1$，于是 (4.1) 给出
\[
  |O_1|=|P_1:P_1\cap P_1|=1.
\]
对所有 $i>1$，有 $P_1\neq P_i$，所以 $P_1\cap P_i<P_1$。由 (4.1)，
\[
  |O_i|=|P_1:P_1\cap P_i|>1,
  \qquad 2\leq i\leq s.
\]
由于 $P_1$ 是 $p$-群，$|P_1:P_1\cap P_i|$ 必为 $p$ 的幂，所以
\[
  p\mid |O_i|,
  \qquad 2\leq i\leq s.
\]
因此
\[
  r=|O_1|+(|O_2|+\cdots+|O_s|)\equiv1\pmod p.
\]

现在证明 (2) 与 (3)。令 $Q$ 是 $G$ 的任意 $p$-子群。假设 $Q$ 不包含于任何 $P_i$（即对任何 $g\in G$，$Q\nleq gPg^{-1}$）。此时对所有 $i$，$Q\cap P_i<Q$，所以由 (4.1)，
\[
  |O_i|=|Q:Q\cap P_i|>1.
\]
因此对所有 $i$，$p\mid |O_i|$，所以 $p$ 整除 $|O_1|+\cdots+|O_s|=r$。这与 $r\equiv1\pmod p$ 矛盾（记住 $r$ 不依赖于 $Q$ 的选择）。此矛盾证明 $Q\leq gPg^{-1}$ 对某个 $g\in G$ 成立。

为说明所有 Sylow $p$-子群在 $G$ 中共轭，令 $Q$ 是 $G$ 的任一 Sylow $p$-子群。由前一段，$Q\leq gPg^{-1}$ 对某个 $g\in G$ 成立。由于 $|gPg^{-1}|=|Q|=p^a$，必有 $gPg^{-1}=Q$。这证明了定理的 (2)。特别地，$S=\Syl_p(G)$，因为每个 Sylow $p$-子群都与 $P$ 共轭，所以 $n_p=r\equiv1\pmod p$，这是 (3) 的第一部分。

最后，由于所有 Sylow $p$-子群都共轭，命题 6 说明
\[
  n_p=|G:N_G(P)|\qquad(P\in\Syl_p(G)),
\]
从而完成 Sylow 定理的证明。

注意，Sylow 定理中的共轭部分结合推论 14 表明，同一个素数 $p$ 的任意两个 Sylow $p$-子群同构。

\medskip
\noindent\textbf{推论 20.} 设 $P$ 是 $G$ 的 Sylow $p$-子群。则下列条件等价：
\begin{enumerate}[label=(\arabic*)]
  \item $P$ 是 $G$ 的唯一 Sylow $p$-子群，即 $n_p=1$；
  \item $P\normal G$；
  \item $P\charac G$；
  \item 由 $p$-幂阶元素生成的所有子群都是 $p$-群；也就是说，若 $X$ 是 $G$ 的任意子集，且对所有 $x\in X$，$|x|$ 都是 $p$ 的幂，则 $\gen{X}$ 是 $p$-群。
\end{enumerate}

\begin{proof}
若 (1) 成立，则对所有 $g\in G$，由于 $gPg^{-1}\in\Syl_p(G)$，有 $gPg^{-1}=P$，即 $P\normal G$。故 (1) 推出 (2)。反过来，若 $P\normal G$ 且 $Q\in\Syl_p(G)$，则由 Sylow 定理，存在 $g\in G$ 使得 $Q=gPg^{-1}=P$。所以 $\Syl_p(G)=\{P\}$，即 (2) 推出 (1)。

由于特征子群正规，(3) 推出 (2)。反过来，若 $P\normal G$，刚才已经证明 $P$ 是 $G$ 的唯一阶为 $p^a$ 的子群，所以 $P\charac G$。因此 (2) 与 (3) 等价。

最后，假设 (1) 成立，并设 $X$ 是 $G$ 的子集，且每个 $x\in X$ 的阶都是 $p$ 的幂。由 Sylow 定理的共轭部分，对每个 $x\in X$，存在某个 $g\in G$ 使得 $x\in gPg^{-1}=P$。因此 $X\subseteq P$，所以 $\gen{X}\leq P$，且 $\gen{X}$ 是 $p$-群。反过来，若 (4) 成立，令 $X$ 为 $G$ 中所有 Sylow $p$-子群的并。若 $P$ 是任一 Sylow $p$-子群，则 $P$ 是 $p$-群 $\gen{X}$ 的子群。由于 $P$ 是 $G$ 的极大阶 $p$-子群，必有 $P=\gen{X}$，于是 (1) 成立。
\end{proof}

\paragraph{例}
设 $G$ 是有限群，$p$ 是素数。
\begin{enumerate}[label=(\arabic*)]
  \item 若 $p$ 不整除 $G$ 的阶，则 $G$ 的 Sylow $p$-子群是平凡群（且 Sylow 定理的所有部分都平凡成立）。若 $|G|=p^a$，则 $G$ 是自身唯一的 Sylow $p$-子群。
  \item 有限阿贝尔群对每个素数 $p$ 都有唯一的 Sylow $p$-子群。这个子群由所有阶为 $p$ 的幂的元素 $x$ 组成，有时称为该阿贝尔群的 $p$-\emph{主分量}。
  \item $S_3$ 有三个 Sylow $2$-子群：$\gen{(1\ 2)}$、$\gen{(2\ 3)}$ 与 $\gen{(1\ 3)}$。它有唯一（因而正规）的 Sylow $3$-子群 $\gen{(1\ 2\ 3)}=A_3$。注意 $3\equiv1\pmod2$。
  \item $A_4$ 有唯一的 Sylow $2$-子群：
  \[
    \{1,(1\ 2)(3\ 4),(1\ 3)(2\ 4),(1\ 4)(2\ 3)\}\cong V_4.
  \]
  它有四个 Sylow $3$-子群：$\gen{(1\ 2\ 3)}$、$\gen{(1\ 2\ 4)}$、$\gen{(1\ 3\ 4)}$ 与 $\gen{(2\ 3\ 4)}$。注意 $4\equiv1\pmod3$。
  \item $S_4$ 有 $n_2=3$ 与 $n_3=4$。由于 $S_4$ 含有一个同构于 $D_8$ 的子群，$S_4$ 的每个 Sylow $2$-子群都同构于 $D_8$。
\end{enumerate}

\paragraph{Sylow 定理的应用}
下面给出 Sylow 定理的一些应用。大多数例子利用 Sylow 定理证明某一特定阶的群不是单群。在讨论从较小群构造较大群的方法（例如形成半直积）之后，我们将能够用这些结果分类某些具体阶 $n$ 的群（正如我们已经对 $n=15$ 所做的）。

由于 Sylow 定理保证了有限群中 $p$-子群的存在，进一步研究素数幂阶群是有价值的。这将在第 6 章进行；那里还会讨论 Sylow 定理的更多应用。

对小阶群而言，Sylow 定理的同余条件常常足以迫使正规子群存在。数值应用 Sylow 定理的第一步，是把群阶分解为素数幂。群阶的最大素因子通常给出 $n_p$ 的最少可能值（例如，对 $n_2$ 的同余条件完全不给出限制），从而限制群 $G$ 的结构。在下面例子中，我们会看到一些情形：Sylow 定理本身并不直接迫使正规子群存在，但一些额外论证（通常涉及对几个不同素数 $p$ 的阶为 $p$ 的元素进行研究）能够证明正规 Sylow 子群存在。

\paragraph{例：阶为 $pq$ 的群，其中 $p,q$ 为素数且 $p<q$}
设 $|G|=pq$，其中 $p,q$ 是素数且 $p<q$。令 $P\in\Syl_p(G)$、$Q\in\Syl_q(G)$。我们证明 $Q\normal G$，并且若 $P$ 也正规，则 $G$ 是循环群。

现在三个条件：$n_q=1+kq$（$k\geq0$）、$n_q\mid p$ 以及 $p<q$，共同迫使 $k=0$。所以 $n_q=1$，即 $Q\normal G$。

由于 $n_p$ 整除素数 $q$，所以只有 $n_p=1$ 或 $n_p=q$ 两种可能。特别地，若 $p\nmid q-1$（即 $q\not\equiv1\pmod p$），则 $n_p$ 不可能等于 $q$，所以 $P\normal G$。

令 $P=\gen{x}$、$Q=\gen{y}$。若 $P\normal G$，则由于 $G/C_G(P)$ 同构于 $\Aut(\ZZ/p\ZZ)$ 的一个子群，而后者阶为 $p-1$，Lagrange 定理以及 $p,q$ 都不整除 $p-1$ 的观察说明 $G=C_G(P)$。此时 $x\in P\leq Z(G)$，所以 $x$ 与 $y$ 交换。（也可直接由第 3.1 节练习 42 得到。）这意味着 $|xy|=pq$（参见第 2.3 节练习），所以此时 $G$ 是循环群：$G\cong\ZZ/pq\ZZ$。

若 $p\mid q-1$，我们将在第 5 章看到存在唯一一个阶为 $pq$ 的非阿贝尔群（其中必有 $n_p=q$）。现在可以证明这个群的存在。令 $Q$ 是 $S_q$ 的一个 Sylow $q$-子群。由第 3 节练习 34，$|N_{S_q}(Q)|=q(q-1)$。由假设 $p\mid q-1$，Cauchy 定理给出 $N_{S_q}(Q)$ 有一个阶为 $p$ 的子群 $P$。由第 3.2 节推论 15，$PQ$ 是一个阶为 $pq$ 的群。由于 $C_{S_q}(Q)=Q$（第 3 节例 2），$PQ$ 是非阿贝尔群。证明 $PQ$ 唯一性的关键是关于 $\Aut(\ZZ/q\ZZ)$ 循环性的定理 17。

\paragraph{例：阶为 $30$ 的群}
令 $G$ 是阶为 $30$ 的群。我们证明 $G$ 有一个同构于 $\ZZ/15\ZZ$ 的正规子群。下一章中我们将用这一信息分类阶为 $30$ 的群。注意，任一阶为 $15$ 的子群必正规（因为指数为 $2$）且循环（由前一结果），所以只需证明存在阶为 $15$ 的子群。最快的方法是引用第 2 节练习 13。这里给出一个替代论证，以说明如何把 Sylow 定理与素数阶元素计数结合起来产生正规子群。

令 $P\in\Syl_5(G)$，$Q\in\Syl_3(G)$。若 $P$ 或 $Q$ 中有一个在 $G$ 中正规，则由第 3 章推论 15，$PQ$ 是阶为 $15$ 的群。还要注意，若 $P$ 或 $Q$ 中有一个正规，则 $P$ 与 $Q$ 都是 $PQ$ 的特征子群；由于 $PQ\normal G$，所以 $P$ 与 $Q$ 都在 $G$ 中正规（第 4 节练习 8(a)）。于是可假设两个 Sylow 子群都不正规。由 Sylow 定理第 3 部分，唯一可能为 $n_5=6$ 与 $n_3=10$。每个阶为 $5$ 的元素都属于某个 Sylow $5$-子群；每个 Sylow $5$-子群含有 $4$ 个非单位元素；并且由 Lagrange 定理，不同 Sylow $5$-子群只在单位元处相交。因此 $G$ 中阶为 $5$ 的元素个数等于一个 Sylow $5$-子群中的非单位元素个数乘以 Sylow $5$-子群个数，即 $4\cdot6=24$。类似地，阶为 $3$ 的元素个数为 $2\cdot10=20$。这荒谬，因为阶为 $30$ 的群不可能含有 $24+20=44$ 个不同元素。故 $P$ 或 $Q$ 中至少有一个正规（因而二者都正规）。

这种计数技巧经常有用（也参见第 6.2 节），并且当 Sylow $p$-子群阶为 $p$ 时尤其有效，因为此时两个不同 Sylow $p$-子群的交只能是单位元。若 Sylow $p$-子群阶为 $p^a$ 且 $a\geq2$，计数元素时需更加小心，因为这时不同 Sylow $p$-子群可能有更多公共元素，即交可能非平凡。

\paragraph{例：阶为 $12$ 的群}
令 $G$ 是阶为 $12$ 的群。我们证明：或者 $G$ 有正规 Sylow $3$-子群，或者 $G\cong A_4$（在后一种情形中，$G$ 有正规 Sylow $2$-子群）。下一章中我们将用这一信息分类阶为 $12$ 的群。

假设 $n_3\neq1$，并取 $P\in\Syl_3(G)$。由于 $n_3\mid4$ 且 $n_3\equiv1\pmod3$，可得 $n_3=4$。不同 Sylow $3$-子群只在单位元相交，每个含有两个阶为 $3$ 的元素，所以 $G$ 含有 $2\cdot4=8$ 个阶为 $3$ 的元素。由于 $|G:N_G(P)|=n_3=4$，有 $N_G(P)=P$。现在 $G$ 通过共轭作用在其四个 Sylow $3$-子群上，所以这个作用给出一个置换表示
\[
  \varphi:G\to S_4
\]
（注意也可以通过左乘作用在 $P$ 的左陪集上，并使用定理 3）。该作用的核 $K$ 是 $G$ 中正规化所有 Sylow $3$-子群的子群。特别地，$K\leq N_G(P)=P$。由于按假设 $P$ 不在 $G$ 中正规，$K=1$，即 $\varphi$ 是单射，于是 $G\cong\varphi(G)$。

由于 $G$ 含有 $8$ 个阶为 $3$ 的元素，而 $S_4$ 中恰有 $8$ 个阶为 $3$ 的元素且全都包含在 $A_4$ 中，可知 $\varphi(G)$ 与 $A_4$ 的交含有阶至少为 $8$ 的子群。由于两个群阶都为 $12$，故 $\varphi(G)=A_4$，从而 $G\cong A_4$。

注意 $A_4$ 确实有 $4$ 个 Sylow $3$-子群（见推论 20 后例 4），所以这样的群 $G$ 确实存在。又令 $V$ 是 $A_4$ 的一个 Sylow $2$-子群。由于 $|V|=4$，它包含 $A_4$ 中剩下的所有元素。特别地，不可能再有另一个 Sylow $2$-子群。因此 $n_2(A_4)=1$，即 $V\normal A_4$（也可直接看出，因为 $V$ 由单位元与 $A_4$ 中三个互不相交换位乘积组成，即 $V$ 是共轭类的并）。

\paragraph{例：阶为 $p^2q$ 的群，其中 $p,q$ 是不同素数}
令 $G$ 是阶为 $p^2q$ 的群。我们证明 $G$ 有一个正规 Sylow 子群（对 $p$ 或 $q$）。下一章中我们将用这一信息分类此阶的一些群（参见第 5.5 节练习 8 到 12）。取 $P\in\Syl_p(G)$ 与 $Q\in\Syl_q(G)$。

先考虑 $p>q$ 的情形。由于 $n_p\mid q$ 且 $n_p\equiv1\pmod p$，必有 $n_p=1$。所以 $P\normal G$。

再考虑 $p<q$ 的情形。若 $n_q=1$，则 $Q\normal G$。否则假设 $n_q>1$，即 $n_q=1+tq$，其中 $t>0$。又 $n_q\mid p^2$，所以 $n_q=p$ 或 $p^2$。由于 $q>p$，不可能有 $n_q=p$，故 $n_q=p^2$。于是
\[
  tq=p^2-1=(p-1)(p+1).
\]
由于 $q$ 是素数，要么 $q\mid p-1$，要么 $q\mid p+1$。前者因 $q>p$ 不可能；故后者成立。由于 $q>p$ 且 $q\mid p+1$，必有 $q=p+1$。这迫使 $p=2$、$q=3$，于是 $|G|=12$。结论由前一例得出。

\paragraph{阶为 $60$ 的群}
我们说明，即使某些给定阶的群可能是单群，Sylow 定理仍可用于揭示该阶群的结构。注意在不同素数之间切换的技巧，以及利用阶小于 $60$ 的群结果来研究阶为 $60$ 的群的归纳过程。

\medskip
\noindent\textbf{命题 21.} 若 $|G|=60$ 且 $G$ 有不止一个 Sylow $5$-子群，则 $G$ 是单群。

\begin{proof}
反设 $|G|=60$、$n_5>1$，但存在正规子群 $H\normal G$，且 $H\neq1,G$。由 Sylow 定理，$n_5$ 唯一可能为 $6$。取 $P\in\Syl_5(G)$，则由于其指数为 $n_5$，有 $|N_G(P)|=10$。

若 $5\mid |H|$，则 $H$ 含有 $G$ 的一个 Sylow $5$-子群；由于 $H$ 正规，它含有该子群的所有 $6$ 个共轭子群。特别地，
\[
  |H|\geq1+6\cdot4=25,
\]
唯一可能是 $|H|=30$。这导致矛盾，因为前面的例子证明了每个阶为 $30$ 的群都有正规（因而唯一）的 Sylow $5$-子群。此论证表明，对 $G$ 的任何真正规子群 $H$，$5$ 都不整除 $|H|$。

若 $|H|=6$ 或 $12$，则 $H$ 有一个正规、因而特征的 Sylow 子群，于是它也在 $G$ 中正规。必要时用这个子群替换 $H$，可假设 $|H|=2,3$ 或 $4$。令 $\overline G=G/H$，则 $|\overline G|=30,20$ 或 $15$。在每种情形下，由前面的结果，$\overline G$ 都有一个阶为 $5$ 的正规子群 $\overline P$。若令 $H_1$ 为 $\overline P$ 在 $G$ 中的完全原像，则 $H_1\normal G$、$H_1\neq G$ 且 $5\mid |H_1|$。这与前一段矛盾，证明完毕。
\end{proof}

\medskip
\noindent\textbf{推论 22.} $A_5$ 是单群。

\begin{proof}
子群 $\gen{(1\ 2\ 3\ 4\ 5)}$ 与 $\gen{(1\ 3\ 2\ 4\ 5)}$ 是 $A_5$ 的不同 Sylow $5$-子群，所以结论立即由命题得出。
\end{proof}

下一个命题说明，阶为 $60$ 的单群在同构意义下唯一。

\medskip
\noindent\textbf{命题 23.} 若 $G$ 是阶为 $60$ 的单群，则 $G\cong A_5$。

\begin{proof}
设 $G$ 是阶为 $60$ 的单群，则 $n_2=3,5$ 或 $15$。取 $P\in\Syl_2(G)$，并令 $N=N_G(P)$，于是 $|G:N|=n_2$。

先观察 $G$ 没有指数小于 $5$ 的真子群 $H$。理由如下：若 $H$ 是 $G$ 中指数为 $4,3$ 或 $2$ 的子群，则由定理 3，$G$ 有一个包含在 $H$ 中的正规子群 $K$，且 $G/K$ 同构于 $S_4$、$S_3$ 或 $S_2$ 的子群。由于 $K\neq G$，单性迫使 $K=1$。但这是不可能的，因为 $60=|G|$ 不整除 $4!$。此论证特别说明 $n_2\neq3$。

若 $n_2=5$，则 $N$ 在 $G$ 中指数为 $5$，所以 $G$ 通过左乘作用在 $N$ 的左陪集集合上，给出一个到 $S_5$ 中的置换表示。如上，该表示的核是真正规子群；由于 $G$ 单，其核为 $1$，故 $G$ 同构于 $S_5$ 的一个子群。把 $G$ 与这个同构副本认同，可假设 $G\leq S_5$。若 $G\not\subseteq A_5$，则 $S_5=GA_5$，由第二同构定理，$A_5\cap G$ 在 $G$ 中指数为 $2$。由于 $G$ 没有（正规）指数为 $2$ 的子群，这是矛盾。故 $G\leq A_5$。由于 $|G|=|A_5|$，$S_5$ 中的 $G$ 的同构副本恰为 $A_5$，如欲证。

最后假设 $n_2=15$。若 $G$ 的每一对不同 Sylow $2$-子群 $P,Q$ 都满足 $P\cap Q=1$，则 $G$ 中出现在 Sylow $2$-子群中的非单位元素数目为 $(4-1)\cdot15=45$。但 $n_5=6$，所以 $G$ 中阶为 $5$ 的元素数目为 $(5-1)\cdot6=24$，合计 $69$ 个元素，矛盾。因此存在不同 Sylow $2$-子群 $P,Q$，使得 $|P\cap Q|=2$。令 $M=N_G(P\cap Q)$。由于 $P$ 与 $Q$ 是阶为 $4$ 的群，故阿贝尔；所以 $P$ 与 $Q$ 都是 $M$ 的子群。又由于 $G$ 单，$M\neq G$。于是 $4$ 整除 $|M|$ 且 $|M|>4$（否则 $P=M=Q$）。唯一可能为 $|M|=12$，即 $M$ 在 $G$ 中指数为 $5$（回忆 $M$ 不能有指数 $3$ 或 $1$）。现在把前一段论证中的 $N$ 换成 $M$，得到 $G\cong A_5$。但这在本情形下导致矛盾，因为 $n_2(A_5)=5$（参见练习）。证明完毕。
\end{proof}

\subsection*{练习}
设 $G$ 是有限群，$p$ 是素数。
\begin{enumerate}[label=\arabic*.]
  \item 证明若 $P\in\Syl_p(G)$，且 $H$ 是包含 $P$ 的 $G$ 的子群，则 $P\in\Syl_p(H)$。给出例子说明一般而言，$G$ 的子群的 Sylow $p$-子群未必是 $G$ 的 Sylow $p$-子群。
  \item 证明若 $H$ 是 $G$ 的子群，且 $Q\in\Syl_p(H)$，则对所有 $g\in G$，$gQg^{-1}\in\Syl_p(gHg^{-1})$。
  \item 用 Sylow 定理证明 Cauchy 定理。（注意，在 Sylow 定理证明中我们只使用了阿贝尔群的 Cauchy 定理——命题 3.21——所以这里的推理不循环。）
  \item 给出 $D_{12}$ 与 $S_3\times S_3$ 的所有 Sylow $2$-子群和 Sylow $3$-子群。
  \item 证明对每个奇素数 $p$，$D_{2n}$ 的 Sylow $p$-子群都是循环且正规的。
  \item 给出 $A_4$ 的所有 Sylow $3$-子群，以及 $S_4$ 的所有 Sylow $3$-子群。
  \item 给出 $S_4$ 的所有 Sylow $2$-子群，并找出 $S_4$ 中元素把其中一个共轭到其他每一个。
  \item 给出 $S_5$ 的两个不同 Sylow $2$-子群，以及 $S_5$ 中一个把其中一个共轭到另一个的元素。
  \item 给出 $\SL_2(\FF_3)$ 的所有 Sylow $3$-子群（参见第 2.1 节练习 9）。
  \item 证明 $\SL_2(\FF_3)$ 中由
  \[
    \begin{pmatrix}0&-1\\1&0\end{pmatrix}
    \quad\text{与}\quad
    \begin{pmatrix}1&1\\1&-1\end{pmatrix}
  \]
  生成的子群是 $\SL_2(\FF_3)$ 的唯一 Sylow $2$-子群（参见第 2.4 节练习 10）。
  \item 说明 $\SL_2(\FF_3)$ 的中心是由 $\pm I$ 组成的阶为 $2$ 的群，其中 $I$ 为单位矩阵。证明
  \[
    \SL_2(\FF_3)/Z(\SL_2(\FF_3))\cong A_4.
  \]
  ［用阶为 $12$ 的群的事实。］
  \item 令 $2n=2^ak$，其中 $k$ 为奇数。证明 $D_{2n}$ 的 Sylow $2$-子群个数为 $k$。［证明若 $P\in\Syl_2(D_{2n})$，则 $N_{D_{2n}}(P)=P$。］
  \item 证明阶为 $56$ 的群对某个整除其阶的素数 $p$ 有正规 Sylow $p$-子群。
  \item 证明阶为 $312$ 的群对某个整除其阶的素数 $p$ 有正规 Sylow $p$-子群。
  \item 证明阶为 $351$ 的群对某个整除其阶的素数 $p$ 有正规 Sylow $p$-子群。
  \item 设 $|G|=pqr$，其中 $p,q,r$ 是素数且 $p<q<r$。证明 $G$ 对 $p,q,r$ 中的某一个有正规 Sylow 子群。
  \item 证明若 $|G|=105$，则 $G$ 有正规 Sylow $5$-子群和正规 Sylow $7$-子群。
  \item 证明阶为 $200$ 的群有正规 Sylow $5$-子群。
  \item 证明若 $|G|=6545$，则 $G$ 不是单群。
  \item 证明若 $|G|=1365$，则 $G$ 不是单群。
  \item 证明若 $|G|=2907$，则 $G$ 不是单群。
  \item 证明若 $|G|=132$，则 $G$ 不是单群。
  \item 证明若 $|G|=462$，则 $G$ 不是单群。
  \item 证明若 $G$ 是阶为 $231$ 的群，则 $Z(G)$ 含有 $G$ 的一个 Sylow $11$-子群，且 Sylow $7$-子群在 $G$ 中正规。
  \item 证明若 $G$ 是阶为 $385$ 的群，则 $Z(G)$ 含有 $G$ 的一个 Sylow $7$-子群，且 Sylow $11$-子群在 $G$ 中正规。
  \item 令 $G$ 是阶为 $105$ 的群。证明若 $G$ 的 Sylow $3$-子群正规，则 $G$ 是阿贝尔群。
  \item 令 $G$ 是阶为 $315$ 的群，并且有正规 Sylow $3$-子群。证明 $Z(G)$ 含有 $G$ 的一个 Sylow $3$-子群，并推出 $G$ 是阿贝尔群。
  \item 令 $G$ 是阶为 $1575$ 的群。证明若 $G$ 的 Sylow $3$-子群正规，则其 Sylow $5$-子群与 Sylow $7$-子群正规。在这种情形下证明 $G$ 是阿贝尔群。
  \item 若 $G$ 是阶小于 $100$ 的非阿贝尔单群，证明 $G\cong A_5$。［排除除 $60$ 外的所有阶。］
  \item 阶为 $168$ 的单群必须有多少个阶为 $7$ 的元素？
  \item 对 $p=2,3,5$，求 $n_p(A_5)$ 与 $n_p(S_5)$。［注意 $A_4\leq A_5$。］
  \item 设 $P$ 是 $H$ 的 Sylow $p$-子群，且 $H$ 是 $K$ 的子群。若 $P\normal H$ 且 $H\normal K$，证明 $P\normal K$。推出若 $P\in\Syl_p(G)$ 且 $H=N_G(P)$，则 $N_G(H)=H$（换言之，Sylow $p$-子群的正规化子是自正规化的）。
  \item 设 $P$ 是 $G$ 的正规 Sylow $p$-子群，$H$ 是 $G$ 的任意子群。证明 $P\cap H$ 是 $H$ 的唯一 Sylow $p$-子群。
  \item 设 $P\in\Syl_p(G)$ 且 $N\normal G$。用 Sylow 定理的共轭部分证明 $P\cap N$ 是 $N$ 的 Sylow $p$-子群。推出 $PN/N$ 是 $G/N$ 的 Sylow $p$-子群（注意这也可由第二同构定理完成，参见第 3.3 节练习 9）。
  \item 设 $P\in\Syl_p(G)$ 且 $H\leq G$。证明对某个 $g\in G$，$gPg^{-1}\cap H$ 是 $H$ 的 Sylow $p$-子群。给出显式例子说明对任何 $h\in H$，$hPh^{-1}\cap H$ 未必是 $H$ 的 Sylow $p$-子群（特别地，不像 $H$ 正规时那样，本题第一部分不能总取 $g=1$）。
  \item 证明若 $N\normal G$，则 $n_p(G/N)\leq n_p(G)$。
  \item 令 $R$ 是 $G$ 的正规 $p$-子群（不必是 Sylow 子群）。
  \begin{enumerate}[label=(\alph*)]
    \item 证明 $R$ 包含在 $G$ 的每个 Sylow $p$-子群中。
    \item 若 $S$ 是 $G$ 的另一个正规 $p$-子群，证明 $RS$ 也是 $G$ 的正规 $p$-子群。
    \item 子群 $O_p(G)$ 定义为由 $G$ 的所有正规 $p$-子群生成的群。证明 $O_p(G)$ 是 $G$ 的唯一最大正规 $p$-子群，并且 $O_p(G)$ 等于 $G$ 的所有 Sylow $p$-子群的交。
    \item 令 $\overline G=G/O_p(G)$。证明 $O_p(\overline G)=1$（即 $\overline G$ 没有非平凡正规 $p$-子群）。
  \end{enumerate}
  \item 用 Sylow 定理证明中的方法说明：若 $n_p\not\equiv1\pmod{p^2}$，则存在 $G$ 的不同 Sylow $p$-子群 $P,Q$，使得
  \[
    |P:P\cap Q|=|Q:P\cap Q|=p.
  \]
  \item 证明 $\GL_n(\FF_p)$ 中严格上三角幺幂矩阵子群（参见第 2.1 节练习 17）是该有限群的 Sylow $p$-子群。［用第 1.4 节中的阶公式求 $\GL_n(\FF_p)$ 的 Sylow $p$-子群阶。］
  \item 证明 $\GL_2(\FF_p)$ 的 Sylow $p$-子群个数为 $p+1$。［给出两个不同的 Sylow $p$-子群。］
  \item 证明 $\SL_2(\FF_4)\cong A_5$（参见第 2.1 节练习 9 中 $\SL_2(\FF_4)$ 的定义）。
  \item 证明正二十面体在 $\RR^3$ 中的刚体运动群同构于 $A_5$。［回忆该群阶为 $60$：第 1.2 节练习 13。］
  \item 证明正十二面体在 $\RR^3$ 中的刚体运动群同构于 $A_5$。（与立方体和正四面体类似，正二十面体与正十二面体互为对偶立体。）［回忆该群阶为 $60$：第 1.2 节练习 12。］
  \item 设 $p$ 是整除有限群 $G$ 的阶的最小素数。若 $P\in\Syl_p(G)$ 且 $P$ 是循环群，证明 $N_G(P)=C_G(P)$。
  \item 设 $p$ 是奇素数。求 $S_{2p}$ 的 Sylow $p$-子群的生成元。说明它是阶为 $p^2$ 的阿贝尔群。
  \item 设 $p$ 是素数。求 $S_{p^2}$ 的 Sylow $p$-子群的生成元。说明它是阶为 $p^{p+1}$ 的非阿贝尔群。
  \item 编写并运行一个计算机程序：
  \begin{enumerate}[label=(\roman*)]
    \item 给出每个小于 $10000$ 的奇数 $n$，其中 $n$ 不是素数幂，且存在某个素因子 $p$，使得 Sylow 定理的同余条件并不迫使所有阶为 $n$ 的群都有 $n_p=1$；
    \item 对 (i) 中的每个 $n$，给出 $n$ 的素数幂分解，并给出整除 $n$ 的每个素数 $p$ 的所有允许的 $n_p$ 取值列表（即 Sylow 定理第 3 部分未排除的取值）。
  \end{enumerate}
  \item 对所有小于 $1000$ 的偶数，执行与前一题相同的过程。解释所得列表长度与所检验整数个数之间的相对关系。
  \item 证明若 $|G|=2^nm$，其中 $m$ 为奇数，且 $G$ 有循环 Sylow $2$-子群，则 $G$ 有一个阶为 $m$ 的正规子群。［用归纳法以及第 2 节练习 11、12。］
  \item 证明若 $U$ 与 $W$ 是 $G$ 的 Sylow $p$-子群 $P$ 的正规子集，则 $U$ 与 $W$ 在 $G$ 中共轭，当且仅当它们在 $N_G(P)$ 中共轭。推出 $P$ 的中心中两个元素在 $G$ 中共轭，当且仅当它们在 $N_G(P)$ 中共轭。（若 $N_P(U)=P$，则称 $P$ 的子集 $U$ 在 $P$ 中正规。）
  \item 设 $P$ 是 $G$ 的 Sylow $p$-子群，$M$ 是包含 $N_G(P)$ 的任意 $G$ 的子群。证明
  \[
    |G:M|\equiv1\pmod p.
  \]
\end{enumerate}

以下一系列练习导向所有整数 $n$ 的分类，这些 $n$ 具有性质：每个阶为 $n$ 的群都是循环群（例如 $n=15$ 就是这样的整数）。这些论证是“奇数阶群可解”定理证明的一个极大简化原型：它们利用真子群的结构（交换性）以及它们嵌入整个群的方式（我们将看到不同极大子群交于单位元），通过计数论证获得矛盾。在证明奇数阶群可解时，人们用归纳把问题化归到最小反例是单群的情形——但在那里每个真子群只是可解的（不像这里是阿贝尔的）。在这种情形下对极大子群结构及嵌入的分析复杂得多，而计数论证大致由特征标理论论证取代（将在第六部分讨论）。

\begin{enumerate}[label=\arabic*.,resume]
  \item 假设 $G$ 是有限单群，且每个真子群都是阿贝尔群。若 $M,N$ 是 $G$ 的两个不同极大子群，证明 $M\cap N=1$。［参见第 3 节练习 23。］
  \item 用前一题证明：若 $G$ 是任一非阿贝尔群且每个真子群都是阿贝尔群，则 $G$ 不是单群。［设 $G$ 是此断言的反例，并用第 3 节练习 24 说明 $G$ 有不止一个极大子群共轭类。仿照第 3 节练习 23 的方法，对位于 $M$ 与 $N$ 的所有共轭中的元素计数，其中 $M,N$ 是不共轭的极大子群；说明这会给出多于 $|G|$ 个元素。］
  \item 证明如下分类：若 $G$ 是阶为 $p_1p_2\cdots p_r$ 的有限群，其中 $p_i$ 是互不相同的素数，且对所有 $i,j$，$p_i$ 都不整除 $p_j-1$，则 $G$ 是循环群。［归纳地，每个真子群都是循环的，所以由前一题，$G$ 不是单群。若 $N$ 是非平凡真正规子群，则 $N$ 循环，且 $G/N$ 作为 $N$ 的自同构作用。用命题 16 证明 $N\leq Z(G)$，并用归纳证明 $G/Z(G)$ 循环，从而由第 3.1 节练习 36，$G$ 是阿贝尔群。］
  \item 证明前一题的逆命题：若整数 $n\geq2$ 满足每个阶为 $n$ 的群都是循环群，则 $n=p_1p_2\cdots p_r$ 是互不相同素数的乘积，且对所有 $i,j$，$p_i$ 不整除 $p_j-1$。［若 $n$ 不是这种形式，利用阶为 $p^2$ 的非循环群以及阶为 $pq$ 的非循环群（其中 $p\mid q-1$）的直积，构造阶为 $n$ 的非循环群。］
  \item 若 $G$ 是有限群且每个真子群都是阿贝尔群，证明 $G$ 可解。
\end{enumerate}

\section{$A_n$ 的单性}

$A_n$（$n\geq5$）的单性有若干证明。最初等的证明涉及说明 $A_n$ 由 $3$-循环生成；然后说明一个正规子群必含有一个 $3$-循环，因而必含有所有 $3$-循环，所以不能是真子群。这里给出一种计算较少的方法。

注意，$A_3$ 是阿贝尔单群，而 $A_4$ 不是单群（$n_2(A_4)=1$）。

\medskip
\noindent\textbf{定理 24.} 对所有 $n\geq5$，$A_n$ 是单群。

\begin{proof}
对 $n$ 作归纳。$n=5$ 的结果已经建立，所以假设 $n\geq6$，并令 $G=A_n$。假设存在 $H\normal G$，且 $H\neq1,G$。

对每个 $i\in\{1,2,\ldots,n\}$，令 $G_i$ 为 $G$ 在其自然作用在 $\{1,2,\ldots,n\}$ 上时 $i$ 的稳定子。因此 $G_i\leq G$ 且 $G_i\cong A_{n-1}$。由归纳假设，$G_i$ 对 $1\leq i\leq n$ 都是单群。

先假设存在某个 $\tau\in H$，满足 $\tau\neq1$ 但 $\tau(i)=i$ 对某个 $i\in\{1,2,\ldots,n\}$ 成立。由于 $\tau\in H\cap G_i$ 且 $H\cap G_i\normal G_i$，由 $G_i$ 的单性，必有 $H\cap G_i=G_i$，即
\[
  G_i\leq H.
\]
由第 1 节练习 2，$\sigma G_i\sigma^{-1}=G_{\sigma(i)}$，所以对所有 $i$，
\[
  \sigma G_i\sigma^{-1}\leq \sigma H\sigma^{-1}=H.
\]
于是
\[
  G_j\leq H,
  \qquad j=1,2,\ldots,n.
\]
任意 $\lambda\in A_n$ 可写成偶数个（设为 $2t$ 个）换位的乘积，即
\[
  \lambda=\lambda_1\lambda_2\cdots\lambda_t,
\]
其中每个 $\lambda_k$ 是两个换位的乘积。由于 $n>4$，每个 $\lambda_k\in G_j$ 对某个 $j$ 成立，因此
\[
  G=\gen{G_1,G_2,\ldots,G_n}\leq H,
\]
这与 $H$ 真子群矛盾。因此，若 $\tau\neq1$ 是 $H$ 的元素，则对所有 $i\in\{1,2,\ldots,n\}$，都有 $\tau(i)\neq i$；即 $H$ 的非单位元素都不固定 $\{1,2,\ldots,n\}$ 的任何元素。

于是若 $\tau_1,
\tau_2\in H$ 且
\[
  \tau_1(i)=\tau_2(i)\text{ 对某个 }i,
\]
则 $\tau_1=\tau_2$，因为此时 $\tau_2^{-1}\tau_1(i)=i$。换言之，
\[
  \tau_1(i)=\tau_2(i)\text{ 对某个 }i\quad\Longrightarrow\quad \tau_1=\tau_2. \tag{4.2}
\]

假设存在 $\tau\in H$，使得 $\tau$ 的循环分解中含有长度至少为 $3$ 的循环，例如
\[
  \tau=(a_1\ a_2\ a_3\ \cdots)(b_1\ b_2\ \cdots)\cdots.
\]
取 $\sigma\in G$，满足 $\sigma(a_1)=a_1$、$\sigma(a_2)=a_2$，但 $\sigma(a_3)\neq a_3$（注意由于 $n\geq5$，这样的 $\sigma\in A_n$ 存在）。由命题 10，
\[
  \tau_1=\sigma\tau\sigma^{-1}
  =(a_1\ a_2\ \sigma(a_3)\ \cdots)(\sigma(b_1)\ \sigma(b_2)\ \cdots)\cdots,
\]
所以 $\tau$ 与 $\tau_1$ 是 $H$ 中不同元素，但 $\tau(a_1)=\tau_1(a_1)=a_2$，与 (4.2) 矛盾。这证明 $H$ 的非单位元素的循环分解中只能出现 $2$-循环。

取 $\tau\in H$ 且 $\tau\neq1$，于是
\[
  \tau=(a_1\ a_2)(a_3\ a_4)(a_5\ a_6)\cdots
\]
（这里用到 $n\geq6$）。令 $\sigma=(a_1\ a_2)(a_3\ a_5)\in G$。则
\[
  \tau_1=\sigma\tau\sigma^{-1}
  =(a_1\ a_2)(a_5\ a_4)(a_3\ a_6)\cdots,
\]
所以 $\tau$ 与 $\tau_1$ 是 $H$ 中不同元素，但 $\tau(a_1)=\tau_1(a_1)=a_2$，再次与 (4.2) 矛盾。由此完成 $A_n$ 单性的证明。
\end{proof}

\subsection*{练习}
设 $G$ 是群，$\Omega$ 是无限集合。
\begin{enumerate}[label=\arabic*.]
  \item 证明对所有 $n\geq5$，$A_n$ 没有指数小于 $n$ 的真子群。
  \item 求对所有 $n\geq5$，$S_n$ 的所有正规子群。
  \item 证明对所有 $n\geq5$，$A_n$ 是 $S_n$ 中唯一的指数小于 $n$ 的真子群。
  \item 证明对每个 $n\geq3$，$A_n$ 由所有 $3$-循环组成的集合生成。
  \item 证明若存在子群链 $G_1\leq G_2\leq\cdots\leq G$，使得 $G=\bigcup_{i=1}^{\infty}G_i$，并且每个 $G_i$ 都是单群，则 $G$ 是单群。
  \item 令 $D$ 为 $S_\Omega$ 中由只移动 $\Omega$ 的有限多个元素的置换组成的子群（见第 3 节练习 17），并令 $A$ 为 $D$ 中所有在其移动的有限点集上作用为偶置换的元素组成的集合。证明 $A$ 是一个无限单群。［证明 $D$ 的任意两个元素都包含在 $D$ 的某个有限单子群中。］
  \item 沿用前一题记号，证明若 $H\normal S_\Omega$ 且 $H\neq1$，则 $A\leq H$；也就是说，$A$ 是 $S_\Omega$ 的唯一非平凡极小正规子群。
  \item 沿用前两题记号，证明 $|D|=|A|=|\Omega|$。推出
  \[
    S_\Omega\cong S_\Delta\quad\Longrightarrow\quad |\Omega|=|\Delta|.
  \]
  ［用 $D$ 由换位生成的事实。你可以假设基数为 $|\Omega|$ 的集合的可数并与有限直积仍具有基数 $|\Omega|$。］
\end{enumerate}

% \chapter{直积、半直积与阿贝尔群}

本章考察由较小的群构造较大群的两种较容易的方法，即直积与半直积。这使我们能够陈述有限生成阿贝尔群基本定理；特别地，该定理完全分类了所有有限阿贝尔群。

\section{直积}

我们从有限多个群以及可数多个群的直积的定义开始（任意多个群的直积将在练习中讨论）。

\noindent\textbf{定义.}
\begin{enumerate}[label=(\arabic*)]
  \item 设群 $G_1,G_2,\ldots,G_n$ 的运算分别为 $*_1,*_2,\ldots,*_n$。它们的\emph{直积}
  \[
    G_1\times G_2\times\cdots\times G_n
  \]
  是所有 $n$ 元组 $(g_1,g_2,\ldots,g_n)$ 组成的集合，其中 $g_i\in G_i$，其运算按分量定义：
  \[
    (g_1,g_2,\ldots,g_n)*(h_1,h_2,\ldots,h_n)
    =(g_1*_1h_1,g_2*_2h_2,\ldots,g_n*_nh_n).
  \]

  \item 类似地，群 $G_1,G_2,\ldots$ 的直积
  \[
    G_1\times G_2\times\cdots
  \]
  是所有序列 $(g_1,g_2,\ldots)$ 组成的集合，其中 $g_i\in G_i$，其运算也按分量定义：
  \[
    (g_1,g_2,\ldots)*(h_1,h_2,\ldots)=(g_1*_1h_1,g_2*_2h_2,\ldots).
  \]
\end{enumerate}

虽然直积各因子中的运算可以不同，但照例，我们把所有抽象群都写成乘法形式。因此，例如上面 (1) 中的运算可简写为
\[
  (g_1,g_2,\ldots,g_n)(h_1,h_2,\ldots,h_n)=(g_1h_1,g_2h_2,\ldots,g_nh_n).
\]

\paragraph{例}
\begin{enumerate}[label=(\arabic*)]
  \item 设 $G_i=\RR$（运算为加法），$i=1,2,\ldots,n$。则 $\RR\times\RR\times\cdots\times\RR$（$n$ 个因子）就是熟悉的 Euclid $n$ 维空间 $\RR^n$，其通常的向量加法为
  \[
    (a_1,a_2,\ldots,a_n)+(b_1,b_2,\ldots,b_n)
    =(a_1+b_1,a_2+b_2,\ldots,a_n+b_n).
  \]

  \item 为了说明构成直积的群（以及相应的运算）可以完全任意，令 $G_1=\ZZ$，$G_2=S_3$，$G_3=\GL_2(\RR)$，其群运算分别为加法、复合以及矩阵乘法。则 $G_1\times G_2\times G_3$ 中的运算由
  \[
    \left(n,\sigma,\begin{pmatrix}a&b\\ c&d\end{pmatrix}\right)
    \left(m,\tau,\begin{pmatrix}p&q\\ r&s\end{pmatrix}\right)
    =
    \left(n+m,\sigma\circ\tau,
    \begin{pmatrix}ap+br&aq+bs\\ cp+dr&cq+ds\end{pmatrix}\right)
  \]
  给出。
\end{enumerate}

\medskip
\noindent\textbf{命题 1.} 若 $G_1,\ldots,G_n$ 是群，则它们的直积是一个群，其阶为
\[
  |G_1||G_2|\cdots |G_n|,
\]
若某个 $G_i$ 是无限群，则直积也是无限群。

\begin{proof}
令 $G=G_1\times G_2\times\cdots\times G_n$。由于 $G$ 的运算按分量定义，而每个因子 $G_i$ 中都满足相同的群公理，所以验证 $G$ 的群公理是直接的。例如，结合律可如下验证。设
\[
  (a_1,a_2,\ldots,a_n),\quad (b_1,b_2,\ldots,b_n),\quad (c_1,c_2,\ldots,c_n)\in G.
\]
则
\begin{align*}
(a_1,\ldots,a_n)[(b_1,\ldots,b_n)(c_1,\ldots,c_n)]
&=(a_1,\ldots,a_n)(b_1c_1,\ldots,b_nc_n)\\
&=(a_1(b_1c_1),\ldots,a_n(b_nc_n))\\
&=((a_1b_1)c_1,\ldots,(a_nb_n)c_n)\\
&=[(a_1,\ldots,a_n)(b_1,\ldots,b_n)](c_1,\ldots,c_n),
\end{align*}
其中第三步使用了每个分量中的结合律。剩下的群公理验证类似：$G$ 的单位元是 $n$ 元组 $(1_1,1_2,\ldots,1_n)$，其中 $1_i$ 是 $G_i$ 的单位元；元素 $(g_1,g_2,\ldots,g_n)$ 的逆元是 $(g_1^{-1},g_2^{-1},\ldots,g_n^{-1})$，其中 $g_i^{-1}$ 是 $g_i$ 在 $G_i$ 中的逆元。$G$ 的阶的公式显然成立。
\end{proof}

若重新排列直积的因子，所得直积同构于原直积（参见练习 7）。

下一个命题表明，直积 $G_1\times G_2\times\cdots\times G_n$ 含有每个 $G_i$ 的一个同构拷贝。可以把这些特定的拷贝看作直积的“坐标轴”；在 $\RR\times\RR$ 的情形中，它们正是 $x$ 轴和 $y$ 轴。不过要注意，不应把这些“坐标轴”看作直积中唯一的 $G_i$ 的拷贝。例如，在 $\RR\times\RR$ 中，过原点的任意直线都是一个同构于 $\RR$ 的子群（而 $\RR\times\RR$ 有无穷多对可作为坐标轴的直线，也就是任一坐标系旋转后的坐标轴）。命题的第二部分表明，存在投影到各个分量上的同态。

\medskip
\noindent\textbf{命题 2.} 设 $G_1,G_2,\ldots,G_n$ 是群，令 $G=G_1\times\cdots\times G_n$。
\begin{enumerate}[label=(\arabic*)]
  \item 对每个固定的 $i$，$G$ 中所有在第 $j$ 个位置（$j\neq i$）为 $G_j$ 的单位元、而在第 $i$ 个位置任取 $G_i$ 元素的元素构成 $G$ 的一个子群，并同构于 $G_i$：
  \[
    G_i\cong \{(1,1,\ldots,1,g_i,1,\ldots,1)\mid g_i\in G_i\},
  \]
  其中 $g_i$ 出现在第 $i$ 个位置。若把 $G_i$ 与这个子群等同起来，则 $G_i\normal G$ 且
  \[
    G/G_i\cong G_1\times\cdots\times G_{i-1}\times G_{i+1}\times\cdots\times G_n.
  \]

  \item 对每个固定的 $i$，定义
  \[
    \pi_i:G\to G_i,\qquad \pi_i((g_1,g_2,\ldots,g_n))=g_i.
  \]
  则 $\pi_i$ 是满同态，并且
  \[
    \ker\pi_i=\{(g_1,\ldots,g_{i-1},1,g_{i+1},\ldots,g_n)\mid g_j\in G_j,\ j\neq i\}
  \]
  同构于 $G_1\times\cdots\times G_{i-1}\times G_{i+1}\times\cdots\times G_n$。

  \item 在 (1) 中的等同之下，若 $x\in G_i$ 且 $y\in G_j$，其中 $i\neq j$，则 $xy=yx$。
\end{enumerate}

\begin{proof}
(1) 因为 $G$ 中的运算按分量定义，由子群判别法立即得到
\[
\{(1,1,\ldots,1,g_i,1,\ldots,1)\mid g_i\in G_i\}
\]
是 $G$ 的子群。此外，映射 $g_i\mapsto (1,1,\ldots,1,g_i,1,\ldots,1)$ 是 $G_i$ 到该子群的同构。以下把 $G_i$ 与它在 $G$ 中的这个同构拷贝等同。

为证明 (1) 的其余部分，考虑映射
\[
  \varphi:G\to G_1\times\cdots\times G_{i-1}\times G_{i+1}\times\cdots\times G_n
\]
定义为
\[
  \varphi(g_1,g_2,\ldots,g_n)=(g_1,\ldots,g_{i-1},g_{i+1},\ldots,g_n),
\]
也就是删去第 $i$ 个分量。映射 $\varphi$ 是同态，因为
\begin{align*}
\varphi((g_1,
\ldots,g_n)(h_1,\ldots,h_n))
&=\varphi(g_1h_1,\ldots,g_nh_n)\\
&=(g_1h_1,\ldots,g_{i-1}h_{i-1},g_{i+1}h_{i+1},\ldots,g_nh_n)\\
&=(g_1,\ldots,g_{i-1},g_{i+1},\ldots,g_n)(h_1,\ldots,h_{i-1},h_{i+1},\ldots,h_n)\\
&=\varphi(g_1,
\ldots,g_n)\varphi(h_1,
\ldots,h_n).
\end{align*}
由于每个位置上的元素可以任意选取，$\varphi$ 是满射。并且
\[
  \ker\varphi=\{(g_1,
\ldots,g_n)
\mid g_j=1\text{ 对所有 }j\neq i\}=G_i.
\]
这证明 $G_i$ 是 $G$ 的正规子群（特别地也再次证明了这个 $G_i$ 的拷贝是子群），而第一同构定理给出 (1) 的最后断言。

(2) 中 $\pi_i$ 是满同态且核如上所述的论证与 (1) 很相似，细节留给读者。

(3) 中，若
\[
  x=(1,
\ldots,1,g_i,1,
\ldots,1),
  \qquad
  y=(1,
\ldots,1,g_j,1,
\ldots,1),
\]
其中所示元素分别出现在第 $i,j$ 个位置，则（设 $i<j$，以便记号清楚）
\[
  xy=(1,
\ldots,1,g_i,1,
\ldots,1,g_j,1,
\ldots,1)=yx.
\]
证明完毕。
\end{proof}

命题 2 的一个推广见练习 2。以后我们将继续把命题 (1) 中的“坐标轴”子群与它们的同构拷贝 $G_i$ 等同。第 $i$ 个这样的子群通常称为 $G$ 的第 $i$ 个\emph{分量}或第 $i$ 个\emph{因子}。例如，在 $\ZZ_n\times\ZZ_m$ 中计算时，可以令 $x$ 为第一因子的生成元，令 $y$ 为第二因子的生成元，并把 $\ZZ_n\times\ZZ_m$ 的元素写成 $x^ay^b$ 的形式。这相当于用 $x$ 和 $y$ 代替形式有序对 $(x,1)$ 与 $(1,y)$（因此 $x^ay^b$ 代替 $(x^a,y^b)$）。

\paragraph{例}
\begin{enumerate}[label=(\arabic*)]
  \item 在命题 2 的记号下，由 (3) 可知若 $x_i\in G_i$，$1\leq i\leq n$，则对所有 $k\in\ZZ$，
  \[
    (x_1x_2\cdots x_n)^k=x_1^kx_2^k\cdots x_n^k.
  \]
  因而 $x_1x_2\cdots x_n$ 的阶是使得 $x_i^k=1$ 对所有 $i$ 成立的最小正整数 $k$。因此
  \[
    |x_1x_2\cdots x_n|=\operatorname{lcm}(|x_1|,|x_2|,\ldots,|x_n|),
  \]
  其中该阶为无限，当且仅当某个 $x_i$ 的阶为无限。

  \item 令 $p$ 为素数，$n\in\ZZ^+$，考虑
  \[
    E_{p^n}=\underbrace{\ZZ_p\times\ZZ_p\times\cdots\times\ZZ_p}_{n\text{ 个因子}}.
  \]
  则 $E_{p^n}$ 是阶为 $p^n$ 的阿贝尔群，并且对所有 $x\in E_{p^n}$ 都有 $x^p=1$。这个群就是第 4.4 节中描述的阶为 $p^n$ 的初等阿贝尔群。

  \item 设 $p$ 为素数。我们说明阶为 $p^2$ 的初等阿贝尔群恰有 $p+1$ 个阶为 $p$ 的子群（特别地，不止两个显然的子群）。令 $E=E_{p^2}$。由于 $E$ 的每个非单位元素都有阶 $p$，每个这样的元素都生成 $E$ 的一个阶为 $p$ 的循环子群。由 Lagrange 定理，不同的阶为 $p$ 的子群交平凡。因此，$E$ 的 $p^2-1$ 个非单位元素被分割为大小为 $p-1$ 的子集（每个这样的子集正是某个阶为 $p$ 的子群的非单位元素集合）。于是阶为 $p$ 的子群个数为
  \[
    \frac{p^2-1}{p-1}=p+1.
  \]
  当 $p=2$ 时，$E$ 是 Klein 四元群；我们已经看到它有 $3$ 个阶为 $2$ 的子群（也参见练习 10 和 11）。
\end{enumerate}

\subsection*{练习}
\begin{enumerate}[label=\arabic*.]
  \item 证明直积的中心是中心的直积：
  \[
    Z(G_1\times G_2\times\cdots\times G_n)=Z(G_1)\times Z(G_2)\times\cdots\times Z(G_n).
  \]
  推出：群的直积是阿贝尔群当且仅当每个因子都是阿贝尔群。

  \item 设 $G_1,G_2,\ldots,G_n$ 是群，令 $G=G_1\times\cdots\times G_n$。令 $I$ 是 $\{1,\ldots,n\}$ 的一个真非空子集，并令 $J=\{1,\ldots,n\}-I$。定义 $G_I$ 为 $G$ 中在所有 $j\in J$ 的位置上都是 $G_j$ 的单位元的元素集合。
  \begin{enumerate}[label=(\alph*)]
    \item 证明 $G_I$ 同构于群 $G_i$（$i\in I$）的直积。
    \item 证明 $G_I$ 是 $G$ 的正规子群，且 $G/G_I\cong G_J$。
    \item 证明 $G\cong G_I\times G_J$。
  \end{enumerate}

  \item 在前一题记号下，令 $I,K$ 是 $\{1,2,\ldots,n\}$ 的任意两个不交非空子集，并令 $G_I,G_K$ 为上面定义的子群。证明对所有 $x\in G_I$ 与所有 $y\in G_K$，都有 $xy=yx$。

  \item 设 $A,B$ 是有限群，$p$ 是素数。证明 $A\times B$ 的任意 Sylow $p$-子群都形如 $P\times Q$，其中 $P\in\Syl_p(A)$ 且 $Q\in\Syl_p(B)$。证明
  \[
    n_p(A\times B)=n_p(A)n_p(B).
  \]
  把这两个结果推广到有限多个有限群的直积。

  \item 给出 $Q_8\times \ZZ_4$ 的一个非正规子群（注意每个因子的每个子群都是正规的）。

  \item 证明 $Q_8\times E_{2^n}$ 的所有子群都是正规的。

  \item 设 $G_1,G_2,
\ldots,G_n$ 是群，令 $\pi\in S_n$。证明映射
  \[
    \varphi_\pi:G_1\times G_2\times\cdots\times G_n
    \longrightarrow G_{\pi^{-1}(1)}\times G_{\pi^{-1}(2)}\times\cdots\times G_{\pi^{-1}(n)}
  \]
  定义为
  \[
    \varphi_\pi(g_1,g_2,\ldots,g_n)=(g_{\pi^{-1}(1)},g_{\pi^{-1}(2)},\ldots,g_{\pi^{-1}(n)})
  \]
  是同构（因此改变直积因子的顺序不改变同构类型）。

  \item 设 $G_1=G_2=\cdots=G_n$，令 $G=G_1\times\cdots\times G_n$。在前题记号下，证明 $\varphi_\pi\in\Aut(G)$。并证明映射 $\pi\mapsto \varphi_\pi$ 是 $S_n$ 到 $\Aut(G)$ 的单同态。（特别地，$\varphi_{\pi_1}\circ\varphi_{\pi_2}=\varphi_{\pi_1\pi_2}$。在 $\varphi_\pi$ 的定义中需要 $\pi^{-1}$ 的原因正是在这里。根本原因是：若 $e_i$ 是第 $i$ 位为 $1$、其余为 $0$ 的 $n$ 元组，则 $S_n$ 按 $\pi\cdot e_i=e_{\pi(i)}$ 作用在 $\{e_1,
\ldots,e_n\}$ 上，这是一个左群作用。若 $n$ 元组 $(g_1,
\ldots,g_n)$ 表示为 $g_1e_1+\cdots+g_ne_n$，则这个左作用扩展为
  \[
    \pi\cdot(g_1e_1+
\cdots+g_ne_n)=g_1e_{\pi(1)}+\cdots+g_ne_{\pi(n)}.
  \]
  右边 $e_{\pi(i)}$ 的系数是 $g_i$，所以 $e_i$ 的系数是 $g_{\pi^{-1}(i)}$，即得到上述公式。若在定义中使用 $\pi$ 而非 $\pi^{-1}$，映射 $\pi\mapsto\varphi_\pi$ 不一定是同态；它对应于右群作用。）

  \item 令所有 $G_i$ 都为域 $F$，并用前一题证明：每行每列恰有一个 $1$ 的 $n\times n$ 矩阵集合是 $\GL_n(F)$ 的一个同构于 $S_n$ 的子群。这些矩阵称为置换矩阵，因为它们只是置换 $n$ 维向量空间 $F\times\cdots\times F$ 的标准基 $e_1,
\ldots,e_n$。

  \item 令 $p$ 为素数。设 $A$ 与 $B$ 是两个阶为 $p$ 的循环群，生成元分别为 $x$ 与 $y$。令 $E=A\times B$，于是 $E$ 是阶为 $p^2$ 的初等阿贝尔群 $E_{p^2}$。证明 $E$ 中不同的阶为 $p$ 的子群为
  \[
    \gen{x},\ \gen{xy},\ \gen{xy^2},\ldots,\gen{xy^{p-1}},\ \gen{y},
  \]
  共 $p+1$ 个。

  \item 令 $p$ 为素数，$n\in\ZZ^+$。求初等阿贝尔群 $E_{p^n}$ 中阶为 $p$ 的子群个数公式。

  \item 设 $A,B$ 是群。假设 $Z(A)$ 含有子群 $Z_1$，$Z(B)$ 含有子群 $Z_2$，且 $Z_1\cong Z_2$。设该同构由 $x_i\mapsto y_i$ 给出，$x_i\in Z_1$。$A$ 与 $B$ 的一个\emph{中心积}是商群
  \[
    (A\times B)/Z,
    \qquad Z=\{(x_i,y_i^{-1})\mid x_i\in Z_1\},
  \]
  记为 $A*B$。它不是唯一的，因为它依赖于 $Z_1,Z_2$ 以及它们之间的同构。（可把 $A*B$ 看作把直积 $A\times B$ 中每个 $x_i\in Z_1$ 与对应的 $y_i\in Z_2$ 识别起来后得到的群。）
  \begin{enumerate}[label=(\alph*)]
    \item 证明 $A$ 和 $B$ 在商群 $A*B$ 中的像分别同构于 $A$ 和 $B$，并且这两个像相交于一个同构于 $Z_1$ 的中心子群。求 $|A*B|$。
    \item 令 $\ZZ_4=\gen{x}$。令 $D_8=\gen{r,s}$ 与 $Q_8=\gen{i,j}$ 由通常的生成元与关系给出。令 $\ZZ_4*D_8$ 为把 $x^2$ 与 $r^2$ 识别起来的中心积，令 $\ZZ_4*Q_8$ 为把 $x^2$ 与 $-1$ 识别起来的中心积。证明 $\ZZ_4*D_8\cong \ZZ_4*Q_8$。
  \end{enumerate}

  \item 给出前一题中构造的群 $\ZZ_4*D_8$ 与 $\ZZ_4*Q_8$ 的展示。

  \item 令 $G=A_1\times A_2\times\cdots\times A_n$，且对每个 $i$，令 $B_i\normal A_i$。证明
  \[
    B_1\times B_2\times\cdots\times B_n\normal G
  \]
  且
  \[
    (A_1\times\cdots\times A_n)/(B_1\times\cdots\times B_n)
    \cong (A_1/B_1)\times\cdots\times(A_n/B_n).
  \]

  \item 下面的练习描述任意多个群的直积。设 $I$ 是任意非空指标集，对每个 $i\in I$ 设 $(G_i,*_i)$ 是一个群。群 $G_i$（$i\in I$）的直积是集合
  \[
    G=\prod_{i\in I}G_i
  \]
  （即 $G_i$ 的笛卡尔积），并按如下方式定义二元运算：若 $\prod a_i$ 与 $\prod b_i$ 是 $G$ 的元素，则
  \[
    \left(\prod_{i\in I}a_i\right)\left(\prod_{i\in I}b_i\right)
    =\prod_{i\in I}(a_i*_i b_i).
  \]
  \begin{enumerate}[label=(\alph*)]
    \item 证明这个二元运算良定义且满足结合律。
    \item 证明元素 $\prod 1_i$ 满足 $G$ 的单位元公理，其中 $1_i$ 是 $G_i$ 的单位元。
    \item 证明元素 $\prod a_i^{-1}$ 是 $\prod a_i$ 的逆元，其中每个分量 $a_i^{-1}$ 在群 $G_i$ 中取逆。
  \end{enumerate}
  因而直积是一个群。（注意，当 $I=\{1,2,
\ldots,n\}$ 时，这个直积定义与正文中的 $n$ 元组定义相同。）

  \item 陈述并证明命题 2 对任意直积的推广。

  \item 令 $I$ 是任意非空指标集，并令 $G_i$ 是群。群 $G_i$ 的\emph{限制直积}或\emph{直和}是直积中除有限多个分量外均为单位元的元素集合，也就是所有 $\prod a_i\in\prod_{i\in I}G_i$，其中除有限多个 $i\in I$ 外都有 $a_i=1_i$。
  \begin{enumerate}[label=(\alph*)]
    \item 证明限制直积是直积的子群。
    \item 证明限制直积是直积的正规子群。
    \item 令 $I=\ZZ^+$，令 $p_i$ 为第 $i$ 个整数素数。若对所有 $i\in\ZZ^+$ 有 $G_i=\ZZ/p_i\ZZ$，证明 $G_i$ 的限制直积中的每个元素都有有限阶，但 $\prod_{i\in\ZZ^+}G_i$ 中有无限阶元素。证明在这个例子中，限制直积是直积的挠子群（参见第 2.1 节练习 6）。
  \end{enumerate}

  \item 对 (a) 到 (e) 的每一项，给出满足指定性质的群的例子：
  \begin{enumerate}[label=(\alph*)]
    \item 一个无限群，其中每个元素的阶为 $1$ 或 $2$；
    \item 一个无限群，其中每个元素都有有限阶，但对每个正整数 $n$ 都有一个阶为 $n$ 的元素；
    \item 一个既有无限阶元素又有阶为 $2$ 的元素的群；
    \item 一个群 $G$，使得每个有限群都同构于 $G$ 的某个子群；
    \item 一个非平凡群 $G$，使得 $G\cong G\times G$。
  \end{enumerate}
\end{enumerate}

\section{有限生成阿贝尔群基本定理}

\noindent\textbf{定义.}
\begin{enumerate}[label=(\arabic*)]
  \item 群 $G$ 称为\emph{有限生成的}，若存在 $G$ 的有限子集 $A$ 使得 $G=\gen{A}$。
  \item 对每个 $r\in\ZZ$ 且 $r\geq 0$，令 $\ZZ^r=\ZZ\times\ZZ\times\cdots\times\ZZ$ 为 $r$ 个 $\ZZ$ 的直积，其中 $\ZZ^0=1$。群 $\ZZ^r$ 称为秩为 $r$ 的\emph{自由阿贝尔群}。
\end{enumerate}

注意，任意有限群 $G$ 当然是有限生成的，只需取 $A=G$ 作为生成集即可。此外，$\ZZ^r$ 由 $e_1,e_2,\ldots,e_r$ 有限生成，其中 $e_i$ 是第 $i$ 位为 $1$、其余为 $0$ 的 $r$ 元组。现在可以陈述（有限生成）阿贝尔群的基本分类定理。

\medskip
\noindent\textbf{定理 3（有限生成阿贝尔群基本定理）.} 设 $G$ 是有限生成阿贝尔群。则：
\begin{enumerate}[label=(\arabic*)]
  \item 存在整数 $r,n_1,n_2,\ldots,n_s$ 使得
  \[
    G\cong \ZZ^r\times \ZZ_{n_1}\times\ZZ_{n_2}\times\cdots\times\ZZ_{n_s},
  \]
  并满足
  \begin{enumerate}[label=(\alph*)]
    \item $r\geq 0$ 且 $n_j\geq 2$ 对所有 $j$ 成立；
    \item $n_{i+1}\mid n_i$，$1\leq i\leq s-1$。
  \end{enumerate}

  \item (1) 中的表达式是唯一的：若
  \[
    G\cong \ZZ^t\times \ZZ_{m_1}\times\ZZ_{m_2}\times\cdots\times\ZZ_{m_u},
  \]
  其中 $t,m_1,
\ldots,m_u$ 也满足 (a) 与 (b)，则 $t=r$，$u=s$，并且对所有 $i$ 有 $m_i=n_i$。
\end{enumerate}

\begin{proof}
我们将在第 12.1 节把这个定理作为一个更一般分类定理的推论推出。对有限群，我们会在第 6.1 节末尾给出另一个证明。
\end{proof}

\noindent\textbf{定义.} 定理 3 中的整数 $r$ 称为 $G$ 的\emph{自由秩}或 \emph{Betti 数}，整数 $n_1,n_2,\ldots,n_s$ 称为 $G$ 的\emph{不变因子}。定理 3(1) 中对 $G$ 的描述称为 $G$ 的\emph{不变因子分解}。

定理 3 断言，阿贝尔群的自由秩以及（有序的）不变因子列表都是唯一确定的，因此两个有限生成阿贝尔群同构当且仅当它们有相同的自由秩与相同的不变因子列表。注意，有限生成阿贝尔群是有限群当且仅当它的自由秩为零。

有限阿贝尔群的阶就是其不变因子的乘积（由命题 1）。若有限阿贝尔群 $G$ 的不变因子为 $n_1,n_2,
\ldots,n_s$，其中 $n_{i+1}\mid n_i$，$1\leq i\leq s-1$，则称 $G$ 是\emph{类型} $(n_1,n_2,
\ldots,n_s)$ 的群。

定理 3 给出了一种有效方法，用来列出给定阶的所有有限阿贝尔群。也就是说，为了找出所有阶为 $n$ 的阿贝尔群（在同构意义下），必须找出所有有限整数序列 $n_1,n_2,
\ldots,n_s$，使得
\begin{enumerate}[label=(\arabic*)]
  \item $n_j\geq 2$ 对所有 $j\in\{1,2,
\ldots,s\}$ 成立；
  \item $n_{i+1}\mid n_i$，$1\leq i\leq s-1$；
  \item $n_1n_2\cdots n_s=n$。
\end{enumerate}
定理 3 说明，这类序列与阶为 $n$ 的有限阿贝尔群的同构类集合之间存在一一对应。

在说明如何对具体的 $n$ 找出所有这类序列之前，先作一些一般性说明。首先注意 $n_1\geq n_2\geq\cdots\geq n_s$，所以 $n_1$ 是最大的不变因子。又由 (3)，每个 $n_i$ 都整除 $n$。若 $p$ 是 $n$ 的任一素因子，则由 (3) 可知 $p$ 必整除某个 $n_i$。于是由 (2)，$p$ 也整除所有 $j\leq i$ 的 $n_j$。因此
\[
  \text{每个 }n\text{ 的素因子都必须整除第一个不变因子 }n_1.
\]
特别地，若 $n$ 是若干互异素数之积（每个素数只出现一次），\footnote{这样的整数称为无平方因子的，因为它们不被任何大于 $1$ 的平方整除。}则 $n\mid n_1$，于是 $n=n_1$。这证明了：若 $n$ 是无平方因子的，则阶为 $n$ 的阿贝尔群只有一种可能的不变因子列表，即 $n_1=n$。

\medskip
\noindent\textbf{推论 4.} 若 $n$ 是互异素数之积，则在同构意义下，唯一的阶为 $n$ 的阿贝尔群是阶为 $n$ 的循环群 $\ZZ_n$。

把 $n$ 分解为素数幂是确定阶为 $n$ 的阿贝尔群所有可能不变因子列表的第一步。

\paragraph{例}
设 $n=180=2^2\cdot 3^2\cdot 5$。如上所述，必须有 $2\cdot3\cdot5\mid n_1$，所以 $n_1$ 的可能值为
\[
  n_1=2^2\cdot3^2\cdot5,
  \quad 2^2\cdot3\cdot5,
  \quad 2\cdot3^2\cdot5,
  \quad \text{或 }2\cdot3\cdot5.
\]
对每一种情况，都要找出所有可能的 $n_2$（满足 $n_2\mid n_1$ 且 $n_1n_2\mid n$）；再对每个得到的 $n_1,n_2$ 找出所有可能的 $n_3$，依此继续，直到得到所有满足 (1)--(3) 的列表。

例如，若 $n_1=2\cdot3^2\cdot5$，则唯一既整除 $n_1$ 又使 $n_1n_2$ 整除 $n$ 的 $n_2$ 是 $n_2=2$。此时 $n_1n_2=n$，所以列表完整：$2\cdot3^2\cdot5,2$。对应的阿贝尔群是 $\ZZ_{90}\times\ZZ_2$。

若 $n_1=2\cdot3\cdot5$，则 $n_2$ 的候选为 $2,3$ 或 $6$。若 $n_2=2$ 或 $3$，则因 $n_3\mid n_2$，必有 $n_3=n_2$（并且由 (3)，列表中还必须有第三项）。这导致矛盾，因为 $n_1n_2n_3$ 将分别被 $2^3$ 或 $3^3$ 整除，而 $n$ 不被这两个数整除。因此，第一项为 $2\cdot3\cdot5$ 的不变因子列表唯一为 $2\cdot3\cdot5,2\cdot3$，对应的阿贝尔群为 $\ZZ_{30}\times\ZZ_6$。

类似地，阶为 $180$ 的所有允许的不变因子列表及相应的阿贝尔群如下：
\[
\begin{array}{c|c}
\text{不变因子} & \text{阿贝尔群}\\ \hline
2^2\cdot3^2\cdot5 & \ZZ_{180}\\
2\cdot3^2\cdot5,
\ 2 & \ZZ_{90}\times\ZZ_2\\
2^2\cdot3\cdot5,
\ 3 & \ZZ_{60}\times\ZZ_3\\
2\cdot3\cdot5,
\ 2\cdot3 & \ZZ_{30}\times\ZZ_6
\end{array}
\]

上面的过程有些临时性，但它表明，确定给定阶 $n$ 的所有阿贝尔群的不变因子列表强烈依赖于 $n$ 的分解。下面的定理（我们将看到它在有限阿贝尔群情形中等价于基本定理）给出一种更系统且计算上快得多的方法来确定给定阶的所有有限阿贝尔群。更具体地说，若
\[
  n=p_1^{\alpha_1}p_2^{\alpha_2}\cdots p_k^{\alpha_k}
\]
是 $n$ 的素因子分解，则它说明阶为 $n$ 的阿贝尔群的所有允许不变因子列表，可以通过分别找出阶为 $p_i^{\alpha_i}$ 的群的允许列表来确定。对素数幂 $p^\alpha$，我们将看到，确定所有允许列表的问题等价于确定 $\alpha$ 的所有分拆（并且与 $p$ 无关）。

\medskip
\noindent\textbf{定理 5.} 设 $G$ 是阶为 $n>1$ 的阿贝尔群，并设 $n$ 分解为互异素数幂的唯一分解为
\[
  n=p_1^{\alpha_1}p_2^{\alpha_2}\cdots p_k^{\alpha_k}.
\]
则：
\begin{enumerate}[label=(\arabic*)]
  \item $G\cong A_1\times A_2\times\cdots\times A_k$，其中 $|A_i|=p_i^{\alpha_i}$；
  \item 对每个 $A\in\{A_1,A_2,
\ldots,A_k\}$，若 $|A|=p^\alpha$，则
  \[
    A\cong \ZZ_{p^{\beta_1}}\times\ZZ_{p^{\beta_2}}\times\cdots\times\ZZ_{p^{\beta_r}},
  \]
  其中 $\beta_1\geq\beta_2\geq\cdots\geq\beta_r\geq1$ 且 $\beta_1+\beta_2+\cdots+\beta_r=\alpha$；
  \item (1) 与 (2) 中的分解唯一。也就是说，若 $G\cong B_1\times\cdots\times B_m$，且 $|B_i|=p_i^{\alpha_i}$，则 $B_i\cong A_i$，并且 $B_i$ 与 $A_i$ 有相同的不变因子。
\end{enumerate}

\noindent\textbf{定义.} 前一定理中出现的整数 $p^{\beta_j}$ 称为 $G$ 的\emph{初等因子}。定理 5(1) 与 5(2) 中对 $G$ 的描述称为 $G$ 的\emph{初等因子分解}。

定理 (1) 中的子群 $A_i$ 是 $G$ 的 Sylow $p_i$-子群。因此 (1) 说明 $G$ 同构于其 Sylow 子群的直积（注意它们是正规子群，因为 $G$ 是阿贝尔群，所以也是唯一的）。第 (1) 部分常称为有限阿贝尔群的\emph{主分解定理}。\footnote{回忆：对阿贝尔群而言，Sylow $p$-子群有时称为 $p$-主分量。}同定理 3 一样，我们稍后会证明这个定理。

注意，对素数 $p$，有 $p^\beta\mid p^\gamma$ 当且仅当 $\beta\leq\gamma$。此外，$p^{\beta_1}\cdots p^{\beta_r}=p^\alpha$ 当且仅当 $\beta_1+
\cdots+
\beta_r=\alpha$。因此，定理 5(2) 中 $A$ 的分解就是 $A$ 的不变因子分解，只是把整数 $p^{\beta_j}$ 的“整除”条件转化为了指数上的“加法”条件。现在可见，$G$ 的初等因子就是当 $p$ 遍历 $|G|$ 的所有素因子时，Sylow $p$-子群的不变因子。

由定理 5，为了找出阶为 $n=p_1^{\alpha_1}p_2^{\alpha_2}\cdots p_k^{\alpha_k}$ 的所有阿贝尔群，必须对每个 $i$ 找出阶为 $p_i^{\alpha_i}$ 的群的所有可能不变因子列表。每个阿贝尔群的初等因子集合由从这 $k$ 个列表中各取一个不变因子集合而得到。阿贝尔群就是以这些初等因子为阶的循环群的直积；不同的初等因子列表给出非同构群。

在素数幂阶 $p^\beta$ 的情形，不变因子的条件化为
\begin{enumerate}[label=(\arabic*)]
  \item $\beta_j\geq 1$ 对所有 $j$ 成立；
  \item $\beta_i\geq\beta_{i+1}$ 对所有 $i$ 成立；
  \item $\beta_1+
\beta_2+
\cdots+
\beta_t=\beta$。
\end{enumerate}
因此，每个不变因子列表就是 $\beta$ 的一个按降序排列的分拆。特别地，阶为 $p^\beta$ 的非同构阿贝尔群个数等于 $\beta$ 的分拆数；这个数与素数 $p$ 无关。例如，阶为 $p^5$ 的阿贝尔群个数由 $5$ 的分拆列表得到：
\[
\begin{array}{c|c}
\text{不变因子中的指数} & \text{阿贝尔群}\\ \hline
5 & \ZZ_{p^5}\\
4,1 & \ZZ_{p^4}\times\ZZ_p\\
3,2 & \ZZ_{p^3}\times\ZZ_{p^2}\\
3,1,1 & \ZZ_{p^3}\times\ZZ_p\times\ZZ_p\\
2,2,1 & \ZZ_{p^2}\times\ZZ_{p^2}\times\ZZ_p\\
2,1,1,1 & \ZZ_{p^2}\times\ZZ_p\times\ZZ_p\times\ZZ_p\\
1,1,1,1,1 & \ZZ_p\times\ZZ_p\times\ZZ_p\times\ZZ_p\times\ZZ_p
\end{array}
\]
所以阶为 $p^5$ 的非同构阿贝尔群恰有 $7$ 个，列表中第一个是循环群 $\ZZ_{p^5}$，最后一个是初等阿贝尔群 $E_{p^5}$。

若 $n=p_1^{\alpha_1}p_2^{\alpha_2}\cdots p_k^{\alpha_k}$，并且 $q_i$ 是 $\alpha_i$ 的分拆数，则阶为 $n$ 的（不同的、非同构的）阿贝尔群个数为 $q_1q_2\cdots q_k$。

\paragraph{例}
若 $n=1800=2^3\cdot3^2\cdot5^2$，可列出如下：
\[
\begin{array}{c|c|c}
\text{阶 }p^\beta & \beta\text{ 的分拆} & \text{阿贝尔群}\\ \hline
2^3 & 3;
\ 2,1;
\ 1,1,1 & \ZZ_8;
\ \ZZ_4\times\ZZ_2;
\ \ZZ_2\times\ZZ_2\times\ZZ_2\\
3^2 & 2;
\ 1,1 & \ZZ_9;
\ \ZZ_3\times\ZZ_3\\
5^2 & 2;
\ 1,1 & \ZZ_{25};
\ \ZZ_5\times\ZZ_5
\end{array}
\]
通过从三行右列中各取一个阿贝尔群并取其直积，得到所有阶为 $1800$ 的阿贝尔群同构类型。以所有可能方式这样做，得到所有同构类型：
\[
\begin{array}{lll}
\ZZ_8\times\ZZ_9\times\ZZ_{25} &
\ZZ_8\times\ZZ_9\times\ZZ_5\times\ZZ_5 &
\ZZ_8\times\ZZ_3\times\ZZ_3\times\ZZ_{25} \\
\ZZ_8\times\ZZ_3\times\ZZ_3\times\ZZ_5\times\ZZ_5 &
\ZZ_4\times\ZZ_2\times\ZZ_9\times\ZZ_{25} &
\ZZ_4\times\ZZ_2\times\ZZ_9\times\ZZ_5\times\ZZ_5 \\
\ZZ_4\times\ZZ_2\times\ZZ_3\times\ZZ_3\times\ZZ_{25} &
\ZZ_4\times\ZZ_2\times\ZZ_3\times\ZZ_3\times\ZZ_5\times\ZZ_5 &
\ZZ_2\times\ZZ_2\times\ZZ_2\times\ZZ_9\times\ZZ_{25} \\
\ZZ_2\times\ZZ_2\times\ZZ_2\times\ZZ_9\times\ZZ_5\times\ZZ_5 &
\ZZ_2\times\ZZ_2\times\ZZ_2\times\ZZ_3\times\ZZ_3\times\ZZ_{25} &
\ZZ_2\times\ZZ_2\times\ZZ_2\times\ZZ_3\times\ZZ_3\times\ZZ_5\times\ZZ_5 .
\end{array}
\]
由基本定理，这就是阶为 $1800$ 的所有阿贝尔群的完整列表；每个阶为 $1800$ 的阿贝尔群都同构于其中恰好一个群，并且列表中任两个群都不同构。

我们强调，$G$ 的初等因子不是 $G$ 的不变因子（而是 $G$ 的某些子群的不变因子）。例如在上面列表第一种情形中，初等因子 $8,9,25$ 不满足不变因子列表的整除条件。

接下来说明如何由有限阿贝尔群的不变因子列表得到初等因子列表，反过来也如何由初等因子列表得到不变因子列表。无论群作为循环群直积以何种形式给出，都可以确定这些不变量。需要如下命题。

\medskip
\noindent\textbf{命题 6.} 设 $m,n\in\ZZ^+$。
\begin{enumerate}[label=(\arabic*)]
  \item $\ZZ_m\times\ZZ_n\cong\ZZ_{mn}$ 当且仅当 $(m,n)=1$。
  \item 若 $n=p_1^{\alpha_1}p_2^{\alpha_2}\cdots p_k^{\alpha_k}$，则
  \[
    \ZZ_n\cong \ZZ_{p_1^{\alpha_1}}\times\ZZ_{p_2^{\alpha_2}}\times\cdots\times\ZZ_{p_k^{\alpha_k}}.
  \]
\end{enumerate}

\begin{proof}
由于 (2) 可由 (1) 与对 $k$ 的归纳容易得到，我们集中证明 (1)。令 $\ZZ_m=\gen{x}$，$\ZZ_n=\gen{y}$，并令 $l=\operatorname{lcm}(m,n)$。注意 $l=mn$ 当且仅当 $(m,n)=1$。设 $x^ay^b$ 是 $\ZZ_m\times\ZZ_n$ 的一个典型元素。由第 1 节例 1，
\[
  (x^ay^b)^l=x^{la}y^{lb}=1^a1^b=1,
\]
因为 $m\mid l$ 且 $n\mid l$。若 $(m,n)\neq1$，则 $\ZZ_m\times\ZZ_n$ 的每个元素的阶至多为 $l$，严格小于 $mn$，所以它不可能同构于 $\ZZ_{mn}$。反之，若 $(m,n)=1$，则
\[
  |xy|=\operatorname{lcm}(|x|,|y|)=mn.
\]
由阶的考虑，$\ZZ_m\times\ZZ_n=\gen{xy}$ 是循环群。证明完毕。
\end{proof}

\paragraph{由不变因子得到初等因子}
设 $G$ 是类型 $(n_1,n_2,
\ldots,n_s)$ 的阿贝尔群，即
\[
  G\cong \ZZ_{n_1}\times\ZZ_{n_2}\times\cdots\times\ZZ_{n_s}.
\]
令
\[
  n=p_1^{\alpha_1}p_2^{\alpha_2}\cdots p_k^{\alpha_k}=n_1n_2\cdots n_s.
\]
把每个 $n_i$ 分解为
\[
  n_i=p_1^{\beta_{i1}}p_2^{\beta_{i2}}\cdots p_k^{\beta_{ik}},
  \qquad \beta_{ij}\geq0.
\]
由上面的命题，
\[
  \ZZ_{n_i}\cong \ZZ_{p_1^{\beta_{i1}}}\times\ZZ_{p_2^{\beta_{i2}}}\times\cdots\times\ZZ_{p_k^{\beta_{ik}}}
\]
对每个 $i$ 成立。若 $\beta_{ij}=0$，则 $\ZZ_{p_j^{\beta_{ij}}}=1$，可从直积中删去而不改变同构类型。因此，$G$ 的初等因子正是所有整数
\[
  p_j^{\beta_{ij}},\qquad 1\leq j\leq k,
\quad 1\leq i\leq s,
\quad \beta_{ij}\neq0.
\]

例如，若 $|G|=2^3\cdot3^2\cdot5^2$ 且 $G$ 的类型为 $(30,30,2)$，则
\[
  G\cong \ZZ_{30}\times\ZZ_{30}\times\ZZ_2.
\]
由于 $\ZZ_{30}\cong \ZZ_2\times\ZZ_3\times\ZZ_5$，
\[
  G\cong \ZZ_2\times\ZZ_3\times\ZZ_5\times\ZZ_2\times\ZZ_3\times\ZZ_5\times\ZZ_2.
\]
因此 $G$ 的初等因子为 $2,3,5,2,3,5,2$，把同素数幂归在一起，可写为 $2,2,2,3,3,5,5$。特别地，$G$ 同构于上面阶 $1800$ 的列表中的最后一个群。

若对每个 $j$ 把所有因子 $\ZZ_{p_j^{\beta_{ij}}}$ 收集在一起，所得直积就是 $G$ 的 Sylow $p_j$-子群 $A_j$。因此，上段中群的 Sylow $2$-子群同构于 $\ZZ_2\times\ZZ_2\times\ZZ_2$，即阶为 $8$ 的初等阿贝尔群。

\paragraph{由任意循环分解得到初等因子}
只要有限阿贝尔群 $G$ 被写成循环群的直积，即使循环分量的阶不是不变因子，上述过程也能得到 $G$ 的初等因子。例如，若 $G=\ZZ_6\times\ZZ_{15}$，列表 $6,15$ 既不是不变因子列表（整除条件失败），也不是初等因子列表（它们不是素数幂）。分解 $6=2\cdot3$ 与 $15=3\cdot5$，得到初等因子 $2,3,3,5$，并且
\[
  G\cong \ZZ_2\times\ZZ_3\times\ZZ_3\times\ZZ_5.
\]

\paragraph{由初等因子得到不变因子}
设 $G$ 是阶为 $n$ 的阿贝尔群，其中 $n=p_1^{\alpha_1}p_2^{\alpha_2}\cdots p_k^{\alpha_k}$，并且已知 $G$ 的初等因子。按以下步骤得到 $G$ 的不变因子：
\begin{enumerate}[label=(\arabic*)]
  \item 首先把所有同一个素数的幂的初等因子放在一起。这样得到 $k$ 个整数列表（每个 $p_j$ 一个列表）。
  \item 在每个列表中，将整数按非增顺序排列。
  \item 在这 $k$ 个列表中，设最长的列表有 $t$ 个整数。通过在每个列表末尾添上适当数量的 $1$，使每个列表长度都等于 $t$。
  \item 对每个 $i\in\{1,2,
\ldots,t\}$，第 $i$ 个不变因子 $n_i$ 是各个列表中第 $i$ 个整数的乘积。
\end{enumerate}
这样排列列表的目的，是确保有整除条件 $n_{i+1}\mid n_i$。

例如，设 $G$ 的初等因子为 $2,3,2,25,3,2$（所以 $|G|=2^3\cdot3^2\cdot5^2$）。分组并把每个列表补到长度 $3$，得到
\[
\begin{array}{c|c|c}
 p=2 & p=3 & p=5\\ \hline
 2&3&25\\
 2&3&1\\
 2&1&1
\end{array}
\]
所以 $G$ 的不变因子为 $2\cdot3\cdot25,
\ 2\cdot3\cdot1,
\ 2\cdot1\cdot1$，即
\[
  G\cong \ZZ_{150}\times\ZZ_6\times\ZZ_2.
\]

下面给出阶为 $1800$ 的阿贝尔群的不变因子分解；该列表中的第 $i$ 个群同构于前面按初等因子得到的第 $i$ 个群：
\[
\begin{array}{llll}
\ZZ_{1800},& \ZZ_{360}\times\ZZ_5,& \ZZ_{600}\times\ZZ_3,& \ZZ_{120}\times\ZZ_{15},\\
\ZZ_{900}\times\ZZ_2,& \ZZ_{180}\times\ZZ_{10},& \ZZ_{300}\times\ZZ_6,& \ZZ_{60}\times\ZZ_{30},\\
\ZZ_{450}\times\ZZ_2\times\ZZ_2,& \ZZ_{90}\times\ZZ_{10}\times\ZZ_2,&
\ZZ_{150}\times\ZZ_6\times\ZZ_2,& \ZZ_{30}\times\ZZ_{30}\times\ZZ_2.
\end{array}
\]

利用基本定理 3 与 5 的唯一性断言，我们可用这些过程判断任意两个有限循环群直积是否同构。例如，若要判断 $\ZZ_6\times\ZZ_{15}\cong\ZZ_{10}\times\ZZ_9$ 是否成立，先看二者阶是否相同（均为 $90$），再比较初等因子：
\[
\begin{array}{l}
\ZZ_6\times\ZZ_{15}\text{ 的初等因子为 }2,3,3,5,
\text{ 且同构于 }\ZZ_2\times\ZZ_3\times\ZZ_3\times\ZZ_5,\\
\ZZ_{10}\times\ZZ_9\text{ 的初等因子为 }2,5,9,
\text{ 且同构于 }\ZZ_2\times\ZZ_5\times\ZZ_9.
\end{array}
\]
初等因子列表不同，所以由定理 5，这两个群不同构。注意 $\ZZ_6\times\ZZ_{15}$ 没有阶为 $9$ 的元素，而 $\ZZ_{10}\times\ZZ_9$ 有这样的元素（参见练习 5）。

上述过程（稍加展开）给出了一个证明：对有限阿贝尔群，定理 3 与定理 5 等价。细节留给读者。

现在可以更好地理解分类定理的威力与局限。一方面，给定任意正整数 $n$，我们能显式描述所有阶为 $n$ 的阿贝尔群，这是一个重要成就。另一方面，判断一个具体群属于阶为 $n$ 的哪一个同构类型所需的信息可能相当多；若 $n$ 被较大的素数幂整除，这个信息量尤其大。

本节以一些以后会用到的术语作结。

\noindent\textbf{定义.}
\begin{enumerate}[label=(\arabic*)]
  \item 若 $G$ 是类型 $(n_1,n_2,
\ldots,n_t)$ 的有限阿贝尔群，则整数 $t$ 称为 $G$ 的\emph{秩}（此处自由秩为 $0$，不会混淆）。
  \item 若 $G$ 是任意群，$G$ 的\emph{指数}是满足 $x^n=1$ 对所有 $x\in G$ 成立的最小正整数 $n$；若不存在这样的正整数，则指数为 $\infty$。
\end{enumerate}

\subsection*{练习}
\begin{enumerate}[label=\arabic*.]
  \item 对 (a)--(e) 的每一项，给出指定阶的非同构阿贝尔群个数，不必列出群：(a) 阶 $100$；(b) 阶 $576$；(c) 阶 $1155$；(d) 阶 $42875$；(e) 阶 $2704$。

  \item 对 (a)--(e) 的每一项，给出指定阶的所有阿贝尔群的不变因子列表：(a) 阶 $270$；(b) 阶 $9801$；(c) 阶 $320$；(d) 阶 $105$；(e) 阶 $44100$。

  \item 对 (a)--(e) 的每一项，给出指定阶的所有阿贝尔群的初等因子列表，并把每个列表与前一题中相应的不变因子列表配对：(a) 阶 $270$；(b) 阶 $9801$；(c) 阶 $320$；(d) 阶 $105$；(e) 阶 $44100$。

  \item 对 (a)--(d) 的每一项，判断列出的阿贝尔群中哪些彼此同构（这里 $\{a_1,a_2,
\ldots,a_k\}$ 表示阿贝尔群 $\ZZ_{a_1}\times\ZZ_{a_2}\times\cdots\times\ZZ_{a_k}$）。
  \begin{enumerate}[label=(\alph*)]
    \item $\{4,9\}$，$\{6,6\}$，$\{8,3\}$，$\{9,4\}$，$\{6,4\}$，$\{64\}$。
    \item $\{2^2,2\cdot3^2\}$，$\{2^2\cdot3,2\cdot3\}$，$\{2^3\cdot3^2\}$，$\{2^2\cdot3^2,2\}$。
    \item $\{5^2\cdot7^2,3^2\cdot5\cdot7\}$，$\{3^2\cdot5^2\cdot7,5\cdot7^2\}$，$\{3\cdot5^2,7^2,3\cdot5\cdot7\}$，$\{5^2\cdot7,3^2\cdot5,7^2\}$。
    \item $\{2^2\cdot5\cdot7,2^3\cdot5^3,2\cdot5^2\}$，$\{2^3\cdot5^3\cdot7,2^3\cdot5^3\}$，$\{2^2\cdot2\cdot7,2^3,5^3,5^3\}$，$\{2\cdot5^3,2^2\cdot5^3,2^3,7\}$。
  \end{enumerate}

  \item 设 $G$ 是类型 $(n_1,n_2,
\ldots,n_t)$ 的有限阿贝尔群。证明：$G$ 含有阶为 $m$ 的元素当且仅当 $m\mid n_1$。推出 $G$ 的指数为 $n_1$。

  \item 证明任意有限群都有有限指数。给出一个有有限指数的无限群例子。指数为 $m$ 的有限群是否必含有阶为 $m$ 的元素？

  \item 设 $p$ 为素数，$A=\gen{x_1}\times\gen{x_2}\times\cdots\times\gen{x_n}$ 是阿贝尔 $p$-群，其中 $|x_i|=p^{\alpha_i}>1$。定义 $p$ 次幂映射
  \[
    \varphi:A\to A,
\qquad \varphi(x)=x^p.
  \]
  \begin{enumerate}[label=(\alph*)]
    \item 证明 $\varphi$ 是同态。
    \item 用给定生成元描述 $\varphi$ 的像与核。
    \item 证明 $\ker\varphi$ 与 $A/\im\varphi$ 的秩都等于 $n$，并证明二者均同构于阶为 $p^n$ 的初等阿贝尔群 $E_{p^n}$。
  \end{enumerate}

  \item 设 $A$ 是有限阿贝尔群（乘法记法），$p$ 是素数。令
  \[
    A^p=\gen{a^p\mid a\in A},
    \qquad A_p=\{x\in A\mid x^p=1\}.
  \]
  \begin{enumerate}[label=(\alph*)]
    \item 证明 $A/A^p\cong A_p$。（提示：证明二者都是初等阿贝尔群且阶相同。）
    \item 证明 $A$ 中阶为 $p$ 的子群个数等于 $A$ 中指数为 $p$ 的子群个数。（提示：化约到 $A$ 是初等阿贝尔 $p$-群的情形。）
  \end{enumerate}

  \item 令 $A=\ZZ_{60}\times\ZZ_{45}\times\ZZ_{12}\times\ZZ_{36}$。求 $A$ 中阶为 $2$ 的元素个数以及 $A$ 中指数为 $2$ 的子群个数。

  \item 设 $n,k$ 是正整数，$A$ 是秩为 $n$ 的自由阿贝尔群（加法记法）。证明 $A/kA$ 同构于 $n$ 个 $\ZZ/k\ZZ$ 的直积，其中 $kA=\{ka\mid a\in A\}$。

  \item 设 $G$ 是非平凡有限阿贝尔群，秩为 $t$。
  \begin{enumerate}[label=(\alph*)]
    \item 证明 $G$ 的秩等于其 Sylow 子群秩的最大值。
    \item 证明 $G$ 可由 $t$ 个元素生成，但没有少于 $t$ 个元素的子集能生成 $G$。
  \end{enumerate}

  \item 设 $n,m$ 为正整数，$d=(n,m)$。令 $\ZZ_n=\gen{x}$，$\ZZ_m=\gen{y}$。令 $A$ 为 $\gen{x}$ 与 $\gen{y}$ 的中心积，在其中识别一个阶为 $d$ 的元素。它有展示
  \[
    \gen{x,y\mid x^n=y^m=1,
\ xy=yx,
\ x^{n/d}=y^{m/d}}.
  \]
  把 $A$ 描述为两个循环群的直积。

  \item 设 $A=\gen{x_1}\times\cdots\times\gen{x_r}$ 是有限阿贝尔群，$|x_i|=n_i$。给出 $A$ 的展示。证明：若任意群 $G$ 含有两两可交换的元素 $g_1,
\ldots,g_r$，且 $g_i^{n_i}=1$，则存在唯一同态 $A\to G$，把 $x_i$ 送到 $g_i$。

  \item 对任意群 $G$，定义 $G$ 的\emph{对偶群} $\widehat G$ 为所有从 $G$ 到 $\CC$ 中单位根乘法群的同态组成的集合。按函数逐点乘法定义 $\widehat G$ 上的群运算：若 $\chi,\psi$ 是同态，则 $(\chi\psi)(g)=\chi(g)\psi(g)$。
  \begin{enumerate}[label=(\alph*)]
    \item 证明这个运算使 $\widehat G$ 成为阿贝尔群。
    \item 若 $G$ 是有限阿贝尔群，证明 $\widehat G\cong G$。（提示：把 $G$ 写成 $\gen{x_1}\times\cdots\times\gen{x_r}$；若 $n_i=|x_i|$，定义 $\chi_i$ 把 $x_i$ 送到 $e^{2\pi i/n_i}$，把 $x_j$（$j\neq i$）送到 $1$。）
  \end{enumerate}
  这个结果常表述为：有限阿贝尔群是自对偶的。它意味着有限阿贝尔群的子群格图上下翻转后相同。不过 $G$ 与其对偶之间没有自然同构（同构依赖于生成元的选择）；常说有限阿贝尔群非典范地同构于其对偶。

  \item 令 $G=\gen{x}\times\gen{y}$，其中 $|x|=8$ 且 $|y|=4$。
  \begin{enumerate}[label=(\alph*)]
    \item 找出所有 $a,b\in G$，使得 $G=\gen{a}\times\gen{b}$（把 $a,b$ 用 $x,y$ 表示）。
    \item 令 $H=\gen{x^2y,y^2}\cong\ZZ_4\times\ZZ_2$。证明不存在 $a,b\in G$ 使得 $G=\gen{a}\times\gen{b}$ 且 $H=\gen{a^2}\times\gen{b^2}$。
  \end{enumerate}

  \item 证明没有有限生成阿贝尔群是可除群（参见第 2.4 节练习 19）。
\end{enumerate}

\section{小阶群表}

此时，我们可以给出大多数小阶群的同构类型表。

下页表中那些不熟悉的非阿贝尔群将在第 5 节关于半直积的内容中构造（那里也会解释相关记号）。这里先给出它们的生成元与关系，即展示。

阶为 $12$ 的群 $\ZZ_3\rtimes\ZZ_4$ 可由生成元与关系
\[
  \gen{x,y\mid x^4=y^3=1,
\ x^{-1}yx=y^{-1}}
\]
描述。也就是说，它有一个正规 Sylow $3$-子群 $\gen{y}$，并且共轭作用下由一个阶为 $4$ 的元素 $x$ 取逆（$x^2$ 中心化 $y$）。

群 $(\ZZ_3\times\ZZ_3)\rtimes\ZZ_2$ 有生成元与关系
\[
  \gen{x,y,z\mid x^2=y^3=z^3=1,
\ yz=zy,
\ x^{-1}yx=y^{-1},
\ x^{-1}zx=z^{-1}}.
\]
也就是说，它有一个同构于 $\ZZ_3\times\ZZ_3$ 的正规 Sylow $3$-子群 $\gen{y,z}$，该子群被一个阶为 $2$ 的元素 $x$ 通过共轭取逆。

阶为 $20$ 的群 $\ZZ_5\rtimes\ZZ_4$ 有生成元与关系
\[
  \gen{x,y\mid x^4=y^5=1,
\ x^{-1}yx=y^{-1}}.
\]
它有一个正规 Sylow $5$-子群 $\gen{y}$，并被阶为 $4$ 的元素 $x$ 通过共轭取逆（$x^2$ 中心化 $y$）。

\begin{center}
\begin{tabular}{c|c|c|c}
\hline
\text{阶} & \text{同构类型数} & \text{阿贝尔群} & \text{非阿贝尔群}\\ \hline
1&1&1&无\\
2&1&$\ZZ_2$&无\\
3&1&$\ZZ_3$&无\\
4&2&$\ZZ_4,
\ \ZZ_2\times\ZZ_2$&无\\
5&1&$\ZZ_5$&无\\
6&2&$\ZZ_6$&$D_6\cong S_3$\\
7&1&$\ZZ_7$&无\\
8&5&$\ZZ_8,
\ \ZZ_4\times\ZZ_2,
\ \ZZ_2\times\ZZ_2\times\ZZ_2$&$D_8,
\ Q_8$\\
9&2&$\ZZ_9,
\ \ZZ_3\times\ZZ_3$&无\\
10&2&$\ZZ_{10}$&$D_{10}$\\
11&1&$\ZZ_{11}$&无\\
12&5&$\ZZ_{12},
\ \ZZ_6\times\ZZ_2$&$A_4,
\ D_{12},
\ \ZZ_3\rtimes\ZZ_4$\\
13&1&$\ZZ_{13}$&无\\
14&2&$\ZZ_{14}$&$D_{14}$\\
15&1&$\ZZ_{15}$&无\\
16&14&$\ZZ_{16},
\ \ZZ_8\times\ZZ_2,
\ \ZZ_4\times\ZZ_4,
\ \ZZ_4\times\ZZ_2\times\ZZ_2,
\ \ZZ_2^4$&未列出\\
17&1&$\ZZ_{17}$&无\\
18&5&$\ZZ_{18},
\ \ZZ_6\times\ZZ_3$&$D_{18},
\ S_3\times\ZZ_3,
\ (\ZZ_3\times\ZZ_3)\rtimes\ZZ_2$\\
19&1&$\ZZ_{19}$&无\\
20&5&$\ZZ_{20},
\ \ZZ_{10}\times\ZZ_2$&$D_{20},
\ \ZZ_5\rtimes\ZZ_4,
\ F_{20}$\\ \hline
\end{tabular}
\end{center}

阶为 $20$ 的群 $F_{20}$ 有生成元与关系
\[
  \gen{x,y\mid x^4=y^5=1,
\ xyx^{-1}=y^2}.
\]
也就是说，它有一个正规 Sylow $5$-子群 $\gen{y}$，该子群被一个阶为 $4$ 的元素 $x$ 通过共轭作用平方。可以验证，这个群出现在 $S_5$ 中 Sylow $5$-子群的正规化子中，例如
\[
  F_{20}=\gen{(2
\ 3
\ 5
\ 4),(1
\ 2
\ 3
\ 4
\ 5)}.
\]
这个群称为阶为 $20$ 的 Frobenius 群。

\subsection*{练习}
\begin{enumerate}[label=\arabic*.]
  \item 证明 $D_{16}$、$\ZZ_2\times D_8$、$\ZZ_2\times Q_8$、$\ZZ_4*D_8$、$QD_{16}$ 与 $M$ 是两两不同构的阶为 $16$ 的非阿贝尔群（其中 $\ZZ_4*D_8$ 在第 1 节练习 12 中描述，而 $QD_{16}$ 与 $M$ 在第 2.5 节练习中描述）。
\end{enumerate}

\section{识别直积}

到目前为止，我们已经看到直积既可用来由“较小”群构造“较大”群，也可用来把有限生成阿贝尔群分解成循环因子。甚至某些以其他形式给出的非阿贝尔群也可能分解成较小群的直积。本节的目的，是给出一个判别准则，用来识别一个群何时是其某些子群的直积，并用若干例子说明。

在此之前，我们引入关于交换子的标准记号与初等结果。这些记号将简化叙述，并会在第 6 章讨论幂零群时再次使用。

\noindent\textbf{定义.} 设 $G$ 是群，$x,y\in G$，且 $A,B$ 是 $G$ 的非空子集。
\begin{enumerate}[label=(\arabic*)]
  \item 定义
  \[
    [x,y]=x^{-1}y^{-1}xy,
  \]
  称为 $x$ 与 $y$ 的\emph{交换子}。
  \item 定义
  \[
    [A,B]=\gen{[a,b]\mid a\in A,
\ b\in B},
  \]
  即由来自 $A$ 与 $B$ 的元素的交换子生成的子群。
  \item 定义
  \[
    G'=\gen{[x,y]\mid x,y\in G},
  \]
  称为 $G$ 的\emph{交换子子群}。
\end{enumerate}

交换子 $[x,y]$ 等于 $1$ 当且仅当 $x$ 与 $y$ 可交换，这解释了该术语。下面命题表明交换子如何衡量 $xy$ 与 $yx$ 在 $G$ 中的差别。

\medskip
\noindent\textbf{命题 7.} 设 $G$ 是群，$x,y\in G$，且 $H\leq G$。则：
\begin{enumerate}[label=(\arabic*)]
  \item $xy=yx[x,y]$；特别地，$xy=yx$ 当且仅当 $[x,y]=1$。
  \item $H\normal G$ 当且仅当 $[H,G]\leq H$。
  \item 对任意 $G$ 的自同构 $\alpha$，有 $\alpha([x,y])=[\alpha(x),\alpha(y)]$；因此 $G'\charac G$，且 $G/G'$ 是阿贝尔群。
  \item $G/G'$ 是 $G$ 的最大阿贝尔商群：若 $H\normal G$ 且 $G/H$ 是阿贝尔群，则 $G'\leq H$；反之，若 $G'\leq H$，则 $H\normal G$ 且 $G/H$ 是阿贝尔群。
  \item 若 $\varphi:G\to A$ 是任意到阿贝尔群 $A$ 的同态，则 $\varphi$ 通过 $G/G'$ 分解；也就是说 $G'\leq\ker\varphi$，并且存在唯一同态 $\bar\varphi:G/G'\to A$ 使得下图交换：
  \[
  \begin{array}{ccc}
  G & \longrightarrow & G/G'\\
  & \searrow_{\varphi} & \downarrow \bar\varphi\\
  & & A
  \end{array}
  \]
\end{enumerate}

\begin{proof}
(1) 由 $[x,y]$ 的定义立即得到。

(2) 按定义，$H\normal G$ 当且仅当对所有 $g\in G$ 与 $h\in H$ 有 $g^{-1}hg\in H$。而对 $h\in H$，$g^{-1}hg\in H$ 当且仅当 $h^{-1}g^{-1}hg\in H$，也就是 $[h,g]\in H$。因此 $H\normal G$ 当且仅当 $[H,G]\leq H$。

(3) 设 $\alpha\in\Aut(G)$。则
\[
  \alpha([x,y])=\alpha(x^{-1}y^{-1}xy)=\alpha(x)^{-1}\alpha(y)^{-1}\alpha(x)\alpha(y)=[\alpha(x),\alpha(y)].
\]
故 $\alpha$ 把交换子的集合双射到自身；由于交换子生成 $G'$，得 $\alpha(G')=G'$，即 $G'\charac G$。又若 $xG'$ 与 $yG'$ 是 $G/G'$ 的任意元素，则
\[
  (xG')(yG')=(xy)G'=(yx[x,y])G'=(yx)G'=(yG')(xG'),
\]
所以 $G/G'$ 是阿贝尔群。

(4) 若 $H\normal G$ 且 $G/H$ 是阿贝尔群，则对任意 $x,y\in G$，有 $(xH)(yH)=(yH)(xH)$，于是
\[
  H=(xH)^{-1}(yH)^{-1}(xH)(yH)=x^{-1}y^{-1}xyH=[x,y]H.
\]
所以 $[x,y]\in H$ 对所有 $x,y$ 成立，从而 $G'\leq H$。

反之，若 $G'\leq H$，则由 (3)，$G/G'$ 是阿贝尔群，因而 $G/G'$ 的每个子群都正规。特别地，$H/G'\normal G/G'$。由格同构定理 $H\normal G$，再由第三同构定理
\[
  G/H\cong (G/G')/(H/G'),
\]
故 $G/H$ 是阿贝尔群。

(5) 这只是用同态语言重述 (4)。
\end{proof}

对 $G$ 取商 $G/G'$ 会把所有交换子压成单位元，使商群中的所有元素都可交换。正如 (4) 所示，强反命题也成立：$G/H$ 是阿贝尔群当且仅当 $G'\leq H$。

下一节我们将给出一个阶为 $96$ 的群，它的交换子子群中有一个元素不能写成任何单个交换子 $[x,y]$。因此 $G'$ 不一定仅由所有单个交换子组成，而是由这些元素生成的群。

\paragraph{例}
\begin{enumerate}[label=(\arabic*)]
  \item 群 $G$ 是阿贝尔群当且仅当 $G'=1$。

  \item 有时无需显式计算交换子即可求出交换子子群。例如，若 $G=D_8$，由于 $Z(D_8)=\gen{r^2}\normal D_8$ 且 $D_8/Z(D_8)$ 是阿贝尔群（Klein 四元群），交换子子群 $D_8'$ 是 $Z(D_8)$ 的子群。又 $D_8$ 非阿贝尔，所以 $D_8'$ 非平凡，唯一可能是 $D_8'=Z(D_8)$。类似地，$Q_8'=Z(Q_8)=\gen{-1}$。更一般地，若 $G$ 是阶为 $p^3$ 的非阿贝尔群，其中 $p$ 是素数，则 $G'=Z(G)$ 且 $|G'|=p$（练习 7）。

  \item 令
  \[
    D_{2n}=\gen{r,s\mid r^n=s^2=1,
\ s^{-1}rs=r^{-1}}.
  \]
  由于 $[r,s]=r^{-2}$，有 $\gen{r^2}\leq D_{2n}'$。另一方面，$\gen{r^2}\normal D_{2n}$，且 $D_{2n}/\gen{r^2}$ 中 $r$ 与 $s$ 的像生成该商群，并且是可交换的阶不超过 $2$ 的元素，所以商群阿贝尔，从而 $D_{2n}'\leq\gen{r^2}$。因此 $D_{2n}'=\gen{r^2}$。若 $n$ 为奇数，则 $\gen{r^2}=\gen r$；若 $n$ 为偶数，则 $\gen{r^2}$ 在 $\gen r$ 中指数为 $2$。因此 $D_{2n}'$ 在 $D_{2n}$ 中的指数分别为 $2$ 或 $4$，取决于 $n$ 是奇数还是偶数。

  \item 由于由 $g\in G$ 给出的共轭是 $G$ 的自同构，命题 (3) 给出 $[a^g,b^g]=[a,b]^g$。也就是说，交换子的共轭仍是交换子。例如，一旦在 $S_n$ 中把某个循环类型的一个元素写成交换子，同一循环类型的所有元素都是交换子。特别地，$S_5$ 中每个 $5$-循环都是交换子：把五边形顶点标为 $1,
\ldots,5$，可视 $D_{10}\leq S_5$。由前例，阶为 $5$ 的元素是 $D_{10}$ 中的交换子，因而也是 $S_5$ 中的交换子。显式地，
  \[
    (1
\ 4
\ 2
\ 5
\ 3)=[(1
\ 2
\ 3
\ 4
\ 5),(2
\ 5)(4
\ 3)].
  \]
\end{enumerate}

下面的结果其实可从命题 3.13 的证明中推出，但这里单独列出以便引用。

\medskip
\noindent\textbf{命题 8.} 设 $H,K$ 是群 $G$ 的子群。集合 $HK$ 中每个元素写成 $hk$（$h\in H,k\in K$）的不同方式数为 $|H\cap K|$。特别地，若 $H\cap K=1$，则 $HK$ 的每个元素都能唯一写成 $hk$，其中 $h\in H,k\in K$。

\begin{proof}
练习。
\end{proof}

本节的主要结果是如下识别定理。

\medskip
\noindent\textbf{定理 9.} 设 $G$ 是群，$H,K$ 是其子群，满足：
\begin{enumerate}[label=(\arabic*)]
  \item $H\normal G$ 且 $K\normal G$；
  \item $H\cap K=1$。
\end{enumerate}
则
\[
  HK\cong H\times K.
\]

\begin{proof}
由假设 (1)，$HK$ 是 $G$ 的子群（参见推论 3.15）。令 $h\in H$，$k\in K$。由于 $H\normal G$，有 $k^{-1}hk\in H$，所以 $h^{-1}(k^{-1}hk)\in H$。类似地，由 $K\normal G$，有 $(h^{-1}k^{-1}h)k\in K$。而
\[
  h^{-1}k^{-1}hk
\]
同时属于 $H$ 与 $K$，故由 $H\cap K=1$ 得 $h^{-1}k^{-1}hk=1$，即 $hk=kh$。所以 $H$ 的每个元素与 $K$ 的每个元素可交换。

由前一命题，$HK$ 的每个元素可唯一写成 $hk$，其中 $h\in H,k\in K$。定义
\[
  \varphi:HK\to H\times K,
\qquad hk\mapsto(h,k).
\]
这是良定义的。若 $h_1,h_2\in H$，$k_1,k_2\in K$，则 $h_2$ 与 $k_1$ 可交换，所以
\[
  (h_1k_1)(h_2k_2)=(h_1h_2)(k_1k_2).
\]
于是
\[
\varphi(h_1k_1h_2k_2)=(h_1h_2,k_1k_2)=(h_1,k_1)(h_2,k_2)=\varphi(h_1k_1)\varphi(h_2k_2),
\]
故 $\varphi$ 是同态。又因为每个元素的表示唯一，$\varphi$ 是双射，故为同构。
\end{proof}

\noindent\textbf{定义.} 若 $G$ 是群，$H,K$ 是 $G$ 的正规子群且 $H\cap K=1$，则称 $HK$ 是 $H$ 与 $K$ 的\emph{内直积}。需要强调时，我们称 $H\times K$ 为 $H$ 与 $K$ 的\emph{外直积}。

由定理 9，内直积与外直积的区别只是记号上的：内直积的元素写成 $hk$，外直积的元素写成有序对 $(h,k)$。此前我们已经在这些形式之间切换过。例如，当 $\ZZ_n=\gen a$、$\ZZ_m=\gen b$ 时，我们写 $x=(a,1)$、$y=(1,b)$，于是 $\ZZ_n\times\ZZ_m$ 的每个元素可写成 $x^ry^s$。

\paragraph{例}
\begin{enumerate}[label=(\arabic*)]
  \item 若 $n$ 是正奇数，则 $D_{4n}\cong D_{2n}\times\ZZ_2$。令
  \[
    D_{4n}=\gen{r,s\mid r^{2n}=s^2=1,
\ srs=r^{-1}}
  \]
  为通常的展示。取 $H=\gen{s,r^2}$，$K=\gen{r^n}$。几何上，若 $D_{4n}$ 是正 $2n$ 边形的对称群，则 $H$ 是由连接顶点 $2i$ 与 $2i+2$ 得到的内接正 $n$ 边形的对称群。令 $\rho=r^2$，则 $H$ 有 $D_{2n}$ 的通常展示。$H$ 在 $D_{4n}$ 中指数为 $2$，所以正规。又 $|r^n|=2$，并且 $r$ 显然中心化 $r^n$；由 $srs=r^{-1}$ 得 $sr^ns=r^{-n}=r^n$，故 $K\leq Z(D_{4n})$，从而 $K\normal D_{4n}$。最后，因为 $r^2$ 的阶为奇数，$r^n\notin H$，所以 $H\cap K=1$。定理 9 给出同构。

  \item 令 $I$ 是 $\{1,2,
\ldots,n\}$ 的子集，令 $G$ 是 $S_n$ 中 $I$ 的集合稳定子：
  \[
    G=\{\sigma\in S_n\mid \sigma(i)\in I\text{ 对所有 }i\in I\}.
  \]
  令 $J=\{1,2,
\ldots,n\}-I$。设 $H$ 为 $I$ 的逐点稳定子，$K$ 为 $J$ 的逐点稳定子：
  \[
    H=\{\sigma\in G\mid \sigma(i)=i\text{ 对所有 }i\in I\},
  \]
  \[
    K=\{\tau\in G\mid \tau(j)=j\text{ 对所有 }j\in J\}.
  \]
  容易看出 $H,K$ 都是 $G$ 的正规子群（事实上它们分别是 $G$ 在 $I$ 与 $J$ 上作用的核）。且 $H\cap K=1$。每个 $G$ 的元素稳定 $I$ 与 $J$，所以其循环分解中的每个循环只涉及 $I$ 的元素或只涉及 $J$ 的元素。因此每个元素都可写成 $\sigma_I\sigma_J$，其中 $\sigma_I\in H$、$\sigma_J\in K$。于是 $G=HK$，由定理 9，$G\cong H\times K$。若 $m=|I|$，则
  \[
    H\cong S_{n-m},
\qquad K\cong S_m,
\qquad G\cong S_m\times S_{n-m},
  \]
  其中约定 $S_0=1$。

  \item 令 $\sigma\in S_n$，令
  \[
    I=\{i\in\{1,2,
\ldots,n\}\mid \sigma(i)=i\}
  \]
  为被 $\sigma$ 逐点固定的集合。若 $C=C_{S_n}(\sigma)$，则由第 4.3 节练习 18，$C$ 稳定 $I$ 及其补集 $J$。由前例，$C$ 同构于 $H\times K$ 的某个子群，其中 $H$ 是固定 $I$ 逐点的置换子群，$K$ 是固定 $J$ 逐点的置换子群。注意 $\sigma\in H$。并且任意固定 $J$ 逐点的置换都与 $\sigma$ 可交换，因为二者移动的整数没有交集。因此 $C$ 包含整个 $K$，且
  \[
    C_{S_n}(\sigma)=C_H(\sigma)\times K\cong C_{S_{n-|I|}}(\sigma)\times S_{|I|}.
  \]
  特别地，若 $\sigma$ 是 $S_n$ 中的一个 $m$-循环，则
  \[
    C_{S_n}(\sigma)=\gen{\sigma}\times S_{n-m}.
  \]
  后者的阶为 $m(n-m)!$，这正是第 4.3 节中计算所得。
\end{enumerate}

\subsection*{练习}
设 $G$ 是群。
\begin{enumerate}[label=\arabic*.]
  \item 证明若 $x,y\in G$，则 $[y,x]=[x,y]^{-1}$。推出对任意子集 $A,B\subseteq G$，有 $[A,B]=[B,A]$。
  \item 证明 $G$ 的子群 $H$ 正规当且仅当 $[G,H]\leq H$。
  \item 设 $a,b,c\in G$。证明
  \[
    [a,bc]=[a,c](c^{-1}[a,b]c),
\qquad
    [ab,c]=(b^{-1}[a,c]b)[b,c].
  \]
  \item 求 $S_4$ 与 $A_4$ 的交换子子群。
  \item 证明对所有 $n\geq5$，$A_n$ 是 $S_n$ 的交换子子群。
  \item 把 $A_5$ 中每种循环类型的一个代表元表示为 $S_5$ 中的交换子。
  \item 证明若 $p$ 是素数且 $P$ 是阶为 $p^3$ 的非阿贝尔群，则 $P'=Z(P)$。
  \item 假设 $x,y\in G$，且 $x,y$ 都与 $[x,y]$ 可交换。证明对所有 $n\in\ZZ^+$，
  \[
    (xy)^n=x^ny^n[y,x]^{n(n-1)/2}.
  \]
  \item 证明若 $p$ 是奇素数且 $P$ 是阶为 $p^3$ 的群，则 $p$ 次幂映射 $x\mapsto x^p$ 是从 $P$ 到 $Z(P)$ 的同态。若 $P$ 非循环，说明该同态的核的阶为 $p^2$ 或 $p^3$。在阶为 $8$ 的非阿贝尔群中，平方映射是否为同态？上述证明中哪里需要 $p$ 为奇数？
  \item 证明有限阿贝尔群是其 Sylow 子群的直积。
  \item 证明若 $G=HK$，其中 $H,K$ 是 $G$ 的特征子群且 $H\cap K=1$，则 $\Aut(G)\cong\Aut(H)\times\Aut(K)$。推出若 $G$ 是有限阶阿贝尔群，则 $\Aut(G)$ 同构于其 Sylow 子群自同构群的直积。
  \item 用定理 4.17 描述有限循环群的自同构群。
  \item 证明 $D_{8n}$ 不同构于 $D_{4n}\times\ZZ_2$。
  \item 设 $F$ 是域，令
  \[
    G=\{(a_{ij})\in\GL_n(F)\mid a_{ij}=0\text{ 若 }i>j,
\ a_{11}=a_{22}=\cdots=a_{nn}\}
  \]
  为所有对角元相同的可逆上三角矩阵组成的群。证明 $G\cong D\times U$，其中 $D$ 是所有非零标量矩阵组成的群，$U$ 是所有对角线上为 $1$ 的上三角矩阵组成的群。
  \item 若 $A,B$ 是 $G$ 的正规子群，且 $G/A$ 与 $G/B$ 都是阿贝尔群，证明 $G/(A\cap B)$ 是阿贝尔群。
  \item 证明若 $K\normal G$，则 $K'\normal G$。
  \item 若 $K\normal G$ 且 $K$ 是循环群，证明 $G'\leq C_G(K)$。（回忆循环群的自同构群是阿贝尔群。）
  \item 设 $K_1,
\ldots,K_n$ 是非阿贝尔单群，令 $G=K_1\times\cdots\times K_n$。证明 $G$ 的每个正规子群都形如 $G_I$，其中 $I\subseteq\{1,2,
\ldots,n\}$。
  \item 群 $H$ 称为\emph{完美群}，若 $H'=H$。
  \begin{enumerate}[label=(\alph*)]
    \item 证明每个非阿贝尔单群都是完美群。
    \item 证明若 $H,K$ 是群 $G$ 的完美子群，则 $\gen{H,K}$ 也是完美群。推广到任意一族完美子群生成的子群。
    \item 证明完美子群的任意共轭仍是完美子群。
    \item 证明任意群 $G$ 有唯一的极大完美子群，且该子群正规。
  \end{enumerate}
  \item 令 $H(F)$ 为域 $F$ 上的 Heisenberg 群（参见第 1.4 节练习 11）。求 $X,Y\in H(F)$ 时交换子 $[X,Y]$ 的显式公式，并证明 $H(F)$ 的交换子子群等于其中心。
\end{enumerate}

\section{半直积}

本节研究两个群 $H$ 与 $K$ 的\emph{半直积}。它推广了直积的概念：在直积中 $H$ 与 $K$ 都正规，而半直积中只要求其中一个因子正规。这个构造使我们在某些情况下能由群 $H$ 与 $K$ 构造一个“较大”的群 $G$，其中分别含有同构于 $H$ 与 $K$ 的子群。在这种情形中，子群 $H$ 将是 $G$ 的正规子群，而 $K$ 不一定正规。因此，即使 $H$ 与 $K$ 都是阿贝尔群，我们也能够构造非阿贝尔群。这个构造会大大扩充我们可用的群例子。与前一节一样，我们随后证明一个识别定理，使我们能把某些熟悉的群分解成较小的“因子”，并由此得到一些分类定理。

作为动机，假设已有一个群 $G$，其中含有子群 $H,K$，满足
\begin{enumerate}[label=(\alph*)]
  \item $H\normal G$（但 $K$ 不一定正规）；
  \item $H\cap K=1$。
\end{enumerate}
此时 $HK$ 仍是 $G$ 的子群（推论 3.15），并且由命题 8，$HK$ 的每个元素都唯一写成 $hk$，其中 $h\in H,k\in K$。因此 $HK$ 与有序对 $(h,k)$ 的集合之间有一个双射，$hk\mapsto(h,k)$。给定 $h_1k_1$ 与 $h_2k_2$，它们的乘积可写为
\begin{align}
(h_1k_1)(h_2k_2)
&=h_1k_1h_2(k_1^{-1}k_1)k_2\notag\\
&=h_1(k_1h_2k_1^{-1})k_1k_2\notag\\
&=h_3k_3, \tag{5.1}
\end{align}
其中 $h_3=h_1(k_1h_2k_1^{-1})\in H$，$k_3=k_1k_2\in K$。这里使用了 $H\normal G$。

半直积的基本思想是把这个构造反过来：从两个抽象群 $H$ 与 $K$ 出发，设法定义一个群，其中含有它们的同构拷贝，并使 (a) 与 (b) 成立。式 (5.1) 中，$k_3$ 只由 $K$ 中的乘法给出，而 $h_3$ 由 $H$ 中的 $h_1$ 与 $k_1h_2k_1^{-1}$ 的乘积给出。由于 $H\normal G$，群 $K$ 通过共轭作用在 $H$ 上：
\[
  k\cdot h=khk^{-1},
\qquad h\in H,
\ k\in K.
\]
于是 (5.1) 可写为
\[
  (h_1k_1)(h_2k_2)=(h_1(k_1\cdot h_2))(k_1k_2). \tag{5.2}
\]
$K$ 在 $H$ 上的共轭作用给出从 $K$ 到 $\Aut(H)$ 的同态 $\varphi$，所以 (5.2) 表明 $HK$ 中的乘法只依赖于 $H$ 中的乘法、$K$ 中的乘法以及同态 $\varphi$。

现在，我们用这个解释来定义一个群：给定群 $H,K$ 以及同态 $\varphi:K\to\Aut(H)$。

\medskip
\noindent\textbf{定理 10.} 设 $H,K$ 是群，$\varphi:K\to\Aut(H)$ 是同态。用 $\cdot$ 表示由 $\varphi$ 确定的 $K$ 在 $H$ 上的左作用。令 $G$ 为所有有序对 $(h,k)$ 组成的集合，其中 $h\in H,k\in K$，并在 $G$ 上定义乘法
\[
  (h_1,k_1)(h_2,k_2)=(h_1(k_1\cdot h_2),k_1k_2).
\]
则：
\begin{enumerate}[label=(\arabic*)]
  \item 这个乘法使 $G$ 成为群，且 $|G|=|H||K|$。
  \item 集合 $\{(h,1)\mid h\in H\}$ 与 $\{(1,k)\mid k\in K\}$ 是 $G$ 的子群，并且映射 $h\mapsto(h,1)$ 与 $k\mapsto(1,k)$ 分别是 $H$ 与 $K$ 到这些子群的同构。
\end{enumerate}
在把 $H,K$ 与 (2) 中的同构拷贝等同后，还有：
\begin{enumerate}[label=(\arabic*),resume]
  \item $H\normal G$；
  \item $H\cap K=1$；
  \item 对所有 $h\in H,k\in K$，有
  \[
    khk^{-1}=k\cdot h=\varphi(k)(h).
  \]
\end{enumerate}

\begin{proof}
利用 $\cdot$ 是 $K$ 在 $H$ 上的作用，可直接验证 $G$ 是群。例如，结合律如下：对所有 $(a,x),(b,y),(c,z)\in G$，
\begin{align*}
((a,x)(b,y))(c,z)
&=(a(x\cdot b),xy)(c,z)\\
&=(a(x\cdot b)((xy)\cdot c),xyz)\\
&=(a(x\cdot b)(x\cdot(y\cdot c)),xyz)\\
&=(a(x\cdot(b(y\cdot c))),xyz)\\
&=(a,x)(b(y\cdot c),yz)\\
&=(a,x)((b,y)(c,z)).
\end{align*}
单位元为 $(1,1)$，而
\[
  (h,k)^{-1}=(k^{-1}\cdot h^{-1},k^{-1}).
\]
$G$ 的阶显然为 $|H||K|$，证明了 (1)。

令 $\bar H=\{(h,1)\mid h\in H\}$，$\bar K=\{(1,k)\mid k\in K\}$。若 $a,b\in H$，则
\[
  (a,1)(b,1)=(a(1\cdot b),1)=(ab,1),
\]
若 $x,y\in K$，则
\[
  (1,x)(1,y)=(1,xy).
\]
这说明 $\bar H,
\bar K$ 是 $G$ 的子群，并且 (2) 中的映射是同构。

显然 $H\cap K=1$。又
\begin{align*}
(1,k)(h,1)(1,k)^{-1}
&=(1,k)(h,1)(1,k^{-1})\\
&=(k\cdot h,k)(1,k^{-1})\\
&=(k\cdot h,k\cdot1,kk^{-1})\\
&=(k\cdot h,1),
\end{align*}
所以在上述等同下 $khk^{-1}=k\cdot h$，即 (5)。最后，$K$ 正规化 $H$，而 $H$ 当然正规化自身；因为 $G=HK$，得 $N_G(H)=G$，即 $H\normal G$。
\end{proof}

\noindent\textbf{定义.} 设 $H,K$ 是群，$\varphi:K\to\Aut(H)$ 是同态。定理 10 中描述的群称为 $H$ 与 $K$ 关于 $\varphi$ 的\emph{半直积}，记为
\[
  H\rtimes_\varphi K.
\]
在没有混淆时，简写为 $H\rtimes K$。

这个记号提醒我们，$H$ 的拷贝是正规的“因子”，并且半直积构造在 $H$ 与 $K$ 中并不对称（与直积不同）。在给出例子之前，先说明半直积何时就是直积。

\medskip
\noindent\textbf{命题 11.} 设 $H,K$ 是群，$\varphi:K\to\Aut(H)$ 是同态。下列条件等价：
\begin{enumerate}[label=(\arabic*)]
  \item $H\rtimes K$ 与 $H\times K$ 之间的恒等集合映射是群同态（因此是同构）；
  \item $\varphi$ 是从 $K$ 到 $\Aut(H)$ 的平凡同态；
  \item $K\normal H\rtimes K$。
\end{enumerate}

\begin{proof}
(1)$\Rightarrow$(2)：在半直积中
\[
  (h_1,k_1)(h_2,k_2)=(h_1(k_1\cdot h_2),k_1k_2).
\]
若恒等集合映射是群同态，则同一乘积也应为 $(h_1h_2,k_1k_2)$。比较第一分量，得 $k_1\cdot h_2=h_2$ 对所有 $h_2\in H,k_1\in K$ 成立，即 $K$ 平凡作用在 $H$ 上。

(2)$\Rightarrow$(3)：若 $\varphi$ 平凡，则 $K$ 平凡作用在 $H$ 上，故由定理 10(5)，$H$ 的元素与 $K$ 的元素可交换。于是 $H$ 正规化 $K$，而 $K$ 正规化自身；由于 $G=HK$，得到 $K\normal G$。

(3)$\Rightarrow$(1)：若 $K\normal H\rtimes K$，则如定理 9 的证明，对所有 $h\in H,k\in K$，有 $[h,k]\in H\cap K=1$。所以 $hk=kh$，$K$ 在 $H$ 上的作用平凡。于是半直积中的乘法与直积中的乘法相同。
\end{proof}


\paragraph{例}
以下所有例子中，$H,K$ 是群，$\varphi:K\to\Aut(H)$ 是同态，$K$ 在 $H$ 上的相应作用用点号表示。令 $G=H\rtimes K$，并且如定理 10 中那样，把 $H,K$ 视为 $G$ 的子群。我们将使用命题 4.16 与命题 4.17 来确定某些特定群 $H$ 的自同构群中的同态 $\varphi$。以下每个例子中，证明 $\varphi$ 是同态都很容易（因为 $K$ 通常是循环群），故细节略去。
\begin{enumerate}[label=(\arabic*)]
  \item 令 $H$ 为任意阿贝尔群（甚至可为无限阶），令 $K=\gen{x}\cong\ZZ_2$ 为阶为 $2$ 的群。定义 $\varphi:K\to\Aut(H)$，把 $x$ 送到 $H$ 上取逆的自同构；于是相应作用为
  \[
    x\cdot h=h^{-1},\qquad h\in H.
  \]
  则 $G$ 含有指数为 $2$ 的子群 $H$，并且
  \[
    xhx^{-1}=h^{-1},\qquad h\in H.
  \]
  特别有趣的是 $H$ 循环的情形：若 $H=\ZZ_n$，则 $G$ 可识别为 $D_{2n}$；若 $H=\ZZ$，则记 $G$ 为 $D_\infty$。

  \item 可以用若干方式推广前例。一种方式是令 $H$ 为任意阿贝尔群，并令 $K=\gen{x}\cong\ZZ_{2n}$ 为阶为 $2n$ 的循环群。仍定义 $\varphi$ 把 $x$ 送到取逆自同构，因此 $x^2$ 在 $H$ 上作为恒等自同构。在 $G$ 中，
  \[
    xhx^{-1}=h^{-1},\qquad x^2hx^{-2}=h,
  \]
  对所有 $h\in H$ 成立。因此 $x^2\in Z(G)$。特别地，若 $H=\ZZ_3$ 且 $K=\ZZ_4$，则 $G$ 是一个阶为 $12$ 的非阿贝尔群，并且不同构于 $A_4$ 或 $D_{12}$（因为它的 Sylow $2$-子群 $K$ 是阶为 $4$ 的循环群）。

  \item 继续前例，令 $H=\gen{h}\cong\ZZ_{2^n}$，令 $K=\gen{x}\cong\ZZ_4$，并在 $G$ 中令 $xhx^{-1}=h^{-1}$。如上，$x^2\in Z(G)$。由于 $x$ 取逆 $h$（即取逆 $H$），$x$ 也取逆 $H$ 中唯一的阶为 $2$ 的子群 $\gen z$，其中 $z=h^{2^{n-1}}$。于是
  \[
    xzx^{-1}=z^{-1}=z,
  \]
  所以 $x$ 中心化 $z$，从而 $z\in Z(G)$。于是 $x^2z\in Z(G)$，故 $\gen{x^2z}\normal G$。令
  \[
    \bar G=G/\gen{x^2z}.
  \]
  因为 $x^2$ 与 $z$ 是两个不同的、可交换的阶为 $2$ 的元素，$x^2z$ 的阶为 $2$，所以 $|\bar G|=|G|/2=2^{n+1}$。在取商形成 $\bar G$ 时，我们把 $x^2$ 与 $h^{2^{n-1}}$ 在商群中识别起来。特别地，当 $n=2$ 时，$x$ 与 $h$ 都有阶 $4$，$x$ 取逆 $h$，且 $h^2=x^2$，所以此时 $\bar G\cong Q_8$。一般地，可以检验 $\bar G$ 有唯一的阶为 $2$ 的子群（即 $\gen{x^2}$），并且它等于 $\bar G$ 的中心。群 $\bar G$ 称为阶为 $2^{n+1}$ 的\emph{广义四元数群}，记为 $Q_{2^{n+1}}$：
  \[
    Q_{2^{n+1}}=
    \gen{h,x\mid h^{2^n}=x^4=1,
    \ x^{-1}hx=h^{-1},
    \ h^{2^{n-1}}=x^2}.
  \]

  \item 令 $H=\QQ$（加法群），令 $K=\gen{x}\cong\ZZ$。定义 $\varphi$ 把 $x$ 送到 $H$ 上“乘以 $2$”的映射，所以 $x$ 在 $h\in H$ 上的作用为 $x\cdot h=2h$。注意乘以 $2$ 是 $H$ 的自同构，因为它有双边逆，即乘以 $1/2$。在群 $G$ 中，$\ZZ\leq\QQ$，并且 $x\ZZ x^{-1}$ 是 $\ZZ$ 的真子群（即 $2\ZZ$）。因此 $x\notin N_G(\ZZ)$，尽管 $x\ZZ x^{-1}\leq\ZZ$（注意 $x^{-1}\ZZ x$ 不包含于 $\ZZ$）。这表明，在无限群中，要证明元素 $g$ 正规化子群 $A$，一般不能只证明 $A$ 被 $g$ 共轭后包含在 $A$ 中（这在有限群中是充分的）。

  \item 对任意群 $H$，令 $K=\Aut(H)$，并令 $\varphi$ 为从 $K$ 到 $\Aut(H)$ 的恒等映射。半直积 $H\rtimes\Aut(H)$ 称为 $H$ 的\emph{全形群}，记为 $\Hol(H)$。下面描述一些全形群；这些同构的验证作为本章末尾的练习给出。
  \begin{enumerate}[label=(\alph*)]
    \item $\Hol(\ZZ_2\times\ZZ_2)\cong S_4$。
    \item 若 $|G|=n$，$\pi:G\to S_n$ 是左正则表示（第 4.2 节），则
    \[
      N_{S_n}(\pi(G))\cong \Hol(G).
    \]
    特别地，由于 $\ZZ_n$ 的一个生成元的左正则表示是 $S_n$ 中的一个 $n$-循环，于是对任意 $n$-循环 $(1\ 2\ \cdots\ n)$，有
    \[
      N_{S_n}\bigl(\gen{(1\ 2\ \cdots\ n)}\bigr)
      \cong \Hol(\ZZ_n)=\ZZ_n\rtimes\Aut(\ZZ_n).
    \]
    注意后一个群的阶为 $n\varphi(n)$。
  \end{enumerate}

  \item 令 $p,q$ 为素数，且 $p<q$，令 $H=\ZZ_q$，$K=\ZZ_p$。我们已经看到，若 $p$ 不整除 $q-1$，则每个阶为 $pq$ 的群都是循环群（见命题 4.16 后的例子）。这与如下事实一致：若 $p$ 不整除 $q-1$，则从 $\ZZ_p$ 到 $\Aut(\ZZ_q)$ 不存在非平凡同态（后者由命题 4.17 是阶为 $q-1$ 的循环群）。现在假设 $p\mid q-1$。由 Cauchy 定理，$\Aut(\ZZ_q)$ 含有一个阶为 $p$ 的子群（它唯一，因为 $\Aut(\ZZ_q)$ 是循环群）。因此存在从 $K$ 到 $\Aut(H)$ 的非平凡同态 $\varphi$。相应的群 $G=H\rtimes K$ 阶为 $pq$，且 $K$ 在 $G$ 中不正规（命题 11）。特别地，$G$ 是非阿贝尔群。我们马上将证明，在同构意义下，$G$ 是唯一的阶为 $pq$ 的非阿贝尔群。若 $p=2$，则 $G$ 必同构于 $D_{2q}$。

  \item 令 $p$ 为奇素数。我们构造两个不同构的阶为 $p^3$ 的非阿贝尔群（稍后将证明任意阶为 $p^3$ 的非阿贝尔群都同构于这两个之一）。

  首先令 $H=\ZZ_p\times\ZZ_p$，令 $K=\ZZ_p$。由命题 4.17，$\Aut(H)\cong\GL_2(\FF_p)$，并且
  \[
    |\GL_2(\FF_p)|=(p^2-1)(p^2-p).
  \]
  由于 $p$ 整除 $|\Aut(H)|$，由 Cauchy 定理，$H$ 有一个阶为 $p$ 的自同构。因此存在从 $K$ 到 $\Aut(H)$ 的非平凡同态 $\varphi$，所以相应的群 $H\rtimes K$ 是阶为 $p^3$ 的非阿贝尔群。更显式地，若 $H=\gen a\times\gen b$，且 $x$ 是 $K$ 的生成元，则令 $x$ 作用在 $a,b$ 上为
  \[
    x\cdot a=ab,
    \qquad
    x\cdot b=b,
  \]
  这定义了 $x$ 在整个 $H$ 上的作用。把 $H$ 看作二维 $\FF_p$-向量空间，以 $a,b$ 为基，则 $x$ 的作用（可用加法记号看作非奇异线性变换）具有矩阵
  \[
    \begin{pmatrix}1&0\\1&1\end{pmatrix}\in\GL_2(\FF_p).
  \]
  所得半直积有展示
  \[
    \gen{x,a,b\mid x^p=a^p=b^p=1,
    \ ab=ba,
    \ xax^{-1}=ab,
    \ xbx^{-1}=b}.
  \]
  事实上，这个群由 $\{x,a\}$ 生成，并称为 $\ZZ/p\ZZ$ 上的 Heisenberg 群（参见练习 25）。

  接着令 $H=\ZZ_{p^2}$，令 $K=\ZZ_p$。再次由命题 4.17，$\Aut(H)\cong\ZZ_{p(p-1)}$，所以 $H$ 有阶为 $p$ 的自同构。因此存在从 $K$ 到 $\Aut(H)$ 的非平凡同态 $\varphi$，从而群 $H\rtimes K$ 是阶为 $p^3$ 的非阿贝尔群。更显式地，若 $H=\gen y$，且 $x$ 是 $K$ 的生成元，则令 $x$ 作用在 $y$ 上为
  \[
    x\cdot y=y^{1+p}.
  \]
  所得半直积有展示
  \[
    \gen{x,y\mid x^p=y^{p^2}=1,
    \ xyx^{-1}=y^{1+p}}.
  \]
  这两个群不同构：前者没有阶为 $p^2$ 的元素（参见练习 25），而后者显然有这样的元素，即 $y$。

  \item 令 $H=Q_8\times(\ZZ_2\times\ZZ_2)=\gen{i,j}\times(\gen a\times\gen b)$，令 $K=\gen y\cong\ZZ_3$。由
  \[
    i\mapsto j\mapsto k=ij,
    \qquad
    a\mapsto b\mapsto ab
  \]
  定义的映射容易看出给出 $H$ 的一个阶为 $3$ 的自同构。令 $\varphi$ 为从 $K$ 到 $\Aut(H)$ 的同态，把 $y$ 送到这个自同构，并令 $G$ 为相应的半直积，使得 $y\in G$ 的作用为
  \[
    y\cdot i=j,
    \quad
    y\cdot j=k,
    \quad
    y\cdot a=b,
    \quad
    y\cdot b=ab.
  \]
  群 $G=H\rtimes K$ 是阶为 $96$ 的非阿贝尔群，并且元素 $i^2a\in G'$，但 $i^2a$ 不能写成任何单个交换子 $[x,y]$，其中 $x,y\in G$（检验后一断言是一个初等计算）。
\end{enumerate}

与直积情形一样，我们现在证明半直积的识别定理。

\medskip
\noindent\textbf{定理 12.} 设 $G$ 是群，$H,K$ 是其子群，满足
\begin{enumerate}[label=(\arabic*)]
  \item $H\normal G$；
  \item $H\cap K=1$。
\end{enumerate}
令 $\varphi:K\to\Aut(H)$ 为把 $k\in K$ 送到 $H$ 上由 $k$ 左共轭给出的自同构的同态。则
\[
  HK\cong H\rtimes_\varphi K.
\]
特别地，若 $G=HK$ 且 $H,K$ 满足 (1)、(2)，则 $G$ 是 $H$ 与 $K$ 的半直积。

\begin{proof}
因为 $H\normal G$，$HK$ 是 $G$ 的子群。由命题 8，$HK$ 的每个元素唯一写成 $hk$，其中 $h\in H,k\in K$。因此映射 $hk\mapsto(h,k)$ 是从 $HK$ 到 $H\rtimes K$ 的集合双射。它是同态这一事实正是本节开头推导半直积定义时的计算。
\end{proof}

\noindent\textbf{定义.} 设 $H$ 是群 $G$ 的子群。若 $K$ 是 $G$ 的子群，且 $G=HK$、$H\cap K=1$，则称 $K$ 是 $H$ 在 $G$ 中的一个\emph{补子群}。

用这个术语，识别半直积的准则就是：某个真正规子群必须存在一个补子群。并非每个群都是其两个真子群的半直积（例如单群），但我们已经看到，半直积的概念大大扩展了我们已知的群的列表。

\paragraph{一些分类}
现在我们应用定理 12 来分类若干阶 $n$ 的群。以下每个论证的基本思想是：
\begin{enumerate}[label=(\alph*)]
  \item 证明每个阶为 $n$ 的群都有真子群 $H,K$，满足定理 12 的假设并且 $G=HK$；
  \item 找出 $H,K$ 的所有可能同构类型；
  \item 对 (b) 中找到的每一对 $H,K$，找出所有可能的同态 $\varphi:K\to\Aut(H)$；
  \item 对 (c) 中找到的每个三元组 $H,K,\varphi$，形成半直积 $H\rtimes K$（因此任意阶为 $n$ 的群 $G$ 都同构于这些显式构造的群之一），并在所有这些半直积中判定哪些彼此同构。这样就得到阶为 $n$ 的群的互不同构类型列表。
\end{enumerate}
为了启动这个过程，我们首先必须在任意阶为 $n$ 的群 $G$ 中找到满足上述条件的子群 $H,K$。当 $n$ “较小”时，通常可由 Sylow 定理做到这一点。为了证明 $H$ 的正规性，我们使用 Sylow 定理的共轭部分，或第 4 章中建立的其他正规性准则（例如推论 4.5）。在下面许多例子中，$|H|$ 与 $|K|$ 互素，因此 $H\cap K=1$ 由 Lagrange 定理得到。

由于 $H,K$ 是 $G$ 的真子群，应当把 $H,K$ 的确定看作按归纳已经完成。在我们讨论的例子中，$H,K$ 的阶足够小，使得我们可由先前结果知道所有可能的同构类型。例如，在大多数情形中，$H,K$ 的阶是素数或素数平方。

在我们的例子中，可能的同态 $\varphi:K\to\Aut(H)$ 相对较少，特别是在考虑某些对称性之后（例如当 $K$ 循环时，把 $K$ 的一个生成元换成另一个）。

最后，在我们的例子中，这个过程得到的半直积数量很少，并且我们会发现它们大多两两不同构。一般来说，判定这些半直积之间的同构可能是更精细的问题，练习 4 就说明了这一点。

我们强调，把给定阶 $n$ 的每个群都“分解”为某个半直积的这种方法，并不适用于任意 $n$。例如，$Q_8$ 不是半直积，因为它没有真子群有补子群（虽然我们看到过它是一个半直积的商群）。经验上，当群阶 $n$ 不被任何素数的高次幂整除时，这个过程通常效果很好。另一方面，在另一个极端，当 $a$ 很大时，阶为 $p^a$ 的群中只有很小比例是非平凡半直积（$p$ 为素数）。

\paragraph{例：阶为 $pq$ 的群，其中 $p<q$ 为素数}
令 $G$ 是阶为 $pq$ 的任意群，令 $P\in\Syl_p(G)$，$Q\in\Syl_q(G)$。在 Sylow 定理应用的例 1 中，我们已经证明 $G\cong Q\rtimes P$，其中 $\varphi:P\to\Aut(Q)$。由于 $P,Q$ 阶为素数，它们都是循环群。群 $\Aut(Q)$ 是阶为 $q-1$ 的循环群。若 $p$ 不整除 $q-1$，则从 $P$ 到 $\Aut(Q)$ 的唯一同态是平凡同态，因此这种情形下唯一的半直积是直积，即 $G$ 是循环群。

现在考虑 $p\mid q-1$ 的情形，并令 $P=\gen y$。由于 $\Aut(Q)$ 是循环群，它含有唯一的阶为 $p$ 的子群，设为 $\gen\gamma$；任意同态 $\varphi:P\to\Aut(Q)$ 都必须把 $y$ 送到 $\gamma$ 的某个幂。因此有 $p$ 个同态 $\varphi_i:P\to\Aut(Q)$，由
\[
  \varphi_i(y)=\gamma^i,
  \qquad 0\leq i\leq p-1
\]
给出。由于 $\varphi_0$ 是平凡同态，$Q\rtimes_{\varphi_0}P\cong Q\times P$，同前。每个 $i\neq0$ 的 $\varphi_i$ 都给出一个非阿贝尔的阶为 $pq$ 的群 $G_i$。容易检验这些群都同构，因为对每个 $i>0$，存在 $P$ 的某个生成元 $y_i$，使得 $\varphi_i(y_i)=\gamma$。因此，除了 $P$ 的（任意）生成元选择之外，这些半直积都相同（参见练习 6；也参见第 4.3 节练习 28）。

\paragraph{例：阶为 $30$ 的群}
由 Sylow 定理后的例子，每个阶为 $30$ 的群 $G$ 都含有一个阶为 $15$ 的子群 $H$。由前例，$H$ 是循环群，并且 $H$ 在 $G$ 中正规（指数为 $2$）。由 Sylow 定理，$G$ 中存在一个阶为 $2$ 的子群 $K$。于是 $G=HK$ 且 $H\cap K=1$，所以 $G\cong H\rtimes K$，其中 $\varphi:K\to\Aut(H)$。由命题 4.16，
\[
  \Aut(\ZZ_{15})\cong(\ZZ/15\ZZ)^\times\cong\ZZ_4\times\ZZ_2.
\]
后一个同构可以直接计算，也可以使用前一节练习 11：把 $H$ 写成
\[
  H=\gen a\times\gen b\cong\ZZ_5\times\ZZ_3,
\]
由于这两个子群在 $H$ 中都是特征子群，有
\[
  \Aut(H)\cong\Aut(\ZZ_5)\times\Aut(\ZZ_3).
\]
特别地，$\Aut(H)$ 恰有三个阶为 $2$ 的元素，它们在 $H=\gen a\times\gen b$ 上的作用如下：
\[
\begin{array}{ccc}
 a\mapsto a, &\qquad& a\mapsto a^{-1}, \qquad a\mapsto a^{-1},\\
 b\mapsto b^{-1}, && b\mapsto b, \qquad b\mapsto b^{-1}.
\end{array}
\]
因此从 $K$ 到 $\Aut(H)$ 有三个非平凡同态，分别把 $K$ 的生成元送到这三个阶为 $2$ 的元素之一（照例，平凡同态给出直积 $H\times K\cong\ZZ_{30}$）。

令 $K=\gen k$。若同态 $\varphi_1:K\to\Aut(H)$ 把 $k$ 送到上面第一个自同构（即 $k\cdot a=a$ 且 $k\cdot b=b^{-1}$），则 $G_1=H\rtimes_{\varphi_1}K$ 容易看出同构于 $\ZZ_5\times D_6$（注意在这个半直积中，$k$ 中心化 $H$ 中阶为 $5$ 的元素 $a$，所以分解为直积 $\gen a\times\gen{b,k}$）。

若 $\varphi_2$ 把 $k$ 送到第二个自同构，则 $G_2=H\rtimes_{\varphi_2}K$ 容易看出同构于 $\ZZ_3\times D_{10}$（注意在这个半直积中，$k$ 中心化 $H$ 中阶为 $3$ 的元素 $b$，所以分解为直积 $\gen b\times\gen{a,k}$）。

若 $\varphi_3$ 把 $k$ 送到第三个自同构，则 $G_3=H\rtimes_{\varphi_3}K$ 容易看出同构于 $D_{30}$。

注意这些群两两不同构，因为它们中心的阶分别为 $30$（阿贝尔情形）、$5$（对 $G_1$）、$3$（对 $G_2$）与 $1$（对 $G_3$）。

我们强调，虽然事后看来，这个过程没有给出任何不能只用直积构造出来的群，但这个论证证明了这是阶为 $30$ 的群的完整同构类型列表。

\paragraph{例：阶为 $12$ 的群}
令 $G$ 是阶为 $12$ 的群，令 $V\in\Syl_2(G)$，$T\in\Syl_3(G)$。由第 4.5 节关于阶为 $12$ 的群的讨论，我们知道 $V$ 或 $T$ 在 $G$ 中正规（为了说明方法，我们不使用第 4 章所得结果的全部力度，即不直接使用“要么 $T\normal G$，要么 $G\cong A_4$”）。由 Lagrange 定理，$V\cap T=1$。因此 $G$ 是一个半直积。注意 $V\cong\ZZ_4$ 或 $\ZZ_2\times\ZZ_2$，而 $T\cong\ZZ_3$。

\emph{情形 1：$V\normal G$。}
我们必须确定从 $T$ 到 $\Aut(V)$ 的所有可能同态。若 $V\cong\ZZ_4$，则 $\Aut(V)\cong\ZZ_2$，从 $T$ 到 $\Aut(V)$ 没有非平凡同态。因此具有正规循环 Sylow $2$-子群的阶为 $12$ 的唯一群是 $\ZZ_{12}$。

于是设 $V\cong\ZZ_2\times\ZZ_2$。此时 $\Aut(V)\cong S_3$，并且 $\Aut(V)$ 有唯一的阶为 $3$ 的子群，设为 $\gen\gamma$。若 $T=\gen y$，则从 $T$ 到 $\Aut(V)$ 有三个可能同态：
\[
  \varphi_i:T\to\Aut(V),
  \qquad
  \varphi_i(y)=\gamma^i,
  \qquad i=0,1,2.
\]
照例，$\varphi_0$ 是平凡同态，给出直积 $\ZZ_2\times\ZZ_2\times\ZZ_3$。同态 $\varphi_1$ 与 $\varphi_2$ 给出同构的半直积，因为二者只差 $T$ 的生成元选择（即 $\varphi_1(y)=\gamma$，而 $\varphi_2(y')=\gamma$，其中 $y'=y^2$ 也是 $T$ 的生成元；也参见练习 6）。这种情形中的唯一非阿贝尔群是 $A_4$。

\emph{情形 2：$T\normal G$。}
我们必须确定从 $V$ 到 $\Aut(T)$ 的所有可能同态。注意 $\Aut(T)=\gen A\cong\ZZ_2$，其中 $A$ 取逆 $T$。若 $V=\gen x\cong\ZZ_4$，则从 $V$ 到 $\Aut(T)$ 恰有两个同态：平凡同态以及把 $x$ 送到 $A$ 的同态。照例，平凡同态给出直积 $\ZZ_3\times\ZZ_4\cong\ZZ_{12}$。非平凡同态给出本节命题 11 后例 2 中讨论的半直积。

最后，设 $V=\gen a\times\gen b\cong\ZZ_2\times\ZZ_2$。从 $V$ 到 $\Aut(T)$ 恰有三个非平凡同态，由指定其核为 $V$ 中三个阶为 $2$ 的子群之一来确定。例如，$\varphi_1(a)=A$ 且 $\varphi_1(b)=A$ 的核为 $\gen{ab}$；也就是说，在这个半直积中，$a,b$ 都通过取逆作用在 $T$ 上，而 $ab$ 中心化 $T$。若 $\varphi_2,\varphi_3$ 的核分别为 $\gen a$ 与 $\gen b$，则容易检验所得三个半直积都同构于 $S_3\times\ZZ_2$，其中 $\ZZ_2$ 直因子是 $\varphi_i$ 的核。例如，
\[
  T\rtimes_{\varphi_1}V=\gen{a,T}\times\gen{ab}.
\]
综上，阶为 $12$ 的群恰有 $5$ 个，其中 $3$ 个非阿贝尔。

\paragraph{例：阶为 $p^3$ 的群，$p$ 为奇素数}
令 $G$ 是阶为 $p^3$ 的群，$p$ 为奇素数，并假设 $G$ 非循环。由前一节练习 9，映射 $x\mapsto x^p$ 是从 $G$ 到 $Z(G)$ 的同态，其核的阶为 $p^2$ 或 $p^3$。在前一种情形中，$G$ 必含有阶为 $p^2$ 的元素；在后一种情形中，每个非单位元素的阶都是 $p$。

\emph{情形 1：$G$ 有阶为 $p^2$ 的元素。}
令 $x$ 是阶为 $p^2$ 的元素，并令 $H=\gen x$。注意由于 $H$ 的指数为 $p$，由推论 4.5，$H\normal G$。若 $E$ 是 $p$ 次幂映射的核，则在此情形中 $E\cong\ZZ_p\times\ZZ_p$，且 $E\cap H=\gen{x^p}$。令 $y$ 是 $E-H$ 中任意元素，并令 $K=\gen y$。由构造，$H\cap K=1$，所以
\[
  G\cong \ZZ_{p^2}\rtimes\ZZ_p,
\]
其中 $\varphi:K\to\Aut(H)$。若 $\varphi$ 是平凡同态，则 $G\cong\ZZ_{p^2}\times\ZZ_p$，因此只需考虑非平凡同态。由命题 4.17，$\Aut(H)\cong\ZZ_{p(p-1)}$ 是循环群，所以含有唯一的阶为 $p$ 的子群，显式地为 $\gen\gamma$，其中
\[
  \gamma(x)=x^{1+p}.
\]
照例，至多差循环群 $K$ 的生成元选择，只有一个非平凡同态 $\varphi:K\to\Aut(H)$，由 $\varphi(y)=\gamma$ 给出；因此在这种情形中，在同构意义下只有一个非阿贝尔群 $H\rtimes K$。这个群已在上面的例 7 中描述。

\emph{情形 2：每个非单位元素的阶都是 $p$。}
在此情形中，令 $H$ 为 $G$ 的任意阶为 $p^2$ 的子群（见第 4.3 节练习 29）。必有 $H\cong\ZZ_p\times\ZZ_p$。令 $K=\gen y$，其中 $y$ 是 $G-H$ 的任意元素。由于 $H$ 的指数为 $p$，$H\normal G$；又因为 $K$ 阶为 $p$ 但不包含于 $H$，$H\cap K=1$。于是
\[
  G\cong(\ZZ_p\times\ZZ_p)\rtimes\ZZ_p,
\]
其中 $\varphi:K\to\Aut(H)$。若 $\varphi$ 平凡，则 $G\cong\ZZ_p\times\ZZ_p\times\ZZ_p$（初等阿贝尔群），所以可假设 $\varphi$ 非平凡。由命题 4.17，
\[
  \Aut(H)\cong\GL_2(\FF_p),
\]
因此 $|\Aut(H)|=(p^2-1)(p^2-p)$。注意 $\Aut(H)$ 的 Sylow $p$-子群阶为 $p$，所以由 Sylow 定理，$\Aut(H)$ 中所有阶为 $p$ 的子群在 $\Aut(H)$ 中共轭。显式地（如上面例 7 所讨论的），$\Aut(H)$ 中每个阶为 $p$ 的子群都共轭于 $\gen\gamma$，其中若 $H=\gen a\times\gen b$，则自同构 $\gamma$ 定义为
\[
  \gamma(a)=ab,
  \qquad
  \gamma(b)=b.
\]
关于二维向量空间 $H$ 的 $\FF_p$-基 $a,b$，这个自同构具有矩阵
\[
  \begin{pmatrix}1&0\\1&1\end{pmatrix}\in\GL_2(\FF_p).
\]
因此（再次引用练习 6），在这种情形中只有一个半直积同构类型。

最后，由于这两个非阿贝尔群的 $p$ 次幂映射的核具有不同阶，它们不同构。这个群的一个展示也在上面的例 7 中给出。

\subsection*{练习}
设 $H,K$ 是群，$\varphi:K\to\Aut(H)$ 是同态，并照例把 $H,K$ 视为 $G=H\rtimes_\varphi K$ 的子群。
\begin{enumerate}[label=\arabic*.]
  \item 证明 $C_K(H)=\ker\varphi$（回忆 $C_K(H)=C_G(H)\cap K$）。
  \item 证明 $C_H(K)=N_H(K)$。
  \item 在命题 11 后例 1 中，证明 $G-H$ 的每个元素都有阶 $2$。证明 $G$ 是阿贝尔群当且仅当 $h^2=1$ 对所有 $h\in H$ 成立。
  \item 令 $p=2$，检验阶为 $p^3$ 的两个非阿贝尔群的构造在此情形中仍有效。证明所得两个群都同构于 $D_8$。
  \item 令 $G=\Hol(\ZZ_2\times\ZZ_2)$。
  \begin{enumerate}[label=(\alph*)]
    \item 证明 $G=H\rtimes K$，其中 $H=\ZZ_2\times\ZZ_2$ 且 $K\cong S_3$。推出 $|G|=24$。
    \item 证明 $G\cong S_4$。（提示：让 $G$ 作用在 $K$ 的左陪集上，得到到 $S_4$ 的同态；用练习 1 证明该表示忠实。）
  \end{enumerate}
  \item 假设 $K$ 是循环群，$H$ 是任意群，$\varphi_1,
\varphi_2$ 是从 $K$ 到 $\Aut(H)$ 的同态，且 $\varphi_1(K)$ 与 $\varphi_2(K)$ 是 $\Aut(H)$ 中共轭子群。若 $K$ 无限，还假设 $\varphi_1,
\varphi_2$ 单射。通过构造显式同构证明
  \[
    H\rtimes_{\varphi_1}K\cong H\rtimes_{\varphi_2}K.
  \]

  \item 本题描述阶为 $56$ 的十三个同构类型。
  \begin{enumerate}[label=(\alph*)]
    \item 证明阶为 $56$ 的阿贝尔群有三个。
    \item 证明每个阶为 $56$ 的群都有正规 Sylow $2$-子群或正规 Sylow $7$-子群。
    \item 构造下列具有正规 Sylow $7$-子群的阶为 $56$ 的非阿贝尔群，其中 Sylow $2$-子群 $S$ 如下：当 $S\cong\ZZ_2\times\ZZ_2\times\ZZ_2$ 时一个群；当 $S\cong\ZZ_4\times\ZZ_2$ 时两个不同构群；当 $S\cong\ZZ_8$ 时一个群；当 $S\cong Q_8$ 时两个不同构群；当 $S\cong D_8$ 时三个不同构群。
    \item 设 $G$ 是阶为 $56$ 且 Sylow $7$-子群非正规的群。证明若 $S$ 是 $G$ 的 Sylow $2$-子群，则 $S\cong\ZZ_2\times\ZZ_2\times\ZZ_2$。
    \item 证明存在唯一一个 Sylow $7$-子群非正规的阶为 $56$ 的群。（存在性可用 $|\GL_3(\FF_2)|=168$，唯一性用练习 6。）
  \end{enumerate}

  \item 构造一个阶为 $75$ 的非阿贝尔群。分类所有阶为 $75$ 的群（共有三个）。

  \item 证明矩阵
  \[
    \begin{pmatrix}0&-1\\1&4\end{pmatrix}
  \]
  是 $\GL_2(\FF_{19})$ 中阶为 $5$ 的元素。利用这个矩阵构造一个阶为 $1805$ 的非阿贝尔群，并给出该群的展示。分类阶为 $1805$ 的群（共有三个同构类型）。

  \item 本题分类阶为 $147$ 的群（共有六个同构类型）。
  \begin{enumerate}[label=(\alph*)]
    \item 证明阶为 $147$ 的阿贝尔群有两个。
    \item 证明每个阶为 $147$ 的群都有正规 Sylow $7$-子群。
    \item 证明 Sylow $7$-子群为循环群的非阿贝尔群唯一。
    \item 令
    \[
      t_1=\begin{pmatrix}2&0\\0&1\end{pmatrix},
\qquad
      t_2=\begin{pmatrix}1&0\\0&2\end{pmatrix}
    \]
    为 $\GL_2(\FF_7)$ 的元素。证明 $P=\gen{t_1,t_2}$ 是 $\GL_2(\FF_7)$ 的 Sylow $3$-子群，且 $P\cong\ZZ_3\times\ZZ_3$。推出 $\GL_2(\FF_7)$ 的每个阶为 $3$ 的子群都共轭于 $P$ 的某个子群。
    \item 由第 1 节例 3，$P$ 有四个阶为 $3$ 的子群，分别为 $P_1=\gen{t_1}$、$P_2=\gen{t_2}$、$P_3=\gen{t_1t_2}$ 与 $P_4=\gen{t_1t_2^2}$。对 $i=1,2,3,4$，令 $G_i=(\ZZ_7\times\ZZ_7)\rtimes_{\varphi_i}\ZZ_3$，其中 $\varphi_i$ 是 $\ZZ_3$ 到 $\Aut(\ZZ_7\times\ZZ_7)$ 的子群 $P_i$ 上的同构。对每个 $i$，用生成元与关系描述 $G_i$。推出 $G_1\cong G_2$。
    \item 证明 $G_1$ 不同构于 $G_3$ 或 $G_4$。（提示：$G_1$ 的中心阶为 $7$，而 $G_3,G_4$ 的中心平凡。）
    \item 证明 $G_3\not\cong G_4$。（提示：$G_3$ 中每个阶为 $7$ 的子群都正规，而 $G_4$ 有不正规的阶为 $7$ 的子群。）
    \item 证明上述六个非同构群就是所有阶为 $147$ 的群。
  \end{enumerate}

  \item 分类阶为 $28$ 的群（共有四个同构类型）。
  \item 分类阶为 $20$ 的群（共有五个同构类型）。
  \item 分类阶为 $4p$ 的群，其中 $p>3$ 为素数。（若 $p\equiv3\pmod4$ 有四个同构类型；若 $p\equiv1\pmod4$ 有五个同构类型。）
  \item 本题分类阶为 $60$ 的群（共有十三个同构类型）。令 $G$ 是阶为 $60$ 的群，$P$ 是 Sylow $5$-子群，$Q$ 是 Sylow $3$-子群。
  \begin{enumerate}[label=(\alph*)]
    \item 证明若 $P$ 在 $G$ 中不正规，则 $G\cong A_5$。
    \item 证明若 $P\normal G$ 但 $Q$ 不正规，则 $G\cong A_4\times\ZZ_5$。
    \item 证明若 $P,Q$ 都正规，则 $G\cong\ZZ_{15}\rtimes T$，其中 $T\cong\ZZ_4$ 或 $\ZZ_2\times\ZZ_2$。说明此时当 $T$ 循环时有六个同构类型（其中一个阿贝尔），当 $T$ 是 Klein 四元群时有五个同构类型（其中一个阿贝尔）。
  \end{enumerate}
  \item 设 $p$ 是奇素数。证明 $\GL_2(\FF_p)$ 中每个阶为 $2$ 的元素都共轭于一个对角元为 $\pm1$ 的对角矩阵。分类阶为 $2p^2$ 的群。
  \item 证明从 $\ZZ_2$ 到 $\Aut(\ZZ_8)$ 恰有 $4$ 个不同同态。证明由此得到的半直积分别是 $\ZZ_8\times\ZZ_2$、$D_{16}$、拟二面体群 $QD_{16}$ 与模群 $M$。
  \item 证明对任意 $n\geq3$，从 $\ZZ_2$ 到 $\Aut(\ZZ_{2^n})$ 恰有 $4$ 个不同同态。证明所得半直积给出 $4$ 个阶为 $2^{n+1}$ 的非同构群。
  \item 证明若 $H$ 是任意群，则存在群 $G$，其中 $H\normal G$，并且 $H$ 的每个自同构都可实现为 $G$ 的某个内自同构在 $H$ 上的限制。
  \item 令 $H$ 是阶为 $n$ 的群，$K=\Aut(H)$，构造 $G=\Hol(H)=H\rtimes K$。令 $G$ 作用在 $K$ 的左陪集上，并令 $\pi:G\to S_n$ 为相应置换表示。
  \begin{enumerate}[label=(\alph*)]
    \item 证明 $H$ 的元素是 $K$ 在 $G$ 中左陪集的一组代表元，并且在此代表元选择下，$\pi$ 限制到 $H$ 是 $H$ 的正则表示。
    \item 证明 $\pi(G)$ 是 $S_n$ 中 $\pi(H)$ 的正规化子。推出任意有限群 $H$ 的正则表示像在 $S_n$ 中的正规化子同构于 $\Hol(H)$。
    \item 推出 $S_n$ 中由一个 $n$-循环生成的群的正规化子同构于 $\Hol(\ZZ_n)$，阶为 $n\varphi(n)$。
  \end{enumerate}
  \item 设 $p$ 是奇素数。证明若 $P$ 是非循环 $p$-群，则 $P$ 含有一个正规子群 $U\cong\ZZ_p\times\ZZ_p$。推出对奇素数 $p$，若一个 $p$-群只含有唯一的阶为 $p$ 的子群，则它是循环群。
  \item 设 $p$ 是奇素数，$P$ 是 $p$-群。证明若 $P$ 的每个子群都正规，则 $P$ 是阿贝尔群。
  \item 设 $F$ 是域，$n$ 是正整数，$G$ 是 $\GL_n(F)$ 中上三角矩阵群。
  \begin{enumerate}[label=(\alph*)]
    \item 证明 $G$ 是半直积 $U\rtimes D$，其中 $U$ 是对角线上全为 $1$ 的上三角矩阵群，$D$ 是 $\GL_n(F)$ 中的对角矩阵群。
    \item 当 $n=2$ 时，描述 $D$ 到 $\Aut(U)$ 的同态。
  \end{enumerate}
  \item 设 $K,L$ 是群，$n$ 是正整数，$\rho:K\to S_n$ 是同态，$H$ 是 $n$ 个 $L$ 的直积。第 1 节练习 8 构造了同态 $S_n\to\Aut(H)$，使 $S_n$ 置换 $H$ 的 $n$ 个因子。复合给出 $K\to\Aut(H)$ 的同态。$L$ 关于 $K$ 的\emph{花环积}是相应的半直积 $H\rtimes K$，记为 $L\wr K$。
  \begin{enumerate}[label=(\alph*)]
    \item 若 $K,L$ 有限且 $\rho$ 是 $K$ 的左正则表示，用 $|K|,|L|$ 表示 $|L\wr K|$。
    \item 令 $p$ 为素数，$K=L=\ZZ_p$，且 $\rho$ 为 $K$ 的左正则表示。证明 $\ZZ_p\wr\ZZ_p$ 是阶为 $p^{p+1}$ 的非阿贝尔群，并同构于 $S_{p^2}$ 的一个 Sylow $p$-子群。
  \end{enumerate}
  \item 设 $n>1$。证明如下分类：阶为 $n$ 的每个群都是阿贝尔群，当且仅当
  \[
    n=p_1^{\alpha_1}p_2^{\alpha_2}\cdots p_r^{\alpha_r},
  \]
  其中 $p_1,
\ldots,p_r$ 为互异素数，$\alpha_i=1$ 或 $2$，且对所有 $i,j$，$p_i\nmid p_j^{\alpha_j}-1$。
  \item 令 $H(\FF_p)$ 为有限域 $\FF_p=\ZZ/p\ZZ$ 上的 Heisenberg 群。证明 $H(\FF_2)\cong D_8$，并且 $H(\FF_p)$ 的指数为 $p$，且同构于例 7 中的第一个非阿贝尔群。
\end{enumerate}
% \chapter{群论中的进一步主题}

\section{$p$-群、幂零群与可解群}

设 $p$ 是素数，$G$ 是阶为 $p^a n$ 的有限群，其中 $p\nmid n$。回忆，一个（有限）$p$-群是指阶为 $p$ 的幂的群。Sylow 定理说明，$p$-群大量出现在 $G$ 的子群中；为了利用这一现象来剖析有限群的结构，我们需要先建立 $p$-群的一些基本性质。下一节中我们将在许多具体情形中应用这些结果。

在给出关于 $p$-群的结果之前，先回忆一个在前面练习中已经出现过的定义。

\medskip
\noindent\textbf{定义.} 群 $G$ 的一个\emph{极大子群}是指 $G$ 的一个真子群 $M$，使得不存在 $G$ 的子群 $H$ 满足
\[
  M < H < G.
\]

由阶数考虑可知，有限群的每个真子群都包含在某个极大子群中。相反，无限群可以有极大子群，也可以没有极大子群。例如，$p\ZZ$ 是 $\ZZ$ 的极大子群，而 $\QQ$（加法群）没有极大子群（参见本节末练习 16）。

我们现在把以后需要的 $p$-群性质集中到一个总括性定理中。

\medskip
\noindent\textbf{定理 1.} 设 $p$ 是素数，$P$ 是阶为 $p^a$ 的群，其中 $a\geq 1$。则：
\begin{enumerate}[label=(\arabic*)]
  \item $P$ 的中心非平凡：$Z(P)\neq 1$。
  \item 若 $H$ 是 $P$ 的非平凡正规子群，则 $H$ 与中心的交非平凡：$H\cap Z(P)\neq 1$。特别地，每个阶为 $p$ 的正规子群都包含在中心中。
  \item 若 $H\normal P$，则对 $|H|$ 的每个因子 $p^b$，$H$ 含有一个阶为 $p^b$ 且在 $P$ 中正规的子群。特别地，对每个 $b\in\{0,1,\ldots,a\}$，$P$ 都有一个阶为 $p^b$ 的正规子群。
  \item 若 $H<P$，则 $H<N_P(H)$，即 $P$ 的每个真子群都是其在 $P$ 中正规化子的真子群。
  \item $P$ 的每个极大子群都具有指标 $p$，并且在 $P$ 中正规。
\end{enumerate}

\begin{proof}
这些结果归根结底依赖类方程，读者可先回顾第 4.3 节。

(1) 是第 4 章定理 8；它也是 (2) 在 $H=P$ 时的特例。我们先证明 (2)，尽管下面的论证只是第 4 章论证的轻微推广，但这里不直接引用第 4 章定理 8。设 $H$ 是 $P$ 的非平凡正规子群。回忆：对 $P$ 的每个共轭类 $C$，由于 $H$ 正规，或者 $C\subseteq H$，或者 $C\cap H=\varnothing$。取 $P$ 的共轭类代表元
\[
  a_1,a_2,\ldots,a_r,
\]
其中 $a_1,\ldots,a_k\in H$，而 $a_{k+1},\ldots,a_r\notin H$。令 $C_i$ 为 $a_i$ 在 $P$ 中的共轭类。于是
\[
  C_i\subseteq H\quad (1\leq i\leq k),
  \qquad
  C_i\cap H=\varnothing\quad (k+1\leq i\leq r).
\]
必要时重新编号 $a_1,\ldots,a_k$，可假设 $a_1,\ldots,a_s$ 代表大小为 $1$ 的共轭类（即它们在 $Z(P)$ 中），而 $a_{s+1},\ldots,a_k$ 代表大小大于 $1$ 的共轭类。由于 $H$ 是这些类的不交并，
\[
  |H|=|H\cap Z(P)|+\sum_{i=s+1}^{k}|P:C_P(a_i)|.
\]
现在 $p$ 整除 $|H|$，并且 $p$ 整除和式中的每一项 $|P:C_P(a_i)|$，所以 $p$ 整除它们的差 $|H\cap Z(P)|$。因此 $H\cap Z(P)\neq 1$。若 $|H|=p$，则由 $H\cap Z(P)\neq 1$ 得 $H\leq Z(P)$。这证明了 (2)。

接着用关于 $a$ 的归纳法证明 (3)。若 $a=1$ 或 $H=1$，结论平凡。设 $a>1$ 且 $H\neq 1$。由 (2)，$H\cap Z(P)\neq 1$；由 Cauchy 定理，$H\cap Z(P)$ 含有一个阶为 $p$ 的（正规）子群 $Z$。用横线表示取商 $P/Z$。这个商群阶为 $p^{a-1}$，且 $\overline H\normal \overline P$。由归纳假设，对每个使 $p^b$ 整除 $|\overline H|$ 的非负整数 $b$，$\overline H$ 中存在阶为 $p^b$ 且在 $\overline P$ 中正规的子群 $\overline K$。若 $K$ 是 $\overline K$ 在 $P$ 中的完全原像，则 $|K|=p^{b+1}$。由这种过程得到的 $H$ 的所有子群，再加上恒等子群，便给出 $|H|$ 的每个因子阶的、在 $P$ 中正规的子群。(3) 的第二个断言是取 $H=P$ 的特例。

(4) 也用关于 $|P|$ 的归纳法证明。若 $P$ 是阿贝尔群，则 $P$ 的所有子群都在 $P$ 中正规，结论平凡。因此可假设 $|P|>p$（事实上由推论 4.9 可知 $|P|>p^2$）。设 $H$ 是 $P$ 的真子群。由于 $Z(P)$ 的每个元素都与 $P$ 的每个元素交换，所以 $Z(P)$ 正规化 $P$ 的每个子群。由 (1)，$Z(P)\neq 1$。若 $Z(P)\nleq H$，则 $H$ 真包含于 $\langle H,Z(P)\rangle$，而后者又包含在 $N_P(H)$ 中，于是 (4) 成立。故可假设 $Z(P)\leq H$。用横线表示取商 $P/Z(P)$。由于 $\overline P$ 的阶比 $P$ 小，由归纳假设，$\overline H$ 真包含于 $N_{\overline P}(\overline H)$。由格同构定理可知，$N_P(H)$ 是 $N_{\overline P}(\overline H)$ 在 $P$ 中的完全原像，因此在这种情形下也得到 $H$ 真包含于其正规化子。

最后证明 (5)。设 $M$ 是 $P$ 的极大子群。由定义 $M<P$，所以由 (4)，$M<N_P(M)$。由极大性，必有 $N_P(M)=P$，即 $M\normal P$。格同构定理说明 $P/M$ 是一个没有非平凡真子群的 $p$-群；但由 (3)，$P/M$ 对每个整除 $|P/M|$ 的阶都有子群。因此唯一可能是 $|P/M|=p$。证明完毕。
\end{proof}

\medskip
\noindent\textbf{定义.}
\begin{enumerate}[label=(\arabic*)]
  \item 对任意（有限或无限）群 $G$，归纳定义如下子群：
  \[
    Z_0(G)=1,
    \qquad
    Z_1(G)=Z(G),
  \]
  并且 $Z_{i+1}(G)$ 是 $G$ 中包含 $Z_i(G)$ 的子群，使得
  \[
    Z_{i+1}(G)/Z_i(G)=Z(G/Z_i(G)),
  \]
  即 $Z_{i+1}(G)$ 是自然投影下 $G/Z_i(G)$ 的中心在 $G$ 中的完全原像。子群链
  \[
    Z_0(G)\leq Z_1(G)\leq Z_2(G)\leq\cdots
  \]
  称为 $G$ 的\emph{上中心列}。“上”这个词表示 $Z_i(G)\leq Z_{i+1}(G)$。
  \item 若存在 $c\in\ZZ$ 使 $Z_c(G)=G$，则称群 $G$ 是\emph{幂零的}。这样的最小 $c$ 称为 $G$ 的\emph{幂零类}。
\end{enumerate}

本节末的一个练习说明，对所有 $i$，$Z_i(G)$ 都是 $G$ 的特征子群（从而正规）。以后我们将自由使用这一事实。

\paragraph{注}
\begin{enumerate}[label=(\arabic*)]
  \item 若 $G$ 是阿贝尔群，则 $G$ 是幂零的（当 $|G|>1$ 时幂零类为 $1$），因为此时 $G=Z(G)=Z_1(G)$。应把幂零群看作位于阿贝尔群与可解群之间的一类结构：
  \[
    \text{循环群}\subset \text{阿贝尔群}\subset \text{幂零群}\subset \text{可解群}\subset \text{所有群}.
  \]
  上述包含关系均为真包含，稍后将验证。

  \item 对任何有限群，由阶数考虑，必存在整数 $n$ 使得
  \[
    Z_n(G)=Z_{n+1}(G)=Z_{n+2}(G)=\cdots.
  \]
  例如，对所有 $n\in\ZZ^+$，$Z_n(S_3)=1$。一旦上中心列中相邻两项相同，这条链就从该处稳定。例如，若 $G=\ZZ_2\times S_3$，则对所有 $n$，
  \[
    Z(G)=Z_1(G)=Z_2(G)=Z_n(G)
  \]
  的阶都是 $2$。按定义，对非幂零群，所有 $Z_n(G)$ 都是 $G$ 的真子群。

  \item 对无限群 $G$，可能所有 $Z_i(G)$ 都是 $G$ 的真子群（因此 $G$ 不是幂零群），但
  \[
    G=\bigcup_{i=0}^{\infty} Z_i(G).
  \]
  满足这种性质的群称为\emph{超幂零群}；它们具有幂零群的一些性质，但并非全部。虽然我们主要处理有限幂零群，但不涉及阶、Sylow 子群等概念的结果也适用于无限群。即使对无限群，处理幂零群的主要技巧之一也是关于幂零类的归纳法。
\end{enumerate}

\medskip
\noindent\textbf{命题 2.} 设 $p$ 是素数，$P$ 是阶为 $p^a$ 的群。则 $P$ 是幂零群，并且幂零类至多为 $a-1$。

\begin{proof}
对每个 $i\geq 0$，$P/Z_i(P)$ 是 $p$-群。因此，由定理 1(1)，
\[
  |P/Z_i(P)|>1\quad\Longrightarrow\quad Z(P/Z_i(P))\neq 1.
\]
于是若 $Z_i(P)\neq P$，则 $|Z_{i+1}(P)|\geq p|Z_i(P)|$，所以 $|Z_{i+1}(P)|\geq p^{i+1}$。特别地，$|Z_a(P)|\geq p^a$，故 $P=Z_a(P)$，所以 $P$ 的幂零类至多为 $a$。

$P$ 的幂零类恰等于 $a$ 的唯一可能情形是 $|Z_i(P)|=p^i$ 对所有 $i$ 成立。此时 $Z_{a-2}(P)$ 在 $P$ 中指标为 $p^2$，故 $P/Z_{a-2}(P)$ 是阿贝尔群（由推论 4.9）。于是 $P/Z_{a-2}(P)$ 等于其中心，从而 $Z_{a-1}(P)=P$，矛盾。因此 $P$ 的幂零类至多为 $a-1$。
\end{proof}

\paragraph{例} $D_8$ 与 $Q_8$ 都是幂零类为 $2$ 的幂零群。更一般地，$D_{2^n}$ 的幂零类为 $n-1$。这可通过归纳证明：对 $n\geq 3$，$|Z(D_{2^n})|=2$ 且 $D_{2^n}/Z(D_{2^n})\cong D_{2^{n-1}}$（细节留作练习）。若 $n$ 不是 $2$ 的幂，则 $D_{2n}$ 不是幂零群（参见练习 10）。

下面给出有限群幂零性的若干等价（也常常更可操作的）刻画。

\medskip
\noindent\textbf{定理 3.} 设 $G$ 是有限群，$p_1,p_2,\ldots,p_s$ 是整除 $|G|$ 的不同素数，且 $P_i\in\Syl_{p_i}(G)$，$1\leq i\leq s$。则以下条件等价：
\begin{enumerate}[label=(\arabic*)]
  \item $G$ 是幂零群；
  \item 若 $H<G$，则 $H<N_G(H)$，即 $G$ 的每个真子群都是其在 $G$ 中正规化子的真子群；
  \item 对 $1\leq i\leq s$，$P_i\normal G$，即每个 Sylow 子群都在 $G$ 中正规；
  \item $G\cong P_1\times P_2\times\cdots\times P_s$。
\end{enumerate}

\begin{proof}
(1) 推出 (2) 的证明与 $p$-群情形相同；所需的唯一事实是：若 $G$ 幂零，则 $G/Z(G)$ 也幂零。因此细节略去（参见练习）。

为证 (2) 推出 (3)，令 $P=P_i$，$N=N_G(P)$。由于 $P\normal N$，由推论 4.20，$P$ 是 $N$ 的特征子群。因为 $P\charac N\normal N_G(N)$，得到 $P\normal N_G(N)$。这意味着 $N_G(N)\leq N$，从而 $N_G(N)=N$。由 (2)，必有 $N=G$，即 $P\normal G$。

接着证明 (3) 推出 (4)。对任意 $t$，$1\leq t\leq s$，用归纳法证明
\[
  P_1P_2\cdots P_t\cong P_1\times P_2\times\cdots\times P_t.
\]
由于每个 $P_i$ 都在 $G$ 中正规，$P_1\cdots P_t$ 是 $G$ 的子群。令 $H=P_1\cdots P_{t-1}$，$K=P_t$。由归纳假设，$H\cong P_1\times\cdots\times P_{t-1}$，特别地 $|H|=|P_1|\cdots |P_{t-1}|$。而 $|K|=|P_t|$，所以 $|H|$ 与 $|K|$ 互素；由 Lagrange 定理，$H\cap K=1$。又 $P_1\cdots P_t=HK$，由直积识别定理（定理 5.9），
\[
  HK\cong H\times K\cong (P_1\times\cdots\times P_{t-1})\times P_t.
\]
取 $t=s$ 得 (4)。

最后证明 (4) 推出 (1)。由第 5.1 节练习 1，
\[
  Z(P_1\times\cdots\times P_s)\cong Z(P_1)\times\cdots\times Z(P_s).
\]
由第 5.1 节练习 14，
\[
  G/Z(G)\cong (P_1/Z(P_1))\times\cdots\times(P_s/Z(P_s)).
\]
因此 (4) 的假设也适用于 $G/Z(G)$。由定理 1，若 $P_i\neq 1$ 则 $Z(P_i)\neq 1$，所以若 $G\neq 1$，则 $|G/Z(G)|<|G|$。由归纳法，$G/Z(G)$ 幂零；再由练习 6，$G$ 幂零。证明完毕。
\end{proof}

由上一定理立即得到有限阿贝尔群基本定理第一部分（第 5.2 节定理 5）。稍后我们还会由中国剩余定理给出另一个证明。

\medskip
\noindent\textbf{推论 4.} 有限阿贝尔群是它的 Sylow 子群的直积。

下面的命题稍后会用来证明有限域的乘法群是循环群（无需使用有限阿贝尔群基本定理）。

\medskip
\noindent\textbf{命题 5.} 若有限群 $G$ 满足：对每个整除 $|G|$ 的正整数 $n$，方程 $x^n=1$ 在 $G$ 中至多有 $n$ 个解，则 $G$ 是循环群。

\begin{proof}
设 $|G|=p_1^{a_1}\cdots p_s^{a_s}$，并令 $P_i\in\Syl_{p_i}(G)$。由于 $p_i^{a_i}\mid |G|$ 且 $P_i$ 的 $p_i^{a_i}$ 个元素都是 $x^{p_i^{a_i}}=1$ 的解，由假设 $P_i$ 必含有 $G$ 中该方程的所有解。于是 $P_i$ 是唯一的（故正规）Sylow $p_i$-子群。由定理 3，$G$ 是其 Sylow 子群的直积。

由定理 1，每个 $P_i$ 有一个指标为 $p_i$ 的正规子群 $M_i$。由于 $|M_i|=p_i^{a_i-1}$，且 $G$ 中至多有 $p_i^{a_i-1}$ 个元素满足 $x^{p_i^{a_i-1}}=1$，由 Lagrange 定理可知，$M_i$ 含有所有满足该方程的元素。因此 $P_i\setminus M_i$ 中的任一元素都满足 $x^{p_i^{a_i}}=1$ 但 $x^{p_i^{a_i-1}}\neq 1$，也就是说它的阶为 $p_i^{a_i}$。故每个 $P_i$ 都是循环群。$G$ 是两两互素阶循环群的直积，因而是循环群。
\end{proof}

下一个命题称为 Frattini 论证。我们将用它给出有限幂零群的另一个刻画，它在下一节也很有用。

\medskip
\noindent\textbf{命题 6（Frattini 论证）.} 设 $G$ 是有限群，$H\normal G$，$P$ 是 $H$ 的 Sylow $p$-子群。则
\[
  G=HN_G(P),
\]
并且 $|G:H|$ 整除 $|N_G(P)|$。

\begin{proof}
由推论 3.15，$HN_G(P)$ 是 $G$ 的子群，且由于 $H\normal G$，有 $HN_G(P)=N_G(P)H$。设 $g\in G$。由于 $P^g\leq H^g=H$，$P$ 与 $P^g$ 都是 $H$ 的 Sylow $p$-子群。由 Sylow 定理应用于 $H$，存在 $x\in H$ 使 $P^g=P^x$。于是 $gx^{-1}\in N_G(P)$，故 $g\in N_G(P)x$。因为 $g$ 任意，得到 $G=N_G(P)H$。

对 $G=N_G(P)H$ 应用第二同构定理，得到
\[
  |G:H|=|N_G(P):N_G(P)\cap H|,
\]
所以 $|G:H|$ 整除 $|N_G(P)|$。证明完毕。
\end{proof}

\medskip
\noindent\textbf{命题 7.} 有限群幂零当且仅当它的每个极大子群都是正规子群。

\begin{proof}
设 $G$ 是有限幂零群，$M$ 是 $G$ 的极大子群。与定理 1 的证明类似，由定理 3(2) 有 $M<N_G(M)$，而 $M$ 的极大性迫使 $N_G(M)=G$，即 $M\normal G$。

反过来，假设有限群 $G$ 的每个极大子群都是正规子群。令 $P$ 是 $G$ 的一个 Sylow $p$-子群。我们证明 $P\normal G$，然后由定理 3(3) 得 $G$ 幂零。若 $P$ 在 $G$ 中不正规，取一个包含 $N_G(P)$ 的极大子群 $M$。由假设 $M\normal G$，于是由 Frattini 论证，$G=MN_G(P)$。但 $N_G(P)\leq M$，故 $MN_G(P)=M$，矛盾。证明完毕。
\end{proof}

\subsection*{换位子与下中心列}

为完整起见，我们给出群的下中心列的定义，并说明它与上中心列的关系。由于以后不会使用这些结果，证明留作（直接的）练习。

回忆，群 $G$ 中两个元素 $x,y$ 的\emph{换位子}定义为
\[
  [x,y]=x^{-1}y^{-1}xy,
\]
而 $G$ 的两个子群 $H,K$ 的换位子子群为
\[
  [H,K]=\langle [h,k]\mid h\in H,\ k\in K\rangle.
\]
换位子及换位子子群的基本性质已在第 5.4 节建立。

\medskip
\noindent\textbf{定义.} 对任意（有限或无限）群 $G$，归纳定义
\[
  G^0=G,
  \qquad
  G^1=[G,G],
  \qquad
  G^{i+1}=[G,G^i].
\]
子群链
\[
  G=G^0\geq G^1\geq G^2\geq\cdots
\]
称为 $G$ 的\emph{下中心列}。“下”表示 $G^i\geq G^{i+1}$。

与上中心列一样，本节末练习中也要求验证对所有 $i$，$G^i$ 都是 $G$ 的特征子群。下一定理给出上、下中心列之间的关系。

\medskip
\noindent\textbf{定理 8.} 群 $G$ 幂零当且仅当存在 $n\geq 0$ 使 $G^n=1$。更精确地，$G$ 的幂零类为 $c$ 当且仅当 $c$ 是使 $G^c=1$ 的最小非负整数。若 $G$ 的幂零类为 $c$，则对所有 $i\in\{0,1,\ldots,c-1\}$，
\[
  Z_i(G)\leq G^{c-i-1}\leq Z_{i+1}(G).
\]

\begin{proof}
这可通过对上中心列或下中心列的长度作直接归纳证明。
\end{proof}

一般而言，上中心列与下中心列的各项不必相同，不过在某些群中它们会相同。

\paragraph{注}
\begin{enumerate}[label=(\arabic*)]
  \item 若 $G$ 是阿贝尔群，我们已经看到 $G'=G^1=1$，故下中心列一项后终止于单位群。
  \item 与上中心列类似，对任意有限群，由阶数考虑，必存在整数 $n$ 使得
  \[
    G^n=G^{n+1}=G^{n+2}=\cdots.
  \]
  对非幂零群，$G^n$ 是 $G$ 的非平凡子群。例如，在第 5.4 节中我们证明 $S_3'=S_3''=A_3$。由于 $S_3$ 不是幂零的，必有 $S_3^i=A_3$（$i\geq 1$）。事实上
  \[
    (123)=[(12),(132)]\in [S_3,S_3']=S_3^2.
  \]
  一旦下中心列中相邻两项相同，这条链就从该处稳定。注意 $S_3$ 给出一个例子：其下中心列有两个不同项，而上中心列的所有项都等于单位群；因此对非幂零群，两种中心列的长度不必相同。
\end{enumerate}

\subsection*{可解群与导出列}

回忆第 3.4 节中，可解群被定义为具有一个列
\[
  1=H_0\normal H_1\normal\cdots\normal H_s=G,
\]
使每个因子 $H_{i+1}/H_i$ 都是阿贝尔群。现在我们用一列递降的特征子群给出可解性的另一个刻画。

\medskip
\noindent\textbf{定义.} 对任意群 $G$，归纳定义子群列
\[
  G^{(0)}=G,
  \qquad
  G^{(1)}=[G,G],
  \qquad
  G^{(i+1)}=[G^{(i)},G^{(i)}]\quad(i\geq 1).
\]
此子群列称为 $G$ 的\emph{导出列}或\emph{换位子列}。这些项也常写作 $G^{(1)}=G'$、$G^{(2)}=G''$ 等。同样，证明每个 $G^{(i)}$ 都是 $G$ 的特征子群留作练习。

重要的是，虽然 $G^{(0)}=G^0$ 且 $G^{(1)}=G^1$，但一般并不成立 $G^{(i)}=G^i$。区别在于：下中心列的第 $i+1$ 项是第 $i$ 项与整个群 $G$ 的换位子，而导出列的第 $i+1$ 项是第 $i$ 项与自身的换位子。因此
\[
  G^{(i)}\leq G^i\quad\text{对所有 }i,
\]
且包含可以是真包含。例如在 $G=S_3$ 中，$G^1=G'=A_3$ 且 $G^2=[S_3,A_3]=A_3$，而 $G^{(2)}=[A_3,A_3]=1$。

\medskip
\noindent\textbf{定理 9.} 群 $G$ 可解当且仅当存在 $n\geq 0$ 使 $G^{(n)}=1$。

\begin{proof}
先设 $G$ 可解，故存在列
\[
  1=H_0\normal H_1\normal\cdots\normal H_s=G,
\]
且每个因子 $H_{i+1}/H_i$ 都是阿贝尔群。我们用归纳法证明 $G^{(i)}\leq H_{s-i}$。$i=0$ 时显然。若 $G^{(i)}\leq H_{s-i}$，则
\[
  G^{(i+1)}=[G^{(i)},G^{(i)}]\leq [H_{s-i},H_{s-i}].
\]
由于 $H_{s-i}/H_{s-i-1}$ 是阿贝尔群，由命题 5.7(4)，$[H_{s-i},H_{s-i}]\leq H_{s-i-1}$。归纳完成。因 $H_0=1$，得到 $G^{(s)}=1$。

反过来，若 $G^{(n)}=1$，则由命题 5.7(4)，取 $H_i=G^{(n-i)}$，便得到一条每个商群都为阿贝尔的正规列。因此 $G$ 可解。
\end{proof}

若 $G$ 可解，使 $G^{(n)}=1$ 的最小非负整数 $n$ 称为 $G$ 的\emph{可解长度}。导出列是具有阿贝尔逐次商的最短列，并且它还具有所有项都在整个群中为特征子群的性质（而不仅仅是每一项在下一项中正规）。其内在定义也使许多场合更容易处理，如下面命题所示。

\medskip
\noindent\textbf{命题 10.} 设 $G,K$ 是群，$H\leq G$，$\varphi:G\to K$ 是满同态。
\begin{enumerate}[label=(\arabic*)]
  \item 对所有 $i\geq 0$，$H^{(i)}\leq G^{(i)}$。特别地，若 $G$ 可解，则 $H$ 可解；也就是说，可解群的子群可解（且 $H$ 的可解长度不超过 $G$ 的可解长度）。
  \item 对所有 $i\geq 0$，$\varphi(G^{(i)})=K^{(i)}$。特别地，可解群的同态像和商群可解（可解长度不超过定义域的可解长度）。
  \item 若 $N\normal G$，且 $N$ 与 $G/N$ 都可解，则 $G$ 可解。
\end{enumerate}

\begin{proof}
(1) 由 $H\leq G$ 可知 $[H,H]\leq [G,G]$，即 $H^{(1)}\leq G^{(1)}$；再归纳得 $H^{(i)}\leq G^{(i)}$。若 $G^{(n)}=1$，则 $H^{(n)}=1$。

(2) 注意由换位子定义
\[
  \varphi([x,y])=[\varphi(x),\varphi(y)].
\]
故归纳得 $\varphi(G^{(i)})\leq K^{(i)}$。由于 $\varphi$ 是满射，$K$ 中的每个换位子都来自 $G$ 中某个换位子的像，于是归纳得到等号。若 $G^{(n)}=1$，则 $K^{(n)}=1$。

(3) 设 $G/N$ 与 $N$ 的可解长度分别为 $n,m$。对自然投影 $\varphi:G\to G/N$ 应用 (2)，得
\[
  \varphi(G^{(n)})=(G/N)^{(n)}=1N,
\]
即 $G^{(n)}\leq N$。于是
\[
  G^{(n+m)}=(G^{(n)})^{(m)}\leq N^{(m)}=1.
\]
由定理 9，$G$ 可解。证明完毕。
\end{proof}

关于有限群可解性，还有一些附加条件如下。

\medskip
\noindent\textbf{定理 11.} 设 $G$ 是有限群。
\begin{enumerate}[label=(\arabic*)]
  \item （Burnside）若 $|G|=p^aq^b$，其中 $p,q$ 是素数，则 $G$ 可解。
  \item （Philip Hall）若对每个整除 $|G|$ 的素数 $p$，将 $|G|$ 写为 $p^a m$ 且 $(p,m)=1$，并且 $G$ 有阶为 $m$ 的子群，则 $G$ 可解。也就是说，若对每个素数 $p$，$G$ 都有一个指标等于 Sylow $p$-子群阶的子群，则 $G$ 可解；这样的子群称为 Sylow $p$-补。
  \item （Feit--Thompson）若 $|G|$ 是奇数，则 $G$ 可解。
  \item （Thompson）若对每一对元素 $x,y\in G$，$\langle x,y\rangle$ 都是可解群，则 $G$ 可解。
\end{enumerate}

Burnside 定理将在第 19 章证明，并由此推出 Philip Hall 的推广。正如第 3.5 节提到的，Feit--Thompson 定理的证明长达 255 页。Thompson 定理最早是作为一篇 475 页论文的推论得到的，而该论文最终也依赖 Feit--Thompson 定理。

\subsection*{有限阿贝尔群基本定理的一个证明}

我们概述一个群论证明：每个有限阿贝尔群都是循环群的直积（即第 5.2 节定理 5 的第 1、2 部分）。有限生成阿贝尔群分类定理（第 5.2 节定理 3）将在第 12 章作为一个更一般定理的推论得到。

由推论 4，只需证明对素数 $p$，任意阿贝尔 $p$-群都是循环群的直积（第 5.5 定理中的整除条件可通过重排因子平凡地实现）。设 $A$ 是阿贝尔 $p$-群。我们对 $|A|$ 作归纳。

若 $E$ 是初等阿贝尔 $p$-群（即对所有 $x\in E$，$x^p=1$），先证明：
\[
  \text{对任意 }x\in E，\text{存在 }M\leq E\text{ 使 }E=M\times\langle x\rangle.
\]
若 $x=1$，取 $M=E$。否则，取 $M$ 为满足 $x\notin M$ 的、阶尽可能大的 $E$ 的子群。若 $M$ 在 $E$ 中指标不是 $p$，令 $\overline E=E/M$。则 $\overline E$ 是初等阿贝尔群，且存在 $\overline y\in \overline E-\langle \overline x\rangle$。由于 $\overline y$ 阶为 $p$，$\overline x\notin \langle\overline y\rangle$。$\langle\overline y\rangle$ 在 $E$ 中的完全原像是不含 $x$ 且阶大于 $M$ 的子群，矛盾。因此 $|E:M|=p$，于是
\[
  E=M\langle x\rangle,
  \qquad
  M\cap\langle x\rangle=1.
\]
由直积识别定理，$E=M\times\langle x\rangle$。

现在令 $\varphi:A\to A$ 由 $\varphi(x)=x^p$ 定义（见第 5.2 节练习 7）。由于 $A$ 阿贝尔，$\varphi$ 是同态。记 $K=\ker\varphi$，$H=\im\varphi$。则 $K=\{x\in A\mid x^p=1\}$，而 $H$ 是 $A$ 中所有 $p$ 次幂组成的子群。注意 $K$ 与 $A/H$ 都是初等阿贝尔群。由第一同构定理，
\[
  |A:H|=|K|.
\]
由归纳假设，
\[
  H=\langle h_1\rangle\times\cdots\times\langle h_r\rangle
  \cong \ZZ_{p^{a_1}}\times\cdots\times\ZZ_{p^{a_r}},
  \qquad a_i\geq 1.
\]
由 $\varphi$ 的定义，存在 $g_i\in A$ 使 $g_i^p=h_i$。令
\[
  A_0=\langle g_1,\ldots,g_r\rangle.
\]
可验证
\begin{enumerate}[label=(\alph*)]
  \item $A_0=\langle g_1\rangle\times\cdots\times\langle g_r\rangle$；
  \item $A_0/H=\langle g_1H\rangle\times\cdots\times\langle g_rH\rangle$ 是阶为 $p^r$ 的初等阿贝尔群；
  \item $H\cap K=\langle h_1^{p^{a_1-1}}\rangle\times\cdots\times\langle h_r^{p^{a_r-1}}\rangle$ 是阶为 $p^r$ 的初等阿贝尔群。
\end{enumerate}
若 $K\leq H$，则 $|K|=|K\cap H|=p^r=|A_0:H|$，比较阶得到 $A_0=A$，定理得证。

否则 $K\nleq H$。用横线表示取商 $A/H$，取 $x\in K-H$，则 $|\overline x|=|x|=p$。由证明开头对初等阿贝尔 $p$-群的断言，应用于 $\overline E=\overline A$，存在 $\overline A$ 的子群 $\overline M$ 使
\[
  \overline A=\overline M\times\langle\overline x\rangle.
\]
若 $M$ 是 $\overline M$ 在 $A$ 中的完全原像，则由于 $x$ 的阶为 $p$ 且 $x\notin M$，有 $\langle x\rangle\cap M=1$。由直积识别定理，
\[
  A=M\times\langle x\rangle.
\]
由归纳假设，$M$ 是循环群的直积，因此 $A$ 也是循环群的直积。证明完成。

有限阿贝尔群分解为循环群直积的唯一性（第 5.5 定理第 3 部分）也可以用 $p$ 次幂映射作归纳证明。这基本上就是第 12.1 节证明有限生成阿贝尔群基本定理唯一性部分时采用的过程。

\subsection*{练习}
\begin{enumerate}[label=\arabic*.]
  \item 证明对所有 $i$，$Z_i(G)$ 是 $G$ 的特征子群。
  \item 证明定理 1 的 (2) 和 (4) 对任意有限幂零群 $G$ 也成立，不必要求 $G$ 是 $p$-群。
  \item 若 $G$ 有限，证明 $G$ 幂零当且仅当对每个整除 $|G|$ 的阶都有一个正规子群；并证明 $G$ 循环当且仅当对每个整除 $|G|$ 的阶都有唯一子群。
  \item 证明有限幂零群的极大子群具有素数指标。
  \item 证明定理 1 的 (2) 和 (4) 对无限幂零群 $G$ 也成立。
  \item 证明若 $G/Z(G)$ 幂零，则 $G$ 幂零。
  \item 证明幂零群的子群和商群都是幂零的（你的证明应适用于无限群）。给出一个明确例子：某群 $G$ 有正规子群 $H$，且 $H$ 与 $G/H$ 都幂零，但 $G$ 不幂零。
  \item 证明若 $p$ 是素数且 $P$ 是阶为 $p^3$ 的非阿贝尔群，则 $|Z(P)|=p$ 且 $P/Z(P)\cong \ZZ_p\times\ZZ_p$。
  \item 证明有限群 $G$ 幂零当且仅当任取 $a,b\in G$ 且 $(|a|,|b|)=1$ 时都有 $ab=ba$。[使用定理 3 的第 (4) 部分。]
  \item 证明 $D_{2n}$ 幂零当且仅当 $n$ 是 $2$ 的幂。[使用练习 9。]
  \item 在额外假设 $G$ 阿贝尔的情形下，通过调用有限阿贝尔群基本定理给出命题 5 的另一个证明。
  \item 求 $A_4$ 和 $S_4$ 的上中心列与下中心列。
  \item 求 $A_n$ 和 $S_n$（$n\geq 5$）的上中心列与下中心列。
  \item 证明对所有 $i$，$G^i$ 是 $G$ 的特征子群。
  \item 证明 $Z_i(D_{2^n})=D_{2^n}^{\,n-1-i}$。
  \item 证明 $\QQ$ 没有极大子群。[回忆第 3.2 节练习 21。]
  \item 证明对所有 $i$，$G^{(i)}$ 是 $G$ 的特征子群。
  \item 证明若 $G'/G''$ 与 $G''/G'''$ 都是循环群，则 $G''=1$。[可假设 $G'''=1$，然后令 $G/G''$ 通过共轭作用于循环群 $G''$。]
  \item 证明不存在一个群，其换位子子群同构于 $S_4$。[使用前一题。]
  \item 设 $p$ 是素数，$P$ 是有限群 $G$ 的 $p$-子群，$N\normal G$ 且 $|N|$ 与 $p$ 互素，令 $\overline G=G/N$。证明：
  \begin{enumerate}[label=(\alph*)]
    \item $N_{\overline G}(\overline P)=\overline{N_G(P)}$；[使用 Frattini 论证。]
    \item $C_{\overline G}(\overline P)=\overline{C_G(P)}$。[使用 (a)。]
  \end{enumerate}
\end{enumerate}

对任意群 $G$，其 Frattini 子群（记为 $\Phi(G)$）定义为 $G$ 所有极大子群的交；若 $G$ 没有极大子群，则规定 $\Phi(G)=G$。下面几题讨论这个重要子群。

\begin{enumerate}[label=\arabic*.,resume]
  \item 证明 $\Phi(G)$ 是 $G$ 的特征子群。
  \item 证明若 $N\normal G$，则 $\Phi(N)\leq \Phi(G)$。给出一个明确例子说明若 $N$ 不正规，则这个包含不一定成立。
  \item 计算 $\Phi(S_3)$、$\Phi(A_4)$、$\Phi(S_4)$、$\Phi(A_5)$ 和 $\Phi(S_5)$。
  \item 称 $G$ 的元素 $x$ 是\emph{非生成元}，若对 $G$ 的每个真子群 $H$，$\langle x,H\rangle$ 仍是 $G$ 的真子群。证明 $\Phi(G)$ 正是 $G$ 的非生成元集合（这里 $|G|>1$）。
  \item 设 $G$ 是有限群。证明 $\Phi(G)$ 幂零。[用 Frattini 论证证明 $\Phi(G)$ 的每个 Sylow 子群都在 $G$ 中正规。]
  \item 设 $p$ 是素数，$P$ 是有限 $p$-群，令 $\overline P=P/\Phi(P)$。
  \begin{enumerate}[label=(\alph*)]
    \item 证明 $\overline P$ 是初等阿贝尔 $p$-群。[证明 $P'\leq \Phi(P)$ 且对所有 $x\in P$，$x^p\in\Phi(P)$。]
    \item 证明若 $N\normal P$ 且 $P/N$ 是初等阿贝尔群，则 $\Phi(P)\leq N$。用同态与交换图表述这一泛性质。
    \item 设 $\overline P$ 是阶为 $p^r$ 的初等阿贝尔群。由练习 24 推出：若 $\overline x_1,\ldots,\overline x_r$ 是 $\overline P$ 在 $\FF_p$ 上的一个基，且 $x_i$ 是陪集 $\overline x_i$ 中任一元素，则 $P=\langle x_1,
\ldots,x_r\rangle$。反过来证明若 $y_1,\ldots,y_s$ 生成 $P$，则 $s\geq r$。推出 Burnside 基定理：$y_1,
\ldots,y_s$ 是 $P$ 的极小生成集当且仅当其像是 $\overline P=P/\Phi(P)$ 的一个基。并推出 $P$ 的任意极小生成集有 $r$ 个元素。
    \item 证明若 $P/\Phi(P)$ 循环，则 $P$ 循环。推出若 $P/P'$ 循环，则 $P$ 循环。
    \item 设 $\alpha$ 是 $P$ 的素数阶 $q$ 的自同构，且 $q\neq p$。证明若 $\alpha$ 固定陪集 $x\Phi(P)$，则它固定该陪集中的某个元素。[把 $\alpha$ 看作该陪集的 $p^a$ 个元素上的一个阶为 $1$ 或 $q$ 的置换。]
    \item 用 (e) 和 (c) 推出：$P$ 的任意阶与 $p$ 互素的非平凡自同构，在 $P/\Phi(P)$ 上诱导非平凡自同构。推出任意阶与 $p$ 互素的 $P$ 的自同构群同构于 $\Aut(\overline P)=\GL_r(\FF_p)$ 的一个子群。
  \end{enumerate}
  \item 推广前一题 (d)：设 $p$ 是素数，$P$ 是 $p$-群，且 $\overline P=P/\Phi(P)$ 是阶为 $p^r$ 的初等阿贝尔群。证明 $P$ 恰有
  \[
    \frac{p^r-1}{p-1}
  \]
  个极大子群。[由于 $P$ 的每个极大子群都含有 $\Phi(P)$，由格同构定理，$P$ 的极大子群与初等阿贝尔群 $\overline P$ 的极大子群一一对应。因此只需证明阶为 $p^r$ 的初等阿贝尔 $p$-群的极大子群个数如上。一个办法是使用阿贝尔群同构于其对偶群的结果（参见第 5.2 节练习 14），所以指标为 $p$ 的子群个数等于阶为 $p$ 的子群个数。]
  \item 证明若 $p$ 是素数，$P=\ZZ_p\times\ZZ_{p^2}$，则 $|\Phi(P)|=p$ 且 $P/\Phi(P)\cong \ZZ_p\times\ZZ_p$。推出 $P$ 有 $p+1$ 个极大子群。
  \item 证明若 $p$ 是素数，$P$ 是阶为 $p^3$ 的非阿贝尔群，则 $\Phi(P)=Z(P)$ 且 $P/\Phi(P)\cong \ZZ_p\times\ZZ_p$。推出 $P$ 有 $p+1$ 个极大子群。
  \item 令 $p$ 是奇素数，$P_1=\ZZ_p\times\ZZ_{p^2}$，$P_2$ 为有一个阶为 $p^2$ 元素的非阿贝尔 $p^3$ 阶群。证明 $P_1$ 与 $P_2$ 有相同的子群格。
  \item 对任意群 $G$，称 $G$ 的极小正规子群为 $G$ 的一个正规子群 $M$，使得包含在 $M$ 中的 $G$ 的正规子群只有 $1$ 和 $M$。证明有限可解群的每个极小正规子群都是某个素数 $p$ 的初等阿贝尔 $p$-群。[若 $M$ 是 $G$ 的极小正规子群，考察它的特征子群 $M'$ 与 $\langle x^p\mid x\in M\rangle$。]
  \item 证明有限可解群的每个极大子群都有素数幂指标。[令 $H$ 是 $G$ 的极大子群，$M$ 是 $G$ 的极小正规子群——参见前一练习。对 $G/M$ 用归纳法，并分别考虑 $M\leq H$ 与 $M\nleq H$ 两种情形。]
  \item 设 $\pi$ 是任意素数集合。有限群 $G$ 的子群 $H$ 称为 Hall $\pi$-子群，若 $|H|$ 的所有素因子都在 $\pi$ 中，且 $|H|$ 与 $|G:H|$ 互素。（注意若 $\pi=\{p\}$，Hall $\pi$-子群就是 Sylow $p$-子群；Hall 子群在第 3.3 节练习 10 中已引入。）证明可解群的 Sylow 定理推广：若 $G$ 是有限可解群，则对每个素数集合 $\pi$，$G$ 有 Hall $\pi$-子群，且任意两个 Hall $\pi$-子群（对同一集合 $\pi$）共轭。[固定 $\pi$，对 $|G|$ 归纳，并同时证明存在性与共轭性。令 $M$ 是 $G$ 的一个极小正规子群，则 $M$ 是某素数 $p$ 的 $p$-群。若 $p\in\pi$，对 $G/M$ 应用归纳假设。若 $p\notin\pi$，化到 $|G|=p^a n$ 的情形，其中 $p^a=|M|$，$n$ 是 $G$ 的 Hall $\pi$-子群的阶。在这种情况下，令 $N/M$ 是 $G/M$ 的一个极小正规子群，则 $N/M$ 是某个素数 $q\neq p$ 的 $q$-群。令 $Q\in\Syl_q(N)$。若 $Q\normal G$，用 $Q$ 代替 $M$ 继续论证；若 $Q$ 在 $G$ 中不正规，用 Frattini 论证说明 $N_G(Q)$ 是 $G$ 的一个 Hall $\pi$-子群，并在这种情形下也证明共轭性。]
  \item 设 $p$ 整除有限可解群 $G$ 的阶。假设 $G$ 没有非平凡的、阶与 $p$ 互素的正规子群。令 $P$ 是 $G$ 的最大正规 $p$-子群（参见第 4.5 节练习 37）。注意上面练习 31 说明 $P\neq1$。证明 $C_G(P)\leq P$，即 $C_G(P)=Z(P)$。[令 $N=C_G(P)$，并用前一练习证明 $N=Z(P)\times H$，其中 $H$ 是 $N$ 的一个 Hall $\pi$-子群——这里 $\pi$ 是 $|N|$ 的所有不同于 $p$ 的素因子集合。证明 $H\normal G$，从而得到所需结论：$H=1$。]
  \item 证明若有限群 $G$ 的每个真子群都是幂零群，则 $G$ 可解。[证明一个极小反例是单群。取不同极大子群 $M,N$，使 $|M\cap N|$ 尽可能大，并用定理 3 的第 (2) 部分证明 $M\cap N=1$。然后应用第 4.5 节练习 53 的方法。]
  \item 设 $p$ 是素数，$V$ 是 $\FF_p$ 上非零有限维向量空间，$\phi\in\GL(V)$ 的阶为 $p$ 的幂。证明存在非零 $v\in V$ 使 $\phi(v)=v$。
  \item 设 $V$ 是 $\FF_2$ 上有限维向量空间，$\phi\in\GL(V)$ 的阶为 $2$（即 $V$ 是非平凡初等阿贝尔 $2$-群，$\phi$ 是 $V$ 的阶为 $2$ 的自同构）。证明映射 $v\mapsto v+\phi(v)$ 是 $V$ 到自身的同态。证明其像中的每个元素都被 $\phi$ 固定。推出 $\phi$ 的不动点子空间维数至少为 $\frac12\dim V$。（注意若 $G$ 是 $V$ 与 $\langle\phi\rangle$ 的半直积，其中 $V\normal G$ 且 $\phi$ 通过共轭作用在 $V$ 上，把每个 $v\in V$ 送到 $\phi(v)$，则 $\phi$ 在 $V$ 上的不动点就是 $C_V(\phi)$，而上面的映射就是换位子映射 $v\mapsto [v,\phi]$。用这种术语，本题是要证明 $|C_V(\phi)|^2\geq |V|$。）
  \item 用前一题证明：若 $P$ 是中心循环的 $2$-群，$M$ 是 $P$ 中指标为 $2$ 的子群，则 $Z(M)$ 的秩至多为 $2$。[阶为 $2$ 的群 $G/M$ 通过共轭作用在 $\FF_2$ 向量空间 $\{z\in Z(M)\mid z^2=1\}$ 上，而该作用的不动点在 $P$ 的中心中。]
\end{enumerate}

\section{中等阶群中的应用}

本节的目的是通过一系列例子展示前面发展出的许多技巧。这些例子大量使用 Sylow 定理，并说明它们如何用于研究有限群。受 H\"older 纲领启发，我们主要讨论如下问题：对某些 $n$，证明每个阶为 $n$ 的群都有真而非平凡的正规子群（即不存在阶为 $n$ 的单群）。多数情况下，一旦得到这一结论我们就停止。不过读者应注意，在达到这一结果的过程中，我们实际上已经获得了关于给定阶 $n$ 的任意群的大量信息；这些信息可以进一步用于分类这些阶的群，但一般需要超出半直积构造的技术。

由于对素数 $p$，我们已经证明除 $\ZZ_p$ 外不存在单 $p$-群，而 $p$-群的结构又可能非常复杂（回忆第 5.3 节的表），所以我们不显式研究 $p$-群本身的结构。相反，我们将把上一节的 $p$-群理论应用于非素数幂阶群的子群。

最后，对某些 $n$（例如 $60,168,360,504,\ldots$），确实存在阶为 $n$ 的单群，因此当然不能强迫这些阶的所有群都非单。与第 4.5 节一样，我们有时可以证明阶为 $n$ 的单群唯一，并剖析其内部结构（Sylow 数等）。我们将把阶为 $168$ 的单群作为另一个检验例子来研究。因此，Sylow 定理将在不同情境中被应用，以展示如何操纵给定阶的群。

本节末我们还将讨论群的存在性问题，特别是有限单群的存在性问题。

对 $n<10000$，存在 $60$ 个奇的非素数幂整数，使得 Sylow 定理的同余条件不能迫使至少一个 Sylow 子群正规；也就是说，对所有素数 $p\mid n$，$n_p$ 都可能大于 $1$。例如，由第 4.5 节的结果，形如 $pq$（$p,q$ 为不同素数）的数不会出现在我们的名单中。相反，对小于 $500$ 的偶数，已经有 $46$ 个可能成为单群阶的候选数。我们许多数值例子来自这些名单，且常用奇数，因为 Sylow 同余条件给出的 $n_p$ 可能值较少。这些例子的目的在于说明已证明结果的用法；其中许多也可由更高级的技术处理（例如 Feit--Thompson 定理说明不存在奇合数阶单群）。

正如第 4.5 节 $n=30$ 的情形中看到的，尽管 Sylow 定理允许 $n_5=6$ 和 $n_3=10$，进一步分析表明阶为 $30$ 的任意群都必须有 $n_5=1$ 且 $n_3=1$。因此，Sylow 定理的同余部分是 Sylow 子群正规的充分条件，但远非必要条件。对于许多 $n$（如 $n=120$），我们可证明不存在阶为 $n$ 的单群，故存在非平凡正规子群，但该子群未必是 Sylow 子群。例如 $S_5$ 与 $\SL_2(\FF_5)$ 都有阶 $120$；$S_5$ 有唯一的非平凡真正规子群 $A_5$，阶为 $60$，而 $\SL_2(\FF_5)$ 有唯一的非平凡真正规子群 $Z(\SL_2(\FF_5))\cong\ZZ_2$，阶为 $2$；二者都不是 Sylow 子群。因此产生正规子群的技术必须足够灵活。本节中我们将考察不同素数的 Sylow 子群、Sylow 子群的交、$p$-子群的正规化子以及许多不那么显然的子群。这里列出的初等方法即使对“中等”阶群也远非穷尽。

\subsection*{一些技巧}

在列出给定（“中等”）阶群中产生正规子群的一些技巧之前，我们注意：所有处理阶为某个具体 $n$ 的群的问题，首先都必须把 $n$ 分解成素数幂，并计算对每个素数 $p\mid n$，$n_p$ 的允许值。最后一步中要熟练进行模 $p$ 计算。我们描述的技巧如下：
\begin{enumerate}[label=(\arabic*)]
  \item 计数元素；
  \item 利用小指标子群；
  \item 置换表示；
  \item 让不同素数 $p$ 的 $p$-子群相互制约；
  \item 研究 Sylow $p$-子群交的正规化子。
\end{enumerate}

\subsection*{计数元素}

设 $G$ 是阶为 $n$ 的群，$p\mid n$，$P\in\Syl_p(G)$。若 $|P|=p$，则 $P$ 的每个非单位元素阶为 $p$，而 $G$ 中每个阶为 $p$ 的元素都在 $P$ 的某个共轭中。由 Lagrange 定理，不同共轭的 $P$ 只在单位元处相交。因此在此情形下，$G$ 中阶为 $p$ 的元素数为
\[
  n_p(p-1).
\]

若对不同素数 $p$ 的 Sylow $p$-子群都具有素数阶，且假设没有一个正规，则有时可推出素数阶元素个数超过 $|G|$。这种矛盾说明至少一个 $n_p$ 必须为 $1$，即某个 Sylow 子群正规。

这正是第 4.5 节用来证明不存在阶为 $30$ 的单群的论证。再如，设 $|G|=105=3\cdot5\cdot7$。若 $G$ 单，则必须有
\[
  n_3=7,
  \qquad n_5=21,
  \qquad n_7=15.
\]
于是
\[
\begin{array}{rcl}
\text{阶为 }3\text{ 的元素数} &=& 7\cdot2=14,\\
\text{阶为 }5\text{ 的元素数} &=& 21\cdot4=84,\\
\text{阶为 }7\text{ 的元素数} &=& 15\cdot6=90,
\end{array}
\]
总计 $188>|G|$，矛盾。

有时计数素数阶元素不会得到“太多”元素。不过，剩下的元素可能少到必须形成某个正规子群。第 4.5 节证明阶为 $12$ 的群中 $n_2=1$ 或 $n_3=1$ 的方法本质上就是如此。当 $G$ 有一个阶为 $p$ 的 Sylow $p$-子群 $P$ 且 $N_G(P)=P$ 时，这种技巧尤其有效。例如，令 $|G|=56$。若 $G$ 单，则 Sylow $7$-子群数唯一可能是 $8$，所以阶为 $7$ 的元素数为
\[
  8\cdot6=48.
\]
于是 $G$ 中还剩 $56-48=8$ 个元素。Sylow $2$-子群含有 $8$ 个元素，且其中无阶为 $7$ 的元素，因此 Sylow $2$-子群至多一个，故 $G$ 有正规 Sylow $2$-子群。

\subsection*{利用小指标子群}

回忆第 4.2 节的结果：若 $G$ 有指标为 $k$ 的子群 $H$，则存在一个同态 $G\to S_k$，其核包含于 $H$。若 $k>1$，该核是 $G$ 的真正规子群；若我们试图证明 $G$ 非单，可反证假设该核为单位群。于是由第一同构定理，$G$ 同构于 $S_k$ 的一个子群，特别地 $|G|$ 整除 $k!$。因此，若 $G$ 是有限单群，且 $k$ 是满足 $|G|\mid k!$ 的最小整数，则 $G$ 不含指标小于 $k$ 的真子群。研究给定阶 $n$ 的群时，应在一开始就计算这个最小允许指标。

在我们考虑的例子中，这个计算通常很容易：$n$ 常常分解为
\[
  n=p_1^{\alpha_1}p_2^{\alpha_2}\cdots p_s^{\alpha_s},
  \qquad p_1<p_2<\cdots<p_s,
\]
并且在我们的例子中，$\alpha_s$ 通常等于 $1$ 或 $2$。在这种情况下，真子群的最小指标必须至少是 $p_s$（分别地，至少是 $2p_s$），而且这常常就是它的精确值。

例如，不存在阶为 $3393$ 的单群。因为
\[
  3393=3^2\cdot 13\cdot29,
\]
而真子群的最小可能指标是 $29$（$3393\nmid 28!$，因为 $29\nmid 28!$）。但任意阶为 $3393$ 的单群必须有 $n_3=13$，因此对 $P\in\Syl_3(G)$，$N_G(P)$ 的指标为 $13$，矛盾。

\subsection*{置换表示}

这是前一方法的细化。若 $G$ 是阶为 $n$ 的单群，且有指标为 $k$ 的真子群，则 $G$ 同构于 $S_k$ 的一个子群。于是可把 $G$ 看作 $S_k$ 的子群，并在 $S_k$ 内计算。有时不仅仅用 Lagrange 定理，而是直接证明 $S_k$ 不含阶为 $n$ 的单子群。两种有用限制是：
\begin{enumerate}[label=(\arabic*)]
  \item 若 $G$ 含有某个阶的元素或子群，则 $S_k$ 也必须含有；
  \item 若 $P\in\Syl_p(G)$ 且 $P$ 也是 $S_k$ 的 Sylow $p$-子群，则 $|N_G(P)|$ 必须整除 $|N_{S_k}(P)|$。
\end{enumerate}

情形 (2) 常发生在 $p$ 为素数且 $k=p$ 或 $p+1$ 时。此时 $p^2\nmid k!$，所以 $G$ 的 Sylow $p$-子群也是 $S_k$ 的 Sylow $p$-子群。$S_k$ 的 Sylow $p$-子群正是由一个 $p$-轮换生成的群，且不同 Sylow $p$-子群交平凡。因此
\[
  \#\Syl_p(S_k)=
  \frac{\text{$p$-轮换个数}}{\text{一个 Sylow $p$-子群中的 $p$-轮换个数}}
  =\frac{k(k-1)\cdots(k-p+1)}{p(p-1)}.
\]
该数给出 $S_k$ 中 Sylow $p$-子群正规化子的指标。因此，当 $k=p$ 或 $p+1$ 时，
\[
  |N_{S_k}(P)|=p(p-1).
\]

例如，若 $G$ 是阶为 $396=2^2\cdot3^2\cdot11$ 的单群，则必须有 $n_{11}=12$。若 $P\in\Syl_{11}(G)$，则 $|G:N_G(P)|=12$，所以 $|N_G(P)|=33$。由于 $G$ 有指标 $12$ 的子群，$G$ 同构于 $S_{12}$ 的子群。但在 $S_{12}$ 中，$P\in\Syl_{11}(S_{12})$ 且 $|N_{S_{12}}(P)|=110$。由 $N_G(P)\leq N_{S_{12}}(P)$，应有 $33\mid110$，不可能。因此不存在阶为 $396$ 的单群。

有时可稍进一步，在 $A_k$ 中而非 $S_k$ 中工作。这个小改进只偶尔有用，且只针对偶数阶群。

\medskip
\noindent\textbf{命题 12.}
\begin{enumerate}[label=(\arabic*)]
  \item 若 $G$ 没有指标为 $2$ 的子群且 $G\leq S_k$，则 $G\leq A_k$。
  \item 若 $P\in\Syl_p(S_k)$，其中 $p$ 是奇素数，则 $P\in\Syl_p(A_k)$ 且
  \[
    |N_{A_k}(P)|=\frac12|N_{S_k}(P)|.
  \]
\end{enumerate}

\begin{proof}
(1) 由第二同构定理：若 $G\nleq A_k$，则 $A_k<G A_k$，故 $GA_k=S_k$。于是
\[
  2=|S_k:A_k|=|GA_k:A_k|=|G:G\cap A_k|,
\]
这说明 $G$ 有指标为 $2$ 的子群。

(2) 若 $P\in\Syl_p(S_k)$ 且 $p$ 为奇素数，则由 (1)（或阶数考虑）$P\leq A_k$，故 $P\in\Syl_p(A_k)$。由 Frattini 论证，
\[
  S_k=N_{S_k}(P)A_k.
\]
因此 $N_{S_k}(P)\nleq A_k$，从而 $N_{S_k}(P)\cap A_k=N_{A_k}(P)$ 是 $N_{S_k}(P)$ 的指标 $2$ 子群。
\end{proof}

例如，不存在阶为 $264$ 的单群。若 $G$ 是阶为 $264=2^3\cdot3\cdot11$ 的单群，则必须有 $n_{11}=12$。如前，$G$ 同构于 $S_{12}$ 的子群。由于 $G$ 单（故无指标 $2$ 子群），$G\leq A_{12}$。令 $P\in\Syl_{11}(G)$。由 $n_{11}=12=|G:N_G(P)|$，有 $|N_G(P)|=22$。而
\[
  |N_{A_{12}}(P)|=\frac12|N_{S_{12}}(P)|=\frac12\cdot11(11-1)=55.
\]
但 $22\nmid55$，与 $N_G(P)\leq N_{A_{12}}(P)$ 矛盾。

最后，我们强调，我们只是略微触及了某些置换表示所能提供的组合信息。在余下的例子中，只要可能，我们还会说明这一技巧的其他应用。

\subsection*{让不同素数的 $p$-子群相互制约}

设 $p,q$ 是不同素数，且每个阶为 $pq$ 的群都是循环群。若 $p<q$，这等价于 $p\nmid q-1$。若 $G$ 有阶为 $q$ 的 Sylow $q$-子群 $Q$，且 $p\mid |N_G(Q)|$，则在 $N_G(Q)$ 中由 Cauchy 定理得到一个阶为 $p$ 的群 $P$ 正规化 $Q$（注意 $P$ 不一定是 $G$ 的 Sylow $p$-子群）。于是 $PQ$ 是群；若 $PQ$ 阿贝尔，则
\[
  Q\leq N_G(P),
  \qquad \text{从而 } q\mid |N_G(P)|.
\]
类似地，若 $G$ 的 Sylow $p$-子群阶为 $p$，且 $q$ 整除某个 Sylow $p$-正规化子的阶，也可作对称论证。这类数值信息可能足以迫使 $N_G(P)=G$（即 $P\normal G$），或者至少迫使 $N_G(P)$ 具有小于置换表示所允许最小指标的指标，从而产生矛盾。

例如，不存在阶为 $1785$ 的单群。若存在，令 $G$ 是阶为
\[
  1785=3\cdot5\cdot7\cdot17
\]
的单群。$n_{17}$ 的唯一可能值为 $35$，所以若 $Q$ 是 Sylow $17$-子群，则 $|G:N_G(Q)|=35$，故 $|N_G(Q)|=3\cdot17$。令 $P$ 是 $N_G(Q)$ 的 Sylow $3$-子群。群 $PQ$ 是阿贝尔的，因为 $3\nmid 17-1$，所以 $Q\leq N_G(P)$，从而 $17\mid |N_G(P)|$。此时 $P\in\Syl_3(G)$。$n_3$ 的允许值为 $7,85,595$；但由于 $17\mid |N_G(P)|$，不能有 $17\mid |G:N_G(P)|=n_3$，故 $n_3=7$。这意味着 $G$ 有一个指标 $7$ 的真子群，与此阶的最小真子群指标为 $17$ 矛盾。另一种说法是：若 $n_3=7$，则
\[
  |N_G(P)|=3\cdot5\cdot17,
\]
并且由 Sylow 定理应用于 $N_G(P)$，有 $Q\normal N_G(P)$。这与 $|N_G(Q)|=3\cdot17$ 矛盾。

还可以细化这一方法，不要求 $P$ 与 $Q$ 都为素数阶。若 $p,q$ 是整除 $|G|$ 的不同素数，$Q\in\Syl_q(G)$，且 $p\mid |N_G(Q)|$，令 $P\in\Syl_p(N_G(Q))$。可在 $N_G(Q)$ 中应用 Sylow 定理，判断是否 $P\normal N_G(Q)$；若是，则可能迫使 $N_G(P)$ 具有小指标。若 $P$ 是整个 $G$ 的 Sylow $p$-子群，则可再用 Sylow 同余条件限制 $|N_G(P)|$。若 $P$ 不是 $G$ 的 Sylow $p$-子群，则 $P$ 包含在某个 $P^*\in\Syl_p(G)$ 中；由于 $P<P^*$，定理 1(4) 给出 $P<N_{P^*}(P)$，从而 $N_G(P)$ 的阶被比 $|P|$ 更高的 $p$ 的幂整除，并且还被 $|Q|$ 整除。

例如，不存在阶为 $3675$ 的单群。若存在，令 $G$ 是阶为
\[
  3675=3\cdot5^2\cdot7^2
\]
的单群。$n_7$ 的唯一可能值为 $15$，故对 $Q\in\Syl_7(G)$，$|G:N_G(Q)|=15$，$|N_G(Q)|=245=5\cdot7^2$。令 $N=N_G(Q)$，$P\in\Syl_5(N)$。由 Sylow 同余条件应用于 $N$，得 $P\normal N$。由于 $|P|=5$，$P$ 不是 $G$ 的 Sylow $5$-子群，故包含在某个 $P^*\in\Syl_5(G)$ 中。由于 $P$ 是 $5$-群 $P^*$ 的指标 $5$ 子群，由定理 1，$P\normal P^*$，也就是说 $P^*\leq N_G(P)$。因此
\[
  \langle N,P^*\rangle\leq N_G(P),
  \qquad \text{从而 } 7^2\cdot5^2\mid |N_G(P)|.
\]
所以 $|G:N_G(P)|$ 整除 $3$。这不可能，因为 $P$ 不正规且 $G$ 没有指标 $3$ 的子群。

\subsection*{研究 Sylow $p$-子群交的正规化子}

前面第一种方法中的计数论证不能直接推广到非素数阶 Sylow 子群的原因之一是：若 $P\in\Syl_p(G)$ 且 $|P|=p^a$、$a\geq2$，则不同共轭的 $P$ 不一定只在单位子群处相交。若不同共轭的 $P$ 确实只交于单位子群，则仍可计数得到 $p$-幂阶元素个数为
\[
  n_p(|P|-1).
\]

假设存在 $R\in\Syl_p(G)$，$R\neq P$ 且 $P\cap R\neq 1$。令 $P_0=P\cap R$。则 $P_0<P$ 且 $P_0<R$，由定理 1，
\[
  P_0<N_P(P_0),
  \qquad
  P_0<N_R(P_0).
\]
可试图由此证明 $P_0$ 在 $G$ 中的正规化子足够大（即指标足够小），从而用前述方法得到矛盾。

一个特别有效的特殊情形是 $|P_0|=p^{a-1}$，也就是两个 Sylow $p$-子群 $R$ 与 $P$ 有很大的交。此时令 $N=N_G(P_0)$。由上述推理（因为 $P_0$ 是 $p$-群 $P$ 与 $R$ 的极大子群），$P_0\normal P$ 且 $P_0\normal R$，即
\[
  N\text{ 有两个不同的 Sylow }p\text{-子群 }P\text{ 与 }R.
\]
特别地，$|N|=p^a k$，其中由 Sylow 定理 $k\geq p+1$。

总之，若 Sylow $p$-子群两两交平凡，则可计数 $p$-幂阶元素；否则存在某个 Sylow 子群交，其正规化子“很大”。对任意群阶，事先未必能判断哪种现象发生，因此非单性证明可能需要分成两个互斥情形并分别推出矛盾。当 Sylow $p$-子群阶为 $p^2$ 时，这种过程尤其适用（例如，在证明阶为 $60$ 的单群同构于 $A_5$ 时，计数 $2$-幂阶元素就使用了这种思路——见第 4.5 节命题 23）。

下面给出一个引理，提供迫使非平凡 Sylow 交的充分条件。

\medskip
\noindent\textbf{引理 13.} 在有限群 $G$ 中，若 $n_p\not\equiv 1\pmod{p^2}$，则存在不同的 Sylow $p$-子群 $P,R$，使得 $P\cap R$ 在 $P$ 和 $R$ 中指标都为 $p$（因而在二者中正规）。

\begin{proof}
这是 Sylow 定理同余部分证明的一个简单细化。令 $P$ 通过共轭作用于集合 $\Syl_p(G)$。设 $O_1,
\ldots,O_s$ 为该作用的轨道，其中 $O_1=\{P\}$。若对每个 $R\in\Syl_p(G)$、$R\neq P$ 都有 $p^2\mid |P:P\cap R|$，则每个 $O_i$ 的大小均被 $p^2$ 整除（$i=2,
\ldots,s$）。于是作为轨道长度之和，$n_p=1+kp^2$，与假设矛盾。因此存在 $R\in\Syl_p(G)$ 使 $|P:P\cap R|=p$。
\end{proof}

例如，不存在阶为 $1053$ 的单群。若存在，令 $G$ 是阶为
\[
  1053=3^4\cdot13
\]
的单群，令 $P\in\Syl_3(G)$。必须有 $n_3=13$。但 $13\not\equiv1\pmod{3^2}$，所以存在 $P,R\in\Syl_3(G)$ 使 $|P\cap R|=3^3$。令 $N=N_G(P\cap R)$，则由上面的论证 $P,R\leq N$。因此 $3^4\mid |N|$ 且 $|N|>3^4$。唯一可能是 $N=G$，即 $P\cap R\normal G$，矛盾。

\subsection*{阶为 $168$ 的单群}

现在我们说明如何用许多技巧来剖析并分类某些单群。我们将分类阶为 $168$ 的单群。由于单群没有非平凡正规子群，这一过程不同于第 5.5 节的方法，但总体思想代表了有限单群研究中的常见方法。

假设存在阶为
\[
  168=2^3\cdot3\cdot7
\]
的单群 $G$。我们先推导它的许多性质：Sylow 子群的数目与结构、共轭类等。所有这些计算只依赖于 $G$ 的阶与单性。随后用这些结果证明 $G$ 唯一，最后证明阶为 $168$ 的单群确实存在。

因为 $|G|\nmid 6!$，有
\begin{enumerate}[label=(\arabic*)]
  \item $G$ 没有指标小于 $7$ 的真子群。
\end{enumerate}
否则 $G$ 作用于该子群的陪集会给出一个（因 $G$ 单而必为单射的）同态 $G\to S_n$，其中 $n\leq6$。

由 $G$ 的单性和 Sylow 定理立刻得到：
\begin{enumerate}[label=(\arabic*),resume]
  \item $n_7=8$，所以 Sylow $7$-子群的正规化子阶为 $21$。特别地，没有阶为 $2$ 的元素正规化 Sylow $7$-子群，且 $G$ 没有阶为 $14$ 的元素。
\end{enumerate}

若 $G$ 有阶为 $21$ 的元素，则 $G$ 的 Sylow $3$-子群正规化子的阶会被 $7$ 整除。于是 $n_3$ 与 $7$ 互素；又 $n_3\mid 8$，便有 $n_3=4$，与 (1) 矛盾。因此：
\begin{enumerate}[label=(\arabic*),resume]
  \item $G$ 没有阶为 $21$ 的元素。
\end{enumerate}

由 Sylow 定理，$n_3=7$ 或 $28$；下面排除前者。若 $n_3=7$，令 $P\in\Syl_3(G)$，并令 $T$ 是 $N_G(P)$（阶为 $24$）的 Sylow $2$-子群。每个 Sylow $3$-子群都正规化 $G$ 的某个 Sylow $7$-子群，因此 $P$ 正规化某个 $R\in\Syl_7(G)$。对每个 $t\in T$，$P=tPt^{-1}$ 也正规化 $tRt^{-1}$。子群 $T$ 通过共轭作用于八个 Sylow $7$-子群；由于由 (2) 没有阶为 $2$ 的元素正规化 Sylow $7$-子群，$T$ 的作用是传递的。于是 $P$ 正规化 $G$ 的每个 Sylow $7$-子群，即 $P$ 包含在所有 Sylow $7$-正规化子的交中。但这个交是 $G$ 的真正规子群，因此必须平凡，矛盾。故
\begin{enumerate}[label=(\arabic*),resume]
  \item $n_3=28$，Sylow $3$-子群的正规化子阶为 $6$。
\end{enumerate}

由于 $n_2=7$ 或 $21$，有 $n_2\not\equiv1\pmod 8$。由引理 13（或相应练习），存在一对不同 Sylow $2$-子群有非平凡交。选取所有这样的对中交的阶最大的 $T_1,T_2$，令 $U=T_1\cap T_2$。接下来证明
\begin{enumerate}[label=(\arabic*),resume]
  \item $U$ 是 Klein 四元群，且 $N_G(U)\cong S_4$。
\end{enumerate}

令 $N=N_G(U)$。由于 $|U|=2$ 或 $4$，且 $N$ 通过共轭置换 $U$ 的非单位元素，$N$ 中阶为 $7$ 的子群将与 $U$ 中某个阶为 $2$ 的元素交换，违背 (2)。故 $7\nmid |N|$。由引理 13 的思想，$N$ 有不止一个 Sylow $2$-子群，故 $|N|=2^a\cdot3$，其中 $a=2$ 或 $3$。令 $P\in\Syl_3(N)$。由于 $P$ 是 $G$ 的 Sylow $3$-子群，由 (4)，$N_N(P)$ 的阶为 $3$ 或 $6$。于是由 Sylow 定理，$N$ 必有四个 Sylow $3$-子群，并且 $N$ 通过共轭传递置换它们。由于任意阶为 $12$ 的群都有正规 Sylow $2$-子群或正规 Sylow $3$-子群（参见第 4.5 节），必有 $|N|=24$。令 $K$ 是 $N$ 作用于其四个 Sylow $3$-子群的共轭作用的核。若 $K=1$，则 $N\cong S_4$。若 $K\neq1$，因为 $K\leq N_N(P)$，$|K|$ 整除 $6$；且 $P$ 不正规化另一个 Sylow $3$-子群，所以 $P\nleq K$。于是 $|K|=2$。但 $N/K$ 是阶为 $12$ 的群，并且有不止一个 Sylow $2$-子群和四个 Sylow $3$-子群，矛盾于上面引用的阶 $12$ 群的性质。故 $N\cong S_4$。由于 $S_4$ 有唯一的非平凡正规 $2$-子群 $V_4$，得到 (5)。

由 $N\cong S_4$，$N$ 含有 $G$ 的 Sylow $2$-子群，而且 $N_N(P)\cong S_3$（由 (4) 也有 $N_G(P)\cong S_3$）。于是得到
\begin{enumerate}[label=(\arabic*),resume]
  \item $G$ 的 Sylow $2$-子群同构于 $D_8$；
  \item $G$ 中 Sylow $3$-子群的正规化子同构于 $S_3$，因此 $G$ 没有阶为 $6$ 的元素。
\end{enumerate}

由 (2) 和 (7)，没有阶为 $2$ 的元素与奇素数阶元素交换。若 $T\in\Syl_2(G)$，则由 (6)，$T\cong D_8$，所以 $Z(T)=\langle z\rangle$，其中 $z$ 阶为 $2$。于是 $T\leq C_G(z)$，而由刚才所述，$|C_G(z)|$ 没有奇素因子，故 $C_G(z)=T$。任意正规化 $T$ 的元素都正规化其中心，因而与 $z$ 交换，所以 Sylow $2$-子群自正规化。故
\begin{enumerate}[label=(\arabic*),resume]
  \item $n_2=21$，且若 $T\in\Syl_2(G)$、$Z(T)=\langle z\rangle$，则 $C_G(z)=T$。
\end{enumerate}

由于 $|C_G(z)|=8$，元素 $z$ 有 $21$ 个共轭。由 (6)，$G$ 有一个阶为 $4$ 的元素共轭类，其大小为 $42$。由 (2) 有 $48$ 个阶为 $7$ 的元素；由 (4) 有 $56$ 个阶为 $3$ 的元素。这些加起来已经占据了所有 $167$ 个非单位元素，因此每个阶为 $2$ 的元素都与 $z$ 共轭，即
\begin{enumerate}[label=(\arabic*),resume]
  \item $G$ 有唯一的阶为 $2$ 元素的共轭类。
\end{enumerate}

保持记号，令 $T\in\Syl_2(G)$，$U\leq T$，并令 $W$ 为 $T$ 中另一个 Klein 四元群。由 Sylow 定理可知 $U$ 与 $W$ 在 $G$ 中不共轭，因为它们在 $N_G(T)=T$ 中不共轭（参见第 4.5 节练习 50）。接下来证明
\begin{enumerate}[label=(\arabic*),resume]
  \item $N_G(W)\cong S_4$。
\end{enumerate}
为此写 $W=\langle z,w\rangle$，其中 $\langle z\rangle=Z(T)$。由于 $w$ 在 $G$ 中与 $z$ 共轭，$C_G(w)=T_0$ 是另一个含有 $W$ 且不同于 $T$ 的 Sylow $2$-子群。因此 $W=T\cap T_0$。因为 $U$ 是 Sylow $2$-子群最大交的任意选择，给出 (5) 的论证也给出 (10)。

现在记录已证明或易于推出的结果。

\medskip
\noindent\textbf{命题 14.} 若 $G$ 是阶为 $168$ 的单群，则：
\begin{enumerate}[label=(\arabic*)]
  \item $n_2=21$，$n_3=28$，$n_7=8$；
  \item $G$ 的 Sylow $2$-子群是二面体群，Sylow $3$-与 $7$-子群是循环群；
  \item $G$ 同构于 $A_7$ 的一个子群，并且没有指标 $\leq6$ 的子群；
  \item $G$ 的共轭类如下：单位元；两个阶为 $7$ 的元素类，各含 $24$ 个元素（可由任一阶为 $7$ 的元素及其逆元代表）；一个阶为 $3$ 的元素类，含 $56$ 个元素；一个阶为 $4$ 的元素类，含 $42$ 个元素；一个阶为 $2$ 的元素类，含 $21$ 个元素。特别地，$G$ 的每个元素阶都是某个素数的幂；
  \item 若 $T\in\Syl_2(G)$，而 $U,W$ 是 $T$ 中两个 Klein 四元群，则 $U,W$ 在 $G$ 中不共轭，且 $N_G(U)\cong N_G(W)\cong S_4$；
  \item $G$ 恰有三个极大子群的共轭类，其中两个同构于 $S_4$，另一个同构于阶为 $21$ 的非阿贝尔群。
\end{enumerate}

以上所有计算都建立在“存在阶为 $168$ 的单群”的假设上。没有导出矛盾并不证明这种群存在，只是强烈提示可能存在。下面说明 $G$ 的内部子群结构如何导出一个几何，使 $G$ 作用在其上，从而证明若阶为 $168$ 的单群存在则唯一，并最终证明它确实存在。

沿用记号，令 $U_1,
\ldots,U_7$ 为 $U$ 的共轭，$W_1,
\ldots,W_7$ 为 $W$ 的共轭。称 $U_i$ 为\emph{点}，$W_j$ 为\emph{直线}。定义一个“关联关系”：
\[
  \text{点 }U_i\text{ 在直线 }W_j\text{ 上}\quad\Longleftrightarrow\quad U_i\text{ 正规化 }W_j.
\]
注意 $U_i$ 正规化 $W_j$ 当且仅当 $U_iW_j\cong D_8$，而这又当且仅当 $W_j$ 正规化 $U_i$。在每个同构于 $S_4$ 的点或线稳定子中，有唯一正规 $4$-群 $V$，以及恰好三个其他（非正规）$4$-群 $A_1,A_2,A_3$。群 $VA_i$ 是该 $S_4$ 的三个 Sylow $2$-子群。因此
\begin{enumerate}[label=(\arabic*),resume]
  \item 每条直线恰含 $3$ 个点，每个点恰在 $3$ 条直线上。
\end{enumerate}
由于 $S_4$ 中任意两个非正规 $4$-群生成整个 $S_4$，从而唯一确定该 $S_4$ 中另外两个 Klein 四元群，得到
\begin{enumerate}[label=(\arabic*),resume]
  \item 直线上任意两个点唯一确定这条直线（也唯一确定其第三个点）。
\end{enumerate}
因为共有 $7$ 个点和 $7$ 条直线，简单计数给出
\begin{enumerate}[label=(\arabic*),resume]
  \item 任意两个点在唯一一条直线上，任意两条直线交于唯一一个点。
\end{enumerate}

这一点线配置因而满足所谓射影平面的公理。现在可直接证明该关联几何唯一，并可用图 \ref{fig:fano} 表示：点为顶点，直线为三角形的六条边和中线以及内切圆（见练习 27）。该关联几何称为\emph{阶为 $2$ 的射影平面}或\emph{Fano 平面}，记为 $\mathcal F$。一般地，一个“阶”为 $N$ 的射影平面有 $N^2+N+1$ 个点，并有同样数目的直线。注意在这一点上射影平面 $\mathcal F$ 的确存在——我们已经显式给出了满足 (11) 到 (13) 的点与直线——尽管群 $G$ 尚未知道存在。

\begin{figure}[htbp]
\centering
\begin{tikzpicture}[scale=2.0]
  \coordinate (A) at (90:1);
  \coordinate (B) at (210:1);
  \coordinate (C) at (330:1);
  \coordinate (D) at ($(B)!0.5!(C)$);
  \coordinate (E) at ($(A)!0.5!(C)$);
  \coordinate (F) at ($(A)!0.5!(B)$);
  \coordinate (O) at (0,0);
  \draw (A)--(B)--(C)--cycle;
  \draw (A)--(D);
  \draw (B)--(E);
  \draw (C)--(F);
  \draw (O) circle (0.5);
  \foreach \P in {A,B,C,D,E,F,O}{\fill (\P) circle (1.2pt);}
\end{tikzpicture}
\caption{图 1}
\label{fig:fano}
\end{figure}

该平面的一个自同构是点和直线的任意置换，要求保持关联关系。图中三角形的六个对称显然给出 $\mathcal F$ 的自同构，但 $\mathcal F$ 还有更多自同构。

每个 $g\in G$ 都通过共轭作用在点和线的集合上，且该作用保持关联关系。只有单位元固定所有点，因此通过此作用，$G$ 同构于 $\Aut(\mathcal F)$ 的一个子群。任何固定同一直线上的两个点以及不在该直线上的第三个点的 $\mathcal F$ 自同构都固定所有点。因此，$\mathcal F$ 的任一自同构由其在任意三个不共线点上的作用唯一确定。容易计算这种三元组有 $168$ 个，所以 $\mathcal F$ 至多有 $168$ 个自同构。这证明
\[
  \text{若单群 }G\text{ 存在，则它唯一，且 }G\cong\Aut(\mathcal F).
\]

分类过程还剩两步：证明 $\mathcal F$ 确有 $168$ 个自同构，并证明 $\Aut(\mathcal F)$ 确实是单群。虽然这两个步骤可以用图论方法完成，我们采用一种遵循“代数群”理论思想的方法。令 $V$ 是 $\FF_2$ 上三维向量空间，即 $V$ 是阶为 $8$ 的初等阿贝尔 $2$-群 $\ZZ_2\times\ZZ_2\times\ZZ_2$。由第 4.4 节命题 17，
\[
  \Aut(V)=\GL(V)\cong \GL_3(\FF_2)
\]
的阶为 $168$。称 $V$ 的七个 $1$ 维子空间（即非平凡循环子群）为点，称七个 $2$ 维子空间（即阶为 $4$ 的子群）为直线，并规定点 $p$ 与线 $L$ 关联当且仅当 $p\subseteq L$。这些点与线容易看出满足上面的 (11)--(13)，因而表示 Fano 平面。由于 $\GL(V)$ 忠实地作用在这些点和线上并保持关联，$|\Aut(\mathcal F)|\geq168$。结合已知上界，得到
\[
  \Aut(\mathcal F)\cong \GL(V)\cong\GL_3(\FF_2),
  \qquad |\Aut(\mathcal F)|=168.
\]

最后证明 $\GL(V)$ 是单群。反设 $H$ 是 $\GL(V)$ 的非平凡真正规子群。令 $\Omega$ 为七个点的集合，$N$ 为 $\GL(V)$ 中某个点的稳定子。由于 $\GL(V)$ 传递作用在 $\Omega$ 上，$N$ 指标为 $7$。所有 $N$ 的共轭的交固定所有点，因此是单位群。于是 $H\nleq N$，故 $\GL(V)=HN$。由 $|H:H\cap N|=|HN:N|$，得 $7\mid |H|$。由于 $\GL(V)$ 同构于 $S_7$ 的子群，且 $S_7$ 的 Sylow $7$-子群正规化子阶为 $42$，$\GL(V)$ 没有正规 Sylow $7$-子群，故由 Sylow 定理 $n_7(\GL(V))=8$。若 $H$ 有正规 Sylow $7$-子群，则它在 $H$ 中特征，因而在 $\GL(V)$ 中正规；所以 $H$ 也没有唯一 Sylow $7$-子群。由于 $n_7(H)\equiv1\pmod7$ 且 $n_7(H)\leq n_7(\GL(V))=8$，必有 $n_7(H)=8$。于是 $8\mid |H|$，从而 $56\mid |H|$；因 $H$ 真，必有 $|H|=56$。通常的计数论证说明 $H$ 有正规（故特征）Sylow $2$-子群，于是该子群在 $\GL(V)$ 中正规。但这会使 $\GL(V)$ 有唯一 Sylow $2$-子群。另一方面，$\GL_3(\FF_2)$ 中上三角矩阵组与下三角矩阵组都是阶为 $8$ 的子群，矛盾。于是证明如下定理。

\medskip
\noindent\textbf{定理 15.} 在同构意义下，阶为 $168$ 的单群唯一，即 $\GL_3(\FF_2)$；它也是射影平面 $\mathcal F$ 的自同构群。

注意，我们也可以把 $W_j$ 称为点、把 $U_i$ 称为直线。这种点线“对偶性”连同阶为 $168$ 单群的唯一性，可用来证明 $G$ 存在一个外自同构，它交换点与直线，即把 $U$ 共轭到 $W$。

许多有限单群族可用类似方法分类。在更一般的背景中，称为建筑的几何结构扮演射影平面的角色（射影平面是 $A_2$ 型建筑的特例）。在这种语境中，子群 $N_G(U)$ 与 $N_G(W)$ 是 $G$ 的抛物子群，而 $U,W$ 分别是它们的幺幂根基。特别地，所有单线性群都可由其抛物子群的结构与交来刻画，等价地，也可由它们在相关建筑上的作用来刻画。

\subsection*{关于群的存在性问题的注记}

如同数学其他领域（例如微分方程理论）中一样，我们可以假设某个数学系统存在（例如某方程的解），并由此推导关于该系统的大量信息。一般来说，如果经过相当努力仍未从初始假设推出矛盾，人们会开始怀疑确有系统满足这些假设。然而，任何数量的一致数据都不能证明存在性。假设我们分析一个阶为 $3^3\cdot7\cdot13\cdot409$ 的假想单群 $G$，类似于上面对阶为 $168$ 的单群的分析。经过一定工作可证明可能的 Sylow 数唯一：
\[
  n_3=7\cdot409,
  \qquad
  n_7=3^3\cdot13\cdot409,
  \qquad
  n_{13}=3^2\cdot7\cdot409,
  \qquad
  n_{409}=3^3\cdot7\cdot13.
\]
还可证明这样的 $G$ 没有阶为 $pq$（$p,q$ 为不同素数）的元素，没有阶为 $9$ 的元素，并且不同 Sylow 子群交平凡。于是计数所有 Sylow $p$-子群中的元素，恰好得到 $|G|$。此时我们会得到 $G$ 的完整子群结构与类方程，可能猜想存在这样的单群。但 Feit--Thompson 定理断言不存在奇合数阶单群。（不过要注意，阶为 $3^3\cdot7\cdot13\cdot409$ 的可能单群的这种配置正是 Feit--Thompson 定理证明中必须处理的情形之一，因此在这个例子中引用该结果实际上有循环之嫌。我们将在第 19.3 节证明不存在这种阶的单群；也见练习 29。）这里的要点是：即便我们拥有和阶 $168$ 情形一样多的数据，也不能在没有新技术的情况下证明存在。

当处理非单群时，我们至少有一种由较小群构造较大群的方法：半直积。尽管这种方法相当受限，它仍体现了非单群可以以某种构造方式由较小群拼合而成。对单群，这一过程完全失效；因此这条技术分界线强化了 H\"older 纲领的重要性：确定单群，并理解这些群如何组合成更大的群。

有限单群研究如阶 $168$ 的讨论所示，使用了许多与非单群研究相同的工具（用以剖析子群结构等），但还需要其他构造技术。正如上面提到的，这些技术常常涉及代数或几何方法，把单群构造成某些本身具有内在兴趣的数学结构的自同构群，从而以迷人的方式把群论与数学和科学的其他领域联系起来。因此，尽管我们在有限群分析中已经走了很远，但这一分支中仍有许多领域只是略作触及。

无限群的分析通常涉及相当不同的方法，下一节将介绍其中的一些。

\subsection*{练习}
\paragraph{计数元素}
\begin{enumerate}[label=\arabic*.]
  \item 证明：对固定 $P\in\Syl_p(G)$，若对所有 $R\in\Syl_p(G)-\{P\}$ 都有 $P\cap R=1$，则任意两个不同 Sylow $p$-子群 $P_1,P_2$ 都有 $P_1\cap P_2=1$。推出在这种情形下，$G$ 中非单位的 $p$-幂阶元素数为 $(|P|-1)|G:N_G(P)|$。
  \item 在群 $S_3\times S_3$ 中展示一对交为单位的 Sylow $2$-子群，并展示另一对交为阶 $2$ 的群的 Sylow $2$-子群。
  \item 证明若 $|G|=380$，则 $G$ 不单。[只需计数奇素数阶元素。]
  \item 证明不存在阶为 $80,351,3875,5313$ 的单群。
  \item 设 $G$ 是阶为 $pm$ 的可解群，其中 $p$ 是不整除 $m$ 的素数，$P\in\Syl_p(G)$。若 $N_G(P)=P$，证明 $G$ 有一个阶为 $m$ 的正规子群。证明中哪里用到了 $G$ 的可解性？
\end{enumerate}

\paragraph{利用小指标子群}
\begin{enumerate}[label=\arabic*.,resume]
  \item 证明不存在阶为 $2205,4125,5103,6545,6435$ 的单群。
\end{enumerate}

\paragraph{置换表示}
\begin{enumerate}[label=\arabic*.,resume]
  \item 证明不存在阶为 $1755$ 或 $5265$ 的单群。[用 Sylow $3$-子群说明 $G\leq S_{13}$，并考察 Sylow $13$-子群的正规化子。]
  \item 证明不存在阶为 $792$ 或 $918$ 的单群。
  \item 证明不存在阶为 $336$ 的单群。
\end{enumerate}

\paragraph{让不同 $p$-子群相互制约}
\begin{enumerate}[label=\arabic*.,resume]
  \item 证明不存在阶为 $4095,4389,5313,6669$ 的单群。
  \item 证明不存在阶为 $4851$ 或 $5145$ 的单群。
  \item 证明不存在阶为 $9555$ 的单群。[令 $Q\in\Syl_{13}(G)$，$P\in\Syl_7(N_G(Q))$。证明 $Q\normal N_G(P)$，并说明为何矛盾。]
\end{enumerate}

\paragraph{Sylow 交的正规化子}
\begin{enumerate}[label=\arabic*.,resume]
  \item 设 $G$ 有不止一个 Sylow $p$-子群。取一对不同 Sylow $p$-子群 $P,Q$，使 $|P\cap Q|$ 在所有这样的对中最大。证明 $N_G(P\cap Q)$ 有不止一个 Sylow $p$-子群，且 $N_G(P\cap Q)$ 的任意两个不同 Sylow $p$-子群都交于 $P\cap Q$。（因此 $|N_G(P\cap Q)|$ 被 $p\cdot |P\cap Q|$ 以及某个不同于 $p$ 的素数整除。注意 $N_G(P\cap Q)$ 的 Sylow $p$-子群不一定是 $G$ 的 Sylow 子群。）
  \item 证明不存在阶为 $144,525,2025,3159$ 的单群。
\end{enumerate}

\paragraph{一般练习}
\begin{enumerate}[label=\arabic*.,resume]
  \item 分类阶为 $105$ 的群。
  \item 证明不存在奇阶 $<10000$ 的非阿贝尔单群。
  \item \begin{enumerate}[label=(\alph*)]
    \item 证明不存在阶为 $420$ 的单群。
    \item 证明除阶 $2,60,168,360$ 外，不存在偶阶 $<500$ 的单群。
  \end{enumerate}
  \item 证明若 $G$ 阶为 $36$，则 $G$ 有正规 Sylow $2$-子群或正规 Sylow $3$-子群。
  \item 证明没有阶为 $6$ 子群的阶 $12$ 群同构于 $A_4$。
  \item 证明没有阶为 $6$ 元素的阶 $24$ 群同构于 $S_4$。
  \item 推广引理 13：证明若 $n_p\not\equiv1\pmod{p^k}$，则存在不同 Sylow $p$-子群 $P,R$，使 $P\cap R$ 在二者中指标至多为 $p^{k-1}$。
  \item 在所有不同 Sylow $p$-子群对中，取 $P,R$ 使 $|P\cap R|$ 最大。证明 $N_G(P\cap R)$ 不是 $p$-群。
  \item 设 $A,B$ 是 $G$ 的 Sylow $p$-子群 $P$ 的正规子集。证明若 $A,B$ 在 $G$ 中共轭，则它们在 $N_G(P)$ 中共轭。
  \item 设 $G$ 是阶为 $pqr$ 的群，其中 $p<q<r$ 为素数。证明 $G$ 的 Sylow $r$-子群正规。
  \item 设 $G$ 是阶为 $p^2qr$ 的单群，其中 $p,q,r$ 为素数。证明 $|G|=60$。
  \item 判断并证明或给出反例：若 $G$ 是阶为 $168$ 的群且有不止一个 Sylow $7$-子群，则 $G$ 单。
  \item 证明若 $\mathcal F$ 是任意满足阶 $168$ 单群小节中 (11)--(13) 的点线集合，则 $\mathcal F$ 的关联图唯一确定，并与图 \ref{fig:fano} 相同（至多重标记点和线）。[取一条直线和一个不在该直线上的点。把这条直线画成一个等边三角形的底边，把该点画成不在底边上的顶点。用公理证明剩余点与直线的关联关系于是如图 \ref{fig:fano} 那样唯一确定。]
  \item 设 $G$ 是阶为 $3^3\cdot7\cdot13\cdot409$ 的单群。对每个 $p\in\{3,7,13,409\}$ 计算 $n_p$ 的所有允许值，并化简到每个 $n_p$ 都只有唯一可能值的情形。
  \item 利用假想阶为 $3^3\cdot7\cdot13\cdot409$ 单群的 Sylow 数信息，证明不存在这样的群。[使用次数为 $819$ 的置换表示。]
  \item 假设 $G$ 是阶为 $720$ 的单群。尽可能多地求出 $G$ 的性质（Sylow 数、Sylow 子群同构类型、共轭类等）。是否存在这样的群？
\end{enumerate}

\section{自由群简介}

本节介绍所谓自由群的基本理论。这将使我们能够精确化前几章中使用过的“生成元”和“关系”的概念。本节结果只依赖同态的基本理论。

由集合 $S$ 生成的自由群 $F(S)$ 的基本思想是：$S$ 中元素不满足任何关系（$S$ 是“无关系”的）。例如，若 $S=\{a,b\}$，则两个生成元 $a,b$ 上的自由群的元素形如 $a,aa,ab,abab,ba b$ 等，称为 $a,b$ 中的字，还包括这些元素的逆；所有这些元素都视为不同。若把相同项合并，可得到熟悉形式如 $a^3b^{-2},aba^{-1}b^2$ 等。这样的元素通过连接它们的字相乘。直观上可先把 $F(S)$ 定义为 $S$ 中所有字的集合，并以连接作为乘法。虽然本质上我们就是这样做，但为了证明连接运算良定义且满足结合律，需要更形式化的构造。毕竟，即使熟悉的记号 $a^n$ 表示 $a\cdot a\cdots a$（$n$ 项），也只有在知道乘积与括号放置无关后才合法。

自由群没有关系这一事实的一个重要性质是：从集合 $S$ 到任意群 $G$ 的任意映射，都唯一地扩张为从群 $F(S)$ 到 $G$ 的同态。基本原因是我们指定了生成元的像，而其他元素的像由同态性质唯一决定；没有关系需要担心，因而生成元的像可任意指定。这通常称为自由群的泛性质，并且事实上刻画了 $F(S)$。

“自由性”的概念出现在许多代数系统中。在线性代数中，当代数系统是向量空间时，$F(S)$ 就是以 $S$ 为基的向量空间。每个向量都是 $S$ 中元素的唯一线性组合（类似于“字”），而从基 $S$ 到另一个向量空间 $V$ 的任意集合映射都唯一扩张为线性变换。

在交换图语言中，自由群的泛性质可写作：给定集合映射 $\varphi:S\to G$，其中 $G$ 是群，存在唯一同态 $\Phi:F(S)\to G$ 使 $\Phi|_S=\varphi$，即图
\[
\begin{array}{ccc}
S & \longrightarrow & F(S)\\
 & \searrow_{\varphi} & \downarrow^{\Phi}\\
 & & G
\end{array}
\]
交换。

下面构造 $F(S)$。构造中的唯一困难是验证字的连接运算良定义且满足结合律。为证明结合律，我们回到最基本层面：字中所有指数都是 $\pm1$。

先为 $S$ 的元素引入逆和单位元。令 $S^{-1}$ 为一个与 $S$ 不交的集合，并有一个从 $S$ 到 $S^{-1}$ 的双射。对每个 $s\in S$，把相应元素记为 $s^{-1}$；对每个 $t\in S^{-1}$，把相应的 $S$ 中元素记为 $t^{-1}$，于是 $(s^{-1})^{-1}=s$。再取一个不在 $S\cup S^{-1}$ 中的单点集 $\{1\}$，令 $1^{-1}=1$，并对任意 $x\in S\cup S^{-1}\cup\{1\}$ 令 $(x^{-1})^{-1}=x$。

$S$ 上的一个\emph{字}定义为一个序列
\[
  (s_1,s_2,s_3,\ldots),
  \qquad s_i\in S\cup S^{-1}\cup\{1\},
\]
且对充分大的 $i$ 有 $s_i=1$。换言之，可把一个字看作由 $S$ 中元素及其逆组成的有限乘积（允许重复）。为保证表达唯一，我们只考虑相邻项之间没有显然消去的字。字 $(s_1,s_2,s_3,\ldots)$ 称为\emph{约化的}，若
\begin{enumerate}[label=(\arabic*)]
  \item 当 $s_i\neq1$ 时，$s_{i+1}\neq s_i^{-1}$；
  \item 若某个 $k$ 有 $s_k=1$，则对所有 $i\geq k$，$s_i=1$。
\end{enumerate}
约化字 $(1,1,1,\ldots)$ 称为空字，记为 $1$。以后把约化字
\[
  (s_1^{\varepsilon_1},s_2^{\varepsilon_2},\ldots,s_n^{\varepsilon_n},1,1,\ldots)
\]
简写为
\[
  s_1^{\varepsilon_1}s_2^{\varepsilon_2}\cdots s_n^{\varepsilon_n},
  \qquad s_i\in S,
  \quad \varepsilon_i=\pm1.
\]
令 $F(S)$ 为 $S$ 上所有约化字的集合，并通过
\[
  s\longmapsto (s,1,1,1,\ldots)
\]
把 $S$ 嵌入 $F(S)$。在此嵌入下把 $S$ 与其像等同。若 $S=\varnothing$，则 $F(S)=\{1\}$。

现在在 $F(S)$ 上定义二元运算。技术困难是保证两个约化字的乘积仍为约化字。定义虽然形式化，但只是“相邻逆元逐次消去”的规则（例如，$ab^{-1}a$ 乘以 $a^{-1}ba$ 应约化为 $aa$）。设
\[
  r_1^{\delta_1}\cdots r_m^{\delta_m},
  \qquad
  s_1^{\varepsilon_1}\cdots s_n^{\varepsilon_n}
\]
为约化字，并先假设 $m\leq n$。令 $k$ 为范围 $1\leq k\leq m+1$ 中使
\[
  s_k^{\varepsilon_k}\neq r_{m-k+1}^{-\delta_{m-k+1}}
\]
成立的最小整数。则这两个约化字的乘积定义为
\[
\begin{cases}
 r_1^{\delta_1}\cdots r_{m-k+1}^{\delta_{m-k+1}}
 s_k^{\varepsilon_k}s_{k+1}^{\varepsilon_{k+1}}\cdots s_n^{\varepsilon_n}, & k\leq m,\\[2mm]
 s_{m+1}^{\varepsilon_{m+1}}\cdots s_n^{\varepsilon_n}, & k=m+1\leq n,\\[2mm]
 1, & k=m+1\text{ 且 }m=n.
\end{cases}
\]
当 $m\geq n$ 时类似地定义乘积。因此无论哪种情形，乘积都是一个约化字，由此得到 $F(S)$ 上的二元运算。

\medskip
\noindent\textbf{定理 16.} $F(S)$ 在上述二元运算下是一个群。

\begin{proof}
容易检查空字 $1$ 是单位元，约化字
\[
  s_1^{\varepsilon_1}s_2^{\varepsilon_2}\cdots s_n^{\varepsilon_n}
\]
的逆为
\[
  s_n^{-\varepsilon_n}s_{n-1}^{-\varepsilon_{n-1}}\cdots s_1^{-\varepsilon_1}.
\]
证明的困难部分是结合律。可对所涉字的“长度”作归纳并分情况讨论；也可如下进行。对每个 $s\in S\cup S^{-1}\cup\{1\}$，定义 $\alpha_s:F(S)\to F(S)$ 为
\[
\alpha_s(s_1^{\varepsilon_1}\cdots s_n^{\varepsilon_n})=
\begin{cases}
 ss_1^{\varepsilon_1}\cdots s_n^{\varepsilon_n}, & s_1^{\varepsilon_1}\neq s^{-1},\\[1mm]
 s_2^{\varepsilon_2}\cdots s_n^{\varepsilon_n}, & s_1^{\varepsilon_1}=s^{-1}.
\end{cases}
\]
由于 $\alpha_{s^{-1}}\circ\alpha_s$ 是 $F(S)$ 到自身的恒等映射，$\alpha_s$ 是 $F(S)$ 的一个置换。令 $A(S)$ 为集合 $F(S)$ 上的对称群中由 $\{\alpha_s\mid s\in S\}$ 生成的子群。容易看出映射
\[
  s_1^{\varepsilon_1}s_2^{\varepsilon_2}\cdots s_n^{\varepsilon_n}
  \longmapsto
  \alpha_{s_1^{\varepsilon_1}}\circ\alpha_{s_2^{\varepsilon_2}}\circ\cdots\circ\alpha_{s_n^{\varepsilon_n}}
\]
是 $F(S)$ 与 $A(S)$ 之间的集合双射，并保持它们的二元运算。由于 $A(S)$ 是群，因而满足结合律，所以 $F(S)$ 也满足结合律。
\end{proof}

自由群的泛性质现在立即得到。

\medskip
\noindent\textbf{定理 17.} 设 $G$ 是群，$S$ 是集合，$\varphi:S\to G$ 是集合映射。则存在唯一的群同态 $\Phi:F(S)\to G$ 使上述图交换。

\begin{proof}
若 $\Phi$ 是同态，则必须满足
\[
  \Phi(s_1^{\varepsilon_1}s_2^{\varepsilon_2}\cdots s_n^{\varepsilon_n})
  =\varphi(s_1)^{\varepsilon_1}\varphi(s_2)^{\varepsilon_2}\cdots\varphi(s_n)^{\varepsilon_n},
\]
这证明唯一性。直接检查该公式确实定义了一个同态，得到存在性。
\end{proof}

\medskip
\noindent\textbf{推论 18.} $F(S)$ 在保持集合 $S$ 上恒等的唯一同构意义下唯一。

\begin{proof}
这由泛性质推出。设 $F(S)$ 与 $F'(S)$ 是两个由 $S$ 生成的自由群。由于 $S$ 包含于二者，由定理的泛性质得到唯一同态
\[
  \Phi:F(S)\to F'(S),
  \qquad
  \Phi':F'(S)\to F(S),
\]
且二者在 $S$ 上都是恒等。复合 $\Phi'\Phi$ 是 $F(S)$ 到自身、在 $S$ 上恒等的同态；由唯一性，它必须是恒等映射。同理 $\Phi\Phi'$ 也是恒等。因此 $\Phi$ 是同构。
\end{proof}

\medskip
\noindent\textbf{定义.} 群 $F(S)$ 称为集合 $S$ 上的\emph{自由群}。若存在集合 $S$ 使 $F=F(S)$，则称 $F$ 是\emph{自由群}；此时称 $S$ 为 $F$ 的一组\emph{自由生成元}（或\emph{自由基}）。$S$ 的基数称为该自由群的\emph{秩}。

现在可以用指数记号简化自由群中的表达式，例如写 $a^3b^{-2}$ 代替形式约化字 $aaab^{-1}b^{-1}$。不过，在 $\{a,b\}$ 上的自由群中，像 $aba$ 这样的表达式不能再简化。下面提到该领域中的一个重要定理。

\medskip
\noindent\textbf{定理 19（Schreier）.} 自由群的子群仍是自由群。

这个定理的证明并不平凡，这里不包括证明。使用覆盖空间可以给出一个漂亮的证明（参见 J.-P. Serre, \emph{Trees}, Springer-Verlag, 1980）。

\subsection*{表示}

设 $G$ 是任意群。则 $G$ 是某个自由群的同态像：取 $S=G$，并取 $\varphi$ 为从集合 $G$ 到群 $G$ 的恒等映射；定理 17 给出一个从 $F(G)$ 到 $G$ 的满同态。更一般地，若 $S$ 是 $G$ 的任意生成子集，即 $G=\langle S\rangle$，则也存在唯一的满同态 $F(S)\to G$，它在 $S$ 上为恒等。注意现在可独立表述“子集生成群”：$G=\langle S\rangle$ 当且仅当把 $S$ 到 $G$ 的恒等映射扩张而得的映射 $\pi:F(S)\to G$ 是满射。

\medskip
\noindent\textbf{定义.} 设 $S$ 是群 $G$ 的子集，且 $G=\langle S\rangle$。
\begin{enumerate}[label=(\arabic*)]
  \item $G$ 的一个\emph{表示}是一个二元组 $(S,R)$，其中 $R$ 是 $F(S)$ 中的一组字，使得 $F(S)$ 中 $\langle R\rangle$ 的正规闭包（即包含 $\langle R\rangle$ 的最小正规子群）等于同态 $\pi:F(S)\to G$ 的核（这里 $\pi$ 扩张 $S$ 到 $S$ 的恒等映射）。$S$ 的元素称为 $G$ 的\emph{生成元}，$R$ 的元素称为 $G$ 的\emph{关系}。
  \item 若存在一个表示 $(S,R)$ 使 $S$ 有限，则称 $G$ 是\emph{有限生成的}；若存在一个表示 $(S,R)$ 使 $S$ 与 $R$ 都有限，则称 $G$ 是\emph{有限表示的}。
\end{enumerate}

注意，若 $(S,R)$ 是一个表示，则映射 $F(S)\to G$ 的核不是 $\langle R\rangle$ 本身，而是由 $R$ 及其所有共轭生成的更大的正规闭包。即使固定集合 $S$，一个群也会有许多不同表示（例如可以加入冗余关系）。若 $G$ 有有限表示，$S=\{s_1,
\ldots,s_n\}$ 且 $R=\{w_1,
\ldots,w_k\}$，则按前几章写作
\[
  G=\langle s_1,s_2,\ldots,s_n\mid w_1=w_2=\cdots=w_k=1\rangle.
\]
若 $w$ 是字 $w_1w_2^{-1}$，则把 $w=1$ 写作 $w_1=w_2$。

\paragraph{例}
\begin{enumerate}[label=(\arabic*)]
  \item 每个有限群都是有限表示的。设 $G=\{g_1,
\ldots,g_n\}$。取 $S=G$，令 $\pi:F(S)\to G$ 扩张 $S$ 上的恒等映射。令 $R_0$ 为所有字 $g_ig_jg_k^{-1}$ 的集合，其中 $g_ig_j=g_k$ 在 $G$ 中成立。显然 $R_0\subseteq\ker\pi$。若 $N$ 是 $R_0$ 在 $F(S)$ 中的正规闭包，并令 $\overline G=F(S)/N$，则 $G$ 是 $\overline G$ 的同态像。而且元素集合 $\{\overline g_i\mid i=1,
\ldots,n\}$ 在乘法下封闭；由于它生成 $\overline G$，必有 $|\overline G|\leq |G|$，从而 $N=\ker\pi$，所以 $(S,R_0)$ 是 $G$ 的一个表示。

  这说明给定有限群 $G$ 的一个表示 $(S,R)$ 的充分条件：
  \begin{enumerate}[label=(\roman*)]
    \item $S$ 必须生成 $G$；
    \item 任意由 $S$ 生成并满足 $R$ 中关系的群，其阶都不超过 $|G|$。
  \end{enumerate}

  \item 阿贝尔群很容易表示。例如
  \[
    \ZZ\cong F(\{a\})=\langle a\rangle,
  \]
  \[
    \ZZ\times\ZZ\cong \langle a,b\mid [a,b]=1\rangle,
  \]
  \[
    \ZZ_n\times\ZZ_m\cong \langle a,b\mid a^n=b^m=[a,b]=1\rangle.
  \]
  回忆 $[a,b]=a^{-1}b^{-1}ab$。

  \item 前几章中一些熟悉的非阿贝尔群有简单表示：
  \[
    D_{2n}=\langle r,s\mid r^n=s^2=1,\ s^{-1}rs=r^{-1}\rangle,
  \]
  \[
    Q_8=\langle i,j\mid i^4=1,
\ j^2=i^2,
\ j^{-1}ij=i^{-1}\rangle.
  \]
  例如，要检查 $D_{2n}$ 的表示，注意这些关系说明该群有一个由 $r$ 生成、阶至多为 $n$ 的正规子群，其商由 $s$ 生成且阶至多为 $2$。因此任何满足这些生成元与关系的群阶至多为 $2n$。由于我们已经知道存在阶为 $2n$ 并满足这些条件的群 $D_{2n}$，故该抽象表示就是 $D_{2n}$。

  \item 如第 1.2 节所述，一般而言，甚至判断给定生成元与关系所定义的群是否为单位群都极其困难。例如下面两个表示中，第一个群是无限群，第二个是单位群：
  \[
  \langle x_1,x_2,x_3,x_4\mid x_2^{-1}x_1x_2=x_1^2,
  x_3^{-1}x_2x_3=x_2^2,
  x_4^{-1}x_3x_4=x_3^2,
  x_1^{-1}x_4x_1=x_4^2\rangle,
  \]
  \[
  \langle x_1,x_2,x_3\mid x_2x_1x_2^{-1}=x_1^2,
  x_3x_2x_3^{-1}=x_2^2,
  x_1x_3x_1^{-1}=x_3^2\rangle.
  \]

  \item 容易看出 $S_n$ 由相邻换位 $(12),(23),\ldots,(n-1\ n)$ 生成，并且满足关系
  \[
    ((i\ i+1)(i+1\ i+2))^3=1,
  \]
  以及
  \[
    [(i\ i+1),(j\ j+1)]=1\quad\text{当 }|i-j|\geq2.
  \]
  可对 $n$ 归纳证明这些关系构成 $S_n$ 的一个表示：
  \[
    S_n\cong\langle t_1,
\ldots,t_{n-1}\mid t_i^2=1,
\ (t_it_{i+1})^3=1,
\ [t_i,t_j]=1\text{ whenever }|i-j|\geq2\rangle.
  \]
\end{enumerate}

如第 1.6 节所述，可用群的表示寻找群间同态或群的自同构。例如在分类阶为 $6$ 的群时，我们证明任意非阿贝尔阶 $6$ 群由一个阶为 $3$ 的元素和一个阶为 $2$、使其取逆的元素生成；因此存在从 $S_3$ 到任意非阿贝尔阶 $6$ 群的满同态（比较阶后为同构）。更一般地，设 $G$ 由生成元 $a,b$ 和关系 $r_1,
\ldots,r_k$ 表示。若 $H$ 中元素 $a',b'$ 满足这些关系，则存在从 $G$ 到 $H$ 的同态。具体地，若 $\pi:F(\{a,b\})\to G$ 是表示同态，则可定义 $\pi':F(\{a,b\})\to H$，使 $\pi'(a)=a'$ 且 $\pi'(b)=b'$。由于 $a',b'$ 满足 $G$ 的关系，有 $\ker\pi\leq\ker\pi'$，所以 $\pi'$ 通过 $\ker\pi$ 分解，从而得到
\[
  G\cong F(\{a,b\})/\ker\pi\longrightarrow H.
\]
特别地，若 $\langle a',b'\rangle=H=G$，则该同态是 $G$ 的自同构。反过来，任一自同构必须把一组生成元送到另一组满足相同关系的生成元。

例如
\[
  D_8=\langle a,b\mid a^2=b^4=1,
\ aba=b^{-1}\rangle.
\]
任意一对元素 $a',b'$，其中 $a'$ 是非中心的阶 $2$ 元素，$b'$ 阶为 $4$，都满足相同关系。$D_8$ 有四个非中心阶 $2$ 元素和两个阶 $4$ 元素，因此 $D_8$ 有 $8$ 个自同构。

同样，$Q_8$ 中任意一对阶为 $4$、且不相等也不互为逆的元素必生成 $Q_8$，并满足例 3 中的关系。容易检查这样的对有 $24$ 个，所以
\[
  |\Aut(Q_8)|=24.
\]

自由对象也可在许多（但不是所有）其他范畴中构造。例如，幺半群是带有二元运算、满足除逆元存在性外所有群公理的集合。幺半群范畴中的自由对象在理论计算机科学中起基础作用，用以刻画机器（Turing 机等）的行为。后面章节还会遇到自由代数（即多项式代数）与自由模。

\subsection*{练习}
\begin{enumerate}[label=\arabic*.]
  \item 设 $F_1,F_2$ 是有限秩自由群。证明 $F_1\cong F_2$ 当且仅当它们有相同秩。若要把证明推广到无限秩（结论仍为真），还需要哪些事实？
  \item 证明若 $|S|>1$，则 $F(S)$ 非阿贝尔。
  \item 证明两个生成元上的自由群的换位子子群不是有限生成的（特别地，有限生成群的子群不必有限生成）。
  \item 证明自由群的每个非单位元素都有无限阶。
  \item 用两个生成元给出 $A_4$ 的一个有限表示。
  \item 用两个生成元给出 $S_4$ 的一个有限表示。
  \item 证明以下是阶为 $8$ 的四元数群的一个表示：
  \[
    Q_8=\langle a,b\mid a^2=b^2,
\ a^{-1}ba=b^{-1}\rangle.
  \]
  \item 用表示求群 $\ZZ_2\times\ZZ_4$ 和 $\ZZ_4\times\ZZ_4$ 的自同构群的阶。
  \item 证明 $\Aut(Q_8)\cong S_4$。
  \item 本题展示 $S_6$ 的一个非内自同构（因此，连同第 4.4 节练习 19，说明 $|\Aut(S_6):\Inn(S_6)|=2$）。令
  \[
  \begin{aligned}
  t_1'&=(12)(34)(56),&
  t_2'&=(14)(25)(36),\\
  t_3'&=(13)(24)(56),&
  t_4'&=(12)(36)(45),\\
  t_5'&=(14)(23)(56).
  \end{aligned}
  \]
  证明 $t_1',
\ldots,t_5'$ 满足：$(t_i')^2=1$；当 $|i-j|\geq2$ 时 $(t_i't_j')^2=1$；以及对 $i=1,2,3,4$，$(t_i't_{i+1}')^3=1$。推出映射
  \[
    (12)\mapsto t_1',\quad (23)\mapsto t_2',\quad (34)\mapsto t_3',\quad
    (45)\mapsto t_4',\quad (56)\mapsto t_5'
  \]
  扩张为 $S_6$ 的一个自同构，且它显然不是内自同构，因为它不把换位送到换位。
  \item 设 $S$ 是集合。表示 $\langle S\mid [s,t]=1\text{ for all }s,t\in S\rangle$ 所给的群称为 $S$ 上的自由阿贝尔群，记为 $A(S)$。证明 $A(S)$ 有如下泛性质：若 $G$ 是任意阿贝尔群，$\varphi:S\to G$ 是任意集合映射，则存在唯一群同态 $\Phi:A(S)\to G$ 使 $\Phi|_S=\varphi$。推出若 $A$ 是基数为 $n$ 的集合上的自由阿贝尔群，则
  \[
    A\cong \underbrace{\ZZ\times\ZZ\times\cdots\times\ZZ}_{n\text{ 个因子}}.
  \]
  \item 设 $S$ 是集合，$c$ 是正整数。表述 $S$ 上幂零类为 $c$ 的自由幂零群的概念，并证明它相对于幂零类 $\leq c$ 的幂零群具有适当的泛性质。
  \item 证明不存在一个由两个元素生成的幂零群 $N$，使得每个由两个元素生成的幂零群都是 $N$ 的同态像（也就是说，前一题中指定类 $c$ 是必要的）。
\end{enumerate}

% \setchapterintro{群论研究的是具有由单个二元运算定义的代数结构的某些对象的一般性质。环的研究则关心具有两个二元运算（称为加法和乘法）并由分配律联系起来的对象。我们首先研究群的结构理论发展中若干基本要点的类比。特别地，我们引入子环、商环、理想（它们是正规子群的类比）以及环同态。随后我们集中讨论由于存在两个二元运算而自然出现的关于一般环的问题。关于乘法逆元的问题导致域的概念，并最终导致某些具体域的构造，例如有限域。对熟悉的整数环 $\ZZ$ 这类环的算术（整除性、最大公因子等）的研究，在第 8 章中导致素元和唯一分解的概念。第 7 章和第 8 章的结果随后将应用于第 9 章中的多项式环。
% 	
% 	第二部分所发展的环的基本理论，是本书其余四个部分的基石。环作用的理论（模）构成本书第三部分。在那里我们将看到，环的结构如何反映在它所作用的对象的结构中，并由此证明一些强有力的分类定理。环的结构理论，特别是多项式环的结构理论，在第四部分中构成域和域上多项式方程理论的基础。在那里，环论、域论和群论之间丰富的相互作用，会导出许多关于域的结构和多项式根的理论的优美结果。第五部分继续研究环以及环论在几何和扩张理论等主题中的应用。第六部分中，对某些特定类型的环（群环）及其作用对象（模）的研究，再次给出深刻的分类定理；这些定理的后果随后被用来给出关于有限群的新结果和新见解。}
%\part{环论}
% 


\chapter{环导论}

\section{基本定义与例子}

\noindent\textbf{定义.}
\begin{enumerate}[label=(\arabic*)]
  \item 环 $R$ 是带有两个二元运算 $+$ 与 $\times$（称为加法与乘法）的集合，满足下列公理：
  \begin{enumerate}[label=(\roman*)]
    \item $(R,+)$ 是 Abel 群；
    \item $\times$ 满足结合律：$(a\times b)\times c=a\times(b\times c)$ 对所有 $a,b,c\in R$ 成立；
    \item $R$ 中满足分配律：对所有 $a,b,c\in R$，
    \[
      (a+b)\times c=(a\times c)+(b\times c),
      \qquad
      a\times(b+c)=(a\times b)+(a\times c).
    \]
  \end{enumerate}
  \item 若乘法交换，则称环 $R$ 是\emph{交换的}。
  \item 若存在元素 $1\in R$，使
  \[
    1\times a=a\times1=a\qquad\text{对所有 }a\in R
  \]
  成立，则称环 $R$ 有\emph{单位元}（或含有 $1$）。
\end{enumerate}

以后通常把 $a\times b$ 简写为 $ab$。$R$ 的加法单位元总记为 $0$，环元素 $a$ 的加法逆元记为 $-a$。

要求 $R$ 在加法下为群是相当自然的；但要求这个群为 Abel 群看起来或许有些人为。一个动机是：如果环 $R$ 有 $1$，那么加法交换性事实上由分配律强制给出。为说明这一点，用两种不同方式计算 $(1+1)(a+b)$，只使用分配律（不假定加法交换）。一方面得到
\[
  (1+1)(a+b)=1(a+b)+1(a+b)=1a+1b+1a+1b=a+b+a+b,
\]
另一方面得到
\[
  (1+1)(a+b)=(1+1)a+(1+1)b=1a+1a+1b+1b=a+a+b+b.
\]
由于 $R$ 在加法下是群，消去可得 $b+a=a+b$，即 $R$ 在加法下必然交换。

域是最重要的环例之一。注意下面的定义只是第 1.4 节中定义的另一种表述。

\medskip
\noindent\textbf{定义.} 带单位元 $1$ 且 $1\neq0$ 的环 $R$ 称为\emph{除环}（或\emph{斜域}），若每个非零元素 $a\in R$ 都有乘法逆元，即存在 $b\in R$ 使 $ab=ba=1$。交换除环称为\emph{域}。

\medskip
\noindent\textbf{例.}
\begin{enumerate}[label=(\arabic*)]
  \item 最简单的环例子是\emph{平凡环}：取 $R$ 为任意交换群（把群运算记为 $+$），并定义乘法 $a\times b=0$ 对所有 $a,b\in R$ 成立。容易看出这给出一个交换环。特别地，若 $R=\{0\}$ 是平凡群，所得环称为\emph{零环}，记为 $R=0$。除零环外，平凡环没有单位元（$R=0$ 是唯一满足 $1=0$ 的环；我们常通过条件 $1\neq0$ 排除它）。虽然平凡环有两个二元运算，乘法并不给加法群增加新的结构，环论在这里不能给出超出 Abel 群论的信息。
  \item 整数环 $\ZZ$ 在通常加法与乘法下是带单位元（整数 $1$）的交换环。环公理和加法群公理一样，最终来自自然数系统的基本公理。注意在乘法下 $\ZZ-\{0\}$ 不是群（事实上，这个环中有乘法逆元的元素很少）。我们很快会回到逆元问题。
  \item 同样，有理数 $\QQ$、实数 $\RR$、复数 $\CC$ 都是带单位元的交换环（事实上它们都是域）。这些环的环公理最终来自 $\ZZ$ 的环公理。当我们从 $\ZZ$ 构造 $\QQ$（第 7.5 节）以及从 $\RR$ 构造 $\CC$（第 13.1 节例 1）时，将验证这一点；这两种构造都是更一般过程的特例。从 $\QQ$ 构造 $\RR$ 以及验证环公理，则通常在基础分析教材中完成。
  \item 商群 $\ZZ/n\ZZ$ 在剩余类的加法和乘法（通常称为“模算术”）下是带单位元（元素 $\bar1$）的交换环。我们已经看到，加法 Abel 群公理来自商群理论的一般原理。稍后我们将证明其余环公理，特别是剩余类乘法的良定义性，类似地来自商环的一般理论。
  \item \textbf{（实 Hamilton 四元数）} 令 $\HH$ 为所有形如
  \[
    a+bi+cj+dk\qquad (a,b,c,d\in\RR)
  \]
  的元素组成的集合（粗略说，即以 $1,i,j,k$ 为“变量”且实系数的多项式）。加法逐分量定义：
  \[
  (a+bi+cj+dk)+(a'+b'i+c'j+d'k)
  =(a+a')+(b+b')i+(c+c')j+(d+d')k.
  \]
  乘法由展开 $(a+bi+cj+dk)(a'+b'i+c'j+d'k)$、使用分配律（注意项的次序）并用关系
  \[
    i^2=j^2=k^2=-1,\\qquad ij=k,
    \quad jk=i,\quad ki=j,
  \]
  \[
    ji=-k,
    \quad kj=-i,
    \quad ik=-j
  \]
  化简来定义，其中实系数与 $i,j,k$ 交换。例如
  \[
  \begin{aligned}
    (1+i+2j)(j+k)
    &=1(j+k)+i(j+k)+2j(j+k)\\
    &=j+k+ij+ik+2j^2+2jk\\
    &=j+k+k+(-j)+2(-1)+2i\\
    &=-2+2i+2k.
  \end{aligned}
  \]
  $\HH$ 是环这一事实可由直接但较长的公理验证证明（乘法结合律尤其繁琐）。Hamilton 四元数是带单位元的非交换环，其中 $1=1+0i+0j+0k$。同样，把上面的 $a,b,c,d$ 取为有理数可定义有理 Hamilton 四元数环。实和有理 Hamilton 四元数都是除环，非零元素的逆元由
  \[
    (a+bi+cj+dk)^{-1}=\frac{a-bi-cj-dk}{a^2+b^2+c^2+d^2}
  \]
  给出。
  \item 一个重要的环类由函数环给出。设 $X$ 是非空集合，$A$ 是任意环。所有集合映射 $f:X\to A$ 的集合 $R$ 在逐点加法和逐点乘法下构成环：
  \[
    (f+g)(x)=f(x)+g(x),\qquad (fg)(x)=f(x)g(x).
  \]
  $R$ 的每条环公理都直接来自 $A$ 的相应公理。$R$ 交换当且仅当 $A$ 交换；$R$ 有 $1$ 当且仅当 $A$ 有 $1$（此时 $R$ 的 $1$ 必然是 $X$ 上恒等取值为 $1$ 的函数）。

  如果 $X$ 与 $A$ 有更多结构，还可形成尊重这些结构的函数环。例如，若 $A=\RR$，且 $X=[0,1]$ 是 $\RR$ 中闭区间，则所有从 $[0,1]$ 到 $\RR$ 的连续函数构成一个带 $1$ 的交换环（这需要基本极限定理保证连续函数之和与积仍连续）。
  \item 没有单位元的环的一个例子是偶整数环 $2\ZZ$，在通常的整数加法和乘法下构成环（偶整数的和与积仍是偶整数）。

  另一个自然出现在分析中的例子如下。函数 $f:\RR\to\RR$ 称为有\emph{紧支撑}，若存在实数 $a,b$（依赖于 $f$），使 $x\notin[a,b]$ 时 $f(x)=0$，即 $f$ 在某个有界区间外为零。所有有紧支撑的函数 $f:\RR\to\RR$ 构成没有单位元的交换环（单位函数不可能有紧支撑）。同样，所有有紧支撑的连续函数 $f:\RR\to\RR$ 也构成没有单位元的交换环。
\end{enumerate}

下一节我们给出从给定环构造“较大”环的三种重要方式（类似上面例 6），从而大大扩充例子列表。在此之前，我们先给出任意环的一些基本性质。整数环 $\ZZ$ 是一个很好的参考例子，尽管它具有比一般环更多的代数结构（例如它交换且有单位元）。不过，它的基本算术在一般环中仍成立，如下所示。

\medskip
\noindent\textbf{命题 1.} 设 $R$ 是环。则：
\begin{enumerate}[label=(\arabic*)]
  \item 对所有 $a\in R$，$0a=a0=0$；
  \item 对所有 $a,b\in R$，$(-a)b=a(-b)=-(ab)$；
  \item 对所有 $a,b\in R$，$(-a)(-b)=ab$；
  \item 若 $R$ 有单位元 $1$，则单位元唯一，且 $-a=(-1)a$。
\end{enumerate}

\begin{proof}
这些都由分配律和加法群中的消去律推出。例如，(1) 来自
\[
  0a=(0+0)a=0a+0a.
\]
(2) 中等式 $(-a)b=-(ab)$ 来自
\[
  ab+(-a)b=(a+(-a))b=0b=0.
\]
其余结论类似，留给读者验证。
\end{proof}

这个命题说明，由于分配律，环的加法结构与乘法结构能很好地相互配合，正如熟悉的整数例子中那样。

不过，与整数不同，一般环可能有许多具有乘法逆元的元素，也可能有非零元素 $a,b$ 其乘积为零。下面给出这些与乘法结构有关的性质的专门名称。

\medskip
\noindent\textbf{定义.} 设 $R$ 是环。
\begin{enumerate}[label=(\arabic*)]
  \item 非零元素 $a\in R$ 称为\emph{零因子}，若存在非零 $b\in R$，使 $ab=0$ 或 $ba=0$。
  \item 假设 $R$ 有单位元 $1\neq0$。元素 $u\in R$ 称为 $R$ 中的\emph{单位}，若存在 $v\in R$，使 $uv=vu=1$。$R$ 中单位全体记为 $R^\times$。
\end{enumerate}

容易看出，环 $R$ 中的单位在乘法下构成群，所以 $R^\times$ 称为 $R$ 的\emph{单位群}。用这个术语，域就是带单位元 $1\neq0$ 的交换环 $F$，其中每个非零元素都是单位，即 $F^\times=F-\{0\}$。

注意零因子不可能是单位。例如若 $a$ 是 $R$ 中单位且 $ab=0$，其中 $b\neq0$，取 $v\in R$ 使 $va=1$，则
\[
  b=1b=(va)b=v(ab)=v0=0,
\]
矛盾。类似地，若 $ba=0$ 对某个非零 $b$ 成立，则 $a$ 不可能是单位。因此域没有零因子。

\medskip
\noindent\textbf{例.}
\begin{enumerate}[label=(\arabic*)]
  \item 整数环 $\ZZ$ 没有零因子，其单位只有 $\pm1$，即 $\ZZ^\times=\{\pm1\}$。注意每个非零整数在较大的环 $\QQ$ 中都有逆元，所以“是单位”这一性质依赖于所处的环。
  \item 设 $n\geq2$。在 $\ZZ/n\ZZ$ 中，与 $n$ 互素的元素 $\bar u$ 是单位（下一章将证明）。另一方面，若非零整数 $a$ 与 $n$ 不互素，则 $\bar a$ 是 $\ZZ/n\ZZ$ 中的零因子。令 $d=(a,n)$，取 $b=n/d$。由假设 $d>1$，故 $0<b<n$，于是 $\bar b\neq\bar0$；但构造表明 $n\mid ab$，即在 $\ZZ/n\ZZ$ 中 $\bar a\bar b=\bar0$。这说明 $\ZZ/n\ZZ$ 中每个非零元素要么是单位，要么是零因子。此外，每个非零元素都是单位当且仅当 $0<a<n$ 的每个整数 $a$ 都与 $n$ 互素；这当且仅当 $n$ 是素数。故 $\ZZ/n\ZZ$ 是域当且仅当 $n$ 是素数。
  \item 若 $R$ 是从闭区间 $[0,1]$ 到 $\RR$ 的所有函数构成的环，则 $R$ 的单位是处处不为零的函数（其逆为函数 $1/f$）。若 $f$ 不是单位且不是零函数，则 $f$ 是零因子：定义
  \[
    g(x)=\begin{cases}0,& f(x)\neq0,\\ 1,& f(x)=0,
    \end{cases}
  \]
  则 $g$ 不是零函数但 $f(x)g(x)=0$ 对所有 $x$ 成立。
  \item 若 $R$ 是从 $[0,1]$ 到 $\RR$ 的所有连续函数构成的环，则单位仍是处处不为零的函数；但现在存在既不是单位也不是零因子的元素。例如 $f(x)=x-\frac12$ 只有一个零点，所以不是单位。另一方面，若 $gf=0$，则 $g$ 在所有 $x\neq\frac12$ 处为零；唯一具有此性质的连续函数是零函数。因此 $f$ 既不是单位也不是零因子。类似地，只有有限（甚至可数）个零点的函数都不是零因子。该环也含有许多零因子。例如令
  \[
    f(x)=\begin{cases}0,&0\leq x\leq\frac12,\\ x-\frac12,&\frac12\leq x\leq1,
    \end{cases}
  \]
  并令 $g(x)=f(1-x)$，则 $f,g$ 是非零连续函数且乘积为零函数。
  \item 设 $D$ 是一个不是 $\QQ$ 中完全平方数的有理数，定义
  \[
    \QQ(\sqrt D)=\{a+b\sqrt D\mid a,b\in\QQ\}\subseteq\CC.
  \]
  该集合显然对减法封闭，恒等式
  \[
    (a+b\sqrt D)(c+d\sqrt D)=(ac+bdD)+(ad+bc)\sqrt D
  \]
  表明它也对乘法封闭。因此 $\QQ(\sqrt D)$ 是 $\CC$ 的子环（若 $D>0$，甚至是 $\RR$ 的子环），特别是带单位元的交换环。容易证明，$D$ 不是平方数这一假设意味着 $\QQ(\sqrt D)$ 的每个元素都能唯一写成 $a+b\sqrt D$ 的形式。该假设还意味着若 $a,b$ 不全为 $0$，则 $a^2-Db^2\neq0$；又因
  \[
    (a+b\sqrt D)(a-b\sqrt D)=a^2-Db^2,
  \]
  若 $a+b\sqrt D\neq0$，则
  \[
    \frac{a-b\sqrt D}{a^2-Db^2}
  \]
  是 $a+b\sqrt D$ 在 $\QQ(\sqrt D)$ 中的逆元。这说明这个交换环中每个非零元素都是单位，即 $\QQ(\sqrt D)$ 是域，称为\emph{二次域}（参见第 13.2 节）。

  有理数 $D$ 可写为 $D=f^2D'$，其中 $f\in\QQ$，$D'$ 是唯一的整数，且不被任何大于 $1$ 的整数的平方整除，即 $D'$ 是 $-1$ 或 $\pm1$ 乘以若干个 $\ZZ$ 中不同素数的乘积。称 $D'$ 为 $D$ 的无平方因子部分。于是 $\sqrt D=f\sqrt{D'}$，从而 $\QQ(\sqrt D)=\QQ(\sqrt{D'})$。因此在定义二次域 $\QQ(\sqrt D)$ 时，不妨假定 $D$ 是无平方因子整数。
\end{enumerate}

具有整数环 $\ZZ$ 的某些特征的环有专门名称。

\medskip
\noindent\textbf{定义.} 带单位元 $1\neq0$ 的交换环称为\emph{整环}，若它没有零因子。

整环中没有零因子给出消去性质。

\medskip
\noindent\textbf{命题 2.} 设 $a,b,c$ 是任意环中的元素，且 $a$ 不是零因子。若 $ab=ac$，则 $a=0$ 或 $b=c$（即若 $a\neq0$，则可以消去 $a$）。特别地，若 $a,b,c$ 是整环中的任意元素，且 $ab=ac$，则 $a=0$ 或 $b=c$。

\begin{proof}
若 $ab=ac$，则 $a(b-c)=0$，故或者 $a=0$，或者 $b-c=0$。第二个断言由第一个断言和整环定义立即得到。
\end{proof}

\medskip
\noindent\textbf{推论 3.} 任意有限整环都是域。

\begin{proof}
设 $R$ 是有限整环，$a$ 是 $R$ 的非零元素。由消去律，映射 $x\mapsto ax$ 是单射。由于 $R$ 有限，此映射也是满射。特别地，存在 $b\in R$ 使 $ab=1$，即 $a$ 是单位。由于 $a$ 是任意非零元素，$R$ 是域。
\end{proof}

Wedderburn 的一个著名结果断言，有限除环必为交换的，即必为域。这个定理的证明在第 13.6 节末的练习中给出提纲。

在第 5 节中，我们会更详细地研究零因子与单位之间的关系。我们将看到，交换环中每个非零且不是零因子的元素，在某个较大的环中都有乘法逆元。这给命题 2 中的消去律提供了另一种视角。

定义了环的概念以后，自然也有子环的概念。

\medskip
\noindent\textbf{定义.} 环 $R$ 的\emph{子环}是 $R$ 的一个加法子群，并且对乘法封闭。

换言之，环 $R$ 的子集 $S$ 是子环，若 $R$ 中的加法和乘法限制到 $S$ 后使 $S$ 成为环。要证明环 $R$ 的一个子集是子环，只需检查它非空，并且对减法和乘法封闭。

\medskip
\noindent\textbf{例.}
\begin{enumerate}[label=(\arabic*)]
  \item $\ZZ$ 是 $\QQ$ 的子环，$\QQ$ 是 $\RR$ 的子环。“是……的子环”这个性质显然具有传递性。
  \item $2\ZZ$ 是 $\ZZ$ 的子环；任意整数 $n$ 的 $n\ZZ$ 也是 $\ZZ$ 的子环。当 $n\geq2$ 时，$\ZZ/n\ZZ$ 不是 $\ZZ$ 的子环（也不是子群）。
  \item 从 $\RR$ 到 $\RR$ 的所有连续函数构成的环，是所有从 $\RR$ 到 $\RR$ 的函数构成的环的子环。从 $\RR$ 到 $\RR$ 的所有可微函数构成的环，是二者的子环。
  \item $S=\ZZ+\ZZ i+\ZZ j+\ZZ k$，即整 Hamilton 四元数，构成实或有理 Hamilton 四元数环的子环：两个这类四元数相乘仍然得到整数系数的四元数。这个环（不是除环）可用于证明若干数论结果。
  \item 若 $R$ 是含有域 $F$ 的单位元的子环，则 $R$ 是整环。反过来也成立：任意整环都包含在某个域中（参见第 5 节）。
\end{enumerate}

\medskip
\noindent\textbf{例（二次整数环）.}
设 $D$ 是无平方因子整数。从加法和乘法可立即看出，子集
\[
  \ZZ[\sqrt D]=\{a+b\sqrt D\mid a,b\in\ZZ\}
\]
构成二次域 $\QQ(\sqrt D)$ 的子环。若 $D\equiv1\pmod4$，则稍大的子集
\[
  \ZZ\!bigg[\frac{1+\sqrt D}{2}\bigg]
  =\left\{a+b\frac{1+\sqrt D}{2}\mid a,b\in\ZZ\right\}
\]
也是子环：对加法封闭显然，而
\[
\left(a+b\frac{1+\sqrt D}{2}\right)
\left(c+d\frac{1+\sqrt D}{2}\right)
=(ac+bd\frac{D-1}{4})+(ad+bc+bd)\frac{1+\sqrt D}{2}
\]
与 $D\equiv1\pmod4$ 一起说明它对乘法封闭。

定义
\[
  \mcO=\mcO_{\QQ(\sqrt D)}=\ZZ[\omega]=\{a+b\omega\mid a,b\in\ZZ\},
\]
其中
\[
  \omega=\begin{cases}
  \sqrt D,&D\equiv2,3\pmod4,\\[2mm]
  \dfrac{1+\sqrt D}{2},&D\equiv1\pmod4.
  \end{cases}
\]
它称为二次域 $\QQ(\sqrt D)$ 的\emph{整数环}。这个术语来自这样一个事实：域 $\QQ(\sqrt D)$ 的子环 $\mcO$ 中的元素有许多类似于有理数域 $\QQ$ 中整数子环 $\ZZ$ 的性质（并且它们正是第 15.3 节中解释的 $\ZZ$ 在 $\QQ(\sqrt D)$ 中的整闭包）。

当 $D=-1$ 时，得到 Gaussian 整数环 $\ZZ[i]$，即复数 $a+bi\in\CC$ 中 $a,b$ 都是整数者。Gauss 在约 1800 年引入这些数，是为了陈述双二次互反律；该定律讨论素数模四次幂之间存在的优美关系。我们很快会看到这个环的代数结构在数论问题中的另一个有用应用。

定义域范数
\[
  N:\QQ(\sqrt D)\to\QQ,
  \qquad N(a+b\sqrt D)=(a+b\sqrt D)(a-b\sqrt D)=a^2-Db^2\in\QQ,
\]
如前所述，当 $a+b\sqrt D\neq0$ 时这个数非零。范数给出域 $\QQ(\sqrt D)$ 中“大小”的度量。例如当 $D=-1$ 时，$a+bi$ 的范数是 $a^2+b^2$，即这个复数作为复平面中向量时长度的平方。我们将在本例和后续例子中使用域范数来建立环 $\mcO$ 的许多性质。

容易验证 $N$ 是乘法性的，即对所有 $\alpha,\beta\in\QQ(\sqrt D)$ 有 $N(\alpha\beta)=N(\alpha)N(\beta)$。在子环 $\mcO$ 上，域范数也容易写成
\[
N(a+b\omega)=(a+b\omega)(a+b\bar\omega)=
\begin{cases}
 a^2-Db^2,&D\equiv2,3\pmod4,\\[1mm]
 a^2+ab+\dfrac{1-D}{4}b^2,&D\equiv1\pmod4,
\end{cases}
\]
其中
\[
  \bar\omega=\begin{cases}
  -\sqrt D,&D\equiv2,3\pmod4,\\[1mm]
  \dfrac{1-\sqrt D}{2},&D\equiv1\pmod4.
  \end{cases}
\]
于是对每个 $\alpha\in\mcO$，$N(\alpha)$ 实际上是整数。

我们可用此范数刻画 $\mcO$ 中的单位。若 $\alpha\in\mcO$ 的域范数 $N(\alpha)=\pm1$，则上式说明 $\alpha^{-1}=\pm\bar\alpha$，仍属于 $\mcO$，故 $\alpha$ 是 $\mcO$ 中的单位。反过来，若 $\alpha$ 是 $\mcO$ 中的单位，设 $\alpha\beta=1$，其中 $\beta\in\mcO$。范数的乘法性给出
\[
  N(\alpha)N(\beta)=N(\alpha\beta)=N(1)=1.
\]
由于 $N(\alpha)$ 和 $N(\beta)$ 都是整数，它们都必须为 $\pm1$。因此
\[
  \alpha\text{ 是 }\mcO\text{ 中的单位，当且仅当 }N(\alpha)=\pm1.
\]
特别地，求整数方程 $x^2-Dy^2=\pm1$（初等数论中称为 Pell 方程）的整数解，本质上等价于求环 $\mcO$ 的单位。

当 $D=-1$ 时，Gaussian 整数 $\ZZ[i]$ 中的单位是满足 $a^2+b^2=\pm1$ 的元素 $a+bi$，其中 $a,b\in\ZZ$，故单位群为 $\{\pm1,\pm i\}$。当 $D=-3$ 时，$\ZZ[(1+\sqrt{-3})/2]$ 中的单位由满足 $a^2+ab+b^2=\pm1$ 的整数 $a,b$ 决定，也即满足 $(2a+b)^2+3b^2=\pm4$；由此容易看出单位群为 $6$ 阶群 $\{\pm1,\pm\rho,\pm\rho^2\}$，其中 $\rho=(-1+\sqrt{-3})/2$。对任何其他 $D<0$，类似地容易看出唯一单位是 $\{\pm1\}$。

相反，当 $D>0$ 时，可以证明单位群 $\mcO^\times$ 总是无限的。例如在环 $\mcO=\ZZ[\sqrt2]$ 中，$1+\sqrt2$ 是一个单位（域范数为 $-1$），且
\[
  \{\pm(1+\sqrt2)^n\mid n\in\ZZ\}
\]
是无限个互异单位（事实上这就是该情形的完整单位群，但证明较困难）。

\subsection*{练习}
在以下各题中，设 $R$ 是带单位元 $1$ 的环。
\begin{enumerate}[label=\arabic*.]
	\item 证明在 $R$ 中有 $(-1)^2=1$。
	\item 证明若 $u$ 是 $R$ 中的单位，则 $-u$ 也是 $R$ 中的单位。
	\item 设 $R$ 是带单位元的环，$S$ 是 $R$ 的含单位元子环。证明：若 $u$ 是 $S$ 中的单位，则 $u$ 是 $R$ 中的单位。举例说明反命题不成立。
	\item 证明任意非空子环族的交仍是子环。
	\item 判定下列各集合中哪些是 $\QQ$ 的子环：
	\begin{enumerate}[label=(\alph*)]
		\item 既约表示时分母为奇数的有理数全体；
		\item 既约表示时分母为偶数的有理数全体；
		\item 非负有理数全体；
		\item 有理数平方全体；
		\item 既约表示时分子为奇数的有理数全体；
		\item 既约表示时分子为偶数的有理数全体。
	\end{enumerate}
	\item 判定下列各集合中哪些是闭区间 $[0,1]$ 到 $\RR$ 的所有函数构成的环的子环：
	\begin{enumerate}[label=(\alph*)]
		\item 所有满足 $f(q)=0$ 对一切 $q\in\QQ\cap[0,1]$ 成立的函数 $f(x)$；
		\item 所有多项式函数；
		\item 只有有限个零点的函数全体，再加上零函数；
		\item 有无穷多个零点的函数全体；
		\item 所有满足 $\lim_{x\to1^-}f(x)=0$ 的函数 $f$；
		\item 函数 $\sin nx$ 与 $\cos mx$ 的所有有理线性组合，其中 $m,n\in\{0,1,2,\ldots\}$。
	\end{enumerate}
	\item 环 $R$ 的中心定义为 $\{z\in R\mid zr=rz\text{ 对所有 }r\in R\}$，即与 $R$ 中每个元素都交换的元素全体。证明环的中心是含单位元的子环。证明除环的中心是域。
	\item 描述实 Hamilton 四元数环 $\HH$ 的中心。证明 $\{a+bi\mid a,b\in\RR\}$ 是 $\HH$ 的子环，且是一个域，但它不包含在 $\HH$ 的中心中。
	\item 对固定元素 $a\in R$，定义 $C(a)=\{r\in R\mid ra=ar\}$。证明 $C(a)$ 是 $R$ 的含 $a$ 子环。证明 $R$ 的中心是所有 $a\in R$ 的子环 $C(a)$ 的交。
	\item 证明若 $D$ 是除环，则对所有 $a\in D$，$C(a)$ 是除环（参见上一题）。
	\item 证明若 $R$ 是整环且某个 $x\in R$ 满足 $x^2=1$，则 $x=\pm1$。
	\item 证明含单位元的域的任意子环都是整环。
	\item 环 $R$ 中元素 $x$ 称为幂零元，若存在 $m\in\ZZ^+$ 使 $x^m=0$。
	\begin{enumerate}[label=(\alph*)]
		\item 若 $n=a^k b$，其中 $a,b$ 是整数，证明 $\bar a\bar b$ 是 $\ZZ/n\ZZ$ 中的幂零元。
		\item 若 $a\in\ZZ$，证明 $\bar a\in\ZZ/n\ZZ$ 是幂零元，当且仅当 $n$ 的每个素因子也整除 $a$。特别地，显式确定 $\ZZ/72\ZZ$ 中的幂零元。
		\item 设 $R$ 是从非空集合 $X$ 到域 $F$ 的所有函数构成的环。证明 $R$ 不含非零幂零元。
	\end{enumerate}
	\item 设 $x$ 是交换环 $R$ 中的幂零元（参见上一题）。
	\begin{enumerate}[label=(\alph*)]
		\item 证明 $x$ 要么为零，要么为零因子。
		\item 证明对所有 $r\in R$，$rx$ 是幂零元。
		\item 证明 $1+x$ 是 $R$ 中的单位。
		\item 推出幂零元与单位之和仍是单位。
	\end{enumerate}
	\item 若环 $R$ 中对所有 $a\in R$ 都有 $a^2=a$，则称 $R$ 为 Boolean 环。证明每个 Boolean 环都是交换环。
	\item 证明唯一的整环型 Boolean 环是 $\ZZ/2\ZZ$。
	\item 设 $R,S$ 是环。证明直积 $R\times S$ 在逐分量加法与乘法下构成环。证明 $R\times S$ 交换当且仅当 $R,S$ 都交换。证明 $R\times S$ 有单位元当且仅当 $R,S$ 都有单位元。
	\item 证明 $\{(r,r)\mid r\in R\}$ 是 $R\times R$ 的子环。
	\item 设 $I$ 为任意非空指标集，并且对每个 $i\in I$，$R_i$ 是环。证明直积 $\prod_{i\in I}R_i$ 在逐分量加法与乘法下构成环。
	\item 设 $R$ 是整数序列 $(a_1,a_2,a_3,\ldots)$ 的全体，其中 $a_i\in\ZZ$ 且除有限多个外均为 $0$（称为无穷多个 $\ZZ$ 的直和）。证明 $R$ 在逐分量加法与乘法下构成环，但没有单位元。
	\item 设 $X$ 是任意非空集合，$\mcP(X)$ 是 $X$ 的所有子集构成的幂集。在 $\mcP(X)$ 上定义
	\[
	A+B=(A-B)\cup(B-A),\qquad A\times B=A\cap B,
	\]
	即加法为对称差，乘法为交。
	\begin{enumerate}[label=(\alph*)]
		\item 证明 $\mcP(X)$ 在这些运算下构成环（$\mcP(X)$ 及其子环常称为集合环）。
		\item 证明这个环是交换的、有单位元，并且是 Boolean 环。
	\end{enumerate}
	\item 给出一个无穷 Boolean 环的例子。
	\item 设 $D$ 是无平方因子整数，$\mcO$ 是二次域 $\QQ(\sqrt D)$ 的整数环。对任意正整数 $f$，证明集合
	\[
	\mcO_f=\ZZ[f\omega]=\{a+bf\omega\mid a,b\in\ZZ\}
	\]
	是 $\mcO$ 的含单位元子环。证明 $[\mcO:\mcO_f]=f$（作为加法阿贝尔群的指数）。反过来，证明 $\mcO$ 的含单位元、且在 $\mcO$ 中作为加法阿贝尔群指数为 $f$ 的任意子环都等于 $\mcO_f$。（环 $\mcO_f$ 称为 $\QQ(\sqrt D)$ 中导子为 $f$ 的 order；整数环 $\mcO$ 称为 $\QQ(\sqrt D)$ 中的极大 order。）
	\item 对 $D=3,5,6,7$，通过在每个二次整数环 $\mcO$ 中展示一个具有无穷乘法阶的显式单位，说明单位群 $\mcO^\times$ 是无穷的。
	\item 设 $I$ 是整 Hamilton 四元数环，定义
	\[
	N:I\to\ZZ,
	\qquad N(a+bi+cj+dk)=a^2+b^2+c^2+d^2
	\]
	（映射 $N$ 称为范数）。
	\begin{enumerate}[label=(\alph*)]
		\item 证明 $N(\alpha)=\alpha\bar\alpha$ 对所有 $\alpha\in I$ 成立，其中若 $\alpha=a+bi+cj+dk$，则 $\bar\alpha=a-bi-cj-dk$。
		\item 证明 $N(\alpha\beta)=N(\alpha)N(\beta)$ 对所有 $\alpha,\beta\in I$ 成立。
		\item 证明 $I$ 中元素是单位当且仅当它的范数为 $+1$。证明 $I^\times$ 同构于 $8$ 阶四元数群。\emph{提示：}非零元素 $\alpha$ 在有理四元数环中的逆元为 $\bar\alpha/N(\alpha)$。
	\end{enumerate}
	\item 设 $K$ 是域。$K$ 上的一个离散赋值是函数 $v:K^\times\to\ZZ$，满足
	\begin{enumerate}[label=(\roman*)]
		\item $v(ab)=v(a)+v(b)$，即 $v$ 是从 $K$ 的非零元素乘法群到 $\ZZ$ 的同态；
		\item $v$ 是满射；
		\item 若 $x,y\in K^\times$ 且 $x+y\neq0$，则 $v(x+y)\geq \min\{v(x),v(y)\}$。
	\end{enumerate}
	集合
	\[
	R=\{x\in K^\times\mid v(x)\geq0\}\cup\{0\}
	\]
	称为 $v$ 的赋值环。
	\begin{enumerate}[label=(\alph*)]
		\item 证明 $R$ 是 $K$ 的含单位元子环。（一般地，若存在域 $K$ 和离散赋值 $v$，使 $R$ 是 $v$ 的赋值环，则称 $R$ 为离散赋值环。）
		\item 证明对每个非零 $x\in K$，$x$ 或 $x^{-1}$ 属于 $R$。
		\item 证明 $x$ 是 $R$ 的单位当且仅当 $v(x)=0$。
	\end{enumerate}
	\item 前一题的一个具体离散赋值环例子如下：$p$ 为素数，$K=\QQ$，并定义
	\[
	v_p:\QQ^\times\to\ZZ,
	\qquad v_p(a/b)=\alpha,
	\]
	其中 $a/b=p^\alpha c/d$，且 $p\nmid c$、$p\nmid d$。证明相应赋值环 $R$ 是所有分母与 $p$ 互素的有理数构成的环。描述这个赋值环的单位。
	\item 设 $R$ 是带 $1\neq0$ 的环。若非零元素 $a$ 满足存在非零 $x\in R$ 使 $ax=0$，则称 $a$ 是左零因子。对称地，若非零元素 $b$ 满足存在非零 $y\in R$ 使 $yb=0$，则称 $b$ 是右零因子（于是零因子是左或右零因子）。若存在 $s\in R$ 使 $su=1$，则称 $u$ 有左逆；若存在 $t\in R$ 使 $vt=1$，则称 $v$ 有右逆。
	\begin{enumerate}[label=(\alph*)]
		\item 证明 $u$ 是单位当且仅当它同时有左逆和右逆（即 $u$ 有双边逆）。
		\item 证明若 $u$ 有右逆，则 $u$ 不是右零因子。
		\item 证明若 $u$ 有不止一个右逆，则 $u$ 是左零因子。
		\item 证明若 $R$ 是有限环，则每个有右逆的元素都是单位（即有双边逆）。
	\end{enumerate}
	\item 设 $A$ 是任意带单位元 $1\neq0$ 的交换环。令 $R$ 为加法群 $A$ 到自身的所有群同态的集合，其中加法定义为逐点加法，乘法定义为函数复合。证明这些运算使 $R$ 成为带单位元的环。证明 $R$ 的单位就是 $A$ 的群自同构（参见第 1.6 节练习 20）。
	\item 设 $A=\ZZ\times\ZZ\times\ZZ\times\cdots$ 为按正整数编号的 $\ZZ$ 的直积（于是 $A$ 在逐分量加法与乘法下成环），并令 $R$ 为上一题中从 $A$ 到自身的所有群同态构成的环。设 $\varphi\in R$ 定义为
	\[
	\varphi(a_1,a_2,a_3,\ldots)=(a_2,a_3,\ldots),
	\]
	设 $\psi\in R$ 定义为
	\[
	\psi(a_1,a_2,a_3,\ldots)=(0,a_1,a_2,a_3,\ldots).
	\]
	\begin{enumerate}[label=(\alph*)]
		\item 证明 $\varphi\psi$ 是 $R$ 的单位元，但 $\psi\varphi$ 不是 $R$ 的单位元（即 $\psi$ 是 $\varphi$ 的右逆，但不是左逆）。
		\item 给出 $\varphi$ 的无穷多个右逆。
		\item 找出非零元素 $\pi\in R$，使 $\varphi\pi=0$ 但 $\pi\varphi\neq0$。
		\item 证明不存在非零元素 $\lambda\in R$ 使 $\lambda\varphi=0$（即 $\varphi$ 是左零因子但不是右零因子）。
	\end{enumerate}
\end{enumerate}



\section{例子：多项式环、矩阵环与群环}

本节介绍三类重要的环：多项式环、矩阵环与群环。随着本书的发展，我们将看到这三类环常常彼此相关。例如在第六部分中，我们会看到复数域上的群 $G$ 的群环 $\CC G$ 是若干个 $\CC$ 上矩阵环的直积。

这些环除了本身有趣以外，还有许多重要应用。在第三部分中，我们将用多项式环证明矩阵的若干分类定理；特别地，它们决定何时一个矩阵相似于对角矩阵。在第六部分中，我们将用群环研究群作用，并证明若干进一步的重要分类定理。

\subsection*{多项式环}
固定一个带单位元的交换环 $R$。我们以一种可能已经熟悉的方式定义多项式环，至少对实系数多项式来说如此。用 Cartesian 积给出的定义见附录 I。令 $x$ 为不定元。形式和
\[
  a_nx^n+a_{n-1}x^{n-1}+\cdots+a_1x+a_0
\]
其中 $n\geq0$ 且每个 $a_i\in R$，称为以 $R$ 中 $a_i$ 为系数、关于 $x$ 的多项式。若 $a_n\neq0$，则该多项式称为次数为 $n$，$a_nx^n$ 称为首项，$a_n$ 称为首系数（零多项式的首系数约定为 $0$）。若 $a_n=1$，多项式称为首一的。所有这类多项式的集合称为以 $R$ 为系数、变量为 $x$ 的多项式环，记为 $R[x]$。

使 $R[x]$ 成为环的加法与乘法，就是初等代数中熟悉的运算：加法逐项进行
\[
\begin{aligned}
 &(a_nx^n+a_{n-1}x^{n-1}+\cdots+a_1x+a_0)\\
 &\quad +(b_nx^n+b_{n-1}x^{n-1}+\cdots+b_1x+b_0)\\
 &= (a_n+b_n)x^n+(a_{n-1}+b_{n-1})x^{n-1}+\cdots+(a_1+b_1)x+(a_0+b_0),
\end{aligned}
\]
这里允许 $a_n$ 或 $b_n$ 为 $0$，以便定义不同次数多项式的和。乘法先对只有一个非零项的多项式定义
\[
  (ax^i)(bx^j)=abx^{i+j},
\]
再用分配律推广到所有多项式（通常称为“展开并合并同类项”）：
\[
\begin{aligned}
 &(a_0+a_1x+a_2x^2+\cdots)(b_0+b_1x+b_2x^2+\cdots)\\
 &=a_0b_0+(a_0b_1+a_1b_0)x+(a_0b_2+a_1b_1+a_2b_0)x^2+\cdots.
\end{aligned}
\]
一般地，积中 $x^k$ 的系数为 $\sum_{i=0}^k a_i b_{k-i}$。这些运算有意义，因为 $R$ 是环，系数的和与积都已定义。容易验证这些加法和乘法确实使 $R[x]$ 成为环。

环 $R$ 以常数多项式的形式出现在 $R[x]$ 中。注意由乘法定义，$R[x]$ 是带单位元的交换环（单位元是 $R$ 的单位元 $1$）。上面假设系数环 $R$ 交换，是因为这是我们主要关心的情形；但即使 $R$ 非交换或没有单位元，上述 $R[x]$ 中的加法与乘法定义仍然有效。若系数环为整数 $\ZZ$（分别为有理数 $\QQ$），则多项式环 $\ZZ[x]$（分别为 $\QQ[x]$）就是初等代数中熟悉的整系数（有理系数）多项式环。

另一个例子是系数在 $\ZZ/3\ZZ$ 中的多项式环 $(\ZZ/3\ZZ)[x]$。这个环由 $x$ 的非负幂构成，系数为 $0,1,2$，系数按模 $3$ 计算。例如若
\[
  p(x)=x^2+2x+1,
  \qquad q(x)=x^3+x+2,
\]
则
\[
  p(x)+q(x)=x^3+x^2,
\]
并且
\[
  p(x)q(x)=x^5+2x^4+2x^3+x^2+2x+2.
\]

系数所在的环会实质性地改变多项式的行为。例如，多项式 $x^2+1$ 在多项式环 $\ZZ[x]$ 中不是平方，但在多项式环 $(\ZZ/2\ZZ)[x]$ 中是平方，因为在这个环中
\[
  (x+1)^2=x^2+2x+1=x^2+1.
\]

\medskip
\noindent\textbf{命题 4.} 设 $R$ 是整环，$p(x),q(x)$ 是 $R[x]$ 的非零元素。则：
\begin{enumerate}[label=(\arabic*)]
  \item $\deg(p(x)q(x))=\deg p(x)+\deg q(x)$；
  \item $R[x]$ 的单位恰好是 $R$ 的单位；
  \item $R[x]$ 是整环。
\end{enumerate}

\begin{proof}
若 $R$ 没有零因子，则 $R[x]$ 也没有零因子：若 $p(x)$ 和 $q(x)$ 的首项分别为 $a_nx^n$ 和 $b_mx^m$，则 $p(x)q(x)$ 的首项为 $a_nb_mx^{n+m}$，而 $a_nb_m\neq0$。这同时证明了 (3) 并验证了 (1)。若 $p(x)$ 是单位，设 $p(x)q(x)=1$，则由 (1)
\[
  \deg p(x)+\deg q(x)=0,
\]
所以 $p(x)$ 与 $q(x)$ 都是 $R$ 的元素。由于它们的积为 $1$，$p(x)$ 是 $R$ 中单位。这证明 (2)。
\end{proof}

若环 $R$ 有零因子，则 $R[x]$ 也有零因子，因为 $R\subseteq R[x]$。此外，若 $f(x)$ 是 $R[x]$ 中的零因子，即存在非零 $g(x)\in R[x]$ 使 $f(x)g(x)=0$，则事实上存在非零 $c\in R$ 使 $cf(x)=0$（参见练习 2）。

若 $S$ 是 $R$ 的子环，则 $S[x]$ 是 $R[x]$ 的子环。例如 $\ZZ[x]$ 是 $\QQ[x]$ 的子环。$R[x]$ 的其他子环例子包括所有关于 $x^2$ 的多项式（即只出现 $x$ 的偶次幂的多项式）以及所有常数项为零的多项式（后一个子环没有单位元）。多项式环，特别是域上的多项式环，将在第 9 章中系统研究。

\subsection*{矩阵环}
固定任意环 $R$，令 $n$ 为正整数。用 $M_n(R)$ 表示所有以 $R$ 中元素为项的 $n\times n$ 矩阵的集合。元素 $(a_{ij})$ 是一个 $n\times n$ 方阵，其第 $i$ 行第 $j$ 列元素为 $a_{ij}\in R$。在通常的实矩阵加法与乘法规则下，矩阵集合成为环。加法逐项进行：$(a_{ij})+(b_{ij})$ 的第 $i,j$ 元为 $a_{ij}+b_{ij}$。矩阵积 $(a_{ij})(b_{ij})$ 的第 $i,j$ 元为 $\sum_{k=1}^n a_{ik}b_{kj}$（矩阵必须是方阵，才能使任意两个元素都可相乘）。直接计算可验证这些运算使 $M_n(R)$ 成为环。当 $R$ 是域时，第三部分中我们将用较少计算的方法证明 $M_n(R)$ 是环。

注意若 $R$ 是任意非平凡环（即使交换），且 $n\geq2$，则 $M_n(R)$ 不交换：若 $ab\neq0$，令 $A$ 为第 $1,1$ 位为 $a$ 其余为 $0$ 的矩阵，$B$ 为第 $1,2$ 位为 $b$ 其余为 $0$ 的矩阵，则 $AB$ 在第 $1,2$ 位有非零项 $ab$，而 $BA$ 是零矩阵。这两个矩阵也说明当 $R$ 非零且 $n\geq2$ 时，$M_n(R)$ 有零因子。

$M_n(R)$ 的元素 $(a_{ij})$ 称为\emph{标量矩阵}，若存在 $a\in R$，使 $a_{ii}=a$ 对所有 $i$ 成立，且当 $i\neq j$ 时 $a_{ij}=0$。标量矩阵全体是 $M_n(R)$ 的子环。这个子环是 $R$ 的一个拷贝：若矩阵 $A$ 主对角线上都是 $a$，矩阵 $B$ 主对角线上都是 $b$，则 $A+B$ 主对角线上都是 $a+b$，而 $AB$ 主对角线上都是 $ab$，其余项为 $0$。若 $R$ 交换，则标量矩阵与 $M_n(R)$ 的所有元素交换。若 $R$ 有 $1$，则主对角线上全为 $1$ 的标量矩阵（$n\times n$ 单位矩阵）是 $M_n(R)$ 的单位元。在这种情形，$M_n(R)$ 的单位就是可逆 $n\times n$ 矩阵，单位群记为 $\GL_n(R)$，称为 $R$ 上 $n$ 次一般线性群。

若 $S$ 是 $R$ 的子环，则 $M_n(S)$ 是 $M_n(R)$ 的子环。例如 $M_n(\ZZ)$ 是 $M_n(\QQ)$ 的子环，$M_n(2\ZZ)$ 又是二者的子环。$M_n(R)$ 的另一个子环例子是上三角矩阵全体
\[
  \{(a_{ij})\mid a_{pq}=0\text{ whenever }p>q\},
\]
即主对角线以下的所有项都为零的矩阵；容易验证上三角矩阵的和与积仍为上三角。

\subsection*{群环}
固定带单位元 $1\neq0$ 的交换环 $R$，令 $G=\{g_1,g_2,\ldots,g_n\}$ 是有限群，群运算写作乘法。定义 $G$ 在 $R$ 上的\emph{群环} $RG$ 为所有形式和
\[
  a_1g_1+a_2g_2+\cdots+a_ng_n,
  \qquad a_i\in R,
\]
组成的集合。若 $g_1$ 是 $G$ 的恒等元，则 $a_1g_1$ 简写为 $a_1$。类似地，对 $g\in G$，元素 $1g$ 简写为 $g$。

加法逐分量定义：
\[
\begin{aligned}
&(a_1g_1+a_2g_2+\cdots+a_ng_n)+(b_1g_1+b_2g_2+\cdots+b_ng_n)\\
&=(a_1+b_1)g_1+(a_2+b_2)g_2+\cdots+(a_n+b_n)g_n.
\end{aligned}
\]
乘法先定义
\[
  (ag_i)(bg_j)=(ab)g_k,
\]
其中 $ab$ 在 $R$ 中相乘，$g_ig_j=g_k$ 是群 $G$ 中的乘积。随后用分配律推广到所有形式和；在乘积
\[
(a_1g_1+\cdots+a_ng_n)(b_1g_1+\cdots+b_ng_n)
\]
中，$g_k$ 的系数为 $\sum_{g_ig_j=g_k}a_ib_j$。直接验证可知这些运算使 $RG$ 成为环（事实上这里不需要 $R$ 交换）。乘法结合律来自 $G$ 中群运算的结合律。$RG$ 交换当且仅当 $G$ 是交换群。

\medskip
\noindent\textbf{例.}
设 $G=D_8$ 是 $8$ 阶二面体群，取通常生成元 $r,s$，满足 $r^4=s^2=1$ 且 $rs=sr^{-1}$，令 $R=\ZZ$。元素
\[
  \alpha=r+r^2-2s,
  \qquad \beta=-3r^2+rs
\]
是 $\ZZ D_8$ 的典型成员。它们的和与积为
\[
  \alpha+\beta=r-2r^2-2s+rs,
\]
并且
\[
\begin{aligned}
\alpha\beta&=(r+r^2-2s)(-3r^2+rs)\\
&=r(-3r^2+rs)+r^2(-3r^2+rs)-2s(-3r^2+rs)\\
&=-3r^3+r^2s-3+r^3s+6r^2s-2r^3\\
&=-3-5r^3+7r^2s+r^3s.
\end{aligned}
\]

环 $R$ 以“常数”形式和的方式出现在 $RG$ 中，即 $G$ 的恒等元的 $R$-倍（注意把 $RG$ 中的加法和乘法限制到这些元素时，正是 $R$ 中的加法和乘法）。这些 $R$ 的元素与 $RG$ 的所有元素交换。$R$ 的单位元也是 $RG$ 的单位元。

群 $G$ 也出现在 $RG$ 中（元素 $g_i$ 作为 $1g_i$ 出现；例如上面的 $r,s\in D_8$ 也属于群环 $\ZZ D_8$）。把环 $RG$ 的乘法限制到 $G$，恰好是 $G$ 的群运算。特别地，$G$ 的每个元素在环 $RG$ 中都有乘法逆元（即它在 $G$ 中的逆元）。这说明 $G$ 是 $RG$ 的单位群的子群。

若 $|G|>1$，则 $RG$ 总有零因子。例如，取 $G$ 中任意阶为 $m>1$ 的元素 $g$，则
\[
  (1-g)(1+g+\cdots+g^{m-1})=1-g^m=0.
\]
若第二个因子非零（例如当 $R$ 没有 $1$ 的非平凡周期时），这便给出零因子；在一般群环中还会出现许多其他零因子。群环是把群论与环论联系起来的基本对象，在第 18 章中将系统使用。

\subsection*{练习}
在以下各题中，设 $R$ 是带单位元的交换环。
\begin{enumerate}[label=\arabic*.]
	\item 令 $p(x)=2x^3-3x^2+4x-5$，$q(x)=7x^3+3-3x^4$。在 (a)、(b)、(c) 中，分别假定两个给定多项式的系数取自指定环（(b)、(c) 中整数系数按模 $n$ 化简），计算 $p(x)+q(x)$ 与 $p(x)q(x)$：
	\begin{enumerate}[label=(\alph*)]
		\item $R=\ZZ$；
		\item $R=\ZZ/2\ZZ$；
		\item $R=\ZZ/3\ZZ$。
	\end{enumerate}
	\item 设 $p(x)=a_nx^n+a_{n-1}x^{n-1}+\cdots+a_1x+a_0\in R[x]$。证明 $p(x)$ 是 $R[x]$ 中的零因子，当且仅当存在非零 $b\in R$ 使 $bp(x)=0$。\emph{提示：}取次数最小的非零多项式 $g(x)=b_mx^m+b_{m-1}x^{m-1}+\cdots+b_0$ 满足 $g(x)p(x)=0$。证明 $b_ma_n=0$，从而 $a_ng(x)$ 是次数小于 $m$ 且与 $p(x)$ 相乘仍为 $0$ 的多项式；推出 $a_ng(x)=0$。类似地对 $i$ 归纳证明 $a_{n-i}g(x)=0$，$i=0,1,\ldots,n$，并推出 $b_mp(x)=0$。
	\item 定义不定元 $x$ 上、系数在 $R$ 中的形式幂级数集合 $R[[x]]$ 为所有形式无穷和
	\[
	\sum_{n=0}^{\infty}a_nx^n=a_0+a_1x+a_2x^2+a_3x^3+\cdots .
	\]
	加法和乘法按实或复幂级数的通常规则定义，也就是把多项式的加法和乘法延伸到“无穷次数多项式”：
	\[
	\sum_{n=0}^\infty a_nx^n+\sum_{n=0}^\infty b_nx^n=\sum_{n=0}^\infty(a_n+b_n)x^n,
	\]
	\[
	\left(\sum_{n=0}^\infty a_nx^n\right)\left(\sum_{n=0}^\infty b_nx^n\right)
	=\sum_{n=0}^\infty\left(\sum_{k=0}^n a_kb_{n-k}\right)x^n.
	\]
	“形式”表示不考虑收敛性，因此形式幂级数不必表示 $\RR$ 上的函数。
	\begin{enumerate}[label=(\alph*)]
		\item 证明 $R[[x]]$ 是带 $1$ 的交换环。
		\item 证明 $1-x$ 是 $R[[x]]$ 中的单位，其逆元为 $1+x+x^2+\cdots$。
		\item 证明 $\sum_{n=0}^{\infty}a_nx^n$ 是 $R[[x]]$ 中的单位，当且仅当 $a_0$ 是 $R$ 中的单位。
	\end{enumerate}
	\item 证明若 $R$ 是整环，则形式幂级数环 $R[[x]]$ 也是整环。
	\item 设 $F$ 是域，定义系数在 $F$ 中的形式 Laurent 级数环 $F((x))$ 为
	\[
	F((x))=\left\{\sum_{n\geq N}a_nx^n\mid a_n\in F,\,N\in\ZZ\right\}.
	\]
	（$F((x))$ 的每个元素都是 $x$ 的幂级数加上 $1/x$ 的多项式，即负次幂项只有有限多个。）
	\begin{enumerate}[label=(\alph*)]
		\item 证明 $F((x))$ 是域。
		\item 定义映射
		\[
		v:F((x))^\times\to\ZZ,
		\qquad v\!\bigg(\sum_{n\geq N}a_nx^n\bigg)=N,
		\]
		其中 $a_N$ 是级数的第一个非零系数（即 $N$ 是该级数在 $0$ 处零点或极点的阶）。证明 $v$ 是 $F((x))$ 上的离散赋值，其离散赋值环是 $F[[x]]$（参见第 7.1 节练习 26）。
	\end{enumerate}
	\item 设 $S$ 是带单位元 $1\neq0$ 的环。令 $n\in\ZZ^+$，并令 $A$ 是一个 $n\times n$ 矩阵，其元素来自 $S$，其 $i,j$ 元记为 $a_{ij}$。设 $E_{ij}$ 是 $M_n(S)$ 中第 $i,j$ 元为 $1$，其余元素全为 $0$ 的矩阵。
	\begin{enumerate}[label=(\alph*)]
		\item 证明 $E_{ij}A$ 是这样一个矩阵：它的第 $i$ 行等于 $A$ 的第 $j$ 行，其余所有行均为零。
		\item 证明 $AE_{ij}$ 是这样一个矩阵：它的第 $j$ 列等于 $A$ 的第 $i$ 列，其余所有列均为零。
		\item 推出 $E_{pq}AE_{rs}$ 是第 $p,s$ 元为 $a_{qr}$、其余元素全为零的矩阵。
	\end{enumerate}
	\item 证明环 $M_n(R)$ 的中心是标量矩阵全体（参见第 7.1 节练习 7）。\emph{提示：}使用上一题。
	\item 设 $S$ 为任意环，$n\geq2$。证明若 $A$ 是 $M_n(S)$ 中任意严格上三角矩阵，则 $A^n=0$（严格上三角矩阵是指主对角线及其下方元素全为零的矩阵）。
	\item 令 $\alpha=r+r^2-2s$，$\beta=-3r^2+rs$ 为本节中描述的整群环 $\ZZ D_8$ 的两个元素。计算 $\ZZ D_8$ 中的下列元素：
	\begin{enumerate}[label=(\alph*)]
		\item $\beta\alpha$；
		\item $\alpha^2$；
		\item $\alpha\beta-\beta\alpha$；
		\item $\beta\alpha\beta$。
	\end{enumerate}
	\item 考虑整群环 $\ZZ S_3$ 中的元素
	\[
	\alpha=3(12)-5(23)+14(123),
	\qquad
	\beta=6(1)+2(23)-7(132),
	\]
	其中 $(1)$ 是 $S_3$ 的恒等置换。计算：
	\begin{enumerate}[label=(\alph*)]
		\item $\alpha+\beta$；
		\item $2\alpha-3\beta$；
		\item $\alpha\beta$；
		\item $\beta\alpha$；
		\item $\alpha^2$。
	\end{enumerate}
	\item 在假定 $\alpha$ 与 $\beta$ 的系数属于 $\ZZ/3\ZZ$ 的情况下（即 $\alpha,\beta\in(\ZZ/3\ZZ)S_3$），重复上一题。
	\item 设 $G=\{g_1,\ldots,g_n\}$ 是有限群。证明元素
	\[
	N=g_1+g_2+\cdots+g_n
	\]
	属于群环 $RG$ 的中心（参见第 7.1 节练习 7）。
	\item 设 $K=\{k_1,\ldots,k_m\}$ 是有限群 $G$ 的一个共轭类。
	\begin{enumerate}[label=(\alph*)]
		\item 证明元素 $\widehat K=k_1+\cdots+k_m$ 属于群环 $RG$ 的中心（参见第 7.1 节练习 7）。\emph{提示：}验证对所有 $g\in G$ 有 $g^{-1}\widehat K g=\widehat K$。
		\item 设 $K_1,
		\ldots,K_r$ 是 $G$ 的共轭类，并对每个 $K_i$ 令 $\widehat K_i$ 为 $RG$ 中 $K_i$ 元素之和。证明元素 $\alpha\in RG$ 属于 $RG$ 的中心，当且仅当
		\[
		\alpha=a_1\widehat K_1+a_2\widehat K_2+\cdots+a_r\widehat K_r
		\]
		其中 $a_1,a_2,\ldots,a_r\in R$。
	\end{enumerate}
\end{enumerate}



\section{环同态与商环}

环同态是从一个环到另一个环、同时保持加法结构和乘法结构的映射。

\medskip
\noindent\textbf{定义.} 设 $R,S$ 是环。
\begin{enumerate}[label=(\arabic*)]
  \item \emph{环同态}是映射 $\varphi:R\to S$，满足
  \begin{enumerate}[label=(\roman*)]
    \item $\varphi(a+b)=\varphi(a)+\varphi(b)$ 对所有 $a,b\in R$ 成立（即 $\varphi$ 是加法群同态）；
    \item $\varphi(ab)=\varphi(a)\varphi(b)$ 对所有 $a,b\in R$ 成立。
  \end{enumerate}
  \item 环同态 $\varphi$ 的\emph{核}记为 $\ker\varphi$，是 $R$ 中映到 $S$ 中 $0$ 的元素全体（即把 $\varphi$ 看作加法群同态时的核）。
  \item 双射环同态称为\emph{同构}。
\end{enumerate}

若语境清楚，我们常把“环同态”简称为“同态”。类似地，若 $A$ 与 $B$ 是环，$A\cong B$ 总表示环同构，除非另有说明。

\medskip
\noindent\textbf{例.}
\begin{enumerate}[label=(\arabic*)]
  \item 映射 $\varphi:\ZZ\to\ZZ/2\ZZ$，把偶数送到 $0$，把奇数送到 $1$，是环同态。它保持加法，因为两个偶数或两个奇数之和为偶数，一个偶数与一个奇数之和为奇数；它保持乘法，因为两个奇数之积为奇数，偶数与任意整数之积为偶数。$\ker\varphi$ 是偶整数集合；$1\in\ZZ/2\ZZ$ 上方的纤维是奇整数集合。
  \item 对 $n\in\ZZ$，映射 $\varphi_n:\ZZ\to\ZZ$，$\varphi_n(x)=nx$ 一般不是环同态，因为 $\varphi_n(xy)=nxy$，而 $\varphi_n(x)\varphi_n(y)=n^2xy$。因此 $\varphi_n$ 是环同态只有在 $n^2=n$，即 $n=0,1$ 时成立。注意 $\varphi_n$ 始终是加法群同态，所以处理环时必须小心检查两个环运算都被保持。$\varphi_0$ 是零同态，$\varphi_1$ 是恒等同态。
  \item 令 $\varphi:\QQ[x]\to\QQ$ 为把有理系数多项式映到其常数项的映射，即 $\varphi(p(x))=p(0)$。这是环同态，因为两个多项式和的常数项是常数项之和，积的常数项是常数项之积。$a\in\QQ$ 上方的纤维由常数项为 $a$ 的所有多项式构成。特别地，核由常数项为 $0$ 的多项式构成。
\end{enumerate}

\medskip
\noindent\textbf{命题 5.} 设 $R,S$ 是环，$\varphi:R\to S$ 是同态。
\begin{enumerate}[label=(\arabic*)]
  \item $\varphi$ 的像是 $S$ 的子环。
  \item $\varphi$ 的核是 $R$ 的子环。进一步，若 $a\in\ker\varphi$，则对每个 $r\in R$ 有 $ra,ar\in\ker\varphi$；也就是说，$\ker\varphi$ 对来自 $R$ 的元素的左右乘法封闭。
\end{enumerate}

\begin{proof}
(1) 若 $s_1,s_2\in\im\varphi$，则 $s_1=\varphi(r_1)$、$s_2=\varphi(r_2)$，其中 $r_1,r_2\in R$。于是
\[
  \varphi(r_1-r_2)=s_1-s_2,
  \qquad \varphi(r_1r_2)=s_1s_2.
\]
故 $s_1-s_2$ 与 $s_1s_2$ 仍在像中，所以像对减法和乘法封闭，是 $S$ 的子环。

(2) 若 $\alpha,\beta\in\ker\varphi$，则 $\varphi(\alpha)=\varphi(\beta)=0$。于是 $\varphi(\alpha-\beta)=0$ 且 $\varphi(\alpha\beta)=0$，所以核对减法和乘法封闭，是 $R$ 的子环。类似地，对任意 $r\in R$，
\[
  \varphi(r\alpha)=\varphi(r)\varphi(\alpha)=\varphi(r)0=0,
  \qquad
  \varphi(\alpha r)=\varphi(\alpha)\varphi(r)=0\varphi(r)=0,
\]
故 $r\alpha,\alpha r\in\ker\varphi$。
\end{proof}

群同态的情形中，同态的纤维自然具有一个群结构，并同构于像；这导致按正规子群取商群的概念。环同态也有类似结果。设 $\varphi:R\to S$ 是环同态，核为 $I$。由于 $R$ 和 $S$ 特别是加法 Abel 群，$\varphi$ 特别是 Abel 群同态；$\varphi$ 的纤维是核 $I$ 的加法陪集 $r+I$。这些纤维具有自然同构于 $\varphi$ 的像的环结构：若 $X$ 是 $a\in S$ 上方纤维，$Y$ 是 $b\in S$ 上方纤维，则 $X+Y$ 是 $a+b$ 上方纤维，$XY$ 是 $ab$ 上方纤维。用核的陪集表示，就是
\[
  (r+I)+(s+I)=(r+s)+I,
  \tag{7.1}
\]
\[
  (r+I)(s+I)=rs+I.
  \tag{7.2}
\]
当 $I$ 是某个同态的核时，加法商群 $R/I$ 也因此具有乘法结构，称为 $R$ 对 $I$ 的\emph{商环}。

现在考虑任意加法子群 $I$ 是否能用 (7.1)、(7.2) 定义商环。由于 $R$ 的加法群 Abel，$I$ 自动正规，所以 $R/I$ 自动是加法 Abel 群。问题在于乘法是否与代表元选择无关。若把代表元 $r,s$ 换成 $r+\alpha,s+\beta$（其中 $\alpha,\beta\in I$），必须有
\[
  (r+\alpha)(s+\beta)+I=rs+I
\]
对所有 $r,s\in R$、$\alpha,\beta\in I$ 成立。取 $r=s=0$ 得 $I$ 必须对乘法封闭；取 $s=0$ 得 $r\beta\in I$，取 $r=0$ 得 $\alpha s\in I$。反过来，如果 $I$ 对 $R$ 中元素的左右乘法封闭，则上式成立。所以这正是必要且充分条件。

\medskip
\noindent\textbf{定义.} 设 $R$ 是环，$I\subseteq R$，$r\in R$。
\begin{enumerate}[label=(\arabic*)]
  \item $rI=\{ra\mid a\in I\}$，$Ir=\{ar\mid a\in I\}$。
  \item 若 $I$ 是 $R$ 的子环，并且 $rI\subseteq I$ 对所有 $r\in R$ 成立，则称 $I$ 是 $R$ 的\emph{左理想}。若改为 $Ir\subseteq I$ 对所有 $r\in R$ 成立，则称 $I$ 为\emph{右理想}。
  \item 同时是左理想和右理想的子集称为\emph{理想}（或强调为\emph{双边理想}）。
\end{enumerate}

在交换环中，左、右和双边理想的概念相同。要证明一个子集 $I$ 是理想，必须证明它非空、对减法封闭，并且对乘以 $R$ 中任意元素封闭（不仅仅是对乘以 $I$ 中元素封闭）。若 $R$ 有 $1$，则 $(-1)a=-a$，于是此时也可检查非空、对加法封闭、对 $R$ 中元素乘法封闭。

\medskip
\noindent\textbf{命题 6.} 设 $R$ 是环，$I$ 是 $R$ 的理想。则加法商群 $R/I$ 在运算
\[
  (r+I)+(s+I)=(r+s)+I,
  \qquad
  (r+I)(s+I)=rs+I
\]
下成为环，称为 $R$ 对 $I$ 的商环。反之，若这些运算对加法子群 $I$ 良定义并使 $R/I$ 成为环，则 $I$ 是 $R$ 的理想。

\begin{proof}
良定义性的必要性和充分性已由上面的代表元讨论证明。若乘法良定义，则商环的每条公理都直接来自 $R$ 中相应公理。例如一个分配律为
\[
\begin{aligned}
 (r+I)[(s+I)+(t+I)]&=(r+I)((s+t)+I)\\
 &=r(s+t)+I=(rs+rt)+I\\
 &=(rs+I)+(rt+I)\\
 &=[(r+I)(s+I)]+[(r+I)(t+I)].
\end{aligned}
\]
其余公理类似。
\end{proof}

\medskip
\noindent\textbf{定理 7.}
\begin{enumerate}[label=(\arabic*)]
  \item \textbf{（环的第一同构定理）} 若 $\varphi:R\to S$ 是环同态，则
  \[
    R/\ker\varphi\cong\varphi(R).
  \]
  \item 若 $I$ 是 $R$ 的理想，则自然投影
  \[
    \pi:R\to R/I,
    \qquad \pi(r)=r+I
  \]
  是满环同态，核为 $I$。因此理想恰好是环同态的核。
\end{enumerate}

\begin{proof}
把第一同构定理用于加法群，得到 $R/\ker\varphi$ 与 $\varphi(R)$ 的加法群同构；该同构把 $r+\ker\varphi$ 送到 $\varphi(r)$。乘法也被保持，因为
\[
  (r+
  \ker\varphi)(s+\ker\varphi)=rs+
  \ker\varphi
\]
映到 $\varphi(rs)=\varphi(r)\varphi(s)$。第二部分由商环乘法定义立即得到；其核正是满足 $r+I=I$ 的元素，即 $r\in I$。
\end{proof}

\medskip
\noindent\textbf{例.}
\begin{enumerate}[label=(\arabic*)]
  \item 在任意环 $R$ 中，$R$ 与 $0$ 都是理想；相应商环为零环和 $R$ 本身。
  \item 在 $\ZZ$ 中，$n\ZZ$ 是理想，商环为 $\ZZ/n\ZZ$。经典的“模 $n$ 化简”就是自然投影 $\ZZ\to\ZZ/n\ZZ$。
  \item 在 $\ZZ[x]$ 中，所有常数项为零的多项式构成理想；更一般地，所有能被固定多项式 $f(x)$ 整除的多项式构成主理想 $(f(x))$。
  \item 若 $R$ 是从集合 $X$ 到环 $A$ 的函数环，固定 $c\in X$，求值映射 $E_c(f)=f(c)$ 是环同态；其核是所有在 $c$ 处为零的函数构成的理想。
  \item 若 $J$ 是 $R$ 的理想，则 $M_n(J)$ 是 $M_n(R)$ 的理想，且
  \[
    M_n(R)/M_n(J)\cong M_n(R/J).
  \]
  \item 群环 $RG$ 有\emph{增广映射}
  \[
    \epsilon\left(\sum a_i g_i\right)=\sum a_i.
  \]
  它是到 $R$ 的环同态；核称为增广理想。
\end{enumerate}

\medskip
\noindent\textbf{例（化简同态）.}
从 $\ZZ$ 到 $\ZZ/n\ZZ$ 的自然投影通常称为“模 $n$ 化简”。这是环同态这一事实在初等数论中有重要后果。例如要解整数方程
\[
  x^2+y^2+z^2=8w+7
\]
或类似 Diophantine 方程时，若整数解存在，那么同一个等式在任意商环 $\ZZ/n\ZZ$ 中也必须成立。选取合适的 $n$ 常能给出矛盾，或至少限制可能的剩余类。比如模 $4$ 时平方只可能为 $0$ 或 $1$，因此许多方程可通过检查有限个模 $4$ 可能性排除。即使有解，此方法也能给出解模 $n$ 的限制；结合第 7.6 节的中国剩余定理，可得到模很大整数的限制，有助于数值搜索。类似地，从 $\ZZ[x]$ 到 $(\ZZ/p\ZZ)[x]$ 的“系数模 $p$ 化简”也是环同态，在第 9 章中将用于判断多项式能否分解。

下面给出环的其余同构定理。每一个都可先用群论中的相应定理得到加法群同构（或第四同构定理中的对应），然后检查该同构保持乘法；检查直接来自商环乘法定义。例如 (2) 中的同构由 $r+I\mapsto r+J$ 给出，显然保持乘法。

\medskip
\noindent\textbf{定理 8.} 设 $R$ 是环。
\begin{enumerate}[label=(\arabic*)]
  \item \textbf{（环的第二同构定理）} 若 $A$ 是子环，$B$ 是 $R$ 的理想，则 $A+B=\{a+b\mid a\in A,b\in B\}$ 是 $R$ 的子环，$A\cap B$ 是 $A$ 的理想，并且
  \[
    (A+B)/B\cong A/(A\cap B).
  \]
  \item \textbf{（环的第三同构定理）} 若 $I\subseteq J$ 是 $R$ 的理想，则 $J/I$ 是 $R/I$ 的理想，并且
  \[
    (R/I)/(J/I)\cong R/J.
  \]
  \item \textbf{（环的第四或格同构定理）} 若 $I$ 是 $R$ 的理想，则对应 $A\mapsto A/I$ 给出包含 $I$ 的 $R$ 的子环与 $R/I$ 的子环之间保持包含关系的双射。并且，包含 $I$ 的子环 $A$ 是 $R$ 的理想，当且仅当 $A/I$ 是 $R/I$ 的理想。
\end{enumerate}

例如取 $R=\ZZ$，$I=12\ZZ$。商环 $\bar R=\ZZ/12\ZZ$ 的理想
\[
  \bar R,
  \quad 2\ZZ/12\ZZ,
  \quad 3\ZZ/12\ZZ,
  \quad 4\ZZ/12\ZZ,
  \quad 6\ZZ/12\ZZ,
  \quad 0
\]
对应于 $\ZZ$ 中包含 $12\ZZ$ 的理想
\[
  \ZZ,
  \quad 2\ZZ,
  \quad 3\ZZ,
  \quad 4\ZZ,
  \quad 6\ZZ,
  \quad 12\ZZ.
\]

若 $I,J$ 是环 $R$ 的理想，则所有 $a+b$（$a\in I,b\in J$）组成的集合不仅是 $R$ 的子环（如第二同构定理），而且是 $R$ 的理想。我们还可以定义两个理想的积。

\medskip
\noindent\textbf{定义.} 设 $I,J$ 是 $R$ 的理想。
\begin{enumerate}[label=(\arabic*)]
  \item 定义 $I$ 与 $J$ 的和为 $I+J=\{a+b\mid a\in I,b\in J\}$。
  \item 定义 $I$ 与 $J$ 的积 $IJ$ 为所有形如 $ab$（$a\in I,b\in J$）的元素的有限和组成的集合。
  \item 对 $n\geq1$，$I$ 的第 $n$ 次幂 $I^n$ 定义为所有形如 $a_1a_2\cdots a_n$（每个 $a_i\in I$）的元素的有限和组成的集合；等价地，$I^1=I$，$I^n=I I^{n-1}$。
\end{enumerate}

容易看出 $I+J$ 是包含 $I$ 与 $J$ 的最小理想，$IJ$ 是包含在 $I\cap J$ 中的理想（但可以严格小于 $I\cap J$，见练习）。注意 $IJ$ 的元素是来自 $I$ 和 $J$ 的乘积的有限和；单纯由乘积 $ab$ 构成的集合一般不对加法封闭，因而一般不是理想。

\medskip
\noindent\textbf{例.}
\begin{enumerate}[label=(\arabic*)]
  \item 在 $\ZZ$ 中令 $I=6\ZZ$，$J=10\ZZ$。则 $I+J$ 由所有 $6x+10y$ 组成，等于 $2\ZZ$；一般地，$m\ZZ+n\ZZ=d\ZZ$，其中 $d=(m,n)$。积 $IJ$ 由所有有限个 $(6x)(10y)$ 的和组成，即 $60\ZZ$。
  \item 令 $I$ 是 $\ZZ[x]$ 中常数项为偶数的多项式理想。多项式 $2$ 与 $x$ 属于 $I$，故 $4=2\cdot2$ 与 $x^2=x\cdot x$ 都在 $I^2$ 中，和 $x^2+4$ 也在 $I^2$ 中。但容易检查，$x^2+4$ 不能写成 $I$ 中两个元素的单个乘积 $p(x)q(x)$。
\end{enumerate}

\subsection*{练习}
在以下各题中，设 $R$ 是带单位元 $1\neq0$ 的环。
\begin{enumerate}[label=\arabic*.]
	\item 证明环 $2\ZZ$ 与 $3\ZZ$ 不同构。
	\item 证明环 $\ZZ[x]$ 与 $\QQ[x]$ 不同构。
	\item 找出 $\ZZ$ 的所有同态像。
	\item 找出从 $\ZZ$ 到 $\ZZ/30\ZZ$ 的所有环同态。分别描述其核与像。
	\item 描述从环 $\ZZ\times\ZZ$ 到 $\ZZ$ 的所有环同态。分别描述其核与像。
	\item 判定下列从 $M_2(\ZZ)$ 到 $\ZZ$ 的映射中哪些是环同态：
	\begin{enumerate}[label=(\alph*)]
		\item $\begin{pmatrix}a&b\\ c&d\end{pmatrix}\mapsto a$（投影到 $(1,1)$ 元）；
		\item $\begin{pmatrix}a&b\\ c&d\end{pmatrix}\mapsto a+d$（矩阵的迹）；
		\item $\begin{pmatrix}a&b\\ c&d\end{pmatrix}\mapsto ad-bc$（矩阵的行列式）。
	\end{enumerate}
	\item 设
	\[
	R=\left\{\begin{pmatrix}a&b\\0&d\end{pmatrix}\mid a,b,d\in\ZZ\right\}
	\]
	是 $M_2(\ZZ)$ 中上三角矩阵构成的子环。证明映射
	\[
	\varphi:R\to\ZZ\times\ZZ,
	\qquad
	\begin{pmatrix}a&b\\0&d\end{pmatrix}\mapsto(a,d)
	\]
	是满同态，并描述其核。
	\item 判定下列哪些是环 $\ZZ\times\ZZ$ 的理想：
	\begin{enumerate}[label=(\alph*)]
		\item $\{(a,a)\mid a\in\ZZ\}$；
		\item $\{(2a,2b)\mid a,b\in\ZZ\}$；
		\item $\{(2a,0)\mid a\in\ZZ\}$；
		\item $\{(a,-a)\mid a\in\ZZ\}$。
	\end{enumerate}
	\item 判定第 7.1 节练习 6 中的哪些集合是从 $[0,1]$ 到 $\RR$ 的所有函数构成的环的理想。
	\item 判定下列哪些是环 $\ZZ[x]$ 的理想：
	\begin{enumerate}[label=(\alph*)]
		\item 常数项是 $3$ 的倍数的所有多项式；
		\item $x^2$ 的系数是 $3$ 的倍数的所有多项式；
		\item 常数项、$x$ 的系数和 $x^2$ 的系数都为 $0$ 的所有多项式；
		\item $\ZZ[x^2]$（即只出现 $x$ 的偶次幂的多项式）；
		\item 系数和为 $0$ 的所有多项式；
		\item 满足 $p'(0)=0$ 的所有多项式 $p(x)$，其中 $p'(x)$ 是 $p(x)$ 对 $x$ 的通常一阶导数。
	\end{enumerate}
	\item 设 $R$ 是闭区间 $[0,1]$ 上所有实值连续函数构成的环。证明映射
	\[
	\varphi:R\to\RR,
	\qquad \varphi(f)=\int_0^1 f(t)\,dt
	\]
	是加法群同态，但不是环同态。
	\item 设 $D$ 是不是 $\ZZ$ 中完全平方数的整数，并令
	\[
	S=\left\{\begin{pmatrix}a&b\\ Db&a\end{pmatrix}\mid a,b\in\ZZ\right\}.
	\]
	\begin{enumerate}[label=(\alph*)]
		\item 证明 $S$ 是 $M_2(\ZZ)$ 的子环。
		\item 若 $D$ 不是 $\ZZ$ 中的完全平方数，证明映射 $\varphi:\ZZ[\sqrt D]\to S$，
		\[
		\varphi(a+b\sqrt D)=\begin{pmatrix}a&b\\Db&a\end{pmatrix}
		\]
		是环同构。
		\item 若 $D\equiv1\pmod4$ 且 $D$ 无平方因子，证明集合
		\[
		\left\{\begin{pmatrix}a&b\\ \frac{D-1}{4}b&a+b\end{pmatrix}\mid a,b\in\ZZ\right\}
		\]
		是 $M_2(\ZZ)$ 的子环，并且同构于二次整数环 $\mcO$。
	\end{enumerate}
	\item 证明环 $M_2(\RR)$ 含有一个同构于 $\CC$ 的子环。
	\item 证明环 $M_4(\RR)$ 含有一个同构于实 Hamilton 四元数环 $\HH$ 的子环。
	\item 设 $X$ 是非空集合，$\mcP(X)$ 是第 7.1 节练习 21 中定义的所有子集构成的 Boolean 环。设 $R$ 是从 $X$ 到 $\ZZ/2\ZZ$ 的所有函数构成的环。对每个 $A\in\mcP(X)$，定义函数
	\[
	\chi_A:X\to\ZZ/2\ZZ,
	\qquad
	\chi_A(x)=\begin{cases}1,&x\in A,\\0,&x\notin A.
	\end{cases}
	\]
	（$\chi_A$ 称为 $A$ 的特征函数，其值在 $\ZZ/2\ZZ$ 中。）证明映射 $\mcP(X)\to R$，$A\mapsto\chi_A$ 是环同构。
	\item 设 $\varphi:R\to S$ 是满环同态。证明 $R$ 的中心的像包含在 $S$ 的中心中（参见第 7.1 节练习 7）。
	\item 设 $R,S$ 是带单位元的非零环，分别记其单位元为 $1_R$ 和 $1_S$。设 $\varphi:R\to S$ 是非零环同态。
	\begin{enumerate}[label=(\alph*)]
		\item 证明若 $\varphi(1_R)\neq1_S$，则 $\varphi(1_R)$ 是 $S$ 中的零因子。推出若 $S$ 是整环，则每个从 $R$ 到 $S$ 的环同态都把 $R$ 的单位元送到 $S$ 的单位元。
		\item 证明若 $\varphi(1_R)=1_S$，则对 $R$ 的每个单位 $u$，$\varphi(u)$ 是 $S$ 的单位，且 $\varphi(u^{-1})=\varphi(u)^{-1}$。
	\end{enumerate}
	\item \begin{enumerate}[label=(\alph*)]
		\item 若 $I,J$ 是 $R$ 的理想，证明 $I\cap J$ 也是 $R$ 的理想。
		\item 证明任意非空理想族的交仍是理想（不要假定该族可数）。
	\end{enumerate}
	\item 证明若 $I_1\subseteq I_2\subseteq\cdots$ 是 $R$ 的理想，则 $\bigcup_{n=1}^{\infty}I_n$ 是 $R$ 的理想。
	\item 设 $I$ 是 $R$ 的理想，$S$ 是 $R$ 的子环。证明 $I\cap S$ 是 $S$ 的理想。举例说明：子环 $S$ 的理想不一定都能写成某个 $R$ 的理想 $I$ 与 $S$ 的交。
	\item 证明 $M_n(R)$ 的每个双边理想都等于 $M_n(J)$，其中 $J$ 是 $R$ 的某个双边理想。\emph{提示：}先用第 7.2 节练习 6(c) 证明 $M_n(R)$ 的一个理想中所有矩阵元素的集合构成 $R$ 的理想。
	\item 设 $a\in R$。
	\begin{enumerate}[label=(\alph*)]
		\item 证明 $\{x\in R\mid ax=0\}$ 是右理想，$\{y\in R\mid ya=0\}$ 是左理想（分别称为 $a$ 的右、左零化子）。
		\item 证明若 $L$ 是 $R$ 的左理想，则 $\{x\in R\mid xa=0\text{ 对所有 }a\in L\}$ 是双边理想（称为 $L$ 的左零化子）。
	\end{enumerate}
	\item 设 $S$ 是 $R$ 的子环，$I$ 是 $R$ 的理想。证明若 $S\cap I=0$，则 $S\cong\bar S$，其中横线表示模 $I$ 到 $R/I$ 中。
	\item 设 $\varphi:R\to S$ 是环同态。
	\begin{enumerate}[label=(\alph*)]
		\item 证明若 $J$ 是 $S$ 的理想，则 $\varphi^{-1}(J)$ 是 $R$ 的理想。把它应用于 $R$ 是 $S$ 的子环且 $\varphi$ 是包含同态的情形，推出若 $J$ 是 $S$ 的理想，则 $J\cap R$ 是 $R$ 的理想。
		\item 证明若 $\varphi$ 满且 $I$ 是 $R$ 的理想，则 $\varphi(I)$ 是 $S$ 的理想。给出一个例子说明若 $\varphi$ 不满，这个结论可能失败。
	\end{enumerate}
	\item 假设 $R$ 是带 $1$ 的交换环。证明二项式定理
	\[
	(a+b)^n=\sum_{k=0}^n {n\choose k}a^kb^{n-k}
	\]
	在 $R$ 中成立，其中二项式系数 ${n\choose k}$ 在 $R$ 中解释为 $R$ 的单位元 $1$ 相加 ${n\choose k}$ 次。
	\item 环 $R$ 的特征是最小正整数 $n$，使 $1+1+\cdots+1=0$（$n$ 次）在 $R$ 中成立；若不存在这样的整数，则称 $R$ 的特征为 $0$。例如，$\ZZ/n\ZZ$ 的特征为 $n$，$\ZZ$ 的特征为 $0$。
	\begin{enumerate}[label=(\alph*)]
		\item 证明映射 $\ZZ\to R$，
		\[
		k\mapsto
		\begin{cases}
			1+1+\cdots+1\ (k\text{ 次}),&k>0,\\
			0,&k=0,\\
			-1-1-\cdots-1\ (-k\text{ 次}),&k<0,
		\end{cases}
		\]
		是环同态，其核为 $n\ZZ$，其中 $n$ 是 $R$ 的特征。（这解释了为什么对没有任何个 $1$ 之和为零的环使用“特征 $0$”而非旧式说法“特征 $\infty$”。）
		\item 确定 $\QQ$、$\ZZ[x]$、$(\ZZ/n\ZZ)[x]$ 的特征。
		\item 证明若 $p$ 是素数且 $R$ 是特征 $p$ 的交换环，则对所有 $a,b\in R$ 有 $(a+b)^p=a^p+b^p$。
	\end{enumerate}
	\item 证明非零 Boolean 环的特征为 $2$（参见第 7.1 节练习 15）。
	\item 证明整环的特征是素数 $p$ 或 $0$（参见练习 26）。
	\item 设 $R$ 是交换环。回忆（第 7.1 节练习 13）元素 $x\in R$ 称为幂零，若 $x^n=0$ 对某个 $n\in\ZZ^+$ 成立。证明所有幂零元构成一个理想，称为 $R$ 的 nilradical，记作 $\mathfrak N(R)$。\emph{提示：}用二项式定理证明 $\mathfrak N(R)$ 对加法封闭。
	\item 证明若 $R$ 是交换环，$\mathfrak N(R)$ 是其 nilradical，则 $R/\mathfrak N(R)$ 中唯一的幂零元是零；即证明 $\mathfrak N(R/\mathfrak N(R))=0$。
	\item 证明矩阵
	\[
	\begin{pmatrix}0&1\\0&0\end{pmatrix},
	\qquad
	\begin{pmatrix}0&0\\1&0\end{pmatrix}
	\]
	是 $M_2(\ZZ)$ 中的幂零元，但它们的和不是幂零元（注意这两个矩阵不交换）。推出非交换环 $M_2(\ZZ)$ 中幂零元全体不是理想。
	\item 设 $\varphi:R\to S$ 是环同态。证明若 $x$ 是 $R$ 中的幂零元，则 $\varphi(x)$ 是 $S$ 中的幂零元。
	\item 假设 $R$ 交换。令 $p(x)=a_nx^n+a_{n-1}x^{n-1}+\cdots+a_1x+a_0\in R[x]$。
	\begin{enumerate}[label=(\alph*)]
		\item 证明 $p(x)$ 是 $R[x]$ 中的单位，当且仅当 $a_0$ 是单位且 $a_1,a_2,\ldots,a_n$ 都是 $R$ 中的幂零元。\emph{提示：}参见第 7.1 节练习 14。
		\item 证明 $p(x)$ 是 $R[x]$ 中的幂零元，当且仅当 $a_0,a_1,\ldots,a_n$ 都是 $R$ 中的幂零元。
	\end{enumerate}
	\item 设 $I,J$ 是 $R$ 的理想。
	\begin{enumerate}[label=(\alph*)]
		\item 证明 $I+J$ 是包含 $I$ 与 $J$ 的最小理想。
		\item 证明 $IJ$ 是包含在 $I\cap J$ 中的理想。
		\item 给出一个例子使 $IJ\neq I\cap J$。
		\item 若 $R$ 交换且 $I+J=R$，证明 $IJ=I\cap J$。
	\end{enumerate}
	\item 设 $I,J,K$ 是 $R$ 的理想。
	\begin{enumerate}[label=(\alph*)]
		\item 证明 $I(J+K)=IJ+IK$ 且 $(I+J)K=IK+JK$。
		\item 若 $J\subseteq I$，证明 $I\cap(J+K)=J+(I\cap K)$。
	\end{enumerate}
	\item 设 $I$ 是 $\ZZ[x]$ 中所有常数项为零的多项式组成的理想。证明
	\[
	I^n=\{a_nx^n+a_{n+1}x^{n+1}+\cdots+a_{n+m}x^{n+m}\mid a_i\in\ZZ,\,m\geq0\}
	\]
	即首个非零项次数至少为 $n$ 的多项式全体。
	\item 若理想 $N$ 满足 $N^n=0$ 对某个 $n\geq1$ 成立，则称 $N$ 为幂零理想。证明在环 $\ZZ/p^n\ZZ$ 中，理想 $p\ZZ/p^n\ZZ$ 是幂零理想。
\end{enumerate}



\section{理想的性质}

本节中 $R$ 总是带单位元 $1\neq0$ 的环。若 $A\subseteq R$，所有包含 $A$ 的理想的交仍是理想，因此存在包含 $A$ 的最小理想。

\medskip
\noindent\textbf{定义.} 对 $A\subseteq R$，记 $(A)$ 为包含 $A$ 的最小理想，称为由 $A$ 生成的理想。由单个元素生成的理想称为\emph{主理想}；由有限集合生成的理想称为\emph{有限生成理想}。

若 $R$ 交换，则
\[
  (A)=\{r_1a_1+\cdots+r_na_n\mid n\geq0,
  \ r_i\in R,
  \ a_i\in A\}.
\]
特别地，$(a)=Ra=\{ra\mid r\in R\}$。在交换环中，$b\in(a)$ 当且仅当 $a$ 整除 $b$；这也等价于 $(b)\subseteq(a)$。

\medskip
\noindent\textbf{例.}
\begin{enumerate}[label=(\arabic*)]
  \item $0=(0)$，$R=(1)$。在 $\ZZ$ 中，每个理想都是主理想 $n\ZZ=(n)$；而 $(m,n)$ 由 $m$ 与 $n$ 的最大公因子生成。
  \item $\ZZ[x]$ 中的理想 $(2,x)$ 不是主理想；它恰是所有常数项为偶数的整系数多项式。
  \item 群环 $RG$ 的增广理想由 $\{g-1\mid g\in G\}$ 生成；若 $G=\langle a\rangle$ 是循环群，则增广理想由 $a-1$ 生成。
\end{enumerate}

\medskip
\noindent\textbf{命题 9.} 设 $I$ 是 $R$ 的理想。
\begin{enumerate}[label=(\arabic*)]
  \item $I=R$ 当且仅当 $I$ 含有单位。
  \item 若 $R$ 是交换环，则 $R$ 是域当且仅当 $R$ 的理想只有 $0$ 与 $R$。
\end{enumerate}

\begin{proof}
若 $I$ 含有单位 $u$，取 $v\in R$ 使 $vu=1$，由于 $I$ 是理想，$1=vu\in I$，故 $I=R$。反过来若 $I=R$，当然含有单位 $1$。若 $R$ 是域且 $I$ 是非零理想，取 $0\neq a\in I$，则 $a$ 是单位，故 $I=R$。反过来，若 $R$ 的理想只有 $0$ 与 $R$，对任意 $0\neq a\in R$，主理想 $(a)$ 非零，故等于 $R$，于是 $1=ra$ 对某个 $r\in R$ 成立；因 $R$ 交换，$a$ 是单位。故 $R$ 是域。
\end{proof}

\medskip
\noindent\textbf{推论 10.} 若 $R$ 是域，则从 $R$ 到任意环的非零环同态都是单射。

\begin{proof}
非零同态的核是 $R$ 的真理想；由命题 9，域 $R$ 的真理想只能是 $0$。
\end{proof}

\medskip
\noindent\textbf{定义.} 真理想 $M$ 称为\emph{极大理想}，若包含 $M$ 的理想只有 $M$ 与 $R$。

\medskip
\noindent\textbf{命题 11.} 在带单位元的环中，每个真理想都包含在某个极大理想中。

\begin{proof}
设 $I$ 是真理想。考虑所有包含 $I$ 的真理想组成的集合，按包含关系偏序。该集合非空，因为含有 $I$。若 $\mathcal C$ 是其中的一条链，则 $\bigcup_{J\in\mathcal C}J$ 是理想：两个元素同属链中某个较大的理想，因而对加法、取负和左右乘法封闭。它仍是真理想，否则若含有 $1$，则 $1$ 已在链中某个成员里，矛盾。由 Zorn 引理，存在极大元；这就是包含 $I$ 的极大理想。
\end{proof}

\medskip
\noindent\textbf{命题 12.} 假设 $R$ 交换。理想 $M$ 是极大理想，当且仅当 $R/M$ 是域。

\begin{proof}
由格同构定理，$R/M$ 的理想与 $R$ 中包含 $M$ 的理想一一对应。于是 $M$ 极大当且仅当 $R/M$ 的理想只有 $0$ 与 $R/M$。由命题 9，这当且仅当 $R/M$ 是域。
\end{proof}

\medskip
\noindent\textbf{定义.} 假设 $R$ 交换。真理想 $P$ 称为\emph{素理想}，若 $ab\in P$ 蕴含 $a\in P$ 或 $b\in P$。

\medskip
\noindent\textbf{命题 13.} 假设 $R$ 交换。理想 $P$ 是素理想，当且仅当 $R/P$ 是整环。

\begin{proof}
在商环中 $(a+P)(b+P)=0$ 当且仅当 $ab\in P$。因此 $R/P$ 没有非零零因子，正等价于：若 $ab\in P$，则 $a\in P$ 或 $b\in P$。
\end{proof}

\medskip
\noindent\textbf{推论 14.} 交换环中的每个极大理想都是素理想。

\begin{proof}
若 $M$ 极大，则 $R/M$ 是域；域是整环，所以由命题 13，$M$ 素。
\end{proof}

\medskip
\noindent\textbf{例.}
\begin{enumerate}[label=(\arabic*)]
  \item 在 $\ZZ$ 中，极大理想为 $p\ZZ$，其中 $p$ 为素数；素理想为 $0$ 和所有 $p\ZZ$。
  \item 在连续函数环 $R=C([0,1],\RR)$ 中，对 $c\in[0,1]$，
  \[
    M_c=\{f\in R\mid f(c)=0\}
  \]
  是极大理想，因为求值映射 $f\mapsto f(c)$ 的商同构于域 $\RR$。
  \item 在非交换环中，极大理想和商为域之间需要更谨慎地区分。若商环是除环，则相应双边理想是极大双边理想；但非交换环也可能有许多一侧理想。
\end{enumerate}

\subsection*{练习}
在以下各题中，设 $R$ 是带单位元 $1\neq0$ 的环。
\begin{enumerate}[label=\arabic*.]
	\item 设 $L_j$ 是 $M_n(R)$ 的左理想，其元素只允许第 $j$ 列任意、其余列全为零；设 $E_{ij}$ 是 $M_n(R)$ 中第 $i,j$ 元为 $1$、其余元为零的矩阵。证明对任意 $i$，有 $L_j=M_n(R)E_{ij}$。（参见第 7.2 节练习 6。）
	\item 假设 $R$ 交换。证明群环 $RG$ 的增广理想由 $\{g-1\mid g\in G\}$ 生成。证明若 $G=\langle a\rangle$ 是循环群，则增广理想由 $a-1$ 生成。
	\item \begin{enumerate}[label=(\alph*)]
		\item 设 $p$ 是素数，$G$ 是阶为 $p^n$ 的阿贝尔群。证明群环 $\FF_pG$ 的 nilradical 是增广理想（参见第 7.3 节练习 29）。\emph{提示：}使用上一题。
		\item 设 $G=\{g_1,\ldots,g_n\}$ 是有限群，且 $R$ 交换。证明若 $r$ 是 $RG$ 的增广理想中任意元素，则
		\[
		r(g_1+\cdots+g_n)=0.
		\]
		\emph{提示：}使用上一题。
	\end{enumerate}
	\item 假设 $R$ 交换。证明 $R$ 是域当且仅当 $0$ 是极大理想。
	\item 证明若 $M$ 是理想且 $R/M$ 是域，则 $M$ 是极大理想（不假定 $R$ 交换）。
	\item 证明 $R$ 是除环当且仅当它的左理想只有 $(0)$ 与 $R$。（把“左”换成“右”也成立。）
	\item 设 $R$ 是带 $1$ 的交换环。证明多项式环 $R[x]$ 中由 $x$ 生成的主理想是素理想，当且仅当 $R$ 是整环。证明 $(x)$ 是极大理想，当且仅当 $R$ 是域。
	\item 设 $R$ 是整环。证明对 $a,b\in R$，$(a)=(b)$ 当且仅当存在 $R$ 的单位 $u$ 使 $a=ub$。
	\item 令 $R$ 为 $[0,1]$ 上所有连续函数组成的环，令 $I$ 为满足 $f(1/3)=f(1/2)=0$ 的所有函数 $f(x)$ 的集合。证明 $I$ 是 $R$ 的理想，但不是素理想。
	\item 假设 $R$ 交换。证明若 $P$ 是 $R$ 的素理想且 $P$ 不含零因子，则 $R$ 是整环。
	\item 假设 $R$ 交换。设 $I,J$ 是 $R$ 的理想，并设 $P$ 是包含 $IJ$ 的素理想（例如 $P$ 包含 $I\cap J$）。证明 $I$ 或 $J$ 包含在 $P$ 中。
	\item 假设 $R$ 交换，且 $I=(a_1,a_2,\ldots,a_n)$、$J=(b_1,b_2,\ldots,b_m)$ 是 $R$ 中两个有限生成理想。证明积理想 $IJ$ 由元素 $a_ib_j$（$i=1,
	\ldots,n$，$j=1,\ldots,m$）有限生成。
	\item 设 $\varphi:R\to S$ 是交换环同态。
	\begin{enumerate}[label=(\alph*)]
		\item 证明若 $P$ 是 $S$ 的素理想，则 $\varphi^{-1}(P)=R$ 或 $\varphi^{-1}(P)$ 是 $R$ 的素理想。把它应用于 $R$ 是 $S$ 的子环且 $\varphi$ 是包含同态的情形，推出若 $P$ 是 $S$ 的素理想，则 $P\cap R$ 要么等于 $R$，要么是 $R$ 的素理想。
		\item 证明若 $M$ 是 $S$ 的极大理想且 $\varphi$ 满，则 $\varphi^{-1}(M)$ 是 $R$ 的极大理想。举例说明若 $\varphi$ 不满则不一定成立。
	\end{enumerate}
	\item 假设 $R$ 交换。设 $x$ 是不定元，$f(x)$ 是 $R[x]$ 中次数 $n\geq1$ 的首一多项式。用横线表示模理想 $(f(x))$ 到商环 $R[x]/(f(x))$。
	\begin{enumerate}[label=(\alph*)]
		\item 证明 $R[x]/(f(x))$ 的每个元素都可写成 $\overline{p(x)}$，其中 $p(x)\in R[x]$ 的次数小于 $n$，即
		\[
		R[x]/(f(x))=\{\overline{a_0+a_1x+\cdots+a_{n-1}x^{n-1}}\mid a_0,
		\ldots,a_{n-1}\in R\}.
		\]
		\emph{提示：}若 $f(x)=x^n+b_{n-1}x^{n-1}+\cdots+b_0$，则在商环中 $\bar x^n=-(b_{n-1}\bar x^{n-1}+\cdots+b_0)$。用它化简 $x$ 的高次幂。
		\item 证明若 $p(x)$ 与 $q(x)$ 是 $R[x]$ 中不同且次数都小于 $n$ 的多项式，则 $\overline{p(x)}\neq\overline{q(x)}$。\emph{提示：}否则 $p(x)-q(x)$ 是首一多项式 $f(x)$ 的 $R[x]$ 倍数。
		\item 若 $f(x)=a(x)b(x)$，且 $a(x),b(x)$ 的次数都小于 $n$，证明 $\overline{a(x)}$ 是 $R[x]/(f(x))$ 中的零因子。
		\item 若 $f(x)=x^n-a$，其中 $a\in R$ 是幂零元，证明 $\bar x$ 在 $R[x]/(f(x))$ 中是幂零元。
		\item 设 $p$ 是素数，$R=\FF_p$，且 $f(x)=x^p-a$，其中 $a\in\FF_p$。证明 $\bar x-a$ 在 $R[x]/(f(x))$ 中是幂零元。\emph{提示：}使用第 7.3 节练习 26(c)。
	\end{enumerate}
	\item 设 $x^2+x+1$ 是多项式环 $\FF_2[x]$ 的元素，并用横线表示通过商环 $E=\FF_2[x]/(x^2+x+1)$。
	\begin{enumerate}[label=(\alph*)]
		\item 证明 $E$ 有 $4$ 个元素：$\bar0,\bar1,\bar x,\bar x+\bar1$。
		\item 写出 $E$ 的 $4\times4$ 加法表，并推出加法群 $E$ 同构于 Klein 四元群。
		\item 写出 $E^\times$ 的 $4\times4$ 乘法表，并证明 $E^\times$ 同构于 $3$ 阶循环群。推出 $E$ 是域。
	\end{enumerate}
	\item 设 $x^4-16$ 是多项式环 $\ZZ[x]$ 的元素，并用横线表示通过商环 $E=\ZZ[x]/(x^4-16)$。
	\begin{enumerate}[label=(\alph*)]
		\item 找一个次数 $\leq3$ 的多项式，使它与 $7x^{13}-11x^9+5x^5-2x^3+3$ 模 $(x^4-16)$ 同余。
		\item 证明 $\bar x-2$ 与 $\bar x+2$ 是 $E$ 中的零因子。
	\end{enumerate}
	\item 设 $x^3-2x+1$ 是多项式环 $\ZZ[x]$ 的元素，并用横线表示通过商环 $E=\ZZ[x]/(x^3-2x+1)$。令
	\[
	p(x)=2x^7-7x^5+4x^3-9x+1,
	\qquad q(x)=(x-1)^4.
	\]
	\begin{enumerate}[label=(\alph*)]
		\item 将 $E$ 中下列每个元素表示成 $\overline{f(x)}$，其中 $f(x)$ 的次数 $\leq2$：$\overline{p(x)}$、$\overline{q(x)}$、$\overline{p(x)}+\overline{q(x)}$、$\overline{p(x)}\,\overline{q(x)}$。
		\item 证明 $E$ 不是整环。
		\item 证明 $\bar x$ 是 $E$ 中的单位。
	\end{enumerate}
	\item 证明若 $R$ 是整环，$R[[x]]$ 是不定元 $x$ 的形式幂级数环，则由 $x$ 生成的主理想是素理想（参见第 7.2 节练习 3）。证明由 $x$ 生成的主理想是极大理想，当且仅当 $R$ 是域。
	\item 设 $R$ 是有限带单位交换环。证明 $R$ 的每个素理想都是极大理想。
	\item 证明没有零因子的非零有限交换环是域（若有单位元，这是推论 3；本题不要假定环有 $1$）。
	\item 证明带单位元 $1\neq0$ 且没有零因子的有限环是域（可以引用 Wedderburn 定理）。
	\item 设 $p\in\ZZ^+$ 是素数，定义 $\FF_p$ 上的四元数为
	\[
	a+bi+cj+dk,
	\qquad a,b,c,d\in\ZZ/p\ZZ,
	\]
	加法逐分量进行，乘法用实四元数中 $i,j,k$ 的相同关系定义。
	\begin{enumerate}[label=(\alph*)]
		\item 证明 $\FF_p$ 上的四元数是整 Hamilton 四元数环的同态像（参见第 7.1 节）。
		\item 证明 $\FF_p$ 上的四元数含有零因子（因而不可能是除环）。\emph{提示：}使用上一题。
	\end{enumerate}
	\item 证明在 Boolean 环中（参见第 7.1 节练习 15），每个素理想都是极大理想。
	\item 证明在 Boolean 环中，每个有限生成理想都是主理想。
	\item 假设 $R$ 交换，且对每个 $a\in R$ 都存在整数 $n>1$（依赖于 $a$）使 $a^n=a$。证明 $R$ 的每个素理想都是极大理想。
	\item 证明交换环 $R$ 的素理想包含每个幂零元（参见第 7.1 节练习 13）。推出 $R$ 的 nilradical（参见第 7.3 节练习 29）包含在 $R$ 的所有素理想之交中。（第 15.2 节将证明 nilradical 等于所有素理想之交。）
	\item 设 $R$ 是带 $1\neq0$ 的交换环。证明若 $a$ 是 $R$ 的幂零元，则对所有 $b\in R$，$1-ab$ 是单位。
	\item 证明若 $R$ 是交换环，$N=(a_1,a_2,\ldots,a_m)$，其中每个 $a_i$ 都是幂零元，则 $N$ 是幂零理想（参见第 7.3 节练习 37）。推出若 $R$ 的 nilradical 是有限生成的，则它是幂零理想。
	\item 设 $p$ 是素数，$G$ 是阶为 $p$ 的幂的有限群（即 $p$-群）。证明群环 $(\ZZ/p\ZZ)G$ 的增广理想是幂零理想。（注意这个环可以是非交换的。）\emph{提示：}使用练习 2。
	\item 设 $I$ 是交换环 $R$ 的理想，定义
	\[
	\rad I=\{r\in R\mid r^n\in I\text{ 对某个 }n\in\ZZ^+\},
	\]
	称为 $I$ 的根。证明 $\rad I$ 是包含 $I$ 的理想，并且 $(\rad I)/I$ 是商环 $R/I$ 的 nilradical，即 $(\rad I)/I=\mathfrak N(R/I)$（参见第 7.3 节练习 29）。
	\item 若交换环 $R$ 的理想 $I$ 满足 $\rad I=I$，则称 $I$ 为根理想。
	\begin{enumerate}[label=(\alph*)]
		\item 证明 $R$ 的每个素理想都是根理想。
		\item 设 $n>1$。证明 $0$ 是 $\ZZ/n\ZZ$ 中的根理想，当且仅当 $n$ 是不同素数一次幂的乘积（即 $n$ 无平方因子）。推出 $(n)$ 是 $\ZZ$ 的根理想，当且仅当 $n$ 是 $\ZZ$ 中不同素数的乘积。
	\end{enumerate}
	\item 设 $I$ 是交换环 $R$ 的理想，并定义
	\[
	\Jac I=\bigcap\{M\mid M\text{ 是 }R\text{ 的极大理想且 }I\subseteq M\},
	\]
	约定 $\Jac R=R$。（当 $I=0$ 时，$\Jac 0$ 称为环 $R$ 的 Jacobson 根；一般 $\Jac I$ 是 $R/I$ 的 Jacobson 根在 $R$ 中的原像。）
	\begin{enumerate}[label=(\alph*)]
		\item 证明 $\Jac I$ 是包含 $I$ 的 $R$ 的理想。
		\item 证明 $\rad I\subseteq\Jac I$，其中 $\rad I$ 是练习 30 中定义的根。
		\item 设 $n>1$。用 $n$ 的素因子分解描述 $\Jac(n\ZZ)$。
	\end{enumerate}
	\item 设 $R$ 是闭区间 $[0,1]$ 上所有连续函数构成的环，并且对每个 $c\in[0,1]$，令 $M_c=\{f\in R\mid f(c)=0\}$（回忆 $M_c$ 已被证明是 $R$ 的极大理想）。
	\begin{enumerate}[label=(\alph*)]
		\item 证明若 $M$ 是 $R$ 的任意极大理想，则存在实数 $c\in[0,1]$ 使 $M=M_c$。
		\item 证明若 $b,c$ 是 $[0,1]$ 中不同点，则 $M_b\neq M_c$。
		\item 证明 $M_c$ 不等于由 $x-c$ 生成的主理想。
		\item 证明 $M_c$ 不是有限生成理想。
	\end{enumerate}
	上一题说明，闭区间 $[0,1]$ 的点与 $[0,1]$ 上连续函数环 $R$ 的极大理想集合之间有双射 $c\mapsto M_c$。对 $\RR$ 的任意子集 $X$，更一般地，对任意完全正则拓扑空间 $X$，映射 $c\mapsto M_c$ 是从 $X$ 到 $R$ 的极大理想集合的单射，其中 $R$ 是 $X$ 上有界连续实值函数环，$M_c$ 是在 $c$ 处为零的函数理想。令 $\beta(X)$ 为 $R$ 的极大理想集合。可以在 $\beta(X)$ 上定义一个拓扑，使得把 $X$ 与其在 $\beta(X)$ 中的像识别后，$X$（带原拓扑）成为 $\beta(X)$ 的子空间；而且 $\beta(X)$ 是紧空间，称为 $X$ 的 Stone--\v Cech 紧化。
	\item 设 $R$ 是从 $\RR$ 到 $\RR$ 的所有连续函数构成的环，并且对每个 $c\in\RR$，令 $M_c=\{f\in R\mid f(c)=0\}$。
	\begin{enumerate}[label=(\alph*)]
		\item 令 $I$ 为 $R$ 中所有紧支撑函数 $f(x)$ 的集合（即当 $|x|$ 充分大时 $f(x)=0$）。证明 $I$ 是 $R$ 的理想，但不是素理想。
		\item 令 $M$ 是包含 $I$ 的极大理想（由 (a)，它真包含 $I$）。证明对任意 $c\in\RR$，$M\neq M_c$（参见上一题）。
	\end{enumerate}
	\item 设 $A=(a_1,a_2,\ldots,a_n)$ 是 $R$ 的非零有限生成理想。证明存在一个理想 $B$，它关于“不包含 $A$”这一性质是极大的。\emph{提示：}使用 Zorn 引理。
	\item 假设 $R$ 交换。证明 $R$ 的素理想集合关于包含关系有极小元（可能是零理想）。\emph{提示：}使用 Zorn 引理。
	\item 若交换环 $R$ 只有一个极大理想，则称 $R$ 为局部环。证明若 $R$ 是局部环且极大理想为 $M$，则 $R-M$ 中每个元素都是单位。反过来，证明若带 $1$ 交换环 $R$ 的非单位全体构成一个理想 $M$，则 $R$ 是以 $M$ 为唯一极大理想的局部环。
	\item 证明所有分母为奇数的有理数构成的环是局部环，其唯一极大理想是由 $2$ 生成的主理想。
	\item 沿用第 7.1 节练习 26 的记号，设 $K$ 是域，$v$ 是 $K$ 上的离散赋值，$R$ 是 $v$ 的赋值环。对每个整数 $k\geq0$，定义
	\[
	A_k=\{r\in R\mid v(r)\geq k\}\cup\{0\}.
	\]
	\begin{enumerate}[label=(\alph*)]
		\item 证明 $A_k$ 是主理想，且 $A_0\supseteq A_1\supseteq A_2\supseteq\cdots$。
		\item 证明若 $I$ 是 $R$ 的任意非零理想，则 $I=A_k$ 对某个 $k\geq0$ 成立。推出 $R$ 是以 $A_1$ 为唯一极大理想的局部环。
	\end{enumerate}
	\item 假设 $R$ 交换。证明下列条件等价（亦参见第 7.1 节练习 13、14）：
	\begin{enumerate}[label=(\roman*)]
		\item $R$ 恰有一个素理想；
		\item $R$ 的每个元素要么是幂零元，要么是单位；
		\item $R/\mathfrak N(R)$ 是域（参见第 7.3 节练习 29）。
	\end{enumerate}
	\item 交换环 $R$ 的真理想 $Q$ 称为 primary ideal，若 $ab\in Q$ 且 $a\notin Q$，则对某个正整数 $n$ 有 $b^n\in Q$。（注意该定义中 $a,b$ 的对称性意味着：若 $Q$ 是 primary ideal 且 $ab\in Q$，其中 $a,b$ 都不在 $Q$ 中，则 $a$ 的某个正幂与 $b$ 的某个正幂都在 $Q$ 中。）证明关于 primary ideal 的下列事实。
	\begin{enumerate}[label=(\alph*)]
		\item $\ZZ$ 的 primary ideal 是 $0$ 和 $(p^n)$，其中 $p$ 是素数，$n$ 是正整数。
		\item $R$ 的每个素理想都是 primary ideal。
		\item $R$ 的理想 $Q$ 是 primary ideal，当且仅当 $R/Q$ 中每个零因子都是 $R/Q$ 的幂零元。
		\item 若 $Q$ 是 primary ideal，则 $\rad(Q)$ 是素理想（参见练习 30）。
	\end{enumerate}
\end{enumerate}



\section{分式环}

本节 $R$ 为交换环。命题 2 说明：若 $a$ 不是零因子且 $ab=ac$，则可以消去 $a$。本节的目标是把这类可消去的元素变成真正的单位。对整环而言，所得的较大环是域，称为分式域或商域；原型是 $\ZZ\subseteq\QQ$。

\medskip
\noindent\textbf{定理 15.} 设 $R$ 是交换环，$D\subseteq R$ 非空，不含 $0$，不含零因子，并且对乘法封闭。则存在带单位元的交换环 $Q$，包含 $R$ 作为子环，并使 $D$ 的每个元素在 $Q$ 中成为单位。进一步：
\begin{enumerate}[label=(\arabic*)]
  \item $Q$ 的每个元素可写为 $rd^{-1}$，其中 $r\in R,d\in D$。若 $D=R-\{0\}$，则 $Q$ 是域。
  \item $Q$ 是使 $D$ 中元素都可逆的最小环：若 $\varphi:R\to S$ 是到带单位交换环 $S$ 的单射同态，并且每个 $\varphi(d)$（$d\in D$）都是 $S$ 的单位，则存在单射同态 $\Phi:Q\to S$ 延拓 $\varphi$。
\end{enumerate}

\begin{proof}
令
\[
  \mathcal F=\{(r,d)\mid r\in R,
  \ d\in D\}.
\]
在 $\mathcal F$ 上定义关系
\[
  (r,d)\sim(s,e)
  \quad\Longleftrightarrow\quad re=sd.
\]
由于 $D$ 中元素不是零因子，容易验证 $\sim$ 是等价关系：反身性和对称性显然；若 $(r,d)\sim(s,e)$ 且 $(s,e)\sim(t,f)$，即 $re=sd$、$sf=te$，则
\[
  (rf)ed=(re)fd=(sd)fd=sdfd=d(sf)d=d(te)d=(td)ed,
\]
消去非零因子 $ed\in D$，得 $rf=td$，故 $(r,d)\sim(t,f)$。

记 $(r,d)$ 的等价类为 $r/d$。在等价类上定义
\[
  \frac ab+
  \frac cd=\frac{ad+bc}{bd},
  \qquad
  \frac ab\frac cd=\frac{ac}{bd}.
\]
若改变代表元，由等价关系和 $R$ 的交换性可直接验证这些定义良定义。于是等价类集合在这些运算下成为交换环；零元为 $0/1$，单位元为 $1/1$，而 $a/b$ 的加法逆元为 $(-a)/b$。

映射
\[
  R\longrightarrow Q,
  \qquad r\longmapsto \frac r1
\]
是环同态。若 $r/1=0/1$，则 $r=0$，所以它是单射，于是可把 $R$ 视为 $Q$ 的子环。对 $d\in D$，元素 $d/1$ 的逆元是 $1/d$，因为 $(d/1)(1/d)=d/d=1/1$。每个 $Q$ 的元素按构造都是 $r/d=(r/1)(d/1)^{-1}$。若 $D=R-\{0\}$ 且 $r/d\neq0$，则 $r\neq0$，故 $r\in D$，并且
\[
  \frac dr
\]
是 $r/d$ 的逆元；因此 $Q$ 是域。

最后证明最小性。设 $\varphi:R\to S$ 满足定理条件。定义
\[
  \Phi\!\left(\frac rd\right)=\varphi(r)\varphi(d)^{-1}.
\]
若 $r/d=s/e$，即 $re=sd$，则
\[
  \varphi(r)\varphi(e)=\varphi(s)\varphi(d),
\]
乘以 $\varphi(d)^{-1}\varphi(e)^{-1}$ 得到 $\varphi(r)\varphi(d)^{-1}=\varphi(s)\varphi(e)^{-1}$，故 $\Phi$ 良定义。它保持加法与乘法由定义直接验证，并显然延拓 $\varphi$。若 $\Phi(r/d)=0$，则 $\varphi(r)\varphi(d)^{-1}=0$，故 $\varphi(r)=0$；由 $\varphi$ 单射，$r=0$，从而 $r/d=0$。所以 $\Phi$ 单射。
\end{proof}

\medskip
\noindent\textbf{定义.} 上述 $Q$ 称为 $R$ 关于分母集合 $D$ 的\emph{分式环}，记为 $D^{-1}R$。若 $R$ 是整环且 $D=R-\{0\}$，则 $Q$ 称为 $R$ 的\emph{分式域}，记为 $\Frac(R)$。

\medskip
\noindent\textbf{推论 16.} 设 $R$ 是整环，$Q=\Frac(R)$。若 $F$ 是任意包含 $R$ 的域，则 $Q$ 同构于 $F$ 的某个子域。

\begin{proof}
取 $D=R-\{0\}$，把包含映射 $R\hookrightarrow F$ 作为定理 15 中的 $\varphi$。由于 $F$ 是域，$D$ 中每个元素的像都是单位。定理 15 给出单射 $Q\to F$，于是 $Q$ 同构于其像，即 $F$ 的一个子域。
\end{proof}

例如 $\Frac(\ZZ)=\QQ$，若 $F$ 是域，则 $\Frac(F[x])$ 记为 $F(x)$，称为 $F$ 上的有理函数域。若 $D=\{1,d,d^2,\ldots\}$，则分式环常记为 $R[1/d]$。

\subsection*{练习}
在以下各题中，设 $R$ 是带单位元 $1\neq0$ 的交换环。
\begin{enumerate}[label=\arabic*.]
	\item 补全定理 15 证明中的所有细节。
	\item 设 $R$ 是整环，$D$ 是 $R$ 的非空乘法封闭子集。证明分式环 $D^{-1}R$ 同构于 $R$ 的分式域的某个子环（因而也是整环）。
	\item 设 $F$ 是域。证明 $F$ 含有唯一的最小子域 $F_0$，并且 $F_0$ 同构于 $\QQ$ 或某个素数 $p$ 的 $\ZZ/p\ZZ$（$F_0$ 称为 $F$ 的素子域）。\emph{提示：}参见第 7.3 节练习 26。
	\item 证明 $\RR$ 的任意子域都必须包含 $\QQ$。
	\item 若 $F$ 是域，证明 $F[[x]]$（系数在 $F$ 中、不定元为 $x$ 的形式幂级数环）的分式域是形式 Laurent 级数域 $F((x))$（参见第 7.2 节练习 3 和 5）。证明幂级数环 $\ZZ[[x]]$ 的分式域真包含于 Laurent 级数域 $\QQ((x))$ 的真子域中。\emph{提示：}考虑 $e^x$ 的级数。
	\item 证明实数域 $\RR$ 含有一个子环 $A$，满足 $1\in A$，并且 $A$ 关于性质 $1/2\notin A$ 在包含关系下极大。\emph{提示：}使用 Zorn 引理。（第 15.3 节练习 13 说明 $\RR$ 是 $A$ 的分式域，因此 $\RR$ 是一个真子环的分式域。）
\end{enumerate}



\section{中国剩余定理}

本节除非另有说明，环均为带 $1\neq0$ 的交换环。环的直积按分量加法与乘法定义。两个理想 $A,B$ 称为\emph{互素}或\emph{余极大}，若
\[
  A+B=R.
\]

\medskip
\noindent\textbf{定理 17（中国剩余定理）.} 设 $A_1,A_2,\ldots,A_k$ 是 $R$ 的理想。映射
\[
  R\longrightarrow R/A_1\times R/A_2\times\cdots\times R/A_k,
  \qquad r\longmapsto(r+A_1,r+A_2,\ldots,r+A_k)
\]
是环同态，核为 $A_1\cap A_2\cap\cdots\cap A_k$。若这些理想两两互素，则此映射满，并且
\[
  A_1\cap A_2\cap\cdots\cap A_k=A_1A_2\cdots A_k,
\]
从而
\[
  R/(A_1A_2\cdots A_k)
  \cong R/A_1\times R/A_2\times\cdots\times R/A_k.
\]

\begin{proof}
该映射显然是环同态；一个元素 $r$ 落入核，当且仅当 $r\in A_i$ 对所有 $i$ 成立，故核为交。

先证 $k=2$。令 $A=A_1$、$B=A_2$。若 $A+B=R$，取 $x\in A$、$y\in B$ 使 $x+y=1$。在 $R/A\times R/B$ 中，$x$ 的像为 $(0,1)$，$y$ 的像为 $(1,0)$。因此任意 $(r+A,s+B)$ 等于 $sy+rx$ 的像，故映射满。总有 $AB\subseteq A\cap B$。若 $c\in A\cap B$，则
\[
  c=c(x+y)=cx+cy.
\]
由于 $c\in B$ 且 $x\in A$，有 $cx\in AB$；由于 $c\in A$ 且 $y\in B$，有 $cy\in AB$。故 $c\in AB$，于是 $A\cap B=AB$。

一般情形由归纳得到。若 $A_1,
\ldots,A_k$ 两两互素，则 $A_1\cdots A_{k-1}$ 与 $A_k$ 互素，并且 $A_1\cap\cdots\cap A_{k-1}=A_1\cdots A_{k-1}$。把二理想情形应用于 $(A_1\cdots A_{k-1},A_k)$，即可同时得到满性和交等于积。
\end{proof}

当 $m,n$ 互素时，取 $R=\ZZ$、$A=m\ZZ$、$B=n\ZZ$，得到经典形式
\[
  \ZZ/mn\ZZ\cong \ZZ/m\ZZ\times\ZZ/n\ZZ.
\]
这说明模互素整数的同余方程组有解且解模乘积唯一。

\medskip
\noindent\textbf{推论 18.} 若
\[
  n=p_1^{\alpha_1}p_2^{\alpha_2}\cdots p_k^{\alpha_k}
\]
是正整数 $n$ 的不同素数幂分解，则
\[
  \ZZ/n\ZZ
  \cong
  \ZZ/p_1^{\alpha_1}\ZZ\times\ZZ/p_2^{\alpha_2}\ZZ\times\cdots\times
  \ZZ/p_k^{\alpha_k}\ZZ
\]
作为环同构。特别地，
\[
  (\ZZ/n\ZZ)^\times
  \cong
  (\ZZ/p_1^{\alpha_1}\ZZ)^\times\times\cdots\times
  (\ZZ/p_k^{\alpha_k}\ZZ)^\times.
\]
比较阶数得到 Euler 函数的乘法性：若 $(a,b)=1$，则 $\varphi(ab)=\varphi(a)\varphi(b)$；且
\[
  \varphi(p^\alpha)=p^{\alpha-1}(p-1).
\]

\begin{proof}
理想 $p_i^{\alpha_i}\ZZ$ 两两互素，故中国剩余定理给出第一同构。环同构把单位送到单位并反映单位，所以诱导单位群同构。最后的 $\varphi$ 公式来自单位群阶数的乘法性，以及模 $p^\alpha$ 中恰有 $p^\alpha-p^{\alpha-1}$ 个与 $p$ 互素的剩余类。
\end{proof}

这个推论也是确定 $(\ZZ/n\ZZ)^\times$ 结构的第一步；接下来只需分别研究素数幂模的单位群。

\subsection*{练习}
在以下各题中，设 $R$ 是带单位元 $1\neq0$ 的环。
\begin{enumerate}[label=\arabic*.]
	\item 若 $e\in R$ 满足 $e^2=e$，称 $e$ 为幂等元。假设 $e$ 是 $R$ 中的幂等元并且 $er=re$ 对所有 $r\in R$ 成立。证明 $Re$ 和 $R(1-e)$ 是 $R$ 的双边理想，并且
	\[
	R\cong Re\times R(1-e).
	\]
	证明 $e$ 与 $1-e$ 分别是子环 $Re$ 与 $R(1-e)$ 的单位元。
	\item 设 $R$ 是有限 Boolean 环，且有单位元 $1\neq0$（参见第 7.1 节练习 15）。证明
	\[
	R\cong \ZZ/2\ZZ\times\cdots\times\ZZ/2\ZZ.
	\]
	\emph{提示：}使用上一题。
	\item 设 $R,S$ 是带单位元的环。证明 $R\times S$ 的每个理想都形如 $I\times J$，其中 $I$ 是 $R$ 的理想，$J$ 是 $S$ 的理想。
	\item 证明若 $R,S$ 是非零环，则 $R\times S$ 永远不是域。
	\item 设 $n_1,n_2,\ldots,n_k$ 两两互素：$(n_i,n_j)=1$ 对所有 $i\neq j$ 成立。
	\begin{enumerate}[label=(\alph*)]
		\item 证明中国剩余定理蕴含：对任意 $a_1,
		\ldots,a_k\in\ZZ$，存在 $x\in\ZZ$ 满足同余方程组
		\[
		x\equiv a_1\pmod {n_1},\qquad x\equiv a_2\pmod {n_2},\qquad \ldots,
		\qquad x\equiv a_k\pmod {n_k},
		\]
		并且解 $x$ 模 $n=n_1n_2\cdots n_k$ 唯一。
		\item 令 $\bar n_i=n/n_i$。由于 $\bar n_i$ 与 $n_i$ 互素，令 $t_i$ 是 $\bar n_i$ 模 $n_i$ 的逆元。证明 (a) 中的解由
		\[
		x\equiv a_1t_1\bar n_1+a_2t_2\bar n_2+\cdots+a_kt_k\bar n_k\pmod n
		\]
		给出。注意，元素 $t_i$ 可由预备章第 2 节所述 Euclidean 算法快速求出（若 $an_i+b\bar n_i=(n_i,\bar n_i)=1$，则 $t_i=b$）；这些 $t_i$ 随即对任意 $a_1,a_2,
		\ldots,a_k$ 快速给出同余方程组的解。
		\item 解同余方程组
		\[
		x\equiv1\pmod8,
		\qquad x\equiv2\pmod{25},
		\qquad x\equiv3\pmod{81},
		\]
		以及同余方程组
		\[
		y\equiv5\pmod8,
		\qquad y\equiv12\pmod{25},
		\qquad y\equiv47\pmod{81}.
		\]
	\end{enumerate}
	\item 设 $f_1(x),f_2(x),\ldots,f_k(x)$ 是次数相同为 $d$ 的整系数多项式。设 $n_1,n_2,\ldots,n_k$ 两两互素，即 $(n_i,n_j)=1$ 对所有 $i\neq j$ 成立。用中国剩余定理证明存在次数为 $d$ 的整系数多项式 $f(x)$，使
	\[
	f(x)\equiv f_1(x)\pmod {n_1},\qquad f(x)\equiv f_2(x)\pmod {n_2},\qquad \ldots,
	\qquad f(x)\equiv f_k(x)\pmod {n_k},
	\]
	即 $f(x)$ 的系数与 $f_i(x)$ 的系数模 $n_i$ 一致。证明若所有 $f_i(x)$ 都是首一多项式，则 $f(x)$ 也可取为首一。\emph{提示：}对每个系数分别在 $\ZZ$ 中应用中国剩余定理。
	\item 设 $m,n$ 是正整数，且 $n\mid m$。证明自然满环投影 $\ZZ/m\ZZ\to\ZZ/n\ZZ$ 在单位群上也是满的：
	\[
	(\ZZ/m\ZZ)^\times\twoheadrightarrow(\ZZ/n\ZZ)^\times.
	\]
	以下四题发展直极限及其“对偶”概念逆极限。在这些题中，$I$ 是带偏序 $\leq$ 的非空指标集（参见附录 I）。对每个 $i\in I$，令 $A_i$ 为加法阿贝尔群。在练习 8 中还假设 $I$ 是有向集：对任意 $i,j\in I$，存在 $k\in I$ 使 $i\leq k$ 且 $j\leq k$。
	\item 假设对每对指标 $i,j$，当 $i\leq j$ 时给定映射 $\rho_{ij}:A_i\to A_j$，并且满足：
	\begin{enumerate}[label=(\roman*)]
		\item 若 $i\leq j\leq k$，则 $\rho_{ik}=\rho_{jk}\circ\rho_{ij}$；
		\item 对所有 $i\in I$，$\rho_{ii}=1$。
	\end{enumerate}
	令 $B$ 为所有 $A_i$ 的不交并。对 $a\in A_i$、$b\in A_j$，定义关系 $\sim$：
	\[
	a\sim b\quad\Longleftrightarrow\quad \text{存在 }k\text{ 使 }i,j\leq k\text{ 且 }\rho_{ik}(a)=\rho_{jk}(b).
	\]
	\begin{enumerate}[label=(\alph*)]
		\item 证明 $\sim$ 是 $B$ 上的等价关系。（等价类集合称为有向系统 $\{A_i\}$ 的直极限或归纳极限，记为 $\varinjlim A_i$。在本题其余部分，令 $A=\varinjlim A_i$。）
		\item 令 $\bar x$ 表示 $x$ 在 $A$ 中的等价类，定义 $\rho_i:A_i\to A$ 为 $\rho_i(a)=\bar a$。证明若每个 $\rho_{ij}$ 都是单射，则每个 $\rho_i$ 也是单射（于是可把每个 $A_i$ 识别为 $A$ 的子集）。
		\item 假设所有 $\rho_{ij}$ 都是群同态。对 $a\in A_i$、$b\in A_j$，定义
		\[
		\bar a+\bar b=\overline{\rho_{ik}(a)+\rho_{jk}(b)},
		\]
		其中 $k$ 是任意满足 $i,j\leq k$ 的指标。证明该运算良定义，并使 $A$ 成为阿贝尔群。推出 (b) 中的映射 $\rho_i:A_i\to A$ 都是群同态。
		\item 若所有 $A_i$ 都是带 $1$ 的交换环，且所有 $\rho_{ij}$ 都是把 $1$ 送到 $1$ 的环同态，证明 $A$ 同样可赋予带 $1$ 的交换环结构，使所有 $\rho_i$ 都成为环同态。
		\item 在 (c) 的假设下，证明直极限有如下泛性质：若 $C$ 是任意阿贝尔群，并且对每个 $i\in I$ 有同态 $\phi_i:A_i\to C$，满足当 $i\leq j$ 时 $\phi_i=\phi_j\circ\rho_{ij}$，则存在唯一同态 $\phi:A\to C$，使 $\phi\circ\rho_i=\phi_i$ 对所有 $i$ 成立。
	\end{enumerate}
	\item 令 $I$ 为实直线上包含固定实数 $p$ 的所有开区间 $U=(a,b)$ 的集合。按反包含排序：$U\leq V$ 若 $V\subseteq U$（注意 $I$ 是有向集）。对每个 $U$，令 $A_U$ 为 $U$ 上所有连续实值函数构成的环。若 $V\subseteq U$，定义限制映射 $\rho_{UV}:A_U\to A_V$，$f\mapsto f|_V$。令 $A=\varinjlim A_U$。在上一题记号下，证明映射 $\rho_U:A_U\to A$ 并非单射，但全都是满射（$A$ 称为点 $p$ 处连续函数芽环）。
	现在发展逆极限。继续假设 $I$ 是偏序集（不一定有向），并且每个 $A_i$ 是群。
	\item 假设对每对指标 $i,j$，当 $i\leq j$ 时给定映射 $\pi_{ji}:A_j\to A_i$，并且满足：
	\begin{enumerate}[label=(\roman*)]
		\item 若 $i\leq j\leq k$，则 $\pi_{ji}\circ\pi_{kj}=\pi_{ki}$；
		\item 对所有 $i\in I$，$\pi_{ii}=1$。
	\end{enumerate}
	令 $P$ 为直积 $\prod_{i\in I}A_i$ 中所有元素 $(a_i)_{i\in I}$ 组成的子集，满足当 $i\leq j$ 时 $\pi_{ji}(a_j)=a_i$（这里 $a_i,a_j$ 分别是直积元素的第 $i$ 与第 $j$ 分量）。集合 $P$ 称为系统 $\{A_i\}$ 的逆极限或射影极限，记为 $\varprojlim A_i$。
	\begin{enumerate}[label=(\alph*)]
		\item 假设所有 $\pi_{ji}$ 都是群同态。证明 $P$ 是直积群的子群（参见第 5.1 节练习 15）。
		\item 假设 (a) 的假设，并令 $I=\ZZ^+$（通常序）。对每个 $i\in I$，令 $\pi_i:P\to A_i$ 是到第 $i$ 分量的投影。证明若每个 $\pi_{ji}$ 都是满射，则每个 $\pi_i$ 也是满射（于是每个 $A_i$ 是 $P$ 的商群）。
		\item 若所有 $A_i$ 都是带 $1$ 的交换环，且所有 $\pi_{ji}$ 都是把 $1$ 送到 $1$ 的环同态，证明 $P$ 同样可赋予带 $1$ 的交换环结构，使所有 $\pi_i$ 都成为环同态。
		\item 在 (a) 的假设下，证明逆极限有如下泛性质：若 $D$ 是任意群，并且对每个 $i\in I$ 有同态 $\eta_i:D\to A_i$，满足当 $i\leq j$ 时 $\eta_i=\pi_{ji}\circ\eta_j$，则存在唯一同态 $\eta:D\to P$，使 $\pi_i\circ\eta=\eta_i$ 对所有 $i$ 成立。
	\end{enumerate}
	\item 设 $p$ 是素数，令 $I=\ZZ^+$，$A_i=\ZZ/p^i\ZZ$，并令 $\pi_{ji}$ 为自然投影
	\[
	\pi_{ji}:a\pmod {p^j}\longmapsto a\pmod {p^i}.
	\]
	逆极限 $\varprojlim \ZZ/p^i\ZZ$ 称为 $p$-进整数环，记为 $\ZZ_p$。
	\begin{enumerate}[label=(\alph*)]
		\item 证明 $\ZZ_p$ 的每个元素都可唯一写成无穷形式和
		\[
		b_0+b_1p+b_2p^2+b_3p^3+\cdots,
		\qquad b_i\in\{0,1,\ldots,p-1\}.
		\]
		描述这些形式和在 $\ZZ_p$ 中加法与乘法所对应的规则。\emph{提示：}把每个 $\ZZ/p^i\ZZ$ 中的最小非负剩余写成 $p$ 进展开，再描述映射 $\pi_{ji}$。特别注意 $\ZZ_p$ 是不可数的。
		\item 证明 $\ZZ_p$ 是含有一份整数环的整环。
		\item 证明 (a) 中的 $b_0+b_1p+b_2p^2+b_3p^3+\cdots$ 是 $\ZZ_p$ 的单位，当且仅当 $b_0\neq0$。
		\item 证明 $p\ZZ_p$ 是 $\ZZ_p$ 的唯一极大理想，且 $\ZZ_p/p\ZZ_p\cong\ZZ/p\ZZ$（其中 $p=0+1p+0p^2+0p^3+
		\cdots$）。证明 $\ZZ_p$ 的每个理想都形如 $p^n\ZZ_p$，其中 $n\geq0$。
		\item 证明若 $a_1\not\equiv0\pmod p$，则直极限 $\ZZ_p$ 中存在元素 $a=(a_i)$，满足 $a_1^{p-1}\equiv1\pmod p$，并且 $\pi_{j1}(a_j)=a_1$ 对所有 $j$ 成立。推出 $\ZZ_p$ 含有 $p-1$ 个不同的 $(p-1)$ 次单位根。
	\end{enumerate}
\end{enumerate}

% \chapter{Euclidean整环、主理想整环与唯一分解整环}

有许多类环比一般的环具有更多代数结构。本章考察三类这样的环：具有除法算法的环（Euclidean整环）、每个理想都是主理想的环（主理想整环），以及元素能分解为素元的环（唯一分解整环）。这类环的主要例子是整数环 $\ZZ$ 和系数在域 $F$ 中的多项式环 $F[x]$。在这里，我们把绪论中关于整数 $\ZZ$ 陈述过的全部定理证明为更一般环论结果的特殊情形。这些结果将在下一章应用到特殊情形 $F[x]$ 中。

本章中所有的环都是交换环。

\section{Euclidean整环}

我们首先定义整环 $R$ 上范数的概念。这里的范数本质上只是 $R$ 中元素“大小”的一种度量。

\noindent\textbf{定义.}
设 $R$ 是整环。任意满足 $N(0)=0$ 的函数
\[
N:R\longrightarrow \ZZ^+\cup\{0\}
\]
称为 $R$ 上的一个\emph{范数}。若对所有 $a\neq0$ 都有 $N(a)>0$，则称 $N$ 为\emph{正范数}。

注意，这里的“范数”概念相当弱；同一个整环 $R$ 可能有若干不同的范数。

\noindent\textbf{定义.}
整环 $R$ 称为\emph{Euclidean整环}（或称它具有\emph{除法算法}），若存在 $R$ 上的一个范数 $N$，使得对任意 $a,b\in R$，$b\neq0$，存在 $q,r\in R$ 满足
\[
a=qb+r,
\qquad r=0 \text{ 或 } N(r)<N(b).
\]
元素 $q$ 称为商，$r$ 称为余数。

除法算法的重要性在于：它允许我们对 $R$ 中两个元素 $a,b$ 进行 Euclidean 算法。通过连续“除法”（这些除法实际上发生在 $R$ 的分式域中），可以写成
\begin{align*}
	a&=q_0b+r_0,\tag{0}\\
	b&=q_1r_0+r_1,\tag{1}\\
	r_0&=q_2r_1+r_2,\tag{2}\\
	&\ \vdots\\
	r_{n-2}&=q_nr_{n-1}+r_n,\tag{$n$}\\
	r_{n-1}&=q_{n+1}r_n.\tag{$n+1$}
\end{align*}
其中 $r_n$ 是最后一个非零余数。这样的 $r_n$ 必然存在，因为只要余数非零，就有
\[
N(b)>N(r_0)>N(r_1)>\cdots>N(r_n),
\]
这是一列递减的非负整数，不可能无限继续。注意这些元素一般并不唯一。

\paragraph{例}
\begin{enumerate}[label=(\arabic*)]
	\setcounter{enumi}{-1}
	\item 域是 Euclidean整环的平凡例子：任意范数都满足条件，例如 $N(a)=0$ 对所有 $a$ 成立。因为对 $b\neq0$，总有 $a=qb+0$，其中 $q=ab^{-1}$。
	
	\item 整数环 $\ZZ$ 是 Euclidean整环，范数为 $N(a)=|a|$。这就是初等算术中的长除法。若 $a,b$ 是非零整数且先设 $b>0$，则半开区间 $[nb,(n+1)b)$，$n\in\ZZ$，划分实数轴，故 $a$ 落在某个区间中，设 $a\in[kb,(k+1)b)$。令 $q=k$，则 $r=a-qb$ 满足 $0\leq r<|b|$。若 $b<0$，对 $-b>0$ 应用刚才的论证即可。若 $a$ 不是 $b$ 的倍数，通常有两个可能的商与余数。例如
	\[
	5=2\cdot2+1=3\cdot2-1.
	\]
	若要求余数非负，则商和余数唯一。
	
	\item 若 $F$ 是域，则多项式环 $F[x]$ 是 Euclidean整环，范数为多项式的次数。其除法算法就是多项式长除法；证明与整数情形类似，将在下一章给出。这里系数必须来自域，因为除法算法最终依赖于能够除以任意非零系数。我们将在第 2 节证明，若 $R$ 不是域，则 $R[x]$ 不是 Euclidean整环。
	
	\item 第 7.1 节中的二次整数环 $\mcO$ 都是整环，并带有由域范数绝对值得到的范数；但一般地，$\mcO$ 并不关于该范数（也不关于任何范数）成为 Euclidean整环。不过 Gaussian 整数环 $\ZZ[i]$ 是 Euclidean整环，其范数为
	\[
	N(a+bi)=a^2+b^2.
	\]
	设 $\alpha=a+bi$，$\beta=c+di\neq0$ 是 $\ZZ[i]$ 中元素。在域 $\QQ(i)$ 中
	\[
	{\alpha\over\beta}=r+si,
	\qquad
	r={ac+bd\over c^2+d^2},\quad s={bc-ad\over c^2+d^2}.
	\]
	取整数 $p,q$ 分别最接近有理数 $r,s$，则 $|r-p|\leq1/2$ 且 $|s-q|\leq1/2$。令
	\[
	\theta=(r-p)+(s-q)i,
	\qquad \gamma=\beta\theta.
	\]
	于是 $\gamma=\alpha-(p+qi)\beta\in\ZZ[i]$，且
	\[
	N(\gamma)=N(\theta)N(\beta)\leq {1\over2}N(\beta)<N(\beta).
	\]
	所以
	\[
	\alpha=(p+qi)\beta+\gamma
	\]
	给出了除法算法。该算法很显式；商 $p+qi$ 可由 $r,s$ 很快确定。不过商不一定唯一。
	
	同样的证明也可用于证明当 $D=-2,-3,-7,-11$ 时，相应的二次整数环关于域范数是 Euclidean整环。稍后我们将看到 $\ZZ[\sqrt{-5}]$ 不是 Euclidean整环；而 $\ZZ[(1+\sqrt{-19})/2]$ 对任何范数都不是 Euclidean整环的证明在本节末给出。
	
	\item 离散赋值环也是 Euclidean整环。设 $K$ 是域，$v:K^\times\to\ZZ$ 是离散赋值，即满足
	\[
	v(ab)=v(a)+v(b),\qquad v \text{ 满射},\qquad
	v(x+y)\geq \min\{v(x),v(y)\}
	\]
	（当 $x+y\neq0$ 时）。集合
	\[
	\{x\in K^\times\mid v(x)\geq0\}\cup\{0\}
	\]
	是 $K$ 的子环，称为赋值环。若整环 $R$ 是某个离散赋值在其分式域上的赋值环，则称 $R$ 是离散赋值环。例如，对固定素数 $p$，所有分母与 $p$ 互素的有理数组成 $\QQ$ 中的离散赋值环。
	
	在离散赋值环中，令 $N(0)=0$，并在非零元素上令 $N=v$，则 $N$ 给出 Euclidean 算法：若 $N(a)<N(b)$，则 $a=0\cdot b+a$；若 $N(a)\geq N(b)$，则 $q=ab^{-1}\in R$，所以 $a=qb+0$。
\end{enumerate}

\medskip
\noindent\textbf{命题 1.}
Euclidean整环中的每个理想都是主理想。更确切地说，若 $I$ 是 Euclidean整环 $R$ 中的任一非零理想，则
\[
I=(d),
\]
其中 $d$ 是 $I$ 中任一具有最小范数的非零元素。

\begin{proof}
	若 $I=(0)$，无须证明。否则取 $I$ 中非零且范数最小的元素 $d$，其存在性来自 $\ZZ$ 的良序性。显然 $(d)\subseteq I$。反过来，取 $a\in I$，由除法算法可写
	\[
	a=qd+r,
	\qquad r=0\text{ 或 }N(r)<N(d).
	\]
	由于 $r=a-qd\in I$，$d$ 的范数最小性迫使 $r=0$。于是 $a=qd\in(d)$。故 $I=(d)$。
\end{proof}

命题 1 可用于证明某些整环不可能是 Euclidean整环：只要找出一个非主理想即可。

\paragraph{例}
\begin{enumerate}[label=(\arabic*)]
	\item 设 $R=\ZZ[x]$。理想 $(2,x)$ 不是主理想，所以 $\ZZ[x]$ 不是 Euclidean整环；尽管 $\QQ[x]$ 是 Euclidean整环。
	
	\item 设 $R=\ZZ[\sqrt{-5}]$，其域范数为 $N(a+b\sqrt{-5})=a^2+5b^2$。考虑理想
	\[
	I=(3,2+\sqrt{-5}).
	\]
	若 $I=(a+b\sqrt{-5})$ 为主理想，则 $3$ 与 $2+\sqrt{-5}$ 都被 $a+b\sqrt{-5}$ 整除。取范数得
	\[
	9=N(\alpha)(a^2+5b^2),
	\]
	从而 $a^2+5b^2$ 只能是 $1,3,9$。值为 $3$ 不可能；值为 $9$ 会推出生成元为 $\pm3$，与 $2+\sqrt{-5}$ 不被 $3$ 整除矛盾；值为 $1$ 会推出 $I=R$，从而 $1=3\gamma+(2+\sqrt{-5})\delta$，再乘以 $2-\sqrt{-5}$ 可推出 $2-\sqrt{-5}$ 被 $3$ 整除，也矛盾。因此 $I$ 不是主理想，$R$ 不是 Euclidean整环。
\end{enumerate}

Euclidean算法的一个基本后果是产生最大公因子。我们先在一般交换环中精确定义最大公因子。

\noindent\textbf{定义.}
设 $R$ 是交换环，$a,b\in R$ 且 $b\neq0$。
\begin{enumerate}[label=(\arabic*)]
	\item 若存在 $x\in R$ 使 $a=bx$，则称 $a$ 是 $b$ 的倍数，或称 $b$ 整除 $a$，记作 $b\mid a$。
	\item $a$ 与 $b$ 的一个\emph{最大公因子}是一个非零元素 $d$，满足：
	\begin{enumerate}[label=(\roman*)]
		\item $d\mid a$ 且 $d\mid b$；
		\item 若 $d'\mid a$ 且 $d'\mid b$，则 $d'\mid d$。
	\end{enumerate}
	最大公因子记为 $\gcd(a,b)$，或简单记作 $(a,b)$。
\end{enumerate}

注意 $b\mid a$ 当且仅当 $a\in(b)$，当且仅当 $(a)\subseteq(b)$。若 $d$ 是 $a,b$ 的公共因子，则 $(d)$ 包含 $a,b$，从而包含由 $a,b$ 生成的理想。因此，最大公因子（若存在）可以看作包含 $a,b$ 所生成理想的唯一最小主理想的一个生成元。

\noindent\textbf{命题 2.}
若 $a,b$ 是交换环 $R$ 中非零元素，且由 $a,b$ 生成的理想是主理想 $(d)$，则 $d$ 是 $a$ 与 $b$ 的最大公因子。

这解释了为什么符号 $(a,b)$ 常同时表示由 $a,b$ 生成的理想以及 $a,b$ 的最大公因子。一个整环若每个二元生成理想 $(a,b)$ 都是主理想，称为 Bezout 整环。

\noindent\textbf{命题 3.}
设 $R$ 是整环。若两个元素 $d,d'$ 生成同一个主理想，即 $(d)=(d')$，则 $d'=ud$，其中 $u$ 是 $R$ 的单位。特别地，若 $d,d'$ 都是 $a,b$ 的最大公因子，则它们相差一个单位因子。

\begin{proof}
	可假设 $d,d'$ 非零。由 $d\in(d')$ 与 $d'\in(d)$，存在 $x,y\in R$ 使
	\[
	d=xd',\qquad d'=yd.
	\]
	故 $d=xyd$，因 $R$ 是整环且 $d\neq0$，得 $xy=1$。因此 $x,y$ 都是单位。第二个断言由最大公因子的相互整除立刻推出。
\end{proof}

\noindent\textbf{定理 4.}
设 $R$ 是 Euclidean整环，$a,b$ 是 $R$ 中非零元素。令 $d=r_n$ 为对 $a,b$ 进行 Euclidean 算法时最后一个非零余数。则
\begin{enumerate}[label=(\arabic*)]
	\item $d$ 是 $a$ 与 $b$ 的最大公因子；
	\item 主理想 $(d)$ 就是由 $a,b$ 生成的理想。特别地，存在 $x,y\in R$ 使
	\[
	d=ax+by.
	\]
\end{enumerate}

\begin{proof}
	由命题 1，$(a,b)$ 是主理想。只需证明 $r_n$ 生成此理想。先从 Euclidean 算法最后一个方程向前追踪整除关系，可得 $r_n\mid r_{n-1}$，继而归纳得到 $r_n\mid b$ 与 $r_n\mid a$。另一方面，从第一个方程开始向后归纳：$r_0=a-q_0b\in(a,b)$，$r_1=b-q_1r_0\in(a,b)$，若 $r_{k-1},r_k\in(a,b)$，则
	\[
	r_{k+1}=r_{k-1}-q_{k+1}r_k\in(a,b).
	\]
	故 $r_n\in(a,b)$。于是 $(a,b)=(r_n)$，结论成立。
\end{proof}

例如，令 $a=2210$，$b=1131$。Euclidean 算法为
\begin{align*}
	2210&=1\cdot1131+1079,\\
	1131&=1\cdot1079+52,\\
	1079&=20\cdot52+39,\\
	52&=1\cdot39+13,\\
	39&=3\cdot13.
\end{align*}
所以 $\gcd(2210,1131)=13$。回代可得
\[
13=(-22)\cdot2210+43\cdot1131.
\]

整数 Euclidean 算法非常快：计算两个整数 $a,b$ 的最大公因子的步数至多为较小数位数的 $5$ 倍。换言之，该算法关于整数大小是对数级的。

由 $\gcd(a,b)=d$ 还可得到一次不定方程
\[
ax+by=N
\]
有整数解当且仅当 $d\mid N$。若 $(x_0,y_0)$ 是一个解，则全体解为
\[
x=x_0+m{b\over d},
\qquad
y=y_0-m{a\over d},
\qquad m\in\ZZ.
\]

本节最后给出另一个有时能用来证明某整环不是 Euclidean整环的判别准则。\footnote{这里以及下一节部分材料采用 J.C. Wilson, \emph{A principal ideal ring that is not a Euclidean ring}, Math. Mag. 46 (1973), pp. 34--38 对 Th. Motzkin 思想的阐述，并使用 Kenneth S. Williams, \emph{Note on non-Euclidean Principal Ideal Domains}, Math. Mag. 48 (1975), pp. 176--177 中的一个简化。}对任意整环，记
\[
\overline R=R^\times\cup\{0\}
\]
若 $u\in R-\overline R$，且对每个 $x\in R$，存在 $z\in\overline R$ 使 $u\mid x-z$，则称 $u$ 为\emph{万能边除子}。也就是说，对 $u$ 有一种弱形式的除法算法：每个 $x$ 都可写成
\[
x=qu+z,
\]
其中 $z$ 是 $0$ 或单位。

\noindent\textbf{命题 5.}
设 $R$ 是不是域的整环。若 $R$ 是 Euclidean整环，则 $R$ 中存在万能边除子。

\begin{proof}
	设 $R$ 关于范数 $N$ 是 Euclidean整环。取 $R-\overline R$ 中范数最小的元素 $u$。对任意 $x\in R$，写
	\[
	x=qu+r,
	\]
	其中 $r=0$ 或 $N(r)<N(u)$。由 $u$ 的选择，余数 $r$ 必在 $\overline R$ 中。因此 $u$ 是万能边除子。
\end{proof}

\paragraph{例}
令
\[
R=\ZZ\left[{1+\sqrt{-19}\over2}\right].
\]
本节末的命题可用来证明它关于任何范数都不是 Euclidean整环。事实上，$R$ 中单位只有 $\pm1$，所以 $\overline R=\{0,\pm1\}$。设 $u$ 是万能边除子。对 $x=2$，$u$ 必整除 $2,1$ 或 $3$ 中的某个非零元素；因此 $u$ 只能是 $\pm2$ 或 $\pm3$。但取 $x=(1+\sqrt{-19})/2$，可以直接检查 $x,x\pm1$ 都不被 $\pm2$ 或 $\pm3$ 整除。因此不存在万能边除子，$R$ 不是 Euclidean整环。

\paragraph{练习}
\begin{enumerate}[label=\arabic*.]
	\item 对下列五组整数 $a$ 与 $b$，求它们的最大公因子 $d$，并把 $d$ 写成 $a$ 与 $b$ 的线性组合 $ax+by$。
	\begin{enumerate}[label=(\alph*)]
		\item $a=20,\ b=13$。
		\item $a=69,\ b=372$。
		\item $a=11391,\ b=5673$。
		\item $a=507885,\ b=60808$。
		\item $a=91442056588823,\ b=779086434385541$（Euclidean 算法对这两个整数只需 $7$ 步）。
	\end{enumerate}
	
	\item 对下列每一对整数 $a$ 与 $n$，证明 $a$ 与 $n$ 互素，并求 $a$ 模 $n$ 的逆元（参见绪论章第 3 节）。
	\begin{enumerate}[label=(\alph*)]
		\item $a=13,\ n=20$。
		\item $a=69,\ n=89$。
		\item $a=1891,\ n=3797$。
		\item $a=6003722857,\ n=77695236973$（Euclidean 算法对这两个整数只需 $3$ 步）。
	\end{enumerate}
	
	\item 设 $R$ 是 Euclidean整环。设 $m$ 是所有非零元素范数的集合中的最小整数。证明 $R$ 中每个范数为 $m$ 的非零元素都是单位。推出：若存在范数为零的非零元素，则它必为单位。
	
	\item 设 $R$ 是 Euclidean整环。
	\begin{enumerate}[label=(\alph*)]
		\item 证明：若 $(a,b)=1$ 且 $a\mid bc$，则 $a\mid c$。更一般地，证明：若 $a\mid bc$ 且 $a,b$ 非零，则
		\[
		{a\over (a,b)}\mid c.
		\]
		\item 考虑 Diophantine 方程
		\[
		ax+by=N,
		\]
		其中 $a,b,N$ 为整数且 $a,b$ 非零。设 $x_0,y_0$ 是一个解：
		\[
		ax_0+by_0=N.
		\]
		证明此方程的全体解为
		\[
		x=x_0+m\,{b\over (a,b)},\qquad
		y=y_0-m\,{a\over (a,b)},
		\]
		其中 $m$ 遍历所有整数。\emph{提示：}若 $x,y$ 是 $ax+by=N$ 的解，证明
		\[
		a(x-x_0)=b(y_0-y)
		\]
		并使用 (a)。
	\end{enumerate}
	
	\item 求下列方程的全部整数解：
	\begin{enumerate}[label=(\alph*)]
		\item $2x+4y=5$。
		\item $17x+29y=31$。
		\item $85x+145y=505$。
	\end{enumerate}
	
	\item （邮票问题）设 $a,b$ 是两个互素正整数。证明：任意充分大的正整数 $N$ 都可以写成 $a$ 与 $b$ 的非负整数线性组合 $ax+by$；也就是说，存在整数 $N_0$，使得对所有 $N\geq N_0$，方程
	\[
	ax+by=N
	\]
	都有非负整数解 $x,y$。事实上，证明整数 $ab-a-b$ 不能写成 $a$ 与 $b$ 的正线性组合，但每个大于 $ab-a-b$ 的整数都是 $a$ 与 $b$ 的正线性组合（因此每个大于 $ab-a-b$ 的“邮资”都可以只用面值为 $a$ 与 $b$ 的邮票得到）。
	
	\item 用 Euclidean 算法求 $\ZZ[i]$ 中理想 $(85,1+13i)$ 的生成元，即求 $85$ 与 $1+13i$ 的最大公因子。对理想 $(47-13i,53+56i)$ 做同样的计算。
	
	\noindent 已知（但证明并不容易）$D=-1,-2,-3,-7,-11,-19,-43,-67,-163$ 是仅有的负整数 $D$，使得 $\mcO$ 中每个理想都是主理想（即按下一节术语，$\mcO$ 是 P.I.D.）。下一题的结果精确判定了哪些 $D<0$ 的二次整数环是 Euclidean整环。
	
	\item 设 $F=\QQ(\sqrt D)$ 是二次域，$\mcO$ 是其对应的二次整数环，$N$ 是第 7.1 节中的域范数。
	\begin{enumerate}[label=(\alph*)]
		\item 设 $D=-1,-2,-3,-7$ 或 $-11$。证明 $\mcO$ 关于 $N$ 是 Euclidean整环。\emph{提示：}修改正文中对 $\ZZ[i]$（$D=-1$）的证明。对 $D=-3,-7,-11$，证明 $F$ 中每个元素都与 $\mcO$ 中某个元素相差一个范数至多为
		\[
		{(1+|D|)^2\over 16|D|}
		\]
		的元素，而这个数对这些 $D$ 都小于 $1$。把 $\mcO$ 中的点画在 $\CC$ 中可能有帮助。
		\item 设 $D=-43,-67$ 或 $-163$。证明 $\mcO$ 关于任何范数都不是 Euclidean整环。\emph{提示：}套用正文中 $D=-19$ 的证明。
	\end{enumerate}
	
	\item 证明二次整数域 $\QQ(\sqrt2)$ 的整数环 $\mcO$ 关于第 7.1 节中域范数绝对值给出的范数是 Euclidean整环。
	
	\item 证明：对 $\ZZ[i]$ 的任意非零理想 $I$，商环 $\ZZ[i]/I$ 是有限的。\emph{提示：}利用 $I=(\alpha)$（其中 $\alpha\neq0$）以及此 Euclidean整环中的除法算法，说明 $I$ 的每个陪集都可由范数小于 $N(\alpha)$ 的元素代表。
	
	\item 设 $R$ 是含 $1$ 的交换环，$a,b$ 是 $R$ 中非零元素。$a$ 与 $b$ 的一个最小公倍数是指 $R$ 中的一个元素 $e$，满足：
	\begin{enumerate}[label=(\roman*)]
		\item $a\mid e$ 且 $b\mid e$；
		\item 若 $a\mid e'$ 且 $b\mid e'$，则 $e\mid e'$。
	\end{enumerate}
	\begin{enumerate}[label=(\alph*)]
		\item 证明：$a$ 与 $b$ 的最小公倍数（若存在）是包含在 $(a)\cap(b)$ 中的唯一最大主理想的一个生成元。
		\item 推出：Euclidean整环中任意两个非零元素都有最小公倍数，并且在乘以单位的意义下唯一。
		\item 证明：在 Euclidean整环中，$a$ 与 $b$ 的最小公倍数为
		\[
		{ab\over (a,b)},
		\]
		其中 $(a,b)$ 是 $a$ 与 $b$ 的最大公因子。
	\end{enumerate}
	
	\item （一种公钥密码）设 $N$ 为正整数。设 $M$ 是与 $N$ 互素的整数，$d$ 是与 $\varphi(N)$ 互素的整数，其中 $\varphi$ 表示 Euler 的 $\varphi$ 函数。证明：若
	\[
	M_1\equiv M^d\pmod N,
	\]
	则
	\[
	M\equiv M_1^{d'}\pmod N,
	\]
	其中 $d'$ 是 $d$ 模 $\varphi(N)$ 的逆元：
	\[
	dd'\equiv1\pmod{\varphi(N)}.
	\]
	
	\noindent\emph{说明：}这个结果是标准公钥密码的基础。设 $N=pq$ 是两个不同大素数的乘积（例如每个约有 $100$ 位数字）。若 $M$ 是一条消息，则 $M_1\equiv M^d\pmod N$ 是 $M$ 的加密版本，而通过计算 $M_1^{d'}\pmod N$ 可以解密（即使 $M,N$ 很大，这些幂也可通过连续平方相当容易地计算）。$N$ 与 $d$ 的值（但不是 $p,q$）公开，因此任何拥有消息 $M$ 的人都可以发送其编码消息 $M^d\pmod N$。要解码，似乎必须确定 $d'$，这又需要确定
	\[
	\varphi(N)=\varphi(pq)=(p-1)(q-1)
	\]
	（不过尚未有人证明不存在其他解码方案）。这种方法作为密码的成功依赖于分解 $N$ 为素数乘积的必要性，而目前没有足够有效的算法可完成这一点（例如最朴素的检查所有小于等于 $\sqrt N$ 的因子的办法，在这里需要约 $10^{100}$ 次计算；即使每秒 $100$ 亿次计算，也大约需要 $300$ 年，而且当然还可以继续增大 $p,q$ 的大小）。
\end{enumerate}


\section{主理想整环（P.I.D.s）}

\noindent\textbf{定义.}
\emph{主理想整环}（Principal Ideal Domain，P.I.D.）是每个理想都是主理想的整环。

命题 1 已经证明每个 Euclidean整环都是主理想整环。因此，所有关于主理想整环的结果自动适用于 Euclidean整环。

\paragraph{例}
\begin{enumerate}[label=(\arabic*)]
	\item 如命题 1 后面所说，整数环 $\ZZ$ 是 P.I.D.。我们在第 7.4 节看到，多项式环 $\ZZ[x]$ 含有非主理想，因此不是 P.I.D.。
	\item 命题 1 后的例 2 说明，二次整数环 $\ZZ[\sqrt{-5}]$ 不是 P.I.D.；事实上，理想 $(3,1+\sqrt{-5})$ 是非主理想。两个非主理想的乘积可以是主理想，例如
	\[
	(3,1+\sqrt{-5})(3,1-\sqrt{-5})=(3)
	\]
	（参见练习 5 以及下面命题 12 之前的例子）。
\end{enumerate}

并非每个主理想整环都是 Euclidean整环。下面将证明，上一节已经证明不是 Euclidean整环的二次整数环
\[
\ZZ\left[{1+\sqrt{-19}\over2}\right]
\]
仍然是一个 P.I.D.。

从理想论的观点看，主理想整环是在域（其理想只有平凡的 $(0)$ 和 $(1)$）之外自然要研究的一类环。Euclidean整环所享有的许多性质在主理想整环中也成立。不过 Euclidean整环相对于主理想整环的一个显著优点是：虽然在两种情形下最大公因子都存在，但在 Euclidean整环中有一个算法来计算它们。因此，依赖于最大公因子存在性的结果往往可以在更大的主理想整环类中证明；然而具体例子的计算（即这些结果的具体应用）在有 Euclidean算法可用时通常更有效。

我们收集上一节中已经证明的关于最大公因子的若干事实。

\medskip
\noindent\textbf{命题 6.}
设 $R$ 是主理想整环，$a,b$ 是 $R$ 中非零元素。令 $d$ 为由 $a,b$ 生成的主理想的一个生成元。则
\begin{enumerate}[label=(\arabic*)]
	\item $d$ 是 $a,b$ 的最大公因子；
	\item $d$ 可写成 $a,b$ 的 $R$-线性组合，也就是说，存在 $x,y\in R$ 使
	\[
	d=ax+by;
	\]
	\item $d$ 在乘以 $R$ 中单位的意义下唯一。
\end{enumerate}

\begin{proof}
	这只是命题 2 与命题 3。
\end{proof}

回忆，极大理想总是素理想，但反过来一般不真。我们在第 7.4 节观察到，$\ZZ$ 的每个非零素理想都是极大理想。如下命题说明，这个有用事实在任意主理想整环中成立。

\medskip
\noindent\textbf{命题 7.}
主理想整环中的每个非零素理想都是极大理想。

\begin{proof}
	设 $(p)$ 是主理想整环 $R$ 中的非零素理想，并设 $I=(m)$ 是包含 $(p)$ 的任一理想。我们必须证明 $I=(p)$ 或 $I=R$。由于 $p\in(m)$，有 $p=rm$ 对某个 $r\in R$ 成立。因为 $(p)$ 是素理想且 $rm\in(p)$，所以 $r$ 或 $m$ 必在 $(p)$ 中。若 $m\in(p)$，则 $(p)=(m)=I$。另一方面，若 $r\in(p)$，写 $r=ps$。此时
	\[
	p=rm=psm,
	\]
	所以 $sm=1$（回忆 $R$ 是整环），故 $m$ 是单位，于是 $I=R$。
\end{proof}

如前所述，若 $F$ 是域，则多项式环 $F[x]$ 是 Euclidean整环，从而也是主理想整环（这将在下一章证明）。其逆命题也成立。直观地说，若 $I$ 是 $R$ 中的一个理想（例如 $\ZZ$ 中的理想 $(2)$），则 $R[x]$ 中的理想 $(I,x)$（例如 $\ZZ[x]$ 中的理想 $(2,x)$）通常比 $I$ 本身多需要一个生成元，因此一般不是主理想。

\medskip
\noindent\textbf{推论 8.}
若 $R$ 是任意交换环，并且多项式环 $R[x]$ 是主理想整环（或 Euclidean整环），则 $R$ 必为域。

\begin{proof}
	假设 $R[x]$ 是主理想整环。因为 $R$ 是 $R[x]$ 的子环，所以 $R$ 必为整环（回忆 $R[x]$ 有单位当且仅当 $R$ 有单位）。理想 $(x)$ 是 $R[x]$ 中的非零素理想，因为 $R[x]/(x)$ 同构于整环 $R$。由命题 7，$(x)$ 是极大理想，于是商环 $R$ 由第 7.4 节命题 12 是域。
\end{proof}

本节最后一个结果将用来证明并非每个 P.I.D. 都是 Euclidean整环，并把主理想性质同 Euclidean 条件的另一个弱化形式联系起来。

\medskip
\noindent\textbf{定义.}
若 $N$ 是正范数，并且对每个非零 $a,b\in R$，或者 $a$ 是理想 $(b)$ 的元素，或者理想 $(a,b)$ 中有一个范数严格小于 $b$ 的范数的非零元素，则称 $N$ 是一个 \emph{Dedekind--Hasse范数}。也就是说，或者 $b$ 在 $R$ 中整除 $a$，或者存在 $s,t\in R$ 使
\[
0<N(sa-tb)<N(b).
\]

注意，如果对某个正范数 $N$，$R$ 是 Euclidean的，则总能用 $s=1$ 满足 Dedekind--Hasse 条件，所以这确实是 Euclidean 条件的一个弱化。

\medskip
\noindent\textbf{命题 9.}
整环 $R$ 是 P.I.D. 当且仅当 $R$ 有一个 Dedekind--Hasse范数。\footnote{Dedekind--Hasse 范数在 $R$ 上推出 $R$ 是 P.I.D.（当 $R$ 是代数整数环时还等价）是 Dedekind 与 Hasse 的经典判别准则，见 \emph{\"Uber eindeutige Zerlegung in Primelemente oder in Primhauptideale in Integrit\"atsbereichen}, Jour. f\"ur die Reine und Angew. Math. 159 (1928), pp. 3--12。其逆命题一般成立这一观察较新，见 John Greene, \emph{Principal Ideal Domains are almost Euclidean}, Amer. Math. Monthly 104 (1997), pp. 154--156。}

\begin{proof}
	令 $I$ 是 $R$ 的任一非零理想，取 $I$ 中满足 $N(b)$ 最小的非零元素 $b$。设 $a$ 是 $I$ 中任一非零元素，则理想 $(a,b)$ 包含于 $I$。于是 Dedekind--Hasse 条件以及 $b$ 的最小性说明 $a\in(b)$，所以 $I=(b)$ 是主理想。反方向将在下一节证明（推论 16）。
\end{proof}

\paragraph{例}
令
\[
R=\ZZ\left[{1+\sqrt{-19}\over2}\right]
\]
为上一节末所考虑的二次整数环。我们说明定义在 $R$ 上的正的域范数
\[
N\left(a+b{1+\sqrt{-19}\over2}\right)=a^2+ab+5b^2
\]
是 Dedekind--Hasse范数；于是由命题 9 与上一节结果可知，$R$ 是 P.I.D. 但不是 Euclidean整环。

设 $\alpha,\beta$ 是 $R$ 的非零元素且 $\alpha/\beta\notin R$。我们必须证明存在 $s,t\in R$ 使
\[
0<N(s\alpha-t\beta)<N(\beta),
\]
由域范数的乘法性，这等价于
\[
0<N\left({\alpha\over\beta}s-t\right)<1. \tag{*}
\]
写
\[
{\alpha\over\beta}={a+b\sqrt{-19}\over c}\in \QQ[\sqrt{-19}],
\]
其中整数 $a,b,c$ 没有公因子且 $c>1$（因为假设 $\beta$ 不整除 $\alpha$）。既然 $a,b,c$ 没有公因子，存在整数 $x,y,z$ 使 $ax+by+cz=1$。把
\[
ay-19bx=cq+r
\]
写成商为 $q$、余数为 $r$ 且 $|r|\leq c/2$ 的形式，并令
\[
s=y+x\sqrt{-19},\qquad t=q-z\sqrt{-19}.
\]
则直接计算给出
\[
0<N\left({\alpha\over\beta}s-t\right)
={ (ay-19bx-cq)^2+19(ax+by+cz)^2\over c^2}
\leq {1\over4}+{19\over c^2},
\]
因此当 $c\geq5$ 时，上述 $s,t$ 满足 $(*)$。

若 $c=2$，则 $a,b$ 中一个为偶数而另一个为奇数（否则 $\alpha/\beta\in R$），这时直接检查可知
\[
s=1,
\qquad
 t={ (a-1)+b\sqrt{-19}\over2}
\]
是 $R$ 的元素并满足 $(*)$。

若 $c=3$，则整数 $a^2+19b^2$ 不能被 $3$ 整除（模 $3$ 后这是 $a^2+b^2$，容易看出它模 $3$ 为 $0$ 当且仅当 $a,b$ 都模 $3$ 为 $0$；但那样 $a,b,c$ 有公因子）。写
\[
a^2+19b^2=3q+r,
\qquad r=1\text{ 或 }2.
\]
于是再次直接检查可知
\[
s=a-b\sqrt{-19},
\qquad t=q
\]
是 $R$ 的元素并满足 $(*)$。

最后，若 $c=4$，则 $a,b$ 不全为偶数。若 $a,b$ 中一个为偶数而另一个为奇数，则 $a^2+19b^2$ 是奇数，所以可写
\[
a^2+19b^2=4q+r,
\qquad 0<r<4,
\]
其中 $q,r\in\ZZ$。这时 $s=a-b\sqrt{-19}$ 与 $t=q$ 满足 $(*)$。若 $a,b$ 都为奇数，则
\[
a^2+19b^2\equiv 1+3 \pmod 8,
\]
所以可写 $a^2+19b^2=8q+4$，其中 $q\in\ZZ$。这时
\[
s={a-b\sqrt{-19}\over2},
\qquad t=q
\]
是 $R$ 中元素并满足 $(*)$。

\paragraph{练习}
\begin{enumerate}[label=\arabic*.]
	\item 证明：在主理想整环中，两个理想 $(a)$ 与 $(b)$ 互素（参见第 7.6 节）当且仅当 $a$ 与 $b$ 的一个最大公因子是 $1$（此时称 $a$ 与 $b$ \emph{互素}或\emph{相对素}）。
	
	\item 证明：P.I.D. 中任意两个非零元素都有最小公倍数（参见第 1 节练习 11）。
	
	\item 证明：P.I.D. 对素理想取商仍是 P.I.D.。
	
	\item 设 $R$ 是整环。证明：若下列两个条件成立，则 $R$ 是主理想整环：
	\begin{enumerate}[label=(\roman*)]
		\item $R$ 中任意两个非零元素 $a,b$ 都有最大公因子，且该最大公因子可写成 $ra+sb$ 的形式，其中 $r,s\in R$；
		\item 若 $a_1,a_2,a_3,\ldots$ 是 $R$ 中非零元素，且对所有 $i$ 都有 $a_{i+1}\mid a_i$，则存在正整数 $N$，使得对所有 $n\geq N$，$a_n$ 都是 $a_N$ 乘以某个单位。
	\end{enumerate}
	
	\item 令 $R$ 为二次整数环 $\ZZ[\sqrt{-5}]$。定义理想
	\[
	I_2=(2,1+\sqrt{-5}),\qquad
	I_3=(3,2+\sqrt{-5}),\qquad
	I_3'=(3,2-\sqrt{-5}).
	\]
	\begin{enumerate}[label=(\alph*)]
		\item 证明 $I_2,I_3,I_3'$ 都是 $R$ 中的非主理想。\emph{提示：}命题 1 后面的例 2 已对 $I_3$ 给出证明。
		\item 通过证明 $I_2^2$ 是由 $2$ 生成的主理想，即
		\[
		I_2^2=(2),
		\]
		说明两个非主理想的乘积可以是主理想。
		\item 类似地证明
		\[
		I_2I_3=(1-\sqrt{-5}),\qquad I_2I_3'=(1+\sqrt{-5})
		\]
		都是主理想。推出主理想 $(6)$ 是四个理想的乘积：
		\[
		(6)=I_2^2I_3I_3'.
		\]
	\end{enumerate}
	
	\item 设 $R$ 是整环，并假设 $R$ 中每个素理想都是主理想。本题证明 $R$ 中每个理想都是主理想，即 $R$ 是 P.I.D.。
	\begin{enumerate}[label=(\alph*)]
		\item 假设 $R$ 中非主理想的集合非空，证明这个集合在包含关系下有一个极大元（按假设它不是素理想）。\emph{提示：}使用 Zorn 引理。
		\item 令 $I$ 是一个关于“非主”性质极大的理想，且取 $a,b\in R$ 满足 $ab\in I$，但 $a\notin I$ 且 $b\notin I$。令 $I_a=(I,a)$ 为由 $I$ 与 $a$ 生成的理想，令 $I_b=(I,b)$ 为由 $I$ 与 $b$ 生成的理想，并定义
		\[
		J=\{r\in R\mid rI_a\subseteq I\}.
		\]
		证明 $I_a=(\alpha)$ 与 $J=(\beta)$ 都是 $R$ 中的主理想，并且
		\[
		I\subseteq I_b\subseteq J,
		\qquad I_aJ=(\alpha\beta)\subseteq I.
		\]
		\item 若 $x\in I$，证明 $x=s\alpha$ 对某个 $s\in J$ 成立。推出 $I=I_aJ$ 是主理想，得到矛盾；从而推出 $R$ 是 P.I.D.。
	\end{enumerate}
	
	\item 一个整环 $R$ 若满足每个由两个元素生成的理想都是主理想（即对每个 $a,b\in R$，都有 $(a,b)=(d)$ 对某个 $d\in R$ 成立），则称为 \emph{Bezout整环}（也参见第 3 节练习 11）。
	\begin{enumerate}[label=(\alph*)]
		\item 证明：整环 $R$ 是 Bezout整环当且仅当 $R$ 中每一对元素 $a,b$ 都有一个最大公因子 $d$，并且 $d$ 可写成 $a,b$ 的 $R$-线性组合：
		\[
		d=ax+by
		\]
		其中 $x,y\in R$。
		\item 证明：Bezout整环的每个有限生成理想都是主理想。\emph{提示：}在第 9.2 节与第 9.3 节的练习中可以看到不是每个理想都主的 Bezout整环。
		\item 设 $F$ 是 Bezout整环 $R$ 的分式域。证明 $F$ 的每个元素都可以写成 $a/b$ 的形式，其中 $a,b\in R$ 且 $a,b$ 互素（参见练习 1）。
	\end{enumerate}
	
	\item 证明：若 $R$ 是主理想整环，$D$ 是 $R$ 的一个乘法封闭子集，则 $D^{-1}R$ 也是 P.I.D.（参见第 7.5 节）。
\end{enumerate}


\section{唯一分解整环（U.F.D.s）}

在整数环 $\ZZ$ 的情形中，还有另一种从初等算术中熟悉的方法来确定两个元素 $a,b$ 的最大公因子，即把 $a,b$ “分解为素数”，再由此容易地确定最大公因子。这也可以推广到一类更大的环，称为唯一分解整环（U.F.D.）——我们马上给出定义。随后我们将证明
\[
\text{每个主理想整环都是唯一分解整环}
\]
从而关于唯一分解整环的每个结果都自动适用于 Euclidean整环和主理想整环。

我们先引入一些术语。

\medskip
\noindent\textbf{定义.}
设 $R$ 是整环。
\begin{enumerate}[label=(\arabic*)]
	\item 设 $r\in R$ 非零且不是单位。若每当 $r=ab$（其中 $a,b\in R$）时，$a,b$ 中至少一个必须是 $R$ 中的单位，则称 $r$ 在 $R$ 中\emph{不可约}。否则称 $r$ \emph{可约}。
	\item 非零元素 $p\in R$ 称为 $R$ 中的\emph{素元}，若由 $p$ 生成的理想 $(p)$ 是素理想。换言之，非零元素 $p$ 是素元，当且仅当 $p$ 不是单位，并且对任意 $a,b\in R$，若 $p\mid ab$，则 $p\mid a$ 或 $p\mid b$。
	\item 若 $R$ 中两个元素 $a,b$ 相差一个单位因子，即 $a=ub$ 对某个单位 $u\in R$ 成立，则称 $a,b$ 在 $R$ 中\emph{相伴}。
\end{enumerate}

\medskip
\noindent\textbf{命题 10.}
整环中的素元总是不可约元。

\begin{proof}
	设 $(p)$ 是非零素理想且 $p=ab$。于是 $ab=p\in(p)$，故由素理想的定义，$a$ 或 $b$ 中有一个（设为 $a$）属于 $(p)$。从而 $a=pr$ 对某个 $r\in R$ 成立。这推出
	\[
	p=ab=prb,
	\]
	所以 $rb=1$，故 $b$ 是单位。这说明 $p$ 不可约。
\end{proof}

一般地，不可约元不一定是素元。例如，考虑二次整数环 $R=\ZZ[\sqrt{-5}]$ 中的元素 $3$。第 1 节中的计算说明 $3$ 在 $R$ 中不可约，但 $3$ 不是素元，因为
\[
(2+\sqrt{-5})(2-\sqrt{-5})=3^2
\]
被 $3$ 整除，而 $2+\sqrt{-5}$ 与 $2-\sqrt{-5}$ 中没有一个在 $R$ 中被 $3$ 整除。

然而，若 $R$ 是主理想整环，则素元与不可约元两个概念相同。特别地，在 $\ZZ$ 与 $F[x]$（其中 $F$ 是域）中二者相同。

\medskip
\noindent\textbf{命题 11.}
在主理想整环中，非零元素是素元当且仅当它不可约。

\begin{proof}
	我们已经证明素元蕴含不可约。反过来，设 $p$ 不可约。我们必须证明 $p$ 是素元，即理想 $(p)$ 是素理想。若 $M$ 是任一包含 $(p)$ 的理想，则由假设 $M=(m)$ 是主理想。由于 $p\in(m)$，有 $p=rm$ 对某个 $r$ 成立。但 $p$ 不可约，所以按定义 $r$ 或 $m$ 是单位。这分别意味着 $(p)=(m)$ 或 $(m)=(1)$。因此包含 $(p)$ 的理想只有 $(p)$ 和 $(1)$，即 $(p)$ 是极大理想。极大理想都是素理想，证明完成。
\end{proof}

\paragraph{例}
命题 11 给出二次整数环 $\ZZ[\sqrt{-5}]$ 不是 P.I.D. 的另一个证明：因为 $3$ 在此环中不可约但不是素元。

整数环 $\ZZ$ 中的不可约元就是初等算术中熟悉的素数（以及它们的负数），而两个整数 $a,b$ 相伴当且仅当 $a=\pm b$。

在整数环 $\ZZ$ 中，任一整数 $n$ 都可写成素数（不一定互异）的乘积。理由如下。若 $n$ 本身不是素数，则由定义可以写成 $n=n_1n_2$，其中 $n_1,n_2$ 都不是单位，即都不是 $\pm1$。$n_1,n_2$ 的绝对值都小于 $n$ 本身。若二者都是素数，则已经把 $n$ 写成素数乘积。若 $n_1$ 或 $n_2$ 中有一个不是素数，则它又可分解成两个（更小的）整数。由于整数的绝对值不能无限下降，我们最终必然只剩下素数因子，于是得到 $n$ 的素数乘积分解。

例如，若 $n=2210$，上述算法如下进行：$n$ 不是素数，因为
\[
n=2\cdot1105.
\]
整数 $2$ 是素数，但 $1105$ 不是：
\[
1105=5\cdot221.
\]
整数 $5$ 是素数，但 $221$ 不是：
\[
221=13\cdot17.
\]
到此算法终止，因为 $13$ 与 $17$ 都是素数。这给出 $2210$ 的素因子分解
\[
2210=2\cdot5\cdot13\cdot17.
\]
类似地，我们得到
\[
1131=3\cdot13\cdot29.
\]
在这些例子中每个素数只出现一次，但一般当然不必如此。

在环 $\ZZ$ 中，不但每个整数 $n$ 都能写成素数乘积，而且这种分解事实上是唯一的：同一个正整数 $n$ 的任意两个素因子分解只差正素因子的书写顺序。这里限制为正整数，是为了避免把 $15$ 的分解 $(3)(5)$ 与 $(-3)(-5)$ 看作本质不同。$\ZZ$ 的这个唯一分解性质（我们很快将证明）对于整数算术非常有用。具有类似性质的一般环有一个名称。

\medskip
\noindent\textbf{定义.}
\emph{唯一分解整环}（Unique Factorization Domain，U.F.D.）是一个整环 $R$，其中每个非零非单位元素 $r\in R$ 都有如下两个性质：
\begin{enumerate}[label=(\roman*)]
	\item $r$ 可写成 $R$ 中有限个不可约元 $p_i$ 的乘积（不一定互异）：
	\[
	r=p_1p_2\cdots p_n;
	\]
	\item (i) 中的分解在相伴意义下唯一：即若
	\[
	r=q_1q_2\cdots q_m
	\]
	是 $r$ 的另一个不可约元分解，则 $m=n$，并且在重新编号之后，对 $i=1,2,\ldots,n$，$p_i$ 与 $q_i$ 相伴。
\end{enumerate}

\paragraph{例}
\begin{enumerate}[label=(\arabic*)]
	\item 域 $F$ 平凡地是唯一分解整环，因为每个非零元素都是单位，所以没有元素需要验证性质 (i) 和 (ii)。
	\item 如上所述，我们很快将证明每个主理想整环都是唯一分解整环（因此特别地，$\ZZ$ 与 $F[x]$，其中 $F$ 是域，都是唯一分解整环）。
	\item 下一章还将证明：只要 $R$ 本身是唯一分解整环，多项式环 $R[x]$ 也是唯一分解整环（这与主理想整环或 Euclidean整环的性质形成对比，后两种性质不会从 $R$ 传递到 $R[x]$）。这个结果连同上一例将说明 $\ZZ[x]$ 是唯一分解整环。
	\item Gaussian 整数的子环
	\[
	R=\ZZ[2i]=\{a+2bi\mid a,b\in\ZZ\},\qquad i^2=-1,
	\]
	是整环但不是唯一分解整环（这类环在第 7.1 节练习 23 中已经引入）。元素 $2$ 与 $2i$ 是不可约元，并且在 $R$ 中不相伴，因为 $i\notin R$；而
	\[
	4=2\cdot2=(-2i)(2i)
	\]
	给出 $R$ 中的两个不同分解。也可直接检查 $2i$ 在 $R$ 中不可约但不是素元（因为 $R/(2i)\cong\ZZ/4\ZZ$）。在更大的 Gaussian 整数环 $\ZZ[i]$（它是唯一分解整环）中，$2$ 与 $2i$ 相伴，因为 $i$ 在这个较大环中是单位。我们将在第 9.3 节末给出 $\ZZ[2i]$ 不是唯一分解整环的稍有不同的证明（该证明无需检查 $2$ 与 $2i$ 不可约）。
	\item 二次整数环 $\ZZ[\sqrt{-5}]$ 是另一个不是唯一分解整环的整环，因为
	\[
	6=2\cdot3=(1+\sqrt{-5})(1-\sqrt{-5})
	\]
	给出 $6$ 分解为不可约元的两个不同分解。$\ZZ[\sqrt{-5}]$ 中的主理想 $(6)$ 可以写成四个非主素理想的乘积：
	\[
	(6)=P_2^2P_3P_3',
	\]
	而元素 $6$ 的两个不同分解可解释为该理想乘积重新组合为主理想乘积的两种方式：$P_2^2=(2)$ 与 $P_3P_3'=(3)$ 的乘积，以及 $P_2P_3=(1+\sqrt{-5})$ 与 $P_2P_3'=(1-\sqrt{-5})$ 的乘积（参见练习 8）。
	
	虽然二次整数环 $\mcO$ 中的元素未必有唯一分解，但有一个定理（推论 16.16）断言 $\mcO$ 中每个理想都能唯一地写成素理想的乘积。理想唯一分解为素理想乘积的性质，对于代数数域的整数环（其中二次整数环就是例子）一般成立，并导向第 16 章考虑的 Dedekind整环概念。历史上，正是数论中代数整数环里的元素分解为不可约元的唯一性失败，最初导致了理想的定义。由此产生的这些环中分解为素理想的唯一性，使理想的元素具有一种“理想的”（在“完美”或“合意”的意义上）行为，这也正是这些如今基本的代数对象术语选择的基础。
\end{enumerate}

唯一分解整环中不可约元的第一个性质是它们也都是素元。也许会认为可以由下面命题以及前面提到的定理（我们马上要证明）“每个主理想整环都是唯一分解整环”推出命题 11；然而命题 11 将用于后一定理的证明中。

\medskip
\noindent\textbf{命题 12.}
在唯一分解整环中，非零元素是素元当且仅当它不可约。

\begin{proof}
	设 $R$ 是唯一分解整环。由命题 10，素元不可约；因此只需证明每个不可约元都是素元。设 $p$ 是 $R$ 中不可约元，并假设 $p\mid ab$，其中 $a,b\in R$；我们必须证明 $p$ 整除 $a$ 或 $b$。说 $p$ 整除 $ab$，就是说 $ab=pc$ 对某个 $c\in R$ 成立。把 $a,b$ 写成不可约元乘积，由这个等式以及 $ab$ 分解为不可约元的唯一性可知，不可约元 $p$ 必与 $a$ 的分解或 $b$ 的分解中出现的某个不可约元相伴。不妨设 $p$ 与 $a$ 的分解中某个不可约元相伴，也就是说，$a$ 可写成
	\[
	a=(up)p_2\cdots p_n,
	\]
	其中 $u$ 是单位，$p_2,\ldots,p_n$ 是某个（可能为空的）不可约元集合。那么 $p\mid a$，因为
	\[
	a=pd,
	\qquad d=up_2\cdots p_n.
	\]
	证明完成。
\end{proof}

在唯一分解整环中，我们现在将互换使用“素元”和“不可约元”这两个术语；不过通常仍说 $\ZZ$ 中的“素数”和 $F[x]$ 中的“不可约多项式”。

我们将用前一命题证明：在唯一分解整环中任意两个非零元素 $a,b$ 都有最大公因子。

\medskip
\noindent\textbf{命题 13.}
设 $a,b$ 是唯一分解整环 $R$ 中两个非零元素，并设
\[
a=u p_1^{e_1}p_2^{e_2}\cdots p_n^{e_n},
\qquad
b=v p_1^{f_1}p_2^{f_2}\cdots p_n^{f_n}
\]
是 $a,b$ 的素因子分解，其中 $u,v$ 是单位，素元 $p_1,p_2,\ldots,p_n$ 互不相同，并且指数 $e_i,f_i\geq0$。则元素
\[
d=p_1^{\min(e_1,f_1)}p_2^{\min(e_2,f_2)}\cdots p_n^{\min(e_n,f_n)}
\]
（若所有指数都是 $0$，则 $d=1$）是 $a,b$ 的最大公因子。

\begin{proof}
	由于 $d$ 中出现的每个素元的指数不大于 $a$ 与 $b$ 的分解中相应指数，$d$ 整除 $a$ 且整除 $b$。为证明 $d$ 是最大公因子，令 $c$ 是 $a,b$ 的任意公共因子，并写 $c=q_1^{g_1}q_2^{g_2}\cdots q_m^{g_m}$ 为 $c$ 的素因子分解。由于每个 $q_i$ 整除 $c$，因而整除 $a$ 与 $b$，由前一命题可知 $q_i$ 必整除某个 $p_j$。特别地，在相伴意义下（即至多相差一个单位），$c$ 中出现的素元必须是 $a,b$ 中出现的素元的子集：
	\[
	\{q_1,q_2,\ldots,q_m\}\subseteq\{p_1,p_2,\ldots,p_n\}.
	\]
	类似地，$c$ 中各素元的指数也不大于 $d$ 中相应指数。因此 $c\mid d$，证明完成。
\end{proof}

\paragraph{例}
在上面的例子中，$a=2210$ 且 $b=1131$，从它们的素因子分解立即得到 $(a,b)=13$。注意，如果 $a,b$ 的素因子分解已知，则上面的命题能立刻给出它们的最大公因子；但找到这些素因子分解在计算上非常耗时。Euclidean算法是确定两个整数最大公因子的最快方法，但遗憾的是，它几乎不给出这些整数的素因子分解信息。

现在我们来到本章所介绍的几类环之间的一个主要联系。

\medskip
\noindent\textbf{定理 14.}
每个主理想整环都是唯一分解整环。特别地，每个 Euclidean整环都是唯一分解整环。

\begin{proof}
	注意第二个断言由第一个断言推出，因为 Euclidean整环是主理想整环。为证明第一个断言，设 $R$ 是主理想整环，$r$ 是 $R$ 中非零且不是单位的元素。我们必须先证明 $r$ 可写成 $R$ 中有限个不可约元的乘积，然后验证该分解在相差单位的意义下唯一。
	
	第一部分证明的方法与确定整数的素因子分解完全类似。假设 $r$ 非零且不是单位。若 $r$ 本身不可约，则完成。若不然，由定义可将 $r$ 写成乘积 $r=r_1r_2$，其中 $r_1,r_2$ 都不是单位。若这两个元素都不可约，则同样完成，因为已经把 $r$ 写成不可约元乘积。否则，至少有一个元素，设为 $r_1$，是可约的，于是可写成两个非单位元素的乘积 $r_1=r_{11}r_{12}$，如此继续。我们必须证明这个过程终止，即必然到达某一步，使得作为 $r$ 的因子得到的所有元素都是不可约元。假设不是这样。由分解 $r=r_1r_2$ 得到理想的真包含
	\[
	(r)\subset (r_1)\subset R.
	\]
	第一个包含是真包含，因为 $r_2$ 不是单位；最后一个包含是真包含，因为 $r_1$ 不是单位。由 $r_1$ 的分解类似得到
	\[
	(r)\subset (r_1)\subset (r_{11})\subset R.
	\]
	若这个分解过程在有限步后不终止，则会得到一个无限上升理想链
	\[
	(r)\subset (r_1)\subset (r_{11})\subset\cdots\subset R,
	\]
	其中所有包含都是真包含，并且选择公理保证存在这样一条无限链（参见附录 I）。
	
	现在证明主理想整环中的任何上升理想链
	\[
	I_1\subseteq I_2\subseteq\cdots\subseteq R
	\]
	最终稳定，即存在正整数 $n$，使得对所有 $k\geq n$，$I_k=I_n$。\footnote{同一个论证可用来证明更一般的陈述：在每个理想都是有限生成的交换环中，上升理想链最终稳定。第 12 章将需要这个结果，并会重复其细节。}特别地，不可能存在所有包含都是真包含的无限上升理想链。令
	\[
	I=\bigcup_{i=1}^{\infty} I_i.
	\]
	容易验证（如第 7.4 节命题 11 的证明中一样）$I$ 是理想。由于 $R$ 是主理想整环，$I$ 主生成，设 $I=(a)$。因为 $I$ 是上述理想的并，$a$ 必属于链中某个理想，设 $a\in I_n$。于是
	\[
	I_n\subseteq I=(a)\subseteq I_n,
	\]
	所以 $I=I_n$，链从 $I_n$ 起稳定。这证明 $R$ 中每个非零非单位元素都有某个不可约元分解。
	
	还需证明上述分解本质唯一。我们对元素 $r$ 的某个分解中不可约因子的个数 $n$ 归纳。若 $n=0$，则 $r$ 是单位。若有 $r=qc$（某个其他分解），其中 $q$ 不可约，则 $q$ 将整除一个单位，因而 $q$ 本身是单位，矛盾。现设 $n\geq1$，且对 $r$ 有两个乘积
	\[
	r=p_1p_2\cdots p_n=q_1q_2\cdots q_m,
	\]
	其中 $p_i,q_j$ 是（不一定互异的）不可约元。于是 $p_1$ 整除右边乘积，由命题 11，$p_1$ 必整除其中一个因子。必要时重新编号，可设 $p_1\mid q_1$。于是 $q_1=p_1u$ 对某个 $u\in R$ 成立，而由于 $q_1$ 不可约，$u$ 实际上必须是单位。因此 $p_1$ 与 $q_1$ 相伴。在整环中可消去 $p_1$（这是合法的），得到
	\[
	p_2\cdots p_n=q_2' q_3\cdots q_m,
	\]
	其中 $q_2'=uq_2$ 仍不可约（与 $q_2$ 相伴）。由关于 $n$ 的归纳假设，左边每个因子与最右边的因子逐一匹配（在相伴意义下），从而也与中间原来的因子匹配（它们只差相伴）。由于 $p_1$ 与 $q_1$（在最初重新编号后）已经证明相伴，归纳步骤和定理证明完成。
\end{proof}

\medskip
\noindent\textbf{推论 15（算术基本定理）.}
整数环 $\ZZ$ 是唯一分解整环。

\begin{proof}
	整数环 $\ZZ$ 是 Euclidean整环，因此由定理是唯一分解整环。
\end{proof}

现在可以完成命题 9 中关于整环 $R$ 是否为 P.I.D. 与是否存在 Dedekind--Hasse范数之间的等价性。

\medskip
\noindent\textbf{推论 16.}
设 $R$ 是 P.I.D.，则 $R$ 上存在一个乘法性的 Dedekind--Hasse范数。

\begin{proof}
	若 $R$ 是 P.I.D.，则 $R$ 是 U.F.D.。定义范数 $N$ 如下：令 $N(0)=0$；若 $u$ 是单位，令 $N(u)=1$；若
	\[
	a=p_1p_2\cdots p_n
	\]
	是 $a$ 的不可约元分解，则令 $N(a)=2^n$。这一定义良好，因为 $a$ 的不可约因子的个数唯一。显然 $N(ab)=N(a)N(b)$，所以 $N$ 是正的且乘法性的。为证明 $N$ 是 Dedekind--Hasse范数，设 $a,b$ 是 $R$ 中非零元素。由假设，由 $a,b$ 生成的理想是主理想，设 $(a,b)=(r)$。若 $a$ 不属于理想 $(b)$，则 $r$ 也不属于 $(b)$，即 $r$ 不被 $b$ 整除。因为 $b=xr$ 对某个 $x\in R$ 成立，所以 $x$ 不是单位，于是
	\[
	N(b)=N(x)N(r)>N(r).
	\]
	因此 $(a,b)$ 中含有一个范数严格小于 $b$ 的范数的非零元素，证明完成。
\end{proof}

\subsection*{Gaussian 整数中的分解}

我们以描述 Gaussian 整数 $\ZZ[i]$ 中的不可约元并给出它对初等数论中 Fermat 一个著名定理的应用，结束对唯一分解整环的讨论。这是特别合适的，因为对 $\ZZ[i]$ 的经典研究开启了环的代数研究。

一般地，设 $\mcO$ 是二次整数环，$N$ 是第 7.1 节中引入的相应域范数。设 $\alpha\in\mcO$ 是一个范数为 $\ZZ$ 中素数 $p$ 的元素。若 $\alpha=\beta\gamma$，其中 $\beta,\gamma\in\mcO$，则
\[
p=N(\alpha)=N(\beta)N(\gamma),
\]
所以 $N(\beta)$ 或 $N(\gamma)$ 中有一个是 $\pm1$，另一个是 $\pm p$。因为我们已经看到，$\mcO$ 中元素的范数为 $\pm1$ 当且仅当它是 $\mcO$ 中的单位，所以 $\alpha$ 的一个因子是单位。因此
\[
\text{若 }N(\alpha)\text{ 是 }\ZZ\text{ 中的 }\pm\text{素数，则 }\alpha\text{ 在 }\mcO\text{ 中不可约。}
\]

设 $\pi$ 是 $\mcO$ 中的素元，并设 $(\pi)$ 是 $\mcO$ 中由 $\pi$ 生成的理想。由于 $(\pi)$ 是 $\mcO$ 中素理想，容易检查 $(\pi)\cap\ZZ$ 是 $\ZZ$ 中素理想（若整数 $a,b$ 满足 $ab\in(\pi)$，则 $a$ 或 $b$ 是 $(\pi)$ 的元素，所以 $a$ 或 $b$ 属于 $(\pi)\cap\ZZ$）。由于 $N(\pi)$ 是 $(\pi)$ 中的非零整数，我们有
\[
(\pi)\cap\ZZ=p\ZZ
\]
对某个整数素数 $p$ 成立。由 $p\in(\pi)$ 可知，$\pi$ 在 $\mcO$ 中整除整数素数 $p$；因此，$\mcO$ 中的素元可以通过确定 $\ZZ$ 中素数在较大环 $\mcO$ 中如何分解来找到。设 $\pi$ 整除 $\mcO$ 中的素数 $p$，比如 $p=\pi\pi'$。则
\[
N(\pi)N(\pi')=N(p)=p^2.
\]
由于 $\pi$ 不是单位，只有两种可能：或者 $N(\pi)=\pm p^2$，或者 $N(\pi)=\pm p$。在前一情形，$N(\pi')=\pm1$，故 $\pi'$ 是单位，而 $p=\pi$（至多相差相伴）在 $\ZZ[i]$ 中不可约。后一情形中，$N(\pi)=N(\pi')=\pm p$，故 $\pi'$ 也不可约，并且 $p=\pi\pi'$ 是恰好两个不可约元的乘积。

现在考虑特殊情形 $D=-1$，即 Gaussian 整数 $\ZZ[i]$。我们已经看到 $\ZZ[i]$ 的单位是 $\pm1$ 与 $\pm i$。第 1 节证明了 $\ZZ[i]$ 是 Euclidean整环，所以它也是主理想整环与唯一分解整环；因此其中不可约元与素元相同，并可通过观察 $\ZZ$ 中素数在较大环 $\ZZ[i]$ 中如何分解来确定。

在此情形，若 $\alpha=a+bi$，则
\[
N(\alpha)=\alpha\overline\alpha=a^2+b^2,
\]
其中 $\overline\alpha=a-bi$ 是 $\alpha$ 的复共轭。由上面的讨论可知，素数 $p$ 在 $\ZZ[i]$ 中恰分解为两个不可约元，当且仅当
\[
p=a^2+b^2
\]
是两个整数平方之和（否则 $p$ 在 $\ZZ[i]$ 中仍不可约）。若 $p=a^2+b^2$，则 $\ZZ[i]$ 中相应的不可约元是 $a\pm bi$。

显然
\[
2=1^2+1^2
\]
是两个平方和，给出分解
\[
2=(1+i)(1-i)=-i(1+i)^2.
\]
不可约元 $1+i$ 与 $1-i=-i(1+i)$ 相伴，并且容易检查这是共轭不可约元 $a+bi$ 与 $a-bi$ 能相伴的唯一情形。

由于任何整数的平方模 $4$ 只可能同余于 $0$ 或 $1$，$\ZZ$ 中一个奇素数若是两个平方和，则必同余于 $1$ 模 $4$。因此，若 $p$ 是 $\ZZ$ 中满足 $p\equiv3\pmod4$ 的素数，则 $p$ 不是两个平方和，并且 $p$ 在 $\ZZ[i]$ 中仍不可约。

现在设 $p$ 是 $\ZZ$ 中满足 $p\equiv1\pmod4$ 的素数。我们将证明 $p$ 在 $\ZZ[i]$ 中不可能不可约；这将说明 $p=(a+bi)(a-bi)$ 在 $\ZZ[i]$ 中分解为两个不同不可约元的乘积，等价地，说明 $p=a^2+b^2$ 是两个平方和。我们先证明如下初等数论结果。

\medskip
\noindent\textbf{引理 17.}
素数 $p\in\ZZ$ 整除某个形如 $n^2+1$ 的整数，当且仅当 $p$ 是 $2$，或 $p$ 是同余于 $1$ 模 $4$ 的奇素数。

\begin{proof}
	$p=2$ 的断言是平凡的，因为 $2\mid 1^2+1$。若 $p$ 是奇素数，注意 $p\mid n^2+1$ 等价于在 $\ZZ/p\ZZ$ 中有 $n^2=-1$。这又等价于 $n$ 的剩余类在乘法群 $(\ZZ/p\ZZ)^\times$ 中有阶 $4$。因此，$p$ 整除某个形如 $n^2+1$ 的整数，当且仅当 $(\ZZ/p\ZZ)^\times$ 中含有阶为 $4$ 的元素。由 Lagrange 定理，若 $(\ZZ/p\ZZ)^\times$ 含有阶为 $4$ 的元素，则
	\[
	|(\ZZ/p\ZZ)^\times|=p-1
	\]
	被 $4$ 整除，即 $p\equiv1\pmod4$。
	
	反过来，设 $p-1$ 被 $4$ 整除。我们先说明 $(\ZZ/p\ZZ)^\times$ 中恰有一个阶为 $2$ 的元素。若 $m^2\equiv1\pmod p$，则 $p$ 整除
	\[
	m^2-1=(m-1)(m+1).
	\]
	因此 $p$ 整除 $m-1$（即 $m\equiv1\pmod p$），或者 $p$ 整除 $m+1$（即 $m\equiv-1\pmod p$），所以 $-1$ 是 $(\ZZ/p\ZZ)^\times$ 中唯一的阶为 $2$ 的剩余类。现在 Abel 群 $(\ZZ/p\ZZ)^\times$ 含有一个阶为 $4$ 的子群 $H$（例如，对子群 $\{\pm1\}$ 取商后含有阶为 $2$ 的子群，其原像就是 $(\ZZ/p\ZZ)^\times$ 中一个阶为 $4$ 的子群）。由于 Klein 四元群有三个阶为 $2$ 的元素，而 $(\ZZ/p\ZZ)^\times$——因此 $H$——只有一个阶为 $2$ 的元素，$H$ 必为四阶循环群。因此 $(\ZZ/p\ZZ)^\times$ 含有阶为 $4$ 的元素，即 $H$ 的一个生成元。
\end{proof}

\noindent\emph{说明：}稍后我们将证明（第 9.5 节推论 19）$(\ZZ/p\ZZ)^\times$ 是循环群；由此可立即看出存在阶为 $4$ 的元素当且仅当 $p-1$ 被 $4$ 整除。

由引理 17，若 $p\equiv1\pmod4$ 是素数，则存在 $n\in\ZZ$ 使得 $p$ 在 $\ZZ$ 中整除 $n^2+1$，因此当然在 $\ZZ[i]$ 中整除
\[
n^2+1=(n+i)(n-i).
\]
若 $p$ 在 $\ZZ[i]$ 中不可约，则 $p$ 将在 $\ZZ[i]$ 中整除 $n+i$ 或 $n-i$。在这种情形下，由于 $p$ 是实数，它也将整除其复共轭，因而整除二者的差 $2i$。这显然不可能。于是我们证明了如下结果。

\medskip
\noindent\textbf{命题 18.}
\begin{enumerate}[label=(\arabic*)]
	\item （Fermat 平方和定理）素数 $p$ 可写成两个整数平方和
	\[
	p=a^2+b^2,
	\qquad a,b\in\ZZ,
	\]
	当且仅当 $p=2$ 或 $p\equiv1\pmod4$。除交换 $a,b$ 以及改变 $a,b$ 的符号外，$p$ 表示为两个平方和的方式唯一。
	\item Gaussian 整数 $\ZZ[i]$ 中的不可约元如下：
	\begin{enumerate}[label=(\alph*)]
		\item $1+i$（它的范数为 $2$）；
		\item $\ZZ$ 中满足 $p\equiv3\pmod4$ 的素数 $p$（它们的范数为 $p^2$）；
		\item 对于 $\ZZ$ 中满足 $p\equiv1\pmod4$ 的素数 $p$，若 $p=a^2+b^2=(a+bi)(a-bi)$，则 $a+bi$ 与 $a-bi$ 是 $p$ 的两个不同不可约因子（二者范数均为 $p$）。
	\end{enumerate}
\end{enumerate}

命题 18 的第一部分是初等数论中 Fermat 的一个著名定理，它还有许多其他证明。

更一般地，整数 $n\in\ZZ$ 能否写成两个整数平方和 $n=A^2+B^2$，等价于 $n$ 是否是 Gaussian 整数中某个元素 $A+Bi$ 的范数，即
\[
n=A^2+B^2=N(A+Bi).
\]
把
\[
A+Bi=\pi_1\pi_2\cdots\pi_k
\]
唯一地（至多差单位）写成不可约元乘积，由命题 18 对 $\ZZ[i]$ 中不可约元的显式描述可知，$n$ 是一个范数当且仅当 $n$ 的那些同余于 $3$ 模 $4$ 的素因子均以偶数次幂出现。进一步地，若 $n$ 满足此条件，则 $A+Bi$ 在 $\ZZ[i]$ 中的唯一分解还允许我们数出 $n$ 表示为两个平方和的方式数，如下推论所示。

\medskip
\noindent\textbf{推论 19.}
设 $n$ 是正整数，并写
\[
n=2^k p_1^{a_1}\cdots p_r^{a_r}q_1^{b_1}\cdots q_s^{b_s},
\]
其中 $p_1,\ldots,p_r$ 是互异的同余于 $1$ 模 $4$ 的素数，$q_1,\ldots,q_s$ 是互异的同余于 $3$ 模 $4$ 的素数。则 $n$ 可写成 $\ZZ$ 中两个平方和，即
\[
n=A^2+B^2,
\qquad A,B\in\ZZ,
\]
当且仅当每个 $b_i$ 都是偶数。进一步地，若 $n$ 满足此条件，则 $n$ 表示为两个平方和的表示数为
\[
4(a_1+1)(a_2+1)\cdots(a_r+1).
\]

\begin{proof}
	推论中的第一句话已经在上面证明。现假设 $b_1,\ldots,b_s$ 都为偶数。对每个同余于 $1$ 模 $4$ 的素数 $p_i$，写
	\[
	p_i=\pi_i\overline{\pi_i},
	\]
	其中 $\pi_i$ 与 $\overline{\pi_i}$ 是命题 18(2)(c) 中的不可约元。若 $N(A+Bi)=n$，则考察范数可知，在相差单位的意义下，$A+Bi$ 在 $\ZZ[i]$ 中分解为不可约元的形式是
	\[
	A+Bi=(1+i)^k
	(\pi_1^{\alpha_{1,1}}\overline{\pi_1}^{\alpha_{1,2}})\cdots
	(\pi_r^{\alpha_{r,1}}\overline{\pi_r}^{\alpha_{r,2}})
	q_1^{b_1/2}\cdots q_s^{b_s/2},
	\]
	其中非负整数 $\alpha_{i,1},\alpha_{i,2}$ 满足
	\[
	\alpha_{i,1}+\alpha_{i,2}=a_i,
	\qquad i=1,2,\ldots,r.
	\]
	由于 $\alpha_{i,1}$ 可取 $0,1,\ldots,a_i$（而后 $\alpha_{i,2}$ 即被确定），所以在相差单位的意义下，范数为 $n$ 的 $\ZZ[i]$ 中元素 $A+Bi$ 共有
	\[
	(a_1+1)(a_2+1)\cdots(a_r+1)
	\]
	个。最后，由于 $\ZZ[i]$ 中有四个单位，推论中的第二句话随之成立。
\end{proof}

\paragraph{例}
由于
\[
493=17\cdot29
\]
且两个素数都同余于 $1$ 模 $4$，$493=A^2+B^2$ 是两个整数平方和。又因为
\[
17=(4+i)(4-i),
\qquad
29=(5+2i)(5-2i),
\]
所以 $A+Bi$ 在 $\ZZ[i]$ 中在相差单位意义下的可能分解是
\begin{align*}
(4+i)(5+2i)&=18+13i,\\
(4+i)(5-2i)&=22-3i,\\
(4-i)(5+2i)&=22+3i,\\
(4-i)(5-2i)&=18-13i.\\
\end{align*}
乘以 $-1$ 会同时改变两个符号；乘以 $i$ 会交换 $A$ 与 $B$ 并引入一个符号变化。因此
\[
493=(\pm18)^2+(\pm13)^2=(\pm22)^2+(\pm3)^2,
\]
其中所有可能的符号选择给出了 $493$ 作为两个平方和的 $16$ 种表示中的 $8$ 种；剩下的 $8$ 种通过交换两个加数得到。

类似地，整数
\[
58000957=73^2\cdot17\cdot29
\]
也恰好能以 $16$ 种方式写成两个平方和，这些表示由上面 $493=A^2+B^2$ 中的每个整数 $A,B$ 都乘以 $73$ 得到。

\subsection*{小结}

综上，我们在含单位交换环的各类之间有如下包含关系：
\[
\text{域}\subset \text{Euclidean整环}\subset \text{P.I.D.s}\subset \text{U.F.D.s}\subset \text{整环},
\]
且所有包含都是真包含。回忆，$\ZZ$ 是不是域的 Euclidean整环；二次整数环
\[
\ZZ\left[{1+\sqrt{-19}\over2}\right]
\]
是主理想整环但不是 Euclidean整环；$\ZZ[x]$ 是唯一分解整环（第 9 章定理 7）但不是主理想整环；而 $\ZZ[\sqrt{-5}]$ 是整环但不是唯一分解整环。

\paragraph{练习}
\begin{enumerate}[label=\arabic*.]
	\item 令 $G=\QQ^\times$ 为非零有理数的乘法群。若 $\alpha=p/q\in G$，其中 $p,q$ 是互素整数，定义映射
	\[
	\varphi:G\longrightarrow G
	\]
	为在 $p$ 与 $q$ 的素数幂分解中互换素数 $2$ 与 $3$（例如
	\[
	\varphi(2^4 3^1 5^1 13^2)=3^4 2^1 5^1 13^2,
	\qquad
	\varphi(3/16)=\varphi(3/2^4)=2/3^4=2/81,
	\]
	并且 $\varphi$ 在所有分子、分母都与 $2$ 和 $3$ 互素的有理数上为恒等）。
	\begin{enumerate}[label=(\alph*)]
		\item 证明 $\varphi$ 是群同构。
		\item 证明群 $G$ 到自身有无穷多个同构。
		\item 证明上述同构都不能扩张为环 $\QQ$ 到自身的同构。事实上，证明 $\QQ$ 的唯一环同构是恒等映射。
	\end{enumerate}
	
	\item 设 $a,b$ 是唯一分解整环 $R$ 中的非零元素。证明 $a,b$ 有最小公倍数（参见第 1 节练习 11），并像命题 13 描述最大公因子那样，用 $a,b$ 的素因子分解描述它。
	
	\item 求整数
	\[
	2130797=17^2\cdot73\cdot101
	\]
	表示为两个平方和的全部表示。
	
	\item 证明：若一个整数是两个有理数平方的和，则它也是两个整数平方的和（例如
	\[
	13=(1/5)^2+(18/5)^2=2^2+3^2
	\]
	）。
	
	\item 令 $R=\ZZ[\sqrt n]$，其中 $n$ 是大于 $3$ 的无平方因子整数。
	\begin{enumerate}[label=(\alph*)]
		\item 证明 $2,\sqrt n$ 与 $1+\sqrt n$ 在 $R$ 中不可约。
		\item 证明 $R$ 不是 U.F.D.。推出当 $D\equiv2,3\pmod4$ 且 $D<-3$ 时，二次整数环 $\mcO$ 不是 U.F.D.（因此也不是 Euclidean整环，也不是 P.I.D.）。\emph{提示：}说明 $\sqrt n$ 或 $1+\sqrt n$ 不是素元。
		\item 给出 $R$ 中一个非主理想的显式例子。\emph{提示：}利用 (b)，考虑包含非素理想 $(\sqrt n)$ 或 $(1+\sqrt n)$ 的极大理想。
	\end{enumerate}
	
	\item
	\begin{enumerate}[label=(\alph*)]
		\item 证明商环 $\ZZ[i]/(1+i)$ 是一个 $2$ 元域。
		\item 设 $q\in\ZZ$ 是满足 $q\equiv3\pmod4$ 的素数。证明商环 $\ZZ[i]/(q)$ 是一个含 $q^2$ 个元素的域。
		\item 设 $p\in\ZZ$ 是满足 $p\equiv1\pmod4$ 的素数，并按命题 18 写成 $p=\pi\overline\pi$。证明中国剩余定理（第 7.6 节定理 17）的假设满足，并且有环同构
		\[
		\ZZ[i]/(p)\cong \ZZ[i]/(\pi)\times \ZZ[i]/(\overline\pi)
		\]
		作为环。证明商环 $\ZZ[i]/(p)$ 的阶为 $p^2$，并推出 $\ZZ[i]/(\pi)$ 与 $\ZZ[i]/(\overline\pi)$ 都是 $p$ 元域。
	\end{enumerate}
	
	\item 设 $\pi$ 是 $\ZZ[i]$ 中的不可约元。
	\begin{enumerate}[label=(\alph*)]
		\item 对任意整数 $n\geq0$，证明 $(\pi^{n+1})=\pi^{n+1}\ZZ[i]$ 是 $(\pi^n)=\pi^n\ZZ[i]$ 中的一个理想，并且乘以 $\pi^n$ 诱导出加法 Abel 群同构
		\[
		\ZZ[i]/(\pi)\cong (\pi^n)/(\pi^{n+1}).
		\]
		\item 证明
		\[
		|\ZZ[i]/(\pi^n)|=|\ZZ[i]/(\pi)|^n.
		\]
		\item 证明：对 $\ZZ[i]$ 中任意非零元素 $\alpha$，商环 $\ZZ[i]/(\alpha)$ 的阶等于 $N(\alpha)$。\emph{提示：}结合 (b)、中国剩余定理以及上一题的结果。
	\end{enumerate}
	
	\item 令 $R$ 为二次整数环 $\ZZ[\sqrt{-5}]$，并定义理想
	\[
	I_2=(2,1+\sqrt{-5}),\qquad
	I_3=(3,2+\sqrt{-5}),\qquad
	I_3'=(3,2-\sqrt{-5}).
	\]
	\begin{enumerate}[label=(\alph*)]
		\item 证明 $2,3,1+\sqrt{-5},1-\sqrt{-5}$ 在 $R$ 中都是不可约元，且任意两个都不相伴，并且
		\[
		6=2\cdot3=(1+\sqrt{-5})(1-\sqrt{-5})
		\]
		是 $6$ 在 $R$ 中分解为不可约元的两个不同分解。
		\item 证明 $I_2,I_3,I_3'$ 是 $R$ 中的素理想。\emph{提示：}对 $I_3$，注意由环的第三同构定理有
		\[
		R/I_3\cong (R/(3))/(I_3/(3)).
		\]
		证明 $R/(3)$ 有 $9$ 个元素，$I_3/(3)$ 有 $3$ 个元素，且 $R/I_3$ 作为加法 Abel 群同构于 $\ZZ/3\ZZ$。推出 $I_3$ 是极大理想（因此是素理想），并且 $R/I_3\cong \ZZ/3\ZZ$ 作为环。
		\item 说明 (a) 中的分解给出理想等式
		\[
		(6)=(2)(3),\qquad
		(6)=(1+\sqrt{-5})(1-\sqrt{-5}).
		\]
		证明这两个理想分解给出了理想 $(6)$ 分解为素理想乘积的同一个分解（参见第 2 节练习 5）。
	\end{enumerate}
	
	\item 假设二次整数环 $\mcO$ 是 P.I.D.。证明 $\mcO$ 上的域范数绝对值 $|N|$（参见第 7.1 节）是 $\mcO$ 上的 Dedekind--Hasse范数。推出：若二次整数环 $\mcO$ 存在任何 Dedekind--Hasse范数，则事实上 $\mcO$ 上的域范数绝对值已经给出了一个 Dedekind--Hasse范数。\emph{提示：}若 $\alpha,\beta\in\mcO$，则 $(\alpha,\beta)=(\gamma)$ 对某个 $\gamma\in\mcO$ 成立。证明若 $\beta$ 不整除 $\alpha$，则
	\[
	0<|N(\gamma)|<|N(\beta)|
	\]
	——使用 $\mcO$ 中的单位恰为范数为 $\pm1$ 的元素这一事实。
	
	\noindent\emph{说明：}如果 $\mcO$ 关于某个范数是 Euclidean整环，并不一定意味着它关于域范数绝对值是 Euclidean整环（虽然对 $D<0$ 这是真的，参见第 1 节练习 8）。一个例子是 $D=69$（参见 D. Clark, \emph{A quadratic field which is Euclidean but not norm-Euclidean}, Manuscripta Math. 83 (1994), pp. 327--330）。
	
	\item （$k$ 阶 Euclidean整环）设 $R$ 是整环，$N:R\to\ZZ^+\cup\{0\}$ 是 $R$ 上的范数。若对任意 $a,b\in R$，$b\neq0$，存在 $q,r\in R$ 满足
	\[
	a=qb+r,
	\qquad r=0\text{ 或 }N(r)<N(b),
	\]
	则 $R$ 关于 $N$ 是 Euclidean的。现在弱化这个条件：假设对任意 $a,b\in R$，$b\neq0$，存在 $q,q',r,r'\in R$ 使得
	\[
	a=qb+r,
	\qquad
	b=q'r+r',
	\qquad r'=0\text{ 或 }N(r')<N(b),
	\]
	即两次除法之后的余数变小。称这样的整环为 \emph{2 阶 Euclidean整环}。
	\begin{enumerate}[label=(\alph*)]
		\item 证明：在 2 阶 Euclidean整环中迭代这些除法会产生 $a$ 与 $b$ 的一个最大公因子，并且该最大公因子是 $a$ 与 $b$ 的线性组合。推出：2 阶 Euclidean整环中的每个有限生成理想都是主理想。（不过存在不是 P.I.D. 的 2 阶 Euclidean整环。）\emph{提示：}模仿定理 4 的证明。
		\item 证明：若一个 2 阶 Euclidean整环中每个非零非单位都可以分解为有限个不可约元的乘积，则它是唯一分解整环。\emph{提示：}先证明不可约元都是素元，如下：若 $p$ 不可约，$p\mid ab$，且 $p$ 不整除 $a$，利用 (a) 写出
		\[
		px+ay=1
		\]
		对某些 $x,y$ 成立。两边乘以 $b$，推出 $p\mid b$，所以 $p$ 是素元。然后沿用定理 14 中唯一性的证明。
		\item 对任意整数 $k\geq1$，作出显然的推广以定义 $k$ 阶 Euclidean整环。证明若把“2 阶 Euclidean”替换为“$k$ 阶 Euclidean”，(a) 与 (b) 仍成立。
	\end{enumerate}
	
	\noindent\emph{说明：}存在 2 阶 Euclidean 但不是 Euclidean的环。也存在关于给定范数不是 Euclidean、但关于该范数是 $k$ 阶 Euclidean的环（例如环 $\ZZ[\sqrt{14}]$ 关于通常范数 $N(a+b\sqrt{14})=|a^2-14b^2|$ 不是 Euclidean的，但关于此范数是 2 阶 Euclidean的）。$k$ 阶 Euclidean 条件还与如下问题有关：由 $R$ 中元素构成的可逆 $n\times n$ 矩阵群 $\operatorname{GL}_n(R)$ 是否由初等矩阵生成（初等矩阵指主对角线上全为 $1$，在主对角线外某一位置有一个 $1$，其余位置全为 $0$ 的矩阵）。
	
	\item （P.I.D. 的刻画）证明：$R$ 是 P.I.D. 当且仅当 $R$ 是同时为 Bezout整环的 U.F.D.（参见第 2 节练习 7）。\emph{提示：}一个方向由定理 14 给出。反过来，令 $a$ 是理想 $I$ 中一个非零元素，使其不可约因子个数最少。通过证明若存在 $b\in I$ 且 $b\notin(a)$，则 $(a,b)=(d)$ 会导致矛盾，从而证明 $I=(a)$。
\end{enumerate}

% \chapter{多项式环}

本章研究多项式环。我们首先汇总前两章中已经出现的若干事实，并在需要时给出引用。基本定义在第 7.2 节已稍详细地给出。为方便起见，除非另有说明，本章中的 $R$ 总表示带单位元 $1\neq0$ 的交换环。

\section{定义与基本性质}

多项式环 $R[x]$ 是由形如
\[
a_nx^n+a_{n-1}x^{n-1}+\cdots+a_1x+a_0
\]
的形式和组成的集合，其中 $n\geq0$ 且 $a_i\in R$。若 $a_n\neq0$，则称该多项式次数为 $n$，$a_nx^n$ 为首项，$a_n$ 为首系数；零多项式的首系数定义为 $0$。若首系数为 $1$，称多项式为首一的。多项式加法按系数逐项定义：
\[
\sum_{i=0}^n a_ix^i+\sum_{i=0}^n b_ix^i
=\sum_{i=0}^n (a_i+b_i)x^i,
\]
其中可把某些高次系数视为 $0$。乘法由
\[
(ax^i)(bx^j)=abx^{i+j}
\]
和分配律确定；一般地，
\[
\left(\sum_{i=0}^n a_i x^i\right)\left(\sum_{j=0}^m b_jx^j\right)
=\sum_{k=0}^{n+m}\left(\sum_{i=0}^k a_i b_{k-i}\right)x^k,
\]
其中在指标超出范围时相应系数视为 $0$。这样 $R[x]$ 是带单位元的交换环，其单位元就是 $R$ 中的 $1$，并且我们把 $R$ 认作 $R[x]$ 中的常数多项式子环。

我们已经知道，如果 $R$ 是整环，则非零多项式乘积的首项等于首项的乘积。下面把第 7.2 节命题 4 重新记录于此。

\noindent\textbf{命题 1.} 设 $R$ 是整环，则：
\begin{enumerate}[label=(\arabic*)]
	\item 若 $p(x),q(x)\neq0$，则
	\[
	\deg(p(x)q(x))=\deg p(x)+\deg q(x).
	\]
	\item $R[x]$ 的单位正是 $R$ 的单位。
	\item $R[x]$ 是整环。
\end{enumerate}

还应回忆：若 $R$ 是整环，则 $R[x]$ 的商域由所有 $p(x)/q(x)$ 组成，其中 $q(x)\neq0$；它称为系数在 $R$ 中的 $x$ 的有理函数域。

\noindent\textbf{命题 2.} 设 $I$ 是环 $R$ 的一个理想，令 $(I)=I[x]$ 表示 $R[x]$ 中由 $I$ 生成的理想，即所有系数都属于 $I$ 的多项式组成的集合。则
\[
R[x]/(I)\cong (R/I)[x].
\]
特别地，若 $I$ 是 $R$ 的素理想，则 $(I)$ 是 $R[x]$ 的素理想。

\begin{proof}
	定义自然映射
	\[
	\varphi:R[x]\longrightarrow (R/I)[x]
	\]
	为把多项式的每个系数模 $I$ 化。加法和乘法的定义说明 $\varphi$ 是环同态。其核正是系数全在 $I$ 中的多项式，即 $I[x]=(I)$，故得到同构。若 $I$ 是素理想，则 $R/I$ 是整环，由命题 1，$(R/I)[x]$ 也是整环，故 $(I)$ 是素理想。
\end{proof}

注意：若 $I$ 是 $R$ 的极大理想，$(I)$ 不一定是 $R[x]$ 的极大理想。不过若 $I$ 在 $R$ 中极大，则 $R[x]$ 中由 $I$ 与 $x$ 共同生成的理想是极大理想。

\paragraph{例}
取 $R=\ZZ$，$I=n\ZZ$。上面的同构写成
\[
\ZZ[x]/n\ZZ[x]\cong (\ZZ/n\ZZ)[x].
\]
自然投影就是把整数系数多项式的系数模 $n$ 化。若 $n$ 合成，商环不一定是整环；若 $n=p$ 是素数，则 $(\ZZ/p\ZZ)[x]$ 是整环，甚至是 Euclidean 整环。因此在 $\ZZ[x]$ 中，系数全被 $p$ 整除的多项式构成一个素理想。

\noindent\textbf{定义.} $R$ 上 $n$ 个变量 $x_1,x_2,\ldots,x_n$ 的多项式环记为
\[
R[x_1,x_2,\ldots,x_n],
\]
并递归地定义为
\[
R[x_1,x_2,\ldots,x_n]=R[x_1,x_2,\ldots,x_{n-1}][x_n].
\]
因此，$n$ 元多项式可视为关于 $x_n$ 的一元多项式，其系数是前 $n-1$ 个变量的多项式。

更具体地，一个非零多项式是有限多个非零单项式项的和，每项形如
\[
a x_1^{d_1}x_2^{d_2}\cdots x_n^{d_n},
\qquad a\in R,
\]
其中 $d_i\geq0$。首一项 $x_1^{d_1}x_2^{d_2}\cdots x_n^{d_n}$ 简称为一个单项式，并称为该项 $a x_1^{d_1}\cdots x_n^{d_n}$ 的单项式部分。指数 $d_i$ 称为该项关于 $x_i$ 的次数，和
\[
d=d_1+d_2+\cdots+d_n
\]
称为该项的次数。有序 $n$ 元组 $(d_1,d_2,\ldots,d_n)$ 称为该项的多次数。非零多项式的次数是其各单项式项次数的最大值。若所有项次数相同，则称该多项式是齐次的，或称为一个形式。若 $f$ 是 $n$ 个变量的非零多项式，则 $f$ 中所有次数为 $k$ 的单项式项之和称为 $f$ 的 $k$ 次齐次分量。次数为 $d$ 的非零多项式可唯一写成
\[
f=f_0+f_1+\cdots+f_d,
\]
其中 $f_k$ 是 $f$ 的 $k$ 次齐次分量。

任意多个变量的多项式环也可类似定义为有限个单项式项的和；等价地，它是所有有限变量多项式环的并。

\paragraph{例}
二元多项式环 $\ZZ[x,y]$ 由所有形如 $a x^i y^j$ 的单项式项的有限和组成，其中 $a\in\ZZ$，该项的次数为 $i+j$。例如
\[
p(x,y)=2x^3+xy-y^2,
\qquad
q(x,y)=-3xy+2y^2+x^2y^3
\]
都是 $\ZZ[x,y]$ 中的元素，次数分别为 $3$ 与 $5$。我们有
\[
p(x,y)+q(x,y)=2x^3-2xy+y^2+x^2y^3
\]
以及
\[
\begin{aligned}
p(x,y)q(x,y)= {}& -6x^4y+4x^3y^2+2x^5y^3-3x^2y^2+5xy^3+x^3y^4 \\
&{}-2y^4-x^2y^5,
\end{aligned}
\]
这是一个次数为 $8$ 的多项式。若把最后这个多项式看成关于 $y$ 的多项式，其系数在 $\ZZ[x]$ 中，则可写成
\[
(-6x^4)y+(4x^3-3x^2)y^2+(2x^5+5x)y^3+(x^3-2)y^4-(x^2)y^5.
\]
若 $f=f(x,y)=p(x,y)q(x,y)$，则 $f$ 的非零齐次分量为
\[
f_4=-3x^2y^2+5xy^3-2y^4,
\]
\[
f_5=-6x^4y+4x^3y^2,
\qquad
f_7=x^3y^4-x^2y^5,
\qquad
f_8=2x^5y^3.
\]

命题 1 的所有陈述对任意多个变量的多项式环仍然成立；有限多个变量由归纳得到，任意多个变量由“有限变量多项式环的并”得到。


\paragraph{练习}
\begin{enumerate}[label=\arabic*.]
	\item 设
	\[
	p(x,y,z)=2x^2y-3xy^3z+4y^2z^5,
	\qquad q(x,y,z)=7x^2+5x^2y^3z^4-3x^2z^3
	\]
	是 $\ZZ[x,y,z]$ 中的多项式。
	\begin{enumerate}[label=(\alph*)]
		\item 把 $p$ 与 $q$ 都写成关于 $x$ 的多项式，其系数在 $\ZZ[y,z]$ 中。
		\item 求 $p$ 与 $q$ 各自的次数。
		\item 分别求 $p$ 与 $q$ 关于三个变量 $x,y,z$ 的次数。
		\item 计算 $pq$，并分别求 $pq$ 关于三个变量 $x,y,z$ 的次数。
		\item 把 $pq$ 写成关于变量 $z$ 的多项式，其系数在 $\ZZ[x,y]$ 中。
	\end{enumerate}
	
	\item 在假定 $p$ 与 $q$ 的系数属于 $\ZZ/3\ZZ$ 的情形下，重复前一题。
	
	\item 设 $R$ 是交换环，$x_1,x_2,\ldots,x_n$ 是 $R$ 上彼此独立的变量。证明：对 $\{1,2,\ldots,n\}$ 的任意置换 $\pi$，有环同构
	\[
	R[x_{\pi(1)},x_{\pi(2)},\ldots,x_{\pi(n)}]\cong R[x_1,x_2,\ldots,x_n].
	\]
	
	\item 证明 $\QQ[x,y]$ 中的理想 $(x)$ 与 $(x,y)$ 都是素理想，但只有后者是极大理想。
	
	\item 证明 $\ZZ[x,y]$ 中的理想 $(x,y)$ 与 $(2,x,y)$ 都是素理想，但只有后者是极大理想。
	
	\item 证明 $(x,y)$ 不是 $\QQ[x,y]$ 中的主理想。
	
	\item 设 $R$ 是含 $1$ 的交换环。证明：$R$ 上多于一个变量的多项式环不是主理想整环。
	
	\item 设 $F$ 是域，并令
	\[
	R=F[x,x^2y,x^3y^2,\ldots,x^ny^{n-1},\ldots]
	\]
	为多项式环 $F[x,y]$ 的子环。
	\begin{enumerate}[label=(\alph*)]
		\item 证明 $R$ 与 $F[x,y]$ 的分式域相同。
		\item 证明 $R$ 含有一个不是有限生成的理想。
	\end{enumerate}
	
	\item 证明：以任意交换环为系数的无限多个变量的多项式环含有不是有限生成的理想。
	
	\item 证明环
	\[
	\ZZ[x_1,x_2,x_3,\ldots]/(x_1x_2,x_3x_4,x_5x_6,\ldots)
	\]
	含有无限多个极小素理想（参见第 7.4 节练习 36）。
	
	\item 证明 $\QQ[x,y]$ 中理想 $I=(x,y^2)$ 的根理想是 $(x,y)$（参见第 7.4 节练习 30）。由此推出 $I$ 是一个 primary 理想，但不是素理想的幂（参见第 7.4 节练习 41）。
	
	\item 设 $R=\QQ[x,y,z]$，并用横线表示到商环 $\QQ[x,y,z]/(xy-z^2)$ 的像。证明
	\[
	P=(\overline{x},\overline{z})
	\]
	是素理想。证明 $\overline{x}\overline{y}\in P^2$，但 $\overline{y}$ 的任何幂都不属于 $P^2$。（这说明 $P$ 是一个素理想，而它的平方不是 primary 理想；参见第 7.4 节练习 41。）
	
	\item 证明：对任意域 $F$，两个环
	\[
	F[x,y]/(y^2-x),\qquad F[x,y]/(y^2-x^2)
	\]
	不同构。
	
	\item 设 $R$ 是整环，$i,j$ 是互素整数。证明 $R[x,y]$ 中的理想 $(x^i-y^j)$ 是素理想。\emph{提示：}考虑环同态 $\varphi:R[x,y]\to R[t]$，定义为 $x\mapsto t^j$，$y\mapsto t^i$。证明 $R[x,y]$ 中任一元素模 $(x^i-y^j)$ 可化为一个关于 $y$ 次数至多为 $j-1$ 的多项式，并观察当 $0\leq s<j$ 时 $\varphi(x^r y^s)$ 的指数互不相同。
	
	\item 设 $p(x_1,x_2,\ldots,x_n)$ 是 $R[x_1,\ldots,x_n]$ 中次数为 $k$ 的齐次多项式。证明：对所有 $\lambda\in R$，有
	\[
	p(\lambda x_1,\lambda x_2,\ldots,\lambda x_n)=\lambda^k p(x_1,x_2,\ldots,x_n).
	\]
	
	\item 证明两个齐次多项式的乘积仍是齐次多项式。
	
	\item 若 $I$ 是 $R[x_1,\ldots,x_n]$ 中的理想，并且只要 $p\in I$，$p$ 的每个齐次分量也都属于 $I$，则称 $I$ 为齐次理想。证明：一个理想是齐次理想，当且仅当它可由齐次多项式生成。\emph{提示：}用次数归纳证明“若”方向。
	
	\item 下面的练习说明在非交换环 $R$ 上处理多项式时，尤其是把多项式看成函数时，必须小心（$R[x]$ 中的环运算仍按交换环情形同样定义）。设 $R$ 是任意环，$\operatorname{Func}(R)$ 为从 $R$ 到自身的所有函数构成的环。若 $p(x)\in R[x]$，令 $\overline p\in\operatorname{Func}(R)$ 是由
	\[
	\overline p(r)=p(r)
	\]
	定义的函数，即通常把 $R[x]$ 中多项式“代入 $r$”所得的函数。
	\begin{enumerate}[label=(\alph*)]
		\item 对固定的 $a\in R$，证明“在 $a$ 处取值”是从 $\operatorname{Func}(R)$ 到 $R$ 的环同态（参见第 7.3 节定理 7 后例 4）。
		\item 证明映射 $\varphi:R[x]\to\operatorname{Func}(R)$，$\varphi(p(x))=\overline p$，一般不是环同态。由此推出，多项式恒等式在把多项式视为函数后不一定给出相应的函数恒等式。\emph{提示：}若 $R=\mathbb H$ 是实 Hamilton 四元数环，证明 $p(x)=x^2+1$ 可分解为 $(x+i)(x-i)$，但 $p(j)=0$ 而 $(j+i)(j-i)\neq0$。
		\item 对固定的 $a\in R$，证明 (a) 与 (b) 中映射的复合，即 $R[x]\to R$ 的“在 $a$ 处取值”映射，是环同态当且仅当 $a$ 属于 $R$ 的中心。
	\end{enumerate}
\end{enumerate}

\section{域上的多项式环 I}

现在更仔细地考虑系数环是域 $F$ 的情形。在 $F[x]$ 上定义范数
\[
N(p(x))=\deg p(x),\qquad N(0)=0.
\]
熟悉的多项式长除法在任意域上都成立。

\noindent\textbf{定理 3.} 设 $F$ 是域，则 $F[x]$ 是 Euclidean 整环。更具体地，若 $a(x),b(x)\in F[x]$ 且 $b(x)\neq0$，则存在唯一的 $q(x),r(x)\in F[x]$ 使
\[
a(x)=q(x)b(x)+r(x),
\qquad r(x)=0\text{ 或 }\deg r(x)<\deg b(x).
\]

\begin{proof}
	若 $a(x)=0$，取 $q=r=0$。否则对 $n=\deg a(x)$ 归纳。设 $\deg b=m$。若 $n<m$，取 $q=0,r=a$。若 $n\geq m$，写
	\[
	a(x)=a_nx^n+\cdots+a_0,
	\qquad b(x)=b_mx^m+\cdots+b_0,
	\]
	其中 $b_m\neq0$。因为 $F$ 是域，可令
	\[
	a'(x)=a(x)-{a_n\over b_m}x^{n-m}b(x),
	\]
	此多项式次数小于 $n$。由归纳假设，
	\[
	a'(x)=q'(x)b(x)+r(x),
	\]
	其中 $r=0$ 或 $\deg r<m$。令
	\[
	q(x)=q'(x)+{a_n\over b_m}x^{n-m}
	\]
	即得存在性。
	
	若另有 $a=q_1b+r_1$ 满足条件，则
	\[
	b(q-q_1)=r_1-r.
	\]
	右边次数小于 $m=\deg b$，而若 $q-q_1\neq0$，左边次数至少为 $m$，矛盾。因此 $q=q_1$，进而 $r=r_1$。
\end{proof}

\noindent\textbf{推论 4.} 若 $F$ 是域，则 $F[x]$ 是主理想整环，也是唯一分解整环。

\begin{proof}
	由上一章 Euclidean 整环、P.I.D. 与 U.F.D. 的关系立即得到。
\end{proof}

回忆第 8.2 节推论 8：若 $R[x]$ 是 P.I.D.（或 Euclidean 整环），则 $R$ 必为域。下一节将看到，U.F.D. 性质不同：只要 $R$ 是 U.F.D.，$R[x]$ 也是 U.F.D.。

\paragraph{例}
\begin{enumerate}[label=(\arabic*)]
	\item 由上面的评注可知，环 $\ZZ[x]$ 不是主理想整环。我们已经见过（第 7.4 节开头例 3）理想 $(2,x)$ 不是这个环中的主理想。
	\item $\QQ[x]$ 是主理想整环，因为系数取在域 $\QQ$ 中。在 $\ZZ[x]$ 这个子环中，由 $2$ 与 $x$ 生成的理想不是主理想；然而在 $\QQ[x]$ 中由它们生成的理想是主理想，事实上它就是整个环（所以生成元可取 $1$），因为 $2$ 是 $\QQ[x]$ 中的单位。
	\item 若 $p$ 是素数，则将 $\ZZ[x]$ 按素理想 $(p)$ 约化得到的环 $\ZZ/p\ZZ[x]$ 是主理想整环，因为其系数取在域 $\ZZ/p\ZZ$ 中。这个例子说明，一个不是主理想整环的环的商环可能是主理想整环。继续跟踪上面理想 $(2,x)$：若 $p=2$，则理想 $(2,x)$ 在商环 $\ZZ/2\ZZ[x]$ 中约化为理想 $(x)$，这是一个真（极大）理想；若 $p\neq2$，则 $2$ 在商环中是单位，所以 $(2,x)$ 约化为整个环 $\ZZ/p\ZZ[x]$。
	\item $\QQ[x,y]$，即有理系数二元多项式环，不是主理想整环，因为这个环是 $\QQ[x][y]$，而 $\QQ[x]$ 不是域（任何正次数元素都不可逆）。练习中要求证明理想 $(x,y)$ 不是这个环中的主理想。
\end{enumerate}

还要注意：若 $F\subseteq E$ 是域扩张，且 $a(x),b(x)\in F[x]$，则在 $F[x]$ 中对 $a$ 除以 $b$ 得到的商和余数，与在更大的环 $E[x]$ 中得到的商和余数相同。这由除法算法的唯一性推出。特别地，$b$ 在 $E[x]$ 中整除 $a$ 当且仅当它在 $F[x]$ 中整除 $a$；把最大公因子规定为首一后，在 $F[x]$ 与 $E[x]$ 中所得相同。


\paragraph{练习}
以下设 $F$ 为域，$x$ 为 $F$ 上的不定元。
\begin{enumerate}[label=\arabic*.]
	\item 设 $f(x)\in F[x]$ 的次数为 $n\geq1$，并用横线表示到商环 $F[x]/(f(x))$ 的像。证明：对每个 $g(x)$，存在唯一的次数至多为 $n-1$ 的多项式 $g_0(x)$，使
	\[
	\overline{g(x)}=\overline{g_0(x)}.
	\]
	等价地，$1,x,\ldots,x^{n-1}$ 是 $F$-向量空间 $F[x]/(f(x))$ 的一组基；特别地，该空间维数为 $n$。\emph{提示：}使用除法算法。
	
	\item 设 $F$ 是有 $q$ 个元素的有限域，$f(x)\in F[x]$ 的次数为 $n\geq1$。证明 $F[x]/(f(x))$ 有 $q^n$ 个元素。\emph{提示：}使用前一题。
	
	\item 设 $f(x)\in F[x]$。证明 $F[x]/(f(x))$ 是域，当且仅当 $f(x)$ 不可约。\emph{提示：}使用第 8.2 节命题 7。
	
	\item 设 $F$ 是有限域。证明 $F[x]$ 含有无限多个素元。（注意在无限域上，一次多项式已经给出多项式环中的无限多个素元。）
	
	\item 设 $F$ 是域，$p(x)\in F[x]$。列出环 $F[x]/(p(x))$ 的所有理想，并用 $p(x)$ 的分解来描述它们。
	
	\item 简要描述下列环的环结构：
	\[
	\text{(a) }\ZZ[x]/(2),\quad
	\text{(b) }\ZZ[x]/(x),\quad
	\text{(c) }\ZZ[x]/(x^2),\quad
	\text{(d) }\ZZ[x,y]/(x^2,y^2,2).
	\]
	证明在最后一个环中，对每个元素 $a$ 都有 $a^2=0$ 或 $1$，并确定满足 $a^2=0$ 的元素。确定这些环各自的特征（参见第 7.3 节练习 26）。
	
	\item 确定环 $\ZZ[x]/(2,x^3+1)$ 的所有理想。
	
	\item 在 $\QQ[x]$ 中求
	\[
	a(x)=x^3-2,
	\qquad b(x)=x+1
	\]
	的最大公因子，并把它写成 $a(x)$ 与 $b(x)$ 的一个 $\QQ[x]$-线性组合。
	
	\item 在 $\QQ[x]$ 中求
	\[
	a(x)=x^5+2x^3+x^2+x+1,
	\qquad b(x)=x^5+x^4+2x^3+2x^2+2x+1
	\]
	的最大公因子，并把它写成 $a(x)$ 与 $b(x)$ 的一个 $\QQ[x]$-线性组合。
	
	\item 在 $\QQ[x]$ 中求
	\[
	a(x)=x^3+4x^2+x-6,
	\qquad b(x)=x^5-6x+5
	\]
	的最大公因子，并把它写成 $a(x)$ 与 $b(x)$ 的一个 $\QQ[x]$-线性组合。
	
	\item 设 $f(x),g(x)$ 是 $\QQ[x]$ 中两个非零多项式，最大公因子为 $d(x)$。
	\begin{enumerate}[label=(\alph*)]
		\item 给定 $h(x)\in\QQ[x]$，证明存在 $a(x),b(x)\in\QQ[x]$ 满足
		\[
		a(x)f(x)+b(x)g(x)=h(x)
		\]
		当且仅当 $h(x)$ 被 $d(x)$ 整除。
		\item 若 $a_0(x),b_0(x)\in\QQ[x]$ 是 (a) 中方程的一组特解，证明该方程的全部解为
		\[
		a(x)=a_0(x)+m(x){g(x)\over d(x)},
		\qquad
		b(x)=b_0(x)-m(x){f(x)\over d(x)},
		\]
		其中 $m(x)$ 遍历 $\QQ[x]$。参见第 8.1 节练习 4。
	\end{enumerate}
	
	\item 设 $F[x,y_1,y_2,\ldots]$ 是域 $F$ 上变量 $x,y_1,y_2,\ldots$ 的多项式环，令
	\[
	I=(x-y_1^2,\,y_1-y_2^2,\,\ldots,\,y_i-y_{i+1}^2,\ldots)
	\]
	为其中的理想。定义
	\[
	R=F[x,y_1,y_2,\ldots]/I,
	\]
	于是，在 $R$ 中每个 $y_{i+1}$ 的平方都是 $y_i$，并且 $y_i^{2^i}=x$（模 $I$），也就是说 $x$ 对每个 $i$ 都有一个 $2^i$ 次根。记 $y_i$ 在 $R$ 中的像为 $x^{1/2^i}$。令 $R_n$ 为由 $F$ 与 $x^{1/2^n}$ 生成的 $R$ 的子环。
	\begin{enumerate}[label=(\alph*)]
		\item 证明 $R_1\subseteq R_2\subseteq\cdots$，并且 $R=\bigcup_{n\geq1}R_n$。
		\item 证明 $R_n$ 同构于 $F$ 上的一元多项式环，因此 $R_n$ 是 P.I.D.。推出 $R$ 是 Bezout 整环（参见第 8.2 节练习 7）。\emph{提示：}先证明
		\[
		S_n=F[x,y_1,\ldots,y_n]/(x-y_1^2,y_1-y_2^2,\ldots,y_{n-1}-y_n^2)
		\]
		同构于多项式环 $F[y_n]$。然后证明 $y_n$ 在 $R_n$ 中满足的任何多项式关系，都给出某个 $N\geq n$ 时 $S_N$ 中的相应关系。
		\item 证明 $R$ 中由 $x,x^{1/2},x^{1/4},\ldots$ 生成的理想不是有限生成的（因此 $R$ 不是 P.I.D.）。
	\end{enumerate}
	
	\item 本题介绍一个非交换环，它是“右”Euclidean 整环（也是“左”主理想整环），但不是“左”Euclidean 整环（也不是“右”主理想整环）。设 $F$ 是特征为 $p$ 的域，且并非每个元素都是 $p$ 次幂，即 $F\neq F^p$（例如有理函数域 $F=\FF_p(t)$）。令 $R=F\{x\}$ 为 $x$ 的“扭曲”多项式环，其元素为 $\sum_{i=0}^n a_i x^i$，其中 $a_i\in F$；加法按项定义，乘法由
	\[
	\left(\sum_{i=0}^n a_i x^i\right)
	\left(\sum_{j=0}^m b_j x^j\right)
	=
	\sum_{k=0}^{n+m}\left(\sum_{i+j=k} a_i b_j^{p^i}\right)x^k
	\]
	定义。这来自规定 $xa=a^p x$（$a\in F$），因而 $x$ 的幂与系数不交换，再自然延拓。令 $N$ 为取多项式次数的范数：$N(f)=\deg f$。
	\begin{enumerate}[label=(\alph*)]
		\item 证明对每个 $a\in F$ 和每个整数 $k\geq0$，有 $x^k a=a^{p^k}x^k$，并证明 $R$ 在上述乘法下是环。\emph{提示：}用特征 $p$ 中 $(a+b)^p=a^p+b^p$，从而 $(a+b)^{p^k}=a^{p^k}+b^{p^k}$。
		\item 证明 $R$ 中两个元素乘积的次数等于两个元素次数之和。证明 $R$ 没有零因子。
		\item 证明 $R$ 关于 $N$ 是“右 Euclidean”的：对任意 $f,g\in R$，$g\neq0$，存在 $q,r\in R$ 使
		\[
		f=qg+r,
		\qquad r=0\text{ 或 }\deg r<\deg g.
		\]
		用此证明 $R$ 的每个左理想都是主理想。
		\item 取 $f=cx$，其中 $c\in F$ 且 $c\notin F^p$，并取 $g=x$。证明不存在 $q,r\in R$ 使
		\[
		f=gq+r,
		\qquad r=0\text{ 或 }\deg r<\deg g.
		\]
		特别地，$R$ 关于 $N$ 不是“左 Euclidean”的。证明由 $x$ 与 $cx$ 生成的 $R$ 的右理想不是主理想。推出 $R$ 关于任何范数都不是“左 Euclidean”的。
	\end{enumerate}
\end{enumerate}

\section{唯一分解整环上的多项式环}

若 $R$ 是整环，它嵌入其分式域 $F$，从而 $R[x]\subseteq F[x]$。由于 $F[x]$ 是 Euclidean 整环，许多计算可以先在 $F[x]$ 中完成，再尝试清去分母回到 $R[x]$。自然问题是：若 $R$ 是 U.F.D.，$R[x]$ 是否也是 U.F.D.？答案是肯定的。关键是 Gauss 引理。

\noindent\textbf{命题 5（Gauss 引理）.} 设 $R$ 是 U.F.D.，$F$ 是其分式域，$p(x)\in R[x]$。若 $p(x)$ 在 $F[x]$ 中可约，则 $p(x)$ 在 $R[x]$ 中也可约。更准确地说，若
\[
p(x)=A(x)B(x)
\]
其中 $A(x),B(x)\in F[x]$ 是非常数多项式，则存在非零 $r,s\in F$，使
\[
rA(x)=a(x)\in R[x],\qquad sB(x)=b(x)\in R[x],
\]
并且
\[
p(x)=a(x)b(x)
\]
是在 $R[x]$ 中的分解。

\begin{proof}
	把 $A,B$ 的系数通分，得
	\[
	d p(x)=a'(x)b'(x),
	\]
	其中 $d\in R$ 非零，$a',b'\in R[x]$。若 $d$ 是单位，结论显然。否则把 $d$ 分解为不可约元乘积 $d=p_1\cdots p_n$。因为 $R$ 是 U.F.D.，不可约元 $p_1$ 是素元，所以 $(p_1)$ 是素理想。由命题 2，$(R/(p_1))[x]$ 是整环。将等式模 $p_1$ 化，得
	\[
	0=\overline{a'}(x)\overline{b'}(x),
	\]
	故两因子之一为零，即 $a'$ 或 $b'$ 的所有系数都被 $p_1$ 整除。于是可把 $p_1$ 从右边一个因子中消去，同时把左边的 $d$ 换成 $d/p_1$。依次对 $d$ 的所有不可约因子重复此过程，最终得到
	\[
	p(x)=a(x)b(x),\qquad a,b\in R[x],
	\]
	且 $a,b$ 分别为 $A,B$ 的 $F$-倍数。
\end{proof}

注意，一般不能要求 $a,b$ 是 $A,B$ 的 $R$-倍数。例如在 $\ZZ[x]$ 中把 $x^2$ 写为 $(2x)(\frac12x)$，虽然可清分母得到 $x\cdot x$，但不是原两因子的整数倍。

\noindent\textbf{推论 6.} 设 $R$ 是 U.F.D.，$F$ 是其分式域，$p(x)\in R[x]$。若 $p(x)$ 的各系数最大公因子为 $1$，则 $p(x)$ 在 $R[x]$ 中不可约当且仅当它在 $F[x]$ 中不可约。特别地，若首一多项式 $p(x)$ 在 $R[x]$ 中不可约，则它在 $F[x]$ 中不可约。

\begin{proof}
	若 $p$ 在 $F[x]$ 中可约，由 Gauss 引理在 $R[x]$ 中可约。反过来，若 $p$ 在 $R[x]$ 中可约，由于系数最大公因子为 $1$，其分解中的两个因子都不能是 $R$ 中的常数非单位，故它在 $F[x]$ 中也可约。
\end{proof}

\noindent\textbf{定理 7.} $R$ 是 U.F.D. 当且仅当 $R[x]$ 是 U.F.D.。

\begin{proof}
	若 $R[x]$ 是 U.F.D.，则常数多项式的唯一分解迫使 $R$ 自身是 U.F.D.。
	
	反之设 $R$ 是 U.F.D.，$F$ 是其分式域。取非零 $p(x)\in R[x]$。令 $d$ 为 $p$ 的系数最大公因子，写
	\[
	p(x)=d p'(x),
	\]
	其中 $p'$ 的系数最大公因子为 $1$。$d$ 在 $R$ 中唯一分解；因此只需处理系数最大公因子为 $1$ 且次数正的情形。
	
	在 $F[x]$ 中，$p(x)$ 唯一分解为不可约多项式乘积。由 Gauss 引理，可清分母得到 $R[x]$ 中的分解；因为 $p$ 的系数最大公因子为 $1$，每个因子的系数最大公因子也为 $1$，由推论 6 它们在 $R[x]$ 中不可约。于是存在不可约分解。
	
	唯一性也由 $F[x]$ 中的唯一性推出。若
	\[
	p(x)=q_1(x)\cdots q_r(x)=q_1'(x)\cdots q_s'(x)
	\]
	是在 $R[x]$ 中的两个不可约分解，则在 $F[x]$ 中两边也是不可约分解，因此 $r=s$ 且适当重排后 $q_i$ 与 $q_i'$ 在 $F[x]$ 中相伴。设 $q=(a/b)q'$，其中 $a,b\in R$ 非零，且 $q,q'$ 的系数最大公因子均为 $1$。则
	\[
	bq=aq'.
	\]
	左边系数最大公因子为 $b$（至多差一个单位），右边为 $a$，故 $a$ 与 $b$ 在 $R$ 中相伴，从而 $q$ 与 $q'$ 在 $R[x]$ 中相伴。
\end{proof}

\noindent\textbf{推论 8.} 若 $R$ 是 U.F.D.，则以 $R$ 为系数的任意多个变量的多项式环仍是 U.F.D.。

\begin{proof}
	有限多个变量由定理 7 归纳得到；任意多个变量的情形由其为有限变量多项式环的并得到。
\end{proof}

\paragraph{例}
\begin{enumerate}[label=(\arabic*)]
	\item $\ZZ[x]$、$\ZZ[x,y]$ 等都是 U.F.D.。$\ZZ[x]$ 是 U.F.D. 但不是 P.I.D. 的例子。
	\item $\QQ[x]$、$\QQ[x,y]$ 等也都是 U.F.D.。
\end{enumerate}

若 $R$ 只是整环而非 U.F.D.，推论 6 的结论可能失败。例如 $R=\ZZ[2i]=\{a+2bi\mid a,b\in\ZZ\}$，其分式域为 $\QQ(i)$。多项式 $x^2+1$ 在 $R[x]$ 中不可约，但在 $\QQ(i)[x]$ 中分解为 $(x-i)(x+i)$。这也说明 $\ZZ[2i]$ 不是 U.F.D.。


\paragraph{练习}
\begin{enumerate}[label=\arabic*.]
	\item 设 $R$ 是整环，分式域为 $F$，且 $p(x)$ 是 $R[x]$ 中的首一多项式。假设
	\[
	p(x)=a(x)b(x),
	\]
	其中 $a(x),b(x)$ 是 $F[x]$ 中次数小于 $p(x)$ 的首一多项式。证明：若 $a(x)\notin R[x]$，则 $R$ 不是唯一分解整环。由此推出 $\ZZ[2\sqrt2]$ 不是 U.F.D.。
	
	\item 证明：若 $f(x),g(x)$ 是有理系数多项式，且乘积 $f(x)g(x)$ 有整数系数，则 $g(x)$ 的任一系数与 $f(x)$ 的任一系数的乘积都是整数。
	
	\item 设 $F$ 是域。证明 $F[x]$ 中 $x$ 的系数等于 $0$ 的所有多项式组成一个子环 $R$，并证明 $R$ 不是 U.F.D.。\emph{提示：}证明
	\[
	x^6=(x^2)^3=(x^3)^2
	\]
	给出 $x^6$ 到不可约元的两个不同分解。
	
	\item 设
	\[
	R=\ZZ+x\QQ[x]\subseteq\QQ[x]
	\]
	为常数项为整数的有理系数多项式集合。
	\begin{enumerate}[label=(\alph*)]
		\item 证明 $R$ 是整环，其单位为 $\pm1$。
		\item 证明 $R$ 中的不可约元正是 $\pm p$（其中 $p$ 为 $\ZZ$ 中素数）以及那些在 $\QQ[x]$ 中不可约且常数项为 $\pm1$ 的多项式 $f(x)$。证明这些不可约元在 $R$ 中都是素元。
		\item 证明 $x$ 不能写成 $R$ 中不可约元的乘积（特别地，$x$ 不是不可约元），并推出 $R$ 不是 U.F.D.。
		\item 证明 $x$ 在 $R$ 中不是素元，并描述商环 $R/(x)$。
	\end{enumerate}
	
	\item 设 $R=\ZZ+x\QQ[x]\subseteq\QQ[x]$ 为前一题中的环。
	\begin{enumerate}[label=(\alph*)]
		\item 设 $f(x),g(x)\in\QQ[x]$ 为两个非零有理系数多项式，$x^r$ 是在 $\QQ[x]$ 中同时整除 $f(x)$ 与 $g(x)$ 的最高幂（即 $r$ 是 $f(x)$ 或 $g(x)$ 中出现的最低次数项的次数）。设 $f_r,g_r$ 分别为 $f(x),g(x)$ 中 $x^r$ 的系数（按定义至少一个非零）。于是 $\ZZ f_r+\ZZ g_r=\ZZ d_r$，其中 $d_r\in\QQ$ 非零（参见第 2.4 节练习 14）。证明存在 $d(x)\in\QQ[x]$，它是 $f(x),g(x)$ 在 $\QQ[x]$ 中的最大公因子，并且其最低次数项为 $d_r x^r$。
		\item 证明 $f(x)=d(x)q_1(x)$ 且 $g(x)=d(x)q_2(x)$，其中 $q_1(x),q_2(x)$ 都属于 $\QQ[x]$ 的子环 $R$。
		\item 证明
		\[
		d(x)=a(x)f(x)+b(x)g(x)
		\]
		其中 $a(x),b(x)\in R$。\emph{提示：}在 Euclidean 整环 $\QQ[x]$ 中 $a(x),b(x)$ 的存在性是直接的；用第 9.2 节练习 11 说明它们可取在 $R$ 中。
		\item 由 (a) 与 (b) 推出在 $\QQ[x]$ 中
		\[
		Rf(x)+Rg(x)=Rd(x),
		\]
		并用它证明 $R$ 是 Bezout 整环（参见第 8.2 节练习 7）。
		\item 证明 (d)、前一题结果以及第 8.3 节练习 11 蕴含 $R$ 必含有非主理想（因而含有非有限生成理想）。事实上证明 $I=x\QQ[x]$ 是 $R$ 的一个非有限生成理想。
	\end{enumerate}
\end{enumerate}

\section{不可约性判别法}

若 $R$ 是 U.F.D.，则任意多变量多项式环仍是 U.F.D.。因此我们关心怎样判定 $R[x]$ 中的多项式是否不可约。一般地，判断高次数或多变量多项式是否分解并不容易。本节给出若干常用判别法。主要讨论一元多项式，系数环通常是 U.F.D.；通过 Gauss 引理，可转化到分式域上讨论。

\noindent\textbf{命题 9.} 设 $F$ 是域，$p(x)\in F[x]$。则 $p(x)$ 有一次因子，当且仅当 $p(x)$ 在 $F$ 中有根，即存在 $a\in F$ 使 $p(a)=0$。

\begin{proof}
	若有一次因子，可取为 $x-a$，显然 $p(a)=0$。反之，若 $p(a)=0$，由 $F[x]$ 中除法算法
	\[
	p(x)=q(x)(x-a)+r,
	\]
	其中 $r$ 为常数。代入 $x=a$ 得 $r=0$，故 $x-a$ 整除 $p(x)$。
\end{proof}

\noindent\textbf{命题 10.} 域 $F$ 上二次或三次多项式可约，当且仅当它在 $F$ 中有根。

\begin{proof}
	次数为 $2$ 或 $3$ 的多项式若可约，必有一个一次因子。结论由命题 9 得到。
\end{proof}

\noindent\textbf{命题 11（有理根判别法）.} 设
\[
p(x)=a_nx^n+a_{n-1}x^{n-1}+\cdots+a_0\in\ZZ[x].
\]
若既约分数 $r/s\in\QQ$ 是 $p(x)$ 的根，则
\[
r\mid a_0,
\qquad s\mid a_n.
\]
特别地，若 $p(x)$ 是首一整数系数多项式，且对所有整除常数项的整数 $d$ 都有 $p(d)\neq0$，则 $p(x)$ 没有有理根。

\begin{proof}
	由 $p(r/s)=0$，乘以 $s^n$ 得
	\[
	a_nr^n+a_{n-1}r^{n-1}s+\cdots+a_0s^n=0.
	\]
	移项可见 $s\mid a_nr^n$，因 $(r,s)=1$，所以 $s\mid a_n$。同理移项解出 $a_0s^n$ 可得 $r\mid a_0$。
\end{proof}

\paragraph{例}
\begin{enumerate}[label=(\arabic*)]
	\item $x^3-3x-1$ 在 $\QQ[x]$ 中不可约。因为次数为 $3$，只需检验有理根；候选根只有 $\pm1$，代入均非零。
	\item 若 $p$ 为素数，则 $x^2-p$ 与 $x^3-p$ 在 $\QQ[x]$ 中不可约。
	\item $x^2+1$ 在 $\FF_2[x]$ 中可约，因为 $1$ 是根，且 $x^2+1=(x+1)^2$。
	\item $x^2+x+1$ 在 $\FF_2[x]$ 中不可约，因为
\[
0^2+0+1=1,
\qquad
1^2+1+1=1.
\]
	\item 类似地，多项式 $x^3+x+1$ 在 $\FF_2[x]$ 中不可约。
\end{enumerate}

上述根判别只适合低次数；四次多项式可能分解为两个二次不可约因子的乘积而没有一次因子。另一类常用方法是把系数模某个理想化。

\noindent\textbf{命题 12.} 设 $I$ 是整环 $R$ 的真理想，$p(x)\in R[x]$ 是非常数首一多项式。若 $p(x)$ 在 $(R/I)[x]$ 中的像不能分解为两个较低次数多项式的乘积，则 $p(x)$ 在 $R[x]$ 中不可约。

\begin{proof}
	若 $p$ 在 $R[x]$ 中可约，则由于 $p$ 是首一多项式，可写成两个首一的、次数较低的非常数多项式 $a(x)b(x)$ 的乘积。把这个等式的系数模 $I$ 化，便在 $(R/I)[x]$ 中给出 $p(x)$ 的像的一个分解。由于两个因子仍是首一的，它们在商环中的次数与原来相同，特别仍为正次数并且都低于 $p(x)$ 的次数。这与假设矛盾。
\end{proof}

该命题给出一种检验方法：找到一个真理想，使约化后的多项式不可约，则原多项式不可约。但反过来不成立；例如 $x^2+1$ 在 $\ZZ[x]$ 中不可约，却在模 $2$ 下可约。还有不可约多项式模每个素数都可约的例子。

\paragraph{例}
\begin{enumerate}[label=(\arabic*)]
	\item $x^2+x+1$ 与 $x^3+x+1$ 在 $\ZZ[x]$ 中不可约，因为它们模 $2$ 后在 $\FF_2[x]$ 中不可约。
	\item $x^2+1$ 在 $\ZZ[x]$ 中不可约，因为模 $3$ 后无根；但它模 $2$ 可约。
	\item 多变量中也可用此法，但要注意某些非单位在商环中会变成单位。例如 $xy+x+y+1=(x+1)(y+1)$，但模 $(x)$ 或模 $(y)$ 后可能“看起来”不可约。
\end{enumerate}

\noindent\textbf{命题 13（Eisenstein 判别法）.} 设 $P$ 是整环 $R$ 的素理想，令
\[
f(x)=x^n+a_{n-1}x^{n-1}+\cdots+a_1x+a_0\in R[x],\qquad n\geq1.
\]
若
\[
a_{n-1},\ldots,a_1,a_0\in P,
\qquad a_0\notin P^2,
\]
则 $f(x)$ 在 $R[x]$ 中不可约。

\begin{proof}
	设 $f(x)=a(x)b(x)$ 是 $R[x]$ 中的非常数分解。因为 $f$ 是首一的，我们可以把 $a(x)$ 与 $b(x)$ 都取为首一多项式。将系数模 $P$ 化后，假设给出
	\[
	x^n=\overline{a}(x)\overline{b}(x)
	\]
	在 $(R/P)[x]$ 中成立。由于 $P$ 是素理想，$R/P$ 是整环，所以 $(R/P)[x]$ 中因式分解的次数相加，且 $x^n$ 的两个首一因子只能是幂 $x^r$ 与 $x^{n-r}$。特别地，$\overline{a}(x)$ 与 $\overline{b}(x)$ 的常数项都为 $0$。这就是说 $a(x),b(x)$ 的常数项都属于 $P$。于是 $f$ 的常数项 $a_0$ 是这两个常数项的乘积，因而属于 $P^2$，这与 $a_0\notin P^2$ 矛盾。
\end{proof}

\noindent\textbf{推论 14（$\ZZ[x]$ 中的 Eisenstein 判别法）.} 设 $p$ 是整数素数，
\[
f(x)=x^n+a_{n-1}x^{n-1}+\cdots+a_1x+a_0\in\ZZ[x].
\]
若 $p\mid a_i$ 对 $0\leq i\leq n-1$ 成立，但 $p^2\nmid a_0$，则 $f(x)$ 在 $\ZZ[x]$ 与 $\QQ[x]$ 中都不可约。

\begin{proof}
	这是命题 13 应用于素理想 $(p)$ 的情形，再结合 Gauss 引理。
\end{proof}

\paragraph{例}
\begin{enumerate}[label=(\arabic*)]
	\item $x^4+10x+5$ 由素数 $5$ 的 Eisenstein 判别法不可约。
	\item 若整数 $a$ 被素数 $p$ 整除但不被 $p^2$ 整除，则 $x^n-a$ 在 $\ZZ[x]$ 中不可约。特别地，$x^n-p$ 对所有 $n\geq1$ 不可约。
	\item $x^4+1$ 不能直接用 Eisenstein 判别法，但
	\[
	(x+1)^4+1=x^4+4x^3+6x^2+4x+2
	\]
	对素数 $2$ 满足 Eisenstein 条件，故 $x^4+1$ 不可约。
	\item 对素数 $p$，分圆多项式
	\[
	\Phi_p(x)=\frac{x^p-1}{x-1}=x^{p-1}+x^{p-2}+\cdots+x+1
	\]
	也可通过替换 $x\mapsto x+1$ 后用 Eisenstein 判别法证明在 $\ZZ[x]$ 中不可约。
\end{enumerate}


\paragraph{练习}
\begin{enumerate}[label=\arabic*.]
	\item 判断下列多项式在所给环中是否不可约。若可约，给出其不可约分解。记 $\FF_p$ 为有限域 $\ZZ/p\ZZ$，其中 $p$ 为素数。
	\begin{enumerate}[label=(\alph*)]
		\item $x^2+x+1$ 在 $\FF_2[x]$ 中。
		\item $x^3+x+1$ 在 $\FF_3[x]$ 中。
		\item $x^4+1$ 在 $\FF_5[x]$ 中。
		\item $x^4+10x^2+1$ 在 $\ZZ[x]$ 中。
	\end{enumerate}
	
	\item 证明下列多项式在 $\ZZ[x]$ 中不可约：
	\begin{enumerate}[label=(\alph*)]
		\item $x^4-4x^3+6$。
		\item $x^6+30x^5-15x^3+6x-120$。
		\item $x^4+4x^3+6x^2+2x+1$。\emph{提示：}把 $x-1$ 代入 $x$。
		\item $\displaystyle { (x+2)^p-2^p\over x}$，其中 $p$ 是奇素数。
	\end{enumerate}
	
	\item 对所有 $n\geq1$，证明多项式
	\[
	(x-1)(x-2)\cdots(x-n)-1
	\]
	在 $\ZZ[x]$ 上不可约。\emph{提示：}若该多项式分解，考察因子在 $x=1,2,\ldots,n$ 处的取值。
	
	\item 对所有 $n\geq1$ 且 $n\neq4$，证明多项式
	\[
	(x-1)(x-2)\cdots(x-n)+1
	\]
	在 $\ZZ[x]$ 上不可约。
	
	\item 找出 $\FF_2[x]$ 中次数不超过 $3$ 的所有首一不可约多项式；再在 $\FF_3[x]$ 中做同样的事。
	
	\item 构造下列各阶的域：
	\[
	\text{(a) }9,\quad \text{(b) }49,\quad \text{(c) }8,\quad \text{(d) }81.
	\]
	可以把它们写成 $F[x]/(f(x))$ 的形式，其中 $F$ 是某个域、$f$ 是某个多项式。\emph{提示：}使用第 9.2 节练习 2 与练习 3。
	
	\item 证明 $\RR[x]/(x^2+1)$ 是一个域，并且同构于复数域 $\CC$。
	
	\item 证明
	\[
	K_1=\FF_{11}[x]/(x^2+1),
	\qquad K_2=\FF_{11}[y]/(y^2+2y+2)
	\]
	都是有 $121$ 个元素的域。证明把 $K_1$ 中元素 $p(\overline{x})$ 送到 $K_2$ 中元素 $p(\overline{y}+1)$ 的映射（其中 $p$ 为 $\FF_{11}$ 系数多项式）是良定义的，并给出从 $K_1$ 到 $K_2$ 的环同构（从而是域同构）。
	
	\item 证明多项式 $x^2-\sqrt2$ 在 $\ZZ[\sqrt2]$ 上不可约（可用 $\ZZ[\sqrt2]$ 是 U.F.D. 这一事实；参见第 8.1 节练习 9）。
	
	\item 证明多项式
	\[
	p(x)=x^4-4x^2+8x+2
	\]
	在二次域
	\[
	F=\QQ(\sqrt{-2})=\{a+b\sqrt{-2}\mid a,b\in\QQ\}
	\]
	上不可约。\emph{提示：}先在唯一分解整环 $\ZZ[\sqrt{-2}]$ 中使用命题 11 的方法（参见第 8.1 节练习 8），证明若 $a\in\ZZ[\sqrt{-2}]$ 是 $p(x)$ 的根，则 $a$ 是 $2$ 在 $\ZZ[\sqrt{-2}]$ 中的因子。推出 $a$ 必为 $\pm1,\pm\sqrt{-2}$ 或 $\pm2$，从而 $p(x)$ 在 $F$ 上没有一次因子。类似地证明 $p(x)$ 不能分解为两个 $F$ 系数二次多项式的乘积。
	
	\item 证明 $x^2+y^2-1$ 在 $\QQ[x,y]$ 中不可约。
	
	\item 证明
	\[
	x^{n-1}+x^{n-2}+\cdots+x+1
	\]
	在 $\ZZ[x]$ 上不可约，当且仅当 $n$ 是素数。
	
	\item 证明：对所有整数 $n\neq1,-3,-5$，多项式 $x^3+nx+2$ 在 $\ZZ[x]$ 上不可约。
	
	\item 把两个多项式 $x^8-1$ 与 $x^6-1$ 分别在下列环上分解为不可约因子乘积：
	\[
	\text{(a) }\ZZ,
	\quad \text{(b) }\ZZ/2\ZZ,
	\quad \text{(c) }\ZZ/3\ZZ.
	\]
	
	\item 设 $F$ 是域。证明作为形式幂级数环 $F[[x]]$ 中的系数多项式，$X^n-x$ 在 $F[[x]][X]$ 中不可约。\emph{提示：}回忆 $F[[x]]$ 是 Euclidean 整环；参见第 7.2 节练习 5 与第 8.1 节例 4。
	
	\item 设 $F$ 是域，$f(x)$ 是 $F[x]$ 中次数为 $n$ 的多项式。多项式
	\[
	g(x)=x^n f(1/x)
	\]
	称为 $f(x)$ 的反多项式。
	\begin{enumerate}[label=(\alph*)]
		\item 用 $f$ 的系数描述 $g$ 的系数。
		\item 证明 $f$ 不可约当且仅当 $g$ 不可约。
	\end{enumerate}
	
	\item 证明 Eisenstein 判别法的以下变形：设 $P$ 是唯一分解整环 $R$ 的素理想，令
	\[
	f(x)=a_nx^n+a_{n-1}x^{n-1}+\cdots+a_1x+a_0\in R[x],\qquad n\geq1.
	\]
	假设 $a_n\notin P$，$a_{n-1},\ldots,a_0\in P$，且 $a_0\notin P^2$。证明 $f(x)$ 在 $F[x]$ 中不可约，其中 $F$ 是 $R$ 的分式域。
	
	\item 证明 $6x^5+14x^3-21x+35$ 与 $18x^5-30x^2+120x+360$ 在 $\QQ[x]$ 中不可约。
	
	\item 设 $F$ 是域，
	\[
	f(x)=a_nx^n+a_{n-1}x^{n-1}+\cdots+a_0\in F[x].
	\]
	$f(x)$ 的导数定义为
	\[
	D_x(f(x))=n a_nx^{n-1}+(n-1)a_{n-1}x^{n-2}+\cdots+a_1,
	\]
	其中 $na=a+\cdots+a$（$n$ 次）。注意 $D_x(f(x))$ 仍是 $F$ 系数多项式。若存在包含 $F$ 的某个域 $E$ 以及 $a\in E$，使 $(x-a)^2$ 在 $E[x]$ 中整除 $f(x)$，则称 $f(x)$ 有重根。我们将在第 13.5 节证明，非零非常数多项式 $f(x)$ 有重根，当且仅当 $f(x)$ 与其导数不互素（这可由 $F[x]$ 中的 Euclidean 算法检测）。用这一判别法判断下列多项式是否有重根：
	\begin{enumerate}[label=(\alph*)]
		\item $x^3-3x-2\in\QQ[x]$。
		\item $x^3+3x+2\in\QQ[x]$。
		\item $x^6-4x^4+6x^3+4x^2-12x+9\in\QQ[x]$。
		\item 证明对任意素数 $p$ 与任意 $a\in\FF_p$，多项式 $x^p-a$ 都有重根。
	\end{enumerate}
	
	\item 证明 $\ZZ/6\ZZ[x]$ 中的多项式 $f(x)=x$ 可分解为
	\[
	x=(3x+4)(4x+3),
	\]
	因而不是不可约多项式。
	\begin{enumerate}[label=(\alph*)]
		\item 证明 $f(x)$ 模 $\ZZ/6\ZZ$ 的两个非平凡理想 $(2)$ 与 $(3)$ 约化后都是不可约多项式，从而说明命题 12 中 $R$ 必须是整环这一条件是必要的。
		\item 证明在 $\ZZ/6\ZZ[x]$ 中任何分解 $f(x)=g(x)h(x)$ 里，$g(x)$ 模 $(2)$ 的约化为 $1$ 或 $x$，而 $h(x)$ 模 $(2)$ 的约化相应为 $x$ 或 $1$；对模 $(3)$ 约化亦然。确定 $f(x)$ 在 $\ZZ/6\ZZ[x]$ 中的全部分解。\emph{提示：}用中国剩余定理。
		\item 证明理想 $(3,x)$ 是 $\ZZ/6\ZZ[x]$ 中的主理想。
		\item 在 $\ZZ/30\ZZ[x]$ 中证明
		\[
		f(x)=x=(10x+21)(15x+16)(6x+25).
		\]
		证明任意两个这些因子的乘积仍具有相同次数。证明 $f(x)$ 对 $\ZZ/30\ZZ$ 中任一素理想约化后都是不可约多项式。确定 $f(x)$ 在 $\ZZ/30\ZZ[x]$ 中的全部分解。\emph{提示：}考虑模 $(2),(3),(5)$ 的约化并用中国剩余定理。
		\item 把 (d) 推广到 $\ZZ/n\ZZ[x]$，其中 $n$ 是 $k$ 个不同素数的乘积。
	\end{enumerate}
\end{enumerate}

\section{域上的多项式环 II}

设 $F$ 是域。本节证明关于一元多项式环 $F[x]$ 的若干进一步结果。

\noindent\textbf{命题 15.} $F[x]$ 中的极大理想正是由不可约多项式 $f(x)$ 生成的理想 $(f(x))$。特别地，
\[
F[x]/(f(x))
\]
是域，当且仅当 $f(x)$ 不可约。

\begin{proof}
	这是第 8.2 节命题 7 对 P.I.D. $F[x]$ 的应用。
\end{proof}

\noindent\textbf{命题 16.} 设 $g(x)$ 是 $F[x]$ 中非常数多项式，并设
\[
g(x)=f_1(x)^{n_1}f_2(x)^{n_2}\cdots f_k(x)^{n_k}
\]
是其不可约分解，其中 $f_i(x)$ 两两不同。则有环同构
\[
F[x]/(g(x))
\cong
F[x]/(f_1^{n_1})\times F[x]/(f_2^{n_2})\times\cdots\times F[x]/(f_k^{n_k}).
\]

\begin{proof}
	由中国剩余定理得到。因为不同不可约多项式 $f_i,f_j$ 在 Euclidean 整环 $F[x]$ 中互素，所以理想 $(f_i^{n_i})$ 与 $(f_j^{n_j})$ 互素。
\end{proof}

一个根 $a$ 对应一次因子 $x-a$。若 $f(x)$ 被 $(x-a)^m$ 整除但不被 $(x-a)^{m+1}$ 整除，则称 $a$ 是重数为 $m$ 的根。

\noindent\textbf{命题 17.} 若多项式 $f(x)$ 在 $F$ 中有根 $a_1,a_2,\ldots,a_k$（可重复计算），则
\[
(x-a_1)(x-a_2)\cdots(x-a_k)
\]
整除 $f(x)$。特别地，域上一元 $n$ 次多项式在该域中至多有 $n$ 个根，即使按重数计算也是如此。

\begin{proof}
	第一部分由命题 9 归纳得到。由于一次因子不可约，第二部分由 $F[x]$ 的唯一分解性得到。
\end{proof}

\noindent\textbf{命题 18.} 域的乘法群的任一有限子群都是循环群。特别地，若 $F$ 是有限域，则 $F^\times$ 是循环群。

\begin{proof}
	用有限生成阿贝尔群基本定理。设有限子群同构于
	\[
	\ZZ/n_1\ZZ\times\ZZ/n_2\ZZ\times\cdots\times\ZZ/n_k\ZZ,
	\qquad n_k\mid n_{k-1}\mid\cdots\mid n_1.
	\]
	每个循环因子中都有 $n_k$ 个元素满足 $x^{n_k}=1$。若 $k>1$，直积中满足 $x^{n_k}=1$ 的元素个数会超过 $n_k$，但域中多项式 $x^{n_k}-1$ 至多有 $n_k$ 个根，矛盾。因此 $k=1$，该群循环。
\end{proof}

\noindent\textbf{推论 19.} 若 $p$ 是素数，则 $(\ZZ/p\ZZ)^\times$ 是循环群。

\begin{proof}
	这是有限域 $\ZZ/p\ZZ$ 的非零元素乘法群。
\end{proof}

\noindent\textbf{推论 20.} 设 $n\geq2$，其素因子分解为
\[
n=p_1^{a_1}p_2^{a_2}\cdots p_r^{a_r}.
\]
则：
\begin{enumerate}[label=(\arabic*)]
	\item
	\[
	(\ZZ/n\ZZ)^\times\cong
	(\ZZ/p_1^{a_1}\ZZ)^\times\times\cdots\times(\ZZ/p_r^{a_r}\ZZ)^\times.
	\]
	\item 对 $a\geq2$，$(\ZZ/2^a\ZZ)^\times$ 是一个 $2$ 阶循环群和一个 $2^{a-2}$ 阶循环群的直积。
	\item 对奇素数 $p$，$(\ZZ/p^a\ZZ)^\times$ 是阶为 $p^{a-1}(p-1)$ 的循环群。
\end{enumerate}

\begin{remark}
这些同构描述了 $n$ 阶循环群 $Z_n$ 的自同构群的群论结构，因为 $\operatorname{Aut}(Z_n)\cong(\ZZ/n\ZZ)^\times$（参见第 4.4 节命题 16）。特别地，对素数 $p$，$p$ 阶循环群的自同构群是 $p-1$ 阶循环群。
\end{remark}

\begin{proof}
	这主要是汇集前面的结果。(1) 中的同构由中国剩余定理得到（见第 7.6 节推论 18）。(2) 中的同构直接来自第 2.3 节练习 22 与练习 23。

	现在设 $p$ 为奇素数。群 $(\ZZ/p^a\ZZ)^\times$ 是阶为 $p^{a-1}(p-1)$ 的阿贝尔群。由第 2.3 节练习 21，这个群的 Sylow $p$-子群是循环的。模 $p$ 约化映射
\[
\varphi: \ZZ/p^a\ZZ\longrightarrow \ZZ/p\ZZ
\]
是环同态，它在单位群上给出满同态
\[
(\ZZ/p^a\ZZ)^\times\longrightarrow (\ZZ/p\ZZ)^\times.
\]
后一个群由推论 19 是 $p-1$ 阶循环群。此映射的核的阶为 $p^{a-1}$，因此对每个素数 $q\neq p$，$(\ZZ/p^a\ZZ)^\times$ 的 Sylow $q$-子群同构地映入循环群 $(\ZZ/p\ZZ)^\times$。于是 $(\ZZ/p^a\ZZ)^\times$ 的所有 Sylow 子群都是循环的，从而 (3) 成立，证明完毕。
\end{proof}


\paragraph{练习}
\begin{enumerate}[label=\arabic*.]
	\item 设 $F$ 是域，$f(x)$ 是 $F[x]$ 中的非常数多项式。用 $f(x)$ 的因式分解描述 $F[x]/(f(x))$ 的 nilradical（参见第 7.3 节练习 29）。
	
	\item 对第 9.4 节练习 6 中构造出的每个域，给出其非零元素乘法群的一个生成元。
	
	\item 设 $p$ 是 $\ZZ$ 中的奇素数，$n$ 是正整数。证明 $x^n-p$ 在 $\ZZ[i]$ 上不可约。\emph{提示：}使用第 8 章命题 18 与 Eisenstein 判别法。
	
	\item 证明 $x^3+12x^2+18x+6$ 在 $\ZZ[i]$ 上不可约。\emph{提示：}使用第 8 章命题 18 与 Eisenstein 判别法。
	
	\item 令 $\varphi$ 表示 Euler $\varphi$-函数。证明恒等式
	\[
	\sum_{d\mid n}\varphi(d)=n,
	\]
	其中求和遍历 $n$ 的所有正因子。\emph{提示：}先观察当 $n=p^m$ 为素数幂时，求和会望远镜相消。写 $n=p^m n'$，其中 $p\nmid n'$。证明
	\[
	\sum_{d\mid n}\varphi(d)=
	\left(\sum_{d''\mid p^m}\varphi(d'')\right)
	\left(\sum_{d'\mid n'}\varphi(d')\right)
	\]
	可由右边展开并使用 $a,b$ 互素时 $\varphi(ab)=\varphi(a)\varphi(b)$ 得到。用归纳完成证明。也可令 $Z$ 为 $n$ 阶循环群，并证明由于 $Z$ 对每个 $d\mid n$ 有唯一的 $d$ 阶子群，$Z$ 中阶为 $d$ 的元素个数为 $\varphi(d)$；于是 $|Z|$ 是所有 $d\mid n$ 的 $\varphi(d)$ 之和。
	
	\item 设 $G$ 是域 $F$ 的乘法群 $F^\times$ 的一个有限子群，阶为 $n$。令 $\varphi$ 为 Euler $\varphi$-函数，并令 $\psi(d)$ 表示 $G$ 中阶为 $d$ 的元素个数。证明对 $n$ 的每个因子 $d$，都有 $\psi(d)=\varphi(d)$。特别地推出 $\psi(n)\geq1$，所以 $G$ 是循环群。\emph{提示：}注意对任意整数 $N\geq1$，多项式 $x^N-1$ 在 $F$ 中至多有 $N$ 个根。推出对任意 $N$ 有 $\sum_{d\mid N}\psi(d)\leq N$。由前一题 $\sum_{d\mid N}\varphi(d)=N$，用归纳证明每个 $d\mid n$ 都有 $\psi(d)\leq\varphi(d)$。再由
	\[
	\sum_{d\mid n}\psi(d)=n=\sum_{d\mid n}\varphi(d)
	\]
	推出每个 $d\mid n$ 都有 $\psi(d)=\varphi(d)$。
	
	\item 证明域的加法群与乘法群从不同构。\emph{提示：}分三种情形考虑：$F$ 有限；$-1\neq1$；以及 $-1=1$。
\end{enumerate}

\section{域上的多变量多项式与 Grobner 基}

本节考虑多变量多项式，介绍一些基本计算工具并指出若干应用。本节结果在第 10 到第 14 章中不需要；更多应用将在第 15 章给出。

我们在第 2 节证明了域 $F$ 上一元多项式环 $F[x]$ 是 Euclidean 整环；推论 8 又说明 $F[x_1,\ldots,x_n]$ 是 U.F.D.。不过由第 8.2 节推论 8 可知，除非 $n=1$，这个环不是 P.I.D.。下面第一个结果说明，这些多项式环中的理想虽然未必是主理想，却总是有限生成的。具有这种性质的一般环有一个专门名称。

\noindent\textbf{定义.} 带单位元的交换环 $R$ 称为 Noether 环，如果 $R$ 的每个理想都是有限生成的。

Noether 环将在第 15 章与第 16 章中更详细地研究。本节发展若干基本理论以及由此得到的算法，用来处理 $F[x_1,\ldots,x_n]$ 中的（有限生成）理想。

正如第 1 节所见，一个 $n$ 元多项式环可看成一个一元多项式环，其系数环是 $n-1$ 元多项式环。循着这个归纳观点，正如我们在定理 7 与推论 8 中所做的那样，可以从下面更一般的结果推出 $F[x_1,x_2,\ldots,x_n]$ 是 Noether 环。

\noindent\textbf{定理 21（Hilbert 基定理）.} 若 $R$ 是 Noether 环，则 $R[x]$ 也是 Noether 环。

\begin{proof}
	设 $I$ 是 $R[x]$ 的理想，令 $L$ 为 $I$ 中所有元素的首系数组成的集合。先证明 $L$ 是 $R$ 的理想。零多项式属于 $I$，所以 $0\in L$。若 $f=ax^d+\cdots$ 与 $g=bx^e+\cdots$ 是 $I$ 中次数分别为 $d,e$、首系数分别为 $a,b\in R$ 的多项式，则对任意 $r\in R$，$ra-b$ 要么为零，要么是多项式 $r x^e f-x^d g$ 的首系数；而这个多项式属于 $I$。因此 $ra-b\in L$，说明 $L$ 是 $R$ 的理想。由于 $R$ 是 Noether 环，$L$ 有有限生成元，设为 $a_1,a_2,\ldots,a_n\in R$。对每个 $i=1,\ldots,n$，取 $f_i\in I$，使其首系数为 $a_i$。设 $e_i=\deg f_i$，并令 $N=\max\{e_1,e_2,\ldots,e_n\}$。

	对每个 $d\in\{0,1,\ldots,N-1\}$，令 $L_d$ 为 $I$ 中次数为 $d$ 的多项式的首系数再加上 $0$ 构成的集合。与 $L$ 的情形相同，每个 $L_d$ 也是 $R$ 的理想，故有限生成。对每个非零 $L_d$，取生成元 $b_{d,1},b_{d,2},\ldots,b_{d,n_d}\in R$，并取 $f_{d,i}\in I$ 为次数 $d$、首系数 $b_{d,i}$ 的多项式。

	我们证明多项式 $f_1,\ldots,f_n$ 连同所有非零 $L_d$ 所对应的多项式 $f_{d,i}$ 生成 $I$，即
\[
I=(\{f_1,\ldots,f_n\}\cup\{f_{d,i}\mid 0\leq d<N,
1\leq i\leq n_d\}).
\]
记右边理想为 $I'$。按构造 $I'\subseteq I$。若 $I'\neq I$，则可取 $f\in I-I'$ 为非零且次数最小的多项式。令 $d=\deg f$，$a$ 为 $f$ 的首系数。

	若 $d\geq N$，因为 $a\in L$，可写成 $a=r_1a_1+\cdots+r_na_n$。于是
\[
g=r_1x^{d-e_1}f_1+\cdots+r_nx^{d-e_n}f_n
\]
属于 $I'$，并且与 $f$ 有相同次数 $d$ 与相同首系数 $a$。故 $f-g\in I$ 的次数小于 $f$。由 $f$ 的最小性，必有 $f-g\in I'$，从而 $f\in I'$，矛盾。

	若 $d<N$，则 $a\in L_d$，可写成 $a=r_1b_{d,1}+\cdots+r_{n_d}b_{d,n_d}$。令 $g=r_1f_{d,1}+\cdots+r_{n_d}f_{d,n_d}$，则 $g\in I'$，并与 $f$ 有相同次数和首系数，于是同样得到矛盾。因此 $I=I'$ 有限生成。由于 $I$ 任意，$R[x]$ 是 Noether 环。
\end{proof}

\noindent\textbf{推论 22.} 若 $F$ 是域，则 $F[x_1,x_2,\ldots,x_n]$ 中每个理想都是有限生成的。

\begin{proof}
	域显然 Noether，多次应用 Hilbert 基定理即可。
\end{proof}

若 $I$ 由一个可能无限的多项式集合 $S$ 生成，则推论 22 说明 $I$ 有限生成；事实上只需要 $S$ 中有限多个多项式即可生成 $I$（参见练习 1）。

正如 Hilbert 基定理的证明所显示的，多项式环 $R[x]$ 中一个理想的所有多项式首系数组成的集合是 $R$ 中一个很有用的理想，可用来理解原理想。这启发我们更一般地研究 $F[x_1,x_2,\ldots,x_n]$ 中的“首项”。为此需要先在单项式上指定一个全序；否则一般不能说一个多项式的哪一项是“首项”。在推论 22 的归纳证明中，我们隐含地选择了这种序：先把多项式看作关于 $x_1$ 的多项式，其系数在 $F[x_2,\ldots,x_n]$ 中；再把首系数看成关于 $x_2$ 的多项式，依此类推。这就是词典单项式序的一个例子。

一个常用次序是词典序：先规定变量顺序，例如
\[
x_1>x_2>\cdots>x_n,
\]
并规定具有指数 $(a_1,a_2,\ldots,a_n)$ 的单项式项
\[
A x_1^{a_1}x_2^{a_2}\cdots x_n^{a_n}
\]
高于具有指数 $(b_1,b_2,\ldots,b_n)$ 的单项式项
\[
B x_1^{b_1}x_2^{b_2}\cdots x_n^{b_n},
\]
如果在两个 $n$ 元组首次出现差异的分量处有 $a_i>b_i$。这类似于字典中单词的排序方式。注意，单项式项的排序只在乘以单位（$F^\times$ 中的元素）意义下定义，并且两个单项式同时乘以同一个非零单项式不会改变其顺序。一般情形可形式化如下。

\noindent\textbf{定义.} 一个单项式序是单项式集合上的良序 $\geq$，并满足：若 $m_1\geq m_2$，则对任意单项式 $m$ 有 $mm_1\geq mm_2$。等价地，可在单项式多次数集合 $\ZZ_{\geq0}^n$ 上给出满足 $\alpha\geq\beta\Rightarrow \alpha+\gamma\geq\beta+\gamma$ 的良序。

容易证明，对任何单项式序，每个单项式 $m$ 都满足 $m\geq1$（参见练习 2）。利用 Hilbert 基定理还不难证明：若单项式上的任意全序满足每个单项式 $m\geq1$ 且 $m_1\geq m_2\Rightarrow mm_1\geq mm_2$，则它必是良序。本书的例子为简明起见主要使用直观的词典序；不过在实际计算中，其他单项式序常有重要优势，练习中还介绍了若干常用单项式序。

固定单项式序后，可定义首项。

\noindent\textbf{定义.} 固定 $F[x_1,\ldots,x_n]$ 上的单项式序。
\begin{enumerate}[label=(\arabic*)]
	\item 非零多项式 $f$ 的首项 $\operatorname{LT}(f)$ 是 $f$ 中按该序最大的单项式项；$0$ 的首项定义为 $0$。$f$ 的多次数 $\alpha(f)$ 定义为其首项的多次数。
	\item 若 $I$ 是理想，$I$ 的首项理想记为 $\operatorname{LT}(I)$，定义为
	\[
	\operatorname{LT}(I)=(\operatorname{LT}(f)\mid f\in I).
	\]
\end{enumerate}

首项依赖于单项式序。例如在 $F[x,y]$ 中，若 $x>y$，则 $\operatorname{LT}(2xy+y^3)=2xy$；若 $y>x$，则首项为 $y^3$。首项理想通常也依赖于选择的序。

由单项式生成的理想称为单项式理想。它们较容易处理：多项式属于一个单项式理想，当且仅当每个单项式项都被某个生成单项式整除。

若 $I=(f_1,\ldots,f_m)$，则
\[
(\operatorname{LT}(f_1),\ldots,\operatorname{LT}(f_m))\subseteq \operatorname{LT}(I),
\]
但包含一般可以是严格的。

\paragraph{例}
\begin{enumerate}[label=(\arabic*)]
	\item 在 $F[x,y]$ 中取词典序 $x>y$。令
\[
f_1=x^3y-xy^2+1,
\qquad f_2=x^2y^2-y^3-1,
\qquad I=(f_1,f_2).
\]
则 $\operatorname{LT}(f_1)=x^3y$（多次数为 $(3,1)$），$\operatorname{LT}(f_2)=x^2y^2$（多次数为 $(2,2)$）。因此 $\operatorname{LT}(I)$ 包含 $x^3y$ 与 $x^2y^2$，所以
\[
(x^3y,x^2y^2)\subseteq\operatorname{LT}(I).
\]
又因为
\[
yf_1-xf_2=y(x^3y-xy^2+1)-x(x^2y^2-y^3-1)=x+y,
\]
故 $g=x+y\in I$，从而首项理想还包含 $\operatorname{LT}(g)=x$。这说明 $\operatorname{LT}(I)$ 严格大于 $(\operatorname{LT}(f_1),\operatorname{LT}(f_2))=(x^3y,x^2y^2)$，因为后者中每个元素总次数至少为 $4$。稍后会看到在此情形下 $\operatorname{LT}(I)=(x,y^4)$。
	\item 对同一个理想 $I$，若使用词典序 $y>x$，则 $f_1,f_2$ 的首项分别为 $-xy^2$ 与 $-y^3$。稍后会看到此时 $\operatorname{LT}(I)=(x^4,y)$，这与使用 $x>y$ 时得到的首项理想不同，并且仍严格大于 $(\operatorname{LT}(f_1),\operatorname{LT}(f_2))$。
	\item 取 $F[x,y]$ 上任意单项式序，设 $f=f(x,y)$ 是非零多项式。主理想 $I=(f)$ 中每个元素的首项都是 $f$ 的首项的倍数，所以在这种情形下
\[
\operatorname{LT}(I)=(\operatorname{LT}(f)).
\]
\end{enumerate}

在一元情形中，首项用于除法算法，把一个多项式 $g$ 模另一个多项式 $f$ 化为唯一余式 $r$，且 $r=0$ 当且仅当 $g\in(f)$。当变量数至少为 $2$ 时，$F[x_1,\ldots,x_n]$ 不是 Euclidean 整环（因为不是 P.I.D.），情形更复杂。在上面第一个例子中，$g=x+y=yf_1-xf_2$ 属于 $I$，但 $f_1,f_2$ 都不整除 $g$；若只尝试用 $f_1,f_2$ 去“除” $g$，会得到非零余式 $g$ 本身，因而误以为 $g\notin I$。原因是 $yf_1$ 与 $xf_2$ 的首项发生了相消，而 $\operatorname{LT}(f_1)$ 与 $\operatorname{LT}(f_2)$ 本身不足以生成 $\operatorname{LT}(I)$。一个理想的生成集若其首项生成了理想中所有元素首项生成的理想，则有如下名称。

\noindent\textbf{定义.} 理想 $I\subseteq F[x_1,\ldots,x_n]$ 的一个 Grobner 基，是 $I$ 的有限生成集合 $\{g_1,\ldots,g_m\}$，满足
\[
\operatorname{LT}(I)=(\operatorname{LT}(g_1),\ldots,\operatorname{LT}(g_m)).
\]
这里的“基”并非向量空间意义下的基，而是一个特殊的生成集。

Grobner 基的重要性质是：任意多项式都可唯一写成 $I$ 中元素加一个“余式”，因此可以有效判定理想成员资格。

\subsection*{一般多项式除法}
固定单项式序，取非零多项式 $g_1,\ldots,g_m\in F[x_1,\ldots,x_n]$。给定 $f$，令商 $q_1,\ldots,q_m$ 和余式 $r$ 初始均为 $0$。反复考察当前被除式的首项：
\begin{enumerate}[label=(\roman*)]
	\item 若其首项可被某个 $\operatorname{LT}(g_i)$ 整除，设
	\[
	\operatorname{LT}(f)=a_i\operatorname{LT}(g_i),
	\]
	则把 $a_i$ 加到 $q_i$ 中，并把 $f$ 替换为 $f-a_i g_i$。
	\item 若其首项不能被任何 $\operatorname{LT}(g_i)$ 整除，则把该首项加到余式 $r$，并从 $f$ 中删去该首项。
\end{enumerate}
由于单项式序是良序，该过程终止，得到
\[
f=q_1g_1+\cdots+q_mg_m+r,
\]
其中每个 $q_i g_i$ 的多次数都不超过 $f$ 的多次数，而余式 $r$ 的任何非零项都不被 $\operatorname{LT}(g_1),\ldots,\operatorname{LT}(g_m)$ 中任何一个整除。一般地，商和余式依赖于除子的顺序。

\paragraph{例}
以下在 $F[x,y]$ 中固定词典序 $x>y$。
\begin{enumerate}[label=(\arabic*)]
	\item 设 $f=x^3y^3+3x^2y^4$，$g=xy^4$。$f$ 的首项 $x^3y^3$ 不被 $g$ 的首项整除，于是将 $x^3y^3$ 加入余式，新的被除式为 $3x^2y^4$。这时 $3x^2y^4$ 被 $xy^4$ 整除，商项为 $3x$，过程终止。得到商 $q=3x$、余式 $r=x^3y^3$，即
\[
x^3y^3+3x^2y^4=(3x)(xy^4)+x^3y^3.
\]
注意若像一元除法那样在第一步因首项不能整除就停止，则会留下整个 $f$ 作为“余式”，尽管 $f$ 的另一部分仍可被 $g$ 整除；这正是一般多项式除法中需要步骤 (ii) 的原因。
	\item 设 $f=x^2+x-y^2+y$，$g_1=xy+1$，$g_2=x+y$。按顺序除以 $g_1,g_2$，第一次迭代中 $x^2$ 不被 $xy$ 整除，但被 $x$ 整除，于是 $q_2$ 中加入 $x$，被除式变为 $-xy+x-y^2+y$。第二次迭代中首项 $-xy$ 被 $xy$ 整除，$q_1=-1$，被除式变为 $x-y^2+y+1$。第三次迭代中首项 $x$ 被 $x$ 整除，$q_2=x+1$，被除式变为 $-y^2+1$。此后 $-y^2$ 与 $1$ 都不能被 $xy$ 或 $x$ 整除，故进入余式。结果为
\[
q_1=-1,
\qquad q_2=x+1,
\qquad r=-y^2+1,
\]
并且
\[
x^2+x-y^2+y=(-1)(xy+1)+(x+1)(x+y)+(-y^2+1).
\]
	\item 仍取 $f=x^2+x-y^2+y$，但交换除子次序，令 $g_1=x+y$、$g_2=xy+1$。容易计算得到不同的商和余式：
\[
q_1=x-y+1,
\qquad q_2=0,
\qquad r=2y,
\]
这说明一般多项式除法的余式通常依赖于除子的顺序。
\end{enumerate}

\noindent\textbf{定理 23.} 固定单项式序，设 $\{g_1,\ldots,g_m\}$ 是非零理想 $I\subseteq F[x_1,\ldots,x_n]$ 的 Grobner 基。则：
\begin{enumerate}[label=(\arabic*)]
	\item 每个多项式 $f$ 可唯一写成
	\[
	f=h+r,
	\]
	其中 $h\in I$，而余式 $r$ 的任何非零单项式项都不被任何 $\operatorname{LT}(g_i)$ 整除。
	\item $h$ 与 $r$ 可由一般多项式除法计算，且与使用 $g_i$ 的顺序无关。
	\item 余式 $r$ 给出 $f$ 在商环 $F[x_1,\ldots,x_n]/I$ 中陪集的唯一代表。特别地，
	\[
	f\in I \quad\Longleftrightarrow\quad r=0.
	\]
\end{enumerate}

\begin{proof}
	一般多项式除法总给出 $f=\sum q_i g_i+r$，令 $h=\sum q_i g_i\in I$。若另有 $f=h'+r'$ 满足同样条件，则
	\[
	r-r'=h'-h\in I.
	\]
	若 $r-r'\neq0$，其首项属于 $\operatorname{LT}(I)=(\operatorname{LT}(g_1),\ldots,\operatorname{LT}(g_m))$，因此应被某个 $\operatorname{LT}(g_i)$ 整除。但 $r,r'$ 的各项都不具有这种可整除性，矛盾。故 $r=r'$，从而 $h=h'$。其余结论立即得到。
\end{proof}

\noindent\textbf{命题 24.} 固定单项式序，设 $I$ 是 $F[x_1,\ldots,x_n]$ 中非零理想。
\begin{enumerate}[label=(\arabic*)]
	\item 若 $g_1,\ldots,g_m\in I$ 满足
	\[
	\operatorname{LT}(I)=(\operatorname{LT}(g_1),\ldots,\operatorname{LT}(g_m)),
	\]
	则 $\{g_1,\ldots,g_m\}$ 是 $I$ 的 Grobner 基。
	\item $I$ 总有 Grobner 基。
\end{enumerate}

\begin{proof}
	(1) 若 $f\in I$，用一般多项式除法除以 $g_i$ 得 $f=\sum q_i g_i+r$。若 $r\neq0$，则 $r\in I$，所以 $\operatorname{LT}(r)\in\operatorname{LT}(I)$，必被某个 $\operatorname{LT}(g_i)$ 整除，与余式性质矛盾。因此 $r=0$，$g_i$ 生成 $I$。
	
	(2) 首项理想 $\operatorname{LT}(I)$ 是单项式理想；由 Hilbert 基定理，它有限生成。选择有限多个 $h_1,\ldots,h_k\in I$，使这些生成元为 $\operatorname{LT}(h_i)$。由 (1) 得到 Grobner 基。
\end{proof}

为了判定给定生成集是否是 Grobner 基，定义 $S$-多项式。若 $f_1,f_2$ 是非零多项式，令 $M$ 为 $\operatorname{LT}(f_1)$ 与 $\operatorname{LT}(f_2)$ 的首一最小公倍单项式，则
\[
S(f_1,f_2)=\frac{M}{\operatorname{LT}(f_1)}f_1-
\frac{M}{\operatorname{LT}(f_2)}f_2.
\]
它的作用是消去首项。

\noindent\textbf{引理 25.} 若 $f_1,
\ldots,f_m$ 在 $F[x_1,\ldots,x_n]$ 中有相同多次数 $\alpha$，并且常数线性组合
\[
h=a_1f_1+\cdots+a_mf_m
\]
的多次数严格小于 $\alpha$，则 $h$ 可写成相邻 $S$-多项式的常数线性组合。

\begin{proof}
	把 $f_i$ 写为 $f_i=c_i f_i'$，其中 $c_i\in F$，$f_i'$ 是多次数为 $\alpha$ 的首一多项式。于是
\[
\begin{aligned}
h={}&\sum a_i c_i f_i' \\
={}& a_1c_1(f_1'-f_2')+(a_1c_1+a_2c_2)(f_2'-f_3')+\cdots \\
&{}+(a_1c_1+\cdots+a_{m-1}c_{m-1})(f_{m-1}'-f_m')
+(a_1c_1+\cdots+a_mc_m)f_m'.
\end{aligned}
\]
注意 $f_{i-1}'-f_i'=S(f_{i-1},f_i)$。由于 $h$ 和每个 $f_{i-1}'-f_i'$ 的多次数都严格小于 $\alpha$，最后一项的系数必须为 $0$，即 $a_1c_1+\cdots+a_mc_m=0$。因此最后一项消失，上式给出所需的相邻 $S$-多项式的常数线性组合。
\end{proof}

\noindent\textbf{命题 26（Buchberger 判别法）.} 设 $R=F[x_1,\ldots,x_n]$ 并固定单项式序。若 $I=(g_1,\ldots,g_m)$ 是非零理想，则 $G=\{g_1,\ldots,g_m\}$ 是 $I$ 的 Grobner 基，当且仅当对所有 $1\leq i<j\leq m$，$S(g_i,g_j)$ 被 $G$ 做一般多项式除法所得余式为 $0$。

\begin{proof}
	若 $G$ 是 Grobner 基，则 $S(g_i,g_j)\in I$，由定理 23 余式为 $0$。
	
	反之设所有 $S$-多项式余式为 $0$。取任意 $f\in I$，在所有表示 $f=\sum h_i g_i$ 中，选取使各项 $h_i g_i$ 的最大多次数最小的一个。若 $f$ 的多次数小于该最大多次数，则最高多次数项必相消。由引理 25，这种相消由 $S$-多项式控制；而每个 $S$-多项式又可由 $G$ 表示且所有项多次数更小，矛盾于最小性。因此 $\operatorname{LT}(f)$ 必属于 $(\operatorname{LT}(g_1),\ldots,\operatorname{LT}(g_m))$。所以 $G$ 是 Grobner 基。
\end{proof}

\subsection*{Buchberger 算法}
从生成集 $G=\{g_1,\ldots,g_m\}$ 出发，计算所有 $S(g_i,g_j)$ 对 $G$ 的余式。若全部为 $0$，则由 Buchberger 判别法 $G$ 是 Grobner 基。若某个余式 $r\neq0$，则把 $r$ 加入 $G$ 并重新计算。该过程在有限步后终止，得到 Grobner 基。

若 Grobner 基中某个 $\operatorname{LT}(g_j)$ 被另一个 $\operatorname{LT}(g_i)$ 整除，则可删去 $g_j$，仍保持 Grobner 基。若每个首项首一且任一首项不被其他首项整除，称为极小 Grobner 基。极小 Grobner 基不唯一，但其元素个数和首项唯一（参见练习 15）。

\paragraph{例}
\begin{enumerate}[label=(\arabic*)]
	\item 在 $F[x,y]$ 中取词典序 $x>y$，考虑本节开头的理想 $I=(f_1,f_2)$，其中 $f_1=x^3y-xy^2+1$，$f_2=x^2y^2-y^3-1$。计算
\[
S(f_1,f_2)=yf_1-xf_2=x+y,
\]
而它除以 $\{f_1,f_2\}$ 的余式就是自身，所以 $\{f_1,f_2\}$ 不是 Grobner 基。令 $f_3=x+y$，增广生成集为 $G'=\{f_1,f_2,f_3\}$。此时 $S(f_1,f_2)=0\bmod G'$，并且
\[
S(f_1,f_3)=f_1-x^2yf_3=-x^2y^2-xy^2+1=0\bmod G',
\]
\[
S(f_2,f_3)=f_2-xyf_3=-xy^2-y^3-1=y^4-y^3-1\bmod G'.
\]
令 $f_4=y^4-y^3-1$，再增广为 $G''=\{f_1,f_2,f_3,f_4\}$。前面的零余式仍为零，且 $S(f_2,f_3)=0\bmod G''$；继续计算得
\[
S(f_1,f_4)=S(f_2,f_4)=S(f_3,f_4)=0\bmod G''.
\]
因此
\[
\{x^3y-xy^2+1,
\ x^2y^2-y^3-1,
\ x+y,
\ y^4-y^3-1\}
\]
是 $I$ 的一个 Grobner 基。于是 $\operatorname{LT}(I)$ 由这些多项式的首项生成，故
\[
\operatorname{LT}(I)=(x^3y,x^2y^2,x,y^4)=(x,y^4),
\]
与前文所说一致。实际上 $x+y$ 与 $y^4-y^3-1$ 的首项已经生成 $\operatorname{LT}(I)$，由命题 24，
\[
I=(x+y,
 y^4-y^3-1)
\]
给出关于此序的极小 Grobner 基。这比原来的描述简单得多。
	\item 对同一个理想 $I$，若使用词典序 $y>x$，则 $S(f_1,f_2)$ 产生余式 $f_3=-x-y$，接着 $S(f_1,f_3)$ 产生余式 $f_4=-x^4-x^3+1$，随后所有余式均为零。得到 Grobner 基
\[
\{x^3y-xy^2+1,
 x^2y^2-y^3-1,
 -x-y,
 -x^4-x^3+1\}.
\]
此时
\[
\operatorname{LT}(I)=(-xy^2,-y^3,-y,-x^4)=(y,x^4),
\]
且
\[
I=(x+y,
 x^4+x^3-1)
\]
给出关于该序的另一个极小 Grobner 基。
\end{enumerate}

在上面第一个例子中，$\{x+y^4-y^3+y-1,\ y^4-y^3-1\}$ 也是 $I$ 的一个极小 Grobner 基（它只是 $\{f_3+f_4,f_4\}$），所以即使固定单项式序，极小 Grobner 基也不唯一。若加强首项整除条件，可得到重要的唯一性。

\noindent\textbf{定义.} 固定单项式序。非零理想 $I$ 的 Grobner 基 $\{g_1,\ldots,g_m\}$ 称为约化 Grobner 基，如果：
\begin{enumerate}[label=(\alph*)]
	\item 每个 $\operatorname{LT}(g_i)$ 首一；
	\item 对 $j\neq i$，$g_j$ 的任何一项都不被 $\operatorname{LT}(g_i)$ 整除。
\end{enumerate}

\noindent\textbf{定理 27.} 对固定的单项式序，$F[x_1,\ldots,x_n]$ 中每个非零理想都有唯一的约化 Grobner 基。

\begin{proof}
	存在性由极小 Grobner 基中逐个把每个生成元除以其余生成元而得。若 $G=\{g_i\}$ 与 $G'=\{g_i'\}$ 是两个约化 Grobner 基，则它们有相同首项集合。重排后可设 $\operatorname{LT}(g_i)=\operatorname{LT}(g_i')$。若 $g_i-g_i'\neq0$，其首项属于 $\operatorname{LT}(I)$，故被某个公共首项整除。但约化条件排除这种可能，矛盾。故 $g_i=g_i'$。
\end{proof}

\noindent\textbf{推论 28.} 两个理想 $I,J\subseteq F[x_1,\ldots,x_n]$ 相等，当且仅当在任一固定单项式序下，它们有相同的约化 Grobner 基。

\paragraph{例}
\begin{enumerate}[label=(\arabic*)]
	\item 考虑 $F[x,y]$ 中理想 $I=(h_1,h_2,h_3)$，其中
\[
h_1=x^2+xy^5+y^4,
\quad h_2=xy^6-xy^3+y^5-y^2,
\quad h_3=xy^5-xy^2.
\]
在词典序 $x>y$ 下，计算得
\[
S(h_1,h_2)=S(h_1,h_3)=0\bmod\{h_1,h_2,h_3\},
\]
而
\[
S(h_2,h_3)=y^5-y^2\bmod\{h_1,h_2,h_3\}.
\]
令 $h_4=y^5-y^2$，则 $S(h_i,h_j)=0\bmod\{h_1,h_2,h_3,h_4\}$ 对所有 $1\leq i<j\leq4$ 成立。因此
\[
x^2+xy^5+y^4,
\quad xy^6-xy^3+y^5-y^2,
\quad xy^5-xy^2,
\quad y^5-y^2
\]
是 $I$ 的 Grobner 基。其首项为 $x^2,xy^6,xy^5,y^5$；删去不必要的生成元得到极小 Grobner 基 $\{x^2+xy^5+y^4,
 y^5-y^2\}$。再将第一个多项式除以第二个，得到余式 $x^2+xy^2+y^4$，所以 $I$ 的约化 Grobner 基为
\[
\{x^2+xy^2+y^4,
 y^5-y^2\}.
\]
	\item 考虑 $J=(h_1,h_2,h_3)$，其中
\[
h_1=xy^3+y^3+1,
\quad h_2=x^3y-x^3+1,
\quad h_3=x+y.
\]
对同一词典序 $x>y$，有 $S(h_1,h_2)=0\bmod\{h_1,h_2,h_3\}$，而 $S(h_1,h_3)=y^4-y^3-1\bmod\{h_1,h_2,h_3\}$。令 $h_4=y^4-y^3-1$，则 $\{h_1,h_2,h_3,h_4\}$ 是 Grobner 基。由首项可化为极小 Grobner 基 $\{x+y,
 y^4-y^3-1\}$，且已是约化 Grobner 基。它与命题 26 后例 1 中理想的约化 Grobner 基相同，因此
\[
(x^3y-xy^2+1,
 x^2y^2-y^3-1)
=(xy^3+y^3+1,
 x^3y-x^3+1,
 x+y),
\]
并且二者都等于理想 $(x+y,
 y^4-y^3-1)$。
\end{enumerate}

\subsection*{Grobner 基与代数方程求解：消元}
设 $S=\{f_1,\ldots,f_m\}$ 是 $n$ 个变量的多项式集合。要求解方程组
\[
f_1=0,
\quad f_2=0,
\ldots,
\quad f_m=0,
\]
等价于研究它们生成的理想 $I=(f_1,\ldots,f_m)$。若 $g_1,\ldots,g_s$ 是同一理想 $I$ 的另一组生成元，则方程组 $g_1=\cdots=g_s=0$ 有相同解集。

若 $I$ 中存在只含某些后面变量的非零多项式，就可以逐步求解。这称为消元。

\noindent\textbf{定义.} 若 $I\subseteq F[x_1,
\ldots,x_n]$ 是理想，则
\[
I_i=I\cap F[x_{i+1},\ldots,x_n]
\]
称为关于词典序 $x_1>\cdots>x_n$ 的第 $i$ 个消元理想。

\noindent\textbf{命题 29（消元）.} 设 $G=\{g_1,\ldots,g_m\}$ 是非零理想 $I\subseteq F[x_1,\ldots,x_n]$ 关于词典序 $x_1>\cdots>x_n$ 的 Grobner 基。则
\[
G\cap F[x_{i+1},\ldots,x_n]
\]
是第 $i$ 个消元理想 $I_i=I\cap F[x_{i+1},\ldots,x_n]$ 的 Grobner 基。特别地，
\[
I\cap F[x_{i+1},\ldots,x_n]=0
\]
当且仅当 $G\cap F[x_{i+1},\ldots,x_n]=0$。

\begin{proof}
	记 $G_i=G\cap F[x_{i+1},\ldots,x_n]$。显然 $G_i\subseteq I_i$。由命题 24，只需证明 $G_i$ 的首项生成 $\operatorname{LT}(I_i)$。若 $f\in I_i$，则 $f\in I$，而 $G$ 是 $I$ 的 Grobner 基，所以 $\operatorname{LT}(f)$ 可由 $G$ 的首项整除。由于 $\operatorname{LT}(f)$ 只含变量 $x_{i+1},\ldots,x_n$，在词典序下对应的 $g_j$ 也只能含这些变量，故 $g_j\in G_i$。结论成立。
\end{proof}

\paragraph{例}
\begin{enumerate}[label=(\arabic*)]
	\item 椭圆
\[
2x^2+2xy+y^2-2x-2y=0
\]
与圆
\[
x^2+y^2=1
\]
有两个交点。对理想
\[
I=(2x^2+2xy+y^2-2x-2y,
x^2+y^2-1)\subseteq\RR[x,y]
\]
按词典序 $x>y$ 求 Grobner 基以消去 $x$，得到
\[
g_1=2x+y^2+5y^3-2,
\qquad g_2=5y^4-4y^3.
\]
由 $g_2=0$ 得 $y=0$ 或 $y=4/5$；代入 $g_1=0$，得交点
\[
(1,0),
\qquad (-3/5,4/5).
\]
若改用词典序 $y>x$ 以消去 $y$，得到 Grobner 基
\[
\{y^2+x^2-1,
\ 2yx-2y+x^2-2x+1,
\ 5x^3-7x^2-x+3\}.
\]
因为 $5x^3-7x^2-x+3=(x-1)^2(5x+3)$，可知 $x=1$ 或 $x=-3/5$，再得到相同解，不过计算较繁。
	\item 上一个例子也可用初等方法求解。现在考虑 $\CC$ 上方程组
\[
x^3-2xy+y^3=0,
\qquad
x^5-2x^2y^2+y^5=0.
\]
对理想 $(x^3-2xy+y^3,
 x^5-2x^2y^2+y^5)$ 按词典序 $x>y$ 计算 Grobner 基，可得
\[
g_1=x^3-2xy+y^3,
\]
\[
g_2=200xy+193y^9+158y^8-45y^7-456y^6+50y^5-100y^4,
\]
\[
g_3=y^{10}-y^8-2y^7+2y^6.
\]
任何原方程组的解都满足 $g_1=g_2=g_3=0$。由于
\[
g_3=y^6(y-1)^2(y^2+2y+2),
\]
有 $y=0,1$ 或 $y=-1\pm i$。分别代入 $g_1,g_2$，得到全部四个解
\[
(x,y)=(0,0),
\quad (1,1),
\quad (-1+i,-1-i),
\quad (-1-i,-1+i).
\]
	\item 考虑 $\CC$ 上方程组
\[
x+y+z=1,
\qquad x^2+y^2+z^2=2,
\qquad x^3+y^3+z^3=3.
\]
关于词典序 $x>y>z$ 的约化 Grobner 基为
\[
\{x+y+z-1,
\ y^2+yz-y+z^2-z-1/2,
\ z^3-z^2-(1/2)z-1/6\}.
\]
因此 $z$ 是多项式 $t^3-t^2-(1/2)t-1/6$ 的根；由对称性，$x,y$ 也都是该多项式的根。对这个三次多项式的每个根，$y$ 有两个对应值，$x$ 随之唯一确定。所得六个解就是该三次多项式三个互异根的所有排列。
\end{enumerate}

这些例子显示，研究多项式方程组 $f_1=0,
 f_2=0,\ldots,f_m=0$ 的解，与研究它们在 $F[x_1,\ldots,x_n]$ 中生成的理想 $I=(f_1,f_2,\ldots,f_m)$ 密切相关。这一基本联系是第 15 章将介绍的“代数几何”的出发点；那里还会给出 Grobner 基的进一步应用。

\subsection*{理想交的计算}
下面说明如何计算多项式环中理想的和、积与交这些基本集合论运算。若 $I=(f_1,\ldots,f_s)$，$J=(h_1,\ldots,h_l)$ 是 $F[x_1,\ldots,x_n]$ 中两个理想，则
\[
I+J=(f_1,\ldots,f_s,h_1,\ldots,h_l),
\]
而
\[
IJ=(f_ih_j\mid 1\leq i\leq s,
1\leq j\leq l).
\]
下面的命题说明如何计算两个理想的交。

\noindent\textbf{命题 30.} 若 $I,J$ 是 $F[x_1,\ldots,x_n]$ 中任意两个理想，则 $tI+(1-t)J$ 是 $F[t,x_1,\ldots,x_n]$ 中的理想，并且
\[
I\cap J=(tI+(1-t)J)\cap F[x_1,\ldots,x_n].
\]
特别地，关于词典序 $t>x_1>\cdots>x_n$，$I\cap J$ 是 $tI+(1-t)J$ 的第一个消元理想。

\begin{proof}
	显然 $tI$ 与 $(1-t)J$ 是 $F[t,x_1,\ldots,x_n]$ 中理想，因此其和也是理想。若 $f\in I\cap J$，则
\[
f=tf+(1-t)f
\]
说明 $f\in (tI+(1-t)J)\cap F[x_1,\ldots,x_n]$。反过来，设
\[
f=tf_1+(1-t)h_1\in F[x_1,\ldots,x_n],
\qquad f_1\in I,
\quad h_1\in J.
\]
则
\[
t(f_1-h_1)=f-h_1\in F[x_1,\ldots,x_n].
\]
因为右边不含 $t$，可得 $f_1-h_1=0$，于是 $f=f_1=h_1\in I\cap J$。由等式可知，关于序 $t>x_1>\cdots>x_n$，$I\cap J$ 正是 $tI+(1-t)J$ 的第一个消元理想。
\end{proof}

若 $I=(f_1,\ldots,f_s)$、$J=(h_1,\ldots,h_l)$，则
\[
tI+(1-t)J=(tf_1,\ldots,tf_s,(1-t)h_1,\ldots,(1-t)h_l).
\]
由命题 29，对该理想按词典序 $t>x_1>\cdots>x_n$ 计算 Grobner 基，其中不含 $t$ 的元素给出 $F[x_1,\ldots,x_n]$ 中 $I\cap J$ 的一个 Grobner 基。

\paragraph{例}
若 $I=(x,y)^2=(x^2,xy,y^2)$，$J=(x)$，在 $F[t,x,y]$ 中按 $t>x>y$ 计算 $tI+(1-t)J$ 的约化 Grobner 基，可得
\[
\{tx-x,\ ty^2,\ x^2,\ xy\}.
\]
去掉含 $t$ 的元素，得
\[
I\cap J=(x^2,xy).
\]


\paragraph{练习}
\begin{enumerate}[label=\arabic*.]
	\item 设 $I$ 是 $F[x_1,\ldots,x_n]$ 中由一个（可能无限的）多项式集合 $S$ 生成的理想。证明 $S$ 中存在有限个多项式就足以生成 $I$。\emph{提示：}由定理 21 写 $I=(f_1,\ldots,f_m)$，再把每个 $f_i\in I$ 用 $S$ 中的多项式表示。
	
	\item 设 $\geq$ 是任一单项式序。
	\begin{enumerate}[label=(\alph*)]
		\item 对任意非零多项式 $f,g$，证明 $\operatorname{LT}(fg)=\operatorname{LT}(f)\operatorname{LT}(g)$，且 $\alpha(fg)=\alpha(f)+\alpha(g)$。
		\item 证明 $\alpha(f+g)\leq\max(\alpha(f),\alpha(g))$；若 $\alpha(f)\neq\alpha(g)$，则等号成立。
		\item 证明对每个单项式 $m$，都有 $m\geq1$。
		\item 证明若 $m_1$ 整除 $m_2$，则 $m_2\geq m_1$。推出一个多项式的首项不整除它的任一较低项。
	\end{enumerate}
	
	\item 证明：若 $\geq$ 是非空集合上的全序或偏序，则下列条件等价：
	\begin{enumerate}[label=(\roman*)]
		\item 每个非空子集都含有最小元。
		\item 不存在无限严格递降序列 $a_1>a_2>a_3>\cdots$（这称为递降链条件，或 D.C.C.）。
	\end{enumerate}
	推出一般多项式除法总在有限步后终止。
	
	\item 设 $\geq$ 是单项式序。对单项式 $m_1,m_2$ 定义 $m_1\geq_g m_2$，若 $\deg m_1>\deg m_2$，或 $\deg m_1=\deg m_2$ 且 $m_1\geq m_2$。
	\begin{enumerate}[label=(\alph*)]
		\item 证明 $\geq_g$ 也是单项式序。（称 $\geq_g$ 是 $\geq$ 的分次化。以次数为首要比较标准的序常称为分次序或次数序，因此本题给出构造分次序的方法。）
		\item 词典序 $x_1>\cdots>x_n$ 的分次化称为分次词典单项式序。证明在分次词典序下 $x_1^2>x_2^2x_3>x_1x_2>x_1^2>x_2^2>x_1$（按原书给出的相应单项式顺序检验），并比较其与普通词典序下的顺序。
	\end{enumerate}
	
	\item 反分次词典单项式序定义如下：先给变量 $\{x_1,x_2,\ldots,x_n\}$ 选定一个次序；对单项式 $m_1,m_2$，若 $\deg m_1>\deg m_2$，或 $\deg m_1=\deg m_2$ 且从 $x_n,x_{n-1},\ldots,x_1$ 的顺序看第一个使两者指数不同的变量在 $m_1$ 中的指数较小，则定义 $m_1\geq m_2$。
	\begin{enumerate}[label=(\alph*)]
		\item 证明反分次词典序是单项式序，并满足 $x_1>x_2>\cdots>x_n$。
		\item 证明在 $F[x_1,x_2]$ 中相对于 $\{x_1,x_2\}$ 的反分次词典序就是满足 $x_1>x_2$ 的分次词典序，但在 $F[x_1,x_2,x_3]$ 中的反分次词典序不是任何词典序的分次化。
		\item 对相对于 $\{x_1,x_2,x_3\}$ 的反分次词典序，写出并验证原书给出的单项式链
		\[
		x_1x_2^2x_3>x_1^2x_3^2>x_1^2x_2>x_2x_3^2>x_1x_2>x_2^2>x_1x_3>x_3^2>x_1>x_2.
		\]
	\end{enumerate}
	
	\item 验证在词典序 $x>y>z$ 下，原书列出的单项式顺序
	\[
	x^3y>x^3z>x^2z^2>xyz>x^2y>xz^2>y^2z^2>y^2z
	\]
	成立；并分别验证对应的分次词典序与相对于 $\{x,y,z\}$ 的反分次词典序下原书列出的两个顺序。
	
	\item 对单项式
	\[
	xz^2,\\ x^2y^2z,
	xy^2z,
	x^3y,
	x^3z^2,
	x^2,
	x^2yz^2,
	x^2z^2
	\]
	分别按词典序 $x>y>z$、对应的分次词典序、以及相对于 $\{x,y,z\}$ 的反分次词典序进行排序。
	
	\item 证明 $F[x_1,\ldots,x_n]$ 上有 $n!$ 个不同的词典单项式序。类似地证明有 $n!$ 个不同的分次词典序和 $n!$ 个不同的反分次词典序。
	
	\item 可以证明 $F[x_1,\ldots,x_n]$ 上任一单项式序都可如下得到。对 $k\geq n$，取 Euclidean 空间 $\RR^n$ 中非零向量 $v_1,
	\ldots,v_k$，两两正交：$v_i\cdot v_j=0$（$i\neq j$），并假设 $v_1$ 的所有坐标非负。定义单项式序为：$m_1>m_2$ 当且仅当对某个 $t\leq k$，有 $v_i\cdot\alpha(m_1)=v_i\cdot\alpha(m_2)$ 对 $i=1,2,\ldots,t-1$ 成立，而 $v_t\cdot\alpha(m_1)>v_t\cdot\alpha(m_2)$。
	\begin{enumerate}[label=(\alph*)]
		\item 令 $k=n$，取 $v_i=(0,\ldots,0,1,0,\ldots,0)$，其中 $1$ 在第 $i$ 个位置。证明该序给出 $x_1>x_2>\cdots>x_n$ 的词典序。
		\item 令 $k=n$，定义 $v_1=(1,1,\ldots,1)$，并令
		\[
		v_i=(1,1,\ldots,1,-n+i-1,0,\ldots,0),
		\qquad 2\leq i\leq n,
		\]
		其中末尾有 $i-2$ 个零。证明该序给出相对于 $\{x_1,\ldots,x_n\}$ 的分次词典序。
	\end{enumerate}
	
	\item 设 $I$ 是由单项式 $m_1,\ldots,m_k$ 生成的单项式理想。证明多项式 $f\in F[x_1,\ldots,x_n]$ 属于 $I$，当且仅当 $f$ 的每个单项式项都被某个 $m_i$ 整除。\emph{提示：}对 $a_1,
	\ldots,a_k\in F[x_1,
	\ldots,x_n]$ 展开 $a_1m_1+\cdots+a_km_k$，注意每个单项式项都被某个 $m_i$ 整除。证明 $x^2yz+3xy^2\in J=(xyz,y^2)\subseteq F[x,y,z]$，但它不属于 $I'=(xz^2,y^2)$。
	
	\item 固定 $R=F[x_1,\ldots,x_n]$ 上的单项式序，并设 $\{g_1,\ldots,g_m\}$ 是理想 $I$ 的 Grobner 基。证明 $h\in\operatorname{LT}(I)$ 当且仅当 $h$ 是若干单项式项的和，且每项都被某个 $\operatorname{LT}(g_i)$ 整除。\emph{提示：}使用前一题。
	
	\item 设 $I$ 是由单项式 $g_1,\ldots,g_m$ 生成的单项式理想。用前一题直接证明 $\{g_1,\ldots,g_m\}$ 是 $I$ 的 Grobner 基。
	
	\item 设 $I$ 是由单项式 $g_1,\ldots,g_m$ 生成的单项式理想。用 Buchberger 判别法证明 $\{g_1,\ldots,g_m\}$ 是 $I$ 的 Grobner 基。
	
	\item 设 $I$ 是 $R=F[x_1,\ldots,x_n]$ 中的单项式理想，且 $\{m_1,
	\ldots,m_k\}$ 是生成 $I$ 的极小单项式集合，即每个 $m_i$ 是单项式且没有真子集生成 $I$。证明这些 $m_i$（$1\leq i\leq k$）唯一。\emph{提示：}使用练习 10。
	
	\item 固定 $R=F[x_1,
	\ldots,x_n]$ 上的单项式序。
	\begin{enumerate}[label=(\alph*)]
		\item 证明 $\{g_1,\ldots,g_m\}$ 是理想 $I$ 的极小 Grobner 基，当且仅当 $\{\operatorname{LT}(g_1),\ldots,\operatorname{LT}(g_m)\}$ 是 $\operatorname{LT}(I)$ 的极小生成集。
		\item 证明 $I$ 的极小 Grobner 基的首项唯一确定，且 $I$ 的任意两个极小 Grobner 基的元素个数相同。\emph{提示：}使用 (a) 与前一题。
	\end{enumerate}
	
	\item 固定 $F[x_1,
	\ldots,x_n]$ 上的单项式序，并设 $G=\{g_1,
	\ldots,g_m\}$ 是非零理想 $I$ 的一组生成元。证明若 $S(g_i,g_j)$ 除以 $G$ 的余式不为零，则理想
	\[
	(\operatorname{LT}(g_1),\ldots,\operatorname{LT}(g_m),\operatorname{LT}(S(g_i,g_j)))
	\]
	严格大于理想 $(\operatorname{LT}(g_1),\ldots,\operatorname{LT}(g_m))$。由此推出命题 26 后描述的计算 Grobner 基的算法在有限步后终止。\emph{提示：}使用练习 1。
	
	\item 固定 $F[x,y]$ 上的词典序 $x>y$。用 Buchberger 判别法证明 $\{x^2y-y^2,\ x^3-xy\}$ 是理想 $I=(x^2y-y^2,x^3-xy)$ 的 Grobner 基。
	
	\item 对 $F[x,y]$ 中由 $x>y$ 定义的词典序，证明 $\{x-y^3,\ y^5-y^6\}$ 是理想 $I=(x-y^3,-x^2+xy^2)$ 的约化 Grobner 基。
	
	\item 固定 $F[x,y]$ 上的词典序 $x>y$。
	\begin{enumerate}[label=(\alph*)]
		\item 证明 $\{x^3-y,\ x^2y-y^2,\ xy^2-y^2,\ y^3-y^2\}$ 是理想 $I=(-x^3+y,x^2y-y^2)$ 的约化 Grobner 基。
		\item 判断多项式 $f=x^6-x^5y$ 是否属于理想 $I$。
	\end{enumerate}
	
	\item 固定 $F[x,y,z]$ 上的词典序 $x>y>z$。证明
	\[
	\{x^2+xy+z,\ xyz+z^2,\ xz^2,\ z^3\}
	\]
	是理想 $I=(x^2+xy+z,xyz+z^2)$ 的约化 Grobner 基；特别地，推出首项理想 $\operatorname{LT}(I)$ 需要四个生成元。
	
	\item 固定 $F[x,y]$ 上的词典序 $x>y$。用 Buchberger 判别法证明 $\{x^2y-y^2,x^3-xy\}$ 是理想 $I=(x^2y-y^2,x^3-xy)$ 的 Grobner 基。
	
	\item 设 $I=(x^2-y,x^2y-z)\subseteq F[x,y,z]$。
	\begin{enumerate}[label=(\alph*)]
		\item 证明对于词典序 $x>y>z$，$\{x^2-y,y^2-z\}$ 是 $I$ 的约化 Grobner 基。
		\item 证明对于词典序 $z>x>y$，$\{x^2-y,z-y^2\}$ 是 $I$ 的约化 Grobner 基（注意这与 (a) 中的多项式本质上相同）。
		\item 证明对于词典序 $z>y>x$，$\{y-x^2,z-x^4\}$ 是 $I$ 的约化 Grobner 基。
	\end{enumerate}
	
	\item 证明 $F[x,y]$ 中理想
	\[
	I=(x^2y+xy^2-2y,
	x^2+xy-x+y^2-2y,
	xy^2-x-y+y^3)
	\]
	与
	\[
	J=(x-y^2,xy-y,x^2-y)
	\]
	相等。
	
	\item 用约化 Grobner 基证明 $F[x,y,z]$ 中理想
	\[
	I=(x^3-yz,
	yz+y)
	\]
	与
	\[
	J=(x^3z+x^3,
	x^3+y)
	\]
	相等。
	
	\item 证明对理想
	\[
	I=(x^2+xy^2,
	x^2-y^3,
	y^3-y^2)
	\]
	按词典序 $x>y$ 的约化 Grobner 基为
	\[
	\{x^2-y^2,
	y^3-y^2,
	xy^2+y^2\}.
	\]
	
	\item 证明理想 $I=(xy+y^2,x^2y+xy^2+x^2)$ 的约化 Grobner 基：对于词典序 $x>y$ 为
	\[
	\{x^2,xy+y^2,y^3\},
	\]
	而对于词典序 $y>x$ 为
	\[
	\{y^2+yx,x^2\}.
	\]
	
	\item 对推论 28 后例子中的理想 $I$，按分次词典单项式序证明
	\[
	\{x^3-y^3,
	x^2+xy^2+y^4,
	x^2y+xy^3+y^2\}
	\]
	是它的约化 Grobner 基。（注意：虽然这给出 $I$ 的三个生成元，而非例中词典序下的两个，但次数更低。）
	
	\item 设
	\[
	I=(x^4-y^4+z^3-1,
	x^3+y^2+z^2-1).
	\]
	证明：对词典序 $x>y>z$，$I$ 的约化 Grobner 基有五个元素（五个生成元中的最高次数为 $12$，最多单项式项数为 $35$）；对词典序 $y>z>x$，有两个元素（最高次数为 $6$，最多项数为 $8$）；而对于反分次词典单项式序，约化 Grobner 基为
	\[
	\{x^3+y^2+z^2-1,
	xy^2+xz^2-x+y^4-z^3+1\}.
	\]
	
	\item 在 $\CC$ 上求解方程组
	\[
	x^2-yz=3,
	\qquad y^2-xz=4,
	\qquad z^2-xy=5.
	\]
	
	\item 对词典序 $x>y$，求理想
	\[
	I=(x^2+xy+y^2-1,
	x^2+4y^2-4)
	\]
	的一个 Grobner 基，并用它求实平面 $\RR^2$ 中椭圆 $x^2+xy+y^2=1$ 与 $x^2+4y^2=4$ 的四个交点。
	
	\item 用 Grobner 基求复数域 $\CC$ 上方程组
	\[
	2x^3+2x^2y^2+3y^3=0,
	\qquad
	3x^5+2x^3y^3+2y^5=0
	\]
	的全部六个解。
	
	\item 用 Grobner 基证明 $F[x,y,z]$ 中
	\[
	(x,z)\cap(y^2,x-yz)=(xy,x-yz).
	\]
	
	\item 用 Grobner 基计算 $F[x,y]$ 中两个理想
	\[
	(x^3y-xy^2+1,
	x^2y^2-y^3-1)
	\]
	与
	\[
	(x^2-y^2,
	x^3+y^3)
	\]
	的交。
	
	\item 以下四题处理环 $R$ 中两个理想 $I,J$ 的理想商。定义两个理想 $I,J$ 的理想商为
	\[
	(I:J)=\{r\in R\mid rJ\subseteq I\}.
	\]
	\begin{enumerate}[label=(\alph*)]
		\item 设 $R$ 是整环，$0\neq f\in R$，$I$ 是 $R$ 的理想。证明：若 $\{g_1,
		\ldots,g_s\}$ 生成理想 $I\cap(f)$，则 $\{g_1/f,
		\ldots,g_s/f\}$ 生成理想商 $(I:(f))$。
		\item 若 $I$ 是交换环 $R$ 的理想，且 $f_1,
		\ldots,f_s\in R$，证明理想商
		\[
		(I:(f_1,
		\ldots,f_s))
		\]
		等于理想
		\[
		\bigcap_{i=1}^s (I:(f_i)).
		\]
	\end{enumerate}
	
	\item 若
	\[
	I=(x^2y+z^3,
	x+y^3-z,
	2y^4z-yz^2-z^3)
	\]
	且
	\[
	J=(x^2y^5,
	x^3z^4,
	y^3z^7)
	\]
	为 $\QQ[x,y,z]$ 中理想，证明
	\[
	(I:J)=(z^2,
	y+z,
	x-z).
	\]
	\emph{提示：}使用前一题与命题 30。
	
	\item 设 $K$ 是 $R$ 中的理想，$I$ 是包含 $K$ 的理想，$J$ 是任一理想。若 $\overline I,\overline J$ 分别表示 $I,J$ 在商环 $R/K$ 中的像，证明
	\[
	(\overline I:\overline J)=\overline{(I:J)},
	\]
	其中右边为理想商 $(I:J)$ 在 $R/K$ 中的像。
	
	\item 令 $K=(y^5-z^4)\subseteq R=\QQ[y,z]$。对下列每一对理想 $I,J$，使用前两题和命题 30，在环 $R/K$ 中验证理想商 $(\overline I:\overline J)$：
	\begin{enumerate}[label=(\roman*)]
		\item $I=(y^3,y^5-z^4)$，$J=(z)$，则 $(\overline I:\overline J)=(\overline{y}^3,
		\overline{z}^3)$。
		\item $I=(y^3,z,y^5-z^4)$，$J=(y)$，则 $(\overline I:\overline J)=(\overline{y}^2,
		\overline{z})$。
		\item $I=(y,y^3,z,y^5-z^4)$，$J=(1)$，则 $(\overline I:\overline J)=(\overline y,
		\overline z)$。
	\end{enumerate}
	
	\item 练习 38--44 给出单项式理想的若干初等性质。由 Hilbert 基定理可知理想有限生成，但这些题不必假设有限生成；有限生成与无限生成时证明相同。这些练习可用于独立证明 Hilbert 基定理（练习 44）。以下 $M,N$ 为单项式理想，分别由单项式集合 $\{m_i\mid i\in I\}$ 与 $\{n_j\mid j\in J\}$ 生成。证明两个单项式理想的和与积仍是单项式理想，并证明
	\[
	M+N=(m_i,n_j\mid i\in I,
	j\in J),
	\qquad MN=(m_i n_j\mid i\in I,
	j\in J).
	\]
	
	\item 证明：若 $\{M_s\mid s\in S\}$ 是任意非空的单项式理想族，并且在包含关系下全序，则
	\[
	\bigcup_{s\in S} M_s
	\]
	是单项式理想。（特别地，任意递增单项式理想列的并仍是单项式理想；参见第 7.3 节练习 19。）
	
	\item 证明两个单项式理想的交仍是单项式理想，并证明
	\[
	M\cap N=(e_{ij}\mid i\in I,
	j\in J),
	\]
	其中 $e_{ij}$ 是 $m_i$ 与 $n_j$ 的最小公倍单项式。\emph{提示：}使用练习 10。
	
	\item 证明：对任意单项式 $n$，理想商
	\[
	(M:(n))=(m_i/d_i\mid i\in I),
	\]
	其中 $d_i$ 是 $m_i$ 与 $n$ 的最大公因单项式（参见练习 34）。证明若 $N$ 有限生成，则两个单项式理想的理想商 $(M:N)$ 仍是单项式理想。
	
	\item 
	\begin{enumerate}[label=(\alph*)]
		\item 证明 $M$ 是单项式素理想，当且仅当 $M=(S)$，其中 $S$ 是 $\{x_1,x_2,\ldots,x_n\}$ 的某个子集。（特别地，单项式素理想只有有限多个，且每个都是有限生成的。）
		\item 证明 $(x_1,
		\ldots,x_n)$ 是唯一的单项式极大理想。
	\end{enumerate}
	
	\item （Dickson 引理——Hilbert 基定理的一个特殊情形）按以下方法证明 $F[x_1,
	\ldots,x_n]$ 中每个单项式理想都是有限生成的。令
	\[
	\mathcal S=\{N\mid N\text{ 是一个非有限生成的单项式理想}\},
	\]
	并反设 $\mathcal S\neq\varnothing$。
	\begin{enumerate}[label=(\alph*)]
		\item 证明 $\mathcal S$ 含有一个极大元素 $M$。\emph{提示：}使用练习 39 与 Zorn 引理。
		\item 证明存在单项式 $x,y\notin M$，使 $xy\in M$。\emph{提示：}使用练习 42(a)。
		\item 对 (b) 中的 $x$，证明 $M$ 含有一个有限生成的单项式理想 $M_0$，使
		\[
		M_0+(x)=M+(x),
		\qquad
		M=M_0+(x)(M:(x)),
		\]
		其中 $(M:(x))$ 是练习 41 中定义的单项式理想，而 $(x)(M:(x))$ 是这两个理想的积。推出 $M$ 有限生成，矛盾，从而证明 $\mathcal S=\varnothing$。\emph{提示：}使用 $M$ 的极大性和前面练习。
	\end{enumerate}
	
	\item 若 $I$ 是 $F[x_1,
	\ldots,x_n]$ 中非零理想，用 Dickson 引理证明 $\operatorname{LT}(I)$ 有限生成。推出 $I$ 有 Grobner 基，并推出 Hilbert 基定理。参见命题 24。
	
	\item （图的 $n$-着色）一个大小为 $N$ 的有限图 $\mathcal G$ 由顶点 $i\in\{1,2,
	\ldots,N\}$ 与若干边 $(i,j)$ 组成。一种 $n$-着色是把 $n$ 种颜色之一赋给每个顶点，使由边相连的顶点颜色不同。设 $F$ 是至少含 $n$ 个元素的域。为每个顶点 $i$ 引入变量 $x_i$，并选取 $F$ 中 $n$ 个不同元素组成集合 $S$ 代表 $n$ 种颜色。则 $\mathcal G$ 的 $n$-着色等价于为每个 $i=1,2,
	\ldots,N$ 赋值 $x_i=a_i$，其中 $a_i\in S$，且若 $(i,j)$ 是边，则 $a_i\neq a_j$。令
	\[
	f(x)=\prod_{a\in S}(x-a)
	\]
	为 $F[x]$ 中以 $S$ 中元素为根的 $n$ 次多项式，则 $x_i=a_i\in S$ 等价于方程 $f(x_i)=0$。条件 $a_i\neq a_j$ 可用
	\[
	g(x_i,x_j)=0,
	\qquad
	g(x_i,x_j)={f(x_i)-f(x_j)\over x_i-x_j}
	\]
	表示，其中 $g\in F[x_i,x_j]$。因此求 $\mathcal G$ 的 $n$-着色等价于求解方程组
	\[
	f(x_i)=0 \quad (i=1,2,
	\ldots,N),
	\qquad
	g(x_i,x_j)=0 \quad ((i,j)\text{ 为边}).
	\]
	（也可以用任意满足 $a_i\neq a_j\Rightarrow g(a_i,a_j)=0$ 的多项式 $g$。）由 Hilbert 零点定理（参见第 15.3 节推论 33）可知，该方程组有解，因而 $\mathcal G$ 有 $n$-着色，除非由这些多项式生成的 $F[x_1,
	\ldots,x_N]$ 中理想不是真理想。这又等价于该理想关于任意单项式序的约化 Grobner 基仅为 $\{1\}$。而当 $n$-着色存在时，像命题 29 后例子那样解这个方程组可给出显式着色。
	
	可选的域 $F$ 与集合 $S$ 有很多。例如，若 $F$ 含有 $n$ 个不同的 $n$ 次单位根，可取 $f(x)=x^n-1$，并可取
	\[
	g(x_i,x_j)={x_i^n-x_j^n\over x_i-x_j}
	=x_i^{n-1}+x_i^{n-2}x_j+\cdots+x_i x_j^{n-2}+x_j^{n-1};
	\]
	或者取 $F=\FF_p$ 中任意 $n$ 个元素的子集，其中 $p\geq n$ 为素数。特别地，当 $n=p$ 时，由 Fermat 小定理可取 $f(x)=x^p-x$，并取
	\[
	g(x_i,x_j)=(x_i-x_j)^{p-1}-1.
	\]
	\begin{enumerate}[label=(\alph*)]
		\item 考虑有 8 个顶点、14 条边
		\[
		(1,3),(1,4),(1,5),(2,4),(2,7),(2,8),(3,4),(3,6),(3,8),(4,5),(5,6),(6,7),(6,8),(7,8)
		\]
		的图 $\mathcal G$ 的可能 3-着色。取 $F=\FF_3$，颜色为 $0,1,2\in\FF_3$，并假设顶点 1 颜色为 $0$。此时
		\[
		f(x)=x(x-1)(x-2)=x^3-x\in\FF_3[x],
		\qquad g(x_i,x_j)=x_i^2+x_i x_j+x_j^2-1.
		\]
		若 $I$ 是由 $x_1$、$x_i^3-x_i$（$2\leq i\leq8$）以及所有边 $(i,j)$ 对应的 $g(x_i,x_j)$ 生成的理想，证明关于词典序 $x_1>x_2>\cdots>x_8$，$I$ 的约化 Grobner 基为
		\[
		\{x_1,
		x_2,
		x_3+x_8,
		x_4+2x_8,
		x_5+x_8,
		x_6,
		x_7+x_8,
		x_8^2+2\}.
		\]
		推出 $\mathcal G$ 有两个不同的 3-着色，它们由顶点 8 的颜色决定（该颜色必须为 $\FF_3$ 中非零元素），并写出这些着色。再证明若向 $\mathcal G$ 加入边 $(3,7)$，则该图不能 3-着色。
		\item 取 $F=\FF_5$，四种颜色为 $1,2,3,4\in\FF_5$，于是 $f(x)=x^4-1$，并可取
		\[
		g(x_i,x_j)=x_i^3+x_i^2x_j+x_i x_j^2+x_j^3.
		\]
		证明有 5 个顶点和 9 条边
		\[
		(1,3),(1,4),(1,5),(2,3),(2,4),(2,5),(3,4),(3,5),(4,5)
		\]
		的图（即去掉一条边的五顶点完全图）可以 4-着色，但不能 3-着色。
		\item 用 Grobner 基证明：有 9 个顶点和 22 条边
		\[
		(1,4),(1,6),(1,7),(1,8),(2,3),(2,4),(2,6),(2,7),(3,5),(3,7),(3,9),(4,5),(4,6),(4,7),(4,9),
		\]
		\[
		(5,6),(5,7),(5,8),(5,9),(6,7),(6,9),(7,8)
		\]
		的图 $\mathcal G$ 在颜色置换意义下恰有四种 4-着色（所以总共有 $96$ 种 4-着色）。证明若加入边 $(1,5)$，则 $\mathcal G$ 不能 4-着色。
	\end{enumerate}
\end{enumerate}

% \setchapterintro{本部分研究称为\emph{模}的数学对象。模的使用由 Emmy Noether 等数学家所推动，它展示了这种结构的力量与优雅。向量空间只是模的特殊情形：当底环为域时，模就是向量空间。若 $R$ 是环，$R$-模 $M$ 的定义与群作用十分相似，其中 $R$ 扮演“作用者”的角色，$M$ 扮演被作用集合的角色；但模的公理还要求 $M$ 本身有阿贝尔群结构。模是环的“表示对象”：环通过作用在模上来体现其结构。随着理论展开，我们会看到环 $R$ 的理想结构与其模的结构相互反映，类似于群的正规子群结构如何反映在置换表示中。}
%\part{模与向量空间}
% 

\chapter{模论导论}




\section{基本定义与例子}

我们从模的定义开始。


\noindent\textbf{定义.} 设 $R$ 是一个环（不要求交换，也不要求有 $1$）。一个\emph{左 $R$-模}，或称 $R$ 上的左模，是一个集合 $M$，连同：
\begin{enumerate}[label=(\arabic*)]
	\item 一个二元运算 $+$，使 $M$ 在 $+$ 下成为阿贝尔群；
	\item 一个 $R$ 在 $M$ 上的作用，即映射 $R\times M\to M$，记为 $(r,m)\mapsto rm$，满足对所有 $r,s\in R$ 与 $m,n\in M$，
	\begin{enumerate}[label=(\alph*)]
		\item $(r+s)m=rm+sm$；
		\item $(rs)m=r(sm)$；
		\item $r(m+n)=rm+rn$。
	\end{enumerate}
\end{enumerate}
若 $R$ 有单位元 $1$，则还要求
\[
1m=m,\qquad m\in M.
\]
“左”表示环元素写在模元素左边。右 $R$-模可类似定义。若 $R$ 交换，左 $R$-模可通过 $mr=rm$ 变为右 $R$-模；若 $R$ 不交换，此定义一般不满足相应公理。除非另有说明，本书中的“模”均指左模，并且在有单位元的环上均指幺模。

当 $R$ 是域 $F$ 时，$R$-模的公理正是 $F$ 上向量空间的公理，因此
\[
F\text{-模} \quad\text{与}\quad F\text{ 上向量空间}
\]
是同一个概念。


\noindent\textbf{定义.} 设 $M$ 是 $R$-模。$M$ 的 $R$-子模是 $M$ 的加法子群 $N$，且对 $R$ 的作用封闭，即
\[
rn\in N,\qquad r\in R,\, n\in N.
\]
因此子模就是在限制运算下本身仍为模的子集。当 $R=F$ 是域时，子模就是向量子空间。每个 $R$-模 $M$ 都有两个显然的子模：$M$ 与 $0$。

\paragraph{例}
\begin{enumerate}[label=(\arabic*)]
	\item 任意环 $R$ 本身是左 $R$-模，作用就是环乘法。此时 $R$ 作为自身左模的子模正是 $R$ 的左理想；作为自身右模的子模正是右理想。若 $R$ 不交换，两种模结构可能不同。
	\item 若 $R=F$ 是域，则 $F^n$ 在逐分量加法与数乘下是 $F$-模，也就是 $F$ 上 $n$ 维向量空间。
	\item 若 $R$ 有 $1$，则
	\[
	R^n=\{(a_1,\ldots,a_n)\mid a_i\in R\}
	\]
	在逐分量加法和左乘下是 $R$-模，称为 $R$ 上秩为 $n$ 的自由模。
	\item 同一个阿贝尔群可有许多不同的 $R$-模结构。若 $M$ 是 $R$-模，$S$ 是 $R$ 的含同一单位元的子环，则 $M$ 自动也是 $S$-模。例如 $\RR$ 同时是 $\RR$-模、$\QQ$-模与 $\ZZ$-模。
	\item 若 $I$ 是 $R$ 的双边理想，且 $am=0$ 对所有 $a\in I$、$m\in M$ 成立，则称 $M$ 被 $I$ 零化。此时 $M$ 可成为 $R/I$-模，定义
	\[
	(r+I)m=rm.
	\]
	若 $R$ 交换且 $I$ 极大，则 $R/I$ 是域，于是 $M$ 是 $R/I$ 上的向量空间。
\end{enumerate}

\paragraph{$\ZZ$-模}
设 $A$ 是任意阿贝尔群，记加法为 $+$。对 $n\in\ZZ$ 与 $a\in A$ 定义
\[
na=\begin{cases}
	a+a+\cdots+a & n>0,\\
	0 & n=0,\\
	-a-a-\cdots-a & n<0.
\end{cases}
\]
这使 $A$ 成为 $\ZZ$-模，并且这是唯一可能的幺 $\ZZ$-模结构。因此
\[
\ZZ\text{-模} \quad\text{就是}\quad \text{阿贝尔群},
\]
而 $\ZZ$-子模就是子群。若 $A$ 是有限阿贝尔群且 $|A|=m$，则 $mx=0$ 对所有 $x\in A$ 成立，所以 $A$ 也是 $\ZZ/m\ZZ$-模。特别地，若 $p$ 为素数且 $px=0$ 对所有 $x\in A$ 成立，则 $A$ 是 $\FF_p$ 上的向量空间；这正是初等阿贝尔 $p$-群的情形。

\paragraph{$F[x]$-模}
设 $F$ 是域，$V$ 是 $F$ 上向量空间，$T:V\to V$ 是线性变换。令 $x$ 在 $V$ 上作用为 $T$，并令
\[
p(x)v=(a_nT^n+a_{n-1}T^{n-1}+\cdots+a_1T+a_0I)(v)
\]
其中 $p(x)=a_nx^n+\cdots+a_0\in F[x]$。这样 $V$ 成为一个 $F[x]$-模。反过来，任一 $F[x]$-模给出一个 $F$-向量空间 $V$，以及由 $x$ 的作用确定的线性变换 $T$。因此有对应
\[
\{F[x]\text{-模}\}
\quad\longleftrightarrow\quad
\{(V,T):V\text{ 是 }F\text{-向量空间},\,T:V\to V\text{ 线性}\}.
\]
在此对应下，$F[x]$-子模正是 $T$-稳定（或 $T$-不变）子空间。这个例子将在第 12 章中用于研究矩阵的有理标准形与 Jordan 标准形。


\noindent\textbf{命题 1（子模判别法）.} 设 $M$ 是 $R$-模。子集 $N\subseteq M$ 是子模当且仅当
\begin{enumerate}[label=(\arabic*)]
	\item $N\neq\emptyset$；
	\item 对所有 $r\in R$ 与 $x,y\in N$，有 $x+ry\in N$。
\end{enumerate}

\begin{proof}
	若 $N$ 是子模，则显然满足两条。反过来，取 $r=-1$，由 $x-y\in N$ 得 $N$ 是加法子群；再取 $x=0$，得 $ry\in N$。故 $N$ 对加法与 $R$ 的作用封闭。
\end{proof}


\noindent\textbf{定义.} 设 $R$ 是带单位元的交换环。一个 $R$-代数是一个带单位元的环 $A$，连同环同态
\[
f:R\longrightarrow A,
\]
把 $1_R$ 映到 $1_A$，并且 $f(R)$ 包含于 $A$ 的中心。

若 $A$ 是 $R$-代数，则 $A$ 自然成为左右 $R$-模，定义
\[
r\cdot a=a\cdot r=f(r)a.
\]
若 $A,B$ 是 $R$-代数，则 $R$-代数同态是保持单位元且满足
\[
\varphi(r\cdot a)=r\cdot\varphi(a)
\]
的环同态。

\paragraph{例}
任意带单位元的环都是 $\ZZ$-代数；若 $R$ 是 $A$ 的中心中含 $1_A$ 的子环，则 $A$ 是 $R$-代数。特别地，$R[x]$、任意多变量多项式环以及群环 $RG$ 都是 $R$-代数。若 $R=F$ 是域，则说 $A$ 是 $F$-代数等价于说 $A$ 的中心中含有一个与 $F$ 同构的子域，且单位元相同。


\paragraph{练习}
在以下练习中，$R$ 是带单位元的环，$M$ 是左 $R$-模。
\begin{enumerate}[label=\arabic*.]
	\item 证明对所有 $m\in M$，有 $0m=0$ 且 $(-1)m=-m$。
	\item 证明 $R^\times$ 与 $M$ 满足第 1.7 节中关于乘法群 $R^\times$ 在集合 $M$ 上作用的两个公理。
	\item 设 $rm=0$，其中 $r\in R$、$m\in M$ 且 $m\neq0$。证明 $r$ 没有左逆，即不存在 $s\in R$ 使 $sr=1$。
	\item 令 $M$ 为例 3 中的模 $R^n$，并令 $I_1,I_2,\ldots,I_n$ 为 $R$ 的左理想。证明下列集合是 $M$ 的子模：
	\begin{enumerate}[label=(\alph*)]
		\item $\{(x_1,x_2,\ldots,x_n)\mid x_i\in I_i\}$；
		\item $\{(x_1,x_2,\ldots,x_n)\mid x_i\in R\text{ 且 }x_1+x_2+\cdots+x_n=0\}$。
	\end{enumerate}
	\item 对 $R$ 的任意左理想 $I$，定义
	\[
	IM=\left\{\sum_{\text{有限}} a_i m_i\mid a_i\in I,\;m_i\in M\right\},
	\]
	即所有形如 $am$（$a\in I,m\in M$）的元素的有限和构成的集合。证明 $IM$ 是 $M$ 的子模。
	\item 证明任意非空子模族的交仍是子模。
	\item 设 $N_1\subseteq N_2\subseteq\cdots$ 是 $M$ 的子模升链。证明 $\bigcup_{i=1}^\infty N_i$ 是 $M$ 的子模。
	\item 若 $R$-模 $M$ 的元素 $m$ 满足 $rm=0$ 对某个非零 $r\in R$ 成立，则称 $m$ 为\emph{挠元素}。挠元素集合记为
	\[
	\Tor(M)=\{m\in M\mid rm=0\text{ 对某个非零 }r\in R\text{ 成立}\}.
	\]
	\begin{enumerate}[label=(\alph*)]
		\item 证明若 $R$ 是整环，则 $\Tor(M)$ 是 $M$ 的子模，称为 $M$ 的\emph{挠子模}。
		\item 给出一个环 $R$ 与一个 $R$-模 $M$ 的例子，使得 $\Tor(M)$ 不是子模。[考虑 $R$-模 $R$ 中的挠元素。]
		\item 若 $R$ 有零因子，证明每个非零 $R$-模都有非零挠元素。
	\end{enumerate}
	\item 若 $N$ 是 $M$ 的子模，定义 $N$ 在 $R$ 中的零化子为
	\[
	\{r\in R\mid rn=0\text{ 对所有 }n\in N\}.
	\]
	证明它是 $R$ 的双边理想。
	\item 若 $I$ 是 $R$ 的右理想，定义 $I$ 在 $M$ 中的零化子为
	\[
	\{m\in M\mid am=0\text{ 对所有 }a\in I\}.
	\]
	证明它是 $M$ 的子模。
	\item 设 $M$ 是阿贝尔群（即 $\ZZ$-模）
	\[
	M=\ZZ/24\ZZ\times \ZZ/15\ZZ\times \ZZ/50\ZZ.
	\]
	\begin{enumerate}[label=(\alph*)]
		\item 求 $M$ 在 $\ZZ$ 中的零化子，即给出这个主理想的一个生成元。
		\item 令 $I=2\ZZ$。把 $I$ 在 $M$ 中的零化子描述为循环群的直积。
	\end{enumerate}
	\item 在前面练习的记号下，证明关于零化子的下列事实：
	\begin{enumerate}[label=(\alph*)]
		\item 设 $N$ 是 $M$ 的子模，$I$ 是它在 $R$ 中的零化子。证明 $I$ 在 $M$ 中的零化子包含 $N$。给出一个例子说明 $I$ 在 $M$ 中的零化子不必等于 $N$。
		\item 设 $I$ 是 $R$ 的右理想，$N$ 是它在 $M$ 中的零化子。证明 $N$ 在 $R$ 中的零化子包含 $I$。给出一个例子说明 $N$ 在 $R$ 中的零化子不必等于 $I$。
	\end{enumerate}
	\item 设 $I$ 是 $R$ 的理想。令 $M'$ 为 $M$ 中所有被 $I$ 的某个幂 $I^k$ 零化的元素 $a$ 的集合，其中 $k$ 可依 $a$ 而变。证明 $M'$ 是 $M$ 的子模。[使用练习 7。]
	\item 设 $z$ 是 $R$ 的中心中的元素，即 $zr=rz$ 对所有 $r\in R$ 成立。证明
	\[
	zM=\{zm\mid m\in M\}
	\]
	是 $M$ 的子模。又设 $R$ 是域上的 $2\times2$ 矩阵环，$e$ 是 $(1,1)$ 位置为 $1$、其余位置为 $0$ 的矩阵。证明当 $M=R$ 作为左 $R$-模看待时，$eR$ 不是左 $R$-子模；在这个例子中矩阵 $e$ 不在 $R$ 的中心。
	\item 若 $M$ 是有限阿贝尔群，则 $M$ 自然是 $\ZZ$-模。这个作用能否扩张为使 $M$ 成为 $\QQ$-模的作用？
	\item 证明在 $F[x]$-模的例子中，对移位算子所描述的子空间 $U_k$ 是该 $T$ 下所有的 $F[x]$-子模。
	\item 设 $T$ 是向量空间 $V$ 上的移位算子，$e_1,\ldots,e_n$ 是 $F[x]$-模例子中描述的通常基向量。若 $m\leq n$，求
	\[
	(a_mx^m+a_{m-1}x^{m-1}+\cdots+a_0)e_n.
	\]
	\item 令 $F=\RR$，$V=\RR^2$，并令 $T:V\to V$ 为绕原点顺时针旋转 $\pi/2$ 弧度的线性变换。证明对此 $T$，$V$ 与 $0$ 是仅有的 $F[x]$-子模。
	\item 令 $F=\RR$，$V=\RR^2$，并令 $T:V\to V$ 为到 $y$ 轴的投影。证明对此 $T$，$V$、$0$、$x$ 轴与 $y$ 轴是仅有的 $F[x]$-子模。
	\item 令 $F=\RR$，$V=\RR^2$，并令 $T:V\to V$ 为绕原点顺时针旋转 $\pi$ 弧度的线性变换。证明对此 $T$，$V$ 的每个子空间都是 $F[x]$-子模。
	\item 设 $n\in\ZZ^+$ 且 $n>1$，令 $R$ 为域 $F$ 上的 $n\times n$ 矩阵环。令 $M$ 为第一列中元素任意、其余列全为零的 $n\times n$ 矩阵集合。证明当 $R$ 作为自身上的左模看待时，$M$ 是 $R$ 的子模；但当 $R$ 作为右 $R$-模看待时，$M$ 不是 $R$ 的子模。
	\item 设 $A$ 是带单位元 $1_A$ 的环，同时是一个幺左 $R$-模，并满足
	\[
	r\cdot(ab)=(r\cdot a)b=a(r\cdot b)
	\]
	对所有 $r\in R$ 与 $a,b\in A$ 成立。证明映射 $f:R\to A$，$f(r)=r\cdot1_A$，是环同态，把 $1_R$ 映到 $1_A$，且 $f(R)$ 包含于 $A$ 的中心。推出 $A$ 是一个 $R$-代数，并且由这个代数结构诱导的 $R$-模结构正是原来的 $R$-模结构。
	\item 令 $A$ 为直积环 $\CC\times\CC$。设 $\tau_1$ 是 $\CC$ 上的恒等映射，$\tau_2$ 是复共轭。对任意一对 $p,q\in\{1,2\}$（可相同），定义
	\[
	f_{p,q}:\CC\longrightarrow \CC\times\CC,
	\qquad f_{p,q}(z)=(\tau_p(z),\tau_q(z)).
	\]
	例如 $f_{2,1}:z\mapsto(\bar z,z)$。
	\begin{enumerate}[label=(\alph*)]
		\item 证明每个 $f_{p,q}$ 都是单的环同态，并且它们在子域 $\RR\subseteq\CC$ 上都相同。推出 $A$ 有四种不同的 $\CC$-代数结构。分别明确给出复数 $z$ 在有序对 $(u,v)\in A$ 上的作用 $z\cdot(u,v)$。
		\item 证明若 $f_{p,q}\neq f_{p',q'}$，则 $A$ 上恒等映射不是从通过 $f_{p,q}$ 看作 $\CC$-代数的 $A$ 到通过 $f_{p',q'}$ 看作 $\CC$-代数的 $A$ 的 $\CC$-代数同态；尽管它是 $\RR$-代数同构。
		\item 证明对任意 $p,q$，存在 $A$ 到自身的一个环同构，使得通过 $f_{p,q}$ 看作 $\CC$-代数的 $A$ 与通过 $f_{1,1}$ 看作 $\CC$-代数的 $A$ 同构，后者是 $A$ 的“自然” $\CC$-代数结构。
	\end{enumerate}
\end{enumerate}
\noindent\textbf{说明.} 在前一题中，$A=\CC\times\CC$ 不是关于任一直接因子分量中的 $\CC$ 副本的 $\CC$-代数，例如子环 $\CC\times0\cong\CC$，因为 $A$ 不是这些 $\CC$ 副本上的幺模；这些子环的 $1$ 与 $A$ 的 $1$ 不同。


\section{商模与模同态}

本节给出商模与模同态的基本理论。


\noindent\textbf{定义.} 设 $M,N$ 是 $R$-模。
\begin{enumerate}[label=(\arabic*)]
	\item 映射 $\varphi:M\to N$ 称为 $R$-模同态，若
	\[
	\varphi(x+y)=\varphi(x)+\varphi(y),
	\qquad
	\varphi(rx)=r\varphi(x)
	\]
	对所有 $x,y\in M$ 与 $r\in R$ 成立。
	\item 若 $R$-模同态既单又满，则称为 $R$-模同构。若存在同构，记为 $M\cong N$。
	\item 核与像照常定义为
	\[
	\ker\varphi=\{m\in M\mid \varphi(m)=0\},
	\qquad
	\varphi(M)=\{\varphi(m)\mid m\in M\}.
	\]
	\item $\Hom_R(M,N)$ 表示所有 $R$-模同态 $M\to N$ 的集合。
\end{enumerate}
任何 $R$-模同态都是加法群同态，但反之不一定成立。核与像都是子模。

\paragraph{例}
\begin{enumerate}[label=(\arabic*)]
	\item 当 $R=\ZZ$ 时，$\ZZ$-模同态正是阿贝尔群同态。
	\item 若 $R$ 是域，则 $R$-模同态就是线性变换。
	\item 若 $I$ 零化 $M$ 与 $N$，则 $R$-模同态 $M\to N$ 自动是 $R/I$-模同态。
\end{enumerate}


\noindent\textbf{命题 2.} 设 $L,M,N$ 是 $R$-模。
\begin{enumerate}[label=(\arabic*)]
	\item 映射 $\varphi:M\to N$ 是 $R$-模同态，当且仅当
	\[
	\varphi(rx+y)=r\varphi(x)+\varphi(y)
	\]
	对所有 $r\in R$ 与 $x,y\in M$ 成立。
	\item 若 $\varphi,\psi\in\Hom_R(M,N)$，定义
	\[
	(\varphi+\psi)(m)=\varphi(m)+\psi(m).
	\]
	则 $\Hom_R(M,N)$ 是阿贝尔群。若 $R$ 交换，再定义
	\[
	(r\varphi)(m)=r\varphi(m),
	\]
	则 $\Hom_R(M,N)$ 是 $R$-模。
	\item 若 $\varphi\in\Hom_R(L,M)$ 且 $\psi\in\Hom_R(M,N)$，则 $\psi\circ\varphi\in\Hom_R(L,N)$。
	\item $\Hom_R(M,M)$ 在逐点加法与复合乘法下是带单位元的环；若 $R$ 交换，则它是 $R$-代数。
\end{enumerate}


\noindent\textbf{定义.} 环 $\Hom_R(M,M)$ 称为 $M$ 的自同态环，记为
\[
\End_R(M).
\]
其元素称为自同态。若 $R$ 交换，有自然映射 $R\to\End_R(M)$，$r\mapsto rI$，即“乘以 $r$”的变换；其像在 $\End_R(M)$ 的中心中。


\noindent\textbf{命题 3.} 设 $N$ 是 $R$-模 $M$ 的子模。商阿贝尔群 $M/N$ 可成为 $R$-模，定义
\[
r(x+N)=rx+N.
\]
自然投影
\[
\pi:M\longrightarrow M/N,
\qquad \pi(x)=x+N
\]
是 $R$-模同态，其核为 $N$。

\begin{proof}
	若 $x+N=y+N$，则 $x-y\in N$，从而 $r(x-y)\in N$，所以 $rx+N=ry+N$，故作用良定义。模公理由 $M$ 中的相应公理推出。自然投影作为阿贝尔群同态的核为 $N$，且
	\[
	\pi(rm)=rm+N=r(m+N)=r\pi(m).
	\]
\end{proof}


\noindent\textbf{定义.} 若 $A,B$ 是 $M$ 的子模，则它们的和为
\[
A+B=\{a+b\mid a\in A,\ b\in B\}.
\]
它是含 $A$ 与 $B$ 的最小子模。


\noindent\textbf{定理 4（模的同构定理）.}
\begin{enumerate}[label=(\arabic*)]
	\item \textbf{第一同构定理：} 若 $\varphi:M\to N$ 是 $R$-模同态，则 $\ker\varphi$ 是 $M$ 的子模，且
	\[
	M/\ker\varphi\cong \varphi(M).
	\]
	\item \textbf{第二同构定理：} 若 $A,B$ 是 $M$ 的子模，则
	\[
	(A+B)/B\cong A/(A\cap B).
	\]
	\item \textbf{第三同构定理：} 若 $A\subseteq B\subseteq M$ 是子模，则
	\[
	(M/A)/(B/A)\cong M/B.
	\]
	\item \textbf{第四（格）同构定理：} 设 $N$ 是 $M$ 的子模。含 $N$ 的 $M$ 的子模与 $M/N$ 的子模之间存在双射
	\[
	A\longleftrightarrow A/N,
	\qquad A\supseteq N,
	\]
	且此对应保持和与交。
\end{enumerate}


\paragraph{练习}
在以下练习中，$R$ 是带单位元的环，$M$ 是左 $R$-模。
\begin{enumerate}[label=\arabic*.]
	\item 用子模判别法证明 $R$-模同态的核与像都是子模。
	\item 证明“作为 $R$-模同构”是任意 $R$-模集合上的等价关系。
	\item 给出一个从一个 $R$-模到另一个 $R$-模的映射的明确例子，它是群同态但不是 $R$-模同态。
	\item 设 $A$ 是任意 $\ZZ$-模，$a\in A$，$n$ 是正整数。证明映射
	\[
	\varphi_a:\ZZ/n\ZZ\longrightarrow A,\qquad \varphi_a(\bar k)=ka
	\]
	是良定义的 $\ZZ$-模同态当且仅当 $na=0$。证明
	\[
	\Hom_{\ZZ}(\ZZ/n\ZZ,A)\cong A_n,
	\qquad A_n=\{a\in A\mid na=0\},
	\]
	其中 $A_n$ 是 $\ZZ$ 的理想 $(n)$ 在 $A$ 中的零化子（参见第 1 节练习 10）。
	\item 列出所有从 $\ZZ/30\ZZ$ 到 $\ZZ/21\ZZ$ 的 $\ZZ$-模同态。
	\item 证明
	\[
	\Hom_{\ZZ}(\ZZ/n\ZZ,\ZZ/m\ZZ)\cong \ZZ/(n,m)\ZZ.
	\]
	\item 设 $z$ 是 $R$ 中心的一个固定元素。证明映射 $m\mapsto zm$ 是 $M$ 到自身的 $R$-模同态。证明若 $R$ 是交换环，则映射
	\[
	R\longrightarrow \End_R(M),\qquad r\mapsto rI
	\]
	是环同态，其中 $I$ 是恒等自同态。
	\item 设 $\varphi:M\to N$ 是 $R$-模同态。证明
	\[
	\varphi(\Tor(M))\subseteq \Tor(N)
	\]
	（参见第 1 节练习 8）。
	\item 设 $R$ 是交换环。证明 $\Hom_R(R,M)$ 与 $M$ 作为左 $R$-模同构。[证明 $\Hom_R(R,M)$ 中每个元素都由其在 $R$ 的单位元上的值唯一确定。]
	\item 设 $R$ 是交换环。证明 $\Hom_R(R,R)$ 与 $R$ 作为环同构。
	\item 设 $A_1,A_2,\ldots,A_n$ 是 $R$-模，且 $B_i$ 是 $A_i$ 的子模，$i=1,2,\ldots,n$。证明
	\[
	(A_1\times\cdots\times A_n)/(B_1\times\cdots\times B_n)
	\cong (A_1/B_1)\times\cdots\times(A_n/B_n).
	\]
	[回忆第 5.1 节练习 14。]
	\item 设 $I$ 是 $R$ 的左理想，$n$ 是正整数。证明
	\[
	R^n/IR^n\cong R/IR\times\cdots\times R/IR\quad(n\text{ 个因子}),
	\]
	其中 $IR^n$ 如第 1 节练习 5 中定义。[使用前一题。]
	\item 设 $I$ 是交换环 $R$ 中的幂零理想（参见第 7.3 节练习 37），$M,N$ 是 $R$-模，$\varphi:M\to N$ 是 $R$-模同态。证明若诱导映射
	\[
	\bar\varphi:M/IM\longrightarrow N/IN
	\]
	是满射，则 $\varphi$ 是满射。
	\item 令 $R=\ZZ[x]$，令 $A=\ZZ[t_1,t_2,\ldots]$ 是独立不定元 $t_1,t_2,\\ldots$ 的多项式环。如下定义 $R$ 在 $A$ 上的作用：
	\begin{enumerate}[label=\arabic*)]
		\item $1\in R$ 作为恒等映射作用在 $A$ 上；
		\item 对 $n\geq1$，令 $x^n\circ 1=t_n$，令 $x^n\circ t_i=t_{n+i}$（$i=1,2,\ldots$），并令 $x^n$ 在 $A$ 中总次数至少为 $2$ 的单项式上作用为 $0$；
		\item 按 $\ZZ$-线性延拓，即使模公理 $2(a)$ 与 $2(c)$ 成立。
	\end{enumerate}
	\begin{enumerate}[label=(\alph*)]
		\item 证明 $x^{p+q}\circ t_i=x^p\circ(x^q\circ t_i)=t_{p+q+i}$，并用此证明在这个作用下 $A$ 是一个幺 $R$-模。
		\item 证明映射 $\varphi:R\to A$，$\varphi(r)=r\circ1_A$，是把环 $R$ 映入环 $A$ 的 $R$-模同态，且把 $1_R$ 映到 $1_A$；但它不是从 $R$ 到 $A$ 的环同态。
	\end{enumerate}
\end{enumerate}


\section{模的生成、直和与自由模}

设 $R$ 是带 $1$ 的环。本节推广子模的和、生成集、有限生成、循环模、直和与自由模。


\noindent\textbf{定义.} 设 $M$ 是 $R$-模，$N_1,
\ldots,N_n$ 是 $M$ 的子模。
\begin{enumerate}[label=(\arabic*)]
	\item $N_1+
	\cdots+N_n$ 是所有有限和
	\[
	a_1+
	\cdots+a_n,
	\qquad a_i\in N_i
	\]
	组成的集合。
	\item 对任意子集 $A\subseteq M$，令
	\[
	RA=\{r_1a_1+
	\cdots+r_ma_m\mid r_i\in R,\ a_i\in A,\ m\geq1\},
	\]
	约定 $R\emptyset=0$。称 $RA$ 为 $A$ 生成的子模。若 $N=RA$，称 $A$ 为 $N$ 的生成集。
	\item 若 $N$ 由有限集生成，则称 $N$ 有限生成。
	\item 若 $N=Ra$ 对某个 $a\in M$ 成立，则称 $N$ 为循环模。
\end{enumerate}
当 $R$ 有 $1$ 时，$A\subseteq RA$。$RA$ 是含 $A$ 的最小子模。循环模当然是有限生成模。

\paragraph{例}
\begin{enumerate}[label=(\arabic*)]
	\item 当 $R=\ZZ$ 时，$Za$ 就是阿贝尔群中由 $a$ 生成的循环子群。有限生成 $\ZZ$-模就是有限生成阿贝尔群。
	\item $R$ 作为自身左模是循环模，因为 $R=R1$。此时 $R$ 的子模是左理想；循环子模对应主理想，有限生成子模对应有限生成理想。
	\item $R^n$ 由标准向量 $e_1,
	\ldots,e_n$ 生成：
	\[
	(s_1,
	\ldots,s_n)=\sum_{i=1}^n s_i e_i.
	\]
	\item 若 $V$ 是由线性变换 $T$ 给出的 $F[x]$-模，则 $V$ 是由 $v$ 生成的循环 $F[x]$-模，当且仅当
	\[
	V=\{p(T)v\mid p(x)\in F[x]\},
	\]
	即 $\{v,Tv,T^2v,
	\ldots\}$ 张成 $V$。
\end{enumerate}


\noindent\textbf{定义.} 设 $M_1,
\ldots,M_k$ 是 $R$-模。它们的直积由所有 $k$ 元组
\[
(m_1,
\ldots,m_k),
\qquad m_i\in M_i
\]
组成，在逐分量加法和逐分量 $R$-作用下成为 $R$-模，记为
\[
M_1\times\cdots\times M_k.
\]
有限多个模的直积也称为外直和，记为
\[
M_1\oplus\cdots\oplus M_k.
\]


\noindent\textbf{命题 5.} 设 $N_1,
\ldots,N_k$ 是 $R$-模 $M$ 的子模。令
\[
\pi:N_1\oplus\cdots\oplus N_k\longrightarrow N_1+
\cdots+N_k,
\qquad (a_1,
\ldots,a_k)\mapsto a_1+
\cdots+a_k.
\]
以下条件等价：
\begin{enumerate}[label=(\arabic*)]
	\item $\pi$ 是同构；
	\item 对每个 $i$，有
	\[
	N_i\cap(N_1+
	\cdots+N_{i-1}+N_{i+1}+
	\cdots+N_k)=0;
	\]
	\item $N_1+
	\cdots+N_k$ 中每个元素可唯一写成 $a_1+
	\cdots+a_k$，其中 $a_i\in N_i$。
\end{enumerate}
若这些条件成立，称 $M=N_1+
\cdots+N_k$ 是 $N_i$ 的内直和，写作
\[
M=N_1\oplus\cdots\oplus N_k.
\]


\noindent\textbf{定义.} $R$-模 $F$ 称为在子集 $A$ 上自由，若每个非零元素 $x\in F$ 都能唯一写成有限和
\[
x=r_1a_1+
\cdots+r_na_n,
\]
其中 $r_i\in R$ 非零且 $a_i\in A$ 两两不同。此时 $A$ 称为 $F$ 的基或自由生成元集。若 $R$ 交换，$|A|$ 称为 $F$ 的秩。


\noindent\textbf{定理 6.} 对任意集合 $A$，存在在 $A$ 上自由的 $R$-模 $F(A)$，并满足如下泛性质：若 $M$ 是任意 $R$-模，$\varphi:A\to M$ 是集合映射，则存在唯一的 $R$-模同态
\[
\Phi:F(A)\longrightarrow M
\]
使 $\Phi(a)=\varphi(a)$ 对所有 $a\in A$ 成立。若 $A=\{a_1,
\ldots,a_n\}$，则
\[
F(A)=Ra_1\oplus\cdots\oplus Ra_n\cong R^n.
\]

\begin{proof}
	若 $A=\emptyset$，取 $F(A)=0$。若 $A$ 非空，令 $F(A)$ 为所有函数 $f:A\to R$，其中除有限多个 $a$ 外 $f(a)=0$。在逐点加法与逐点 $R$-乘法下它是 $R$-模。把 $a\in A$ 识别为在 $a$ 处取值 $1$、其他处取值 $0$ 的函数。于是 $F(A)$ 的每个元素都唯一写成有限形式和 $\sum r_i a_i$。给定 $\varphi:A\to M$，定义
	\[
	\Phi\left(\sum r_i a_i\right)=\sum r_i\varphi(a_i),
	\]
	由唯一表示可知良定义且为 $R$-模同态。唯一性来自 $A$ 生成 $F(A)$。
\end{proof}


\noindent\textbf{推论 7.}
\begin{enumerate}[label=(\arabic*)]
	\item 若 $F_1,F_2$ 都在同一集合 $A$ 上自由，则存在唯一同构 $F_1\cong F_2$，且在 $A$ 上为恒等。
	\item 若 $F$ 是以 $A$ 为基的自由 $R$-模，则 $F\cong F(A)$，并满足同样的泛性质。
\end{enumerate}
当 $R=\ZZ$ 时，集合 $A$ 上的自由模称为 $A$ 上的自由阿贝尔群。


\paragraph{练习}
在以下练习中，$R$ 是带单位元的环，$M$ 是左 $R$-模。
\begin{enumerate}[label=\arabic*.]
	\item 证明若集合 $A$ 与 $B$ 有相同基数，则自由模 $F(A)$ 与 $F(B)$ 同构。
	\item 假设 $R$ 交换。证明 $R^n\cong R^m$ 当且仅当 $n=m$，即两个有限秩自由 $R$-模同构当且仅当它们的秩相同。[对 $R$ 的极大理想 $I$ 应用第 2 节练习 12。可使用如下事实：若 $F$ 是域，则 $F^n\cong F^m$ 当且仅当 $n=m$，即有限维向量空间同构当且仅当维数相同；这将在第 11.1 节证明。]
	\item 证明第 1 节练习 18 和 19 中的 $F[x]$-模都是循环模。
	\item 若 $R$-模 $M$ 的每个元素 $m$ 都满足 $rm=0$ 对某个非零 $r\in R$ 成立（$r$ 可依 $m$ 而变），则称 $M$ 为\emph{挠模}，即在第 1 节练习 8 的记号下 $M=\Tor(M)$。证明每个有限阿贝尔群都是挠 $\ZZ$-模。给出一个无限阿贝尔群的例子，它是挠 $\ZZ$-模。
	\item 设 $R$ 是整环。证明每个有限生成的挠 $R$-模都有非零零化子，即存在非零 $r\in R$ 使 $rm=0$ 对所有 $m\in M$ 成立；这里 $r$ 不依赖于 $m$（模的零化子见第 1 节练习 9）。给出一个挠 $R$-模的例子，其零化子是零理想。
	\item 证明若 $M$ 是由 $n$ 个元素生成的有限生成 $R$-模，则 $M$ 的每个商模可由 $n$ 个或更少元素生成。推出循环模的商仍是循环模。
	\item 设 $N$ 是 $M$ 的子模。证明若 $M/N$ 与 $N$ 都有限生成，则 $M$ 也有限生成。
	\item 令 $S$ 为整数序列 $(a_1,a_2,a_3,\ldots)$ 的集合，其中除有限多个 $a_i$ 外皆为 $0$（即无限多个 $\ZZ$ 的直和）。回忆 $S$ 在逐分量加法和乘法下是一个环，且没有乘法单位元（参见第 7.1 节练习 20）。证明 $S$ 作为自身上的模不是有限生成的。
	\item 称 $R$-模 $M$ 为\emph{不可约}的，如果 $M\neq0$ 且 $0$ 与 $M$ 是它仅有的子模。证明 $M$ 不可约当且仅当 $M\neq0$ 且 $M$ 是循环模，并且任一非零元素都是生成元。确定所有不可约 $\ZZ$-模。
	\item 假设 $R$ 交换。证明 $R$-模 $M$ 不可约当且仅当 $M$ 作为 $R$-模同构于 $R/I$，其中 $I$ 是 $R$ 的极大理想。[由前一题，若 $M$ 不可约，则对任意固定非零元素 $m\in M$，有自然映射 $R\to M$，$r\mapsto rm$。]
	\item 证明若 $M_1$ 与 $M_2$ 是不可约 $R$-模，则任意非零 $R$-模同态 $M_1\to M_2$ 都是同构。推出若 $M$ 不可约，则 $\End_R(M)$ 是除环（此结果称为 Schur 引理）。[考虑核与像。]
	\item 设 $R$ 是交换环，$A,B,M$ 是 $R$-模。证明下列 $R$-模同构：
	\begin{enumerate}[label=(\alph*)]
		\item $\Hom_R(A\times B,M)\cong \Hom_R(A,M)\times\Hom_R(B,M)$；
		\item $\Hom_R(M,A\times B)\cong \Hom_R(M,A)\times\Hom_R(M,B)$。
	\end{enumerate}
	\item 设 $R$ 是交换环，$F$ 是有限秩自由 $R$-模。证明 $R$-模同构
	\[
	\Hom_R(F,R)\cong F.
	\]
	\item 设 $R$ 是交换环，$F$ 是秩为 $n$ 的自由 $R$-模。证明
	\[
	\Hom_R(F,M)\cong M\times\cdots\times M\quad(n\text{ 个因子}).
	\]
	[使用第 2 节练习 9 与练习 12。]
	\item 若 $e\in R$ 满足 $e^2=e$ 且 $er=re$ 对所有 $r\in R$ 成立，则称 $e$ 为\emph{中心幂等元}。若 $e$ 是 $R$ 中的中心幂等元，证明
	\[
	M=eM\oplus(1-e)M.
	\]
	[回忆第 1 节练习 14。]
	
	以下两题建立模的中国剩余定理（参见第 7.6 节）。
	\item 对 $R$ 的任意理想 $I$，令 $IM$ 为第 1 节练习 5 中定义的子模。令 $A_1,\ldots,A_k$ 为环 $R$ 中任意理想。证明映射
	\[
	M\longrightarrow M/A_1M\times\cdots\times M/A_kM,
	\qquad m\longmapsto(m+A_1M,\ldots,m+A_kM)
	\]
	是 $R$-模同态，其核为
	\[
	A_1M\cap A_2M\cap\cdots\cap A_kM.
	\]
	\item 在前一题的记号下，再假设理想 $A_1,\ldots,A_k$ 两两互素，即 $A_i+A_j=R$ 对所有 $i\neq j$ 成立。证明
	\[
	M/(A_1\cdots A_k)M\cong M/A_1M\times\cdots\times M/A_kM.
	\]
	[参见第 7.6 节中环的中国剩余定理的证明。]
	\item 设 $R$ 是主理想整环，$M$ 是被非零真理想 $(a)$ 零化的 $R$-模。设
	\[
	a=p_1^{\alpha_1}p_2^{\alpha_2}\cdots p_k^{\alpha_k}
	\]
	是 $a$ 在 $R$ 中分解为互异素元幂的唯一分解。令 $M_i$ 为 $p_i^{\alpha_i}$ 在 $M$ 中的零化子，即
	\[
	M_i=\{m\in M\mid p_i^{\alpha_i}m=0\},
	\]
	称为 $M$ 的 $p_i$-primary 分量。证明
	\[
	M=M_1\oplus M_2\oplus\cdots\oplus M_k.
	\]
	\item 证明若 $M$ 是阶为 $a=p_1^{\alpha_1}p_2^{\alpha_2}\cdots p_k^{\alpha_k}$ 的有限阿贝尔群，则把它看作 $\ZZ$-模时，$M$ 被 $(a)$ 零化；$M$ 的 $p_i$-primary 分量是 $M$ 的唯一 Sylow $p_i$-子群；并且 $M$ 同构于它的 Sylow 子群的直积。
	\item 设 $I$ 是非空指标集，并对每个 $i\in I$，令 $M_i$ 为 $R$-模。$M_i$ 的直积定义为其作为阿贝尔群的直积（参见第 5.1 节练习 15），并以逐分量乘法给出 $R$ 的作用。$M_i$ 的直和定义为阿贝尔群的限制直积（参见第 5.1 节练习 17），并以逐分量乘法给出 $R$ 的作用。换言之，$M_i$ 的直和是直积 $\prod_{i\in I}M_i$ 的子集，由所有只有有限多个分量 $m_i$ 非零的元素 $\prod_{i\in I}m_i$ 构成；$R$ 在直积或直和上的作用为
	\[
	r\prod_{i\in I}m_i=\prod_{i\in I}rm_i.
	\]
	直和记为 $\bigoplus_{i\in I}M_i$。
	\begin{enumerate}[label=(\alph*)]
		\item 证明 $M_i$ 的直积是 $R$-模，$M_i$ 的直和是其直积的子模。
		\item 若 $R=\ZZ$，$I=\ZZ^+$，且 $M_i$ 是阶为 $i$ 的循环群，证明 $M_i$ 的直和与直积不同构。[观察挠元素。]
	\end{enumerate}
	\item 设 $I$ 是非空指标集，并对每个 $i\in I$，令 $N_i$ 是 $M$ 的子模。证明下列条件等价：
	\begin{enumerate}[label=(\roman*)]
		\item 由所有 $N_i$ 生成的 $M$ 的子模同构于这些 $N_i$ 的直和；
		\item 若 $\{i_1,i_2,\ldots,i_k\}$ 是 $I$ 的任意有限子集，则
		\[
		N_{i_1}\cap(N_{i_2}+\cdots+N_{i_k})=0;
		\]
		\item 若 $\{i_1,i_2,\ldots,i_k\}$ 是 $I$ 的任意有限子集，则
		\[
		N_{i_1}+\cdots+N_{i_k}=N_{i_1}\oplus\cdots\oplus N_{i_k};
		\]
		\item 对由 $N_i$ 生成的子模中的每个元素 $x$，存在唯一的一族元素 $a_i\in N_i$（$i\in I$），使得除有限多个 $a_i$ 外均为零，且 $x$ 是这些 $a_i$ 的有限和。
	\end{enumerate}
	\item 设 $R$ 是主理想整环，$M$ 是挠 $R$-模（见练习 4），$p$ 是 $R$ 中素元；不要假设 $M$ 有限生成，因此它不必有非零零化子（参见练习 5）。$M$ 的 $p$-primary 分量定义为 $M$ 中所有被 $p$ 的某个正幂零化的元素的集合。
	\begin{enumerate}[label=(\alph*)]
		\item 证明 $p$-primary 分量是子模。[见第 1 节练习 13。]
		\item 当 $M$ 有非零零化子时，证明这里的 $p$-primary 分量定义与练习 18 中的定义一致。
		\item 证明 $M$ 是其所有 $p$-primary 分量的（可能无限）直和，其中 $p$ 遍历 $R$ 的所有素元。
	\end{enumerate}
	\item 证明任意自由 $R$-模的直和仍是自由的。
	\item （自由模的任意直积不必自由）对每个正整数 $i$，令 $M_i$ 为自由 $\ZZ$-模 $\ZZ$，并令
	\[
	M=\prod_{i\in\ZZ^+}M_i.
	\]
	每个 $M$ 的元素可唯一写成 $(a_1,a_2,a_3,\ldots)$，其中 $a_i\in\ZZ$。令 $N$ 为 $M$ 中所有只有有限多个非零 $a_i$ 的元构成的子模。假设 $M$ 是有基 $\mathcal B$ 的自由 $\ZZ$-模。
	\begin{enumerate}[label=(\alph*)]
		\item 证明 $N$ 是可数的。
		\item 证明存在 $\mathcal B$ 的可数子集 $\mathcal B_1$，使 $N$ 包含于由 $\mathcal B_1$ 生成的子模 $N_1$ 中。还要证明 $N_1$ 是可数的。
		\item 令 $\overline M=M/N_1$。证明 $\overline M$ 是自由 $\ZZ$-模。推出若 $x$ 是 $\overline M$ 中非零元素，则只有有限多个不同的正整数 $k$ 使得 $x=km$ 对某个 $m\in\overline M$ 成立。
		\item 令
		\[
		S=\{(b_1,b_2,b_3,\ldots)\mid b_i=\pm i!\text{ 对所有 }i\}.
		\]
		证明 $S$ 不可数。推出存在 $s\in S$ 使 $s\notin N_1$。
		\item 证明 $M$ 是自由的假设导致矛盾：由 (d) 取 $s\in S$ 且 $s\notin N_1$。证明对每个正整数 $k$，存在 $m\in\overline M$ 使 $s=km$，与 (c) 矛盾。[用到 $N\subseteq N_1$。]
	\end{enumerate}
	\item 在直接极限的构造中（第 7.6 节练习 8），证明若所有 $A_i$ 都是 $R$-模且映射 $\rho_{ij}$ 都是 $R$-模同态，则直接极限 $A=\varinjlim A_i$ 可以自然地赋予 $R$-模结构，使得映射 $\rho_i:A_i\to A$ 全都是 $R$-模同态。验证相应的泛性质，即对所有与 $\rho_{ij}$ 相容的 $R$-模同态 $\varphi_i:A_i\to C$，第 7.6 节练习 8(e) 的结论仍成立。
	\item 对逆极限作出与前一题相应的分析，证明 $R$-模的逆极限是 $R$-模并满足适当的泛性质（参见第 7.6 节练习 10）。
	\item （非交换环上的自由模不必有唯一的秩）令 $M$ 为练习 24 中的 $\ZZ$-模 $\ZZ\times\ZZ\times\cdots$，令 $R$ 为其自同态环 $R=\End_{\ZZ}(M)$（参见第 7.1 节练习 29 和 30）。定义 $\varphi_1,\varphi_2\in R$ 为
	\[
	\varphi_1(a_1,a_2,a_3,\ldots)=(a_1,a_3,a_5,\ldots),
	\qquad
	\varphi_2(a_1,a_2,a_3,\ldots)=(a_2,a_4,a_6,\ldots).
	\]
	\begin{enumerate}[label=(\alph*)]
		\item 证明 $\{\varphi_1,\varphi_2\}$ 是左 $R$-模 $R$ 的自由基。[定义
		\[
		\psi_1(a_1,a_2,\ldots)=(a_1,0,a_2,0,\ldots),
		\quad
		\psi_2(a_1,a_2,\ldots)=(0,a_1,0,a_2,\ldots).
		\]
		验证 $\varphi_i\psi_i=1$，$\varphi_1\psi_2=0=\varphi_2\psi_1$，且 $\psi_1\varphi_1+\psi_2\varphi_2=1$。用这些关系证明 $\varphi_1,
		\varphi_2$ 作为左 $R$-模的元素线性无关且生成 $R$。]
		\item 用 (a) 证明 $R\cong R^2$，并推出 $R\cong R^n$ 对所有 $n\in\ZZ^+$ 成立。
	\end{enumerate}
\end{enumerate}


\section{模的张量积}

本节讨论张量积。张量积首先作为“扩张标量”的工具出现。设 $R\subseteq S$ 是带单位元的环，若 $N$ 是 $S$-模，则通过限制作用，$N$ 自动成为 $R$-模。反过来，若从一个 $R$-模 $N$ 出发，希望把标量从 $R$ 扩张到 $S$，一般不能直接在 $N$ 上定义 $S$-模结构。张量积 $S\otimes_R N$ 给出最自然的目标。

构造如下：取集合 $S\times N$ 上的自由阿贝尔群，并模去由以下元素生成的子群：
\[
\begin{aligned}
	&(s_1+s_2,n)-(s_1,n)-(s_2,n),\\
	&(s,n_1+n_2)-(s,n_1)-(s,n_2),\\
	&(sr,n)-(s,rn),
\end{aligned}
\]
其中 $s,s_1,s_2\in S$，$n,n_1,n_2\in N$，$r\in R$。所得商群记为 $S\otimes_R N$，$(s,n)$ 的像记为 $s\otimes n$。于是有关系
\[
\begin{aligned}
	(s_1+s_2)\otimes n&=s_1\otimes n+s_2\otimes n,\\
	s\otimes(n_1+n_2)&=s\otimes n_1+s\otimes n_2,\\
	sr\otimes n&=s\otimes rn.
\end{aligned}
\]
它自然成为左 $S$-模，定义
\[
s\left(\sum s_i\otimes n_i\right)=\sum ss_i\otimes n_i.
\]
称 $S\otimes_R N$ 为由 $N$ 通过从 $R$ 到 $S$ 扩张标量得到的 $S$-模。自然映射
\[
t:N\longrightarrow S\otimes_R N,
\qquad n\mapsto 1\otimes n
\]
是 $R$-模同态，但一般不单。


\noindent\textbf{定理 8（扩张标量的泛性质）.} 设 $R\subseteq S$，$N$ 是左 $R$-模，$t:N\to S\otimes_R N$ 由 $t(n)=1\otimes n$ 定义。若 $L$ 是左 $S$-模，$\varphi:N\to L$ 是 $R$-模同态，则存在唯一的 $S$-模同态
\[
\Phi:S\otimes_R N\longrightarrow L
\]
使 $\varphi=\Phi\circ t$。反过来，任意 $S$-模同态 $\Phi:S\otimes_R N\to L$ 与 $t$ 复合给出 $R$-模同态 $N\to L$。


\noindent\textbf{推论 9.} 若 $t:N\to S\otimes_R N$ 是上述映射，则
\[
\ker t
\]
等于所有在任意 $R$-模同态 $N\to L$（其中 $L$ 是 $S$-模）下都映为 $0$ 的元素。特别地，$N$ 可嵌入某个 $S$-模中当且仅当 $t$ 是单射。

\paragraph{例}
\begin{enumerate}[label=(\arabic*)]
	\item $\QQ\otimes_\ZZ \ZZ\cong \QQ$，对应于把 $\ZZ$ 中的标量扩张到 $\QQ$。
	\item $\QQ\otimes_\ZZ (\ZZ/n\ZZ)=0$，因为有理向量空间中没有非零有限阶元素。
	\item 一般地，若 $d=(m,n)$，则
	\[
	\ZZ/m\ZZ\otimes_\ZZ \ZZ/n\ZZ\cong \ZZ/d\ZZ.
	\]
	\item 若 $I$ 是 $R$ 的双边理想，$N$ 是左 $R$-模，则
	\[
	(R/I)\otimes_R N\cong N/IN.
	\]
\end{enumerate}

现在给出一般张量积。设 $M$ 是右 $R$-模，$N$ 是左 $R$-模。在 $M\times N$ 上自由生成的阿贝尔群中，模去由
\[
\begin{aligned}
	&(m_1+m_2,n)-(m_1,n)-(m_2,n),\\
	&(m,n_1+n_2)-(m,n_1)-(m,n_2),\\
	&(mr,n)-(m,rn)
\end{aligned}
\]
生成的子群，得到阿贝尔群 $M\otimes_R N$，其元素称为张量。简单张量记为 $m\otimes n$。


\noindent\textbf{定理 10（张量积的泛性质）.} 设 $M$ 是右 $R$-模，$N$ 是左 $R$-模。若 $D$ 是阿贝尔群，映射
\[
\varphi:M\times N\longrightarrow D
\]
满足双加性并满足 $\varphi(mr,n)=\varphi(m,rn)$，则称为 $R$-平衡映射。张量积 $M\otimes_R N$ 连同自然映射
\[
t(m,n)=m\otimes n
\]
满足：对任意 $R$-平衡映射 $\varphi:M\times N\to D$，存在唯一群同态
\[
\Phi:M\otimes_R N\longrightarrow D
\]
使 $\varphi=\Phi\circ t$。


\noindent\textbf{推论 11.} $R$-平衡映射 $M\times N\to D$ 与群同态 $M\otimes_R N\to D$ 自然一一对应。

若 $R$ 交换，左右模结构相同，此时平衡映射通常称为双线性映射，并且 $M\otimes_R N$ 也自然成为 $R$-模。


\noindent\textbf{定义.} 设 $R$ 与 $S$ 是带单位元的环。若阿贝尔群 $M$ 同时是左 $S$-模与右 $R$-模，且
\[
(sm)r=s(mr),
\]
则称 $M$ 为 $(S,R)$-双模。


\noindent\textbf{定理 13（两个同态的张量积）.} 设 $M,M'$ 是右 $R$-模，$N,N'$ 是左 $R$-模，且
\[
\varphi:M\to M',
\qquad
\psi:N\to N'
\]
是 $R$-模同态。则存在唯一群同态
\[
\varphi\otimes\psi:M\otimes_R N\longrightarrow M'\otimes_R N'
\]
使
\[
(\varphi\otimes\psi)(m\otimes n)=\varphi(m)\otimes\psi(n).
\]
若还带有适当的双模结构，则这是相应的模同态；并且张量积与复合相容。


\noindent\textbf{定理 14（张量积的结合性）.} 若 $M$ 是右 $R$-模，$N$ 是 $(R,T)$-双模，$L$ 是左 $T$-模，则有自然同构
\[
(M\otimes_R N)\otimes_T L
\cong
M\otimes_R (N\otimes_T L),
\]
把 $(m\otimes n)\otimes l$ 映到 $m\otimes(n\otimes l)$。

若 $R$ 交换，且 $M,N,L$ 是左 $R$-模，则
\[
(M\otimes_R N)\otimes_R L\cong M\otimes_R(N\otimes_R L).
\]
因此多重张量积可不带歧义地写作 $M_1\otimes_R\cdots\otimes_R M_n$。


\noindent\textbf{定义.} 设 $R$ 交换。映射
\[
\varphi:M_1\times\cdots\times M_n\longrightarrow L
\]
称为 $R$ 上 $n$-多重线性映射，若它对每一个变量分别都是 $R$-模同态。


\noindent\textbf{推论 16.} $M_1\otimes_R\cdots\otimes_R M_n$ 是关于 $n$-重线性映射的泛对象：$n$-重线性映射 $M_1\times\cdots\times M_n\to L$ 与 $R$-模同态
\[
M_1\otimes_R\cdots\otimes_R M_n\to L
\]
自然一一对应。


\noindent\textbf{定理 17（直和的张量积）.} 设 $M,M'$ 是右 $R$-模，$N,N'$ 是左 $R$-模。则有自然同构
\[
(M\oplus M')\otimes_R N
\cong
(M\otimes_R N)\oplus(M'\otimes_R N),
\]
以及
\[
M\otimes_R(N\oplus N')
\cong
(M\otimes_R N)\oplus(M\otimes_R N').
\]


\noindent\textbf{推论 18（自由模的扩张标量）.} 若 $N\cong R^n$ 是自由 $R$-模，则
\[
S\otimes_R N\cong S^n
\]
作为左 $S$-模。


\noindent\textbf{推论 19.} 若 $R$ 交换，$M\cong R^s$、$N\cong R^t$ 是自由模，基分别为 $m_1,
\ldots,m_s$ 与 $n_1,
\ldots,n_t$，则 $M\otimes_R N$ 是秩为 $st$ 的自由模，基为
\[
m_i\otimes n_j,
\qquad 1\leq i\leq s,
\quad 1\leq j\leq t.
\]


\noindent\textbf{命题 20.} 若 $R$ 交换，$M,N$ 是 $R$-模，则有自然同构
\[
M\otimes_R N\cong N\otimes_R M,
\]
把 $m\otimes n$ 映到 $n\otimes m$。但当 $M=N$ 时，一般不能推出 $m\otimes n=n\otimes m$。


\noindent\textbf{命题 21.} 若 $R$ 交换，$A,B$ 是 $R$-代数，则在 $A\otimes_R B$ 上定义
\[
(a\otimes b)(a'\otimes b')=aa'\otimes bb'
\]
是良定义的，并使 $A\otimes_R B$ 成为 $R$-代数。


\paragraph{练习}
设 $R$ 是带单位元的环。
\begin{enumerate}[label=\arabic*.]
	\item 设 $f:R\to S$ 是环同态，且 $f(1_R)=1_S$。验证 $sr=sf(r)$ 定义了 $S$ 上的右 $R$-作用，并使 $S$ 成为一个 $(S,R)$-双模。
	\item 证明元素“$2\otimes1$”在 $\ZZ\otimes_{\ZZ}\ZZ/2\ZZ$ 中为 $0$，但在 $2\ZZ\otimes_{\ZZ}\ZZ/2\ZZ$ 中非零。
	\item 证明 $\CC\otimes_{\RR}\CC$ 与 $\CC\otimes_{\CC}\CC$ 都是左 $\RR$-模，但它们作为 $\RR$-模不同构。
	\item 证明 $\QQ\otimes_{\ZZ}\QQ$ 与 $\QQ\otimes_{\QQ}\QQ$ 作为左 $\QQ$-模同构。[证明二者都是 $\QQ$ 上的一维向量空间。]
	\item 设 $A$ 是阶为 $n$ 的有限阿贝尔群，$p^k$ 是素数 $p$ 整除 $n$ 的最高幂。证明
	\[
	\ZZ/p^k\ZZ\otimes_{\ZZ} A
	\]
	同构于 $A$ 的 Sylow $p$-子群。
	\item 若 $R$ 是任意整环，商域为 $Q$，证明
	\[
	(Q/R)\otimes_R(Q/R)=0.
	\]
	\item 若 $R$ 是任意整环，商域为 $Q$，$N$ 是左 $R$-模，证明张量积 $Q\otimes_R N$ 的每个元素均可写为简单张量
	\[
	(1/d)\otimes n
	\]
	的形式，其中 $0\neq d\in R$，$n\in N$。
	\item 设 $R$ 是整环，商域为 $Q$，$N$ 是任意 $R$-模。令 $U=R^\times$ 表示 $R$ 中所有非零元素的集合，定义 $U^{-1}N$ 为有序对 $(u,n)$ 的等价类集合，其中 $u\in U$、$n\in N$，等价关系为
	\[
	(u,n)\sim(u',n')\quad\Longleftrightarrow\quad u'n=un'\text{ 在 }N\text{ 中成立}.
	\]
	\begin{enumerate}[label=(\alph*)]
		\item 证明 $U^{-1}N$ 在加法
		\[
		(u_1,n_1)+(u_2,n_2)=(u_1u_2,u_2n_1+u_1n_2)
		\]
		下成为阿贝尔群。证明 $r(u,n)=(u,rn)$ 定义了 $R$ 在 $U^{-1}N$ 上的作用，使其成为 $R$-模。[这是第 15 章第 4 节一般讨论的局部化的例子；也参见第 7 章第 5 节。]
		\item 证明从 $Q\times N$ 到 $U^{-1}N$ 的映射
		\[
		(a/b,n)\longmapsto(b,an)
		\]
		（其中 $a\in R,b\in U,n\in N$）是 $R$-平衡映射，从而诱导同态 $f:Q\otimes_RN\to U^{-1}N$。证明从 $U^{-1}N$ 到 $Q\otimes_RN$ 的映射
		\[
		g((u,n))=(1/u)\otimes n
		\]
		良定义且是 $f$ 的逆同态。推出 $Q\otimes_RN\cong U^{-1}N$ 作为 $R$-模。
		\item 由 (b) 推出，在 $Q\otimes_RN$ 中 $(1/d)\otimes n=0$ 当且仅当存在非零 $r\in R$ 使 $rn=0$。
		\item 若 $A$ 是阿贝尔群，证明 $\QQ\otimes_{\ZZ}A=0$ 当且仅当 $A$ 是挠阿贝尔群，即每个元素都有有限阶。
	\end{enumerate}
	\item 设 $R$ 是整环，商域为 $Q$，$N$ 是任意 $R$-模。令 $Q\otimes_RN$ 为由 $N$ 从 $R$ 到 $Q$ 的标量扩张所得的模。证明 $R$-模同态
	\[
	\iota:N\longrightarrow Q\otimes_RN,
	\qquad n\longmapsto 1\otimes n
	\]
	的核是 $N$ 的挠子模（参见第 1 节练习 8）。[使用前一题。]
	\item 设 $R$ 交换，$N\cong R^n$ 是秩为 $n$ 的自由 $R$-模，基为 $e_1,
	\ldots,e_n$。
	\begin{enumerate}[label=(\alph*)]
		\item 对任意非零 $R$-模 $M$，证明 $M\otimes_R N$ 的每个元素都可唯一写成
		\[
		\sum_{i=1}^n m_i\otimes e_i,
		\qquad m_i\in M.
		\]
		推出若 $\sum_{i=1}^n m_i\otimes e_i=0$，则所有 $m_i=0$。
		\item 证明若 $\sum m_i\otimes n_i=0$ 在 $M\otimes_R N$ 中成立，其中 $n_i$ 仅假设为 $R$-线性无关，则不一定推出所有 $m_i=0$。[考虑 $R=\ZZ$、$n=1$、$M=\ZZ/2\ZZ$ 与元素 $1\otimes2$。]
	\end{enumerate}
	\item 设 $\{e_1,e_2\}$ 是 $V=\RR^2$ 的一组基。证明 $V\otimes_{\RR}V$ 中的元素
	\[
	e_1\otimes e_2+e_2\otimes e_1
	\]
	不能写为简单张量 $v\otimes w$，其中 $v,w\in\RR^2$。
	\item 设 $V$ 是域 $F$ 上的向量空间，$v,v'$ 是 $V$ 的非零元素。证明
	\[
	v\otimes v'=v'\otimes v\quad\text{在 }V\otimes_FV\text{ 中}
	\]
	当且仅当 $v=av'$ 对某个 $a\in F$ 成立。
	\item 证明通常的点积
	\[
	(a_1,\ldots,a_n)\cdot(b_1,\ldots,b_n)=a_1b_1+\cdots+a_nb_n
	\]
	是从 $\RR^n\times\RR^n$ 到 $\RR$ 的双线性映射。
	\item 设 $I$ 是任意非空指标集，并对每个 $i\in I$，令 $N_i$ 是左 $R$-模。令 $M$ 为右 $R$-模。证明群同构
	\[
	M\otimes_R\left(\bigoplus_{i\in I}N_i\right)
	\cong \bigoplus_{i\in I}(M\otimes_R N_i),
	\]
	其中任意模族的直和定义见第 3 节练习 20。[使用与两个模直和情形相同的论证，并注意何处用到了直和假设；参见下一题。]
	\item 证明张量积一般不与直积交换。[考虑从 $\ZZ$ 到 $\QQ$ 的标量扩张作用于模 $M_i=\ZZ/2^i\ZZ$（$i=1,2,\ldots$）的直积。]
	\item 设 $R$ 交换，$I,J$ 是 $R$ 的理想，于是 $R/I$ 与 $R/J$ 自然是 $R$-模。
	\begin{enumerate}[label=(\alph*)]
		\item 证明 $R/I\otimes_RR/J$ 的每个元素都可写为形如
		\[
		(1\bmod I)\otimes(r\bmod J)
		\]
		的简单张量。
		\item 证明存在 $R$-模同构
		\[
		R/I\otimes_RR/J\cong R/(I+J),
		\]
		它把 $(r\bmod I)\otimes(r'\bmod J)$ 映到 $rr'\bmod(I+J)$。
	\end{enumerate}
	\item 令 $I=(2,x)$ 为环 $R=\ZZ[x]$ 中由 $2$ 与 $x$ 生成的理想。环 $\ZZ/2\ZZ=R/I$ 自然是同时被 $2$ 与 $x$ 零化的 $R$-模。
	\begin{enumerate}[label=(\alph*)]
		\item 证明映射
		\[
		\varphi:I\times I\longrightarrow \ZZ/2\ZZ,
		\]
		定义为
		\[
		\varphi(a_0+a_1x+\cdots+a_nx^n,
		b_0+b_1x+\cdots+b_mx^m)=b_1\bmod2
		\]
		是 $R$-双线性的。
		\item 证明存在 $R$-模同态 $I\otimes_R I\to\ZZ/2\ZZ$，把 $p(x)\otimes q(x)$ 映到 $p(0)q'(0)$，其中 $q'$ 是 $q$ 的通常多项式导数。
		\item 证明在 $I\otimes_R I$ 中，$2\otimes x\neq x\otimes2$。
	\end{enumerate}
	\item 设 $I$ 是整环 $R$ 中的主理想。证明 $R$-模 $I\otimes_RI$ 没有非零挠元素，即若 $0\neq r\in R$ 且 $rm=0$（$m\in I\otimes_RI$），则 $m=0$。
	\item 令 $I=(2,x)$ 为环 $R=\ZZ[x]$ 中由 $2$ 与 $x$ 生成的理想，如练习 17。证明 $I\otimes_RI$ 中的非零元素
	\[
	2\otimes x-x\otimes2
	\]
	是挠元素。进一步证明它同时被 $2$ 与 $x$ 零化，并且由 $2\otimes x-x\otimes2$ 生成的 $I\otimes_RI$ 的子模同构于 $R/I$。
	\item 令 $I=(2,x)$ 为环 $R=\ZZ[x]$ 中由 $2$ 与 $x$ 生成的理想。证明 $I\otimes_RI$ 中的元素
	\[
	2\otimes2+x\otimes x
	\]
	不是简单张量，即不能写成 $a\otimes b$，其中 $a,b\in I$。
	\item 设 $R$ 交换，$I,J$ 是 $R$ 的理想。
	\begin{enumerate}[label=(\alph*)]
		\item 证明存在从 $I\otimes_RJ$ 到乘积理想 $IJ$ 的满 $R$-模同态，把 $i\otimes j$ 映到元素 $ij$。
		\item 给出一个例子说明 (a) 中的映射不必单（参见练习 17）。
	\end{enumerate}
	\item 设 $M$ 同时是左、右 $R$-模，并且 $rm=mr$ 对所有 $r\in R$、$m\in M$ 成立。证明任意 $r_1,r_2\in R$，元素 $r_1r_2$ 与 $r_2r_1$ 在 $M$ 上的作用相同。（这解释了在 $R$-代数定义中假设 $R$ 交换是相当自然的。）
	\item 验证命题 19 中的乘法确实使 $A\otimes_RB$ 成为一个 $R$-代数。
	\item 证明把环 $\RR$ 从 $\ZZ$ 扩张标量到 Gaussian 整数 $\ZZ[i]$ 得到的环同构于 $\CC$：
	\[
	\ZZ[i]\otimes_{\ZZ}\RR\cong\CC
	\]
	作为环同构。
	\item 设 $R$ 是交换环 $S$ 的子环，$x$ 是 $S$ 上的不定元。证明 $S[x]$ 与 $S\otimes_RR[x]$ 作为 $S$-代数同构。
	\item 设 $S$ 是含 $R$ 的交换环，且 $1_S=1_R$。令 $x_1,
	\ldots,x_n$ 是 $S$ 上独立的不定元。证明对多项式环 $R[x_1,
	\ldots,x_n]$ 的任意理想 $I$，有 $S$-代数同构
	\[
	S\otimes_R(R[x_1,
	\ldots,x_n]/I)
	\cong S[x_1,
	\ldots,x_n]/IS[x_1,
	\ldots,x_n].
	\]
	
	下一题说明本节末引入的环 $\CC\otimes_{\RR}\CC$ 同构于 $\CC\times\CC$。也可通过练习 26 与第 9.5 节命题 16 证明这一点，因为 $\CC\cong\RR[x]/(x^2+1)$。环 $\CC\times\CC$ 也在第 1 节练习 23 中讨论过。
	\item
	\begin{enumerate}[label=(\alph*)]
		\item 在命题 21 后的例子 $A=\CC\otimes_{\RR}\CC$ 中，写出两个元素
		\[
		a\cdot1+b\cdot e_2+c\cdot e_3+d\cdot e_4
		\quad\text{与}\quad
		a'\cdot1+b'\cdot e_2+c'\cdot e_3+d'\cdot e_4
		\]
		的乘法公式，其中 $1=1\otimes1$ 是 $A$ 的单位元。
		\item 令
		\[
		E_1=\frac12(1\otimes1+i\otimes i),\qquad
		E_2=\frac12(1\otimes1-i\otimes i).
		\]
		证明 $E_1E_2=0$、$E_1+E_2=1$，且 $E_j^2=E_j$（$j=1,2$）。$E_1,E_2$ 称为 $A$ 中的正交幂等元。推出 $A$ 作为环同构于两个主理想的直积：
		\[
		A\cong AE_1\times AE_2
		\]
		（参见第 7.6 节练习 1）。
		\item 证明映射
		\[
		\varphi:\CC\times\CC\longrightarrow\CC\times\CC,
		\qquad \varphi(z_1,z_2)=(z_1z_2,z_1\overline{z_2})
		\]
		是 $\RR$-双线性的，其中 $\overline{z_2}$ 表示 $z_2$ 的复共轭。
		\item 令 $\Phi$ 为由 (c) 中 $\varphi$ 得到的 $\RR$-模同态 $A\to\CC\times\CC$。证明 $\Phi(E_1)=(0,1)$ 且 $\Phi(E_2)=(1,0)$。还要证明 $\Phi$ 是 $\CC$-线性的，其中 $\CC$ 的作用在 $A$ 中作用于左张量因子，在 $\CC\times\CC$ 中作用于两个因子。推出 $\Phi$ 是满射。证明 $\Phi$ 是 $\CC$-代数同构。
	\end{enumerate}
\end{enumerate}


\section{正合列、投射模、内射模与平坦模}

研究代数对象 $B$（例如群、环或模）结构的一个基本结果是第一同构定理。它把 $B$ 的子对象（分别为正规子群、理想或子模）与 $B$ 的可能同态像联系起来。我们已经看到许多例子：由对这些“较小”组成部分的理解来分析 $B$ 的结构，例如通过确定二面体群 $D_8$ 的中心以及相应的中心商群来分析 $D_8$ 的结构。

在这些例子中，我们通常先给定 $B$，再构造同态 $\varphi$（往往通过指定 $\ker\varphi\leq B$ 来隐式给出），并研究 $\ker\varphi$ 与商 $B/\ker\varphi$。现在我们更详细地考虑反方向的问题：是否可以先指定这些“较小组成部分”。更确切地说，给定两个模 $A$ 与 $C$，是否存在一个模 $B$，它含有 $A$ 的一个同构副本，并且所得商模 $B/A$ 同构于 $C$。此时称 $B$ 是 $C$ 被 $A$ 的一个\emph{扩张}。于是自然要问：给定 $A$ 与 $C$，这样的 $B$ 有多少个；以及 $B$ 的性质在多大程度上由 $A$ 与 $C$ 的相应性质决定。群和环的情形当然也有类似问题。这就是第 3.4 节中首先对群讨论过的扩张问题。本节主要关心一个环 $R$ 上的左模；在必要时会指出对于其他结构，特别是非交换群，需要作怎样的修改。与前一节一样，本节中所有环都含有单位元 $1$。

我们先引入一个非常方便的记号。说 $A$ 同构于 $B$ 的一个子模，就是说存在一个单同态
\[
\psi:A\hookrightarrow B
\]
于是 $A\cong \psi(A)\leq B$。说 $C$ 同构于所得商，就是说存在一个满同态
\[
\varphi:B\twoheadrightarrow C
\]
且 $\ker\varphi=\psi(A)$。特别地，这给出一对同态
\[
A\xrightarrow{\psi}B\xrightarrow{\varphi}C
\]
并且 $\operatorname{Im}\psi=\ker\varphi$。满足这一性质的一对同态有如下名称。

\noindent\textbf{定义.}
\begin{enumerate}[label=(\arabic*)]
\item 一对同态
\[
X\xrightarrow{\alpha}Y\xrightarrow{\beta}Z
\]
称为（在 $Y$ 处）\emph{正合}，若
\[
\operatorname{Im}\alpha=\ker\beta.
\]
\item 一个同态列
\[
\cdots\longrightarrow X_{n-1}\longrightarrow X_n\longrightarrow X_{n+1}\longrightarrow\cdots
\]
称为\emph{正合列}，若它在每个处于一对同态之间的 $X_n$ 处都是正合的。
\end{enumerate}

用这个术语，上面的 $A\xrightarrow{\psi}B\xrightarrow{\varphi}C$ 在 $B$ 处正合。我们还可以用这个术语表示 $\psi$ 是单射、$\varphi$ 是满射这一事实。

\noindent\textbf{命题 22.} 设 $A,B,C$ 是某个环 $R$ 上的 $R$-模。则
\begin{enumerate}[label=(\arabic*)]
\item 序列
\[
0\longrightarrow A\xrightarrow{\psi}B
\]
在 $A$ 处正合，当且仅当 $\psi$ 是单射。
\item 序列
\[
B\xrightarrow{\varphi}C\longrightarrow0
\]
在 $C$ 处正合，当且仅当 $\varphi$ 是满射。
\end{enumerate}

\begin{proof}
唯一确定的同态 $0\to A$ 在 $A$ 中的像是 $0$。这正是 $\psi$ 的核，当且仅当 $\psi$ 是单射。类似地，唯一确定的零同态 $C\to0$ 的核是整个 $C$，它等于 $\varphi$ 的像，当且仅当 $\varphi$ 是满射。
\end{proof}

\noindent\textbf{推论 23.} 序列
\[
0\longrightarrow A\xrightarrow{\psi}B\xrightarrow{\varphi}C\longrightarrow0
\]
正合，当且仅当 $\psi$ 是单射、$\varphi$ 是满射，并且
\[
\operatorname{Im}\psi=\ker\varphi,
\]
也就是说，$B$ 是 $C$ 被 $A$ 的一个扩张。

\noindent\textbf{定义.} 正合列
\[
0\longrightarrow A\xrightarrow{\psi}B\xrightarrow{\varphi}C\longrightarrow0
\]
称为一个\emph{短正合列}。

用这个记号，扩张问题可以简洁地表述为：给定模 $A$ 与 $C$，确定所有短正合列
\[
0\longrightarrow A\longrightarrow B\longrightarrow C\longrightarrow0.
\tag{10.9}
\]
我们下面会看到，正合列记号在分析 $A$ 与 $C$ 的性质多大程度上决定 $B$ 的相应性质时也极为方便。若 $A,B,C$ 是用乘法记号书写的群，则序列 (10.9) 写作
\[
1\longrightarrow A\longrightarrow B\longrightarrow C\longrightarrow1,
\tag{10.9'}
\]
其中 $1$ 表示平凡群。命题 22 与推论 23 在作显然的记号改变后仍成立。

注意，任意正合列都可以写成一连串短正合列；因为说
\[
X\longrightarrow Y\longrightarrow Z
\]
在 $Y$ 处正合，等价于说
\[
0\longrightarrow \alpha(X)\longrightarrow Y\longrightarrow Y/\ker\beta\longrightarrow0
\]
是短正合列。

\paragraph{例}
\begin{enumerate}[label=(\arabic*)]
\item 给定模 $A$ 与 $C$，总可以构造它们的直和 $B=A\oplus C$，并有短正合列
\[
0\longrightarrow A\xrightarrow{\iota} A\oplus C\xrightarrow{\pi} C\longrightarrow0,
\]
其中 $\iota(a)=(a,0)$，$\pi(a,c)=c$。特别地，$C$ 被 $A$ 的扩张至少存在一个。

\item 作为前例的特殊情形，取两个 $\ZZ$-模 $A=\ZZ$ 与 $C=\ZZ/n\ZZ$，得到
\[
0\longrightarrow \ZZ\longrightarrow \ZZ\oplus(\ZZ/n\ZZ)\longrightarrow \ZZ/n\ZZ\longrightarrow0,
\]
这是 $\ZZ/n\ZZ$ 被 $\ZZ$ 的一个扩张。

$\ZZ/n\ZZ$ 被 $\ZZ$ 的另一个扩张由短正合列
\[
0\longrightarrow \ZZ\xrightarrow{\,n\,}\ZZ\xrightarrow{\pi}\ZZ/n\ZZ\longrightarrow0
\]
给出，其中 $n$ 表示乘以 $n$ 的映射 $x\mapsto nx$，$\pi$ 表示自然投影。注意，上面两个短正合列中间的模并不同构，尽管相应的 $A$ 项与 $C$ 项同构。因此，$\ZZ/n\ZZ$ 被 $\ZZ$ 的扩张至少有两种“本质上不同”或“不等价”的方式。

\item 若 $\varphi:B\to C$ 是任意同态，则可以形成正合列
\[
0\longrightarrow \ker\varphi\xrightarrow{\iota}B\xrightarrow{\varphi}\operatorname{Im}\varphi\longrightarrow0,
\]
其中 $\iota$ 是包含映射。特别地，若 $\varphi$ 满射，则序列 $\varphi:B\to C$ 可以扩张为一个以 $A=\ker\varphi$ 为左端项的短正合列。

\item 前例的一个特别重要的情形如下。设 $M$ 是 $R$-模，$S$ 是 $M$ 的一组生成元。令 $F(S)$ 是集合 $S$ 上的自由 $R$-模。由定理 6 得到的唯一 $R$-模同态 $F(S)\to M$（在 $S$ 上为恒等）是满射。若 $K$ 是它的核，则有短正合列
\[
0\longrightarrow K\longrightarrow F(S)\longrightarrow M\longrightarrow0.
\]
更一般地，当 $M$ 是任意群（可能非阿贝尔）时，上述短正合列（若用乘法记号则两端为 $1$）描述了 $M$ 的一个表示，其中 $K$ 是由定义 $M$ 的关系所生成的自由群 $F(S)$ 的正规子群（参见第 6.3 节）。

\item Klein 四元群被二阶循环群 $\ZZ_2$ 的两个“不等价”扩张 $G$ 为
\[
1\longrightarrow \ZZ_2\longrightarrow D_8\longrightarrow \ZZ_2\times\ZZ_2\longrightarrow1,
\]
以及
\[
1\longrightarrow \ZZ_2\longrightarrow Q_8\longrightarrow \ZZ_2\times\ZZ_2\longrightarrow1.
\]
在这两种情形中，$\ZZ_2$ 都被单射地映到 $G$ 的中心中，而 $\varphi$ 是 $G$ 到 $G/Z(G)$ 的自然投影。

Klein 四元群被 $\ZZ_2$ 的另外两个不等价扩张出现在 $G$ 为阿贝尔群 $\ZZ_2\times\ZZ_2\times\ZZ_2$ 或 $\ZZ_2\times\ZZ_4$ 时，只需选取适当的映射 $\psi$ 与 $\varphi$。
\end{enumerate}

上述例子 2 与 5 表明，对于固定的 $A$ 与 $C$，一般可能存在多个 $C$ 被 $A$ 的扩张。为了区分不同扩张，我们定义两个正合列之间的同态（与同构）。先回忆：包含若干同态的图称为交换的，若图中任意具有相同起点与终点的同态复合都相等；也就是说，沿图中同态路径所定义的复合映射只依赖于起点和终点，而不依赖于所选择的路径。

\noindent\textbf{定义.} 设
\[
0\longrightarrow A\longrightarrow B\longrightarrow C\longrightarrow0
\]
和
\[
0\longrightarrow A'\longrightarrow B'\longrightarrow C'\longrightarrow0
\]
是两个模的短正合列。
\begin{enumerate}[label=(\arabic*)]
\item 一个\emph{短正合列同态}是模同态三元组 $\alpha,\beta,\gamma$，使得下图交换：
\[
\begin{array}{ccccccccc}
0&\longrightarrow&A&\longrightarrow&B&\longrightarrow&C&\longrightarrow&0\\
 &&\alpha\downarrow&&\beta\downarrow&&\gamma\downarrow&&\\
0&\longrightarrow&A'&\longrightarrow&B'&\longrightarrow&C'&\longrightarrow&0.
\end{array}
\]
若 $\alpha,\beta,\gamma$ 都是同构，则称该同态为短正合列的同构；此时相应的扩张 $B$ 与 $B'$ 称为同构扩张。

\item 若 $A=A'$、$C=C'$，并且存在如 (1) 中的同构且它在 $A$ 与 $C$ 上均为恒等映射（即 $\alpha$ 与 $\gamma$ 是恒等映射），则称两个正合列\emph{等价}。此时相应的扩张 $B$ 与 $B'$ 称为等价扩张。
\end{enumerate}

若 $B$ 与 $B'$ 是同构扩张，则特别地，$B$ 与 $B'$ 作为 $R$-模同构；但还更多：存在 $B$ 与 $B'$ 之间的 $R$-模同构，它限制为 $A$ 到 $A'$ 的同构，并在商 $C$ 与 $C'$ 上诱导同构。对于给定的 $A$ 与 $C$，两个 $C$ 被 $A$ 的扩张等价这一条件还要更强：必须存在 $B$ 与 $B'$ 之间的 $R$-模同构，它在 $A$ 上限制为恒等映射，并在 $C$ 上诱导恒等映射。同构扩张的概念衡量在允许 $A$ 与 $C$ 被同构替换时，$C$ 被 $A$ 的扩张有多少种；等价扩张的概念衡量当 $A$ 与 $C$ 被刚性固定时，有多少种不同扩张。

乘法群短正合列 (10.9') 之间的同态与同构也类似定义。

容易验证，短正合列同态的复合仍是短正合列同态。同样，若三元组 $\alpha,\beta,\gamma$ 是同构（或等价），则 $\alpha^{-1},\beta^{-1},\gamma^{-1}$ 给出反向的同构（或等价）。因此，在任意短正合列集合上，“同构”（或“等价”）是等价关系。

\paragraph{例}
\begin{enumerate}[label=(\arabic*)]
\item 设 $m,n$ 是大于 $1$ 的整数。假设 $n\mid m$，令 $k=m/n$。可以从上面例 2 的 $\ZZ$-模短正合列定义到另一个短正合列的映射：
\[
\begin{array}{ccccccccc}
0&\longrightarrow&\ZZ&\xrightarrow{\,n\,}&\ZZ&\longrightarrow&\ZZ/n\ZZ&\longrightarrow&0\\
 &&\alpha\downarrow&&\beta\downarrow&&\gamma\downarrow&&\\
0&\longrightarrow&\ZZ/k\ZZ&\xrightarrow{\,\iota\,}&\ZZ/m\ZZ&\longrightarrow&\ZZ/n\ZZ&\longrightarrow&0,
\end{array}
\]
其中 $\alpha$ 与 $\beta$ 是自然投影，$\gamma$ 是恒等映射，$\iota$ 把 $a\bmod k$ 映到 $na\bmod m$，而右边的投影是 $\ZZ/m\ZZ$ 到其商 $(\ZZ/m\ZZ)/(n\ZZ/m\ZZ)\cong\ZZ/n\ZZ$ 的自然投影。容易验证这是短正合列同态。

\item 对短正合列
\[
0\longrightarrow \ZZ\xrightarrow{\,n\,}\ZZ\longrightarrow\ZZ/n\ZZ\longrightarrow0
\]
把每个模都用 $x\mapsto -x$ 映到自身，得到该正合列到自身的一个同构。这个同构不是正合列等价，因为它在第一个 $\ZZ$ 上不是恒等映射。

\item 推论 23 后例 1 与例 2 中的短正合列不同构，因为中间的扩张模不是同构的 $\ZZ$-模（阿贝尔群）。同样，例 5 中的两个扩张 $D_8$ 与 $Q_8$ 也不同构（因而不等价），尽管两列的两端项 $A$ 与 $C$ 相同。

\item 考虑映射
\[
\begin{array}{ccccccccc}
0&\longrightarrow&\ZZ/2\ZZ&\xrightarrow{\psi}&\ZZ/2\ZZ\oplus\ZZ/2\ZZ&\xrightarrow{\varphi}&\ZZ/2\ZZ&\longrightarrow&0\\
 &&\Vert&&\beta\downarrow&&\Vert&&\\
0&\longrightarrow&\ZZ/2\ZZ&\xrightarrow{\psi'}&\ZZ/2\ZZ\oplus\ZZ/2\ZZ&\xrightarrow{\varphi'}&\ZZ/2\ZZ&\longrightarrow&0,
\end{array}
\]
其中 $\psi$ 把 $\ZZ/2\ZZ$ 单射到直和的第一分量，$\varphi$ 投影到第二分量；而 $\psi'$ 嵌入到第二分量，$\varphi'$ 投影到第一分量。若 $\beta$ 是交换直和两个分量的映射，则图交换，于是得到两个正合列的一个等价；这个等价同构本身并不是恒等同构。

\item 我们展示两个同构但不等价的 $\ZZ$-模扩张。对 $i=1,2$ 定义
\[
0\longrightarrow \ZZ/2\ZZ\xrightarrow{\psi}\ZZ/4\ZZ\oplus\ZZ/2\ZZ\xrightarrow{\varphi_i}\ZZ/2\ZZ\oplus\ZZ/2\ZZ\longrightarrow0,
\]
其中两列中的 $\psi$ 都由 $1\mapsto (2,0)$ 给出，$\varphi_1$ 定义为
\[
\varphi_1(a\bmod4,b\bmod2)=(a\bmod2,b\bmod2),
\]
而
\[
\varphi_2(a\bmod4,b\bmod2)=(b\bmod2,a\bmod2).
\]
容易看出这两列都是短正合列。

这两个正合列之间的一个显然同构由三元组 $\operatorname{id},\operatorname{id},\gamma$ 给出，其中 $\gamma:\ZZ/2\ZZ\oplus\ZZ/2\ZZ\to\ZZ/2\ZZ\oplus\ZZ/2\ZZ$ 是交换两个直接因子的映射。

现在说明这两个同构正合列不等价。由于 $\varphi_1(0,1)=(0,1)$，任意从第一列到第二列的等价 $\operatorname{id},\beta,\operatorname{id}$ 必须把 $(0,1)\in \ZZ/4\ZZ\oplus\ZZ/2\ZZ$ 映到 $\ZZ/4\ZZ\oplus\ZZ/2\ZZ$ 中在 $\varphi_2$ 下映到 $(0,1)$ 的元素，即 $(1,0)$ 或 $(3,0)$。但这是不可能的，因为同构 $\beta$ 不能把二阶元素映到四阶元素。

换言之，涉及同一个扩张模 $B$ 的等价，是 $B$ 的自同构，且它在 $\psi(A)$ 与 $B/\psi(A)$ 上都限制为恒等。任意这样的 $B=\ZZ/4\ZZ\oplus\ZZ/2\ZZ$ 的自同构必须固定陪集 $(0,1)+\psi(A)$，因为这是唯一含有二阶元素的非恒等陪集。因此，把这个陪集送到 $C$ 中不同元素的映射给出不等价扩张。特别地，仍以 $A=\ZZ/2\ZZ$、$B=\ZZ/4\ZZ\oplus\ZZ/2\ZZ$、$C=\ZZ/2\ZZ\oplus\ZZ/2\ZZ$ 为模，还有第三个不等价扩张，它把陪集 $(0,1)+\psi(A)$ 映到 $C$ 中的元素 $(1,1)$。

类似推理表明，若 $B\cong D_8$，则 $\ZZ_2\times\ZZ_2$ 被 $\ZZ_2$ 的群扩张也有三个不等价但同构的情形（参见练习）。
\end{enumerate}

短正合列同态中的 $\alpha,\beta,\gamma$ 并非彼此独立。下一个结果给出这三个同态之间的一些关系。

\noindent\textbf{命题 24（短五引理）.} 设 $\alpha,\beta,\gamma$ 是短正合列的同态
\[
\begin{array}{ccccccccc}
0&\longrightarrow&A&\longrightarrow&B&\longrightarrow&C&\longrightarrow&0\\
 &&\alpha\downarrow&&\beta\downarrow&&\gamma\downarrow&&\\
0&\longrightarrow&A'&\longrightarrow&B'&\longrightarrow&C'&\longrightarrow&0.
\end{array}
\]
则：
\begin{enumerate}[label=(\arabic*)]
\item 若 $\alpha$ 与 $\gamma$ 单射，则 $\beta$ 也单射。
\item 若 $\alpha$ 与 $\gamma$ 满射，则 $\beta$ 也满射。
\item 若 $\alpha$ 与 $\gamma$ 同构，则 $\beta$ 也同构（于是两个序列同构）。
\end{enumerate}

\noindent\textbf{说明.} 这些结果对（可能非阿贝尔的）群短正合列也成立，正如证明所显示的那样。

\begin{proof}
我们证明 (1)，把 (2) 的证明留作练习；(3) 立即由 (1) 与 (2) 推出。设 $\alpha$ 与 $\gamma$ 单射，并设 $b\in B$ 满足 $\beta(b)=0$。记第一短正合列中的同态为 $\psi:A\to B$ 与 $\varphi:B\to C$。由于 $\beta(b)=0$，其在商 $C'$ 中的像也为 $0$。由图的交换性可知 $\gamma(\varphi(b))=0$。因为 $\gamma$ 单射，得到 $\varphi(b)=0$，即 $b\in\ker\varphi$。由第一列正合，$b$ 属于 $\psi$ 的像，所以 $b=\psi(a)$ 对某个 $a\in A$ 成立。再由图交换，$\alpha(a)$ 在 $B'$ 中的像等于 $\beta(\psi(a))=\beta(b)=0$。但 $\alpha$ 与 $A'\to B'$ 都是单射，因此 $a=0$。最后 $b=\psi(a)=\psi(0)=0$，故 $\beta$ 确为单射。
\end{proof}

我们已经看到，$C$ 被 $A$ 的扩张至少有一个，即直接和 $B=A\oplus C$。在这种情况下，模 $B$ 除了子模 $A$ 外，还含有一个同构于 $C$ 的子模 $C'$（即 $C'=0\oplus C$），而这个与 $A$ 互补的子模把 $B$ “分裂”为直和。在群的情形中，若 $B$ 中正规子群有一个子群补 $C'$，则 $B$ 是一个半直积（参见第 5 章第 5 节）。在模的语境中，$B$ 是直和这一事实反映了底层群结构是阿贝尔的；对阿贝尔群，半直积就是直积。在任一情形中，相应的短正合列都称为“分裂”的。

\noindent\textbf{定义.}
\begin{enumerate}[label=(\arabic*)]
\item 设 $R$ 是环，
\[
0\longrightarrow A\xrightarrow{\psi}B\xrightarrow{\varphi}C\longrightarrow0
\]
是 $R$-模短正合列。若 $B$ 中存在 $\psi(A)$ 的一个 $R$-模补，则称该序列\emph{分裂}。此时，直到同构，$B=A\oplus C$；更精确地说，$B=\psi(A)\oplus C'$ 对某个子模 $C'$ 成立，并且 $\varphi$ 把 $C'$ 同构地映到 $C$，即 $\varphi(C')\cong C$。
\item 若
\[
1\longrightarrow A\xrightarrow{\psi}B\xrightarrow{\varphi}C\longrightarrow1
\]
是群短正合列，若 $B$ 中存在 $\psi(A)$ 的一个子群补，则称该序列\emph{分裂}。此时，直到同构，$B=A\rtimes C$；更精确地说，$B=\psi(A)\rtimes C'$ 对某个子群 $C'$ 成立，并且 $\varphi$ 把 $C'$ 同构地映到 $C$。
\end{enumerate}
在这两种情形中，扩张 $B$ 称为 $C$ 被 $A$ 的一个\emph{分裂扩张}。

扩张是否分裂的问题，就是 $B$ 中是否存在一个由 $\varphi$ 同构到 $C$ 的 $\psi(A)$ 的补。因此分裂扩张的概念也可以等价地用同态语言表述。

\noindent\textbf{命题 25.} $R$-模短正合列
\[
0\longrightarrow A\xrightarrow{\psi}B\xrightarrow{\varphi}C\longrightarrow0
\]
分裂，当且仅当存在 $R$-模同态 $\mu:C\to B$ 使
\[
\varphi\circ\mu=\operatorname{id}_C.
\]
类似地，群短正合列
\[
1\longrightarrow A\xrightarrow{\psi}B\xrightarrow{\varphi}C\longrightarrow1
\]
分裂，当且仅当存在群同态 $\mu:C\to B$ 使 $\varphi\circ\mu$ 是 $C$ 上的恒等映射。

\begin{proof}
这直接来自定义：若给定 $\mu$，令 $C'=\mu(C)\leq B$；若给定补 $C'$，则定义 $\mu=\varphi^{-1}:C\to C'\hookrightarrow B$。
\end{proof}

\noindent\textbf{定义.} 在命题 25 的记号下，任意集合映射 $\mu:C\to B$，若满足 $\varphi\circ\mu=\operatorname{id}_C$，称为 $\varphi$ 的一个\emph{截面}。若 $\mu$ 是命题 25 中的同态，则称 $\mu$ 为该序列的一个\emph{分裂同态}。

注意，$\varphi$ 的截面不过是给商 $B/\ker\varphi\cong C$ 在 $B$ 中选择一组陪集代表。若该截面是（分裂）同态，则这组陪集代表在 $B$ 中构成子模（分别为子群），于是该子模（分别为子群）给出 $\psi(A)$ 的一个补。

\paragraph{例}
\begin{enumerate}[label=(\arabic*)]
\item 分裂短正合列
\[
0\longrightarrow A\longrightarrow A\oplus C\longrightarrow C\longrightarrow0
\]
有显然的分裂同态 $\mu(c)=(0,c)$。

\item $\ZZ/n\ZZ$ 被 $\ZZ$ 的扩张
\[
0\longrightarrow \ZZ\longrightarrow \ZZ\oplus(\ZZ/n\ZZ)\longrightarrow \ZZ/n\ZZ\longrightarrow0
\]
是分裂的；分裂同态把 $\ZZ/n\ZZ$ 同构地映到直和的第二因子。另一方面，$\ZZ$-模正合列
\[
0\longrightarrow \ZZ\xrightarrow{\,n\,}\ZZ\longrightarrow\ZZ/n\ZZ\longrightarrow0
\]
不分裂，因为不存在从 $\ZZ/n\ZZ$ 到 $\ZZ$ 的非零同态。

\item $D_8$ 与 $Q_8$ 都不是 $\ZZ_2\times\ZZ_2$ 被 $\ZZ_2$ 的分裂扩张，因为在这两个群中，中心都没有子群补（第 2.5 节给出了这些群的子群结构）。

\item 群 $D_8$ 是 $\ZZ_2$ 被 $\ZZ_4$ 的分裂扩张，即存在分裂短正合列
\[
1\longrightarrow \ZZ_4\longrightarrow D_8\longrightarrow\ZZ_2\longrightarrow1.
\]
用 $D_8$ 的通常生成元 $r,s$，这就是
\[
1\longrightarrow\langle r\rangle\longrightarrow D_8\longrightarrow\langle s\rangle\longrightarrow1.
\]
这里包含映射把 $\langle r\rangle$ 嵌入 $D_8$，而投影 $\pi:r^as^b\mapsto s^b$ 给出到商 $D_8/\langle r\rangle\cong\ZZ_2$ 的映射。分裂同态 $\mu$ 把 $\langle s\rangle$ 同构地映到 $D_8$ 中的补 $\langle s\rangle$。等价地，$D_8$ 是正规子群 $\langle r\rangle\cong\ZZ_4$ 与 $\langle s\rangle\cong\ZZ_2$ 的半直积。

另一方面，$Q_8$ 也是 $\ZZ_2$ 被 $\ZZ_4$ 的一个扩张（例如 $\langle i\rangle\cong\ZZ_4$，其商同构于 $\ZZ_2$），但 $Q_8$ 不是 $\ZZ_2$ 被 $\ZZ_4$ 的分裂扩张：$Q_8$ 的任意四阶循环子群都没有补。

第 5.5 节包含更多群的分裂扩张例子。
\end{enumerate}

命题 25 表明，$C$ 被 $A$ 的扩张 $B$ 是分裂扩张，当且仅当从 $B$ 到商 $C$ 的投影 $\varphi:B\to C$ 有一个分裂同态 $\mu$。下一个命题特别表明，对模来说，这等价于序列左端的 $\psi$ 有一个分裂同态。

\noindent\textbf{命题 26.} 设
\[
0\longrightarrow A\xrightarrow{\psi}B\xrightarrow{\varphi}C\longrightarrow0
\]
是模的短正合列（分别地，设
\[
1\longrightarrow A\xrightarrow{\psi}B\xrightarrow{\varphi}C\longrightarrow1
\]
是群的短正合列）。则 $B=\psi(A)\oplus C'$ 对某个 $B$ 的子模 $C'$ 成立且 $\varphi(C')\cong C$（分别地，$B=\psi(A)\times C'$ 对某个 $B$ 的子群 $C'$ 成立且 $\varphi(C')\cong C$），当且仅当存在同态 $\lambda:B\to A$ 使
\[
\lambda\circ\psi=\operatorname{id}_A.
\]

\begin{proof}
证明与命题 25 类似。若给定 $\lambda$，定义 $C'=\ker\lambda\leq B$。若给定 $C'$，则由 $B=\psi(A)\oplus C'$（或群情形的直积）定义
\[
\lambda:B=\psi(A)\oplus C'\longrightarrow A,
\qquad \lambda(\psi(a)+c')=a.
\]
在群情形中，此时 $C'=\ker\lambda$ 在 $B$ 中正规，所以 $C'$ 是 $\psi(A)$ 的正规补，从而 $B$ 是 $\psi(A)$ 与 $C'$ 的直积（参见第 5.4 节定理 9）。
\end{proof}

命题 26 表明，对一般群扩张，在序列左端存在分裂同态 $\lambda$ 比扩张分裂这一条件更强：在这种情况下扩张群是直积，而不仅仅是半直积。在模的语境中，这两个概念等价，这再次反映了其底层群的阿贝尔性质，在那里半直积总是直积。

\subsection*{模与 $\Hom_R(D,-)$}

设 $R$ 是带单位元的环，并设 $R$-模 $M$ 是 $N$ 被 $L$ 的扩张，对应短正合列为
\[
0\longrightarrow L\xrightarrow{\psi}M\xrightarrow{\varphi}N\longrightarrow0.
\]
自然要问：$L$ 与 $N$ 的性质是否推出扩张 $M$ 的相应性质。我们首先考虑如下情形：从某个固定 $R$-模 $D$ 到 $L$ 或 $N$ 的 $R$-模同态，是否蕴含也存在从 $D$ 到 $M$ 的 $R$-模同态。

若给定从 $D$ 到 $L$ 的同态，则得到到 $M$ 的同态十分直接：若 $f\in\Hom_R(D,L)$，则复合
\[
f'=\psi\circ f
\]
是从 $D$ 到 $M$ 的 $R$-模同态。图示为
\[
\begin{array}{ccc}
D&\xrightarrow{f'}&M\\
&\searrow_f&\uparrow\psi\\[-2mm]
&&L.
\end{array}
\]
换言之，与 $\psi$ 复合诱导映射
\[
\psi':\Hom_R(D,L)\longrightarrow\Hom_R(D,M),
\qquad f\longmapsto f'=\psi\circ f.
\]
回忆命题 2，$\Hom_R(D,L)$ 与 $\Hom_R(D,M)$ 都是阿贝尔群。

\noindent\textbf{命题 27.} 设 $D,L,M$ 是 $R$-模，$\psi:L\to M$ 是 $R$-模同态。则映射
\[
\psi':\Hom_R(D,L)\longrightarrow\Hom_R(D,M),
\qquad f\longmapsto \psi\circ f
\]
是阿贝尔群同态。若 $\psi$ 是单射，则 $\psi'$ 也是单射；即若
\[
0\longrightarrow L\xrightarrow{\psi}M
\]
正合，则
\[
0\longrightarrow \Hom_R(D,L)\xrightarrow{\psi'}\Hom_R(D,M)
\]
也正合。

\begin{proof}
$\psi'$ 是同态是直接的。若 $\psi$ 单射，则 $D$ 到 $L$ 的两个不同同态 $f,g$ 给出 $D$ 到 $M$ 的两个不同复合 $\psi\circ f$ 与 $\psi\circ g$；这就是说 $\psi'$ 也单射。
\end{proof}

尽管从到子模 $L$ 的同态得到到 $M$ 的同态很直接，对于到商 $N$ 的同态情形则不那么显然。更确切地说，给定 $R$-模同态 $f:D\to N$，问题是是否存在 $R$-模同态 $F:D\to M$，它把 $f$ 延拓或提升到 $M$，即使下图交换：
\[
\begin{array}{ccc}
&&D\\
&F\swarrow&\downarrow f\\[-1mm]
M&\xrightarrow{\varphi}&N.
\end{array}
\]
与前面一样，与 $\varphi$ 复合诱导阿贝尔群同态
\[
\varphi':\Hom_R(D,M)\longrightarrow\Hom_R(D,N),
\qquad F\longmapsto F'=\varphi\circ F.
\]
用 $\varphi'$ 的语言，$f$ 到 $N$ 的同态可提升到 $M$ 中的同态，当且仅当 $f$ 在 $\varphi'$ 的像中。

一般来说，不一定能把从 $D$ 到 $N$ 的同态提升为从 $D$ 到 $M$ 的同态。例如，考虑前面非分裂正合列
\[
0\longrightarrow\ZZ\xrightarrow{\,2\,}\ZZ\longrightarrow\ZZ/2\ZZ\longrightarrow0.
\]
令 $D=\ZZ/2\ZZ$，令 $f$ 是 $D$ 到 $N$ 的恒等映射。从 $D$ 到 $M=\ZZ$ 的任意同态 $F$ 都必须把 $D$ 映到 $0$（因为 $\ZZ$ 没有二阶元素），所以 $\pi\circ F$ 把 $D$ 映到 $N$ 中的 $0$，特别地 $\pi\circ F\neq f$。用映射 $\varphi'$ 的语言，这说明若
\[
M\xrightarrow{\varphi}N\longrightarrow0
\]
正合，则
\[
\Hom_R(D,M)\xrightarrow{\varphi'}\Hom_R(D,N)\longrightarrow0
\]
不一定正合。

上述结果可以整齐地概括为下列定理的一部分。

\noindent\textbf{定理 28.} 设 $D,L,M,N$ 是 $R$-模。若
\[
0\longrightarrow L\xrightarrow{\psi}M\xrightarrow{\varphi}N\longrightarrow0
\]
正合，则相应序列
\[
0\longrightarrow \Hom_R(D,L)\xrightarrow{\psi'}\Hom_R(D,M)\xrightarrow{\varphi'}\Hom_R(D,N)
\tag{10.10}
\]
正合。

同态 $f:D\to N$ 可以提升为同态 $F:D\to M$，当且仅当 $f\in\Hom_R(D,N)$ 位于 $\varphi'$ 的像中。一般地，$\varphi':\Hom_R(D,M)\to\Hom_R(D,N)$ 不必满射；$\varphi'$ 满射当且仅当每个从 $D$ 到 $N$ 的同态都可提升为从 $D$ 到 $M$ 的同态，此时序列 (10.10) 可扩张为一个短正合列。

序列 (10.10) 对所有 $R$-模 $D$ 都正合，当且仅当序列
\[
0\longrightarrow L\xrightarrow{\psi}M\xrightarrow{\varphi}N
\]
正合。

\begin{proof}
第一部分中唯一尚未证明的是 (10.10) 在 $\Hom_R(D,M)$ 处的正合性，即 $\ker\varphi'=\operatorname{Im}\psi'$。设 $F:D\to M$ 是 $\Hom_R(D,M)$ 中属于 $\ker\varphi'$ 的元素，即作为 $D$ 到 $N$ 的同态有 $\varphi\circ F=0$。若 $d\in D$，则 $\varphi(F(d))=0$，所以 $F(d)\in\ker\varphi$。由给定短正合列正合，$\ker\varphi=\operatorname{Im}\psi$，故存在 $l\in L$ 使 $F(d)=\psi(l)$。由于 $\psi$ 单射，这个 $l$ 唯一，于是得到良定义映射 $F':D\to L$，由 $F'(d)=l$ 给出。容易验证 $F'$ 是同态，即 $F'\in\Hom_R(D,L)$。又 $\psi(F'(d))=\psi(l)=F(d)$，所以 $F=\psi'(F')$，从而 $F$ 在 $\psi'$ 的像中。这证明 $\ker\varphi'\subseteq\operatorname{Im}\psi'$。反过来，若 $F$ 在 $\psi'$ 的像中，则 $F=\psi'(F')$ 对某个 $F'\in\Hom_R(D,L)$ 成立。于是任意 $d\in D$ 有
\[
\varphi(F(d))=\varphi(\psi(F'(d)))=0
\]
因为 $\varphi\circ\psi=0$。故 $F\in\ker\varphi'$，得到反包含。

最后的断言中，注意在证明 (10.10) 正合时并未用到 $\varphi$ 的满射性，因此“若”方向已经证明。反过来，假设 (10.10) 对所有 $R$-模 $D$ 都正合。一般地，对任意左 $R$-模 $X$，有 $\Hom_R(R,X)\cong X$，同构由把同态映到它在 $1\in R$ 上的值给出。取 $D=R$，立即得到
\[
0\longrightarrow L\xrightarrow{\psi}M\xrightarrow{\varphi}N
\]
的正合性。
\end{proof}

由定理 28，序列
\[
0\longrightarrow \Hom_R(D,L)\xrightarrow{\psi'}\Hom_R(D,M)\xrightarrow{\varphi'}\Hom_R(D,N)\longrightarrow0
\tag{10.11}
\]
一般不是短正合列，因为 $\varphi'$ 不一定满射。序列 (10.11) 是否正合，精确地衡量了从 $D$ 到 $M$ 的同态在多大程度上由从 $D$ 到 $L$ 与从 $D$ 到 $N$ 的同态对唯一确定。更确切地说，序列 (10.11) 正合当且仅当存在由 $F\mapsto(g,f)$ 给出的双射，其中 $F:D\to M$ 是同态，$g:D\to L$ 与 $f:D\to N$ 是同态，并且 $F|_{\psi(L)}=\psi'(g)$、$f=\varphi'(F)$。

序列 (10.11) 正合的一个情形是原序列
\[
0\longrightarrow L\longrightarrow M\longrightarrow N\longrightarrow0
\]
为分裂正合列，也就是 $M=L\oplus N$。此时序列 (10.11) 也是分裂短正合列，正如下列命题的第一部分所示。

\noindent\textbf{命题 29.} 设 $D,L,N$ 是 $R$-模。则
\begin{enumerate}[label=(\arabic*)]
\item
\[
\Hom_R(D,L\oplus N)\cong\Hom_R(D,L)\oplus\Hom_R(D,N),
\]
\item
\[
\Hom_R(L\oplus N,D)\cong\Hom_R(L,D)\oplus\Hom_R(N,D).
\]
\end{enumerate}

\begin{proof}
令 $\pi_1:L\oplus N\to L$ 为到 $L$ 的自然投影，$\pi_2$ 为到 $N$ 的自然投影。若 $f\in\Hom_R(D,L\oplus N)$，则复合 $\pi_1\circ f$ 与 $\pi_2\circ f$ 分别给出 $\Hom_R(D,L)$ 与 $\Hom_R(D,N)$ 中的元素。这定义了从 $\Hom_R(D,L\oplus N)$ 到 $\Hom_R(D,L)\oplus\Hom_R(D,N)$ 的映射，容易验证它是同态。反过来，给定 $f_1\in\Hom_R(D,L)$ 与 $f_2\in\Hom_R(D,N)$，定义 $f\in\Hom_R(D,L\oplus N)$ 为
\[
f(d)=(f_1(d),f_2(d)).
\]
这给出反向同态，并且与前一映射互为逆映射，从而证明 (1)。 (2) 的证明类似，留作练习。
\end{proof}

命题 29 的结果可由归纳立即推广到任意有限个 $R$-模的直和。通常说 $\Hom$ 在任一变量上与有限直和交换。对于无限直和情形，情况更复杂：若把 $L\oplus N$ 换成任意直和，则 (1) 仍成立，但右端的直和需换成直积（练习 13 表明必须使用直积）；(2) 在两边直和均换成直积时仍成立。

这个命题表明，若
\[
0\longrightarrow L\xrightarrow{\psi}M\xrightarrow{\varphi}N\longrightarrow0
\]
是 $R$-模的分裂短正合列，则对每个 $R$-模 $D$，
\[
0\longrightarrow\Hom_R(D,L)\xrightarrow{\psi'}\Hom_R(D,M)\xrightarrow{\varphi'}\Hom_R(D,N)\longrightarrow0
\]
也是阿贝尔群的分裂短正合列。练习 14 说明其逆命题成立：若对每个 $R$-模 $D$，上述 $\Hom$ 序列都正合，则原序列是分裂短正合列；这还推出，若原 $\Hom$ 序列对每个 $D$ 都正合，则事实上它对每个 $D$ 都分裂正合。

命题 29 从 $L,M,N$ 的角度给出了序列 (10.11) 正合的一种情形。下一个结果采用略微不同的视角，刻画那些模 $D$：对它们而言，定理 28 中的序列 (10.10) 总能扩张为短正合列。

\noindent\textbf{命题 30.} 设 $P$ 是 $R$-模。下列条件等价：
\begin{enumerate}[label=(\arabic*)]
\item 对任意 $R$-模 $L,M,N$，若
\[
0\longrightarrow L\longrightarrow M\longrightarrow N\longrightarrow0
\]
是短正合列，则
\[
0\longrightarrow\Hom_R(P,L)\longrightarrow\Hom_R(P,M)\longrightarrow\Hom_R(P,N)\longrightarrow0
\]
也是短正合列。
\item 对任意 $R$-模 $M,N$，若 $M\to N\to0$ 正合，则每个从 $P$ 到 $N$ 的 $R$-模同态都能提升为到 $M$ 的 $R$-模同态；也就是说，给定 $f\in\Hom_R(P,N)$，存在提升 $F\in\Hom_R(P,M)$ 使下图交换：
\[
\begin{array}{ccc}
&&P\\
&F\swarrow&\downarrow f\\[-1mm]
M&\longrightarrow&N\longrightarrow0.
\end{array}
\]
\item 若 $P$ 是 $R$-模 $M$ 的商，则 $P$ 同构于 $M$ 的一个直和因子；也就是说，每个短正合列
\[
0\longrightarrow L\longrightarrow M\longrightarrow P\longrightarrow0
\]
都分裂。
\item $P$ 是某个自由 $R$-模的直和因子。
\end{enumerate}

\begin{proof}
(1) 与 (2) 的等价正是定理 28 中一项结果的重述。

现在假设 (2) 成立，并令
\[
0\longrightarrow L\xrightarrow{\psi}M\xrightarrow{\varphi}P\longrightarrow0
\]
正合。由 (2)，从 $P$ 到 $P$ 的恒等映射可提升为同态 $\mu:P\to M$，使得 $\varphi\circ\mu=\operatorname{id}_P$。因此 $\mu$ 是该序列的分裂同态，这证明 (3)。

每个模 $P$ 都是某个自由模的商，例如以 $P$ 的元素集合为基的自由模。因此存在正合列
\[
0\longrightarrow\ker\varphi\longrightarrow \mathcal F\xrightarrow{\varphi}P\longrightarrow0
\]
其中 $\mathcal F$ 是自由 $R$-模（参见推论 23 后例 4）。若 (3) 成立，则此序列分裂，所以 $\mathcal F$ 同构于 $\ker\varphi$ 与 $P$ 的直和，从而证明 (4)。

最后证明 (4) 蕴含 (2)。设 $P$ 是某个集合 $S$ 上自由模的直和因子，即
\[
\mathcal F(S)=P\oplus K,
\]
并给定 (2) 中从 $P$ 到 $N$ 的同态 $f$。令 $\pi$ 表示从 $\mathcal F(S)$ 到 $P$ 的自然投影，于是 $f\circ\pi$ 是从 $\mathcal F(S)$ 到 $N$ 的同态。对任意 $s\in S$，定义
\[
n_s=f\circ\pi(s)\in N,
\]
并取 $m_s\in M$ 使 $\varphi(m_s)=n_s$（这由 $\varphi$ 满射得到）。由自由模的泛性质（第 3 节定理 6），存在唯一 $R$-模同态 $F':\mathcal F(S)\to M$，满足 $F'(s)=m_s$。于是对所有 $s\in S$ 有
\[
\varphi\circ F'(s)=\varphi(m_s)=n_s=f\circ\pi(s),
\]
从而 $\varphi\circ F'=f\circ\pi$ 在 $\mathcal F(S)$ 上成立。定义 $F:P\to M$ 为
\[
F(d)=F'((d,0)).
\]
由于 $F$ 是包含 $P\hookrightarrow\mathcal F(S)$ 与同态 $F'$ 的复合，它是 $R$-模同态，并且
\[
\varphi\circ F(d)=\varphi\circ F'((d,0))=f\circ\pi((d,0))=f(d).
\]
故 $F$ 是所需提升。这证明 (4) 蕴含 (2)，从而完成证明。
\end{proof}

\noindent\textbf{定义.} 若 $R$-模 $P$ 满足命题 30 中任一等价条件，则称 $P$ 为\emph{投射模}。

命题 30 的第三个陈述可重述为：任何投影到 $P$ 上的模 $M$，都含有 $P$ 的一个同构副本作为直和因子。这解释了“投射”这一术语。

下列结果由命题 30（及其证明）立即得到。

\noindent\textbf{推论 31.} 自由模是投射模。有限生成模是投射模，当且仅当它是某个有限生成自由模的直和因子。每个模都是某个投射模的商。

若固定 $D$，则对任意 $R$-模 $X$，有一个相应的阿贝尔群 $\Hom_R(D,X)$。进一步，$R$-模同态 $\alpha:X\to Y$ 诱导阿贝尔群同态
\[
\alpha':\Hom_R(D,X)\longrightarrow\Hom_R(D,Y),
\qquad \alpha'(f)=\alpha\circ f.
\]
换言之，映射 $\Hom_R(D,-)$ 是从 $R$-模范畴到阿贝尔群范畴的协变函子（参见附录 II）。定理 28 表明，把这个函子作用于正合列
\[
0\longrightarrow L\xrightarrow{\psi}M\xrightarrow{\varphi}N\longrightarrow0
\]
的各项，会得到正合列
\[
0\longrightarrow\Hom_R(D,L)\xrightarrow{\psi'}\Hom_R(D,M)\xrightarrow{\varphi'}\Hom_R(D,N).
\]
这称为 $\Hom_R(D,-)$ 是\emph{左正合函子}。由命题 30，函子 $\Hom_R(D,-)$ 正合（即总把短正合列变成短正合列），当且仅当 $D$ 是投射的。总结为：

\noindent\textbf{推论 32.} 若 $D$ 是 $R$-模，则函子 $\Hom_R(D,-)$ 从 $R$-模范畴到阿贝尔群范畴是左正合的。它正合当且仅当 $D$ 是投射 $R$-模。

注意，若 $\Hom_R(D,-)$ 把短正合列变为短正合列，则它也把任意长度的正合列变为正合列，因为任意正合列都可以分解成一连串短正合列。

如我们所见，函子 $\Hom_R(D,-)$ 一般在右端不正合。衡量这类函子偏离正合的程度，会引出第 17 章中讨论的“同调代数”概念。

\paragraph{例}
\begin{enumerate}[label=(\arabic*)]
\item 我们将在第 11.1 节看到，若 $R=F$ 是域，则每个 $F$-模都是投射的（尽管那里只证明有限生成模情形）。

\item 由推论 31，$\ZZ$ 是投射 $\ZZ$-模。这也可以直接看出：设 $f:\ZZ\to N$ 是同态，且 $M\xrightarrow{\varphi}N\to0$ 正合。同态 $f$ 完全由 $n=f(1)$ 决定。取 $m\in M$ 使 $\varphi(m)=n$，先定义 $F(1)=m$，再按加性延拓到整个 $\ZZ$，便得到 $f$ 的提升 $F:\ZZ\to M$。

由命题 30 的第一条，由于 $\ZZ$ 是投射的，若
\[
0\longrightarrow L\longrightarrow M\longrightarrow N\longrightarrow0
\]
是 $\ZZ$-模正合列，则
\[
0\longrightarrow\Hom_\ZZ(\ZZ,L)\longrightarrow\Hom_\ZZ(\ZZ,M)\longrightarrow\Hom_\ZZ(\ZZ,N)\longrightarrow0
\]
也是正合列。这也可直接由阿贝尔群同构 $\Hom_\ZZ(\ZZ,M)\cong M$ 看出；在这个同构下，上面两个正合列本质上相同。

\item 自由 $\ZZ$-模没有非零有限阶元素，因此非零有限阿贝尔群不可能同构于自由模的子模。由推论 31，可知非零有限阿贝尔群都不是投射 $\ZZ$-模。

\item 特别地，对 $n\geq2$，$\ZZ$-模 $\ZZ/n\ZZ$ 不是投射的。由定理 28 必然可以找到一个短正合列，使对其作用函子 $\Hom_\ZZ(\ZZ/n\ZZ,-)$ 后右端不再正合。一个这样的序列就是前面例 2 的正合列
\[
0\longrightarrow\ZZ\xrightarrow{\,n\,}\ZZ\longrightarrow\ZZ/n\ZZ\longrightarrow0.
\]
注意 $\Hom_\ZZ(\ZZ/n\ZZ,\ZZ)=0$，因为不存在从 $\ZZ/n\ZZ$ 到 $\ZZ$ 的非零 $\ZZ$-模同态。又容易看出 $\Hom_\ZZ(\ZZ/n\ZZ,\ZZ/n\ZZ)\cong\ZZ/n\ZZ$：每个同态 $f$ 由 $f(1)=a\in\ZZ/n\ZZ$ 唯一确定，反之给定任意 $a$ 都有唯一同态 $f_a$ 满足 $f_a(1)=a$，映射 $f_a\mapsto a$ 即给出同构。作用 $\Hom_\ZZ(\ZZ/n\ZZ,-)$ 后得到
\[
0\longrightarrow0\longrightarrow0\longrightarrow\ZZ/n\ZZ\longrightarrow0,
\]
它在唯一非零项处不正合。

\item 由于 $\QQ/\ZZ$ 是挠 $\ZZ$-模，它不是自由模的子模，因而不是投射的。也注意到正合列
\[
0\longrightarrow\ZZ\longrightarrow\QQ\longrightarrow\QQ/\ZZ\longrightarrow0
\]
不分裂，因为 $\QQ$ 不含同构于 $\QQ/\ZZ$ 的子模。

\item $\ZZ$-模 $\QQ$ 不是投射的（参见练习）。

\item 我们将在第 12 章看到，有限生成 $\ZZ$-模投射当且仅当它自由。

\item 令 $R$ 为交换环 $\ZZ/2\ZZ\times\ZZ/2\ZZ$，其加法与乘法按分量进行。若 $P_1$ 与 $P_2$ 分别是由 $(1,0)$ 与 $(0,1)$ 生成的主理想，则
\[
R=P_1\oplus P_2,
\]
故由命题 30，$P_1$ 与 $P_2$ 都是投射 $R$-模。二者都不是自由模，因为任意自由模的阶都应是 $4$ 的倍数。

\item 两个投射模的直和仍是投射模（参见练习 3）。

\item 我们将在第 VI 部分看到，若 $F$ 是任意域且 $n\in\ZZ^+$，则所有 $n\times n$ 矩阵构成的环 $R=M_n(F)$ 具有这样的性质：每个 $R$-模都是投射的。我们还会看到，若 $G$ 是阶为 $n$ 的有限群，并且域 $F$ 中 $n\neq0$，则群环 $FG$ 也具有每个模都是投射的性质。
\end{enumerate}

\subsection*{内射模与 $\Hom_R(-,D)$}

若
\[
0\longrightarrow L\xrightarrow{\psi}M\xrightarrow{\varphi}N\longrightarrow0
\]
是 $R$-模短正合列，则除了考虑从某个 $R$-模 $D$ 到 $L$ 或 $N$ 的映射，以及它们在多大程度上决定从 $D$ 到 $M$ 的映射外，还可以考虑“对偶”问题：从 $L$ 或 $N$ 到 $D$ 的映射。此时，从 $N$ 到 $D$ 的映射很容易处理：一个 $R$-模映射 $N\to D$ 立即通过与 $\varphi$ 复合给出一个 $M\to D$ 的映射。容易验证这定义了阿贝尔群的单同态
\[
\varphi':\Hom_R(N,D)\longrightarrow\Hom_R(M,D),
\qquad f\longmapsto f'=f\circ\varphi.
\]
换言之，若
\[
M\xrightarrow{\varphi}N\longrightarrow0
\]
正合，则
\[
0\longrightarrow\Hom_R(N,D)\xrightarrow{\varphi'}\Hom_R(M,D)
\]
正合。注意，这里同态群上的诱导映射方向与原来的映射方向相反。

另一方面，给定 $R$-模同态 $f:L\to D$，不一定能把 $f$ 扩张为从 $M$ 到 $D$ 的映射 $F$；也就是说，给定 $f$，不一定存在 $F$ 使图
\[
\begin{array}{ccc}
L&\xrightarrow{\psi}&M\\
&f\searrow&\downarrow F\\[-1mm]
&&D
\end{array}
\]
交换。例如，考虑 $\ZZ$-模正合列
\[
0\longrightarrow\ZZ\xrightarrow{\,2\,}\ZZ\longrightarrow\ZZ/2\ZZ\longrightarrow0,
\]
其中 $\psi$ 是乘以 $2$，$\varphi$ 是自然投影。取 $D=\ZZ/2\ZZ$，令 $f:\ZZ\to\ZZ/2\ZZ$ 为第一项 $\ZZ$ 的模 $2$ 化简。第二项 $\ZZ$ 到 $\ZZ/2\ZZ$ 的非零同态只有一个（即模 $2$ 化简），但此 $F$ 不能提升 $f$，因为 $F\circ\psi(\ZZ)=F(2\ZZ)=0$，所以 $F\circ\psi\neq f$。

与 $\psi$ 复合诱导阿贝尔群同态
\[
\psi':\Hom_R(M,D)\longrightarrow\Hom_R(L,D),
\]
用 $\psi'$ 的语言，同态 $f\in\Hom_R(L,D)$ 可提升为从 $M$ 到 $D$ 的同态，当且仅当 $f$ 位于 $\psi'$ 的像中。上面的例子表明，若
\[
0\longrightarrow L\xrightarrow{\psi}M
\]
正合，则
\[
\Hom_R(M,D)\xrightarrow{\psi'}\Hom_R(L,D)\longrightarrow0
\]
不一定正合。

这些结果可概括为定理 28 的对偶版本。

\noindent\textbf{定理 33.} 设 $D,L,M,N$ 是 $R$-模。若
\[
0\longrightarrow L\xrightarrow{\psi}M\xrightarrow{\varphi}N\longrightarrow0
\]
正合，则相应序列
\[
0\longrightarrow\Hom_R(N,D)\xrightarrow{\varphi'}\Hom_R(M,D)\xrightarrow{\psi'}\Hom_R(L,D)
\tag{10.12}
\]
正合。

同态 $f:L\to D$ 可提升为同态 $F:M\to D$，当且仅当 $f\in\Hom_R(L,D)$ 位于 $\psi'$ 的像中。一般地，$\psi':\Hom_R(M,D)\to\Hom_R(L,D)$ 不一定满射；$\psi'$ 满射当且仅当每个从 $L$ 到 $D$ 的同态都可以提升为从 $M$ 到 $D$ 的同态，在这种情形下序列 (10.12) 可扩张为短正合列。

序列 (10.12) 对所有 $R$-模 $D$ 都正合，当且仅当序列
\[
L\xrightarrow{\psi}M\xrightarrow{\varphi}N\longrightarrow0
\]
正合。

\begin{proof}
第一部分中剩下的只是证明 (10.12) 在 $\Hom_R(M,D)$ 处正合。这个证明与定理 28 中对应结果的证明非常相似，留作练习。注意这里并不需要 $\psi$ 的单射性，这证明了最后一句中的“若”方向。

现在假设序列 (10.12) 对所有 $R$-模 $D$ 都正合。首先证明 $\varphi:M\to N$ 是满射。取
\[
D=N/\varphi(M).
\]
若 $\pi_1:N\to N/\varphi(M)$ 是自然投影同态，则由 $\pi_1$ 的定义，$\pi_1\circ\varphi(M)=0$。也就是说，$\varphi'(\pi_1)=0$。因为对所有模 $D$，$\varphi'$ 都单射，所以 $\pi_1=0$，于是 $N=\varphi(M)$，故 $\varphi$ 满射。

接着证明 $\varphi\circ\psi=0$，这将推出 $\operatorname{Im}\psi\subseteq\ker\varphi$。取 $D=N$，并注意 $N$ 上恒等映射 $\operatorname{id}_N$ 属于 $\Hom_R(N,N)$，所以 $\varphi'(\operatorname{id}_N)\in\Hom_R(M,N)$。由 (10.12) 在 $D=N$ 时的正合性，$\varphi'(\operatorname{id}_N)\in\ker\psi'$，因此
\[
\psi'(\varphi'(\operatorname{id}_N))=0.
\]
这正是 $\operatorname{id}_N\circ\varphi\circ\psi=0$，即 $\varphi\circ\psi=0$。

最后证明 $\ker\varphi\subseteq\operatorname{Im}\psi$。取 $D=M/\psi(L)$，令 $\pi_2:M\to M/\psi(L)$ 为自然投影。由于按定义 $\pi_2(\psi(L))=0$，故 $\psi'(\pi_2)=0$。正合性推出 $\pi_2$ 在 $\varphi'$ 的像中，即存在 $f\in\Hom_R(N,M/\psi(L))$ 使
\[
\pi_2=\varphi'(f)=f\circ\varphi.
\]
若 $m\in\ker\varphi$，则
\[
\pi_2(m)=f(\varphi(m))=0,
\]
这意味着 $m\in\psi(L)$，因为 $\pi_2$ 是到商 $M/\psi(L)$ 的投影。因此 $\ker\varphi\subseteq\operatorname{Im}\psi$，证明完成。
\end{proof}

由定理 33，序列
\[
0\longrightarrow\Hom_R(N,D)\xrightarrow{\varphi'}\Hom_R(M,D)\xrightarrow{\psi'}\Hom_R(L,D)\longrightarrow0
\]
一般不是短正合列，因为 $\psi'$ 不必满射。这个序列是否正合，精确地衡量了从 $M$ 到 $D$ 的同态在多大程度上由从 $L$ 与 $N$ 到 $D$ 的同态对唯一确定。

命题 29 的第二部分说明，当原短正合列
\[
0\longrightarrow L\longrightarrow M\longrightarrow N\longrightarrow0
\]
为分裂正合列时，上述序列是正合的。事实上，在这种情形中，对每个 $R$-模 $D$，
\[
0\longrightarrow\Hom_R(N,D)\longrightarrow\Hom_R(M,D)\longrightarrow\Hom_R(L,D)\longrightarrow0
\]
也是阿贝尔群的分裂短正合列。练习 14 表明逆命题成立：若对每个 $R$-模 $D$，这个 $\Hom$ 序列都正合，则原序列是分裂短正合列。

也有命题 30 前三部分的对偶版本，它描述那些 $R$-模 $D$：对它们而言，定理 33 中的序列 (10.12) 总能扩张为短正合列。

\noindent\textbf{命题 34.} 设 $Q$ 是 $R$-模。下列条件等价：
\begin{enumerate}[label=(\arabic*)]
\item 对任意 $R$-模 $L,M,N$，若
\[
0\longrightarrow L\longrightarrow M\longrightarrow N\longrightarrow0
\]
是短正合列，则
\[
0\longrightarrow\Hom_R(N,Q)\longrightarrow\Hom_R(M,Q)\longrightarrow\Hom_R(L,Q)\longrightarrow0
\]
也是短正合列。
\item 对任意 $R$-模 $L,M$，若 $0\to L\xrightarrow{\psi}M$ 正合，则每个从 $L$ 到 $Q$ 的 $R$-模同态都可提升为从 $M$ 到 $Q$ 的 $R$-模同态；即给定 $f\in\Hom_R(L,Q)$，存在提升 $F\in\Hom_R(M,Q)$ 使下图交换：
\[
\begin{array}{ccc}
0\longrightarrow&L&\xrightarrow{\psi}M\\
&f\downarrow&\swarrow F\\[-1mm]
&Q&
\end{array}
\]
\item 若 $Q$ 是 $R$-模 $M$ 的子模，则 $Q$ 是 $M$ 的直和因子；也就是说，每个短正合列
\[
0\longrightarrow Q\longrightarrow M\longrightarrow N\longrightarrow0
\]
都分裂。
\end{enumerate}

\begin{proof}
(1) 与 (2) 的等价是定理 33 的一部分。现在假设 (2) 成立，且
\[
0\longrightarrow Q\xrightarrow{\psi}M\longrightarrow N\longrightarrow0
\]
正合。取 $L=Q$ 且 $f$ 为 $Q$ 到自身的恒等映射，由 (2) 得存在同态 $F:M\to Q$，满足 $F\circ\psi=1$。于是 $F$ 是该序列的分裂同态，从而证明 (3)。(3) 蕴含 (2) 的证明在练习中给出。
\end{proof}

\noindent\textbf{定义.} 若 $R$-模 $Q$ 满足命题 34 的任一等价条件，则称 $Q$ 为\emph{内射模}。

命题 34 的第三个陈述可重述为：任何把 $Q$ 嵌入其中的模 $M$，都含有 $Q$ 的一个同构副本作为直和因子。这解释了“内射”这一术语。

若固定 $D$，则对任意 $R$-模 $X$，有阿贝尔群 $\Hom_R(X,D)$。进一步，$R$-模同态 $\alpha:X\to Y$ 诱导阿贝尔群同态
\[
\alpha':\Hom_R(Y,D)\longrightarrow\Hom_R(X,D),
\qquad \alpha'(f)=f\circ\alpha,
\]
它“反转”箭头方向。换言之，映射 $\Hom_R(-,D)$ 是从 $R$-模范畴到阿贝尔群范畴的反变函子（参见附录 II）。定理 33 表明，把这个函子作用于正合列
\[
0\longrightarrow L\xrightarrow{\psi}M\xrightarrow{\varphi}N\longrightarrow0
\]
的各项，得到正合列
\[
0\longrightarrow\Hom_R(N,D)\xrightarrow{\varphi'}\Hom_R(M,D)\xrightarrow{\psi'}\Hom_R(L,D).
\]
这称为 $\Hom_R(-,D)$ 是左正合（反变）函子。注意，$\Hom_R(-,D)$ 与前面考虑的 $\Hom_R(D,-)$ 都是左正合的；前者反转原短正合列中映射的方向，后者保持方向。

由命题 34，函子 $\Hom_R(-,D)$ 正合，即总把短正合列变为短正合列（因而把任意长度正合列变为正合列），当且仅当 $D$ 是内射的。总结为下列与推论 32 对偶的结果。

\noindent\textbf{推论 35.} 若 $D$ 是 $R$-模，则函子 $\Hom_R(-,D)$ 从 $R$-模范畴到阿贝尔群范畴是左正合的。它正合当且仅当 $D$ 是内射 $R$-模。

我们已经看到，$R$-模投射当且仅当它是自由 $R$-模的直和因子。要给出内射 $R$-模如此简单的刻画并不容易。下一个结果给出 $Q$ 为内射 $R$-模的一个判别法（这是 Baer 的结果，Baer 大约在 1940 年引入内射模的概念），并由此在 $R=\ZZ$ 时（更一般地，在 $R$ 为 P.I.D. 时）给出内射模的刻画。回忆，$\ZZ$-模 $A$（即加法记号的阿贝尔群）称为\emph{可除的}，若对每个非零整数 $n$ 都有 $A=nA$。例如，$\QQ$ 与 $\QQ/\ZZ$ 都是可除的（参见第 2.4 节练习 18、19 与第 3.1 节练习 15）。

\noindent\textbf{命题 36.} 设 $Q$ 是 $R$-模。
\begin{enumerate}[label=(\arabic*)]
\item （Baer 判别法）模 $Q$ 内射，当且仅当对 $R$ 的每个左理想 $I$，任意 $R$-模同态 $g:I\to Q$ 都可以扩张为 $R$-模同态 $G:R\to Q$。
\item 若 $R$ 是 P.I.D.，则 $Q$ 内射当且仅当对每个非零 $r\in R$ 都有 $rQ=Q$。特别地，$\ZZ$-模内射当且仅当它是可除的。若 $R$ 是 P.I.D.，则内射 $R$-模的商模仍是内射的。
\end{enumerate}

\begin{proof}
若 $Q$ 内射，且 $g:I\to Q$ 是从 $R$ 的非零左理想 $I$ 到 $Q$ 的 $R$-模同态，则由命题 34(2) 作用于正合列 $0\to I\to R$，$g$ 可扩张为从 $R$ 到 $Q$ 的 $R$-模同态。这证明了 (1) 的“只若”方向。

反过来，假设每个同态 $g:I\to Q$ 都可提升为同态 $G:R\to Q$。要证明 $Q$ 内射，必须证明：若 $0\to L\to M$ 正合且 $f:L\to Q$ 是 $R$-模同态，则存在提升 $F:M\to Q$ 扩张 $f$。令 $\mathcal S$ 为所有对 $(f',L')$ 组成的集合，其中 $L'$ 是 $M$ 中含 $L$ 的子模，$f':L'\to Q$ 是 $f$ 的提升。规定
\[
(f',L')\leq(f'',L'')
\]
若 $L'\subseteq L''$ 且 $f''$ 在 $L'$ 上等于 $f'$。这把 $\mathcal S$ 部分有序。由于 $\mathcal S\neq\emptyset$，由 Zorn 引理，$\mathcal S$ 有极大元 $(F,M')$。映射 $F:M'\to Q$ 是 $f$ 的一个提升。只需证明 $M'=M$。

若存在 $m\in M$ 不在 $M'$ 中，令
\[
I=\{r\in R\mid rm\in M'\}.
\]
容易验证 $I$ 是 $R$ 的左理想，并且映射
\[
g:I\to Q,
\qquad g(x)=F(xm)
\]
是 $R$-模同态。由假设，存在 $g$ 的提升 $G:R\to Q$。考虑 $M$ 的子模 $M'+Rm$，定义
\[
F':M'+Rm\longrightarrow Q,
\qquad F'(m'+rm)=F(m')+G(r).
\]
若 $m_1+r_1m=m_2+r_2m$，则 $(r_1-r_2)m=m_2-m_1$，所以 $r_1-r_2\in I$。于是
\[
G(r_1-r_2)=g(r_1-r_2)=F((r_1-r_2)m)=F(m_2-m_1),
\]
从而 $F(m_1)+G(r_1)=F(m_2)+G(r_2)$。故 $F'$ 良定义，并且立即可见它是扩张 $f$ 到 $M'+Rm$ 的 $R$-模同态。这与 $M'$ 的极大性矛盾，所以 $M'=M$。于是 (1) 得证。

证明 (2)。设 $R$ 是 P.I.D.。$R$ 的任意非零理想 $I$ 都形如 $I=(r)$，其中 $r\in R$ 非零。从 $I$ 到 $Q$ 的 $R$-模同态 $f$ 完全由 $f(r)=q\in Q$ 决定。这个同态可扩张为 $F:R\to Q$，当且仅当存在 $q'\in Q$ 使 $F(1)=q'$ 且
\[
q=f(r)=F(r)=rq'.
\]
因此 Baer 判别法对 $Q$ 成立，当且仅当 $rQ=Q$。这证明了 (2) 的前两个断言。最后一个断言是因为，若 $Q$ 对所有非零 $r\in R$ 满足 $rQ=Q$，则 $Q$ 的商模也满足相同性质。
\end{proof}

\paragraph{例}
\begin{enumerate}[label=(\arabic*)]
\item 因为 $\ZZ$ 不可除，$\ZZ$ 不是内射 $\ZZ$-模。这也可由短正合列
\[
0\longrightarrow\ZZ\xrightarrow{\,2\,}\ZZ\longrightarrow\ZZ/2\ZZ\longrightarrow0
\]
不分裂看出。
\item 有理数群 $\QQ$ 是内射 $\ZZ$-模。
\item 内射 $\ZZ$-模 $\QQ$ 的商 $\QQ/\ZZ$ 是内射 $\ZZ$-模。
\item 直接和的可除 $\ZZ$-模仍可除，因此内射 $\ZZ$-模的直接和仍内射。例如 $\QQ\oplus\QQ/\ZZ$ 是内射 $\ZZ$-模。（亦见练习 4。）
\item 我们将在第 12 章看到，非零有限生成 $\ZZ$-模都不是内射的。
\item 设环 $R$ 是整环。若 $R$-模 $A$ 满足对每个非零 $r\in R$ 都有 $rA=A$，则称 $A$ 是\emph{可除 $R$-模}。命题 36 的证明说明，在这种情形中，内射 $R$-模必可除。
\item 我们将在第 11.1 节看到，若 $R=F$ 是域，则每个 $F$-模都是内射的。
\item 我们将在第 VI 部分看到，若 $F$ 是任意域且 $n\in\ZZ^+$，则 $R=M_n(F)$ 具有这样的性质：每个 $R$-模都是内射的（也都是投射的）。还将看到，若 $G$ 是阶为 $n$ 的有限群且 $n\neq0$ 于域 $F$ 中，则群环 $FG$ 也具有每个模都是内射的（且投射的）性质。
\end{enumerate}

\noindent\textbf{推论 37.} 每个 $\ZZ$-模都是某个内射 $\ZZ$-模的子模。

\begin{proof}
设 $M$ 是 $\ZZ$-模，令 $A$ 是 $M$ 的任意一组 $\ZZ$-模生成元。令 $\mathcal F=F(A)$ 为集合 $A$ 上的自由 $\ZZ$-模。由定理 6，存在从 $\mathcal F$ 到 $M$ 的满 $\ZZ$-模同态；若 $K$ 是其核，则 $K$ 是 $\mathcal F$ 的 $\ZZ$-子模，并且我们可把 $M$ 识别为 $\mathcal F/K$。令 $\mathcal Q$ 为集合 $A$ 上的自由 $\QQ$-模。则 $\mathcal Q$ 是若干个 $\QQ$ 的直接和，因此可除，从而由命题 36 是内射 $\ZZ$-模，并且包含 $\mathcal F$。于是 $K$ 也是 $\mathcal Q$ 的 $\ZZ$-子模，所以商 $\mathcal Q/K$ 仍由命题 36 是内射的。由于
\[
M=\mathcal F/K\subseteq \mathcal Q/K,
\]
故 $M$ 包含于一个内射 $\ZZ$-模中。
\end{proof}

推论 37 可用来证明下列适用于任意 $R$-模的更一般版本。该定理是定理 6 与推论 31 中“每个 $R$-模都是某个投射 $R$-模的商”这一结果的内射类似物。

\noindent\textbf{定理 38.} 设 $R$ 是带单位元的环，$M$ 是 $R$-模。则 $M$ 包含于某个内射 $R$-模中。

\begin{proof}
证明见练习 15 至 17 的提示。
\end{proof}

可以证明比定理 38 更强的结果：存在含有 $M$ 的一个极小内射 $R$-模 $H$，其意义是任意将 $M$ 单射到内射 $R$-模 $Q$ 的映射都经由 $H$ 因子化。更精确地说，若 $M\hookrightarrow Q$ 且 $Q$ 是内射 $R$-模，则存在单射 $\iota:H\hookrightarrow Q$，其在 $M$ 上限制为恒等映射；利用 $\iota$ 把 $H$ 识别为 $Q$ 的子模，则
\[
M\subseteq H\subseteq Q.
\]
这个模 $H$ 称为 $M$ 的\emph{内射包}或\emph{内射包络}。$M$ 的内射包关于把 $M$ 包含进内射 $R$-模中的泛性质，应与定理 6 中 $M$ 的一组生成元 $A$ 上自由模 $F(A)$ 关于从 $M$ 出发的同态的泛性质比较。例如，$\ZZ$ 的内射包是 $\QQ$，任意域的内射包是它自身（参见练习）。

\subsection*{平坦模与 $D\otimes_R-$}

现在考虑 $R$-模扩张
\[
0\longrightarrow L\xrightarrow{\psi}M\xrightarrow{\varphi}N\longrightarrow0
\]
在张量积下的行为。

设 $D$ 是右 $R$-模。对左 $R$-模的任意同态 $f:X\to Y$，由定理 13 得到阿贝尔群同态
\[
1\otimes f:D\otimes_R X\longrightarrow D\otimes_R Y.
\]
若此外 $D$ 是 $(S,R)$-双模（例如 $S=R$ 交换且 $D$ 赋予第 4 节中的标准 $(R,R)$-双模结构），则 $1\otimes f$ 是左 $S$-模同态。换言之，
\[
D\otimes_R-:X\longmapsto D\otimes_RX
\]
是从左 $R$-模范畴到阿贝尔群范畴的协变函子；当 $D$ 是 $(S,R)$-双模时，它是到左 $S$-模范畴的协变函子（参见附录 II）。类似地，若 $D$ 是左 $R$-模，则 $-\otimes_RD$ 是从右 $R$-模范畴到阿贝尔群范畴的协变函子；若 $D$ 是 $(R,S)$-双模，则是到右 $S$-模范畴的协变函子。注意，与 $\Hom$ 不同，张量积在两个变量中都是协变的。因此我们集中讨论 $D\otimes_R-$，而把 $-\otimes_RD$ 所需的小修改留作练习。

我们已经看到一些例子，其中由单射 $\psi:L\hookrightarrow M$ 诱导的映射
\[
1\otimes\psi:D\otimes_RL\longrightarrow D\otimes_RM
\]
不再是单射。例如，$\ZZ$-模单射 $\ZZ\hookrightarrow\QQ$ 诱导从
\[
\ZZ/2\ZZ\otimes_\ZZ\ZZ\cong\ZZ/2\ZZ
\]
到
\[
\ZZ/2\ZZ\otimes_\ZZ\QQ=0
\]
的零映射。另一方面，设 $\varphi:M\to N$ 是满 $R$-模同态。张量积 $D\otimes_RN$ 作为阿贝尔群由简单张量 $d\otimes n$ 生成，其中 $d\in D$、$n\in N$。$\varphi$ 的满性给出 $n=\varphi(m)$ 对某个 $m\in M$ 成立，于是
\[
(1\otimes\varphi)(d\otimes m)=d\otimes\varphi(m)=d\otimes n.
\]
这说明 $1\otimes\varphi:D\otimes_RM\to D\otimes_RN$ 是满同态。由此得到下列定理的大部分。

\noindent\textbf{定理 39.} 设 $D$ 是右 $R$-模，$L,M,N$ 是左 $R$-模。若
\[
0\longrightarrow L\xrightarrow{\psi}M\xrightarrow{\varphi}N\longrightarrow0
\]
正合，则相应阿贝尔群序列
\[
D\otimes_RL\xrightarrow{1\otimes\psi}D\otimes_RM\xrightarrow{1\otimes\varphi}D\otimes_RN\longrightarrow0
\tag{10.13}
\]
正合。

若 $D$ 是 $(S,R)$-双模，则 (10.13) 是左 $S$-模的正合列。特别地，若 $S=R$ 是交换环，则 (10.13) 关于标准 $R$-模结构是 $R$-模正合列。映射 $1\otimes\psi$ 一般不是单射，也就是说 (10.13) 一般不能扩张为短正合列。

序列 (10.13) 对所有右 $R$-模 $D$ 都正合，当且仅当
\[
L\xrightarrow{\psi}M\xrightarrow{\varphi}N\longrightarrow0
\]
正合。

\begin{proof}
第一部分还需证明 (10.13) 在 $D\otimes_RM$ 处正合。由于 $\varphi\circ\psi=0$，有
\[
(1\otimes\varphi)\circ(1\otimes\psi)=0,
\]
从而
\[
\operatorname{Im}(1\otimes\psi)\subseteq\ker(1\otimes\varphi).
\]
因此有自然投影
\[
\pi:(D\otimes_RM)/\operatorname{Im}(1\otimes\psi)
\longrightarrow
(D\otimes_RM)/\ker(1\otimes\varphi)=D\otimes_RN.
\]
由两个投影复合得到的映射
\[
D\otimes_RM\longrightarrow(D\otimes_RM)/\operatorname{Im}(1\otimes\psi)
\xrightarrow{\pi}D\otimes_RN
\]
就是 $1\otimes\varphi$。我们将证明 $\pi$ 是同构，这就说明 $1\otimes\varphi$ 的核正是第一投影的核，即 $\operatorname{Im}(1\otimes\psi)$，从而得到 (10.13) 在 $D\otimes_RM$ 处正合。

定义映射
\[
\pi':D\times N\longrightarrow (D\otimes_RM)/\operatorname{Im}(1\otimes\psi)
\]
如下：对 $d\in D$ 与 $n\in N$，取任意 $m\in M$ 使 $\varphi(m)=n$，令
\[
\pi'(d,n)=d\otimes m.
\]
这是良定义的：若另一个元素 $m'$ 也映到 $n$，则 $m'-m\in\ker\varphi=\operatorname{Im}\psi$，所以 $m'=m+\psi(l)$ 对某个 $l\in L$ 成立；于是
\[
d\otimes\psi(l)=(1\otimes\psi)(d\otimes l)
\]
在商中为零。容易验证 $\pi'$ 是平衡映射，所以诱导同态
\[
\widetilde\pi:D\otimes_RN\longrightarrow (D\otimes_RM)/\operatorname{Im}(1\otimes\psi)
\]
且 $\widetilde\pi(d\otimes n)=d\otimes m$。于是
\[
\widetilde\pi\circ\pi(d\otimes m)=\widetilde\pi(d\otimes\varphi(m))=d\otimes m
\]
而同样 $\pi\circ\widetilde\pi=1$。故 $\pi$ 与 $\widetilde\pi$ 互为逆同构，证明 (10.13) 正合。注意，在此证明中并未用到 $\psi$ 的单射性。

最后，若 (10.13) 对每个右 $R$-模 $D$ 都正合，则取 $D=R$。一般地，对任意左 $R$-模 $X$ 有
\[
R\otimes_RX\cong X
\]
（见推论 9 后例 1），于是得到 $L\to M\to N\to0$ 的正合性。
\end{proof}

由定理 39，序列
\[
0\longrightarrow D\otimes_RL\xrightarrow{1\otimes\psi}D\otimes_RM\xrightarrow{1\otimes\varphi}D\otimes_RN\longrightarrow0
\]
一般不正合，因为 $1\otimes\psi$ 不必单射。然而若原短正合列
\[
0\longrightarrow L\xrightarrow{\psi}M\xrightarrow{\varphi}N\longrightarrow0
\]
分裂，则由于张量积与直和交换（定理 17），可知
\[
0\longrightarrow D\otimes_RL\longrightarrow D\otimes_RM\longrightarrow D\otimes_RN\longrightarrow0
\]
也是分裂短正合列。

由定理 39 立即得到下列刻画，它描述那些模 $D$：对其而言，(10.13) 总可扩张为短正合列。

\noindent\textbf{命题 40.} 设 $A$ 是右 $R$-模。下列条件等价：
\begin{enumerate}[label=(\arabic*)]
\item 对任意左 $R$-模 $L,M,N$，若
\[
0\longrightarrow L\xrightarrow{\psi}M\xrightarrow{\varphi}N\longrightarrow0
\]
是短正合列，则
\[
0\longrightarrow A\otimes_RL\xrightarrow{1\otimes\psi}A\otimes_RM\xrightarrow{1\otimes\varphi}A\otimes_RN\longrightarrow0
\]
也是短正合列。
\item 对任意左 $R$-模 $L,M$，若
\[
0\longrightarrow L\xrightarrow{\psi}M
\]
正合（即 $\psi$ 单射），则
\[
0\longrightarrow A\otimes_RL\xrightarrow{1\otimes\psi}A\otimes_RM
\]
是阿贝尔群的正合列（即 $1\otimes\psi$ 单射）。
\end{enumerate}

\noindent\textbf{定义.} 若右 $R$-模 $A$ 满足命题 40 的任一等价条件，则称 $A$ 为\emph{平坦模}。

对固定右 $R$-模 $D$，定理 39 的第一部分称为函子 $D\otimes_R-$ 的\emph{右正合性}。

\noindent\textbf{推论 41.} 若 $D$ 是右 $R$-模，则函子 $D\otimes_R-$ 从左 $R$-模范畴到阿贝尔群范畴是右正合的。若 $D$ 是 $(S,R)$-双模（例如 $S=R$ 交换且 $D$ 赋予标准 $R$-模结构），则 $D\otimes_R-$ 是从左 $R$-模范畴到左 $S$-模范畴的右正合函子。该函子正合当且仅当 $D$ 是平坦 $R$-模。

我们已经见到一些平坦模。

\noindent\textbf{推论 42.} 自由模是平坦的；更一般地，投射模是平坦的。

\begin{proof}
要证明自由 $R$-模 $F$ 平坦，只需证明：对任意 $R$-模单射 $\psi:L\to M$，诱导映射
\[
1\otimes\psi:F\otimes_RL\longrightarrow F\otimes_RM
\]
也是单射。先假设 $F\cong R^n$ 是有限生成自由模。此时
\[
F\otimes_RL\cong R^n\otimes_RL\cong L^n,
\]
因为 $R\otimes_RL\cong L$ 且张量积与直和交换。类似地 $F\otimes_RM\cong M^n$，在这些同构下，$1\otimes\psi$ 就是由每个分量中的包含 $\psi$ 所诱导的自然映射 $L^n\to M^n$。特别地，$1\otimes\psi$ 单射，所以有限生成自由模平坦。

现在设 $F$ 是任意自由模，且元素 $\sum f_i\otimes l_i\in F\otimes_RL$ 在 $1\otimes\psi$ 下映到 $0$。这意味着元素 $\sum f_i\otimes\psi(l_i)$ 在 $F\otimes_RM$ 中为 $0$，因而它可写成第 4 节公式 (6) 中那些定义张量积的生成关系的有限和。由于这些项有限，所得等式中所有第一坐标都落在 $F$ 的某个有限生成自由子模 $F'$ 中。于是这个等式说明 $\sum f_i\otimes l_i\in F'\otimes_RL$ 在 $F'\otimes_RM$ 中映到 $0$。由已证的有限生成自由情形，$F'$ 平坦，故 $\sum f_i\otimes l_i=0$ 于 $F'\otimes_RL$，从而也为 $0$ 于 $F\otimes_RL$。所以 $1\otimes\psi$ 单射，$F$ 平坦。

若 $P$ 是投射模，则由命题 30，$P$ 是某个自由模 $F$ 的直和因子，设 $F=P\oplus P'$。若 $\psi:L\to M$ 单射，则由 $F$ 平坦，
\[
1\otimes\psi:F\otimes_RL\longrightarrow F\otimes_RM
\]
单射。由于 $F=P\oplus P'$ 且张量积与直和交换，这说明
\[
1\otimes\psi:(P\otimes_RL)\oplus(P'\otimes_RL)
\longrightarrow (P\otimes_RM)\oplus(P'\otimes_RM)
\]
单射，因此限制到第一直和因子也给出 $P\otimes_RL\to P\otimes_RM$ 的单射。故 $P$ 平坦。
\end{proof}

\paragraph{例}
\begin{enumerate}[label=(\arabic*)]
\item 因为 $\ZZ$ 是投射 $\ZZ$-模，它是平坦的。定理 39 前面的例子表明 $\ZZ/2\ZZ$ 不是平坦 $\ZZ$-模。

\item $\ZZ$-模 $\QQ$ 是平坦 $\ZZ$-模，如下所示。设 $\psi:L\to M$ 是 $\ZZ$-模单射。$\QQ\otimes_\ZZ L$ 的每个元素都可写成 $(1/d)\otimes l$ 的形式，其中 $d$ 是非零整数，$l\in L$（第 4 节练习 7）。若 $(1/d)\otimes l$ 在 $1\otimes\psi$ 的核中，则 $(1/d)\otimes\psi(l)$ 在 $\QQ\otimes_\ZZ M$ 中为 $0$。由第 4 节练习 8，这意味着存在非零整数 $c$ 使 $c\psi(l)=0$ 于 $M$。于是 $\psi(cl)=0$，由 $\psi$ 的单射性，$cl=0$ 于 $L$。因此
\[
(1/d)\otimes l=(1/(cd))\otimes(cl)=0
\]
于 $\QQ\otimes_\ZZ L$。这证明 $1\otimes\psi$ 单射。

\item $\ZZ$-模 $\QQ/\ZZ$ 是内射的（由命题 36），但不是平坦的：由乘以 $2$ 给出的单射 $\psi(z)=2z$ 从 $\ZZ$ 到 $\ZZ$，与 $\QQ/\ZZ$ 张量后不再保持单射。事实上，在
\[
(\QQ/\ZZ)\otimes_\ZZ \ZZ\longrightarrow(\QQ/\ZZ)\otimes_\ZZ\ZZ
\]
中，非零元素 $(\tfrac12+\ZZ)\otimes1$ 位于核中。用 $\QQ/\ZZ\cong(\QQ/\ZZ)\otimes_\ZZ\ZZ$ 识别，这就是说乘以 $2$ 的映射以 $1/2$ 为核中的元素。

\item 平坦模的直和仍平坦（练习 5）。特别地，$\QQ\oplus\ZZ$ 平坦。这个模既非投射也非内射：$\QQ$ 非投射（练习 8），$\ZZ$ 非内射（命题 36；另参见练习 3、4）。
\end{enumerate}

我们以 $\Hom$ 与张量积之间的一个重要关系结束本节。

\noindent\textbf{定理 43（伴随结合律）.} 设 $R,S$ 是环，$A$ 是右 $R$-模，$B$ 是 $(R,S)$-双模，$C$ 是右 $S$-模。则有阿贝尔群同构
\[
\Hom_S(A\otimes_RB,C)
\cong
\Hom_R(A,\Hom_S(B,C)).
\]
这里的同态群均指右模同态。注意 $\Hom_S(B,C)$ 具有右 $R$-模结构（参见练习）。若 $R=S$ 交换，则这是在标准 $R$-模结构下的 $R$-模同构。

\begin{proof}
设 $\varphi:A\otimes_RB\to C$ 是同态。对固定 $a\in A$，定义从 $B$ 到 $C$ 的映射 $\Phi(a)$ 为
\[
\Phi(a)(b)=\varphi(a\otimes b).
\]
容易验证 $\Phi(a)$ 是右 $S$-模同态，并且映射
\[
\Phi:A\longrightarrow\Hom_S(B,C),
\qquad a\longmapsto \Phi(a)
\]
是右 $R$-模同态。令 $f(\varphi)=\Phi$，得到从 $\Hom_S(A\otimes_RB,C)$ 到 $\Hom_R(A,\Hom_S(B,C))$ 的群同态。

反过来，设 $\Phi:A\to\Hom_S(B,C)$ 是同态。定义 $A\times B\to C$ 的映射为
\[
(a,b)\longmapsto \Phi(a)(b).
\]
这是 $R$-平衡映射，因此诱导同态
\[
\varphi:A\otimes_RB\to C.
\]
令 $g(\Phi)=\varphi$，得到反向群同态。两个构造互为逆映射，从而得到定理中的同构。
\end{proof}

作为定理 43 的第一个应用，我们在 $S=R$ 为交换环的情形给出定理 39 第一部分——张量积右正合性——的另一种证明。若
\[
0\longrightarrow L\longrightarrow M\longrightarrow N\longrightarrow0
\]
正合，则由定理 33，对每个 $R$-模 $E$，序列
\[
0\longrightarrow\Hom_R(N,E)\longrightarrow\Hom_R(M,E)\longrightarrow\Hom_R(L,E)
\]
正合。再由定理 28，对所有 $D$ 与 $E$，
\[
0\longrightarrow\Hom_R(D,\Hom_R(N,E))\longrightarrow\Hom_R(D,\Hom_R(M,E))\longrightarrow\Hom_R(D,\Hom_R(L,E))
\]
正合。由伴随结合律，这意味着对任意 $D$ 与所有 $E$，序列
\[
0\longrightarrow\Hom_R(D\otimes_RN,E)\longrightarrow\Hom_R(D\otimes_RM,E)\longrightarrow\Hom_R(D\otimes_RL,E)
\]
正合。于是由定理 33 的第二部分，得到
\[
D\otimes_RL\longrightarrow D\otimes_RM\longrightarrow D\otimes_RN\longrightarrow0
\]
对所有 $D$ 正合，这正是张量积的右正合性。

作为定理 43 的第二个应用，我们证明交换环 $R$ 上两个投射模的张量积仍投射（另一个更直接的证明见练习 9）。

\noindent\textbf{推论 44.} 若 $R$ 交换，则两个投射 $R$-模的张量积仍是投射 $R$-模。

\begin{proof}
设 $P_1$ 与 $P_2$ 是投射模。由推论 32，$\Hom_R(P_2,-)$ 是从 $R$-模范畴到 $R$-模范畴的正合函子。再由同一推论，复合函子
\[
\Hom_R(P_1,\Hom_R(P_2,-))
\]
是正合的。由定理 43，这意味着 $\Hom_R(P_1\otimes_RP_2,-)$ 是 $R$-模上的正合函子。再次由推论 32，$P_1\otimes_RP_2$ 是投射的。
\end{proof}

\paragraph{本节总结}
函子 $\Hom_R(A,-)$、$\Hom_R(-,A)$ 与 $A\otimes_R-$ 都把左 $R$-模映到阿贝尔群；函子 $-\otimes_RA$ 把右 $R$-模映到阿贝尔群。若 $R$ 交换，则四个函子都把 $R$-模映到 $R$-模。
\begin{enumerate}[label=(\arabic*)]
\item 设 $A$ 是左 $R$-模。函子 $\Hom_R(A,-)$ 是协变左正合函子；模 $A$ 投射当且仅当 $\Hom_R(A,-)$ 正合（即也右正合）。
\item 设 $A$ 是左 $R$-模。函子 $\Hom_R(-,A)$ 是反变左正合函子；模 $A$ 内射当且仅当 $\Hom_R(-,A)$ 正合。
\item 设 $A$ 是右 $R$-模。函子 $A\otimes_R-$ 是协变右正合函子；模 $A$ 平坦当且仅当 $A\otimes_R-$ 正合（即也左正合）。
\item 设 $A$ 是左 $R$-模。函子 $-\otimes_RA$ 是协变右正合函子；模 $A$ 平坦当且仅当 $-\otimes_RA$ 正合。
\item 投射模都是平坦模。$\ZZ$-模 $\QQ/\ZZ$ 是内射的但不平坦；$\ZZ$-模 $\ZZ\oplus\QQ$ 是平坦的，但既不投射也不内射。
\end{enumerate}



\paragraph{练习}
设 $R$ 是带单位元的环。
\begin{enumerate}[label=\arabic*.]
	\item 设
	\[
	\begin{array}{ccccc}
		A & \xrightarrow{\psi} & B & \xrightarrow{\varphi} & C\\
		\alpha\downarrow & & \beta\downarrow & & \gamma\downarrow\\
		A'&\xrightarrow{\psi'}&B'&\xrightarrow{\varphi'}&C'
	\end{array}
	\]
	是群的交换图，且两行正合。证明：
	\begin{enumerate}[label=(\alph*)]
		\item 若 $\varphi$ 与 $\alpha$ 满射，且 $\beta$ 单射，则 $\gamma$ 单射。[若 $c\in\ker\gamma$，取 $b\in B$ 使 $\varphi(b)=c$。证明 $\varphi'(\beta(b))=0$，从而 $\beta(b)=\psi'(a')$ 对某个 $a'\in A'$ 成立。再取 $a\in A$ 使 $\alpha(a)=a'$，并证明 $\beta(\psi(a))=\beta(b)$。推出 $b=\psi(a)$，故 $c=\varphi(b)=0$。]
		\item 若 $\psi'$、$\alpha$ 与 $\gamma$ 都单射，则 $\beta$ 单射。
		\item 若 $\varphi$、$\alpha$ 与 $\gamma$ 都满射，则 $\beta$ 满射。
		\item 若 $\beta$ 单射，$\alpha$ 与 $\gamma$ 满射，则 $\gamma$ 单射。
		\item 若 $\beta$ 满射，$\gamma$ 与 $\psi'$ 单射，则 $\alpha$ 满射。
	\end{enumerate}
	\item 设
	\[
	\begin{array}{ccccccc}
		A&\longrightarrow&B&\longrightarrow&C&\longrightarrow&D\\
		\alpha\downarrow&&\beta\downarrow&&\gamma\downarrow&&\delta\downarrow\\
		A'&\longrightarrow&B'&\longrightarrow&C'&\longrightarrow&D'
	\end{array}
	\]
	是群的交换图，且两行正合。证明：
	\begin{enumerate}[label=(\alph*)]
		\item 若 $\alpha$ 满射，$\beta$ 与 $\delta$ 单射，则 $\gamma$ 单射。
		\item 若 $\delta$ 单射，$\alpha$ 与 $\gamma$ 满射，则 $\beta$ 满射。
	\end{enumerate}
	\item 设 $P_1,P_2$ 是 $R$-模。证明 $P_1\oplus P_2$ 是投射 $R$-模当且仅当 $P_1$ 与 $P_2$ 都是投射的。
	\item 设 $Q_1,Q_2$ 是 $R$-模。证明 $Q_1\oplus Q_2$ 是内射 $R$-模当且仅当 $Q_1$ 与 $Q_2$ 都是内射的。
	\item 设 $A_1,A_2$ 是 $R$-模。证明 $A_1\oplus A_2$ 是平坦 $R$-模当且仅当 $A_1$ 与 $A_2$ 都是平坦的。更一般地，证明任意直和 $\bigoplus A_i$ 是平坦的当且仅当每个 $A_i$ 都平坦。[使用张量积与任意直和交换的事实。]
	\item 证明环 $R$ 的下列条件等价：
	\begin{enumerate}[label=(\roman*)]
		\item 每个 $R$-模都是投射的；
		\item 每个 $R$-模都是内射的。
	\end{enumerate}
	\item 设 $A$ 是非零有限阿贝尔群。
	\begin{enumerate}[label=(\alph*)]
		\item 证明 $A$ 不是投射 $\ZZ$-模。
		\item 证明 $A$ 不是内射 $\ZZ$-模。
	\end{enumerate}
	\item 设 $Q$ 是非零可除 $\ZZ$-模。证明 $Q$ 不是投射 $\ZZ$-模。推出有理数群 $\QQ$ 不是投射 $\ZZ$-模。[先证明若 $F$ 是任意自由模，则 $\bigcap_{n\geq1}nF=0$（用 $F$ 的一组基证明）。再反设 $Q$ 投射，并由命题 30(4) 得到矛盾。]
	\item 假设 $R$ 交换且有 $1$。
	\begin{enumerate}[label=(\alph*)]
		\item 证明两个自由 $R$-模的张量积仍是自由的。[使用张量积与直和交换的事实。]
		\item 用 (a) 证明两个投射 $R$-模的张量积仍是投射的。
	\end{enumerate}
	\item 设 $R$ 与 $S$ 是带单位元的环，$M,N$ 是左 $R$-模。再假设 $M$ 是 $(R,S)$-双模。
	\begin{enumerate}[label=(\alph*)]
		\item 对 $s\in S$ 与 $\varphi\in\Hom_R(M,N)$，定义 $(s\varphi):M\to N$ 为
		\[
		(s\varphi)(m)=\varphi(ms).
		\]
		证明 $s\varphi$ 是左 $R$-模同态，并且 $S$ 在 $\Hom_R(M,N)$ 上的这个作用使其成为左 $S$-模。
		\item 令 $S=R$ 且 $M=R$，把 $R$ 作为通常的 $(R,R)$-双模。对每个 $n\in N$，定义 $\varphi_n:R\to N$ 为 $\varphi_n(r)=rn$，即 $\varphi_n$ 是把 $1_R$ 映到 $n$ 的唯一 $R$-模同态。证明 $\varphi_n\in\Hom_R(R,N)$。用 (a) 证明映射 $n\mapsto\varphi_n$ 是左 $R$-模同构：
		\[
		N\cong\Hom_R(R,N).
		\]
		\item 推出若 $N$ 是自由（相应地，投射、内射、平坦）左 $R$-模，则 $\Hom_R(R,N)$ 也是自由（相应地，投射、内射、平坦）左 $R$-模。
	\end{enumerate}
	\item 设 $R$ 与 $S$ 是带单位元的环，$M,N$ 是左 $R$-模。再假设 $N$ 是 $(R,S)$-双模。
	\begin{enumerate}[label=(\alph*)]
		\item 对 $s\in S$ 与 $\varphi\in\Hom_R(M,N)$，定义 $(\varphi s):M\to N$ 为
		\[
		(\varphi s)(m)=\varphi(m)s.
		\]
		证明 $\varphi s$ 是左 $R$-模同态，并且 $S$ 在 $\Hom_R(M,N)$ 上的这个作用使其成为右 $S$-模。推出对任意 $R$-模 $M$，$\Hom_R(M,R)$ 是右 $R$-模，称为 $M$ 的\emph{对偶模}。
		\item 令 $N=R$，把它作为通常的 $(R,R)$-双模。证明在 (a) 定义的作用下，映射 $r\mapsto\varphi_r$ 是右 $R$-模同构：
		\[
		\Hom_R(R,R)\cong R,
		\]
		其中 $\varphi_r$ 是把 $1_R$ 映到 $r$ 的同态。推出若 $M$ 是有限生成自由左 $R$-模，则 $\Hom_R(M,R)$ 是同秩的自由右 $R$-模（也参见练习 13）。
		\item 证明若 $M$ 是有限生成投射 $R$-模，则其对偶模 $\Hom_R(M,R)$ 也是投射的。
	\end{enumerate}
	\item 设 $A$ 是 $R$-模，$I$ 是任意非空指标集，并对每个 $i\in I$，令 $B_i$ 是 $R$-模。证明下列阿贝尔群同构；当 $R$ 交换时，也证明它们是 $R$-模同构。（任意模族的直和与直积见第 3 节练习 20。）
	\begin{enumerate}[label=(\alph*)]
		\item $\Hom_R\left(\bigoplus_{i\in I}B_i,A\right)\cong\prod_{i\in I}\Hom_R(B_i,A)$；
		\item $\Hom_R\left(A,\prod_{i\in I}B_i\right)\cong\prod_{i\in I}\Hom_R(A,B_i)$。
	\end{enumerate}
	\item
	\begin{enumerate}[label=(\alph*)]
		\item 证明具有可数基的自由 $\ZZ$-模的对偶模不是自由的。[使用前一题与第 3 节练习 24。也见第 11.3 节练习 5。]
		\item 证明具有可数基的自由 $\ZZ$-模的对偶模也不是投射的。[可使用任意自由 $\ZZ$-模的子模都是自由的这一事实。]
	\end{enumerate}
	\item 设
	\[
	0\longrightarrow L\xrightarrow{\psi}M\xrightarrow{\varphi}N\longrightarrow0
	\]
	是 $R$-模序列。
	\begin{enumerate}[label=(\alph*)]
		\item 证明对所有 $R$-模 $D$，相伴序列
		\[
		0\longrightarrow\Hom_R(D,L)\xrightarrow{\psi'}\Hom_R(D,M)
		\xrightarrow{\varphi'}\Hom_R(D,N)\longrightarrow0
		\]
		都是阿贝尔群短正合列，当且仅当原序列是分裂短正合列。[为了证明序列分裂，取 $D=N$，并证明 $\Hom_R(N,N)$ 中恒等映射到 $\Hom_R(N,M)$ 的提升给出 $\varphi$ 的一个分裂同态。]
		\item 证明对所有 $R$-模 $D$，相伴序列
		\[
		0\longrightarrow\Hom_R(N,D)\xrightarrow{\varphi'}\Hom_R(M,D)
		\xrightarrow{\psi'}\Hom_R(L,D)\longrightarrow0
		\]
		都是阿贝尔群短正合列，当且仅当原序列是分裂短正合列。
	\end{enumerate}
	\item 设 $M$ 是左 $R$-模，其中 $R$ 是带 $1$ 的环。
	\begin{enumerate}[label=(\alph*)]
		\item 证明 $\Hom_{\ZZ}(R,M)$ 在作用
		\[
		(r\varphi)(r')=\varphi(r'r)
		\]
		下是左 $R$-模（见练习 10）。
		\item 设 $0\to A\xrightarrow{\psi}B$ 是 $R$-模的正合序列。证明若从 $A$ 到 $M$ 的每个同态 $f$ 都可提升为从 $B$ 到 $M$ 的同态 $F$，且 $f=F\circ\psi$，则从 $A$ 到 $\Hom_{\ZZ}(R,M)$ 的每个同态 $f'$ 也可提升为从 $B$ 到 $\Hom_{\ZZ}(R,M)$ 的同态 $F'$，且 $f'=F'\circ\psi$。[给定 $f'$，证明 $f(a)=f'(a)(1_R)$ 定义了 $A$ 到 $M$ 的同态。若 $F$ 是 $f$ 到 $B$ 的相应提升，证明 $F'(b)(r)=F(rb)$ 定义了 $B$ 到 $\Hom_{\ZZ}(R,M)$ 的同态，并提升 $f'$。]
		\item 证明若 $Q$ 是内射 $R$-模，则 $\Hom_{\ZZ}(R,Q)$ 也是内射 $R$-模。
	\end{enumerate}
	\item 本题证明定理 38：每个左 $R$-模 $M$ 都包含在一个内射左 $R$-模中。
	\begin{enumerate}[label=(\alph*)]
		\item 证明 $M$ 包含于一个内射 $\ZZ$-模 $Q$ 中。[把 $M$ 看成 $\ZZ$-模；使用推论 37。]
		\item 证明
		\[
		\Hom_R(R,M)\subseteq\Hom_{\ZZ}(R,M)\subseteq\Hom_{\ZZ}(R,Q).
		\]
		\item 利用 $R$-模同构 $M\cong\Hom_R(R,M)$（练习 10）以及前一题，推出 $M$ 包含于一个内射模中。
	\end{enumerate}
	\item 本题完成命题 34 的证明。设 $Q$ 是 $R$-模，并具有性质：每个短正合列
	\[
	0\longrightarrow Q\longrightarrow M_1\longrightarrow N\longrightarrow0
	\]
	都分裂。又设序列 $0\to L\xrightarrow{\psi}M$ 正合。证明每个 $R$-模同态 $f:L\to Q$ 都可提升为 $R$-模同态 $F:M\to Q$，满足 $f=F\circ\psi$。[由前一题，$Q$ 包含于一个内射 $R$-模中。使用分裂性质以及练习 4（注意练习 4 可用命题 34(2) 作为内射模的定义来证明）。]
	\item 证明 $\ZZ$-模 $\ZZ$ 的内射包是 $\QQ$。[设 $H$ 是 $\ZZ$ 的内射包，证明 $\QQ$ 包含 $H$ 的一个同构副本。利用 $H$ 的可除性证明 $1/n\in H$ 对所有非零整数 $n$ 成立，从而推出 $H=\QQ$。]
	\item 若 $F$ 是域，证明 $F$ 的内射包就是 $F$。
	\item 证明交换环 $R$ 上不定元 $x$ 的多项式环 $R[x]$ 是平坦 $R$-模。
	\item 设 $R$ 与 $S$ 是带单位元的环，$M$ 是右 $R$-模，$N$ 是 $(R,S)$-双模。若 $M$ 在 $R$ 上平坦，且 $N$ 作为 $S$-模平坦，证明 $M\otimes_RN$ 作为右 $S$-模平坦。
	\item 设 $R$ 是交换环，$M,N$ 是平坦 $R$-模。证明 $M\otimes_RN$ 是平坦 $R$-模。[使用前一题。]
	\item 证明由环同态 $f:R\to S$（$f(1_R)=1_S$，参见第 4 节推论 12 后例 6）把平坦右 $R$-模 $M$ 改变基环得到的右 $S$-模 $M\otimes_RS$ 是平坦 $S$-模。
	\item 证明 $A$ 是平坦 $R$-模当且仅当对任意左 $R$-模 $L,M$，其中 $L$ 有限生成，若 $\psi:L\to M$ 单射，则
	\[
	1\otimes\psi:A\otimes_RL\longrightarrow A\otimes_RM
	\]
	也是单射。[使用推论 42 证明中的技巧。]
	\item （平坦性判别法）本题 (a)--(c) 证明：$A$ 是平坦 $R$-模当且仅当对每个有限生成理想 $I$，由包含 $I\hookrightarrow R$ 诱导的映射
	\[
	A\otimes_RI\longrightarrow A\otimes_RR\cong A
	\]
	仍为单射，等价地 $A\otimes_RI\cong AI\subseteq A$。
	\begin{enumerate}[label=(\alph*)]
		\item 证明若 $A$ 平坦，则 $A\otimes_RI\to A\otimes_RR$ 单射。
		\item 若对每个有限生成理想 $I$，$A\otimes_RI\to A\otimes_RR$ 单射，证明对每个理想 $I$ 也有 $A\otimes_RI\to A\otimes_RR$ 单射。证明若 $K$ 是有限生成自由模 $F$ 的任意子模，则 $A\otimes_RK\to A\otimes_RF$ 单射。证明对任意自由模 $F$ 也成立。[参见推论 42 的证明。]
		\item 在 (b) 的假设下，设 $L,M$ 是 $R$-模，且 $L\xrightarrow{\psi}M$ 单射。证明 $A\otimes_RL\xrightarrow{1\otimes\psi}A\otimes_RM$ 单射，并推出 $A$ 平坦。[把 $M$ 写成自由模 $F$ 的商，得到短正合列
		\[
		0\to K\to F\to M\to0.
		\]
		令 $J=j^{-1}(\psi(L))$，$t:J\to F$ 为自然包含。证明相应交换图具有正合行；张量 $A$ 后的诱导图也交换且行正合。用 (b) 证明 $1\otimes t$ 单射，再用练习 1 推出 $1\otimes\psi$ 单射。]
		\item （商模的平坦性判别法）设 $A=F/K$，其中 $F$ 平坦（例如自由）且 $K$ 是 $F$ 的 $R$-子模。证明 $A$ 平坦当且仅当
		\[
		FI\cap K=KI
		\]
		对每个有限生成理想 $I$ 成立。[用 (a) 证明 $F\otimes_RI\cong FI$，并观察 $K\otimes_RI$ 的像是 $KI$；把正合列 $0\to K\to F\to A\to0$ 与 $I$ 张量，证明 $A\otimes_RI\cong FI/KI$，再应用平坦性判别法。]
	\end{enumerate}
	\item 设 $R$ 是 P.I.D.。本题证明 $A$ 是平坦 $R$-模当且仅当 $A$ 是无挠 $R$-模，即若 $a\in A$ 非零且 $r\in R$，$ra=0$ 蕴含 $r=0$。
	\begin{enumerate}[label=(\alph*)]
		\item 设 $A$ 平坦。对固定 $r\in R$，考虑由乘以 $r$ 定义的映射 $\psi_r:R\to R$，$\psi_r(x)=rx$。若 $r\neq0$，证明 $\psi_r$ 单射。由 $A$ 的平坦性推出从 $A$ 到 $A$ 的映射 $a\mapsto ra$ 单射，从而 $A$ 无挠。
		\item 设 $A$ 无挠。若 $I$ 是 $R$ 的非零理想，则 $I=rR$ 对某个非零 $r\in R$ 成立。证明 (a) 中的映射 $\psi_r$ 诱导 $R\cong I$ 的 $R$-模同构，且 $\psi_r$ 与包含 $t:I\to R$ 的复合 $R\xrightarrow{\psi_r}I\xrightarrow{t}R$ 是乘以 $r$。证明复合
		\[
		A\otimes_RR\xrightarrow{1\otimes\psi_r}A\otimes_RI
		\xrightarrow{1\otimes t}A\otimes_RR
		\]
		在 $A\otimes_RR=A$ 的识别下对应于 $a\mapsto ra$，并且由于 $A$ 无挠此复合单射。证明 $1\otimes\psi_r$ 是同构，并推出 $1\otimes t$ 单射。用前一题推出 $A$ 平坦。
	\end{enumerate}
	\item 设 $M,A,B$ 是 $R$-模。
	\begin{enumerate}[label=(\alph*)]
		\item 设 $f:A\to M$ 与 $g:B\to M$ 是 $R$-模同态。证明
		\[
		X=\{(a,b)\mid a\in A,
		b\in B,
		f(a)=g(b)\}
		\]
		是直和 $A\oplus B$ 的 $R$-子模，称为 $f$ 与 $g$ 的\emph{拉回}或\emph{纤维积}，并且存在以 $\pi_1,
		\pi_2$ 为到第一、第二分量的自然投影的交换图。
		\item 设 $f':M\to A$ 与 $g':M\to B$ 是 $R$-模同态。证明 $A\oplus B$ 对子模
		\[
		\{(f'(m),-g'(m))\mid m\in M\}
		\]
		的商 $Y$ 是 $R$-模，称为 $f'$ 与 $g'$ 的\emph{推出}或\emph{纤维和}，并且存在以 $\pi_1,
		\pi_2$ 为由到第一、第二分量的映射诱导到商上的自然映射的交换图。
	\end{enumerate}
	\item
	\begin{enumerate}[label=(\alph*)]
		\item （Schanuel 引理）若
		\[
		0\to K\to P\to M\to0,
		\qquad
		0\to K'\to P'\to M\to0
		\]
		是 $R$-模正合列，且 $P,P'$ 投射，证明
		\[
		P\oplus K'\cong P'\oplus K
		\]
		作为 $R$-模同构。[设 $X$ 是前一题中相应同态的纤维积，证明存在正合列 $0\to\ker\pi\to X\to P\to0$，且 $\ker\pi\cong K'$。推出 $X\cong P\oplus K'$。类似地证明 $X\cong P'\oplus K$。]
		\item 若
		\[
		0\to M\to Q\to L\to0,
		\qquad
		0\to M\to Q'\to L'\to0
		\]
		是 $R$-模正合列，且 $Q,Q'$ 内射，证明
		\[
		Q\oplus L'\cong Q'\oplus L
		\]
		作为 $R$-模同构。
	\end{enumerate}
\end{enumerate}
\noindent 两个 $R$-模 $M,N$ 称为\emph{投射等价}，若存在投射模 $P,P'$ 使 $M\oplus P\cong N\oplus P'$。类似地，若存在内射模 $Q,Q'$ 使 $M\oplus Q\cong N\oplus Q'$，则称 $M,N$ \emph{内射等价}。前一题表明其中的 $K$ 与 $K'$ 投射等价，$L$ 与 $L'$ 内射等价。


% \chapter{向量空间}

本章回顾任意域 $F$ 上有限维向量空间的基本理论（若干无限维向量空间理论见练习）。由于证明与实向量空间中相应论证完全相同，我们的处理十分简略。大体上，我们只包括本书其他部分将要用到的结果，因此像 Gauss--Jordan 消元、行阶梯形、求子空间的基的方法、矩阵的初等性质等基本主题，或者不予讨论，或者放在练习中讨论。因此，读者应把本章看作线性代数的复习，并看作域论与 Galois 理论的预备。特征多项式与特征值将在下一章中以更大的背景重新回顾与处理。

\section{定义与基本理论}

向量空间的术语与模的术语略有不同。也就是说，当环 $R$ 是域时，我们在上一章定义的 $R$-模的许多性质有不同的名称。下表给出这些新术语的词典（其中许多可能已经为读者所熟悉）。每一个相应的向量空间性质的定义都与模论中的定义完全相同，只是额外假定环 $R$ 是一个域（所以这些定义不再重复）。
\[
\begin{array}{c|c}
\text{$R$ 为任意环时的术语} & \text{$R$ 为域时的术语}\\ \hline
M\text{ 是 }R\text{-模} & M\text{ 是 }R\text{ 上的向量空间}\\
m\text{ 是 }M\text{ 的元素} & m\text{ 是 }M\text{ 中的向量}\\
a\text{ 是环元素} & a\text{ 是标量}\\
N\text{ 是 }M\text{ 的子模} & N\text{ 是 }M\text{ 的子空间}\\
M/N\text{ 是商模} & M/N\text{ 是商空间}\\
M\text{ 是秩为 }n\text{ 的自由模} & M\text{ 是维数为 }n\text{ 的向量空间}\\
M\text{ 是有限生成模} & M\text{ 是有限维向量空间}\\
M\text{ 是非零循环模} & M\text{ 是 }1\text{ 维向量空间}\\
\varphi:M\to N\text{ 是 }R\text{-模同态} & \varphi:M\to N\text{ 是线性变换}\\
M\text{ 与 }N\text{ 作为 }R\text{-模同构} & M\text{ 与 }N\text{ 作为向量空间同构}\\
A\subseteq M\text{ 生成 }M & A\subseteq M\text{ 张成 }M\\
M=RA & M\text{ 的每个元素都是 }A\text{ 中元素的线性组合，即 }M=\operatorname{Span}(A)
\end{array}
\]

本章余下部分中，$F$ 是一个域，$V$ 是 $F$ 上的向量空间。

我们将要证明的关于向量空间的最初结果之一，是向量空间都是自由 $F$-模，也就是说，它们有基。虽然我们的论证只处理有限维空间的情形，但任意向量空间的相应结果将在练习中作为 Zorn 引理的应用来证明。读者可先复习上一章关于自由模的一节，尤其是其中与同态有关的性质。

\noindent\textbf{定义.}
\begin{enumerate}[label=(\arabic*)]
\item $V$ 的子集 $S$ 称为一组\emph{线性无关向量}，若等式
\[
a_1v_1+a_2v_2+\cdots+a_nv_n=0,
\qquad a_1,a_2,\ldots,a_n\in F,
\quad v_1,v_2,\ldots,v_n\in S,
\]
蕴含 $a_1=a_2=\cdots=a_n=0$。
\item 向量空间 $V$ 的一组\emph{基}是一个有序的线性无关向量集，并且它张成 $V$。特别地，即使两组基只是彼此重排，也仍被看作不同的基。这有时称为\emph{有序基}。
\end{enumerate}

\paragraph{例}
\begin{enumerate}[label=(\arabic*)]
\item 变量 $x$ 的、系数来自域 $F$ 的多项式空间 $V=F[x]$ 特别地是 $F$ 上向量空间。元素 $1,x,x^2,\ldots$ 按定义线性无关（也就是说，一个多项式为 $0$ 当且仅当它的所有系数都为 $0$）。因为这些元素按定义也张成 $V$，它们是 $V$ 的一组基。
\item 一个线性齐次常系数微分方程（例如 $y''-3y'+2y=0$）在 $\mathbb C$ 上的解全体构成 $\mathbb C$ 上向量空间，因为微分是线性算子。这个向量空间中的元素线性无关，是指它们作为函数线性无关。例如，容易看出 $e^t$ 与 $e^{2t}$ 是方程 $y''-3y'+2y=0$ 的解（关于 $t$ 微分）。它们作为函数线性无关，因为 $ae^t+be^{2t}=0$ 蕴含 $a+b=0$（令 $t=0$）以及 $ae+be^2=0$（令 $t=1$），而这两个方程的唯一解是 $a=b=0$。微分方程中的一个定理说明，这些元素张成该方程的解集，因此它们是这个空间的一组基。
\end{enumerate}

\noindent\textbf{命题 1.} 设集合 $A=\{v_1,v_2,\ldots,v_n\}$ 张成向量空间 $V$，但 $A$ 的任何真子集都不张成 $V$。则 $A$ 是 $V$ 的一组基。特别地，$F$ 上任何有限生成（即有限张成）的向量空间都是自由 $F$-模。

\begin{proof}
只需证明 $v_1,v_2,\ldots,v_n$ 线性无关。设
\[
a_1v_1+a_2v_2+\cdots+a_nv_n=0
\]
且并非所有 $a_i$ 都为 $0$。重新排列后，可假定 $a_1\neq0$，于是
\[
v_1=-\frac1{a_1}(a_2v_2+\cdots+a_nv_n).
\]
由此可知，$\{v_2,v_3,\ldots,v_n\}$ 也张成 $V$，因为任何 $v_1,v_2,\ldots,v_n$ 的线性组合，都可利用上式写成 $v_2,v_3,\ldots,v_n$ 的线性组合。这是矛盾。
\end{proof}

\paragraph{例}
设 $F$ 是域，考虑 $F[x]/(f(x))$，其中
\[
f(x)=x^n+a_{n-1}x^{n-1}+\cdots+a_1x+a_0.
\]
理想 $(f(x))$ 是向量空间 $F[x]$ 的子空间，商 $F[x]/(f(x))$ 也是 $F$ 上向量空间。由 Euclidean 算法，每个多项式 $a(x)\in F[x]$ 都可唯一写成
\[
a(x)=q(x)f(x)+r(x),
\]
其中 $r(x)\in F[x]$ 且 $0\leq \deg r(x)\leq n-1$（或 $r=0$）。因为 $q(x)f(x)\in (f(x))$，所以商中的每个元素都由一个次数至多为 $n-1$ 的多项式 $r(x)$ 表示。两个不同的这种多项式不可能在商中相同，因为那将意味着它们的差（一个次数至多为 $n-1$ 的非零多项式）被 $f(x)$（次数为 $n$）整除。由此可知元素
\[
\overline{1},\overline{x},\overline{x}^2,\ldots,\overline{x}^{n-1}
\]
（横线一如往常表示这些元素在商中的像）作为 $F$ 上向量空间张成 $F[x]/(f(x))$，并且这些元素的任何真子集都不张成它，因此这些元素给出 $F[x]/(f(x))$ 的一组基。

\noindent\textbf{推论 2.} 设有限集 $A$ 张成向量空间 $V$。则 $A$ 含有 $V$ 的一组基。

\begin{proof}
取 $A$ 的一个张成 $V$ 的子集 $B$，使得 $B$ 的任何真子集都不再张成 $V$（显然存在这样的子集）。由命题 1，$B$ 是 $V$ 的一组基。
\end{proof}

\noindent\textbf{定理 3（替换定理）.} 设 $A=\{a_1,a_2,\ldots,a_n\}$ 是 $V$ 的一组含 $n$ 个元素的基，且 $\{b_1,b_2,\ldots,b_m\}$ 是 $V$ 中一组线性无关向量。则可给 $a_1,a_2,\ldots,a_n$ 适当排序，使得对每个 $k\in\{1,2,\ldots,m\}$，集合
\[
\{b_1,b_2,\ldots,b_k,a_{k+1},a_{k+2},\ldots,a_n\}
\]
都是 $V$ 的一组基。换言之，元素 $b_1,b_2,\ldots,b_m$ 可逐次替换基 $A$ 中的元素，并仍保持为基。特别地，$m\leq n$。

\begin{proof}
对 $k$ 归纳。若 $k=0$，无事可证，因为 $A$ 已给定为 $V$ 的一组基。现在假定
\[
\{b_1,b_2,\ldots,b_k,a_{k+1},a_{k+2},\ldots,a_n\}
\]
是 $V$ 的一组基。特别地，这是一组张成集，因此 $b_{k+1}$ 是线性组合：
\begin{equation}
 b_{k+1}=\beta_1b_1+\cdots+\beta_kb_k+\alpha_{k+1}a_{k+1}+\cdots+\alpha_na_n.
\tag{11.1}
\end{equation}
并非所有的 $\alpha_i$ 都为 $0$，否则 $b_{k+1}$ 将是 $b_1,b_2,\ldots,b_k$ 的线性组合，这与这些元素线性无关矛盾。必要时重新排列后，可假定 $\alpha_{k+1}\neq0$。于是由上式可把 $a_{k+1}$ 解为 $b_{k+1}$ 与 $b_1,b_2,\ldots,b_k,a_{k+2},\ldots,a_n$ 的线性组合，从而
\[
\operatorname{Span}\{b_1,b_2,\ldots,b_k,b_{k+1},a_{k+2},\ldots,a_n\}
=
\operatorname{Span}\{b_1,b_2,\ldots,b_k,a_{k+1},a_{k+2},\ldots,a_n\}.
\]
因此前者张成 $V$。还须证明 $b_1,\ldots,b_k,b_{k+1},a_{k+2},\ldots,a_n$ 线性无关。若
\begin{equation}
\gamma_1b_1+\cdots+\gamma_kb_k+\gamma_{k+1}b_{k+1}
+\alpha_{k+2}a_{k+2}+\cdots+\alpha_na_n=0,
\tag{11.2}
\end{equation}
则把式 (11.1) 中 $b_{k+1}$ 的表达式代入，得到旧集合
\[
\{b_1,b_2,\ldots,b_k,a_{k+1},a_{k+2},\ldots,a_n\}
\]
的一个等于 $0$ 的线性组合，其中 $a_{k+1}$ 的系数为 $\gamma_{k+1}\alpha_{k+1}$。由于这个旧集合按归纳假设是基，所有系数均为 $0$；特别地 $\gamma_{k+1}\alpha_{k+1}=0$，而 $\alpha_{k+1}\neq0$，故 $\gamma_{k+1}=0$。于是式 (11.2) 化为
\[
\gamma_1b_1+\cdots+\gamma_kb_k+
\alpha_{k+2}a_{k+2}+\cdots+\alpha_na_n=0.
\]
再由归纳假设，其他所有系数也都为 $0$。因此
$\{b_1,b_2,\ldots,b_k,b_{k+1},a_{k+2},\ldots,a_n\}$ 是 $V$ 的一组基，归纳完成。
\end{proof}

\noindent\textbf{推论 4.}
\begin{enumerate}[label=(\arabic*)]
\item 若 $V$ 有一组含 $n$ 个元素的有限基，则任意线性无关向量集至多有 $n$ 个元素，任意张成集至少有 $n$ 个元素。
\item 若 $V$ 有某个有限基，则 $V$ 的任意两组基有相同基数。
\end{enumerate}

\begin{proof}
(1) 这是定理 3 的最后结论与推论 2 的重述。

(2) 由 (1) 立即得到，因为一组基既是张成集又是线性无关集。
\end{proof}

\noindent\textbf{定义.} 若 $V$ 是有限生成 $F$-模（即有有限基），则任意一组基的基数称为 $V$ 的\emph{维数}，记为 $\dim_F V$；当 $F$ 从上下文中明确时，也简记为 $\dim V$。此时称 $V$ 在 $F$ 上\emph{有限维}。若 $V$ 不是有限生成的，则称 $V$ \emph{无限维}（记作 $\dim V=\infty$）。

\paragraph{例}
\begin{enumerate}[label=(\arabic*)]
\item 方程 $y''-3y'+2y=0$ 的解空间在 $\mathbb C$ 上维数为 $2$（例如以 $e^t,e^{2t}$ 为基）。一般地，微分方程中的一个定理说，$\mathbb C$ 上一个 $n$ 阶线性齐次常系数微分方程的解空间，是 $\mathbb C$ 上维数为 $n$ 的向量空间。
\item 上面考虑的由非零多项式 $f(x)$ 生成的商 $F[x]/(f(x))$ 在 $F$ 上维数为 $n=\deg f(x)$。空间 $F[x]$ 及其子空间 $(f(x))$ 都是 $F$ 上无限维向量空间。
\end{enumerate}

\noindent\textbf{推论 5（扩充引理）.} 若 $A$ 是有限维空间 $V$ 中的一组线性无关向量，则存在一组 $V$ 的基包含 $A$。

\begin{proof}
这也立即来自定理 3，因为我们可用 $A$ 中的元素逐次替换 $V$ 的任意给定基中的元素（这样的基因 $V$ 有限维而存在）。
\end{proof}

\noindent\textbf{定理 6.} 若 $V$ 是 $F$ 上 $n$ 维向量空间，则 $V\cong F^n$。特别地，$F$ 上任意两个同维有限维向量空间同构。

\begin{proof}
取 $V$ 的基 $v_1,v_2,\ldots,v_n$。定义映射
\[
\varphi:F^n\longrightarrow V
\]
为
\[
\varphi(a_1,a_2,\ldots,a_n)=a_1v_1+a_2v_2+\cdots+a_nv_n.
\]
映射 $\varphi$ 显然是 $F$-线性的；因为 $v_i$ 张成 $V$，它是满射；因为 $v_i$ 线性无关，它是单射。因此 $\varphi$ 是同构。
\end{proof}

\paragraph{例}
\begin{enumerate}[label=(\arabic*)]
\item 设 $\mathbb F$ 是含 $q$ 个元素的有限域，$W$ 是 $\mathbb F$ 上 $k$ 维向量空间。我们说明 $W$ 的不同有序基的数目为
\[
(q^k-1)(q^k-q)(q^k-q^2)\cdots(q^k-q^{k-1}).
\]
每组基都可如下构造。任意非零向量 $w_1$ 都可作为基的第一个元素。因为 $W\cong\mathbb F^k$，所以 $|W|=q^k$，故 $w_1$ 有 $q^k-1$ 种选择。任何不在 $w_1$ 张成的一维空间中的向量都与 $w_1$ 线性无关，因而可选作第二个基向量 $w_2$。一维空间同构于 $\mathbb F$，所以有 $q$ 个元素，因此 $w_2$ 有 $q^k-q$ 种选择。继续这样做，在第 $i$ 步，任何不在 $w_1,w_2,\ldots,w_{i-1}$ 张成的 $(i-1)$ 维空间中的向量，都与 $w_1,w_2,\ldots,w_{i-1}$ 线性无关，因而可选作第 $i$ 个基向量 $w_i$。一个 $(i-1)$ 维空间同构于 $\mathbb F^{i-1}$，所以有 $q^{i-1}$ 个元素。因此 $w_i$ 有 $q^k-q^{i-1}$ 种选择。这个过程在选定 $w_k$ 时终止，因为此时我们在 $k$ 维空间中已有 $k$ 个线性无关向量，因而得到一组基。
\item 设 $\mathbb F$ 是含 $q$ 个元素的有限域，$V$ 是 $\mathbb F$ 上 $n$ 维向量空间。对每个 $k\in\{1,2,\ldots,n\}$，我们说明 $V$ 中 $k$ 维子空间的数目为
\[
\frac{(q^n-1)(q^n-q)\cdots(q^n-q^{k-1})}
{(q^k-1)(q^k-q)\cdots(q^k-q^{k-1})}.
\]
任意 $k$ 维空间由 $k$ 个线性无关向量张成。与前一例相同地论证，上式分子是在 $n$ 维空间中选取 $k$ 个线性无关向量的方法数。两组 $k$ 个线性无关向量张成同一空间 $W$，当且仅当它们都是 $k$ 维空间 $W$ 的基。为了得到不同 $k$ 维子空间的数目，必须除以重复次数，即固定的 $k$ 维空间的基的数目。这个出现在分母中的因子正是例 1 中计算出的数目。
\end{enumerate}

接下来，我们证明一个重要关系，它联系子空间的维数、相应商空间的维数以及整个空间的维数。

\noindent\textbf{定理 7.} 设 $V$ 是 $F$ 上向量空间，$W$ 是 $V$ 的子空间。则 $V/W$ 是向量空间，并且
\[
\dim V=\dim W+\dim V/W,
\]
其中若一边为无限，则两边均为无限。

\begin{proof}
设 $W$ 在 $F$ 上维数为 $m$，$V$ 在 $F$ 上维数为 $n$，并令 $w_1,w_2,\ldots,w_m$ 为 $W$ 的一组基。由推论 5，这些 $V$ 中的线性无关元素可扩充为 $V$ 的一组基
\[
w_1,w_2,\ldots,w_m,v_{m+1},\ldots,v_n.
\]
从 $V$ 到 $V/W$ 的自然满投影把每个 $w_i$ 送到 $0$。没有非平凡的 $v_i$ 线性组合被送到 $0$，因为这会意味着这个线性组合是 $W$ 的元素，与 $v_i$ 的选择矛盾。因此这个投影映射的像 $V/W$ 同构于 $V$ 中由 $v_i$ 张成的子空间，故 $\dim V/W=n-m$。这证明了有限维情形。若任一边无限，则容易构造无限多个线性无关向量，从而证明另一边也无限。
\end{proof}

\noindent\textbf{推论 8.} 设 $\varphi:V\to U$ 是 $F$ 上向量空间的线性变换。则 $\ker\varphi$ 是 $V$ 的子空间，$\varphi(V)$ 是 $U$ 的子空间，并且
\[
\dim V=\dim\ker\varphi+\dim\varphi(V).
\]

\begin{proof}
这立即来自定理 7。注意，定理 7 的证明实际上就是本推论的特殊情形，其中 $U$ 是商 $V/W$，而 $\varphi$ 是自然投影同态。
\end{proof}

\noindent\textbf{推论 9.} 设 $\varphi:V\to W$ 是同一有限维数的两个向量空间之间的线性变换。则下列条件等价：
\begin{enumerate}[label=(\arabic*)]
\item $\varphi$ 是同构；
\item $\varphi$ 单，即 $\ker\varphi=0$；
\item $\varphi$ 满，即 $\varphi(V)=W$；
\item $\varphi$ 把 $V$ 的一组基送到 $W$ 的一组基。
\end{enumerate}

\begin{proof}
这些条件的等价性由推论 8 通过计算维数立即得到。
\end{proof}

\noindent\textbf{定义.} 若 $\varphi:V\to U$ 是 $F$ 上向量空间的线性变换，$\ker\varphi$ 有时称为 $\varphi$ 的\emph{零空间}，$\ker\varphi$ 的维数称为 $\varphi$ 的\emph{零化度}。$\varphi(V)$ 的维数称为 $\varphi$ 的\emph{秩}。若 $\ker\varphi=0$，称这个变换\emph{非奇异}。

\paragraph{例}
设 $F$ 是含 $q$ 个元素的有限域，$V$ 是 $F$ 上 $n$ 维向量空间。回忆，一般线性群 $GL(V)$ 是从 $V$ 到 $V$ 的所有非奇异线性变换构成的群（群运算为复合）。我们说明这个群的阶为
\[
|GL(V)|=(q^n-1)(q^n-q)(q^n-q^2)\cdots(q^n-q^{n-1}).
\]
为此，固定 $V$ 的一组基 $v_1,\ldots,v_n$。一个线性变换非奇异当且仅当它把这组基送到 $V$ 的另一组基。此外，若 $w_1,\ldots,w_n$ 是 $V$ 的任意一组基，则由 10.3 节定理 6，存在唯一线性变换把 $v_i$ 送到 $w_i$，$1\leq i\leq n$。因此，从 $V$ 到自身的非奇异线性变换数等于 $V$ 的不同基的数目。这个数在上面的例 1（取 $k=n$）中已经计算出来，正是 $GL(V)$ 的阶。

\paragraph{练习}
\begin{enumerate}[label=\arabic*.]
\item 令 $V=\mathbb R^n$，并令 $(a_1,a_2,\ldots,a_n)$ 为 $V$ 中固定向量。证明 $V$ 中所有满足 $a_1x_1+a_2x_2+\cdots+a_nx_n=0$ 的元素 $(x_1,x_2,\ldots,x_n)$ 构成 $V$ 的一个子空间。确定这个子空间的维数并找出一组基。
\item 令 $V$ 为变量 $x$ 中次数至多为 $5$ 的有理系数多项式全体。证明 $V$ 是 $\mathbb Q$ 上维数为 $6$ 的向量空间，且 $1,x,x^2,\ldots,x^5$ 为一组基。证明 $1,1+x,1+x+x^2,\ldots,1+x+x^2+x^3+x^4+x^5$ 也是 $V$ 的一组基。
\item 设 $\varphi$ 是线性变换 $\varphi:\mathbb R^4\to \mathbb R$，满足
\[
\varphi((1,0,0,0))=1,
\quad \varphi((1,-1,0,0))=0,
\quad \varphi((1,-1,1,0))=1,
\quad \varphi((1,-1,1,-1))=0.
\]
求 $\varphi((a,b,c,d))$。
\item 证明闭区间 $[a,b]$ 上所有实值函数构成的空间是 $\mathbb R$ 上无限维向量空间，其中 $a<b$。
\item 证明闭区间 $[a,b]$ 上所有连续实值函数构成的空间是 $\mathbb R$ 上无限维向量空间，其中 $a<b$。
\item 设 $V$ 为有限维向量空间。若 $\varphi$ 是 $V$ 到自身的任意线性变换，证明存在整数 $m$，使得 $\varphi^m$ 的像与 $\varphi^m$ 的核的交为 $\{0\}$。
\item 设 $\varphi$ 是 $n$ 维向量空间 $V$ 到自身的线性变换，且满足 $\varphi^2=0$。证明 $\varphi$ 的像包含于 $\varphi$ 的核，并由此证明 $\varphi$ 的秩至多为 $n/2$。
\item 设 $V$ 是 $F$ 上向量空间，$\varphi$ 是向量空间 $V$ 到自身的线性变换。若非零元素 $v\in V$ 满足 $\varphi(v)=\lambda v$，其中 $\lambda\in F$，则称 $v$ 为 $\varphi$ 的一个具有特征值 $\lambda$ 的特征向量。证明对任意固定的 $\lambda\in F$，$\varphi$ 的具有特征值 $\lambda$ 的所有特征向量连同 $0$ 构成 $V$ 的一个子空间。
\item 设 $V$ 是 $F$ 上向量空间，$\varphi$ 是向量空间 $V$ 到自身的线性变换。设对 $i=1,2,\ldots,k$，$v_i\in V$ 是 $\varphi$ 的特征向量，其特征值为 $\lambda_i\in F$（参见前一题），且所有特征值 $\lambda_i$ 两两不同。证明 $v_1,v_2,\ldots,v_k$ 线性无关。\emph{[提示：对 $k$ 归纳；写出 $v_i$ 之间的一个线性相关关系，并作用 $\varphi$，得到另一个涉及特征值的线性相关关系；再减去第一个线性关系的适当倍数，得到包含更少元素的线性相关关系。]} 推出 $n$ 维向量空间上的任意线性变换至多有 $n$ 个不同特征值。
\item 在以下练习中，令 $V$ 为域 $F$ 上任意维数的向量空间。证明任意向量空间 $V$ 都有一组基（约定空集是零空间的基）。\emph{[提示：令 $\mathcal S$ 为 $V$ 的所有线性无关向量子集构成的集合，按包含关系偏序；对 $\mathcal S$ 应用 Zorn 引理，并证明 $\mathcal S$ 的一个极大元素是一组基。]}
\item 改进上一题中的论证，证明 $V$ 的任意线性无关向量集都包含于 $V$ 的一组基中。
\item 若 $F$ 是有限域或可数域，$V$ 是 $F$ 上无限维向量空间，且有基 $B$，证明 $V$ 的基数等于 $B$ 的基数。在这种情况下推出 $V$ 的任意两组基有相同基数。
\item 证明作为 $\mathbb Q$ 上向量空间，对所有 $n\in\mathbb Z^+$，有 $\mathbb R^n\cong\mathbb R$（特别地，这意味着 $\mathbb R^n$ 与 $\mathbb R$ 作为加法阿贝尔群同构）。
\item 设 $A$ 是无限维空间 $V$ 的一组基。证明 $V$ 同构于由集合 $A$ 指标的域 $F$ 的若干副本的直和。证明由 $A$ 指标的 $F$ 的若干副本的直积是 $F$ 上向量空间，且其维数严格大于 $V$ 的维数（关于无限多个模的直和与直积的定义，见 10.3 节练习）。
\end{enumerate}

\section{线性变换的矩阵}

本节始终设 $V,W$ 是同一域 $F$ 上向量空间，$\mathcal B=\{v_1,v_2,\ldots,v_n\}$ 是 $V$ 的一组（有序）基，$\mathcal E=\{w_1,w_2,\ldots,w_m\}$ 是 $W$ 的一组（有序）基，且 $\varphi\in\Hom(V,W)$ 是从 $V$ 到 $W$ 的线性变换。对每个 $j\in\{1,2,\ldots,n\}$，把 $v_j$ 在 $\varphi$ 下的像写成基 $\mathcal E$ 的线性组合：
\begin{equation}
\varphi(v_j)=\sum_{i=1}^m a_{ij}w_i. \tag{11.3}
\end{equation}
令 $M_{\mathcal E}^{\mathcal B}(\varphi)=(a_{ij})$ 为 $m\times n$ 矩阵，其 $(i,j)$ 元为 $a_{ij}$（也就是说，把上面计算 $\varphi(v_j)$ 时 $w_i$ 的系数作为该矩阵第 $j$ 列）。矩阵 $M_{\mathcal E}^{\mathcal B}(\varphi)$ 称为 $\varphi$ 关于基 $\mathcal B,\mathcal E$ 的矩阵。出现在 $M$ 下标处的是定义域的基，上标处的是陪域的基。给定此矩阵后，可如下恢复线性变换 $\varphi$：若要计算 $v\in V$ 的 $\varphi(v)$，先把 $v$ 写成基 $\mathcal B$ 的线性组合：
\[
v=\sum_{i=1}^n\alpha_i v_i,
\qquad \alpha_i\in F,
\]
再计算 $m\times n$ 与 $n\times1$ 矩阵的乘积
\[
M_{\mathcal E}^{\mathcal B}(\varphi)
\begin{pmatrix}\alpha_1\\ \alpha_2\\ \vdots\\ \alpha_n\end{pmatrix}
=
\begin{pmatrix}\beta_1\\ \beta_2\\ \vdots\\ \beta_m\end{pmatrix}.
\]
则 $v$ 在 $\varphi$ 下的像为
\[
\varphi(v)=\sum_{i=1}^m\beta_iw_i.
\]
换言之，$\varphi(v)$ 关于基 $\mathcal E$ 的坐标列向量，是由矩阵 $M_{\mathcal E}^{\mathcal B}(\varphi)$ 左乘 $v$ 关于基 $\mathcal B$ 的坐标列向量得到的（有时记为 $[\varphi(v)]_{\mathcal E}=M_{\mathcal E}^{\mathcal B}(\varphi)[v]_{\mathcal B}$）。

\noindent\textbf{定义.} 上面与线性变换 $\varphi$ 关联的 $m\times n$ 矩阵 $A=(a_{ij})$ 称为关于基 $\mathcal B,\mathcal E$ \emph{表示}线性变换 $\varphi$。类似地，$\varphi$ 是关于基 $\mathcal B,\mathcal E$ 由 $A$ 表示的线性变换。

\paragraph{例}
\begin{enumerate}[label=(\arabic*)]
\item 令 $V=\mathbb R^3$，取标准基 $\mathcal B=\{(1,0,0),(0,1,0),(0,0,1)\}$；令 $W=\mathbb R^2$，取标准基 $\mathcal E=\{(1,0),(0,1)\}$。令 $\varphi$ 为线性变换
\[
\varphi(x,y,z)=(x+2y,x+y+z).
\]
因为 $\varphi(1,0,0)=(1,1)$，$\varphi(0,1,0)=(2,1)$，$\varphi(0,0,1)=(0,1)$，所以矩阵 $A=M_{\mathcal E}^{\mathcal B}(\varphi)$ 为
\[
\begin{pmatrix}1&2&0\\ 1&1&1\end{pmatrix}.
\]
\item 令 $V=W$ 为微分方程 $y''-3y'+2y=0$ 在 $\mathbb C$ 上的二维解空间，取 $\mathcal B=\mathcal E$ 为基 $v_1=e^t,v_2=e^{2t}$。因为方程系数为常数，容易检查若 $y$ 是解，则它的导数 $y'$ 也是解。因此映射 $\varphi=d/dt$（关于 $t$ 微分）是从 $V$ 到自身的线性变换。由于
\[
\varphi(v_1)=\frac{d(e^t)}{dt}=e^t=v_1,
\qquad
\varphi(v_2)=\frac{d(e^{2t})}{dt}=2e^{2t}=2v_2,
\]
相对于这些基的矩阵是对角矩阵
\[
\begin{pmatrix}1&0\\0&2\end{pmatrix}.
\]
\item 令 $V=W=\mathbb Q^3=\{(x,y,z)\mid x,y,z\in\mathbb Q\}$ 为有理数域 $F=\mathbb Q$ 上通常的三维有序三元组向量空间，并设 $\varphi$ 为线性变换
\[
\varphi(x,y,z)=(9x+4y+5z,-4x-3z,-6x-4y-2z),
\qquad x,y,z\in\mathbb Q,
\]
从 $V$ 到自身。对定义域与值域均取标准基 $e_1=(1,0,0)$，$e_2=(0,1,0)$，$e_3=(0,0,1)$。因为
\[
\varphi(1,0,0)=(9,-4,-6),\quad
\varphi(0,1,0)=(4,0,-4),\quad
\varphi(0,0,1)=(5,-3,-2),
\]
所以相对于这些基表示此线性变换的矩阵是
\[
A=\begin{pmatrix}
9&4&5\\
-4&0&-3\\
-6&-4&-2
\end{pmatrix}.
\]
\end{enumerate}

\noindent\textbf{定理 10.} 设 $V$ 是 $F$ 上维数为 $n$ 的向量空间，$W$ 是 $F$ 上维数为 $m$ 的向量空间，基分别为 $\mathcal B,\mathcal E$。则映射
\[
\Hom_F(V,W)\longrightarrow M_{m\times n}(F),
\qquad
\varphi\longmapsto M_{\mathcal E}^{\mathcal B}(\varphi)
\]
是一个向量空间同构。特别地，在固定基的选择后，线性变换与它们的相应矩阵之间存在一一对应。

\begin{proof}
矩阵 $M_{\mathcal E}^{\mathcal B}(\varphi)$ 的列由式 (11.3) 中 $\varphi$ 在基 $\mathcal B$ 上的作用决定。这特别说明映射 $\varphi\mapsto M_{\mathcal E}^{\mathcal B}(\varphi)$ 是 $F$-线性的，因为 $\varphi$ 是 $F$-线性的。此映射满，因为给定一个矩阵 $M$，先按式 (11.3) 在一组基上定义 $\varphi$，再按线性性延拓，便得到一个矩阵为 $M$ 的线性变换。此映射单，因为两个线性变换若在一组基上一致，则它们相同。
\end{proof}

注意，不同的基选择会给出不同的同构。因此，正如向量空间没有自然的基选择一样，$\Hom_F(V,W)$ 与 $M_{m\times n}(F)$ 之间也没有自然同构。

\noindent\textbf{推论 11.} $\Hom_F(V,W)$ 的维数是 $(\dim V)(\dim W)$。

\begin{proof}
$M_{m\times n}(F)$ 的维数为 $mn$。
\end{proof}

\noindent\textbf{定义.} 一个 $m\times n$ 矩阵 $A$ 称为\emph{非奇异}的，若 $Ax=0$ 且 $x\in F^n$ 蕴含 $x=0$。

矩阵与线性变换中“非奇异”一词的联系如下：设 $A=M_{\mathcal E}^{\mathcal B}(\varphi)$ 是与线性变换 $\varphi$ 关联的矩阵（相对于某些基 $\mathcal B,\mathcal E$）。则不依赖于基的选择，$m\times n$ 矩阵 $A$ 非奇异当且仅当线性变换 $\varphi$ 是从 $n$ 维空间 $V$ 到 $m$ 维空间 $W$ 的非奇异线性变换（参见练习）。

现在假设 $U,V,W$ 都是 $F$ 上有限维向量空间，其有序基分别为 $\mathcal D,\mathcal B,\mathcal E$，其中 $\mathcal B,\mathcal E$ 如前所述，并设 $\mathcal D=\{u_1,u_2,\ldots,u_k\}$。假设 $\psi:U\to V$ 与 $\varphi:V\to W$ 是线性变换。它们的复合 $\varphi\circ\psi$ 是从 $U$ 到 $W$ 的线性变换，因此可计算它相对于相应基的矩阵：$M_{\mathcal E}^{\mathcal D}(\varphi\circ\psi)$ 由计算
\[
\varphi\circ\psi(u_j)=\sum_{i=1}^m\gamma_{ij}w_i
\]
并把系数 $\gamma_{ij}$ 放在 $M_{\mathcal E}^{\mathcal D}(\varphi\circ\psi)$ 的第 $j$ 列得到。另一方面，分别计算 $\psi$ 与 $\varphi$ 的矩阵：
\[
\psi(u_j)=\sum_{p=1}^n\alpha_{pj}v_p,
\qquad
\varphi(v_p)=\sum_{i=1}^m\beta_{ip}w_i,
\]
于是 $M_{\mathcal B}^{\mathcal D}(\psi)=(\alpha_{pj})$，$M_{\mathcal E}^{\mathcal B}(\varphi)=(\beta_{ip})$。利用这些系数，可把 $\gamma$ 写成 $\alpha$ 与 $\beta$ 的表达式：
\[
\begin{aligned}
\varphi\circ\psi(u_j)
&=\varphi\left(\sum_{p=1}^n\alpha_{pj}v_p\right)\\
&=\sum_{p=1}^n\alpha_{pj}\varphi(v_p)\\
&=\sum_{p=1}^n\alpha_{pj}\sum_{i=1}^m\beta_{ip}w_i\\
&=\sum_{p=1}^n\sum_{i=1}^m\alpha_{pj}\beta_{ip}w_i.
\end{aligned}
\]
交换上面二重和的求和顺序可知，$w_i$ 的系数为
\[
\gamma_{ij}=\sum_{p=1}^n\alpha_{pj}\beta_{ip}.
\]
计算 $\varphi$ 与 $\psi$ 的矩阵积（按这个顺序）得到
\[
(\beta_{ij})(\alpha_{ij})=(\delta_{ij}),
\qquad
\delta_{ij}=\sum_{p=1}^n\beta_{ip}\alpha_{pj}.
\]
比较这两个和，并使用域乘法的交换性，可知对所有 $i,j$，有 $\gamma_{ij}=\delta_{ij}$。这个计算证明了如下结果。

\noindent\textbf{定理 12.} 用上述记号，
\[
M_{\mathcal E}^{\mathcal D}(\varphi\circ\psi)
=M_{\mathcal E}^{\mathcal B}(\varphi)M_{\mathcal B}^{\mathcal D}(\psi),
\]
也就是说，相对于相容的基选择，表示线性变换 $\varphi$ 与 $\psi$ 的矩阵的乘积，是表示复合线性变换 $\varphi\circ\psi$ 的矩阵。

\noindent\textbf{推论 13.} 矩阵乘法满足结合律与分配律（只要维数使得乘积有定义）。一个 $n\times n$ 矩阵 $A$ 非奇异当且仅当它可逆。

\begin{proof}
设 $A,B,C$ 是矩阵，使得乘积 $(AB)C$ 与 $A(BC)$ 有定义，并令 $S,T,R$ 表示相关线性变换。由定理 12，与 $AB$ 对应的线性变换为复合 $S\circ T$，所以与 $(AB)C$ 对应的线性变换为复合 $(S\circ T)\circ R$。类似地，与 $A(BC)$ 对应的线性变换为 $S\circ(T\circ R)$。由于函数复合满足结合律，这两个线性变换相同，故由定理 10 得 $(AB)C=A(BC)$。分配律类似证明。也可以通过对矩阵进行直接（虽然冗长的）计算证明这些结果。

若 $A$ 可逆，则 $Ax=0$ 蕴含 $x=A^{-1}Ax=A^{-1}0=0$，所以 $A$ 非奇异。反过来，若 $A$ 非奇异，固定 $V$ 的基 $\mathcal B,\mathcal E$，令 $\varphi$ 为关于这些基由 $A$ 表示的从 $V$ 到自身的线性变换。由推论 9，$\varphi$ 是 $V$ 到自身的同构，因而有逆 $\varphi^{-1}$。令 $B$ 为 $\varphi^{-1}$ 关于基 $\mathcal E,\mathcal B$（注意顺序）表示的矩阵。则
\[
AB=M_{\mathcal E}^{\mathcal B}(\varphi)M_{\mathcal B}^{\mathcal E}(\varphi^{-1})
=M_{\mathcal E}^{\mathcal E}(\varphi\circ\varphi^{-1})
=M_{\mathcal E}^{\mathcal E}(1)=I.
\]
同理 $BA=I$，所以 $B$ 是 $A$ 的逆。
\end{proof}

\noindent\textbf{推论 14.}
\begin{enumerate}[label=(\arabic*)]
\item 若 $\mathcal B$ 是 $n$ 维空间 $V$ 的一组基，则映射 $\varphi\mapsto M_{\mathcal B}^{\mathcal B}(\varphi)$ 是 $\Hom_F(V,V)$ 到 $M_n(F)$ 的环同构与向量空间同构。
\item $GL(V)\cong GL_n(F)$，其中 $\dim V=n$。特别地，若 $F$ 是有限域，则有限群 $GL_n(F)$ 的阶（等于 $|GL(V)|$）由第 1 节末尾的公式给出。
\end{enumerate}

\begin{proof}
(1) 在定理 10 中已经看到此映射是 $F$ 上向量空间同构。推论 13 说明 $M_n(F)$ 在矩阵乘法下是环，而定理 12 说明乘法在此映射下保持，因此它也是环同构。

(2) 由 (1) 立即得到，因为环同构把单位送到单位。
\end{proof}

\noindent\textbf{定义.} 若 $A$ 是一个以 $F$ 为系数的 $m\times n$ 矩阵，则 $A$ 的\emph{行秩}（分别地，\emph{列秩}）是 $A$ 的线性无关行（分别地，列）的最大个数（其中把 $A$ 的行或列分别看作仿射 $n$ 空间、$m$ 空间中的向量）。

矩阵的秩与相关线性变换的秩之间的关系如下：线性变换 $\varphi$ 的秩等于矩阵 $M_{\mathcal E}^{\mathcal B}(\varphi)$ 的列秩（参见练习）。我们还将看到，任意矩阵的行秩与列秩相同。

现在考虑从向量空间到自身的同一线性变换，在两种不同基选择下得到的两个矩阵之间的关系（线性变换从一个向量空间 $V$ 到另一个向量空间 $W$ 的一般陈述见练习）。

\noindent\textbf{定义.} 两个 $n\times n$ 矩阵 $A$ 与 $B$ 称为\emph{相似}，若存在可逆（即非奇异）$n\times n$ 矩阵 $P$ 使得 $P^{-1}AP=B$。从向量空间 $V$ 到自身的两个线性变换 $\varphi$ 与 $\psi$ 称为\emph{相似}，若存在从 $V$ 到 $V$ 的非奇异线性变换 $T$ 使得 $T^{-1}\varphi T=\psi$。

设 $\mathcal B$ 与 $\mathcal E$ 是同一向量空间 $V$ 的两组基，且 $\varphi\in\Hom_F(V,V)$。令 $I$ 为从 $V$ 到 $V$ 的恒等映射，并令 $P=M_{\mathcal B}^{\mathcal E}(I)$ 为其相应矩阵（换言之，把基 $\mathcal E$ 的元素用基 $\mathcal B$ 表示，注意顺序，并把所得坐标作为矩阵 $P$ 的列）。注意，若 $\mathcal B\neq\mathcal E$，则 $P$ 不是恒等矩阵。于是
\[
P^{-1}M_{\mathcal B}^{\mathcal B}(\varphi)P
=M_{\mathcal E}^{\mathcal E}(\varphi).
\]
若 $[v]_{\mathcal B}$ 是 $v\in V$ 关于基 $\mathcal B$ 的 $n\times1$ 坐标矩阵，$[v]_{\mathcal E}$ 类似定义，则
\[
[v]_{\mathcal B}=P[v]_{\mathcal E}.
\]
矩阵 $P$ 称为从 $\mathcal B$ 到 $\mathcal E$ 的\emph{过渡矩阵}或\emph{换基矩阵}，这个对 $M_{\mathcal B}^{\mathcal B}(\varphi)$ 的相似作用称为换基。这说明，同一线性变换相对于两组不同基的矩阵是相似的。

反过来，设 $A$ 与 $B$ 是由一个非奇异矩阵 $P$ 联系的相似 $n\times n$ 矩阵。令 $\mathcal B$ 为 $n$ 维向量空间 $V$ 的一组基。利用给定矩阵 $A$ 按式 (11.3) 定义从 $V$（基为 $\mathcal B$）到 $V$（仍以 $\mathcal B$ 为基）的线性变换 $\varphi$，即
\[
\varphi(v_j)=\sum_{i=1}^na_{ij}v_i.
\]
于是按 $\varphi$ 的定义，$A=M_{\mathcal B}^{\mathcal B}(\varphi)$。再用 $P$ 的第 $i$ 列作为 $w_i$ 关于基 $\mathcal B$ 的坐标，定义 $V$ 的一组新基 $\mathcal E$（因此按定义 $P=M_{\mathcal B}^{\mathcal E}(I)$）。于是
\[
B=P^{-1}AP=P^{-1}M_{\mathcal B}^{\mathcal B}(\varphi)P=M_{\mathcal E}^{\mathcal E}(\varphi)
\]
就是 $\varphi$ 关于基 $\mathcal E$ 的矩阵。这说明任意两个相似的 $n\times n$ 矩阵都以这种方式出现为同一线性变换关于两种不同基选择的表示矩阵。

注意，从 $V$ 到自身的线性变换的换基，就是由群 $GL(V)$ 中某个元素对 $V$ 的非奇异线性变换作共轭。特别地，“相似”关系是一个等价关系，其等价类就是 $GL(V)$ 在 $\Hom_F(V,V)$ 上按共轭作用的轨道。若 $\varphi\in GL(V)$（即 $\varphi$ 是可逆线性变换），则 $\varphi$ 的相似类正是群 $GL(V)$ 中 $\varphi$ 的共轭类。

\paragraph{例}
令 $V=\mathbb Q^3$，并令 $\varphi$ 为前面考虑过的线性变换
\[
\varphi(x,y,z)=(9x+4y+5z,-4x-3z,-6x-4y-2z),
\qquad x,y,z\in\mathbb Q.
\]
关于标准基 $\mathcal B$，$b_1=(1,0,0)$，$b_2=(0,1,0)$，$b_3=(0,0,1)$，我们已经看到表示此线性变换的矩阵为
\[
A=M_{\mathcal B}^{\mathcal B}(\varphi)=
\begin{pmatrix}
9&4&5\\
-4&0&-3\\
-6&-4&-2
\end{pmatrix}.
\]
现在取 $V$ 的基 $\mathcal E$：
\[
e_1=(2,-1,-2),\qquad e_2=(1,0,-1),\qquad e_3=(3,-2,-2)
\]
（我们马上会看到这确实是一组基）。由于
\[
\begin{aligned}
\varphi(e_1)&=\varphi(2,-1,-2)=(4,-2,-4)=2e_1+0e_2+0e_3,\\
\varphi(e_2)&=\varphi(1,0,-1)=(4,-1,-4)=1e_1+2e_2+0e_3,\\
\varphi(e_3)&=\varphi(3,-2,-2)=(9,-6,-6)=0e_1+0e_2+3e_3,
\end{aligned}
\]
所以关于这组基表示 $\varphi$ 的矩阵为
\[
B=M_{\mathcal E}^{\mathcal E}(\varphi)=
\begin{pmatrix}
2&1&0\\
0&2&0\\
0&0&3
\end{pmatrix}.
\]
把基 $\mathcal E$ 的元素用基 $\mathcal B$ 表示，有
\[
e_1=2b_1-b_2-2b_3,
\qquad e_2=b_1-b_3,
\qquad e_3=3b_1-2b_2-2b_3.
\]
故
\[
P=M_{\mathcal B}^{\mathcal E}(I)=
\begin{pmatrix}
2&1&3\\
-1&0&-2\\
-2&-1&-2
\end{pmatrix},
\qquad
P^{-1}=\begin{pmatrix}
-2&-1&-2\\
2&2&1\\
1&0&1
\end{pmatrix},
\]
它把 $A$ 共轭为 $B$，也就是说，$P^{-1}AP=B$，这很容易检查。（顺便注意，因为 $P$ 可逆，所以这证明了 $\mathcal E$ 确为 $V$ 的一组基。）

我们顺便指出，表示此线性变换 $\varphi$ 的矩阵 $B$ 比表示 $\varphi$ 的矩阵 $A$ 简单得多。研究给定线性变换的最简单可能矩阵表示（以及选择哪组基来实现它），就是下一章将讨论的典范形理论。

\paragraph{向量空间张量积上的线性变换}
为方便起见，我们在向量空间的特殊情形下重述 10.4 节的推论 18 与 19。

\noindent\textbf{命题 15.} 设 $F$ 是域 $K$ 的子域。若 $W$ 是 $F$ 上 $m$ 维向量空间，基为 $w_1,\ldots,w_m$，则 $K\otimes_F W$ 是 $K$ 上 $m$ 维向量空间，基为
\[
1\otimes w_1,\ldots,1\otimes w_m.
\]

\noindent\textbf{命题 16.} 设 $V$ 与 $W$ 是域 $F$ 上有限维向量空间，基分别为 $v_1,\ldots,v_n$ 与 $w_1,\ldots,w_m$。则 $V\otimes_F W$ 是 $F$ 上维数为 $nm$ 的向量空间，其一组基为
\[
v_i\otimes w_j,
\qquad 1\leq i\leq n,
\quad 1\leq j\leq m.
\]

\noindent\textbf{注记.} 若 $v$ 与 $w$ 分别是 $V$ 与 $W$ 的非零元素，则由命题可知 $v\otimes w$ 是 $V\otimes_F W$ 的非零元素，因为总可构造 $V$ 与 $W$ 的基，使得其第一个基向量分别为 $v,w$。在两个 $R$-模的张量积 $M\otimes_R N$ 中，当 $R$ 不是域时，一般而言要判定两个非零元素的张量积 $m\otimes n$ 何时为零，要困难得多。

现在令 $V,W,X,Y$ 是 $F$ 上有限维向量空间，并令
\[
\varphi:V\to X,
\qquad
\psi:W\to Y
\]
为线性变换。我们计算线性变换
\[
\varphi\otimes\psi:V\otimes W\to X\otimes Y
\]
的矩阵。令 $\mathcal B_1=\{v_1,\ldots,v_n\}$ 与 $\mathcal B_2=\{w_1,\ldots,w_m\}$ 分别为 $V$ 与 $W$ 的（有序）基；令 $\mathcal E_1=\{x_1,\ldots,x_r\}$ 与 $\mathcal E_2=\{y_1,\ldots,y_s\}$ 分别为 $X$ 与 $Y$ 的（有序）基。令 $\mathcal B=\{v_i\otimes w_j\}$ 与 $\mathcal E=\{x_i\otimes y_j\}$ 为由命题 16 给出的 $V\otimes W$ 与 $X\otimes Y$ 的基；我们稍后指定其顺序。设
\[
\varphi(v_i)=\sum_{p=1}^r\alpha_{pi}x_p,
\qquad
\psi(w_j)=\sum_{q=1}^s\beta_{qj}y_q.
\]
则
\begin{align}
(\varphi\otimes\psi)(v_i\otimes w_j)
&=\varphi(v_i)\otimes\psi(w_j)\notag\\
&=\left(\sum_{p=1}^r\alpha_{pi}x_p\right)
\otimes
\left(\sum_{q=1}^s\beta_{qj}y_q\right)\notag\\
&=\sum_{p=1}^r\sum_{q=1}^s\alpha_{pi}\beta_{qj}(x_p\otimes y_q).
\tag{11.8}
\end{align}
鉴于 (11.8) 中求和的顺序，我们把基 $\mathcal E$ 排成 $r$ 个有序列，第 $p$ 个列表为
\[
x_p\otimes y_1,
 x_p\otimes y_2,
 \ldots,
 x_p\otimes y_s,
\]
并类似地排列基 $\mathcal B$。于是方程 (11.8) 决定了 $\varphi\otimes\psi$ 的相应矩阵的列。所得矩阵 $M_{\mathcal E}^{\mathcal B}(\varphi\otimes\psi)$ 是一个 $r\times n$ 分块矩阵，其第 $(p,q)$ 块是 $s\times m$ 矩阵 $\alpha_{pq}M_{\mathcal E_2}^{\mathcal B_2}(\psi)$。换言之，$\varphi\otimes\psi$ 的矩阵由 $\varphi$ 的矩阵出发，把其中每个元素都乘以 $\psi$ 的矩阵得到。这种矩阵有一个名称。

\noindent\textbf{定义.} 设 $A=(a_{ij})$ 与 $B$ 分别为 $r\times n$ 与 $s\times m$ 矩阵，系数在任意交换环中。$A$ 与 $B$ 的 \emph{Kronecker 积}或\emph{张量积}，记为 $A\otimes B$，是一个 $rs\times nm$ 矩阵，它由 $r\times n$ 个分块组成，其中第 $(i,j)$ 块是 $s\times m$ 矩阵 $a_{ij}B$。

用这个术语，我们有：

\noindent\textbf{命题 17.} 设 $\varphi:V\to X$ 与 $\psi:W\to Y$ 是有限维向量空间之间的线性变换。则表示 $\varphi$ 与 $\psi$ 的矩阵的 Kronecker 积，是 $\varphi\otimes\psi$ 的一个矩阵表示。

\paragraph{例}
令 $V=X=\mathbb R^3$，二者都取基 $v_1,v_2,v_3$；令 $W=Y=\mathbb R^2$，二者都取基 $w_1,w_2$。设 $\varphi:\mathbb R^3\to\mathbb R^3$ 由
\[
\varphi(av_1+bv_2+cv_3)=cv_1+2av_2-3bv_3
\]
给出，$\psi:\mathbb R^2\to\mathbb R^2$ 由
\[
\psi(aw_1+bw_2)=(a+3b)w_1+(4b-2a)w_2
\]
给出。相对于所选基，$\varphi$ 与 $\psi$ 的矩阵分别为
\[
\begin{pmatrix}
0&0&1\\
2&0&0\\
0&-3&0
\end{pmatrix},
\qquad
\begin{pmatrix}
1&3\\
-2&4
\end{pmatrix}.
\]
因此相对于有序基
\[
\mathcal B=\{v_1\otimes w_1,
 v_1\otimes w_2,
 v_2\otimes w_1,
 v_2\otimes w_2,
 v_3\otimes w_1,
 v_3\otimes w_2\}
\]
有
\[
M_{\mathcal B}^{\mathcal B}(\varphi\otimes\psi)=
\begin{pmatrix}
0&0&0&0&1&3\\
0&0&0&0&-2&4\\
2&6&0&0&0&0\\
-4&8&0&0&0&0\\
0&0&-3&-9&0&0\\
0&0&6&-12&0&0
\end{pmatrix},
\]
这是通过依次用 $\varphi$ 的矩阵中的各个元素乘以 $\psi$ 的 $2\times2$ 矩阵得到的。

\paragraph{练习}
\begin{enumerate}[label=\arabic*.]
\item 令 $V$ 为变量 $x$ 中次数至多为 $5$ 的有理系数多项式全体。求从基 $1,x,x^2,\ldots,x^5$ 到基 $1,1+x,1+x+x^2,\ldots,1+x+x^2+x^3+x^4+x^5$ 的过渡矩阵。
\item 令 $V$ 为上一题的向量空间。令 $\varphi=d/dx$ 为 $V$ 到自身的线性变换，即对多项式通常关于 $x$ 求导。求 $\varphi$ 关于上一题中两组基的矩阵。
\item 令 $V$ 为变量 $x$ 中次数至多为 $n$ 的 $F$ 系数多项式全体。求从基 $1,x,x^2,\ldots,x^n$ 到元素 $1,x-\lambda,\ldots,(x-\lambda)^{n-1},(x-\lambda)^n$ 的过渡矩阵，其中 $\lambda$ 是 $F$ 的固定元素。推出这些元素是 $V$ 的一组基。
\item 令 $\varphi$ 是 $\mathbb R^2$ 到自身的线性变换，它把平面绕原点逆时针旋转角 $\theta$。证明 $\varphi$ 关于 $\mathbb R^2$ 标准基的矩阵为
\[
\begin{pmatrix}\cos\theta&-\sin\theta\\ \sin\theta&\cos\theta\end{pmatrix}.
\]
\item 证明：$m\times n$ 矩阵 $A$ 非奇异，当且仅当线性变换 $\varphi$ 是从 $n$ 维空间 $V$ 到 $m$ 维空间 $W$ 的非奇异线性变换，其中 $A=M_{\mathcal E}^{\mathcal B}(\varphi)$，且这与基 $\mathcal B,\mathcal E$ 的选择无关。
\item 证明若 $\varphi\in\Hom_F(F^n,F^m)$，且 $\mathcal B,\mathcal E$ 分别为 $F^n,F^m$ 的自然基，则 $\varphi$ 的值域等于 $M_{\mathcal E}^{\mathcal B}(\varphi)$ 的列向量集合的张成空间。推出 $\varphi$ 的秩（作为线性变换）等于 $M_{\mathcal E}^{\mathcal B}(\varphi)$ 的列秩。
\item 证明任意两个相似矩阵有相同的行秩与相同的列秩。
\item 设 $V$ 是 $F$ 上 $n$ 维向量空间，$\varphi$ 是向量空间 $V$ 到自身的线性变换。
\begin{enumerate}[label=(\alph*)]
\item 证明若 $V$ 有一组由 $\varphi$ 的特征向量组成的基（参见第 1 节练习 8），则 $\varphi$ 关于这组基（定义域和值域都取这组基）的矩阵是对角矩阵，对角元为相应特征值。
\item 若 $A$ 是关于 $V$ 的一组给定基（定义域和值域都取这组基）表示 $\varphi$ 的 $n\times n$ 矩阵，证明 $A$ 与对角矩阵相似，当且仅当 $V$ 有一组由 $\varphi$ 的特征向量组成的基。
\end{enumerate}
\item 若 $W$ 是向量空间 $V$ 的子空间，并且在 $\varphi$ 下稳定（即 $\varphi(W)\subseteq W$），证明 $\varphi$ 诱导 $W$ 上的线性变换 $\varphi|_W$，以及商向量空间 $V/W$ 上的线性变换 $\overline\varphi$。若 $\varphi|_W$ 与 $\overline\varphi$ 非奇异，证明 $\varphi$ 非奇异。若 $V$ 有限维，证明逆命题成立，并给出 $V$ 无限维时的一个反例。
\item 设 $V$ 是 $n$ 维向量空间，$\varphi$ 是 $V$ 到自身的线性变换。设 $W$ 是 $V$ 的一个在 $\varphi$ 下稳定的 $m$ 维子空间。
\begin{enumerate}[label=(\alph*)]
\item 证明存在 $V$ 的一组基，使得 $\varphi$ 的矩阵形如
\[
\begin{pmatrix}A&B\\0&C\end{pmatrix},
\]
其中 $A$ 是 $m\times m$ 矩阵，$B$ 是 $m\times(n-m)$ 矩阵，$C$ 是 $(n-m)\times(n-m)$ 矩阵（这样的矩阵称为分块上三角矩阵）。
\item 证明若存在 $\varphi$-不变子空间 $W'$，使得 $V=W\oplus W'$，则 $W$ 与 $W'$ 的基给出 $V$ 的一组基，相对于这组基，$\varphi$ 的矩阵为分块对角矩阵
\[
\begin{pmatrix}A&0\\0&C\end{pmatrix},
\]
其中 $A$ 是 $m\times m$ 矩阵，$C$ 是 $(n-m)\times(n-m)$ 矩阵。
\item 反过来证明，若存在 $V$ 的一组基，使得相对于这组基 $\varphi$ 的矩阵如 (b) 中那样分块对角，则存在维数分别为 $m$ 与 $n-m$ 的 $\varphi$-不变子空间 $W,W'$，并且 $V=W\oplus W'$。
\end{enumerate}
\item 设 $\varphi$ 是有限维向量空间 $V$ 到自身的线性变换，且 $\varphi^2=\varphi$。
\begin{enumerate}[label=(\alph*)]
\item 证明 $\operatorname{image}\varphi\cap\ker\varphi=0$。
\item 证明 $V=\operatorname{image}\varphi\oplus\ker\varphi$。
\item 证明存在 $V$ 的一组基，使得 $\varphi$ 关于这组基的矩阵是对角矩阵，其对角元全为 $0$ 或 $1$。
\end{enumerate}
满足 $\varphi^2=\varphi$ 的线性变换称为幂等线性变换。本题证明幂等线性变换只是投影到某个子空间上的映射。
\item 令 $V=\mathbb R^2$，$v_1=(1,0)$，$v_2=(0,1)$，所以 $v_1,v_2$ 是 $V$ 的一组基。令 $\varphi$ 是 $V$ 到自身的线性变换，其关于这组基的矩阵为
\[
\begin{pmatrix}2&1\\0&2\end{pmatrix}.
\]
证明若 $W$ 是由 $v_1$ 生成的子空间，则 $W$ 在 $\varphi$ 的作用下稳定。证明不存在 $\varphi$-不变子空间 $W'$ 使得 $V=W\oplus W'$。
\item 设 $V$ 是维数为 $n$ 的向量空间，$W$ 是维数为 $m$ 的向量空间，均在域 $F$ 上。设 $A$ 是表示线性变换 $\varphi:V\to W$ 的 $m\times n$ 矩阵，对应于 $V$ 的基 $\mathcal B_1$ 与 $W$ 的基 $\mathcal E_1$。类似地，设 $B$ 是相同线性变换 $\varphi$ 对应于 $V$ 的基 $\mathcal B_2$ 与 $W$ 的基 $\mathcal E_2$ 的 $m\times n$ 矩阵。令 $P=M_{\mathcal B_1}^{\mathcal B_2}(I)$，其中 $I$ 表示从 $V$ 到 $V$ 的恒等映射；令 $Q=M_{\mathcal E_1}^{\mathcal E_2}(I)$，其中 $I$ 表示从 $W$ 到 $W$ 的恒等映射。证明 $Q^{-1}=M_{\mathcal E_2}^{\mathcal E_1}(I)$，并且 $Q^{-1}AP=B$，从而给出同一线性变换相对于不同基选择的表示矩阵之间的一般关系。
\end{enumerate}

以下练习回顾 Gauss--Jordan 消元过程。这是解决许多涉及向量空间问题的最快计算方法之一：解线性方程组、求矩阵逆、计算行列式、确定一组向量的张成空间、判定线性无关或线性相关关系，以及确定线性变换的核与像。我们把有关定义一并列在这里，以便与练习相配合。

矩阵的一次初等行变换是指以下三种操作之一：交换两行；用某个非零标量乘以一行；把一行的某个标量倍加到另一行。若矩阵 $A$ 可通过有限次初等行变换化为矩阵 $C$，则称 $A$ 可行化简为 $C$。

一个 $m\times n$ 矩阵称为\emph{化简行阶梯形}，若：
\begin{enumerate}[label=(\alph*)]
\item 第 $i$ 行的第一个非零项 $a_{i j_i}$ 是 $1$，并且对应的第 $j_i$ 列中所有其他项都为 $0$；
\item $j_1<j_2<\cdots<j_r$，其中 $r$ 是非零行的数目，也就是说，每一行开头零的个数严格增加（因此称为阶梯形）。
\end{enumerate}
增广矩阵 $(A\mid C)$ 称为化简行阶梯形，若它的系数矩阵 $A$ 是化简行阶梯形。例如，下面的第一个矩阵是化简行阶梯形，而第二个矩阵不是：
\[
\begin{pmatrix}
1&0&2&0&2\\
0&1&1&0&18\\
0&0&0&1&-7\\
0&0&0&0&0
\end{pmatrix},
\qquad
\begin{pmatrix}
1&-1&0\\
0&0&1\\
0&0&0
\end{pmatrix}.
\]
在化简行阶梯形矩阵中，每个非零行的第一个非零项所在的列称为主元列，其余列称为非主元列。

\begin{enumerate}[label=\arabic*.,resume]
\item 证明若 $A$ 可行化简为 $C$，则 $C$ 也可行化简为 $A$。证明关系“$A\sim C$ 当且仅当 $A$ 可行化简为 $C$”是等价关系。\emph{[观察初等行变换都是可逆的。]} 同一等价类中的矩阵称为行等价。
\item 证明两个行等价矩阵的行秩相同。\emph{[只需证明两个相差一次初等行变换的矩阵。]}
\item 证明任意增广矩阵都可经过一系列初等行变换化为化简行阶梯形矩阵。\emph{[对矩阵的列数或行数归纳。]}
\item 设 $A$ 与 $C$ 是两个化简行阶梯形矩阵。证明若 $A$ 与 $C$ 行等价，则 $A=C$。
\item 证明一个化简行阶梯形矩阵的行秩等于其非零行数。
\item 证明下列矩阵的化简行阶梯形分别是上文紧接练习 16 前给出的两个矩阵：
\[
\begin{pmatrix}
1&4&8&-1&-1\\
2&3&9&0&-5\\
-2&2&-2&1&14\\
4&1&11&-13&-4
\end{pmatrix},
\qquad
\begin{pmatrix}
-3&3&1\\
1&-1&0\\
2&-2&-3
\end{pmatrix}.
\]
\item （线性方程组的求解）设 $(A'\mid C')$ 是增广矩阵 $(A\mid C)$ 的化简行阶梯形。显然，$A'$ 的零行数至少与 $(A'\mid C')$ 的零行数一样多。
\begin{enumerate}[label=(\alph*)]
\item 证明若 $A'$ 的零行数严格大于 $(A'\mid C')$ 的零行数，则方程组 $AX=C$ 无解。
\end{enumerate}
由 (a)，可假设 $A'$ 与 $(A'\mid C')$ 有相同数目的非零行 $r$（于是 $n\geq r$）。
\begin{enumerate}[label=(\alph*),resume]
\item 证明若 $r=n$，则方程组 $AX=C$ 恰有一个解。
\item 证明若 $r<n$，则方程组 $AX=C$ 有无穷多个解。事实上，证明与 $(A'\mid C')$ 的非主元列对应的 $n-r$ 个变量可任意指定，而与主元列对应的其余 $r$ 个变量随后被唯一确定。
\end{enumerate}
\item 求下列线性方程组的解：
\begin{enumerate}[label=(\alph*)]
\item $-3x+3y+z=5,\quad x-y=0,\quad 2x-2y=-3$。
\item $x-2y+z=5,\quad x-4y+6z=10,\quad 4x-11y+11z=12$。
\item $x-2y+z=5,\quad y-2z=17,\quad 2x-3y=27$。
\item $x+y-3z+2u=2,\quad 3x-2y+5z+u=1,\quad 6x+y-4z+3u=7,\quad 2x+2y-6z=4$。
\item $x+y+4z+8u-w=-1,\quad x+2y+3z+9u-5w=-2,\quad -2y+2z-2u+v+14w=3,\quad x+4y+z+11u-13w=-4$。
\end{enumerate}
\item 设 $A$ 与 $B$ 是两个行等价的 $m\times n$ 矩阵。
\begin{enumerate}[label=(\alph*)]
\item 证明齐次方程 $AX=0$ 的解集与齐次方程 $BX=0$ 的解集相同。\emph{[只需证明两个矩阵相差一次初等行变换的情形。]}
\item 证明把 $A$ 的列视为 $F^m$ 中向量时，$A$ 的列向量满足的任意线性相关关系，也被 $B$ 的列向量满足。
\item 由 (b) 推出，$A$ 的线性无关列的最大数目等于 $B$ 的线性无关列的最大数目。
\end{enumerate}
\item 设 $A'$ 是化简行阶梯形矩阵。
\begin{enumerate}[label=(\alph*)]
\item 证明 $A'$ 的非零行线性无关。证明 $A'$ 的主元列线性无关，且非主元列线性依赖于主元列。（注意主元元素的作用。）
\item 证明化简行阶梯形矩阵中的线性无关列数等于线性无关行数，即这样的矩阵的行秩与列秩相同。
\end{enumerate}
\item 使用前两题以及练习 15，证明任意矩阵的行秩等于列秩。
\item （逆矩阵的计算）设 $A$ 是 $n\times n$ 矩阵。
\begin{enumerate}[label=(\alph*)]
\item 证明 $A$ 有一个以 $B_1,B_2,\ldots,B_n$ 为列的逆矩阵 $B$，当且仅当方程组 $AB_1=e_1,AB_2=e_2,\ldots,AB_n=e_n$ 都有解，其中 $e_i$ 是标准基列向量。
\item 证明 $A$ 可逆当且仅当 $A$ 与 $n\times n$ 恒等矩阵行等价。
\item 证明 $A$ 有逆矩阵 $B$ 当且仅当增广矩阵 $(A\mid I)$ 可经行变换化为 $(I\mid B)$，其中 $I$ 是 $n\times n$ 恒等矩阵。
\end{enumerate}
\item 用行化简求下列矩阵的逆：
\[
A=\begin{pmatrix}-7&-1&-4\\7&1&3\\1&0&0\end{pmatrix},
\qquad
B=\begin{pmatrix}1&1&0&2\\0&2&1&-1\\0&2&0&0\\-1&1&1&0\end{pmatrix}.
\]
\item （计算张成空间、线性无关性与线性相关关系）设 $V$ 是有基 $e_1,e_2,\ldots,e_m$ 的 $m$ 维向量空间，$v_1,v_2,\ldots,v_n$ 是 $V$ 中向量。设 $A$ 是 $m\times n$ 矩阵，其列是 $v_i$ 关于基 $e_1,\ldots,e_m$ 的坐标；设 $A'$ 是 $A$ 的化简行阶梯形。
\begin{enumerate}[label=(\alph*)]
\item 设 $B$ 是任意与 $A$ 行等价的矩阵。令 $w_1,w_2,\ldots,w_n$ 是坐标列为 $B$ 的列的向量。证明 $v_1,\ldots,v_n$ 满足的任意线性关系
\begin{equation}
x_1v_1+x_2v_2+\cdots+x_nv_n=0
\tag{11.5}
\end{equation}
在把 $v_i$ 替换为 $w_i$ 后仍成立。
\item 证明由 $A'$ 的主元列给出的坐标的向量线性无关，而由非主元列给出坐标的向量线性依赖于它们。
\item （判断向量是否线性无关）证明 $v_1,v_2,\ldots,v_n$ 线性无关，当且仅当 $A'$ 有 $n$ 个非零行（即秩为 $n$）。
\item （确定向量之间的线性相关关系）由 (c)，$v_1,\ldots,v_n$ 线性相关当且仅当 $A'$ 有非主元列。证明定义线性相关关系 (11.5) 的解由 $A'$ 定义的线性方程给出；与非主元列对应的变量可任意指定，其余变量随后被唯一确定。
\item （确定向量组张成的子空间）证明由 $v_1,\ldots,v_n$ 张成的子空间 $W$ 的维数为 $r$，其中 $r$ 是 $A'$ 的非零行数；并且 $W$ 的一组基由原向量中与 $A'$ 的主元列对应的 $v_{j_i}$ 给出，$i=1,2,\ldots,r$。
\end{enumerate}
\item 令 $V=\mathbb R^5$，取标准基，并考虑向量
\[
\begin{aligned}
v_1&=(1,1,3,-2,3),&
 v_2&=(0,1,0,-1,0),&
 v_3&=(2,3,6,-5,6),\\
v_4&=(0,3,1,-3,1),&
 v_5&=(2,-1,-1,-1,-1).
\end{aligned}
\]
\begin{enumerate}[label=(\alph*)]
\item 证明以这些向量为列的矩阵
\[
A=\begin{pmatrix}
1&0&2&0&2\\
1&1&3&3&-1\\
3&0&6&1&-1\\
-2&-1&-5&-3&-1\\
3&0&6&1&-1
\end{pmatrix}
\]
的化简行阶梯形为
\[
A'=\begin{pmatrix}
1&0&2&0&2\\
0&1&1&0&18\\
0&0&0&1&-7\\
0&0&0&0&0\\
0&0&0&0&0
\end{pmatrix},
\]
其中第 $1,2,4$ 列为主元列，其余两列为非主元列。
\item 推出这些向量线性相关；由 $v_1,
\ldots,v_5$ 张成的子空间 $W$ 是 $3$ 维的，并且
\[
v_1=(1,1,3,-2,3),\quad v_2=(0,1,0,-1,0),\quad v_4=(0,3,1,-3,1)
\]
是 $W$ 的一组基。
\item 由 (a) 推出任意线性关系
\[
x_1v_1+x_2v_2+x_3v_3+x_4v_4+x_5v_5=0
\]
的系数满足
\[
x_1+2x_3+2x_5=0,
\qquad
x_2+x_3+18x_5=0,
\qquad
x_4-7x_5=0.
\]
由此推出与非主元列对应的第 $3$ 与第 $5$ 个变量 $x_3,x_5$ 可任意指定，其余变量唯一确定为
\[
x_1=-2x_3-2x_5,
\qquad
x_2=-x_3-18x_5,
\qquad
x_4=7x_5.
\]
特别地，证明
\[
-2v_1-v_2+v_3=0,
\qquad
-2v_1-18v_2+7v_4+v_5=0,
\]
分别对应于 $(x_3=1,x_5=0)$ 与 $(x_3=0,x_5=1)$。
\end{enumerate}
\item 对下列每组 $\mathbb R^4$ 中向量，判断它是否线性无关；若线性相关，求一个显式的线性相关关系。
\begin{enumerate}[label=(\alph*)]
\item $(1,-4,3,0),(0,-1,4,-3),(1,-1,1,-1),(2,2,-1,-3)$。
\item $(1,-2,4,1),(2,-3,9,-1),(1,0,6,-5),(2,-5,7,5)$。
\item $(1,-2,0,1),(2,-2,0,0),(-1,3,0,-2),(-2,1,0,1)$。
\item $(0,1,1,0),(1,0,1,1),(2,2,2,0),(0,-1,1,1)$。
\end{enumerate}
\item 对下列每组 $\mathbb R^4$ 中向量，求其张成的子空间并给出该子空间的一组基。
\begin{enumerate}[label=(\alph*)]
\item $(1,-2,5,3),(2,3,1,-4),(3,8,-3,-5)$。
\item $(2,-5,3,0),(0,-2,5,-3),(1,-1,1,-1),(-3,2,-1,2)$。
\item $(1,-2,0,1),(2,-2,0,0),(-1,3,0,-2),(-2,1,0,1)$。
\item $(1,1,0,-1),(1,2,3,0),(2,3,3,-1),(1,2,2,-2),(2,3,2,-3),(1,3,4,-3)$。
\end{enumerate}
\item （计算线性变换的像与核）设 $V$ 是有基 $e_1,\ldots,e_n$ 的 $n$ 维向量空间，$W$ 是有基 $f_1,\ldots,f_m$ 的 $m$ 维向量空间。设 $\varphi:V\to W$ 是线性变换，相对于这些基的矩阵为 $A=(a_{ij})$，其中
\[
\varphi(e_j)=\sum_{i=1}^m a_{ij}f_i,
\qquad j=1,2,\ldots,n.
\]
设 $A'$ 是与 $A$ 行等价的化简行阶梯形矩阵。
\begin{enumerate}[label=(\alph*)]
\item （确定线性变换的像）证明 $\varphi(V)$ 的维数为 $r$，其中 $r$ 是 $A'$ 的非零行数；并且 $\varphi(V)$ 的一组基由 $\varphi(e_{j_i})$ 给出，其中 $j_i$ 是 $A'$ 中主元列的指标。
\item （确定线性变换的核）$\ker\varphi$ 中的元素正是关于基 $e_1,\ldots,e_n$ 的坐标 $(x_1,\ldots,x_n)$ 满足方程 $AX=0$ 的向量。
\begin{enumerate}[label=(\roman*)]
\item 证明 $\varphi$ 单当且仅当 $A'$ 有 $n$ 个非零行（即秩为 $n$）。
\item 由 (i)，$\ker\varphi$ 非平凡当且仅当 $A'$ 有非主元列。证明与非主元列对应的变量可任意指定，其余变量被唯一确定，从而给出 $\ker\varphi$ 的元素。特别地，证明让一个非主元变量等于 $1$，其他非主元变量等于 $0$，并解出所有主元变量，所得向量的坐标给出 $\ker\varphi$ 的一组基。推出 $\dim\ker\varphi=n-r$，其中 $r$ 是 $A$ 的秩。
\end{enumerate}
\end{enumerate}
\item 令 $V=\mathbb R^5$，$W=\mathbb R^4$，都取标准基。设 $\varphi:V\to W$ 定义为
\[
\varphi(x,y,z,u,v)=(x+2y+3z+4u+4v,-2x-4y+2v,x+2y+u-2v,x+2y-v).
\]
\begin{enumerate}[label=(\alph*)]
\item 证明 $\varphi$ 对应矩阵为
\[
A=\begin{pmatrix}
1&2&3&4&4\\
-2&-4&0&0&2\\
1&2&0&1&-2\\
1&2&0&0&-1
\end{pmatrix},
\]
且与 $A$ 行等价的化简行阶梯形为
\[
A'=\begin{pmatrix}
1&2&0&0&-1\\
0&0&1&0&3\\
0&0&0&1&-1\\
0&0&0&0&0
\end{pmatrix},
\]
其中第 $1,3,4$ 列为主元列，其余两列为非主元列。
\item 推出 $\varphi$ 的像是 $3$ 维的，且 $V$ 的第 $1,3,4$ 个基元素的像，即 $(1,-2,1,1),(3,0,0,0),(4,0,1,0)$，构成 $\varphi(V)$ 的一组基。
\item 由 (a) 推出 $\ker\varphi$ 中的向量 $(x,y,z,u,v)$ 满足
\[
x+2y-v=0,
\qquad
z+3v=0,
\qquad
u-v=0.
\]
由此推出与非主元列对应的第 $2$ 与第 $5$ 个变量 $y,v$ 可任意指定，其余变量唯一确定为
\[
x=-2y+v,
\qquad
z=-3v,
\qquad
u=v.
\]
证明 $(-2,1,0,0,0)$ 与 $(1,0,-3,1,1)$ 构成 $\ker\varphi$ 的一组基，它们分别对应于 $(y=1,v=0)$ 与 $(y=0,v=1)$。
\end{enumerate}
\item 设 $\varphi$ 是 $\mathbb R^4$ 到自身的线性变换，其相对于标准基的矩阵为
\[
A=\begin{pmatrix}
-1&0&3&0\\
-2&2&1&-1\\
-1&1&0&-3\\
1&-2&-1&1
\end{pmatrix}.
\]
求 $\varphi$ 的像与核的一组基。
\item 设 $\varphi:\mathbb R^4\to\mathbb R^2$ 是线性变换，满足
\[
\varphi((1,0,0,0))=(1,-1),
\quad \varphi((1,-1,0,0))=(0,0),
\]
\[
\varphi((1,-1,1,0))=(1,-1),
\quad \varphi((1,-1,1,-1))=(0,0).
\]
求 $\varphi$ 的像与核的一组基。
\item 令 $V$ 为所有 $2\times2$ 实矩阵构成的向量空间，并令 $\varphi:V\to\mathbb R$ 把矩阵 $A\in V$ 送到 $A$ 的对角元之和（即 $A$ 的迹）。
\begin{enumerate}[label=(\alph*)]
\item 证明四个矩阵
\[
\begin{pmatrix}1&0\\0&0\end{pmatrix},
\begin{pmatrix}0&1\\0&0\end{pmatrix},
\begin{pmatrix}0&0\\1&0\end{pmatrix},
\begin{pmatrix}0&0\\0&1\end{pmatrix}
\]
构成 $V$ 的一组基。
\item 证明 $\varphi$ 是线性变换，并求 $\varphi$ 关于 (a) 中 $V$ 的基以及 $\mathbb R$ 的标准基的矩阵。求 $\ker\varphi$ 的维数与一组基。
\end{enumerate}
\item 令 $V$ 为 $\mathbb Q$ 上由变量 $x$ 中次数至多为 $5$ 的多项式构成的 $6$ 维向量空间。定义 $\varphi:V\to V$ 为
\[
\varphi(f)=x^2f''-6xf'+12f,
\]
其中 $f''$ 与 $f'$ 分别表示通常关于 $x$ 的二阶与一阶导数。
\begin{enumerate}[label=(\alph*)]
\item 证明 $\varphi$ 是 $V$ 到自身的线性变换。
\item 求 $\varphi$ 的像与核的一组基。
\end{enumerate}
\item 令 $V$ 为域 $F$ 上变量 $x$ 中次数至多为 $6$ 的多项式构成的 $7$ 维向量空间。令 $\varphi$ 是 $V$ 到自身的线性变换，定义为 $\varphi(f)=f'$，其中 $f'$ 表示通常关于 $x$ 的导数。对下列域，求 $\varphi$ 的像与核的一组基：
\begin{enumerate}[label=(\alph*)]
\item $F=\mathbb R$；
\item $F=\mathbb F_2$，二元素有限域（注意例如在该域上 $(x^2)'=2x=0$）；
\item $F=\mathbb F_3$；
\item $F=\mathbb F_5$。
\end{enumerate}
\item 设 $A,B$ 是方阵。证明它们的 Kronecker 积的迹等于它们的迹的乘积：
\[
\operatorname{tr}(A\otimes B)=\operatorname{tr}(A)\operatorname{tr}(B).
\]
（回忆方阵的迹是其对角元之和。）
\item 设 $F$ 是 $K$ 的子域，$\psi:V\to W$ 是有限维 $F$-向量空间之间的线性变换。
\begin{enumerate}[label=(\alph*)]
\item 证明 $1\otimes\psi$ 是从 $K\otimes_F V$ 到 $K\otimes_F W$ 的 $K$-线性变换（这里 $1$ 表示 $K$ 到自身的恒等映射）。
\item 设 $\mathcal B=\{v_1,\ldots,v_n\}$ 与 $\mathcal E=\{w_1,\ldots,w_m\}$ 分别为 $V$ 与 $W$ 的基。证明 $1\otimes\psi$ 关于基 $\{1\otimes v_1,\ldots,1\otimes v_n\}$ 与 $\{1\otimes w_1,\ldots,1\otimes w_m\}$ 的矩阵，与 $\psi$ 关于 $\mathcal B$ 与 $\mathcal E$ 的矩阵相同。
\end{enumerate}
\end{enumerate}

\section{对偶向量空间}

\noindent\textbf{定义.}
\begin{enumerate}[label=(\arabic*)]
\item 对任意 $F$ 上向量空间 $V$，令
\[
V^*=\Hom_F(V,F)
\]
为从 $V$ 到 $F$ 的线性变换空间，称为 $V$ 的\emph{对偶空间}。$V^*$ 的元素称为\emph{线性泛函}。
\item 若 $\mathcal B=\{v_1,v_2,\ldots,v_n\}$ 是有限维空间 $V$ 的一组基，对每个 $i\in\{1,2,\ldots,n\}$ 定义 $v_i^*\in V^*$，由其在基 $\mathcal B$ 上的作用给出：
\begin{equation}
v_i^*(v_j)=
\begin{cases}
1,& i=j,\\
0,& i\neq j,
\end{cases}
\qquad 1\leq j\leq n.
\tag{11.6}
\end{equation}
\end{enumerate}

\noindent\textbf{命题 18.} 用上述记号，$\{v_1^*,v_2^*,\ldots,v_n^*\}$ 是 $V^*$ 的一组基。特别地，若 $V$ 有限维，则 $V^*$ 与 $V$ 有相同维数。

\begin{proof}
注意，因为 $V$ 有限维，
\[
\dim V^*=\dim\Hom_F(V,F)=\dim V=n
\]
（推论 11），所以既然有 $n$ 个 $v_i^*$，只需证明它们线性无关。若
\[
a_1v_1^*+a_2v_2^*+\cdots+a_nv_n^*=0
\quad\text{于 }\Hom_F(V,F),
\]
则把这个元素作用于 $v_i$，并利用上面的式 (11.6)，得到 $a_i=0$。因为 $i$ 任意，这些元素线性无关。
\end{proof}

\noindent\textbf{定义.} $V^*$ 的基 $\{v_1^*,v_2^*,\ldots,v_n^*\}$ 称为 $\{v_1,v_2,\ldots,v_n\}$ 的\emph{对偶基}。

后面的练习说明，若 $V$ 无限维，则总有 $\dim V<\dim V^*$。对任意维数的空间，$V^*$ 是 $V$ 的“代数”对偶空间。若 $V$ 还有额外结构，例如连续结构（即拓扑），则也可定义其他类型的对偶空间（例如 $V$ 的连续对偶，要求线性泛函为连续映射）。阅读其他著作时必须注意（特别是分析书籍），在使用“对偶空间”与“线性泛函”时隐含了哪些限定词。

\paragraph{例}
令 $[a,b]$ 是 $\mathbb R$ 中闭区间，$V$ 是所有连续函数 $f:[a,b]\to\mathbb R$ 构成的实向量空间。若 $a<b$，则 $V$ 无限维。对每个 $g\in V$，函数
\[
\varphi_g:V\to\mathbb R,
\qquad
\varphi_g(f)=\int_a^b f(t)g(t)\,dt
\]
是 $V$ 上的线性泛函。

\noindent\textbf{定义.} $V^*$ 的对偶，即 $V^{**}$，称为 $V$ 的\emph{双对偶}或\emph{二重对偶}。

注意，对有限维空间 $V$，有 $\dim V=\dim V^*$，且 $\dim V^*=\dim V^{**}$，所以 $V$ 与 $V^{**}$ 是同构向量空间。对无限维空间，$\dim V<\dim V^{**}$（见练习），所以 $V$ 与 $V^{**}$ 不可能同构。在有限维空间的情形，有一种自然的，即与基无关或无坐标的方式来展示向量空间与其二重对偶之间的同构。更一般地说，基本思想如下：若 $X$ 是任意集合，$S$ 是从 $X$ 到域 $F$ 的函数集合，我们通常想到选定 $f\in S$ 并随着 $x$ 遍历 $X$ 来计算 $f(x)$。也可以反过来想：固定 $X$ 中一点 $x$，并随着 $f$ 遍历 $S$ 来计算 $f(x)$。后一个过程称为在 $x$ 处求值，它说明对每个 $x\in X$，有一个函数 $E_x:S\to F$ 定义为 $E_x(f)=f(x)$。这给出一个映射 $x\mapsto E_x$，从 $X$ 到 $S$ 上的 $F$ 值函数集合。若 $S$ “分离点”，即对 $X$ 中不同点 $x,y$，存在某个 $f\in S$ 使 $f(x)\neq f(y)$，则映射 $x\mapsto E_x$ 是单射。下一个定理的证明把这种“角色互换”过程应用到 $X=V$ 与 $S=V^*$ 的情形；证明 $E_x$ 是 $S$ 上的线性 $F$ 值函数，即 $E_x$ 属于 $V^*$ 的对偶空间；并证明映射 $x\mapsto E_x$ 是从 $V$ 到 $V^{**}$ 的线性变换。注意整个过程中没有提到“基”这个词（虽然知道 $V^{**}$ 的维数很方便，而这是通过选基得到的）。特别地，证明并不以熟悉的短语“取 $V$ 的一组基……”开头。

\noindent\textbf{定理 19.} 存在从 $V$ 到 $V^{**}$ 的自然单射线性变换。若 $V$ 有限维，则此线性变换是同构。

\begin{proof}
令 $v\in V$。定义映射（在 $v$ 处求值）
\[
E_v:V^*\to F,
\qquad
E_v(f)=f(v).
\]
则
\[
E_v(f+ag)=(f+ag)(v)=f(v)+ag(v)=E_v(f)+aE_v(g),
\]
所以 $E_v$ 是从 $V^*$ 到 $F$ 的线性变换。因此 $E_v\in\Hom_F(V^*,F)=V^{**}$。这定义了自然映射
\[
\Phi:V\to V^{**},
\qquad
\Phi(v)=E_v.
\]
映射 $\Phi$ 是线性的：对 $v,w\in V$ 与 $a\in F$，对每个 $f\in V^*$，
\[
E_{v+aw}(f)=f(v+aw)=f(v)+af(w)=E_v(f)+aE_w(f),
\]
所以
\[
\Phi(v+aw)=E_{v+aw}=E_v+aE_w=\Phi(v)+a\Phi(w).
\]
为证明 $\Phi$ 单，令 $v$ 是 $V$ 中任一非零向量。由扩充引理，存在包含 $v$ 的一组基 $\mathcal B$。定义线性变换 $f:V\to F$，使它把 $v$ 送到 $1$，并把 $\mathcal B-\{v\}$ 中每个元素都送到 $0$。则 $f\in V^*$ 且
\[
E_v(f)=f(v)=1.
\]
因此 $\Phi(v)=E_v$ 在 $V^{**}$ 中非零。这证明 $\ker\Phi=0$，即 $\Phi$ 单。

若 $V$ 有有限维数 $n$，则由命题 18，$V^*$ 以及 $V^{**}$ 也都有维数 $n$。在这种情况下，$\Phi$ 是从 $V$ 到同维有限维向量空间的单射线性变换，因而是同构。
\end{proof}

令 $V,W$ 是 $F$ 上有限维向量空间，基分别为 $\mathcal B,\mathcal E$，对偶基分别为 $\mathcal B^*,\mathcal E^*$。固定某个 $\varphi\in\Hom_F(V,W)$。则对每个 $f\in W^*$，复合 $f\circ\varphi$ 是从 $V$ 到 $F$ 的线性变换，即 $f\circ\varphi\in V^*$。因此映射 $f\mapsto f\circ\varphi$ 定义了从 $W^*$ 到 $V^*$ 的函数。我们把这个诱导在对偶空间上的函数记为 $\varphi^*$。

\noindent\textbf{定理 20.} 用上述记号，$\varphi^*$ 是从 $W^*$ 到 $V^*$ 的线性变换，并且 $M_{\mathcal B^*}^{\mathcal E^*}(\varphi^*)$ 是矩阵 $M_{\mathcal E}^{\mathcal B}(\varphi)$ 的转置（回忆矩阵 $(a_{ij})$ 的转置是矩阵 $(a_{ji})$）。

\begin{proof}
映射 $\varphi^*$ 线性，因为
\[
(f+ag)\circ\varphi=(f\circ\varphi)+a(g\circ\varphi).
\]
定义 $\varphi$ 的方程（由它的矩阵）为
\[
\varphi(v_j)=\sum_{i=1}^m a_{ij}w_i,
\qquad 1\leq j\leq n.
\]
为计算 $\varphi^*$ 的矩阵，由 $\varphi^*$ 与 $w_k^*$ 的定义，有
\[
\varphi^*(w_k^*)(v_j)
=(w_k^*\circ\varphi)(v_j)
=w_k^*\left(\sum_{i=1}^m a_{ij}w_i\right)
=a_{kj}.
\]
另一方面，
\[
\left(\sum_{i=1}^n a_{ki}v_i^*\right)(v_j)=a_{kj}
\]
对所有 $j$ 成立。这说明下面两个线性泛函在 $V$ 的一组基上一致，因而它们是 $V^*$ 的同一个元素：
\[
\varphi^*(w_k^*)=\sum_{i=1}^n a_{ki}v_i^*.
\]
这就把 $\varphi^*$ 关于基 $\mathcal E^*$ 与 $\mathcal B^*$ 的矩阵确定为 $\varphi$ 的矩阵的转置。
\end{proof}

\noindent\textbf{推论 21.} 对任意矩阵 $A$，$A$ 的行秩等于 $A$ 的列秩。

\begin{proof}
令 $\varphi:V\to W$ 为一个线性变换，它关于 $V,W$ 的某些固定基的矩阵为 $A$。由定理 20，$\varphi^*:W^*\to V^*$ 关于对偶基的矩阵是 $A$ 的转置。$A$ 的列秩是 $\varphi$ 的秩；而 $A$ 的行秩（即 $A$ 的转置的列秩）是 $\varphi^*$ 的秩（参见第 2 节练习 6）。因此只需证明 $\varphi$ 与 $\varphi^*$ 有相同的秩。现在
\[
f\in\ker\varphi^*
\Longleftrightarrow
\varphi^*(f)=0
\Longleftrightarrow
f\circ\varphi(v)=0\quad\text{对所有 }v\in V,
\]
\[
\Longleftrightarrow
\varphi(V)\subseteq\ker f
\Longleftrightarrow
f\in\Ann(\varphi(V)),
\]
其中 $\Ann(S)$ 是下面练习 3 中描述的 $S$ 的零化子。因此 $\Ann(\varphi(V))=\ker\varphi^*$。由练习 3，
\[
\dim\Ann(\varphi(V))=\dim W-\dim\varphi(V).
\]
由推论 8，
\[
\dim\ker\varphi^*=\dim W^*-\dim\varphi^*(W^*).
\]
由于 $W$ 与 $W^*$ 维数相同，得到
\[
\dim\varphi(V)=\dim\varphi^*(W^*),
\]
这正是所需结论。
\end{proof}

\paragraph{练习}
\begin{enumerate}[label=\arabic*.]
\item 设 $V$ 是有限维向量空间。证明定理 20 中的映射 $\varphi\mapsto\varphi^*$ 给出 $\End(V)$ 与 $\End(V^*)$ 之间的环同构。
\item 令 $V$ 为变量 $x$ 中次数至多为 $5$ 的有理系数多项式全体，并取基 $1,x,x^2,\ldots,x^5$。证明下列映射都是 $V$ 的对偶空间中的元素，并把它们表示为对偶基的线性组合：
\begin{enumerate}[label=(\alph*)]
\item $E:V\to\mathbb Q$，定义为 $E(p(x))=p(3)$，即在 $x=3$ 处取值。
\item $\varphi:V\to\mathbb Q$，定义为 $\varphi(p(x))=\int_0^1p(t)\,dt$。
\item $\varphi:V\to\mathbb Q$，定义为 $\varphi(p(x))=\int_0^1t^2p(t)\,dt$。
\item $\varphi:V\to\mathbb Q$，定义为 $\varphi(p(x))=p'(5)$，其中 $p'(x)$ 是 $p(x)$ 关于 $x$ 的通常导数。
\end{enumerate}
\item 设 $S$ 是有限维空间 $V$ 的对偶空间 $V^*$ 的任意子集。定义 $S$ 在 $V$ 中的零化子为
\[
\Ann(S)=\{v\in V\mid f(v)=0\text{ 对所有 }f\in S\}.
\]
\begin{enumerate}[label=(\alph*)]
\item 证明 $\Ann(S)$ 是 $V$ 的子空间。
\item 设 $W_1,W_2$ 是 $V^*$ 的子空间。证明
\[
\Ann(W_1+W_2)=\Ann(W_1)\cap\Ann(W_2),
\qquad
\Ann(W_1\cap W_2)=\Ann(W_1)+\Ann(W_2).
\]
\item 设 $W_1,W_2$ 是 $V^*$ 的子空间。证明 $W_1=W_2$ 当且仅当 $\Ann(W_1)=\Ann(W_2)$。
\item 证明 $S$ 的零化子等于由 $S$ 张成的 $V^*$ 的子空间的零化子。
\item 假设 $V$ 有限维，基为 $v_1,\ldots,v_n$。若 $S=\{v_1^*,\ldots,v_k^*\}$，其中 $k\leq n$，证明 $\Ann(S)$ 是由 $v_{k+1},\ldots,v_n$ 张成的子空间。
\item 假设 $V$ 有限维。证明若 $W^*$ 是 $V^*$ 的任意子空间，则
\[
\dim\Ann(W^*)=\dim V-\dim W^*.
\]
\end{enumerate}
\item 若 $V$ 是无限维向量空间，且 $A$ 是其一组基，证明 $A^*=\{v^*\mid v\in A\}$ 不张成 $V^*$。
\item 若 $V$ 是无限维向量空间，且 $A$ 是其一组基，证明 $V^*$ 同构于由 $A$ 指标的域 $F$ 的若干副本的直积。由此推出 $\dim V^*>\dim V$。\emph{[使用第 1 节练习 14。]}
\end{enumerate}

\section{行列式}

虽然本书中行列式主要用于域上的向量空间，但行列式理论可以在任意带 $1$ 的交换环上发展。本节设 $R$ 是带 $1$ 的交换环。

\noindent\textbf{定义.}
\begin{enumerate}[label=(\arabic*)]
\item 映射
\[
\varphi:V_1\times V_2\times\cdots\times V_n\to W
\]
称为\emph{多重线性}，若在固定除第 $i$ 个变量以外的所有变量后，所得映射 $V_i\to W$ 是 $R$-模同态。若 $V_i=V$，称其为 $V$ 上的 $n$-重线性函数；若又 $W=R$，称其为 $n$-重线性形式。
\item $V$ 上的 $n$-重线性函数 $\varphi$ 称为\emph{交错}，若当某两个相邻变量相等时 $\varphi$ 的值为 $0$。若任意交换两个变量都不改变函数值，则称它\emph{对称}。
\end{enumerate}

\paragraph{例}
点积是 $\mathbb R^3$ 上的对称双线性形式。三重积
\[
(u,v,w)\longmapsto u\cdot(v\times w)
\]
是 $\mathbb R^3$ 上的交错三重线性形式。

\noindent\textbf{命题 22.} 设 $\varphi$ 是 $V$ 上的 $n$-重线性交错函数。则：
\begin{enumerate}[label=(\arabic*)]
\item 交换两个相邻变量会使 $\varphi$ 的值变为相反数；
\item 对任意 $\sigma\in S_n$，
\[
\varphi(v_{\sigma(1)},\ldots,v_{\sigma(n)})
=\epsilon(\sigma)\varphi(v_1,\ldots,v_n),
\]
其中 $\epsilon(\sigma)$ 是置换 $\sigma$ 的符号；
\item 若任意两个变量相等，则 $\varphi(v_1,\ldots,v_n)=0$；
\item 若把 $v_i$ 替换为 $v_i+av_j$，其中 $j\neq i$，则 $\varphi$ 的值不变。
\end{enumerate}

\begin{proof}
对 (1)，由交错性与多重线性性，
\[
0=\varphi(\ldots,v_i+v_{i+1},v_i+v_{i+1},\ldots)
\]
展开后，含有两个相邻相等变量的两项为 $0$，故
\[
\varphi(\ldots,v_i,v_{i+1},\ldots)+
\varphi(\ldots,v_{i+1},v_i,\ldots)=0.
\]
这证明交换相邻两个变量会改变符号。任意置换都可写成相邻换位的乘积，且符号就是所用换位次数的奇偶性，因此 (2) 成立。(3) 来自 (2)：若两个变量相等，交换这两个变量不改变输入，却把值乘以 $-1$，故值为其相反数；在一般交换环上也可通过把这两个相等变量用相邻换位移到相邻位置，然后用交错性得到 $0$。(4) 由多重线性性，替换后增加的项含有两个相等的变量 $v_j$，由 (3) 为 $0$。
\end{proof}

\noindent\textbf{命题 23.} 设 $\varphi$ 是 $V$ 上的 $n$-重线性交错函数，并且对 $1\leq j\leq n$，
\[
w_j=a_{1j}v_1+a_{2j}v_2+\cdots+a_{nj}v_n.
\]
则
\[
\varphi(w_1,\ldots,w_n)
=\sum_{\sigma\in S_n}\epsilon(\sigma)
 a_{\sigma(1),1}a_{\sigma(2),2}\cdots a_{\sigma(n),n}
 \varphi(v_1,
\ldots,v_n).
\]

\begin{proof}
由多重线性性，把每个 $w_j$ 的表达式代入并展开，得到对所有 $i_1,
\ldots,i_n\in\{1,\ldots,n\}$ 的求和项
\[
a_{i_1,1}a_{i_2,2}\cdots a_{i_n,n}
\varphi(v_{i_1},v_{i_2},\ldots,v_{i_n}).
\]
若某两个指标相等，则由命题 22(3)，相应项为 $0$。剩下的项正是指标 $(i_1,
\ldots,i_n)$ 是 $1,
2,
\ldots,n$ 的一个排列的项，即 $i_j=\sigma(j)$。对这些项，由命题 22(2) 得
\[
\varphi(v_{\sigma(1)},\ldots,v_{\sigma(n)})
=\epsilon(\sigma)\varphi(v_1,
\ldots,v_n),
\]
从而得到所述公式。
\end{proof}

\noindent\textbf{定义.} $R$ 上的一个 $n\times n$ \emph{行列式函数}是任意函数
\[
\det:M_{n\times n}(R)\to R
\]
满足：
\begin{enumerate}[label=(\arabic*)]
\item 它作为列向量的函数是 $n$-重线性交错形式；
\item $\det I=1$，其中 $I$ 是 $n\times n$ 恒等矩阵。
\end{enumerate}

\noindent\textbf{定理 24.} $R$ 上存在唯一的 $n\times n$ 行列式函数，并且对矩阵 $A=(a_{ij})$，它由公式
\[
\det(a_{ij})=
\sum_{\sigma\in S_n}\epsilon(\sigma)
 a_{\sigma(1),1}a_{\sigma(2),2}\cdots a_{\sigma(n),n}
\]
计算。

\begin{proof}
先用右边的公式定义一个函数。多重线性性来自每一列中各项只一次出现且求和可分配；若两列相等，则成对的置换项相消，因而该函数是交错的；对恒等矩阵，只有恒等置换给出非零项，故其值为 $1$。因此存在性成立。

为证明唯一性，令 $e_i$ 是第 $i$ 个位置为 $1$、其余位置为 $0$ 的列向量。若 $A_1,A_2,
\ldots,A_n$ 是 $A$ 的列，则
\[
A_1=a_{11}e_1+a_{21}e_2+\cdots+a_{n1}e_n,
\qquad
A_2=a_{12}e_1+a_{22}e_2+\cdots+a_{n2}e_n,
\]
依此类推。由命题 23，
\[
\det A=
\sum_{\sigma\in S_n}\epsilon(\sigma)
 a_{\sigma(1),1}a_{\sigma(2),2}\cdots a_{\sigma(n),n}
 \det(e_1,e_2,
\ldots,e_n).
\]
由行列式函数公理 (2)，$\det(e_1,e_2,
\ldots,e_n)=1$，所以 $\det A$ 的值只能是所给公式。这证明唯一性。
\end{proof}

\noindent\textbf{推论 25.} 行列式也是 $M_{n\times n}(R)$ 的行向量的 $n$-重线性函数，并且对任意 $n\times n$ 矩阵 $A$，
\[
\det A=\det(A^t).
\]

\begin{proof}
第一句话是第二句话的直接结果，所以只需证明矩阵与其转置有相同行列式。对 $A=(a_{ij})$，计算 $A^t$ 的行列式得
\[
\det(A^t)=\sum_{\sigma\in S_n}\epsilon(\sigma)
 a_{1,\sigma(1)}a_{2,\sigma(2)}\cdots a_{n,\sigma(n)}.
\]
由于 $1,
\ldots,n$ 每个数在 $\sigma(1),
\ldots,
\sigma(n)$ 中恰出现一次，可把乘积重新排列成
\[
a_{\sigma^{-1}(1),1}a_{\sigma^{-1}(2),2}\cdots a_{\sigma^{-1}(n),n}.
\]
又因为符号同态满足 $\epsilon(\sigma^{-1})=\epsilon(\sigma)$，当 $\sigma$ 遍历 $S_n$ 时 $\sigma^{-1}$ 也遍历 $S_n$，故上式的和与 $\det A$ 的公式相同。
\end{proof}

\noindent\textbf{定理 26（Cramer 法则）.} 若 $A_1,A_2,
\ldots,A_n$ 是 $n\times n$ 矩阵 $A$ 的列，且
\[
B=\beta_1A_1+\cdots+\beta_nA_n,
\]
则
\[
\beta_i\det A=
\det(A_1,
\ldots,A_{i-1},B,A_{i+1},
\ldots,A_n).
\]

\begin{proof}
这立即来自命题 22(3)：把第 $i$ 列替换为给定的线性组合 $B$ 并利用关于列的线性性展开，除含 $A_i$ 的项外，其他项都因出现两列相同而为 $0$。
\end{proof}

\noindent\textbf{推论 27.} 若 $R$ 是整环，则对 $A\in M_n(R)$，$\det A=0$ 当且仅当 $A$ 的列向量在自由 $R$-模 $R^n$ 中线性相关。等价地，$\det A=0$ 当且仅当 $A$ 的行向量线性相关。

\begin{proof}
若列向量线性相关，则某一列是其余列的线性组合（在整环的分式域中，或先取一个非零系数解出），由交错性与多重线性性得行列式为 $0$。反过来，若 $\det A=0$，把 $R$ 嵌入其分式域 $K$，则 $A$ 在 $K$ 上作为矩阵行列式也为 $0$，所以列向量在 $K^n$ 中线性相关。清除分母得到在 $R^n$ 中的非平凡线性相关关系。行向量情形由推论 25 得到。
\end{proof}

\noindent\textbf{定理 28.} 对矩阵 $A,B\in M_n(R)$，
\[
\det(AB)=(\det A)(\det B).
\]

\begin{proof}
把 $A$ 固定，并把函数 $D(B)=\det(AB)$ 看作 $B$ 的列向量的函数。由于 $AB$ 的每一列是 $A$ 乘以 $B$ 的相应列，$D$ 对 $B$ 的列是 $n$-重线性的；若 $B$ 有两列相同，则 $AB$ 有相应两列相同，所以 $D(B)=0$。因此 $D$ 是 $B$ 的列的交错 $n$-重线性形式。由行列式函数唯一性的证明，必有
\[
D(B)=D(I)\det B.
\]
而 $D(I)=\det(AI)=\det A$，故 $\det(AB)=(\det A)(\det B)$。
\end{proof}

\noindent\textbf{定义.} 设 $A=(a_{ij})$ 是 $n\times n$ 矩阵。对每个 $i,j$，令 $A_{ij}$ 为从 $A$ 中删去第 $i$ 行与第 $j$ 列所得的 $(n-1)\times(n-1)$ 矩阵。元素
\[
(-1)^{i+j}\det(A_{ij})
\]
称为 $A$ 的 $(i,j)$-\emph{余因子}。

\noindent\textbf{定理 29（沿第 $i$ 行的余因子展开公式）.} 若 $A=(a_{ij})$ 是 $n\times n$ 矩阵，则对每个固定的 $i$，
\[
\det A=
(-1)^{i+1}a_{i1}\det A_{i1}+
(-1)^{i+2}a_{i2}\det A_{i2}+\cdots+
(-1)^{i+n}a_{in}\det A_{in}.
\]

\begin{proof}
由行列式公式，把和按置换 $\sigma$ 使得 $\sigma(j)=i$ 的那个列指标 $j$ 分组。固定 $j$ 后，含有因子 $a_{ij}$ 的所有项，其余因子正是子矩阵 $A_{ij}$ 的行列式展开项，并且把 $i$ 移到第 $j$ 个位置引入符号 $(-1)^{i+j}$。于是得到所述展开公式。
\end{proof}

\noindent\textbf{定理 30（逆矩阵的余因子公式）.} 令 $A=(a_{ij})$ 为 $n\times n$ 矩阵。设 $B$ 为 $A$ 的余因子矩阵的转置，即 $B=(\beta_{ij})$，其中
\[
\beta_{ij}=(-1)^{i+j}\det A_{ji}.
\]
则
\[
AB=BA=(\det A)I.
\]
此外，$\det A$ 是 $R$ 的单位当且仅当 $A$ 是 $M_n(R)$ 中的单位；此时
\[
A^{-1}=(\det A)^{-1}B.
\]

\begin{proof}
矩阵 $AB$ 的 $(i,j)$ 元为
\[
\sum_{k=1}^n a_{ik}\beta_{kj}
=
\sum_{k=1}^n a_{ik}(-1)^{k+j}\det A_{jk}.
\]
若 $i=j$，这正是沿第 $i$ 行的余因子展开，等于 $\det A$。若 $i\neq j$，它是把矩阵 $A$ 的第 $j$ 行替换为第 $i$ 行所得矩阵沿第 $j$ 行的余因子展开；该矩阵有两行相同，因此行列式为 $0$。故 $AB=(\det A)I$。同理使用列展开，或对转置应用上式，得 $BA=(\det A)I$。

若 $\det A$ 是 $R$ 的单位，则上式给出 $A$ 的逆为 $(\det A)^{-1}B$。反过来，若 $A$ 可逆，则 $I=AA^{-1}$，所以由定理 28，
\[
1=\det I=\det A\cdot \det(A^{-1}),
\]
从而 $\det A$ 是 $R$ 的单位。
\end{proof}

\paragraph{练习}
\begin{enumerate}[label=\arabic*.]
\item 陈述并证明方阵 $A$ 沿第 $j$ 列的余因子展开公式。
\item 设 $F$ 为域，$A_1,A_2,\ldots,A_n$ 是 $F^n$ 中列向量。令 $A$ 为以 $A_i$ 为第 $i$ 列的矩阵。证明这些向量构成 $F^n$ 的一组基，当且仅当 $\det A\neq0$。
\item 设 $R$ 为任意带 $1$ 的交换环，$V$ 为 $R$-模，且 $x_1,x_2,\ldots,x_n\in V$。假设对某个 $A=(a_{ij})\in M_n(R)$ 有
\[
\sum_{j=1}^n a_{ij}x_j=0,
\qquad i=1,2,\ldots,n.
\]
证明 $(\det A)x_i=0$，对所有 $i=1,2,\ldots,n$ 成立。
\item （矩阵行列式的计算）本题概述如何用 Gauss--Jordan 消元计算行列式。这是计算大行列式最有效的一般方法。设 $A$ 是 $n\times n$ 矩阵。
\begin{enumerate}[label=(\alph*)]
\item 证明初等行变换对行列式有如下影响：
\begin{enumerate}[label=(\roman*)]
\item 交换两行会改变行列式的符号；
\item 把一行的倍数加到另一行不改变行列式；
\item 把某一行乘以 $F$ 中非零元素 $u$，会把行列式乘以 $u$。
\end{enumerate}
\item 证明 $\det A\neq0$ 当且仅当 $A$ 与 $n\times n$ 恒等矩阵行等价。假设 $A$ 可经行化简化为恒等矩阵，过程中共使用 $s$ 次如 (i) 的行交换，并用非零元素 $u_1,u_2,\ldots,u_r$ 乘行，如 (iii) 所述。证明
\[
\det A=(-1)^s(u_1u_2\cdots u_r)^{-1}.
\]
\end{enumerate}
\item 用行化简计算下列矩阵的行列式：
\[
A=\begin{pmatrix}5&4&-6\\-2&0&2\\3&4&-2\end{pmatrix},
\qquad
B=\begin{pmatrix}1&2&-4&4\\2&-1&4&8\\1&0&1&-2\\0&1&-2&3\end{pmatrix}.
\]
\item （Minkowski 判别准则）设 $A$ 是实数项 $n\times n$ 矩阵，所有对角元为正，所有非对角元为负，并且每一行的元素和均为正。证明 $\det A\neq0$。\emph{[提示：考虑方程组 $AX=0$，并假设存在非平凡解 $(x_1,\ldots,x_n)$。若 $x_i$ 具有最大的绝对值，证明第 $i$ 个方程导致矛盾。]}
\end{enumerate}

\section{张量代数、对称代数与外代数}

本节中 $R$ 是带 $1$ 的交换环，并假设每个 $R$-模上的左右 $R$-作用相同。主要应用情形是 $R=F$ 为域。

设 $M$ 是 $R$-模。对 $k\geq1$，定义
\[
T^k(M)=\underbrace{M\otimes_R M\otimes_R\cdots\otimes_R M}_{k\text{ 个因子}},
\qquad
T^0(M)=R.
\]
$T^k(M)$ 的元素称为 $k$-张量。令
\[
T(M)=R\oplus T^1(M)\oplus T^2(M)\oplus\cdots
=\bigoplus_{k=0}^{\infty}T^k(M).
\]

\noindent\textbf{定理 31.} 若 $M$ 是交换环 $R$ 上任意 $R$-模，则：
\begin{enumerate}[label=(\arabic*)]
\item $T(M)$ 是一个包含 $M$ 的 $R$-代数，其乘法由
\[
(m_1\otimes\cdots\otimes m_i)(m_1'\otimes\cdots\otimes m_j')
=m_1\otimes\cdots\otimes m_i\otimes m_1'\otimes\cdots\otimes m_j'
\]
并再由分配律扩张定义。因此 $T^i(M)T^j(M)\subseteq T^{i+j}(M)$。
\item 若 $A$ 是任意 $R$-代数，$\varphi:M\to A$ 是 $R$-模同态，则存在唯一的 $R$-代数同态 $\Phi:T(M)\to A$，使得 $\Phi|_M=\varphi$。
\end{enumerate}

\begin{proof}
乘法的定义在单张量上给出，并由张量积的泛性质和双线性性延拓到直和。结合律来自张量积括号的自然结合性，单位元是 $T^0(M)=R$ 中的 $1$，所以得到一个 $R$-代数；$M=T^1(M)$ 自然包含其中。

为证明 (2)，设 $\varphi:M\to A$ 为 $R$-模同态。对每个 $k$，由张量积的泛性质，存在唯一 $R$-模同态
\[
\Phi_k:T^k(M)\to A,
\qquad
\Phi_k(m_1\otimes\cdots\otimes m_k)=\varphi(m_1)\cdots\varphi(m_k),
\]
且 $k=0$ 时取 $\Phi_0:R\to A$ 为结构同态。把这些 $\Phi_k$ 在直和上相加，得到唯一的 $R$-模同态 $\Phi:T(M)\to A$。由定义可见它保持乘法与单位，因而是 $R$-代数同态；唯一性由 $T(M)$ 中单张量由 $M$ 的元素乘积生成而得。
\end{proof}

\noindent\textbf{定义.} 环 $T(M)$ 称为 $M$ 的\emph{张量代数}。

\noindent\textbf{命题 32.} 设 $V$ 是域 $F$ 上有限维向量空间，基为 $\mathcal B=\{v_1,\ldots,v_n\}$。则 $k$-张量
\[
v_{i_1}\otimes v_{i_2}\otimes\cdots\otimes v_{i_k},
\qquad v_{i_j}\in\mathcal B,
\]
构成 $T^k(V)$ 在 $F$ 上的一组基（约定当 $k=0$ 时基向量为 $F$ 中元素 $1$）。特别地，
\[
\dim_F T^k(V)=n^k.
\]

\begin{proof}
这立即来自第 2 节命题 16。
\end{proof}

定理 31 与命题 32 说明，空间 $T(V)$ 可看作 $F$ 上非交换变量 $v_1,
\ldots,v_n$ 的非交换多项式代数（乘法即连接单词）。若 $V$ 是无限维的，也可用同样方式理解，只是变量集合由一组基指标化。

\paragraph{例}
\begin{enumerate}[label=(\arabic*)]
\item 令 $R=\mathbb Z$，$M=\mathbb Q/\mathbb Z$。由 10.4 节推论 12 后的例子，$(\mathbb Q/\mathbb Z)\otimes_{\mathbb Z}(\mathbb Q/\mathbb Z)=0$。因此
\[
T(\mathbb Q/\mathbb Z)=\mathbb Z\oplus(\mathbb Q/\mathbb Z),
\]
且乘法使两个正次数元素的乘积为 $0$。
\item 令 $R=\mathbb Z$，$M=\mathbb Z/n\mathbb Z$。由 10.4 节推论 12 后的例子，$(\mathbb Z/n\mathbb Z)\otimes_{\mathbb Z}(\mathbb Z/n\mathbb Z)\cong\mathbb Z/n\mathbb Z$。因此对所有 $i>0$，$T^i(M)\cong M$，所以
\[
T(M)=\mathbb Z\oplus M\oplus M\oplus\cdots.
\]
\end{enumerate}

\noindent\textbf{定义.}
\begin{enumerate}[label=(\arabic*)]
\item 环 $S$ 称为\emph{分次环}，若
\[
S=S_0\oplus S_1\oplus S_2\oplus\cdots
\]
并且 $S_iS_j\subseteq S_{i+j}$。$S_k$ 称为 $k$ 次齐次分量。
\item 分次环 $S$ 的理想 $I$ 称为\emph{分次理想}，若
\[
I=\bigoplus_{k\geq0}(I\cap S_k).
\]
\item 分次环之间的同态若保持次数，则称为分次环同态。
\end{enumerate}

\paragraph{例}
多项式环 $S=R[x_1,x_2,\ldots,x_n]$ 是交换环 $R$ 上的分次环。这里 $S_0=R$，$k$ 次齐次分量是所有次数为 $k$ 的单项式的 $R$-线性组合构成的子群。

由 $x_1,
\ldots,x_n$ 生成的理想 $I$ 是分次理想：每个常数项为零的多项式都可唯一写成次数 $k\geq1$ 的齐次多项式之和，而这些齐次分量都仍有零常数项，因而属于 $I$。更一般地，一个理想是分次理想当且仅当它可由齐次多项式生成（参见 9.1 节练习 17）。并不是分次环的每个理想都是分次理想。例如，在分次环 $\mathbb Z[x]$ 中，由 $1+x$ 生成的主理想 $J$ 不是分次理想：$1+x\in J$ 而 $1\notin J$，所以 $1+x$ 不能写成每个分量都属于 $J$ 的齐次多项式之和。

\noindent\textbf{命题 33.} 设 $S$ 是分次环，$I$ 是 $S$ 的分次理想，并令 $I_k=I\cap S_k$。则 $S/I$ 自然成为分次环，且其 $k$ 次齐次分量同构于
\[
S_k/I_k.
\]

\begin{proof}
因为 $I$ 是分次理想，每个陪集可唯一写成各齐次分量的陪集之和：
\[
S/I\cong\bigoplus_{k\geq0}S_k/I_k
\]
作为加法群。又由 $S_iS_j\subseteq S_{i+j}$ 以及 $I$ 是理想且分次，乘法满足
\[
(S_i/I_i)(S_j/I_j)\subseteq S_{i+j}/I_{i+j},
\]
并且良定义。因此 $S/I$ 是自然的分次环。
\end{proof}

\subsection*{对称代数}

命题 33 的第一个应用是构造张量代数的交换商。任何从 $M$ 到交换 $R$-代数的映射都必须把 $m_1m_2$ 与 $m_2m_1$ 识别，因而对应的 $R$-代数同态必须通过下面定义的商。这给出定理 31 的“交换化”版本。

\noindent\textbf{定义.} $R$-模 $M$ 的\emph{对称代数}是 $R$-代数 $T(M)$ 对由所有
\[
m_1\otimes m_2-m_2\otimes m_1,
\qquad m_1,m_2\in M
\]
生成的理想 $C(M)$ 取商所得的 $R$-代数，记为
\[
S(M)=T(M)/C(M).
\]
这是一个交换 $R$-代数。

理想 $C(M)$ 由次数为 $2$ 的齐次元素生成，所以它是分次理想。由命题 33，对称代数是分次环。其 $k$ 次齐次分量记为 $S^k(M)$，称为 $M$ 的第 $k$ 个\emph{对称幂}。由于生成元都在次数 $\geq2$ 中，$C(M)\cap M=0$，所以 $M=T^1(M)$ 在 $S(M)$ 中的像同构于 $M$；我们把 $M$ 与其像相同看待，因而 $S^1(M)=M$。类似地 $S^0(M)=R$，所以对称代数也是 $R$-代数。

下一定理的第一部分说明，$M$ 的第 $k$ 个对称幂中的元素可看作有限个单张量 $m_1\otimes\cdots\otimes m_k$ 的和，其中因子的顺序重排被识别。还回忆，$k$-重线性映射 $\varphi:M\times\cdots\times M\to N$ 称为对称的，若对 $1,
2,
\ldots,k$ 的任意置换 $\sigma$，都有
\[
\varphi(m_1,
\ldots,m_k)=\varphi(m_{\sigma(1)},
\ldots,m_{\sigma(k)}).
\]

\noindent\textbf{定理 34.} 设 $M$ 是交换环 $R$ 上的 $R$-模，$S(M)$ 是其对称代数。
\begin{enumerate}[label=(\arabic*)]
\item $S^k(M)$ 等于 $M\otimes\cdots\otimes M$（$k$ 个因子）模掉由所有
\[
m_1\otimes\cdots\otimes m_k
-
 m_{\sigma(1)}\otimes\cdots\otimes m_{\sigma(k)}
\]
生成的子模，其中 $\sigma\in S_k$。
\item 若 $\varphi:M\times\cdots\times M\to N$ 是对称 $k$-重线性映射，则存在唯一 $R$-模同态
\[
\Phi:S^k(M)\to N
\]
使得 $\varphi=\Phi\circ t$，其中
\[
t(m_1,
\ldots,m_k)=m_1\otimes\cdots\otimes m_k\pmod{C(M)}.
\]
\item 若 $A$ 是交换 $R$-代数，$\varphi:M\to A$ 是 $R$-模同态，则存在唯一 $R$-代数同态
\[
\Phi:S(M)\to A
\]
使得 $\Phi|_M=\varphi$。
\end{enumerate}

\begin{proof}
(1) 因为 $C(M)$ 由所有二次交换子生成，故在 $k$ 次分量中的部分由把相邻两个因子交换所产生的差生成。任意置换都可写成相邻换位的乘积，因此模掉这些差正好识别所有由因子重排得到的单张量。

(2) 由张量积的泛性质，对称 $k$-重线性映射 $\varphi$ 首先唯一诱导出 $T^k(M)$ 上的 $R$-模同态。由于 $\varphi$ 对因子置换不变，上述在 (1) 中被模掉的所有生成元都落入核中，因此该同态唯一地通过商 $S^k(M)$ 分解。

(3) 由张量代数的泛性质，$\varphi$ 唯一诱导 $R$-代数同态 $T(M)\to A$。由于 $A$ 交换，所有 $m_1\otimes m_2-m_2\otimes m_1$ 都被送到 $0$，所以 $C(M)$ 包含于核中。于是同态唯一地通过商 $S(M)$ 分解。唯一性来自 $S(M)$ 由 $M$ 作为 $R$-代数生成。
\end{proof}

\noindent\textbf{推论 35.} 若 $V$ 是域 $F$ 上 $n$ 维向量空间，则 $S(V)$ 作为分次 $F$-代数同构于 $F$ 上 $n$ 个变量的多项式环
\[
F[x_1,
\ldots,x_n].
\]
在此同构下，$S^k(V)$ 对应所有 $k$ 次齐次多项式构成的空间。特别地，
\[
\dim_F S^k(V)=\binom{n+k-1}{k}.
\]

\begin{proof}
令 $\mathcal B=\{v_1,
\ldots,v_n\}$ 为 $V$ 的一组基。由命题 32，$T^k(V)$ 的一组基对应于 $\mathcal B^k$ 中有序 $k$ 元组。定理 34(1) 把仅差一个置换的 $k$ 元组识别，因此 $S^k(V)$ 的基对应于从 $n$ 个基向量中允许重复地取 $k$ 个并不计顺序的方式，也就是次数为 $k$ 的单项式。对应 $v_i\mapsto x_i$，得到分次代数同构。次数为 $k$ 的单项式数为 $\binom{n+k-1}{k}$，故维数公式成立。
\end{proof}

\subsection*{外代数}

在交错映射中，任何含有两个相同参数的值必须为 $0$。因此与交错多重线性映射相配的商，是把所有 $m\otimes m$ 置为零所得的商。

\noindent\textbf{定义.} $R$-模 $M$ 的\emph{外代数}是 $R$-代数 $T(M)$ 对由所有
\[
m\otimes m,
\qquad m\in M
\]
生成的理想 $A(M)$ 取商所得的 $R$-代数，记为
\[
\bigwedge(M)=T(M)/A(M).
\]
元素 $m_1\otimes\cdots\otimes m_k$ 的像记为
\[
m_1\wedge m_2\wedge\cdots\wedge m_k.
\]
外代数中的乘法称为\emph{楔积}或\emph{外积}。外代数的 $k$ 次齐次分量
\[
\bigwedge^k(M)
\]
称为 $M$ 的第 $k$ 个\emph{外幂}。

由商的定义，乘法在如下意义下是交错的：若某个简单积 $m_1\wedge\cdots\wedge m_k$ 中 $m_i=m_{i+1}$，则它在 $\bigwedge(M)$ 中为 $0$。于是
\[
0=(m+m')\wedge(m+m')
=(m\wedge m)+(m\wedge m')+(m'\wedge m)+(m'\wedge m')
=(m\wedge m')+(m'\wedge m),
\]
所以在单张量上乘法也是反交换的：
\[
m\wedge m'=-m'\wedge m
\]
对所有 $m,m'\in M$ 成立。不过这种反交换性不推广到任意乘积；也就是说，一般不一定有 $ab=-ba$ 对所有 $a,b\in\bigwedge(M)$ 成立（参见练习 4）。

\noindent\textbf{定理 36.} 设 $M$ 是交换环 $R$ 上的 $R$-模，$\bigwedge(M)$ 是其外代数。
\begin{enumerate}[label=(\arabic*)]
\item $\bigwedge^k(M)$ 等于 $M\otimes\cdots\otimes M$（$k$ 个因子）模掉由所有含有两个相等因子的单张量生成的子模。特别地，若某两个 $m_i=m_j$ 且 $i\neq j$，则
\[
m_1\wedge\cdots\wedge m_k=0.
\]
\item 若 $\varphi:M\times\cdots\times M\to N$ 是交错 $k$-重线性映射，则存在唯一 $R$-模同态
\[
\Phi:\bigwedge^k(M)\to N
\]
使得 $\varphi=\Phi\circ t$，其中
\[
t(m_1,
\ldots,m_k)=m_1\wedge\cdots\wedge m_k.
\]
\end{enumerate}

\begin{proof}
理想 $A(M)$ 的 $k$ 次齐次部分由所有单张量生成，这些单张量中某两个相邻因子相等。由于在外代数中交换相邻两个因子会引入负号，通过一系列相邻交换与符号改变，可以把任意含有两个相等因子的简单 $k$-张量中的这两个相等因子移到相邻位置。因此由含任意两个相等因子的单张量生成的子模与 $A(M)$ 的 $k$ 次部分相同，这证明 (1)。

(2) 与定理 34 类似。由张量积的泛性质，交错 $k$-重线性映射先诱导 $T^k(M)$ 上的同态；由于任何含有两个相等因子的单张量都被交错映射送到 $0$，$A^k(M)$ 包含于核中，所以该同态唯一地通过 $\bigwedge^k(M)$ 分解。
\end{proof}

\noindent\textbf{注记.} 外代数也满足类似于定理 34(3) 的泛性质：若 $A$ 是一个 $R$-代数，并且 $a^2=0$ 对所有来自 $M$ 的像成立，则从 $M$ 到 $A$ 的 $R$-模同态唯一延拓为从 $\bigwedge(M)$ 到 $A$ 的 $R$-代数同态（见练习）。

\paragraph{例}
\begin{enumerate}[label=(\arabic*)]
\item 若 $V$ 是 $F$ 上一维向量空间，基元素为 $v$，则 $\bigwedge^k(V)$ 由形如
\[
a_1v\wedge a_2v\wedge\cdots\wedge a_kv=a_1a_2\cdots a_k(v\wedge v\wedge\cdots\wedge v)
\]
的有限和组成。因为 $v\wedge v=0$，所以
\[
\bigwedge^0(V)=F,
\qquad
\bigwedge^1(V)=V,
\qquad
\bigwedge^k(V)=0\quad(k\geq2).
\]
\item 若 $V$ 是二维向量空间，基为 $v,v'$，则
\[
\bigwedge(V)=F\oplus V\oplus F(v\wedge v')\oplus0\oplus0\oplus\cdots.
\]
事实上，$\bigwedge^2(V)$ 由 $v\wedge v'$ 张成，而所有更高外幂为零。
\end{enumerate}

\noindent\textbf{推论 37.} 设 $V$ 是域 $F$ 上有限维向量空间，基为 $\mathcal B=\{v_1,
\ldots,v_n\}$。则向量
\[
v_{i_1}\wedge v_{i_2}\wedge\cdots\wedge v_{i_k},
\qquad
1\leq i_1<i_2<\cdots<i_k\leq n,
\]
构成 $\bigwedge^k(V)$ 的一组基；当 $k>n$ 时，$\bigwedge^k(V)=0$（当 $k=0$ 时，基向量为 $F$ 中元素 $1$）。特别地，
\[
\dim_F\bigwedge^k(V)=\binom nk.
\]

\begin{proof}
如定理 36 的证明所示，模掉 $A^k(V)$ 后，任何简单 $k$-张量中的项的顺序都可在引入符号改变的情况下重新排列。因此，推论中所列的 $k$-张量（下标递增且没有重复项）生成 $\bigwedge^k(V)$。为证明这些向量线性无关，只需对每个给定的 $v_{i_1}\wedge\cdots\wedge v_{i_k}$ 构造一个交错 $k$-重线性函数 $V^k\to F$，它在该生成元上取值 $1$，在所有其他生成元上取值 $0$。这样的函数可在命题 32 给出的 $T^k(V)$ 的基上定义：若输入是 $v_{i_1},
\ldots,v_{i_k}$ 的同一集合的某个排列，则取相应置换的符号；否则取 $0$。它是交错多重线性函数，由定理 36 的泛性质给出所需线性泛函。这证明线性无关。若 $k>n$，任何 $k$ 个基向量中必有重复，因此外积为零。
\end{proof}

推论 37 的结果对任意秩为 $n$ 的自由 $R$-模也成立。特别地，若 $M$ 是自由 $R$-模，则 $\bigwedge^k(M)$ 是秩为 $\binom nk$ 的自由 $R$-模。

\paragraph{例}
令 $R=\mathbb Z[x,y]$，$I=(x,y)$。虽然 $\bigwedge^2(R)=0$，但 $\bigwedge^2(I)\neq0$。定义
\[
\varphi(ax+by,cx+dy)=ad-bc\pmod{(x,y)}
\]
从 $I\times I$ 到 $R/I\cong\mathbb Z$。这个映射是良定义的交错 $R$-双线性映射，并且 $\varphi(x,y)=1$。由 $\bigwedge^2(I)$ 关于这类映射的泛性质，可知 $x\wedge y\neq0$。这个例子说明，当 $R$-模不是自由模时，推论 37 的结果不能无条件推广。

\noindent\textbf{注记.} 理想 $I$ 是一个秩为 $1$（但非自由）的 $R$-模。一个模的秩是指把它张量到分式域后得到的向量空间的维数。这个例子表明，即使 $I$ 的秩为 $1$，其第二外幂也可以非零。

\subsection*{张量代数的同态}

若 $\varphi:M\to N$ 是任意 $R$-模同态，则在第 $k$ 个张量幂上诱导映射
\[
T^k(\varphi):T^k(M)\to T^k(N),
\qquad
m_1\otimes\cdots\otimes m_k
\mapsto
\varphi(m_1)\otimes\cdots\otimes\varphi(m_k).
\]
直接可见，此映射把理想 $C(M)$ 与 $A(M)$ 的每个齐次分量的生成元送到相应理想的生成元。因此 $\varphi$ 在商上诱导 $R$-模同态
\[
S^k(\varphi):S^k(M)\to S^k(N),
\qquad
\bigwedge^k(\varphi):\bigwedge^k(M)\to\bigwedge^k(N).
\]
此外，这三个映射都是环同态（因而是分次 $R$-代数同态）。

特别重要的情形是 $M=V$ 为域 $F$ 上 $n$ 维向量空间，且 $\varphi:V\to V$ 是自同态。此时由推论 37，$\bigwedge^n(\varphi)$ 把一维空间 $\bigwedge^n(V)$ 映到自身。令 $v_1,
\ldots,v_n$ 是 $V$ 的一组基。则 $\bigwedge^n(V)$ 的基为 $v_1\wedge\cdots\wedge v_n$，因此存在标量 $D(\varphi)$，使得
\[
\bigwedge^n(\varphi)(v_1\wedge\cdots\wedge v_n)
=
\varphi(v_1)\wedge\cdots\wedge\varphi(v_n)
=
D(\varphi)\,v_1\wedge\cdots\wedge v_n.
\]
将 $\varphi(v_i)$ 以基 $v_j$ 展开，左边的系数作为 $\varphi$ 矩阵列向量的交错多重线性函数，满足第 4 节中行列式函数的两个公理。于是由定理 24 的唯一性，得到以下刻画。

\noindent\textbf{命题 38.} 若 $\varphi$ 是 $n$ 维向量空间 $V$ 上的自同态，则
\[
\bigwedge^n(\varphi)(w)=(\det\varphi)w,
\qquad w\in\bigwedge^n(V).
\]

注意，命题 38 把自同态 $\varphi$ 的行列式刻画为其在最高外幂上的作用标量。由于它以基无关方式定义，行列式在换基下不变。

若 $M$ 是 $R$-模 $N$ 的直和因子，则可证明 $T(M)$（分别地，$S(M)$ 与 $\bigwedge(M)$）是 $T(N)$（分别地，$S(N)$ 与 $\bigwedge(N)$）的 $R$-子代数（参见练习）。当 $R=F$ 是域时，$N$ 的每个子空间 $M$ 都是 $N$ 的直和因子，因此 $M$ 的相应代数是 $N$ 的相应代数的子代数。

\subsection*{对称张量与交错张量}

对称代数与外代数在某些情形下也可用对称张量与交错张量（定义如下）来定义，这会把这些代数识别为张量代数的子代数，而不是商代数。

对任意 $R$-模 $M$，对称群 $S_k$ 通过置换因子自然左作用在 $M\times\cdots\times M$ 上，从而也作用在 $T^k(M)$ 上：
\[
\sigma(m_1\otimes\cdots\otimes m_k)
=m_{\sigma^{-1}(1)}\otimes\cdots\otimes m_{\sigma^{-1}(k)}.
\]
使用 $\sigma^{-1}$ 是为了得到左群作用。

\noindent\textbf{定义.}
\begin{enumerate}[label=(\arabic*)]
\item 若 $z\in T^k(M)$ 满足 $\sigma z=z$ 对所有 $\sigma\in S_k$ 成立，则称 $z$ 是\emph{对称 $k$-张量}。
\item 若 $z\in T^k(M)$ 满足 $\sigma z=\epsilon(\sigma)z$ 对所有 $\sigma\in S_k$ 成立，则称 $z$ 是\emph{交错 $k$-张量}。
\end{enumerate}

\paragraph{例}
若 $m,n\in M$，则 $m\otimes n+n\otimes m$ 是对称 $2$-张量，而 $m\otimes n-n\otimes m$ 是交错 $2$-张量（当 $2$ 在 $R$ 中不是零因子时，这与通常直观一致）。

\noindent\textbf{命题 39.} 设 $\sigma\in S_k$，$\epsilon(\sigma)$ 为其符号。则：
\begin{enumerate}[label=(\arabic*)]
\item 对任意 $w\in S^k(M)$，有 $\sigma w=w$；
\item 对任意 $w\in\bigwedge^k(M)$，有 $\sigma w=\epsilon(\sigma)w$。
\end{enumerate}

\begin{proof}
第一条立即来自定理 34(1)。在定理 36 的证明中我们已经看到，在 $\bigwedge^k(M)$ 中，交换任意两个相邻因子会使简单积变号。由于相邻换位生成 $S_k$，且 $\epsilon$ 是群同态，故第二条对简单积成立；再由 $R$-线性性，对所有 $w\in\bigwedge^k(M)$ 成立。
\end{proof}

由命题 39，对称群 $S_k$ 在对称 $k$-张量的子模与商模 $S^k(M)$ 上都平凡作用；类似地，$S_k$ 在交错 $k$-张量的子模上与在 $\bigwedge^k(M)$ 上以同样方式作用。现在我们说明，当 $k!$ 是 $R$ 中的单位时，这些相应的子模与商模同构（这里 $k!$ 是 $R$ 中的 $1$ 自加 $k!$ 次）。

对任意 $k$-张量 $z\in T^k(M)$，定义
\[
\operatorname{Sym}(z)=\sum_{\sigma\in S_k}\sigma z,
\qquad
\operatorname{Alt}(z)=\sum_{\sigma\in S_k}\epsilon(\sigma)\sigma z.
\]
它们分别称为 $z$ 的\emph{对称化}与\emph{斜对称化}。

\noindent\textbf{命题 40.} 假设 $k!$ 是环 $R$ 中的单位，$M$ 是 $R$-模。则：
\begin{enumerate}[label=(\arabic*)]
\item 映射 $(1/k!)\operatorname{Sym}$ 诱导 $R$-模同构
\[
\frac1{k!}\operatorname{Sym}:S^k(M)\xrightarrow{\sim}\{\text{对称 }k\text{-张量}\};
\]
\item 映射 $(1/k!)\operatorname{Alt}$ 诱导 $R$-模同构
\[
\frac1{k!}\operatorname{Alt}:\bigwedge^k(M)\xrightarrow{\sim}\{\text{交错 }k\text{-张量}\}.
\]
\end{enumerate}

\begin{proof}
相应映射都是满的 $R$-同态：若 $z$ 是对称张量，则 $\operatorname{Sym}(z)=k!z$，故 $(1/k!)\operatorname{Sym}(z)=z$；若 $z$ 是交错张量，则 $\operatorname{Alt}(z)=k!z$，故 $(1/k!)\operatorname{Alt}(z)=z$。

对称情形中，$C^k(M)$ 中的生成元是交换两个因子之差，而对这样的差，对称化为 $0$，故 $(1/k!)\operatorname{Sym}$ 通过商 $S^k(M)$。若 $z$ 被送到 $0$，则 $\operatorname{Sym}(z)=0$，并且
\[
z-\frac1{k!}\operatorname{Sym}(z)
=\frac1{k!}\sum_{\sigma\in S_k}(z-\sigma z)
\]
属于 $C^k(M)$，所以 $z$ 在 $S^k(M)$ 中为 $0$。故映射单。

交错情形类似。$A^k(M)$ 中被模掉的生成元在斜对称化下为 $0$，所以映射通过商 $\bigwedge^k(M)$。若 $\operatorname{Alt}(z)=0$，则
\[
z-\frac1{k!}\operatorname{Alt}(z)
=\frac1{k!}\sum_{\sigma\in S_k}(z-
\epsilon(\sigma)\sigma z)
\]
属于 $A^k(M)$（这一事实也在练习中要求验证），因而 $z$ 在外幂中为 $0$。故映射单。
\end{proof}

若 $k!$ 在 $R$ 中不可逆，则一般不能把第 $k$ 个外幂直接等同于交错 $k$-张量空间。这个差异在正特征或一般交换环上尤其重要。文献中（尤其是在 $\mathbb R$ 与 $\mathbb C$ 等域上发展理论时）常常隐含 $k!$ 可逆的情形。在这种情形下，两个交错张量 $z,w$ 的乘法定义为先在 $T(M)$ 中取积 $zw=z\otimes w$，再把所得张量投影到交错张量子模中。注意，两个交错张量的简单乘积不必交错（例如，一个交错张量的平方是对称张量）。

\paragraph{例}
设 $V$ 是域 $F$ 上向量空间，且 $k!\neq0$ 于 $F$ 中。$T^k(V)$ 中 $A^k(V)$ 的向量空间补空间有很多（例如把 $A^k(V)$ 的一组基扩充为 $T^k(V)$ 的一组基）。这些补空间依赖于 $T^k(V)$ 的基选择，因此仅从向量空间角度看彼此无法区分。$T^k(V)$ 上由 $S_k$ 的作用给出的附加结构，挑出了 $A^k(V)$ 的一个唯一补空间，即命题 40 中的交错张量子空间。

\paragraph{练习}
在这些练习中，$R$ 是带 $1$ 的交换环，$M$ 是 $R$-模；$F$ 是域，$V$ 是 $F$ 上有限维向量空间。
\begin{enumerate}[label=\arabic*.]
\item 证明若 $M$ 是循环 $R$-模，则 $T(M)=S(M)$，即张量代数 $T(M)$ 是交换的。
\item 补全命题 33 的证明细节：
\[
S/I\cong\bigoplus_{k\geq0}S_k/I_k.
\]
先证明 $I_iI_j\subseteq I_{i+j}$；用此证明乘法
\[
(S_i/I_i)(S_j/I_j)\subseteq S_{i+j}/I_{i+j}
\]
良定义，然后检查环公理，并验证命题 33 证明中的陈述。
\item 证明对 $\mathbb Z$-模 $\mathbb Z$，映射 $\operatorname{Sym}_2$ 的像由形如 $a(1\otimes1)$ 的 $2$-张量组成，其中 $a$ 为偶整数。特别地，推出对称张量 $1\otimes1\in\mathbb Z\otimes_{\mathbb Z}\mathbb Z$ 不包含在映射 $\operatorname{Sym}$ 的像中。
\item 证明
\[
m\wedge n_1\wedge n_2\wedge\cdots\wedge n_k
=(-1)^k(n_1\wedge n_2\wedge\cdots\wedge n_k\wedge m).
\]
特别地，对所有 $x,y,z\in M$，有
\[
x\wedge(y\wedge z)=(y\wedge z)\wedge x.
\]
\item 证明若 $M$ 是秩为 $n$ 的自由 $R$-模，则 $\bigwedge^i(M)$ 是秩为 $\binom ni$ 的自由 $R$-模，$i=0,1,2,\ldots$。
\item 若 $A$ 是任意满足 $a^2=0$（对所有 $a\in A$）的 $R$-代数，且 $\varphi:M\to A$ 是 $R$-模同态，证明存在唯一的 $R$-代数同态
\[
\Phi:\bigwedge(M)\to A
\]
使得 $\Phi|_M=\varphi$。
\item 令 $R=\mathbb Z[x,y]$，$I=(x,y)$。
\begin{enumerate}[label=(\alph*)]
\item 证明若 $ax+by=a'x+b'y$ 于 $R$ 中成立，则对某个多项式 $f(x,y)\in R$，有
\[
a'=a+yf,
\qquad
b'=b-xf.
\]
\item 证明推论 37 后例子中的映射
\[
\varphi(ax+by,cx+dy)=ad-bc\pmod{(x,y)}
\]
是良定义的交错 $R$-双线性映射 $I\times I\to\mathbb Z=R/I$。
\end{enumerate}
\item 设 $R$ 是整环，$F$ 是其分式域。
\begin{enumerate}[label=(\alph*)]
\item 把 $F$ 看作 $R$-模，证明 $\bigwedge^2F=0$。
\item 设 $I$ 是 $F$ 的任意 $R$-子模（例如 $R$ 中任意理想）。证明当 $i\geq2$ 时，$\bigwedge^i I$ 是挠 $R$-模，即对每个 $x\in\bigwedge^i I$，存在非零 $r\in R$ 使得 $rx=0$。
\item 给出一个整环 $R$ 以及 $F$ 中的一个 $R$-模 $I$，使得对每个 $i\geq0$ 都有 $\bigwedge^i I\neq0$。\emph{[参见推论 37 后的例子。]}
\end{enumerate}
\item 令 $R=\mathbb Z[G]$ 为阶为 $2$ 的群 $G=\{1,\sigma\}$ 的群环。令 $M=\mathbb Ze_1+\mathbb Ze_2$ 为秩为 $2$ 的自由 $\mathbb Z$-模，基为 $e_1,e_2$。定义
\[
\sigma(e_1)=e_1+2e_2,
\qquad
\sigma(e_2)=-e_2.
\]
证明这使 $M$ 成为 $R$-模，并证明 $R$-模 $\bigwedge^2M$ 是一个 $2$ 阶群，生成元为 $e_1\wedge e_2$。
\item 对任意 $k$-张量 $z$，证明
\[
z-\frac1{k!}\operatorname{Alt}(z)
=
\frac1{k!}\sum_{\sigma\in S_k}\bigl(z-\epsilon(\sigma)\sigma z\bigr),
\]
并用它证明命题 40 中 $R$-模同态 $(1/k!)\operatorname{Alt}$ 的核是 $A^k(M)$。
\item 证明 $\operatorname{Alt}_k$ 的像是 $T^k(V)$ 中最大的子空间，并且在该子空间上，对称群 $S_k$ 中每个置换 $\sigma$ 都按标量 $\epsilon(\sigma)$ 作用。
\item 
\begin{enumerate}[label=(\alph*)]
\item 证明若 $f(x,y)$ 是 $V$ 上交错双线性映射，即对所有 $x\in V$ 有 $f(x,x)=0$，则对所有 $x,y\in V$ 有
\[
f(x,y)=-f(y,x).
\]
\item 假设 $F$ 中 $-1\neq1$。证明 $f(x,y)$ 是 $V$ 上交错双线性映射，当且仅当对所有 $x,y\in V$，有 $f(x,y)=-f(y,x)$。
\item 假设 $F$ 中 $-1=1$。证明 $V$ 上每个交错双线性形式 $f(x,y)$ 都是对称的，即对所有 $x,y\in V$ 有 $f(x,y)=f(y,x)$。证明存在一个对称双线性映射不是交错的。\emph{[一种方法：通过维数计数证明 $C^2(V)\subset A^2(V)$ 且 $C^2(V)\neq A^2(V)$。或者构造一个显式的对称但非交错的映射。]}
\end{enumerate}
\item 设 $F$ 为任意满足 $-1\neq1$ 的域，$V$ 为 $F$ 上向量空间。证明
\[
V\otimes_F V=S^2(V)\oplus\bigwedge^2(V),
\]
即每个 $2$-张量都可唯一写成一个对称张量与一个交错张量之和。
\item 证明若 $M$ 是 $R$-模 $N$ 的直和因子，则 $T(M)$（分别地，$S(M)$ 与 $\bigwedge(M)$）是 $T(N)$（分别地，$S(N)$ 与 $\bigwedge(N)$）的 $R$-子代数。
\end{enumerate}

% 
% ============================================
% chapter12_full_zh.tex — 第12章：主理想整环上的模
% 按原书第12章 PDF 逐节重译校订；正文、证明、例子与练习均保留
% ============================================

\chapter{主理想整环上的模}

本章的主要目的是证明有限生成模在一类特别好的环上，即在主理想整环（Principal Ideal Domain, P.I.D.）上的结构定理。这个定理是环的理想结构（对 P.I.D. 来说特别简单）反映在它的模结构中的一个例子。如果把这个结果应用于 P.I.D. 为整数环 $\mathbb Z$ 的情形，那么就得到有限生成阿贝尔群基本定理的一个证明（我们在第 5 章已经不加证明地考察过它）。如果把这个结构定理应用于 P.I.D. 为域 $F$ 上关于 $x$ 的多项式环 $F[x]$ 的情形，那么就得到所谓矩阵有理标准形与 Jordan 标准形的基本结果。在进入证明以前，我们先简要讨论这两个重要应用。

我们在第 5 章已经讨论过如下结果：任意有限生成阿贝尔群同构于若干循环阿贝尔群的直和，而这些循环阿贝尔群要么是 $\mathbb Z$，要么是某个正整数 $n\neq 0$ 对应的 $\mathbb Z/n\mathbb Z$。还要回忆，阿贝尔群与 $\mathbb Z$-模是同一件事。由于 $\mathbb Z$ 的理想恰好是平凡理想 $(0)$ 以及由正整数 $n$ 生成的主理想 $(n)=n\mathbb Z$，我们看到，用模的语言来说，有限生成阿贝尔群基本定理断言任意有限生成 $\mathbb Z$-模都是形如 $\mathbb Z/I$ 的模（其中 $I$ 是 $\mathbb Z$ 的理想；这些正是循环 $\mathbb Z$-模）的直和，并且当直和写成一种特定形式时还带有唯一性陈述。注意 $\mathbb Z$ 的理想结构与它的有限生成模（也就是有限生成阿贝尔群）的结构之间的对应关系。

P.I.D. 上有限生成模基本定理说，当主理想整环 $\mathbb Z$ 被任意 P.I.D. 取代时，同样的结果仍成立。特别地，我们已经在第 10 章看到，关于域 $F$ 上多项式环 $F[x]$ 的模与一个 $F$-向量空间 $V$ 加上一个固定的线性变换 $T$ 是同一件事（其中元素 $x$ 在 $V$ 上按线性变换 $T$ 作用）。在这个情形中，基本定理将说明这样的向量空间是形如 $F[x]/I$ 的模的直和，其中 $I$ 是 $F[x]$ 的理想，因而或者是平凡理想 $(0)$，或者是由某个非零多项式 $f(x)$ 生成的主理想 $(f(x))$（这些就是循环 $F[x]$-模）；同样地，当直和写成一种特定形式时还带有唯一性陈述。如果把这翻译回向量空间和线性变换的语言，我们便能得到关于线性变换 $T$ 的信息。

例如，设 $V$ 是 $F$ 上 $n$ 维向量空间并且我们为 $V$ 选定一组基。那么给出 $V$ 到自身的线性变换 $T$，等价于给出一个系数在 $F$ 中的 $n\times n$ 矩阵 $A$（如果为 $V$ 选取另一组基，则得到 $T$ 的另一个矩阵 $B$，它与 $A$ 相似，即 $B$ 具有 $P^{-1}AP$ 的形式，其中 $P$ 是定义换基的可逆矩阵）。我们将看到，在这种情形下，基本定理蕴含（假设域 $F$ 含有给定线性变换 $T$ 的所有“特征值”）存在 $V$ 的一组基，使得 $T$ 的相应矩阵尽可能接近对角矩阵，因而有一个特别简单的形式。这就是 Jordan 标准形。有理标准形是 $T$ 的矩阵的另一种简单形式（它不要求 $T$ 的特征值属于 $F$）。这样，我们就能给出域 $F$ 上任意 $n\times n$ 矩阵的标准形，即找到与给定 $n\times n$ 矩阵相似而又特别简单的矩阵（例如几乎对角）。

\paragraph{例}
令
\[
V=\mathbb Q^3=\{(x,y,z)\mid x,y,z\in\mathbb Q\}
\]
为有理数域 $F=\mathbb Q$ 上通常的三维有序三元组向量空间，并设 $T$ 是线性变换
\[
T(x,y,z)=(9x+4y+5z,-4x-3z,-6x-4y-2z),\qquad x,y,z\in\mathbb Q.
\]
如果取 $V$ 的标准基 $e_1=(1,0,0)$，$e_2=(0,1,0)$，$e_3=(0,0,1)$，则表示这个线性变换的矩阵为
\[
A=\begin{pmatrix}
9&4&5\\
-4&0&-3\\
-6&-4&-2
\end{pmatrix}.
\]
我们将看到，这个矩阵 $A$ 的 Jordan 标准形是更简单的矩阵
\[
B=\begin{pmatrix}
2&1&0\\
0&2&0\\
0&0&3
\end{pmatrix},
\]
它通过改取基 $f_1=(2,-1,-2)$，$f_2=(1,0,-1)$，$f_3=(3,-2,-2)$ 得到，因为在这种情况下
\[
\begin{aligned}
T(f_1)&=T(2,-1,-2)=(4,-2,-4)=2f_1+0f_2+0f_3,\\
T(f_2)&=T(1,0,-1)=(4,-1,-4)=1f_1+2f_2+0f_3,\\
T(f_3)&=T(3,-2,-2)=(9,-6,-6)=0f_1+0f_2+3f_3.
\end{aligned}
\]
所以相对于这组基表示 $T$ 的矩阵的列分别是 $(2,0,0)$、$(1,2,0)$ 与 $(0,0,3)$，也就是说 $T$ 相对于这组基的矩阵是 $B$。特别地，$A$ 与更简单的矩阵 $B$ 相似。事实上这个线性变换 $T$ 不能对角化（即不存在一组 $V$ 的基使相应矩阵为对角矩阵），所以矩阵 $B$ 已经是 $T$ 可能达到的最接近对角矩阵的形式。

下面第一节给出若干一般定义，并陈述和证明任意 P.I.D. 上的基本定理；然后再回到标准形的应用（对阿贝尔群的应用出现在第 5 章）。这些应用可以独立于一般证明来阅读。对 Euclidean 整环有效的一种替代证明（因而特别适用于环 $\mathbb Z$ 和 $F[x]$），按行、列运算的方式给出，并且在计算上有用；它在练习中加以概述。

\section{基本理论}

我们首先描述若干一般的有限性条件。令 $R$ 为环，$M$ 为左 $R$-模。

\noindent\textbf{定义.}
\begin{enumerate}[label=(\arabic*)]
\item 左 $R$-模 $M$ 称为 \emph{Noether $R$-模}，或称满足子模上的升链条件（ascending chain condition on submodules，简称 A.C.C. on submodules），如果不存在子模的无限严格递增链；也就是说，每当
\[
M_1\subseteq M_2\subseteq M_3\subseteq\cdots
\]
是 $M$ 的子模升链时，存在正整数 $m$，使得对所有 $k\geq m$ 都有 $M_k=M_m$（因此这条链在第 $m$ 阶段稳定：$M_m=M_{m+1}=M_{m+2}=\cdots$）。
\item 环 $R$ 称为 \emph{Noether 环}，如果 $R$ 作为它自身上的左模是 Noether 的；也就是说，如果 $R$ 中不存在左理想的无限严格递增链。
\end{enumerate}

在一个（可能非交换的）环 $R$ 中，也可以类似地表述右理想和双边理想上的 A.C.C. 概念。对非交换环而言，这些性质不必相互相关。

\noindent\textbf{定理 1.} 设 $R$ 为环，$M$ 为左 $R$-模。下列条件等价：
\begin{enumerate}[label=(\arabic*)]
\item $M$ 是 Noether $R$-模；
\item $M$ 的每个非空子模集合在包含关系下含有极大元；
\item $M$ 的每个子模都是有限生成的。
\end{enumerate}

\begin{proof}
$[(1)\Rightarrow(2)]$ 假设 $M$ 是 Noether 的，并令 $\Sigma$ 是 $M$ 的任意非空子模族。任选 $M_1\in\Sigma$。如果 $M_1$ 已是 $\Sigma$ 中的极大元，则 (2) 成立；于是设 $M_1$ 不是极大元。则存在某个 $M_2\in\Sigma$ 使 $M_1\subsetneq M_2$。如果 $M_2$ 是 $\Sigma$ 中的极大元，则 (2) 成立；否则又存在 $M_3\in\Sigma$ 严格包含 $M_2$。如此继续下去可见，若 (2) 失败，则可借助选择公理构造出 $\Sigma$ 中元素的无限严格递增链
\[
M_1\subsetneq M_2\subsetneq M_3\subsetneq\cdots,
\]
这与 (1) 矛盾。

$[(2)\Rightarrow(3)]$ 假设 (2) 成立，令 $N$ 是 $M$ 的任意子模。令 $\Sigma$ 为 $N$ 的所有有限生成子模的集合。由于 $\{0\}\in\Sigma$，这个集合非空。由 (2)，$\Sigma$ 含有极大元 $N'$。若 $N'\neq N$，取 $x\in N-N'$。由于 $N'\in\Sigma$，子模 $N'$ 按假设有限生成，因此由 $N'$ 与 $x$ 生成的子模也有限生成，并且严格包含 $N'$。这与 $N'$ 的极大性矛盾。因此 $N=N'$ 有限生成。

$[(3)\Rightarrow(1)]$ 假设 (3) 成立，并令
\[
M_1\subseteq M_2\subseteq M_3\subseteq\cdots
\]
是 $M$ 的子模链。令
\[
N=\bigcup_{i=1}^{\infty}M_i,
\]
注意 $N$ 是一个子模。由 (3)，$N$ 有有限个生成元，设为 $a_1,a_2,\ldots,a_n$。因为每个 $a_i\in N$，每个 $a_i$ 都属于这条链中的某个子模，设 $a_i\in M_{j_i}$。令
\[
m=\max\{j_1,j_2,\ldots,j_n\}.
\]
则所有 $a_i$ 都属于 $M_m$，所以它们生成的模包含在 $M_m$ 中，即 $N\subseteq M_m$。于是 $M_m=N=M_k$ 对所有 $k\geq m$ 成立，这证明了 (1)。
\end{proof}

\noindent\textbf{推论 2.} 若 $R$ 是 P.I.D.，则 $R$ 的每个非空理想集合有极大元，并且 $R$ 是 Noether 环。

\begin{proof}
P.I.D. $R$ 在定理 1 中对 $M=R$ 满足条件 (3)，因为 $R$ 的每个理想都是主理想，因而由一个元素生成。
\end{proof}

回忆：即使 $M$ 本身是有限生成 $R$-模，$M$ 的子模也未必有限生成，所以 $M$ 是 Noether $R$-模这个条件通常强于 $M$ 是有限生成 $R$-模这个条件。

在转向本章主要结果之前，我们需要一个关于“线性相关”的结果。

\noindent\textbf{命题 3.} 设 $R$ 是整环，$M$ 是秩 $n<\infty$ 的自由 $R$-模。则 $M$ 中任意 $n+1$ 个元素都是 $R$-线性相关的；也就是说，对任意 $y_1,y_2,\ldots,y_{n+1}\in M$，存在不全为零的 $r_1,r_2,\ldots,r_{n+1}\in R$，使得
\[
r_1y_1+r_2y_2+\cdots+r_{n+1}y_{n+1}=0.
\]

\begin{proof}
最快的证明是把 $R$ 嵌入它的分式域 $F$（因为 $R$ 是整环），并观察到：由于
\[
M\cong R\oplus R\oplus\cdots\oplus R\quad(n\text{ 次}),
\]
我们得到
\[
F\otimes_R M\cong F\oplus F\oplus\cdots\oplus F.
\]
后者是 $F$ 上的 $n$ 维向量空间，所以 $M$ 中任意 $n+1$ 个元素在张量到 $F$ 后都是 $F$-线性相关的。把这些标量的分母清除（例如同乘所有分母的乘积）便得到这 $n+1$ 个元素之间的一个 $R$-线性相关关系。

也可如下证明。令 $e_1,\ldots,e_n$ 是自由 $R$-模 $M$ 的一组基，并令 $y_1,\ldots,y_{n+1}$ 是 $M$ 的任意 $n+1$ 个元素。对 $1\leq i\leq n+1$，把
\[
y_i=a_{1i}e_1+a_{2i}e_2+\cdots+a_{ni}e_n
\]
写在基 $e_1,e_2,\ldots,e_n$ 下。令 $A$ 为 $(n+1)\times(n+1)$ 矩阵，它的第 $i,j$ 个条目为 $a_{ij}$（$1\leq i\leq n$，$1\leq j\leq n+1$），最后一行全为 $0$，所以当然 $\det A=0$。由于 $R$ 是整环，第 11.4 节推论 27 说明 $A$ 的列 $R$-线性相关。列之间的任一线性相关关系正给出 $y_i$ 之间的相应线性相关关系，从而完成证明。
\end{proof}

若 $R$ 是任意整环，$M$ 是任意 $R$-模，回忆
\[
\operatorname{Tor}(M)=\{x\in M\mid rx=0\text{ 对某个非零 }r\in R\}
\]
是 $M$ 的子模（称为 $M$ 的\emph{挠子模}）。如果 $N$ 是 $\operatorname{Tor}(M)$ 的任意子模，则称 $N$ 为 $M$ 的一个挠子模（因此 $\operatorname{Tor}(M)$ 是 $M$ 的所有挠子模的并，即 $M$ 的极大挠子模）。如果 $\operatorname{Tor}(M)=0$，则称 $M$ 是\emph{无挠}的。

对 $M$ 的任意子模 $N$，$N$ 的零化子是 $R$ 的理想
\[
\operatorname{Ann}(N)=\{r\in R\mid rn=0\text{ 对所有 }n\in N\}.
\]
注意，如果 $N$ 不是 $M$ 的挠子模，则 $\operatorname{Ann}(N)=(0)$。容易看出，如果 $N,L$ 是 $M$ 的子模且 $N\subseteq L$，则 $\operatorname{Ann}(L)\subseteq \operatorname{Ann}(N)$。若 $R$ 是 P.I.D. 且 $N\subseteq L\subseteq M$，并且 $\operatorname{Ann}(N)=(a)$、$\operatorname{Ann}(L)=(b)$，则 $a\mid b$。特别地，$M$ 的任一元素 $x$ 的零化子整除 $M$ 的零化子（当 $R=\mathbb Z$ 时，这由 Lagrange 定理蕴含）。

\noindent\textbf{定义.} 对任意整环 $R$，$R$-模 $M$ 的\emph{秩}是 $M$ 中 $R$-线性无关元素个数的最大值。

前一命题说明：在整环上，一个有限秩自由 $R$-模 $M$ 的子模的秩被 $M$ 的秩所界定。这个秩的概念与此前同名术语一致。如果环 $R=F$ 是域，则 $R$-模 $M$ 的秩就是 $M$ 作为 $F$-向量空间的维数，并且任意极大 $F$-线性无关集都是 $M$ 的一组基。然而，对一般整环，秩为 $n$ 的 $R$-模 $M$ 不必有“基”，也就是说，即使 $M$ 无挠，它也不必是自由 $R$-模。因此使用秩的概念时需要谨慎，尤其要注意 $M$ 的挠元素。练习 1--6 与 20 给出秩的另一种刻画，并给出一些（无挠）秩 $1$ 但非自由的 $R$-模例子。

下面这个重要结果说明：如果 $N$ 是 P.I.D. 上有限秩自由模的子模，则 $N$ 仍是有限秩自由模；而且可以为两个模选择生成元，使它们之间具有一种简单关系。

\noindent\textbf{定理 4.} 设 $R$ 是主理想整环，$M$ 是有限秩 $n$ 的自由 $R$-模，$N$ 是 $M$ 的子模。则
\begin{enumerate}[label=(\arabic*)]
\item $N$ 是自由模，秩为 $m$，且 $m\leq n$；
\item 存在 $M$ 的一组基 $y_1,y_2,\ldots,y_n$，使得 $a_1y_1,a_2y_2,\ldots,a_my_m$ 是 $N$ 的一组基，其中 $a_1,a_2,\ldots,a_m$ 是 $R$ 中非零元素，并满足整除关系
\[
a_1\mid a_2\mid\cdots\mid a_m.
\]
\end{enumerate}

\begin{proof}
当 $N=\{0\}$ 时定理平凡，故设 $N\neq\{0\}$。对每个 $R$-模同态 $\varphi:M\to R$，$N$ 的像 $\varphi(N)$ 是 $R$ 的子模，即 $R$ 的理想。由于 $R$ 是 P.I.D.，该理想必为主理想，记为 $\varphi(N)=(a_\varphi)$，其中 $a_\varphi\in R$。令
\[
\Sigma=\{(a_\varphi)\mid \varphi\in\operatorname{Hom}_R(M,R)\}
\]
为以这种方式由 $M$ 到 $R$ 的 $R$-模同态得到的主理想集合。集合 $\Sigma$ 非空，因为取零同态便得到 $(0)\in\Sigma$。由推论 2，$\Sigma$ 至少有一个极大元；也就是说，存在至少一个同态 $v:M\to R$，使得主理想 $v(N)=(a_v)$ 不真包含在 $\Sigma$ 的任何其他元素中。令 $a_1=a_v$，并取 $y\in N$ 使 $v(y)=a_1$。

先说明 $a_1$ 非零。令 $x_1,x_2,\ldots,x_n$ 是自由模 $M$ 的任意一组基，并令 $\pi_i\in\operatorname{Hom}_R(M,R)$ 为相对于这组基投到第 $i$ 个坐标的自然投影同态。因为 $N\neq\{0\}$，存在某个 $i$ 使 $\pi_i(N)\neq0$；这特别说明 $\Sigma$ 不仅仅含有零理想。由于 $(a_1)$ 是 $\Sigma$ 的极大元，得到 $a_1\neq0$。

接下来说明，对每个 $\varphi\in\operatorname{Hom}_R(M,R)$，元素 $a_1$ 都整除 $\varphi(y)$。令 $d$ 是由 $a_1$ 和 $\varphi(y)$ 生成的主理想的生成元。因此 $d$ 是 $a_1$ 与 $\varphi(y)$ 的公因子，并且存在 $r_1,r_2\in R$ 使
\[
d=r_1a_1+r_2\varphi(y).
\]
考虑同态 $\psi=r_1v+r_2\varphi:M\to R$。则
\[
\psi(y)=(r_1v+r_2\varphi)(y)=r_1a_1+r_2\varphi(y)=d.
\]
所以 $d\in\psi(N)$，从而 $(d)\subseteq\psi(N)$。另一方面，因为 $d$ 是 $a_1$ 的因子，也有 $(a_1)\subseteq(d)$。于是
\[
(a_1)\subseteq(d)\subseteq\psi(N).
\]
由 $(a_1)$ 在 $\Sigma$ 中的极大性，只能有等号：$(a_1)=(d)=\psi(N)$。特别地，$(a_1)=(d)$ 说明 $a_1\mid\varphi(y)$，因为 $d$ 整除 $\varphi(y)$。

把这个结论应用于投影同态 $\pi_i$，可知 $a_1$ 整除所有 $\pi_i(y)$。写
\[
\pi_i(y)=a_1b_i\qquad(1\leq i\leq n)
\]
并定义
\[
y_1=\sum_{i=1}^n b_ix_i.
\]
注意 $a_1y_1=y$。因为
\[
a_1=v(y)=v(a_1y_1)=a_1v(y_1),
\]
且 $a_1$ 是整环 $R$ 中的非零元素，所以
\[
v(y_1)=1.
\]

现在验证这个元素 $y_1$ 可以取为 $M$ 的一组基中的一个元素，而且 $a_1y_1$ 可以取为 $N$ 的一组基中的一个元素；也就是证明
\[
\text{(a)}\quad M=Ry_1\oplus\ker v,
\qquad
\text{(b)}\quad N=Ra_1y_1\oplus (N\cap\ker v).
\]
为证明 (a)，取任意 $x\in M$，写成
\[
x=v(x)y_1+(x-v(x)y_1).
\]
因为
\[
v(x-v(x)y_1)=v(x)-v(x)v(y_1)=v(x)-v(x)\cdot1=0,
\]
可知 $x-v(x)y_1\in\ker v$。这说明 $x$ 可以写成 $Ry_1$ 中元素与 $\ker v$ 中元素之和，故 $M=Ry_1+\ker v$。若 $ry_1\in\ker v$，则
\[
0=v(ry_1)=rv(y_1)=r,
\]
所以这个交中的元素只能为 $0$，于是直和成立。

为证明 (b)，注意对每个 $x'\in N$，按 $a_1$ 作为 $v(N)$ 的生成元的定义，$v(x')$ 被 $a_1$ 整除。写 $v(x')=ba_1$（$b\in R$），则上面 (a) 中使用的分解给出
\[
x'=v(x')y_1+(x'-v(x')y_1)=ba_1y_1+(x'-ba_1y_1),
\]
其中第二项既在 $\ker v$ 中，也在 $N$ 中。这说明 $N=Ra_1y_1+(N\cap\ker v)$。该和为直和则是 (a) 中直和性的特例。

现在对 $N$ 的秩 $m$ 作归纳来证明 (1)。如果 $m=0$，则 $N$ 是挠模；但自由模无挠，故 $N=0$，(1) 平凡成立。于是设 $m>0$。由 (b) 中的直和容易看出 $N\cap\ker v$ 的秩为 $m-1$（参见练习 3）。由归纳假设，$N\cap\ker v$ 是秩 $m-1$ 的自由 $R$-模。再由 (b) 的直和性可知，把 $a_1y_1$ 添加到 $N\cap\ker v$ 的任一组基中，得到 $N$ 的一组基，所以 $N$ 也是自由的（秩为 $m$），这证明了 (1)。

最后对 $M$ 的秩 $n$ 作归纳证明 (2)。把 (1) 应用于子模 $\ker v$，可知这个子模是自由的；又由 (a) 的直和可知它的秩为 $n-1$。由归纳假设应用于模 $\ker v$（它扮演 $M$ 的角色）及其子模 $\ker v\cap N$（它扮演 $N$ 的角色），可知存在 $\ker v$ 的一组基 $y_2,y_3,\ldots,y_n$，使得 $a_2y_2,a_3y_3,\ldots,a_my_m$ 是 $N\cap\ker v$ 的一组基，并且 $a_2,a_3,\ldots,a_m\in R$ 满足
\[
a_2\mid a_3\mid\cdots\mid a_m.
\]
由 (a) 与 (b) 的直和性，$y_1,y_2,\ldots,y_n$ 是 $M$ 的一组基，而 $a_1y_1,a_2y_2,\ldots,a_my_m$ 是 $N$ 的一组基。为完成归纳，只需说明 $a_1\mid a_2$。定义同态 $\varphi:M\to R$，在 $M$ 的这组基上令
\[
\varphi(y_1)=\varphi(y_2)=1,
\qquad
\varphi(y_i)=0\quad(i>2).
\]
于是对这个同态 $\varphi$，有 $a_1=\varphi(a_1y_1)$，所以 $a_1\in\varphi(N)$，从而 $(a_1)\subseteq\varphi(N)$。由 $(a_1)$ 在 $\Sigma$ 中的极大性，得到 $(a_1)=\varphi(N)$。又因为 $a_2=\varphi(a_2y_2)\in\varphi(N)$，所以 $a_2\in(a_1)$，即 $a_1\mid a_2$。定理证明完毕。
\end{proof}

回忆：对任意环 $R$（不必交换，也不必有 $1$），左 $R$-模 $C$ 称为循环 $R$-模，如果存在 $x\in C$ 使 $C=Rx$。此时可定义 $R$-模同态
\[
\pi:R\longrightarrow C
\]
为 $\pi(r)=rx$；由于 $C=Rx$，这个同态是满射。第一同构定理给出左 $R$-模同构
\[
R/\ker\pi\cong C.
\]
若 $R$ 是 P.I.D.，则 $\ker\pi$ 是主理想 $(a)$，所以循环 $R$-模 $C$ 具有 $R/(a)$ 的形式，其中 $(a)=\operatorname{Ann}(C)$。

循环模是最简单的模（因为只需要一个生成元）。基本定理的存在性部分断言：P.I.D. 上任意有限生成模同构于有限多个循环模的直和。

\noindent\textbf{定理 5（基本定理，存在性：不变因子形式）.} 设 $R$ 是 P.I.D.，$M$ 是有限生成 $R$-模。
\begin{enumerate}[label=(\arabic*)]
\item $M$ 同构于有限多个循环模的直和。更准确地说，
\[
M\cong R^r\oplus R/(a_1)\oplus R/(a_2)\oplus\cdots\oplus R/(a_m)
\]
其中 $r\geq0$ 是整数，$a_1,a_2,\ldots,a_m$ 是 $R$ 中非零且非单位的元素，并满足整除关系
\[
a_1\mid a_2\mid\cdots\mid a_m.
\]
\item $M$ 无挠当且仅当 $M$ 自由。
\item 在 (1) 的分解中，
\[
\operatorname{Tor}(M)\cong R/(a_1)\oplus R/(a_2)\oplus\cdots\oplus R/(a_m).
\]
特别地，$M$ 是挠模当且仅当 $r=0$；在这种情形下，$M$ 的零化子是理想 $(a_m)$。
\end{enumerate}

\begin{proof}
由于 $M$ 按假设可由有限集生成，取 $M$ 的一组生成元 $x_1,x_2,\ldots,x_n$，并要求其基数最小。令 $R^n$ 为秩 $n$ 的自由 $R$-模，基为 $b_1,b_2,\ldots,b_n$。定义同态
\[
\pi:R^n\longrightarrow M
\]
由 $\pi(b_i)=x_i$（对所有 $i$）给出。由于 $x_1,\ldots,x_n$ 生成 $M$，$\pi$ 自动满射。由模的第一同构定理，
\[
R^n/\ker\pi\cong M.
\]
现在对 $R^n$ 及其子模 $\ker\pi$ 应用定理 4，可以选取 $R^n$ 的另一组基 $y_1,y_2,\ldots,y_n$，使得 $a_1y_1,a_2y_2,\ldots,a_my_m$ 是 $\ker\pi$ 的一组基，其中 $a_1,a_2,\ldots,a_m\in R$ 且
\[
a_1\mid a_2\mid\cdots\mid a_m.
\]
这给出
\[
M\cong R^n/\ker\pi=(Ry_1\oplus Ry_2\oplus\cdots\oplus Ry_n)/(Ra_1y_1\oplus Ra_2y_2\oplus\cdots\oplus Ra_my_m).
\]

为了识别右侧商模，我们使用自然满的 $R$-模同态
\[
Ry_1\oplus\cdots\oplus Ry_n\longrightarrow R/(a_1)\oplus R/(a_2)\oplus\cdots\oplus R/(a_m)\oplus R^{n-m}
\]
它把
\[
\alpha_1y_1+\cdots+\alpha_ny_n
\]
送到
\[
(\alpha_1\bmod(a_1),\ldots,\alpha_m\bmod(a_m),\alpha_{m+1},\ldots,\alpha_n).
\]
这个同态的核显然正是那些 $\alpha_i$ 被 $a_i$ 整除（$i=1,2,\ldots,m$）的元素集合，也就是
\[
Ra_1y_1\oplus Ra_2y_2\oplus\cdots\oplus Ra_my_m
\]
（参见练习 7）。于是得到
\[
M\cong R/(a_1)\oplus R/(a_2)\oplus\cdots\oplus R/(a_m)\oplus R^{n-m}.
\]
若某个 $a_i$ 是 $R$ 中单位，则 $R/(a_i)=0$，因此可从直和中删去那些初始的单位因子对应的项。这样得到 (1) 中的分解（其中 $r=n-m$）。

由于对 $R$ 中任意非零元素 $a$，$R/(a)$ 都是挠 $R$-模，(1) 立即蕴含 $M$ 是无挠模当且仅当 $M\cong R^r$，这就是 (2)。而 (3) 也直接由定义得到，因为 $R/(a)$ 的零化子显然是理想 $(a)$，并且整除关系使各循环挠和的总零化子为最大不变因子的理想 $(a_m)$。
\end{proof}

我们很快将证明定理 5 中分解的唯一性：如果还有
\[
M\cong R^{r'}\oplus R/(b_1)\oplus R/(b_2)\oplus\cdots\oplus R/(b_{m'})
\]
其中 $r'\geq0$，$b_1,b_2,\ldots,b_{m'}$ 是非零非单位元素并满足
\[
b_1\mid b_2\mid\cdots\mid b_{m'},
\]
则 $r=r'$，$m=m'$，并且 $(a_i)=(b_i)$（也就是 $a_i$ 与 $b_i$ 只差一个单位因子）对所有 $i$ 成立。正是整除条件 $a_1\mid a_2\mid\cdots\mid a_m$ 给出这种唯一性。

\noindent\textbf{定义.} 定理 5 中的整数 $r$ 称为 $M$ 的\emph{自由秩}或 \emph{Betti 数}；元素 $a_1,a_2,\ldots,a_m\in R$（只差乘以 $R$ 中单位的歧义）称为 $M$ 的\emph{不变因子}。

注意，在证明 $M$ 的不变因子的唯一性之前，严格地说我们应称它们为 $M$ 的一组不变因子（自由秩同理），意思是能给出上述定理 (1) 中分解的任意一组元素。

利用中国剩余定理，可以把定理 5 中的循环模进一步分解，使 $M$ 成为若干循环模的直和，而这些循环模的零化子尽可能简单（即为 $(0)$ 或由 $R$ 中素元幂生成）。这给出另一种分解；我们也会看到它是唯一的。现在描述它。

设 $a$ 是主理想整环 $R$ 中的非零元素。由于 $R$ 也是唯一因子分解整环，可写
\[
a=up_1^{\alpha_1}p_2^{\alpha_2}\cdots p_s^{\alpha_s},
\]
其中 $p_i$ 是 $R$ 中两两不同的素元，$u$ 是单位。这个分解在差一个单位意义下唯一，所以理想 $(p_i^{\alpha_i})$（$i=1,\ldots,s$）唯一确定。若 $i\neq j$，则
\[
(p_i^{\alpha_i})+(p_j^{\alpha_j})=R,
\]
因为这两个理想的和由一个最大公因子生成，而对不同素元 $p_i,p_j$ 来说最大公因子为 $1$。换言之，理想 $(p_i^{\alpha_i})$ 两两互素。所有这些理想的交就是 $(a)$，因为 $a$ 是 $p_1^{\alpha_1},p_2^{\alpha_2},\ldots,p_s^{\alpha_s}$ 的最小公倍元。因此中国剩余定理（定理 7.17）说明
\[
R/(a)\cong R/(p_1^{\alpha_1})\oplus R/(p_2^{\alpha_2})\oplus\cdots\oplus R/(p_s^{\alpha_s})
\]
既是环同构，也是 $R$-模同构。

把这应用到定理 5 中的模，可将 $M$ 的每个不变因子 $a_i$ 对应的直和项 $R/(a_i)$ 写成若干循环模的直和，这些循环模的零化子为 $a_i$ 的素幂因子。于是得到：

\noindent\textbf{定理 6（基本定理，存在性：初等因子形式）.} 设 $R$ 是 P.I.D.，$M$ 是有限生成 $R$-模。则 $M$ 是有限多个循环模的直和，其零化子或者为 $(0)$，或者由 $R$ 中素元的幂生成；也就是说，
\[
M\cong R^r\oplus R/(p_1^{\alpha_1})\oplus R/(p_2^{\alpha_2})\oplus\cdots\oplus R/(p_t^{\alpha_t}),
\]
其中 $r\geq0$ 为整数，$p_1^{\alpha_1},\ldots,p_t^{\alpha_t}$ 是 $R$ 中（不必互异的）素元的正幂。

我们是用 $M$ 的不变因子的素幂因子证明定理 6 的。事实上，我们将看到，把 $M$ 分解为定理 6 中零化子为 $(0)$ 或素幂的循环模直和也是唯一的；也就是说，整数 $r$ 与理想 $(p_1^{\alpha_1}),\ldots,(p_t^{\alpha_t})$ 由 $M$ 唯一确定。这些素幂有名称：

\noindent\textbf{定义.} 设 $R$ 是 P.I.D.，$M$ 是如定理 6 中的有限生成 $R$-模。素幂 $p_1^{\alpha_1},\ldots,p_t^{\alpha_t}$（只差乘以 $R$ 中单位的歧义）称为 $M$ 的\emph{初等因子}。

设 $M$ 是主理想整环 $R$ 上的有限生成挠模。如果对定理 6 的分解中出现的不同素元 $p_1,p_2,\ldots,p_n$，把所有对应同一素元 $p_i$ 的循环因子聚在一起，我们特别得到 $M$ 可写成直和
\[
M=N_1\oplus N_2\oplus\cdots\oplus N_n,
\]
其中 $N_i$ 由 $M$ 中被素元 $p_i$ 的某个幂零化的所有元素组成。这个结论对不一定有限生成的 $R$-模也成立：

\noindent\textbf{定理 7（Primary 分解定理）.} 设 $R$ 是 P.I.D.，$M$ 是非零挠 $R$-模（不必有限生成），且具有非零零化子 $a$。设 $a$ 分解为 $R$ 中不同素元的幂的乘积：
\[
a=up_1^{\alpha_1}p_2^{\alpha_2}\cdots p_n^{\alpha_n},
\]
并令
\[
N_i=\{x\in M\mid p_i^{\alpha_i}x=0\},\qquad 1\leq i\leq n.
\]
则 $N_i$ 是 $M$ 的子模，零化子为 $(p_i^{\alpha_i})$，并且它是所有被 $p_i$ 的某个幂零化的元素构成的子模。我们有
\[
M=N_1\oplus N_2\oplus\cdots\oplus N_n.
\]
如果 $M$ 有限生成，则每个 $N_i$ 是有限多个循环模的直和，其零化子均为 $p_i^{\alpha_i}$ 的因子。

\begin{proof}
当 $M$ 在 $R$ 上有限生成时，这些结论已经由定理 6 及其后的分组讨论得到。在一般情形下，显然每个 $N_i$ 是 $M$ 的子模，且其零化子整除 $(p_i^{\alpha_i})$。由于 $R$ 是 P.I.D.，当 $i\neq j$ 时理想 $(p_i^{\alpha_i})$ 与 $(p_j^{\alpha_j})$ 互素，因此可以把中国剩余定理证明中的论证轻微修改后应用于模，得到 $M$ 的这个直和分解。利用这个直和分解，容易看出 $N_i$ 的零化子恰为 $(p_i^{\alpha_i})$。
\end{proof}

\noindent\textbf{定义.} 前一定理中的子模 $N_i$ 称为 $M$ 的 $p_i$-\emph{primary 分量}。

用这个术语来说，有限生成模 $M$ 的初等因子正是 $\operatorname{Tor}(M)$ 的 primary 分量的不变因子。

现在证明关于基本定理中分解的唯一性陈述。注意，如果 $M$ 是交换环 $R$ 上的任意模，$a\in R$，则
\[
aM=\{am\mid m\in M\}
\]
是 $M$ 的子模。还回忆，在主理想整环 $R$ 中，非零素理想都是极大理想，所以 $R$ 对非零素理想的商是域。

\noindent\textbf{引理 8.} 设 $R$ 是 P.I.D.，$p$ 是 $R$ 中素元，$F=R/(p)$。
\begin{enumerate}[label=(\arabic*)]
\item 若 $M=R^r$，则 $M/pM\cong F^r$。
\item 若 $M=R/(a)$，其中 $a\in R$ 非零，则
\[
M/pM\cong
\begin{cases}
F,& p\mid a,\\
0,& p\nmid a.
\end{cases}
\]
\item 若 $M=R/(a_1)\oplus R/(a_2)\oplus\cdots\oplus R/(a_k)$，且每个 $a_i$ 都被 $p$ 整除，则 $M/pM\cong F^k$。
\end{enumerate}

\begin{proof}
(1) 从 $R^r$ 到 $(R/(p))^r$ 有自然映射，把 $(a_1,\ldots,a_r)$ 送到 $(a_1\bmod(p),\ldots,a_r\bmod(p))$。这显然是满的 $R$-模同态，它的核由所有坐标都被 $p$ 整除的 $r$ 元组组成，也就是 $pR^r$。故 $R^r/pR^r\cong(R/(p))^r=F^r$。

(2) 这由同构定理得到。先注意 $p(R/(a))$ 是理想 $(p)$ 在商 $R/(a)$ 中的像，因此等于
\[
((p)+(a))/(a).
\]
理想 $(p)+(a)$ 由 $p$ 与 $a$ 的一个最大公因子生成；所以当 $p\mid a$ 时它是 $(p)$，当 $p\nmid a$ 时它是 $R=(1)$。因此当 $p\mid a$ 时 $pM=(p)/(a)$，当 $p\nmid a$ 时 $pM=R/(a)=M$。若 $p\mid a$，则
\[
M/pM=(R/(a))/((p)/(a))\cong R/(p)=F;
\]
若 $p\nmid a$，则 $M/pM=M/M=0$。

(3) 由 (2) 对直和逐项应用得到，正如定理 5(1) 证明中识别直和商的方式一样。
\end{proof}

\noindent\textbf{定理 9（基本定理，唯一性）.} 设 $R$ 是 P.I.D.。
\begin{enumerate}[label=(\arabic*)]
\item 两个有限生成 $R$-模 $M_1$ 与 $M_2$ 同构，当且仅当它们有相同的自由秩和相同的不变因子列表。
\item 两个有限生成 $R$-模 $M_1$ 与 $M_2$ 同构，当且仅当它们有相同的自由秩和相同的初等因子列表。
\end{enumerate}

\begin{proof}
如果 $M_1$ 与 $M_2$ 有相同的自由秩和相同的不变因子列表，或者有相同的自由秩和相同的初等因子列表，则它们显然同构。

反过来，设 $M_1\cong M_2$。任一 $M_1$ 与 $M_2$ 之间的同构都把 $M_1$ 的挠部分映到 $M_2$ 的挠部分，所以
\[
\operatorname{Tor}(M_1)\cong\operatorname{Tor}(M_2).
\]
于是
\[
R^{r_1}\cong M_1/\operatorname{Tor}(M_1)
\cong M_2/\operatorname{Tor}(M_2)\cong R^{r_2},
\]
其中 $r_1$、$r_2$ 分别是 $M_1$、$M_2$ 的自由秩。令 $p$ 为 $R$ 中任意非零素元。由 $R^{r_1}\cong R^{r_2}$ 可得
\[
R^{r_1}/pR^{r_1}\cong R^{r_2}/pR^{r_2}.
\]
由引理 8(1)，这说明 $F^{r_1}\cong F^{r_2}$，其中 $F=R/pR$ 是域。因此两个 $F$-向量空间同构，维数相等，故 $r_1=r_2$。所以 $M_1$ 与 $M_2$ 有相同的自由秩。

于是只需证明 $M_1$ 与 $M_2$ 有相同的不变因子列表和初等因子列表。为此只需要处理同构的挠模 $\operatorname{Tor}(M_1)$ 与 $\operatorname{Tor}(M_2)$；也就是说，不妨假设 $M_1,M_2$ 都是挠 $R$-模。

先证明它们有相同的初等因子。只需说明：对任一固定素元 $p$，作为 $p$ 的幂的初等因子在 $M_1$ 与 $M_2$ 中相同。如果 $M_1\cong M_2$，则 $M_1$ 的 $p$-primary 子模（即初等因子为 $p$ 的幂的循环因子的直和）同构于 $M_2$ 的 $p$-primary 子模，因为这些子模正是所有被 $p$ 的某个幂零化的元素组成的子模。因此归结为证明：若两个零化子为 $p$ 的幂的模 $M_1,M_2$ 同构，则它们有相同的初等因子。

对 $M_1$ 的零化子中 $p$ 的幂指数作归纳（这与 $M_2$ 中的指数相同，因为两者同构）。若该幂指数为 $0$，则 $M_1=M_2=0$，结论成立。否则 $M_1$ 与 $M_2$ 都有非平凡初等因子。设 $M_1$ 的初等因子为
\[
\underbrace{p,p,\ldots,p}_{m\text{ 次}},\quad p^{\alpha_1},p^{\alpha_2},\ldots,p^{\alpha_s},
\]
其中 $2\leq\alpha_1\leq\alpha_2\leq\cdots\leq\alpha_s$；也就是说，$M_1$ 是若干循环模的直和，生成元可记为 $x_1,x_2,\ldots,x_m,x_{m+1},\ldots,x_{m+s}$，其零化子分别为 $(p),\ldots,(p),(p^{\alpha_1}),\ldots,(p^{\alpha_s})$。那么子模 $pM_1$ 的初等因子为
\[
p^{\alpha_1-1},p^{\alpha_2-1},\ldots,p^{\alpha_s-1},
\]
因为 $pM_1$ 是以 $px_1,\ldots,px_m,px_{m+1},\ldots,px_{m+s}$ 生成的循环模直和，而这些生成元的零化子为 $(1),\ldots,(1),(p^{\alpha_1-1}),\ldots,(p^{\alpha_s-1})$。

类似地，设 $M_2$ 的初等因子为
\[
\underbrace{p,p,\ldots,p}_{n\text{ 次}},\quad p^{\beta_1},p^{\beta_2},\ldots,p^{\beta_t},
\]
其中 $2\leq\beta_1\leq\beta_2\leq\cdots\leq\beta_t$，则 $pM_2$ 的初等因子为
\[
p^{\beta_1-1},p^{\beta_2-1},\ldots,p^{\beta_t-1}.
\]
由于 $M_1\cong M_2$，也有 $pM_1\cong pM_2$；并且 $pM_1$ 的零化子中 $p$ 的幂指数比 $M_1$ 的少 $1$。由归纳假设，$pM_1$ 与 $pM_2$ 的初等因子相同，于是 $s=t$ 且 $\alpha_i-1=\beta_i-1$，即 $\alpha_i=\beta_i$（$1\leq i\leq s$）。最后，因为
\[
M_1/pM_1\cong M_2/pM_2,
\]
由引理 8(3) 得到 $F^{m+s}\cong F^{n+t}$，从而 $m+s=n+t$；已知 $s=t$，故 $m=n$。这证明了 $M_1$ 与 $M_2$ 的初等因子集合相同。

现在证明 $M_1$ 与 $M_2$ 必有相同的不变因子。设
\[
a_1\mid a_2\mid\cdots\mid a_m
\]
是 $M_1$ 的不变因子。把这些元素分解为素幂因子，得到 $M_1$ 的一组初等因子。注意整除关系意味着 $a_m$ 是这些初等因子中每个不同素元所对应的最大素幂的乘积；去掉 $a_m$ 的这些因子后，$a_{m-1}$ 是剩余初等因子中每个不同素元所对应最大素幂的乘积；依此类推。若
\[
b_1\mid b_2\mid\cdots\mid b_n
\]
是 $M_2$ 的不变因子，同样地通过分解得到 $M_2$ 的一组初等因子。上面已经证明 $M_1$ 与 $M_2$ 的初等因子相同，因此由刚才的重组规则可知它们的不变因子也相同。
\end{proof}

\noindent\textbf{推论 10.} 设 $R$ 是 P.I.D.，$M$ 是有限生成 $R$-模。
\begin{enumerate}[label=(\arabic*)]
\item $M$ 的初等因子是 $M$ 的不变因子的素幂因子。
\item $M$ 的最大不变因子是 $M$ 的初等因子中互异素元对应的最大素幂的乘积；次大的不变因子是剩余初等因子中互异素元对应最大素幂的乘积；依此类推。
\end{enumerate}

\begin{proof}
(1) 中的过程给出一组初等因子，而初等因子由定理 9 唯一，所以 (1) 中的过程给出的就是初等因子的集合。(2) 同理。
\end{proof}

\noindent\textbf{推论 11（有限生成阿贝尔群基本定理）.} 见定理 5.3 与定理 5.5。

\begin{proof}
在定理 5、6 与 9 中取 $R=\mathbb Z$ 即可（不过注意，第 5 章中为了计算方便，不变因子的列出顺序与这里相反）。
\end{proof}

推论 10 中在初等因子与不变因子之间转换的过程，在第 5 章有限生成阿贝尔群的情形中有较详细的描述。还应注意，如果一个有限生成模 $M$ 写成形如 $R/(a)$ 的循环模的直和，则出现的理想 $(a)$ 一般并不唯一，除非附加某些条件（例如不变因子的整除条件，或者在初等因子情形中要求 $a$ 是素元的幂）。要判断两个模是否同构，必须先把它们写成这种标准（或规范）形式。
\subsection*{练习}

\begin{enumerate}[label=\arabic*.]
\item 设 $M$ 是整环 $R$ 上的模。
\begin{enumerate}[label=(\alph*)]
\item 设 $x$ 是 $M$ 中非零挠元素。证明 $x$ 与 $0$ ``线性相关''。推出 $\operatorname{Tor}(M)$ 的秩为 $0$，因而任意挠 $R$-模的秩为 $0$。
\item 证明 $M$ 的秩等于无挠商模 $M/\operatorname{Tor}M$ 的秩。
\end{enumerate}

\item 设 $M$ 是整环 $R$ 上的模。
\begin{enumerate}[label=(\alph*)]
\item 设 $M$ 的秩为 $n$，且 $x_1,x_2,\ldots,x_n$ 是 $M$ 中任意极大线性无关集。令 $N=Rx_1+\cdots+Rx_n$。证明 $N\cong R^n$，且 $M/N$ 是挠 $R$-模。等价地，$x_1,\ldots,x_n$ 线性无关，并且对任意 $y\in M$，存在非零 $r\in R$，使得 $ry$ 可写为 $r_1x_1+\cdots+r_nx_n$。
\item 反过来，若 $M$ 含有自由秩为 $n$ 的子模 $N$（即 $N\cong R^n$），并且 $M/N$ 是挠 $R$-模，证明 $M$ 的秩为 $n$。提示：令 $y_1,\ldots,y_{n+1}$ 为 $M$ 中任意 $n+1$ 个元素。利用 $M/N$ 为挠模，取非零 $r_i\in R$ 使 $r_i y_i$ 都可写成 $N$ 的一组基的线性组合，再仿照命题 3 的证明说明 $r_i y_i$，从而 $y_i$，线性相关。
\end{enumerate}

\item 设 $R$ 为整环，$A,B$ 分别为秩 $m,n$ 的 $R$-模。证明 $A\oplus B$ 的秩为 $m+n$。提示：使用上一题。

\item 设 $R$ 为整环，$M$ 为 $R$-模，$N$ 为 $M$ 的子模。设 $M$ 的秩为 $n$，$N$ 的秩为 $r$，商模 $M/N$ 的秩为 $s$。证明 $n=r+s$。提示：取 $x_1,\ldots,x_s\in M$，其像在 $M/N$ 中构成极大线性无关集；再取 $x_{s+1},\ldots,x_{s+r}$ 为 $N$ 中极大线性无关集。证明这些元素在 $M$ 中线性无关，并且任意 $y\in M$ 的某个非零倍数可写成它们的线性组合，然后使用练习 2。

\item 令 $R=\mathbb Z[x]$，令 $M=(2,x)$ 为由 $2$ 与 $x$ 生成的理想，把它看作 $R$ 的子模。证明 $\{2,x\}$ 不是 $M$ 的基。提示：找出这两个元素之间的一个非平凡 $R$-线性相关关系。证明 $M$ 的秩为 $1$，但 $M$ 不是秩 $1$ 的自由模（参见练习 2）。

\item 证明：若 $R$ 是整环，$M$ 是 $R$ 的任意非主理想，则 $M$ 是无挠、秩 $1$ 的 $R$-模，但不是自由 $R$-模。

\item 设 $R$ 为任意环，$A_1,A_2,\ldots,A_m$ 为 $R$-模，$B_i$ 为 $A_i$ 的子模，$1\leq i\leq m$。证明
\[
(A_1\oplus\cdots\oplus A_m)/(B_1\oplus\cdots\oplus B_m)
\cong (A_1/B_1)\oplus\cdots\oplus(A_m/B_m).
\]

\item 设 $R$ 为 P.I.D.，$B$ 为挠 $R$-模，$p$ 为 $R$ 中素元。证明：若存在非零 $b\in B$ 满足 $pb=0$，则 $\operatorname{Ann}(B)\subseteq (p)$。

\item 给出一个整环 $R$ 与非零挠 $R$-模 $M$ 的例子，使得 $\operatorname{Ann}(M)=0$。证明：若 $N$ 是有限生成挠 $R$-模，则 $\operatorname{Ann}(N)\ne0$。

\item 设 $p$ 为 P.I.D. $R$ 中的素元，$N$ 为 $R$-模。证明 $N$ 的 $p$-primary 分量是 $N$ 的子模，并证明 $N$ 是其各 $p$-primary 分量的直和（分量个数不必有限）。

\item 设 $R$ 为 P.I.D.，$a\in R$ 非零，$M=R/(a)$。对 $R$ 的任意素元 $p$，证明
\[
p^{k-1}M/p^kM\cong
\begin{cases}
R/(p),& 1\leq k\leq n,\\
0,& k>n,
\end{cases}
\]
其中 $n$ 是 $a$ 中 $p$ 的幂指数。

\item 设 $R$ 为 P.I.D.，$p$ 为 $R$ 中素元。
\begin{enumerate}[label=(\alph*)]
\item 设 $M$ 是有限生成挠 $R$-模。用上一题证明
\[
p^{k-1}M/p^kM\cong F^{n_k},
\]
其中 $F=R/(p)$，$n_k$ 是 $M$ 的初等因子中形如 $p^a$ 且 $a\geq k$ 的个数。
\item 设 $M_1,M_2$ 是同构的有限生成挠 $R$-模。用 (a) 证明：对每个 $k\geq0$，$M_1$ 与 $M_2$ 有相同数目的初等因子 $p^a$，其中 $a\geq k$。证明这推出 $M_1$ 与 $M_2$ 有相同的初等因子集合。
\end{enumerate}

\item 若 $M$ 是 P.I.D. $R$ 上的有限生成模，描述 $M/\operatorname{Tor}(M)$ 的结构。

\item 设 $R$ 为 P.I.D.，$M$ 为挠 $R$-模。证明 $M$ 不可约（参见第 10.3 节练习 9--11）当且仅当对每个非零 $m\in M$，都有 $M=Rm$，且 $m$ 的零化子是某个非零素理想 $(p)$。

\item 证明：若 $R$ 是 Noether 环，则 $R^n$ 是 Noether $R$-模。提示：固定 $R^n$ 的一组基。若 $M$ 是 $R^n$ 的子模，说明 $M$ 中元素的第一坐标构成 $R$ 的子模，因而有限生成。取 $m_1,
\ldots,m_k\in M$，使它们的第一坐标生成该子模。证明 $M$ 的任意元素都可写成这些 $m_i$ 的 $R$-线性组合加上一个第一坐标为 $0$ 的 $M$ 中元素。再证明 $M\cap R^{n-1}$ 是 $R^{n-1}$ 的子模，其中 $R^{n-1}$ 表示 $R^n$ 中第一坐标为 $0$ 的元素集合，最后对 $n$ 归纳。

\medskip
以下一组练习给出定理 5 在 $R$ 为 Euclidean 整环时的矩阵证明，使用行、列初等变换。这尤其适用于 $R=\mathbb Z$ 与 $R=F[x]$，并且在计算上很有用。设 $R$ 为 Euclidean 整环，$M$ 为 $R$-模。

\item 证明 $M$ 有限生成当且仅当存在某个整数 $n$ 和满的 $R$-同态 $\varphi:R^n\to M$。（这对任意环 $R$ 都成立。）

设 $\varphi:R^n\to M$ 为满 $R$-模同态。由练习 15，$\ker\varphi$ 有限生成。若 $x_1,
\ldots,x_n$ 是 $R^n$ 的基，$y_1,
\ldots,y_m$ 是 $\ker\varphi$ 的生成元，则
\[
y_i=a_{i1}x_1+a_{i2}x_2+\cdots+a_{in}x_n,
\qquad i=1,2,\ldots,m,
\]
其中 $a_{ij}\in R$。因此同态 $\varphi$（也就是 $M$ 的模结构）由 $R^n$ 的生成元选择和矩阵 $A=(a_{ij})$ 决定。这样的矩阵称为\emph{关系矩阵}。

\item
\begin{enumerate}[label=(\alph*)]
\item 证明在 $R^n$ 的基中交换 $x_i$ 与 $x_j$，会交换相应关系矩阵的第 $i$ 列和第 $j$ 列。
\item 证明：对任意 $a\in R$，把基元素 $x_j$ 替换为 $x_j-ax_i$ 后仍得到 $R^n$ 的一组基；相应关系矩阵等于原关系矩阵，只是把 $a$ 倍的第 $j$ 列加到第 $i$ 列上。提示：
\[
\cdots+a_ix_i+\cdots+a_jx_j+\cdots
=\cdots+(a_i+aa_j)x_i+\cdots+a_j(x_j-ax_i)+\cdots.
\]
\end{enumerate}

\item
\begin{enumerate}[label=(\alph*)]
\item 证明交换生成元 $y_i,y_j$ 会交换关系矩阵的第 $i$ 行和第 $j$ 行。
\item 证明：对任意 $a\in R$，把 $y_j$ 替换为 $y_j-ay_i$ 仍给出 $\ker\varphi$ 的一组生成元；相应关系矩阵等于原矩阵，只是把 $-a$ 倍的第 $i$ 行加到第 $j$ 行上。
\end{enumerate}

\item 由前两题可知，可通过重新选择 $R^n$ 与 $\ker\varphi$ 的生成元，对给定关系矩阵作初等行、列变换。若所有关系矩阵都是零矩阵，则 $\ker\varphi=0$，$M\cong R^n$。否则，令 $a_1$ 为某个初始关系矩阵中所有非零元素的最大公因子。
\begin{enumerate}[label=(\alph*)]
\item 证明可用初等行、列变换使 $a_1$ 出现在关系矩阵左上角，并且所得矩阵中 $a_1$ 整除所有元素。
\item 证明存在如下形式的关系矩阵
\[
\begin{pmatrix}
 a_1&0&\cdots&0\\
 0&a_{22}&\cdots&a_{2n}\\
 \vdots&\vdots&\ddots&\vdots\\
 0&a_{m2}&\cdots&a_{mn}
\end{pmatrix},
\]
其中 $a_1$ 整除所有条目。
\item 令 $a_2$ 为 (b) 中除 $a_1$ 之外所有条目的最大公因子。证明存在如下形式的关系矩阵
\[
\begin{pmatrix}
 a_1&0&0&\cdots&0\\
 0&a_2&0&\cdots&0\\
 0&0&a_{33}&\cdots&a_{3n}\\
 \vdots&\vdots&\vdots&\ddots&\vdots
\end{pmatrix},
\]
其中 $a_1\mid a_2$，且 $a_2$ 整除其余所有条目。
\item 证明存在形如
\[
\begin{pmatrix}
D&0\\
0&0
\end{pmatrix}
\]
的关系矩阵，其中 $D$ 为对角矩阵，非零对角元为 $a_1,a_2,\ldots,a_k$，$k\leq n$，并满足
\[
a_1\mid a_2\mid\cdots\mid a_k.
\]
推出
\[
M\cong R/(a_1)\oplus R/(a_2)\oplus\cdots\oplus R/(a_k)\oplus R^{n-k}.
\]
\end{enumerate}

若 $n$ 不是生成 $M$ 所需的最小生成元个数，则上述某些初始 $a_i$ 可能是单位；对应直和项为零。去掉这些无关因子即可得到 $M$ 的不变因子。并且，在上面的约化过程中跟踪列变换，就可得到不变因子分解中各循环子模的 $R$-生成元。

\item 设 $R$ 为整环，$F$ 为其分式域，$M$ 为任意 $R$-模。证明 $M$ 的秩等于 $F\otimes_R M$ 作为 $F$-向量空间的维数。

\item 证明 P.I.D. 上的有限生成模是投射模当且仅当它是自由模。

\item 设 $R$ 为不是域的 P.I.D.。证明没有有限生成 $R$-模是内射模。提示：用第 10.5 节练习 4 分别考虑挠模与自由模。
\end{enumerate}



\section{有理标准形}

现在把上一节的结果应用到 P.I.D. 为 $F[x]$ 的特殊情形，其中 $F[x]$ 是域 $F$ 上关于 $x$ 的多项式环。

令 $V$ 是 $F$ 上维数为 $n$ 的有限维向量空间，$T$ 是 $V$ 的一个固定线性变换（即 $V$ 到自身的线性变换）。正如第 10 章所见，可以把 $V$ 看作 $F[x]$-模，其中元素 $x$ 在 $V$ 上按线性变换 $T$ 作用（于是 $x$ 中任意多项式按 $T$ 中相同多项式作用）。由于 $V$ 按假设是 $F$ 上有限维的，所以按定义它是有限生成 $F$-模，因而当然是有限生成 $F[x]$-模；所以前一节的分类定理适用。

任何非零自由 $F[x]$-模（同构于若干个 $F[x]$ 的直和）作为 $F$-向量空间都是无限维的，因此如果 $V$ 在 $F$ 上有限维，那么它作为 $F[x]$-模事实上必为挠模（即自由秩为 $0$）。于是由基本定理，$V$ 作为 $F[x]$-模同构于有限多个循环挠 $F[x]$-模的直和。我们将看到，这个 $V$ 的分解允许我们选取 $V$ 的一组基，使得线性变换 $T$ 的矩阵表示具有一种特定的简单形式。当使用 $V$ 的不变因子分解时，得到 $T$ 的矩阵的有理标准形；本节分析这种形式。当使用初等因子分解（并且 $F$ 含有 $T$ 的所有特征值）时，得到下一节讨论的 Jordan 标准形，也就是前面提到的尽可能接近对角矩阵的 $T$ 的矩阵。基本定理的唯一性部分保证有理标准形和 Jordan 标准形是唯一的（因此称为标准形）。

这些标准形的一个重要用途，是分类 $V$ 的不同线性变换。特别地，它们使我们能够判定两个矩阵何时表示同一个线性变换，即判定给定的两个 $n\times n$ 矩阵何时相似。

注意，这又是被作用空间的结构（例如 $V$ 的不变因子分解）用来获得作用在其上的代数对象（这里是线性变换）的重要信息的一个例子。第 18 章将在群作用在向量空间上的情形中讨论这一点（这称为群的表示论）。

在详细描述有理标准形之前，我们先引入一些线性代数概念。

\noindent\textbf{定义.}
\begin{enumerate}[label=(\arabic*)]
\item 如果存在非零向量 $v\in V$ 使 $T(v)=\lambda v$，则 $F$ 中元素 $\lambda$ 称为线性变换 $T$ 的\emph{特征值}。在这种情形下，$v$ 称为 $T$ 的一个对应于特征值 $\lambda$ 的\emph{特征向量}。
\item 如果 $A$ 是一个系数在 $F$ 中的 $n\times n$ 矩阵，$\lambda\in F$，并且存在非零 $n\times1$ 列向量 $v$ 使 $Av=\lambda v$，则称 $\lambda$ 为 $A$ 的特征值，$v$ 为对应特征向量。
\item 如果 $\lambda$ 是线性变换 $T$ 的特征值，则集合
\[
\{v\in V\mid T(v)=\lambda v\}
\]
称为 $T$ 对应于特征值 $\lambda$ 的\emph{特征空间}。类似地，如果 $\lambda$ 是 $n\times n$ 矩阵 $A$ 的特征值，则满足 $Av=\lambda v$ 的 $n\times1$ 矩阵 $v$ 的集合称为 $A$ 对应于特征值 $\lambda$ 的特征空间。
\end{enumerate}

注意，如果固定 $V$ 的一组基 $\mathcal B$，则 $V$ 的任意线性变换 $T$ 都有一个对应的 $n\times n$ 矩阵 $A$。反过来，若 $A$ 是任意 $n\times n$ 矩阵，则由 $T(v)=Av$ 定义的映射（右边的 $v$ 指 $v$ 在固定基 $\mathcal B$ 下的坐标列向量）是 $V$ 的线性变换。于是 $v$ 是 $T$ 的特征值 $\lambda$ 的特征向量，当且仅当 $v$ 在 $\mathcal B$ 下的坐标列向量是 $A$ 的特征值 $\lambda$ 的特征向量。换言之，线性变换 $T$ 的特征值与 $T$ 在任意固定基下矩阵 $A$ 的特征值相同。

\noindent\textbf{定义.} $V$ 到自身的线性变换的行列式定义为表示该线性变换的任意矩阵的行列式（注意这与所选基无关）。

\noindent\textbf{命题 12.} 下列条件等价：
\begin{enumerate}[label=(\arabic*)]
\item $\lambda$ 是 $T$ 的特征值；
\item $\lambda I-T$ 是 $V$ 的奇异线性变换；
\item $\det(\lambda I-T)=0$。
\end{enumerate}

\begin{proof}
$\lambda$ 是 $T$ 的特征值且有对应特征向量 $v$，当且仅当 $v$ 是 $\lambda I-T$ 的核中的非零向量。因此 (1) 与 (2) 等价。由关于行列式的结果，线性变换奇异当且仅当其行列式为零，故 (2) 与 (3) 等价。
\end{proof}

\noindent\textbf{定义.} 令 $x$ 为 $F$ 上的不定元。多项式
\[
\det(xI-T)
\]
称为 $T$ 的\emph{特征多项式}，记为 $c_T(x)$。如果 $A$ 是系数在 $F$ 中的 $n\times n$ 矩阵，则 $\det(xI-A)$ 称为 $A$ 的特征多项式，记为 $c_A(x)$。

由展开行列式容易看出，$T$ 或 $A$ 的特征多项式都是次数 $n=\dim V$ 的首一多项式。命题 12 说明 $T$（或 $A$）的特征值集合正是 $T$（或 $A$）特征多项式的根的集合。特别地，$T$ 至多有 $n$ 个不同的特征值。

我们已经看到，$V$ 通过线性变换 $T$ 看作 $F[x]$-模时是挠 $F[x]$-模。令 $m(x)\in F[x]$ 是生成 $F[x]$ 中 $\operatorname{Ann}(V)$ 的唯一首一多项式。等价地，$m(x)$ 是使 $V$ 被零化的最低次数首一多项式（即 $m(T)$ 是零线性变换），并且若 $f(x)\in F[x]$ 是任意零化 $V$ 的多项式，则 $m(x)$ 整除 $f(x)$。由于选择 $V$ 的一组基可给出所有 $n\times n$ 矩阵环与 $V$ 到自身线性变换集合的同构，对 $F$ 上任意 $n\times n$ 矩阵 $A$，同样存在唯一首一最低次数多项式 $m(x)$ 使 $m(A)$ 为零矩阵。

\noindent\textbf{定义.} 生成 $F[x]$ 中理想 $\operatorname{Ann}(V)$ 的唯一首一多项式称为 $T$ 的\emph{最小多项式}，记为 $m_T(x)$。在矩阵 $A$ 上取值为零矩阵的最低次数首一多项式称为 $A$ 的\emph{最小多项式}，记为 $m_A(x)$。

容易看出（参见练习 5），这些最小多项式的次数至多为 $n^2$，其中 $n=\dim V$。我们很快将证明 $T$ 的最小多项式整除 $T$ 的特征多项式（这是 Cayley--Hamilton 定理），矩阵情形同理；所以实际上这些多项式的次数至多为 $n$。

现在描述线性变换 $T$（以及矩阵 $A$）的有理标准形。由定理 5，有 $F[x]$-模同构
\[
V\cong F[x]/(a_1(x))\oplus F[x]/(a_2(x))\oplus\cdots\oplus F[x]/(a_m(x)),\tag{12.1}
\]
其中 $a_1(x),a_2(x),\ldots,a_m(x)$ 是 $F[x]$ 中次数至少为 $1$ 的多项式，并满足整除条件
\[
a_1(x)\mid a_2(x)\mid\cdots\mid a_m(x).
\]
这些不变因子只在差一个 $F[x]$ 中单位的意义下确定；由于 $F[x]$ 的单位恰是 $F$ 的非零元素（非零常数多项式），我们可规定这些多项式为首一，从而使它们唯一。

由于 $V$ 的零化子是理想 $(a_m(x))$（定理 5(3)），立即得到：

\noindent\textbf{命题 13.} 最小多项式 $m_T(x)$ 是 $V$ 的最大不变因子。$V$ 的所有不变因子都整除 $m_T(x)$。

我们下面将看到如何计算不仅 $T$ 的最小多项式，也计算其他不变因子。

现在为 (12.1) 中 $V$ 的每个直和因子选取一组基，使得 $T$ 的矩阵非常简单。回忆：等式 (12.1) 左侧的线性变换 $T$ 就是右侧每个因子上元素 $x$ 的乘法作用。

在第 11 章命题 1 后的例子中已经看到，若
\[
a(x)=x^k+b_{k-1}x^{k-1}+\cdots+b_1x+b_0
\]
是 $F[x]$ 中任意首一多项式，并令 $\bar x=x\bmod(a(x))$，则元素
\[
1,\bar x,\bar x^2,\ldots,\bar x^{k-1}
\]
给出 $F[x]/(a(x))$ 的 $F$-向量空间基。关于这组基，乘以 $x$ 的线性变换按如下方式作用：
\[
1\mapsto\bar x,
\quad \bar x\mapsto\bar x^2,
\quad\ldots,
\quad \bar x^{k-2}\mapsto\bar x^{k-1},
\]
而
\[
\bar x^{k-1}\mapsto \bar x^k=-b_0-b_1\bar x-\cdots-b_{k-1}\bar x^{k-1},
\]
最后一个等式来自 $a(\bar x)=0$。因此关于这组基，乘以 $x$ 的矩阵是
\[
\begin{pmatrix}
0&0&0&\cdots&0&-b_0\\
1&0&0&\cdots&0&-b_1\\
0&1&0&\cdots&0&-b_2\\
\vdots&\vdots&\vdots& &\vdots&\vdots\\
0&0&0&\cdots&1&-b_{k-1}
\end{pmatrix}.
\]

\noindent\textbf{定义.} 上述矩阵称为首一多项式 $a(x)$ 的\emph{伴随矩阵}，记为 $C(a)$。

\noindent\textbf{定义.} 若一个块对角矩阵的对角块为首一多项式 $a_1(x),\ldots,a_m(x)$ 的伴随矩阵 $C(a_1),\ldots,C(a_m)$，并且
\[
a_1(x)\mid a_2(x)\mid\cdots\mid a_m(x),
\]
则称该矩阵处于\emph{有理标准形}。

\noindent\textbf{定理 14（线性变换的有理标准形）.} 设 $V$ 是域 $F$ 上有限维向量空间，$T$ 是 $V$ 上线性变换。则存在 $V$ 的一组基，使得 $T$ 关于这组基的矩阵处于有理标准形。这个有理标准形由 $V$ 作为 $F[x]$-模的不变因子唯一确定。

\begin{proof}
把 $V$ 看作由 $T$ 定义的 $F[x]$-模。由定理 5，$V$ 分解为循环挠模 $F[x]/(a_i(x))$ 的直和，其中 $a_1\mid\cdots\mid a_m$。在每个因子 $F[x]/(a_i(x))$ 上取基 $1,\bar x,\ldots,\bar x^{d_i-1}$。上面的计算说明，在这一基上元素 $x$ 的作用（即 $T$ 的作用）由伴随矩阵 $C(a_i)$ 给出。把这些基依直和顺序合在一起，得到 $V$ 的一组基；在这组基下 $T$ 的矩阵是这些伴随矩阵的块对角直和，因而处于有理标准形。唯一性来自定理 9 中不变因子的唯一性。
\end{proof}

\noindent\textbf{定理 15.} 设 $S,T$ 是 $V$ 的线性变换。下列条件等价：
\begin{enumerate}[label=(\arabic*)]
\item $S$ 与 $T$ 相似，即存在 $V$ 的一个非奇异线性变换 $P$，使得 $S=P^{-1}TP$；
\item $S$ 与 $T$ 使 $V$ 成为同构的 $F[x]$-模；
\item $S$ 与 $T$ 有相同的不变因子。
\end{enumerate}

\begin{proof}
若 $S=P^{-1}TP$，则 $P$ 给出两个 $F[x]$-模之间的同构：在第一个模中 $x$ 按 $S$ 作用，在第二个模中 $x$ 按 $T$ 作用；等式 $TP=PS$ 正是 $P$ 与 $x$ 的作用可交换。反过来，若存在 $F[x]$-模同构 $P$，则 $P$ 作为 $F$-线性映射非奇异，并满足 $P(Sv)=T(Pv)$，故 $S=P^{-1}TP$。条件 (2) 与 (3) 的等价性正是定理 9 对 P.I.D. $F[x]$ 的应用。
\end{proof}

\noindent\textbf{定理 16（矩阵的有理标准形）.} 设 $A$ 是域 $F$ 上的 $n\times n$ 矩阵。则 $A$ 相似于一个有理标准形矩阵，也就是说，存在 $F$ 上可逆 $n\times n$ 矩阵 $P$，使 $P^{-1}AP$ 处于有理标准形。这个有理标准形唯一。

\begin{proof}
令 $V=F^n$，让 $A$ 定义 $V$ 上的线性变换 $T(v)=Av$。由定理 14 可选 $V$ 的一组基，使 $T$ 的矩阵为有理标准形。若 $P$ 的列是这组新基向量在标准基下的坐标，则 $P^{-1}AP$ 就是该有理标准形。唯一性由不变因子的唯一性给出。
\end{proof}

\noindent\textbf{定义.} 域 $F$ 上 $n\times n$ 矩阵的不变因子，定义为由该矩阵定义的线性变换使 $F^n$ 成为 $F[x]$-模后所得的不变因子。

\noindent\textbf{定理 17.} 设 $A$ 与 $B$ 是域 $F$ 上的 $n\times n$ 矩阵。则 $A$ 与 $B$ 相似当且仅当它们有相同的不变因子。

\begin{proof}
这是定理 15 直接应用到 $V=F^n$ 的矩阵形式。相似矩阵给出同一个线性变换在不同基下的矩阵；而同一不变因子列表也由基本定理唯一地确定相应 $F[x]$-模的同构类，进而确定相似类。
\end{proof}

\noindent\textbf{推论 18.} 设 $A,B$ 是域 $F$ 上的两个 $n\times n$ 矩阵，并设 $F$ 是域 $K$ 的子域。
\begin{enumerate}[label=(\arabic*)]
\item 若 $A$ 与 $B$ 在 $F$ 上相似，则它们在 $K$ 上相似。
\item 若 $A$ 与 $B$ 在 $K$ 上相似，并且二者在 $K[x]$ 中的不变因子都事实上属于 $F[x]$，则 $A$ 与 $B$ 在 $F$ 上相似。
\end{enumerate}

\begin{proof}
(1) 若 $B=P^{-1}AP$ 且 $P$ 的条目在 $F$ 中，则把 $P$ 看作 $K$ 上矩阵同样给出 $K$ 上相似。 (2) 若在 $K$ 上相似，则由定理 17 它们在 $K[x]$ 中有相同的不变因子。若这些不变因子属于 $F[x]$，则它们也是作为 $F$ 上矩阵的不变因子；再由定理 17，$A$ 与 $B$ 在 $F$ 上相似。
\end{proof}

\noindent\textbf{引理 19.} 设 $a(x)\in F[x]$ 是任意首一多项式，$C(a)$ 是它的伴随矩阵。则 $C(a)$ 的最小多项式和特征多项式都等于 $a(x)$。更一般地，如果 $M$ 是块对角矩阵，对角块为 $A_1,\ldots,A_k$，则 $M$ 的特征多项式为这些 $A_i$ 的特征多项式之积，最小多项式为这些 $A_i$ 的最小多项式的最小公倍式。

\begin{proof}
伴随矩阵 $C(a)$ 表示 $F[x]/(a)$ 中“乘以 $x$”的线性变换；整个循环模的零化子正是 $(a)$，所以最小多项式为 $a$。其特征多项式可直接由 $xI-C(a)$ 的行列式展开得到，也可由后面的命题 20 得到。对于块对角矩阵，行列式按块相乘给出特征多项式的乘积；一个多项式零化整个块对角矩阵当且仅当它零化每个块，所以最小多项式是各块最小多项式的最小公倍式。
\end{proof}

\noindent\textbf{命题 20.} 设 $A$ 的不变因子为 $a_1(x),a_2(x),\ldots,a_m(x)$。则
\[
c_A(x)=a_1(x)a_2(x)\cdots a_m(x).
\]
特别地，$m_A(x)$ 整除 $c_A(x)$。

\begin{proof}
矩阵 $A$ 相似于有理标准形，其对角块是 $C(a_1),\ldots,C(a_m)$。相似矩阵有相同特征多项式；而块对角矩阵的特征多项式为各块特征多项式的乘积，且每个伴随矩阵 $C(a_i)$ 的特征多项式为 $a_i$。因此 $c_A=a_1\cdots a_m$。又最小多项式是最大不变因子 $a_m$，而 $a_m\mid a_1\cdots a_m$，故 $m_A\mid c_A$。
\end{proof}

下面给出计算矩阵不变因子的一种矩阵方法。命题 20 中的关系在确定矩阵 $A$ 的不变因子时经常很有用，特别是对低阶矩阵而言（参见练习 3、4 以及下面的例子）。以下结果依赖上一节练习 16--19；它在一般情形中计算不变因子，其证明在练习中概述。

设 $A$ 是域 $F$ 上的 $n\times n$ 矩阵。于是 $xI-A$ 是一个条目在 $F[x]$ 中的 $n\times n$ 矩阵。下列三种运算称为初等行、列运算：
\begin{enumerate}[label=(\alph*)]
\item 交换两行或两列；
\item 把一行或一列的一个 $F[x]$ 倍加到另一行或另一列上；
\item 把任一行或任一列乘以 $F[x]$ 中的一个单位，也就是乘以 $F$ 中一个非零元素。
\end{enumerate}

\noindent\textbf{定理 21.} 设 $A$ 是域 $F$ 上的 $n\times n$ 矩阵。利用上述三种初等行、列运算，条目在 $F[x]$ 中的矩阵 $xI-A$ 可以化为对角形（称为 $A$ 的 Smith 标准形）
\[
\begin{pmatrix}
1&&&&&\\
&\ddots&&&&\\
&&1&&&\\
&&&a_1(x)&&\\
&&&&\ddots&\\
&&&&&a_m(x)
\end{pmatrix},
\]
其中 $a_1(x),a_2(x),\ldots,a_m(x)$ 是 $F[x]$ 中首一非零元素，次数至少为 $1$，并满足
\[
a_1(x)\mid a_2(x)\mid\cdots\mid a_m(x).
\]
这些元素 $a_1(x),\ldots,a_m(x)$ 正是 $A$ 的不变因子。

\begin{proof}
参见练习。
\end{proof}

\paragraph{不变因子分解算法：化为有理标准形}
如上一节末尾的练习所述，跟踪把 $xI-A$ 对角化所需的运算，可以显式给出一个矩阵 $P$，使得 $P^{-1}AP$ 为有理标准形。等价地，若 $V$ 是一个给定的 $F[x]$-模，并且有向量空间基 $[e_1,e_2,\ldots,e_n]$，那么 $P$ 给出使 $V$ 分解为循环 $F[x]$-模直和的不变因子分解所需的换基。特别地，若 $A$ 是由 $x$ 定义的 $F[x]$-模 $V$ 的线性变换 $T$ 在这组基下的矩阵（即 $T(e_j)=xe_j=\sum_{i=1}^n a_{ij}e_i$，其中 $A=(a_{ij})$），那么矩阵 $P$ 给出 $V$ 的换基，使得 $T$ 的矩阵为有理标准形。

我们先在一般背景下描述算法：给定一个具有向量空间基 $[e_1,e_2,\ldots,e_n]$ 的 $F[x]$-模 $V$，求它的不变因子分解（证明在练习中概述）。然后描述怎样把给定的 $n\times n$ 矩阵 $A$ 化为有理标准形（这时不再显式提及底层向量空间和相应线性变换）。下面例 2 与例 3 给出该算法的具体数值例子。

\paragraph{不变因子分解算法}
设 $V$ 是具有向量空间基 $[e_1,e_2,\ldots,e_n]$ 的 $F[x]$-模（于是这些元素尤其是 $V$ 作为 $F[x]$-模的生成元）。设 $T$ 是由 $x$ 定义的 $V$ 到自身的线性变换，并设 $A$ 是 $T$ 关于 $V$ 的这组基的 $n\times n$ 矩阵，即
\[
T(e_j)=xe_j=\sum_{i=1}^n a_{ij}e_i,
\qquad A=(a_{ij}).
\]
\begin{enumerate}[label=(\arabic*)]
\item 使用以下三种初等行、列运算，在 $F[x]$ 上把矩阵 $xI-A$ 对角化，并记录所用的行运算：
\begin{enumerate}[label=(\alph*)]
\item 交换两行或两列（第 $i$ 行与第 $j$ 行交换记为 $R_i\leftrightarrow R_j$；列同理记为 $C_i\leftrightarrow C_j$）；
\item 把一行或一列的一个 $F[x]$ 倍加到另一行或列上（若把 $p(x)$ 倍第 $j$ 行加到第 $i$ 行，记为 $R_i+p(x)R_j\mapsto R_i$；列同理记为 $C_i+p(x)C_j\mapsto C_i$）；
\item 把任一行或任一列乘以 $F[x]$ 中的一个单位，也就是乘以 $F$ 中的一个非零元素（若把第 $i$ 行乘以 $u\in F^\times$，记为 $uR_i$；列同理记为 $uC_i$）。
\end{enumerate}
\item 从 $F[x]$-模生成元 $[e_1,e_2,\ldots,e_n]$ 出发，对 (1) 中使用的每一个行运算，按下列规则改变生成元集合：
\begin{enumerate}[label=(\alph*)]
\item 若第 $i$ 行与第 $j$ 行交换，则交换第 $i$ 个与第 $j$ 个生成元；
\item 若把 $p(x)$ 倍第 $j$ 行加到第 $i$ 行，则从第 $j$ 个生成元中减去 $p(x)$ 倍第 $i$ 个生成元（注意指标）；
\item 若第 $i$ 行乘以单位 $u\in F$，则把第 $i$ 个生成元除以 $u$。
\end{enumerate}
\item 当 $xI-A$ 已经被对角化为定理 21 中的形式时，$V$ 的生成元 $[e_1,e_2,\ldots,e_n]$ 已经成为原来 $e_1,e_2,\ldots,e_n$ 的 $F[x]$-线性组合。利用 $xe_j=T(e_j)=\sum_{i=1}^n a_{ij}e_i$，把这些元素写成 $e_1,e_2,\ldots,e_n$ 的 $F$-线性组合。对角化以后，前 $n-m$ 个这样的线性组合为 $0$（这给计算提供一个有用的数值检验），其余 $m$ 个线性组合非零，也就是说，$V$ 的生成元具有形式
\[
[0,\ldots,0,f_1,\ldots,f_m],
\]
并且它们精确对应于定理 21 中的对角元素。元素 $f_1,\ldots,f_m$ 是 $V$ 的不变因子分解中循环因子的 $F[x]$-模生成元（零化子分别为 $(a_1(x)),\ldots,(a_m(x))$）：
\[
V=F[x]f_1\oplus F[x]f_2\oplus\cdots\oplus F[x]f_m,
\qquad
F[x]f_i\cong F[x]/(a_i(x)),
\]
其中 $i=1,2,\ldots,m$。这给出 $F[x]$-模 $V$ 的不变因子分解。
\item 每个循环因子的相应向量空间基由元素
\[
f_i,Tf_i,T^2f_i,\ldots,T^{\deg a_i(x)-1}f_i
\]
给出。
\item 把 (4) 中算出的向量空间基的第 $k$ 个元素写成原向量空间基 $[e_1,e_2,\ldots,e_n]$ 的线性组合，并把所得坐标作为一个 $n\times n$ 矩阵 $P$ 的第 $k$ 列。则 $P^{-1}AP$ 处于有理标准形（对角块为多项式 $a_i(x)$ 的伴随矩阵）。这是线性变换 $T$ 关于 (4) 中向量空间基的矩阵。
\end{enumerate}

现在描述把给定 $n\times n$ 矩阵 $A$ 化为有理标准形的算法，即求 $n\times n$ 矩阵 $P$，使 $P^{-1}AP$ 处于有理标准形。这不过是把上面的算法应用到向量空间 $V=F^n$（$n\times1$ 列向量空间）及其标准基 $[e_1,e_2,\ldots,e_n]$，并令 $T$ 为由 $A$ 和这组基定义的线性变换。下面省略对底层向量空间与相应线性变换的显式提及，所以算法纯粹是矩阵论的。

\paragraph{把 $n\times n$ 矩阵化为有理标准形}
设 $A$ 是域 $F$ 上的 $n\times n$ 矩阵。
\begin{enumerate}[label=(\arabic*)]
\item 使用上述三种初等行、列运算，在 $F[x]$ 上把矩阵 $xI-A$ 对角化，并记录所用的行运算。定义 $d_1,\ldots,d_m$ 分别为对角线上出现的首一非常数多项式 $a_1(x),\ldots,a_m(x)$ 的次数。
\item 从 $n\times n$ 单位矩阵 $P'$ 出发，对 (1) 中使用的每一个行运算，按以下规则改变矩阵 $P'$：
\begin{enumerate}[label=(\alph*)]
\item 若 $R_i\leftrightarrow R_j$，则交换 $P'$ 的第 $i$ 列与第 $j$ 列（即对 $P'$ 作 $C_i\leftrightarrow C_j$）；
\item 若 $R_i+p(x)R_j\mapsto R_i$，则从 $P'$ 的第 $j$ 列中减去矩阵 $p(A)$ 乘以 $P'$ 的第 $i$ 列所得的列向量（即对 $P'$ 作 $C_j-p(A)C_i\mapsto C_j$；注意指标）；
\item 若 $uR_i$，则把 $P'$ 的第 $i$ 列的所有元素除以 $u$（即对 $P'$ 作 $u^{-1}C_i$）。
\end{enumerate}
\item 当 $xI-A$ 被对角化为定理 21 中的形式时，矩阵 $P'$ 的前 $n-m$ 列为 $0$（这给计算提供一个有用的数值检验），其余 $m$ 列非零。对每个 $i=1,2,\ldots,m$，把 $P'$ 的第 $i$ 个非零列向量依次乘以 $A^0=I,A^1,A^2,\ldots,A^{d_i-1}$，并按此顺序把所得列向量作为 $n\times n$ 矩阵 $P$ 接下来的 $d_i$ 列。则 $P^{-1}AP$ 处于有理标准形（其对角块为 (1) 中多项式 $a_1(x),\ldots,a_m(x)$ 的伴随矩阵）。
\end{enumerate}

在线性变换（或矩阵）的标准形理论中，特征多项式扮演有限阿贝尔群的阶的角色，而最小多项式扮演指数的角色（毕竟它们是同一类不变量：一个是主理想整环 $\mathbb Z$ 上模的不变量，另一个是主理想整环 $F[x]$ 上模的不变量），因此我们可以直接解决与第 5 章中有限阿贝尔群问题相对应的问题。特别包括：
\begin{enumerate}[label=(\Alph*)]
\item 确定给定矩阵的有理标准形（类似于把有限阿贝尔群分解为循环群直积）；
\item 判定两个给定矩阵是否相似（类似于判定两个给定有限阿贝尔群是否同构）；
\item 确定域 $F$ 上具有给定特征多项式的矩阵的所有相似类（类似于确定给定阶的所有阿贝尔群）；
\item 确定域 $F$ 上具有给定最小多项式的所有 $n\times n$ 矩阵相似类（类似于确定给定指数、秩至多为 $n$ 的所有阿贝尔群）。
\end{enumerate}

\paragraph{例}
\begin{enumerate}[label=(\arabic*)]
\item 我们求下列有理矩阵的有理标准形，并判定它们是否相似：
\[
A=\begin{pmatrix}2&-2&14\\0&3&-7\\0&0&2\end{pmatrix},\qquad
B=\begin{pmatrix}1&-44&-14\\0&85&28\\0&0&3\end{pmatrix},\qquad
C=\begin{pmatrix}1&2&2\\0&2&-1\\0&0&3\end{pmatrix}.
\]
直接计算表明三个矩阵有相同的特征多项式
\[
c_A(x)=c_B(x)=c_C(x)=(x-2)^2(x-3).
\]
由于最小多项式与特征多项式有相同根，最小多项式只可能是 $(x-2)(x-3)$ 或 $(x-2)^2(x-3)$。很快可算得
\[
(A-2I)(A-3I)=0,
\]
而 $(B-2I)(B-3I)\neq0$（其 $(1,1)$ 条目非零），且 $(C-2I)(C-3I)\neq0$（其 $(1,2)$ 条目非零）。因此
\[
m_A(x)=(x-2)(x-3),\qquad m_B(x)=m_C(x)=(x-2)^2(x-3).
\]
于是 $B$ 与 $C$ 不再有其他不变因子。由于 $A$ 的不变因子整除最小多项式且乘积为特征多项式，$A$ 的不变因子为
\[
x-2,\qquad (x-2)(x-3)=x^2-5x+6.
\]
（对 $2\times2$ 和 $3\times3$ 矩阵，特征多项式与最小多项式已足以确定所有不变因子，参见练习 3、4。）因此 $B$ 与 $C$ 相似，而二者都不与 $A$ 相似。有理标准形为（注意 $(x-2)^2(x-3)=x^3-7x^2+16x-12$）
\[
\begin{pmatrix}2&0&0\\0&0&-6\\0&1&5\end{pmatrix},
\qquad
\begin{pmatrix}0&0&12\\1&0&-16\\0&1&7\end{pmatrix},
\qquad
\begin{pmatrix}0&0&12\\1&0&-16\\0&1&7\end{pmatrix}.
\]

\item 上例的有理标准形只是通过确定矩阵的特征多项式与最小多项式得到的。如前所述，对 $2\times2$ 与 $3\times3$ 矩阵这已经足够，因为这些信息足以确定所有不变因子。然而对更大矩阵一般并不足够（参见下一例），还需要更多计算来确定不变因子。本例再次用上面两种算法计算例 (1) 中矩阵 $A$ 的有理标准形。对这个小矩阵来说这在计算上更麻烦（这一点会很明显），但即使在此情形它也有一个优点：它还显式算出矩阵 $P$，使 $P^{-1}AP$ 为有理标准形。

\noindent I.（不变因子分解）我们在 $\mathbb Q[x]$ 中用行、列运算把矩阵
\[
xI-A=\begin{pmatrix}
x-2&2&-14\\
0&x-3&7\\
0&0&x-2
\end{pmatrix}
\]
化为对角形。采用不变因子分解算法中的记号，$R_i\leftrightarrow R_j$ 表示交换第 $i$ 行与第 $j$ 行；若把 $a$ 倍第 $j$ 行加到第 $i$ 行，记作 $R_i+aR_j\mapsto R_i$；第 $i$ 行乘以 $u$ 简记作 $uR_i$（列运算同理，用 $C$ 代替 $R$）。在下面的化简中，最开始两个运算相当临时，只是为了使计算中各处都是整数。该化简把 $xI-A$ 化为对角矩阵，其非单位对角元为
\[
x-2,\qquad x^2-5x+6,
\]
因而确定了该矩阵的不变因子，这与例 (1) 所得一致。令 $V$ 为 $\mathbb Q$ 上三维向量空间，基为 $e_1,e_2,e_3$，并令 $T$ 为相应线性变换（它定义 $x$ 在 $V$ 上的作用），即
\[
xe_1=T(e_1)=2e_1,
\qquad
xe_2=T(e_2)=-2e_1+3e_2,
\qquad
xe_3=T(e_3)=14e_1-7e_2+2e_3.
\]
化简中用到的行运算为
\[
R_1+R_2\mapsto R_1,
\quad -R_1,
\quad R_2+(x-3)R_1\mapsto R_2,
\quad R_2-7R_3\mapsto R_2,
\quad R_2\leftrightarrow R_3.
\]
从 $V$ 的基 $[e_1,e_2,e_3]$ 出发，并按正文中的规则改变它，得到
\[
\begin{aligned}
[e_1,e_2,e_3]
&\mapsto [e_1,e_2-e_1,e_3]
\mapsto [-e_1,e_2-e_1,e_3]\\
&\mapsto [-e_1-(x-3)(e_2-e_1),e_2-e_1,e_3]\\
&\mapsto [-e_1-(x-3)(e_2-e_1),e_2-e_1,e_3+7(e_2-e_1)]\\
&\mapsto [-e_1-(x-3)(e_2-e_1),e_3+7(e_2-e_1),e_2-e_1].
\end{aligned}
\]
利用上面对 $x$ 作用的公式，可见这些最后的元素是
\[
[0,-7e_1+7e_2+e_3,-e_1+e_2],
\]
分别对应于 $xI-A$ 对角化形式中的元素 $1,x-2,x^2-5x+6$。因此
\[
f_1=-7e_1+7e_2+e_3,
\qquad
f_2=-e_1+e_2
\]
是 $V$ 作为 $\mathbb Q[x]$-模的不变因子分解中两个循环因子的 $\mathbb Q[x]$-模生成元。相应的 $\mathbb Q$-向量空间基为 $f_1$ 以及 $f_2,xf_2=Tf_2$，即
\[
-7e_1+7e_2+e_3,
\qquad
-e_1+e_2,
\qquad
T(-e_1+e_2)=-4e_1+3e_2.
\]
于是矩阵
\[
P=\begin{pmatrix}
-7&-1&-4\\
7&1&3\\
1&0&0
\end{pmatrix}
\]
把 $A$ 共轭到其有理标准形：
\[
P^{-1}AP=\begin{pmatrix}2&0&0\\0&0&-6\\0&1&5\end{pmatrix},
\]
这很容易直接检验。

\noindent II.（直接把 $A$ 化为有理标准形）利用对角化 $xI-A$ 时使用的行运算，按上面的算法确定矩阵 $P'$。这里 $d_1=1,d_2=2$，分别对应于 $P'$ 的第二与第三个非零列。因此 $P$ 的列分别为
\[
\begin{pmatrix}-7\\7\\1\end{pmatrix},
\qquad
\begin{pmatrix}-1\\1\\0\end{pmatrix},
\qquad
A\begin{pmatrix}-1\\1\\0\end{pmatrix}=\begin{pmatrix}-4\\3\\0\end{pmatrix},
\]
从而再次得到上面的矩阵 $P$。

\item 对 $3\times3$ 矩阵 $A$ 来说，仅为确定有理标准形（等价地，不变因子）而执行上面冗长计算并非必要，正如例 (1) 所示。然而对 $n\times n$ 矩阵，若 $n\geq4$，一般来说特征多项式与最小多项式的计算并不足以确定所有不变因子，因此前一例中更广泛的计算可能变得必要。比如考虑矩阵
\[
D=\begin{pmatrix}
1&0&0&-4\\
2&-1&0&-4\\
1&1&0&-2\\
0&0&1&3
\end{pmatrix}.
\]
简短计算表明 $D$ 的特征多项式为 $(x-1)^4$。可能的最小多项式为 $x-1,(x-1)^2,(x-1)^3,(x-1)^4$。显然 $D-I\neq0$，另一次简短计算给出 $(D-I)^2=0$，所以 $D$ 的最小多项式为 $(x-1)^2$。于是有两种可能的不变因子集合：
\[
x-1,x-1,(x-1)^2
\qquad\text{和}\qquad
(x-1)^2,(x-1)^2.
\]
为确定 $D$ 的不变因子，把前一例的方法应用于 $4\times4$ 矩阵
\[
xI-D=\begin{pmatrix}
x-1&0&0&4\\
-2&x+1&0&4\\
-1&-1&x&2\\
0&0&-1&x-3
\end{pmatrix}.
\]
由初等行、列运算从此矩阵得到的对角矩阵为
\[
\begin{pmatrix}
1&0&0&0\\
0&1&0&0\\
0&0&(x-1)^2&0\\
0&0&0&(x-1)^2
\end{pmatrix},
\]
这说明 $D$ 的不变因子为 $(x-1)^2,(x-1)^2$。一组把 $xI-D$ 对角化的初等行、列运算为
\[
\begin{gathered}
R_1\leftrightarrow R_3,
\quad -R_1,
\quad R_2+2R_1\mapsto R_2,
\quad R_3-(x-1)R_1\mapsto R_3,\\
C_3+(x-1)C_1\mapsto C_3,
\quad C_4+2C_1\mapsto C_4,
\quad R_2\leftrightarrow R_4,
\quad -R_2,\\
R_3+2R_2\mapsto R_3,
\quad R_4-(x+1)R_2\mapsto R_4,
\quad C_3+2C_2\mapsto C_3,
\quad C_4+(x-3)C_2\mapsto C_4.
\end{gathered}
\]
I.（不变因子分解）若在此情形中 $e_1,e_2,e_3,e_4$ 是 $V$ 的一组基，则像前一例一样使用这些行运算可见，对应于上述因子的 $V$ 的生成元为
\[
(x-1)e_1-2e_2-e_3=0,
\qquad
-2e_1+(x+1)e_2-e_4=0,
\qquad
e_1,
\qquad
e_2.
\]
因此在此情形中 $V$ 的不变因子分解中两个直因子的向量空间基为 $e_1,Te_1$ 以及 $e_2,Te_2$，其中 $T$ 是由 $D$ 定义的线性变换，即
\[
e_1,
\quad e_1+2e_2+e_3,
\quad e_2,
\quad 2e_1-e_2+e_4.
\]
相应的换基矩阵为
\[
P=\begin{pmatrix}
1&1&0&2\\
0&2&1&-1\\
0&1&0&0\\
0&0&0&1
\end{pmatrix},
\]
使得 $P^{-1}DP$ 为有理标准形，这也可直接检验。

II.（直接把 $D$ 化为有理标准形）与例 (2) 类似，由对角化 $xI-D$ 所用的行运算确定算法中的矩阵 $P'$。这里 $d_1=2,d_2=2$，分别对应于 $P'$ 的第三与第四个非零列。把这些列分别乘以 $I,D$，即可再次得到上面的矩阵 $P$。

\item 本例确定所有系数在 $\mathbb Q$ 中、特征多项式为 $(x^4-1)(x^2-1)$ 的矩阵 $A$ 的相似类。首先注意，具有次数 $6$ 的特征多项式的矩阵必为 $6\times6$ 矩阵。多项式 $(x^4-1)(x^2-1)$ 在 $\mathbb Q[x]$ 中分解为
\[
(x-1)^2(x+1)^2(x^2+1).
\]
由于 $A$ 的最小多项式 $m_A(x)$ 与 $c_A(x)$ 有相同根，可知 $(x-1)(x+1)(x^2+1)$ 整除 $m_A(x)$。设某个 $A$ 的不变因子为 $a_1(x),\ldots,a_m(x)$，于是 $a_m(x)=m_A(x)$，$a_i(x)\mid a_{i+1}(x)$（特别地，所有不变因子都整除 $m_A(x)$），且
\[
a_1(x)a_2(x)\cdots a_m(x)=(x^4-1)(x^2-1).
\]
容易看出，在这些限制下唯一允许的列表为
\begin{enumerate}[label=(\alph*)]
\item $(x-1)(x+1),\ (x-1)(x+1)(x^2+1)$；
\item $x-1,\ (x-1)(x+1)^2(x^2+1)$；
\item $x+1,\ (x-1)^2(x+1)(x^2+1)$；
\item $(x-1)^2(x+1)^2(x^2+1)$。
\end{enumerate}
现在可以容易写出相应的伴随矩阵直和，从而得到 $4$ 个相似类的代表。下一节将看到即使在 $M_6(\mathbb C)$ 中仍然只有 $4$ 个相似类。

\item 本例求所有满足 $A^6=I$ 且条目在 $\mathbb Q$ 中的 $3\times3$ 矩阵 $A$ 的相似类。对每个这样的 $A$，其最小多项式整除 $x^6-1$，而在 $\mathbb Q[x]$ 中该多项式的完全分解为
\[
x^6-1=(x-1)(x+1)(x^2-x+1)(x^2+x+1).
\]
反过来，若 $B$ 是任意 $3\times3$ 矩阵且其最小多项式整除 $x^6-1$，则 $B^6=I$。对 $B$ 的最小多项式的唯一限制是其次数至多为 $3$（由 Cayley--Hamilton 定理）。因此这种矩阵 $A$ 的最小多项式的唯一可能为
\begin{enumerate}[label=(\alph*)]
\item $x-1$；
\item $x+1$；
\item $x^2-x+1$；
\item $x^2+x+1$；
\item $(x-1)(x+1)$；
\item $(x-1)(x^2-x+1)$；
\item $(x-1)(x^2+x+1)$；
\item $(x+1)(x^2-x+1)$；
\item $(x+1)(x^2+x+1)$。
\end{enumerate}
在有理标准形的限制下，它们给出以下允许的不变因子列表：
\begin{enumerate}[label=(\roman*)]
\item $x-1,x-1,x-1$；
\item $x+1,x+1,x+1$；
\item $x-1,(x-1)(x+1)$；
\item $x+1,(x-1)(x+1)$；
\item $(x-1)(x^2-x+1)$；
\item $(x-1)(x^2+x+1)$；
\item $(x+1)(x^2-x+1)$；
\item $(x+1)(x^2+x+1)$。
\end{enumerate}
注意，如果最小多项式为 $x^2+x+1$ 或 $x^2-x+1$，则不存在合适的不变因子集合。现在可以写出相应的有理标准形；例如 (i) 是 $I$，(ii) 是 $-I$，而 (iii) 是
\[
\begin{pmatrix}
1&0&0\\
0&0&1\\
0&1&0
\end{pmatrix}.
\]
还要注意，这个结果也可表述为：任意条目在 $\mathbb Q$ 中且阶（当然是乘法阶）整除 $6$ 的 $3\times3$ 矩阵，都相似于这 $8$ 个矩阵之一；所以本例确定了群 $GL_3(\mathbb Q)$ 中阶为 $1,2,3,6$ 的所有元素（差一个相似）。
\end{enumerate}

\subsection*{练习}

\begin{enumerate}[label=\arabic*.]
\item 证明相似的线性变换（或 $n\times n$ 矩阵）有相同的特征多项式和相同的最小多项式。

\item 设 $M$ 如引理 19。证明 $M$ 的最小多项式是 $A_1,
\ldots,A_k$ 的最小多项式的最小公倍式。

\item 证明 $F$ 上两个不是标量矩阵的 $2\times2$ 矩阵相似当且仅当它们有相同的特征多项式。

\item 证明两个 $3\times3$ 矩阵相似当且仅当它们有相同的特征多项式和相同的最小多项式。给出一个 $4\times4$ 矩阵的明确反例，说明此命题对 $4\times4$ 矩阵不成立。

\item 直接利用 $n$ 维 $F$-向量空间 $V$ 到自身的全体线性变换构成维数为 $n^2$ 的 $F$-向量空间这一事实，证明线性变换 $T$ 的最小多项式次数至多为 $n^2$。

\item 证明 $n\times n$ 矩阵 $A$ 的特征多项式的常数项为 $(-1)^n\det A$，$x^{n-1}$ 的系数为 $A$ 的对角线元素之和的负数。对角线元素之和称为 $A$ 的迹，记为 $\operatorname{tr}(A)$。证明 $\det A$ 是 $A$ 的特征值之积，$\operatorname{tr}(A)$ 是 $A$ 的特征值之和。

\item 求矩阵
\[
\begin{pmatrix}
0&1&0&0\\
0&0&1&0\\
0&0&0&1\\
1&0&0&0
\end{pmatrix}
\]
的特征值。

\item 验证伴随矩阵
\[
\begin{pmatrix}
0&0&0&\cdots&0&-a_0\\
1&0&0&\cdots&0&-a_1\\
0&1&0&\cdots&0&-a_2\\
\vdots&\vdots&\vdots& &\vdots&\vdots\\
0&0&0&\cdots&1&-a_{n-1}
\end{pmatrix}
\]
的特征多项式为
\[
x^n+a_{n-1}x^{n-1}+\cdots+a_1x+a_0.
\]

\item 求下列矩阵的有理标准形：
\[
\begin{pmatrix}0&-1&-1\\0&0&0\\-1&0&0\end{pmatrix},\qquad
\begin{pmatrix}c&0&-1\\0&c&1\\-1&1&c\end{pmatrix},\qquad
\begin{pmatrix}
422&465&15&-30\\
-420&-463&-15&30\\
840&930&32&-60\\
-140&-155&-5&12
\end{pmatrix}.
\]

\item 求 $\mathbb Q$ 上所有最小多项式为 $(x+2)^2(x-1)$ 的 $6\times6$ 矩阵相似类。给出所有不变因子列表并写出若干相应矩阵即可。

\item 求 $\mathbb C$ 上所有特征多项式为 $(x^4-1)(x^2-1)$ 的 $6\times6$ 矩阵相似类。

\item 求 $\mathbb F_2$ 上所有满足 $A^6=I$ 的 $3\times3$ 矩阵 $A$ 的相似类（与正文中在 $\mathbb Q$ 上的答案比较）。同样，求所有满足 $B^{20}=I$ 的 $4\times4$ 矩阵 $B$ 的相似类。

\item 证明：在 $\mathbb Q[x]$ 中给定特征多项式时，$\mathbb Q$ 上 $3\times3$ 矩阵的相似类数等于任意 $\mathbb Q$ 的扩域上的相似类数。给出例子说明此结论对 $4\times4$ 矩阵一般不成立。

\item 求特征多项式为 $x^2(x^2+1)^2$ 的线性变换的所有可能有理标准形。

\item 在相似意义下求所有精确阶为 $4$ 的 $2\times2$ 有理矩阵（即 $M_2(\mathbb Q)$ 中的矩阵，阶指乘法阶）。若矩阵条目来自 $\mathbb C$，也做同样问题。

\item 证明在 $\mathbb F_{19}[x]$ 中
\[
x^5-1=(x-1)(x^2-4x+1)(x^2+5x+1).
\]
用此在相似意义下求所有 $\mathbb F_{19}$ 上乘法阶为 $5$ 的 $2\times2$ 矩阵。

\item 求 $GL_3(\mathbb F_2)$ 的共轭类代表元。提示：与第 6 章定理 15 和命题 14 比较。

\item 设 $V$ 是 $\mathbb Q$ 上有限维向量空间，$T$ 是 $V$ 上非奇异线性变换，并满足
\[
T^{-1}=T^2+T.
\]
证明 $\dim V$ 被 $3$ 整除。若 $\dim V=3$，证明所有这样的变换 $T$ 都相似。

\item 令 $V$ 为无限维实向量空间
\[
\mathbb R^\infty=\{(a_0,a_1,a_2,
\ldots)\mid a_0,a_1,a_2,
\ldots\in\mathbb R\}.
\]
定义 $T:V\to V$ 为
\[
T(a_0,a_1,a_2,
\ldots)=(0,a_0,a_1,a_2,
\ldots).
\]
证明 $T$ 没有特征向量。

\item 设 $\ell$ 为素数，
\[
\Phi_\ell(x)=\frac{x^\ell-1}{x-1}=x^{\ell-1}+x^{\ell-2}+\cdots+x+1\in\mathbb Z[x]
\]
为第 $\ell$ 个分圆多项式，它在 $\mathbb Q$ 上不可约（见推论 9.14 后的例 4）。本题确定 $\Phi_\ell(x)$ 模任意素数 $p$ 后因子的最小次数，特别确定它何时模 $p$ 不可约。（这实际上确定了 $\Phi_\ell(x)$ 模 $p$ 的完整分解，参见第 13.6 节练习 8。）
\begin{enumerate}[label=(\alph*)]
\item 证明若 $p=\ell$，则 $\Phi_\ell(x)$ 在 $\mathbb F_\ell[x]$ 中被 $x-1$ 整除。
\item 设 $p\ne\ell$，令 $f$ 为 $p$ 在 $\mathbb F_\ell^\times$ 中的阶，即 $f$ 是使 $p^f\equiv1\pmod\ell$ 的最小正整数。证明 $m=f$ 是群 $GL_m(\mathbb F_p)$ 含有阶为 $\ell$ 的元素 $A$ 的第一个 $m$ 值。提示：使用第 11.1 节末尾 $GL_m(\mathbb F_p)$ 阶数公式。
\item 证明 $\Phi_\ell(x)$ 在 $\mathbb F_p[x]$ 中不被任何次数小于 $f$ 的多项式整除。提示：考虑这种因子的伴随矩阵并使用 (b)。令 $m_A(x)\in\mathbb F_p[x]$ 为 (b) 中矩阵 $A$ 的最小多项式，推出 $m_A(x)$ 不可约、次数为 $f$，并整除 $\Phi_\ell(x)$。
\item 特别地，证明 $\Phi_\ell(x)$ 模 $p$ 不可约当且仅当 $\ell-1$ 是使 $p$ 模 $\ell$ 同余于 $1$ 的最小幂，即 $p$ 是模 $\ell$ 的原根。
\end{enumerate}

\item 证明定理 21 前描述的前两种初等行、列操作不改变矩阵行列式，第三种初等操作只使行列式乘以一个单位。由定理 21 推出 $A$ 的特征多项式与 $A$ 的不变因子乘积只差一个单位。由于二者按定义都是首一多项式，推出它们相等（这给出命题 20 的另一证明）。

\medskip
以下练习勾勒定理 21 的证明。它们在 Euclidean 整环 $F[x]$ 上显式执行上一节练习 16--19 中描述的构造。设 $V$ 为 $n$ 维向量空间，基为 $v_1,v_2,\ldots,v_n$，$T$ 是由矩阵 $A$ 和此基定义的线性变换，即
\[
T(v_j)=\sum_{i=1}^n a_{ij}v_i,
\qquad j=1,2,
\ldots,n,
\]
其中 $A=(a_{ij})$。令 $F[x]^n$ 为秩 $n$ 的自由 $F[x]$-模，基为 $\xi_1,\xi_2,
\ldots,\xi_n$。有自然满 $F[x]$-模同态
\[
\varphi:F[x]^n\to V,
\qquad \xi_i\mapsto v_i.
\]
由于 $x$ 在 $V$ 上按 $T$ 作用，故
\[
x(v_j)=\sum_{i=1}^n a_{ij}v_i.
\]

\item 证明元素
\[
\nu_j=-a_{1j}\xi_1-\cdots-a_{j-1,j}\xi_{j-1}+(x-a_{jj})\xi_j-a_{j+1,j}\xi_{j+1}-\cdots-a_{nj}\xi_n
\]
对 $j=1,2,
\ldots,n$ 都属于 $\ker\varphi$。

\item
\begin{enumerate}[label=(\alph*)]
\item 证明 $x\xi_j=\nu_j+f_j$，其中 $f_j\in F\xi_1+\cdots+F\xi_n$ 是由 $\xi_1,
\ldots,
\xi_n$ 张成的 $F$-向量空间中的元素。
\item 证明
\[
F[x]\xi_1+\cdots+F[x]\xi_n
=(F[x]\nu_1+\cdots+F[x]\nu_n)+(F\xi_1+\cdots+F\xi_n).
\]
\end{enumerate}

\item 证明 $\nu_1,
\nu_2,
\ldots,
\nu_n$ 生成 $\ker\varphi$。提示：用上一题说明 $\ker\varphi$ 的任意元素可写成由 $\nu_1,
\ldots,
\nu_n$ 生成的模中的元素加上形如 $b_1\xi_1+\cdots+b_n\xi_n$ 的元素，其中 $b_i\in F$；再证明这种元素在 $\ker\varphi$ 中当且仅当所有 $b_i=0$，因为 $v_1,
\ldots,v_n$ 是 $V$ 的 $F$-基。

\item 证明 $\ker\varphi$ 的生成元 $\nu_1,
\ldots,
\nu_n$ 对应的关系矩阵为
\[
\begin{pmatrix}
 x-a_{11}&-a_{21}&\cdots&-a_{n1}\\
 -a_{12}&x-a_{22}&\cdots&-a_{n2}\\
 \vdots&\vdots&\ddots&\vdots\\
 -a_{1n}&-a_{2n}&\cdots&x-a_{nn}
\end{pmatrix}=xI-A^t.
\]
推出定理 21 以及求 $A$ 的不变因子的算法可由上一节练习 16--19 得到（注意，对此关系矩阵的行、列操作正是对定理 21 中矩阵作列、行操作，这解释了不变因子算法为什么要跟踪所用的行操作）。
\end{enumerate}



\section{Jordan 标准形}

继续沿用上一节记号：$F$ 是域，$F[x]$ 是 $F$ 上多项式环，$V$ 是 $F$ 上维数为 $n$ 的有限维向量空间，$T$ 是固定线性变换，由此把 $V$ 变为 $F[x]$-模；$A$ 是 $F$ 上的 $n\times n$ 矩阵。回忆，一旦固定 $V$ 的一组基，任意线性变换 $T$ 都定义一个矩阵 $A$；反之任意矩阵 $A$ 都定义一个线性变换 $T$。

上一节使用 P.I.D. $F[x]$ 上有限生成模基本定理的不变因子形式，得到线性变换和矩阵的有理标准形。本节使用基本定理的初等因子形式得到 Jordan 标准形。我们将看到，处于这个标准形的矩阵尽可能接近对角矩阵，因此它们比有理标准形中的矩阵更简单（但会失去某些“有理性”结果）。

模的初等因子是不变因子的素幂因子（推论 10）。对 $F[x]$-模 $V$ 来说，不变因子是次数至少为 $1$ 的首一多项式 $a_1(x),a_2(x),\ldots,a_m(x)$，并满足 $a_1\mid a_2\mid\cdots\mid a_m$；相应初等因子就是这些多项式的不可约因子的幂。和不变因子一样，我们要求它们为首一，从而唯一。

为了得到最简单的初等因子，我们假设多项式 $a_1(x),a_2(x),\ldots,a_m(x)$ 全部分解为一次因子，也就是说，$V$ 的初等因子都是形如 $(x-\lambda)^k$ 的线性多项式的幂。由于初等因子的乘积是特征多项式，这等价于假设域 $F$ 含有线性变换 $T$（或表示它的矩阵 $A$）的所有特征值。

在这个对 $F$ 的假设下，由定理 6 立即可知，$V$ 是有限多个循环 $F[x]$-模
\[
F[x]/(x-\lambda)^k
\]
的直和，其中 $\lambda\in F$ 是 $T$ 的某个特征值，这些直和因子对应于 $V$ 的初等因子。

现在为 $V$ 的每个初等因子直和因子选择一组向量空间基，使得 $T$ 的相应矩阵特别简单。回忆，由 $F[x]$-模结构的定义，$V$ 上的线性变换 $T$ 就是每个直和因子 $F[x]/(x-\lambda)^k$ 上元素 $x$ 的乘法作用。

考虑商模 $F[x]/(x-\lambda)^k$ 中的元素
\[
(\bar x-\lambda)^{k-1},(\bar x-\lambda)^{k-2},\ldots,\bar x-\lambda,1.
\]
把这些多项式展开为 $x$ 的多项式，可见它们与 $F$-基 $\bar x^{k-1},\bar x^{k-2},\ldots,\bar x,1$ 之间的转换矩阵是一个对角元全为 $1$ 的上三角矩阵。因此这个矩阵可逆（行列式为 $1$），所列元素确为 $F[x]/(x-\lambda)^k$ 的一组 $F$-基。关于这组基，乘以 $x$ 的线性变换按特别简单的方式作用（注意 $x=\lambda+(x-\lambda)$，且 $(\bar x-\lambda)^k=0$）：
\[
(\bar x-\lambda)^{k-1}\mapsto \lambda(\bar x-\lambda)^{k-1},
\]
\[
(\bar x-\lambda)^{k-2}\mapsto \lambda(\bar x-\lambda)^{k-2}+(\bar x-\lambda)^{k-1},
\]
依此类推，最后
\[
1\mapsto \lambda\cdot1+(\bar x-\lambda).
\]
因此关于这组基，乘以 $x$ 的矩阵为
\[
J_k(\lambda)=
\begin{pmatrix}
\lambda&1&0&\cdots&0\\
0&\lambda&1&\cdots&0\\
\vdots&\vdots&\ddots&\ddots&\vdots\\
0&0&\cdots&\lambda&1\\
0&0&\cdots&0&\lambda
\end{pmatrix}.
\]

\noindent\textbf{定义.} 上述 $k\times k$ 矩阵称为特征值为 $\lambda$ 的 $k\times k$ \emph{初等 Jordan 矩阵}，或称特征值为 $\lambda$、大小为 $k$ 的 \emph{Jordan 块}。

把这应用到 $V$ 的初等因子分解中的每个循环因子上，就得到 $V$ 的一组向量空间基，使得 $T$ 的矩阵是对应于 $V$ 初等因子的 Jordan 块的直和，即对角线上是 Jordan 块的块对角矩阵。这个矩阵由 $F[x]$-模 $V$ 的初等因子唯一决定（至多差一个对角块的排列）；反过来，由定理 9，初等因子列表也唯一决定 $V$ 的 $F[x]$-模同构类。

\noindent\textbf{定义.}
\begin{enumerate}[label=(\arabic*)]
\item 若一个矩阵是块对角矩阵，并且对角线上各块均为 Jordan 块，则称该矩阵处于\emph{Jordan 标准形}。
\item 线性变换 $T$ 的一个 \emph{Jordan 标准形}是一个表示 $T$ 且处于 Jordan 标准形的矩阵。
\end{enumerate}

上面的讨论已经证明了任意满足条件的线性变换 $T$ 都有 Jordan 标准形。与有理标准形情形一样，由初等因子的唯一性可知，Jordan 标准形在对角线上 Jordan 块排列次序以外唯一。因此称它为 $T$ 的 Jordan 标准形。总结如下：

\noindent\textbf{定理 22（线性变换的 Jordan 标准形）.} 设 $V$ 是域 $F$ 上有限维向量空间，$T$ 是 $V$ 上线性变换。假设 $F$ 含有 $T$ 的所有特征值。
\begin{enumerate}[label=(\arabic*)]
\item 存在 $V$ 的一组基，使得 $T$ 的矩阵处于 Jordan 标准形；也就是说，它是块对角矩阵，对角块是 $V$ 的初等因子所对应的 Jordan 块。
\item $T$ 的 Jordan 标准形在对角线上 Jordan 块的排列次序以外唯一。
\end{enumerate}

相应的矩阵形式如下：

\noindent\textbf{定理 23（矩阵的 Jordan 标准形）.} 设 $A$ 是域 $F$ 上 $n\times n$ 矩阵，并假设 $F$ 含有 $A$ 的所有特征值。
\begin{enumerate}[label=(\arabic*)]
\item 矩阵 $A$ 相似于一个处于 Jordan 标准形的矩阵；也就是说，存在 $F$ 上可逆 $n\times n$ 矩阵 $P$，使得 $P^{-1}AP$ 是块对角矩阵，其对角块是 $A$ 的初等因子对应的 Jordan 块。
\item $A$ 的 Jordan 标准形在对角线上 Jordan 块的排列次序以外唯一。
\end{enumerate}

Jordan 标准形与对角矩阵的差异仅在于第一超对角线上可能出现一些 $1$（且只有当存在大小大于 $1$ 的 Jordan 块时才会出现）。因此它接近对角矩阵。下面结果尤其说明，矩阵的 Jordan 标准形确实是相似意义下尽可能接近对角矩阵的形式。

\noindent\textbf{推论 24.}
\begin{enumerate}[label=(\arabic*)]
\item 如果矩阵 $A$ 相似于对角矩阵 $D$，则 $D$ 就是 $A$ 的 Jordan 标准形。
\item 两个对角矩阵相似当且仅当它们的对角元只差一个排列。
\end{enumerate}

\begin{proof}
第一断言由 Jordan 标准形唯一性立即得到，因为对角矩阵本身已经处于 Jordan 标准形（所有 Jordan 块都为大小 $1$）。Jordan 标准形的唯一性也给出 (2)。
\end{proof}

下一推论给出判断矩阵能否对角化的准则。

\noindent\textbf{推论 25.} 若 $A$ 是 $F$ 上 $n\times n$ 矩阵，且 $F$ 含有 $A$ 的所有特征值，则 $A$ 在 $F$ 上相似于对角矩阵，当且仅当 $A$ 的最小多项式没有重根。

\begin{proof}
若 $A$ 相似于对角矩阵，则对角矩阵的最小多项式没有重根（它的根正是对角线上出现的不同元素）。相似矩阵有相同最小多项式，所以 $A$ 的最小多项式没有重根。

反过来，假设 $A$ 的最小多项式没有重根，并令 $B$ 是 $A$ 的 Jordan 标准形。矩阵 $B$ 是以初等 Jordan 矩阵为对角块的块对角矩阵。由上一节末尾的练习，$B$ 的最小多项式是各 Jordan 块最小多项式的最小公倍式。大小为 $k$、特征值为 $\lambda$ 的 Jordan 块的最小多项式是 $(x-\lambda)^k$（这也可直接从该 Jordan 块给出一个零化子为 $((x-\lambda)^k)$ 的循环 $F[x]$-子模看出）。因此如果最小多项式无重根，所有 Jordan 块的大小都必须为 $1$。于是 $B$ 是对角矩阵，故 $A$ 可对角化。
\end{proof}

有理标准形与 Jordan 标准形之间的转换，就是第 5.2 节有限阿贝尔群情形中从不变因子到初等因子以及反向的同一算法。简言之，初等因子是不变因子的素幂因子；把每个不变因子分解为不同一次因子的幂，所得线性多项式幂的集合就是初等因子。反过来，给定初等因子列表时，先按不同特征值分成若干列，在每列内按次数非降排列；再给较短的列前面补常数多项式 $1$，使所有列长度相同；然后把各列第 $i$ 个多项式相乘，得到第 $i$ 个不变因子。

\paragraph{初等因子分解算法：化为 Jordan 标准形}
定理 21 给出计算任意矩阵 $A$ 不变因子的过程。把这些不变因子分解为线性因子的幂便得到 $A$ 的初等因子，从而决定 $A$ 的 Jordan 标准形。

更具体地说，不变因子分解算法从 $V$ 的基 $e_1,\ldots,e_n$ 出发，产生 $V$ 的若干元素 $f_1,\ldots,f_m$，它们是 $V$ 的不变因子分解中循环 $F[x]$-因子的 $F[x]$-模生成元，对应的零化子为 $(a_1(x)),\ldots,(a_m(x))$。由于初等因子分解由对每个循环模 $F[x]/(a_i(x))$ 应用中国剩余定理得到，这也给出初等因子分解中各循环因子的生成元。若
\[
a(x)=(x-\lambda_1)^{\alpha_1}(x-\lambda_2)^{\alpha_2}\cdots(x-\lambda_s)^{\alpha_s}
\]
且 $f$ 是对应于不变因子 $a(x)$ 的循环因子生成元，则
\[
\frac{a(x)}{(x-\lambda_1)^{\alpha_1}}f,
\ldots,
\frac{a(x)}{(x-\lambda_s)^{\alpha_s}}f
\]
分别给出对应初等因子 $(x-\lambda_i)^{\alpha_i}$ 的循环因子生成元。若
\[
g_i=\frac{a(x)}{(x-\lambda_i)^{\alpha_i}}f,
\]
则对应循环因子的向量空间基为
\[
(T-\lambda_i)^{\alpha_i-1}g_i,
(T-\lambda_i)^{\alpha_i-2}g_i,
\ldots,
(T-\lambda_i)g_i,g_i.
\]
把这些向量在原基下的坐标作为矩阵 $P$ 的列，则 $P^{-1}AP$ 为 Jordan 标准形。

\paragraph{例}
设 $A,B,C$ 是上一节例 1 中的矩阵，并取 $F=\mathbb Q$。注意 $\mathbb Q$ 含有这些矩阵的所有特征值。由于已经求出不变因子，可立即得到初等因子：$A$ 的初等因子为 $x-2,x-2,x-3$；$B$ 与 $C$ 的初等因子为 $(x-2)^2$ 和 $x-3$。因此 $A$ 相似于对角矩阵，而 $B,C$ 不可对角化。

\paragraph{例}
对上一节矩阵 $A$ 的显式基，已得到不变因子分解中的生成元
\[
f_1=-7e_1+7e_2+e_3,
\qquad f_2=-e_1+e_2,
\]
对应不变因子 $x-2$ 与 $(x-2)(x-3)$。由上述算法，初等因子分解中的生成元为
\[
f_1,\quad (x-3)f_2,\quad (x-2)f_2,
\]
分别对应 $x-2,x-2,x-3$。直接计算得这些向量为 $-7e_1+7e_2+e_3$、$-e_1$、$-2e_1+e_2$。以它们为列组成矩阵 $P$，则 $P^{-1}AP$ 为相应 Jordan 标准形。

\paragraph{例}
若一个 $4\times4$ 矩阵的不变因子为 $(x-1)^2,(x-1)^2$，相应 $\mathbb Q[x]$-模生成元为 $f_1=e_1$ 与 $f_2=e_2$，则这些不变因子本身就是初等因子。相应 Jordan 基为 $(T-I)f_1,f_1,(T-I)f_2,f_2$；在原书给定的具体矩阵中，这些向量为 $2e_2+e_3,e_1,2e_1-e_2+e_4,e_2$，从而给出两个大小为 $2$、特征值为 $1$ 的 Jordan 块。

当扩张底域时，相似类的数目可能增加（由推论 18(2)，永远不会减少）。例如在 $\mathbb C$ 上，特征多项式 $(x^4-1)(x^2-1)$ 的矩阵可进一步按一次因子 $x-i$ 与 $x+i$ 分解 Jordan 标准形；而在 $\mathbb Q$ 上这些因子合并为不可约二次因子 $x^2+1$，只能在有理标准形中体现。
\subsection*{练习}

\begin{enumerate}[label=\arabic*.]
\item 设向量空间 $V$ 是若干循环 $F[x]$-模的直和，其零化子分别为
\[
(x+1)^2,
\qquad (x-1)(x^2+1)^2,
\qquad (x^4-1),
\qquad (x+1)(x^2-1).
\]
求 $V$ 的不变因子和初等因子。

\item 证明：若 $\lambda_1,
\ldots,\lambda_n$ 是 $n\times n$ 矩阵 $A$ 的特征值，则对任意 $k\geq0$，$\lambda_1^k,
\ldots,
\lambda_n^k$ 是 $A^k$ 的特征值。

\item 用上面例 2 的方法，求显式矩阵 $P_1,P_2$，使 $P_1^{-1}BP_1$ 与 $P_2^{-1}CP_2$ 为 Jordan 标准形。用此显式构造一个矩阵 $Q$，它把 $B$ 共轭为 $C$，从而直接证明这两个矩阵相似。

\item 证明矩阵
\[
\begin{pmatrix}
9&4&5\\
-4&0&-3\\
-6&-4&-2
\end{pmatrix}
\]
的 Jordan 标准形正是本章开头给出的形式。显式求出一个矩阵 $P$，使它把该矩阵共轭到其 Jordan 标准形。说明为什么该矩阵不能对角化。

\item 求矩阵
\[
\begin{pmatrix}
1&0&0\\
0&0&-2\\
0&1&3
\end{pmatrix}
\]
的 Jordan 标准形。

\item 判断下列矩阵中哪些相似：
\[
\begin{pmatrix}-1&4&-4\\2&-1&3\\0&-4&3\end{pmatrix},\quad
\begin{pmatrix}-3&-4&0\\2&3&0\\8&8&1\end{pmatrix},\quad
\begin{pmatrix}-3&2&-4\\2&1&0\\3&-1&3\end{pmatrix},\quad
\begin{pmatrix}-1&4&-4\\0&-3&2\\0&-4&3\end{pmatrix}.
\]

\item 求下列矩阵的 Jordan 标准形：
\[
\begin{pmatrix}5&4&1\\-1&0&0\\-3&-4&1\end{pmatrix},\qquad
\begin{pmatrix}3&4&2\\-2&-3&-1\\-4&-4&-3\end{pmatrix}.
\]

\item 证明矩阵
\[
A=\begin{pmatrix}5&6&0\\-3&-4&0\\-2&0&1\end{pmatrix},
\qquad
B=\begin{pmatrix}3&-1&2\\-10&6&-14\\-6&3&-7\end{pmatrix}
\]
相似。证明 $A$ 与 $B$ 都可对角化，并求显式矩阵 $P_1,P_2$，使 $P_1^{-1}AP_1$ 与 $P_2^{-1}BP_2$ 为对角形式。

\item 证明矩阵
\[
A=\begin{pmatrix}-8&-10&-1\\7&9&1\\3&2&0\end{pmatrix},
\qquad
B=\begin{pmatrix}-3&2&-4\\4&-1&4\\4&-2&5\end{pmatrix}
\]
都以 $(x-1)^2(x+1)$ 为特征多项式，但其中一个可对角化而另一个不可对角化。求两个矩阵的 Jordan 标准形。

\item 求复数域上所有 $2\times2$、$3\times3$ 与 $4\times4$ 矩阵的 Jordan 标准形。

\item 验证矩阵
\[
A=\begin{pmatrix}
1&0&0&0\\
0&1&0&0\\
-2&-2&0&1\\
-2&0&-1&-2
\end{pmatrix}
\]
的特征多项式在 $\mathbb Q$ 上分解为一次因子。求 $A$ 在 $\mathbb Q$ 上的有理标准形与 Jordan 标准形。

\item 求矩阵
\[
\begin{pmatrix}
1&2&0&0\\
0&1&2&0\\
0&0&1&2\\
0&0&0&1
\end{pmatrix}
\]
的 Jordan 标准形。

\item 求矩阵
\[
\begin{pmatrix}
3&0&-2&-3\\
4&-8&14&-15\\
2&-4&7&-7\\
0&2&-4&3
\end{pmatrix}
\]
的 Jordan 标准形。

\item 证明矩阵
\[
A=\begin{pmatrix}
2&0&0&0\\
-4&-1&-4&0\\
2&1&3&0\\
-2&4&9&1
\end{pmatrix},
\qquad
B=\begin{pmatrix}
5&0&-4&-7\\
3&-8&15&-13\\
2&-4&7&-7\\
1&2&-5&1
\end{pmatrix}
\]
相似。

\item 证明矩阵
\[
A=\begin{pmatrix}
0&1&1&1\\
1&0&1&1\\
1&1&0&1\\
1&1&1&0
\end{pmatrix},
\qquad
B=\begin{pmatrix}
5&2&-8&-8\\
-6&-3&8&8\\
-3&-1&3&4\\
3&1&-4&-5
\end{pmatrix}
\]
都以 $(x-3)(x+1)^3$ 为特征多项式。判断它们是否相似，并求每个矩阵的 Jordan 标准形。

\item 求矩阵
\[
\begin{pmatrix}
1&1&1&1\\
0&1&0&-1\\
0&0&1&1\\
0&0&0&1
\end{pmatrix}
\]
的 Jordan 标准形，并求一个矩阵 $P$，使它把该矩阵共轭到其 Jordan 标准形。

\item 证明任意矩阵 $A$ 都与其转置 $A^t$ 相似。

\item 求特征多项式为 $(x-2)^3(x-3)^2$ 的线性变换的所有可能 Jordan 标准形。

\item 证明：所有特征多项式为 $f(x)$ 的 $n\times n$ 矩阵彼此相似，当且仅当 $f(x)$ 在 $F[x]$ 中的唯一因子分解没有重因子。

\item 证明下列两个 $M_p(\mathbb F_p)$ 中的矩阵相似：
\[
\begin{pmatrix}
0&0&0&\cdots&0&1\\
1&0&0&\cdots&0&0\\
0&1&0&\cdots&0&0\\
\vdots&\vdots&\vdots& &\vdots&\vdots\\
0&0&0&\cdots&1&0
\end{pmatrix}
\quad\text{和}\quad
\begin{pmatrix}
1&1&0&\cdots&0&0\\
0&1&1&\cdots&0&0\\
0&0&1&\cdots&0&0\\
\vdots&\vdots&\vdots& &\vdots&\vdots\\
0&0&0&\cdots&1&1\\
0&0&0&\cdots&0&1
\end{pmatrix}.
\]

\item 证明若 $A^2=A$，则 $A$ 相似于一个只以 $0$ 和 $1$ 为对角元的对角矩阵。

\item 证明复数域上满足 $A^3=A$ 的 $n\times n$ 矩阵 $A$ 可对角化。这个结论在任意域 $F$ 上是否成立？

\item 设 $A$ 是 $\mathbb Q$ 上 $2\times2$ 矩阵，满足 $A^3=I$ 但 $A\ne I$。把 $A$ 写成有理标准形；再把它视为 $\mathbb C$ 上矩阵，写出 Jordan 标准形。

\item 证明不存在 $\mathbb Q$ 上 $3\times3$ 矩阵 $A$ 满足 $A^8=I$ 但 $A^4\ne I$。

\item 求 $\mathbb Q$ 上所有条目都等于 $1$ 的 $n\times n$ 矩阵的 Jordan 标准形。

\item 求 $\mathbb F_p$ 上所有条目都等于 $1$ 的 $n\times n$ 矩阵的 Jordan 标准形（答案取决于 $p$ 是否整除 $n$）。

\item 求 $\mathbb Q$ 上 $n\times n$ 矩阵的 Jordan 标准形，其中所有非对角线条目均为 $1$，而主对角线条目均为 $0$。

\item 求 $\mathbb F_p$ 上 $n\times n$ 矩阵的 Jordan 标准形，其中所有非对角线条目均为 $1$，而主对角线条目均为 $0$。

与同一特征值 $\lambda$ 对应的、所有初等因子为 $(x-\lambda)$ 的幂的循环子模的直和，称为 $T$ 关于特征值 $\lambda$ 的\emph{广义特征空间}。这是 $F[x]$ 中素元 $p=x-\lambda$ 的 primary 分量，包含所有被某个 $(T-\lambda)$ 的幂零化的 $V$ 中元素。

\item 设 $V_i$ 是 $T$ 对特征值 $\lambda_i$ 的广义特征空间。对任意 $k\geq0$，证明 $T-\lambda_i$ 在子空间 $(T-\lambda_i)^kV_i$ 上的零化度，等于 $T-\lambda_i$ 在 $(T-\lambda_i)^kV$ 上的零化度，并且等于 $T$ 中特征值为 $\lambda_i$、大小大于 $k$ 的 Jordan 块个数（当 $k=0$ 时给出 Jordan 块的个数）。

\item 设 $\lambda$ 是域 $F$ 上有限维向量空间 $V$ 上线性变换 $T$ 的特征值。令
\[
r_k=\dim_F (T-\lambda)^kV
\]
为 $(T-\lambda)^k$ 在 $V$ 上的秩。证明对任意 $k\geq1$，
\[
r_{k-1}-2r_k+r_{k+1}
\]
是 $T$ 中对应于 $\lambda$ 且大小为 $k$ 的 Jordan 块个数（使用第 1 节练习 12）。这给出一种通过计算矩阵 $(A-\lambda I)^k$ 的秩来确定 Jordan 标准形的有效方法。

\item 设 $N$ 是域 $F$ 上 $n\times n$ 矩阵。若某个幂 $N^k$ 为零矩阵，则称 $N$ 为\emph{幂零矩阵}。证明任意幂零矩阵都相似于一个块对角矩阵，其每个块在第一超对角线上全为 $1$，其余条目全为 $0$。

\item 证明若 $N$ 是 $n\times n$ 幂零矩阵，则事实上 $N^n=0$。

\item 设 $A$ 是严格上三角的 $n\times n$ 矩阵（主对角线及其下方条目均为 $0$）。证明 $A$ 是幂零矩阵。

\item 证明 $n\times n$ 幂零矩阵的迹为 $0$。

\item 对 $0\leq i\leq n$，令 $d_i$ 为 $xI-A$ 的所有 $i\times i$ 子式行列式的最大公因子（把 $0\times0$ 子式定义为 $1$）。证明 $A$ 的 Smith 标准形对角线上第 $i$ 个元素为 $d_i/d_{i-1}$。这给出 $A$ 的不变因子。提示：证明这些最大公因子在初等行、列操作下不变。

\item 令 $V=\mathbb C^n$ 为通常的 $n$ 维复向量空间，$T$ 为线性变换
\[
T(a_1,a_2,
\ldots,a_n)=(0,a_1,a_2,
\ldots,a_{n-1}).
\]
求 $T$ 的 Jordan 标准形。

\item 设 $J$ 是复数域上大小为 $n$、特征值为 $\lambda$ 的 Jordan 块。
\begin{enumerate}[label=(\alph*)]
\item 证明若 $\lambda\ne0$，则 $J^2$ 的 Jordan 标准形是大小为 $n$、特征值为 $\lambda^2$ 的 Jordan 块。
\item 若 $\lambda=0$，证明 $J^2$ 的 Jordan 标准形有两个块，特征值均为 $0$；当 $n$ 为偶数时，两个块大小为 $n/2,n/2$，当 $n$ 为奇数时，大小为 $(n-1)/2,(n+1)/2$。
\end{enumerate}

\item 求矩阵 $A\in M_n(\mathbb C)$ 有平方根的充要条件，即存在 $B\in M_n(\mathbb C)$ 使 $A=B^2$。提示：设 $B$ 为 Jordan 标准形，并使用上一题考虑 $B^2$ 的 Jordan 标准形。

\item 设 $J$ 是特征为 $2$ 的域 $F$ 上大小为 $n$、特征值为 $\lambda$ 的 Jordan 块。求 $J^2$ 的 Jordan 标准形。求 $A\in M_n(F)$ 有平方根的充要条件，即存在 $B\in M_n(F)$ 使 $A=B^2$。

\medskip
以下练习考察矩阵函数（幂级数）并介绍 Jordan 标准形在常微分方程理论中的一些应用。以下矩阵均为 $K$ 上 $n\times n$ 矩阵，其中 $K=\mathbb R$ 或 $\mathbb C$。令
\[
G(x)=\sum_{k=0}^{\infty}\alpha_kx^k
\]
为 $K$ 系数幂级数，令 $G_N(x)=\sum_{k=0}^N\alpha_kx^k$。对 $A\in M_n(K)$，把 $A$ 代入 $G_N$ 得 $G_N(A)$。若对每个 $i,j$，$G_N(A)$ 的 $(i,j)$ 条目收敛到矩阵 $C$ 的 $(i,j)$ 条目，则称
\[
G(A)=\sum_{k=0}^{\infty}\alpha_kA^k
\]
收敛到 $C$，记作 $G(A)=C$。定义
\[
\|A\|=\sum_{i,j=1}^n |a_{ij}|.
\]

\item 证明对所有 $A,B\in M_n(K)$ 和 $\alpha\in K$，有：
\begin{enumerate}[label=(\alph*)]
\item $\|A+B\|\leq \|A\|+\|B\|$；
\item $\|AB\|\leq \|A\|\,\|B\|$；
\item $\|\alpha A\|=|\alpha|\,\|A\|$。
\end{enumerate}

\item 设 $R$ 为实或复幂级数 $G(x)$ 的收敛半径（若 $G(x)$ 对所有 $x\in K$ 收敛，则 $R=\infty$）。
\begin{enumerate}[label=(\alph*)]
\item 证明若 $\|A\|<R$，则 $G(A)$ 收敛。
\item 推出对所有矩阵 $A$，下列幂级数都收敛：
\[
\sin(A)=A-\frac{A^3}{3!}+\frac{A^5}{5!}+\cdots+(-1)^k\frac{A^{2k+1}}{(2k+1)!}+\cdots,
\]
\[
\cos(A)=I-\frac{A^2}{2!}+\frac{A^4}{4!}+\cdots+(-1)^k\frac{A^{2k}}{(2k)!}+\cdots,
\]
\[
\exp(A)=I+A+\frac{A^2}{2!}+\frac{A^3}{3!}+\cdots+\frac{A^k}{k!}+\cdots,
\]
其中 $I$ 是 $n\times n$ 单位矩阵。
\end{enumerate}

为应用于微分方程，引入变量 $t$。对 $A\in M_n(K)$，矩阵 $At$ 由 $A$ 中每个条目乘以 $t$ 得到。因此在收敛点上，$t\mapsto G(At)$ 是从 $K$ 的子集到 $M_n(K)$ 的函数。特别地，$\sin(At)$、$\cos(At)$ 与 $\exp(At)$ 对所有 $t\in K$ 收敛。

\item 设 $P$ 为非奇异 $n\times n$ 矩阵。
\begin{enumerate}[label=(\alph*)]
\item 证明
\[
PG(At)P^{-1}=G(PAtP^{-1})=G(PAP^{-1}t).
\]
这说明在相差一个基变换的意义下，只需对标准形矩阵计算 $G(At)$。提示：先对部分和取极限得到第一个等式；第二个等式由标量矩阵 $tI$ 与任意矩阵可交换立即得出。
\item 证明若 $A$ 是矩阵 $A_1,A_2,\ldots,A_m$ 的直和，则 $G(At)$ 是 $G(A_1t),G(A_2t),\ldots,G(A_mt)$ 的直和。
\item 证明若 $Z$ 是对角矩阵，对角元为 $z_1,z_2,
\ldots,z_n$，则 $G(Zt)$ 是对角矩阵，对角元为 $G(z_1t),G(z_2t),
\ldots,G(z_nt)$。
\end{enumerate}

矩阵 $\exp(A)$ 常记为 $e^A$。下面三题导出初等 Jordan 矩阵 $J$ 的 $\exp(Jt)$ 公式。

\item 证明若 $A$ 与 $B$ 可交换，则
\[
\exp(A+B)=\exp(A)\exp(B).
\]
提示：把 $A,B$ 看作可交换不定元，比较左边幂级数与右边两个幂级数的乘积。

\item 用上一题证明：若 $M$ 为任意矩阵，$\lambda\in K$，则
\[
\exp(\lambda It+M)=e^{\lambda t}\exp(M).
\]

\item 令 $N$ 为第一超对角线上全为 $1$、其余条目全为 $0$ 的 $r\times r$ 矩阵。计算如下幂零矩阵的指数：若
\[
Nt=\begin{pmatrix}
0&t&0&\cdots&0\\
0&0&t&\cdots&0\\
\vdots&\vdots&\ddots&\ddots&\vdots\\
0&0&\cdots&0&t\\
0&0&\cdots&0&0
\end{pmatrix},
\]
则
\[
\exp(Nt)=
\begin{pmatrix}
1&t&\frac{t^2}{2!}&\cdots&\frac{t^{r-1}}{(r-1)!}\\
0&1&t&\cdots&\frac{t^{r-2}}{(r-2)!}\\
\vdots&\vdots&\ddots&\ddots&\vdots\\
0&0&\cdots&1&t\\
0&0&\cdots&0&1
\end{pmatrix}.
\]
推出若 $J$ 是特征值为 $\lambda$ 的 $r\times r$ 初等 Jordan 矩阵，则
\[
\exp(Jt)=
\begin{pmatrix}
 e^{\lambda t}&te^{\lambda t}&\frac{t^2}{2!}e^{\lambda t}&\cdots&\frac{t^{r-1}}{(r-1)!}e^{\lambda t}\\
0&e^{\lambda t}&te^{\lambda t}&\cdots&\frac{t^{r-2}}{(r-2)!}e^{\lambda t}\\
\vdots&\vdots&\ddots&\ddots&\vdots\\
0&0&\cdots&e^{\lambda t}&te^{\lambda t}\\
0&0&\cdots&0&e^{\lambda t}
\end{pmatrix}.
\]
提示：由于 $Nt$ 幂零，$\exp(Nt)$ 是 $Nt$ 的多项式。计算 $\exp(Jt)$ 时写 $Jt=\lambda It+Nt$，再用练习 44。

若 $A\in M_n(K)$，$P$ 是基变换矩阵，使 $P^{-1}AP$ 为 Jordan 标准形。设 $P^{-1}AP$ 是初等 Jordan 矩阵 $J_1,
\ldots,J_m$ 的直和。由前面练习（取 $t=1$），先写 $E=\exp(P^{-1}AP)$ 为 $\exp(J_1),
\ldots,
\exp(J_m)$ 的直和，再变回原基，即得 $\exp(A)=PEP^{-1}$。

\item 对本节例 3 中给出的 $4\times4$ 矩阵 $D$ 与 $P$：
\[
D=\begin{pmatrix}
1&2&-4&4\\
2&-1&4&-8\\
1&0&1&-2\\
0&1&-2&3
\end{pmatrix},
\qquad
P=\begin{pmatrix}
0&1&2&0\\
2&0&-2&1\\
1&0&0&0\\
0&0&1&0
\end{pmatrix},
\]
证明
\[
E=\begin{pmatrix}
e&e&0&0\\
0&e&0&0\\
0&0&e&e\\
0&0&0&e
\end{pmatrix},
\qquad
\exp(D)=\begin{pmatrix}
e&2e&-4e&4e\\
2e&-e&4e&-8e\\
e&0&e&-2e\\
0&e&-2e&3e
\end{pmatrix}.
\]

\item 计算下列每个矩阵的指数：
\begin{enumerate}[label=(\alph*)]
\item 本节例 2 中的矩阵 $A$；
\item 练习 4 中的矩阵（在那里已经求了 Jordan 标准形和一个基变换矩阵）；
\item 练习 16 中的矩阵。
\end{enumerate}

\item 证明 $\exp(0)=I$（其中 $0$ 为零矩阵，$I$ 为单位矩阵）。推出对所有 $A\in M_n(K)$，$\exp(A)$ 非奇异，且逆矩阵为 $\exp(-A)$。

\item 证明
\[
\det(\exp(A))=e^{\operatorname{tr}(A)},
\]
其中 $\operatorname{tr}(A)$ 是 $A$ 的迹。

\item 固定 $A\in M_n(K)$。证明映射
\[
K\longrightarrow GL_n(K),
\qquad
 t\longmapsto \exp(At)
\]
是群同态（这里 $K$ 是域的加法群）。注意这推广了熟悉的指数映射 $K\to K^\times$，即 $n=1$ 情形。子群 $\{\exp(At)\mid t\in K\}$ 称为 $GL_n(K)$ 的一个\emph{单参数子群}。这些子群和指数映射在 Lie 群理论中很重要。

设 $G(x)$ 是收敛半径为无穷大的幂级数，并固定矩阵 $A\in M_n(K)$。矩阵 $G(At)$ 的条目是变量 $t$ 的函数。令 $c_{ij}(t)$ 是 $G(At)$ 的 $(i,j)$ 条目。$G(At)$ 对 $t$ 的导数记为 $\frac{d}{dt}G(At)$，定义为各条目逐项求导所得矩阵。

\item 建立下列导数性质：
\begin{enumerate}[label=(\alph*)]
\item 若 $G(x)=\sum_{k=0}^\infty \alpha_kx^k$，则
\[
\frac{d}{dt}G(At)=A\sum_{k=1}^\infty k\alpha_k(At)^{k-1}.
\]
\item 若 $v$ 是条目为常数的 $n\times1$ 矩阵，则
\[
\frac{d}{dt}\bigl(G(At)v\bigr)=\left(\frac{d}{dt}G(At)\right)v.
\]
\end{enumerate}

\item 由上一题 (a) 推出
\[
\frac{d}{dt}\exp(At)=A\exp(At).
\]

现在令 $y_1(t),
\ldots,y_n(t)$ 为实变量 $t$ 的可微函数，它们满足以下常系数一阶线性微分方程组：
\[
\begin{aligned}
y_1'&=a_{11}y_1+a_{12}y_2+\cdots+a_{1n}y_n,\\
y_2'&=a_{21}y_1+a_{22}y_2+\cdots+a_{2n}y_n,\\
&\ \vdots\\
y_n'&=a_{n1}y_1+a_{n2}y_2+\cdots+a_{nn}y_n.
\end{aligned}\tag{$*$}
\]
令 $A=(a_{ij})$，则可写为 $y'=Ay$。一个条目为 $t$ 的函数、其列向量为方程组 $(*)$ 的线性无关解的 $n\times n$ 矩阵，称为 $(*)$ 的\emph{基本矩阵}。由微分方程理论，$(*)$ 的解向量空间为 $n$ 维，其基本矩阵的列构成全部解空间的一组基。

\item 证明 $\exp(At)$ 是 $(*)$ 的基本矩阵。并证明：若 $C$ 是 $n\times1$ 常数列向量，其条目为 $y_1(0),
\ldots,y_n(0)$，则
\[
y(t)=\exp(At)C
\]
是满足初值条件 $y(0)=C$ 的 $(*)$ 的特解。注意这推广了 $1$ 维方程 $y'=ay$ 的结果：$e^{at}$ 是解空间的一组基，且满足 $y(0)=c$ 的唯一解是 $y=ce^{at}$。提示：使用前面练习。

\item 证明若 $M$ 是 $(*)$ 的基本矩阵，$Q\in M_n(K)$ 非奇异，则 $MQ$ 也是 $(*)$ 的基本矩阵。提示：$MQ$ 的列是 $M$ 的列的线性组合。

现在用前两题解决具体微分方程组：给定 $(*)$ 中的矩阵 $A$，求基变换矩阵 $P$ 使 $B=P^{-1}AP$ 为 Jordan 标准形。则 $\exp(At)=P\exp(Bt)P^{-1}$ 是 $(*)$ 的基本矩阵；由上一题，$P\exp(Bt)$ 也是基本矩阵。这样可用练习 45 后的方法计算 $\exp(Bt)$，且不必求 $P^{-1}$。

例如，若 $A=D$ 且 $P$ 为练习 46 中矩阵，则 $B=P^{-1}AP$ 由两个特征值为 $1$ 的 $2\times2$ Jordan 块组成。因此 $y'=Ay$ 的一个基本矩阵为
\[
P\exp(Bt)=
\begin{pmatrix}
0&1&2&0\\
2&0&-2&1\\
1&0&0&0\\
0&0&1&0
\end{pmatrix}
\begin{pmatrix}
e^t&te^t&0&0\\
0&e^t&0&0\\
0&0&e^t&te^t\\
0&0&0&e^t
\end{pmatrix}
=
\begin{pmatrix}
0&e^t&2e^t&2te^t\\
2e^t&2te^t&-2e^t&e^t(1-2t)\\
e^t&te^t&0&0\\
0&0&e^t&te^t
\end{pmatrix}.
\]
由此得方程组
\[
\begin{aligned}
y_1'&=y_1+2y_2-4y_3+4y_4,\\
y_2'&=2y_1-y_2+4y_3-8y_4,\\
y_3'&=y_1+y_3-2y_4,\\
y_4'&=y_2-2y_3+3y_4
\end{aligned}
\]
的一般解
\[
\begin{pmatrix}y_1\\y_2\\y_3\\y_4\end{pmatrix}
=\alpha_1\begin{pmatrix}0\\2e^t\\e^t\\0\end{pmatrix}
+\alpha_2\begin{pmatrix}e^t\\2te^t\\te^t\\0\end{pmatrix}
+\alpha_3\begin{pmatrix}2e^t\\-2e^t\\0\\e^t\end{pmatrix}
+\alpha_4\begin{pmatrix}2te^t\\e^t(1-2t)\\0\\te^t\end{pmatrix},
\]
其中 $\alpha_1,
\ldots,
\alpha_4\in K$ 任意。

\item 在 (a)--(c) 各部分中，对给定系数矩阵 $A$ 的方程组 $(*)$ 求一个基本矩阵。
\begin{enumerate}[label=(\alph*)]
\item $A$ 为练习 47(a) 中的矩阵。
\item $A$ 为练习 47(b) 中的矩阵。
\item $A$ 为练习 47(c) 中的矩阵。
\end{enumerate}

\item 考虑方程组 $(*)$，其系数矩阵 $A$ 为练习 46 中列出的矩阵 $D$，其基本矩阵已在上一题前算出。求满足初始条件 $y_i(0)=1$（$i=1,2,3,4$）的 $(*)$ 的特解。

接下来研究 $(*)$ 的一个特殊情形。给定常系数 $n$ 阶线性微分方程
\[
y^{(n)}+a_{n-1}y^{(n-1)}+\cdots+a_1y'+a_0y=0,\tag{$**$}
\]
其中 $y^{(k)}$ 表示 $y$ 的第 $k$ 阶导数，$y^{(0)}=y$。令 $y_i=y^{(i-1)}$，$1\leq i\leq n$，可把它化为一阶线性方程组。$(**)$ 的解空间的一组基可由该线性系统的基本矩阵得到：基本矩阵每一列的第 $(1,1)$ 项给出 $(**)$ 的一个解，第一行的 $n$ 个函数组成 $(**)$ 的解空间基。

\item 证明由 $(**)$ 得到的一阶系统的系数矩阵 $A$ 是多项式
\[
x^n+a_{n-1}x^{n-1}+\cdots+a_1x+a_0
\]
的伴随矩阵的转置。

\item 用上述方法求下列微分方程的解空间基：
\begin{enumerate}[label=(\alph*)]
\item $y'''-3y'+2y=0$；
\item $y''''+4y'''+6y''+4y'+y=0$。
\end{enumerate}

微分方程组
\[
\begin{aligned}
y_1'&=F_1(y_1,y_2,
\ldots,y_n),\\
y_2'&=F_2(y_1,y_2,
\ldots,y_n),\\
&\ \vdots\\
y_n'&=F_n(y_1,y_2,
\ldots,y_n)
\end{aligned}
\]
称为\emph{自治系统}，简记为 $y'=F(y)$，其中 $F=(F_1,
\ldots,F_n)$。自治意味着右端不显含时间变量 $t$。常系数线性系统 $(*)$ 是其中每个 $F_i$ 都为线性函数的特殊情形。

常数解 $y=c$（其中 $c=(c_1,
\ldots,c_n)^t$）满足 $y'=F(y)$ 当且仅当
\[
0=F_i(c_1,
\ldots,c_n),
\qquad i=1,2,
\ldots,n.
\]
这些常数解称为\emph{稳态}。若每个解都趋于稳态 $c$，即对任意解 $y(t)$，都有 $\lim_{t\to\infty}y_i(t)=c_i$，则称该稳态\emph{全局渐近稳定}。在线性自治系统中，解空间为向量空间，唯一常数解是零解。下一题给出零解全局渐近稳定的一个充分条件。

\item 证明：若 $A$ 的所有特征值实部均为负，则 $(*)$ 的零解 $y_i(t)=0$（$i=1,
\ldots,n$）全局渐近稳定。若不熟悉复指数函数的行为，可假设所有特征值均为实数（因此为负实数）。利用解的显式形式说明它们均在 $t\to\infty$ 时趋于零。
\end{enumerate}

% \setchapterimage{前面几部分发展了群、环和域这些基本代数结构的理论。接下来的两章研究域的性质，特别是那些由求解方程（例如简单方程 $x^2+1=0$）而产生的域，以及在进行“算术”（加、减、乘、除）时自然出现的域。优美的 Galois 理论把域的结构与某些相关群的结构联系起来，是基本的代数工具之一。它的应用包括尺规作图问题、有限域，以及 Abel 关于一般五次方程不可用根式求解的著名定理。
% }
%\part{域理论与Galois理论}
% % ============================================
% chapter13_full_zh.tex — 第13章：域论
% 与第12章格式统一；正文、证明、例子与练习均保留
% ============================================

前面几部分发展了群、环和域这些基本代数结构的理论。接下来的两章研究域的性质，特别是那些由求解方程（例如简单方程 $x^2+1=0$）而产生的域，以及在进行“算术”（加、减、乘、除）时自然出现的域。优美的 Galois 理论把域的结构与某些相关群的结构联系起来，是基本的代数工具之一。它的应用包括尺规作图问题、有限域，以及 Abel 关于一般五次方程不可用根式求解的著名定理。

\chapter{域论}

\section{域扩张的基本理论}

回忆：域 $F$ 是一个带恒等元的交换环，并且其中每个非零元素都有逆。等价地，$F$ 的非零元素集合 $F^\times=F-\{0\}$ 在乘法下是一个阿贝尔群。

与任意域 $F$ 相联系的第一个不变量之一是它的特征，其定义如下：如果 $1_F$ 表示 $F$ 的恒等元，那么 $F$ 含有由 $1_F$ 生成的加法子群中的元素
\[
1_F,\quad 1_F+1_F,\quad 1_F+1_F+1_F,\ldots,
\]
这些元素未必全都不同。对正整数 $n$，令 $n\cdot 1_F=1_F+\cdots+1_F$（$n$ 次）。于是有两种可能：或者所有元素 $n\cdot 1_F$ 都互不相同，或者存在某个正整数 $n$ 使得 $n\cdot 1_F=0$。

\noindent\textbf{定义.} 域 $F$ 的\emph{特征}，记为 $\operatorname{ch}(F)$，定义为满足 $p\cdot 1_F=0$ 的最小正整数 $p$，如果这样的 $p$ 存在；若不存在这样的 $p$，则定义为 $0$。

容易看到
\[
n\cdot 1_F+m\cdot 1_F=(m+n)\cdot 1_F
\]
以及
\[
(n\cdot 1_F)(m\cdot 1_F)=mn\cdot 1_F \tag{13.1}
\]
对正整数 $m$ 和 $n$ 成立。由此可知，域的特征或者是 $0$，或者是一个素数 $p$（因此上面的定义中选用字母 $p$），因为若 $n=ab$ 是合数并且 $n\cdot 1_F=0$，则
\[
ab\cdot 1_F=(a\cdot 1_F)(b\cdot 1_F)=0,
\]
而由于 $F$ 是域，$a\cdot 1_F$ 或 $b\cdot 1_F$ 中有一个为 $0$，所以这样的最小整数必然是素数。还可推出，若 $n\cdot 1_F=0$，则 $n$ 被 $p$ 整除。

\noindent\textbf{命题 1.} 域 $F$ 的特征 $\operatorname{ch}(F)$ 或为 $0$，或为素数 $p$。若 $\operatorname{ch}(F)=p$，则对任意 $a\in F$，
\[
p\cdot a=\underbrace{a+a+\cdots+a}_{p\text{ 次}}=0.
\]

\begin{proof}
	只有第二个断言尚未证明；它立即来自 $F$ 中显然的等式
	\[
	p\cdot a=p\cdot(1_Fa)=(p\cdot 1_F)(a).
	\]
\end{proof}

\paragraph{注记：} 这种特征的概念也适用于任意整环，并且一个整环的特征与其分式域的特征相同。

\paragraph{例}
\begin{enumerate}[label=(\arabic*)]
	\item 域 $\mathbb Q$ 和 $\mathbb R$ 的特征都为 $0$：$\operatorname{ch}(\mathbb Q)=\operatorname{ch}(\mathbb R)=0$。整环 $\mathbb Z$ 的特征也为 $0$。
	\item 对任意素数 $p$，有限域 $\mathbb F_p=\mathbb Z/p\mathbb Z$ 的特征为 $p$。
	\item 系数在域 $\mathbb F_p$ 中的变量 $x$ 的多项式整环 $\mathbb F_p[x]$ 的特征为 $p$，它的分式域 $\mathbb F_p(x)$（系数在 $\mathbb F_p$ 中的关于 $x$ 的有理函数域）的特征也为 $p$。
\end{enumerate}

若对正整数 $n$ 定义 $(-n)\cdot 1_F=-(n\cdot 1_F)$，并定义 $0\cdot 1_F=0$，则由等式 (13.1) 得到一个自然环同态
\[
\varphi:\mathbb Z\longrightarrow F,
\qquad
n\longmapsto n\cdot 1_F,
\]
并且我们可通过注意到 $\ker(\varphi)=\operatorname{ch}(F)\mathbb Z$ 来解释 $F$ 的特征。对核取商给出 $\mathbb Z$ 或 $\mathbb Z/p\mathbb Z$ 到 $F$ 的一个嵌入（取决于 $\operatorname{ch}(F)=0$ 还是 $\operatorname{ch}(F)=p$）。由于 $F$ 是域，我们看到，取决于 $F$ 的特征，$F$ 含有一个同构于 $\mathbb Q$（$\mathbb Z$ 的分式域）或同构于 $\mathbb F_p=\mathbb Z/p\mathbb Z$（$\mathbb Z/p\mathbb Z$ 的分式域）的子域；并且在任一情形中，这个子域都是 $F$ 中含有 $1_F$ 的最小子域（即 $F$ 中由 $1_F$ 生成的域）。

\noindent\textbf{定义.} 域 $F$ 的\emph{素子域}是由 $F$ 的乘法恒等元 $1_F$ 生成的子域。它（同构于）$\mathbb Q$（若 $\operatorname{ch}(F)=0$）或 $\mathbb F_p$（若 $\operatorname{ch}(F)=p$）。

\paragraph{注记：} 我们通常把域 $F$ 的恒等元 $1_F$ 简记为 $1$。于是，在特征为 $p$ 的域中，有 $p\cdot 1=0$，这常常简写为 $p=0$（例如，在特征为 $2$ 的域中 $2=0$）。然而应当记住，这只是一个简写说法——元素“$p$”实际上是 $p\cdot 1_F$，并不是 $F$ 中一个不同的元素。鉴于命题 1 的第二个断言，这种记法是有用的。

\paragraph{例}
\begin{enumerate}[label=(\arabic*)]
	\item $\mathbb Q$ 和 $\mathbb R$ 的素子域都是 $\mathbb Q$。
	\item 域 $\mathbb F_p(x)$ 的素子域同构于 $\mathbb F_p$，由常数多项式给出。
\end{enumerate}

\noindent\textbf{定义.} 若 $K$ 是含有子域 $F$ 的域，则称 $K$ 是 $F$ 的\emph{扩域}（或简称为 $F$ 的\emph{扩张}），记为 $K/F$，也用图
\[
\begin{array}{c}
	K\\
	|\\
	F
\end{array}
\]
表示。特别地，每个域 $F$ 都是它的素子域的扩张。域 $F$ 有时称为该扩张的\emph{基域}。

域扩张的记号 $K/F$ 是“$K$ 在 $F$ 上”的简写，而不是 $K$ 对 $F$ 的商。

若 $K/F$ 是任意域扩张，则 $K$ 中定义的乘法使 $K$ 成为 $F$ 上的向量空间。特别地，每个域 $F$ 都可看作其素域上的向量空间。

\noindent\textbf{定义.} 域扩张 $K/F$ 的\emph{次数}（或\emph{相对次数}，或\emph{指数}），记为 $[K:F]$，是 $K$ 作为 $F$ 上向量空间的维数（即 $[K:F]=\dim_F K$）。若 $[K:F]$ 有限，则称该扩张为\emph{有限}的；否则称为\emph{无限}的。

一类重要的域扩张，是试图在给定域 $F$ 上求解方程而得到的扩张。例如，若 $F=\mathbb R$ 是实数域，那么简单方程 $x^2+1=0$ 在 $F$ 中没有解。于是产生一个问题：是否存在某个含有 $\mathbb R$ 的更大的域，使得这个方程在其中有解？正是这个问题促使 Gauss 引入了复数 $\mathbb C=\mathbb R+\mathbb R i$，其中 $i$ 被定义为满足 $i^2+1=0$。然后按照初等代数中熟悉的通常规则在 $\mathbb C$ 中定义加法和乘法，并检查这样定义的 $\mathbb C$ 事实上是一个域，即每个非零元素都有逆。

给定任意域 $F$ 和任意多项式 $p(x)\in F[x]$，可以提出类似问题：是否存在 $F$ 的扩张 $K$，含有方程 $p(x)=0$ 的一个解（即含有 $p(x)$ 的一个根）？注意，这里可假设多项式 $p(x)$ 在 $F[x]$ 中不可约，因为 $p(x)$ 的任意因子的一个根当然也是 $p(x)$ 的一个根。答案是肯定的，并且几乎立即来自我们关于多项式环 $F[x]$ 的工作。我们先回忆关于域同态的以下有用结果（第 7 章推论 10），它来自域 $F$ 的理想只有 $0$ 和 $F$ 这一事实。

\noindent\textbf{命题 2.} 设 $\varphi:F\longrightarrow F'$ 是域同态。则 $\varphi$ 或恒等于 $0$，或是单射；因此 $\varphi$ 的像或者是 $0$，或者同构于 $F$。

\noindent\textbf{定理 3.} 设 $F$ 是域，$p(x)\in F[x]$ 是不可约多项式。则存在一个含有 $F$ 的一个同构拷贝的域 $K$，使得 $p(x)$ 在 $K$ 中有根。把 $F$ 与这个同构拷贝等同起来，便说明存在 $F$ 的一个扩张，使得 $p(x)$ 有根。

\begin{proof}
	考虑多项式环 $F[x]$ 对由 $p(x)$ 生成的理想取商得到的商
	\[
	K=F[x]/(p(x)).
	\]
	由于假设 $p(x)$ 是 P.I.D. $F[x]$ 中的不可约多项式，理想 $(p(x))$ 是极大理想。因此 $K$ 实际上是一个域（这是第 7 章命题 12）。从 $F[x]$ 到商 $F[x]/(p(x))$ 的典范投影 $\pi$ 限制在 $F\subset F[x]$ 上，给出同态
	\[
	\varphi=\pi|_F:F\longrightarrow K,
	\]
	它并非恒等于 $0$，因为它把 $F$ 的恒等元 $1$ 映到 $K$ 的恒等元 $1$。因此由上面的命题，$\varphi(F)\cong F$ 是包含在 $K$ 中的 $F$ 的一个同构拷贝。我们把 $F$ 与它在 $K$ 中的同构像等同起来，并把 $F$ 看作 $K$ 的一个子域。若 $\bar{x}=\pi(x)$ 表示 $x$ 在商 $K$ 中的像，则
	\[
	\begin{aligned}
		p(\bar{x})&=\overline{p(x)} &&\text{（因为 }\pi\text{ 是同态）}\
		&=p(x)\pmod{p(x)} &&\text{在 }F[x]/(p(x))\text{ 中}\
		&=0 &&\text{在 }F[x]/(p(x))\text{ 中},
	\end{aligned}
	\]
	所以 $K$ 的确含有多项式 $p(x)$ 的一个根。因此 $K$ 是 $F$ 的一个扩张，且 $p(x)$ 在其中有根。
\end{proof}

稍后我们将使用这个结果来构造含有 $p(x)$ 的所有根的 $F$ 的扩张（这就是分裂域的概念，是 Galois 理论中一个核心的研究对象）。

为了更充分地理解上面构造的域 $K=F[x]/(p(x))$，给出这个域中元素的简单表示是有用的。由于 $F$ 是 $K$ 的子域，我们特别可以询问 $K$ 作为 $F$ 上向量空间的一组基。

\noindent\textbf{定理 4.} 设 $p(x)\in F[x]$ 是域 $F$ 上次数为 $n$ 的不可约多项式，设 $K$ 为域 $F[x]/(p(x))$。令
\[
\theta=x\bmod (p(x))\in K.
\]
则元素
\[
1,\theta,\theta^2,\ldots,\theta^{n-1}
\]
构成 $K$ 作为 $F$ 上向量空间的一组基，所以扩张次数为 $n$，即 $[K:F]=n$。因此
\[
K=\{a_0+a_1\theta+a_2\theta^2+\cdots+a_{n-1}\theta^{n-1}\mid a_0,a_1,\ldots,a_{n-1}\in F\}
\]
由所有 $\theta$ 的次数 $<n$ 的多项式组成。

\begin{proof}
	设 $a(x)\in F[x]$ 是任意系数在 $F$ 中的多项式。由于 $F[x]$ 是 Euclidean 整环（这是第 9 章定理 3），我们可以用 $p(x)$ 除 $a(x)$：
	\[
	a(x)=q(x)p(x)+r(x),\qquad q(x),r(x)\in F[x],\quad \deg r(x)<n.
	\]
	由于 $q(x)p(x)$ 属于理想 $(p(x))$，可知
	\[
	a(x)\equiv r(x)\pmod{p(x)},
	\]
	这说明 $F[x]/(p(x))$ 中的每个剩余类都由一个次数小于 $n$ 的多项式表示。因此，$1,x,x^2,\ldots,x^{n-1}$ 在商中的像 $1,\theta,\theta^2,\ldots,\theta^{n-1}$ 张成这个商作为 $F$ 上的向量空间。还需证明这些元素线性无关，从而形成商在 $F$ 上的一组基。
	
	如果 $K$ 中的元素 $1,\theta,\theta^2,\ldots,\theta^{n-1}$ 不是线性无关的，那么存在一个线性组合
	\[
	b_0+b_1\theta+b_2\theta^2+\cdots+b_{n-1}\theta^{n-1}=0
	\]
	在 $K$ 中成立，其中 $b_0,b_1,\ldots,b_{n-1}\in F$ 且不全为 $0$。这等价于
	\[
	b_0+b_1x+b_2x^2+\cdots+b_{n-1}x^{n-1}\equiv 0\pmod{p(x)},
	\]
	即
	\[
	p(x)\text{ 整除 } b_0+b_1x+b_2x^2+\cdots+b_{n-1}x^{n-1}
	\]
	在 $F[x]$ 中成立。但这是不可能的，因为 $p(x)$ 的次数为 $n$，而右边非零多项式的次数 $<n$。这证明 $1,\theta,\theta^2,\ldots,\theta^{n-1}$ 是 $K$ 在 $F$ 上的一组基，所以按照定义 $[K:F]=n$。定理最后的断言显然成立。
\end{proof}

这个定理给出了域 $F[x]/(p(x))$ 中元素的一种简单描述：它们是 $\theta$ 的次数 $<n$ 的多项式，其中 $\theta$ 是一个满足 $p(\theta)=0$ 的元素（在 $K$ 中）。剩下的只是说明如何对这种形式写出的元素作加法和乘法。商 $F[x]/(p(x))$ 中的加法就是通常的多项式加法。商 $F[x]/(p(x))$ 中多项式 $a(x)$ 与 $b(x)$ 的乘法，是先在 $F[x]$ 中求乘积 $a(x)b(x)$，再通过用 $p(x)$ 除 $a(x)b(x)$ 并求余式，找到陪集 $a(x)b(x)+(p(x))$ 的次数 $<n$ 的代表元（如上面的证明所示）。

这也可以容易地用 $\theta$ 来完成，如下：我们可以假设 $p(x)$ 是首一多项式（因为把它乘以一个常数不会改变它的根，也不会改变它生成的理想），设
\[
p(x)=x^n+p_{n-1}x^{n-1}+\cdots+p_1x+p_0.
\]
于是，在 $K$ 中，由于 $p(\theta)=0$，有
\[
\theta^n=-(p_{n-1}\theta^{n-1}+\cdots+p_1\theta+p_0),
\]
即 $\theta^n$ 是 $\theta$ 的较低次幂的线性组合。两边同乘以 $\theta$，并再用这些较低次幂替换右边出现的 $\theta^n$，可见 $\theta^{n+1}$ 也是 $\theta$ 的次数 $<n$ 的多项式。类似地，$\theta$ 的任意正幂都可以写成 $\theta$ 的次数 $<n$ 的多项式，因此 $\theta$ 的任意多项式都可以写成 $\theta$ 的次数 $<n$ 的多项式。现在，在 $K$ 中作乘法很容易：只要把两个次数 $<n$ 的 $\theta$ 多项式的乘积再写成另一个次数 $<n$ 的 $\theta$ 多项式。我们把这总结为：

\noindent\textbf{推论 5.} 设 $K$ 如定理 4 中所述，令 $a(\theta),b(\theta)\in K$ 为 $\theta$ 的两个次数 $<n$ 的多项式。则 $K$ 中的加法简单地定义为通常的多项式加法，而 $K$ 中的乘法定义为
\[
a(\theta)b(\theta)=r(\theta),
\]
其中 $r(x)$ 是在 $F[x]$ 中用 $p(x)$ 除多项式 $a(x)b(x)$ 后得到的余式（次数 $<n$）。

由上面已经证明的结果，这样在 $\theta$ 的次数 $<n$ 的多项式上定义的加法和乘法使 $K$ 成为一个域，所以也可以除以非零元素；这一点从运算定义本身看并不是立即显然的。

定理 4 中多项式 $p(x)$ 在 $F$ 上不可约这一点也很重要。一般来说，推论 5 中的加法和乘法（对任意多项式 $p(x)$ 都可用同样方式定义）若 $p(x)$ 不可约，并不会使 $\theta$ 的次数 $<n$ 的多项式构成一个域。事实上，一般情况下这个集合甚至不是整环（其结构由第 9 章命题 16 给出）。为了描述含有一般多项式 $f(x)$ 在 $F$ 上的一个根 $\theta$ 的域，需要把 $f(x)$ 在 $F[x]$ 中分解为不可约因子，然后对 $f(x)$ 的一个以 $\theta$ 为根的不可约因子 $p(x)$ 应用上面的结果。我们将在以下各节进一步考虑这一点。

\paragraph{例}
\begin{enumerate}[label=(\arabic*)]
	\item 若把这个构造应用到特殊情形 $F=\mathbb R$ 且 $p(x)=x^2+1$，则得到域
	\[
	\mathbb R[x]/(x^2+1),
	\]
	这是 $\mathbb R$ 的一个次数为 $2$ 的扩张，其中 $x^2+1$ 有一个根。这个域的元素形如 $a+b\theta$，其中 $a,b\in\mathbb R$。加法定义为
	\[
	(a+b\theta)+(c+d\theta)=(a+c)+(b+d)\theta. \tag{13.2a}
	\]
	作乘法时使用事实 $\theta^2+1=0$，即 $K$ 中 $\theta^2=-1$。（或者，注意到 $x^2$ 除以 $x^2+1$ 在 $\mathbb R[x]$ 中的余式也是 $-1$。）于是
	\begin{equation}
		\begin{aligned}
			(a+b\theta)(c+d\theta)&=ac+(ad+bc)\theta+bd\theta^2\\
			&=ac+(ad+bc)\theta+bd(-1)\\
			&=(ac-bd)+(ad+bc)\theta.
		\end{aligned}
		\tag{13.2b}
	\end{equation}
	这些公式，差别只在于把 $\theta$ 改写为 $i$，正是 $\mathbb C$ 中加法和乘法的公式。换句话说，映射
	\[
	\varphi:\mathbb R[x]/(x^2+1)\longrightarrow \mathbb C,
	\qquad
	a+bx\longmapsto a+bi
	\]
	是一个同态。由于它是双射（例如作为实数上的向量空间之间的映射），它是一个同构。注意，与其先假定 $\mathbb C$ 的存在（并伴随相当冗长地验证它事实上是一个域），我们本可以通过这个同构来定义 $\mathbb C$。于是它是一个域这一事实就是定理 4 的结果。
	
	\item 现在取 $F=\mathbb Q$ 为有理数域，并再次取 $p(x)=x^2+1$（当然，它仍在 $\mathbb Q$ 上不可约）。则同样的构造给出 $\mathbb Q$ 的一个次数为 $2$ 的域扩张 $\mathbb Q(i)$，它含有 $x^2+1$ 的一个根 $i$；加法和乘法公式与上面的 (13.2a) 和 (13.2b) 相同，只是现在 $a$ 与 $b$ 是 $\mathbb Q$ 的元素。
	
	\item 取 $F=\mathbb Q$ 且 $p(x)=x^2-2$，例如由 Eisenstein 判别法可知它在 $\mathbb Q$ 上不可约。于是我们得到 $\mathbb Q$ 的一个次数为 $2$ 的域扩张，它含有 $2$ 的一个平方根 $\theta$，记为 $\mathbb Q(\theta)$。若用 $\sqrt2$ 表示 $\theta$，则这个域的元素形如
	\[
	a+b\sqrt2,
	\qquad a,b\in\mathbb Q,
	\]
	加法定义为
	\[
	(a+b\sqrt2)+(c+d\sqrt2)=(a+c)+(b+d)\sqrt2,
	\]
	乘法定义为
	\[
	(a+b\sqrt2)(c+d\sqrt2)=(ac+2bd)+(ad+bc)\sqrt2.
	\]
	
	\item 令 $F=\mathbb Q$ 且 $p(x)=x^3-2$，它同样由 Eisenstein 判别法不可约。把 $p(x)$ 的一个根记为 $\theta$，我们得到域
	\[
	\mathbb Q[x]/(x^3-2)\cong\{a+b\theta+c\theta^2\mid a,b,c\in\mathbb Q\},
	\]
	其中 $\theta^3=2$，这是一个次数为 $3$ 的扩张。要在这个域中求例如 $1+\theta$ 的逆，可以如下进行：由 $\mathbb Q[x]$ 中的 Euclidean 算法，存在多项式 $a(x)$ 和 $b(x)$ 使得
	\[
	a(x)(1+x)+b(x)(x^3-2)=1
	\]
	（因为 $p(x)=x^3-2$ 不可约，它与每个次数较低的多项式都互素）。在商域中，这个方程说明 $a(\theta)$ 是 $1+\theta$ 的逆。在此情形下，一个简单计算表明可取 $a(x)=\frac13(x^2-x+1)$（且 $b(x)=-\frac13$），所以
	\[
	(1+\theta)^{-1}=\frac{\theta^2-\theta+1}{3}.
	\]
	
	\item 一般地，若 $\theta\in K$ 是不可约多项式
	\[
	p(x)=p_nx^n+p_{n-1}x^{n-1}+\cdots+p_1x+p_0
	\]
	的根，则可由
	\[
	\theta(p_n\theta^{n-1}+p_{n-1}\theta^{n-2}+\cdots+p_1)=-p_0
	\]
	计算 $\theta^{-1}\in K$，即
	\[
	\theta^{-1}=-\frac{1}{p_0}(p_n\theta^{n-1}+p_{n-1}\theta^{n-2}+\cdots+p_1)\in K
	\]
	（注意 $p_0\neq 0$，因为 $p(x)$ 不可约）。
	
	\paragraph{注记：} 在 $\mathbb C$ 的情形或例 (3) 中，这类扩张里逆的求法可能以“分母有理化”之名在初等代数中为人熟知。最后两个例子指出了一个比初等代数中临时性的做法更为一般的过程。
	
	\item 取 $F=\mathbb F_2$，即含有两个元素的有限域，并取 $p(x)=x^2+x+1$，我们先前已经检验过它在 $\mathbb F_2$ 上不可约。这里得到 $\mathbb F_2$ 的一个次数为 $2$ 的扩张
	\[
	\mathbb F_2[x]/(x^2+x+1)\cong\{a+b\theta\mid a,b\in\mathbb F_2\},
	\]
	其中 $\theta^2=-\theta-1=\theta+1$。这个域 $\mathbb F_2(\theta)$（含有四个元素）中的乘法定义为
	\[
	\begin{aligned}
		(a+b\theta)(c+d\theta)&=ac+(ad+bc)\theta+bd\theta^2\\
		&=ac+(ad+bc)\theta+bd(\theta+1)\\
		&=(ac+bd)+(ad+bc+bd)\theta.
	\end{aligned}
	\]
	
	\item 令 $F=k(t)$ 为域 $k$ 上变量 $t$ 的有理函数域（例如 $k=\mathbb Q$ 或 $k=\mathbb F_p$）。令 $p(x)=x^2-t\in F[x]$。则 $p(x)$ 不可约（它在 $k[t]$ 中关于素元 $(t)$ 是 Eisenstein 的）。若把一个根记为 $\theta$，则相应的次数为 $2$ 的域扩张 $F(\theta)$ 由元素
	\[
	\{a(t)+b(t)\theta\mid a(t),b(t)\in F\}
	\]
	组成，其中系数 $a(t)$ 与 $b(t)$ 是系数在 $k$ 中的 $t$ 的有理函数，并且 $\theta^2=t$。
\end{enumerate}

设 $F$ 是域 $K$ 的一个子域，且 $\alpha\in K$ 是 $K$ 的一个元素。那么 $K$ 中同时含有 $F$ 与 $\alpha$ 的子域组成的集合非空（例如 $K$ 本身就是这样的域）。由于子域的交仍是子域，于是存在唯一的含有 $F$ 与 $\alpha$ 的极小子域（即所有具有这个性质的子域的交）。若把 $\alpha$ 替换为 $K$ 中的一族元素 $\alpha,\beta,\ldots$，同样的说法也成立。

\noindent\textbf{定义.} 设 $K$ 是域 $F$ 的一个扩张，并令 $\alpha,\beta,\ldots\in K$ 是 $K$ 的一族元素。那么 $K$ 中同时含有 $F$ 和元素 $\alpha,\beta,
\ldots$ 的最小子域，记为 $F(\alpha,\beta,\ldots)$，称为由 $\alpha,\beta,\ldots$ 在 $F$ 上生成的域。

\noindent\textbf{定义.} 若域 $K$ 由单个元素 $\alpha$ 在 $F$ 上生成，即 $K=F(\alpha)$，则称 $K$ 是 $F$ 的\emph{单扩张}，元素 $\alpha$ 称为该扩张的\emph{本原元}。

稍后我们将刻画域 $F$ 的哪些扩张是单扩张。特别地，我们将证明，特征为 $0$ 的域的每个有限扩张都是单扩张。

当 $\alpha$ 是某个不可约多项式 $p(x)$ 的根时，由 $\alpha$ 在 $F$ 上生成的单扩张 $F(\alpha)$ 与定理 3 中构造的域之间的联系由下述定理给出：

\noindent\textbf{定理 6.} 设 $F$ 是域，$p(x)\in F[x]$ 是不可约多项式。设 $K$ 是 $F$ 的扩域，其中含有 $p(x)$ 的一个根 $\alpha$：$p(\alpha)=0$。令 $F(\alpha)$ 表示由 $\alpha$ 在 $F$ 上生成的 $K$ 的子域。则
\[
F(\alpha)\cong F[x]/(p(x)).
\]

\paragraph{注记：} 这个定理说的是，$F$ 上任意一个含有 $p(x)$ 的根的域，都含有一个同构于定理 3 中构造的 $F$ 的扩张的子域，并且这个域（在同构意义下）是含有这样一个根的 $F$ 的最小扩张。这个结果与定理 3 的区别在于，定理 6 假设某个域 $K$ 中存在 $p(x)$ 的根 $\alpha$，而定理 3 的要点则是证明确实存在这样的扩域 $K$。

\begin{proof}
	有一个自然同态
	\[
	\varphi:F[x]\longrightarrow F(\alpha)\subseteq K,
	\qquad
	a(x)\longmapsto a(\alpha),
	\]
	它通过把 $F$ 按恒等映射送到 $F$，把 $x$ 送到 $\alpha$，再扩张为环同态而得到（即关于 $x$ 的多项式 $a(x)$ 映到关于 $\alpha$ 的多项式 $a(\alpha)$）。由于假设 $p(\alpha)=0$，元素 $p(x)$ 属于 $\varphi$ 的核，所以得到一个诱导同态（仍记为 $\varphi$）
	\[
	\varphi:F[x]/(p(x))\longrightarrow F(\alpha).
	\]
	但由于 $p(x)$ 不可约，左边的商是一个域，并且 $\varphi$ 不是零映射（例如它在 $F$ 上是恒等映射），所以 $\varphi$ 是左边这个域到其像上的一个同构。因为这个像于是是 $F(\alpha)$ 的一个子域，且含有 $F$ 并含有 $\alpha$，由 $F(\alpha)$ 的定义，该映射必为满射，从而证明了定理。
\end{proof}

结合推论 5，当 $\alpha$ 是不可约多项式 $p(x)$ 的一个根时，这就确定了域 $F(\alpha)$：

\noindent\textbf{推论 7.} 若在定理 6 中 $p(x)$ 的次数为 $n$，则
\[
F(\alpha)=\{a_0+a_1\alpha+a_2\alpha^2+\cdots+a_{n-1}\alpha^{n-1}\mid a_0,a_1,\ldots,a_{n-1}\in F\}\subseteq K.
\]

描述由多于一个元素生成的域要更复杂；我们将在下一节回到这个问题。

\paragraph{例}
\begin{enumerate}[label=(\arabic*)]
	\item 在上面的例 (3) 中，我们已经确定了由元素 $\sqrt2\in\mathbb R$ 在 $\mathbb Q$ 上生成的域 $\mathbb Q(\sqrt2)$；我们富有暗示性地把方程 $x^2-2=0$ 的抽象解 $\theta$ 记为符号 $\sqrt2$，而这个符号在域 $\mathbb R$ 中有独立的含义（即 $2$ 的正平方根）。
	
	\item 方程 $x^2-2=0$ 在 $\mathbb R$ 中还有另一个解，即 $-\sqrt2$，$2$ 的负平方根。由这个解在 $\mathbb Q$ 上生成的域由元素
	\[
	\{a+b(-\sqrt2)\mid a,b\in\mathbb Q\}
	\]
	组成，并且它也同构于上面例 (3) 中的域（因而也同构于刚才考虑的域，其同构显式地由 $a+b\sqrt2\mapsto a-b\sqrt2$ 给出）。作为 $\mathbb R$ 的子集，它与例 (1) 中的元素集合相同。
	
	\item 类似地，若用符号 $\sqrt[3]{2}$ 表示 $\mathbb R$ 中 $2$ 的（正）立方根，则由 $\sqrt[3]{2}$ 在 $\mathbb Q$ 上于 $\mathbb R$ 中生成的域由元素
	\[
	\{a+b\sqrt[3]{2}+c(\sqrt[3]{2})^2\mid a,b,c\in\mathbb Q\}
	\]
	组成，并同构于上面例 (4) 中构造的域。
	
	\item 方程 $x^3-2=0$ 在 $\mathbb R$ 中没有更多解，但在 $\mathbb C$ 中有另外两个解，分别为
	\[
	\sqrt[3]{2}\left(\frac{-1+i\sqrt3}{2}\right)
	\quad\text{和}\quad
	\sqrt[3]{2}\left(\frac{-1-i\sqrt3}{2}\right)
	\]
	（这里 $\sqrt3$ 表示 $3$ 的正实平方根），这可以容易地检验。由这两个元素中的任意一个在 $\mathbb Q$ 上生成的域都是 $\mathbb C$ 的子域（但不是 $\mathbb R$ 的子域），并且二者都同构于前一例中构造的域（也同构于更早的例 (4)）。
\end{enumerate}

如定理 6 所表明的，不可约多项式 $p(x)$ 的根在代数上是不可区分的，其意义是，把不可约多项式的任意一个根添入而得到的域都是同构的。在上面最后两个例子中，把 $x^3-2=0$ 的三个可能的（复）根中任一个添入 $\mathbb Q$ 所得的域在代数上全都同构。这些域不是由它们的代数性质区分的，而是由它们的元素是否为实数来区分；这涉及连续运算。

同一个不可约多项式的不同根具有相同代数性质这一事实还可以稍作推广，如下：

令 $\varphi:F\to F'$ 是域同构。映射 $\varphi$ 诱导一个环同构（仍记为 $\varphi$）
\[
\varphi:F[x]\longrightarrow F'[x]
\]
它通过把 $\varphi$ 作用到 $F[x]$ 中多项式的系数上来定义。令 $p(x)\in F[x]$ 是不可约多项式，并令 $p'(x)\in F'[x]$ 是把映射 $\varphi$ 作用到 $p(x)$ 的系数上所得的多项式，即 $p(x)$ 在 $\varphi$ 下的像。同构 $\varphi$ 把极大理想 $(p(x))$ 映到理想 $(p'(x))$，所以这个理想也极大，从而说明 $p'(x)$ 在 $F'[x]$ 中也不可约。下述定理说明，把 $p(x)$ 的一个根添入 $F$ 所得的域，与把 $p'(x)$ 的一个根添入 $F'$ 所得的域具有相同的代数结构（即同构）：

\noindent\textbf{定理 8.} 设 $\varphi:F\to F'$ 是域同构。设 $p(x)\in F[x]$ 是不可约多项式，$p'(x)\in F'[x]$ 是把映射 $\varphi$ 作用到 $p(x)$ 的系数上所得的不可约多项式。令 $\alpha$ 为 $p(x)$ 的一个根（在 $F$ 的某个扩张中），令 $\beta$ 为 $p'(x)$ 的一个根（在 $F'$ 的某个扩张中）。则存在同构
\[
\sigma:F(\alpha)\longrightarrow F'(\beta),
\qquad
\alpha\longmapsto \beta,
\]
它把 $\alpha$ 映到 $\beta$ 并延拓 $\varphi$，即其限制在 $F$ 上等于同构 $\varphi$。

\begin{proof}
	如上所述，同构 $\varphi$ 诱导从 $F[x]$ 到 $F'[x]$ 的一个自然同构，它把极大理想 $(p(x))$ 映到极大理想 $(p'(x))$。对这些理想取商，得到域同构
	\[
	F[x]/(p(x))\cong F'[x]/(p'(x)).
	\]
	由定理 6，左边的域同构于 $F(\alpha)$；同样由该定理，右边的域同构于 $F'(\beta)$。复合这些同构，得到同构 $\sigma$。显然这个同构限制到 $F$ 上就是 $\varphi$，证明完成。
\end{proof}

这个延拓定理在我们稍后考虑 Galois 理论时将相当有用。它可由图
\[
\begin{array}{ccc}
	\sigma:
	F(\alpha)&\longrightarrow&F'(\beta)\\
	&\alpha\longmapsto\beta&\\[-2pt]
	\varphi:\quad F&\longrightarrow&F'
\end{array}
\]
表示。

\subsection*{练习}
\begin{enumerate}[label=\arabic*.]
	\item 证明 $p(x)=x^3+9x+6$ 在 $\mathbb Q[x]$ 中不可约。令 $\theta$ 为 $p(x)$ 的一个根。求 $\mathbb Q(\theta)$ 中 $1+\theta$ 的逆。
	
	\item 证明 $x^3-2x-2$ 在 $\mathbb Q$ 上不可约，并令 $\theta$ 为其一个根。在 $\mathbb Q(\theta)$ 中计算
	\[
	(1+\theta)(1+\theta+\theta^2)
	\quad\text{和}\quad
	\frac{1+\theta}{1+\theta+\theta^2}.
	\]
	
	\item 证明 $x^3+x+1$ 在 $\mathbb F_2$ 上不可约，并令 $\theta$ 为其一个根。计算 $\mathbb F_2(\theta)$ 中 $\theta$ 的各次幂。
	
	\item 直接证明映射
	\[
	a+b\sqrt2\longmapsto a-b\sqrt2
	\]
	是 $\mathbb Q(\sqrt2)$ 到自身的一个同构。
	
	\item 设 $\alpha$ 是 $\mathbb Z[x]$ 中某个首一多项式的有理根。证明 $\alpha$ 是整数。
	
	\item 证明：若 $\alpha$ 是
	\[
	a_nx^n+a_{n-1}x^{n-1}+\cdots+a_1x+a_0
	\]
	的一个根，则 $a_n\alpha$ 是下列首一多项式的一个根：
	\[
	x^n+a_{n-1}x^{n-1}+a_na_{n-2}x^{n-2}+\cdots+a_n^{n-2}a_1x+a_n^{n-1}a_0.
	\]
	
	\item 证明当 $n\neq -1,3,5$ 时，$x^3-nx+2$ 不可约。
	
	\item 证明 $x^5-ax-1\in\mathbb Z[x]$ 不可约，除非 $a=0,2$ 或 $-1$。前两个值对应线性因子，第三个值对应分解
	\[
	(x^2-x+1)(x^3+x^2-1).
	\]
\end{enumerate}
\section{代数扩张}

设 $F$ 是一个域，设 $K$ 是 $F$ 的一个扩张。

\noindent\textbf{定义.} 元素 $\alpha\in K$ 称为在 $F$ 上是\emph{代数的}，如果 $\alpha$ 是某个非零多项式 $f(x)\in F[x]$ 的根。如果 $\alpha$ 在 $F$ 上不是代数的（即，它不是任意一个系数在 $F$ 中的非零多项式的根），则称 $\alpha$ 在 $F$ 上是\emph{超越的}。扩张 $K/F$ 称为\emph{代数扩张}，如果 $K$ 的每个元素都在 $F$ 上是代数的。

注意，如果 $\alpha$ 在域 $F$ 上是代数的，那么它也在 $F$ 的任意扩域 $L$ 上是代数的（若以 $\alpha$ 为根的 $f(x)$ 的系数在 $F$ 中，则这些系数也在 $L$ 中）。

\noindent\textbf{命题 9.} 设 $\alpha$ 在 $F$ 上是代数的。则存在唯一的首一不可约多项式
\[
m_{\alpha,F}(x)\in F[x]
\]
以 $\alpha$ 为根。一个多项式 $f(x)\in F[x]$ 以 $\alpha$ 为根，当且仅当 $m_{\alpha,F}(x)$ 在 $F[x]$ 中整除 $f(x)$。

\begin{proof}
	设 $g(x)\in F[x]$ 是所有以 $\alpha$ 为根的多项式中次数最小的一个。把 $g(x)$ 乘以一个常数，我们可以假设 $g(x)$ 是首一的。若 $g(x)$ 在 $F[x]$ 中可约，例如
	\[
	g(x)=a(x)b(x),
	\]
	其中 $a(x),b(x)\in F[x]$ 的次数都小于 $g(x)$ 的次数，则在 $K$ 中有
	\[
	g(\alpha)=a(\alpha)b(\alpha).
	\]
	由于 $K$ 是域，或者 $a(\alpha)=0$，或者 $b(\alpha)=0$，这与 $g(x)$ 的次数的最小性矛盾。因此 $g(x)$ 是一个以 $\alpha$ 为根的首一不可约多项式。
	
	现在设 $f(x)\in F[x]$ 是任意一个以 $\alpha$ 为根的多项式。由 $F[x]$ 中的 Euclidean 算法，存在多项式 $q(x),r(x)\in F[x]$ 使得
	\[
	f(x)=q(x)g(x)+r(x),\qquad \deg r(x)<\deg g(x).
	\]
	于是，在 $K$ 中
	\[
	f(\alpha)=q(\alpha)g(\alpha)+r(\alpha).
	\]
	由于 $\alpha$ 同时是 $f(x)$ 和 $g(x)$ 的根，得到 $r(\alpha)=0$；除非 $r(x)=0$，这将与 $g(x)$ 的次数最小性矛盾。因此 $g(x)$ 整除 $F[x]$ 中任意以 $\alpha$ 为根的多项式 $f(x)$；特别地，它将整除任意另一个以 $\alpha$ 为根的首一不可约多项式。这证明了 $m_{\alpha,F}(x)=g(x)$ 的唯一性，并完成命题的证明。
\end{proof}

\noindent\textbf{推论 10.} 若 $L/F$ 是域扩张，并且 $\alpha$ 同时在 $F$ 和 $L$ 上是代数的，则 $m_{\alpha,L}(x)$ 在 $L[x]$ 中整除 $m_{\alpha,F}(x)$。

\begin{proof}
	这立即来自把命题 9 的第二个断言应用于 $L$，因为 $m_{\alpha,F}(x)$ 是 $L[x]$ 中以 $\alpha$ 为根的多项式。
\end{proof}

\noindent\textbf{定义.} 命题 9 中的多项式 $m_{\alpha,F}(x)$（如果域 $F$ 已由上下文明确，也简记为 $m_\alpha(x)$）称为 $\alpha$ 在 $F$ 上的\emph{极小多项式}。$m_\alpha(x)$ 的次数称为 $\alpha$ 的\emph{次数}。

注意，由命题可知，$F$ 上一个以 $\alpha$ 为根的首一多项式是 $\alpha$ 在 $F$ 上的极小多项式，当且仅当它在 $F$ 上不可约。练习 20 给出了计算 $\alpha$ 在 $F$ 上的极小多项式的一种方法；而 Gröbner 基理论可以用来计算 $F(\alpha)$ 中其他元素的极小多项式（参见第 15.1 节命题 10 与练习 48）。

\noindent\textbf{命题 11.} 设 $\alpha$ 在域 $F$ 上是代数的，设 $F(\alpha)$ 是由 $\alpha$ 在 $F$ 上生成的域。则
\[
F(\alpha)\cong F[x]/(m_\alpha(x)),
\]
因而特别地
\[
[F(\alpha):F]=\deg m_\alpha(x)=\deg \alpha,
\]
也就是说，$\alpha$ 在 $F$ 上的次数就是它生成的扩张的次数。

\begin{proof}
	这立即来自定理 6。
\end{proof}

\paragraph{例}
\begin{enumerate}[label=(\arabic*)]
	\item $\sqrt2$ 在 $\mathbb Q$ 上的极小多项式是 $x^2-2$，并且 $\sqrt2$ 在 $\mathbb Q$ 上的次数为 $2$：
	\[
	[\mathbb Q(\sqrt2):\mathbb Q]=2.
	\]
	\item $\sqrt[3]{2}$ 在 $\mathbb Q$ 上的极小多项式是 $x^3-2$，并且 $\sqrt[3]{2}$ 在 $\mathbb Q$ 上的次数为 $3$：
	\[
	[\mathbb Q(\sqrt[3]{2}):\mathbb Q]=3.
	\]
	\item 类似地，对任意 $n>1$，多项式 $x^n-2$ 在 $\mathbb Q$ 上不可约，因为它满足 Eisenstein 判别法。把这个多项式的一个根记为 $\sqrt[n]{2}$（照常，若我们想把这个根看作 $\mathbb R$ 中的元素，就保留这个符号表示 $2$ 的正 $n$ 次根；而一般地，这个符号表示代数上不可区分的抽象解中的任意一个），则
	\[
	[\mathbb Q(\sqrt[n]{2}):\mathbb Q]=n.
	\]
	\item 一个元素 $\alpha$ 的极小多项式和次数依赖于基域。例如，在 $\mathbb R$ 上，元素 $\sqrt[n]{2}$ 的次数为 $1$，其极小多项式为
	\[
	m_{\alpha,\mathbb R}(x)=x-\sqrt[n]{2}.
	\]
	\item 考虑 $\mathbb Q$ 上的多项式
	\[
	p(x)=x^3-3x-1,
	\]
	它在 $\mathbb Q$ 上不可约，因为它是一个没有有理根的三次多项式（参见第 9 章命题 11）。因此，对 $p(x)$ 的任意根 $\alpha$，都有 $[\mathbb Q(\alpha):\mathbb Q]=3$。为了以后引用，我们注意到，在实数上快速画出这个函数的图像可见，该图像在区间 $[0,2]$ 中恰好与 $x$ 轴相交一次；也就是说，恰有一个实数 $\alpha$，$0<\alpha<2$，满足
	\[
	\alpha^3-3\alpha-1=0.
	\]
\end{enumerate}

\noindent\textbf{命题 12.} 元素 $\alpha$ 在 $F$ 上是代数的，当且仅当单扩张 $F(\alpha)/F$ 是有限的。更精确地，若 $\alpha$ 是 $F$ 上一个次数为 $n$ 的扩张中的元素，则 $\alpha$ 满足一个在 $F$ 上次数至多为 $n$ 的多项式；而若 $\alpha$ 满足一个在 $F$ 上次数为 $n$ 的多项式，则 $F(\alpha)$ 在 $F$ 上的次数至多为 $n$。

\begin{proof}
	若 $\alpha$ 在 $F$ 上是代数的，则扩张 $F(\alpha)/F$ 的次数就是 $\alpha$ 在 $F$ 上的极小多项式的次数。因此该扩张是有限的；若 $\alpha$ 满足一个次数为 $n$ 的多项式，则该扩张次数至多为 $n$。反过来，假设 $\alpha$ 是 $F$ 上次数为 $n$ 的一个扩张中的元素（例如，若 $[F(\alpha):F]=n$）。那么 $F(\alpha)$ 中的 $n+1$ 个元素
	\[
	1,\alpha,\alpha^2,
	\ldots,\alpha^n
	\]
	在 $F$ 上线性相关，也就是说，存在不全为 $0$ 的 $b_0,b_1,b_2,\ldots,b_n\in F$，使得
	\[
	b_0+b_1\alpha+b_2\alpha^2+\cdots+b_n\alpha^n=0.
	\]
	因此 $\alpha$ 是一个系数在 $F$ 中的非零多项式（次数 $\leq n$）的根，这证明 $\alpha$ 在 $F$ 上是代数的，也证明了命题的第二个断言。
\end{proof}

\noindent\textbf{推论 13.} 若扩张 $K/F$ 是有限的，则它是代数扩张。

\begin{proof}
	若 $\alpha\in K$，则子域 $F(\alpha)$ 特别是 $K$ 作为 $F$ 上向量空间的一个子空间。因此 $[F(\alpha):F]\leq [K:F]$，所以由命题可知 $\alpha$ 在 $F$ 上是代数的。
\end{proof}

\paragraph{注记：} 下面我们将证明这个结果的某种逆命题（定理 17），但注意到存在无限的代数扩张（稍后我们将看到一个例子），所以这个推论的字面逆命题不为真。

\paragraph{例：（特征不为 $2$ 的域的二次扩张）}
设 $F$ 是一个特征不为 $2$ 的域（例如，任意特征为 $0$ 的域，如 $\mathbb Q$），设 $K$ 是 $F$ 的一个次数为 $2$ 的扩张，$[K:F]=2$。令 $\alpha$ 为 $K$ 中任意一个不属于 $F$ 的元素。由上面的命题，$\alpha$ 满足一个在 $F$ 上次数至多为 $2$ 的方程。这个方程不可能次数为 $1$，因为按假设 $\alpha$ 不是 $F$ 的元素。因此 $\alpha$ 的极小多项式是一个首一二次多项式
\[
m_\alpha(x)=x^2+bx+c,\qquad b,c\in F.
\]
由于 $F\subset F(\alpha)\subset K$，并且 $F(\alpha)$ 作为 $F$ 上向量空间已经是 $2$ 维的，我们有
\[
K=F(\alpha).
\]

这个二次方程的根可以由二次公式确定；该公式在任意特征不为 $2$ 的域上都是有效的（这个公式像初等代数中一样由配方得到）：
\[
\alpha=\frac{-b\pm\sqrt{b^2-4c}}{2}
\]
（要求 $F$ 的特征不为 $2$ 的原因是我们必须除以 $2$）。这里 $b^2-4c$ 在 $F$ 中不是平方，因为 $\alpha$ 不是 $F$ 的元素；符号 $\sqrt{b^2-4c}$ 表示 $K$ 中方程
\[
x^2-(b^2-4c)=0
\]
的一个根（见下一段末尾）。注意，在这里并没有一个类似于在 $\mathbb R$ 中选取 $2$ 的正平方根的自然选择；这些根在代数上是不可区分的。

现在 $F(\alpha)=F(\sqrt{b^2-4c})$，理由如下：由上面的公式，$\alpha$ 是右边这个域的元素，所以
\[
F(\alpha)\subset F(\sqrt{b^2-4c}).
\]
反过来，
\[
\sqrt{b^2-4c}=\pm(b+2\alpha)
\]
说明 $\sqrt{b^2-4c}$ 是 $F(\alpha)$ 的元素，于是得到反向包含
\[
F(\sqrt{b^2-4c})\subset F(\alpha),
\]
并且顺便说明方程 $x^2-(b^2-4c)=0$ 在 $K$ 中确实有一个解。

由此可知，$F$ 的任意次数为 $2$ 的扩张 $K$ 都具有形式 $F(\sqrt D)$，其中 $D$ 是 $F$ 中一个不是平方的元素；反过来，每个这样的扩张都是 $F$ 的次数为 $2$ 的扩张。因此，域 $F$ 的次数为 $2$ 的扩张称为 $F$ 的\emph{二次扩张}。

假设 $F$ 是域 $K$ 的子域，而 $K$ 又是域 $L$ 的子域。于是有三个相关的扩张次数——$K$ 和 $L$ 作为 $F$ 上向量空间的维数，以及 $L$ 作为 $K$ 上向量空间的维数。

\noindent\textbf{定理 14.} 设 $F\subset K\subset L$ 是域。则
\[
[L:F]=[L:K][K:F],
\]
也就是说，扩张次数是可乘的；其中，如果等式一边为无限，则另一边也为无限。用图表示为
\[
\begin{array}{ccccc}
	& & [L:F] & & \\
	F & \subset & K & \subset & L\\
	\underbrace{\hspace{2.0cm}}_{[K:F]} & & & \underbrace{\hspace{2.0cm}}_{[L:K]}
\end{array}
\]

\begin{proof}
	首先假设 $[L:K]=m$ 且 $[K:F]=n$ 都有限。令 $\alpha_1,\alpha_2,\ldots,\alpha_m$ 是 $L$ 在 $K$ 上的一组基，令 $\beta_1,\beta_2,\ldots,\beta_n$ 是 $K$ 在 $F$ 上的一组基。那么 $L$ 的每个元素都可写成线性组合
	\[
	a_1\alpha_1+a_2\alpha_2+\cdots+a_m\alpha_m,
	\]
	其中 $a_1,\ldots,a_m$ 是 $K$ 的元素，因此又是 $\beta_1,\ldots,\beta_n$ 的 $F$-线性组合：
	\[
	a_i=b_{i1}\beta_1+b_{i2}\beta_2+\cdots+b_{in}\beta_n,
	\qquad i=1,2,
	\ldots,m, \tag{13.3}
	\]
	其中 $b_{ij}$ 是 $F$ 的元素。把这些表达式代入上面 $a_i$ 的位置可见，$L$ 的每个元素都可写成
	\[
	\sum_{\substack{i=1,2,\ldots,m\\ j=1,2,\ldots,n}} b_{ij}\alpha_i\beta_j
	\]
	的线性组合，系数在 $F$ 中。因此这 $mn$ 个元素 $\alpha_i\beta_j$ 张成 $L$ 作为 $F$ 上的向量空间。
	
	现在假设在 $L$ 中有一个线性关系
	\[
	\sum_{\substack{i=1,2,\ldots,m\\ j=1,2,\ldots,n}} b_{ij}\alpha_i\beta_j=0,
	\]
	其中 $b_{ij}$ 在 $F$ 中。按上面的公式 (13.3) 定义元素 $a_i\in K$，则这个线性关系可写成
	\[
	a_1\alpha_1+a_2\alpha_2+\cdots+a_m\alpha_m=0.
	\]
	由于 $\alpha_i$ 是 $L$ 在 $K$ 上的一组基，所有系数 $a_i$，$i=1,2,\ldots,m$，都必须为 $0$；也就是说，
	\[
	b_{i1}\beta_1+b_{i2}\beta_2+\cdots+b_{in}\beta_n=0,
	\qquad i=1,2,\ldots,m,
	\]
	是在 $K$ 中的等式。由于现在 $\beta_j$，$j=1,2,\ldots,n$，构成 $K$ 在 $F$ 上的一组基，这推出所有 $i,j$ 都有 $b_{ij}=0$。因此元素 $\alpha_i\beta_j$ 在 $F$ 上线性无关，所以形成 $L$ 在 $F$ 上的一组基，并且
	\[
	[L:F]=mn=[L:K][K:F],
	\]
	如所要证明。
	
	若 $[K:F]$ 为无限，则 $K$ 中有无限多个在 $F$ 上线性无关的元素，因而 $L$ 中也有这些元素，所以 $[L:F]$ 也是无限的。类似地，若 $[L:K]$ 为无限，则 $L$ 中有无限多个在 $K$ 上线性无关的元素，当然也在 $F$ 上线性无关，所以 $[L:F]$ 仍为无限。最后，若 $[L:K]$ 和 $[K:F]$ 都有限，则上面的证明说明 $[L:F]$ 有限；因此 $[L:F]$ 为无限意味着 $[L:K]$ 与 $[K:F]$ 中至少有一个为无限。证明完成。
\end{proof}

\paragraph{注记：} 注意这个结果与第一部分中关于群的阶的结果相似。与涉及群的图表一样，我们将经常在域图中标出扩张的相对次数。

扩张次数的可乘性在计算中极为有用。一个特别的应用如下：

\noindent\textbf{推论 15.} 设 $L/F$ 是有限扩张，设 $K$ 是 $L$ 的任意含有 $F$ 的子域，$F\subset K\subset L$。则 $[K:F]$ 整除 $[L:F]$。

\begin{proof}
	这是立即的。
\end{proof}

\paragraph{例}
\begin{enumerate}[label=(\arabic*)]
	\item 元素 $\sqrt2$ 不属于域 $\mathbb Q(\alpha)$，其中 $\alpha$ 是 $x^3-3x-1$ 在 $0$ 与 $2$ 之间的实根；因为我们已经确定
	\[
	[\mathbb Q(\sqrt2):\mathbb Q]=2,
	\qquad [\mathbb Q(\alpha):\mathbb Q]=3,
	\]
	而 $2$ 不整除 $3$。注意，直接证明 $\sqrt2$ 不能写成 $1,\alpha,\alpha^2$ 的有理线性组合并不那么容易。
	\item 照常令 $\sqrt[6]{2}$ 表示 $2$ 的正实六次根。则
	\[
	[\mathbb Q(\sqrt[6]{2}):\mathbb Q]=6.
	\]
	由于 $(\sqrt[6]{2})^3=\sqrt2$，我们有
	\[
	\mathbb Q(\sqrt2)\subset \mathbb Q(\sqrt[6]{2}),
	\]
	并且由扩张次数的可乘性，
	\[
	[\mathbb Q(\sqrt[6]{2}):\mathbb Q(\sqrt2)]=3.
	\]
	这给出域图
	\[
	\begin{array}{ccccc}
		\mathbb Q & \subset & \mathbb Q(\sqrt2) & \subset & \mathbb Q(\sqrt[6]{2})\\
		& 2 & & 3 & \\
		\multicolumn{5}{c}{[\mathbb Q(\sqrt[6]{2}):\mathbb Q]=6.}
	\end{array}
	\]
	特别地，这说明 $\sqrt[6]{2}$ 在 $\mathbb Q(\sqrt2)$ 上的极小多项式次数为 $3$。因此它就是多项式
	\[
	x^3-\sqrt2.
	\]
	注意，直接说明这个多项式在 $\mathbb Q(\sqrt2)$ 上不可约并非完全平凡。
\end{enumerate}

由定理 14，有限扩张的有限扩张仍为有限。接下来的结果使用这一点来说明，由有限多个代数元素生成的扩张是有限的（推广命题 12）。

\noindent\textbf{定义.} 扩张 $K/F$ 称为\emph{有限生成的}，如果存在 $K$ 中的元素 $\alpha_1,\alpha_2,\ldots,\alpha_k$，使得
\[
K=F(\alpha_1,\alpha_2,\ldots,\alpha_k).
\]

回忆：在域 $K$ 中由一组元素在 $F$ 上生成的域，是 $K$ 中含有这些元素和 $F$ 的最小子域。下一个引理将说明，对于有限生成扩张，这个域可以通过一系列单扩张递归地得到。

\noindent\textbf{引理 16.} $F(\alpha,\beta)=(F(\alpha))(\beta)$，即，由 $\alpha$ 和 $\beta$ 在 $F$ 上生成的域，就是由 $\beta$ 在 $\alpha$ 生成的域 $F(\alpha)$ 上生成的域。

\begin{proof}
	这来自所讨论的域的最小性。域 $F(\alpha,\beta)$ 含有 $F$ 和 $\alpha$，因此含有域 $F(\alpha)$；又由于它也含有 $\beta$，由域 $(F(\alpha))(\beta)$ 的最小性得到包含
	\[
	(F(\alpha))(\beta)\subset F(\alpha,\beta).
	\]
	由于域 $(F(\alpha))(\beta)$ 含有 $F$、$\alpha$ 和 $\beta$，由 $F(\alpha,\beta)$ 的最小性得到反向包含
	\[
	F(\alpha,\beta)\subset (F(\alpha))(\beta),
	\]
	这证明了引理。
\end{proof}

由引理，我们有
\[
K=F(\alpha_1,\alpha_2,\ldots,\alpha_k)
=(F(\alpha_1,\alpha_2,\ldots,\alpha_{k-1}))(\alpha_k),
\]
于是通过迭代可见，$K$ 是这样得到的：先取由 $\alpha_1$ 在 $F$ 上生成的域 $F_1$，然后取由 $\alpha_2$ 在 $F_1$ 上生成的域 $F_2$（这一点很重要），依次进行，直到 $F_k=K$。这给出一列域：
\[
F=F_0\subset F_1\subset F_2\subset \cdots\subset F_k=K,
\]
其中
\[
F_{i+1}=F_i(\alpha_{i+1}),\qquad i=0,1,\ldots,k-1.
\]

现在假设元素 $\alpha_1,\alpha_2,\ldots,\alpha_k$ 在 $F$ 上都是代数的，次数分别为 $n_1,n_2,\ldots,n_k$（于是先验地它们在 $F$ 的任意扩域上也都是代数的）。那么这列扩张是命题 11 中所考虑的那种单扩张。相对扩张次数 $[F_{i+1}:F_i]$ 等于 $\alpha_{i+1}$ 在 $F_i$ 上的极小多项式次数，该次数至多为 $n_{i+1}$（并且当且仅当 $\alpha_{i+1}$ 在 $F$ 上的极小多项式在 $F_i$ 上仍不可约时，等于 $n_{i+1}$）。由扩张次数的可乘性，我们看到
\[
[K:F]=[F_k:F_{k-1}][F_{k-1}:F_{k-2}]\cdots [F_1:F_0]
\]
也是有限的，并且 $\leq n_1n_2\cdots n_k$。

这也给出了 $F(\alpha_1,\alpha_2,\ldots,\alpha_k)$ 中元素的描述。为简单起见，考虑域 $F(\alpha,\beta)$ 的情形，其中 $\alpha$ 和 $\beta$ 都在 $F$ 上代数。那么这个域的元素具有形式
\[
b_0+b_1\beta+b_2\beta^2+\cdots+b_{d-1}\beta^{d-1},
\]
其中
\[
d=[F(\alpha)(\beta):F(\alpha)]
\]
是 $\beta$ 在 $F(\alpha)$ 上的次数（它可能严格小于 $\beta$ 在 $F$ 上的次数），而系数 $b_0,b_1,\ldots,b_{d-1}$ 是 $F(\alpha)$ 的元素。系数 $b_i\in F(\alpha)$，$i=0,\ldots,d-1$，具有形式
\[
a_{0i}+a_{1i}\alpha+a_{2i}\alpha^2+\cdots+a_{n-1,i}\alpha^{n-1},
\]
其中
\[
n=[F(\alpha):F]
\]
是 $\alpha$ 在 $F$ 上的次数，且 $a_{ij}$ 是 $F$ 的元素。因此 $F(\alpha,\beta)$ 的元素具有形式
\[
\sum_{\substack{i=0,1,\ldots,n-1\\ j=0,1,\ldots,d-1}} a_{ij}\alpha^i\beta^j,
\qquad a_{ij}\in F.
\]
由于
\[
[F(\alpha,\beta):F]=[F(\alpha,\beta):F(\alpha)][F(\alpha):F]=dn,
\]
元素 $\alpha^i\beta^j$ 事实上构成 $F(\alpha,\beta)$ 在 $F$ 上的一组基。

在实践中，由代数元素 $\alpha$ 生成的域 $F(\alpha)$ 是通过把元素 $\alpha$ 添到 $F$ 中，然后对所得集合在加法和乘法下“闭合”而得到的；这相当于添入 $\alpha$ 的幂 $\alpha^2,\alpha^3,
\ldots$，并取这些元素的（系数来自 $F$ 的）线性组合。当 $\alpha$ 的某个幂是它较低次幂的线性组合时，这个过程终止，而这等价于知道 $\alpha$ 的极小多项式。前面的讨论表明，类似的过程给出由两个元素生成的域 $F(\alpha,\beta)$；递归地，也给出由任意有限个代数元素生成的域。这特别说明，对代数元素而言，关于加法和乘法“闭合”也会关于除法闭合（参见上面推论 5 后的例 5）。若元素不是代数的，则还必须关于取逆“闭合”。这个过程中的困难在于确定相对扩张的次数——例如上面 $F(\alpha,\beta)$ 在 $F(\alpha)$ 上的次数 $d$，对此我们只有先验的上界（即 $\beta$ 在 $F$ 上的次数）。

这类似于在一个群 $G$ 中对一组元素进行“闭合”以确定它们生成的子群。

\paragraph{例}
\begin{enumerate}[label=(\arabic*)]
	\item 扩张 $\mathbb Q(\sqrt[3]{2},\sqrt2)$ 就是扩张 $\mathbb Q(\sqrt[3]{2})$，因为 $\sqrt2$ 已经是这个域的元素。换言之，$\sqrt2$ 在 $\mathbb Q(\sqrt[3]{2})$ 上的次数 $d$ 为 $1$，它严格小于 $\sqrt2$ 在 $\mathbb Q$ 上的次数。稍后我们会看到不那么显然地发生这种情形的例子。
	\item 考虑由 $\sqrt2$ 和 $\sqrt3$ 在 $\mathbb Q$ 上生成的域 $\mathbb Q(\sqrt2,\sqrt3)$。由于 $\sqrt3$ 在 $\mathbb Q$ 上次数为 $2$，扩张
	\[
	\mathbb Q(\sqrt2,\sqrt3)/\mathbb Q(\sqrt2)
	\]
	的次数至多为 $2$，并且它恰好为 $2$ 当且仅当 $x^2-3$ 在 $\mathbb Q(\sqrt2)$ 上不可约。由于这个多项式次数为 $2$，它可约当且仅当它有一个根，即当且仅当 $\sqrt3\in\mathbb Q(\sqrt2)$。假设
	\[
	\sqrt3=a+b\sqrt2,
	\qquad a,b\in\mathbb Q.
	\]
	两边平方得到
	\[
	3=(a^2+2b^2)+2ab\sqrt2.
	\]
	若 $ab\neq0$，则可由这个等式把 $\sqrt2$ 用 $a$ 和 $b$ 表出，从而推出 $\sqrt2$ 是有理数，这是不可能的。若 $b=0$，则有 $\sqrt3=a$ 是有理数，矛盾。最后，若 $a=0$，则有 $\sqrt3=b\sqrt2$；两边乘以 $\sqrt2$ 可见 $\sqrt6$ 是有理数，又矛盾。这说明 $\sqrt3\notin\mathbb Q(\sqrt2)$，从而证明
	\[
	[\mathbb Q(\sqrt2,\sqrt3):\mathbb Q]=4.
	\]
	这个域中的元素（由对 $1,\sqrt2,\sqrt3$ “闭合”得到）包括 $1,\sqrt2,\sqrt3,\sqrt6$；由上面的计算，这些元素构成该域的一组基：
	\[
	\mathbb Q(\sqrt2,\sqrt3)
	=\{a+b\sqrt2+c\sqrt3+d\sqrt6\mid a,b,c,d\in\mathbb Q\}.
	\]
\end{enumerate}

现在我们可以刻画域 $F$ 的有限扩张：

\noindent\textbf{定理 17.} 扩张 $K/F$ 是有限的，当且仅当 $K$ 由有限多个在 $F$ 上代数的元素生成。更精确地，由有限多个在 $F$ 上次数为 $n_1,n_2,\ldots,n_k$ 的代数元素生成的域，是次数 $\leq n_1n_2\cdots n_k$ 的代数扩张。

\begin{proof}
	若 $K/F$ 是次数为 $n$ 的有限扩张，令 $\alpha_1,\alpha_2,\ldots,\alpha_n$ 是 $K$ 作为 $F$ 上向量空间的一组基。由推论 15，$[F(\alpha_i):F]$ 整除 $[K:F]=n$，对 $i=1,2,\ldots,n$ 成立；所以命题 12 说明每个 $\alpha_i$ 都在 $F$ 上代数。由于 $K$ 显然由 $\alpha_1,\alpha_2,\ldots,\alpha_n$ 在 $F$ 上生成，我们看到 $K$ 由有限多个在 $F$ 上代数的元素生成。反过来的方向已经在上面证明。定理的第二个断言立即来自推论 13 和上面的计算。
\end{proof}

上面的第一个例子说明，定理中给出的扩张次数的不等式可能是严格的。我们指出，对这个次数的确定有帮助的信息常常可以通过确定子域并应用推论 15 得到。

\noindent\textbf{推论 18.} 设 $\alpha$ 和 $\beta$ 在 $F$ 上是代数的。则
\[
\alpha\pm\beta,
\quad \alpha\beta,
\quad \alpha/\beta\ (\beta\neq0),
\]
（特别地，当 $\alpha\neq0$ 时 $\alpha^{-1}$）都在 $F$ 上是代数的。

\begin{proof}
	所有这些元素都属于扩张 $F(\alpha,\beta)$，而由定理该扩张在 $F$ 上有限；于是由推论 13，它们都是代数的。
\end{proof}

\noindent\textbf{推论 19.} 设 $L/F$ 是任意扩张。则 $L$ 中所有在 $F$ 上代数的元素构成 $L$ 的一个子域 $K$。

\begin{proof}
	这立即来自前一个推论。
\end{proof}

\paragraph{例}
\begin{enumerate}[label=(\arabic*)]
	\item 考虑扩张 $\mathbb C/\mathbb Q$，令 $\overline{\mathbb Q}$ 表示 $\mathbb C$ 中所有在 $\mathbb Q$ 上代数的元素所构成的子域。特别地，元素 $\sqrt[n]{2}$（$2$ 在 $\mathbb R$ 中的正 $n$ 次根）全都属于 $\overline{\mathbb Q}$，所以对所有整数 $n>1$ 都有
	\[
	[\overline{\mathbb Q}:\mathbb Q]\geq n.
	\]
	因此 $\overline{\mathbb Q}$ 是 $\mathbb Q$ 的一个无限代数扩张，称为\emph{代数数域}。
	\item 考虑域 $\overline{\mathbb Q}\cap\mathbb R$，即 $\mathbb R$ 中所有在 $\mathbb Q$ 上代数的元素构成的子域。域 $\mathbb Q$ 是可数的。$\mathbb Q[x]$ 中任意给定次数 $n$ 的多项式的个数因此也是可数的（因为这样的多项式由从 $\mathbb Q$ 中指定 $n+1$ 个系数决定）。由于这些多项式在 $\mathbb R$ 中最多有 $n$ 个根，在 $\mathbb R$ 中次数为 $n$ 的代数元素的个数是可数的。最后，$\mathbb R$ 中所有代数元素的集合是可数个（由 $n$ 指标化的）可数集的并，因此是可数的。由于 $\mathbb R$ 是不可数的，可知存在（事实上有很多）不是代数的元素，即在 $\mathbb Q$ 上超越的元素。因此 $\mathbb R$ 中代数元素的子域 $\overline{\mathbb Q}\cap\mathbb R$ 是 $\mathbb R$ 的一个真子域，所以 $\overline{\mathbb Q}$ 也是 $\mathbb C$ 的一个真子域。
	
	一般来说，证明一个给定实数不是代数的是极其困难的。例如，已知（这些是定理）
	\[
	\pi=3.14159\ldots
	\quad\text{和}\quad
	e=2.71828\ldots
	\]
	是 $\mathbb R$ 的超越元素。甚至证明这些元素不是有理数也并不太容易。
\end{enumerate}

\noindent\textbf{定理 20.} 若 $K$ 在 $F$ 上代数，且 $L$ 在 $K$ 上代数，则 $L$ 在 $F$ 上代数。

\begin{proof}
	令 $\alpha$ 为 $L$ 的任意元素。则 $\alpha$ 在 $K$ 上是代数的，所以 $\alpha$ 满足某个多项式方程
	\[
	\alpha^n+a_{n-1}\alpha^{n-1}+\cdots+a_1\alpha+a_0=0,
	\]
	其中系数 $a_0,a_1,\ldots,a_{n-1}$ 在 $K$ 中。考虑由 $\alpha$ 和这个多项式的系数在 $F$ 上生成的域
	\[
	F(\alpha,a_0,a_1,\ldots,a_{n-1}).
	\]
	由于 $K/F$ 是代数的，元素 $a_0,a_1,\ldots,a_{n-1}$ 都在 $F$ 上代数，所以由定理 17，扩张
	\[
	F(a_0,a_1,\ldots,a_{n-1})/F
	\]
	是有限的。由上面的方程可见，$\alpha$ 在这个域上生成一个次数至多为 $n$ 的扩张，因为它在这个域上的极小多项式整除上面的多项式。因此
	\[
	\begin{aligned}
		&[F(\alpha,a_0,a_1,\ldots,a_{n-1}):F] \\
		&\qquad = [F(\alpha,a_0,\ldots,a_{n-1}):F(a_0,
		\ldots,a_{n-1})]
		[F(a_0,\ldots,a_{n-1}):F]
	\end{aligned}
	\]
	也是有限的，并且 $F(\alpha,a_0,a_1,\ldots,a_{n-1})/F$ 是代数扩张。特别地，元素 $\alpha$ 在 $F$ 上是代数的，这证明 $L$ 在 $F$ 上代数。
\end{proof}

由域 $K$ 中有限集合 $\alpha_1,\alpha_2,\ldots,\alpha_k$ 生成的子域 $F(\alpha_1,\alpha_2,\ldots,\alpha_k)$ 包含每个域 $F(\alpha_i)$，$i=1,2,\ldots,k$。由定义，它也是 $K$ 中含有这些域的最小子域。

\noindent\textbf{定义.} 设 $K_1$ 和 $K_2$ 是域 $K$ 的两个子域。则 $K_1$ 与 $K_2$ 的\emph{复合域}，记为 $K_1K_2$，是 $K$ 中含有 $K_1$ 和 $K_2$ 二者的最小子域。类似地，$K$ 的任意一族子域的复合域，是含有所有这些子域的最小子域。

注意，复合域 $K_1K_2$ 也可描述为 $K$ 中所有含有 $K_1$ 和 $K_2$ 二者的子域的交；对多于两个域的复合同样如此。这类似于一个群的子集生成的子群（参见第 2.4 节）。

\paragraph{例}
两个域 $\mathbb Q(\sqrt2)$ 和 $\mathbb Q(\sqrt[3]{2})$ 的复合域是域 $\mathbb Q(\sqrt[6]{2})$。这是因为这个域含有上述两个子域（$(\sqrt[6]{2})^3=\sqrt2$ 且 $(\sqrt[6]{2})^2=\sqrt[3]{2}$）；反过来，任何同时含有 $\sqrt2$ 和 $\sqrt[3]{2}$ 的域都含有它们的商，即 $\sqrt[6]{2}$。

现在假设 $K_1$ 和 $K_2$ 是包含在 $K$ 中的 $F$ 的有限扩张。令 $\alpha_1,\alpha_2,\ldots,\alpha_n$ 是 $K_1$ 在 $F$ 上的一组基，令 $\beta_1,\beta_2,\ldots,\beta_m$ 是 $K_2$ 在 $F$ 上的一组基（于是 $[K_1:F]=n$ 且 $[K_2:F]=m$）。显然这些元素给出复合域 $K_1K_2$ 在 $F$ 上的生成元：
\[
K_1K_2=F(\alpha_1,\alpha_2,\ldots,\alpha_n,
\beta_1,\beta_2,\ldots,\beta_m).
\]
由于 $\alpha_1,\alpha_2,\ldots,\alpha_n$ 是 $K_1$ 的一组 $F$-基，任意一个 $\alpha$ 的幂 $\alpha_i^k$ 都是 $\alpha$ 们的系数在 $F$ 中的线性组合；对 $\beta$ 们也有类似的说法。于是，所有形如
\[
\sum_{\substack{i=1,2,\ldots,n\\ j=1,2,\ldots,m}} a_{ij}\alpha_i\beta_j
\]
且系数在 $F$ 中的线性组合所成的集合，在乘法和加法下闭合，因为两个这样的元素相乘时，$\alpha$ 们与 $\beta$ 们的任意较高次幂都可以由线性表达式替换。因此元素 $\alpha_i\beta_j$，$i=1,2,\ldots,n$ 且 $j=1,2,\ldots,m$，张成复合扩张 $K_1K_2$ 在 $F$ 上的向量空间。特别地，
\[
[K_1K_2:F]
\leq mn.
\]
我们把这总结如下：

\noindent\textbf{命题 21.} 设 $K_1$ 和 $K_2$ 是包含在 $K$ 中的域 $F$ 的两个有限扩张。则
\[
[K_1K_2:F]
\leq [K_1:F][K_2:F],
\]
并且等号成立当且仅当其中一个域的一组 $F$-基在另一个域上仍线性无关。若 $\alpha_1,\alpha_2,\ldots,\alpha_n$ 和 $\beta_1,\beta_2,\ldots,\beta_m$ 分别是 $K_1$ 和 $K_2$ 在 $F$ 上的基，则元素 $\alpha_i\beta_j$（$i=1,2,\ldots,n$，$j=1,2,\ldots,m$）张成 $K_1K_2$ 在 $F$ 上的向量空间。

\begin{proof}
	由
	\[
	K_1K_2=F(\alpha_1,\alpha_2,\ldots,\alpha_n,
	\beta_1,\beta_2,\ldots,\beta_m)=K_1(\beta_1,\beta_2,\ldots,\beta_m),
	\]
	如上可见 $\beta_1,\beta_2,\ldots,\beta_m$ 张成 $K_1K_2$ 在 $K_1$ 上的向量空间。因此
	\[
	[K_1K_2:K_1]\leq m=[K_2:F],
	\]
	并且等号成立当且仅当这些元素在 $K_1$ 上线性无关。由于
	\[
	[K_1K_2:F]=[K_1K_2:K_1][K_1:F],
	\]
	这证明了命题。
\end{proof}

由命题（以及它的证明），我们有下面的图：
\[
\begin{array}{ccccc}
	& & K_1K_2 & & \\
	& {\scriptstyle \leq m}\nearrow & & \nwarrow{\scriptstyle \leq n} & \\
	K_1 & & & & K_2\\
	& {\scriptstyle n}\nwarrow & & \nearrow{\scriptstyle m} & \\
	& & F & &
\end{array}
\]

稍后我们会有例子说明命题中的不等式可以是严格的。一个可以确定有等号成立的有用情形如下：

\noindent\textbf{推论 22.} 在命题 21 中，假设 $[K_1:F]=n$，$[K_2:F]=m$，其中 $n$ 与 $m$ 互素：$(n,m)=1$。则
\[
[K_1K_2:F]=[K_1:F][K_2:F]=nm.
\]

\begin{proof}
	一般地，扩张次数 $[K_1K_2:F]$ 被 $n$ 和 $m$ 同时整除，因为 $K_1$ 和 $K_2$ 是 $K_1K_2$ 的子域；因此它被它们的最小公倍数整除。在这里，由于 $(n,m)=1$，这意味着 $[K_1K_2:F]$ 被 $nm$ 整除。结合命题中的不等式
	\[
	[K_1K_2:F]\leq nm
	\]
	即得推论。
\end{proof}

\paragraph{例}
两个域 $\mathbb Q(\sqrt2)$ 与 $\mathbb Q(\sqrt[3]{2})$ 的复合域在 $\mathbb Q$ 上次数为 $6$；我们先前通过实际计算复合域 $\mathbb Q(\sqrt[6]{2})$ 已经确定了这一点。

\subsection*{练习}
\begin{enumerate}[label=\arabic*.]
	\item 设 $\mathbb F$ 是一个特征为 $p$ 的有限域。证明 $|\mathbb F|=p^n$，其中 $n$ 是某个正整数。
	
	\item 设 $g(x)=x^2+x-1$，$h(x)=x^3-x+1$。通过把 $f(x)$ 的一个根添加到域 $F$ 中，得到含有 $4,8,9$ 和 $27$ 个元素的域，其中 $f(x)=g(x)$ 或 $h(x)$，且 $F=\mathbb F_2$ 或 $\mathbb F_3$。写出含有 $4$ 个元素和 $9$ 个元素的域的乘法表，并说明它们的非零元素构成一个循环群。
	
	\item 确定元素 $1+i$ 在 $\mathbb Q$ 上的极小多项式。
	
	\item 确定 $2+\sqrt3$ 和 $1+\sqrt[3]{2}+\sqrt[3]{4}$ 在 $\mathbb Q$ 上的次数。
	
	\item 令 $F=\mathbb Q(i)$。证明 $x^3-2$ 和 $x^3-3$ 在 $F$ 上不可约。
	
	\item 直接从定义证明，域 $F(\alpha_1,\alpha_2,\ldots,\alpha_n)$ 是域 $F(\alpha_1),F(\alpha_2),\ldots,F(\alpha_n)$ 的复合域。
	
	\item 证明
	\[
	\mathbb Q(\sqrt2+\sqrt3)=\mathbb Q(\sqrt2,\sqrt3)
	\]
	（一个包含是显然的；另一个包含可考虑 $(\sqrt2+\sqrt3)^2$ 等）。推出
	\[
	[\mathbb Q(\sqrt2+\sqrt3):\mathbb Q]=4.
	\]
	找出一个由 $\sqrt2+\sqrt3$ 满足的不可约多项式。
	
	\item 设 $F$ 是特征不为 $2$ 的域。设 $D_1$ 和 $D_2$ 是 $F$ 中的元素，二者在 $F$ 中都不是平方。证明：若 $D_1D_2$ 在 $F$ 中不是平方，则 $F(\sqrt{D_1},\sqrt{D_2})$ 在 $F$ 上次数为 $4$；否则在 $F$ 上次数为 $2$。当 $F(\sqrt{D_1},\sqrt{D_2})$ 在 $F$ 上次数为 $4$ 时，该域称为 $F$ 的\emph{双二次扩张}。
	
	\item 设 $F$ 是特征不为 $2$ 的域。设 $a,b$ 是域 $F$ 中的元素，并且 $b$ 在 $F$ 中不是平方。证明存在 $F$ 中的某些 $m,n$ 使得
	\[
	\sqrt{a+\sqrt b}=\sqrt m+\sqrt n
	\]
	的充要条件是 $a^2-b$ 在 $F$ 中是平方。用这点判定域
	\[
	\mathbb Q(\sqrt{a+\sqrt b})\qquad (a,b\in\mathbb Q)
	\]
	何时在 $\mathbb Q$ 上是双二次的。
	
	\item 确定扩张
	\[
	\mathbb Q(\sqrt{3+2\sqrt2})
	\]
	在 $\mathbb Q$ 上的次数。
	
	\item
	\begin{enumerate}[label=(\alph*)]
		\item 令 $\sqrt{3+4i}$ 表示复数 $3+4i$ 的位于第一象限的平方根，令 $\sqrt{3-4i}$ 表示 $3-4i$ 的位于第四象限的平方根。证明
		\[
		[\mathbb Q(\sqrt{3+4i}+\sqrt{3-4i}):\mathbb Q]=1.
		\]
		\item 确定扩张
		\[
		\mathbb Q(\sqrt{1+\sqrt{-3}}+\sqrt{1-\sqrt{-3}})
		\]
		在 $\mathbb Q$ 上的次数。
	\end{enumerate}
	
	\item 假设扩张 $K/F$ 的次数是素数 $p$。证明 $K$ 中任意含有 $F$ 的子域 $E$ 或者是 $K$，或者是 $F$。
	
	\item 假设
	\[
	F=\mathbb Q(\alpha_1,\alpha_2,\ldots,\alpha_n),
	\]
	其中 $\alpha_i^2\in\mathbb Q$，$i=1,2,\ldots,n$。证明 $\sqrt[3]{2}\notin F$。
	
	\item 证明：若 $[F(\alpha):F]$ 是奇数，则
	\[
	F(\alpha)=F(\alpha^2).
	\]
	
	\item 如果 $-1$ 在域 $F$ 中不能表示为平方和，则称域 $F$ 是\emph{形式实}的。设 $F$ 是形式实域，设 $f(x)\in F[x]$ 是奇次数不可约多项式，设 $\alpha$ 是 $f(x)$ 的一个根。证明 $F(\alpha)$ 也是形式实的。[选取一个次数最小的反例 $\alpha$。证明对某些 $p_i(x),g(x)\in F[x]$ 有
	\[
	-1+f(x)g(x)=(p_1(x))^2+\cdots+(p_m(x))^2,
	\]
	其中 $g(x)$ 的次数为奇数且 $<\deg f$。证明 $g$ 的某个根 $\beta$ 在 $F$ 上次数为奇数，并且 $F(\beta)$ 不是形式实的，从而违背 $\alpha$ 的最小性。]
	
	\item 设 $K/F$ 是代数扩张，设 $R$ 是一个包含在 $K$ 中且含有 $F$ 的环。证明 $R$ 是 $K$ 的含有 $F$ 的一个子域。
	
	\item 设 $f(x)$ 是域 $F$ 上次数为 $n$ 的不可约多项式。设 $g(x)$ 是 $F[x]$ 中任意多项式。证明复合多项式
	\[
	f(g(x))
	\]
	的每个不可约因子的次数都被 $n$ 整除。
	
	\item 设 $k$ 是域，设 $k(x)$ 是系数在 $k$ 中的关于 $x$ 的有理函数域。令 $t\in k(x)$ 为有理函数
	\[
	t=\frac{P(x)}{Q(x)},
	\]
	其中 $P(x),Q(x)\in k[x]$ 互素，且 $Q(x)\neq0$。则 $k(x)$ 是 $k(t)$ 的扩张；为了计算其次数，需要计算 $x$ 满足的、系数在 $k(t)$ 中的极小多项式。
	\begin{enumerate}[label=(\alph*)]
		\item 证明变量 $X$ 中、系数在 $k(t)$ 中的多项式
		\[
		P(X)-tQ(X)
		\]
		在 $k(t)$ 上不可约，并且以 $x$ 为根。[由 Gauss 引理，该多项式在 $(k(t))[X]$ 中不可约当且仅当它在 $(k[t])[X]$ 中不可约。然后注意到
		\[
		(k[t])[X]=(k[X])[t].
		\]]
		\item 证明 $P(X)-tQ(X)$ 作为变量 $X$ 中、系数在 $k(t)$ 中的多项式的次数，是 $P(x)$ 与 $Q(x)$ 的次数的最大值。
		\item 证明
		\[
		[k(x):k(t)]=\left[k(x):k\!\bigg(\frac{P(x)}{Q(x)}\bigg)\right]
		=\max(\deg P(x),\deg Q(x)).
		\]
	\end{enumerate}
	
	\item 设 $K$ 是 $F$ 的一个次数为 $n$ 的扩张。
	\begin{enumerate}[label=(\alph*)]
		\item 对任意 $\alpha\in K$，证明由左乘作用在 $K$ 上的 $\alpha$ 是 $K$ 的一个 $F$-线性变换。
		\item 证明 $K$ 同构于 $F$ 上 $n\times n$ 矩阵环的一个子域；因此，$F$ 上 $n\times n$ 矩阵环含有 $F$ 的每个次数 $\leq n$ 的扩张的一个同构拷贝。
	\end{enumerate}
	
	\item 说明：若前一练习中考虑的线性变换“乘以 $\alpha$”的矩阵为 $A$，则 $\alpha$ 是 $A$ 的特征多项式的一个根。这给出了确定 $F$ 的次数为 $n$ 的扩张中元素 $\alpha$ 所满足的次数为 $n$ 的方程的有效步骤。用这个步骤得到 $\sqrt[3]{2}$ 以及 $1+\sqrt[3]{2}+\sqrt[3]{4}$ 所满足的次数为 $3$ 的首一多项式。
	
	\item 设 $K=\mathbb Q(\sqrt D)$，其中 $D$ 是某个无平方因子的整数。令 $\alpha=a+b\sqrt D$ 是 $K$ 的一个元素。使用 $1,\sqrt D$ 作为 $K$ 作为 $\mathbb Q$ 上向量空间的一组基，说明前面练习中考虑的 $K$ 上线性变换“乘以 $\alpha$”具有矩阵
	\[
	\begin{pmatrix}
		a & bD\\
		b & a
	\end{pmatrix}.
	\]
	直接证明映射
	\[
	a+b\sqrt D\longmapsto
	\begin{pmatrix}
		a & bD\\
		b & a
	\end{pmatrix}
	\]
	是域 $K$ 到 $\mathbb Q$ 上 $2\times2$ 矩阵环的一个子域的同构。
	
	\item 设 $K_1$ 和 $K_2$ 是包含在域 $K$ 中的域 $F$ 的两个有限扩张。证明 $F$-代数 $K_1\otimes_F K_2$ 是一个域，当且仅当
	\[
	[K_1K_2:F]=[K_1:F][K_2:F].
	\]
\end{enumerate}
\section{经典的直尺与圆规作图}

作为我们已经得到的关于代数扩张的结果，特别是关于扩张次数的可乘性的一个简单应用，我们可以（以否定的方式）回答希腊人提出的下列几何问题：
\begin{enumerate}[label=\Roman*.]
	\item \emph{（倍立方体）}是否可能仅用直尺与圆规作出一个立方体，使其体积恰好为给定立方体体积的两倍？
	\item \emph{（三等分角）}是否可能仅用直尺与圆规三等分任意给定角 $\theta$？
	\item \emph{（化圆为方）}是否可能仅用直尺与圆规作出一个正方形，使其面积恰好为给定圆的面积？
\end{enumerate}
为回答这些问题，我们必须把用圆规与直尺作出长度这件事翻译成代数术语。设 $1$ 表示一个固定的给定单位距离。于是任意距离都由它的长度 $a\in\mathbb R$ 决定，这使我们可以把几何距离看作实数 $\mathbb R$ 的元素。用给定的单位距离 $1$ 来定义坐标轴上的尺度，我们便可以构造通常的 Cartesian 平面 $\mathbb R^2$，并把我们的所有作图都看作发生在 $\mathbb R^2$ 中。于是，从给定距离 $1$ 出发，一个点 $(x,y)\in\mathbb R^2$ 可作，当且仅当它的坐标 $x$ 与 $y$ 是 $\mathbb R$ 的可作元素。上面的问题于是归结为确定 $\mathbb R$ 中某些特定长度能否由一个固定单位距离出发用圆规与直尺作图得到。这样的实数的全体连同它们的负数将称为 $\mathbb R$ 的\emph{可作元素}；我们将不区分可作长度与可作实数。

每个直尺与圆规作图都由以下四种类型的一系列操作组成：（1）用一条直线连接两个给定点，（2）求两条直线的交点，（3）以给定半径和圆心作圆，以及（4）求一条直线与一个圆的交点，或两个圆的交点。

如果给定两个长度 $a$ 和 $b$，则可用直尺与圆规作出长度 $a\pm b$、$ab$ 和 $a/b$，这是几何中的初等事实（前两个是显然的，后两个由平行线的构造给出（图 1））。

\begin{center}
	\begin{tikzpicture}[scale=0.85,>=stealth]
		% left construction: product
		\coordinate (A) at (0,0);
		\coordinate (B) at (3.6,0);
		\coordinate (C) at (4.4,1.6);
		\coordinate (D) at (1.5,1.05);
		\draw (A)--(B)--(C)--cycle;
		\draw (A)--(D);
		\draw (1.5,0)--(D);
		\draw[<->] (0,-0.45)--node[below]{$1$}(1.5,-0.45);
		\draw[<->] (0,-1.05)--node[below]{$a$}(3.6,-1.05);
		\draw[<->] (-0.15,0.05)--node[above,sloped]{$b$}(1.5,1.05);
		\draw[<->] (-0.4,0.18)--node[above,sloped]{$ab$}(4.15,1.95);
		% right construction: quotient
		\begin{scope}[xshift=6.4cm]
			\coordinate (A) at (0,0);
			\coordinate (B) at (2.2,0);
			\coordinate (C) at (3.2,1.45);
			\coordinate (D) at (0.9,1.0);
			\coordinate (E) at (0.35,0.4);
			\draw (A)--(B)--(C)--cycle;
			\draw (A)--(D);
			\draw (E)--(2.25,1.05);
			\draw[<->] (0,-0.45)--node[below]{$1$}(1.6,-0.45);
			\draw[<->] (0,-1.05)--node[below]{$b$}(3.2,-1.05);
			\draw[<->] (-0.2,0.12)--node[left]{$\frac ab$}(0.75,1.28);
			\draw[<->] (-0.55,0.25)--node[left]{$a$}(0.25,2.0);
		\end{scope}
	\end{tikzpicture}
	
	\textbf{图 1}
\end{center}

如果给定 $a$，则作出 $\sqrt a$ 也是一个初等几何作图：按图 2 所示，作直径为 $1+a$ 的圆，并在直径上作垂线。则 $\sqrt a$ 就是这条垂线的长度。

\begin{center}
	\begin{tikzpicture}[scale=1.0,>=stealth]
		\draw (0,0) arc (180:0:2.7 and 1.55);
		\draw (0,0)--(5.4,0);
		\coordinate (P) at (3.75,0);
		\coordinate (Q) at (3.75,1.45);
		\draw (P)--(Q);
		\draw[<->] ($(P)+(-0.28,0)$)--node[left]{$\sqrt a$}($(Q)+(-0.28,0)$);
		\draw[<->] (0,-0.18)--node[below]{$a$}(3.75,-0.18);
		\draw[<->] (3.75,-0.18)--node[below]{$1$}(5.4,-0.18);
	\end{tikzpicture}\qquad \textbf{图 2}
\end{center}

由此可知，直尺与圆规作图给出了实数中的加法、减法、乘法和除法（除以非零元素）这些全部代数运算，所以可作元素的全体是 $\mathbb R$ 的一个子域。还可以对可作元素开平方。我们现在将看到，这些本质上是仅有的可能操作。

由给定长度 $1$ 出发，通过这些操作可以作出所有有理数 $\mathbb Q$。因此我们可以作出所有坐标为有理数的点 $(x,y)\in\mathbb R^2$。通过开平方，我们可以作出 $\mathbb R$ 的附加元素，所以从 $1$ 出发可作的 $\mathbb R$ 的元素全体形成一个严格大于 $\mathbb Q$ 的域。

连接两个坐标在某个域 $F$ 中的点的直线的通常公式（“两点式”）给出直线的方程，其形式为
\[
ax+by-c=0,
\]
其中 $a,b,c\in F$。联立求解两个这样的方程以确定两条这样的直线的交点，所得解也在 $F$ 中。因此，如果两个点的坐标属于域 $F$，则仅用直尺作图不会产生坐标不也在 $F$ 中的额外点。

圆规作图（上面的类型 (3) 或 (4)）定义的是由一个圆与一条直线或另一个圆相交所得到的点。圆心为 $(h,k)$、半径为 $r$ 的圆有方程
\[
(x-h)^2+(y-k)^2=r^2,
\]
所以，当我们考虑圆规作图对域 $F$ 的元素的影响时，我们是在考虑这样的一个方程与一个线性方程
\[
ax+by-c=0
\]
的联立解，其中 $a,b,c,h,k,r\in F$，或者两个二次方程的联立解。

在线性方程与圆的方程的情形中，例如从线性方程中解出 $y$，并代入圆方程，会得到关于 $x$ 的二次方程（而 $y$ 由 $x$ 线性给出）。因此交点的坐标至多位于 $F$ 的一个二次扩张中。

在两个圆相交的情形中，设为
\[
(x-h)^2+(y-k)^2=r^2
\]
和
\[
(x-h')^2+(y-k')^2=r'^2,
\]
用第二个方程从第一个方程中相减，说明考虑下面两个方程会得到相同的交点：
\[
(x-h)^2+(y-k)^2=r^2
\]
以及
\[
2(h'-h)x+2(k'-k)y=r^2-h^2-k^2-r'^2+h'^2+k'^2,
\]
这正是前面所考虑类型的一个圆与一条直线的交点（事实上，这条直线连接两个交点）。

由此可知，如果给定一族可作元素，则可以作出由这些元素生成的 $\mathbb R$ 的子域 $F$ 中的所有元素；并且对 $F$ 的元素所作的任意直尺与圆规操作，所产生的元素至多位于 $F$ 的一个二次扩张中。由于二次扩张的次数为 $2$，且扩张次数是可乘的，因此如果 $\alpha\in\mathbb R$ 是由域 $F$ 中的元素经过一系列（有限个）直尺与圆规操作得到的，那么 $\alpha$ 属于 $F$ 的一个次数为 $2$ 的幂的扩张 $K$ 中：对某个 $m$，有
\[
[K:F]=2^m.
\]
由于 $[F(\alpha):F]$ 整除这个扩张次数，所以它也必须是 $2$ 的幂。

\noindent\textbf{命题 23.} 若元素 $\alpha\in\mathbb R$ 是由一个域 $F\subset\mathbb R$ 经由一系列圆规与直尺作图得到的，则对某个整数 $k\geq 0$，有
\[
[F(\alpha):F]=2^k.
\]

\noindent\textbf{定理 24.} 经典希腊问题：（I）倍立方体，（II）三等分角，以及（III）化圆为方，均不可能完成。

\begin{proof}
	（I）倍立方体等价于从单位 $1$ 出发在实数中作出 $\sqrt[3]{2}$。由于
	\[
	[\mathbb Q(\sqrt[3]{2}):\mathbb Q]=3
	\]
	不是 $2$ 的幂，这是不可能的。
	
	（II）如果一个角 $\theta$ 可作，那么在 $\mathbb R^2$ 中确定从原点出发、与正 $x$ 轴成角 $\theta$ 且距离为 $1$ 的点，说明 $\cos\theta$（该点的 $x$-坐标）可作（于是 $\sin\theta$ 也可作）。反过来，如果 $\cos\theta$，从而 $\sin\theta$，可作，则具有这些坐标的点给出角 $\theta$。
	
	于是三等分角 $\theta$ 的问题等价于如下问题：给定 $\cos\theta$，作出 $\cos(\theta/3)$。
	
	为说明这并非总是可能的（它当然有时是可能的，例如 $\theta=180^\circ$），考虑 $\theta=60^\circ$。于是 $\cos\theta=\frac12$。由余弦的三倍角公式：
	\[
	\cos\theta=4\cos^3(\theta/3)-3\cos(\theta/3),
	\]
	代入 $\theta=60^\circ$，可见 $\beta=\cos20^\circ$ 满足方程
	\[
	4\beta^3-3\beta-1/2=0,
	\]
	或者
	\[
	8\beta^3-6\beta-1=0.
	\]
	这可写成
	\[
	(2\beta)^3-3(2\beta)-1=0.
	\]
	令 $\alpha=2\beta$。则 $\alpha$ 是一个介于 $0$ 与 $2$ 之间的实数，并满足方程
	\[
	\alpha^3-3\alpha-1=0.
	\]
	但我们在上一节已经考察过这个方程，并确定 $[\mathbb Q(\alpha):\mathbb Q]=3$；同上可见 $\alpha$ 不可作。
	
	（III）化圆为方等价于确定实数
	\[
	\pi=3.14159\ldots
	\]
	是否可作。如前所述，甚至证明这个数不是有理数也是一个困难问题。事实上它是超越的（我们将不加证明地假设这一点），所以 $[\mathbb Q(\pi):\mathbb Q]$ 甚至不是有限的，更不可能是 $2$ 的幂，这说明用直尺与圆规化圆为方是不可能的。
\end{proof}

\paragraph{注记：} 上面的证明说明 $\cos20^\circ$ 和 $\sin20^\circ$ 不能作出。问题随之出现：哪些整数角（以度为单位测量）是可作的？角 $1^\circ$ 和 $2^\circ$ 不可作，因为否则正弦和余弦的加法公式将给出 $20^\circ$ 的可作性。另一方面，初等几何作图（作正五边形得到 $72^\circ$ 的角，作等边三角形得到 $60^\circ$ 的角）连同加法公式和半角公式说明 $\cos3^\circ$ 与 $\sin3^\circ$ 是可作的。由此可知，一个整数度角的三角函数可作，当且仅当这个角是 $3^\circ$ 的倍数。明确地，
\[
\cos 3^\circ
=\frac18(\sqrt3+1)\sqrt{5+\sqrt5}
+\frac1{16}(\sqrt6-\sqrt2)(\sqrt5-1)
\]
并且
\[
\sin 3^\circ
=\frac1{16}(\sqrt6+\sqrt2)(\sqrt5-1)
-\frac18(\sqrt3-1)\sqrt{5+\sqrt5},
\]
这说明它们是由 $\mathbb Q$ 经过连续开平方和域运算得到的。

在第 14.5 节讨论分圆域之后，我们将考虑另一个经典几何问题：“哪些正 $n$ 边形可以用直尺与圆规作出？”（参见命题 14.29）。

这里我们一直谨慎地考虑使用直尺而不是带刻度的尺子来作图，区别在于带刻度的尺子上有标记。如果使用带刻度的尺子，则可以作出许多附加的代数元素。例如，假设 $\theta$ 是一个给定角，并且单位距离 $1$ 被标在尺子上。如图 3 所示，作一个半径为 $1$、圆心角为 $\theta$ 的圆，然后滑动尺子，直到圆上点 $A$ 与 $B$ 之间的距离为 $1$。于是一些初等几何说明（参见练习），所示角 $\alpha$ 为 $\theta/3$，即这个构造（归功于 Archimedes）三等分 $\theta$。特别地，定理 24 中第二个经典问题（三等分角）可以用带刻度的尺子和圆规解决。

\begin{center}
	\begin{tikzpicture}[scale=1.1,>=stealth]
		\draw (-0.5,0)--(6.2,0);
		\coordinate (O) at (3.9,0);
		\draw (O) ++(-1.4,0) arc (180:0:1.4);
		\draw (O)--++(58:1.4);
		\draw (O)--++(-1.4,0);
		\coordinate (A) at (0.25,0);
		\coordinate (B) at (2.05,0.58);
		\coordinate (C) at (4.6,1.25);
		\draw (A)--(C);
		\draw (B)--(O);
		\draw[<->] (A)++(12:0.0)--node[above]{$1$}++(22:1.3);
		\draw (A) ++(0.55,0) arc (0:22:0.55);
		\node at (0.95,0.16) {$\alpha$};
		\draw (O) ++(0.7,0) arc (0:58:0.7);
		\node at (4.55,0.22) {$\theta$};
		\node at (1.92,0.36) {$B$};
		\node[below] at (A) {$A$};
		\node[below] at (O) {$1$};
	\end{tikzpicture}\qquad \textbf{图 3}
\end{center}

定理 24 中第一个经典问题（倍立方体），即作出 $\sqrt[3]{2}$，也可以用带刻度的尺子和圆规解决。下面给出对任意给定的正实数 $k<1$ 作出 $k^{1/3}$ 的构造。这个构造是 J. H. Conway 告诉我们的。

作半径为 $1$ 的圆，并以点 $A=(k,0)$ 为圆心，作出点
\[
B=(0,\sqrt{1-k^2}).
\]
把这个距离三等分，作出点
\[
(0,\tfrac13\sqrt{1-k^2}),
\]
并画出连接此点与 $A$ 的直线。滑动带有标记单位长度 $1$ 的尺子，使它经过点 $B$，并使从交点 $C$ 到它与 $x$-轴的交点 $D$ 的距离为 $1$，如图 4 所示。

于是 $A$ 与 $D$ 之间的距离为 $2k^{1/3}$，而 $B$ 与 $C$ 之间的距离为 $2k^{2/3}$（参见练习）。

\begin{center}
	\begin{tikzpicture}[scale=1.0,>=stealth]
		\coordinate (B) at (0,3.0);
		\coordinate (A) at (1.65,0.85);
		\coordinate (C) at (3.35,1.85);
		\coordinate (D) at (5.2,0.85);
		\coordinate (E) at (0,0);
		\coordinate (F) at (0,0.85);
		\draw (E)--(D);
		\draw (E)--(B);
		\draw (B)--(D);
		\draw (A)--(B);
		\draw (-0.55,0.0)--(4.3,2.45);
		\node[above] at (B) {$B$};
		\node[below] at (A) {$A$};
		\node[right] at (C) {$C$};
		\node[right] at (D) {$D$};
		\node at (0.8,1.03) {$k$};
		\node at (0.96,2.04) {$1$};
		\node at (4.25,1.45) {$1$};
		\node[left] at (0,1.9) {$\sqrt{1-k^2}$};
		\node[left] at (0,0.43) {$\frac13\sqrt{1-k^2}$};
	\end{tikzpicture}\qquad \textbf{图 4}
\end{center}

\subsection*{练习}

\begin{enumerate}[label=\arabic*.]
	\item 证明不可能作出正九边形。
	
	\item 证明 Archimedes 的构造确实三等分角 $\theta$。\,[注意图 5 中的等腰三角形，以证明 $\beta=\gamma=2\alpha$。]
	
	\begin{center}
		\begin{tikzpicture}[scale=1.05,>=stealth]
			\draw (-0.5,0)--(6.2,0);
			\coordinate (O) at (3.9,0);
			\draw (O) ++(-1.4,0) arc (180:0:1.4);
			\draw (O)--++(58:1.4);
			\coordinate (A) at (0.25,0);
			\coordinate (B) at (2.05,0.58);
			\coordinate (C) at (4.6,1.25);
			\draw (A)--(C);
			\draw (B)--(O);
			\draw[dashed] (B)--(O);
			\draw[<->] (A)--node[above]{$1$}(B);
			\draw[<->] (O)--node[below]{$1$}(5.15,0);
			\draw (A) ++(0.55,0) arc (0:22:0.55);
			\node at (0.95,0.16) {$\alpha$};
			\draw (B) ++(0.45,0.13) arc (18:-20:0.45);
			\node at (2.45,0.43) {$\beta$};
			\node at (4.25,1.06) {$\gamma$};
			\draw (O) ++(0.7,0) arc (0:58:0.7);
			\node at (4.55,0.22) {$\theta$};
		\end{tikzpicture}\qquad \textbf{图 5}
	\end{center}
	
	\item 证明文中所示 Conway 的构造确实作出了 $2k^{1/3}$ 和 $2k^{2/3}$。\,[一种方法：令 $(x,y)$ 为点 $C$ 的坐标，令 $a$ 为从 $B$ 到 $C$ 的距离，$b$ 为从 $A$ 到 $D$ 的距离；用相似三角形证明
	\[
	\text{(a)}\quad \frac{y}{1}=\frac{\sqrt{1-k^2}}{1+a},
	\]
	\[
	\text{(b)}\quad \frac{x}{a}=\frac{b+k}{1+a},
	\qquad
	\text{(c)}\quad \frac{y}{x-k}=\frac{\sqrt{1-k^2}}{3k},
	\]
	并且还证明
	\[
	\text{(d)}\quad (1-k^2)+(b+k)^2=(1+a)^2;
	\]
	解这些方程求出 $a$ 和 $b$。]
	
	\item 正七边形的构造等价于 $\cos(2\pi/7)$ 的可作性。我们以后将看到（第 14.5 节和第 14.7 节练习 2），
	\[
	\alpha=2\cos(2\pi/7)
	\]
	满足方程
	\[
	x^3+x^2-2x-1=0.
	\]
	利用这一点证明正七边形不能由直尺与圆规作出。
	
	\item 利用 $\alpha=2\cos(2\pi/5)$ 满足方程
	\[
	x^2+x-1=0
	\]
	这一事实，推出正五边形可以用直尺与圆规作出。
\end{enumerate}
\section{分裂域与代数闭包}

设 $F$ 是一个域。

若 $f(x)$ 是 $F[x]$ 中的任意多项式，则我们已经在第 2 节看到，存在一个域 $K$，它可以（通过把 $F$ 与 $F$ 的一个同构副本等同起来）看作 $F$ 的一个扩张，并且在其中 $f(x)$ 有一个根 $\alpha$。这等价于说 $f(x)$ 在 $K[x]$ 中有一个线性因子 $x-\alpha$（这是第 9 章命题 9）。

\noindent\textbf{定义.} 域 $F$ 的扩域 $K$ 称为多项式 $f(x)\in F[x]$ 的一个\emph{分裂域}，如果 $f(x)$ 在 $K[x]$ 中完全分解为线性因子（或称在 $K[x]$ 中完全分裂），并且 $f(x)$ 不在任何含有 $F$ 的 $K$ 的真子域上完全分解为线性因子。


若 $f(x)$ 的次数为 $n$，则 $f(x)$ 在 $F$ 中至多有 $n$ 个根（第 9 章命题 17），并且当且仅当 $f(x)$ 在 $F[x]$ 中完全分裂时，$f(x)$ 在 $F$ 中恰有 $n$ 个根（按重数计）。

\noindent\textbf{定理 25.} 对任意域 $F$，若 $f(x)\in F[x]$，则存在 $F$ 的一个扩张 $K$，它是 $f(x)$ 的分裂域。


\begin{proof}
	我们首先对 $f(x)$ 的次数 $n$ 用归纳法证明：存在 $F$ 的一个扩张 $E$，在它上面 $f(x)$ 完全分解为线性因子。若 $n=1$，则取 $E=F$。现在设 $n>1$。若 $f(x)$ 在 $F$ 上的不可约因子全都为次数 $1$，则 $F$ 就是 $f(x)$ 的分裂域，因而可取 $E=F$。否则，$f(x)$ 在 $F[x]$ 中至少有一个不可约因子，比如 $p(x)$，次数至少为 $2$。由定理 3，存在 $F$ 的一个扩张 $E_1$，它含有 $p(x)$ 的一个根 $\alpha$。在 $E_1$ 上，多项式 $f(x)$ 有线性因子 $x-\alpha$。$f(x)$ 的其余因子 $f_1(x)$ 的次数为 $n-1$，所以由归纳假设，存在 $E_1$ 的一个扩张 $E$，它含有 $f_1(x)$ 的全部根。由于 $\alpha\in E$，$E$ 是 $F$ 的一个扩张，且含有 $f(x)$ 的全部根。现在令 $K$ 为 $E$ 的所有含有 $F$ 且同时含有 $f(x)$ 全部根的子域的交。则 $K$ 是一个域，并且是 $f(x)$ 的分裂域。
\end{proof}

我们很快将看到，$f(x)$ 的任意两个分裂域都是同构的（这推广了定理 8），所以我们常常（滥用语言地）说一个多项式的\emph{这个}分裂域。

\noindent\textbf{定义.} 若 $K$ 是 $F$ 的代数扩张，并且 $K$ 是 $F$ 上一族多项式 $f(x)\in F[x]$ 的分裂域，则称 $K$ 是 $F$ 的一个\emph{正规扩张}。


我们通常使用术语“分裂域”而不是“正规扩张”（也见第 14.9 节）。


	\begin{enumerate}[label=(\arabic*)]
		\item $x^2-2$ 在 $\mathbb Q$ 上的分裂域正是 $\mathbb Q(\sqrt2)$，因为两个根是 $\pm\sqrt2$，而 $-\sqrt2\in\mathbb Q(\sqrt2)$。
		
		\item $(x^2-2)(x^2-3)$ 的分裂域是由 $\sqrt2$ 和 $\sqrt3$ 在 $\mathbb Q$ 上生成的域 $\mathbb Q(\sqrt2,\sqrt3)$，因为该多项式的根是 $\pm\sqrt2,\pm\sqrt3$。我们已经看到，这是 $\mathbb Q$ 上次数为 $4$ 的扩张，并且有如下已知子域图：
		
		\begin{center}
			\begin{tikzpicture}[scale=1.0,>=stealth]
				\node (top) at (0,3) {$\mathbb Q(\sqrt2,\sqrt3)$};
				\node (l) at (-2.4,1.5) {$\mathbb Q(\sqrt2)$};
				\node (m) at (0,1.5) {$\mathbb Q(\sqrt6)$};
				\node (r) at (2.4,1.5) {$\mathbb Q(\sqrt3)$};
				\node (bot) at (0,0) {$\mathbb Q$};
				\draw (top)--node[left]{$2$}(l);
				\draw (top)--node[right]{$2$}(m);
				\draw (top)--node[right]{$2$}(r);
				\draw (l)--node[left]{$2$}(bot);
				\draw (m)--node[right]{$2$}(bot);
				\draw (r)--node[right]{$2$}(bot);
			\end{tikzpicture}
		\end{center}
		
		\item $x^3-2$ 在 $\mathbb Q$ 上的分裂域不只是 $\mathbb Q(\sqrt[3]{2})$，因为如前面指出的，这个多项式在 $\mathbb C$ 中的三个根是
		\[
		\sqrt[3]{2},\qquad
		\sqrt[3]{2}\left(\frac{-1+i\sqrt3}{2}\right),\qquad
		\sqrt[3]{2}\left(\frac{-1-i\sqrt3}{2}\right),
		\]
		而后两个根不是 $\mathbb Q(\sqrt[3]{2})$ 的元素，因为这个域的元素形如 $a+b\sqrt[3]{2}+c\sqrt[3]{4}$，其中 $a,b,c$ 为有理数，并且所有这样的数都是实数。
		
		该多项式的分裂域 $K$ 由把这三个根全部添加到 $\mathbb Q$ 中得到。注意，由于 $K$ 含有上面前两个根，所以它含有它们的商
		\[
		\frac{-1+\sqrt{-3}}{2},
		\]
		从而 $K$ 含有元素 $\sqrt{-3}$。另一方面，任意含有 $\sqrt[3]{2}$ 和 $\sqrt{-3}$ 的域都含有上面三个根。因此
		\[
		K=\mathbb Q(\sqrt[3]{2},\sqrt{-3})
		\]
		是 $x^3-2$ 在 $\mathbb Q$ 上的分裂域。因为 $\sqrt{-3}$ 满足方程 $x^2+3=0$，这个扩张在 $\mathbb Q(\sqrt[3]{2})$ 上的次数至多为 $2$；又由于我们上面观察到 $\mathbb Q(\sqrt[3]{2})$ 不是分裂域，所以次数必为 $2$。于是
		\[
		[\mathbb Q(\sqrt[3]{2},\sqrt{-3}):\mathbb Q]=6.
		\]
		注意，我们也可以在最后略有不同地进行：指出 $\mathbb Q(\sqrt{-3})$ 是 $K$ 的一个子域，于是指数 $[\mathbb Q(\sqrt{-3}):\mathbb Q]=2$ 整除 $[K:\mathbb Q]$。由于这个扩张次数还被 $3$ 整除（因为 $\mathbb Q(\sqrt[3]{2})\subset K$），次数被 $6$ 整除，因此必为 $6$。
		
		这给出了如下已知子域图：
		\begin{center}
			\begin{tikzpicture}[scale=1.0]
				\node (top) at (0,3.4) {$\mathbb Q(\sqrt[3]{2},\sqrt{-3})$};
				\node (qneg) at (-2.6,1.6) {$\mathbb Q(\sqrt{-3})$};
				\node (q1) at (-0.5,1.9) {$\mathbb Q(\theta_1)$};
				\node (q2) at (1.2,1.65) {$\mathbb Q(\theta_2)$};
				\node (q3) at (3.0,1.45) {$\mathbb Q(\theta_3)$};
				\node (bot) at (-1.1,0) {$\mathbb Q$};
				\draw (top)--(qneg);
				\draw (top)--(q1);
				\draw (top)--(q2);
				\draw (top)--(q3);
				\draw (qneg)--(bot);
				\draw (q1)--(bot);
				\draw (q2)--(bot);
				\draw (q3)--(bot);
			\end{tikzpicture}
		\end{center}
		其中
		\[
		\theta_1=\sqrt[3]{2},\qquad
		\theta_2=\sqrt[3]{2}\left(\frac{-1+i\sqrt3}{2}\right),\qquad
		\theta_3=\sqrt[3]{2}\left(\frac{-1-i\sqrt3}{2}\right).
		\]
		
		\item 在计算分裂域时必须小心。多项式 $x^4+4$ 在 $\mathbb Q$ 上的分裂域比起初可能猜想的更小。事实上，这个多项式在 $\mathbb Q$ 上分解为
		\begin{align*}
			x^4+4 &= x^4+4x^2+4-4x^2 \\
			&=(x^2+2)^2-4x^2 \\
			&=(x^2+2x+2)(x^2-2x+2),
		\end{align*}
		其中这两个因子不可约（又由 Eisenstein 判别法）。用二次公式求解这两个因子的根，我们得到四个根
		\[
		\pm1\pm i,
		\]
		所以这个多项式的分裂域只是域 $\mathbb Q(i)$，它是 $\mathbb Q$ 的次数为 $2$ 的扩张。
	\end{enumerate}


一般地，若 $f(x)\in F[x]$ 是次数为 $n$ 的多项式，则向 $F$ 添加 $f(x)$ 的一个根会生成次数至多为 $n$ 的扩张 $F_1$（并且当且仅当 $f(x)$ 不可约时次数等于 $n$）。在 $F_1$ 上，多项式 $f(x)$ 现在至少有一个线性因子，所以 $f(x)$ 的任何另一个根都在 $F_1$ 上满足一个次数至多为 $n-1$ 的方程。把这样一个根添加到 $F_1$ 中，于是得到 $F_1$ 的一个次数至多为 $n-1$ 的扩张，等等。利用扩张次数的可乘性，这证明了

\noindent\textbf{命题 26.} 一个次数为 $n$ 的多项式在 $F$ 上的分裂域在 $F$ 上的次数至多为 $n!$。


如上面的例子所示，分裂域的次数可能小于 $n!$。以后将用 Galois 理论证明，$\mathbb Q$ 上次数为 $n$ 的“通常”多项式（在一个明确定义的意义下）的分裂域次数为 $n!$，所以这可看作“典型”情形（虽然我们将考虑的大多数有趣例子的分裂域次数更小）。

\paragraph{例（$x^n-1$ 的分裂域：分圆域）}
考虑多项式 $x^n-1$ 在 $\mathbb Q$ 上的分裂域。这个多项式的根称为 $n$ 次单位根。
	
	回忆每个非零复数 $a+bi\in\mathbb C$ 都能唯一写成
	\[
	re^{i\theta}=r(\cos\theta+i\sin\theta),\qquad r>0,
	\quad 0\leq \theta<2\pi,
	\]
	这不过是在复平面中用极坐标表示点 $a+bi$：$r$ 是 $(a,b)$ 到原点的距离，$\theta$ 是与正实轴所成的角。
	
	在 $\mathbb C$ 上，方程 $x^n=1$ 有 $n$ 个不同的解，即元素
	\[
	e^{2\pi k i/n}=\cos\left(\frac{2\pi k}{n}\right)+i\sin\left(\frac{2\pi k}{n}\right)
	\]
	其中 $k=0,1,\ldots,n-1$。在几何上，这些点是在复平面中半径为 $1$ 的圆上从点 $(1,0)$（对应 $k=0$）开始的 $n$ 个等距点（见图 6）。这些数全都是 $n$ 次单位根这一事实是立即的，因为
	\[
	\left(e^{2\pi k i/n}\right)^n=e^{(2\pi k i/n)n}=e^{2\pi k i}=1.
	\]
	于是 $\mathbb C$ 含有 $x^n-1$ 的一个分裂域，而且我们常常把 $x^n-1$ 在 $\mathbb Q$ 上的分裂域看作在 $\mathbb C$ 中由上面的数在 $\mathbb Q$ 上生成的域。
	
	\begin{center}
		\begin{tikzpicture}[scale=1.0,>=stealth]
			\draw[->] (-2.3,0)--(2.7,0);
			\draw[->] (0,-2.3)--(0,2.7);
			\draw (0,0) circle (1.7);
			\foreach \ang in {0,45,72,108,144,216,252,315}{
				\fill ({1.7*cos(\ang)},{1.7*sin(\ang)}) circle (1.5pt);
			}
			\draw[->] (0,0)--(1.7,0) node[right]{$1$};
			\draw[->] (0,0)--({1.7*cos(40)},{1.7*sin(40)}) node[right]{$e^{2\pi i/n}$};
			\draw[->] (0,0)--({1.7*cos(70)},{1.7*sin(70)});
			\draw[->] (0,0)--({1.7*cos(110)},{1.7*sin(110)});
			\draw[->] (0,0)--({1.7*cos(-40)},{1.7*sin(-40)});
			\node at (1.05,2.15) {$\mathbb C$};
			\node at (-0.95,0.95) {$\cdots$};
			\node at (0.95,-1.05) {$\cdots$};
		\end{tikzpicture}
		
		\textbf{图 6}
	\end{center}
	
	在 $x^n-1$ 的任意抽象分裂域 $K/\mathbb Q$ 中，$n$ 次单位根全体在乘法下形成一个群，因为若 $\alpha^n=1$ 且 $\beta^n=1$，则 $(\alpha\beta)^n=1$，所以 $K^\times$ 的这个子集在乘法下封闭。由此可知这是一个循环群（第 9 章命题 18）；我们将看到在 $K$ 中有 $n$ 个不同的根，所以它的阶为 $n$。
	
	\noindent\textbf{定义.} 全体 $n$ 次单位根所成循环群的一个生成元称为一个\emph{本原 $n$ 次单位根}。

	
	令 $\zeta_n$ 表示一个本原 $n$ 次单位根。于是其他本原 $n$ 次单位根是元素 $\zeta_n^a$，其中 $1\leq a<n$ 是与 $n$ 互素的整数，因为这些正是阶为 $n$ 的循环群的其他生成元。特别地，恰有 $\varphi(n)$ 个本原 $n$ 次单位根，其中 $\varphi(n)$ 表示 Euler $\varphi$-函数。
	
	在 $\mathbb C$ 上，我们可以直接看出这一切：令
	\[
	\zeta_n=e^{2\pi i/n}
	\]
	（即从 $1$ 出发逆时针方向遇到的第一个 $n$ 次单位根）。于是所有其他单位根都是 $\zeta_n$ 的幂：
	\[
	e^{2\pi k i/n}=\zeta_n^k,
	\]
	所以 $\zeta_n$ 是 $n$ 次单位根乘法群的一个可能生成元。当我们在 $\mathbb C$ 中看单位根时，通常用 $\zeta_n$ 表示这个本原 $n$ 次单位根的选择。对若干小的 $n$，$\mathbb C$ 中的本原单位根为
	\[
	\zeta_1=1,
	\]
	\[
	\zeta_2=-1,
	\]
	\[
	\zeta_3=\frac{-1+i\sqrt3}{2},
	\]
	\[
	\zeta_4=i,
	\]
	\[
	\zeta_5=\frac{\sqrt5-1}{4}+i\left(\frac{\sqrt{10+2\sqrt5}}{4}\right),
	\]
	\[
	\zeta_6=\frac{1+i\sqrt3}{2},
	\]
	\[
	\zeta_8=\frac{\sqrt2}{2}+i\frac{\sqrt2}{2}
	\]
	（这些公式来自 $n$ 边形的初等几何；无论如何，通过把它们升到相应幂次可以直接验证）。
	
	$x^n-1$ 在 $\mathbb Q$ 上的分裂域是域 $\mathbb Q(\zeta_n)$，这个域有一个名称：
	
	\noindent\textbf{定义.} 域 $\mathbb Q(\zeta_n)$ 称为 $n$ 次单位根的\emph{分圆域}。

	
	确定这个扩张的次数需要分析 $\zeta_n$ 在 $\mathbb Q$ 上的极小多项式，并将推迟到后面（第 6 节）。我们事实上已经考虑过一个重要特例：$n=p$ 为素数。在这种情形下，我们有因式分解
	\[
	x^p-1=(x-1)(x^{p-1}+x^{p-2}+\cdots+x+1),
	\]
	并且由于 $\zeta_p\ne 1$，可知 $\zeta_p$ 是多项式
	\[
	\Phi_p(x)=\frac{x^p-1}{x-1}=x^{p-1}+x^{p-2}+\cdots+x+1
	\]
	的一个根；我们在第 9.4 节证明过这个多项式不可约。于是 $\Phi_p(x)$ 是 $\zeta_p$ 在 $\mathbb Q$ 上的极小多项式，所以
	\[
	[\mathbb Q(\zeta_p):\mathbb Q]=p-1.
	\]
	我们以后将看到，一般地 $[\mathbb Q(\zeta_n):\mathbb Q]=\varphi(n)$，其中 $\varphi(n)$ 是 $n$ 的 Euler phi-函数（因此 $\varphi(p)=p-1$）。


\paragraph{例（$x^p-2$ 的分裂域，$p$ 为素数）}
令 $p$ 为素数，并考虑 $x^p-2$ 的分裂域。若 $\alpha$ 是这个方程的一个根，即 $\alpha^p=2$，则 $(\zeta\alpha)^p=2$，其中 $\zeta$ 是任意 $p$ 次单位根。因此这个方程的解为
	\[
	\zeta\sqrt[p]{2},\qquad \zeta\text{ 是一个 }p\text{ 次单位根},
	\]
	其中照例，如果我们想把这些元素看作复数，符号 $\sqrt[p]{2}$ 表示 $2$ 的正实 $p$ 次根；而如果抽象地看待这些根，则它表示 $x^p=2$ 的任意一个解。由于两个解 $\zeta_p\sqrt[p]{2}$ 与 $\sqrt[p]{2}$ 的比正是 $\zeta_p$，其中 $\zeta_p$ 为一个本原 $p$ 次单位根，$x^p-2$ 在 $\mathbb Q$ 上的分裂域含有 $\mathbb Q(\sqrt[p]{2},\zeta_p)$。另一方面，上述所有根都在这个域中，所以分裂域恰为
	\[
	\mathbb Q(\sqrt[p]{2},\zeta_p).
	\]
	
	这个域含有 $p$ 次单位根的分圆域，并由 $\sqrt[p]{2}$ 在它上面生成，所以它是在该分圆域上的次数至多为 $p$ 的扩张。于是它在 $\mathbb Q$ 上的次数 $\leq p(p-1)$。由于 $\mathbb Q(\sqrt[p]{2})$ 和 $\mathbb Q(\zeta_p)$ 都是子域，所以这个扩张在 $\mathbb Q$ 上的次数同时被 $p$ 和 $p-1$ 整除。由于这两个数互素，扩张次数被 $p(p-1)$ 整除，所以必有
	\[
	[\mathbb Q(\sqrt[p]{2},\zeta_p):\mathbb Q]=p(p-1)
	\]
	（这是推论 22）。特别注意，我们已经证明 $x^p-2$ 在 $\mathbb Q(\zeta_p)$ 上仍不可约，这一点完全不是显然的。我们有如下已知子域图：
	\begin{center}
		\begin{tikzpicture}[scale=1.0]
			\node (top) at (0,3) {$\mathbb Q(\sqrt[p]{2},\zeta_p)$};
			\node (left) at (-2.2,1.4) {$\mathbb Q(\zeta_p)$};
			\node (right) at (2.2,1.4) {$\mathbb Q(\sqrt[p]{2})$};
			\node (bot) at (0,0) {$\mathbb Q$};
			\draw (top)--node[left]{$p$}(left);
			\draw (top)--node[right]{$p-1$}(right);
			\draw (left)--node[left]{$p-1$}(bot);
			\draw (right)--node[right]{$p$}(bot);
		\end{tikzpicture}
	\end{center}
	
	特例 $p=3$ 正是上面的例 3，在那里我们只是把 $3$ 次单位根显式写了出来。


现在回到证明如下问题：多项式 $f(x)$ 在域 $F$ 上的分裂域如何构造并无差别。如定理 8 中那样，把结果表述为两个域之间的任意同构 $\varphi:F\overset{\sim}{\longrightarrow}F'$ 是方便的。

\noindent\textbf{定理 27.} 设 $\varphi:F\overset{\sim}{\longrightarrow}F'$ 是域的同构。设 $f(x)\in F[x]$ 是多项式，且 $f'(x)\in F'[x]$ 是把 $\varphi$ 作用到 $f(x)$ 的系数上所得到的多项式。设 $E$ 是 $f(x)$ 在 $F$ 上的分裂域，$E'$ 是 $f'(x)$ 在 $F'$ 上的分裂域。则同构 $\varphi$ 可扩张为一个同构 $\sigma:E\overset{\sim}{\longrightarrow}E'$，即 $\sigma$ 限制到 $F$ 上就是同构 $\varphi$：
	\[
	\begin{array}{ccc}
		\sigma:E&\longrightarrow&E'\\
		& & \\
		\varphi:F&\longrightarrow&F'.
	\end{array}
	\]


\begin{proof}
	我们对 $f(x)$ 的次数 $n$ 归纳。如定理 8 前的讨论所说，回忆从一个域 $F$ 到另一个域 $F'$ 的同构 $\varphi$ 在多项式环 $F[x]$ 与 $F'[x]$ 之间诱导一个自然同构。特别地，若 $f(x)$ 和 $f'(x)$ 在这个同构下相互对应，则 $f(x)$ 在 $F[x]$ 中的不可约因子对应于 $f'(x)$ 在 $F'[x]$ 中的不可约因子。
	
	若 $f(x)$ 的全部根都在 $F$ 中，则 $f(x)$ 在 $F[x]$ 中完全分裂，并且 $f'(x)$ 在 $F'[x]$ 中完全分裂（其线性因子就是 $f(x)$ 的线性因子的像）。因此 $E=F$ 且 $E'=F'$，在这种情形可取 $\sigma=\varphi$。这说明对 $n=1$ 以及 $f(x)$ 的全部不可约因子都次数为 $1$ 的情形，结果成立。
	
	现在假设归纳地，定理已对任意域 $F$、同构 $\varphi$ 以及次数小于 $n$ 的多项式 $f(x)\in F[x]$ 得到证明。令 $p(x)$ 为 $f(x)$ 在 $F[x]$ 中次数至少为 $2$ 的一个不可约因子，令 $p'(x)$ 为 $f'(x)$ 在 $F'[x]$ 中相应的不可约因子。若 $\alpha\in E$ 是 $p(x)$ 的一个根，且 $\beta\in E'$ 是 $p'(x)$ 的一个根，则由定理 8，我们可以把 $\varphi$ 扩张为一个同构 $\sigma':F(\alpha)\overset{\sim}{\longrightarrow}F'(\beta)$：
	\[
	\begin{array}{ccc}
		\sigma':F(\alpha)&\longrightarrow&F'(\beta)\\
		& & \\
		\varphi:F&\longrightarrow&F'.
	\end{array}
	\]
	令 $F_1=F(\alpha)$，$F'_1=F'(\beta)$，于是有同构 $\sigma':F_1\overset{\sim}{\longrightarrow}F'_1$。在 $F_1$ 上我们有 $f(x)=(x-\alpha)f_1(x)$，其中 $f_1(x)$ 次数为 $n-1$，并且 $f'(x)=(x-\beta)f'_1(x)$。域 $E$ 是 $f_1(x)$ 在 $F_1$ 上的一个分裂域：$f_1(x)$ 的所有根都在 $E$ 中，并且若它们包含在任何更小的含有 $F_1$ 的扩张 $L$ 中，那么由于 $F_1$ 含有 $\alpha$，$L$ 也将含有 $f(x)$ 的全部根，这将与 $E$ 作为 $f(x)$ 在 $F$ 上的分裂域的极小性矛盾。类似地，$E'$ 是 $f'_1(x)$ 在 $F'_1$ 上的一个分裂域。由于 $f_1(x)$ 与 $f'_1(x)$ 的次数都小于 $n$，由归纳假设，存在一个映射 $\sigma:E\overset{\sim}{\longrightarrow}E'$，它扩张同构 $\sigma':F_1\overset{\sim}{\longrightarrow}F'_1$。这给出扩张后的图：
	\[
	\begin{array}{ccccc}
		\sigma:E&\longrightarrow&E'\\
		& & \\
		\sigma':F_1&\longrightarrow&F'_1\\
		& & \\
		\varphi:F&\longrightarrow&F'.
	\end{array}
	\]
	于是，如图所示，$\sigma$ 限制到 $F_1$ 上就是同构 $\sigma'$，所以特别地，$\sigma$ 限制到 $F$ 上就是 $\sigma'$ 限制到 $F$ 上，也就是 $\varphi$，这说明 $\sigma$ 是 $\varphi$ 的一个扩张，证明完成。
\end{proof}

\noindent\textbf{推论 28（分裂域的唯一性）.} 域 $F$ 上一个多项式 $f(x)\in F[x]$ 的任意两个分裂域都是同构的。


\begin{proof}
	取 $\varphi$ 为从 $F$ 到自身的恒等映射，并取 $E$ 和 $E'$ 为 $f(x)(=f'(x))$ 的两个分裂域。
\end{proof}

如前所述，这个结果证明了“$f(x)$ 在 $F$ 上的分裂域”这一术语的合理性，因为任意两个这样的域都是同构的。分裂域在代数元素的研究中扮演自然的角色（如果你要添加一个多项式的一个根，为什么不添加全部根呢？），因此在 Galois 理论中占有特别重要的地位。

我们以讨论 $F$ 的那些含有 $F$ 上所有多项式的所有根的域扩张来结束本节。

\noindent\textbf{定义.} 若 $\overline F$ 是 $F$ 的代数扩张，并且每个 $f(x)\in F[x]$ 都在 $\overline F$ 上完全分裂（因而可说 $\overline F$ 含有 $F$ 上全部代数元素），则称域 $\overline F$ 是 $F$ 的一个\emph{代数闭包}。


\noindent\textbf{定义.} 若每个以 $K$ 中元素为系数的多项式都在 $K$ 中有一个根，则称域 $K$ 是\emph{代数闭的}。


代数闭域的存在并不显然，给定域 $F$ 的代数闭包存在也不显然（我们很快证明这一点）。

注意，若 $K$ 是代数闭的，则事实上每个 $f(x)\in K[x]$ 的全部根都在 $K$ 中，因为按定义 $f(x)$ 在 $K$ 中有一个根 $\alpha$，因而在 $K[x]$ 中有因子 $x-\alpha$。$f(x)$ 的余下因子于是仍是 $K[x]$ 中的一个多项式，所以也有一个根，因而又有一个线性因子，依此类推，所以 $f(x)$ 必完全分裂。因此，若 $K$ 是代数闭的，则 $K$ 本身就是 $K$ 的一个代数闭包；反过来则显然，所以 $K=\overline K$ 当且仅当 $K$ 是代数闭的。

下一个结果说明，“取代数闭包”的过程事实上一步就停止——取一个代数闭包的代数闭包不会给出更大的域：这个域已经是代数闭的（记号上：$\overline{\overline F}=\overline F$）。

\noindent\textbf{命题 29.} 设 $\overline F$ 是 $F$ 的一个代数闭包。则 $\overline F$ 是代数闭的。


\begin{proof}
	设 $f(x)$ 是 $\overline F[x]$ 中的一个多项式，并令 $\alpha$ 是 $f(x)$ 的一个根。于是 $\alpha$ 生成 $\overline F$ 的一个代数扩张 $\overline F(\alpha)$，而 $\overline F$ 在 $F$ 上代数。由定理 20，$\overline F(\alpha)$ 在 $F$ 上代数，所以特别地，其元素 $\alpha$ 在 $F$ 上代数。但于是 $\alpha\in\overline F$，说明 $\overline F$ 是代数闭的。
\end{proof}

给定一个域 $F$，我们已经证明了如何构造（有限）扩张，使其含有任意给定多项式 $f(x)\in F[x]$ 的全部根。直观地，$F$ 的一个代数闭包由所有这些域“生成”。这里的困难在于“生成”在哪里进行，因为它们并不全是某一个给定域的子域。对于有限个多项式 $f_1(x),\ldots,f_k(x)$，我们可以把它们的分裂域看作乘积多项式 $f_1(x)\cdots f_k(x)$ 的分裂域的子域；但同样的想法用于无限多个多项式时，需要大量“记账”式的等同以及 Zorn 引理的一次应用。

我们将换一种方式构造 $F$ 的一个代数闭包：先构造一个含有 $F$ 的代数闭域。证明使用 Artin 的一个巧妙想法：通过为每个多项式引入一个单独变量，非常简洁地解决了构造一个含有适当多项式根的域时的“记账”问题（这最终也依赖于 Zorn 引理）。

\noindent\textbf{命题 30.} 对任意域 $F$，存在一个含有 $F$ 的代数闭域 $K$。


\begin{proof}
	对每个以 $F$ 中元素为系数的非常数首一多项式 $f=f(x)$，令 $x_f$ 表示一个不定元，并考虑由这些变量 $x_f$ 在 $F$ 上生成的多项式环 $F[\ldots,x_f,\ldots]$。在这个多项式环中，考虑由多项式 $f(x_f)$ 生成的理想 $I$。若这个理想不是真理想，则 $1$ 是这个理想的一个元素，因而有一个关系
	\[
	g_1 f_1(x_{f_1})+g_2 f_2(x_{f_2})+\cdots+g_n f_n(x_{f_n})=1,
	\]
	其中 $g_i$，$i=1,2,\ldots,n$，是关于 $x_f$ 的多项式。对 $i=1,2,\ldots,n$，令 $x_{f_i}=x_i$，并令 $x_{n+1},\ldots,x_m$ 为多项式 $g_j$，$j=1,2,\ldots,n$ 中出现的其余变量。于是上面的关系写为
	\[
	g_1(x_1,x_2,\ldots,x_m)f_1(x_1)+\cdots+g_n(x_1,x_2,\ldots,x_m)f_n(x_n)=1.
	\]
	令 $F'$ 为 $F$ 的一个有限扩张，它含有 $f_i(x)$ 的一个根 $\alpha_i$，$i=1,2,\ldots,n$。在上面的多项式方程中令 $x_i=\alpha_i$，$i=1,2,\ldots,n$，并令 $x_{n+1}=\cdots=x_m=0$，比如说，就会推出 $F'$ 中 $0=1$，这显然不可能。
	
	既然理想 $I$ 是真理想，它包含于某个极大理想 $M$ 中（这正是使用 Zorn 引理之处）。于是商
	\[
	K_1=F[\ldots,x_f,\ldots]/M
	\]
	是一个含有 $F$（的一个同构副本）的域。由构造，每个多项式 $f$ 在 $K_1$ 中都有一个根，即 $x_f$ 的像，因为 $f(x_f)\in I\subseteq M$。我们已经构造了一个域 $K_1$，在其中每个系数来自 $F$ 的多项式都有一个根。
	
	用 $K_1$ 代替 $F$ 进行同样构造，得到一个含有 $K_1$ 的域 $K_2$，在其中所有系数来自 $K_1$ 的多项式都有一个根。如此继续，我们得到一列域
	\[
	F=K_0\subseteq K_1\subseteq K_2\subseteq\cdots\subseteq K_j\subseteq K_{j+1}\subseteq\cdots
	\]
	其中每个系数在 $K_j$ 中的多项式在 $K_{j+1}$ 中都有一个根，$j=0,1,\ldots$。令
	\[
	K=\bigcup_{j\geq 0}K_j
	\]
	为这些域的并。则 $K$ 显然是一个含有 $F$ 的域。由于 $K$ 是这些域 $K_j$ 的并，$K[x]$ 中任意多项式 $h(x)$ 的全部系数都属于某个 $K_N$，只要 $N$ 充分大。但于是 $h(x)$ 在 $K_{N+1}$ 中有一个根，所以在 $K$ 中有一个根。因此 $K$ 是代数闭的，证明完成。
\end{proof}

现在我们用含有 $F$ 的代数闭域来构造 $F$ 的一个代数闭包：

\noindent\textbf{命题 31.} 设 $K$ 是一个代数闭域，且 $F$ 是 $K$ 的一个子域。则 $K$ 中在 $F$ 上代数的元素全体 $\overline F$ 是 $F$ 的一个代数闭包。$F$ 的代数闭包在同构意义下唯一。


\begin{proof}
	按定义，$\overline F$ 是 $F$ 的一个代数扩张。每个多项式 $f(x)\in F[x]$ 在 $K$ 上完全分解为线性因子 $x-\alpha$（甚至 $K[x]$ 中的每个多项式也是如此）。但每个 $\alpha$ 都是 $f(x)$ 的一个根，所以在 $F$ 上代数，因而是 $\overline F$ 的一个元素。于是所有线性因子 $x-\alpha$ 的系数都在 $\overline F$ 中，即 $f(x)$ 在 $\overline F[x]$ 中完全分裂，$\overline F$ 是 $F$ 的一个代数闭包。
	
	代数闭包的唯一性（在同构意义下）鉴于分裂域的唯一性是自然的；它沿着同样的思路并配合 Zorn 引理的一次应用来证明，这里省略。
\end{proof}

我们以后将用 Galois 理论证明下面的结果（也存在使用复分析的纯解析证明）。

\noindent\textbf{定理（代数学基本定理）.} 域 $\mathbb C$ 是代数闭的。


由命题 31，我们立即得到：

\noindent\textbf{推论 32.} 域 $\mathbb C$ 含有它任一子域的代数闭包。特别地，$\overline{\mathbb Q}$，即 $\mathbb C$ 中在 $\mathbb Q$ 上代数的复数全体，是 $\mathbb Q$ 的一个代数闭包。


这些考虑的要点是：涉及域 $F$ 上代数元素的所有计算，都可以看作发生在一个（大的）域中，即 $\overline F$。类似地，我们可以有意义地谈论任意一族代数扩张的复合，只要把它们全都看作代数闭包的子域。在 $\mathbb Q$ 或 $\mathbb Q$ 的有限扩张的情形下，我们可以把所有计算都看作发生在 $\mathbb C$ 中。

\subsection*{练习}

\begin{enumerate}[label=\arabic*.]
	\item 求 $x^4-2$ 在 $\mathbb Q$ 上的分裂域及其在 $\mathbb Q$ 上的次数。
	\item 求 $x^4+2$ 在 $\mathbb Q$ 上的分裂域及其在 $\mathbb Q$ 上的次数。
	\item 求 $x^4+x^2+1$ 在 $\mathbb Q$ 上的分裂域及其在 $\mathbb Q$ 上的次数。
	\item 求 $x^6-4$ 在 $\mathbb Q$ 上的分裂域及其在 $\mathbb Q$ 上的次数。
	\item 设 $K$ 是 $F$ 的一个有限扩张。证明：$K$ 是 $F$ 上的一个分裂域，当且仅当 $F[x]$ 中每个在 $K$ 中有根的不可约多项式都在 $K[x]$ 中完全分裂。[使用定理 8 和定理 27。]
	\item 设 $K_1$ 和 $K_2$ 是包含在域 $K$ 中的 $F$ 的有限扩张，并假设二者都是 $F$ 上的分裂域。
	\begin{enumerate}[label=(\alph*)]
		\item 证明它们的复合 $K_1K_2$ 是 $F$ 上的一个分裂域。
		\item 证明 $K_1\cap K_2$ 是 $F$ 上的一个分裂域。[使用前一练习。]
	\end{enumerate}
\end{enumerate}
\section{可分与不可分扩张}

设 $F$ 是一个域，设 $f(x)\in F[x]$ 是一个多项式。在 $f(x)$ 的一个分裂域上，我们有分解
\[
f(x)=(x-\alpha_1)^{n_1}(x-\alpha_2)^{n_2}\cdots(x-\alpha_k)^{n_k},
\]
其中 $\alpha_1,\alpha_2,\ldots,\alpha_k$ 是分裂域中的互异元素，并且对所有 $i$ 有 $n_i\geq 1$。回忆：如果 $n_i>1$，则称 $\alpha_i$ 为一个\emph{重根}；如果 $n_i=1$，则称它为一个\emph{单根}。整数 $n_i$ 称为根 $\alpha_i$ 的\emph{重数}。

\noindent\textbf{定义.} 域 $F$ 上的一个多项式称为\emph{可分的}，如果它没有重根（即，它的所有根都互异）。一个不可分的多项式称为\emph{不可分的}。

注意，如果一个多项式 $f(x)$ 在一个分裂域中有互异的根，则 $f(x)$ 在任意分裂域中都有互异的根（因为这等价于 $f(x)$ 分解成互异的线性因子，而 $f(x)$ 的任意两个分裂域之间存在一个在 $F$ 上的同构，并且它在根的集合上是双射），所以我们不必指明含有 $f(x)$ 全部根的域。

\paragraph{例}
\begin{enumerate}[label=(\arabic*)]
	\item 多项式 $x^2-2$ 在 $\mathbb Q$ 上是可分的，因为它的两个根 $\pm\sqrt2$ 互异。对任意 $n\geq 2$，多项式 $(x^2-2)^n$ 是不可分的，因为它有重根 $\pm\sqrt2$，每个重数都是 $n$。
	
	\item 多项式 $x^2-t(=x^2+t)$ 在域 $F=\mathbb F_2(t)$（系数来自 $\mathbb F_2$ 的 $t$ 的有理函数域）上如我们以前所见是不可约的，但不是可分的。若 $\sqrt t$ 表示某个扩域中的一个根（注意 $\sqrt t\notin F$），则
	\[
	(x-\sqrt t)^2=x^2-2x\sqrt t+t=x^2+t=x^2-t,
	\]
	因为 $F$ 是特征 $2$ 的域。因此这个不可约多项式只有一个根（重数为 $2$），所以它在 $F$ 上不可分。
\end{enumerate}

有一个简单判别准则可以检查一个多项式是否有重根。

\noindent\textbf{定义.} 多项式
\[
f(x)=a_nx^n+a_{n-1}x^{n-1}+\cdots+a_1x+a_0\in F[x]
\]
的\emph{导数}定义为多项式
\[
D_xf(x)=na_nx^{n-1}+(n-1)a_{n-1}x^{n-2}+\cdots+2a_2x+a_1\in F[x].
\]

这个公式无非是微积分中熟悉的多项式导数通常公式。它是纯代数的，因此可以应用于任意域 $F$ 上的多项式，而在这样的域上，导数的解析概念（涉及极限——一种连续运算）可能并不存在。

通常的（微积分）导数公式在这种情形下同样成立，例如和与乘积的导数公式：
\[
D_x(f(x)+g(x))=D_xf(x)+D_xg(x)
\]
以及
\[
D_x(f(x)g(x))=f(x)D_xg(x)+(D_xf(x))g(x).
\]
这些公式可以直接从多项式的定义证明，不需要任何极限运算，并且留作练习。

下一个命题说明，$f(x)$ 的可分性可以在 $f(x)$ 的系数所在的域中用 Euclidean 算法来判定，而不必通过一个分裂域并分解 $f(x)$。

\noindent\textbf{命题 33.} 一个多项式 $f(x)$ 有重根 $\alpha$，当且仅当 $\alpha$ 也是 $D_xf(x)$ 的根，也就是说，$f(x)$ 和 $D_xf(x)$ 都被 $\alpha$ 的极小多项式整除。特别地，$f(x)$ 可分，当且仅当它与它的导数互素：
\[
(f(x),D_xf(x))=1.
\]

\begin{proof}
	首先假设 $\alpha$ 是 $f(x)$ 的一个重根。则在一个分裂域上，
	\[
	f(x)=(x-\alpha)^ng(x)
	\]
	其中 $n\geq 2$，$g(x)$ 是某个多项式。取导数得到
	\[
	D_xf(x)=n(x-\alpha)^{n-1}g(x)+(x-\alpha)^nD_xg(x),
	\]
	这说明（由于 $n\geq2$）$D_xf(x)$ 以 $\alpha$ 为根。
	
	反过来，假设 $\alpha$ 同时是 $f(x)$ 和 $D_xf(x)$ 的根。于是可以写
	\[
	f(x)=(x-\alpha)h(x)
	\]
	其中 $h(x)$ 是某个多项式；取导数得到
	\[
	D_xf(x)=h(x)+(x-\alpha)D_xh(x).
	\]
	由于按假设 $D_xf(\alpha)=0$，把 $\alpha$ 代入最后这个等式，说明 $h(\alpha)=0$。因此 $h(x)=(x-\alpha)h_1(x)$，其中 $h_1(x)$ 是某个多项式，并且
	\[
	f(x)=(x-\alpha)^2h_1(x),
	\]
	说明 $\alpha$ 是 $f(x)$ 的一个重根。
	
	与被 $\alpha$ 的极小多项式整除的等价性来自命题 9。最后一个断言于是显然成立（令 $\alpha$ 表示 $f(x)$ 和 $D_xf(x)$ 的一个公共因子的任意一个根）。
\end{proof}

\paragraph{例}
\begin{enumerate}[label=(\arabic*)]
	\item $\mathbb F_p$ 上的多项式 $x^p-x$ 的导数是 $px^{p-1}-1=-1$，因为该域的特征是 $p$。在这种情形下，导数根本没有根，因此这个多项式没有重根，所以是可分的。
	
	\item 多项式 $x^n-1$ 的导数是 $nx^{n-1}$。在任意特征不整除 $n$ 的域上（包括特征 $0$），这个多项式只有根 $0$（重数为 $n-1$），而 $0$ 不是 $x^n-1$ 的根。因此 $x^n-1$ 是可分的，并且有 $n$ 个互异的 $n$ 次单位根。我们已经通过在 $\mathbb C$ 上显式给出 $n$ 个互异解，直接在 $\mathbb Q$ 上看到了这一点。
	
	\item 如果 $F$ 的特征为 $p$ 且 $p$ 整除 $n$，那么在 $F$ 上互异的 $n$ 次单位根少于 $n$ 个：在这种情形下，由于在 $F$ 中有 $n=0$，导数恒等于 $0$。事实上，在这种情形下，$x^n-1$ 的每个根都是重根。
\end{enumerate}

\noindent\textbf{推论 34.} 特征 $0$ 的域（例如 $\mathbb Q$）上的每个不可约多项式都是可分的。这样一个域上的多项式可分，当且仅当它是互异不可约多项式的乘积。

\begin{proof}
	设 $F$ 是特征 $0$ 的域，设 $p(x)\in F[x]$ 是次数为 $n$ 的不可约多项式。则导数 $D_xp(x)$ 是次数为 $n-1$ 的多项式。在 $F[x]$ 中，不计常数倍，$p(x)$ 的因子只有 $1$ 和 $p(x)$，所以 $D_xp(x)$ 必定与 $p(x)$ 互素。这说明特征 $0$ 的域上的任意不可约多项式都是可分的。推论的第二个断言于是显然成立，因为互异不可约多项式从不共有零点（由命题 9）。
\end{proof}

推论证明中在特征 $p$ 中可能失效的点，是导数 $D_xp(x)$ 的次数为 $n-1$ 这一说法。在特征 $p$ 中，$x^p$ 的任意幂 $x^{pm}$ 的导数恒等于 $0$：
\[
D_x(x^{pm})=pmx^{pm-1}=0,
\]
所以导数次数可能降低不止一阶。不过，如果不可约多项式 $p(x)$ 的导数 $D_xp(x)$ 非零，那么仍然像以前一样可以得出 $p(x)$ 必须可分。

由导数的定义显然可知，如果 $p(x)$ 是一个导数为 $0$ 的多项式，那么 $p(x)$ 中 $x$ 的每个指数都必须是 $p$ 的倍数，其中 $p$ 是 $F$ 的特征：
\[
p(x)=a_mx^{mp}+a_{m-1}x^{(m-1)p}+\cdots+a_1x^p+a_0.
\]
令
\[
p_1(x)=a_mx^m+a_{m-1}x^{m-1}+\cdots+a_1x+a_0,
\]
我们看到 $p(x)$ 是 $x^p$ 的一个多项式，即 $p(x)=p_1(x^p)$。

现在我们证明一个关于在特征 $p$ 的域中取 $p$ 次幂的简单但重要的结果。

\noindent\textbf{命题 35.} 设 $F$ 是特征 $p$ 的域。则对任意 $a,b\in F$，
\[
(a+b)^p=a^p+b^p,
\qquad
(ab)^p=a^pb^p.
\]
换言之，由 $\varphi(a)=a^p$ 定义的 $p$ 次幂映射，是从 $F$ 到 $F$ 的一个单射域同态。

\begin{proof}
	把 $(a+b)^n$ 展开成二项式的二项式定理对任意正整数 $n$ 在任意交换环上都成立（由标准归纳证明）：
	\[
	(a+b)^n=a^n+\binom n1a^{n-1}b+\cdots+\binom ni a^{n-i}b^i+\cdots+b^n.
	\]
	应当注意，二项式系数
	\[
	\binom ni=\frac{n!}{i!(n-i)!}
	\]
	是整数（回忆，对任意环中的元素 $a$，$m\in\mathbb Z$ 时 $ma$ 是有定义的），在这里它们是素子域中的元素。
	
	若 $p$ 是素数，则当 $i=1,2,\ldots,p-1$ 时，二项式系数 $\binom pi$ 全部被 $p$ 整除，因为对这些 $i$，数 $i!$ 和 $(p-i)!$ 只包含小于 $p$ 的因子，因此都与 $p$ 互素，所以不能约去表达式
	\[
	\frac{p!}{i!(p-i)!}
	\]
	分子中的因子 $p$。于是，在特征 $p$ 的域上，$(a+b)^p$ 展开式中的所有中间项都是 $0$，这给出命题的第一个等式。第二个等式是显然的，$\varphi$ 是单射也显然。
\end{proof}

\noindent\textbf{定义.} 命题 35 中的映射称为 $F$ 的 \emph{Frobenius 自同态}。

\noindent\textbf{推论 36.} 设 $\mathbb F$ 是特征 $p$ 的有限域。则 $\mathbb F$ 中每个元素都是 $\mathbb F$ 中的 $p$ 次幂（记号为 $\mathbb F=\mathbb F^p$）。

\begin{proof}
	$\mathbb F$ 的 Frobenius 自同态的单射性蕴含它在 $\mathbb F$ 有限时也是满射，这正是推论的断言。
\end{proof}

现在我们证明推论 34 对有限域的类似结论。

设 $\mathbb F$ 是一个有限域，并设 $p(x)\in\mathbb F[x]$ 是一个系数在 $\mathbb F$ 中的不可约多项式。如果 $p(x)$ 不可分，那么我们已经看到 $p(x)=q(x^p)$，其中 $q(x)\in\mathbb F[x]$。令
\[
q(x)=a_mx^m+a_{m-1}x^{m-1}+\cdots+a_1x+a_0.
\]
由推论 36，每个 $a_i$，$i=1,2,\ldots,m$，都是 $\mathbb F$ 中的一个 $p$ 次幂，比如 $a_i=b_i^p$。于是由命题 35，
\begin{align*}
	p(x)&=q(x^p)=a_m(x^p)^m+a_{m-1}(x^p)^{m-1}+\cdots+a_1x^p+a_0 \\
	&=b_m^p(x^p)^m+b_{m-1}^p(x^p)^{m-1}+\cdots+b_1^px^p+b_0^p \\
	&=(b_mx^m)^p+(b_{m-1}x^{m-1})^p+\cdots+(b_1x)^p+(b_0)^p \\
	&=(b_mx^m+b_{m-1}x^{m-1}+\cdots+b_1x+b_0)^p,
\end{align*}
这说明 $p(x)$ 是 $\mathbb F[x]$ 中某个多项式的 $p$ 次幂，这与 $p(x)$ 不可约矛盾。这证明了：

\noindent\textbf{命题 37.} 有限域 $\mathbb F$ 上的每个不可约多项式都是可分的。$\mathbb F[x]$ 中一个多项式可分，当且仅当它是 $\mathbb F[x]$ 中互异不可约多项式的乘积。

这个结果证明的重要部分，是特征 $p$ 的域 $\mathbb F$ 中每个元素都是 $\mathbb F$ 中的一个 $p$ 次幂。这提示了下面的定义：

\noindent\textbf{定义.} 特征 $p$ 的域 $K$ 称为\emph{完美的}，如果 $K$ 的每个元素都是 $K$ 中的一个 $p$ 次幂，即 $K=K^p$。任意特征 $0$ 的域也称为完美的。

用这个定义，我们看到我们已经证明了，完美域上的每个不可约多项式都是可分的。不难看出，如果 $K$ 不是完美的，那么存在不可分的不可约多项式。

\paragraph{例：有限域的存在性与唯一性}
设 $n>0$ 是任意正整数，并考虑 $\mathbb F_p$ 上多项式 $x^{p^n}-x$ 的分裂域。我们已经看到这个多项式是可分的，因此恰有 $p^n$ 个根。设 $\alpha$ 和 $\beta$ 是这个多项式的任意两个根，即 $\alpha^{p^n}=\alpha$ 且 $\beta^{p^n}=\beta$。于是
\[
(\alpha\beta)^{p^n}=\alpha\beta,
\qquad
(\alpha^{-1})^{p^n}=\alpha^{-1},
\]
并且由命题 35 还得到
\[
(\alpha+\beta)^{p^n}=\alpha^{p^n}+\beta^{p^n}=\alpha+\beta.
\]
因此，由 $x^{p^n}-x$ 在 $\mathbb F_p$ 上的 $p^n$ 个互异根组成的集合 $\mathbb F$ 在其分裂域中对加法、乘法和取逆封闭。于是 $\mathbb F$ 是一个子域，因而事实上必定就是分裂域。由于元素个数是 $p^n$，我们有
\[
[\mathbb F:\mathbb F_p]=n,
\]
这说明对任意 $n>0$，都存在 $\mathbb F_p$ 上次数为 $n$ 的有限域。

现在设 $\mathbb F$ 是任意一个特征为 $p$ 的有限域。若 $\mathbb F$ 在它的素子域 $\mathbb F_p$ 上的维数为 $n$，则 $\mathbb F$ 恰有 $p^n$ 个元素。由于乘法群 $\mathbb F^\times$（事实上是循环的）阶为 $p^n-1$，对每个 $\alpha\neq0$ 的 $\alpha\in\mathbb F$ 有 $\alpha^{p^n-1}=1$，所以对每个 $\alpha\in\mathbb F$ 有 $\alpha^{p^n}=\alpha$。但这意味着 $\alpha$ 是 $x^{p^n}-x$ 的一个根，因此 $\mathbb F$ 包含在这个多项式的一个分裂域中。由于我们已经看到这个分裂域有 $p^n$ 阶，这说明 $\mathbb F$ 是 $x^{p^n}-x$ 的一个分裂域。由于分裂域在同构意义下唯一，这证明了任意阶 $p^n$ 的有限域都存在，并且在同构意义下唯一。我们把 $p^n$ 阶有限域记为 $\mathbb F_{p^n}$。

我们以后还会进一步考虑有限域。

现在我们进一步研究特征 $p$ 的域上不可分不可约多项式的结构。我们在上面已经看到，如果 $p(x)$ 是一个不可分的不可约多项式，则它的导数 $D_xp(x)$ 恒等于 $0$，所以 $p(x)=p_1(x^p)$，其中 $p_1(x)$ 是某个多项式。多项式 $p_1(x)$ 本身可能可分，也可能不可分。如果不可分，则它也同样是 $x^p$ 的一个多项式，$p_1(x)=p_2(x^p)$，于是 $p(x)$ 是 $x^{p^2}$ 的一个多项式：$p(x)=p_2(x^{p^2})$。这样继续下去，我们看到存在唯一确定的 $p$ 的幂 $p^k$，使得 $p(x)=p_k(x^{p^k})$，其中 $p_k(x)$ 的导数非零。显然 $p_k(x)$ 不可约，因为 $p_k(x)$ 的任意分解，在把 $x$ 替换为 $x^{p^k}$ 后，都会立即给出不可约多项式 $p(x)$ 的一个分解。于是 $p_k(x)$ 是可分的。我们把这总结如下：

\noindent\textbf{命题 38.} 设 $p(x)$ 是特征 $p$ 的域 $F$ 上的一个不可约多项式。则存在唯一整数 $k\geq0$ 和唯一不可约可分多项式 $p_{\mathrm{sep}}(x)\in F[x]$，使得
\[
p(x)=p_{\mathrm{sep}}(x^{p^k}).
\]

\noindent\textbf{定义.} 设 $p(x)$ 是特征 $p$ 的域上的一个不可约多项式。上一个命题中 $p_{\mathrm{sep}}(x)$ 的次数称为 $p(x)$ 的\emph{可分次数}，记为 $\deg_s p(x)$。命题中的整数 $p^k$ 称为 $p(x)$ 的\emph{不可分次数}，记为 $\deg_i p(x)$。

从定义和命题可知，$p(x)$ 可分，当且仅当它的不可分次数为 $1$，当且仅当它的次数等于它的可分次数。此外，在关系式 $p(x)=p_{\mathrm{sep}}(x^{p^k})$ 中计算次数，我们得到
\[
\deg p(x)=\deg_s p(x)\deg_i p(x).
\]

\paragraph{例}
\begin{enumerate}[label=(\arabic*)]
	\item 上面考虑过的 $F=\mathbb F_2(t)$ 上的多项式 $p(x)=x^2-t$ 的导数为 $0$，所以不是可分的（正如我们先前所判定的）。这里 $p_{\mathrm{sep}}(x)=x-t$，不可分次数为 $2$。
	
	\item $F=\mathbb F_2(t)$ 上的多项式 $p(x)=x^{2^m}-t$ 是不可约的，其可分多项式部分相同，但不可分次数为 $2^m$。
	
	\item $F=\mathbb F_p(t)$ 上的多项式 $(x^{p^2}-t)(x^p-t)$ 有（两个）不可分不可约因子，所以是不可分的。这个多项式不能写成 $f_{\mathrm{sep}}(x^{p^k})$ 的形式，其中 $f_{\mathrm{sep}}(x)$ 可分，这正是我们上面限定在不可约多项式的原因。这个例子还说明，对于一般多项式，没有类似的分解来定义可分次数和不可分次数。
\end{enumerate}

可分性的概念推广到由这些多项式的根所生成的域。

\noindent\textbf{定义.} 若域 $K$ 的每个元素都是 $F$ 上某个可分多项式的根（等价地，$K$ 的每个元素在 $F$ 上的极小多项式都是可分的），则称 $K$ 在 $F$ 上是\emph{可分的}（或\emph{可分代数的}）。一个不可分的域称为\emph{不可分的}。

我们已经看到，对完美域的有限扩张而言，可分性问题很简单，因为对这些域，代数元素的极小多项式不可约，从而可分。

\noindent\textbf{推论 39.} 完美域的每个有限扩张都是可分的。特别地，$\mathbb Q$ 或有限域的每个有限扩张都是可分的。

我们将在发展一些 Galois 理论之后进一步考虑可分和不可分扩张，特别是定义扩张 $K/F$ 的可分次数和不可分次数。

\subsection*{练习}

\begin{enumerate}[label=\arabic*.]
	\item 证明多项式的导数 $D_x$ 对任意两个多项式 $f(x)$ 和 $g(x)$ 满足
	\[
	D_x(f(x)+g(x))=D_x(f(x))+D_x(g(x))
	\]
	以及
	\[
	D_x(f(x)g(x))=D_x(f(x))g(x)+D_x(g(x))f(x).
	\]
	
	\item 找出 $\mathbb F_2$ 上所有次数为 $1$、$2$ 和 $4$ 的不可约多项式，并证明它们的乘积是 $x^{16}-x$。
	
	\item 证明 $d$ 整除 $n$，当且仅当 $x^d-1$ 整除 $x^n-1$。[注意：如果 $n=qd+r$，则
	\[
	x^n-1=(x^{qd+r}-x^r)+(x^r-1).
	\]]
	
	\item 设 $a>1$ 是整数。证明对任意正整数 $n,d$，$d$ 整除 $n$，当且仅当 $a^d-1$ 整除 $a^n-1$（参见前一练习）。特别地，推出
	\[
	\mathbb F_{p^d}\subseteq\mathbb F_{p^n}
	\]
	当且仅当 $d$ 整除 $n$。
	
	\item 对任意素数 $p$ 和任意非零 $a\in\mathbb F_p$，证明 $x^p-x+a$ 在 $\mathbb F_p$ 上不可约且可分。[关于不可约性：一种做法是先证明若 $\alpha$ 是一个根，则 $\alpha+1$ 也是一个根。另一种做法是假设它可约并计算导数。]
	
	\item 证明
	\[
	x^{p^n-1}-1=\prod_{a\in\mathbb F_{p^n}^{\times}}(x-a).
	\]
	推出
	\[
	\prod_{a\in\mathbb F_{p^n}^{\times}}a=(-1)^{p^n},
	\]
	所以有限域中所有非零元素的乘积在 $p=2$ 时为 $+1$，在 $p$ 为奇数时为 $-1$。当 $p$ 为奇素数且 $n=1$ 时，推出 Wilson 定理：
	\[
	(p-1)!\equiv -1\pmod p.
	\]
	
	\item 设 $K$ 是一个特征为 $p$ 且不是完美域的域：$K\neq K^p$。证明存在 $K$ 上不可约的不可分多项式。推出存在 $K$ 的不可分有限扩张。
	
	\item 证明对任意多项式 $f(x)\in\mathbb F_p[x]$，有 $f(x)^p=f(x^p)$。
	
	\item 说明二项式系数 $\binom{pn}{pi}$ 是 $(1+x)^{pn}$ 的展开式中 $x^{pi}$ 的系数。在 $\mathbb F_p$ 上工作，说明这也是 $(1+x^p)^n$ 中 $(x^p)^i$ 的系数，从而证明
	\[
	\binom{pn}{pi}\equiv \binom ni \pmod p.
	\]
	
	\item 设
	\[
	f(x_1,x_2,\ldots,x_n)\in\mathbb Z[x_1,x_2,\ldots,x_n]
	\]
	是变量 $x_1,x_2,\ldots,x_n$ 中的整数系数多项式。对任意素数 $p$，证明多项式
	\[
	f(x_1,x_2,\ldots,x_n)^p-f(x_1^p,x_2^p,\ldots,x_n^p)
	\in\mathbb Z[x_1,x_2,\ldots,x_n]
	\]
	的所有系数都可被 $p$ 整除。
	
	\item 设 $K[x]$ 是域 $K$ 上的一个多项式环，且 $F$ 是 $K$ 的一个子域。如果 $F$ 是完美域，并且 $f(x)\in F[x]$ 在 $F[x]$ 中没有重复的不可约因子，证明 $f(x)$ 在 $K[x]$ 中也没有重复的不可约因子。
\end{enumerate}
\section{分圆多项式与分圆扩张}

本节的目的是证明第 4 节中引入的、由 $\mathbb Q$ 上 $n$ 次单位根生成的分圆扩张
\[
\mathbb Q(\zeta_n)/\mathbb Q
\]
的次数是 $\varphi(n)$，其中 $\varphi$ 表示 Euler 的 $\varphi$-函数（即与 $n$ 互素的整数 $a$，$1\leq a<n$，的个数，也即群 $(\mathbb Z/n\mathbb Z)^\times$ 的阶）。

\noindent\textbf{定义.} 令 $\mu_n$ 表示 $\mathbb Q$ 上 $n$ 次单位根的群。

那么，如我们已经观察到的，作为群（右边用乘法，左边用加法），有
\[
\mathbb Z/n\mathbb Z\cong \mu_n,
\]
对一个固定的本原 $n$ 次单位根，显式地由映射
\[
a\longmapsto (\zeta_n)^a
\]
给出。本原 $n$ 次单位根由与 $n$ 互素的剩余类给出，所以恰有 $\varphi(n)$ 个本原 $n$ 次单位根。

若 $d$ 是 $n$ 的一个因子且 $\zeta$ 是一个 $d$ 次单位根，则 $\zeta$ 也是一个 $n$ 次单位根，因为
\[
\zeta^n=(\zeta^d)^{n/d}=1.
\]
因此
\[
\mu_d\subseteq \mu_n \qquad \text{对所有 } d\mid n.
\]
反过来，群 $\mu_n$ 中任意元素的阶都是 $n$ 的一个因子，所以如果 $\zeta$ 是一个 $n$ 次单位根，同时又是某个较小的 $d$ 的 $d$ 次单位根，则 $d\mid n$。

\noindent\textbf{定义.} 定义第 $n$ 个\emph{分圆多项式} $\Phi_n(x)$ 为以本原 $n$ 次单位根为根的多项式：
\[
\Phi_n(x)=\prod_{\substack{\zeta\text{ primitive}\\ \zeta\in\mu_n}}(x-\zeta)
=\prod_{\substack{1\leq a<n\\ (a,n)=1}}(x-\zeta_n^a)
\]
（它的次数是 $\varphi(n)$）。

多项式 $x^n-1$ 的根恰好是 $n$ 次单位根，所以我们有分解
\[
x^n-1=\prod_{\substack{\zeta^n=1\\ \text{i.e. }\zeta\in\mu_n}}(x-\zeta).
\]
如果我们把因子 $(x-\zeta)$ 按 $\zeta$ 在 $\mu_n$ 中的阶为 $d$ 的情形分组（即 $\zeta$ 是一个本原 $d$ 次单位根），就得到
\[
x^n-1=\prod_{d\mid n}\ \prod_{\substack{\zeta\in\mu_d\\ \zeta\text{ primitive}}}(x-\zeta).
\]
内层乘积按定义就是 $\Phi_d(x)$，所以我们有分解
\[
x^n-1=\prod_{d\mid n}\Phi_d(x). \tag{13.4}
\]
顺便指出，比较次数给出恒等式
\[
n=\sum_{d\mid n}\varphi(d).
\]

这个分解允许我们递归地计算任意 $n$ 的 $\Phi_n(x)$：显然
\[
\Phi_1(x)=x-1\qquad\text{且}\qquad \Phi_2(x)=x+1.
\]
然后
\[
x^3-1=\Phi_1(x)\Phi_3(x)=(x-1)\Phi_3(x)
\]
给出
\[
\Phi_3(x)=x^2+x+1.
\]
类似地，
\[
x^4-1=\Phi_1(x)\Phi_2(x)\Phi_4(x)=(x-1)(x+1)\Phi_4(x)
\]
给出
\[
\Phi_4(x)=x^2+1
\]
（在这些情形下，这些多项式也可以从显式的单位根直接得到）。继续按这种方式可对任意 $n$ 计算 $\Phi_n(x)$。还要注意，对素数 $p$，我们重新得到我们的多项式
\[
\Phi_p(x)=x^{p-1}+x^{p-2}+\cdots+x+1.
\]
对若干小的 $n$ 值，这些多项式是
\[
\begin{aligned}
	\Phi_5(x)&=x^4+x^3+x^2+x+1,\\
	\Phi_6(x)&=x^2-x+1,\\
	\Phi_7(x)&=x^6+x^5+x^4+x^3+x^2+x+1,\\
	\Phi_8(x)&=x^4+1,\\
	\Phi_9(x)&=x^6+x^3+1,\\
	\Phi_{10}(x)&=x^4-x^3+x^2-x+1,\\
	\Phi_{11}(x)&=x^{10}+x^9+\cdots+x+1,\\
	\Phi_{12}(x)&=x^4-x^2+1.
\end{aligned}
\]

对上面计算的所有值，$\Phi_n(x)$ 都是一个（首一的）整数系数多项式。事实总是如此：

\noindent\textbf{引理 40.} 分圆多项式 $\Phi_n(x)$ 是 $\mathbb Z[x]$ 中次数为 $\varphi(n)$ 的首一多项式。

\begin{proof}
	显然 $\Phi_n(x)$ 是首一的并且次数为 $\varphi(n)$。我们必须证明其系数属于 $\mathbb Z$。我们对 $n$ 用归纳法。结论对 $n=1$（以及 $n\leq12$）成立。归纳地假设对所有 $1\leq d<n$ 有 $\Phi_d(x)\in\mathbb Z[x]$。那么
	\[
	x^n-1=f(x)\Phi_n(x),
	\]
	其中
	\[
	f(x)=\prod_{\substack{d\mid n\\ d\neq n}}\Phi_d(x)
	\]
	是首一的并且系数在 $\mathbb Z$ 中。由于 $f(x)$ 显然在 $F[x]$ 中整除 $x^n-1$，其中 $F=\mathbb Q(\zeta_n)$ 是 $n$ 次单位根的域，并且 $f(x)$ 和 $x^n-1$ 都有系数在 $\mathbb Q$ 中，由除法算法（参见第 9.2 节末尾的注记），$f(x)$ 在 $\mathbb Q[x]$ 中整除 $x^n-1$。由 Gauss 引理，$f(x)$ 在 $\mathbb Z[x]$ 中整除 $x^n-1$，所以 $\Phi_n(x)\in\mathbb Z[x]$。
\end{proof}

顺便说一句，虽然上面例子中计算出的 $\Phi_n(x)$ 的所有系数都是 $0,\pm1$，但已知存在具有任意大系数的分圆多项式。

\noindent\textbf{定理 41.} 分圆多项式 $\Phi_n(x)$ 是 $\mathbb Z[x]$ 中次数为 $\varphi(n)$ 的不可约首一多项式。

\begin{proof}
	我们必须证明 $\Phi_n(x)$ 不可约。若不然，则有分解
	\[
	\Phi_n(x)=f(x)g(x)\qquad\text{其中 } f(x),g(x)\text{ 是 }\mathbb Z[x]\text{ 中首一多项式}
	\]
	其中我们取 $f(x)$ 为 $\Phi_n(x)$ 的一个不可约因子。设 $\zeta$ 是 $1$ 的一个本原 $n$ 次根并且是 $f(x)$ 的一个根（于是 $f(x)$ 是 $\zeta$ 在 $\mathbb Q$ 上的极小多项式），并令 $p$ 表示任意一个不整除 $n$ 的素数。那么 $\zeta^p$ 仍然是 $1$ 的一个本原 $n$ 次根，因而是 $f(x)$ 或 $g(x)$ 的一个根。
	
	假设 $g(\zeta^p)=0$。那么 $\zeta$ 是 $g(x^p)$ 的一个根，并且因为 $f(x)$ 是 $\zeta$ 的极小多项式，$f(x)$ 必须在 $\mathbb Z[x]$ 中整除 $g(x^p)$，也就是说
	\[
	g(x^p)=f(x)h(x),\qquad h(x)\in\mathbb Z[x].
	\]
	如果把这个等式模 $p$ 约化，就得到
	\[
	\overline{g}(x^p)=\overline{f}(x)\overline{h}(x) \qquad\text{在 }\mathbb F_p[x]\text{ 中}.
	\]
	由上一节的注记，
	\[
	\overline{g}(x^p)=(\overline{g}(x))^p,
	\]
	所以我们有等式
	\[
	(\overline{g}(x))^p=\overline{f}(x)\overline{h}(x)
	\]
	在 U.F.D. $\mathbb F_p[x]$ 中成立。于是 $\overline{f}(x)$ 和 $\overline{g}(x)$ 在 $\mathbb F_p[x]$ 中有一个公共因子。
	
	现在，从 $\Phi_n(x)=f(x)g(x)$，模 $p$ 约化可知
	\[
	\overline{\Phi}_n(x)=\overline{f}(x)\overline{g}(x),
	\]
	于是由上面可知 $\overline{\Phi}_n(x)\in\mathbb F_p[x]$ 有重根。但那么 $x^n-1$ 在 $\mathbb F_p$ 上也会有重根，因为它以 $\overline{\Phi}_n(x)$ 为因子。这是一个矛盾，因为我们已经在上一节看到，在任何特征不整除 $n$ 的域上，$x^n-1$ 有 $n$ 个互异的根。
	
	因此 $\zeta^p$ 必须是 $f(x)$ 的一个根。由于这适用于 $f(x)$ 的每一个根 $\zeta$，推出对每个与 $n$ 互素的整数 $a$，$\zeta^a$ 都是 $f(x)$ 的根：把
	\[
	a=p_1p_2\cdots p_k
	\]
	写成不整除 $n$ 的（不必互异的）素数的乘积，则 $\zeta^{p_1}$ 是 $f(x)$ 的根，因而 $(\zeta^{p_1})^{p_2}$ 也是 $f(x)$ 的根，等等。但这意味着每个本原 $n$ 次单位根都是 $f(x)$ 的根，即 $f(x)=\Phi_n(x)$，说明 $\Phi_n(x)$ 不可约。
\end{proof}

\noindent\textbf{推论 42.} $n$ 次单位根的分圆域在 $\mathbb Q$ 上的次数是 $\varphi(n)$：
\[
[\mathbb Q(\zeta_n):\mathbb Q]=\varphi(n).
\]

\begin{proof}
	由定理，$\Phi_n(x)$ 是任意本原 $n$ 次单位根 $\zeta_n$ 的极小多项式。
\end{proof}

\paragraph{例}
$8$ 次单位根的分圆域 $\mathbb Q(\zeta_8)$ 在 $\mathbb Q$ 上的次数是 $\varphi(8)=4$。此域包含 $4$ 次单位根，即
\[
\mathbb Q(i)\subset \mathbb Q(\zeta_8),
\]
并且还包含元素
\[
\zeta_8+\zeta_8^7=\sqrt2
\]
（回忆第 4 节中显式的单位根）。于是
\[
\mathbb Q(\zeta_8)=\mathbb Q(i,\sqrt2).
\]

在练习中概述的分圆多项式的一个有趣数论应用是证明：对任意 $n$，存在无穷多个与 $1$ 模 $n$ 同余的素数。对于素数 $\ell$，$\Phi_\ell(x)$（它在 $\mathbb Z[x]$ 中不可约）在 $\mathbb F_p[x]$ 中的完全分解在下面练习 8 中描述。

在发展出一些 Galois 理论之后，我们将回到分圆域的例子。

\subsection*{练习}

\begin{enumerate}[label=\arabic*.]
	 \item 设 $m$ 和 $n$ 是互素的正整数。设 $\zeta_m$ 是一个本原 $m$ 次单位根，并设 $\zeta_n$ 是一个本原 $n$ 次单位根。证明 $\zeta_m\zeta_n$ 是一个本原 $mn$ 次单位根。
	 
	 \item 设 $\zeta_n$ 是一个本原 $n$ 次单位根，并设 $d$ 是 $n$ 的一个因子。证明 $\zeta_n^d$ 是一个本原 $(n/d)$ 次单位根。
	 
	 \item 证明若一个域包含 $n$ 次单位根且 $n$ 为奇数，则它也包含 $2n$ 次单位根。
	 
	 \item 证明若 $n=p^km$，其中 $p$ 是素数且 $m$ 与 $p$ 互素，则在一个特征为 $p$ 的域上，恰有 $m$ 个互异的 $n$ 次单位根。
	 
	 \item 证明在 $\mathbb Q$ 的任意有限扩张 $K$ 中，单位根只有有限多个。
	 
	 \item 证明对奇数 $n$，$n>1$，有
	 \[
	 \Phi_{2n}(x)=\Phi_n(-x).
	 \]
	 
	 \item 使用第 14.3 节中指出的 Möbius 反演公式证明
	 \[
	 \Phi_m(x)=\prod_{d\mid m}(x^d-1)^{\mu(m/d)}.
	 \]
	 
	 \item 设 $\ell$ 是素数，并设
	 \[
	 \Phi_\ell(x)=\frac{x^\ell-1}{x-1}=x^{\ell-1}+x^{\ell-2}+\cdots+x+1\in\mathbb Z[x]
	 \]
	 是第 $\ell$ 个分圆多项式，它由定理 41 在 $\mathbb Z$ 上不可约。这个练习确定 $\Phi_\ell(x)$ 模任意素数 $p$ 的分解。令 $\zeta$ 表示任意固定的本原 $\ell$ 次单位根。
	 \begin{enumerate}[label=(\alph*)]
	 	\item 证明若 $p=\ell$，则
	 	\[
	 	\Phi_\ell(x)=(x-1)^{\ell-1}\in\mathbb F_\ell[x].
	 	\]
	 	
	 	\item 假设 $p\neq\ell$，并令 $f$ 表示 $p$ 模 $\ell$ 的阶，即 $f$ 是使得 $p^f\equiv1\pmod\ell$ 的 $p$ 的最小幂。使用 $\mathbb F_{p^n}^{\times}$ 是循环群这一事实，证明 $n=f$ 是使得 $\zeta\in\mathbb F_{p^n}$ 的 $p$ 的最小幂 $p^n$。推出 $\zeta$ 在 $\mathbb F_p$ 上的极小多项式次数为 $f$。
	 	
	 	\item 证明对任意不能被 $\ell$ 整除的整数 $a$，有
	 	\[
	 	\mathbb F_p(\zeta)=\mathbb F_p(\zeta^a).
	 	\]
	 	[一个包含关系是显然的。对于另一个，注意 $\zeta=(\zeta^a)^b$，其中 $b$ 是 $a$ 模 $\ell$ 的乘法逆元。] 利用 (b) 推出，在 $\mathbb F_p[x]$ 中，$\Phi_\ell(x)$ 是 $(\ell-1)/f$ 个互异的次数为 $f$ 的不可约多项式的乘积。
	 	
	 	\item 特别地，证明在 $\mathbb F_p[x]$ 中看，
	 	\[
	 	\Phi_7(x)=x^6+x^5+\cdots+x+1
	 	\]
	 	当 $p=7$ 时为 $(x-1)^6$；当 $p\equiv1\pmod7$ 时是互异线性因子的乘积；当 $p\equiv6\pmod7$ 时是 $3$ 个不可约二次多项式的乘积；当 $p\equiv2,4\pmod7$ 时是 $2$ 个不可约三次多项式的乘积；当 $p\equiv3,5\pmod7$ 时不可约。
	 \end{enumerate}
	 
	 \item 设 $A$ 是一个 $\mathbb C$ 上的 $n\times n$ 矩阵，并且对某个整数 $k\geq1$ 有 $A^k=I$。证明 $A$ 可对角化。证明矩阵
	 \[
	 A=\begin{pmatrix}1&\alpha\\0&1\end{pmatrix}
	 \]
	 其中 $\alpha$ 是特征为 $p$ 的域中的一个元素，满足 $A^p=I$，并且若 $\alpha\neq0$ 则不能对角化。
	 
	 \item 令 $\varphi$ 表示有限域 $\mathbb F_{p^n}$ 上的 Frobenius 映射 $x\mapsto x^p$。证明 $\varphi$ 给出 $\mathbb F_{p^n}$ 到自身的一个同构（这样的同构称为一个自同构）。证明 $\varphi^n$ 是恒等映射，并且没有更低的 $\varphi$ 的幂是恒等映射。
	 
	 \item 令 $\varphi$ 表示有限域 $\mathbb F_{p^n}$ 上的 Frobenius 映射 $x\mapsto x^p$，如前一个练习中那样。把 $\varphi$ 看作 $n$ 维 $\mathbb F_p$-向量空间 $\mathbb F_{p^n}$ 的一个 $\mathbb F_p$-线性变换，求 $\varphi$ 在 $\mathbb F_p$ 上的有理标准形。
	 
	 \item 令 $\varphi$ 表示有限域 $\mathbb F_{p^n}$ 上的 Frobenius 映射 $x\mapsto x^p$，如前一个练习中那样。把 $\varphi$ 看作 $n$ 维 $\mathbb F_p$-向量空间 $\mathbb F_{p^n}$ 的一个 $\mathbb F_p$-线性变换，求 $\varphi$ 的 Jordan 标准形（在一个包含全部特征值的域上）。
	 
	 \item （有限除环上的 Wedderburn 定理）本练习概述 Wedderburn 定理的一个证明（依照 Witt），该定理断言一个有限除环 $D$ 是域（即，是交换的）。
	 \begin{enumerate}[label=(\alph*)]
	 	\item 令 $Z$ 表示 $D$ 的中心（即，$D$ 中与 $D$ 的每个元素交换的元素）。证明 $Z$ 是一个包含某个素数 $p$ 的 $\mathbb F_p$ 的域。若 $Z=\mathbb F_q$，证明 $D$ 的阶为 $q^n$，其中 $n$ 是某个整数 [$D$ 是 $Z$ 上的向量空间]。
	 	
	 	\item $D$ 的非零元素 $D^\times$ 构成一个乘法群。对任意 $x\in D^\times$，证明 $D$ 中与 $x$ 交换的元素构成一个包含 $Z$ 的除环。说明这个除环的阶为 $q^m$，其中 $m$ 是某个整数，并且若 $x$ 不是 $Z$ 的元素，则 $m<n$。
	 	
	 	\item 证明群 $D^\times$ 的类方程（定理 4.7）是
	 	\[
	 	q^n-1=(q-1)+\sum_{i=1}^r\frac{q^n-1}{|C_{D^\times}(x_i)|},
	 	\]
	 	其中 $x_1,x_2,\ldots,x_r$ 是 $D^\times$ 中不包含在 $D^\times$ 的中心里的互异共轭类的代表元。由 (b) 推出对每个 $i$，
	 	\[
	 	|C_{D^\times}(x_i)|=q^{m_i}-1
	 	\]
	 	其中某个 $m_i<n$。
	 	
	 	\item 证明由于
	 	\[
	 	\frac{q^n-1}{q^{m_i}-1}
	 	\]
	 	是整数（即，指数 $|D^\times:C_{D^\times}(x_i)|$），所以 $m_i$ 整除 $n$（参见第 5 节练习 4）。推出 $\Phi_n(x)$ 整除 $(x^n-1)/(x^{m_i}-1)$，从而整数 $\Phi_n(q)$ 整除 $(q^n-1)/(q^{m_i}-1)$，$i=1,2,\ldots,r$。
	 	
	 	\item 证明 (c) 和 (d) 推出
	 	\[
	 	\Phi_n(q)=\prod_{\substack{\zeta\text{ primitive}}}(q-\zeta)
	 	\]
	 	整除 $q-1$。证明对任意单位根 $\zeta\neq1$，有 $|q-\zeta|>q-1$（复绝对值）[注意，$1$ 是单位圆上到实直线上的点 $q$ 最近的点]。推出 $n=1$，即 $D=Z$ 是一个域。
	 \end{enumerate}
	 
	 下面的练习给出一个证明：对任意正整数 $m$，存在无穷多个素数 $p$ 满足 $p\equiv1\pmod m$。这是算术级数中素数的 Dirichlet 定理的一个特殊情形；该定理更一般地断言，对任意与 $m$ 互素的 $a$，存在无穷多个素数 $p$ 满足 $p\equiv a\pmod m$。
	 
	 \item 给定任意次数至少为一的首一多项式 $P(x)\in\mathbb Z[x]$，说明整数
	 \[
	 P(1),\ P(2),\ P(3),\ \ldots,\ P(n),\ \ldots
	 \]
	 有无穷多个互异的素因子。
	 
	 [假设 $p_1,p_2,\ldots,p_k$ 是整除各个值 $P(n)$，$n=1,2,\ldots$ 的唯一素数。令 $N$ 是满足 $P(N)=a\neq0$ 的整数。证明
	 \[
	 Q(x)=a^{-1}P(N+ap_1p_2\cdots p_kx)
	 \]
	 是 $\mathbb Z[x]$ 的一个元素，并且对 $n=1,2,\ldots$ 有
	 \[
	 Q(n)\equiv1\pmod {p_1p_2\cdots p_k}.
	 \]
	 推出存在某个整数 $M$ 使得 $Q(M)$ 有一个不同于 $p_1,p_2,\ldots,p_k$ 的素因子，因而 $P(N+ap_1p_2\cdots p_kM)$ 有一个不同于 $p_1,p_2,\ldots,p_k$ 的素因子。]
	 
	 \item 设 $p$ 是一个不整除 $m$ 的奇素数，并令 $\Phi_m(x)$ 为第 $m$ 个分圆多项式。假设 $a\in\mathbb Z$ 满足
	 \[
	 \Phi_m(a)\equiv0\pmod p.
	 \]
	 证明 $a$ 与 $p$ 互素，并且 $a$ 在 $(\mathbb Z/p\mathbb Z)^\times$ 中的阶恰好为 $m$。[因为
	 \[
	 x^m-1=\prod_{d\mid m}\Phi_d(x)=\Phi_m(x)\prod_{\substack{d\mid m\\ d<m}}\Phi_d(x),
	 \]
	 我们首先看到 $a^m-1\equiv0\pmod p$，即 $a^m\equiv1\pmod p$。若 $a\bmod p$ 的阶小于 $m$，则对某个整除 $m$ 的 $d$ 有 $a^d\equiv1\pmod p$，因此对某个 $d<m$ 有 $\Phi_d(a)\equiv0\pmod p$。但于是 $x^m-1$ 模 $p$ 会有 $a$ 作为重根，矛盾。]
	 
	 \item 设 $a\in\mathbb Z$。说明如果 $p$ 是整除 $\Phi_m(a)$ 的奇素数，则或者 $p$ 整除 $m$，或者 $p\equiv1\pmod m$。
	 
	 \item 证明存在无穷多个素数 $p$ 满足 $p\equiv1\pmod m$。
	 
\end{enumerate}
	
% \chapter{Galois 理论}

\section{基本定义}

在前一章中，我们证明了域 $F$ 的一个有限扩张的存在性，这个扩张含有给定多项式 $f(x)$ 的所有根，其中 $f(x)$ 的系数属于 $F$。Galois 理论（以 Évariste Galois，1811--1832 命名）的主要思想，是考察 $f(x)$ 的根的置换群同它的分裂域的代数结构之间的关系。这个联系由下一节的基本定理给出。它可以看作数学中一个重要思想的又一个（极其优美的）应用：一个（在我们这里是代数的）对象作用在另一个对象上，会给出关于二者的结构信息。

本节中我们引入所关心对象的术语和基本性质。令 $K$ 为一个域。

\noindent\textbf{定义.}
\begin{enumerate}[label=(\arabic*)]
\item $K$ 到它自身的一个同构称为 $K$ 的一个\emph{自同构}。$K$ 的自同构全体记为 $\operatorname{Aut}(K)$。若 $a\in K$，我们将把 $\sigma(a)$ 记作 $\sigma a$。
\item 自同构 $\sigma\in\operatorname{Aut}(K)$ 称为\emph{固定}元素 $a\in K$，如果 $\sigma a=a$。如果 $F$ 是 $K$ 的一个子集（例如一个子域），则称自同构 $\sigma$ \emph{固定} $F$，如果它固定 $F$ 的所有元素，即对所有 $a\in F$ 都有 $\sigma a=a$。
\end{enumerate}

注意，任何域至少有一个自同构，即恒等映射，记为 $1$，有时称为平凡自同构。域 $K$ 的素域由 $1\in K$ 生成，并且由于任意自同构 $\sigma$ 都把 $1$ 送到 $1$（也把 $0$ 送到 $0$），即 $\sigma(1)=1$，所以对素域中的所有 $a$ 都有 $\sigma a=a$。因此，域 $K$ 的任意自同构都固定它的素子域。特别地，我们看到 $\mathbb Q$ 和 $\mathbb F_p$ 只有平凡自同构：
\[
\operatorname{Aut}(\mathbb Q)=\{1\},\qquad \operatorname{Aut}(\mathbb F_p)=\{1\}.
\]

\noindent\textbf{定义.} 设 $K/F$ 是域扩张。记 $\operatorname{Aut}(K/F)$ 为 $K$ 的所有固定 $F$ 的自同构的集合。

注意，如果 $F$ 是 $K$ 的素子域，则 $\operatorname{Aut}(K)=\operatorname{Aut}(K/F)$，因为 $K$ 的每个自同构都自动固定 $F$。

如果 $\sigma$ 和 $\tau$ 是 $K$ 的自同构，则复合 $\sigma\tau$（以及复合 $\tau\sigma$，它未必相同）是有定义的，并且仍然是 $K$ 的自同构。

\noindent\textbf{命题 1.} $\operatorname{Aut}(K)$ 在复合下是一个群，并且 $\operatorname{Aut}(K/F)$ 是一个子群。

\begin{proof}
显然 $\operatorname{Aut}(K)$ 是一个群。如果 $\sigma$ 和 $\tau$ 是固定 $F$ 的 $K$ 的自同构，则 $\sigma\tau$ 和 $\sigma^{-1}$ 在 $F$ 上也都是恒等映射，这说明 $\operatorname{Aut}(K/F)$ 是子群。
\end{proof}

下面的命题对于确定代数扩张的自同构极为有用。

\noindent\textbf{命题 2.} 设 $K/F$ 是一个域扩张，并设 $\alpha\in K$ 在 $F$ 上代数。则对任意 $\sigma\in\operatorname{Aut}(K/F)$，$\sigma\alpha$ 是 $\alpha$ 在 $F$ 上的极小多项式的一个根；也就是说，$\operatorname{Aut}(K/F)$ 置换不可约多项式的根。等价地，任何系数在 $F$ 中且以 $\alpha$ 为根的多项式，也以 $\sigma\alpha$ 为根。

\begin{proof}
设 $\alpha$ 满足方程
\[
\alpha^n+a_{n-1}\alpha^{n-1}+\cdots+a_1\alpha+a_0=0,
\]
其中 $a_0,a_1,\ldots,a_{n-1}$ 是 $F$ 的元素。施以自同构 $\sigma$，得到（利用 $\sigma$ 是加法同态）
\[
\sigma(\alpha^n)+\sigma(a_{n-1}\alpha^{n-1})+\cdots+\sigma(a_1\alpha)+\sigma(a_0)=\sigma(0)=0.
\]
再利用 $\sigma$ 也是乘法同态，这变为
\[
(\sigma(\alpha))^n+\sigma(a_{n-1})(\sigma(\alpha))^{n-1}+\cdots+
\sigma(a_1)\sigma(\alpha)+\sigma(a_0)=0.
\]
由假设，$\sigma$ 固定 $F$ 的所有元素，所以 $\sigma(a_i)=a_i$，$i=0,1,\ldots,n-1$。因此
\[
(\sigma\alpha)^n+a_{n-1}(\sigma\alpha)^{n-1}+\cdots+a_1(\sigma\alpha)+a_0=0.
\]
这正是说 $\sigma\alpha$ 是 $F$ 上同一个多项式的根。命题得证。
\end{proof}

\paragraph{例}
\begin{enumerate}[label=(\arabic*)]
\item 令 $K=\mathbb Q(\sqrt2)$。如果 $\tau\in\operatorname{Aut}(\mathbb Q(\sqrt2))=\operatorname{Aut}(\mathbb Q(\sqrt2)/\mathbb Q)$，则 $\tau(\sqrt2)=\pm\sqrt2$，因为这是 $\sqrt2$ 的极小多项式的两个根。由于 $\tau$ 固定 $\mathbb Q$，这完全确定了 $\tau$：
\[
\tau(a+b\sqrt2)=a\pm b\sqrt2.
\]
映射 $\sqrt2\mapsto\sqrt2$ 只是 $\mathbb Q(\sqrt2)$ 的恒等自同构 $1$。映射
\[
\sigma:\sqrt2\mapsto-\sqrt2
\]
是推论 13.7 后的例 2 中考虑的同构。因此
\[
\operatorname{Aut}(\mathbb Q(\sqrt2))=\operatorname{Aut}(\mathbb Q(\sqrt2)/\mathbb Q)=\{1,\sigma\}
\]
是由 $\sigma$ 生成的 $2$ 阶循环群。

\item 令 $K=\mathbb Q(\sqrt[3]{2})$。如前所述，如果 $\tau\in\operatorname{Aut}(K/\mathbb Q)$，则 $\tau$ 完全由它在 $\sqrt[3]{2}$ 上的作用确定，因为
\[
\tau(a+b\sqrt[3]{2}+c(\sqrt[3]{2})^2)=a+b\tau\sqrt[3]{2}+c(\tau\sqrt[3]{2})^2.
\]
由于 $\tau\sqrt[3]{2}$ 必须是 $x^3-2$ 的根，而此方程的另外两个根不属于 $K$（回忆这个多项式的分裂域在 $\mathbb Q$ 上次数为 $6$），唯一可能是 $\tau\sqrt[3]{2}=\sqrt[3]{2}$，即 $\tau=1$。因此
\[
\operatorname{Aut}(\mathbb Q(\sqrt[3]{2})/\mathbb Q)=1
\]
是平凡群。
\end{enumerate}

一般地，如果 $K$ 在 $F$ 上由某个元素集合生成，那么 $\operatorname{Aut}(K/F)$ 中的任意自同构 $\sigma$ 完全由它在这些生成元上的作用决定。如果 $K/F$ 有限，那么 $K$ 在 $F$ 上由有限多个代数元素生成，所以由命题可知，固定 $F$ 的 $K$ 的自同构数目有限，即 $\operatorname{Aut}(K/F)$ 是有限群。特别地，有限扩张的自同构可以看作有限多个方程的根的置换（然而并非每个置换都给出自同构，如上面的例 2 所说明）。正是对方程根的置换的研究，把 Galois 引向了我们正在描述的理论。

我们已经把一个群 $\operatorname{Aut}(K/F)$，即固定 $F$ 的 $K$ 的自同构群，与每个域扩张 $K/F$（等价地，与 $K$ 的子域 $F$）联系起来。也可以反过来，从一个自同构群联系到一个域扩张。

\noindent\textbf{命题 3.} 设 $H\leq \operatorname{Aut}(K)$ 是 $K$ 的自同构群的一个子群。那么 $K$ 中被 $H$ 的所有元素固定的元素集合 $F$ 是 $K$ 的一个子域。

\begin{proof}
令 $h\in H$，并令 $a,b\in F$。于是由定义 $h(a)=a$、$h(b)=b$，所以
\[
h(a\pm b)=h(a)\pm h(b)=a\pm b,
\]
\[
h(ab)=h(a)h(b)=ab,
\]
并且
\[
h(a^{-1})=h(a)^{-1}=a^{-1}.
\]
因此 $F$ 封闭，故为 $K$ 的子域。
\end{proof}

注意，在这个命题中，$H$ 实际上是 $\operatorname{Aut}(K)$ 的一个子群并不重要——$K$ 中被 $\operatorname{Aut}(K)$ 的一个子集的所有元素固定的元素集合也仍是 $K$ 的一个子域。

\noindent\textbf{定义.} 如果 $H$ 是 $K$ 的自同构群的一个子群，则 $K$ 中被 $H$ 的所有元素固定的子域称为 $H$ 的\emph{固定域}。

\noindent\textbf{命题 4.} 上面定义的由域到群、由群到域的对应是反包含的，即
\begin{enumerate}[label=(\arabic*)]
\item 如果 $F_1\subseteq F_2\subseteq K$ 是 $K$ 的两个子域，则
\[
\operatorname{Aut}(K/F_2)\leq \operatorname{Aut}(K/F_1),
\]
并且
\item 如果 $H_1\leq H_2\leq \operatorname{Aut}(K)$ 是两个自同构子群，其相应固定域分别为 $F_1$ 和 $F_2$，则 $F_2\subseteq F_1$。
\end{enumerate}

\begin{proof}
固定 $F_2$ 的任意 $K$ 的自同构也固定它的子域 $F_1$，这给出 (1)。第二个断言类似证明。
\end{proof}

\paragraph{例}
\begin{enumerate}[label=(\arabic*)]
\item 设 $K=\mathbb Q(\sqrt2)$，如上例 1。则 $\operatorname{Aut}(\mathbb Q(\sqrt2))=\operatorname{Aut}(\mathbb Q(\sqrt2)/\mathbb Q)=\{1,\sigma\}$ 的固定域是 $\mathbb Q(\sqrt2)$ 中满足
\[
\sigma(a+b\sqrt2)=a+b\sqrt2
\]
的元素集合，因为每个元素都被恒等自同构固定。这就是方程
\[
a-b\sqrt2=a+b\sqrt2,
\]
它等价于 $b=0$，所以 $\operatorname{Aut}(\mathbb Q(\sqrt2)/\mathbb Q)$ 的固定域就是 $\mathbb Q$。

\item 现在设 $K=\mathbb Q(\sqrt[3]{2})$，如上例 2。在此情形中 $\operatorname{Aut}(K)=1$，所以 $K$ 的每个元素都被固定，即 $\operatorname{Aut}(\mathbb Q(\sqrt[3]{2})/\mathbb Q)$ 的固定域是 $\mathbb Q(\sqrt[3]{2})$。
\end{enumerate}

给定 $K$ 的一个子域 $F$，对应的群是固定 $F$ 的 $K$ 的自同构集合。给定 $K$ 的一个自同构群，则通过取这些自同构的固定域来定义相应的扩张。在上面的第一个例子中，从 $\mathbb Q(\sqrt2)$ 的子域 $\mathbb Q$ 出发，得到群 $\{1,\sigma\}$；再从群 $\{1,\sigma\}$ 出发，又得到子域 $\mathbb Q$，所以二者之间存在一种“对偶”。然而在第二个例子中，从 $\mathbb Q(\sqrt[3]{2})$ 的子域 $\mathbb Q$ 出发，只得到平凡群；再从平凡群出发，得到整个域 $\mathbb Q(\sqrt[3]{2})$。

考察这两个例子提示我们：在第二个例子中，自同构“不够多”，不能迫使固定域为 $\mathbb Q$，而只能得到整个 $\mathbb Q(\sqrt[3]{2})$。这又似乎是因为 $x^3-2$ 的其他根，也就是自同构下 $\sqrt[3]{2}$ 的唯一可能像，并不是 $\mathbb Q(\sqrt[3]{2})$ 的元素。（不过即使它们属于该域，我们也还需要检查由此定义的附加映射是否真的是自同构。）现在我们精确地定义带有“足够多”自同构的域（这将导向 Galois 扩张的定义）。正如甚至从这两个例子就可能猜到的那样（我们将在下一节证明），这些域与分裂域有关。

我们首先考察分裂域情形中自同构群的大小。设 $F$ 是一个域，并设 $E$ 是 $F$ 上 $f(x)\in F[x]$ 的分裂域。主要工具是关于同构延拓存在性的定理 13.27，它断言：$F$ 与 $F'$ 的任意同构 $\varphi:F\xrightarrow{\sim}F'$ 都能延拓为同构 $\sigma:E\xrightarrow{\sim}E'$，其中 $E'$ 是 $f'(x)=\varphi(f(x))\in F'[x]$ 的分裂域。

我们现在用对 $[E:F]$ 的归纳证明，这样的延拓数目至多为 $[E:F]$，并且如果 $f(x)$ 在 $F$ 上可分，则等号成立。如果 $[E:F]=1$，则 $E=F$、$E'=F'$、$\sigma=\varphi$，延拓的数目为 $1$。如果 $[E:F]>1$，那么 $f(x)$ 至少有一个次数 $>1$ 的不可约因子 $p(x)$，并有 $f'(x)$ 的相应不可约因子 $p'(x)$。令 $\alpha$ 为 $p(x)$ 的一个固定根。如果 $\sigma$ 是 $\varphi$ 到 $E$ 的任意延拓，那么 $\sigma$ 限制在 $E$ 的子域 $F(\alpha)$ 上，是 $F(\alpha)$ 到 $E'$ 的某个子域的同构 $\tau$。同构 $\tau$ 完全由它在 $\alpha$ 上的作用，即由 $\tau\alpha$，决定，因为 $\alpha$ 在 $F$ 上生成 $F(\alpha)$。正如在命题 2 中一样，我们看到 $\tau\alpha$ 必须是 $p'(x)$ 的某个根 $\beta$。于是有图
\[
\begin{array}{ccccc}
\sigma:& E & \longrightarrow & E'& \\
& | && |& \\
\tau:& F(\alpha)& \longrightarrow & F'(\beta)& \\
& | && |& \\
\varphi:& F & \longrightarrow & F'& .
\end{array}
\]
反过来，对 $p'(x)$ 的任意根 $\beta$，都存在给出这种图的延拓 $\tau$ 和 $\sigma$（这就是定理 13.8 和定理 13.27）。因此，为了计算延拓 $\sigma$ 的数目，只需要计算这种图的可能数目。

$\varphi$ 到同构 $\tau$ 的延拓数等于 $p'(x)$ 的不同根的数目。由于 $p(x)$ 和 $p'(x)$ 的次数都等于 $[F(\alpha):F]$，我们看到 $\varphi$ 到 $\tau$ 的延拓数至多为 $[F(\alpha):F]$；如果 $p(x)$ 的根互异，则等号成立。

由于 $E$ 也是 $f(x)$ 在 $F(\alpha)$ 上的分裂域，$E'$ 是 $f'(x)$ 在 $F'(\beta)$ 上的分裂域，并且 $[E:F(\alpha)]<[E:F]$，我们可以把归纳假设应用于这些域扩张。由归纳可知，$\tau$ 到 $\sigma$ 的延拓数至多为 $[E:F(\alpha)]$；如果 $f(x)$ 的根互异，则等号成立。

由
\[
[E:F]=[E:F(\alpha)][F(\alpha):F]
\]
可知，$\varphi$ 到 $\sigma$ 的延拓数至多为 $[E:F]$。如果 $p(x)$ 和 $f(x)$ 都有互异根，则等号成立；这等价于 $f(x)$ 有互异根，因为 $p(x)$ 是 $f(x)$ 的因子。归纳证明完成。

在特殊情形 $F=F'$ 且 $\varphi$ 是恒等映射时，有 $f(x)=f'(x)$ 且 $E=E'$，于是 $E$ 到 $E'$ 的、限制到 $F$ 上为 $\varphi$ 的同构，正是固定 $F$ 的 $E$ 的自同构。我们把这表述如下：

\noindent\textbf{命题 5.} 设 $E$ 是多项式 $f(x)\in F[x]$ 在 $F$ 上的分裂域。则
\[
|\operatorname{Aut}(E/F)|\leq [E:F],
\]
如果 $f(x)$ 在 $F$ 上可分，则等号成立。

\paragraph{注记：} 虽然我们主要关心的是计算固定 $F$ 的 $E$ 的自同构（即上面 $F=F'$、$\varphi=1$ 的情形），但是仍然必须考虑更一般的 $\varphi$（以及不同的域 $F'$），因为证明中的归纳步骤涉及同一个多项式 $p(x)$ 的两个根对应的域 $F(\alpha)$ 与 $F(\beta)$。

可以修改上面的证明，更一般地证明对任意有限扩张 $K/F$ 都有 $|\operatorname{Aut}(K/F)|\leq [K:F]$（我们将在下一节从稍微不同的观点证明这一点）。这给出具有“足够多”自同构的域扩张的概念。

\noindent\textbf{定义.} 设 $K/F$ 是有限扩张。如果 $|\operatorname{Aut}(K/F)|=[K:F]$，则称 $K$ 在 $F$ 上是 \emph{Galois} 的，并称 $K/F$ 是一个 \emph{Galois 扩张}。如果 $K/F$ 是 Galois 的，则自同构群 $\operatorname{Aut}(K/F)$ 称为 $K/F$ 的 \emph{Galois 群}，记为 $\operatorname{Gal}(K/F)$。

\paragraph{注记：} 对任意 $K/F$，扩张 $K/F$ 的 Galois 群有时被定义为自同构群 $\operatorname{Aut}(K/F)$。我们采用上面的定义，是为了使记号 $\operatorname{Gal}(K/F)$ 强调扩张 $K/F$ 有最大可能数目的自同构。

\noindent\textbf{推论 6.} 如果 $K$ 是可分多项式 $f(x)$ 在 $F$ 上的分裂域，则 $K/F$ 是 Galois 的。

我们将在下一节看到其逆命题也为真，这将完全刻画 Galois 扩张。

还要注意，推论 6 蕴含任意多项式在 $\mathbb Q$ 上的分裂域都是 Galois 的，因为 $f(x)$ 的分裂域显然与 $f(x)$ 的不可约因子的乘积（即去掉重因子后得到的多项式）的分裂域相同，而后者是可分的（推论 13.34）。

\noindent\textbf{定义.} 如果 $f(x)$ 是 $F$ 上的可分多项式，则 $f(x)$ 在 $F$ 上的 \emph{Galois 群}是 $f(x)$ 的分裂域在 $F$ 上的 Galois 群。

\paragraph{例}
\begin{enumerate}[label=(\arabic*)]
\item 扩张 $\mathbb Q(\sqrt2)/\mathbb Q$ 是 Galois 的，其 Galois 群为
\[
\operatorname{Gal}(\mathbb Q(\sqrt2)/\mathbb Q)=\{1,\sigma\}\cong \mathbb Z/2\mathbb Z,
\]
其中 $\sigma$ 是自同构
\[
\sigma: \mathbb Q(\sqrt2)\longrightarrow \mathbb Q(\sqrt2),\qquad
 a+b\sqrt2\longmapsto a-b\sqrt2.
\]

\item 更一般地，任意特征不为 $2$ 的域 $F$ 的二次扩张 $K$ 都是 Galois 的。这来自推论 13.13 后关于二次扩张的讨论，该讨论说明 $F$ 的任意 $2$ 次扩张 $K$（当 $F$ 的特征不为 $2$ 时）具有形式 $F(\sqrt D)$，因而是 $x^2-D$ 的分裂域（因为若 $\sqrt D\in K$，则 $-\sqrt D\in K$）。

\item 扩张 $\mathbb Q(\sqrt[3]{2})/\mathbb Q$ 不是 Galois 的，因为它的自同构群只有 $1$ 个元素。

\item 扩张 $\mathbb Q(\sqrt2,\sqrt3)$ 在 $\mathbb Q$ 上是 Galois 的，因为它是多项式 $(x^2-2)(x^2-3)$ 的分裂域。任意自同构 $\sigma$ 完全由它在生成元 $\sqrt2$ 和 $\sqrt3$ 上的作用决定，而这两个元素必须分别被送到 $\pm\sqrt2$ 和 $\pm\sqrt3$。因此自同构唯一可能为映射
\[
\begin{array}{c|c|c|c}
\sqrt2\mapsto\sqrt2 & \sqrt2\mapsto-\sqrt2 & \sqrt2\mapsto\sqrt2 & \sqrt2\mapsto-\sqrt2\\
\sqrt3\mapsto\sqrt3 & \sqrt3\mapsto\sqrt3 & \sqrt3\mapsto-\sqrt3 & \sqrt3\mapsto-\sqrt3
\end{array}
\]
由于 Galois 群阶为 $4$，所有这些元素事实上都是 $\mathbb Q(\sqrt2,\sqrt3)$ 在 $\mathbb Q$ 上的自同构。

定义自同构 $\sigma$ 和 $\tau$ 为
\[
\sigma:\left\{
\begin{array}{l}
\sqrt2\mapsto-\sqrt2,\\
\sqrt3\mapsto\sqrt3,
\end{array}\right.
\qquad
\tau:\left\{
\begin{array}{l}
\sqrt2\mapsto\sqrt2,\\
\sqrt3\mapsto-\sqrt3,
\end{array}\right.
\]
或者更明确地，定义为
\[
\sigma: a+b\sqrt2+c\sqrt3+d\sqrt6\longmapsto a-b\sqrt2+c\sqrt3-d\sqrt6
\]
和
\[
\tau: a+b\sqrt2+c\sqrt3+d\sqrt6\longmapsto a+b\sqrt2-c\sqrt3-d\sqrt6
\]
（因为例如
\[
\sigma(\sqrt6)=\sigma(\sqrt2\sqrt3)=\sigma(\sqrt2)\sigma(\sqrt3)=(-\sqrt2)(\sqrt3)=-\sqrt6
\]
）。
于是 $\sigma^2(\sqrt2)=\sigma(\sigma\sqrt2)=\sigma(-\sqrt2)=\sqrt2$，并且显然 $\sigma^2(\sqrt3)=\sqrt3$。因此 $\sigma^2=1$ 是恒等自同构。类似地，$\tau^2=1$。自同构 $\sigma\tau$ 可直接计算：
\[
\sigma\tau(\sqrt2)=\sigma(\tau(\sqrt2))=\sigma(\sqrt2)=-\sqrt2
\]
且
\[
\sigma\tau(\sqrt3)=\sigma(\tau(\sqrt3))=\sigma(-\sqrt3)=-\sqrt3.
\]
因此 $\sigma\tau$ 是 Galois 群中剩下的非平凡自同构。由于这个自同构显然在 Galois 群中也有阶 $2$，我们有
\[
\operatorname{Gal}(\mathbb Q(\sqrt2,\sqrt3)/\mathbb Q)=\{1,\sigma,\tau,\sigma\tau\},
\]
即 Galois 群同构于 Klein 四元群。

与 $\operatorname{Gal}(\mathbb Q(\sqrt2,\sqrt3)/\mathbb Q)$ 的每个子群相联系的是 $\mathbb Q(\sqrt2,\sqrt3)$ 的相应固定子域。例如，与 $\{1,\sigma\tau\}$ 相应的子域是被映射
\[
\sigma\tau:a+b\sqrt2+c\sqrt3+d\sqrt6\longmapsto a-b\sqrt2-c\sqrt3+d\sqrt6
\]
固定的元素集合，这就是元素 $a+d\sqrt6$ 的集合，即域 $\mathbb Q(\sqrt6)$。类似地可以确定 Galois 群的其他子群的固定域：
\[
\begin{array}{c|c}
\text{子群} & \text{固定域}\\ \hline
\{1\} & \mathbb Q(\sqrt2,\sqrt3)\\
\{1,\sigma\} & \mathbb Q(\sqrt3)\\
\{1,\sigma\tau\} & \mathbb Q(\sqrt6)\\
\{1,\tau\} & \mathbb Q(\sqrt2)\\
\{1,\sigma,\tau,\sigma\tau\} & \mathbb Q
\end{array}
\]

\item $x^3-2$ 在 $\mathbb Q$ 上的分裂域是 $6$ 次 Galois 扩张。这个方程的根为
\[
\sqrt[3]{2},\quad \rho\sqrt[3]{2},\quad \rho^2\sqrt[3]{2},
\]
其中
\[
\rho=\frac{-1+\sqrt{-3}}{2}
\]
是一个本原三次单位根。因此分裂域可写为 $\mathbb Q(\sqrt[3]{2},\rho\sqrt[3]{2})$。任意自同构把这两个元素各自送到 $x^3-2$ 的某个根，给出 $9$ 种可能；但由于 Galois 群阶为 $6$，并非每个这样的映射都是域的自同构。

为了确定 Galois 群，我们使用一组更方便的生成元，即 $\sqrt[3]{2}$ 和 $\rho$。任意自同构 $\sigma$ 把 $\sqrt[3]{2}$ 送到 $\sqrt[3]{2}$、$\rho\sqrt[3]{2}$、$\rho^2\sqrt[3]{2}$ 之一，并把 $\rho$ 送到
\[
\rho\quad\text{或}\quad \rho^2=\frac{-1-\sqrt{-3}}{2},
\]
因为它们是分圆多项式 $\Phi_3(x)=x^2+x+1$ 的根。由于 $\sigma$ 完全由它在这两个元素上的作用决定，这只给出 $6$ 种可能，因而每一种可能事实上都是自同构。为了明确给出这些自同构，令 $\sigma$ 和 $\tau$ 为
\[
\sigma:\left\{
\begin{array}{l}
\sqrt[3]{2}\mapsto \rho\sqrt[3]{2},\\
\rho\mapsto \rho,
\end{array}\right.
\qquad
\tau:\left\{
\begin{array}{l}
\sqrt[3]{2}\mapsto \sqrt[3]{2},\\
\rho\mapsto \rho^2=-1-\rho.
\end{array}\right.
\]
与前面一样，这些映射可以在 $\mathbb Q(\sqrt[3]{2},\rho)$ 的元素上明确给出，而这些元素是基
\[
\{1,\sqrt[3]{2},(\sqrt[3]{2})^2,\rho,\rho\sqrt[3]{2},\rho(\sqrt[3]{2})^2\}
\]
的线性组合。例如
\[
\sigma(\rho\sqrt[3]{2})=(\rho)(\rho\sqrt[3]{2})=\rho^2\sqrt[3]{2}=(-1-\rho)\sqrt[3]{2}
=-\sqrt[3]{2}-\rho\sqrt[3]{2},
\]
并且可以类似地确定 $\sigma$ 在其他基元素上的作用。这给出
\[
\begin{aligned}
\sigma:&\ a+b\sqrt[3]{2}+c\sqrt[3]{4}+d\rho+e\rho\sqrt[3]{2}+f\rho\sqrt[3]{4} \\
&\longmapsto
 a-e\sqrt[3]{2}+(f-c)\sqrt[3]{4}+d\rho+(b-e)\rho\sqrt[3]{2}-c\rho\sqrt[3]{4}.
\end{aligned}
\tag{14.1}
\]
Galois 群的其他元素是
\[
1:\left\{
\begin{array}{l}
\sqrt[3]{2}\mapsto\sqrt[3]{2},\\
\rho\mapsto\rho,
\end{array}\right.
\qquad
\sigma^2:\left\{
\begin{array}{l}
\sqrt[3]{2}\mapsto\rho^2\sqrt[3]{2},\\
\rho\mapsto\rho,
\end{array}\right.
\]
计算 $\sigma\tau$ 得到
\[
\sigma\tau:\left\{
\begin{array}{l}
\sqrt[3]{2}\mapsto \sqrt[3]{2}\mapsto \rho\sqrt[3]{2},\\
\rho\mapsto \rho^2\mapsto \rho^2,
\end{array}\right.
\]
即
\[
\sigma\tau:\left\{
\begin{array}{l}
\sqrt[3]{2}\mapsto \rho\sqrt[3]{2},\\
\rho\mapsto \rho^2.
\end{array}\right.
\]
所以 $\sigma\tau=\tau\sigma^2$。类似地可算得 $\sigma^3=\tau^2=1$。因此
\[
\operatorname{Gal}(\mathbb Q(\sqrt[3]{2},\rho)/\mathbb Q)=\langle \sigma,\tau
angle\cong S_3
\]
是 $3$ 个字母上的对称群。或者（并且计算更少地）说，由于
\[
G=\operatorname{Gal}(\mathbb Q(\sqrt[3]{2},\rho)/\mathbb Q)
\]
作为 $x^3-2$ 的 $3$ 个根的置换来作用，$G$ 是 $S_3$ 的一个子群，所以既然它的阶为 $6$，就必须是 $S_3$。上面的计算明确地识别了 $G$ 中的自同构，并给出 $G$ 与 $S_3$ 的一个显式同构。

和前面的例子一样，我们可以确定 Galois 群任意子群的固定域。例如，考虑由 $\sigma$ 生成的子群 $\{1,\sigma,\sigma^2\}$ 的固定域。它们正是被 $\sigma$ 固定的元素（见等式 (14.1) 中的显式公式），因为如果一个元素被 $\sigma$ 固定，那么也被 $\sigma^2$ 固定。（一般地，某个子群的固定域就是被该子群的一组生成元固定的域。）被 $\sigma$ 固定的元素是满足
\[
a=a,\\qquad b=-e,
\qquad c=f-c,
\qquad d=d,
\qquad e=b-e,
\qquad f=-c
\]
的元素，这等价于 $b=c=f=e=0$。因此 $\{1,\sigma,\sigma^2\}$ 的固定域是 $\mathbb Q(\rho)$。

\paragraph{注记：} 这个例子说明，在由生成元的作用确定 Galois 群时必须谨慎。如上所述，并不是每个把 $\sqrt[3]{2}$ 和 $\rho\sqrt[3]{2}$ 送到 $x^3-2$ 的根的映射，都产生域的自同构（例如，映射
\[
\sqrt[3]{2}\mapsto \rho\sqrt[3]{2},\qquad
\rho\sqrt[3]{2}\mapsto \rho\sqrt[3]{2}
\]
显然不可能是自同构，因为它显然不是单射）。关键在于生成元之间可能存在（有时非常微妙的）代数关系，而这些关系必须被自同构保持。例如，这里的两个生成元之商是 $\rho$，它在上面的映射下被送到 $1$，而不是 $\rho$ 的极小多项式的根。换句话说，这些生成元之商满足一个二次方程，而该映射不保持这一性质。

另一个（不那么平凡的）例子，可与第 2 节中关于 $x^8-2$ 的分裂域的讨论比较。

\item 如例 3 中所述，域 $\mathbb Q(\sqrt[4]{2})$ 在 $\mathbb Q$ 上不是 Galois 的，因为任意自同构都由它把 $\sqrt[4]{2}$ 送到何处决定，而在四种可能 $\{\pm\sqrt[4]{2},\pm i\sqrt[4]{2}\}$ 中，只有两个是该域的元素（两个实根）。

注意我们有
\[
\begin{array}{ccccc}
\mathbb Q & \subset & \mathbb Q(\sqrt2) & \subset & \mathbb Q(\sqrt[4]{2})\\
& 2 && 2 &
\end{array}
\]
其中 $\mathbb Q(\sqrt2)/\mathbb Q$ 和 $\mathbb Q(\sqrt[4]{2})/\mathbb Q(\sqrt2)$ 都由例 2 是 Galois 扩张，因为二者都是二次扩张。这说明一个 Galois 扩张的 Galois 扩张未必是 Galois 的。

\item 推论 6 表明，命题 13.37 之后构造的有限域扩张 $\mathbb F_{p^n}/\mathbb F_p$ 是 Galois 的，因为 $\mathbb F_{p^n}$ 是可分多项式 $x^{p^n}-x$ 在 $\mathbb F_p$ 上的分裂域。由此可知，这个扩张的自同构群阶为 $n$。命题 13.35 中的单射同态
\[
\sigma:\mathbb F_{p^n}\longrightarrow \mathbb F_{p^n},\qquad \alpha\longmapsto \alpha^p
\]
在此情形中是满射，因为 $\mathbb F_{p^n}$ 是有限的，所以它是同构。这给出 $\mathbb F_{p^n}$ 的一个自同构，称为 \emph{Frobenius 自同构}，我们记为 $\sigma_p$。迭代 $\sigma_p$，有
\[
\sigma_p^2(\alpha)=\sigma_p(\sigma_p(\alpha))=(\alpha^p)^p=\alpha^{p^2}.
\]
类似地，
\[
\sigma_p^i(\alpha)=\alpha^{p^i},\qquad i=0,1,2,\ldots
\]
由于 $\alpha^{p^n}=\alpha$，我们看到 $\sigma_p^n=1$ 是恒等自同构。$\sigma_p$ 的更低正幂不可能为恒等，因为这将意味着对某个 $i<n$，所有 $\alpha\in\mathbb F_{p^n}$ 都满足 $\alpha^{p^i}=\alpha$，这是不可能的，因为这个方程只有 $p^i$ 个根。由此可知 $\sigma_p$ 在 Galois 群中阶为 $n$，也就是说 $\operatorname{Gal}(\mathbb F_{p^n}/\mathbb F_p)$ 是 $n$ 阶循环群，以 Frobenius 自同构 $\sigma_p$ 为生成元。

\item 第 13.5 节中考虑的不可分扩张 $\mathbb F_2(x)$ 在 $\mathbb F_2(t)$ 上（其中 $x^2-t=0$）不是 Galois 的。这个二次扩张的任意自同构由它在 $x$ 上的作用决定，而 $x$ 必须被送到方程 $X^2-t$ 的一个根。我们已经看到这个方程只有一个根（重数为 $2$），因为我们在特征 $2$ 的域中。因此这个扩张只有平凡自同构。注意，$\mathbb F_2(x)$ 是 $X^2-t$ 在 $\mathbb F_2(t)$ 上的分裂域，所以这个例子说明推论 6 中的可分性条件是必要的。
\end{enumerate}

\subsection*{练习}

\begin{enumerate}[label=\arabic*.]
\item
\begin{enumerate}[label=(\alph*)]
\item 证明：如果域 $K$ 在 $F$ 上由元素 $\alpha_1,\ldots,\alpha_n$ 生成，那么固定 $F$ 的 $K$ 的自同构 $\sigma$ 由 $\sigma(\alpha_1),\ldots,\sigma(\alpha_n)$ 唯一决定。特别地，证明一个自同构固定 $K$ 当且仅当它固定 $K$ 的一组生成元。
\item 设 $G\leq \operatorname{Gal}(K/F)$ 是扩张 $K/F$ 的 Galois 群的一个子群，并设 $\sigma_1,\ldots,\sigma_k$ 是 $G$ 的生成元。证明子域 $E/F$ 被 $G$ 固定当且仅当它被生成元 $\sigma_1,\ldots,\sigma_k$ 固定。
\end{enumerate}

\item 令 $\tau$ 为映射 $\tau:\mathbb C\to\mathbb C$，定义为 $\tau(a+bi)=a-bi$（复共轭）。证明 $\tau$ 是 $\mathbb C$ 的自同构。

\item 求复共轭在 $\mathbb C$ 上的固定域。

\item 证明 $\mathbb Q(\sqrt2)$ 与 $\mathbb Q(\sqrt3)$ 不同构。

\item 明确确定扩张 $\mathbb Q(\sqrt[4]{2})/\mathbb Q(\sqrt2)$ 的自同构。

\item 设 $k$ 是域。
\begin{enumerate}[label=(\alph*)]
\item 证明：对固定的 $a,b\in k$、$a\neq0$，由
\[
\varphi(f(t))=f(at+b)
\]
定义的映射 $\varphi:k[t]\to k[t]$ 是 $k[t]$ 的一个自同构，并且在 $k$ 上为恒等。
\item 反过来，设 $\varphi$ 是 $k[t]$ 的一个自同构，并且在 $k$ 上为恒等。证明存在 $a,b\in k$ 且 $a\neq0$，使得 $\varphi(f(t))=f(at+b)$，如 (a) 中所述。
\end{enumerate}

\item 本练习确定 $\operatorname{Aut}(\mathbb R/\mathbb Q)$。
\begin{enumerate}[label=(\alph*)]
\item 证明任意 $\sigma\in\operatorname{Aut}(\mathbb R/\mathbb Q)$ 把平方送到平方，并把正实数送到正实数。推出对任意 $a,b\in\mathbb R$，若 $a<b$，则 $\sigma a<\sigma b$。
\item 证明若
\[
-\frac1m<a-b<\frac1m
\]
则
\[
-\frac1m<\sigma a-\sigma b<\frac1m
\]
对每个正整数 $m$ 成立。推出 $\sigma$ 是 $\mathbb R$ 上的连续映射。
\item 证明任何在 $\mathbb Q$ 上为恒等的 $\mathbb R$ 上连续映射都是恒等映射，因此 $\operatorname{Aut}(\mathbb R/\mathbb Q)=1$。
\end{enumerate}

\item 证明固定 $k$ 的有理函数域 $k(t)$ 的自同构，恰好是由
\[
t\longmapsto \frac{at+b}{ct+d}
\]
确定的分式线性变换，其中 $a,b,c,d\in k$ 且 $ad-bc\neq0$（于是 $f(t)\in k(t)$ 被送到
\[
f\left(\frac{at+b}{ct+d}\right)
\]
）（参见第 13.2 节练习 18）。

\item 求 $k(t)$ 的自同构 $t\mapsto t+1$ 的固定域。

\item 设 $K$ 是域 $F$ 的一个扩张。令 $\varphi:K\to K'$ 是 $K$ 与域 $K'$ 的一个同构，它把 $F$ 映到 $K'$ 的子域 $F'$。证明映射
\[
\sigma\longmapsto \varphi\sigma\varphi^{-1}
\]
定义一个群同构
\[
\operatorname{Aut}(K/F)\cong \operatorname{Aut}(K'/F').
\]
\end{enumerate}

\section{Galois 理论基本定理}

在上一节考虑的 Galois 扩张 $\operatorname{Gal}(\mathbb Q(\sqrt2,\sqrt3)/\mathbb Q)$ 中，Galois 群的子群图与相应固定域的图之间有很强的相似性：
\[
\begin{tikzpicture}[scale=0.95, baseline=(current bounding box.center), every node/.style={inner sep=2pt}]
\node (top) at (0,2) {$\{1\}$};
\node (l) at (-2,0) {$\{1,\tau\}$};
\node (m) at (0,0) {$\{1,\sigma\tau\}$};
\node (r) at (2,0) {$\{1,\sigma\}$};
\node (bot) at (0,-2) {$\{1,\sigma,\tau,\sigma\tau\}$};
\draw (top)-- node[above left] {$2$} (l);
\draw (top)-- node[right] {$2$} (m);
\draw (top)-- node[above right] {$2$} (r);
\draw (l)-- node[below left] {$2$} (bot);
\draw (m)-- node[right] {$2$} (bot);
\draw (r)-- node[below right] {$2$} (bot);
\end{tikzpicture}
\]
以及相应固定域的图
\[
\begin{tikzpicture}[scale=0.95, baseline=(current bounding box.center), every node/.style={inner sep=2pt}]
\node (top) at (0,2) {$\mathbb Q(\sqrt2,\sqrt3)$};
\node (l) at (-2,0) {$\mathbb Q(\sqrt2)$};
\node (m) at (0,0) {$\mathbb Q(\sqrt6)$};
\node (r) at (2,0) {$\mathbb Q(\sqrt3)$};
\node (bot) at (0,-2) {$\mathbb Q$};
\draw (top)-- node[above left] {$2$} (l);
\draw (top)-- node[right] {$2$} (m);
\draw (top)-- node[above right] {$2$} (r);
\draw (l)-- node[below left] {$2$} (bot);
\draw (m)-- node[right] {$2$} (bot);
\draw (r)-- node[below right] {$2$} (bot);
\end{tikzpicture}
\]
（由于这种对应的反包含性质，我们已经把子群格倒置过来）。注意，这也是该扩张的所有已知子域的图，并且在这个情形中，每个子域本身也是 $\mathbb Q$ 上的 Galois 扩张。

类似地，Galois 群的子群图
\[
\begin{tikzpicture}[scale=0.95, baseline=(current bounding box.center), every node/.style={inner sep=2pt}]
\node (top) at (0,2.1) {$1$};
\node (a) at (-1.8,0.2) {$\langle\sigma\rangle$};
\node (t) at (-0.2,0.8) {$\langle\tau\rangle$};
\node (ta) at (1.1,0.8) {$\langle\tau\sigma\rangle$};
\node (ta2) at (2.4,0.8) {$\langle\tau\sigma^2\rangle$};
\node (bot) at (0,-1.4) {$\langle\sigma,\tau\rangle$};
\draw (top)--(a);
\draw (top)--(t);
\draw (top)--(ta);
\draw (top)--(ta2);
\draw (a)--(bot);
\draw (t)--(bot);
\draw (ta)--(bot);
\draw (ta2)--(bot);
\end{tikzpicture}
\]
与 $x^3-2$ 的分裂域的已知子域图也有很强的相似性：
\[
\begin{tikzpicture}[scale=0.95, baseline=(current bounding box.center), every node/.style={inner sep=2pt}]
\node (top) at (0,2.1) {$\mathbb Q(\sqrt[3]{2},\rho)$};
\node (rho) at (-1.9,0.1) {$\mathbb Q(\rho)$};
\node (a) at (-0.25,0.75) {$\mathbb Q(\sqrt[3]{2})$};
\node (b) at (1.35,0.75) {$\mathbb Q(\rho\sqrt[3]{2})$};
\node (c) at (3.1,0.75) {$\mathbb Q(\rho^2\sqrt[3]{2})$};
\node (bot) at (0,-1.4) {$\mathbb Q$};
\draw (top)--(rho);
\draw (top)--(a);
\draw (top)--(b);
\draw (top)--(c);
\draw (rho)--(bot);
\draw (a)--(bot);
\draw (b)--(bot);
\draw (c)--(bot);
\end{tikzpicture}
\]
第二个图中的子域正是第一个图中子群的固定域。注意，在这对图中，只有由 $\sigma$ 生成的子群 $\langle\sigma\rangle$ 是 $S_3$ 中的正规子群，而且子域 $\mathbb Q(\rho)$ 是唯一在 $\mathbb Q$ 上 Galois 的子域。

Galois 理论基本定理断言，在上面两个例子中观察到的关系不是偶然的，并且对任意 Galois 扩张都成立。在证明它之前，我们先发展一些关于群的特征的预备结果；域自同构将给出这种特征的特殊例子。

\noindent\textbf{定义.} 群 $G$ 的一个取值于域 $L$ 中的\emph{特征}\footnote{这是线性特征的定义。更一般的特征将在第 18 章中研究。} $\chi$ 是从 $G$ 到 $L$ 的乘法群的同态：
\[
\chi:G\longrightarrow L^\times,
\]
即对所有 $g_1,g_2\in G$，有 $\chi(g_1g_2)=\chi(g_1)\chi(g_2)$，并且对所有 $g\in G$，$\chi(g)$ 是 $L$ 的一个非零元素。

\noindent\textbf{定义.} $G$ 的特征 $\chi_1,\chi_2,\ldots,\chi_n$ 称为在 $L$ 上\emph{线性无关}，如果它们作为 $G$ 上的函数线性无关；也就是说，不存在作为函数在 $G$ 上成立的非平凡关系
\begin{equation}
 a_1\chi_1+a_2\chi_2+\cdots+a_n\chi_n=0
 \qquad(a_1,\ldots,a_n\in L\text{ 不全为 }0),
\tag{14.2}
\end{equation}
即对所有 $g\in G$，都有 $a_1\chi_1(g)+a_2\chi_2(g)+\cdots+a_n\chi_n(g)=0$。

\noindent\textbf{定理 7（特征的线性无关性）.} 如果 $\chi_1,\chi_2,\ldots,\chi_n$ 是 $G$ 的互异特征，取值于 $L$，那么它们在 $L$ 上线性无关。

\begin{proof}
设这些特征线性相关。在所有如 (14.2) 的线性相关关系中，选取一个非零系数个数 $m$ 最少的关系。必要时重新编号，可设这 $m$ 个非零系数为 $a_1,a_2,\ldots,a_m$：
\[
 a_1\chi_1+a_2\chi_2+\cdots+a_m\chi_m=0.
\]
于是对任意 $g\in G$，有
\begin{equation}
 a_1\chi_1(g)+a_2\chi_2(g)+\cdots+a_m\chi_m(g)=0.
\tag{14.3}
\end{equation}
取 $g_0\in G$ 使得 $\chi_1(g_0)\neq\chi_m(g_0)$（这样的 $g_0$ 存在，因为 $\chi_1\neq\chi_m$）。由于 (14.3) 对 $G$ 的每个元素成立，特别地有
\[
 a_1\chi_1(g_0g)+a_2\chi_2(g_0g)+\cdots+a_m\chi_m(g_0g)=0,
\]
即
\begin{equation}
 a_1\chi_1(g_0)\chi_1(g)+a_2\chi_2(g_0)\chi_2(g)+\cdots+a_m\chi_m(g_0)\chi_m(g)=0.
\tag{14.4}
\end{equation}
把方程 (14.3) 乘以 $\chi_m(g_0)$，再从方程 (14.4) 中相减，得到
\[
\begin{aligned}
&[\chi_m(g_0)-\chi_1(g_0)]a_1\chi_1(g)+[\chi_m(g_0)-\chi_2(g_0)]a_2\chi_2(g)+\cdots\\
&\qquad +[\chi_m(g_0)-\chi_{m-1}(g_0)]a_{m-1}\chi_{m-1}(g)=0,
\end{aligned}
\]
这对所有 $g\in G$ 成立。但是第一个系数非零，于是这给出了一个非零系数更少的关系，矛盾。
\end{proof}

现在考虑域 $K$ 到域 $L$ 中的一个单同态 $\sigma$，称为 $K$ 到 $L$ 中的一个\emph{嵌入}。特别地，$\sigma$ 是乘法群 $G=K^\times$ 到乘法群 $L^\times$ 的同态，因此 $\sigma$ 可以看作 $K^\times$ 上取值于 $L$ 的一个特征。还要注意，这个特征包含了把 $\sigma$ 仅仅看作 $K$ 上函数时所有有用的取值信息，因为 $K^\times$ 中没有考虑的 $K$ 的唯一点是 $0$，而我们知道 $\sigma$ 把 $0$ 送到 $0$。

\noindent\textbf{推论 8.} 如果 $\sigma_1,\sigma_2,\ldots,\sigma_n$ 是域 $K$ 到域 $L$ 中的互异嵌入，那么它们作为 $K$ 上的函数线性无关。特别地，域 $K$ 的互异自同构作为 $K$ 上的函数线性无关。

下面我们用推论 8 来证明域 $K$ 的自同构群的子群阶数与其固定域上的扩张次数之间的基本关系。

\noindent\textbf{定理 9.} 设 $G=\{\sigma_1=1,\sigma_2,\ldots,\sigma_n\}$ 是域 $K$ 的自同构群的一个子群，令 $F$ 为固定域。那么
\[
[K:F]=n=|G|.
\]

\begin{proof}
先设 $n>[K:F]$，并令 $\omega_1,\omega_2,\ldots,\omega_m$ 是 $K$ 在 $F$ 上的一组基（$m=[K:F]$）。则方程组
\[
\begin{array}{rcl}
\sigma_1(\omega_1)x_1+\sigma_2(\omega_1)x_2+\cdots+\sigma_n(\omega_1)x_n&=&0,\\
\vdots&&\\
\sigma_1(\omega_m)x_1+\sigma_2(\omega_m)x_2+\cdots+\sigma_n(\omega_m)x_n&=&0
\end{array}
\]
是 $m$ 个方程、$n$ 个未知数 $x_1,x_2,\ldots,x_n$ 的方程组。由于未知数多于方程数，它在 $K$ 中有非平凡解 $\beta_1,\beta_2,\ldots,\beta_n$。

令 $a_1,a_2,\ldots,a_m$ 是 $F$ 中任意 $m$ 个元素。按照定义，$F$ 被 $\sigma_1,\ldots,\sigma_n$ 固定，所以这些元素中的每一个都被每个 $\sigma_i$ 固定，即 $\sigma_i(a_j)=a_j$，$i=1,2,\ldots,n$，$j=1,2,\ldots,m$。将上面的第一个方程乘以 $a_1$，第二个方程乘以 $a_2$，依次直到最后一个方程乘以 $a_m$，得到
\[
\begin{array}{rcl}
\sigma_1(a_1\omega_1)\beta_1+\sigma_2(a_1\omega_1)\beta_2+\cdots+\sigma_n(a_1\omega_1)\beta_n&=&0,\\
\vdots&&\\
\sigma_1(a_m\omega_m)\beta_1+\sigma_2(a_m\omega_m)\beta_2+\cdots+\sigma_n(a_m\omega_m)\beta_n&=&0.
\end{array}
\]
把这些方程相加，就看到存在 $K$ 中不全为零的元素 $\beta_1,\ldots,\beta_n$，使得
\[
\sigma_1(a_1\omega_1+\cdots+a_m\omega_m)\beta_1+\cdots+
\sigma_n(a_1\omega_1+\cdots+a_m\omega_m)\beta_n=0
\]
对 $F$ 中所有 $a_1,\ldots,a_m$ 都成立。由于 $\omega_1,\ldots,\omega_m$ 是 $K$ 的 $F$-基，$K$ 的任意元素 $a$ 都可写成 $a_1\omega_1+\cdots+a_m\omega_m$，所以上式意味着
\[
\sigma_1(a)\beta_1+\cdots+\sigma_n(a)\beta_n=0
\]
对所有 $a\in K$ 成立。这就是说互异自同构 $\sigma_1,\ldots,\sigma_n$ 在 $K$ 上线性相关，与推论 8 矛盾。

我们已经证明 $n\leq [K:F]$。注意，到目前为止我们还没有用到 $\sigma_1,\sigma_2,\ldots,\sigma_n$ 是一个群的元素这一事实。

现在设 $n<[K:F]$。那么 $K$ 中存在多于 $n$ 个在 $F$ 上线性无关的元素，设为 $\alpha_1,
\ldots,\alpha_{n+1}$。方程组
\begin{equation}
\sigma_i(\alpha_1)x_1+\sigma_i(\alpha_2)x_2+\cdots+\sigma_i(\alpha_{n+1})x_{n+1}=0,
\qquad i=1,2,\ldots,n,
\tag{14.5}
\end{equation}
有 $n$ 个方程和 $n+1$ 个未知数 $x_1,
\ldots,x_{n+1}$，因此在 $K$ 中有一个解 $\beta_1,
\ldots,\beta_{n+1}$，其中并非所有 $\beta_i$ 都为 $0$。如果解中的所有元素 $\beta_1,
\ldots,\beta_{n+1}$ 都属于 $F$，那么第一个方程（回忆 $\sigma_1=1$ 是恒等自同构）就与 $\alpha_1,
\alpha_2,
\ldots,
\alpha_{n+1}$ 在 $F$ 上线性无关矛盾。因此至少有一个 $\beta_i$ 不属于 $F$。

在 (14.5) 的所有非平凡解 $(\beta_1,
\ldots,
\beta_{n+1})$ 中，选取非零 $\beta_i$ 个数 $r$ 最少的一个。必要时重新编号，可设 $\beta_1,
\ldots,
\beta_r$ 非零。把方程除以 $\beta_r$，还可以假设 $\beta_r=1$。我们已经看到 $\beta_1,
\ldots,
\beta_{r-1},1$ 中至少有一个不属于 $F$（这特别说明 $r>1$），设 $\beta_1\notin F$。于是方程组写作
\begin{equation}
\sigma_1(\alpha_1)\beta_1+\cdots+\sigma_1(\alpha_{r-1})\beta_{r-1}+\sigma_1(\alpha_r)=0,
\tag{14.6}
\end{equation}
或更简略地写为
\begin{equation}
\sigma_i(\alpha_1)\beta_1+\cdots+\sigma_i(\alpha_{r-1})\beta_{r-1}+\sigma_i(\alpha_r)=0,
\qquad i=1,2,\ldots,n.
\tag{14.7}
\end{equation}
由于 $\beta_1\notin F$，存在一个自同构 $\sigma_{k_0}$（$k_0\in\{1,2,
\ldots,n\}$），使得 $\sigma_{k_0}\beta_1\neq\beta_1$。若对 (14.6) 中的方程施以自同构 $\sigma_{k_0}$，得到方程组
\begin{equation}
\sigma_{k_0}\sigma_j(\alpha_1)\sigma_{k_0}(\beta_1)+\cdots+
\sigma_{k_0}\sigma_j(\alpha_{r-1})\sigma_{k_0}(\beta_{r-1})+
\sigma_{k_0}\sigma_j(\alpha_r)=0,
\tag{14.8}
\end{equation}
其中 $j=1,2,
\ldots,n$。但是元素
\[
\sigma_{k_0}\sigma_1,
\sigma_{k_0}\sigma_2,
\ldots,
\sigma_{k_0}\sigma_n
\]
与元素
\[
\sigma_1,
\sigma_2,
\ldots,
\sigma_n
\]
相同，只是顺序可能不同，因为这些元素组成一个群。换句话说，若由 $\sigma_{k_0}\sigma_j=\sigma_i$ 定义指标 $i$，则 $i$ 和 $j$ 都遍历集合 $\{1,2,
\ldots,n\}$。于是 (14.8) 中的方程可写为
\begin{equation}
\sigma_i(\alpha_1)\sigma_{k_0}(\beta_1)+\cdots+
\sigma_i(\alpha_{r-1})\sigma_{k_0}(\beta_{r-1})+
\sigma_i(\alpha_r)=0.
\tag{14.8'}
\end{equation}
现在从 (14.7) 中的方程减去 (14.8') 中的方程，得到方程组
\[
\sigma_i(\alpha_1)(\beta_1-\sigma_{k_0}\beta_1)+\cdots+
\sigma_i(\alpha_{r-1})(\beta_{r-1}-\sigma_{k_0}\beta_{r-1})=0
\]
对 $i=1,2,
\ldots,n$ 成立。但这是方程组 (14.5) 的一个解，并且
\[
 x_1=\beta_1-\sigma_{k_0}(\beta_1)\neq0
\]
（由 $k_0$ 的选择），因此它是非平凡的，同时非零 $x_i$ 的个数少于 $r$。这给出矛盾，从而证明完成。
\end{proof}

这个结果的第一个应用，是证明命题 5 中的不等式对任意有限扩张 $K/F$ 都成立。

\noindent\textbf{推论 10.} 设 $K/F$ 是任意有限扩张。则
\[
|\operatorname{Aut}(K/F)|\leq [K:F],
\]
且等号成立当且仅当 $F$ 是 $\operatorname{Aut}(K/F)$ 的固定域。换言之，$K/F$ 是 Galois 扩张，当且仅当 $F$ 是 $\operatorname{Aut}(K/F)$ 的固定域。

\begin{proof}
令 $F_1$ 为 $\operatorname{Aut}(K/F)$ 的固定域，于是
\[
F\subseteq F_1\subseteq K.
\]
由定理 9，$[K:F_1]=|\operatorname{Aut}(K/F)|$。因此
\[
[K:F]=|\operatorname{Aut}(K/F)|[F_1:F],
\]
这证明了推论。
\end{proof}

\noindent\textbf{推论 11.} 设 $G$ 是域 $K$ 的自同构群的一个有限子群，令 $F$ 为它的固定域。那么固定 $F$ 的 $K$ 的每个自同构都属于 $G$，即
\[
\operatorname{Aut}(K/F)=G,
\]
从而 $K/F$ 是 Galois 扩张，并且 Galois 群为 $G$。

\begin{proof}
按照定义，$F$ 被 $G$ 的所有元素固定，所以 $G\leq\operatorname{Aut}(K/F)$（问题在于是否还有不在 $G$ 中的固定 $F$ 的 $K$ 的自同构，即这个包含是否为真包含）。因此 $|G|\leq|\operatorname{Aut}(K/F)|$。由定理 9，$|G|=[K:F]$，由前一推论，$|\operatorname{Aut}(K/F)|\leq [K:F]$。这给出
\[
[K:F]=|G|\leq |\operatorname{Aut}(K/F)|\leq [K:F],
\]
所以处处取等号，推论得证。
\end{proof}

\noindent\textbf{推论 12.} 如果 $G_1\neq G_2$ 是域 $K$ 的自同构群的两个互异有限子群，那么它们的固定域也互异。

\begin{proof}
设 $F_1$ 为 $G_1$ 的固定域，$F_2$ 为 $G_2$ 的固定域。若 $F_1=F_2$，则按照定义 $F_1$ 被 $G_2$ 固定。由前一推论，任意固定 $F_1$ 的自同构都属于 $G_1$，所以 $G_2\leq G_1$。同理 $G_1\leq G_2$，因此 $G_1=G_2$。
\end{proof}

由以上推论可见，对 $\operatorname{Aut}(K)$ 的互异有限子群取固定域，给出 $K$ 的互异子域，而且 $K$ 在这些子域上都是 Galois 的。此外，扩张次数由子群的阶给出。我们已经在上面对域 $K=\mathbb Q(\sqrt2,\sqrt3)$ 与 $K=\mathbb Q(\sqrt[3]{2},\rho)$ 明确看到了这一点。基本定理的一部分断言，这些就是 $K$ 的所有子域。

下一个结果给出命题 5 的逆命题，并刻画 Galois 扩张。

\noindent\textbf{定理 13.} 扩张 $K/F$ 是 Galois 的，当且仅当 $K$ 是 $F$ 上某个可分多项式的分裂域。进一步地，如果这是成立的，那么系数在 $F$ 中且在 $K$ 中有一个根的每个不可约多项式都是可分的，并且它的所有根都在 $K$ 中（因此特别地，$K/F$ 是可分扩张）。

\begin{proof}
命题 5 已经证明了一个可分多项式的分裂域是 Galois 的。

现在证明如果 $K/F$ 是 Galois，则 $F[x]$ 中每个在 $K$ 中有根的不可约多项式 $p(x)$ 都在 $K$ 中完全分裂。令 $G=\operatorname{Gal}(K/F)$。设 $\alpha\in K$ 是 $p(x)$ 的一个根，并考虑元素
\begin{equation}
\alpha,\sigma_2(\alpha),\ldots,\sigma_n(\alpha)\in K,
\tag{14.9}
\end{equation}
其中 $\{1,\sigma_2,\ldots,\sigma_n\}$ 是 $\operatorname{Gal}(K/F)$ 的元素。令
\[
\alpha,\alpha_2,\ldots,\alpha_r
\]
表示 (14.9) 中的互异元素。若 $\tau\in G$，则因为 $G$ 是群，元素
\[
\{\tau,\tau\sigma_2,\ldots,\tau\sigma_n\}
\]
与元素 $\{1,\sigma_2,\ldots,\sigma_n\}$ 相同，只是顺序可能不同。因此，把 $\tau\in G$ 作用到 (14.9) 中的元素上只是置换这些元素，特别地，把 $\tau$ 作用到 $\alpha,\alpha_2,\alpha_3,
\ldots,
\alpha_r$ 上也只是置换这些元素。于是多项式
\[
f(x)=(x-\alpha)(x-\alpha_2)\cdots(x-\alpha_r)
\]
的系数被 $G$ 的所有元素固定，因为 $G$ 的元素只是置换这些因子。因此这些系数属于 $G$ 的固定域；由推论 10，这个固定域就是 $F$。于是 $f(x)\in F[x]$。

由于 $p(x)$ 不可约并且以 $\alpha$ 为根，$p(x)$ 是 $\alpha$ 在 $F$ 上的极小多项式，因此它整除任何一个系数在 $F$ 中且以 $\alpha$ 为根的多项式（这是命题 13.9）。于是 $p(x)$ 在 $F[x]$ 中整除 $f(x)$；另一方面，由命题 2，$f(x)$ 显然在 $K[x]$ 中整除 $p(x)$。所以
\[
p(x)=f(x).
\]
特别地，这说明 $p(x)$ 是可分的，并且它的所有根都属于 $K$（事实上它们就在元素 $\alpha,
\sigma_2\alpha,
\ldots,
\sigma_n\alpha$ 之中），这证明了定理的最后一个陈述。

为了完成证明，设 $K/F$ 是 Galois 的，令 $\omega_1,\omega_2,
\ldots,
\omega_n$ 是 $K/F$ 的一组基。设 $p_i(x)$ 为 $\omega_i$ 在 $F$ 上的极小多项式，$i=1,2,
\ldots,n$。由刚刚证明的结果，$p_i(x)$ 是可分的，并且它的所有根都在 $K$ 中。令 $g(x)$ 为乘积 $p_1(x)\cdots p_n(x)$ 去掉所有重因子后得到的多项式（“无平方部分”）。那么这两个多项式的分裂域相同，并且这个分裂域就是 $K$（所有根都在 $K$ 中，所以 $K$ 含有分裂域；而 $\omega_1,
\omega_2,
\ldots,
\omega_n$ 都在这些根之中，所以分裂域含有 $K$）。因此 $K$ 是可分多项式 $g(x)$ 的分裂域。
\end{proof}

\noindent\textbf{定义.} 设 $K/F$ 是 Galois 扩张。若 $\alpha\in K$，则 $\sigma\alpha$（其中 $\sigma\in\operatorname{Gal}(K/F)$）称为 $\alpha$ 在 $F$ 上的\emph{共轭}（或 Galois 共轭）。如果 $E$ 是 $K$ 的一个含 $F$ 的子域，则域 $\sigma(E)$ 称为 $E$ 在 $F$ 上的\emph{共轭域}。

定理的证明表明，在 Galois 扩张 $K/F$ 中，任意元素 $\alpha\in K$ 在 $F$ 上的极小多项式的其他根，恰好是 $\alpha$ 在 $K/F$ 的 Galois 群作用下的互异共轭。

这个定理的第二个陈述还说明，如果我们能够找到哪怕一个 $F$ 上不可约多项式，它在 $K$ 中有一个根但不是所有根都在 $K$ 中，那么 $K$ 就不是 $F$ 上的 Galois 扩张。这以很强的意义证明了先前例子中关于 Galois 扩张的直觉：Galois 扩张是具有“足够多”不可约多项式根的扩张（即，如果它含有一个根，就含有所有根）。

最后注意，我们现在已经有了 Galois 扩张 $K/F$ 的四种刻画：
\begin{enumerate}[label=(\arabic*)]
\item $F$ 上可分多项式的分裂域；
\item $F$ 恰好是被 $\operatorname{Aut}(K/F)$ 固定的元素集合的域（一般而言，固定域可能大于 $F$）；
\item 满足 $[K:F]=|\operatorname{Aut}(K/F)|$ 的域（原始定义）；
\item 有限、正规且可分的扩张。
\end{enumerate}

\noindent\textbf{定理 14（Galois 理论基本定理）.} 设 $K/F$ 是 Galois 扩张，并令 $G=\operatorname{Gal}(K/F)$。那么存在一个双射
\[
\left\{\begin{array}{c}
\text{含 }F\text{ 的 }K\text{ 的}\text{子域 }E
\end{array}\right\}
\longleftrightarrow
\left\{\begin{array}{c}
G\text{ 的}\text{子群 }H
\end{array}\right\}
\]
由对应
\[
E\longmapsto \{G\text{ 中固定 }E\text{ 的元素}\},
\qquad
H\longmapsto H\text{ 的固定域}
\]
给出，并且这两个对应互为逆。在此对应下：
\begin{enumerate}[label=(\arabic*)]
\item （反包含）若 $E_1,E_2$ 分别对应 $H_1,H_2$，则 $E_1\subseteq E_2$ 当且仅当 $H_2\subseteq H_1$。
\item $[K:E]=|H|$，并且 $[E:F]=|G:H|$，即 $H$ 在 $G$ 中的指数：
\[
\begin{array}{c}
K\\[-2mm]
|\,|H|\\[-1mm]
E\\[-2mm]
|\,|G:H|\\[-1mm]
F.
\end{array}
\]
\item $K/E$ 总是 Galois 的，其 Galois 群为
\[
\operatorname{Gal}(K/E)=H.
\]
\item $E$ 在 $F$ 上是 Galois 的，当且仅当 $H$ 是 $G$ 的正规子群。若是这种情形，则 Galois 群同构于商群
\[
\operatorname{Gal}(E/F)\cong G/H.
\]
更一般地，即使 $H$ 未必在 $G$ 中正规，固定 $F$ 的 $E$ 的同构（到含有 $K$ 的一个固定代数闭包中）与 $G$ 中陪集 $\{\sigma H\}$ 一一对应。
\item 若 $E_1,E_2$ 分别对应 $H_1,H_2$，则交 $E_1\cap E_2$ 对应由 $H_1,H_2$ 生成的群 $\langle H_1,H_2\rangle$，而复合域 $E_1E_2$ 对应交 $H_1\cap H_2$。因此，含 $F$ 的 $K$ 的子域格与 $G$ 的子群格是“对偶的”（一个格图就是另一个格图上下倒置）。
\end{enumerate}

\begin{proof}
给定 $G$ 的任意子群 $H$，由推论 12，我们得到唯一的固定域 $E=K^H$。这说明上述从右到左的对应是单射。

如果 $K$ 是可分多项式 $f(x)\in F[x]$ 的分裂域，那么对 $K$ 的任意含 $F$ 的子域 $E$，也可以把 $f(x)$ 看作 $E[x]$ 中的元素。于是 $K$ 也是 $f(x)$ 在 $E$ 上的分裂域，所以扩张 $K/E$ 是 Galois 的。由推论 10，$E$ 是 $\operatorname{Aut}(K/E)\leq G$ 的固定域，这说明含 $F$ 的 $K$ 的每个子域都作为 $G$ 的某个子群的固定域出现。因此上述从右到左的对应是满射，从而是双射。这两个对应互为逆，因为由推论 10，固定 $E$ 的自同构正是 $\operatorname{Aut}(K/E)$。

我们已经在命题 4 中看到 Galois 对应是反包含的，这给出 (1)。

若 $E=K^H$ 是 $H$ 的固定域，则定理 9 给出 $[K:E]=|H|$，并且 $[K:F]=|G|$。两式相除得到 $[E:F]=|G:H|$，这证明 (2)。

推论 11 立即给出 (3)。

设 $E=K^H$ 是子群 $H$ 的固定域。每个 $\sigma\in G=\operatorname{Gal}(K/F)$ 限制到 $E$ 上，给出 $E$ 与 $K$ 的子域 $\sigma(E)$ 之间的一个嵌入 $\sigma|_E$。反过来，令
\[
\tau:E\xrightarrow{\sim}\tau(E)\subseteq \overline F
\]
是 $E$（到包含 $K$ 的一个固定代数闭包 $\overline F$ 中）的任意固定 $F$ 的嵌入。那么事实上 $\tau(E)$ 包含在 $K$ 中：若 $\alpha\in E$ 的 $F$ 上极小多项式为 $m_\alpha(x)$，则 $\tau(\alpha)$ 是 $m_\alpha(x)$ 的另一个根，而由定理 13，$K$ 包含所有这些根。如上，$K$ 是 $f(x)$ 在 $E$ 上的分裂域，因而也是 $\tau f(x)$（这与 $f(x)$ 相同，因为 $f(x)$ 的系数在 $F$ 中）在 $\tau(E)$ 上的分裂域。定理 13.27 关于延拓同构的结果于是说明我们可以把 $\tau$ 延拓为一个同构 $\sigma$：
\[
\begin{array}{ccc}
K&\xrightarrow{\sigma}&K\\
\cup&&\cup\\[-1mm]
E&\xrightarrow{\tau}&\tau(E).
\end{array}
\]
由于 $\tau$ 固定 $F$，所以 $\sigma$ 也固定 $F$。因此 $E$ 的每个固定 $F$ 的嵌入都是某个固定 $F$ 的 $K$ 的自同构 $\sigma$ 限制到 $E$ 上所得，换句话说，每个这样的嵌入都具有 $\sigma|_E$ 的形式，其中 $\sigma\in G$。

两个自同构 $\sigma,
\sigma'\in G$ 限制到 $E$ 上给出同一个嵌入，当且仅当 $\sigma^{-1}\sigma'$ 在 $E$ 上是恒等映射。但这时 $\sigma^{-1}\sigma'\in H$（即 $\sigma'\in\sigma H$），因为由 (3)，固定 $E$ 的 $K$ 的自同构正是 $H$ 中的元素。因此 $E$ 的互异嵌入与 $G$ 中的陪集 $\sigma H$ 双射。特别地，
\[
|\operatorname{Emb}(E/F)|=[G:H]=[E:F],
\]
其中 $\operatorname{Emb}(E/F)$ 表示固定 $F$ 的 $E$ 的嵌入（到一个固定代数闭包中）的集合。注意 $\operatorname{Emb}(E/F)$ 包含自同构集合 $\operatorname{Aut}(E/F)$。

扩张 $E/F$ 是 Galois 的，当且仅当 $|\operatorname{Aut}(E/F)|=[E:F]$。由上式，这当且仅当 $E$ 的每个嵌入实际上都是 $E$ 的自同构，即当且仅当对每个 $\sigma\in G$ 都有 $\sigma(E)=E$。

若 $\sigma\in G$，则固定域 $\sigma(E)$ 的 $G$ 的子群为 $\sigma H\sigma^{-1}$，即
\[
\sigma(E)=K^{\sigma H\sigma^{-1}}.
\]
为了看出这一点，注意若 $\sigma a\in\sigma(E)$，则
\[
(\sigma h\sigma^{-1})(\sigma a)=\sigma(ha)=\sigma a,
\qquad h\in H,
\]
因为 $h$ 固定 $a\in E$；这说明 $\sigma H\sigma^{-1}$ 固定 $\sigma(E)$。固定 $\sigma(E)$ 的群的阶等于 $K$ 在 $\sigma(E)$ 上的次数。这个次数等于 $K$ 在 $E$ 上的次数，因为这两个域同构；因此等于 $H$ 的阶。于是 $\sigma H\sigma^{-1}$ 正是固定 $\sigma(E)$ 的群，因为我们已经证明了包含关系并且二者阶数相同。

由于已经证明的 Galois 对应是双射，包含 $F$ 的 $K$ 的两个子域相等当且仅当它们在 $G$ 中的固定子群相等。因此，对所有 $\sigma\in G$ 都有 $\sigma(E)=E$，当且仅当对所有 $\sigma\in G$ 都有 $\sigma H\sigma^{-1}=H$，换句话说，$E$ 在 $F$ 上 Galois 当且仅当 $H$ 是 $G$ 的正规子群。

我们已经把 $E$ 在 $F$ 上的嵌入识别为 $G$ 中 $H$ 的陪集集合，并且当 $H$ 在 $G$ 中正规时看到这些嵌入都是自同构。由自同构复合所定义的群运算可知，在这种情形中，陪集群 $G/H$ 被识别为 Galois 扩张 $E/F$ 的自同构群。因此，当 $H$ 正规于 $G$ 时，$G/H\cong\operatorname{Gal}(E/F)$，这完成了 (4) 的证明。

设 $H_1$ 是固定子域 $E_1$ 的 $G$ 中元素组成的子群，$H_2$ 是固定子域 $E_2$ 的 $G$ 中元素组成的子群。$H_1\cap H_2$ 中任意元素同时固定 $E_1$ 和 $E_2$，因而固定复合域 $E_1E_2$ 中的每个元素，因为这个域中的元素由 $E_1$ 与 $E_2$ 中元素经代数运算得到。反过来，如果自同构 $\sigma$ 固定复合域 $E_1E_2$，则特别地 $\sigma$ 固定 $E_1$，即 $\sigma\in H_1$；并且 $\sigma$ 固定 $E_2$，即 $\sigma\in H_2$，所以 $\sigma\in H_1\cap H_2$。这证明复合域 $E_1E_2$ 对应交 $H_1\cap H_2$。类似地，交 $E_1\cap E_2$ 对应由 $H_1$ 与 $H_2$ 生成的群 $\langle H_1,H_2\rangle$。定理得证。
\end{proof}

\paragraph{例（$\mathbb Q(\sqrt2,\sqrt3)$ 与 $\mathbb Q(\sqrt[3]{2},\rho)$）}
我们已经在本节开头看到了这个定理的例子。现在我们看到，先前给出的域 $\mathbb Q(\sqrt2,\sqrt3)$ 与 $\mathbb Q(\sqrt[3]{2},\rho)$ 的子域图，确实列出了这两个域的所有子域。

由于 Klein 四元群的每个子群都是正规子群，$\mathbb Q(\sqrt2,\sqrt3)$ 的所有子域都是 $\mathbb Q$ 上的 Galois 扩张。

类似地，由于 $S_3$ 的唯一非平凡正规子群是 $3$ 阶子群，我们看到 $K=\mathbb Q(\sqrt[3]{2},\rho)$ 的子域中只有 $\mathbb Q(\rho)$ 在 $\mathbb Q$ 上是 Galois 的，其 Galois 群同构于 $S_3/\langle\sigma\rangle$，即 $2$ 阶循环群。例如，$\mathbb Q(\rho)$ 的非平凡自同构由把 $\langle\sigma\rangle$ 的非平凡陪集中的任意元素（比如 $\tau$）限制到 $\mathbb Q(\rho)$ 上诱导出来。这一点可由先前给出的这些自同构的显式描述看出：陪集中的每个元素 $\tau,\tau\sigma,\tau\sigma^2$ 都把 $\rho$ 映到 $\rho^2$。$\operatorname{Gal}(K/\mathbb Q)$ 的元素限制到那些（非 Galois 的）三次子域上，一般不给出这些域的自同构，而是给出这些域之间的同构，这与定理 (4) 一致。

\paragraph{例（$\mathbb Q(\sqrt2+\sqrt3)$）}
考虑域 $\mathbb Q(\sqrt2+\sqrt3)$。它显然是 Galois 扩张 $\mathbb Q(\sqrt2,\sqrt3)$ 的一个子域。因此，$\sqrt2+\sqrt3$ 在 $\mathbb Q$ 上的极小多项式的其他根，就是 $\sqrt2+\sqrt3$ 在 Galois 群作用下的互异共轭。这些共轭是
\[
\pm\sqrt2\pm\sqrt3,
\]
它们显然互异。因此极小多项式为
\[
[x-(\sqrt2+\sqrt3)][x-(\sqrt2-\sqrt3)][x-(-\sqrt2+\sqrt3)][x-(-\sqrt2-\sqrt3)],
\]
快速计算得这个多项式为 $x^4-10x^2+1$。于是这个多项式不可约，并且
\[
\mathbb Q(\sqrt2,\sqrt3)=\mathbb Q(\sqrt2+\sqrt3),
\]
这可由次数考虑得到，也可由如下事实得到：在 $\{1,\sigma,\tau,\sigma\tau\}$ 中，只有自同构 $1$ 固定 $\sqrt2+\sqrt3$，所以这个域的固定群同 $\mathbb Q(\sqrt2,\sqrt3)$ 的固定群相同。

\paragraph{例（$x^8-2$ 的分裂域）}
$x^8-2$ 在 $\mathbb Q$ 上的分裂域由 $\theta=\sqrt[8]{2}$（任取一个固定的 $8$ 次根，例如实根）和一个本原 $8$ 次单位根 $\zeta=\zeta_8$ 生成。回忆第 13.6 节，
\[
\mathbb Q(\zeta)=\mathbb Q(i,\sqrt2).
\]
由于 $\theta^4=\sqrt2$，我们看到分裂域由 $\theta$ 与 $i$ 生成。子域 $\mathbb Q(\theta)$ 在 $\mathbb Q$ 上次数为 $8$（因为 $x^8-2$ 由 Eisenstein 判别法不可约），并且这个域中的所有元素都是实数。因此 $i\notin\mathbb Q(\theta)$；又因为 $i$ 至多生成这个域的一个二次扩张，分裂域
\[
\mathbb Q(\sqrt[8]{2},\zeta)=\mathbb Q(\sqrt[8]{2},i)
\]
在 $\mathbb Q$ 上次数为 $16$。

Galois 群由它在生成元 $\theta$ 和 $i$ 上的作用决定，这给出可能性
\[
\theta\longmapsto \zeta^a\theta,
\qquad
 i\longmapsto \pm i,
\qquad a=0,1,2,\ldots,7.
\]
由于我们已经看到该扩张次数为 $16$，而这样的映射至多有 $16$ 种可能，故事实上上面每一种映射都是 $\mathbb Q(\sqrt[8]{2},i)$ 在 $\mathbb Q$ 上的一个自同构。

定义两个自同构
\[
\sigma:\begin{cases}
\theta\longmapsto \zeta\theta,\\
i\longmapsto i,
\end{cases}
\qquad
\tau:\begin{cases}
\theta\longmapsto \theta,\\
i\longmapsto -i,
\end{cases}
\]
其中 $\tau$ 是由复共轭诱导的映射。由于
\[
\zeta=\zeta_8=\frac{\sqrt2}{2}+i\frac{\sqrt2}{2}=\frac12(1+i)\sqrt2
\]
并且 $\theta^4=\sqrt2$，我们可以从下面图 1 中 $\zeta$ 的幂的显式表达式轻易计算出 $\zeta$ 的像。

\begin{center}
\begin{tikzpicture}[scale=1.35, baseline=(current bounding box.center)]
\draw[->] (-1.35,0)--(1.35,0);
\draw[->] (0,-1.35)--(0,1.35) node[above] {$i$};
\draw (0,0) circle (1);
\foreach \ang/\lab in {0/{$1$},45/{$\zeta$},90/{$i$},135/{$\zeta^3$},180/{$-1$},225/{$\zeta^5$},270/{$-i$},315/{$\zeta^7$}}
  \node at (\ang:1.22) {\lab};
\fill (0,0) circle (0.01);
\node at (0,-1.65) {图 1};
\end{tikzpicture}
\end{center}

利用这些显式值，我们发现
\[
\sigma:\begin{cases}
\theta\longmapsto \zeta\theta,\\
i\longmapsto i,\\
\zeta\longmapsto -\zeta=\zeta^5.
\end{cases}
\]
注意，我们还关心元素 $\zeta$ 上的作用，是因为在计算自同构的复合时需要它；例如，在计算
\[
\sigma^2(\theta)=\sigma(\zeta\theta)=\sigma(\zeta)\sigma(\theta)=(-\zeta)(\zeta\theta)=-\zeta^2\theta=-i\theta
\]
时就需要它。

我们同样可以计算出如下自同构：
\[
\begin{array}{c@{\quad}c}
\sigma:\begin{cases}\theta\mapsto\zeta\theta\\ i\mapsto i\\ \zeta\mapsto\zeta^5\end{cases}
&
\tau\sigma:\begin{cases}\theta\mapsto\zeta^7\theta\\ i\mapsto -i\\ \zeta\mapsto\zeta^3\end{cases}\\[5mm]
\sigma^2:\begin{cases}\theta\mapsto -i\theta\\ i\mapsto i\\ \zeta\mapsto\zeta\end{cases}
&
\tau\sigma^2:\begin{cases}\theta\mapsto i\theta\\ i\mapsto -i\\ \zeta\mapsto -\zeta\end{cases}\\[5mm]
\sigma^3:\begin{cases}\theta\mapsto\zeta^7\theta\\ i\mapsto i\\ \zeta\mapsto\zeta^5\end{cases}
&
\tau\sigma^3:\begin{cases}\theta\mapsto\zeta\theta\\ i\mapsto -i\\ \zeta\mapsto\zeta^3\end{cases}\\[5mm]
\sigma^4:\begin{cases}\theta\mapsto-\theta\\ i\mapsto i\\ \zeta\mapsto\zeta\end{cases}
&
\tau\sigma^4:\begin{cases}\theta\mapsto-\theta\\ i\mapsto -i\\ \zeta\mapsto-\zeta\end{cases}\\[5mm]
\sigma^5:\begin{cases}\theta\mapsto\zeta^5\theta\\ i\mapsto i\\ \zeta\mapsto\zeta^5\end{cases}
&
\tau\sigma^5:\begin{cases}\theta\mapsto\zeta^3\theta\\ i\mapsto -i\\ \zeta\mapsto\zeta^3\end{cases}\\[5mm]
\sigma^6:\begin{cases}\theta\mapsto i\theta\\ i\mapsto i\\ \zeta\mapsto\zeta\end{cases}
&
\tau\sigma^6:\begin{cases}\theta\mapsto-i\theta\\ i\mapsto -i\\ \zeta\mapsto-\zeta\end{cases}\\[5mm]
\sigma^7:\begin{cases}\theta\mapsto\zeta^3\theta\\ i\mapsto i\\ \zeta\mapsto-\zeta\end{cases}
&
\tau\sigma^7:\begin{cases}\theta\mapsto\zeta^5\theta\\ i\mapsto -i\\ \zeta\mapsto\zeta^3\end{cases}.
\end{array}
\]
由于这些元素（连同 $1$ 与 $\tau$）穷尽了所有可能性，它们就是 Galois 群。特别地，我们看到 $\sigma$ 与 $\tau$ 生成 Galois 群。为了确定这些元素满足的关系，首先显然有 $\tau^2=1$ 和 $(\sigma^4)^2=1$，所以
\[
\sigma^8=\tau^2=1.
\]
另外，计算得
\[
\sigma\tau:\begin{cases}
\theta\longmapsto\zeta\theta,\\
i\longmapsto -i,\\
\zeta\longmapsto\zeta^3,
\end{cases}
\]
所以
\[
\sigma\tau=\tau\sigma^3.
\]
不难证明这些关系完全定义了这个群，即
\[
\operatorname{Gal}(\mathbb Q(\sqrt[8]{2},i)/\mathbb Q)
=\langle\sigma,\tau\mid \sigma^8=\tau^2=1,\ \sigma\tau=\tau\sigma^3\rangle.
\]
这样的群称为\emph{拟二面体群}（回忆 $16$ 阶二面体群会满足关系 $\sigma\tau=\tau\sigma^7$，而不是 $\sigma\tau=\tau\sigma^3$），并且它是 $S_8$ 的子群，因为 Galois 群是 $x^8-2$ 的 $8$ 个根的置换群的一个子群。

这个例子再次说明，从生成元的作用确定 Galois 群时必须小心。上面我们先计算了 Galois 扩张的次数，以确定 Galois 群中的元素个数。如果我们直接从原生成元 $\theta=\sqrt[8]{2}$ 与 $\zeta=\zeta_8$ 出发，可能会（错误地）断定 Galois 群中共有 $32$ 个元素，因为第一个生成元可被送到 $x^8-2$ 的任意 $8$ 个根之一，而第二个生成元可被送到其极小多项式 $\Phi_4(x)=x^4+1$ 的任意 $4$ 个可能根之一。问题如前所述，在于这些选择并不独立。这里的原因由代数关系
\[
\theta^4=\sqrt2=\zeta+\zeta^7
\]
给出，它说明不能独立指定 $\theta$ 与 $\zeta$ 的像——它们的像必须仍然满足这个代数关系。这个关系或许足够微妙，可以作为一个警告：不要轻率地断言某些映射是自同构。我们还要指出，一般而言必须给出映射是自同构的理由。例如可以使用扩张定理，或者像这里一样使用次数考虑。

确定这个群 $G$ 的子群格是一个直接的问题。其格如下：
\[
\begin{tikzpicture}[scale=0.85, baseline=(current bounding box.center), every node/.style={inner sep=2pt}]
\node (G) at (0,4.2) {$G$};
\node (a2t) at (-3.2,3.0) {$\langle\sigma^2,\tau\rangle$};
\node (a) at (0,3.0) {$\langle\sigma\rangle$};
\node (a2ta3) at (3.2,3.0) {$\langle\sigma^2,\tau\sigma^3\rangle$};
\node (a4ta6) at (-4.4,1.75) {$\langle\sigma^4,\tau\sigma^6\rangle$};
\node (a4t) at (-2.2,1.75) {$\langle\sigma^4,\tau\rangle$};
\node (a2) at (0,1.75) {$\langle\sigma^2\rangle$};
\node (ta3) at (2.2,1.75) {$\langle\tau\sigma^3\rangle$};
\node (ta) at (4.4,1.75) {$\langle\tau\sigma\rangle$};
\node (ta2) at (-5.2,0.6) {$\langle\tau\sigma^2\rangle$};
\node (ta6) at (-3.5,0.6) {$\langle\tau\sigma^6\rangle$};
\node (ta4) at (-1.8,0.6) {$\langle\tau\sigma^4\rangle$};
\node (t) at (-0.2,0.6) {$\langle\tau\rangle$};
\node (a4) at (1.2,0.6) {$\langle\sigma^4\rangle$};
\node (one) at (0,-0.7) {$1$};
\foreach \x in {a2t,a,a2ta3}{\draw (G)--(\x);}
\draw (a2t)--(a4ta6); \draw (a2t)--(a4t); \draw (a2t)--(a2);
\draw (a)--(a2);
\draw (a2ta3)--(a2); \draw (a2ta3)--(ta3); \draw (a2ta3)--(ta);
\draw (a4ta6)--(ta2); \draw (a4ta6)--(ta6); \draw (a4ta6)--(a4);
\draw (a4t)--(ta4); \draw (a4t)--(t); \draw (a4t)--(a4);
\draw (a2)--(a4);
\draw (ta3)--(a4); \draw (ta)--(a4);
\foreach \x in {ta2,ta6,ta4,t,a4}{\draw (\x)--(one);}
\end{tikzpicture}
\]

利用基本定理中 (2)——即在 $\mathbb Q$ 上的扩张次数等于固定子群的指数——对上面许多子群确定相应子域（根据基本定理，这给出了 $\mathbb Q(\sqrt[8]{2},i)$ 的所有子域）是相当简单的。于是只需找出一个具有正确次数、并被相应子群固定的子域。还要记住，如果一个子域被某个子群的生成元固定，那么它就被该子群固定。例如，从自同构 $\sigma$ 的显式描述可见，$\mathbb Q(i)$ 被 $\sigma$ 生成的群固定。由于这是指数为 $2$ 的子群，而 $\mathbb Q(i)$ 在 $\mathbb Q$ 上次数为 $2$，它必定就是全部固定域。上面大多数子群的固定域都可以这样简单地确定。

对右边的 $4$ 阶子群（即由 $\tau\sigma^3$ 与由 $\tau\sigma$ 生成的子群），也许不那么容易看出如何确定相应的固定域。对由 $\tau\sigma^3$ 生成的子群 $H$，可作如下处理：元素 $\theta^2=\sqrt[4]{2}$ 显然被 $\sigma^4$ 固定。由上面的图，$\sigma^4$ 是 $H$ 的一个指数为 $2$ 的正规子群，其陪集代表可取 $1,\tau\sigma^3$。考虑元素
\[
\alpha=(1+\tau\sigma^3)\theta^2=\theta^2+\tau\sigma^3\theta^2.
\]
那么 $\alpha$ 被 $\sigma^4$ 固定（我们在 $4$ 阶交换群 $H$ 中，所以 $\sigma^4$ 与 $1$、$\tau\sigma^3$ 交换，并且已经知道 $\theta^2$ 被 $\sigma^4$ 固定）。但是（这正是关键）$\alpha$ 也被 $\tau\sigma^3$ 固定：
\[
\tau\sigma^3\alpha=\tau\sigma^3(1+\tau\sigma^3)\theta^2=[\tau\sigma^3+(\tau\sigma^3)^2]\theta^2
=(\tau\sigma^3+\sigma^4)\theta^2,
\]
而最后这个表达式就是 $\alpha$，因为 $\sigma^4\theta^2=\theta^2$。因此 $\alpha$ 属于 $H$ 的固定域。显式地，
\[
\alpha=\sqrt[4]{2}+i\sqrt[4]{2}=(1+i)\sqrt[4]{2}.
\]
快速检查可见 $\alpha$ 不被自同构 $\sigma^2$ 固定，因此由上面的子群图可知，域 $\mathbb Q(\alpha)$ 的固定子群不可能大于 $H$，从而正是 $H$，这就给出了固定域。回忆一般地，如果 $E$ 是 $H$ 的固定域，那么 $\tau H\tau^{-1}$ 的固定域是 $\tau(E)$，这也给出了 $\langle\tau\sigma\rangle$ 的固定域。对 $H=\langle\tau\sigma^3\rangle$，有 $\tau H\tau^{-1}=\langle\tau\sigma\rangle$，其固定域由
\[
\tau(\alpha)=(1-i)\sqrt[4]{2}
\]
给出。

一般而言，人们试图确定被 Galois 群中给定子群 $H$ 固定的元素（参见练习，那里指出了上面的元素从何而来），并尝试生成足够大的域以得到全部固定域。在我们的情形中，我们能够用单个生成元做到这一点。稍后我们将看到，$\mathbb Q$ 的任意有限扩张都是单扩张，因此会有这种类型的单个生成元；不过一般而言可能难以直接产生它。

元素 $\alpha$ 是多项式
\[
x^4+8
\]
的根；这个多项式因此必定不可约，因为我们已经确定了它的一个根生成 $\mathbb Q$ 上次数为 $4$ 的扩张。

类似地，可以补全 $\mathbb Q(\sqrt[8]{2},i)$ 的子域图。为了强调它与上面子群图的关系，我们把它倒置过来（$\theta=\sqrt[8]{2}$）：
\[
\begin{tikzpicture}[scale=0.74, baseline=(current bounding box.center), every node/.style={inner sep=1.5pt}]
\node (Q) at (0,4.4) {$\mathbb Q$};
\node (sqrt2) at (-3.6,3.15) {$\mathbb Q(\sqrt2)$};
\node (i) at (0,3.15) {$\mathbb Q(i)$};
\node (sqrtm2) at (3.6,3.15) {$\mathbb Q(\sqrt{-2})$};
\node (iq4) at (-5.0,1.9) {$\mathbb Q(i\sqrt[4]{2})$};
\node (q4) at (-3.0,1.9) {$\mathbb Q(\sqrt[4]{2})$};
\node (isqrt2) at (-0.8,1.9) {$\mathbb Q(i,\sqrt2)$};
\node (onep) at (1.6,1.9) {$\mathbb Q((1+i)\sqrt[4]{2})$};
\node (onem) at (4.5,1.9) {$\mathbb Q((1-i)\sqrt[4]{2})$};
\node (ztheta) at (-5.8,0.65) {$\mathbb Q(\zeta\theta)$};
\node (z3theta) at (-4.05,0.65) {$\mathbb Q(\zeta^3\theta)$};
\node (z2theta) at (-2.15,0.65) {$\mathbb Q(\zeta^2\theta)$};
\node (theta) at (-0.45,0.65) {$\mathbb Q(\theta)$};
\node (iq4full) at (1.6,0.65) {$\mathbb Q(i,\sqrt[4]{2})$};
\node (bottom) at (0,-0.8) {$\mathbb Q(i,\theta)$};
\draw (Q)--(sqrt2); \draw (Q)--(i); \draw (Q)--(sqrtm2);
\draw (sqrt2)--(iq4); \draw (sqrt2)--(q4); \draw (sqrt2)--(isqrt2);
\draw (i)--(isqrt2);
\draw (sqrtm2)--(isqrt2); \draw (sqrtm2)--(onep); \draw (sqrtm2)--(onem);
\draw (iq4)--(ztheta); \draw (iq4)--(z3theta); \draw (q4)--(z2theta); \draw (q4)--(theta); \draw (isqrt2)--(iq4full); \draw (onep)--(iq4full); \draw (onem)--(iq4full);
\foreach \x in {ztheta,z3theta,z2theta,theta,iq4full}{\draw (\x)--(bottom);}
\end{tikzpicture}
\]
注意，群 $\langle\sigma^4\rangle$ 是 $G$ 的正规子群（事实上是 $G$ 的中心），并且商群 $G/\langle\sigma^4\rangle\cong D_8$，所以相应固定域 $\mathbb Q(i,\sqrt[4]{2})$ 在 $\mathbb Q$ 上是 Galois 的，Galois 群为 $D_8$。由于它是 Galois 的，它是某个分裂域，显然是 $x^4-2$ 的分裂域。这个域的子域格于是立即由上面的格给出。

我们以这个 Galois 扩张的下列有趣方面结束这个例子。容易验证
\[
\langle\sigma^2,\tau\rangle\cong D_8,
\qquad
\langle\sigma\rangle\cong\mathbb Z/8\mathbb Z,
\qquad
\langle\sigma^2,\tau\sigma^3\rangle\cong Q_8,
\]
其中 $D_8$ 是 $8$ 阶二面体群，$Q_8$ 是 $8$ 阶四元数群。因此，域 $\mathbb Q(\sqrt[8]{2},i)$ 在它的三个二次子域
\[
\mathbb Q(\sqrt2),
\qquad
\mathbb Q(i),
\qquad
\mathbb Q(\sqrt{-2})
\]
上分别是次数为 $8$ 的 Galois 扩张，其 Galois 群分别为二面体群、循环群和四元数群。于是 $8$ 阶群的 $5$ 种可能类型中有三种（并且两个非交换的类型都）作为这个扩张中的 Galois 群出现。

我们将在下面各节考虑更多例子和应用。

\subsection*{练习}

\begin{enumerate}[label=\arabic*.]
\item 求元素 $\sqrt2+\sqrt5$ 在 $\mathbb Q$ 上的极小多项式。

\item 求元素 $1+\sqrt[3]{2}+\sqrt[3]{4}$ 在 $\mathbb Q$ 上的极小多项式。

\item 求 $(x^2-2)(x^2-3)(x^2-5)$ 的 Galois 群。求这个多项式的分裂域的所有子域。

\item 设 $p$ 为素数。求 $x^p-2$ 的 Galois 群的元素。

\item 证明：当 $p$ 为素数时，$x^p-2$ 的 Galois 群同构于矩阵群
\[
\left\{\begin{pmatrix}a&b\\0&1\end{pmatrix}\mid a,b\in\mathbb F_p,
\ a\neq0\right\}.
\]

\item 令 $K=\mathbb Q(\sqrt[8]{2},i)$，并令 $F_1=\mathbb Q(i)$、$F_2=\mathbb Q(\sqrt2)$、$F_3=\mathbb Q(\sqrt{-2})$。证明
\[
\operatorname{Gal}(K/F_1)\cong\mathbb Z_8,
\qquad
\operatorname{Gal}(K/F_2)\cong D_8,
\qquad
\operatorname{Gal}(K/F_3)\cong Q_8.
\]

\item 求 $x^8-2$ 的分裂域中所有在 $\mathbb Q$ 上 Galois 的子域。

\item 设 $K$ 是 $F$ 的一个次数为 $p^n$ 的 Galois 扩张，其中 $p$ 是素数且 $n\geq1$。证明在 $K$ 中存在 $F$ 的次数为 $p$ 和 $p^{n-1}$ 的 Galois 扩张。

\item 给出域 $F_1,F_2,F_3$ 的例子，使得 $\mathbb Q\subset F_1\subset F_2\subset F_3$，$[F_3:\mathbb Q]=8$，并且每个域都在其所有子域上是 Galois 的，唯一例外是 $F_2$ 不在 $\mathbb Q$ 上 Galois。

\item 求 $x^8-3$ 在 $\mathbb Q$ 上的分裂域的 Galois 群。

\item 设 $f(x)\in\mathbb Z[x]$ 是不可约四次多项式，并且它在 $\mathbb Q$ 上的分裂域的 Galois 群为 $S_4$（这样的四次多项式有很多，参见第 6 节）。令 $\theta$ 是 $f(x)$ 的一个根，并设 $K=\mathbb Q(\theta)$。证明 $K$ 是 $\mathbb Q$ 的次数为 $4$ 的扩张，并且没有真子域。是否存在次数为 $4$ 且没有真子域的 $\mathbb Q$ 的 Galois 扩张？

\item 求 $x^4-14x^2+9$ 在 $\mathbb Q$ 上的分裂域的 Galois 群。

\item 证明：如果某个三次多项式在 $\mathbb Q$ 上的分裂域的 Galois 群是 $3$ 阶循环群，那么该三次多项式的所有根都是实数。

\item 证明 $\mathbb Q(\sqrt{2+\sqrt2})$ 是一个循环四次域，即它是一个次数为 $4$ 且 Galois 群为循环群的 Galois 扩张。

\item （二二次扩张）设 $F$ 是一个特征不为 $2$ 的域。
\begin{enumerate}[label=(\alph*)]
\item 若 $K=F(\sqrt{D_1},\sqrt{D_2})$，其中 $D_1,D_2\in F$，并且 $D_1$、$D_2$、$D_1D_2$ 都不是 $F$ 中的平方，证明 $K/F$ 是 Galois 扩张，且 $\operatorname{Gal}(K/F)$ 同构于 Klein 四元群。
\item 反过来，设 $K/F$ 是 Galois 扩张，并且 $\operatorname{Gal}(K/F)$ 同构于 Klein 四元群。证明 $K=F(\sqrt{D_1},\sqrt{D_2})$，其中 $D_1,D_2\in F$，并且 $D_1$、$D_2$、$D_1D_2$ 都不是 $F$ 中的平方。
\end{enumerate}

\item
\begin{enumerate}[label=(\alph*)]
\item 证明 $x^4-2x^2-2$ 在 $\mathbb Q$ 上不可约。
\item 说明这个四次多项式的根为
\[
\alpha_1=\sqrt{1+\sqrt3},\quad
\alpha_2=\sqrt{1-\sqrt3},\quad
\alpha_3=-\sqrt{1+\sqrt3},\quad
\alpha_4=-\sqrt{1-\sqrt3}.
\]
\item 令 $K_1=\mathbb Q(\alpha_1)$、$K_2=\mathbb Q(\alpha_2)$。说明 $K_1\neq K_2$，且 $K_1\cap K_2=\mathbb Q(\sqrt3)=F$。
\item 证明 $K_1$、$K_2$ 和 $K_1K_2$ 都在 $F$ 上 Galois，并且 $\operatorname{Gal}(K_1K_2/F)$ 是 Klein 四元群。明确写出 $\operatorname{Gal}(K_1K_2/F)$ 的元素。确定 Galois 群的所有子群，并给出它们对应的含 $F$ 的 $K_1K_2$ 的固定子域。
\item 证明 $x^4-2x^2-2$ 在 $\mathbb Q$ 上的分裂域次数为 $8$，其 Galois 群为二面体群。
\end{enumerate}

下面两个练习指出了在给定域中构造子域元素的一种方法，并且在许多计算中相当有用。

\item 设 $K/F$ 是任意有限扩张，$\alpha\in K$。令 $L$ 是一个含有 $K$ 的 $F$ 的 Galois 扩张，并令 $H\leq\operatorname{Gal}(L/F)$ 是对应于 $K$ 的子群。定义 $\alpha$ 从 $K$ 到 $F$ 的\emph{范数}为
\[
N_{K/F}(\alpha)=\prod_\sigma \sigma(\alpha),
\]
其中乘积取遍 $K$ 到 $F$ 的一个代数闭包中的所有嵌入（所以由 Galois 理论基本定理，它取遍 $\operatorname{Gal}(L/F)$ 中 $H$ 的一组陪集代表）。这是 $\alpha$ 的 Galois 共轭的乘积。特别地，如果 $K/F$ 是 Galois 的，那么这是
\[
\prod_{\sigma\in\operatorname{Gal}(K/F)}\sigma(\alpha).
\]
\begin{enumerate}[label=(\alph*)]
\item 证明 $N_{K/F}(\alpha)\in F$。
\item 证明 $N_{K/F}(\alpha\beta)=N_{K/F}(\alpha)N_{K/F}(\beta)$，因此范数是从 $K$ 到 $F$ 的乘法映射。
\item 令 $K=F(\sqrt D)$ 是 $F$ 的一个二次扩张。说明
\[
N_{K/F}(a+b\sqrt D)=a^2-Db^2.
\]
\item 设
\[
m_\alpha(x)=x^d+a_{d-1}x^{d-1}+\cdots+a_1x+a_0\in F[x]
\]
是 $\alpha\in K$ 在 $F$ 上的极小多项式。令 $n=[K:F]$。证明 $d$ 整除 $n$，并且 $\alpha$ 的 $d$ 个互异 Galois 共轭在上面的乘积中都重复 $n/d$ 次，从而推出
\[
N_{K/F}(\alpha)=(-1)^n a_0^{n/d}.
\]
\end{enumerate}

\item 在前一题的记号下，定义 $\alpha$ 从 $K$ 到 $F$ 的\emph{迹}为
\[
\operatorname{Tr}_{K/F}(\alpha)=\sum_\sigma \sigma(\alpha),
\]
即 Galois 共轭的和。
\begin{enumerate}[label=(\alph*)]
\item 证明 $\operatorname{Tr}_{K/F}(\alpha)\in F$。
\item 证明
\[
\operatorname{Tr}_{K/F}(\alpha+\beta)=\operatorname{Tr}_{K/F}(\alpha)+\operatorname{Tr}_{K/F}(\beta),
\]
因此迹是从 $K$ 到 $F$ 的加法映射。
\item 令 $K=F(\sqrt D)$ 是 $F$ 的一个二次扩张。说明
\[
\operatorname{Tr}_{K/F}(a+b\sqrt D)=2a.
\]
\item 设 $m_\alpha(x)$ 如前一题所述。证明
\[
\operatorname{Tr}_{K/F}(\alpha)=-a_{d-1}.
\]
\end{enumerate}

\item 在前面各题的记号下，证明对基域 $F$ 中所有 $a$，有
\[
N_{K/F}(a\alpha)=a^nN_{K/F}(\alpha),
\qquad
\operatorname{Tr}_{K/F}(a\alpha)=a\operatorname{Tr}_{K/F}(\alpha).
\]
特别地，对所有 $a\in F$，证明
\[
N_{K/F}(a)=a^n,
\qquad
\operatorname{Tr}_{K/F}(a)=na.
\]

\item 在前面各题的记号下，更一般地证明
\[
\prod_\sigma (x-\sigma(\alpha))=(m_\alpha(x))^{n/d}.
\]

\item 利用特征的线性无关性证明：对任意 Galois 扩张 $K/F$，存在 $\alpha\in K$ 使得 $\operatorname{Tr}_{K/F}(\alpha)\neq0$。

\item 设 $K/F$ 是 Galois 扩张，$\sigma$ 是 Galois 群的一个元素。
\begin{enumerate}[label=(\alph*)]
\item 设 $\alpha\in K$ 具有形式
\[
\alpha=\frac{\beta}{\sigma\beta}
\]
其中 $0\neq\beta\in K$。证明 $N_{K/F}(\alpha)=1$。
\item 设 $\alpha\in K$ 具有形式
\[
\alpha=\beta-\sigma\beta
\]
其中 $\beta\in K$。证明 $\operatorname{Tr}_{K/F}(\alpha)=0$。
\end{enumerate}

下一题和后面的练习 26 建立 Hilbert 定理 90 的乘法形式与加法形式。这些是第一上同调群消失的例子，如第 17.3 节将讨论的那样。

\item （Hilbert 定理 90）设 $K$ 是 $F$ 的 Galois 扩张，其 Galois 群为由 $\sigma$ 生成的 $n$ 阶循环群。设 $\alpha\in K$ 满足 $N_{K/F}(\alpha)=1$。证明 $\alpha$ 具有形式
\[
\alpha=\frac{\beta}{\sigma\beta}
\]
其中 $0\neq\beta\in K$。[利用特征的线性无关性证明存在某个 $\theta\in K$ 使得
\[
\beta=\theta+\alpha\sigma(\theta)+(\alpha\sigma\alpha)\sigma^2(\theta)+\cdots+(
\alpha\sigma\alpha\cdots\sigma^{n-2}\alpha)\sigma^{n-1}(\theta)
\]
非零。利用 $\alpha$ 到 $F$ 的范数为 $1$ 这一事实计算 $\beta/\sigma\beta$。]

\item 证明 Pythagoras 方程 $a^2+b^2=1$ 的有理数解 $a,b\in\mathbb Q$ 具有形式
\[
a=\frac{s^2-t^2}{s^2+t^2},
\qquad
b=\frac{2st}{s^2+t^2}
\]
其中 $s,t\in\mathbb Q$；并由此说明任意整数边直角三角形的边长为
\[
(m^2-n^2,\ 2mn,\ m^2+n^2)
\]
其中 $m,n$ 是整数。[注意 $a^2+b^2=1$ 等价于 $N_{\mathbb Q(i)/\mathbb Q}(a+ib)=1$，然后把上面的 Hilbert 定理 90 用于 $\beta=s+it$。]

\item 将前一题推广，确定方程
\[
a^2+Db^2=1
\]
的所有有理数解，其中 $D\in\mathbb Z$，$D>0$，并且 $D$ 不是 $\mathbb Z$ 中的完全平方。

\item （加法 Hilbert 定理 90）设 $K$ 是 $F$ 的 Galois 扩张，其 Galois 群为由 $\sigma$ 生成的 $n$ 阶循环群。设 $\alpha\in K$ 满足 $\operatorname{Tr}_{K/F}(\alpha)=0$。证明 $\alpha$ 具有形式
\[
\alpha=\beta-\sigma\beta
\]
其中 $\beta\in K$。[由前一练习，取 $\theta\in K$ 满足 $\operatorname{Tr}_{K/F}(\theta)\neq0$，令
\[
\beta=\frac{1}{\operatorname{Tr}_{K/F}(\theta)}
\bigl[\alpha\sigma(\theta)+(\alpha+\sigma\alpha)\sigma^2(\theta)+\cdots+
(\alpha+\sigma\alpha+\cdots+\sigma^{n-2}\alpha)\sigma^{n-1}(\theta)\bigr]
\]
并计算 $\beta-\sigma\beta$。]

\item 令
\[
\alpha=\sqrt{(2+\sqrt2)(3+\sqrt3)}
\]
（为具体起见取正实平方根），并考虑扩张 $E=\mathbb Q(\alpha)$。
\begin{enumerate}[label=(\alph*)]
\item 说明 $\alpha^2=(2+\sqrt2)(3+\sqrt3)$ 不是 $F=\mathbb Q(\sqrt2,\sqrt3)$ 中的平方。[若 $\alpha^2=c^2$，$c\in F$，则对固定 $\mathbb Q(\sqrt2)$ 的自同构 $\varphi\in\operatorname{Gal}(F/\mathbb Q)$ 有
\[
\alpha\varphi(\alpha)=(2+\sqrt2)^2(6)=(c\varphi c)^2.
\]
由于 $c\varphi c=N_{F/\mathbb Q(\sqrt2)}(c)\in\mathbb Q(\sqrt2)$，推出这蕴含 $\sqrt6\in\mathbb Q(\sqrt2)$，矛盾。]
\item 由 (a) 推出 $[E:\mathbb Q]=8$。证明 $\alpha$ 在 $\mathbb Q$ 上的极小多项式的根是 $8$ 个元素
\[
\pm\sqrt{(2\pm\sqrt2)(3\pm\sqrt3)}.
\]
\item 令
\[
\beta=\sqrt{(2-\sqrt2)(3+\sqrt3)}.
\]
证明 $\alpha\beta=\sqrt2(3+\sqrt3)\in F$，因此 $\beta\in E$。类似地说明其他根也都属于 $E$，从而 $E$ 是 $\mathbb Q$ 的 Galois 扩张。说明 Galois 群的元素恰好是由把 $\alpha$ 映到 (b) 中八个元素之一所确定的映射。
\item 令 $\sigma\in\operatorname{Gal}(E/\mathbb Q)$ 为把 $\alpha$ 映到 $\beta$ 的自同构。说明由于 $\sigma(\alpha^2)=\beta^2$，有 $\sigma(\sqrt2)=-\sqrt2$ 且 $\sigma(\sqrt3)=\sqrt3$。由 $\alpha\beta=\sqrt2(3+\sqrt3)$ 推出 $\sigma(\alpha\beta)=-\alpha\beta$，从而 $\sigma(\beta)=-\alpha$。说明 $\sigma$ 是 $\operatorname{Gal}(E/\mathbb Q)$ 中的一个 $4$ 阶元素。
\item 类似地说明由
\[
\tau(\alpha)=\sqrt{(2+\sqrt2)(3-\sqrt3)}
\]
定义的映射 $\tau$ 是 $\operatorname{Gal}(E/\mathbb Q)$ 中的一个 $4$ 阶元素。证明 $\sigma$ 与 $\tau$ 生成 Galois 群，$\sigma^4=\tau^4=1$，$\sigma^2=\tau^2$，且 $\sigma\tau=\tau\sigma^3$。
\item 推出 $\operatorname{Gal}(E/\mathbb Q)\cong Q_8$，即 $8$ 阶四元数群。
\end{enumerate}

\item 设 $f(x)\in F[x]$ 是域 $F$ 上次数为 $n$ 的不可约多项式，$L$ 是 $f(x)$ 在 $F$ 上的分裂域，$\alpha$ 是 $f(x)$ 在 $L$ 中的一个根。若 $K$ 是 $F$ 的包含在 $L$ 中的任意 Galois 扩张，说明多项式 $f(x)$ 在 $K$ 上分解为 $m$ 个不可约多项式的乘积，每个不可约因子的次数为 $d$，其中
\[
m=[F(\alpha)\cap K:F],
\qquad
 d=[K(\alpha):K]
\]
（也参见第 4 节练习 4 中的推广）。[若 $H$ 是 $L/F$ 的 Galois 群中对应于 $K$ 的子群，则 $f(x)$ 在 $K$ 上的因子对应于 $H$ 在 $f(x)$ 的根上的轨道。然后使用第 4.1 节练习 9。]

\item 设 $k$ 是域，$k(t)$ 是变量 $t$ 中的有理函数域。定义 $k(t)$ 到自身的映射 $\sigma$ 与 $\tau$ 为
\[
\sigma f(t)=f\left(\frac{-1}{1-t}\right),
\qquad
\tau f(t)=f\left(\frac1t\right),
\qquad f(t)\in k(t).
\]
\begin{enumerate}[label=(\alph*)]
\item 证明 $\sigma$ 与 $\tau$ 是 $k(t)$ 的自同构（参见第 1 节练习 8），并且它们生成的群 $G=\langle\sigma,\tau\rangle$ 同构于 $S_3$。
\item 证明元素
\[
s=\frac{(t^2-t+1)^3}{t^2(t-1)^2}
\]
被 $G$ 的所有元素固定。
\item 证明 $k(s)$ 正是 $G$ 在 $k(t)$ 中的固定域。[计算扩张次数。]
\end{enumerate}

\item 证明前一题中由 $\tau$ 生成的自同构子群的固定域是
\[
k\left(t+\frac1t\right).
\]
证明由自同构 $\tau\sigma^2$（它把 $t$ 映到 $1-t$）生成的子群的固定域是 $k(t(1-t))$。求由 $\tau\sigma$ 生成的子群的固定域，以及由 $\sigma$ 生成的子群的固定域。

\item 设 $K$ 是 $F$ 的次数为 $n$ 的有限扩张。令 $\alpha$ 是 $K$ 的一个元素。
\begin{enumerate}[label=(\alph*)]
\item 证明 $\alpha$ 通过在 $K$ 上的左乘作用，给出 $K$ 的一个 $F$-线性变换 $T_\alpha$。
\item 证明 $\alpha$ 在 $F$ 上的极小多项式同线性变换 $T_\alpha$ 的极小多项式相同。
\item 证明迹 $\operatorname{Tr}_{K/F}(\alpha)$ 是由 $T_\alpha$ 定义的 $n\times n$ 矩阵的迹（这说明“迹”这个词的两种用法是一致的）。证明范数 $N_{K/F}(\alpha)$ 是 $T_\alpha$ 的行列式。
\end{enumerate}
\end{enumerate}

\section{有限域}

有限域 $\mathbb F$ 的特征为某个素数 $p$，因而是 $\mathbb F_p$ 上的一个有限维向量空间。如果维数为 $n$，即 $[\mathbb F:\mathbb F_p]=n$，则 $\mathbb F$ 恰有 $p^n$ 个元素。我们已经看到（在命题 13.37 之后），$\mathbb F$ 同构于多项式 $x^{p^n}-x$ 的分裂域，因此在同构意义下唯一。我们把阶为 $p^n$ 的有限域记为 $\mathbb F_{p^n}$。

域 $\mathbb F_{p^n}$ 在 $\mathbb F_p$ 上是 Galois 的，它的 Galois 群是由 Frobenius 自同构生成的 $n$ 阶循环群
\[
\operatorname{Gal}(\mathbb F_{p^n}/\mathbb F_p)=\langle \sigma_p\rangle\cong \mathbb Z/n\mathbb Z,
\]
其中
\[
\sigma_p(\alpha)=\alpha^p.
\]
这正是推论 6 后例 7 的内容。由基本定理，$\mathbb F_{p^n}$ 的每个子域对应于 $\mathbb Z/n\mathbb Z$ 的一个子群。因此，对 $n$ 的每个正因子 $d$，恰有一个在 $\mathbb F_p$ 上次数为 $d$ 的子域，即由阶为 $n/d$ 的子群 $\langle \sigma_p^d\rangle$ 固定的域，并且没有其他子域。这个域同构于 $\mathbb F_{p^d}$，即阶为 $p^d$ 的唯一有限域。

由于 Galois 群是 Abel 群，每个子群都是正规子群，所以每个子域 $\mathbb F_{p^d}$（其中 $d$ 是 $n$ 的一个因子）都在 $\mathbb F_p$ 上是 Galois 的（这也可由这些域本身就是分裂域这一事实直接看出）。此外，Galois 群 $\operatorname{Gal}(\mathbb F_{p^d}/\mathbb F_p)$ 由 $\sigma_p$ 在商群
\[
\operatorname{Gal}(\mathbb F_{p^n}/\mathbb F_p)/\langle \sigma_p^d\rangle
\]
中的像生成。如果仍把这个元素记作 $\sigma_p$，我们就重新得到扩张 $\mathbb F_{p^d}/\mathbb F_p$ 的 Frobenius 自同构。（不过要注意，$\sigma_p$ 在 $\operatorname{Gal}(\mathbb F_{p^n}/\mathbb F_p)$ 中阶为 $n$，而在 $\operatorname{Gal}(\mathbb F_{p^d}/\mathbb F_p)$ 中阶为 $d$。）

我们把这些结果概括为如下命题。

\noindent\textbf{命题 15.} 任意有限域都同构于某个 $\mathbb F_{p^n}$，其中 $p$ 是素数且 $n\geq 1$。域 $\mathbb F_{p^n}$ 是多项式 $x^{p^n}-x$ 在 $\mathbb F_p$ 上的分裂域，其 Galois 群是由 Frobenius 自同构 $\sigma_p$ 生成的 $n$ 阶循环群。$\mathbb F_{p^n}$ 的子域全都在 $\mathbb F_p$ 上是 Galois 的，并且同 $n$ 的因子 $d$ 一一对应；它们是域 $\mathbb F_{p^d}$，即 $\sigma_p^d$ 的固定域。

任意有限域的有限扩张的相应陈述，都是命题 15 的容易推论，并在练习中列出。

作为一个初等应用，我们得到整数多项式 $x^4+1\in\mathbb Z[x]$ 的如下结果。

\noindent\textbf{推论 16.} 不可约多项式 $x^4+1\in\mathbb Z[x]$ 模每个素数 $p$ 都可约。

\begin{proof}
考虑素数 $p$ 上的多项式 $x^4+1\in\mathbb F_p[x]$。如果 $p=2$，则 $x^4+1=(x+1)^4$，因此该多项式可约。现在假设 $p$ 为奇素数。由于 $p$ 模 $8$ 同余于 $1,3,5$ 或 $7$，而这些数的平方都模 $8$ 同余于 $1$，所以 $p^2-1$ 可被 $8$ 整除。因此 $x^{p^2-1}-1$ 可被 $x^8-1$ 整除。于是有整除关系
\[
x^4+1\mid x^8-1\mid x^{p^2-1}-1\mid x^{p^2}-x,
\]
这说明 $x^4+1$ 的所有根都是 $x^{p^2}-x$ 的根。（等价地，这些根被 Frobenius 自同构的平方 $\sigma_p^2$ 固定。）由于 $x^{p^2}-x$ 的根就是域 $\mathbb F_{p^2}$，由 $x^4+1$ 的任意一个根生成的扩张在 $\mathbb F_p$ 上的次数至多为 $2$，这意味着 $x^4+1$ 不可能在 $\mathbb F_p$ 上不可约。
\end{proof}

乘法群 $\mathbb F_{p^n}^{\times}$ 显然是一个域的乘法群的有限子群。由命题 9.18，它是循环群。如果 $\theta$ 是任意一个生成元，那么显然
\[
\mathbb F_{p^n}=\mathbb F_p(\theta).
\]
这证明了如下结果。

\noindent\textbf{命题 17.} 有限域 $\mathbb F_{p^n}$ 是单扩张。特别地，对每个 $n\geq 1$，在 $\mathbb F_p$ 上都存在次数为 $n$ 的不可约多项式。

上面我们已经把有限域 $\mathbb F_{p^n}$ 描述为多项式 $x^{p^n}-x$ 的分裂域。由前一命题，这个域也可以描述为 $\mathbb F_p[x]$ 的一个商，即由 $\theta$ 的极小多项式所定义的商。由于 $\theta$ 必然是 $x^{p^n}-x$ 的根，我们看到 $\theta$ 的极小多项式是 $x^{p^n}-x$ 的一个次数为 $n$ 的因子。

反过来，设 $p(x)$ 是任意一个次数为 $d$ 的不可约多项式，并且 $p(x)$ 整除 $x^{p^n}-x$。如果 $\alpha$ 是 $p(x)$ 的一个根，那么扩张 $\mathbb F_p(\alpha)$ 是 $\mathbb F_{p^n}$ 的一个次数为 $d$ 的子域。因此 $d$ 是 $n$ 的一个因子，并且由命题 15，这个扩张是 Galois 的（事实上就是扩张 $\mathbb F_{p^d}$），所以特别地，$p(x)$ 的所有根都包含在 $\mathbb F_p(\alpha)$ 中。

域 $\mathbb F_{p^n}$ 的元素正是 $x^{p^n}-x$ 的根。如果按照这些根在 $\mathbb F_p$ 上的极小多项式的次数 $d$，把这个多项式的因子 $x-\alpha$ 分组，就得到：

\noindent\textbf{命题 18.} 多项式 $x^{p^n}-x$ 正好是 $\mathbb F_p[x]$ 中所有不同的不可约多项式的乘积，其中这些不可约多项式的次数 $d$ 取遍 $n$ 的所有因子。

这个命题可以用来递归地产生 $\mathbb F_p$ 上的不可约多项式。例如，$\mathbb F_2$ 上的不可约二次多项式是
\[
\frac{x^4-x}{x(x-1)}
\]
的因子，于是得到唯一的多项式 $x^2+x+1$。类似地，这个域上的不可约三次多项式是
\[
\frac{x^8-x}{x(x-1)}=x^6+x^5+x^4+x^3+x^2+x+1
\]
的因子，它分解为两个三次多项式
\[
x^3+x+1\qquad\text{和}\qquad x^3+x^2+1.
\]
不可约四次多项式由 $x^{16}-x$ 除以 $x(x-1)$ 以及上面的不可约二次多项式 $x^2+x+1$ 后，再分解为不可约四次多项式而得：
\[
\frac{x^{16}-x}{x(x-1)(x^2+x+1)}
=(x^4+x^3+x^2+x+1)(x^4+x^3+1)(x^4+x+1).
\]

这给出了一种确定 $\mathbb F_p$ 上所有给定次数不可约多项式之乘积的方法。实践中，存在用于模 $p$ 分解多项式的有效算法，它们可以给出各个不可约多项式（见练习）。掌握不可约多项式之所以重要，是因为它们给出有限域 $\mathbb F_{p^n}$ 的一种便于显式计算的表示，即把它表示为商
\[
\mathbb F_p[x]/(f(x)),
\]
其中 $f(x)$ 是次数为 $n$ 的不可约多项式。

还要注意，由于有限域 $\mathbb F_{p^n}$ 在同构意义下唯一，$\mathbb F_p[x]$ 分别模任意两个次数为 $n$ 的不可约多项式所得的商域全都同构。如果 $f_1(x)$ 和 $f_2(x)$ 都是次数为 $n$ 的不可约多项式，那么 $f_2(x)$ 在域
\[
\mathbb F_{p^n}\cong \mathbb F_p[x]/(f_1(x))
\]
中完全分裂。如果把 $f_2(x)$ 的一个根记为 $\alpha(x)$（强调它在 $\mathbb F_p[x]/(f_1(x))$ 中是一个次数 $<n$ 的关于 $x$ 的多项式），则同构由
\[
\mathbb F_p[x]/(f_2(x))\longrightarrow \mathbb F_p[x]/(f_1(x)),
\qquad
x\longmapsto \alpha(x)
\]
给出（我们把第一个域中 $f_2(x)$ 的一个根映到第二个域中 $f_2(x)$ 的一个根）。例如，如果
\[
f_1(x)=x^4+x^3+1,
\qquad
f_2(x)=x^4+x+1
\]
是上面确定的 $\mathbb F_2$ 上两个不可约四次多项式，那么一个简单计算验证
\[
\alpha(x)=x^3+x^2
\]
是 $f_2(x)$ 在 $\mathbb F_{16}=\mathbb F_2[x]/(x^4+x^3+1)$ 中的一个根。于是有
\[
\mathbb F_2[x]/(x^4+x+1)\longrightarrow \mathbb F_2[x]/(x^4+x^3+1)\quad(\cong \mathbb F_{16}),
\qquad
x\longmapsto x^3+x^2.
\]

如果我们假定一个初等数论结果，就可以给出次数为 $n$ 的不可约多项式的个数公式。定义 Möbius $\mu$-函数为
\[
\mu(n)=
\begin{cases}
1, & n=1,\\
0, & \text{如果 } n \text{ 有平方因子},\\
(-1)^r, & \text{如果 } n \text{ 有 } r \text{ 个不同素因子}.
\end{cases}
\]
如果 $f(n)$ 是定义在所有正整数 $n$ 上的函数，并且
\[
F(n)=\sum_{d\mid n} f(d),\qquad n=1,2,\ldots,
\]
那么 Möbius 反演公式说，可以由 $F(n)$ 恢复 $f(n)$：
\[
f(n)=\sum_{d\mid n}\mu(d)F\left(\frac{n}{d}\right),\qquad n=1,2,\ldots.
\]
这是一个初等数论结果，我们在此直接采用。定义
\[
\psi(n)=\text{$\mathbb F_p[x]$ 中次数为 $n$ 的不可约多项式的个数}.
\]
在命题 18 中比较次数得到
\[
p^n=\sum_{d\mid n} d\psi(d).
\]
对 $f(n)=n\psi(n)$ 应用 Möbius 反演公式，得到
\[
n\psi(n)=\sum_{d\mid n}\mu(d)p^{n/d},
\]
从而给出 $\mathbb F_p$ 上次数为 $n$ 的不可约多项式的个数公式：
\[
\psi(n)=\frac1n\sum_{d\mid n}\mu(d)p^{n/d}.
\]
例如，在 $p=2,n=4$ 的情形中，
\[
\psi(4)=\frac14[\mu(1)2^4+\mu(2)2^2+\mu(4)2^1]
=\frac14(16-4+0)=3,
\]
与我们上面直接确定的结果一致。

我们已经看到
\[
\mathbb F_{p^m}\subseteq \mathbb F_{p^n}\quad\text{当且仅当}\quad m\mid n.
\]
特别地，给定任意两个有限域 $\mathbb F_{p^{n_1}}$ 和 $\mathbb F_{p^{n_2}}$，存在第三个有限域含有它们（的同构副本），即 $\mathbb F_{p^{n_1n_2}}$。这给出了这些域上的一个偏序，并允许我们考虑它们的并。由于这些域给出了 $\mathbb F_p$ 的全部有限扩张，我们看到所有 $\mathbb F_{p^n}$ 的并是 $\mathbb F_p$ 的一个代数闭包，在同构意义下唯一：
\[
\overline{\mathbb F_p}=\bigcup_{n\geq 1}\mathbb F_{p^n}.
\]
这给出了 $\mathbb F_p$ 的代数闭包的一个简单描述。

\subsection*{练习}

\begin{enumerate}
\item 在 $\mathbb Z[x]$ 和 $\mathbb F_2[x]$ 中把 $x^8-x$ 分解成不可约多项式的乘积。

\item 写出 $\mathbb F_4$ 和 $\mathbb F_9$ 的乘法表。

\item 证明代数闭域必定是无限域。

\item 构造含有 $16$ 个元素的有限域，并找出它的乘法群的一个生成元。有多少个这样的生成元？

\item 给出 $x^3-x+1$ 与 $x^3-x-1$ 在 $\mathbb F_3$ 上的分裂域之间的一个显式同构。

\item 设 $K=\mathbb Q(\theta)=\mathbb Q(\sqrt{D_1},\sqrt{D_2})$，其中 $D_1,D_2\in\mathbb Z$，是一个双二次扩张，并设
\[
\theta=a+b\sqrt{D_1}+c\sqrt{D_2}+d\sqrt{D_1D_2},
\]
其中 $a,b,c,d\in\mathbb Z$。证明 $\theta$ 在 $\mathbb Q$ 上的极小多项式 $m_\theta(x)$ 在 $\mathbb Q$ 上不可约且次数为 $4$，但模每个素数 $p$ 都可约。特别地，证明多项式 $x^4-10x^2+1$ 在 $\mathbb Z[x]$ 中不可约，但模每个素数都可约。[使用有限域上不存在双二次扩张这一事实。]

\item 证明对每个素数 $p$，$2,3,6$ 中至少有一个是 $\mathbb F_p$ 中的平方。推出多项式
\[
x^6-11x^4+36x^2-36=(x^2-2)(x^2-3)(x^2-6)
\]
模每个素数 $p$ 都有根，但在 $\mathbb Z$ 中没有根。

\item 确定多项式 $x^p-x-a$ 在 $\mathbb F_p$ 上的分裂域，其中 $a\neq 0$ 且 $a\in\mathbb F_p$。显式说明其 Galois 群是循环群。[证明 $\alpha\mapsto \alpha+1$ 是一个自同构。] 这样的扩张称为 Artin--Schreier 扩张（参见第 7 节练习 9）。

\item 令 $q=p^m$ 为素数 $p$ 的一个幂，并令 $\mathbb F_q=\mathbb F_{p^m}$ 为含有 $q$ 个元素的有限域。令 $\sigma_q=\sigma_p^m$ 为 Frobenius 自同构 $\sigma_p$ 的第 $m$ 次幂，称为 $q$-Frobenius 自同构。
\begin{enumerate}[label=(\alph*)]
\item 证明 $\sigma_q$ 固定 $\mathbb F_q$。
\item 证明 $\mathbb F_q$ 的任意次数为 $n$ 的有限扩张都是 $x^{q^n}-x$ 在 $\mathbb F_q$ 上的分裂域，因此在同构意义下唯一。
\item 证明 $\mathbb F_q$ 的任意次数为 $n$ 的有限扩张都是循环扩张，并以 $\sigma_q$ 为生成元。
\item 证明 $\mathbb F_q$ 的唯一次数为 $n$ 的扩张的子域与 $n$ 的因子 $d$ 一一对应。
\end{enumerate}

\item 证明 $n$ 整除 $\varphi(p^n-1)$。[观察 $\varphi(p^n-1)$ 是一个阶为 $p^n-1$ 的循环群的自同构群的阶。]

\item 证明 $x^{p^n}-x+1$ 只有在 $n=1$ 或 $n=p=2$ 时才可能在 $\mathbb F_p$ 上不可约。[注意如果 $\alpha$ 是一个根，则对任意 $a\in\mathbb F_{p^n}$，$\alpha+a$ 也是根。证明这蕴含 $\mathbb F_p(\alpha)$ 含有 $\mathbb F_{p^n}$，且 $[\mathbb F_p(\alpha):\mathbb F_{p^n}]=p$。]

\paragraph{Berlekamp 分解算法.} 下面的练习概述 Berlekamp 分解算法，用于分解 $\mathbb F_p[x]$ 中的多项式。这个算法的效率基于用 Euclidean 算法在 $\mathbb F_p[x]$ 中计算最大公因子的效率，以及用矩阵行化简算法解线性方程组的效率。

令 $f(x)\in\mathbb F_p[x]$ 是次数为 $n$ 的首一多项式，并设
\[
f(x)=p_1(x)p_2(x)\cdots p_k(x),
\]
其中 $p_1(x),p_2(x),\ldots,p_k(x)$ 是 $\mathbb F_p[x]$ 中不同的首一不可约多项式的幂。

\item 证明为了把 $f(x)$ 写成 $\mathbb F_p[x]$ 中不可约多项式的乘积，只要确定因子 $p_1(x),\ldots,p_k(x)$ 即可。[如果 $p(x)=q(x)^N\in\mathbb F_p[x]$，其中 $q(x)$ 是首一不可约多项式，说明可通过检查 $p$ 次幂并与导数计算最大公因子来确定 $q(x)$。]

\item 令 $g(x)\in\mathbb F_p[x]$ 为任意次数 $<n$ 的多项式。记 $R(h(x))$ 为 $h(x)$ 除以 $f(x)$ 后的余式。证明下面条件等价：
\begin{enumerate}[label=(\alph*)]
\item $R(g(x^p))=g(x)$。
\item $f(x)$ 整除
\[
[g(x)-0][g(x)-1]\cdots[g(x)-(p-1)].
\]
[使用 $g(x^p)=g(x)^p$ 以及 $x^p-x$ 在 $\mathbb F_p[x]$ 中的分解。]
\item 对 $i=1,2,\ldots,k$，$p_i(x)$ 整除 (b) 中的乘积。
\item 对每个 $i=1,2,\ldots,k$，存在 $s_i\in\mathbb F_p$ 使得 $p_i(x)$ 整除 $g(x)-s_i$，即
\[
g(x)=s_i\pmod{p_i(x)}.
\]
\end{enumerate}

\item 证明满足前一练习中等价条件的次数 $<n$ 的多项式 $g(x)$ 构成 $\mathbb F_p$ 上一个维数为 $k$ 的向量空间 $V$。[使用中国剩余定理处理 13(d) 中 $s_i$ 的 $p^k$ 种可能选择。]

\item 设
\[
g(x)=b_0+b_1x+\cdots+b_{n-1}x^{n-1}\in V.
\]
对 $j=0,1,\ldots,n-1$，令
\[
R(x^{jp})=a_{0,j}+a_{1,j}x+\cdots+a_{n-1,j}x^{n-1},
\]
并令 $A$ 为 $n\times n$ 矩阵
\[
A=\begin{pmatrix}
a_{0,0}&a_{0,1}&\cdots&a_{0,n-1}\\
a_{1,0}&a_{1,1}&\cdots&a_{1,n-1}\\
\vdots&\vdots& &\vdots\\
a_{n-1,0}&a_{n-1,1}&\cdots&a_{n-1,n-1}
\end{pmatrix}.
\]
证明练习 13 中条件 (a) 对 $g(x)\in V$ 来说等价于
\[
(A-I)B=0,
\tag{**}
\]
其中 $B$ 是以 $b_0,b_1,\ldots,b_{n-1}$ 为项的列矩阵。推出矩阵 $A-I$ 的秩为 $n-k$。注意，这已经足以判断 $f(x)$ 是否不可约，而无需实际确定因子。

\item 令 $g_1(x),g_2(x),\ldots,g_k(x)$ 为 (**) 的一组解基（也就是 $V$ 的一组基），其中可取 $g_1(x)=1$。从 $w(x)=f(x)$ 开始，对 $f(x)$ 已经算出的每个因子，计算最大公因子
\[
(w(x),g_i(x)-s),\qquad i=2,3,\ldots,k,\\quad s\in\mathbb F_p.
\]
注意由练习 13(d)，$f(x)$ 的每个因子 $p_i(x)$ 都整除某个这样的最大公因子。当已经确定了 $k$ 个两两互素的因子时，过程终止。

证明这个过程确实给出所有因子 $p_1(x),p_2(x),\ldots,p_k(x)$，也就是说，可以由这个过程分离出各个因子 $p_1(x),p_2(x),\ldots,p_k(x)$。证明如下：如果并非如此，则对两个因子，比如 $p_1(x)$ 与 $p_2(x)$，对每个 $i=1,2,\ldots,k$ 都存在 $s_i\in\mathbb F_p$，使得 $g_i(x)-s_i$ 同时被 $p_1(x)$ 和 $p_2(x)$ 整除。由中国剩余定理，选择 $g(x)\in V$ 满足
\[
g(x)=0\pmod{p_1(x)},\qquad g(x)=1\pmod{p_2(x)}.
\]
把 $g(x)=\sum_{i=1}^k c_i g_i(x)$ 写成 $V$ 的基的线性组合，并令 $s=\sum_{i=1}^k c_i s_i\in\mathbb F_p$。证明 $s=0\pmod{p_1(x)}$，从而 $s=0$；而 $s=1\pmod{p_2(x)}$，从而 $s=1$，矛盾。

\item 本练习按照前面练习概述的 Berlekamp 分解算法，确定
\[
f(x)=x^5+x^2+4x+6
\]
在 $\mathbb F_7[x]$ 中的分解。
\begin{enumerate}[label=(\alph*)]
\item 证明
\[
x^7=x^2+3x^3+6x^4\pmod{f(x)}.
\]
类似地计算 $x^{14},x^{21}$ 和 $x^{28}$ 模 $f(x)$（注意 $x^{14}$ 最容易通过把 $x^7$ 的结果平方再约化来计算，等等），并说明此时练习 15 中的矩阵 $A$ 为
\[
A=\begin{pmatrix}
0&6&0&4&2\\
0&0&3&0&2\\
0&1&6&3&1\\
0&3&4&3&1\\
0&6&1&4&1
\end{pmatrix}.
\]
\item 说明 $A-I$ 的简化行阶梯形矩阵为
\[
\begin{pmatrix}
1&0&0&6&6\\
0&1&0&1&0\\
0&0&1&1&6\\
0&0&0&0&0\\
0&0&0&0&0
\end{pmatrix}.
\]
推出 $k=2$（所以 $f(x)$ 正好是两个因子的乘积，每个因子都是不可约多项式的幂），并且 $g_1(x)=1$ 与
\[
g_2(x)=x^4+5x^3+x^2+x
\]
给出练习 15 中 (**) 的解空间的一组基。
\item 按照练习 16 中的过程，证明
\[
(f(x),g_2(x)-1)=x^2+3x+5=p_1(x),
\]
且 $f(x)/p_1(x)=x^3+4x^2+4x+4=p_2(x)$，这给出 $\mathbb F_7[x]$ 中整除 $f(x)$ 的不可约多项式的幂。说明这两个因子都不是 $\mathbb F_7[x]$ 中的 $7$ 次幂，并且每个都与它的导数互素，从而推出二者都是不可约多项式，得到 $f(x)$ 到不可约多项式的完整分解
\[
f(x)=(x^2+3x+5)(x^3+4x^2+4x+4)\in\mathbb F_7[x].
\]
\end{enumerate}
\end{enumerate}

\section{复合扩张与单扩张}

现在我们考虑同 Galois 扩张取复合的效果。第一个结果说明，把一个 Galois 扩张“上滑”仍得到一个 Galois 扩张。

\noindent\textbf{命题 19.} 设 $K/F$ 是 Galois 扩张，$F'/F$ 是任意扩张。则 $KF'/F'$ 是 Galois 扩张，其 Galois 群
\[
\operatorname{Gal}(KF'/F')\cong \operatorname{Gal}(K/K\cap F')
\]
同构于 $\operatorname{Gal}(K/F)$ 的一个子群。图示为：
\[
\begin{tikzpicture}[scale=1.0, baseline=(current bounding box.center), every node/.style={inner sep=2pt}]
\node (top) at (0,2.1) {$KF'$};
\node (left) at (-1.6,0.7) {$K$};
\node (right) at (1.6,0.7) {$F'$};
\node (mid) at (0,-0.8) {$K\cap F'$};
\node (bot) at (0,-2.1) {$F$};
\draw (top)--(left);
\draw (top)--(right);
\draw (left)--(mid);
\draw (right)--(mid);
\draw (mid)--(bot);
\end{tikzpicture}
\]

\begin{proof}
如果 $K/F$ 是 Galois 的，则 $K$ 是 $F[x]$ 中某个可分多项式 $f(x)$ 的分裂域。于是 $KF'/F'$ 是把 $f(x)$ 视作 $F'[x]$ 中多项式时的分裂域，因此这个扩张是 Galois 的。由于 $K/F$ 是 Galois 的，固定 $F$ 的 $K$ 的每个嵌入都是 $K$ 的自同构，所以映射
\[
\varphi:\operatorname{Gal}(KF'/F')\longrightarrow \operatorname{Gal}(K/F),
\qquad
\sigma\longmapsto \sigma|_K
\]
由把自同构 $\sigma$ 限制到子域 $K$ 来定义，并且是良定义的。它显然是同态，其核为
\[
\ker\varphi=\{\sigma\in\operatorname{Gal}(KF'/F')\mid \sigma|_K=1\}.
\]
由于 $\operatorname{Gal}(KF'/F')$ 中的元素在 $F'$ 上平凡，核中的元素在 $K$ 和 $F'$ 上都平凡，因此在它们的复合上也平凡，所以核只含恒等自同构。因此 $\varphi$ 是单射。

令 $H$ 表示 $\varphi$ 在 $\operatorname{Gal}(K/F)$ 中的像，令 $K_H$ 表示 $K$ 中对应的、含有 $F$ 的固定子域。由于 $H$ 中的每个元素都固定 $F'$，$K_H$ 含有 $K\cap F'$。另一方面，复合 $K_HF'$ 被 $\operatorname{Gal}(KF'/F')$ 固定（任意 $\sigma\in\operatorname{Gal}(KF'/F')$ 固定 $F'$，并且在 $K_H\subseteq K$ 上按它的限制 $\sigma|_K\in H$ 作用，而后者按定义固定 $K_H$）。由基本定理可知 $K_HF'=F'$，所以 $K_H\subseteq F'$，这给出反向包含 $K_H\subseteq K\cap F'$。因此 $K_H=K\cap F'$，再由基本定理得到 $H=\operatorname{Gal}(K/K\cap F')$，证明完成。
\end{proof}

\noindent\textbf{推论 20.} 设 $K/F$ 是 Galois 扩张，$F'/F$ 是任意有限扩张。则
\[
[KF':F]=\frac{[K:F][F':F]}{[K\cap F':F]}.
\]

\begin{proof}
这由命题以及等式 $[KF':F']=[K:K\cap F']$ 得出，而这个等式由命题中 Galois 群的阶给出。
\end{proof}

例子 $F=\mathbb Q$，$K=\mathbb Q(\sqrt[3]{2})$，$F'=\mathbb Q(\rho\sqrt[3]{2})$，其中 $\rho$ 是一个本原三次单位根，说明如果两个扩张中没有一个是 Galois 的，那么推论 20 的公式一般不成立。

\noindent\textbf{命题 21.} 令 $K_1$ 和 $K_2$ 为域 $F$ 的 Galois 扩张。则
\begin{enumerate}[label=(\arabic*)]
\item 交 $K_1\cap K_2$ 在 $F$ 上是 Galois 的；
\item 复合 $K_1K_2$ 在 $F$ 上是 Galois 的。其 Galois 群同构于直积
\[
\operatorname{Gal}(K_1/F)\times\operatorname{Gal}(K_2/F)
\]
的子群
\[
H=\{(\sigma,\tau)\mid \sigma|_{K_1\cap K_2}=\tau|_{K_1\cap K_2}\},
\]
也就是限制到交 $K_1\cap K_2$ 上相等的那些元素组成的子群。
\end{enumerate}
\[
\begin{tikzpicture}[scale=1.0, baseline=(current bounding box.center), every node/.style={inner sep=2pt}]
\node (top) at (0,2.0) {$K_1K_2$};
\node (left) at (-1.5,0.7) {$K_1$};
\node (right) at (1.5,0.7) {$K_2$};
\node (mid) at (0,-0.8) {$K_1\cap K_2$};
\node (bot) at (0,-2.1) {$F$};
\draw (top)--(left);
\draw (top)--(right);
\draw (left)--(mid);
\draw (right)--(mid);
\draw (mid)--(bot);
\end{tikzpicture}
\]

\begin{proof}
(1) 设 $p(x)$ 是 $F[x]$ 中一个不可约多项式，并且在 $K_1\cap K_2$ 中有一个根 $\alpha$。由于 $\alpha\in K_1$ 且 $K_1/F$ 是 Galois 的，$p(x)$ 的所有根都在 $K_1$ 中。同理，所有根都在 $K_2$ 中，因此 $p(x)$ 的所有根都在 $K_1\cap K_2$ 中。由定理 13 中的判别可以容易地推出 $K_1\cap K_2$ 是 Galois 的。

(2) 如果 $K_1$ 是可分多项式 $f_1(x)$ 的分裂域，而 $K_2$ 是可分多项式 $f_2(x)$ 的分裂域，那么复合就是多项式 $f_1(x)f_2(x)$ 的无平方因子部分的分裂域，因此在 $F$ 上是 Galois 的。

映射
\[
\varphi:\operatorname{Gal}(K_1K_2/F)\longrightarrow
\operatorname{Gal}(K_1/F)\times\operatorname{Gal}(K_2/F),
\qquad
\sigma\longmapsto (\sigma|_{K_1},\sigma|_{K_2})
\]
显然是同态。其核由同时在 $K_1$ 与 $K_2$ 上平凡的元素 $\sigma$ 组成，因此在复合上平凡，所以这个映射是单射。它的像位于子群 $H$ 中，因为
\[
(\sigma|_{K_1})|_{K_1\cap K_2}=\sigma|_{K_1\cap K_2}=(\sigma|_{K_2})|_{K_1\cap K_2}.
\]

子群 $H$ 的阶可以这样计算：对每个 $\sigma\in\operatorname{Gal}(K_1/F)$，有 $|\operatorname{Gal}(K_2/K_1\cap K_2)|$ 个 $\tau\in\operatorname{Gal}(K_2/F)$，其在 $K_1\cap K_2$ 上的限制等于 $\sigma|_{K_1\cap K_2}$。因此
\[
\begin{aligned}
|H|&=|\operatorname{Gal}(K_1/F)|\cdot |\operatorname{Gal}(K_2/K_1\cap K_2)|\\
&=|\operatorname{Gal}(K_1/F)|\frac{|\operatorname{Gal}(K_2/F)|}{|\operatorname{Gal}(K_1\cap K_2/F)|}.
\end{aligned}
\]
由推论 20 和上图可知，$H$ 与 $\operatorname{Gal}(K_1K_2/F)$ 的阶都等于
\[
[K_1K_2:F]=\frac{[K_1:F][K_2:F]}{[K_1\cap K_2:F]}.
\]
因此 $\varphi$ 的像恰好就是 $H$，证明完成。
\end{proof}

\noindent\textbf{推论 22.} 令 $K_1$ 和 $K_2$ 为域 $F$ 的 Galois 扩张，并且 $K_1\cap K_2=F$。则
\[
\operatorname{Gal}(K_1K_2/F)\cong \operatorname{Gal}(K_1/F)\times\operatorname{Gal}(K_2/F).
\]
反过来，如果 $K$ 在 $F$ 上是 Galois 的，且
\[
G=\operatorname{Gal}(K/F)=G_1\times G_2
\]
是两个子群 $G_1$ 和 $G_2$ 的直积，那么 $K$ 是 $F$ 的两个 Galois 扩张 $K_1$ 和 $K_2$ 的复合，并且 $K_1\cap K_2=F$。

\begin{proof}
第一部分立即由命题推出。对于第二部分，令 $K_1$ 为 $G_1\subset G$ 的固定域，令 $K_2$ 为 $G_2\subset G$ 的固定域。那么 $K_1\cap K_2$ 对应于子群 $G_1G_2$，在此情形中它就是整个 $G$，所以 $K_1\cap K_2=F$。复合 $K_1K_2$ 对应于子群 $G_1\cap G_2$，这里它是恒等子群，所以 $K_1K_2=K$，证明完成。
\end{proof}

\noindent\textbf{推论 23.} 令 $E/F$ 为任意有限可分扩张。则 $E$ 包含在某个在 $F$ 上 Galois 的扩张 $K$ 中，并且 $K$ 是极小的，意思是在 $K$ 的一个固定代数闭包中，任何含有 $E$ 的 $F$ 的 Galois 扩张都含有 $K$。

\begin{proof}
存在一个含有 $E$ 的 $F$ 的 Galois 扩张，例如取 $E$ 在 $F$ 上一组基的极小多项式的分裂域的复合（由于 $E$ 在 $F$ 上可分，这些多项式都是可分的）。于是所有含有 $E$ 的 $F$ 的 Galois 扩张的交就是这个域 $K$。
\end{proof}

\noindent\textbf{定义.} 前一推论中含有 $E$ 的 $F$ 的 Galois 扩张 $K$ 称为 $E$ 在 $F$ 上的 \emph{Galois 闭包}。

在 Galois 扩张中工作通常更简单（例如在按推论 20 计算次数时）。可分扩张的 Galois 闭包的存在性，经常可用来把计算化归为对 Galois 扩张的考虑。

回忆：如果 $K=F(\theta)$ 对某个元素 $\theta$ 成立，则称 $F$ 的扩张 $K$ 是\emph{单}的，在这种情形下 $\theta$ 称为 $K$ 的一个本原元。

\noindent\textbf{命题 24.} 令 $K/F$ 为有限扩张。则 $K=F(\theta)$ 当且仅当含有 $F$ 的 $K$ 的子域只有有限多个。

\begin{proof}
首先假设 $K=F(\theta)$ 是单的。令 $E$ 为 $K$ 的一个含有 $F$ 的子域：$F\subseteq E\subseteq K$。令 $f(x)\in F[x]$ 为 $\theta$ 在 $F$ 上的极小多项式，令 $g(x)\in E[x]$ 为 $\theta$ 在 $E$ 上的极小多项式。则 $g(x)$ 在 $E[x]$ 中整除 $f(x)$。令 $E'$ 为 $g(x)$ 的系数在 $F$ 上生成的域。那么 $E'\subseteq E$，并且显然 $\theta$ 在 $E'$ 上的极小多项式仍为 $g(x)$。于是
\[
[K:E]=\deg g(x)=[K:E']
\]
蕴含 $E=E'$。因此，含有 $F$ 的 $K$ 的子域是由 $f(x)$ 的首一因子的系数生成的子域，从而这种子域只有有限多个。

反过来，假设含有 $F$ 的 $K$ 的子域只有有限多个。如果 $F$ 是有限域，那么我们已经看到 $K$ 是一个单扩张（命题 17）。因此可以假设 $F$ 是无限的。因为 $K$ 在 $F$ 上有限生成，显然只需证明 $F(\alpha,\beta)$ 由一个元素生成。考虑子域
\[
F(\alpha+c\beta),\qquad c\in F.
\]
由于 $F$ 中 $c$ 有无限多个选择，而这种子域只有有限多个，存在 $F$ 中的 $c,c'$，$c\neq c'$，使得
\[
F(\alpha+c\beta)=F(\alpha+c'\beta).
\]
于是 $\alpha+c\beta$ 和 $\alpha+c'\beta$ 都属于 $F(\alpha+c\beta)$，取它们的差可知 $(c-c')\beta\in F(\alpha+c\beta)$。因此 $\beta\in F(\alpha+c\beta)$，于是也有 $\alpha\in F(\alpha+c\beta)$。所以
\[
F(\alpha,\beta)\subseteq F(\alpha+c\beta),
\]
而反向包含显然成立，于是
\[
F(\alpha,\beta)=F(\alpha+c\beta),
\]
证明完成。
\end{proof}

\noindent\textbf{定理 25（本原元定理）.} 如果 $K/F$ 是有限且可分的，则 $K/F$ 是单扩张。特别地，特征为 $0$ 的域的任意有限扩张都是单扩张。

\begin{proof}
令 $L$ 为 $K$ 在 $F$ 上的 Galois 闭包。则含有 $F$ 的 $K$ 的任意子域，都由基本定理对应于 Galois 群 $\operatorname{Gal}(L/F)$ 的一个子群。由于这样的子群只有有限多个，前一命题说明 $K/F$ 是单的。最后一句话则由特征为 $0$ 的域的任意有限扩张都是可分的推出。
\end{proof}

如命题的证明所表明的，一个扩张的本原元可以取为该扩张的生成元的一个简单线性组合。在 Galois 扩张的情形中，只需确定一个不被 Galois 群的任何非平凡元素固定的线性组合，因为此时由基本定理，这个线性组合不可能属于任何真子域。

\paragraph{例}
\begin{enumerate}[label=(\arabic*)]
\item 元素 $\sqrt2+\sqrt3$ 生成域 $\mathbb Q(\sqrt2,\sqrt3)$，正如我们已经看到的那样（它不被这个域的四个 Galois 自同构中的任何一个固定）。
\item 有理函数域 $\overline{\mathbb F}_p(x,y)$，其中变量 $x$ 与 $y$ 取系数于 $\mathbb F_p$ 的代数闭包 $\overline{\mathbb F}_p$，不是子域
\[
F=\overline{\mathbb F}_p(x^p,y^p)
\]
的一个单扩张。容易看出
\[
[\overline{\mathbb F}_p(x,y):\overline{\mathbb F}_p(x^p,y^p)]=p^2,
\]
并且子域
\[
F(x+cy),\qquad c\in\overline{\mathbb F}_p
\]
都在 $\overline{\mathbb F}_p(x^p,y^p)$ 上次数为 $p$（注意 $(x+cy)^p=x^p+c^py^p\in\overline{\mathbb F}_p(x^p,y^p)$）。如果其中任意两个子域相等，那么正如命题 24 的证明中一样，我们会得到
\[
\overline{\mathbb F}_p(x,y)=F(x+cy),
\]
这由次数考虑是不可能的。因此存在无限多个这样的子域，并且这个扩张不是单扩张。
\end{enumerate}

\subsection*{练习}
\begin{enumerate}[label=\arabic*.]
\item 求域 $\mathbb Q(\sqrt{1+\sqrt2})$ 在 $\mathbb Q$ 上的 Galois 闭包。

\item 找出 $\mathbb Q(\sqrt2,\sqrt3,\sqrt5)$ 在 $\mathbb Q$ 上的一个本原生成元。

\item 令 $F$ 为包含在 $\mathbb Q$ 上 $n\times n$ 矩阵环中的一个域。证明 $[F:\mathbb Q]\leq n$。（注意，由第 13.2 节练习 19，$\mathbb Q$ 上 $n\times n$ 矩阵环确实含有次数为 $n$ 的 $\mathbb Q$ 上域。）

\item 设 $f(x)\in F[x]$ 是域 $F$ 上次数为 $n$ 的不可约多项式，$L$ 为 $f(x)$ 在 $F$ 上的分裂域，$\alpha$ 为 $f(x)$ 在 $L$ 中的一个根。如果 $K$ 是 $F$ 的任意 Galois 扩张，证明 $f(x)$ 在 $K$ 上分解为 $m$ 个次数均为 $d$ 的不可约多项式的乘积，其中
\[
d=[K(\alpha):K]=[(L\cap K)(\alpha):L\cap K],
\]
并且
\[
m=n/d=[F(\alpha)\cap K:F].
\]
[先证明 $f(x)$ 在 $K$ 上的分解与它在 $L\cap K$ 上的分解相同。然后，若 $H$ 是 $L$ 在 $F$ 上的 Galois 群中对应于 $L\cap K$ 的子群，则 $f(x)$ 在 $L\cap K$ 上的因子对应于 $H$ 在 $f(x)$ 的根上的轨道。使用第 4.1 节练习 9。]

\item 令 $p$ 为素数，$F$ 为一个域。令 $K$ 为 $F$ 的一个 Galois 扩张，其 Galois 群是一个 $p$-群（即次数 $[K:F]$ 是 $p$ 的幂）。这样的扩张称为一个 $p$-扩张（注意，$p$-扩张按定义是 Galois 的）。
\begin{enumerate}[label=(\alph*)]
\item 设 $L$ 是 $K$ 的一个 $p$-扩张。证明 $L$ 在 $F$ 上的 Galois 闭包是 $F$ 的一个 $p$-扩张。
\item 举例说明：如果 $[K:F]$ 是 $p$ 的幂但 $K/F$ 不是 Galois 的，则 (a) 不必成立。
\end{enumerate}

\item 通过明确展示无限多个中间子域，证明 $\overline{\mathbb F}_p(x,y)/\overline{\mathbb F}_p(x^p,y^p)$ 不是单扩张。

\item 设 $F\subseteq K\subseteq L$，且 $\theta\in L$，其中 $p(x)=m_{\theta,F}(x)$。证明作为 $K$-代数，
\[
K\otimes_F F(\theta)\cong K[x]/(p(x)).
\]

\item 设 $K_1$ 与 $K_2$ 是域 $F$ 的两个代数扩张，它们都包含在特征为零的域 $L$ 中。证明 $F$-代数 $K_1\otimes_F K_2$ 没有非零幂零元。[使用前一练习。]
\end{enumerate}

\section{分圆扩张与 $\mathbb Q$ 上的 Abel 扩张}

我们已经确定，$n$ 次单位根的分圆域 $\mathbb Q(\zeta_n)$ 是 $\mathbb Q$ 的一个 Galois 扩张，其次数为 $\varphi(n)$，其中 $\varphi$ 表示 Euler $\varphi$-函数。这个域的任一自同构都由它在本原 $n$ 次单位根 $\zeta_n$ 上的作用唯一决定。这个元素必须被映到另一个本原 $n$ 次单位根（回忆这些根就是不可约分圆多项式 $\Phi_n(x)$ 的根）。因此
\[
\sigma(\zeta_n)=\zeta_n^a
\]
对某个整数 $a$ 成立，其中 $1\leq a<n$，且 $a$ 与 $n$ 互素。由于这样的整数 $a$ 恰有 $\varphi(n)$ 个，于是事实上这些映射中的每一个都确实是 $\mathbb Q(\zeta_n)$ 的一个自同构。还要注意，对任意与 $n$ 互素的整数 $a$，我们都可以用同一个公式定义 $\sigma_a$，并且 $\sigma_a$ 只依赖于 $a$ 模 $n$ 的剩余类。

\noindent\textbf{定理 26.} $n$ 次单位根的分圆域 $\mathbb Q(\zeta_n)$ 的 Galois 群同构于乘法群 $(\mathbb Z/n\mathbb Z)^\times$。这个同构明确由映射
\[
(\mathbb Z/n\mathbb Z)^\times\longrightarrow \operatorname{Gal}(\mathbb Q(\zeta_n)/\mathbb Q),
\qquad
 a\pmod n\longmapsto \sigma_a
\]
给出，其中 $\sigma_a$ 是由
\[
\sigma_a(\zeta_n)=\zeta_n^a
\]
定义的自同构。

\begin{proof}
上面的讨论表明，对任意 $a\pmod n$，$\sigma_a$ 都是一个自同构，所以此映射良定义。它是同态，因为
\[
(\sigma_a\sigma_b)(\zeta_n)=\sigma_a(\zeta_n^b)=(\zeta_n^a)^b=\zeta_n^{ab},
\]
这说明 $\sigma_a\sigma_b=\sigma_{ab}$。由上面的讨论，这个映射是双射，因为我们知道每个 Galois 自同构都形如 $\sigma_a$，其中 $a\pmod n$ 唯一确定。因此这个映射是同构。
\end{proof}

\paragraph{例}
\begin{enumerate}[label=(\arabic*)]
\item 域 $\mathbb Q(\zeta_5)$ 在 $\mathbb Q$ 上是 Galois 的，其 Galois 群
\[
(\mathbb Z/5\mathbb Z)^\times\cong \mathbb Z/4\mathbb Z.
\]
这是我们第一个次数为 $4$ 且 Galois 群为循环群的 $\mathbb Q$ 的 Galois 扩张例子。按照上面的记号，Galois 群的元素是 $\{\sigma_1=1,\sigma_2,\sigma_3,\sigma_4\}$。这个循环群的一个生成元是
\[
\sigma_2:\zeta_5\longmapsto \zeta_5^2
\]
（因为 $2$ 在 $(\mathbb Z/5\mathbb Z)^\times$ 中有阶 $4$）。

这里只有一个非平凡子域，它是 $\mathbb Q$ 的一个二次扩张，即子群 $\{1,\sigma_4=\sigma_{-1}\}$ 的固定域。这个子域中的一个元素由
\[
\alpha=\zeta_5+\sigma_{-1}\zeta_5=\zeta_5+\zeta_5^{-1}
\]
给出，因为这个元素显然被 $\sigma_{-1}$ 固定。元素 $\zeta_5$ 满足
\[
\zeta_5^4+\zeta_5^3+\zeta_5^2+\zeta_5+1=0.
\]
于是注意到
\[
\begin{aligned}
\alpha^2+\alpha-1
&=(\zeta_5^2+2+\zeta_5^{-2})+(\zeta_5+\zeta_5^{-1})-1\\
&=\zeta_5^2+2+\zeta_5^3+\zeta_5+\zeta_5^4-1=0.
\end{aligned}
\]
显式解出 $\alpha$，可知由 $\alpha$ 生成的 $\mathbb Q$ 的二次扩张就是 $\mathbb Q(\sqrt 5)$：
\[
\mathbb Q(\zeta_5+\zeta_5^{-1})=\mathbb Q(\sqrt 5).
\]
一般可以证明（这并不完全平凡）：若 $p$ 是奇素数，则域 $\mathbb Q(\zeta_p)$ 含有二次域 $\mathbb Q(\sqrt{\pm p})$，其中当 $p\equiv 1\pmod 4$ 时取正号，当 $p\equiv 3\pmod 4$ 时取负号（参见第 13.6 节中的练习）。

\item 对 $\mathbb Q(\zeta_{13})$，并且更一般地，对奇素数 $p$，我们可以像前一个例子那样为 $\mathbb Q(\zeta_p)$ 的任意子域构造一个本原元。$\mathbb Q(\zeta_p)$ 在 $\mathbb Q$ 上的一组基为
\[
1,\zeta_p,\zeta_p^2,\ldots,\zeta_p^{p-2}.
\]
由于
\[
\zeta_p^{p-1}+\zeta_p^{p-2}+\cdots+\zeta_p+1=0,
\]
我们看到这些元素
\[
\zeta_p,\zeta_p^2,\ldots,\zeta_p^{p-2},\zeta_p^{p-1}
\]
也构成一组基。选择这组基的原因在于，Galois 群 $\operatorname{Gal}(\mathbb Q(\zeta_p)/\mathbb Q)$ 中的任意 $\sigma$ 都只是置换这些基元素，因为它们正是本原 $p$ 次单位根。注意，正是在这一点上我们需要 $p$ 是素数；一般地，本原 $n$ 次单位根并不构成 $n$ 次单位根分圆域在 $\mathbb Q$ 上的一组基（例如，本原 $4$ 次单位根 $\pm i$ 就不是线性无关的）。

令 $H$ 为 $\mathbb Q(\zeta_p)$ 在 $\mathbb Q$ 上的 Galois 群的任意子群，并令
\[
\alpha_H=\sum_{\sigma\in H}\sigma\zeta_p, \tag{14.10}
\]
即 $\zeta_p$ 在 $H$ 的元素下的共轭之和。对任意 $\tau\in H$，当 $\sigma$ 遍历 $H$ 时，元素 $\tau\sigma$ 也遍历 $H$。因此 $\tau\alpha=\alpha$，所以 $\alpha$ 位于 $H$ 的固定域中。若现在 $\tau$ 不是 $H$ 的元素，那么 $\tau\alpha$ 是若干基元素之和（回忆这里任意自同构都置换这些基元素），其中一项是 $\tau(\zeta_p)$。若有 $\tau\alpha=\alpha$，则由于这些元素构成一组基，我们必须有 $\tau(\zeta_p)=\sigma(\zeta_p)$，其中 $\sigma\zeta_p$ 是式 (10) 中的一项。但这意味着 $\tau\sigma^{-1}=1$，因为这个自同构在 $\zeta_p$ 上是恒等的。于是 $\tau=\sigma\in H$，矛盾。这说明 $\alpha$ 不被任何不属于 $H$ 的自同构固定，所以 $\mathbb Q(\alpha)$ 正是 $H$ 的固定域。

作为一个具体例子，考虑 $\mathbb Q(\zeta_{13})$ 的子域；它们对应于
\[
(\mathbb Z/13\mathbb Z)^\times\cong \mathbb Z/12\mathbb Z
\]
的子群。这个循环群的一个生成元是自同构 $\sigma=\sigma_2$，它把 $\zeta_{13}$ 映到 $\zeta_{13}^2$。非平凡子群对应于 $12$ 的非平凡因子，因此分别有阶 $2,3,4,6$，其生成元分别为 $\sigma^6,\sigma^4,\sigma^3,\sigma^2$。对应的固定域在 $\mathbb Q$ 上的次数分别为 $6,4,3,2$。令 $\zeta=\zeta_{13}$，这些固定域的生成元由
\[
\zeta+\sigma^6\zeta=\zeta+\zeta^{26}=\zeta+\zeta^{-1},
\]
\[
\zeta+\sigma^4\zeta+\sigma^8\zeta=\zeta+\zeta^{24}+\zeta^{28}=\zeta+\zeta^3+\zeta^9,
\]
\[
\zeta+\sigma^3\zeta+\sigma^6\zeta+\sigma^9\zeta=\zeta+\zeta^8+\zeta^{12}+\zeta^5,
\]
以及
\[
\zeta+\sigma^2\zeta+\sigma^4\zeta+\sigma^6\zeta+\sigma^8\zeta+\sigma^{10}\zeta
=\zeta+\zeta^4+\zeta^3+\zeta^{12}+\zeta^9+\zeta^{10}
\]
给出。这个扩张的子域格如下：
\[
\begin{tikzpicture}[scale=1.0, every node/.style={inner sep=2pt}, baseline=(current bounding box.center)]
\node (top) at (0,3.6) {$\mathbb Q(\zeta)$};
\node (mid) at (0,2.2) {$\mathbb Q(\zeta+\zeta^{-1})$};
\node (left) at (-3.0,1.4) {$\mathbb Q(\zeta+\zeta^3+\zeta^9)$};
\node (right) at (3.0,1.0) {$\mathbb Q(\zeta+\zeta^5+\zeta^8+\zeta^{12})$};
\node (low) at (0,-0.3) {$\mathbb Q(\zeta+\zeta^3+\zeta^4+\zeta^9+\zeta^{10}+\zeta^{12})$};
\node (bot) at (3.0,-1.8) {$\mathbb Q$};
\draw (top)--node[left] {$2$} (mid);
\draw (top)--node[above left] {$3$} (left);
\draw (mid)--node[above] {$2$} (right);
\draw (mid)--node[right] {$3$} (low);
\draw (left)--node[below left] {$2$} (low);
\draw (right)--node[right] {$3$} (bot);
\draw (low)--node[below] {$2$} (bot);
\end{tikzpicture}
\]
式 (10) 中构造的元素以及它们的共轭称为 $\zeta$ 的\emph{周期}，它们在分圆域的算术研究中很有用。对它们的组合性质的研究称为\emph{分圆法}。
\end{enumerate}

设
\[
n=p_1^{a_1}p_2^{a_2}\cdots p_k^{a_k}
\]
是 $n$ 分解为互异素数幂的分解。由于
\[
\zeta_{p_2^{a_2}\cdots p_k^{a_k}}
\]
是一个本原 $p_1^{a_1}$ 次单位根，域 $K_1=\mathbb Q(\zeta_{p_1^{a_1}})$ 是 $\mathbb Q(\zeta_n)$ 的一个子域。类似地，每个域
\[
K_i=\mathbb Q(\zeta_{p_i^{a_i}}),\qquad i=1,2,\ldots,k
\]
都是 $\mathbb Q(\zeta_n)$ 的子域。这些域的复合含有乘积
\[
\zeta_{p_1^{a_1}}\zeta_{p_2^{a_2}}\cdots \zeta_{p_k^{a_k}},
\]
它是一个本原 $n$ 次单位根，因此复合域就是 $\mathbb Q(\zeta_n)$。由于扩张次数 $[K_i:\mathbb Q]$ 等于 $\varphi(p_i^{a_i})$，$i=1,2,\ldots,k$，并且
\[
\varphi(n)=\varphi(p_1^{a_1})\varphi(p_2^{a_2})\cdots\varphi(p_k^{a_k}),
\]
这些域 $K_i$ 的复合的次数正好是各 $K_i$ 的次数的乘积。由命题 21（以及从命题中两个域的情形到这里 $k$ 个域的简单归纳）可知所有这些域的交恰为 $\mathbb Q$。于是推论 22 表明 $\mathbb Q(\zeta_n)$ 的 Galois 群是这些子域 $K_i$ 在 $\mathbb Q$ 上的 Galois 群的直积。我们把这总结为下面的推论。

\noindent\textbf{推论 27.} 设
\[
n=p_1^{a_1}p_2^{a_2}\cdots p_k^{a_k}
\]
是正整数 $n$ 分解为互异素数幂的分解。则分圆域 $\mathbb Q(\zeta_{p_i^{a_i}})$，$i=1,2,\ldots,k$，只在域 $\mathbb Q$ 中相交，并且它们的复合是分圆域 $\mathbb Q(\zeta_n)$。我们有
\[
\operatorname{Gal}(\mathbb Q(\zeta_n)/\mathbb Q)
\cong
\operatorname{Gal}(\mathbb Q(\zeta_{p_1^{a_1}})/\mathbb Q)	imes
\operatorname{Gal}(\mathbb Q(\zeta_{p_2^{a_2}})/\mathbb Q)	imes\cdots\times
\operatorname{Gal}(\mathbb Q(\zeta_{p_k^{a_k}})/\mathbb Q),
\]
而在定理 26 的同构下，这就是中国剩余定理：
\[
(\mathbb Z/n\mathbb Z)^\times
\cong
(\mathbb Z/p_1^{a_1}\mathbb Z)^\times\times
(\mathbb Z/p_2^{a_2}\mathbb Z)^\times\times\cdots\times
(\mathbb Z/p_k^{a_k}\mathbb Z)^\times.
\]

\begin{proof}
唯一尚未证明的断言，是把 Galois 群的这个同构与中国剩余定理在群 $(\mathbb Z/n\mathbb Z)^\times$ 上的陈述相识别；这很简单，留给练习。
\end{proof}

由定理 26，Galois 群 $\operatorname{Gal}(\mathbb Q(\zeta_n)/\mathbb Q)$ 特别是一个 Abel 群。

\noindent\textbf{定义.} 扩张 $K/F$ 称为一个 \emph{Abel 扩张}，如果 $K/F$ 是 Galois 的，并且 $\operatorname{Gal}(K/F)$ 是一个 Abel 群。

由于 Abel 群的所有子群和商群仍然是 Abel 群，由 Galois 理论基本定理可知，一个 $F$ 的 Abel 扩张中任意包含 $F$ 的子域又是 $F$ 的 Abel 扩张。由上一节关于复合扩张的结果，我们还看到 Abel 扩张的复合仍然是 Abel 扩张（因为复合的 Galois 群同构于这些 Galois 群的直积的一个子群，因而是 Abel 群）。

判定哪些群可以作为 $\mathbb Q$ 的 Galois 扩张的 Galois 群，是一个开放问题。利用上面的结果，我们可以看到，每个 Abel 群都作为 $\mathbb Q$ 的某个扩张的 Galois 群出现；事实上，它作为某个分圆域的子域的 Galois 群出现。

令 $n=p_1p_2\cdots p_k$ 是互异素数的乘积。于是由中国剩余定理，
\[
(\mathbb Z/n\mathbb Z)^\times
\cong
(\mathbb Z/p_1\mathbb Z)^\times\times(\mathbb Z/p_2\mathbb Z)^\times\times\cdots\times(\mathbb Z/p_k\mathbb Z)^\times
\cong Z_{p_1-1}\times Z_{p_2-1}\times\cdots\times Z_{p_k-1}. \tag{14.11}
\]

现在设 $G$ 为任意有限 Abel 群。由有限生成 Abel 群基本定理，
\[
G\cong Z_{n_1}\times Z_{n_2}\times\cdots\times Z_{n_k}
\]
对某些整数 $n_1,n_2,\ldots,n_k$ 成立。我们把下列事实看作已知：给定任意整数 $m$，存在无穷多个素数 $p$ 满足 $p\equiv 1\pmod m$（参见第 13.6 节后的练习，它用分圆多项式给出一个证明）。利用这个结果，选择互异素数 $p_1,p_2,\ldots,p_k$，使得
\[
p_1\equiv 1\pmod {n_1},
\]
\[
p_2\equiv 1\pmod {n_2},
\]
\[
\vdots
\]
\[
p_k\equiv 1\pmod {n_k},
\]
并令 $n=p_1p_2\cdots p_k$，如上所述。

由构造可知，$n_i$ 整除 $p_i-1$，$i=1,2,\ldots,k$，所以群 $Z_{p_i-1}$ 有一个阶为 $(p_i-1)/n_i$ 的子群 $H_i$，且它对这个子群的商是阶为 $n_i$ 的循环群。因此，式 (11) 中 $(\mathbb Z/n\mathbb Z)^\times$ 对子群
\[
H_1\times H_2\times\cdots\times H_k
\]
的商同构于群 $G$。

由定理 26 和 Galois 理论基本定理，我们看到 $\mathbb Q(\zeta_{p_1p_2\cdots p_k})$ 有一个在 $\mathbb Q$ 上 Galois 的子域，其 Galois 群为 $G$。我们把这一点总结为下面的推论。

\noindent\textbf{推论 28.} 设 $G$ 为任意有限 Abel 群。则存在某个分圆域的子域 $K$，使得
\[
\operatorname{Gal}(K/\mathbb Q)\cong G.
\]

这个结果有一个逆命题（其证明超出我们的范围），即著名的 Kronecker--Weber 定理：

\noindent\textbf{定理（Kronecker--Weber）.} 设 $K$ 是 $\mathbb Q$ 的一个有限 Abel 扩张。则 $K$ 包含在 $\mathbb Q$ 的某个分圆扩张中。

$\mathbb Q$ 的 Abel 扩张是“最容易的”Galois 扩张（至少就其 Galois 群的结构而言），而上面的结果表明，它们可以由 $\mathbb Q$ 的分圆扩张来分类。若以 $\mathbb Q$ 的其他有限扩张为基域，则描述 Abel 扩张要困难得多。对 $\mathbb Q$ 的任意有限扩张 $F$ 的 Abel 扩张的研究称为类域论。存在一种用与 $F$ 相关的不变量来分类 $F$ 的 Abel 扩张的理论，它极大地推广了 $\mathbb Q$ 上分圆域的结果。然而一般地，Abel 扩张的构造远不如分圆域情形那样显式。一个可以作出这种描述的情形是虚二次域（$D$ 为正时的 $\mathbb Q(\sqrt{-D})$）的 Abel 扩张，在那里 Abel 扩张可以通过添入某些椭圆函数的值来构造（这类似于添入单位根，而单位根是指数函数 $e^x$ 在某些 $x$ 处的值）。研究这类 Abel 扩张的算术以及寻找非 Abel 扩张的类似结果，是当前数学研究中丰富而迷人的领域。

我们以正 $n$ 边形能否用直尺圆规作出的问题结束对分圆域的讨论。

回忆（见第 13.3 节），元素 $\alpha$ 在 $\mathbb Q$ 上可作，当且仅当域 $\mathbb Q(\alpha)$ 包含在一个域 $K$ 中，而 $K$ 由一串二次扩张得到：
\[
\mathbb Q=K_0\subset K_1\subset\cdots\subset K_i\subset K_{i+1}\subset\cdots\subset K_m=K, \tag{14.12}
\]
其中
\[
[K_{i+1}:K_i]=2,
\qquad i=0,1,\ldots,m-1.
\]

在 $\mathbb R^2$ 中作正 $n$ 边形显然等价于作出 $n$ 次单位根，因为 $n$ 次单位根就是复平面 $\mathbb C$ 中单位圆上以点 $1$ 为一个顶点的正 $n$ 边形的顶点。

作出 $\zeta_n$ 等价于作出 $\zeta_n$ 在 $\mathbb R^2$ 中的第一坐标 $x$，即 $\zeta_n$ 的实部。由于 $\zeta_n$ 的复共轭正是 $\zeta_n^{-1}$，$\zeta_n$ 的实部为
\[
x=\frac12(\zeta_n+\zeta_n^{-1}).
\]
注意 $\zeta_n$ 在 $\mathbb Q(x)$ 上满足二次方程
\[
\zeta_n^2-2x\zeta_n+1=0.
\]
由于 $\mathbb Q(x)$ 只含有实数，故 $[\mathbb Q(\zeta_n):\mathbb Q(x)]=2$，从而 $\mathbb Q(x)$ 是 $\mathbb Q$ 的一个次数为 $\varphi(n)/2$ 的扩张。

于是，如果正 $n$ 边形可以由直尺圆规作出，那么 $\varphi(n)$ 必须是 $2$ 的幂。反过来，如果 $\varphi(n)=2^m$ 是 $2$ 的幂，则 Galois 群 $\operatorname{Gal}(\mathbb Q(\zeta_n)/\mathbb Q)$ 是一个阶为 $2$ 的幂的 Abel 群，所以 $\operatorname{Gal}(\mathbb Q(x)/\mathbb Q)$ 也有同样性质。由 Abel 群基本定理很容易看出，任意阶为 $2^m$ 的 Abel 群 $G$ 有一列子群
\[
G=G_m>G_{m-1}>\cdots>G_{i+1}>G_i>\cdots>G_0=1
\]
且
\[
[G_{i+1}:G_i]=2,
\qquad i=0,1,2,\ldots,m-1.
\]
把这应用到群 $G=\operatorname{Gal}(\mathbb Q(x)/\mathbb Q)$，并取子群 $G_i$，$i=0,1,\ldots,m-1$ 的固定域，由 Galois 理论基本定理得到如式 (12) 的一列二次扩张。

我们得出结论：正 $n$ 边形可以用直尺圆规作出，当且仅当 $\varphi(n)$ 是 $2$ 的幂。把 $n$ 分解成素数幂以计算 $\varphi(n)$，我们看到这意味着
\[
n=2^k p_1\cdots p_r
\]
是 $2$ 的一个幂与互异奇素数 $p_i$ 的乘积，其中每个 $p_i-1$ 都是 $2$ 的幂。初等练习可知，若素数 $p$ 满足 $p-1$ 是 $2$ 的幂，则 $p$ 必须形如
\[
p=2^{2^s}+1
\]
对某个整数 $s$ 成立。这样的素数称为 Fermat 素数。最初几个是
\[
3=2^1+1,
\]
\[
5=2^2+1,
\]
\[
17=2^4+1,
\]
\[
257=2^8+1,
\]
\[
65537=2^{16}+1.
\]
（但 $2^{32}+1$ 不是素数，因为它被 $641$ 整除。）目前尚不知道是否有无穷多个 Fermat 素数。我们把这总结为下面的命题。

\noindent\textbf{命题 29.} 正 $n$ 边形可以用直尺圆规作出，当且仅当
\[
n=2^k p_1\cdots p_r
\]
是 $2$ 的一个幂与互异 Fermat 素数的乘积。

上面的证明实际上给出了一种通过一串平方根来作正 $n$ 边形的程序。例如，正 $17$ 边形的作图（Gauss 于 1796 年、19 岁时解决）需要构造 $\mathbb Q(\zeta_{17})$ 中次数为 $2,4,8,16$ 的子域。这些子域可以通过形成 $\zeta_{17}$ 的周期来构造，如同上面 $13$ 次单位根的例子一样。在这种情形下，$\mathbb Q(\zeta_{17})$ 由一列二次扩张得到这一事实，反映为周期可以逐次“对半分”（也就是说，如果 $H_1<H_2$ 是子群且 $[H_2:H_1]=2$，那么 $H_1$ 的周期满足一个二次方程，其系数涉及 $H_2$ 的周期）。例如，Galois 群中指数为 $2$ 的子群（由 $\sigma^2$ 生成）的周期为（$\zeta=\zeta_{17}$）
\[
\eta_1=\zeta+\zeta^2+\zeta^4+\zeta^8+\zeta^9+\zeta^{13}+\zeta^{15}+\zeta^{16}
\]
和
\[
\eta_2=\zeta^3+\zeta^5+\zeta^6+\zeta^7+\zeta^{10}+\zeta^{11}+\zeta^{12}+\zeta^{14},
\]
它们把整个 Galois 群的周期“对半分”，并满足
\[
\eta_1+\eta_2=-1
\]
（由 $\zeta_{17}$ 满足的极小多项式得到）以及
\[
\eta_1\eta_2=-4
\]
（这需要计算；由 Galois 理论可知它必为有理数，因为这个乘积被 Galois 群的所有元素固定）。因此这两个周期是二次方程
\[
x^2+x-4=0
\]
的根，我们可以显式求解它。类似地，指数为 $4$ 的子群（由 $\sigma^4$ 生成）的周期自然把这些周期再对半分，因此在这些周期上是二次的，依此类推。这样可以用迭代平方根显式确定 $\zeta_{17}$。例如，可以求得
\[
8(\zeta+\zeta^{-1})=16\cos\frac{2\pi}{17}
\]
（这足以作出正 $17$ 边形）显式地等于
\[
-1+\sqrt{17}+\sqrt{2(17-\sqrt{17})}
+2\sqrt{17+3\sqrt{17}-\sqrt{2(17-\sqrt{17})}-2\sqrt{2(17+\sqrt{17})}}.
\]

一个相对简单的正 $17$ 边形作图（由 J. H. Conway 向我们展示）在练习中给出。

虽然我们已经看到，一般不能仅用连续开平方来求出 $\zeta_n$，但按定义可以通过连续开高次根得到 $\zeta_n$（即取 $1$ 的一个 $n$ 次根）。对于一般 $n$ 次方程的解，这种情形不再成立；一般不能用根式确定其解，我们将在下面几节中看到这一点。

\subsection*{练习}
\begin{enumerate}[label=\arabic*.]
\item 求正文中给出的 $\mathbb Q(\zeta_{13})$ 的各子域的本原生成元所满足的极小多项式。

\item 求由 $\zeta_8$ 的周期生成的 $\mathbb Q(\zeta_8)$ 的子域，并且特别证明并非每个子域都有这样的周期作为本原元。

\item 求 $5$ 次单位根 $\zeta_5$ 的周期 $\alpha=\zeta_5+\zeta_5^{-1}$ 所满足的二次方程。求 $\zeta_5$ 在 $\mathbb Q(\alpha)$ 上所满足的二次方程，并利用它显式求出 $5$ 次单位根。

\item 设 $\sigma_a\in\operatorname{Gal}(\mathbb Q(\zeta_n)/\mathbb Q)$ 表示 $n$ 次单位根分圆域中的自同构，它把 $\zeta_n$ 映到 $\zeta_n^a$，其中 $a$ 与 $n$ 互素且 $\zeta_n$ 是一个本原 $n$ 次单位根。证明对每个 $n$ 次单位根 $\zeta$，都有
\[
\sigma_a(\zeta)=\zeta^a.
\]

\item 设 $p$ 为素数，令 $\epsilon_1,\epsilon_2,\ldots,\epsilon_{p-1}$ 表示本原 $p$ 次单位根。置
\[
P_n=\epsilon_1^n+\epsilon_2^n+\cdots+\epsilon_{p-1}^n,
\]
即所有 $\epsilon_i$ 的 $n$ 次幂之和。证明若 $p$ 不整除 $n$，则 $P_n=-1$；若 $p$ 整除 $n$，则 $P_n=p-1$。〔一种方法：由 $\Phi_p(x)$ 得 $P_1=-1$；证明若 $p$ 不整除 $n$，则 $P_n$ 是 $P_1$ 的一个 Galois 共轭，因此也等于 $-1$。〕

\item 令 $\zeta_n$ 表示一个本原 $n$ 次单位根，令 $K=\mathbb Q(\zeta_n)$ 为相应分圆域。令 $a$ 表示 $\zeta_n$ 从 $K$ 到 $\mathbb Q$ 的迹（参见第 14.2 节练习 18）。证明：若 $n=1$，则 $a=1$；若 $n$ 被某个素数的平方整除，则 $a=0$；若 $n$ 是 $r$ 个互异素数的乘积，则 $a=(-1)^r$。

\item 证明复共轭限制到 $n$ 次单位根分圆域 $\mathbb Q(\zeta_n)$ 上就是自同构 $\sigma_{-1}\in\operatorname{Gal}(\mathbb Q(\zeta_n)/\mathbb Q)$。证明
\[
K^+=\mathbb Q(\zeta_n+\zeta_n^{-1})
\]
是 $K=\mathbb Q(\zeta_n)$ 中实元素组成的子域，称为 $K$ 的极大实子域。

\item 令 $K_n=\mathbb Q(\zeta_{2^{n+2}})$ 为 $2^{n+2}$ 次单位根的分圆域，$n\geq 0$。置
\[
\alpha_n=\zeta_{2^{n+2}}+\zeta_{2^{n+2}}^{-1},
\qquad
K_n^+=\mathbb Q(\alpha_n),
\]
即 $K_n$ 的极大实子域。
\begin{enumerate}[label=(\alph*)]
\item 证明对所有 $n\geq 0$，
\[
[K_n:\mathbb Q]=2^{n+1},\qquad [K_n:K_n^+]=2,
\]
\[
[K_n^+:\mathbb Q]=2^n,
\qquad [K_{n+1}^+:K_n^+]=2.
\]

\item 用 $\alpha_n$ 表示 $\zeta_{2^{n+2}}$ 在 $K_n^+$ 上所满足的二次方程。

\item 证明对 $n\geq 0$，
\[
\alpha_{n+1}^2=2+\alpha_n,
\]
并由此证明
\[
\alpha_n=\sqrt{2+\sqrt{2+\cdots+\sqrt2}}
\]
（共 $n$ 次），从而给出可作的 $2^{n+2}$ 次单位根的一个显式公式。
\end{enumerate}

\item 沿用前一题的记号。
\begin{enumerate}[label=(\alph*)]
\item 证明 $K_n^+$ 是 $\mathbb Q$ 的次数为 $2^n$ 的循环扩张。〔使用 Abel 群的一个显式同构
\[
(\mathbb Z/2^{n+2}\mathbb Z)^\times\cong \mathbb Z/2\mathbb Z\times\mathbb Z/2^n\mathbb Z
\]
（即 $(\mathbb Z/2^{n+2}\mathbb Z)^\times$ 同构于一个阶为 $2$ 的循环群与一个阶为 $2^n$ 的循环群的直积；参见第 2.3 节练习 22 和 23）。〕

\item 证明 $K_n$ 是 $K_{n-1}^+$ 的二二次扩张，并且三个中间域中的两个是 $K_n^+$ 和 $K_{n-1}$。证明 $K_{n-1}^+$ 与 $K_n$ 之间剩下的那个中间域是 $\mathbb Q$ 的次数为 $2^n$ 的循环扩张。
\end{enumerate}

\item 证明 $\mathbb Q(\sqrt[3]{2})$ 不是 $\mathbb Q$ 上任一分圆域的子域。

\item 证明本原 $n$ 次单位根构成 $n$ 次单位根分圆域在 $\mathbb Q$ 上的一组基，当且仅当 $n$ 是无平方因子的（即 $n$ 不被任一素数的平方整除）。

\item 设 $\sigma_p$ 表示有限域 $\mathbb F_q$（其中 $q=p^n$）的 Frobenius 自同构
\[
x\longmapsto x^p.
\]
把 $\mathbb F_q$ 看作 $\mathbb F_p$ 上的 $n$ 维向量空间 $V$，于是可以把 $\sigma_p$ 看作 $V$ 到自身的线性变换。求 $\sigma_p$ 的特征多项式，并证明：线性变换 $\sigma_p$ 在 $\mathbb F_p$ 上可对角化，当且仅当 $n$ 整除 $p-1$；它在 $\mathbb F_p$ 的代数闭包上可对角化，当且仅当 $(n,p)=1$。

\item 设
\[
n=p_1^{a_1}p_2^{a_2}\cdots p_k^{a_k}
\]
为 $n$ 的素因子分解，令 $\zeta_n$ 是一个本原 $n$ 次单位根。对每个 $i=1,2,\ldots,k$，定义 $d_i$ 由
\[
n=p_i^{a_i}d_i
\]
给出，并令
\[
\zeta_{p_i^{a_i}}=\zeta_n^{d_i},
\]
于是 $\zeta_{p_i^{a_i}}$ 是一个特定的本原 $p_i^{a_i}$ 次单位根。令 $\sigma_a\in\operatorname{Gal}(\mathbb Q(\zeta_n)/\mathbb Q)$ 为把 $\zeta_n$ 映到 $\zeta_n^a$ 的自同构，其中 $a$ 与 $n$ 互素。
\begin{enumerate}[label=(\alph*)]
\item 证明对 $i=1,2,\ldots,k$，$\sigma_a$ 把 $\zeta_{p_i^{a_i}}$ 映到 $\zeta_{p_i^{a_i}}^a$，并给出 $\mathbb Q(\zeta_{p_i^{a_i}})/\mathbb Q$ 的一个只依赖于 $a\pmod {p_i^{a_i}}$ 的自同构，我们可将它记为 $\sigma_{a\,({\rm mod}\,p_i^{a_i})}$。

\item 证明映射
\[
\sigma_a\longmapsto
\bigl(\sigma_{a\,({\rm mod}\,p_1^{a_1})},\ldots,
\sigma_{a\,({\rm mod}\,p_k^{a_k})}\bigr)
\]
就是推论 27 中对应于 $(\mathbb Z/n\mathbb Z)^\times$ 的中国剩余定理的同构。
\end{enumerate}

以下练习 14 到 18 确定与一个本原 $17$ 次单位根相关的周期，并给出练习 17 中指出的正 $17$ 边形的简单几何作图的证明。令
\[
\zeta=\zeta_{17}=\cos\frac{2\pi}{17}+i\sin\frac{2\pi}{17}
\]
为 $\mathbb C$ 中一个固定的本原 $17$ 次单位根。

\item 定义 $\zeta$ 的周期如下：
\[
\eta_1=\zeta+\zeta^2+\zeta^4+\zeta^8+\zeta^9+\zeta^{13}+\zeta^{15}+\zeta^{16},
\qquad
\eta_3'=\zeta^6+\zeta^7+\zeta^{10}+\zeta^{11},
\]
\[
\eta_2=\zeta^3+\zeta^5+\zeta^6+\zeta^7+\zeta^{10}+\zeta^{11}+\zeta^{12}+\zeta^{14},
\qquad
\eta_4'=\zeta^3+\zeta^5+\zeta^{12}+\zeta^{14},
\]
\[
\eta_1'=\zeta+\zeta^4+\zeta^{13}+\zeta^{16},
\qquad
\eta_1''=\zeta+\zeta^{16},
\]
\[
\eta_2'=\zeta^2+\zeta^8+\zeta^9+\zeta^{15},
\qquad
\eta_2''=\zeta^4+\zeta^{13}.
\]
\begin{enumerate}[label=(\alph*)]
\item 证明所有这些周期都是实数，并且
\[
\eta_1''=2\cos\frac{2\pi}{17}.
\]
证明作为实数，这些周期近似为
\[
\eta_1\approx 1.562,
\quad
\eta_1'\approx 2.049,
\quad
\eta_3'\approx -2.906,
\quad
\eta_1''\approx 1.865,
\]
\[
\eta_2\approx -2.562,
\quad
\eta_2'\approx -0.488,
\quad
\eta_4'\approx 0.344,
\quad
\eta_2''\approx 0.185.
\]

\item 证明 $\eta_1$ 和 $\eta_2$ 是方程
\[
x^2+x-4=0
\]
的根。

\item 证明 $\eta_1'$ 和 $\eta_2'$ 是方程
\[
x^2-\eta_1x-1=0
\]
的根，并且 $\eta_3'$ 和 $\eta_4'$ 是方程
\[
x^2-\eta_2x-1=0
\]
的根。

\item 证明 $\eta_1''$ 和 $\eta_2''$ 是方程
\[
x^2-\eta_1'x+\eta_2'=0
\]
的根。
\end{enumerate}

\item 证明若 $\tan 2\theta=a$，其中 $0<2\theta<\pi/2$，则 $\tan\theta$ 满足方程
\[
x^2-\frac{2}{a}x-1=0.
\]

\item 令 $C$ 为 $\mathbb R^2$ 中以点 $(h,k)$ 和 $(0,1)$ 为直径的圆。证明这个圆与 $x$ 轴相交，当且仅当 $h^2-4k\geq 0$；并且在这种情形中，两个截距是方程
\[
x^2-hx+k=0
\]
的根。

\item （正 $17$ 边形的作图）画一个以原点 $(0,0)$ 为圆心、半径为 $2$ 的圆。
\begin{enumerate}[label=(\alph*)]
\item 连接点 $(4,0)$ 与点 $(0,1)$，并作直线 $\ell_1$，它平分这条直线与 $y$ 轴之间的角。作直线 $\ell_2$ 垂直于图 2 中的 $\ell_1$。

\begin{center}
\begin{tikzpicture}[scale=0.75, every node/.style={inner sep=2pt}]
\draw[->] (-3.2,0)--(4.7,0) node[right] {$x$};
\draw[->] (0,-2.4)--(0,2.7) node[above] {$y$};
\draw (0,0) circle (2);
\fill (4,0) circle (1.5pt) node[below] {$(4,0)$};
\fill (0,1) circle (1.5pt) node[left] {$(0,1)$};
\draw (4,0)--(0,1);
\draw[dashed] (0,-2.2)--(0,2.5);
\draw[thick] (-2.6,-1.5)--(2.1,1.2) node[right] {$\ell_1$};
\draw[thick] (-1.0,2.3)--(1.4,-1.9) node[right] {$\ell_2$};
\node at (3.8,-1.5) {图 2};
\end{tikzpicture}
\end{center}

\item 以 $\ell_1$ 与 $x$ 轴的交点为圆心，并取到 $(0,1)$ 的距离为半径作圆 $C_1$，令 $A=(s,0)$ 为 $C_1$ 与 $x$ 轴右侧的交点。类似地，令 $B=(t,0)$ 表示 $x$ 轴与圆 $C_2$ 的右侧交点，其中 $C_2$ 的圆心是 $\ell_2$ 与 $x$ 轴的交点，半径等于到 $(0,1)$ 的距离，如图 3 所示。

\begin{center}
\begin{tikzpicture}[scale=0.75, every node/.style={inner sep=2pt}]
\draw[->] (-3.2,0)--(4.7,0) node[right] {$x$};
\draw[->] (0,-2.4)--(0,2.7) node[above] {$y$};
\draw (0,0) circle (2);
\fill (0,1) circle (1.5pt) node[left] {$(0,1)$};
\coordinate (O1) at (0.8,0);
\coordinate (O2) at (-1.3,0);
\draw[thick] (O1) circle (1.28);
\draw[thick] (O2) circle (1.64);
\draw[dashed] (O1)--(0,1);
\draw[dashed] (O2)--(0,1);
\fill (2.08,0) circle (1.5pt) node[below] {$A=(s,0)$};
\fill (0.34,0) circle (1.5pt) node[below] {$B=(t,0)$};
\node at (1.4,1.4) {$C_1$};
\node at (-2.4,1.2) {$C_2$};
\node at (3.8,-1.5) {图 3};
\end{tikzpicture}
\end{center}

\item 在点 $A$ 处作 $x$ 轴的垂线，并把从 $(0,0)$ 到 $B$ 的距离 $t$ 标出，从而作出点 $(s,t)$。作以 $(s,t)$ 和 $(0,1)$ 为直径的圆，令 $P$ 为这个圆与 $x$ 轴的右侧交点。过 $P$ 作 $x$ 轴的垂线，它与半径为 $2$ 的圆相交于正 $17$ 边形的第二个顶点，而第一个顶点取为 $(2,0)$；于是如图 4 所示，用直尺圆规作出了正 $17$ 边形。

\begin{center}
\begin{tikzpicture}[scale=0.75, every node/.style={inner sep=2pt}]
\draw[->] (-3.2,0)--(4.7,0) node[right] {$x$};
\draw[->] (0,-2.4)--(0,2.7) node[above] {$y$};
\draw (0,0) circle (2);
\fill (2,0) circle (1.5pt) node[below right] {$(2,0)$};
\fill (0,1) circle (1.5pt) node[left] {$(0,1)$};
\fill (1.6,0) circle (1.5pt) node[below] {$P$};
\draw[dashed] (1.6,-0.1)--(1.6,2.1);
\fill (1.6,1.2) circle (1.5pt) node[right] {第二顶点};
\fill (1.1,0.6) circle (1.2pt) node[right] {$(s,t)$};
\draw[thick] (0.55,0.8) circle (0.78);
\draw[dashed] (0,1)--(1.1,0.6);
\draw[dashed] (1.1,0.6)--(1.6,0);
\node at (3.8,-1.5) {图 4};
\end{tikzpicture}
\end{center}
\end{enumerate}

\item 沿用前一练习的记号。
\begin{enumerate}[label=(\alph*)]
\item 证明 $\ell_1$ 与 $x$ 轴交于点 $(\eta_1/2,0)$，并且 $\ell_2$ 与 $x$ 轴交于点 $(\eta_2/2,0)$。

\item 证明 $C_1$ 是以点 $(\eta_1,-1)$ 和 $(0,1)$ 为直径的圆。证明 $s=\eta_1'$。类似地，证明 $C_2$ 是以点 $(\eta_2,-1)$ 和 $(0,1)$ 为直径的圆，并且 $t=\eta_4'$。

\item 证明 $P$ 的坐标为 $(\eta_1'',0)$，从而前一题中的作图确实用直尺圆规作出了正 $17$ 边形。
\end{enumerate}
\end{enumerate}

\section{多项式的 Galois 群}

回忆：可分多项式 $f(x)\in F[x]$ 的 Galois 群，定义为 $f(x)$ 的分裂域在 $F$ 上的 Galois 群。

若 $K$ 是 $F$ 的 Galois 扩张，则 $K$ 是 $F$ 上某个可分多项式 $f(x)$ 的分裂域。任意自同构 $\sigma\in \operatorname{Gal}(K/F)$ 都把 $f(x)$ 的一个不可约因子的一个根映到同一不可约因子的另一个根，并且 $\sigma$ 由它在这些根上的作用唯一确定（因为这些根在 $F$ 上生成 $K$）。如果固定 $f(x)$ 的根 $\alpha_1,\ldots,\alpha_n$ 的一个标号，则每个 $\sigma\in\operatorname{Gal}(K/F)$ 都定义了 $\alpha_1,\\ldots,\alpha_n$ 的一个唯一置换，因而定义了下标集合 $\{1,2,\ldots,n\}$ 的一个唯一置换（这个置换依赖于所固定的根的标号）。于是得到一个单射
\[
\operatorname{Gal}(K/F)\hookrightarrow S_n
\]
把 Galois 群嵌入 $n$ 个字母上的对称群；这显然是一个同态（两个群的运算都是复合）。因此我们可以把 Galois 群看成对称群的子群。由于由基本定理知分裂域的次数等于 Galois 群的阶，这也从群论方面解释了为什么 $F$ 上一个 $n$ 次多项式的分裂域次数至多为 $n!$（命题 13.26）。

一般地，如果 $f(x)$ 分解为不可约因子的乘积
\[
f(x)=f_1(x)\cdots f_k(x),
\]
其中 $\deg f_i=n_i$，$i=1,2,\ldots,k$，则 Galois 群只在每个不可约因子的根之间作置换，因此
\[
\operatorname{Gal}(K/F)\leq S_{n_1}\times\cdots\times S_{n_k}.
\]
如果 $f(x)$ 是不可约的，那么给定 $f(x)$ 的任意两个根，存在 $f(x)$ 的 Galois 群 $G$ 中的一个自同构把第一个根映到第二个根（这来自扩张定理 13.27）。这样的群称为在根上传递；也就是说，可以通过作用某个 $G$ 的元素从任意给定的根到达任意另一个根。Galois 群必须在根的块（即不可约因子的根）上传递这一事实，常常有助于减少 $G$ 的可能结构（参见下面关于四次多项式的 Galois 群的讨论）。

\paragraph{例}
\begin{enumerate}[label=(\arabic*)]
\item 考虑双二次扩张 $\mathbb Q(\sqrt2,\sqrt3)$ 在 $\mathbb Q$ 上的扩张，它是 $(x^2-2)(x^2-3)$ 的分裂域。把根标号为
\[
\alpha_1=\sqrt2,\qquad \alpha_2=-\sqrt2,\qquad \alpha_3=\sqrt3,
\qquad \alpha_4=-\sqrt3.
\]
Galois 群的元素为 $\{1,\sigma,\tau,\sigma\tau\}$，其中 $\sigma$ 把 $\sqrt2$ 映到 $-\sqrt2$ 并固定 $\sqrt3$，而 $\tau$ 固定 $\sqrt2$ 并把 $\sqrt3$ 映到 $-\sqrt3$。作为在这些根上的置换，$\sigma$ 交换前两个根并固定后两个根，而 $\tau$ 固定前两个根并交换后两个根；也就是说，作为 $S_4$ 的元素，
\[
\sigma=(12),\qquad \tau=(34),\qquad \sigma\tau=(12)(34).
\]
因此
\[
\operatorname{Gal}(\mathbb Q(\sqrt2,\sqrt3)/\mathbb Q)
\cong \{1,(12),(34),(12)(34)\}\leq S_4,
\]
这把该 Galois 群识别为 $S_4$ 中的 Klein 四元子群。注意，如果改变上面根的标号，就会得到 Galois 群作为 $S_4$ 子群的另一个（同构的）表示；例如，若交换第二个根与第三个根，则得到子群 $\{1,(13),(24),(13)(24)\}$。

\item $x^3-2$ 的 Galois 群作用在三个根
\[
\sqrt[3]{2},\qquad \rho\sqrt[3]{2},\qquad \rho^2\sqrt[3]{2}
\]
上，其中 $\rho$ 是本原三次单位根。按这个顺序排列，先前定义的生成元 $\sigma$ 与 $\tau$ 给出置换
\[
\sigma=(123),\qquad \tau=(23),
\]
于是
\[
\{1,\sigma,\sigma^2,\tau,\sigma\tau,\sigma\tau^2\}
=\{1,(123),(132),(23),(13),(12)\}=S_3.
\]
在这个情形中，Galois 群是三个字母上的整个对称群。
\end{enumerate}

回忆，每个有限群都同构于某个对称群 $S_n$ 的子群。判断每个有限群是否都能作为 $\mathbb Q$ 上某个多项式的 Galois 群，是一个开放问题。我们在上一节已经看到，每个有限 Abel 群都是 $\mathbb Q$ 上某个 Galois 群（来自某个分圆域的子域）。下面我们将显式确定低次数（$\leq 4$）多项式的 Galois 群；这特别会表明 $S_4$ 的每个子群都作为 Galois 群出现。

我们首先引入若干定义，并说明 $n$ 次“一般”多项式的 Galois 群是 $S_n$（因此上面第二个例子应视为“典型”的）。

\noindent\textbf{定义.} 令 $x_1,x_2,\ldots,x_n$ 为不定元。\emph{基本对称函数} $s_1,s_2,\ldots,s_n$ 定义为
\[
\begin{aligned}
s_1&=x_1+x_2+\cdots+x_n,\\
s_2&=x_1x_2+x_1x_3+\cdots+x_2x_3+x_2x_4+\cdots+x_{n-1}x_n,\\
&\ \vdots\\
s_n&=x_1x_2\cdots x_n;
\end{aligned}
\]
也就是说，$x_1,x_2,\ldots,x_n$ 的第 $i$ 个对称函数 $s_i$ 是所有一次取 $i$ 个 $x_j$ 的乘积之和。

\noindent\textbf{定义.} \emph{$n$ 次一般多项式}是多项式
\[
(x-x_1)(x-x_2)\cdots(x-x_n),
\]
其根是不定元 $x_1,x_2,\ldots,x_n$。

容易用归纳看出，$n$ 次一般多项式的系数由根的基本对称函数给出：
\begin{equation}
(x-x_1)(x-x_2)\cdots(x-x_n)=x^n-s_1x^{n-1}+s_2x^{n-2}+\cdots+(-1)^n s_n.\tag{14.13}
\end{equation}

对任意域 $F$，扩张 $F(x_1,x_2,\ldots,x_n)$ 是 $F(s_1,s_2,\ldots,s_n)$ 上的 Galois 扩张，因为它是 $n$ 次一般多项式的分裂域。

若 $\sigma\in S_n$ 是 $\{1,2,\ldots,n\}$ 的任意置换，则 $\sigma$ 通过置换变量 $x_1,x_2,\ldots,x_n$ 的下标作用在 $F(x_1,x_2,\ldots,x_n)$ 中的有理函数上。显然这给出 $F(x_1,x_2,\ldots,x_n)$ 的一个自同构。把 $\sigma\in S_n$ 与这个自同构相识别，便把 $S_n$ 识别为 $\operatorname{Aut}(F(x_1,x_2,\ldots,x_n))$ 的一个子群。

基本对称函数 $s_1,s_2,\ldots,s_n$ 在下标的任意置换下都保持不变（这就是称它们为对称的原因），这说明子域 $F(s_1,s_2,\ldots,s_n)$ 包含在 $S_n$ 的固定域中。由 Galois 理论基本定理，$S_n$ 的固定域在 $F(x_1,x_2,\ldots,x_n)$ 中的指数恰为 $n!$。另一方面，因为 $F(x_1,x_2,\ldots,x_n)$ 是式 (14.13) 的 $n$ 次多项式在 $F(s_1,s_2,\ldots,s_n)$ 上的分裂域，我们有
\begin{equation}
[F(x_1,x_2,\ldots,x_n):F(s_1,s_2,\ldots,s_n)]\leq n!.\tag{14.14}
\end{equation}
于是实际上必有等号，并且 $F(s_1,s_2,\ldots,s_n)$ 正是 $S_n$ 的固定域。这证明了如下结果。

\noindent\textbf{命题 30.} 对称群 $S_n$ 作用在 $n$ 个变量的有理函数域 $F(x_1,x_2,\ldots,x_n)$ 上时，它的固定域是基本对称函数生成的有理函数域
\[
F(s_1,s_2,\ldots,s_n).
\]

\noindent\textbf{定义.} 若有理函数 $f(x_1,x_2,\ldots,x_n)$ 在变量 $x_1,x_2,\ldots,x_n$ 的任意置换下都不变，则称它是\emph{对称的}。

\noindent\textbf{推论 31（对称函数基本定理）.} 变量 $x_1,x_2,\ldots,x_n$ 中的任意对称函数都是基本对称函数 $s_1,s_2,\ldots,s_n$ 的有理函数。

\begin{proof}
对称函数属于上面 $S_n$ 的固定域，因而是 $s_1,\ldots,s_n$ 的有理函数。
\end{proof}

这个推论解释了为什么这些函数称为基本对称函数。

\paragraph{注记}
如果 $f(x_1,\ldots,x_n)$ 是 $x_1,x_2,\ldots,x_n$ 中的一个对称多项式，那么可以看出 $f$ 实际上是 $s_1,s_2,\ldots,s_n$ 的一个多项式；这加强了推论的结论。事实上，如果 $R$ 是任意含 $1$ 的交换环，则系数属于 $R$ 的对称多项式，都是系数属于 $R$ 的基本对称函数多项式。练习中给出的把对称多项式写成基本对称函数多项式的算法，隐含着这个事实的证明。

\paragraph{例}
\begin{enumerate}[label=(\arabic*)]
\item 表达式 $(x_1-x_2)^2$ 关于 $x_1,x_2$ 对称。我们有
\[
(x_1-x_2)^2=(x_1+x_2)^2-4x_1x_2=s_1^2-4s_2,
\]
这是基本对称函数的多项式。

\item 多项式 $x_1^2+x_2^2+x_3^2$ 关于 $x_1,x_2,x_3$ 对称，在此情形中
\[
x_1^2+x_2^2+x_3^2=(x_1+x_2+x_3)^2-2(x_1x_2+x_1x_3+x_2x_3)=s_1^2-2s_2.
\]

\item 多项式 $x_1^2x_2^2+x_1^2x_3^2+x_2^2x_3^2$ 是对称的。由于
\[
\begin{aligned}
(x_1x_2+x_1x_3+x_2x_3)^2
&=x_1^2x_2^2+x_1^2x_3^2+x_2^2x_3^2 \\
&\quad+2(x_1^2x_2x_3+x_2^2x_1x_3+x_3^2x_1x_2)\\
&=x_1^2x_2^2+x_1^2x_3^2+x_2^2x_3^2+2x_1x_2x_3(x_1+x_2+x_3),
\end{aligned}
\]
所以
\[
x_1^2x_2^2+x_1^2x_3^2+x_2^2x_3^2=s_2^2-2s_1s_3.
\]
\end{enumerate}

现在假设我们从一般多项式
\[
x^n-s_1x^{n-1}+s_2x^{n-2}+\cdots+(-1)^n s_n
\]
出发，把它看作域 $F(s_1,s_2,\ldots,s_n)$ 上的多项式，其中 $s_i$，$i=1,2,\ldots,n$，被看作不定元。如果定义这个多项式的根为 $x_1,x_2,\ldots,x_n$，则 $s_i$ 恰好是根 $x_1,\ldots,x_n$ 的基本对称函数。此外，这些根本身也是不定元，也就是说它们之间在 $F$ 上没有多项式关系。事实上，若 $p(t_1,\ldots,t_n)$ 是 $n$ 个变量中系数在 $F$ 中的一个非零多项式，并且 $p(x_1,\ldots,x_n)=0$，则取 $S_n$ 中所有置换 $\sigma$ 的乘积
\[
\prod_{\sigma\in S_n}p(t_{\sigma(1)},\ldots,t_{\sigma(n)}),
\]
得到一个非零对称多项式，它在 $(x_1,
\ldots,x_n)$ 处为 $0$。这会给出 $s_1,\ldots,s_n$ 之间在 $F$ 上的非零多项式关系，矛盾。反过来，如果一个多项式 $f(x)$ 的根在 $F$ 上是独立的不定元，则 $f(x)$ 的系数也是如此；参见第 9 节开头。因此，用“不定根”定义一般多项式，与用“不定系数”定义一般多项式是等价的。从这个观点出发，上述结果可以表述为：

\noindent\textbf{定理 32.} 一般多项式
\[
x^n-s_1x^{n-1}+s_2x^{n-2}+\cdots+(-1)^n s_n
\]
在域 $F(s_1,s_2,\ldots,s_n)$ 上可分，并且其 Galois 群为 $S_n$。

这个结果说，如果 $n$ 次多项式的系数之间没有关系（这就是上面说 $s_i$ 是不定元的含义），那么这个多项式在由其系数生成的域上的 Galois 群就是完整的对称群 $S_n$。粗略地说，这意味着“泛型”的 $n$ 次多项式会有 $S_n$ 作为 Galois 群。不过要注意，在有限域上每个多项式的 Galois 群都是循环群（有限域的所有扩张都是循环的），所以这种意义下的“泛型”多项式并不存在。在 $\mathbb Q$ 上，可以把“泛型”多项式的概念精确化；于是大多数多项式确有完整的对称群作为 Galois 群。

当 $n\geq 5$ 时，$S_n$ 只有一个非平凡正规子群，即指数为 $2$ 的子群 $A_n$。因此一般地，$F(x_1,x_2,\ldots,x_n)$ 中含有 $F(s_1,s_2,\ldots,s_n)$ 的正规子域只有一个，并且它是 $2$ 次扩张。

\noindent\textbf{定义.} 定义 $x_1,x_2,\ldots,x_n$ 的\emph{判别式} $D$ 为
\[
D=\prod_{i<j}(x_i-x_j)^2.
\]
多项式的判别式定义为其根的判别式。

判别式 $D$ 是 $x_1,
\ldots,x_n$ 中的对称函数，因而属于 $K=F(s_1,s_2,\ldots,s_n)$。

当我们最初定义交错群 $A_n$ 时已经看到，一个置换 $\sigma\in S_n$ 属于子群 $A_n$ 当且仅当它固定乘积
\[
\sqrt D=\prod_{i<j}(x_i-x_j)\in \mathbb Z[x_1,x_2,\ldots,x_n].
\]
因此（由基本定理），若 $F$ 的特征不为 $2$，则 $\sqrt D$ 生成 $A_n$ 的固定域，并且生成 $K$ 的一个二次扩张。这证明了如下命题。

\noindent\textbf{命题 33.} 若 $\operatorname{ch}(F)\neq 2$，则置换 $\sigma\in S_n$ 属于 $A_n$ 当且仅当它固定判别式 $D$ 的平方根。

现在我们考虑域 $F$ 上小次数（$\leq 4$）的可分多项式的 Galois 群，并假设 $F$ 的特征不是 $2$ 或 $3$。注意，在 $\mathbb Q$ 上或在有限域上（更一般地，在任意完美域上），任意多项式 $f(x)$ 的分裂域与把 $f(x)$ 的不可约因子各取一次后的乘积的分裂域相同；后者是可分多项式。

若多项式
\[
f(x)=x^n+a_{n-1}x^{n-1}+\cdots+a_1x+a_0
\]
的根为 $\alpha_1,\alpha_2,\ldots,\alpha_n$，则 $f(x)$ 的判别式为
\[
D=\prod_{i<j}(\alpha_i-\alpha_j)^2.
\]
注意，$D=0$ 当且仅当 $f(x)$ 不可分，也就是当且仅当根 $\alpha_1,\ldots,\alpha_n$ 不是两两不同。回忆，在完美域上（例如 $\mathbb Q$ 或有限域），这蕴含 $f(x)$ 可约，因为完美域上的每个不可约多项式都是可分的。

判别式 $D$ 关于 $f(x)$ 的根是对称的，因而被 $f(x)$ 的 Galois 群的所有自同构固定。由基本定理可知 $D\in F$。一般地，判别式可写成 $f(x)$ 的系数的多项式（由推论 31），对于较高次数这些表达式相当复杂（下面给出 $n\leq 4$ 的公式）。最后注意，由于
\[
\sqrt D=\prod_{i<j}(\alpha_i-\alpha_j),
\]
所以 $\sqrt D$ 总是包含在 $f(x)$ 的分裂域中。

若 $f(x)$ 的根两两不同，固定根的某个排序，并按上面方式把 $f(x)$ 的 Galois 群看成 $S_n$ 的子群。

\noindent\textbf{命题 34.} 多项式 $f(x)\in F[x]$ 的 Galois 群是 $A_n$ 的子群，当且仅当判别式 $D\in F$ 是 $F$ 中某个元素的平方。

\begin{proof}
这就是命题 33 在当前情形下的重述。Galois 群包含在 $A_n$ 中，当且仅当 Galois 群的每个元素都固定
\[
\sqrt D=\prod_{i<j}(\alpha_i-\alpha_j),
\]
也就是当且仅当 $\sqrt D\in F$。
\end{proof}

这个性质，加上 $D=0$ 判定重根的出现，正是 $D$ 称为判别式的原因。

\paragraph{二次多项式}
考虑根为 $\alpha,\beta$ 的多项式 $x^2+ax+b$。它的判别式为 $(\alpha-\beta)^2$，可写成根的基本对称函数的多项式。我们在上面例 (1) 中已经做过：
\[
D=s_1^2-4s_2=(-a)^2-4b=a^2-4b,
\]
这就是通常的二次判别式。

该多项式可分当且仅当 $a^2-4b\neq0$。Galois 群是 $S_2$ 的子群；它是平凡群（在这里也就是 $A_2$）当且仅当 $a^2-4b$ 是 $F$ 中的平方，这完全决定了可能的 Galois 群。

注意，这只是重新表述了先前通过显式求根得到的结果：如果多项式可约（即 $D$ 是 $F$ 中平方），则 Galois 群平凡（分裂域就是 $F$）；而如果多项式不可约，则分裂域为二次扩张 $F(\sqrt D)$，Galois 群同构于 $\mathbb Z/2\mathbb Z$。

\paragraph{三次多项式}
设三次多项式为
\begin{equation}
f(x)=x^3+ax^2+bx+c.\tag{14.15}
\end{equation}
若作代换 $x=y-a/3$，它变为三次多项式
\begin{equation}
g(y)=y^3+py+q,\tag{14.16}
\end{equation}
其中
\begin{equation}
p=\frac{1}{3}(3b-a^2),\qquad q=\frac{1}{27}(2a^3-9ab+27c).\tag{14.17}
\end{equation}
这两个多项式的分裂域相同，因为它们的根只相差常数 $a/3\in F$；并且由于判别式公式只涉及根的差，所以二者也有相同的判别式。

令 (14.16) 中多项式的根为 $\alpha,\beta,\gamma$。我们先用 $p$ 和 $q$ 计算这个多项式的判别式。注意
\[
g(y)=(y-\alpha)(y-\beta)(y-\gamma),
\]
故微分得
\[
D_yg(y)=(y-\alpha)(y-\beta)+(y-\alpha)(y-\gamma)+(y-\beta)(y-\gamma).
\]
于是
\[
D_yg(\alpha)=(\alpha-\beta)(\alpha-\gamma),
\quad
D_yg(\beta)=(\beta-\alpha)(\beta-\gamma),
\quad
D_yg(\gamma)=(\gamma-\alpha)(\gamma-\beta).
\]
取乘积得到
\[
D=[(\alpha-\beta)(\alpha-\gamma)(\beta-\gamma)]^2=-D_yg(\alpha)D_yg(\beta)D_yg(\gamma).
\]
由于 $D_yg(y)=3y^2+p$，我们有
\[
\begin{aligned}
-D&=(3\alpha^2+p)(3\beta^2+p)(3\gamma^2+p)\\
&=27\alpha^2\beta^2\gamma^2
+9p(\alpha^2\beta^2+\alpha^2\gamma^2+\beta^2\gamma^2)
+3p^2(\alpha^2+\beta^2+\gamma^2)+p^3.
\end{aligned}
\]
根的基本对称函数中的相应表达式已在上面的例 (2) 和 (3) 中确定。这里 $s_1=0$，$s_2=p$，$s_3=-q$。于是得到
\begin{equation}
D=-4p^3-27q^2.\tag{14.18}
\end{equation}
这也是 (14.15) 中 $f(x)$ 的判别式。用 $a,b,c$ 通过 (14.17) 表示 $D$，得到
\begin{equation}
D=a^2b^2-4b^3-4a^3c-27c^2+18abc.\tag{14.18'}
\end{equation}

\paragraph{三次多项式的 Galois 群}
\begin{enumerate}[label=\alph*.]
\item 如果三次多项式 $f(x)$ 可约，那么它或者分解为三个一次因子，或者分解为一个一次因子与一个不可约二次因子。在第一种情形中 Galois 群平凡；在第二种情形中 Galois 群阶为 $2$。

\item 如果三次多项式 $f(x)$ 不可约，则 $f(x)$ 的一个根生成 $F$ 上次数为 $3$ 的扩张，所以分裂域在 $F$ 上的次数被 $3$ 整除。由于 Galois 群是 $S_3$ 的子群，只有两种可能：$A_3$ 或 $S_3$。Galois 群是 $A_3$（即 $3$ 阶循环群）当且仅当 (14.18) 中的判别式 $D$ 是平方。
\end{enumerate}

明确地说，如果 $D$ 是 $F$ 中某元素的平方，那么不可约三次多项式 $f(x)$ 的分裂域可由向 $F$ 添入 $f(x)$ 的任意一个根得到。所得域是 $F$ 上次数为 $3$ 的 Galois 扩张，其 Galois 群是 $3$ 阶循环群。如果 $D$ 不是 $F$ 中元素的平方，那么 $f(x)$ 的分裂域在 $F$ 上次数为 $6$，因而对 $f(x)$ 的任一根 $\theta$，这个分裂域就是 $F(\theta,\sqrt D)$。这个扩张在 $F$ 上是 Galois 的，其 Galois 群为 $S_3$（可取生成元 $\sigma$ 与 $\tau$，其中 $\sigma$ 把 $\theta$ 映到 $f(x)$ 的另一个根并固定 $\sqrt D$，而 $\tau$ 把 $\sqrt D$ 映到 $-\sqrt D$ 并固定 $\theta$）。

我们看到，在两种情形中，不可约三次多项式 $f(x)$ 的分裂域都通过向 $F$ 添入 $\sqrt D$ 与 $f(x)$ 的一个根而得到。

在下一节引入 Lagrange 预解式后，我们将给出 (14.16) 的根的显式公式（Cardano 公式）。

\paragraph{四次多项式}
令四次多项式为
\[
f(x)=x^4+ax^3+bx^2+cx+d,
\]
它在代换 $x=y-a/4$ 下变为四次多项式
\[
g(y)=y^4+py^2+qy+r,
\]
其中
\[
p=\frac{1}{8}(-3a^2+8b),
\qquad
q=\frac{1}{8}(a^3-4ab+8c),
\]
\[
r=\frac{1}{256}(-3a^4+16a^2b-64ac+256d).
\]
令 $g(y)$ 的根为 $\alpha_1,\alpha_2,\alpha_3,\alpha_4$，并令 $G$ 表示 $g(y)$（或 $f(x)$）的分裂域的 Galois 群。

首先假设 $g(y)$ 可约。如果 $g(y)$ 分解为一个一次因子与一个三次因子，那么 $G$ 就是这个三次因子的 Galois 群，已在上面确定。再假设 $g(y)$ 分解为两个不可约二次因子。则分裂域是扩张 $F(\sqrt{D_1},\sqrt{D_2})$，其中 $D_1,D_2$ 是这两个二次因子的判别式。如果 $D_1$ 与 $D_2$ 不只差一个平方因子，那么这是一个双二次扩张，$G$ 同构于 $S_4$ 的 Klein 四元子群；如果 $D_1$ 是 $D_2$ 乘以一个平方，则这是二次扩张，$G\cong \mathbb Z/2\mathbb Z$。

于是只剩下 $g(y)$ 不可约的情形。在这种情形中回忆，Galois 群在根上传递，也就是说，可以通过作用某个 Galois 群元素从任意给定的根得到任意另一个根。考察各种可能可知，$S_4$ 的传递子群，因而我们的 Galois 群 $G$ 的唯一可能性，是
\[
S_4,
\qquad A_4,
\]
\[
D_8=\{1,(1324),(12)(34),(1423),(13)(24),(14)(23),(12),(34)\}
\quad\text{及其共轭},
\]
\[
V=\{1,(12)(34),(13)(24),(14)(23)\},
\]
\[
C=\{1,(1234),(13)(24),(1432)\}\quad\text{及其共轭}.
\]
这里 $D_8$ 是二面体群，是 $S_4$ 的一个 Sylow $2$-子群，在 $S_4$ 中有 $3$ 个（同构的）共轭子群；$V$ 是 Klein 四元群，在 $S_4$ 中正规；$C$ 是一个循环群，在 $S_4$ 中有 $3$ 个（同构的）共轭子群。

考虑 $g(y)$ 的分裂域中的元素
\[
\theta_1=(\alpha_1+\alpha_2)(\alpha_3+\alpha_4),
\]
\[
\theta_2=(\alpha_1+\alpha_3)(\alpha_2+\alpha_4),
\qquad
\theta_3=(\alpha_1+\alpha_4)(\alpha_2+\alpha_3).
\]
这些元素被 $S_4$ 中的置换在彼此之间置换。$\theta_1$ 在 $S_4$ 中的稳定子是二面体群 $D_8$；$\theta_2$ 与 $\theta_3$ 在 $S_4$ 中的稳定子是阶为 $8$ 的共轭二面体子群。$S_4$ 中稳定所有三个元素的子群是这些子群的交，即 Klein 四元群 $V$。

由于 $S_4$ 只是置换 $\theta_1,\theta_2,\theta_3$，这些 $\theta$ 的基本对称函数被 $S_4$ 的所有元素固定，因而属于 $F$。对称函数中的一个初等计算显示，这些基本对称函数分别为 $2p$、$p^2-4r$ 与 $-q^2$，这说明 $\theta_1,\theta_2,\theta_3$ 是
\[
h(x)=x^3-2px^2+(p^2-4r)x+q^2
\]
的根。这个多项式称为四次多项式 $g(y)$ 的\emph{预解三次方程}。由于
\[
\theta_1-\theta_2=\alpha_1\alpha_3+\alpha_2\alpha_4-\alpha_1\alpha_2-\alpha_3\alpha_4
=-(\alpha_1-\alpha_4)(\alpha_2-\alpha_3),
\]
并且类似地
\[
\theta_1-\theta_3=-(\alpha_1-\alpha_3)(\alpha_2-\alpha_4),
\qquad
\theta_2-\theta_3=-(\alpha_1-\alpha_2)(\alpha_3-\alpha_4),
\]
我们看到预解三次方程的判别式与四次多项式 $g(y)$ 的判别式相同，因而也与四次多项式 $f(x)$ 的判别式相同。利用三次多项式的判别式公式，我们可以容易地用 $p,q,r$ 计算判别式：
\[
D=16p^4r-4p^3q^2-128p^2r^2+144pq^2r-27q^4+256r^3.
\]
由此可以给出 $D$ 用 $a,b,c,d$ 表示的公式：
\[
\begin{aligned}
D={}&-128b^2d^2-4a^3c^3+16b^4d-4b^3c^2-27a^4d^2+18abc^3\\
&+144a^2bd^2-192acd^2+a^2b^2c^2-4a^2b^3d-6a^2c^2d\\
&+144bc^2d+256d^3-27c^4-80ab^2cd+18a^3bcd.
\end{aligned}
\]

预解三次方程的分裂域是四次多项式分裂域的一个子域，所以预解三次方程的 Galois 群是 $G$ 的一个商。因此，知道 Galois 群在预解三次方程 $h(x)$ 的根上的作用，可以给出关于 $g(y)$ 的 Galois 群的信息，如下：

\paragraph{四次多项式的 Galois 群}
\begin{enumerate}[label=\alph*.]
\item 首先假设预解三次方程不可约。如果 $D$ 不是平方，则 $G$ 不包含在 $A_4$ 中，并且预解三次方程的 Galois 群为 $S_3$，这蕴含 $g(y)$ 的分裂域次数被 $6$ 整除。唯一可能于是是 $G=S_4$。

\item 如果预解三次方程不可约且 $D$ 是平方，则 $G$ 是 $A_4$ 的子群，并且 $3$ 整除 $G$ 的阶（预解三次方程的 Galois 群为 $A_3$）。唯一可能是 $G=A_4$。

\item[\textup{c1}.] 现在剩下预解三次方程可约的情形。第一种可能是 $h(x)$ 有 $F$ 中的三个根（也就是完全分裂）。由于每个 $\theta_1,\theta_2,\theta_3$ 都在 $F$ 中，$G$ 的每个元素都固定这三个元素，这意味着 $G\leq V$。唯一可能是 $G=V$。

\item[\textup{c2}.] 如果 $h(x)$ 分解为一个一次因子与一个二次因子，那么 $\theta_1,\theta_2,\theta_3$ 中恰有一个在 $F$ 中，设为 $\theta_1$。则 $G$ 稳定 $\theta_1$，但不稳定 $\theta_2$ 和 $\theta_3$，所以 $G\leq D_8$ 且 $G\nleq V$。这留下两种可能：$G=D_8$ 或 $G=C$。区分它们的一种方法是观察到 $F(\sqrt D)$ 是 $G$ 中落在 $A_4$ 内的元素的固定域。对正在考虑的两种情形，有 $D_8\cap A_4=V$，而 $C\cap A_4=\{1,(13)(24)\}$。第一个群在 $g(y)$ 的根上传递，第二个不是。因此第一种情形当且仅当 $g(y)$ 在 $F(\sqrt D)$ 上不可约时发生。于是，通过在 $F(\sqrt D)$ 中分解 $g(y)$，可以完全确定 $G$，从而在所有情形中完全确定 Galois 群。（参见下面练习以及下一节；在那里说明，在 $\mathbb Q$ 上，如果 $D$ 不是两个平方之和，则 Galois 群不可能是 $4$ 阶循环群；特别地，若 $D<0$，则不可能。）
\end{enumerate}

下一节末尾我们将给出四次多项式根的显式公式。

\paragraph{代数基本定理}
本节最后给出代数基本定理的两个证明。我们需要关于域 $\mathbb C$ 的两个事实：
\begin{enumerate}[label=(\alph*)]
\item 每个实系数奇数次多项式在实数中有根。等价地，$\mathbb R$ 没有非平凡的奇数次有限扩张。
\item 复系数二次多项式在 $\mathbb C$ 中有根。等价地，$\mathbb C$ 没有二次扩张。
\end{enumerate}

第一个结果来自微积分中的介值定理，因为奇数次首一多项式 $f(x)\in\mathbb R[x]$ 的图像在很大的负值处为负，而在很大的正值处为正，因此必在某处与 $x$ 轴相交。它与第二个陈述的等价性如下：$\mathbb R$ 的任意有限扩张都是单扩张，若这个扩张为奇数次，则本原元的极小多项式为奇数次，因而既在 $\mathbb R$ 上不可约又在 $\mathbb R$ 中有根，故次数只能是 $1$。

第二个结果可由直接计算得到。由二次公式，只需说明每个复数 $a=a+bi$（$a,b\in\mathbb R$）在 $\mathbb C$ 中都有平方根。写 $a=re^{i\theta}$，其中 $r\geq0$，$\theta\in[0,2\pi)$，则 $\sqrt r e^{i\theta/2}$ 是 $a$ 的平方根。更显式地，令 $c\in\mathbb R$ 是实数
\[
\frac{a+\sqrt{a^2+b^2}}{2}
\]
的平方根，令 $d\in\mathbb R$ 是实数
\[
\frac{-a+\sqrt{a^2+b^2}}{2}
\]
的平方根，并选取两个平方根的符号，使得 $cd$ 与 $b$ 同号。展开即可看出 $(c+di)^2=a+bi$。

\noindent\textbf{定理 35（代数基本定理）.} 每个次数为 $n$ 的多项式 $f(x)\in\mathbb C[x]$ 在 $\mathbb C$ 中恰有 $n$ 个根（按重数计）。等价地，$\mathbb C$ 是代数闭的。

\begin{proof}[证明一]
只需证明每个多项式 $f(x)\in\mathbb C[x]$ 在 $\mathbb C$ 中有根。令 $\tau$ 表示复共轭自同构。如果 $f(x)$ 在 $\mathbb C$ 中没有根，则共轭多项式 $\overline f(x)=\tau f(x)$（由对 $f(x)$ 的系数作用 $\tau$ 得到）也没有根，因为其根是 $f(x)$ 的根的共轭。乘积 $f(x)\overline f(x)$ 的系数在复共轭下不变，因此有实系数。所以只需证明实系数多项式在 $\mathbb C$ 中有根。

设 $f(x)$ 是次数为 $n$ 的实系数多项式，并写 $n=2^km$，其中 $m$ 为奇数。我们对 $k$ 归纳证明 $f(x)$ 在 $\mathbb C$ 中有根。若 $k=0$，则 $f(x)$ 次数为奇数，由上面事实 (a)，$f(x)$ 在 $\mathbb R$ 中有根，故结论成立。现在设 $k\geq1$。令 $\alpha_1,\alpha_2,\ldots,\alpha_n$ 为 $f(x)$ 的根，并设
\[
K=\mathbb R(\alpha_1,\alpha_2,\ldots,\alpha_n,i).
\]
则 $K$ 是 $\mathbb R$ 的一个 Galois 扩张，包含 $\mathbb C$ 以及 $f(x)$ 的所有根。对任意 $t\in\mathbb R$，考虑多项式
\[
L_t=\prod_{1\leq i<j\leq n}\left[x-(\alpha_i+\alpha_j+t\alpha_i\alpha_j)\right].
\]
$K/\mathbb R$ 的任意自同构都置换这个乘积中的因子，所以 $L_t$ 的系数被 $\operatorname{Gal}(K/\mathbb R)$ 的所有元素固定。因此 $L_t$ 是实系数多项式。$L_t$ 的次数为
\[
\frac{n(n-1)}{2}=2^{k-1}m(2^km-1)=2^{k-1}m',
\]
其中 $m'$ 是奇数（因为 $k\geq1$）。这个次数中 $2$ 的幂次小于 $k$，所以由归纳假设，$L_t$ 在 $\mathbb C$ 中有根。于是对每个 $t\in\mathbb R$，都有某个 $i,j$（$1\leq i<j\leq n$）使 $\alpha_i+\alpha_j+t\alpha_i\alpha_j\in\mathbb C$。由于 $t$ 有无限多个选择而 $i,j$ 只有有限多种取值，存在某一对 $i,j$（不妨为 $i=1,j=2$）以及两个不同实数 $s,t$，使得
\[
\alpha_1+\alpha_2+s\alpha_1\alpha_2\in\mathbb C,
\qquad
\alpha_1+\alpha_2+t\alpha_1\alpha_2\in\mathbb C.
\]
因为 $s\neq t$，所以 $\alpha_1\alpha_2\in\mathbb C$，从而 $\alpha_1+\alpha_2\in\mathbb C$。于是 $\alpha_1$ 与 $\alpha_2$ 是系数在 $\mathbb C$ 中的二次多项式
\[
x^2-(\alpha_1+\alpha_2)x+\alpha_1\alpha_2
\]
的根，故由事实 (b) 它们属于 $\mathbb C$。这完成了证明。
\end{proof}

\begin{proof}[证明二]
第二个证明仍然使用上面的事实 (a) 和 (b)，但用一个涉及 Galois 群的 Sylow $2$-子群的幂零性的简单群论论证，取代上面关于多项式 $L_t$ 的计算。

令 $f(x)$ 是实系数次数为 $n$ 的多项式，并令 $K$ 为 $f(x)$ 在 $\mathbb R$ 上的分裂域。则 $K(i)$ 是 $\mathbb R$ 的一个 Galois 扩张。令 $G$ 表示它的 Galois 群，令 $P_2$ 表示 $G$ 的 Sylow $2$-子群。$P_2$ 的固定域是 $\mathbb R$ 的一个奇数次扩张，因此由 (a) 是平凡的。

于是 $\operatorname{Gal}(K(i)/\mathbb C)$ 是一个 $2$-群。由于 $2$-群有所有可能阶数的子群（回忆这对任意有限 $p$-群都成立，参见定理 6.1），若这个群非平凡，就会存在 $\mathbb C$ 的一个二次扩张，这与 (b) 矛盾。证明完成。
\end{proof}

代数基本定理最早由 Gauss 在 1816 年严格证明（他 1798 年的博士论文中给出了一个使用几何考虑的证明，需要某些拓扑上的论证加以补充）。上面第一个证明本质上归功于 Laplace 在 1795 年的证明（这也是把多项式记作 $L_t$ 的原因）。Laplace 的证明当时被认为不能接受，是因为他假设了多项式的分裂域存在（也就是假设根存在于某个域中），而这一点在当时尚未建立。优雅的第二个证明是 Artin 的一个简化。


\subsection*{练习}
\begin{enumerate}[label=\arabic*.]
	\item 证明有重根的三次多项式必有一次因子。四次多项式是否也如此？
	\item 确定下列多项式的 Galois 群：
	\begin{enumerate}[label=(\alph*)]
		\item $x^3-x^2-4$；
		\item $x^3-2x+4$；
		\item $x^3-x+1$；
		\item $x^3+x^2-2x-1$。
	\end{enumerate}
	\item 证明对任意 $a,b\in\mathbb F_{p^n}$，若 $x^3+ax+b$ 不可约，则 $-4a^3-27b^2$ 是 $\mathbb F_{p^n}$ 中的平方。
	\item 确定 $x^4-25$ 的 Galois 群。
	\item 确定 $x^4+4$ 的 Galois 群。
	\item 确定 $x^4+3x^3-3x-2$ 的 Galois 群。
	\item 确定 $x^4+2x^2+x+3$ 的 Galois 群。
	\item 确定 $x^4+8x+12$ 的 Galois 群。
	\item 确定 $x^4+4x-1$ 的 Galois 群（参见练习 19）。
	\item 确定 $x^5+x-1$ 的 Galois 群。
	\item 设 $F/\mathbb Q$ 是次数为 $4$ 但不是 Galois 的扩张。证明 $F$ 的 Galois 闭包的 Galois 群或者为 $S_4$，或者为 $A_4$，或者为 $8$ 阶二面体群 $D_8$。证明其 Galois 群为二面体群当且仅当 $F$ 含有 $\mathbb Q$ 的一个二次扩张。
	\item 证明 $\mathbb Q$ 的 $4$ 次扩张 $F$ 可由不可约双二次式 $x^4+ax^2+b$ 的一个根生成，当且仅当 $F$ 含有 $\mathbb Q$ 的一个二次扩张。
	\item 
	\begin{enumerate}[label=(\alph*)]
		\item 设 $\pm\alpha,\pm\beta$ 为 $f(x)=x^4+ax^2+b\in\mathbb Z[x]$ 的根。\footnote{参见 Luise-Charlotte Kappe 和 Bette Warren, “An Elementary Test for the Galois Group of a Quartic Polynomial,” \emph{Amer. Math. Monthly}, 96 (1989), pp. 133--137.} 证明 $f(x)$ 不可约当且仅当 $\alpha^2$ 与 $\alpha\pm\beta$ 都不在 $\mathbb Q$ 中。
		\item 假设 $f(x)$ 不可约，令 $G$ 为 $f(x)$ 的 Galois 群。证明：
		\begin{enumerate}[label=(\roman*)]
			\item $G\cong V$（Klein 四元群）当且仅当 $b$ 是 $\mathbb Q$ 中平方；等价地，$\alpha\beta\in\mathbb Q$。
			\item $G\cong C_4$ 当且仅当 $b(a^2-4b)$ 是 $\mathbb Q$ 中平方；等价地，$\mathbb Q(\alpha\beta)=\mathbb Q(\alpha^2)$。
			\item $G\cong D_8$ 当且仅当 $b$ 与 $b(a^2-4b)$ 都不是 $\mathbb Q$ 中平方；等价地，$\alpha\beta\notin\mathbb Q(\alpha^2)$。
		\end{enumerate}
	\end{enumerate}
	\item 证明对任意不同的奇素数 $p,q$，多项式 $x^4-px^2+q\in\mathbb Q[x]$ 不可约，且 Galois 群为 $8$ 阶二面体群。\footnote{同上。}
	\item 证明对每个素数 $p$，$x^4+px+p\in\mathbb Q[x]$ 不可约；当 $p\ne3,5$ 时其 Galois 群为 $S_4$；当 $p=3$ 时为 $8$ 阶二面体群；当 $p=5$ 时为 $4$ 阶循环群。\footnote{同上。}
	\item 确定 $x^4+8x^2+8x+4$ 在 $\mathbb Q$ 上的 Galois 群。确定其分裂域的哪些子域在 $\mathbb Q$ 上是 Galois 的；对每个这样的子域，给出一个 $f(x)\in\mathbb Q[x]$，使该子域是其在 $\mathbb Q$ 上的分裂域。
	\item 将 $x^4-7$ 在 $\mathbb Q$ 上的 Galois 群显式写成根上的置换群。
	\item 设 $\theta$ 为 $x^3-3x+1$ 的一个根。证明该多项式的分裂域为 $\mathbb Q(\theta)$，且 Galois 群为 $3$ 阶循环群。特别地，另外两个根可写为 $a+b\theta+c\theta^2$，其中 $a,b,c\in\mathbb Q$。显式求出这两个根。
	\item 设 $f(x)$ 是 $\mathbb Q[x]$ 中不可约四次多项式，判别式为 $D$。令 $K$ 为其分裂域，并把它看作 $\mathbb C$ 的子域。
	\begin{enumerate}[label=(\alph*)]
		\item 证明 $\mathbb Q(\sqrt D)\subseteq K$。
		\item 令 $\tau$ 为复共轭，$\tau_K$ 为复共轭在 $K$ 上的限制。证明 $\tau_K\in\operatorname{Gal}(K/\mathbb Q)$，并且其阶为 $1$ 或 $2$，取决于 $K$ 中元素是否全为实数。
		\item 证明若 $D<0$，则 $K/\mathbb Q$ 不可能是 $4$ 次循环扩张。
		\item 一般地证明：若 $D<0$ 为无平方因子的整数，则 $\mathbb Q(\sqrt D)$ 不是任何循环四次域的子域（也参见第7节练习19）。
	\end{enumerate}
	\item 确定 $(x^3-2)(x^3-3)$ 在 $\mathbb Q$ 上的 Galois 群。确定所有包含 $\mathbb Q(\rho)$ 的子域，其中 $\rho$ 是本原三次单位根。
	\item 设 $G\leq S_n$，并设 $\sigma_1,
	\ldots,\sigma_k$ 生成 $G$。若函数 $f(x_1,
	\ldots,x_n)$ 被每个 $\sigma_i$ 固定，证明它被 $G$ 固定。
	\item（Newton 公式）设 $f(x)$ 是 $n$ 次首一多项式，根为 $\alpha_1,
	\ldots,\alpha_n$。令 $s_i$ 为根的 $i$ 次基本对称函数，并约定 $s_i=0$（$i>n$）。令 $p_i=\alpha_1^i+\cdots+\alpha_n^i$。证明 Newton 公式：
	\[
	p_1-s_1=0,
	\]
	\[
	p_2-s_1p_1+2s_2=0,
	\]
	\[
	p_3-s_1p_2+s_2p_1-3s_3=0,
	\]
	一般地，若 $1\leq k\leq n$，则
	\[
	p_k-s_1p_{k-1}+s_2p_{k-2}-\cdots+(-1)^{k-1}s_{k-1}p_1+(-1)^k k s_k=0,
	\]
	若 $k>n$，则
	\[
	p_k-s_1p_{k-1}+s_2p_{k-2}-\cdots+(-1)^ns_n p_{k-n}=0.
	\]
	\item 
	\begin{enumerate}[label=(\alph*)]
		\item 若 $x+y+z=1$，$x^2+y^2+z^2=2$，$x^3+y^3+z^3=3$，求 $x^4+y^4+z^4$。
		\item 一般地证明 $x,y,z$ 不是有理数，但对每个正整数 $n$，$x^n+y^n+z^n$ 是有理数。
	\end{enumerate}
	\item 证明特征 $0$ 域上的 $n\times n$ 矩阵 $A$ 幂零当且仅当对所有 $k\ge0$，$\operatorname{tr}(A^k)=0$。
	\item 证明特征 $0$ 域上的两个 $n\times n$ 矩阵 $A,B$ 有相同特征多项式，当且仅当对所有 $k\ge0$，$\operatorname{tr}(A^k)=\operatorname{tr}(B^k)$。
	\item 利用任意两个 $n\times n$ 矩阵 $A,B$ 满足 $\operatorname{tr}(AB)=\operatorname{tr}(BA)$，证明在特征 $0$ 域上 $AB$ 与 $BA$ 有相同特征多项式（该结论对任意特征的域也成立）。
	\item 设 $f(x)$ 是 $n$ 次首一多项式，根为 $\alpha_1,
	\ldots,\alpha_n$。
	\begin{enumerate}[label=(\alph*)]
		\item 证明 $f(x)$ 的判别式 $D$ 是 Vandermonde 行列式的平方：
		\[
		D=\left|\begin{matrix}
			1&\alpha_1&\alpha_1^2&\cdots&\alpha_1^{n-1}\\
			1&\alpha_2&\alpha_2^2&\cdots&\alpha_2^{n-1}\\
			\vdots&\vdots&\vdots&&\vdots\\
			1&\alpha_n&\alpha_n^2&\cdots&\alpha_n^{n-1}
		\end{matrix}\right|^2
		=\prod_{i>j}(\alpha_i-\alpha_j)^2.
		\]
		\item 取上述 Vandermonde 矩阵，左乘其转置并取行列式，证明
		\[
		D=\det\begin{pmatrix}
			p_0&p_1&p_2&\cdots&p_{n-1}\\
			p_1&p_2&p_3&\cdots&p_n\\
			\vdots&\vdots&\vdots&&\vdots\\
			p_{n-1}&p_n&p_{n+1}&\cdots&p_{2n-2}
		\end{pmatrix},
		\]
		其中 $p_i=\alpha_1^i+\cdots+\alpha_n^i$，它可由 Newton 公式用 $f(x)$ 的系数计算。这给出计算判别式的有效方法。
	\end{enumerate}
	\item 设 $\alpha$ 为不可约多项式 $f(x)\in F[x]$ 的根，$K=F(\alpha)$。令 $D$ 为 $f(x)$ 的判别式。证明
	\[
	D=(-1)^{n(n-1)/2}N_{K/F}(f'(\alpha)),
	\]
	其中 $f'(x)=D_xf(x)$ 为导数。
	\item 设 $F$ 是域，
	\[
	f(x)=a_nx^n+a_{n-1}x^{n-1}+\cdots+a_0,
	\quad
	g(x)=b_mx^m+b_{m-1}x^{m-1}+\cdots+b_0
	\]
	为 $F[x]$ 中两个多项式。
	\begin{enumerate}[label=(\alph*)]
		\item 证明 $f(x)$ 与 $g(x)$ 有公共根（等价地，有公共因子）的充要条件是存在次数至多 $m-1$ 的多项式 $a(x)$ 与次数至多 $n-1$ 的多项式 $b(x)$，使
		\[
		a(x)f(x)=b(x)g(x).
		\]
		\item 将 $a(x),b(x)$ 写出并比较 $a(x)f(x)=b(x)g(x)$ 的系数，得到关于其系数的 $n+m$ 个线性方程。证明该线性方程组有非零解当且仅当 Sylvester 行列式 $R(f,g)$ 为零。该行列式称为 $f,g$ 的\emph{结式}，它是一个 $(n+m)\times(n+m)$ 行列式，前 $m$ 行由 $f$ 的系数移位组成，后 $n$ 行由 $g$ 的系数移位组成。
	\end{enumerate}
	\item 
	\begin{enumerate}[label=(\alph*)]
		\item 沿用前题记号，写出 Sylvester 矩阵 $R$ 与列向量 $(x^{n+m-1},x^{n+m-2},\ldots,x,1)^t$ 的矩阵关系，使右边为由 $x^{m-1}f(x),\ldots,f(x),x^{n-1}g(x),\ldots,g(x)$ 组成的列向量。
		\item 令 $R'$ 为 $R$ 的余因子矩阵的转置，如第11章定理30，则 $R'R=R(f,g)I$。利用上面的矩阵等式证明存在 $r(x),s(x)\in F[x]$ 使
		\[
		R(f,g)=r(x)f(x)+s(x)g(x).
		\]
		即两个多项式的结式是它们在 $F[x]$ 中的线性组合。
	\end{enumerate}
	\item 把 $f(x)$ 与 $g(x)$ 看作一般多项式。设 $f$ 的根为 $x_1,
	\ldots,x_n$，$g$ 的根为 $y_1,
	\ldots,y_m$。
	\begin{enumerate}[label=(\alph*)]
		\item 展开行列式证明 $R(f,g)$ 对 $a_i$ 的系数齐次次数为 $m$，对 $b_j$ 的系数齐次次数为 $n$。
		\item 证明 $R(f,g)$ 等于 $a_n^m b_m^n$ 乘以一个关于 $x_1,
		\ldots,x_n$ 和 $y_1,
		\ldots,y_m$ 的对称函数。
		\item 因为当 $x_i=y_j$ 时 $f,g$ 有公共根，证明 $R(f,g)$ 被每个 $x_i-y_j$ 整除。由次数考虑推出
		\[
		R(f,g)=a_n^m b_m^n\prod_{i=1}^n\prod_{j=1}^m(x_i-y_j).
		\]
		\item 证明该乘积也可写成
		\[
		R(f,g)=a_n^m\prod_{i=1}^n g(x_i)=(-1)^{nm}b_m^n\prod_{j=1}^m f(y_j).
		\]
		这给出把 $g$ 在 $f$ 的根上取值的乘积与把 $f$ 在 $g$ 的根上取值的乘积之间的互反关系。
	\end{enumerate}
	\item 现在取特殊情形 $g(x)=f'(x)$，其中
	\[
	f(x)=x^n+a_{n-1}x^{n-1}+\cdots+a_0
	\]
	根为 $\alpha_1,
	\ldots,\alpha_n$。利用
	\[
	R(f,f')=\prod_{i=1}^n f'(\alpha_i)
	\]
	证明
	\[
	D=(-1)^{n(n-1)/2}R(f,f'),
	\]
	其中 $D$ 为 $f(x)$ 的判别式。
	\item 
	\begin{enumerate}[label=(\alph*)]
		\item 证明奇素数 $p$ 的 $p$ 次单位根分圆多项式 $\Phi_p(x)$ 的判别式为
		\[
		(-1)^{(p-1)/2}p^{p-2}.
		\]
		\item 证明对奇素数 $p$，
		\[
		\mathbb Q\bigl(\sqrt{(-1)^{(p-1)/2}p}\bigr)\subseteq \mathbb Q(\zeta_p).
		\]
	\end{enumerate}
	\item 利用前一题证明 $\mathbb Q$ 的每个二次扩张都包含在某个分圆扩张中（Kronecker-Weber 定理的一个特例）。
	\item 证明多项式 $x^n+px+q$ 的判别式为
	\[
	(-1)^{n(n-1)/2}n^nq^{n-1}+(-1)^{(n-1)(n-2)/2}(n-1)^{n-1}p^n.
	\]
	\item 证明
	\[
	x^n+nx^{n-1}+n(n-1)x^{n-2}+\cdots+n(n-1)\cdots(3)(2)x+n!
	\]
	的判别式为
	\[
	(-1)^{n(n-1)/2}(n!)^n.
	\]
	\item 下面练习 37--43 给出两种把对称函数写成基本对称函数的方法。设 $f(x_1,
	\ldots,x_n)$ 是关于 $x_1,
	\ldots,x_n$ 对称的多项式。
	\begin{enumerate}[label=(\alph*)]
		\item 证明每个多项式都可写成齐次多项式之和。证明若 $f$ 对称，则每个齐次部分也对称。
		\item 证明基本对称函数中的单项式 $Bs_1^{a_1}s_2^{a_2}\cdots s_n^{a_n}$ 作为 $x_1,
		\ldots,x_n$ 的多项式，是次数 $a_1+2a_2+\cdots+na_n$ 的齐次多项式。
	\end{enumerate}
	因此写出 $f$ 时可假设 $f$ 齐次。
	\item 取字典序 $x_1>x_2>\cdots>x_n$。若 $f$ 中含有单项式 $Ax_1^{a_1}\cdots x_n^{a_n}$，则由对称性也含有其所有置换。取其中字典序最大的单项式，于是 $a_1\ge a_2\ge\cdots\ge a_n\ge0$。
	\begin{enumerate}[label=(\alph*)]
		\item 证明基本对称函数中的单项式
		\[
		As_1^{a_1-a_2}s_2^{a_2-a_3}\cdots s_n^{a_n}
		\]
		有相同的字典序首项。
		\item 证明从 $f$ 中减去这个单项式后，得到 $0$ 或者一个同次数、所有项字典序更小的对称多项式。
		\item 证明反复执行此过程会终止，并把 $f$ 表成基本对称函数的多项式。
	\end{enumerate}
	\item 利用练习38的算法证明关于 $x_1,
	\ldots,x_n$ 的对称多项式可唯一地表示为基本对称函数的多项式。
	\item 用练习38的方法把下列对称函数写成基本对称函数的多项式：
	\begin{enumerate}[label=(\alph*)]
		\item $(x_1-x_2)^2$；
		\item $x_1^2+x_2^2+x_3^2$；
		\item $x_1^2x_2^2+x_1^2x_3^2+x_2^2x_3^2$。
	\end{enumerate}
	\item 用练习38的方法把 $\sum_{i\ne j}x_i^2x_j$ 写成基本对称函数的多项式。
	\item 已知对称多项式可唯一写成基本对称函数的多项式。设 $f$ 为次数 $M$ 的齐次对称多项式，$N$ 为 $f$ 中任一变量出现的最大次数。
	\begin{enumerate}[label=(\alph*)]
		\item 由约束
		\[
		a_1+2a_2+\cdots+na_n=M,
		\qquad
		a_1+a_2+\cdots+a_n\le N
		\]
		确定可能出现的单项式 $A_is_1^{a_1}s_2^{a_2}\cdots s_n^{a_n}$。
		\item 由于恒等式
		\[
		f(x_1,
		\ldots,x_n)=\sum A_is_1^{a_1}s_2^{a_2}\cdots s_n^{a_n}
		\]
		对任意代入都成立，每次给 $x_i$ 赋值会给出关于 $A_i$ 的线性关系；足够多的代入即可确定 $A_i$。
	\end{enumerate}
	\item 证明
	\[
	(x_1+x_2-x_3-x_4)(x_1+x_3-x_2-x_4)(x_1+x_4-x_2-x_3)
	\]
	关于 $x_1,x_2,x_3,x_4$ 对称，并用前一过程证明它等于
	\[
	s_1^3-4s_1s_2+8s_3.
	\]
	\item 把下列式子表示为基本对称函数：
	\begin{enumerate}[label=(\alph*)]
		\item $\sum_{i\ne j}x_i^2x_j$；
		\item $\sum_{i,j,k\ \mathrm{distinct}}x_i^2x_jx_k$；
		\item $\sum_{i,j,k\ \mathrm{distinct}}x_i^2x_j^2x_k^2$。
	\end{enumerate}
	\item 设 $\alpha_1,
	\alpha_2,
	\alpha_3,
	\alpha_4$ 为 $\mathbb Q$ 上四次多项式 $f(x)$ 的根。证明
	\[
	\alpha_1\alpha_2+\alpha_3\alpha_4,
	\quad
	\alpha_1\alpha_3+\alpha_2\alpha_4,
	\quad
	\alpha_1\alpha_4+\alpha_2\alpha_3
	\]
	被 $f(x)$ 的 Galois 群置换。推出这些元素是某个有理系数三次多项式的根（这也常称为 $f(x)$ 的预解三次方程）。
	\item 若 $f(x)=x^3+px+q\in\mathbb Z[x]$ 不可约，证明其判别式 $D=-4p^3-27q^2$ 是不等于 $0,\pm1$ 的整数。
	\item 证明每个有限群都可以作为某个形如 $F(x_1,x_2,
	\ldots,x_n)/E$ 的域扩张的 Galois 群。
	\item 设 $F$ 是特征 $0$ 的域，并且每个三次多项式在 $F$ 中都有根。设 $f(x)$ 是 $F$ 上不可约四次多项式，其判别式在 $F$ 中为平方。确定 $f(x)$ 的 Galois 群。
	\item 本题确定 $f(x)=x^6-2x^3-2$ 在 $\mathbb Q$ 上的分裂域 $K$。
	\begin{enumerate}[label=(\alph*)]
		\item 证明 $f(x)$ 在 $\mathbb Q$ 上不可约，其根是 $1\pm\sqrt3$ 的三个三次根。
		\item 证明 $K$ 含有三次单位根域 $\mathbb Q(\sqrt{-3})$，且含有 $\mathbb Q(\sqrt3)$，因而含有双二次域 $F=\mathbb Q(i,\sqrt3)$。取 (a) 中两个根的乘积证明 $K$ 含有 $\theta$，并推出 $K$ 是 $L=\mathbb Q(\theta,i,\sqrt3)$ 的扩张。
		\item 证明 $[L:\mathbb Q]=12$，且 $K$ 由 $L$ 添入 $L$ 中某元素的三次根得到，所以 $[K:\mathbb Q]=12$ 或 $36$。
		\item 证明若 $[K:\mathbb Q]=12$，则 $K=\mathbb Q(\theta,i,\sqrt3)$，且 $\operatorname{Gal}(K/\mathbb Q)$ 同构于 $2$ 阶循环群与 $S_3$ 的直积。证明若 $[K:\mathbb Q]=12$，则 $K$ 中唯一的实三次子域为 $\mathbb Q(\theta)$。
		\item 取 (a) 中两个实根的商，证明 $\sqrt[3]{2+\sqrt3}$ 与 $\sqrt[3]{2-\sqrt3}$ 都属于 $K$。令
		\[
		\alpha=\sqrt[3]{2+\sqrt3}+\sqrt[3]{2-\sqrt3}.
		\]
		证明 $\alpha$ 是不可约三次方程 $x^3-3x-4$ 的实根，其判别式为 $-2^2 3^4$。推出 $\mathbb Q(\alpha)$ 的 Galois 闭包含有 $\mathbb Q(i)$，所以 $\mathbb Q(\alpha)\ne\mathbb Q(\theta)$。
		\item 由 (e) 推出 $G=\operatorname{Gal}(K/\mathbb Q)$ 阶为 $36$。显式确定 $G$ 的所有元素，并证明 $G\cong S_3\times S_3$。
	\end{enumerate}
	\item 证明 $x^6-4x^3+1$ 在 $\mathbb Q$ 上的 Galois 群同构于 $12$ 阶二面体群。（提示：两个实根互为倒数。）
	\item（任意域上不可约三次多项式 Galois 群的判别准则）设 $K$ 是域，
	\[
	f(x)=x^3+ax^2+bx+c\in K[x]
	\]
	不可约，因此 $f$ 的 Galois 群为 $S_3$ 或 $A_3$。
	\begin{enumerate}[label=(\alph*)]
		\item 证明 $f$ 的 Galois 群为 $A_3$ 当且仅当二次多项式
		\[
		g(x)=x^2+(ab-3c)x+(b^3+a^3c-6abc+9c^2)
		\]
		在 $K$ 中有根。若 $\alpha,\beta,\gamma$ 是 $f$ 的根，证明 Galois 群为 $A_3$ 当且仅当
		\[
		\theta=\alpha\beta^2+\beta\gamma^2+\gamma\alpha^2
		\]
		属于 $K$，且 $\theta$ 是 $g(x)$ 的根。证明 $g(x)$ 的判别式等于 $f(x)$ 的判别式。
		\item 若 $\operatorname{ch}(K)\ne2$，证明 $f(x)$ 的 Galois 群为 $A_3$ 当且仅当 $f(x)$ 的判别式是 $K$ 中平方。
		\item 若 $\operatorname{ch}(K)=2$，证明 $f(x)$ 的判别式总是平方。证明可取 $f(x)$ 为 $x^3+px+q$ 的形式，且 $f(x)$ 的 Galois 群为 $A_3$ 当且仅当二次式
		\[
		x^2+qx+(p^3+q^2)
		\]
		在 $K$ 中有根；等价地，$(p^3+q^2)/q^2$ 属于 Artin-Schreier 映射 $x\mapsto x^2-x$ 的像。
		\item 若 $K=\mathbb F_2(t)$，其中 $t$ 在 $\mathbb F_2$ 上超越，证明
		\[
		x^3+t^2x+t^3,
		\quad x^3+(t^2+t+1)x+(t^2+t+1),
		\]
		\[
		x^3+(t^2+t+1)x+(t^3+t^2+t)
		\]
		的 Galois 群为 $A_3$，而 $x^3+t^2x+t$ 与 $x^3+x+t$ 的 Galois 群为 $S_3$。
	\end{enumerate}
	\item 本题证明 Sturm 定理，用来确定实多项式 $f(x)\in\mathbb R[x]$ 在区间 $[a,b]$ 中的实根数。可先除去重根，假设所有实根都是单根。对 $f_0(x)=f(x)$ 和 $f_1(x)=f'(x)$ 使用 Euclidean 算法，每一步取负余式，得多项式序列
	\[
	f(x),f'(x),f_2(x),\ldots,f_n(x)
	\]
	满足
	\[
	f_{i-1}(x)=q_i(x)f_i(x)-f_{i+1}(x),
	\quad i=1,
	\ldots,n-1,
	\]
	其中 $f_n(x)\in\mathbb R$ 为非零常数。
	\begin{enumerate}[label=(\alph*)]
		\item 证明相邻多项式 $f_i(x),f_{i+1}(x)$ 无公共零点。
		\item 若 $f_i(c)=0$，证明 $f_{i-1}(c)$ 与 $f_{i+1}(c)$ 一正一负。
		\item 对实数 $a$，令 $V(a)$ 表示 Sturm 序列
		\[
		f(a),f'(a),f_2(a),\ldots,f_n(a)
		\]
		中忽略零项后的符号变化次数。若 $a<\beta$ 且两个 Sturm 序列中元素均非零，证明除非存在 $a<c<\beta$ 且某个 $i$ 满足 $f_i(c)=0$，否则两个 Sturm 序列符号完全相同，因而 $V(a)=V(\beta)$。
		\item 若 $f_j(c)=0$，证明存在充分小区间 $(a,\beta)$ 包含 $c$，使 $f_j(x)$ 在该区间内除 $c$ 外无零点。
		\item 若 $j\ge1$，证明 $f_{j-1}(a),f_j(a),f_{j+1}(a)$ 与 $f_{j-1}(\beta),f_j(\beta),f_{j+1}(\beta)$ 中的符号变化次数相同。
		\item 若 $j=0$，证明 $f(a),f'(a)$ 中的符号变化次数比 $f(\beta),f'(\beta)$ 中的符号变化次数多 $1$。
		\item 证明 Sturm 定理：若 $f$ 的实根全为单根，且 $f(a),f(b)$ 均非零，则 $[a,b]$ 中 $f$ 的实零点数为 $V(a)-V(b)$。
		\item 若 $f(x)=x^5+px+q\in\mathbb R[x]$ 有单根，证明上述序列为
		\[
		f(x),\quad 5x^4+p,
		\quad (-4p/5)x+q,
		\quad -D/(256p^4),
		\]
		其中 $D=256p^5+3125q^4$ 是 $f$ 的判别式。推出：当 $p>0$ 时 $f$ 恰有一个实根；当 $p<0$ 时，$f$ 有 $1$ 个或 $3$ 个实根，取决于 $D>0$ 或 $D<0$。
	\end{enumerate}
\end{enumerate}

\section{可解扩张与根式扩张：五次方程不可解}

我们现在研究用根式求多项式的根的问题，也就是说，用加、减、乘、除以及抽取 $n$ 次根这些代数运算来表示根。二次多项式的根的公式，是初等代数中熟悉的二次公式；下面我们还将推导出三次和四次多项式根的类似公式。然而，对于次数 $\geq 5$ 的多项式，我们将看到这样的公式是不可能存在的——这就是 Abel 关于一般五次方程不可解的定理。其原因相当简单：我们将看到，一个多项式可由根式求解，当且仅当它的 Galois 群是可解群（这也解释了术语的来源），而当 $n\geq 5$ 时，群 $S_n$ 不是可解的。

我们首先讨论简单根式扩张，即通过把 $F$ 中某个元素 $a$ 的 $n$ 次根添入域 $F$ 而得到的扩张。由于当 $a\in F$ 时，多项式 $x^n-a$ 的所有根只相差 $n$ 次单位根的因子，添入其中一个根得到 Galois 扩张，当且仅当这个域含有 $n$ 次单位根。简单根式扩张在基域 $F$ 已经含有相应单位根时表现最好。对 $a\in F$，符号 $\sqrt[n]{a}$ 将用来表示多项式 $x^n-a\in F[x]$ 的任意一个根。

\noindent\textbf{定义.} 扩张 $K/F$ 称为\emph{循环扩张}，如果它是 Galois 扩张，并且它的 Galois 群是循环群。

\noindent\textbf{命题 36.} 设 $F$ 是一个特征不整除 $n$ 且含有 $n$ 次单位根的域。那么对 $a\in F$，扩张 $F(\sqrt[n]{a})$ 在 $F$ 上是次数整除 $n$ 的循环扩张。

\begin{proof}
扩张 $K=F(\sqrt[n]{a})$ 在 $F$ 上是 Galois 的，因为当 $F$ 含有 $n$ 次单位根时，它是 $x^n-a$ 的分裂域。对任意 $\sigma\in\operatorname{Gal}(K/F)$，$\sigma(\sqrt[n]{a})$ 是这个多项式的另一个根，因此
\[
\sigma(\sqrt[n]{a})=\zeta_\sigma\sqrt[n]{a}
\]
对某个 $n$ 次单位根 $\zeta_\sigma$ 成立。这给出映射
\[
\operatorname{Gal}(K/F)\longrightarrow \mu_n,
\qquad
\sigma\longmapsto \zeta_\sigma,
\]
其中 $\mu_n$ 表示 $n$ 次单位根的群。由于 $F$ 含有 $\mu_n$，每个 $n$ 次单位根都被 $\operatorname{Gal}(K/F)$ 的每个元素固定。因此
\[
\begin{aligned}
\sigma\tau(\sqrt[n]{a})
&=\sigma(\zeta_\tau\sqrt[n]{a})\\
&=\zeta_\tau\sigma(\sqrt[n]{a})\\
&=\zeta_\tau\zeta_\sigma\sqrt[n]{a}
=\zeta_\sigma\zeta_\tau\sqrt[n]{a},
\end{aligned}
\]
这说明 $\zeta_{\sigma\tau}=\zeta_\sigma\zeta_\tau$，于是上面的映射是同态。它的核正是固定 $\sqrt[n]{a}$ 的自同构，也就是恒等自同构。因此 $\operatorname{Gal}(K/F)$ 嵌入阶为 $n$ 的循环群 $\mu_n$ 中，命题得证。
\end{proof}

现在设 $K$ 是域 $F$ 上次数为 $n$ 的任意循环扩张，其中 $F$ 的特征不整除 $n$，且 $F$ 含有 $n$ 次单位根。令 $\sigma$ 为循环群 $\operatorname{Gal}(K/F)$ 的一个生成元。

\noindent\textbf{定义.} 对 $\alpha\in K$ 和任意 $n$ 次单位根 $\zeta$，定义 \emph{Lagrange 预解式} $(\alpha,\zeta)\in K$ 为
\[
(\alpha,\zeta)=\alpha+\zeta\sigma(\alpha)+\zeta^2\sigma^2(\alpha)+\cdots+\zeta^{n-1}\sigma^{n-1}(\alpha).
\]

如果把自同构 $\sigma$ 作用到 $(\alpha,\zeta)$ 上，则得到
\[
\sigma(\alpha,\zeta)=\sigma\alpha+\zeta\sigma^2(\alpha)+\zeta^2\sigma^3(\alpha)+\cdots+\zeta^{n-1}\sigma^n(\alpha),
\]
因为 $\zeta$ 是基域 $F$ 的元素，所以被 $\sigma$ 固定。我们在 $\mu_n$ 中有 $\zeta^n=1$，在 $\operatorname{Gal}(K/F)$ 中有 $\sigma^n=1$，所以这可写成
\[
\begin{aligned}
\sigma(\alpha,\zeta)
&=\sigma\alpha+\zeta\sigma^2(\alpha)+\zeta^2\sigma^3(\alpha)+\cdots+\zeta^{-1}\alpha\\
&=\zeta^{-1}\bigl(\alpha+\zeta\sigma(\alpha)+\zeta^2\sigma^2(\alpha)+\cdots+\zeta^{n-1}\sigma^{n-1}(\alpha)\bigr)\\
&=\zeta^{-1}(\alpha,\zeta). \qquad\text{(14.19)}
\end{aligned}
\]
于是
\[
\sigma(\alpha,\zeta)^n=(\zeta^{-1})^n(\alpha,\zeta)^n=(\alpha,\zeta)^n,
\]
所以 $(\alpha,\zeta)^n$ 被 $\operatorname{Gal}(K/F)$ 固定，因此对任意 $\alpha\in K$，$(\alpha,\zeta)^n$ 都是 $F$ 的元素。

令 $\zeta$ 为本原 $n$ 次单位根。由自同构 $1,\sigma,\ldots,\sigma^{n-1}$ 的线性无关性（定理 7），存在 $\alpha\in K$ 使得 $(\alpha,\zeta)\neq 0$。迭代 (19) 得
\[
\sigma^i(\alpha,\zeta)=\zeta^{-i}(\alpha,\zeta),\qquad i=0,1,\ldots,
\]
从而当 $i<n$ 时，$\sigma^i$ 不固定 $(\alpha,\zeta)$。因此这个元素不能属于 $K$ 的任何真子域，所以 $K=F((\alpha,\zeta))$。由于上面已经证明 $(\alpha,\zeta)^n=a\in F$，所以
\[
F(\sqrt[n]{a})=F((\alpha,\zeta))=K.
\]
这证明了命题 36 的如下逆命题。

\noindent\textbf{命题 37.} 域 $F$ 上任意次数为 $n$ 的循环扩张，其中 $F$ 的特征不整除 $n$ 且 $F$ 含有 $n$ 次单位根，都具有形式 $F(\sqrt[n]{a})$，其中 $a\in F$。

\paragraph{注记}
上面两个命题构成所谓 Kummer 理论的一部分。若一个群 $G$ 的每个元素 $g\in G$ 都满足 $g^n=1$，则称 $G$ 的指数为 $n$。设 $F$ 是一个特征不整除 $n$ 且含有 $n$ 次单位根的域。如果取元素 $a_1,\\ldots,a_k\in F^\times$，则如命题 36 那样可见扩张
\[
F(\sqrt[n]{a_1},\sqrt[n]{a_2},\ldots,\sqrt[n]{a_k}) \qquad\text{(14.20)}
\]
是 $F$ 的 Abel 扩张，且它的 Galois 群的指数为 $n$。反过来，指数为 $n$ 的任意 Abel 扩张都具有这种形式。

记 $(F^\times)^n$ 为乘法群 $F^\times$ 中所有非零元素的 $n$ 次幂所组成的子群。商群 $F^\times/(F^\times)^n$ 是一个指数为 $n$ 的 Abel 群。(20) 中扩张的 Galois 群同构于 $F^\times/(F^\times)^n$ 中由元素 $a_1,\ldots,a_k$ 生成的群，并且两个如 (20) 的扩张相等，当且仅当它们在 $F^\times/(F^\times)^n$ 中对应的群相等。

因此，$F^\times/(F^\times)^n$ 的有限生成子群分类了含有 $n$ 次单位根且特征不整除 $n$ 的域上的指数为 $n$ 的 Abel 扩张。这样的扩张称为 \emph{Kummer 扩张}。

这些结果推广了上面 $k=1$ 的情形，并且可以用类似的方法证明。

为简单起见，我们现在考虑基域 $F$ 的特征为 $0$ 的情形。正如前面命题中一样，这些结果也适用于那些特征不整除将要抽取的根的阶数的域。

\noindent\textbf{定义.}
\begin{enumerate}[label=(\arabic*)]
\item 若 $\alpha$ 在 $F$ 上代数，并且 $\alpha$ 是某个域 $K$ 的元素，而 $K$ 可通过一串简单根式扩张得到：
\[
F=K_0\subset K_1\subset\cdots\subset K_i\subset K_{i+1}\subset\cdots\subset K_s=K, \qquad\text{(14.21)}
\]
其中
\[
K_{i+1}=K_i(\sqrt[n_i]{a_i})
\]
对某个 $a_i\in K_i$ 成立，$i=0,1,\ldots,s-1$。这里 $\sqrt[n_i]{a_i}$ 表示多项式 $x^{n_i}-a_i$ 的某个根。则称 $\alpha$ \emph{可由根式表示}，或\emph{可用根式求出}。这样的域 $K$ 称为 $F$ 的一个\emph{根式扩张}。
\item 若多项式 $f(x)\in F[x]$ 的所有根都可用根式求出，则称 $f(x)$ \emph{可由根式求解}。
\end{enumerate}

这给出了直观说法的精确定义：$\alpha$ 由连续的代数运算（加、减、乘和除）以及连续的抽根运算得到。例如，在第 5 节末尾遇到的元素
\[
-1+\sqrt{17}+\sqrt{2(17-\sqrt{17})}
+2\sqrt{17+3\sqrt{17}-\sqrt{2(17-\sqrt{17})}-2\sqrt{2(17+\sqrt{17})}}
\]
用于构造正 $17$ 边形；它可由根式表示，并且属于域 $K_4$，其中
\[
\begin{array}{ll}
K_0=\mathbb Q, & \\
K_1=K_0(\sqrt{a_0}), & a_0=17,\\
K_2=K_1(\sqrt{a_1}), & a_1=2(17-\sqrt{17}),\\
K_3=K_2(\sqrt{a_2}), & a_2=2(17+\sqrt{17}),\\
K_4=K_3(\sqrt{a_3}), & a_3=17+3\sqrt{17}-\sqrt{2(17-\sqrt{17})}-2\sqrt{2(17+\sqrt{17})}.
\end{array}
\]
这些扩张中的每一个都是根式扩张。除了平方根以外不需要其他根这一事实，反映了正 $17$ 边形可用直尺和圆规作出。

在考虑根式扩张时，总可以添入单位根，因为根据定义，单位根本身就是根式。这是有用的，因为这样一来循环扩张变成根式扩张，反过来也成立。特别地有：

\noindent\textbf{引理 38.} 若 $\alpha$ 包含在如 (21) 中的根式扩张 $K$ 中，则 $\alpha$ 包含在一个在 $F$ 上 Galois 的根式扩张中，并且其中每个扩张 $K_{i+1}/K_i$ 都是循环扩张。

\begin{proof}
令 $L$ 为 $K$ 在 $F$ 上的 Galois 闭包。对任意 $\sigma\in\operatorname{Gal}(L/F)$，有子域链
\[
F=\sigma K_0\subset \sigma K_1\subset\cdots\subset \sigma K_i\subset \sigma K_{i+1}\subset\cdots\subset \sigma K_s=\sigma K,
\]
其中 $\sigma K_{i+1}/\sigma K_i$ 仍然是简单根式扩张，因为它由元素 $\sigma(\sqrt[n_i]{a_i})$ 生成，而这个元素是 $\sigma(K_i)$ 上方程 $x^{n_i}-\sigma(a_i)$ 的一个根。容易看出两个根式扩张的复合仍是根式扩张：若 $K'$ 是另一个带有子域 $K'_i$ 的根式扩张，先把 $K'_1$ 与域 $K_0,K_1,\ldots,K_s$ 逐步复合，再把这些域与 $K'_2$ 复合，依此类推，使得这个过程中每个单独的扩张都是简单根式扩张。因此所有共轭域 $\sigma(K)$（$\sigma\in\operatorname{Gal}(L/F)$）的复合仍然是根式扩张。由于这个域正是 $L$，我们看到 $\alpha$ 包含在一个 Galois 根式扩张中。

现在对这个 Galois 根式扩张 $K/F$ 中简单根式扩张所涉及的所有根 $\sqrt[n_i]{a_i}$，把相应的 $n_i$ 次单位根都添入 $F$，得到域 $F'$，然后形成 $F'$ 与根式扩张的复合：
\[
F\subseteq F'=F'K_0\subseteq F'K_1\subseteq\cdots\subseteq F'K_i\subseteq F'K_{i+1}\subseteq\cdots\subseteq F'K_s=F'K.
\]
域 $F'K$ 是 $F$ 的 Galois 扩张，因为它是两个 Galois 扩张的复合。从 $F$ 到 $F'=F'K_0$ 的扩张可以写成一条子域链，使得每个单独的扩张都是循环的（这对任意 Abel 扩张都成立）。每个扩张 $F'K_{i+1}/F'K_i$ 是一个简单根式扩张；又由于我们现在在基域中有相应的单位根，命题 36 表明从 $F'$ 到 $F'K$ 的这些单独扩张中的每一个都是循环扩张。因此 $F'K/F$ 是一个在 $F$ 上 Galois 的根式扩张，并有循环中间扩张，证明完成。
\end{proof}

回忆第 3.4 节（也参见第 6.1 节），有限群 $G$ 称为\emph{可解的}，如果存在一条子群链
\[
1=G_s\leq G_{s-1}\leq\cdots\leq G_{i+1}\leq G_i\leq\cdots\leq G_0=G \qquad\text{(14.22)}
\]
使得 $G_i/G_{i+1}$ 是循环群，$i=0,1,\ldots,s-1$。我们已经证明，可解群的子群和商群都是可解的，并且如果 $H\leq G$ 且 $G/H$ 都是可解的，那么 $G$ 是可解的。

现在我们证明 Galois 关于用根式求多项式根与多项式 Galois 群之间的基本联系。我们继续在特征 $0$ 的域 $F$ 上工作；不过容易看出，这个证明也适用于特征不整除 Galois 群的阶以及所涉及根式阶数的任意域。

\noindent\textbf{定理 39.} 多项式 $f(x)$ 可由根式求解，当且仅当它的 Galois 群是可解群。

\begin{proof}
首先假设 $f(x)$ 可由根式求解。那么 $f(x)$ 的每个根都包含在引理中的某个扩张中。这样的扩张的复合 $L$ 由命题 21 仍是同类型的扩张。令 $G_i$ 为对应于子域 $K_i$ 的子群，$i=0,1,\ldots,s-1$。由于
\[
\operatorname{Gal}(K_{i+1}/K_i)=G_i/G_{i+1},
\qquad i=0,1,\ldots,s-1,
\]
得到 Galois 群 $G=\operatorname{Gal}(L/F)$ 是可解群。域 $L$ 含有 $f(x)$ 的分裂域，因此 $f(x)$ 的 Galois 群是可解群 $G$ 的商群，所以也是可解的。

反过来，假设 $f(x)$ 的 Galois 群 $G$ 是可解群，令 $K$ 为 $f(x)$ 的分裂域。取 $G$ 的如 (22) 中的子群链的固定域，得到一条域链
\[
F=K_0\subset K_1\subset\cdots\subset K_i\subset K_{i+1}\subset\cdots\subset K_s=K,
\]
其中 $K_{i+1}/K_i$，$i=0,1,\ldots,s-1$，是次数为 $n_i$ 的循环扩张。令 $F'$ 为 $F$ 上所有阶为 $n_i$ 的单位根（$i=0,1,\ldots,s-1$）生成的分圆域，并形成复合域 $K_i'=F'K_i$。我们得到一列扩张
\[
F\subseteq F'=F'K_0\subseteq F'K_1\subseteq\cdots\subseteq F'K_i\subseteq F'K_{i+1}\subseteq\cdots\subseteq F'K_s=F'K.
\]
扩张 $F'K_{i+1}/F'K_i$ 是次数整除 $n_i$ 的循环扩张，$i=0,1,\ldots,s-1$（由命题 19）。由于现在基域中有适当的单位根，命题 37 表明这些循环扩张中的每一个都是简单根式扩张。因此 $f(x)$ 的每个根都包含在根式扩张 $F'K$ 中，从而 $f(x)$ 可由根式求解。
\end{proof}

\noindent\textbf{推论 40.} 当 $n\geq 5$ 时，一般 $n$ 次方程不能由根式求解。

\begin{proof}
当 $n\geq 5$ 时，群 $S_n$ 不是可解群，这是我们在第 4 章中证明过的。由定理 32 和定理 39，推论立即得到。
\end{proof}

这个推论说明，对于五次多项式的根，不存在类似于二次多项式根的二次公式那样的根式公式。为了给出一个具体的、次数为 $5$、有理系数且其根不能用根式表示的多项式，我们必须展示一个次数为 $5$ 的有理系数多项式，其 Galois 群为 $S_5$（或者 $A_5$，它也不是可解的）（也参见练习 21，它给出了五次方程可解性的一个判据）。

\paragraph{例}
考虑多项式
\[
f(x)=x^5-6x+3\in\mathbb Q[x].
\]
这个多项式由素数 $3$ 的 Eisenstein 判别法不可约。于是这个多项式的分裂域 $K$ 的次数被 $5$ 整除，因为把 $f(x)$ 的一个根添入 $\mathbb Q$ 会生成一个次数为 $5$ 的扩张。因此 Galois 群 $G$ 是 $S_5$ 的一个子群，且阶被 $5$ 整除，所以含有阶为 $5$ 的元素。$S_5$ 中阶为 $5$ 的元素只能是 $5$-轮换，因此 $G$ 含有一个 $5$-轮换。

由于
\[
f(-2)=-17,
\qquad f(0)=3,
\qquad f(1)=-2,
\qquad f(2)=23,
\]
可知 $f(x)$ 在区间 $(-2,0)$、$(0,1)$ 和 $(1,2)$ 中各有一个实根。由中值定理，如果存在 $4$ 个实根，那么导数 $f'(x)=5x^4-6$ 至少有 $3$ 个实零点，但事实并非如此。因此这些就是全部实根（这也可由 Descartes 符号法则容易得到）。由代数基本定理，$f(x)$ 在 $\mathbb C$ 中有 $5$ 个根。因此 $f(x)$ 有两个非实复根。令 $\tau$ 表示 $\mathbb C$ 中复共轭自同构。由于 $f(x)$ 的系数为实数，这两个复根必须被 $\tau$ 互换（因为它们不是实数，所以不被固定）。因此复共轭在 $K$ 上的限制固定 $f(x)$ 的三个根并交换另外两个根。作为 $G$ 的元素，$\tau|_K$ 因而是一个换位。

现在很容易验证，任意 $5$-轮换与任意换位都生成整个 $S_5$。于是 $G=S_5$，所以 $x^5-6x+3$ 的根不能由根式表示。

如这个例子所表明的，通过了解自同构在 $G$ 作为 $S_n$ 的子群时的\emph{循环型}，可以获得关于 Galois 群的大量信息。实践中，这是确定次数 $\geq 5$ 的多项式的 Galois 群最有效的方法（当然，次数越大就越困难，即便只是因为 $S_n$ 的可能子群极为众多）。我们将在下一节描述这个过程。

由定理 39，任意次数 $n\leq 4$ 的多项式都可由根式求解，因为这些 $n$ 的 $S_n$ 是可解群。当 $n=2$ 时，这就是熟悉的二次公式。当 $n=3$ 时，这个公式称为 Cardano 公式（以 Geronimo Cardano（1501--1576）命名），而 $n=4$ 的公式可以化归到三次情形。公式在任意特征不等于 $2,3$ 的域 $F$ 上都成立，这两个特征正是所需根式阶数和可能的 Galois 群阶数（它们是 $S_3$ 和 $S_4$ 的子群）所涉及的特征。为简单起见，我们将在 $\mathbb Q$ 上推导这些公式。

\paragraph{三次方程的根式解：Cardano 公式}
由定理 39 的证明，以及 $S_3$ 的一个如式 (22) 的合成列由 $1\leq A_3\leq S_3$ 给出这一事实，我们应该预期三次方程
\[
f(x)=x^3+ax^2+bx+c
\]
的解（等价地，在代换 $x=y-a/3$ 下，
\[
g(y)=y^3+py+q,
\]
其中
\[
p=\frac13(3b-a^2),
\qquad
q=\frac1{27}(2a^3-9ab+27c)
\]
）会涉及添入三次单位根，以及形成含有这些单位根的 Lagrange 预解式。

令 $\rho$ 表示本原三次单位根，所以 $\rho^2+\rho+1=0$。设 $g(y)$ 的根为 $\alpha,\beta,\gamma$，于是
\[
\alpha+\beta+\gamma=0,
\]
这也是我们把 $f(x)$ 改写成 $g(x)$ 的原因之一。在域 $\mathbb Q(\sqrt D)$ 上，其中 $D$ 是判别式（上一节中计算过），$g(y)$ 的 Galois 群是 $A_3$，即阶为 $3$ 的循环群。如果再添入 $\rho$，则这个扩张是次数 $3$ 的根式扩张，其生成元由一个 Lagrange 预解式给出，如命题 37 的证明中那样。因此考虑元素
\[
(\alpha,1)=\alpha+\beta+\gamma=0,
\]
\[
\theta_1=(\alpha,\rho)=\alpha+\rho\beta+\rho^2\gamma,
\qquad
\theta_2=(\alpha,\rho^2)=\alpha+\rho^2\beta+\rho\gamma.
\]
注意这些预解式的和为
\[
\theta_1+\theta_2=3\alpha, \qquad\text{(14.23)}
\]
因为 $1+\rho+\rho^2=0$。同样，
\[
\rho^2\theta_1+\rho\theta_2=3\beta,
\qquad
\rho\theta_1+\rho^2\theta_2=3\gamma. \qquad\text{(14.23')}
\]
我们也已经在命题 37 之前一般地说明，这些预解式的三次幂必定属于 $\mathbb Q(\sqrt D,\rho)$。展开 $\theta_1^3$ 得
\[
\alpha^3+\beta^3+\gamma^3
+3\rho(\alpha^2\beta+\beta^2\gamma+\alpha\gamma^2)
+3\rho^2(\alpha\beta^2+\beta\gamma^2+\alpha^2\gamma)
+6\alpha\beta\gamma. \qquad\text{(14.24)}
\]
我们有
\[
\begin{aligned}
\sqrt D&=(\alpha-\beta)(\alpha-\gamma)(\beta-\gamma)\\
&=(\alpha^2\beta+\beta^2\gamma+\alpha\gamma^2)
-(\alpha\beta^2+\beta\gamma^2+\alpha^2\gamma).
\end{aligned}
\]
利用这个等式，可把 (24) 写为
\[
\alpha^3+\beta^3+\gamma^3
+3\rho\Bigl[\frac12(S+\sqrt D)\Bigr]
+3\rho^2\Bigl[\frac12(S-\sqrt D)\Bigr]
+6\alpha\beta\gamma, \qquad\text{(14.24')}
\]
其中为简便起见，我们用 $S$ 表示表达式
\[
(\alpha^2\beta+\beta^2\gamma+\alpha\gamma^2)
+(\alpha\beta^2+\beta\gamma^2+\alpha^2\gamma).
\]
由于 $S$ 在根中是对称的，(24') 中的每一个表达式都是 $\alpha,\beta,\gamma$ 中的对称多项式，因此是基本对称函数 $s_1=0$、$s_2=p$ 和 $s_3=-q$ 的多项式。经过一个简短计算得到
\[
\alpha^3+\beta^3+\gamma^3=-3q,
\qquad
S=3q,
\]
所以由 (24') 得（$\rho+\rho^2=-1$ 且 $\rho-\rho^2=\sqrt{-3}$）
\[
\begin{aligned}
\theta_1^3
&=-3q+\frac32\rho(3q+\sqrt D)+\frac32\rho^2(3q-\sqrt D)-6q\\
&=-\frac{27}{2}q+\frac32\sqrt{-3D}. \qquad\text{(14.25)}
\end{aligned}
\]
同样地，得到
\[
\theta_2^3=-\frac{27}{2}q-\frac32\sqrt{-3D}. \qquad\text{(14.25')}
\]
式 (25) 和 (23) 本质上已经给出了三次方程的解。然而还有一个小问题，即需要从 (25) 中的表达式抽取三次根来得到 $\theta_1$ 和 $\theta_2$。每个表达式有三个可能的三次根，这似乎会在 (23) 中产生总共 $9$ 个表达式。事实并非如此，因为 $\theta_1$ 和 $\theta_2$ 并非相互独立（添入其中一个已经给出含所有根的 Galois 扩张）。类似于上面但更容易的计算表明
\[
\theta_1\theta_2=-3p, \qquad\text{(14.26)}
\]
这说明 $\theta_1$ 的三次根的选择决定了 $\theta_2$。利用 $D=-4p^3-27q^2$，我们得到 Cardano 的显式公式如下。

令
\[
A=\sqrt[3]{-\frac{27}{2}q+\frac32\sqrt{-3D}},
\qquad
B=\sqrt[3]{-\frac{27}{2}q-\frac32\sqrt{-3D}},
\]
其中三次根的选择满足 $AB=-3p$。则方程
\[
y^3+py+q=0
\]
的根为
\[
\alpha=\frac{A+B}{3},
\qquad
\beta=\frac{\rho^2A+\rho B}{3},
\qquad
\gamma=\frac{\rho A+\rho^2B}{3}, \qquad\text{(14.27)}
\]
其中
\[
\rho=-\frac12+\frac12\sqrt{-3}.
\]

\paragraph{例}
\begin{enumerate}[label=(\arabic*)]
\item 考虑三次方程 $x^3-x+1=0$。这个三次方程的判别式为
\[
D=-4(-1)^3-27(1)^2=-23,
\]
它不是有理数的平方，所以这个多项式的 Galois 群是 $S_3$。代入上面的公式得到
\[
A=\sqrt[3]{-\frac{27}{2}+\frac32\sqrt{69}},
\qquad
B=\sqrt[3]{-\frac{27}{2}-\frac32\sqrt{69}},
\]
其中选择 $A$ 为实三次根；再由 $AB=3$ 可知 $B$ 也是实数。三次方程的根由 (27) 给出，我们看到有一个实根和两个共轭复根（当然，我们本来无需解出根也能确定这一点）。

\item 考虑方程
\[
x^3+x^2-2x-1=0.
\]
令 $x=s-1/3$，方程变为
\[
s^3-\frac73s-\frac7{27}=0.
\]
乘以 $27$ 以清除分母，并令 $y=3s$，可见 $y$ 满足三次方程
\[
y^3-21y-7=0.
\]
这个三次方程的判别式 $D$ 为
\[
D=-4(-21)^3-27(-7)^2=3^6 7^2,
\]
这说明这个（在 $7$ 处 Eisenstein 的）三次方程的 Galois 群是 $A_3$。代入上面的公式得到
\[
A=3\sqrt[3]{\frac72+\frac{21}{2}\sqrt{-3}},
\qquad
B=3\sqrt[3]{\frac72-\frac{21}{2}\sqrt{-3}},
\]
并且我们三次方程的根可用 $A$ 和 $B$，通过上面的公式表示。这个三次方程出现在试图把一个本原 $7$ 次单位根 $\zeta_7$ 用根式表示的过程中，类似于前面给出的其他小阶单位根的显式公式（参见练习）。

在这个情形中，$g(-5)=-27$，$g(-1)=13$，$g(0)=-7$，$g(5)=13$，所以这个三次方程有 $3$ 个实根。上面给出的这些根的表达式是复数的共轭和。我们稍后将看到这是必要的，即不可能仅用涉及实数的根式来解出这些实根。
\end{enumerate}

一个有理系数三次方程要么有一个实根和两个共轭虚根，要么有三个实根。这两种情形可以由判别式的符号区分：

设在第一种情形中，根为 $a$ 和 $b\pm ic$，其中 $a,b,c$ 为实数且 $c\neq 0$。则
\[
\begin{aligned}
\sqrt D&=[a-(b+ic)][a-(b-ic)][(b+ic)-(b-ic)]\\
&=2ic[(a-b)^2+c^2]
\end{aligned}
\]
是纯虚数，所以判别式 $D$ 为负。于是上面 $A$ 和 $B$ 的公式中可以把二者都取为实数。(27) 中第一个根于是为实数，另外两个根是共轭复根。

如果三个根全为实数，那么显然 $\sqrt D$ 是实数，所以 $D\geq 0$ 是非负实数。如果 $D=0$，则三次方程有重根。若 $D>0$（有时称为 \emph{Casus irreducibilis}），根的公式涉及非实数的根式，如例 2。现在我们证明，对不可约三次方程而言这是必要的。练习概述了如下推广的证明：如果不可约多项式 $f(x)\in\mathbb Q[x]$ 的所有根都是实数，并且其中一个根可以用实根式表示，那么 $f(x)$ 的次数是 $2$ 的幂，$f(x)$ 的 Galois 群是 $2$-群，并且 $f(x)$ 的根可以用直尺和圆规作出。

假设不可约三次多项式 $f(x)$ 有三个实根，并且能够仅用涉及实数的根式表示其中一个根。那么这个三次多项式的分裂域将包含在根式扩张
\[
\mathbb Q=K_0\subset K_1=\mathbb Q(\sqrt D)\subset\cdots\subset K_i\subset K_{i+1}\subset\cdots\subset K_s=K
\]
中，其中每个域 $K_i$，$i=0,1,\ldots,s$，都包含在实数域 $\mathbb R$ 中，并且 $s\geq 2$，因为二次扩张 $\mathbb Q(\sqrt D)$ 不可能含有不可约三次方程的根。我们以 $\mathbb Q(\sqrt D)$ 开始这个根式扩张，是因为在这个域上，多项式的 Galois 群是次数为 $3$ 的循环群。

注意对任意域 $F$，扩张 $F(\sqrt[mn]{a})$ 可以由两个较小的简单根式扩张得到：令
\[
F_1=F(\sqrt[n]{a})
\]
并令 $b=\sqrt[n]{a}\in F_1$，则
\[
F(\sqrt[mn]{a})=F_1(\sqrt[m]{b}).
\]
因此我们总可以假设根式扩张具有形式 $F(\sqrt[p]{a})$，其中 $p$ 是素数。

现在设 $F$ 是实数域 $\mathbb R$ 的一个子域，且 $a\in F$。令 $p$ 为素数，令 $\alpha=\sqrt[p]{a}$ 表示 $a$ 的一个实 $p$ 次根。那么 $[F(\sqrt[p]{a}):F]$ 必须是 $1$ 或 $p$，理由如下。$\alpha$ 在 $F$ 上的所有共轭都与 $\alpha$ 相差一个 $p$ 次单位根因子。因此 $\alpha$ 在 $F$ 上的极小多项式的常数项是 $\alpha^d\zeta$，其中 $d=[F(\sqrt[p]{a}):F]$ 是极小多项式的次数，$\zeta$ 是某个 $p$ 次单位根。由于 $\alpha$ 是实数且 $\alpha^d\zeta\in F$ 是实数，得到 $\zeta=\pm1$，所以 $\alpha^d\in F$。于是若 $d\neq p$，则 $\alpha^d\in F$ 和 $\alpha^p=a\in F$ 蕴含 $\alpha\in F$，所以 $d=1$。

因此，对于上面的根式扩张，我们可以假设 $[K_{i+1}:K_i]$ 是素数 $p_i$，且 $K_{i+1}=K_i(\sqrt[p_i]{a_i})$，其中 $a_i\in K_i$，$i=0,1,\ldots,s-1$。换言之，原来的实根式扩张塔可以细化为一个塔，使得每个相邻根式扩张的次数都是素数。

如果任意一个含有 $\sqrt D$ 的域含有 $f(x)$ 的一个根，那么它就含有 $f(x)$ 的分裂域，因而含有这个三次多项式的所有根。我们假设 $s$ 取成使得 $K_{s-1}$ 不含这个三次多项式的任何根。

考虑扩张 $K_s/K_{s-1}$。域 $K_s$ 含有三次多项式 $f(x)$ 的所有根，而域 $K_{s-1}$ 不含这些根中的任何一个。因此 $f(x)$ 在 $K_{s-1}$ 上不可约，所以 $[K_s:K_{s-1}]$ 被 $3$ 整除。由于我们已经化归到扩张次数为素数的情形，推出这个扩张次数正好是 $3$，且扩张 $K_s/K_{s-1}$ 是 Galois 的（因为它是 $f(x)$ 在 $K_{s-1}$ 上的分裂域）。又因为 $K_s=K_{s-1}(\sqrt[3]{a})$，其中 $a\in K_{s-1}$，Galois 扩张 $K_s$ 也必须含有 $a$ 的另外两个三次根。这意味着 $K_s$ 含有 $\rho$，一个本原三次单位根。这与 $K_s$ 是 $\mathbb R$ 的子域这一假设矛盾，因而说明不可能只用实根式表示这个三次方程的根。

\paragraph{四次方程的根式解}
现在考虑四次多项式
\[
f(x)=x^4+ax^3+bx^2+cx+d,
\]
它在代换 $x=y-a/4$ 下变为四次多项式
\[
g(y)=y^4+py^2+qy+r,
\]
其中
\[
p=\frac18(-3a^2+8b),
\qquad
q=\frac18(a^3-4ab+8c),
\qquad
r=\frac1{256}(-3a^4+16a^2b-64ac+256d).
\]
令 $g(y)$ 的根为 $\alpha_1,\alpha_2,\alpha_3,\alpha_4$。上一节为这个四次多项式引入的预解三次方程为
\[
h(x)=x^3-2px^2+(p^2-4r)x+q^2,
\]
其根为
\[
\theta_1=(\alpha_1+\alpha_2)(\alpha_3+\alpha_4),
\]
\[
\theta_2=(\alpha_1+\alpha_3)(\alpha_2+\alpha_4),
\qquad
\theta_3=(\alpha_1+\alpha_4)(\alpha_2+\alpha_3).
\]
$f(x)$（或 $g(y)$）的分裂域在预解三次方程 $h(x)$ 的分裂域上的 Galois 群是 Klein 四元群。这样的扩张是双二次扩张，这意味着可以用涉及预解三次方程的根 $\theta_1,\theta_2,\theta_3$ 的表达式的平方根来解出 $\alpha_1,\alpha_2,\alpha_3,\alpha_4$。在这里显然有
\[
(\alpha_1+\alpha_2)(\alpha_3+\alpha_4)=\theta_1,
\qquad
(\alpha_1+\alpha_2)+(\alpha_3+\alpha_4)=0,
\]
这给出
\[
\alpha_1+\alpha_2=\sqrt{-\theta_1},
\qquad
\alpha_3+\alpha_4=-\sqrt{-\theta_1}.
\]
同样地，
\[
\alpha_1+\alpha_3=\sqrt{-\theta_2},
\qquad
\alpha_2+\alpha_4=-\sqrt{-\theta_2},
\]
\[
\alpha_1+\alpha_4=\sqrt{-\theta_3},
\qquad
\alpha_2+\alpha_3=-\sqrt{-\theta_3}.
\]
一个容易的计算表明
\[
\sqrt{-\theta_1}\sqrt{-\theta_2}\sqrt{-\theta_3}=-q,
\]
所以两个平方根的选择决定第三个。由于 $\alpha_1+\alpha_2+\alpha_3+\alpha_4=0$，把上面左侧方程相加得到 $2\alpha_1$，同样可以解出 $g(y)$ 的其他根。我们得到
\[
\begin{aligned}
2\alpha_1&=\sqrt{-\theta_1}+\sqrt{-\theta_2}+\sqrt{-\theta_3},\\
2\alpha_2&=\sqrt{-\theta_1}-\sqrt{-\theta_2}-\sqrt{-\theta_3},\\
2\alpha_3&=-\sqrt{-\theta_1}+\sqrt{-\theta_2}-\sqrt{-\theta_3},\\
2\alpha_4&=-\sqrt{-\theta_1}-\sqrt{-\theta_2}+\sqrt{-\theta_3}.
\end{aligned}
\]
这把四次方程的求解化归为求解相应的预解三次方程。

\subsection*{练习}

\begin{enumerate}[label=\arabic*.]
\item 使用 Cardano 公式求解方程 $x^3+x^2-2=0$。特别地，证明这个方程有实根
\[
\frac13\left(\sqrt[3]{26+15\sqrt3}+\sqrt[3]{26-15\sqrt3}-1\right).
\]
直接证明这个三次方程的根为 $1,-1\pm i$。通过证明
\[
\sqrt[3]{26+15\sqrt3}=2+\sqrt3,
\qquad
\sqrt[3]{26-15\sqrt3}=2-\sqrt3
\]
从而
\[
\sqrt[3]{26+15\sqrt3}+\sqrt[3]{26-15\sqrt3}=4,
\]
解释这一点。

\item 令 $\zeta_7$ 为本原 $7$ 次单位根，并令 $\alpha=\zeta+\zeta^{-1}$。
\begin{enumerate}[label=(\alph*)]
\item 证明 $\zeta_7$ 是 $\mathbb Q(\alpha)$ 上二次多项式 $z^2-\alpha z+1$ 的一个根。
\item 利用 $\zeta_7$ 的极小多项式证明 $\alpha$ 是三次多项式 $x^3+x^2-2x-1$ 的一个根。
\item 利用 (a) 和 (b)，以及正文中 (b) 里的三次方程的显式解，把 $\zeta_7$ 用根式表示，类似于前面给出的其他小阶单位根的显式表达式。（这个表达式的复杂性解释了为什么我们先前没有把 $\zeta_7$ 列入显式单位根的表中。）
\end{enumerate}

\item 设 $F$ 是特征不等于 $2$ 的域。陈述并证明关于 $\alpha,\beta\in F$ 的一个充要条件，使得 $F(\sqrt\alpha)=F(\sqrt\beta)$。用它判断是否有
\[
\mathbb Q(\sqrt{1-\sqrt2})=\mathbb Q(i,\sqrt2).
\]

\item 令 $K=\mathbb Q(\sqrt[n]{a})$，其中 $a\in\mathbb Q$、$a>0$，并假设 $[K:\mathbb Q]=n$（即 $x^n-a$ 不可约）。令 $E$ 为 $K$ 的任意子域，并令 $[E:\mathbb Q]=d$。证明
\[
E=\mathbb Q(\sqrt[d]{a}).
\]
[考虑 $N_{K/E}(\sqrt[n]{a})\in E$。]

\item 令 $K$ 如前一题。证明：若 $n$ 为奇数，则 $K$ 没有非平凡的、在 $\mathbb Q$ 上 Galois 的子域；若 $n$ 为偶数，则 $K$ 中唯一非平凡的、在 $\mathbb Q$ 上 Galois 的子域是 $\mathbb Q(\sqrt a)$。

\item 令 $L$ 为前两题中 $K$ 的 Galois 闭包（即 $x^n-a$ 的分裂域）。证明
\[
[L:\mathbb Q]=n\varphi(n)
\quad\text{或}\quad
\frac12 n\varphi(n).
\]
[注意 $\mathbb Q(\zeta_n)\cap K$ 是 $\mathbb Q$ 的 Galois 扩张。]

\item （循环扩张的 Kummer 生成元）设 $F$ 是特征不整除 $n$ 且含有 $n$ 次单位根的域，并令 $K$ 是次数 $d$ 整除 $n$ 的循环扩张。于是对某个非零 $a\in F$，$K=F(\sqrt[n]{a})$。令 $\sigma$ 为循环群 $\operatorname{Gal}(K/F)$ 的生成元。
\begin{enumerate}[label=(\alph*)]
\item 证明
\[
\sigma(\sqrt[n]{a})=\zeta\sqrt[n]{a}
\]
对某个本原 $d$ 次单位根 $\zeta$ 成立。
\item 假设 $K=F(\sqrt[n]{a})=F(\sqrt[n]{b})$。利用 (a) 证明
\[
\frac{\sigma(\sqrt[n]{a})}{\sqrt[n]{a}}
=\biggl(\frac{\sigma(\sqrt[n]{b})}{\sqrt[n]{b}}\biggr)^i
\]
对某个与 $d$ 互素的整数 $i$ 成立。推出 $\sigma$ 固定元素
\[
\frac{\sqrt[n]{a}}{(\sqrt[n]{b})^i},
\]
所以这个元素属于 $F$。
\item 证明
\[
K=F(\sqrt[n]{a})=F(\sqrt[n]{b})
\]
当且仅当存在 $c,d\in F$ 和整数 $i$，使得
\[
a=b^i c^n,
\qquad
b=a^i d^n,
\]
也就是说，当且仅当 $a$ 与 $b$ 在模 $n$ 次幂的 $F^\times$ 中生成同一个子群。
\end{enumerate}

\item 设 $p,q,r$ 是 $\mathbb Z$ 中的素数，且 $q\neq r$。令 $\sqrt[p]{q}$ 表示 $x^p-q$ 的任意一个根，令 $\sqrt[p]{r}$ 表示 $x^p-r$ 的任意一个根。证明
\[
\mathbb Q(\sqrt[p]{q})\neq \mathbb Q(\sqrt[p]{r}).
\]

\item （Artin--Schreier 扩张）设 $F$ 是特征为 $p$ 的域，且 $K$ 是 $F$ 的次数为 $p$ 的循环扩张。证明 $K=F(\alpha)$，其中 $\alpha$ 是某个多项式
\[
x^p-x-a
\]
的根，$a\in F$。[注意 $\operatorname{Tr}_{K/F}(-1)=0$，因为 $F$ 的特征为 $p$，所以由第 2 节练习 26，有 $-1=\alpha-\sigma\alpha$，其中 $\sigma$ 是 $\operatorname{Gal}(K/F)$ 的一个生成元。证明 $a=\alpha^p-\alpha$ 是 $F$ 的元素。] 注意，因为 $F$ 含有 $p$ 次单位根（即 $1$），这就完成了所有含 $p$ 次单位根的域上素数次数 $p$ 的循环扩张在所有特征下的描述。

\item 令 $K=\mathbb Q(\zeta_p)$ 为素数 $p$ 的 $p$ 次单位根的分圆域，并令 $G=\operatorname{Gal}(K/\mathbb Q)$。令 $\zeta$ 表示任意一个 $p$ 次单位根。证明
\[
\sum_{\sigma\in G}\sigma(\zeta)
\]
（$\zeta$ 从 $K$ 到 $\mathbb Q$ 的迹）等于 $-1$ 或 $p-1$，其情形分别取决于 $\zeta$ 是否为本原 $p$ 次单位根。

\item （经典 Gauss 和）令 $K=\mathbb Q(\zeta_p)$ 为奇素数 $p$ 的 $p$ 次单位根的分圆域，并视作 $\mathbb C$ 的子域，令 $G=\operatorname{Gal}(K/\mathbb Q)$。令 $H$ 表示循环群 $G$ 中指数为 $2$ 的子群。定义
\[
\eta_0=\sum_{\tau\in H}\tau(\zeta_p),
\qquad
\eta_1=\sum_{\tau\in \sigma H}\tau(\zeta_p),
\]
其中 $\sigma$ 是 $\operatorname{Gal}(K/\mathbb Q)$ 的一个生成元（这是 $\zeta_p$ 关于 $H$ 的两个周期，即关于 $G$ 中 $H$ 的两个陪集的共轭和，参见第 5 节）。
\begin{enumerate}[label=(\alph*)]
\item 证明 $\sigma(\eta_0)=\eta_1$，$\sigma(\eta_1)=\eta_0$，并且
\[
\eta_0=\sum_{a=\text{square}}\zeta_p^a,
\qquad
\eta_1=\sum_{b\neq\text{square}}\zeta_p^b,
\]
其中和分别取遍 $(\mathbb Z/p\mathbb Z)^\times$ 中的平方和非平方。[观察 $H$ 是 $(\mathbb Z/p\mathbb Z)^\times$ 中的平方子群。]
\item 证明
\[
\eta_0+\eta_1=(\zeta_p,1)=-1,
\qquad
\eta_0-\eta_1=(\zeta_p,-1),
\]
其中 $(\zeta_p,1)$ 和 $(\zeta_p,-1)$ 是 $\zeta_p$ 的两个 Lagrange 预解式。
\item 令
\[
g=\sum_{i=0}^{p-1}\zeta_p^{i^2}
\]
（经典 Gauss 和）。证明
\[
g=(\zeta_p,-1)=\sum_{i=0}^{p-2}(-1)^i\sigma^i(\zeta_p).
\]
\item 证明若 $\tau\in H$，则 $\tau g=g$；若 $\tau\notin H$，则 $\tau g=-g$。特别地推出 $[\mathbb Q(g):\mathbb Q]=2$。回忆复共轭是 $K$ 上的自同构 $\sigma_{-1}$（参见第 5 节练习 7）。推出：若 $-1$ 是模 $p$ 的平方（即 $p\equiv1\pmod4$），则 $\overline g=g$；若 $-1$ 不是模 $p$ 的平方（即 $p\equiv3\pmod4$），则 $\overline g=-g$，其中 $\overline g$ 表示 $g$ 的复共轭。
\item 证明 $g\overline g=p$。[单位根的复共轭是其倒数。于是
\[
\overline g=\sum_{j=0}^{p-2}(-1)^j(\sigma^j(\zeta_p))^{-1}
\]
给出
\[
\begin{aligned}
g\overline g
&=\sum_{i,j=0}^{p-2}(-1)^i(-1)^j\frac{\sigma^i(\zeta_p)}{\sigma^j(\zeta_p)}
=\sum_{i,j=0}^{p-2}(-1)^{i-j}\sigma^j\left[\frac{\sigma^{i-j}(\zeta_p)}{\zeta_p}\right]\\
&=\sum_{k=0}^{p-2}(-1)^k\sum_{j=0}^{p-2}\sigma^j\left[\frac{\sigma^k(\zeta_p)}{\zeta_p}\right],
\end{aligned}
\]
其中 $k=i-j$。若 $k=0$，元素 $\sigma^k(\zeta_p)/\zeta_p$ 为 $1$；若 $k\neq0$，它是本原 $p$ 次单位根。利用前一题推出内层和在 $k=0$ 时为 $p-1$，否则为 $-1$。]
\item 推出
\[
g^2=(-1)^{(p-1)/2}p
\]
并且
\[
\mathbb Q\left(\sqrt{(-1)^{(p-1)/2}p}\right)
\]
是 $\mathbb Q(\zeta_p)$ 的唯一二次子域。（也参见第 6 节练习 33。）
\end{enumerate}

\item 令 $L$ 为 $\mathbb Q$ 的有限扩张 $\mathbb Q(\alpha)$ 的 Galois 闭包。对任意整除 $\operatorname{Gal}(L/\mathbb Q)$ 的阶的素数 $p$，证明存在 $L$ 的一个子域 $F$，使得 $[L:F]=p$ 且 $L=F(\alpha)$。

\item 设 $F$ 是实数域 $\mathbb R$ 的一个子域。令 $a\in F$，且令 $K=F(\sqrt[n]{a})$，其中 $\sqrt[n]{a}$ 表示 $a$ 的一个实 $n$ 次根。证明：若 $L$ 是 $F$ 的任意 Galois 扩张且包含在 $K$ 中，则 $[L:F]\leq2$。

\item 本题说明在一般情形下，用根式表示实根时必须使用复数，并推广了三次方程的 Casus irreducibilis。令 $f(x)\in\mathbb Q[x]$ 是不可约多项式，且它的所有根都是实数。进一步假设 $f(x)$ 的一个根 $\alpha$ 可用实根式表示（即存在一条实域的根式扩张
\[
\mathbb Q=K_0\subset K_1\subset\cdots\subset K_m\subset\mathbb R,
\]
其中
\[
K_{i+1}=K_i(\sqrt[n_i]{a_i}),
\qquad i=1,2,\ldots,m-1,
\]
对某些整数 $n_i$ 和某些 $a_i\in K_i$ 成立，并且 $\alpha\in K_m$）。证明 $f(x)$ 的 Galois 群是 $2$-群。特别地推出 $f(x)$ 的次数是 $2$ 的幂，并且这样的多项式的实根可完全用实根式表示，当且仅当这些根可由直尺和圆规作出。[论证类似于三次方程的情形。令 $L\subset\mathbb R$ 为 $\mathbb Q(\alpha)$ 的 Galois 闭包，并假设 $\operatorname{Gal}(L/\mathbb Q)$ 的阶被某个奇素数 $p$ 整除。由练习 12，令 $F$ 为 $L$ 的一个子域，使得 $[L:F]=p$ 且 $L=F(\alpha)$，并考虑复合域 $K_i'=FK_i$，$i=0,1,\ldots,m$。这些仍是实根式扩张，并且由正文中 Casus irreducibilis 的论证，可以假设每个 $[K_{i+1}':K_i']$ 都是素数。由于 $\alpha\notin F=FK_0$，存在整数 $s$，使得 $\alpha\notin K_{s-1}'$ 且 $\alpha\in K_s'$。由于扩张都是素数次数，我们有 $K_s'=K_{s-1}'(\alpha)$。因为 $L=F(\alpha)$ 是次数 $p$ 的 Galois 扩张，$K_s'$ 是 $K_{s-1}'$ 的次数 $p$ 的 Galois 扩张，这与前一题矛盾。]

\item （特征 $2$ 中三次方程的“Cardano 公式”）设
\[
f(x)=x^3+px+q
\]
是特征为 $2$ 的域上的不可约三次多项式。令 $\rho$ 为本原三次单位根，令 $\theta,\theta'$ 为二次多项式
\[
x^2+qx+(p^3+q^2)
\]
的根（参见第 6 节练习 50）。令 $\theta_1$ 和 $\theta_2$ 分别为 $p\rho+\theta$ 和 $p\rho+\theta'$ 的三次根，并且选择三次根使得 $\theta_1\theta_2=p$。证明 $f(x)$ 的根为
\[
\alpha=\theta_1+\theta_2,
\qquad
\beta=\rho\alpha+\theta_1,
\qquad
\gamma=\rho\alpha+\theta_2=\alpha+\beta.
\]

\item 令 $a$ 为非零有理数。
\begin{enumerate}[label=(\alph*)]
\item 判断扩张 $\mathbb Q(\sqrt{ai})$（$i^2=-1$）何时在 $\mathbb Q$ 上次数为 $4$。
\item 当 $K=\mathbb Q(\sqrt{ai})$ 在 $\mathbb Q$ 上次数为 $4$ 时，证明 $K$ 在 $\mathbb Q$ 上是 Galois 扩张，且 Galois 群为 Klein 四元群。在这种情形下，确定 $K$ 中所含的 $\mathbb Q$ 的二次扩张。
\end{enumerate}

\item 令 $D\in\mathbb Z$ 为无平方因子的整数，且令 $a\in\mathbb Q$ 为非零有理数。证明
\[
\mathbb Q(\sqrt{a\sqrt D})
\]
不可能是 $\mathbb Q$ 上次数为 $4$ 的循环扩张。

\item 令 $D\in\mathbb Z$ 为无平方因子的整数，且令 $a\in\mathbb Q$ 为非零有理数。证明：若 $\mathbb Q(\sqrt{a\sqrt D})$ 在 $\mathbb Q$ 上是 Galois 扩张，则 $D=-1$。

\item 令 $D\in\mathbb Z$ 为无平方因子的整数，并令 $K=\mathbb Q(\sqrt D)$。
\begin{enumerate}[label=(\alph*)]
\item 证明：若 $D=s^2+t^2$ 是两个有理平方数之和，则存在一个含有 $K$ 的扩张 $L/\mathbb Q$，它在 $\mathbb Q$ 上是 Galois 扩张，并且 Galois 群为阶 $4$ 的循环群。[考虑扩张 $\mathbb Q(\sqrt{D+s\sqrt D})$。]（还要注意：$D$ 是两个有理平方数之和，当且仅当 $D$ 也是两个整数平方数之和，所以不妨假设 $s$ 和 $t$ 是整数。）
\item 反过来证明：若 $K$ 能嵌入如 (a) 中的次数为 $4$ 的循环扩张 $L$ 中，则 $D$ 是两个平方数之和。[一种方法是：(i) 先观察 $L$ 在 $K$ 上是二次的，所以对某些 $a,b\in\mathbb Q$，有 $L=K(\sqrt{a+b\sqrt D})$；(ii) 证明 $L$ 含有二次子域 $\mathbb Q(\sqrt{a^2-b^2D})$，如果 $L/\mathbb Q$ 是循环的，则这个子域必须是 $\mathbb Q(\sqrt D)$，然后使用练习 7。]
\item 特别地推出 $\mathbb Q(\sqrt3)$ 不是任何 $\mathbb Q$ 上次数为 $4$ 的循环扩张的子域。同样推出，当 $D<0$ 为无平方因子的整数时，域 $\mathbb Q(\sqrt D)$ 从不包含在 $\mathbb Q$ 上次数为 $4$ 的循环扩张中（这给出第 6 节练习 19 的另一个证明）。
\end{enumerate}

\item 设 $p$ 为素数。证明 $S_p$ 中任意阶被 $p$ 整除的可解子群，都包含在 $S_p$ 的一个 Sylow $p$-子群的正规化子中（一个阶为 $p(p-1)$ 的 Frobenius 群）。推出：$\mathbb Q[x]$ 中次数为 $p$ 的不可约多项式 $f(x)$ 可由根式求解，当且仅当它的 Galois 群包含在阶为 $p(p-1)$ 的 Frobenius 群中。[令 $G\leq S_p$ 为阶被 $p$ 整除的可解子群。则 $G$ 含有一个 $p$-轮换，因此在 $\{1,2,\ldots,p\}$ 上传递。令 $H<G$ 为 $G$ 中元素 $1$ 的稳定子，所以 $H$ 在 $G$ 中指数为 $p$。证明 $H$ 不含 $G$ 的非平凡正规子群（注意 $H$ 的共轭子群是其他点的稳定子）。令 $G^{(n-1)}$ 为 $G$ 的导出列中最后一个非平凡子群。证明 $H\cap G^{(n-1)}=1$，并推出 $|G^{(n-1)}|=p$，所以 $G$ 的 Sylow $p$-子群（同时也是 $S_p$ 的 Sylow $p$-子群）在 $G$ 中正规。]

\item （五次方程可解性的判据）由前一题，$\mathbb Q[x]$ 中次数为 $5$ 的不可约多项式 $f(x)$ 可由根式求解，当且仅当它的 Galois 群（看作 $S_5$ 的子群）包含在阶为 $20$ 的 Frobenius 群中。已知这等价于一个相关的 $6$ 次多项式 $g(x)$ 有有理根（参见 Dummit, \emph{Solving Solvable Quintics}, Math. Comp., 57 (1991), pp. 387--401）。如果五次多项式具有一般形式（经过平移使 $x^4$ 的系数为零）
\[
f(x)=x^5+px^3+qx^2+rx+s,
\qquad p,q,r,s\in\mathbb Q,
\]
那么相关的 $6$ 次多项式为
\[
\begin{aligned}
g(x)=&\ x^6+8rx^5+(2pq^2-6p^2r+40r^2-50qs)x^4\\
&+(-2q^4+21pq^2r-40p^2r^2+160r^3-15p^2qs-400qrs+125ps^2)x^3\\
&+(p^2q^4-8q^4r+9p^4r^2-136p^2r^3+625q^2s^2+400r^4-6p^3q^2r\\
&\qquad +76pq^2r^2-50pq^3s-1400qr^2s+500prs^2+90p^2qrs)x^2\\
&+(-108p^5s^2+32p^4r^3-256p^2r^4-3125s^4+512r^5-2pq^6+3q^4r^2\\
&\qquad -58q^5s+2750q^2rs^2-31p^3q^3s-500pr^2s^2+19p^2q^4r\\
&\qquad -51p^3q^2r^2+76pq^2r^3-2400qr^3s-325p^2q^2s^2+525p^3rs^2\\
&\qquad +625pqs^3+117p^4qrs+105pq^3rs+260p^2qr^2s)x\\
&+(q^8+256r^6+17q^4r^3-27p^7s^2-4p^6r^3+48p^4r^4-192p^2r^5\\
&\qquad +3125p^2s^4-9375rs^4-1600qr^4s-99p^5rs^2-125pq^4s^2\\
&\qquad -124q^5rs+3250q^2r^2s^2-2000pr^3s^2-13pq^6r+p^5q^2r^2\\
&\qquad +65p^2q^4r^2-128p^3q^2r^3-16pq^2r^4-4p^5q^3s-12p^2q^5s\\
&\qquad -150p^4q^2s^2+1200p^3r^2s^2+18p^6qrs+12p^3q^3rs+196p^4qr^2s\\
&\qquad +590pq^3r^2s-160p^2qr^3s-725p^2q^2rs^2-1250pqrs^3).
\end{aligned}
\]
在特殊情形 $f(x)=x^5+Ax+B$ 中，这个多项式简化为
\[
g(x)=x^6+8Ax^5+40A^2x^4+160A^3x^3+400A^4x^2+(512A^5-3125B^4)x-9375AB^4+256A^6.
\]
\begin{enumerate}[label=(\alph*)]
\item 用这个判据证明多项式 $x^5-5x+12$ 在 $\mathbb Q$ 上的 Galois 群是阶为 $10$ 的二面体群。[证明相关的六次多项式是
\[
x^6-40x^5+1000x^4-20000x^3+250000x^2-66400000x+976000000,
\]
并且有 $x=40$ 作为有理根。也参见第 6 节练习 35。]
\item 用这个判据证明 $x^5-x-1$ 不可由根式求解。
\end{enumerate}
\end{enumerate}

\section{计算 $\mathbb Q$ 上的 Galois 群}

在第 6 节中确定次数 $\leq 4$ 的多项式的 Galois 群，以及在上一节中确定多项式 $x^5-6x+3$ 的 Galois 群时，我们观察到：把自同构看作 $S_n$ 中的元素时，可以从它们的\emph{循环型}得到关于 Galois 群的有用信息。这对于计算 $\mathbb Q$ 上多项式的 Galois 群非常有用；现在我们简要描述其理论根据。

设 $f(x)$ 是有理系数多项式。在确定 $f(x)$ 的 Galois 群时，我们可以假设 $f(x)$ 是可分的，并且具有整数系数。于是 $f(x)$ 的判别式 $D$ 是一个整数，并且非零。

对任意素数 $p$，考虑 $f(x)$ 模 $p$ 的约化
\[
\overline f(x)\in \mathbb F_p[x].
\]
如果 $p$ 整除 $D$，那么约化多项式 $\overline f(x)$ 的判别式在 $\mathbb F_p$ 中满足 $\overline D=0$，所以它不是可分的。

如果 $p$ 不整除 $D$，那么 $\overline f(x)$ 是 $\mathbb F_p$ 上的可分多项式，并且我们可以把 $\overline f(x)$ 分解成互不相同的不可约因子的乘积
\[
\overline f(x)=\overline f_1(x)\overline f_2(x)\cdots\overline f_k(x)
\qquad\text{in }\mathbb F_p[x].
\]
令 $n_i$ 为 $\overline f_i(x)$ 的次数，$i=1,2,\ldots,k$。

这种约化的重要性由代数数论中的如下定理给出；它是研究 $\mathbb Q$ 的有限扩张中算术的一个初等后果（这里我们将它作为已知）。

\noindent\textbf{定理.} 对任意不整除 $f(x)\in\mathbb Z[x]$ 的判别式 $D$ 的素数 $p$，约化
\[
\overline f(x)=f(x)\pmod p
\]
在 $\mathbb F_p$ 上的 Galois 群，作为置换群，同构于 $f(x)$ 在 $\mathbb Q$ 上的 Galois 群的一个子群。

定理中“作为置换群同构”的意思是：不仅 $f(x)$ 模 $p$ 的约化 $\overline f(x)$ 的 Galois 群同构于 $f(x)$ 的 Galois 群的一个子群，而且可以对 $\overline f(x)$ 的根和 $f(x)$ 的根作某种排序（依赖于 $p$），使得在这个同构下，相应自同构作为这些根上的置换的作用是相同的。特别地，$f(x)$ 的 Galois 群中存在与 $\overline f(x)$ 的自同构具有相同循环型的自同构。

$\overline f(x)$ 的 Galois 群是一个\emph{循环}群，因为 $\mathbb F_p$ 的每个有限扩张都是循环扩张。令 $\sigma$ 为这个 $\mathbb F_p$ 上 Galois 群的一个生成元（例如 Frobenius 自同构）。$\overline f_1(x)$ 的根被 Galois 群在它们之间置换，并且给定这些根中的任意两个，都有一个 Galois 自同构把第一个根送到第二个根（回忆：当这种情形发生时，称群在根上是\emph{传递的}）。同样地，Galois 群在每个因子 $\overline f_i(x)$ 的根上都是传递的，$i=1,2,\ldots,k$。由于这些因子两两互素，我们还看到，Galois 群中的任何元素都不会把一个因子的根映到另一个因子的根。

通过给 $\overline f(x)$ 的 $n$ 个根作标号，把 $\sigma$ 看作 $S_n$ 中的元素，并考虑 $\sigma$ 的循环分解。由于 $\sigma$ 在每个因子 $\overline f_i(x)$ 的根之间置换，所以这个循环分解是 $k$ 个互不相同的置换的乘积。由刚才的观察，$\sigma$ 在 $\overline f_1(x)$ 的根上的作用必须是长度为 $n_i$ 的循环；否则 $\sigma$ 的幂不可能在 $\overline f_1(x)$ 的根上传递。同样地，$\sigma$ 在 $\overline f_i(x)$ 的根上的作用给出长度为 $n_i$ 的循环，$i=1,2,\ldots,k$。

由此可见，生成 $\overline f(x)$ 的 Galois 群的自同构 $\sigma$ 的循环分解为
\[
(n_1,n_2,\ldots,n_k),
\]
其中 $n_1,n_2,\ldots,n_k$ 是 $f(x)$ 模 $p$ 约化后的不可约因子的次数。这给出如下结果。

\noindent\textbf{推论 41.} 对任意不整除 $f(x)\in\mathbb Z[x]$ 的判别式的素数 $p$，$f(x)$ 在 $\mathbb Q$ 上的 Galois 群含有一个循环分解为
\[
(n_1,n_2,\ldots,n_k)
\]
的元素，其中 $n_1,n_2,\ldots,n_k$ 是 $f(x)$ 模 $p$ 约化后各不可约因子的次数。

\paragraph{例}
考虑多项式 $x^5-x-1$。这个多项式的判别式为
\[
2869=19\cdot151,
\]
所以我们在不同于 $19,151$ 的素数处约化。模 $2$ 约化时，多项式 $x^5-x-1$ 分解为
\[
(x^2+x+1)(x^3+x^2+1)\pmod 2,
\]
所以 Galois 群有一个 $(2,3)$-循环。把这个元素立方，我们看到 Galois 群含有一个换位。

模 $3$ 约化时，多项式是不可约的，理由如下：$x^5-x-1$ 模 $3$ 没有根，所以如果它在模 $3$ 下可约，那么它就有一个不可约二次因子，从而会与 $x^9-x$ 有公共因子（$x^9-x$ 是 $\mathbb F_3$ 上所有次数 $1$ 和 $2$ 的不可约多项式的乘积），于是会与 $x^4-1$ 或 $x^4+1$ 有公共因子，从而会与 $x^5-x$ 或 $x^5+x$ 有公共因子，从而会与 $-1$ 或 $2x+1$ 有公共因子，而这显然不可能。这同时说明 $x^5-x-1$ 在 $\mathbb Z[x]$ 中不可约，并且在它的 Galois 群中有一个 $5$-循环。

由于 $S_5$ 由任意一个 $5$-循环和任意一个换位生成，所以 $x^5-x-1$ 的 Galois 群是 $S_5$（因此，特别地，这个多项式不能由根式求解，参见第 7 节练习 21）。

上面例子中的论证指出了如何构造以 $S_n$ 为 Galois 群的多项式。我们使用如下事实：$S_n$ 的一个传递子群若含有一个换位和一个 $n-1$-循环，则它就是 $S_n$。令 $f_1$ 是 $\mathbb F_2$ 上次数为 $n$ 的不可约多项式。令 $f_2\in\mathbb F_3[x]$ 是一个次数为 $2$ 的不可约多项式与奇数次数的不可约多项式的乘积（例如，当 $n$ 为偶数时，取一个次数为 $n-3$ 的不可约多项式再乘以 $x$；当 $n$ 为奇数时，取一个次数为 $n-2$ 的不可约多项式）。令 $f_3\in\mathbb F_5[x]$ 是 $x$ 与一个次数为 $n-1$ 的不可约多项式的乘积。最后，令 $f(x)\in\mathbb Z[x]$ 是任意满足
\[
\begin{aligned}
f(x)&\equiv f_1(x)\pmod 2,\\
&\equiv f_2(x)\pmod 3,\\
&\equiv f_3(x)\pmod 5
\end{aligned}
\]
的多项式。

$f(x)$ 模 $2$ 的约化说明 $f(x)$ 在 $\mathbb Z[x]$ 中不可约，因此 Galois 群在 $f(x)$ 的 $n$ 个根上传递。把 $f(x)$ 模 $3$ 的分解给出的元素取适当的奇数次幂，说明 Galois 群含有一个换位。模 $5$ 的分解说明 Galois 群含有一个 $n-1$-循环，因此 Galois 群是 $S_n$。

\noindent\textbf{命题 42.} 对每个 $n\in\mathbb Z^+$，存在无穷多个多项式 $f(x)\in\mathbb Z[x]$，其在 $\mathbb Q$ 上的 Galois 群为 $S_n$。

对于多项式 $f(x)\in\mathbb Z[x]$ 模 $p$ 的分解，有极其有效的算法（参见第 3 节练习 12 到 17），所以以上推论给出了一个确定 Galois 群中元素某些循环型的有效过程。在使用推论 41 时，应当小心不要假设某个特定循环就是 Galois 群中的元素。例如，一个分解可能推出存在一个 $(2,2)$ 循环，譬如 $(1\ 2)(3\ 4)$；另一个分解可能推出存在一个换位。然而不能推出这个换位必然是 $(1\ 2)$（也不能推出它必然是 $(3\ 4)$ 或 $(1\ 3)$，等等）。选择 $(1\ 2)(3\ 4)$ 来表示第一个循环，就固定了根的一种特定排序，而这未必是那个换位表现为 $(1\ 2)$ 时所对应的排序。

推论 41 在判断 Galois 群何时很大（例如 $S_n$）时特别有效，因为一个含有足够多循环型的传递群必须是 $S_n$（例如，一个含有一个换位和一个 $n-1$-循环的 $S_n$ 的传递子群就是 $S_n$，如上所用）。用这种方法最难确定的是小的 Galois 群（例如阶为 $n$ 的循环群），因为一次又一次的分解只会产生阶整除 $n$ 的元素，而人们不能确定是否还有某个尚未出现的素数 $p$ 会给出一个与 Galois 群为循环群这一假设不相容的循环型。如果人们能够“永远计算下去”，至少可以在如下意义下确定 Galois 群元素循环型的精确分布：设 Galois 群 $G\leq S_n$ 的阶为 $N$，并且 $G$ 中有 $n_T$ 个元素具有循环型 $T$（例如 $(2,2)$-循环、换位等），从而循环型 $T$ 在 $G$ 中的“密度”为
\[
d_T=\frac{n_T}{N}.
\]
那么可以在素数集合上定义一种密度（使得谈论“$1/2$ 的素数”等有意义），并且有如下结果（它依赖于代数数论中的 Tchebotarov 密度定理）。

\noindent\textbf{定理.} 使得 $f(x)$ 模 $p$ 分裂为类型 $T$ 的素数 $p$ 的密度恰好为 $d_T$。

这说明，如果我们知道 $f(x)$ 模每个素数的分解，至少就能确定 $G$ 中具有给定循环型的元素数目。不幸的是，即使这样也不足以确定 $G$（直到同构意义下）：已知存在非同构的群，它们对所有循环型都含有相同数目的元素（在 $S_8$ 中有两个非同构的 $96$ 阶群，二者的循环型分布均为：$1$ 个 $1$-循环，$6$ 个 $(2,2)$-循环，$13$ 个 $(2,2,2,2)$-循环，$32$ 个 $(3,3)$-循环，$12$ 个 $(4,4)$-循环，$32$ 个 $(2,6)$-循环）。这样的例子有无穷多个（阶为 $p^3$ 的初等 Abel 群的正则表示与阶为 $p^3$、指数为 $p$ 的非 Abel 群的正则表示，给出 $S_{p^3}$ 中两个非同构的群；对任意素数 $p$，它们的所有非恒等元素都是 $p^2$ 个 $p$-循环的乘积）。

在实际中，人们用 $f(x)$ 模小素数的分解来形成对可能的 Galois 群的判断（基于前面的结果）。然后再试图证明这确实是 Galois 群——这常常是一个困难的问题。对于小次数多项式，存在确定性的算法，部分基于预解多项式的计算。

这些预解多项式类似于前面各节中用来确定四次多项式 Galois 群的三次预解式。这些预解多项式有有理系数，并且它们的根是 $f(x)$ 的根的某些组合（类似于三次预解式中的组合 $(\alpha_1+\alpha_2)(\alpha_3+\alpha_4)$）。然后人们确定这些预解多项式的分解，以获得关于 $f(x)$ 的 Galois 群的信息——例如，线性因子的存在意味着 Galois 群落在 $S_n$ 中所选根的组合的稳定子中（例如，在我们的三次预解式情形中，是阶为 $8$ 的二面体群）。然而应当注意，与四次多项式的三次预解式情形不同，一般构造出来的预解多项式的次数远大于 $f(x)$ 的次数。这种计算技巧的有效性也强烈依赖于对 $S_n$ 的可能传递子群的显式了解。对 $n=2,3,\ldots,8$，$S_n$ 的传递子群的同构类数目分别为
\[
1,2,5,5,16,7,50.
\]
Galois 群的计算有很大研究兴趣，部分动机来自确定哪些群能作为 $\mathbb Q$ 上的 Galois 群这一问题。

我们用一些较容易的例子来说明这些技巧（取自 L. Soicher, \emph{The Computation of Galois Groups}, Master's Thesis, Concordia University, Montreal, 1981）。

\paragraph{例}
\begin{enumerate}
\item $S_5$ 的传递子群共有 $5$ 个同构类，分别由群 $Z_5$、$D_{10}$、$F_{20}$、所谓阶为 $20$ 的 Frobenius 群（$x^5-2$ 的 Galois 群，在 $S_5$ 中由 $(1\ 2\ 3\ 4\ 5)$ 和 $(2\ 3\ 5\ 4)$ 生成）、$A_5$ 和 $S_5$ 给出。这些群的循环型分布如下：
\[
\begin{array}{c|rrrrrrr}
\text{循环型:} & 1 & 2 & (2,2) & 3 & (2,3) & 4 & 5\\ \hline
Z_5    & 1 &   &    &    &    &    & 4\\
D_{10} & 1 &   & 5  &    &    &    & 4\\
F_{20} & 1 &   & 5  &    &    & 10 & 4\\
A_5    & 1 &   & 15 & 20 &    &    & 24\\
S_5    & 1 &10 & 15 & 20 & 20 & 30 & 24
\end{array}
\]

由这些信息，$x^5-x-1$ 的不可约性（给出在 $5$ 个根上的传递性）以及循环型 $(2,3)$ 立刻说明 $x^5-x-1$ 的 Galois 群是 $S_5$。

现在考虑多项式 $x^5+15x+12$。判别式为 $2^{10}3^45^5$，所以 Galois 群不包含在 $A_5$ 中。只有两种可能：$S_5$ 或 $F_{20}$。通过把这个多项式模若干小素数分解，并把所得循环型分布与上表比较，可以容易地判断哪一种更可能。然而这不能证明所猜测的 Galois 群实际上正确。为了判定 $S_5$ 和 $F_{20}$ 中哪一个是正确的，可以计算次数为 $15$ 的预解多项式 $R(x)$，其根是 $S_5$ 作用下
\[
(\alpha_1+\alpha_2-\alpha_3-\alpha_4)^2
\]
的互异置换值，其中 $\alpha_1,\alpha_2,\alpha_3,\alpha_4$ 是 $f(x)$ 的 $4$ 个根。按定义，$S_5$ 在 $R(x)$ 的根上传递；而用上面给出的 $F_{20}$ 的显式生成元不难检验，$F_{20}$ 在这 $15$ 个值上不是传递的。于是 $R(x)$ 在 $\mathbb Q$ 上可约，当且仅当这个五次多项式的 Galois 群是 $F_{20}$。对 $x^5+15x+12$，人们发现预解多项式 $R(x)$ 分解为一个次数为 $5$ 的多项式和一个次数为 $10$ 的多项式，因此这个五次多项式的 Galois 群是 $F_{20}$。也可以使用上一节练习 21（参见练习 6），它也是基于相关预解多项式的计算。

\item 考虑多项式
\[
x^7-14x^5+56x^3-56x+22.
\]
判别式计算为 $2^67^{10}$，所以 Galois 群包含在 $A_7$ 中。

把该多项式在 $3$ 与 $193$ 之间且不等于 $7$ 的 $42$ 个素数处作分解，得到循环型分布：$1$ 个 $1$-循环（$2.38\%$），$30$ 个 $(3,3)$-循环（$71.43\%$），$11$ 个 $7$-循环（$26.19\%$）。$S_7$ 的传递子群有 $7$ 个同构类，其中 $4$ 个包含在 $A_7$ 中。在这些群中，有一个不含 $(3,3)$-循环，于是剩下三种可能：$A_7$、$GL_3(\mathbb F_2)$ 或 $F_{21}$，其中 $F_{21}$ 是阶为 $21$ 的 Frobenius 群（在 $S_7$ 中由 $(1\ 2\ 3\ 4\ 5\ 6\ 7)$ 和 $(2\ 3\ 5)(4\ 7\ 6)$ 生成）。这三个群的循环型分布如下：
\[
\begin{array}{c|rrrrrrrrr}
\text{循环型:} & 1 & 2 & (2,2) & 3 & (2,2,3) & (3,3) & (2,4) & 5 & 7\\ \hline
F_{21}           & 1 &   &     &    &     & 14  &     &     & 6\\
GL_3(\mathbb F_2)& 1 &   & 21  &    &     & 56  & 42  &     & 48\\
A_7              & 1 &21 &105  & 70 & 210 & 280 & 630 & 504 & 720
\end{array}
\]

由此可见，这个多项式的 Galois 群极有可能是阶为 $21$ 的 Frobenius 群。事实确实如此（验证需要计算一个次数为 $35$ 的预解式，并在 $\mathbb Z$ 上分解它——共有三个因子，次数分别为 $7,7$ 和 $21$）。
\end{enumerate}

\subsection*{练习}

\begin{enumerate}
\item 设 $p$ 是素数。利用推论 41 以及这个多项式的 Galois 群是 Klein 四元群这一事实，证明多项式 $x^4+1$ 模 $p$ 要么分裂为两个不可约二次因子，要么分裂为 $4$ 个线性因子。

\item （参见第 6 节练习 48。）
\begin{enumerate}
\item 令 $K$ 为 $x^6-2x^3-2$ 的分裂域。证明如果 $[K:\mathbb Q]=12$，那么
\[
K=\mathbb Q(\sqrt[3]{2},i,\sqrt3),
\]
并且 $K$ 在双二次域 $F=\mathbb Q(i,\sqrt3)$ 上由
\[
\alpha=\sqrt[3]{1+\sqrt3}
\qquad\text{和}\qquad
\beta=\sqrt[3]{1-\sqrt3}
\]
生成。证明如果是这种情形，那么 $\operatorname{Gal}(K/\mathbb Q)$ 中阶为 $3$ 的元素都属于 $\operatorname{Gal}(K/F)$。推出 $\operatorname{Gal}(K/\mathbb Q)$ 中任意阶为 $3$ 的元素都把 $\alpha$ 映到 $1+\sqrt3$ 的另一个立方根，把 $\beta$ 映到 $1-\sqrt3$ 的另一个立方根；并且如果它在 $\alpha$ 或 $\beta$ 上为恒等，则它在整个 $K$ 上为恒等。

\item 证明 $f(x)$ 在 $\mathbb F_{13}$ 上分解为不可约因子的分解是多项式
\[
(x-7)(x-8)(x-11)(x^3+3),
\]
并使用推论 41 证明 $[K:\mathbb Q]=36$。

\item 已知 $G=\operatorname{Gal}(K/\mathbb Q)$ 的阶为 $36$，显式确定 $G$ 的所有元素，并且特别证明 $G$ 同构于 $S_3\times S_3$。
\end{enumerate}

\item 证明 $x^5+20x+16$ 的 Galois 群是 $A_5$。

\item 证明 $x^5+x^4-4x^3-3x^2+3x+1$ 的 Galois 群是阶为 $5$ 的循环群。[证明这是 $\zeta_{11}+\zeta_{11}^{-1}$ 的极小多项式。]

\item 证明 $x^5+11x+44$ 的 Galois 群是二面体群 $D_{10}$（参见第 7 节练习 21）。

\item 证明 $x^5+15x+12$ 的 Galois 群是 $F_{20}$，即阶为 $20$ 的 Frobenius 群（参见第 7 节练习 21）。

\item 证明 $x^6+24x-20$ 的 Galois 群是 $A_6$。

\item 证明 $x^7+7x^4+14x+3$ 的 Galois 群是 $A_7$。

\item 确定一个次数为 $7$、Galois 群为阶 $7$ 的循环群的多项式。

\item 确定 $x^7-7x+3$ 的可能 Galois 群。
\end{enumerate}

\section{超越扩张、不可分扩张、无限 Galois 群}

本节汇集关于任意扩张 $E/F$ 的一些结果。这些结果补充前面各节的结果，并完成任意（可能无限）扩张如何分解的基本图景。由于本节主要意在作为一个概览，所有证明均不列出；凡是这些证明可以由读者容易补出的地方，我们或者在正文中，或者（带提示地）在练习中指出。

本节始终设 $E/F$ 是域扩张。回忆：$E$ 中不在 $F$ 上代数的元素称为在 $F$ 上超越。要记住，涉及超越元的扩张总是无限次的。当我们考虑对 $t$ 赋值时，通常保留说法“$t$ 是 $F$ 上的一个不定元”。然而，从域论观点看，超越元与不定元这两个术语是同义的（所以 $\mathbb R$ 的子域 $\mathbb Q(\pi)$ 与域 $\mathbb Q(t)$ 同构）。

\noindent\textbf{定义.}
\begin{enumerate}[label=(\arabic*)]
\item $E$ 的一个子集 $\{a_1,a_2,\ldots,a_n\}$ 称为在 $F$ 上\emph{代数独立}，如果不存在非零多项式
\[
 f(x_1,x_2,\ldots,x_n)\in F[x_1,x_2,\ldots,x_n]
\]
使得 $f(a_1,a_2,\ldots,a_n)=0$。$E$ 的任意子集 $S$ 称为在 $F$ 上代数独立，如果 $S$ 的每个有限子集都是代数独立的。$S$ 的元素称为 $F$ 上的相互独立超越元。
\item $E/F$ 的一个\emph{超越基}是 $E$ 中一个极大（相对于包含关系）的在 $F$ 上代数独立的子集。
\end{enumerate}

注意，如果 $E/F$ 是代数的，那么空集是 $E$ 的唯一代数独立子集。特别地，代数独立集的元素必定是超越元。此外，容易验证：$S\subseteq E$ 是 $F$ 上代数独立集，当且仅当每个 $s\in S$ 都在 $F(S-\{s\})$ 上超越。还容易看出，$S$ 是 $E/F$ 的超越基，当且仅当 $S$ 是 $F$ 上的一个代数独立超越元集合，并且 $E$ 在 $F(S)$ 上代数。

\noindent\textbf{定理.} 扩张 $E/F$ 有超越基，并且 $E/F$ 的任意两个超越基有相同基数。

\begin{proof}
第一断言是一个标准的 Zorn 引理论证。第二断言的证明使用同证明向量空间任意两个基有相同基数时相同的“替换引理”思想。
\end{proof}

\noindent\textbf{定义.} $E/F$ 的一个超越基的基数称为 $E/F$ 的\emph{超越次数}。

代数扩张恰好是超越次数为 $0$ 的扩张。

这个定理的一个特殊情形是 $E$ 在 $F$ 上有限生成，即
\[
E=F(a_1,a_2,\ldots,a_n)
\]
对 $E$ 中某些（不一定代数独立的）元素 $a_1,\ldots,a_n$ 成立。显然可以给 $a_1,\ldots,a_n$ 重新编号，使得 $a_1,\ldots,a_m$ 是相互独立超越元，而 $a_{m+1},\ldots,a_n$ 在 $F(a_1,\ldots,a_m)$ 上代数（于是 $E$ 是后一个域的有限扩张）。在这种情形下，$E$ 称为 $F$ 上 $m$ 个变量的\emph{函数域}。这类域在代数几何中作为 $m$ 维曲面上的函数域而起基本作用。例如，当 $F=\mathbb C$ 且 $m=1$ 时，这些域在分析中作为紧 Riemann 曲面上的亚纯函数域出现。

注意，如果 $S_1$ 和 $S_2$ 是 $E/F$ 的超越基，那么未必有 $F(S_1)=F(S_2)$。例如，如果 $t$ 在 $\mathbb Q$ 上超越，那么 $\{t\}$ 和 $\{t^2\}$ 都是 $\mathbb Q(t)/\mathbb Q$ 的超越基，但（正如我们马上将看到的）$\mathbb Q(t^2)$ 是 $\mathbb Q(t)$ 的真子域。

现在我们看到，如果 $x_1,x_2,\ldots,x_n$ 是 $F$ 上的不定元，并且
\begin{equation}
 f(x)=(x-x_1)(x-x_2)\cdots (x-x_n) \tag{14.28}
\end{equation}
是次数 $n$ 的一般多项式，那么 $x_i$ 中的 $n$ 个基本对称函数 $s_1,s_2,\ldots,s_n$ 也在 $F$ 上相互独立超越。这是因为 $x_1,\ldots,x_n$ 是 $E=F(x_1,\ldots,x_n)$ 在 $F$ 上的一个超越基（所以超越次数为 $n$），并且 $E$ 在 $F(s_1,\ldots,s_n)$ 上代数（次数为 $n!$）。这个定理迫使 $s_1,\ldots,s_n$ 也是这个扩张的一个超越基（特别地，它们是相互独立超越元）。因此，$F$ 上次数 $n$ 的一般多项式也可以等价地定义为：取 $a_1,\ldots,a_n$ 为任意相互独立超越元（或不定元），并令
\begin{equation}
 f(x)=x^n+a_1x^{n-1}+\cdots+a_n, \tag{14.29}
\end{equation}
其中 $f$ 的根记为 $x_1,\ldots,x_n$（并且 $s_i=(-1)^ia_i$）。

\noindent\textbf{定义.} 如果扩张 $E/F$ 有一个超越基 $S$ 使得 $E=F(S)$，则称 $E/F$ 为\emph{纯超越扩张}。

在前面的讨论中，$F(x_1,\ldots,x_n)$ 与 $F(s_1,\ldots,s_n)$ 都是 $F$ 上的纯超越扩张。作为一个练习（如下），可以证明 $\mathbb Q(t,\sqrt{t^3-t})$ 不是 $\mathbb Q$ 的纯超越扩张，尽管除了 $\mathbb Q$ 中的元素以外，它不含任何在 $\mathbb Q$ 上代数的元素（也就是说，把一般扩张分解为一个纯超越扩张后接一个代数扩张的过程，一般不能反过来，使得代数部分先出现）。

如果 $E$ 是 $F$ 的一个纯超越扩张，其超越次数 $n=1$ 或 $2$，并且 $L$ 是中间域 $F\subseteq L\subseteq E$，且有相同的超越次数，那么 $L$ 仍然是 $F$ 的一个纯超越扩张（Lüroth（$n=1$），Castelnuovo（$n=2$））。然而，如果超越次数 $\geq 3$，这个结果不成立，尽管构造 $L$ 不是纯超越的例子很困难。对于超越次数 $1$ 的扩张，中间域由下面的定理描述。

\noindent\textbf{定理.} 设 $t$ 在 $F$ 上超越。
\begin{enumerate}[label=(\arabic*)]
\item （Lüroth）如果 $F\subseteq K\subseteq F(t)$，那么 $K=F(r)$，其中 $r\in F(t)$。特别地，$F(t)$ 中包含 $F$ 的每个非平凡扩张都是 $F$ 上的纯超越扩张。
\item 如果 $P=P(t)$，$Q=Q(t)$ 是 $F[t]$ 中非零互素多项式，且不全为常数，则
\[
 [F(t):F(P/Q)]=\max(\deg P,\deg Q).
\]
\end{enumerate}

\begin{proof}
(2) 的证明在第 13.2 节练习 18 中概述。
\end{proof}

由这个定理的 (2) 可知，$F(P/Q)=F(t)$ 当且仅当 $P,Q$ 是次数 $\leq 1$ 的非零互素多项式（且不全为常数）。因此
\[
 F(r)=F(t)
\]
当且仅当
\[
 r=\frac{at+b}{ct+d},
\]
其中 $a,b,c,d\in F$ 且 $ad-bc\neq 0$（称为 $t$ 的\emph{分式线性变换}）。对任意 $r\in F(t)-F$，映射 $t\mapsto r$ 都扩张为 $F(t)$ 到自身中的一个嵌入，并在 $F$ 上为恒等。这个嵌入是满的（即，是 $F(t)$ 的一个自同构），恰好对于分式线性变换成立。此外，映射
\[
\operatorname{GL}_2(F)\longrightarrow \operatorname{Aut}(F(t)/F),
\qquad
A=\begin{pmatrix}a&c\\ b&d\end{pmatrix}\longmapsto \sigma_A,
\]
其中 $\sigma_A$ 表示 $F(t)$ 的自同构，它由把 $t$ 映到 $(at+b)/(ct+d)$ 来定义，是一个满同态，其核由标量矩阵组成。因此
\[
 \operatorname{Aut}(F(t)/F)\cong \operatorname{PGL}_2(F),
\]
其中
\[
 \operatorname{PGL}_2(F)=\operatorname{GL}_2(F)/\{\lambda I\mid \lambda\in F^\times\}
\]
给出了这个超越扩张的自同构群（参见第 1 节练习 8）。

当 $\mathbb F$ 是阶为 $q$ 的有限域时，$\operatorname{Aut}(\mathbb F(t)/\mathbb F)\cong \operatorname{PGL}_2(\mathbb F)$ 是一个阶为 $q(q-1)(q+1)$ 的有限群。由推论 11，如果 $K$ 是 $\operatorname{Aut}(\mathbb F(t)/\mathbb F)$ 的固定域，那么 $\mathbb F(t)$ 是 $K$ 上的 Galois 扩张，其 Galois 群等于 $\operatorname{Aut}(\mathbb F(t)/\mathbb F)$。特别地，在这种情形下，$\operatorname{Aut}(\mathbb F(t)/\mathbb F)$ 的固定域不是 $\mathbb F$。

这也提供了 Galois 对应的更多例子；对较小的 $q$ 值可以完全写出。例如，如果 $q=|\mathbb F|=2$，则 $\operatorname{PGL}_2(\mathbb F)$ 是阶为 $6$ 的非 Abel 群，因而同构于 $S_3$，并且有如下子群格：
\[
\begin{tikzpicture}[scale=1.0, baseline=(current bounding box.center), every node/.style={inner sep=2pt}]
\node (one) at (0,2.0) {$1$};
\node (a) at (-2.7,0.6) {$\left\langle\begin{pmatrix}0&1\\1&1\end{pmatrix}\right\rangle$};
\node (b) at (-0.8,0.6) {$\left\langle\begin{pmatrix}0&1\\1&0\end{pmatrix}\right\rangle$};
\node (c) at (1.2,0.6) {$\left\langle\begin{pmatrix}1&0\\1&1\end{pmatrix}\right\rangle$};
\node (d) at (3.1,0.6) {$\left\langle\begin{pmatrix}1&1\\0&1\end{pmatrix}\right\rangle$};
\node (g) at (-0.9,-1.1) {$\operatorname{PGL}_2(\mathbb F_2)$};
\draw (one)--(a); \draw (one)--(b); \draw (one)--(c); \draw (one)--(d);
\draw (a)--(g); \draw (b)--(g); \draw (c)--(g); \draw (d)--(g);
\node at (4.2,-0.8) {图 5};
\end{tikzpicture}
\]

域 $\mathbb F_2(t)$ 在 $\operatorname{Aut}(\mathbb F_2(t)/\mathbb F_2)$ 的固定域 $K$ 上是次数 $6$ 的，并且子域格 $K\subseteq L\subseteq \mathbb F_2(t)$ 与 $S_3$ 的子群格对偶。一个循环子群 $(\sigma)$ 的固定域容易（通过前面的定理）由寻找 $t$ 中一个被 $\sigma$ 固定且满足 $[\mathbb F(t):\mathbb F(r)]=|\sigma|$ 的有理函数 $r$ 得到。例如，如果
\[
\sigma:t\longmapsto \frac{1}{1+t},
\]
则 $\sigma$ 有阶 $3$。有理函数
\[
 r=t+\sigma(t)+\sigma^2(t)=\frac{t^3+t+1}{t(t+1)}
\]
被 $\sigma$ 固定，并且 $[\mathbb F(t):\mathbb F(r)]=3$（由定理的 (2)）。由于 $\mathbb F(r)$ 包含在 $(\sigma)$ 的固定域中，并且 $\mathbb F(t)$ 在该固定域上的次数为 $3$，所以 $\mathbb F(r)$ 就是 $(\sigma)$ 的固定域。用这种方式可以明确描述图 6 中所示的所有含 $K$ 的 $\mathbb F(t)$ 的子域格。
\[
\begin{tikzpicture}[scale=1.0, baseline=(current bounding box.center), every node/.style={inner sep=2pt}]
\node (top) at (0,3) {$\mathbb F_2(t)$};
\node (l) at (-3.4,1.2) {$\mathbb F_2\!\left(\dfrac{t^3+t+1}{t^2+t}\right)$};
\node (m1) at (-0.8,1.5) {$\mathbb F_2\!\left(\dfrac{t^2+1}{t}\right)$};
\node (m2) at (1.2,1.5) {$\mathbb F_2(t^2+t)$};
\node (m3) at (3.4,1.5) {$\mathbb F_2\!\left(\dfrac{t^2}{t+1}\right)$};
\node (bot) at (-0.4,-0.6) {$\mathbb F_2\!\left(\dfrac{(t^3+t+1)(t^3+t^2+1)}{(t^2+t)^2}\right)$};
\draw (top)--(l); \draw (top)--(m1); \draw (top)--(m2); \draw (top)--(m3);
\draw (l)--(bot); \draw (m1)--(bot); \draw (m2)--(bot); \draw (m3)--(bot);
\node at (4.5,-0.1) {图 6};
\end{tikzpicture}
\]

$\mathbb Q$ 的纯超越扩张在把有限群实现为 $\mathbb Q$ 上的 Galois 群的问题中起重要作用。我们描述 Hilbert 的一个深刻结果，它是这一研究领域的基本结果。若 $a_1,a_2,\ldots,a_n$ 是域 $F$ 上的相互独立不定元，则可以把 $a_1,\ldots,a_n$ 在 $F$ 的任意元素上求值（或特化），即把 $F$ 中的值代入“变量” $a_1,a_2,\ldots,a_n$。如果 $E$ 是 $F(a_1,\ldots,a_n)$ 的一个 Galois 扩张，那么 $E$ 作为一个多项式的分裂域得到，而这个多项式的系数在 $F[a_1,\ldots,a_n]$ 中。把 $a_1,\ldots,a_n$ 特化到 $F$ 中，就把这个多项式映到一个系数在 $F$ 中的多项式。$E$ 的特化就是所得特化多项式的分裂域。

\noindent\textbf{定理（Hilbert）.} 设 $x_1,x_2,\ldots,x_n$ 是 $\mathbb Q$ 上的相互独立超越元，设 $E=\mathbb Q(x_1,\ldots,x_n)$，并设 $G$ 是 $E$ 的一个有限自同构群，其固定域为 $K$。如果 $K$ 是 $\mathbb Q$ 的一个纯超越扩张，且有超越基 $a_1,a_2,\ldots,a_n$，那么存在无穷多个把 $a_1,\ldots,a_n$ 特化到 $\mathbb Q$ 中的方式，使得 $E$ 特化为 $\mathbb Q$ 的一个 Galois 扩张，且 Galois 群同构于 $G$。

Hilbert 定理给出一个充分条件，保证特化扩张不会坍缩。一般而言，特化扩张的 Galois 群是 $G$ 的一个子群（参见命题 19），并且可能是 $G$ 的真子群。还已知固定域 $K$ 不一定总是 $\mathbb Q$ 的纯超越扩张。一个这样的例子出现在 $G$ 为阶 $47$ 的循环群时。

这个定理可以用来给出命题 42 的另一个证明：

\noindent\textbf{推论.} 对所有 $n$，$S_n$ 都是 $\mathbb Q$ 上的一个 Galois 群。

\begin{proof}
我们已经证明，$S_n$ 以明显方式作用在 $\mathbb Q(x_1,\ldots,x_n)$ 上时，它的固定域是 $\mathbb Q$ 的纯超越扩张（以基本对称函数为超越基），所以 Hilbert 定理立即推出此推论。
\end{proof}

$K$ 是 $\mathbb Q$ 上纯超越扩张这个假设对 Hilbert 定理的证明是关键的。每个有限群都同构于某个 $S_n$ 的子群，因此对某个 $n$ 作用在 $\mathbb Q(x_1,\ldots,x_n)$ 上。然而，即使对于 $S_n$ 的子群 $A_n$，在这个明显作用下它的固定域是否为 $\mathbb Q$ 的纯超越扩张也尚不知道（尽管由其他方法可知，对所有 $n$，$A_n$ 都是 $\mathbb Q$ 上的 Galois 群）。因此这一研究领域中有许多重要的开放问题。

还应注意，当基域 $\mathbb Q$ 被任意域 $F$ 取代时，Hilbert 定理并不成立（例如设 $F$ 是代数闭域）。特别地，如前所述，第 6 节中的一般多项式 $f(x)$ 对任意 $F$ 在 $F(a_1,\ldots,a_n)$ 上有 Galois 群 $S_n$，但当 $F$ 是有限域时，从它的分裂域得到的特化扩张总是循环的。

下面我们扩展第 13.5 节中描述的不可分扩张理论。设 $p$ 为素数，$F$ 是特征 $p$ 的域。

\noindent\textbf{定义.} 代数扩张 $E/F$ 称为\emph{纯不可分}，如果对每个 $\alpha\in E$，$\alpha$ 在 $F$ 上的极小多项式只有一个不同的根。

容易看出下面条件等价：
\begin{enumerate}[label=(\arabic*)]
\item $E/F$ 是纯不可分的；
\item 如果 $\alpha\in E$ 在 $F$ 上可分，则 $\alpha\in F$；
\item 如果 $\alpha\in E$，则对某个 $n$（依赖于 $\alpha$）有 $\alpha^{p^n}\in F$，并且
\[
 m_{\alpha,F}(x)=x^{p^n}-\alpha^{p^n}.
\]
\end{enumerate}

下面这个容易的命题描述可分扩张与纯不可分扩张的复合。

\noindent\textbf{命题.} 如果 $E_1$ 和 $E_2$ 是 $E$ 的子域，并且二者都是 $F$ 的可分（或者二者都是纯不可分）扩张，那么它们的复合 $E_1E_2$ 是 $F$ 的可分（相应地，纯不可分）扩张。

\begin{proof}
练习。
\end{proof}

由此可立即得到如下结果。

\noindent\textbf{命题.} 设 $E/F$ 是代数扩张。则存在唯一的域 $E_{\mathrm{sep}}$，满足
\[
F\subseteq E_{\mathrm{sep}}\subseteq E,
\]
并且 $E_{\mathrm{sep}}$ 在 $F$ 上可分，而 $E$ 在 $E_{\mathrm{sep}}$ 上纯不可分。域 $E_{\mathrm{sep}}$ 是 $E$ 中在 $F$ 上可分的元素组成的集合。

$E_{\mathrm{sep}}/F$ 的次数称为 $E/F$ 的\emph{可分次数}，$E/E_{\mathrm{sep}}$ 的次数称为 $E/F$ 的\emph{不可分次数}（通常分别记为 $[E:F]_s$ 和 $[E:F]_i$）。这两个次数的乘积就是（通常的）次数。上述命题立即给出如下推论。

\noindent\textbf{推论.} 可分次数（相应地，不可分次数）具有乘法性。

当 $E$ 由 $F$ 上不可约多项式 $p(x)\in F[x]$ 的一个根生成时，扩张 $E/F$ 的可分次数和不可分次数，与第 13.5 节中定义的多项式 $p(x)$ 的可分次数和不可分次数相同。

这个命题断言，任意代数扩张都可以分解为一个可分扩张后接一个纯不可分扩张。本节末练习 3 概述了一个例子，说明这种分解一般不能反过来，即存在一个扩张，它不是某个纯不可分扩张上的可分扩张。我们将马上叙述扩张满足什么条件时，可以把它分解为纯不可分子扩张后接可分子扩张。

现在我们知道，任意扩张 $E/F$ 都可以分解为 $F$ 的一个纯超越扩张 $F(S)$，后接 $F(S)$ 的一个可分扩张 $E_1$，再后接一个纯不可分扩张 $E/E_1$。在某些情形下，位于“顶部”的代数扩张中的不可分性可以通过明智选择超越基而去除：

\noindent\textbf{命题.} 如果 $E$ 是完美域 $F$ 的一个有限生成扩张，那么存在 $E/F$ 的一个超越基 $T$，使得 $E$ 是 $F(T)$ 的可分（代数）扩张。

命题中所描述的超越基 $T$ 称为\emph{分离超越基}。本节末练习 4 用一个非平凡例子说明这一点。

回忆：如果扩张 $E/F$ 是 $F[x]$ 中某个（可能无限的）多项式集合的分裂域，则称 $E/F$ 是\emph{正规}的（特别地，正规扩张是代数的，但不一定有限或可分）。我们先前使用了同义词“分裂域”，这里在任意代数扩张的语境下重新引入“正规”这一术语，因为它在文献中经常使用，尤其是在把一个域嵌入到代数闭包的语境中。虽然下面一组等价条件可以从前面各节中提取出来，但读者应写出完整证明，并检查这些论证对无限扩张和不可分扩张都成立：

\noindent\textbf{命题.} 设 $E/F$ 是任意代数扩张，设 $\Omega$ 是 $E$ 的一个代数闭包。以下条件等价：
\begin{enumerate}[label=(\arabic*)]
\item $E/F$ 是正规扩张（即，它是 $F[x]$ 中某个多项式集合在 $F$ 上的分裂域）；
\item 只要 $\sigma:E\to \Omega$ 是一个嵌入且 $\sigma|_F$ 是恒等映射，就有 $\sigma(E)=E$；
\item 只要不可约多项式 $f(x)\in F[x]$ 在 $E$ 中有一个根，它就在 $E$ 中有所有根。
\end{enumerate}

一般地，正规扩张 $E/F$ 到 $E$ 的一个代数闭包中、且扩张 $F$ 的恒等嵌入的任意嵌入，都是 $E$ 的一个自同构，即为 $\operatorname{Aut}(E/F)$ 的一个元素。此外，只要可分次数有限，这样的自同构数目就等于 $E/F$ 的可分次数：
\[
\text{若 }E/F\text{ 正规且 }[E:F]_s\text{ 有限，则}\qquad
|\operatorname{Aut}(E/F)|=[E:F]_s.
\]
如果 $[E:F]_s$ 是无限的，我们马上会看到 $|\operatorname{Aut}(E/F)|$ 也无限，但二者未必有相同的基数。

如果 $E/F$ 是可分次数有限的正规扩张，令 $E_0$ 为 $\operatorname{Aut}(E/F)$ 的固定域。由推论 11，$E/E_0$ 是一个（可分）Galois 扩张，其次数等于 $|\operatorname{Aut}(E/F)|$。于是 $E_0/F$ 必须是纯不可分的（次数等于 $[E:F]_i$），也就是说，对正规扩张而言，扩张中的可分部分和纯不可分部分可以颠倒。更精确地，我们容易得到如下命题。

\noindent\textbf{命题.} 如果 $E/F$ 正规且 $[E:F]_s<\infty$，那么
\[
E=E_{\mathrm{sep}}E_{\mathrm{pi}},
\]
其中 $E_{\mathrm{pi}}$ 是 $F$ 的一个纯不可分扩张（$E_{\mathrm{pi}}$ 由 $E$ 中所有在 $F$ 上纯不可分的元素组成），并且
\[
E_{\mathrm{sep}}\cap E_{\mathrm{pi}}=F.
\]

最后，我们说明 Galois 理论如何推广到无限扩张。

\noindent\textbf{定义.} 扩张 $E/F$ 称为\emph{Galois 扩张}，如果它是代数、正规且可分的。在这种情形下，$\operatorname{Aut}(E/F)$ 称为这个扩张的 Galois 群，记为 $\operatorname{Gal}(E/F)$。

对无限扩张而言，Galois 群的所有子群组成的集合与 $E$ 中包含 $F$ 的所有子域组成的集合之间未必存在双射，下面的例子说明了这一点。

令 $E$ 为 $\mathbb R$ 的子域，它由给 $\mathbb Q$ 添入所有正有理数的平方根而得到。容易看出，$E$ 也可以描述为多项式集合 $x^2-p$ 的分裂域，其中 $p$ 遍历 $\mathbb Z^+$ 中所有素数。注意 $E$ 是 $\mathbb Q$ 的一个（可数）无限 Galois 扩张。由于 $E$ 的每个自同构 $\sigma$ 都由它在素数平方根上的作用确定，并且 $\sigma$ 或固定或取负每一个这样的平方根，$\sigma^2$ 是恒等自同构。因此 $\operatorname{Aut}(E)$ 是一个无限初等 Abel $2$-群。于是 $\operatorname{Aut}(E)$ 是 $\mathbb F_2$ 上的无限维向量空间。由对偶空间一节（第 11.3 节）的一个练习，$\operatorname{Aut}(E)$ 到 $\mathbb F_2$ 的非零同态的数目是不可数的，从而它们的核（余维数为 $1$ 的子空间）不可数且彼此不同。因此 $\operatorname{Aut}(E)$ 有不可数多个指数为 $2$ 的子群，而 $\mathbb Q$ 只有可数多个二次扩张。

基本问题在于，$\operatorname{Gal}(E/F)$ 的许多（多数）子群并不（以双射方式）对应于 $E$ 中包含 $F$ 的子域。为了挑出 $\operatorname{Gal}(E/F)$ 的“正确”子群集合，我们必须在这个群上引入一个拓扑（称为 Krull 拓扑）。拓扑空间中（拓扑）闭集族的公理，正是挑出相关子群的记账装置（这些公理列在第 15.2 节）。有限扩张的 Galois 理论迫使某些有限指数子群成为闭集，而这些闭集又决定整个群上的拓扑（正如我们预期的，因为 $E$ 中 $F$ 的每个扩张都是有限扩张的复合）。此外，扩张 $E/F$ 的 Galois 群是有限群 $\operatorname{Gal}(K/F)$ 的集合的逆极限，其中 $K$ 遍历 $E$ 中所有 $F$ 的有限 Galois 扩张（参见第 7.6 节练习 10）。

\noindent\textbf{定理（Krull）.} 设 $E/F$ 是 Galois 扩张，Galois 群为 $G$。通过如下方式拓扑化 $G$：取 $G$ 中有限扩张 $F$ 在 $E$ 中的固定子群，以及这些子群的所有左、右陪集，作为闭集的基。则在这个（“Krull”）拓扑下，$G$ 的闭子群与 $E$ 中包含 $F$ 的子域一一对应，并且相应的格是对偶的。$G$ 的闭正规子群对应于 $E$ 中 $F$ 的正规扩张。

当前研究的一个重要方向，是把某些域扩张（例如 $\overline F/F$，其中 $\overline F$ 是 $F$ 的代数闭包）的 Galois 群描述为拓扑群。当 $F=\mathbb Q$ 时，对后一个群所知甚少（特别地，它的有限指数正规子群，即哪些有限群能够作为 $\mathbb Q$ 上的 Galois 群，尚不清楚）。如果 $E$ 是有限域 $\mathbb F_p$ 的代数闭包，则这个扩张的 Galois 群是拓扑循环群 $\widehat{\mathbb Z}$，并以 Frobenius 自同构为拓扑生成元。群 $\widehat{\mathbb Z}$ 是一个不可数群（特别地，它不同构于 $\mathbb Z$），并且具有这样的性质：每个有限指数闭子群都是正规子群，且商群为循环群。注意 $\widehat{\mathbb Z}$ 也必定有非平凡的无限闭子群（不同于 $\mathbb Z$），因为 $E$ 含有在 $\mathbb F_p$ 上无限的真子域（例如，对任意素数 $q$，所有 $\mathbb F_p$ 的 $q$ 幂次数扩张的复合；如第 7.6 节练习 11 所述，这个 $\mathbb F_p$ 的 Galois 扩张的 Galois 群是 $q$-进整数 $\mathbb Z_q$）。

\subsection*{练习}

\begin{enumerate}[label=\arabic*.]
\item 证明每个纯不可分扩张都是正规扩张。

\item 设 $p$ 为素数，令 $K=\mathbb F_p(x,y)$，其中 $x$ 和 $y$ 是 $\mathbb F_p$ 上的相互独立超越元。令
\[
 F=\mathbb F_p(x^p-x,y^p-x).
\]
\begin{enumerate}[label=(\alph*)]
\item 证明 $[K:F]=p^2$，并且 $K/F$ 的可分次数和不可分次数都等于 $p$。
\item 证明存在 $K$ 的一个含 $F$ 的子域 $E$，使得 $E$ 是 $F$ 的次数 $p$ 的纯不可分扩张（于是 $K$ 是 $E$ 的次数 $p$ 的可分扩张）。[令 $s=x^p-x\in F$ 且 $t=y^p-x\in F$，并考虑 $s-t$。]
\end{enumerate}

\item 设 $p$ 为奇素数，设 $s$ 和 $t$ 是 $\mathbb F_p$ 上的相互独立超越元，并令 $F$ 为域 $\mathbb F_p(s,t)$。令 $\beta$ 为 $x^2-sx+t=0$ 的一个根，令 $\alpha$ 为 $x^p-\beta=0$ 的一个根（在 $F$ 的某个代数闭包中）。置 $E=F(\beta)$ 且 $K=F(\alpha)$。
\begin{enumerate}[label=(\alph*)]
\item 证明 $E$ 是 $F$ 的次数 $2$ 的 Galois 扩张，并且 $K$ 是 $E$ 的次数 $p$ 的纯不可分扩张。
\item 证明 $K$ 不是 $F$ 的正规扩张。[如果它是正规扩张，把 $\beta$ 在 $F$ 上共轭，以说明 $K$ 会含有 $s$ 的一个 $p$ 次根，并进而也含有 $t$ 的一个 $p$ 次根，因此 $[K:F]\geq p^2$，矛盾。]
\item 证明不存在域 $K_0$ 使得 $F\subseteq K_0\subseteq K$，且 $K_0/F$ 纯不可分、$K/K_0$ 可分。[如果存在这样的域，使用练习 1 以及两个正规扩张的复合仍为正规这一事实，说明 $K$ 会是 $F$ 的正规扩张。]
\end{enumerate}

\item 在前一题的记号下，证明 $\alpha,s$ 是 $K$ 在 $\mathbb F_p$ 上的一个分离超越基。

\item 设 $p$ 为素数，设 $t$ 在 $\mathbb F_p$ 上超越，并令 $K$ 为从 $\mathbb F_p(t)$ 添入 $t$ 的所有 $p$ 幂次根而得到的域。证明 $K$ 在 $\mathbb F_p$ 上超越次数为 $1$，并且没有分离超越基。

\item 证明：如果 $t$ 在 $\mathbb Q$ 上超越，那么 $\mathbb Q(t,\sqrt{t^3-t})$ 不是 $\mathbb Q$ 的纯超越扩张。（这是一个所谓椭圆函数域的例子。）

\item 令 $k$ 为含 $4$ 个元素的域，$t$ 为 $k$ 上的超越元，$F=k(t^4+t)$ 且 $K=k(t)$。
\begin{enumerate}[label=(\alph*)]
\item 说明 $[K:F]=4$。
\item 说明 $K$ 在 $F$ 上可分。
\item 说明 $K$ 在 $F$ 上是 Galois 的。
\item 描述 Galois 群的子群格和对应的 $K$ 的子域格，并把每个子域写成 $k(r)$ 的形式，其中 $r$ 是某个有理函数。
\end{enumerate}

\item 设 $p$ 为奇素数，$k$ 为特征 $p$ 的代数闭域，设 $t$ 在 $k$ 上超越。假设 $F$ 是 $k(t)$ 的一个次数 $2$ 的域扩张。说明 $F$ 可写成 $k(t,y)$ 的形式，其中 $y\in F$，$y^2\in k(t)$，且 $y$ 在 $k$ 上超越。若 $y^2=4t^3-t-1$，求 $[F:k(y)]$，并把 $k(t)\cap k(y)$ 描述成 $k(r)$，其中 $r\in k(t)$。

\item 设 $t$ 在 $\mathbb F_3$ 上超越，令 $K=\mathbb F_3(t)$，$G=\operatorname{Aut}(K/\mathbb F_3)$，并令 $F$ 为 $G$ 的固定域。
\begin{enumerate}[label=(\alph*)]
\item 证明 $G\cong S_4$，并推出存在唯一的域 $E$，满足 $F\subseteq E\subseteq K$ 且 $[E:F]=2$。[回忆 $G\cong \operatorname{PGL}_2(\mathbb F_3)$；说明 $\operatorname{GL}_2(\mathbb F_3)$ 置换 $\mathbb F_3$ 上二维向量空间中的 $4$ 条直线，而这个置换表示的核为标量矩阵。]
\item 完成 $K$ 中包含 $E$ 的子域格的描述：
\[
\begin{tikzpicture}[scale=0.9, baseline=(current bounding box.center), every node/.style={inner sep=2pt}]
\node (K) at (0,4.0) {$K$};
\node (L1) at (-4.4,2.0) {$E(t^3-t)$};
\node (L2) at (-2.6,1.6) {$**$};
\node (L3) at (-1.1,1.8) {$**$};
\node (L4) at (0.2,2.1) {$E\!\left(\dfrac{t^2+t}{(t-1)^3}\right)$};
\node (L5) at (1.5,2.4) {$**$};
\node (L6) at (3.2,2.8) {$E\!\left(\dfrac{t^2+1}{t+1}\right)$};
\node (L7) at (4.8,2.4) {$**$};
\node (M) at (2.0,0.8) {$**$};
\node (E) at (-0.5,-1.1) {$E\!\left(\dfrac{(t^6+t^4+t^2+1)^2}{(t^3+t)^3}\right)$};
\draw (K)--(L1); \draw (K)--(L2); \draw (K)--(L3); \draw (K)--(L4); \draw (K)--(L5); \draw (K)--(L6); \draw (K)--(L7);
\draw (L1)--(E); \draw (L2)--(E); \draw (L3)--(E); \draw (L4)--(E);
\draw (L5)--(M); \draw (L6)--(M); \draw (L7)--(M); \draw (M)--(E);
\end{tikzpicture}
\]
把每个子域写成 $E(r)$ 的形式，其中 $r$ 是某个有理函数。（$A_4$ 的子群格见第 3.5 节。）
\end{enumerate}

\item 证明一个域的真纯超越扩张绝不代数闭。

\item 设 $S$ 是域 $F$ 上的一组相互独立超越元，设 $\Omega$ 是 $F(S)$ 的一个代数闭包。证明 $S$ 上的任意置换都可扩张为 $\operatorname{Aut}(F(S)/F)$ 的一个元素。证明 $F(S)$ 的任意这种自同构都可扩张为 $\Omega$ 的一个自同构。推出 $\mathbb C$ 有无穷多个自同构。

\item 设 $K$ 是 $\mathbb C$ 的一个子域，它在性质“$\sqrt2\notin K$”下极大。
\begin{enumerate}[label=(\alph*)]
\item 说明这样的域 $K$ 存在。
\item 说明 $\mathbb C$ 在 $K$ 上代数。
\item 证明 $\mathbb C$ 中 $K$ 的每个有限扩张都是 Galois 的，且 Galois 群为循环 $2$-群。
\item 推出 $[\mathbb C:K]$ 是可数的（且不是有限的）。
\end{enumerate}

\item 令 $K$ 为 $\mathbb C$ 的某个自同构的固定域。证明 $\mathbb C$ 中 $K$ 的每个有限扩张都是循环扩张。

\item 令 $K_n$ 为
\[
(x^2-p_1)(x^2-p_2)\cdots(x^2-p_n)
\]
在 $\mathbb Q$ 上的分裂域，其中 $p_1,\ldots,p_n$ 是前 $n$ 个素数。证明 $K_n/\mathbb Q$ 的 Galois 群是阶为 $2^n$ 的初等 Abel $2$-群。

\item 令 $K_0=\mathbb Q$，并对 $n\geq 0$ 定义域 $K_{n+1}$ 为从 $K_n$ 添入 $K_n$ 上所有三次多项式的所有根得到的扩张。令 $K$ 为子域 $K_n$，$n\geq 0$，的并。证明 $K$ 是 $\mathbb Q$ 的 Galois 扩张。证明 $K$ 上的每个三次多项式都在 $K$ 上完全分裂。证明 $K$ 有非平凡代数扩张。

\item 令 $F$ 为 $\mathbb Q$ 上所有不可约三次多项式的分裂域的复合。证明 $F$ 不含 $\mathbb Q$ 的所有二次扩张。

\item 令 $K_0=\mathbb Q$，并对 $n\geq 0$ 定义域 $K_{n+1}$ 为从 $K_n$ 添入 $K_n$ 的元素的所有根式而得到的扩张。令 $K$ 为子域 $K_n$，$n\geq 0$，的并。证明 $K$ 是 $\mathbb Q$ 的 Galois 扩张。证明 $K$ 没有非平凡可解 Galois 扩张。证明 $K$ 有非平凡 Galois 扩张。

\item 令 $F_0=\mathbb Q$，并对 $n\geq 0$ 定义域 $F_{n+1}$ 为从 $F_n$ 添入 $F_n$ 的元素的所有实根式而得到的扩张。令 $F$ 为子域 $F_n$，$n\geq 0$，的并。令 $K^+$ 为前一练习中域 $K$ 上复共轭限制的固定域（$K$ 的最大实子域）。证明 $F\neq K^+$。

\item 本练习证明：若 $K/F$ 是域的 Galois 扩张，则
\[
\operatorname{Gal}(K/F)\cong \varprojlim \operatorname{Gal}(L/F),
\]
其中逆极限取遍 $K$ 中包含的 $F$ 的所有有限 Galois 扩张 $L$。
\begin{enumerate}[label=(\alph*)]
\item 说明 $K$ 是这些域 $L$ 的并。
\item 证明映射
\[
\varphi:\operatorname{Gal}(K/F)\longrightarrow \varprojlim \operatorname{Gal}(L/F)
\]
由把 $\operatorname{Gal}(K/F)$ 中的 $\sigma$ 映到 $(\ldots,\sigma|_L,\ldots)$ 来定义，其中 $\sigma|_L$ 是 $\sigma$ 到 $L$ 的限制，是一个同态。
\item 说明 $\varphi$ 是单射。
\item 如果 $(\ldots,\sigma_L,\ldots)\in \varprojlim \operatorname{Gal}(L/F)$，定义 $\sigma\in \operatorname{Gal}(K/F)$ 为：若 $a\in L$，则 $\sigma(a)=\sigma_L(a)$。证明 $\sigma$ 是良定义的自同构，并推出 $\varphi$ 是满射。
\end{enumerate}
\end{enumerate}


% \setchapterintro{在本部分中，我们继续研究环与模，首先集中讨论交换环。交换代数本身是一个重要主题，同时也是代数学其他领域的基本基础。为说明所引入代数主题的重要性，本章在发展环论的同时，也并行给出仿射代数几何的导论。每一节先介绍基本的代数理论，然后给出这些思想在几何中的应用，并说明如何用第 9 章的 Gröbner 基理论进行计算。这样做有两个目的：一是给出代数技巧在代数几何这一重要数学分支中的应用；二是从几何问题的来源说明这些代数概念的动机。
% 	
% 	几何与代数之间的这种联系展示了两个数学领域之间丰富的相互作用，也再次说明，一个数学思想圈中的结果与结构如何为另一个思想圈提供洞见。
% 	
% 	第 16 章将继续讨论交换环中的若干基本结构，最终给出 Dedekind 整环以及这类环上模的结构定理；这是第 12 章中主理想整环上模结构定理的推广。第 17 章将介绍同调代数的一些基本技术，这些技术继续处理第 10 章中一些序列不再正合时所产生的问题。本章中的群上同调既作为同调代数较深入的应用，也作为这一重要主题的相对自洽的介绍。}
%\part{交换环，代数几何与同调代数引论}
% % chapter 15

% ============================================
% chapter15_part1.tex — 第15章：交换环与代数几何（第一部分）
% 根据 15 part1.pdf 翻译整理；练习部分逐题完整保留
% 范围：Part V 引言、15.1、15.2、15.3 至本 PDF 结尾
% ============================================





\chapter{交换环与代数几何}

本章中，$R$ 始终表示含有 $1\ne0$ 的交换环。

\section{Noether 环与仿射代数集}

本节更详细地研究 Noether 环。Noether 环是主理想整环的自然推广，在第 12 章已作过简要介绍。注意，当 $R$ 被看成它自身上的左模时，$R$ 的 $R$-子模正好是 $R$ 的理想。因此第 12 章第 1 节中的定义可在这里改写如下。

\noindent\textbf{定义.} 交换环 $R$ 称为 \emph{Noether 环}，或称满足理想上的升链条件（理想 A.C.C.），若在 $R$ 中不存在无限严格上升的理想链；换言之，若
\[
I_1\subseteq I_2\subseteq I_3\subseteq\cdots
\]
是 $R$ 的一个上升理想链，则存在正整数 $m$，使得对所有 $k\ge m$ 都有 $I_k=I_m$。

\noindent\textbf{命题 1.} 若 $R$ 是 Noether 环且 $I$ 是 $R$ 的一个理想，则商环 $R/I$ 也是 Noether 环。Noether 环的任意同态像仍为 Noether 环。

\begin{proof}
	商环 $R/I$ 中的任意无限上升理想链，都由格同构定理对应于 $R$ 中包含 $I$ 的无限上升理想链。因此第一断言成立；第二断言由第一同构定理推出。
\end{proof}

\noindent\textbf{定理 2.} 下列条件等价：
\begin{enumerate}[label=(\arabic*)]
	\item $R$ 是 Noether 环；
	\item $R$ 的任意非空理想集合在包含关系下都有极大元；
	\item $R$ 的每个理想都是有限生成的。
\end{enumerate}

\begin{proof}
	证明与第 12.1 节定理 1 在 $R$-模 $M=R$ 的特殊情形完全相同，此时子模即为理想。
\end{proof}

\paragraph{例}
每个主理想整环都是 Noether 环，因为它满足定理 2 的条件 (3)。特别地，$\mathbb Z$、域 $k$ 上的多项式环 $k[x]$、Gaussian 整数环 $\mathbb Z[i]$ 都是 Noether 环。环 $\mathbb Z[x_1,x_2,\ldots]$ 不是 Noether 环，因为理想 $(x_1,x_2,\ldots)$ 不能由有限个元素生成。第 7.4 节练习 33(d) 表明，闭区间 $[0,1]$ 上的实值连续函数环也不是 Noether 环。

Noether 环可以有任意长的上升理想链，也可以有无限长的下降理想链。例如
\[
(2)\supset (4)\supset (8)\supset\cdots
\]
是 $\mathbb Z$ 中的无限下降链。因此 Noether 环不必满足理想下降链条件（D.C.C.）。不过我们将在第 16 章看到，满足理想 D.C.C. 的交换环必定也满足 A.C.C.，因而是 Noether 环；这样的环称为 Artinian 环。

\noindent\textbf{定理 3（Hilbert 基定理）.} 若 $R$ 是 Noether 环，则多项式环 $R[x]$ 也是 Noether 环。

Hilbert 基定理说明，可以从已知的 Noether 环出发构造更大的 Noether 环，这类似于第 9.3 节定理 7：若 $R$ 是 U.F.D.，则 $R[x]$ 也是 U.F.D.。

\noindent\textbf{推论 4.} 若 $k$ 是域，则有限多个变量的多项式环
\[
k[x_1,x_2,\ldots,x_n]
\]
是 Noether 环。

回忆：若 $k$ 是域，$k$-代数是含有 $k$ 作为中心子环且二者恒等元相同的环。

\noindent\textbf{定义.}
\begin{enumerate}[label=(\arabic*)]
	\item 若环 $R$ 作为环由 $k$ 与有限多个元素 $r_1,r_2,\ldots,r_n\in R$ 生成，则称 $R$ 是\emph{有限生成的 $k$-代数}。
	\item 若 $R,S$ 是 $k$-代数，则映射 $\psi:R\to S$ 称为 \emph{$k$-代数同态}，若 $\psi$ 是环同态且在 $k$ 上为恒等映射。
\end{enumerate}

若 $R$ 是 $k$-代数，则 $R$ 同时是环和 $k$ 上的向量空间。必须区分元素生成 $R$ 的意义。例如，$k[x_1,\ldots,x_n]$ 是有限生成 $k$-代数，因为 $x_1,\ldots,x_n$ 是环生成元；但当 $n>0$ 时，它作为 $k$-向量空间是无限维的。

\noindent\textbf{推论 5.} 环 $R$ 是有限生成 $k$-代数，当且仅当存在某个满的 $k$-代数同态
\[
\varphi:k[x_1,x_2,\ldots,x_n]\longrightarrow R
\]
从有限多个变量的多项式环映到 $R$，且在 $k$ 上为恒等映射。因此任意有限生成 $k$-代数都是 Noether 环。

\begin{proof}
	若 $R$ 由 $r_1,\ldots,r_n$ 作为 $k$-代数生成，则定义 $\varphi(x_i)=r_i$，且 $\varphi(a)=a$ 对 $a\in k$ 成立，唯一地扩张为满环同态。反之，若有这样的满同态，则 $x_i$ 的像生成 $R$。最后，由推论 4 与命题 1，有限生成 $k$-代数作为 Noether 环的商仍是 Noether 环。
\end{proof}

\paragraph{例}
若 $k$-代数 $R$ 作为 $k$-向量空间有限维，例如 $R=k[x]/(f(x))$，其中 $f\in k[x]$ 非零，则 $R$ 是有限生成 $k$-代数。此时理想都是 $k$-子空间，因此任意上升或下降理想链至多有 $\dim_k R+1$ 个不同项，所以 $R$ 同时满足 A.C.C. 与 D.C.C.。

\subsection*{仿射代数集}

代数几何的基本思想，是把几何问题转化为诸如 $k[x_1,\ldots,x_n]$ 这类环中理想的代数问题。由于这些环是 Noether 环，许多问题可化为有限多个代数方程的研究，这也是 Hilbert 基定理最初的主要动机之一。

记 $\mathbb A^n$ 为域 $k$ 上的仿射 $n$-空间，即 $k$ 中元素的 $n$ 元组集合。多项式 $f\in k[x_1,\ldots,x_n]$ 可看作 $\mathbb A^n$ 上的 $k$ 值函数
\[
f:(a_1,
\ldots,a_n)\longmapsto f(a_1,\ldots,a_n)\in k.
\]
这些函数构成一个函数环，记为 $k[\mathbb A^n]$，称为 $\mathbb A^n$ 的坐标环。

若 $S\subseteq k[\mathbb A^n]$，定义
\[
Z(S)=\{(a_1,\ldots,a_n)\in\mathbb A^n\mid f(a_1,\ldots,a_n)=0\text{ 对所有 }f\in S\},
\]
其中 $Z(0)=\mathbb A^n$。

\noindent\textbf{定义.} 若 $V\subseteq\mathbb A^n$ 是某些多项式的公共零点集，即 $V=Z(S)$ 对某个 $S\subseteq k[\mathbb A^n]$ 成立，则称 $V$ 是一个\emph{仿射代数集}（或简称代数集）。此时 $V=Z(S)$ 称为 $S$ 在 $\mathbb A^n$ 中的\emph{轨迹}。

若 $S=\{f\}$ 或 $S=\{f_1,\ldots,f_m\}$，则写作 $Z(f)$ 或 $Z(f_1,\ldots,f_m)$。

\paragraph{例}
\begin{enumerate}[label=(\arabic*)]
	\item 若 $n=1$，则单个多项式 $f\in k[x]$ 的轨迹是 $f$ 在 $k$ 中的根集合。$\mathbb A^1$ 中的代数集正是 $\varnothing$、有限集以及整个 $k$。
	\item 任一点集 $\{(a_1,\ldots,a_n)\}$ 是代数集，因为它是
	\[
	Z(x_1-a_1,x_2-a_2,\ldots,x_n-a_n).
	\]
	更一般地，任意有限子集都是仿射代数集。
	\item 直线、平面等线性代数集是线性多项式的零点集。例如 $\mathbb A^2$ 中直线 $ax+by=c$ 是 $f(x,y)=ax+by-c$ 的轨迹；$\mathbb A^3$ 中 $x$ 轴为 $Z(y,z)$，$xy$ 平面为 $Z(z)$。
	\item 非常数多项式 $f$ 的零点集 $Z(f)$ 称为 $\mathbb A^n$ 中的超曲面。圆锥曲线与二次曲面都是这类例子。
\end{enumerate}

对 $S,T\subseteq k[\mathbb A^n]$，有下列基本性质：
\begin{enumerate}[label=(\arabic*)]
	\item 若 $S\subseteq T$，则 $Z(T)\subseteq Z(S)$。
	\item $Z(S)=Z(I)$，其中 $I=(S)$ 是由 $S$ 生成的理想。
	\item 两个代数集的交仍为代数集，且 $Z(S)\cap Z(T)=Z(S\cup T)$；任意交亦然。
	\item 两个代数集的并仍为代数集，若 $I,J$ 是理想，则 $Z(I)\cup Z(J)=Z(IJ)$。
	\item $Z(0)=\mathbb A^n$，$Z(1)=\varnothing$。
\end{enumerate}

因此可以看作有一个反变映射
\[
Z:\{k[\mathbb A^n]\text{ 的理想}\}\longrightarrow\{\mathbb A^n\text{ 中的仿射代数集}\}.
\]
由于 $k[x_1,\ldots,x_n]$ 是 Noether 环，每个理想有限生成，所以每个仿射代数集都是有限多个超曲面的交。

若 $A\subseteq\mathbb A^n$，定义
\[
I(A)=\{f\in k[x_1,\ldots,x_n]\mid f(a_1,\ldots,a_n)=0\text{ 对所有 }(a_1,\ldots,a_n)\in A\}.
\]
这是一个理想，并且是所有在 $A$ 上恒为零的函数构成的唯一最大理想。于是有反变对应
\[
I:\{\mathbb A^n\text{ 的子集}\}\longrightarrow\{k[\mathbb A^n]\text{ 的理想}\}.
\]

\paragraph{例}
\begin{enumerate}[label=(\arabic*)]
	\item 在欧氏平面中，$x$ 轴的理想是 $\mathbb R[x,y]$ 中的 $(y)$。
	\item 点 $(a_1,
	\ldots,a_n)$ 的消没理想是
	\[
	I((a_1,
	\ldots,a_n))=(x_1-a_1,
	\ldots,x_n-a_n),
	\]
	这是极大理想。
	\item 若 $V=Z(x^3-y^2)\subseteq\mathbb A^2$，且 $k$ 无限，则
	\[
	I(V)=(x^3-y^2)\subseteq k[x,y].
	\]
	若 $k$ 有限，则可能有不属于 $(x^3-y^2)$ 但在 $V$ 上为零的多项式。
\end{enumerate}

映射 $I$ 满足：若 $A\subseteq B$，则 $I(B)\subseteq I(A)$；$I(A\cup B)=I(A)\cap I(B)$；$I(\varnothing)=k[x_1,\ldots,x_n]$；若 $k$ 无限，则 $I(\mathbb A^n)=0$。此外，$A\subseteq Z(I(A))$，$I\subseteq I(Z(I))$；若 $V=Z(I)$ 是代数集，则 $V=Z(I(V))$。

\noindent\textbf{定义.} 若 $V\subseteq\mathbb A^n$ 是仿射代数集，则商环
\[
k[\mathbb A^n]/I(V)
\]
称为 $V$ 的\emph{坐标环}，记为 $k[V]$。

若 $V=Z(xy-1)\subseteq\mathbb R^2$ 是双曲线 $y=1/x$，则
\[
\mathbb R[V]=\mathbb R[x,y]/(xy-1)\cong \mathbb R[x,1/x].
\]

设 $V\subseteq\mathbb A^n$，$W\subseteq\mathbb A^m$ 是两个仿射代数集。

\noindent\textbf{定义.} 映射 $\varphi:V\to W$ 称为代数集的\emph{态射}（或多项式映射、正则映射），若存在多项式 $\varphi_1,\ldots,\varphi_m\in k[x_1,
\ldots,x_n]$，使得
\[
\varphi((a_1,
\ldots,a_n))=(\varphi_1(a_1,
\ldots,a_n),
\ldots,
\varphi_m(a_1,
\ldots,a_n)).
\]
若存在态射 $\psi:W\to V$ 使得 $\varphi\circ\psi=1_W$ 且 $\psi\circ\varphi=1_V$，则称 $\varphi$ 是代数集同构。

任一态射 $\varphi:V\to W$ 给出 $k$-代数同态
\[
\widetilde\varphi:k[W]\longrightarrow k[V],\qquad f\longmapsto f\circ\varphi.
\]
反过来，任一 $k$-代数同态 $k[W]\to k[V]$ 由一个唯一的态射 $V\to W$ 诱导。

\noindent\textbf{定理 6.} 设 $V\subseteq\mathbb A^n$ 与 $W\subseteq\mathbb A^m$ 是仿射代数集。则有双射对应
\[
\{V\to W\text{ 的代数集态射}\}
\longleftrightarrow
\{k[W]\to k[V]\text{ 的 }k\text{-代数同态}\}.
\]
更具体地：
\begin{enumerate}[label=(\arabic*)]
	\item 每个态射 $\varphi:V\to W$ 诱导 $\widetilde\varphi(f)=f\circ\varphi$；
	\item 每个 $k$-代数同态都由唯一态射诱导；
	\item 若 $\varphi:V\to W$ 与 $\psi:W\to U$ 是态射，则 $\widetilde{\psi\circ\varphi}=\widetilde\varphi\circ\widetilde\psi$；
	\item $\varphi$ 是同构当且仅当 $\widetilde\varphi$ 是 $k$-代数同构。
\end{enumerate}

\paragraph{例}
若 $k$ 是无限域，令 $V=\mathbb A^1$，$W=Z(x^3-y^2)=\{(a^2,a^3)\mid a\in k\}$。映射 $\varphi(a)=(a^2,a^3)$ 是从 $V$ 到 $W$ 的态射且为双射。其坐标环同态为
\[
\widetilde\varphi:k[x,y]/(x^3-y^2)\longrightarrow k[x],
\qquad x\mapsto x^2,
\quad y\mapsto x^3.
\]
其像为 $k[x^2,x^3]$，不是 $k[x]$，所以 $\varphi$ 不是代数集同构，尽管它是集合双射。

\noindent\textbf{推论 7.} 设 $\varphi:V\to W$ 是仿射代数集之间的映射。则 $\varphi$ 是态射，当且仅当对每个 $f\in k[W]$，复合函数 $f\circ\varphi$ 是 $k[V]$ 的元素。若 $\varphi$ 是态射，则对 $v\in V,w\in W$，有 $\varphi(v)=w$ 当且仅当
\[
\widetilde\varphi^{-1}(I(\{v\}))=I(\{w\}).
\]

\subsection*{仿射代数集与 $k$-代数中的计算}

第 9.6 节的 Gröbner 基理论在仿射代数集计算中非常有用，尤其是在坐标环 $k[\mathbb A^n]/I(V)$ 中计算时。若 $G$ 是 $I(V)$ 的 Gröbner 基，则每个多项式 $f$ 除以 $G$ 后的唯一余式给出了商环中余类 $\bar f$ 的唯一代表。

更一般地，若
\[
\Phi:k[y_1,
\ldots,y_m]/J\longrightarrow k[x_1,
\ldots,x_n]/I
\]
是 $k$-代数同态，且 $\Phi(y_i)=\varphi_i+I$，则可用消元 Gröbner 基计算 $\ker\Phi$ 与 $\operatorname{im}\Phi$。

\noindent\textbf{命题 8.} 沿用上述记号。令
\[
R=k[y_1,
\ldots,y_m,x_1,
\ldots,x_n],
\]
令 $A$ 为由 $y_i-
\varphi_i$ 与 $I$ 的生成元生成的理想。设 $G$ 为 $A$ 关于字典序
\[
x_1>\cdots>x_n>y_1>\cdots>y_m
\]
的约化 Gröbner 基。则：
\begin{enumerate}[label=(\alph*)]
	\item $\ker\Phi$ 是 $A\cap k[y_1,
	\ldots,y_m]$ 模 $J$；$G$ 中属于 $k[y_1,
	\ldots,y_m]$ 的元素（模 $J$）生成 $\ker\Phi$。
	\item 对 $f\in k[x_1,
	\ldots,x_n]$，$\bar f$ 属于 $\operatorname{im}\Phi$ 当且仅当 $f$ 除以 $G$ 后的余式是某个 $h\in k[y_1,
	\ldots,y_m]$；此时 $\Phi(h)=\bar f$。
\end{enumerate}

\noindent\textbf{推论 9.} 同态 $\Phi$ 满，当且仅当对每个 $i=1,
\ldots,n$，约化 Gröbner 基 $G$ 中含有形如
\[
x_i-h_i,
\qquad h_i\in k[y_1,
\ldots,y_m]
\]
的多项式。

\paragraph{例}
令 $\Phi:\mathbb Q[u,v]\to\mathbb Q[x]$ 由 $\Phi(u)=x^2+x$ 与 $\Phi(v)=x^3$ 定义。对理想 $A=(u-x^2-x,v-x^3)$，按字典序 $x>u>v$ 的约化 Gröbner 基为
\[
g_1=x^2+x-u,
\quad
g_2=ux+x-u-v,
\quad
g_3=vx-x-u^2+u+2v,
\quad
g_4=u^3-3uv-v^2-v.
\]
故 $\ker\Phi=(g_4)$，且 $\Phi$ 不满。

\noindent\textbf{命题 10.} 设 $\alpha$ 是不可约多项式 $p(x)\in k[x]$ 的根，且 $\beta=f(\alpha)\in k(\alpha)$。令 $G$ 为 $k[x,y]$ 中理想 $(p,y-f)$ 关于字典序 $x>y$ 的约化 Gröbner 基。则 $\beta$ 在 $k$ 上的最小多项式是 $G\cap k[y]$ 中的首一多项式。

\paragraph{例}
取 $k=\mathbb Q$，令 $\beta=1+\sqrt[3]{2}+3\sqrt[3]{4}\in\mathbb Q(\sqrt[3]{2})$。理想
\[
(x^3-2,
y-(1+x+3x^2))
\]
在 $\mathbb Q[x,y]$ 中关于字典序 $x>y$ 的约化 Gröbner 基包含
\[
y^3-3y^2-15y-93,
\]
因此 $\beta$ 的最小多项式为 $y^3-3y^2-15y-93$。

\subsection*{第15.1节练习}

设 $R$ 是含 $1\ne0$ 的交换环，$k$ 是域。

\begin{enumerate}[label=\arabic*.]
	\item 证明 Hilbert 基定理的逆命题：若多项式环 $R[x]$ 是 Noether 环，则 $R$ 是 Noether 环。
	
	\item 对下列每个环，给出一个显式的无限上升理想链，以证明它们不是 Noether 环：
	\begin{enumerate}[label=(\alph*)]
		\item 闭区间 $[0,1]$ 上的实值连续函数环；
		\item 从任意无限集合 $X$ 到 $\mathbb Z/2\mathbb Z$ 的所有函数构成的环。
	\end{enumerate}
	
	\item 证明变量 $x$ 上 $k$ 的有理函数域 $k(x)$ 不是有限生成 $k$-代数。（回忆：$k(x)$ 是多项式环 $k[x]$ 的分式域。注意 $k(x)$ 是 $k$ 上有限生成的域扩张。）
	
	\item 证明：若 $R$ 是 Noether 环，则形式幂级数环 $R[[x]]$ 也是 Noether 环（参见第 7.2 节练习 3）。[仿照 Hilbert 基定理的证明。]
	
	\item（Fitting 引理）设 $M$ 是 Noether $R$-模，$\theta:M\to M$ 是 $M$ 的 $R$-模自同态。证明对充分大的 $n$，有
	\[
	\ker(\theta^n)\cap\operatorname{image}(\theta^n)=0.
	\]
	证明若 $\theta$ 是满射，则 $\theta$ 是同构。[观察 $\ker(\theta)\subseteq\ker(\theta^2)\subseteq\cdots$。]
	
	\item 设
	\[
	0\longrightarrow M'\longrightarrow M\longrightarrow M''\longrightarrow0
	\]
	是 $R$-模的正合列。证明 $M$ 是 Noether $R$-模，当且仅当 $M'$ 与 $M''$ 都是 Noether $R$-模。
	
	\item 证明 Noether $R$-模的子模、商模以及有限直和仍是 Noether $R$-模。
	
	\item 若 $R$ 是 Noether 环，证明 $M$ 是 Noether $R$-模，当且仅当 $M$ 是有限生成 $R$-模。（因此，Noether 环上有限生成模的任意子模也有限生成。）
	
	\item 对域 $k$，证明多项式环 $k[x]$ 中任何含有 $k$ 的子环都是 Noether 环。举例说明这样的子环不一定是 U.F.D.。[若 $k\subseteq R\subseteq k[x]$ 且 $y\in R-k$，证明 $k[x]$ 是有限生成 $k[y]$-模；然后使用前两个练习。第二部分考虑 $k[x^2,x^3]$。]
	
	\item 证明多项式环 $k[x,y]$ 的子环
	\[
	k[x,x^2y,x^3y^2,
	\ldots,x^iy^{i-1},
	\ldots]
	\]
	不是 Noether 环，因此不是有限生成 $k$-代数。（所以 Noether 环的子环未必 Noether，有限生成 $k$-代数的子代数未必有限生成。）
	
	\item 设 $R$ 是交换环，且所有素理想都是有限生成的。本练习证明 $R$ 是 Noether 环。
	\begin{enumerate}[label=(\alph*)]
		\item 假设 $R$ 中非有限生成理想的集合非空，证明该集合在包含关系下有极大元 $I$，并且 $R/I$ 是 Noether 环。
		\item 证明存在包含 $I$ 的有限生成理想 $J_1,J_2$，使得 $J_1J_2\subseteq I$，且 $J_1J_2$ 有限生成。[观察 $I$ 不是素理想。]
		\item 证明 $I/J_1J_2$ 是 $J_1/J_1J_2$ 的有限生成 $R/I$-子模。[使用练习 8。]
		\item 说明 (c) 蕴含矛盾：$I$ 本应作为 $R$-模有限生成；由此推出 $R$ 是 Noether 环。
	\end{enumerate}
	
	\item 设 $R$ 是 Noether 环，$S$ 是有限生成 $R$-代数。若 $T\subseteq S$ 是一个 $R$-代数，并且 $S$ 是有限生成 $T$-模，证明 $T$ 是有限生成 $R$-代数。[若 $s_1,
	\ldots,s_n$ 生成 $S$ 作为 $R$-代数，且 $s'_1,
	\ldots,s'_m$ 生成 $S$ 作为 $T$-模，说明元素 $s_i$ 与 $s_js_i$ 可写成 $s'_j$ 的有限 $T$-线性组合。若 $T_0$ 是由这些线性组合系数生成的 $R$-子代数，证明 $S$（从而 $T_0$）是有限生成 $T_0$-模，并推出 $T$ 是有限生成 $R$-代数。]
	
	\item 验证映射 $Z$ 与 $I$ 的性质 (1) 到 (10)。
	
	\item 证明任意域 $k$ 上 $\mathbb A^1$ 中的仿射代数集正是 $\varnothing$、$k$ 以及 $k$ 的有限子集。
	
	\item 若 $k=\mathbb F_2$ 且 $V=\{(0,0),(1,1)\}\subseteq\mathbb A^2$，证明 $I(V)$ 是乘积理想 $\mathfrak m_1\mathfrak m_2$，其中
	\[
	\mathfrak m_1=(x,y),
	\qquad
	\mathfrak m_2=(x-1,y-1).
	\]
	
	\item 设 $V$ 是 $\mathbb A^n$ 中的有限代数集。若 $V$ 有 $m$ 个点，证明 $k[V]$ 作为 $k$-代数同构于 $k^m$。[使用中国剩余定理。]
	
	\item 若 $k$ 是有限域，证明 $\mathbb A^n$ 的每个子集都是仿射代数集。
	
	\item 若 $k=\mathbb F_q$ 是含 $q$ 个元素的有限域，证明
	\[
	I(\mathbb A^1)=(x^q-x)\subseteq k[x].
	\]
	
	\item 对每个非常数多项式 $f\in k[x]$，用 $f$ 在 $k[x]$ 中的唯一分解描述 $Z(f)\subseteq\mathbb A^1$，再据此描述 $I(Z(f))$。推出 $I(Z(f))=(f)$ 当且仅当 $f$ 是 $k[x]$ 中互异线性因子的乘积。
	
	\item 若 $f,g$ 是 $k[x,y]$ 中不相伴的不可约多项式（即互不整除），证明 $Z((f,g))$ 在 $\mathbb A^2$ 中要么为空，要么有限。[若 $(f,g)\ne(1)$，令 $R=k[x]$，$F=k(x)$，使用 Gauss 引理证明 $f,g$ 在 $F[y]$ 中互素，从而说明 $(f,g)$ 含有 $k[x]$ 中非零多项式；类似地也含有 $k[y]$ 中非零多项式。]
	
	\item 把每个 $2\times2$ 矩阵
	\[
	\begin{pmatrix}a&b\\ c&d\end{pmatrix}
	\]
	视为 $\mathbb A^4$ 中的点 $(a,b,c,d)$。证明行列式为 $1$ 的矩阵群 $SL_2(k)$ 是 $\mathbb A^4$ 中的代数集。
	
	\item 证明 $SL_n(k)$ 是 $\mathbb A^{n^2}$ 中的仿射代数集。[推广前一题。]
	
	\item 设 $V$ 是 $\mathbb R^2$ 中任意一条直线（某个非零线性多项式 $ax+by-c$ 的零点集）。证明 $\mathbb R[V]$ 作为 $\mathbb R$-代数同构于多项式环 $\mathbb R[x]$，并给出从 $\mathbb A^1$ 到 $V$ 的相应同构。
	
	\item 令 $V=Z(xy-z)\subseteq\mathbb A^3$。证明 $V\cong\mathbb A^2$，给出显式同构 $\varphi$ 及相应的 $k$-代数同构 $\widetilde\varphi:k[V]\to k[\mathbb A^2]$，以及它们的逆。问：$V=Z(xy-z^2)$ 是否同构于 $\mathbb A^2$？
	
	\item 设 $V\subseteq\mathbb A^n$ 是仿射代数集，$f\in k[V]$。$f$ 的图像是 $\mathbb A^{n+1}$ 中的点集
	\[
	\{(a_1,
	\ldots,a_n,f(a_1,
	\ldots,a_n))\}.
	\]
	证明 $f$ 的图像是与 $V$ 同构的仿射代数集。\linebreak\ [一个方向的态射把 $(a_1,
	\ldots,a_n)$ 送到 $(a_1,
	\ldots,a_n,f(a_1,
	\ldots,a_n))$。]
	
	\item 令
	\[
	V=Z(xz-y^2,
	yz-x^3,
	z^2-x^2y)\subseteq\mathbb A^3.
	\]
	\begin{enumerate}[label=(\alph*)]
		\item 证明映射 $\varphi:\mathbb A^1\to V$，$\varphi(t)=(t^3,t^4,t^5)$ 是满态射。[关于满性，若 $(x,y,z)\ne(0,0,0)$，令 $t=y/x$。]
		\item 显式描述对应的 $k$-代数同态 $\widetilde\varphi:k[V]\to k[\mathbb A^1]$。
		\item 证明 $\varphi$ 不是同构。
	\end{enumerate}
	
	\item 设 $\varphi:V\to W$ 是仿射代数集的态射。若 $W'$ 是 $W$ 的仿射代数子集，证明原像
	\[
	V'=\varphi^{-1}(W')
	\]
	是 $V$ 的仿射代数子集。若 $W'=Z(I)$，证明对相应的同态 $\widetilde\varphi:k[W]\to k[V]$，有
	\[
	V'=Z(\widetilde\varphi(I)).
	\]
	
	\item 证明若 $V,W$ 是仿射代数集，则 $V\times W$ 也是仿射代数集，并且
	\[
	k[V\times W]\cong k[V]\otimes_k k[W].
	\]
	
	\item 若 $R=\mathbb Z$，$M=\mathbb Z/n\mathbb Z$，证明 $\operatorname{Ass}_R(M)$ 由所有 $(p)$ 组成，其中 $p$ 遍历 $n$ 的素因子。
	
	\item 若 $M$ 是某族子模 $M_i$ 的并，证明
	\[
	\operatorname{Ass}_R(M)=\bigcup_i \operatorname{Ass}_R(M_i).
	\]
	
	\item 假设 $\operatorname{Ann}(m)=P$，即 $Rm\cong R/P$。证明若 $0\ne m'\in Rm$，则 $\operatorname{Ann}(m')=P$。推出 $\operatorname{Ass}_R(R/P)=\{P\}$。[观察 $R/P$ 是整环。]
	
	\item 设 $M$ 是 $R$-模，$P$ 是所有形如 $\operatorname{Ann}(m)$ 的理想中的极大元，其中 $m\in M$。证明 $P$ 是素理想。[若 $P=\operatorname{Ann}(m)$ 且 $ab\in P$，证明 $bm\ne0$ 蕴含 $\operatorname{Ann}(m)\subseteq\operatorname{Ann}(bm)$，再用 $P$ 的极大性推出 $a\in\operatorname{Ann}(bm)=P$。]
	
	\item 若 $R$ 是 Noether 环且 $M\ne0$ 是 $R$-模，证明 $\operatorname{Ass}_R(M)\ne\varnothing$。[使用练习 32。]
	
	\item 若 $L$ 是 $M$ 的子模，$N\cong M/L$，证明有包含关系
	\[
	\operatorname{Ass}_R(N)\subseteq\operatorname{Ass}_R(M)
	\subseteq
	\operatorname{Ass}_R(L)\cup\operatorname{Ass}_R(N),
	\]
	并说明这两个包含都可能是真包含。[若 $Rm\cong R/P$，说明 $Rm\cap L=0$ 蕴含 $P\in\operatorname{Ass}_R(N)$，而 $Rm\cap L\ne0$ 蕴含 $P\in\operatorname{Ass}_R(L)$（由练习 31）。第二部分考虑 $n\mathbb Z\subseteq\mathbb Z$。]
	
	\item 设 $M$ 是 $R$-模，$S$ 是 $\operatorname{Ass}_R(M)$ 中素理想的一个子集。证明存在 $M$ 的子模 $N$，使得
	\[
	\operatorname{Ass}_R(N)=S,
	\qquad
	\operatorname{Ass}_R(M/N)=\operatorname{Ass}_R(M)-S.
	\]
	[考虑 $M$ 的子模 $N'$ 的集合，使得 $\operatorname{Ass}_R(N')\subseteq S$。用练习 30 与 Zorn 引理取极大子模 $N$。若 $P\in\operatorname{Ass}_R(M/N)$，有子模 $M'/N\cong R/P$。用前一题说明 $\operatorname{Ass}_R(M')\subseteq\operatorname{Ass}_R(N)\cup\{P\}$，再由 $N$ 的极大性推出 $P\in\operatorname{Ass}_R(M)-S$；最后再次用前一题得到等号。]
	
	\item 设 $M$ 由 $m_1,
	\ldots,m_n$ 生成。$M$ 的 Fitting 理想 $\mathcal F_R(M)$（层级 $0$）定义为由所有 $n\times n$ 关系矩阵 $A=(r_{ij})$ 的行列式生成的理想，其中 $r_{i1}m_1+\cdots+r_{in}m_n=0$。证明：
	\begin{enumerate}[label=(\alph*)]
		\item $\mathcal F_R(M)$ 也等于所有 $p\times n$ 关系矩阵（$p\ge1$）的所有 $n\times n$ 子式生成的理想。
		\item 对固定的关系矩阵 $A$，若 $A'$ 由 $A$ 经过任意初等行变换或列变换得到，则 $A$ 的所有 $n\times n$ 子式生成的理想与 $A'$ 的所有 $n\times n$ 子式生成的理想相同。
	\end{enumerate}
	
	\item 设 $m_1,
	\ldots,m_n$ 与 $m'_1,
	\ldots,m'_{n'}$ 是 $M$ 的两组 $R$-模生成元。令 $\mathcal F$ 为用 $m_1,
	\ldots,m_n$ 计算得到的 Fitting 理想，$\mathcal F'$ 为用 $m_1,
	\ldots,m_n,m'_1,
	\ldots,m'_{n'}$ 计算得到的 Fitting 理想。
	\begin{enumerate}[label=(\alph*)]
		\item 证明 $m'_s=a_{s1}m_1+\cdots+a_{sn}m_n$，并推出
		\[
		(-a_{s1},
		\ldots,-a_{sn},0,
		\ldots,0,1,0,
		\ldots,0)
		\]
		是 $m_1,
		\ldots,m_n,m'_1,
		\ldots,m'_{n'}$ 的一个关系。
		\item 若 $A=(r_{ij})$ 是一个 $n\times n$ 关系矩阵，其行是 $m_1,
		\ldots,m_n$ 的关系系数，证明 $\det A=\det A'$，其中 $A'$ 是一个 $(n+n')\times(n+n')$ 关系矩阵，其行是 $m_1,
		\ldots,m_n,m'_1,
		\ldots,m'_{n'}$ 的关系系数。推出 $\mathcal F\subseteq\mathcal F'$。[用 (a) 构造块上三角矩阵 $A'$，其左上角为 $A$，右下角为 $n'\times n'$ 单位矩阵。]
		\item 证明 $\mathcal F'\subseteq\mathcal F$，从而 $\mathcal F'=\mathcal F$。[使用前一练习。]
		\item 由 (c) 推出 $M$ 的 Fitting 理想 $\mathcal F_R(M)$ 是 $M$ 的不变量，与计算时选取的生成元无关。
	\end{enumerate}
	
	\item 本题中所有模都假设有限生成。
	\begin{enumerate}[label=(\alph*)]
		\item 若 $M$ 可由 $n$ 个元素生成，证明
		\[
		\operatorname{Ann}(M)^n\subseteq \mathcal F_R(M)\subseteq\operatorname{Ann}(M).
		\]
		[若 $A$ 是 $n\times n$ 关系矩阵，令 $X$ 为以 $m_1,
		\ldots,m_n$ 为条目的列矩阵，则 $AX=0$。左乘 $A$ 的伴随矩阵，推出 $\det A$ 零化 $M$。]
		\item 若 $M=M_1\times M_2$，证明
		\[
		\mathcal F_R(M)=\mathcal F_R(M_1)\mathcal F_R(M_2).
		\]
		\item 若
		\[
		M=(R/I_1)\times\cdots\times(R/I_n),
		\]
		证明 $\mathcal F_R(M)=I_1I_2\cdots I_n$。
		\item 若 $R=\mathbb Z$ 且 $M$ 是有限生成阿贝尔群，证明：若 $M$ 无限，则 $\mathcal F_{\mathbb Z}(M)=0$；若 $M$ 有限，则 $\mathcal F_{\mathbb Z}(M)=|M|\mathbb Z$。
		\item 若 $I$ 是 $R$ 的理想，证明 $\mathcal F_R(M)$ 在商环 $R/I$ 中的像等于 $\mathcal F_{R/I}(M/IM)$。
		\item 证明 $\mathcal F_R(M/IM)\subseteq(\mathcal F_R(M),I)\subseteq R$。
		\item 若 $\varphi:M\to M'$ 是满 $R$-模同态，证明 $\mathcal F_R(M)\subseteq\mathcal F_R(M')$。
		\item 若
		\[
		0\longrightarrow L\longrightarrow M\longrightarrow N\longrightarrow0
		\]
		是短正合列，证明
		\[
		\mathcal F_R(L)\mathcal F_R(N)\subseteq\mathcal F_R(M).
		\]
		\item 设 $R=k[x,y,z]$。令 $M=R/(x,y^2,yz,z^2)$，令 $L=(x,y,z)/(x,y^2,yz,z^2)$ 为 $M$ 的子模。证明 $\mathcal F_R(M)=(x,y^2,yz,z^2)$，而 $\mathcal F_R(L)=(x,y,z)^2$。（这说明一般地，子模 $L$ 的 Fitting 理想不必包含 $M$ 的 Fitting 理想。）
	\end{enumerate}
	
	\item 设 $M$ 是 $R$-模，且 $\varphi:R^n\to M$ 是满 $R$-模同态（即 $M$ 可由 $n$ 个元素生成）。令 $L=\ker\varphi$。证明由包含 $L\subseteq R^n$ 诱导的 $R$-模同态
	\[
	\bigwedge^n L\longrightarrow \bigwedge^n(R^n)\cong R
	\]
	的像为 Fitting 理想 $\mathcal F_R(M)$。
	
	\item 设 $R,S$ 是交换环，$\varphi:R\to S$ 是环同态，$M$ 是有限生成 $R$-模，且 $M'=S\otimes_R M$ 是把标量从 $R$ 扩张到 $S$ 得到的 $S$-模。证明 $M'$ 在 $S$ 上的 Fitting 理想 $\mathcal F_S(M')$ 是 $M$ 在 $R$ 上的 Fitting 理想 $\mathcal F_R(M)$ 到 $S$ 的扩张。
	
	\item 设 $\{g_1,
	\ldots,g_m\}$ 是 $k[x_1,
	\ldots,x_n]$ 中理想 $I$ 的 Gröbner 基。证明：不能被任何 $LT(g_i)$ 整除的单项式 $m$ 给出商空间
	\[
	k[x_1,
	\ldots,x_n]/I
	\]
	的 $k$-向量空间基。
	
	\item 令
	\[
	I=(x^3y-x^2+1,
	x^2y^2-y^3-1)
	\]
	如命题 9.26 后的例 1。
	\begin{enumerate}[label=(\alph*)]
		\item 用前一题证明 $\{1,y,y^2,y^3\}$ 是 $k$-向量空间 $k[x,y]/I$ 的一组基。
		\item 计算 (a) 中基向量的 $4\times4$ 乘法表。
	\end{enumerate}
	
	\item 设 $K[x_1,
	\ldots,x_n]$ 是域 $K$ 上的 $n$ 元多项式环，$k$ 是 $K$ 的子域。若 $I$ 是 $k[x_1,
	\ldots,x_n]$ 的理想，令 $I'$ 为 $I$ 在 $K[x_1,
	\ldots,x_n]$ 中生成的理想。
	\begin{enumerate}[label=(\alph*)]
		\item 若 $G$ 是 $I$ 关于某个单项式序的 Gröbner 基，证明 $G$ 也是 $I'$ 关于同一单项式序的 Gröbner 基。[使用 Buchberger 判别准则。]
		\item 证明 $k[x_1,
		\ldots,x_n]/I$ 作为 $k$-向量空间的维数等于 $K[x_1,
		\ldots,x_n]/I'$ 作为 $K$-向量空间的维数。[一种方法：使用 (a) 与练习 41。]
		\item 证明 $I=k[x_1,
		\ldots,x_n]$ 当且仅当 $I'=K[x_1,
		\ldots,x_n]$。
	\end{enumerate}
	
	\item 令
	\[
	V=Z(x^3-x^2z-y^2z),\qquad W=Z(x^2+y^2-z^2)
	\]
	在 $\mathbb C^3$ 中。则 $I(V)=(x^3-x^2z-y^2z)$，$I(W)=(x^2+y^2-z^2)$（参见第 15.3 节练习 23）。证明
	\[
	\varphi((a,b,c))=(a^2c-b^2c,
	2abc,
	-a^3)
	\]
	定义了从 $V$ 到 $W$ 的态射。
	
	\item 令 $V=Z(x^3+y^3+7z^3)\subseteq\mathbb C^3$。则 $I(V)=(x^3+y^3+7z^3)$。
	\begin{enumerate}[label=(\alph*)]
		\item 证明
		\[
		\varphi(x)=x(x^3-8),
		\quad
		\varphi(y)=-y(x^3+1),
		\quad
		\varphi(z)=z(x^3-1)
		\]
		定义了从 $k[V]$ 到自身的 $\mathbb C$-代数同态。
		\item 令 $\varphi:V\to V$ 为对应的态射。注意 $(-2,1,1)\in V$，计算 $\varphi((-2,1,1))\in V$。
		\item 证明 $V$ 上存在无穷多个点 $(a,b,c)$，其中 $a,b,c\in\mathbb Z$ 且 $a,b,c$ 的最大公因数为 $1$。
	\end{enumerate}
	
	\item 令
	\[
	V=Z(xz+y^2+z^2,
	xy-xz+yz-2z^2)\subseteq\mathbb C^3,
	\]
	以及
	\[
	W=Z(u^3-uv^2+v^3)\subseteq\mathbb C^2
	\]
	如推论 9 后例 2。证明映射
	\[
	\varphi((a,b))=(-2a^2+ab,
	ab-b^2,
	a^2-ab)
	\]
	定义了从 $W$ 到 $V$ 的态射。证明相应的 $\mathbb C$-代数同态 $k[V]\to k[W]$ 的核由 $x^2-3y^2+yz$ 生成。
	
	\item 定义 $\Phi:\mathbb Q[u,v,w]\to\mathbb Q[x,y]$ 为
	\[
	\Phi(u)=x^2+y,
	\quad
	\Phi(v)=x+y^2,
	\quad
	\Phi(w)=x-y.
	\]
	证明 $x$ 与 $y$ 都不在 $\Phi$ 的像中。证明
	\[
	f=2x^3-4xy-2y^3-4y
	\]
	在 $\Phi$ 的像中，并求出一个映到 $f$ 的 $\mathbb Q[u,v,w]$ 中的多项式。证明 $\ker\Phi$ 由
	\[
	u^2-2uv-2uw^2+4uw+v^2-2vw^2-4vw+w^4+3w^2
	\]
	生成。
	
	\item 设 $\alpha$ 是不可约多项式 $p(x)\in k[x]$ 的根，且
	\[
	\beta={f(\alpha)\over g(\alpha)}
	\]
	其中 $f(x),g(x)\in k[x]$ 且 $g(\alpha)\ne0$。
	\begin{enumerate}[label=(\alph*)]
		\item 证明存在 $a,b\in k[x]$ 使得 $ag+bp=1$，并证明 $\beta=h(\alpha)$，其中 $h=af$。
		\item 证明 $k[x,y]$ 中理想 $(p,y-h)$ 与 $(p,gy-f)$ 相等。
		\item 推出 $\beta$ 的最小多项式是 $G\cap k[y]$ 中的首一多项式，其中 $G$ 是理想 $(p,gy-f)$ 关于字典序 $x>y$ 的约化 Gröbner 基。
		\item 求
		\[
		{3-\sqrt2+\sqrt[3]{4}\over 1+3\sqrt2-3\sqrt[3]{4}}
		\]
		在 $\mathbb Q$ 上的最小多项式。
	\end{enumerate}
\end{enumerate}

\section{根式与仿射簇}

由于多项式 $f$ 的零点与它的各个幂 $f^2,f^3,\ldots$ 的零点相同，一般说来，在环 $k[x_1,x_2,\ldots,x_n]$ 中会有许多不同的理想，其零点轨迹定义同一个仿射 $n$-空间中的代数集 $V$。这引出理想的根式的概念；这个概念可以在任意交换环中定义。

\noindent\textbf{定义.} 设 $I$ 是交换环 $R$ 中的理想。
\begin{enumerate}[label=(\arabic*)]
	\item $I$ 的\emph{根式}，记为 $\operatorname{rad}I$，是 $R$ 中某个幂属于 $I$ 的元素全体，即
	\[
	\operatorname{rad}I=\{a\in R\mid a^k\in I\text{ 对某个 }k\geq 1\}.
	\]
	\item 零理想的根式称为 $R$ 的\emph{幂零根基}（nilradical）。
	\item 若 $I=\operatorname{rad}I$，则称 $I$ 为\emph{根式理想}。
\end{enumerate}
注意，$a\in R$ 属于 $R$ 的幂零根基，当且仅当 $a$ 的某个幂等于 $0$；所以 $R$ 的幂零根基正是 $R$ 中所有幂零元素的集合。

\noindent\textbf{命题 11.} 设 $I$ 是交换环 $R$ 中的理想。则 $\operatorname{rad}I$ 是包含 $I$ 的理想，并且 $(\operatorname{rad}I)/I$ 是 $R/I$ 的幂零根基。特别地，$R/I$ 没有非零幂零元素，当且仅当 $I=\operatorname{rad}I$ 是根式理想。

\begin{proof}
	显然 $I\subseteq \operatorname{rad}I$。按定义，商环 $R/I$ 的幂零根基由商环中某个幂为 $0$ 的元素组成。由环的格同构定理，这个元素集合对应于 $R$ 中某个幂属于 $I$ 的元素，也就是 $\operatorname{rad}I$。因此只需证明任意交换环 $R$ 的幂零根基 $N$ 是理想。
	
	显然 $0\in N$。若 $a\in N$ 且 $r\in R$，则 $a^n=0$ 对某个 $n\geq1$ 成立；由于 $R$ 交换，有
	\[
	(ra)^n=r^na^n=0,
	\]
	所以 $ra\in N$。若 $a,b\in N$，设 $a^n=0$ 且 $b^m=0$。由于二项式定理在交换环中成立（参见第 7.3 节练习 25），
	\[
	(a+b)^{n+m}=\sum_{i=0}^{n+m} r_i a^i b^{n+m-i}
	\]
	其中 $r_i$ 是环中相应的二项式系数。这个和式中的每一项，要么 $i\geq n$，从而 $a^i=0$；要么 $n+m-i\geq m$，从而 $b^{n+m-i}=0$。因此 $(a+b)^{n+m}=0$，于是 $a+b$ 也是幂零元，即 $a+b\in N$。
\end{proof}

\noindent\textbf{命题 12.} 真理想 $I$ 的根式等于所有包含 $I$ 的素理想的交。特别地，幂零根基是 $R$ 中所有素理想的交。

\begin{proof}
	传递到 $R/I$，由命题 11 可知，只需证明 $I=0$ 的情形。在这个情形中，命题断言 $R$ 的幂零根基 $N$ 是 $R$ 中所有素理想的交。令 $N'$ 表示 $R$ 中所有素理想的交。
	
	设 $a$ 是 $R$ 中任意幂零元素，并令 $P$ 是任意素理想。由于 $a^k=0$ 对某个 $k$ 成立，存在最小正整数 $n$ 使得 $a^n\in P$。于是 $a^{n-1}a\in P$；因为 $P$ 是素理想，或者 $a^{n-1}\in P$，或者 $a\in P$。前者与 $n$ 的极小性矛盾，故 $a\in P$。由于 $P$ 任意，得到 $a\in N'$，所以 $N\subseteq N'$。
	
	反向包含通过证明若 $a\notin N$ 则 $a\notin N'$ 来完成。若 $a$ 不是幂零元，令 $\mathcal S$ 为所有不含 $a$ 的任何正幂的真理想组成的族。该族非空，因为 $0\in\mathcal S$。并且如果一条链 $I_1\subseteq I_2\subseteq\cdots$ 中没有任何理想含有 $a$ 的正幂，则这些理想的并也不含有 $a$ 的正幂，故 $\mathcal S$ 中的链有上界。由 Zorn 引理，$\mathcal S$ 有极大元 $P$。我们证明 $P$ 实际上是素理想。若 $x,y\notin P$ 而 $xy\in P$，则由 $P$ 的极大性，存在正整数 $n,m$ 使得
	\[
	a^n\in (x)+P,
	\qquad
	a^m\in (y)+P.
	\]
	于是
	\[
	a^{n+m}\in (xy)+P=P,
	\]
	这与 $P\in\mathcal S$ 矛盾。因此 $P$ 是不含 $a$ 的素理想，从而 $a\notin N'$。这证明了 $N'\subseteq N$，命题得证。
\end{proof}

注意，在 Noether 环中，前一证明中对 Zorn 引理的使用可以用定理 2 避开。

\noindent\textbf{推论 13.} 素理想（从而极大理想）都是根式理想。

\begin{proof}
	若 $P$ 是素理想，则 $P$ 显然是所有包含 $P$ 的素理想的交，所以由命题 12，$P=\operatorname{rad}P$。
\end{proof}

\paragraph{例}
\begin{enumerate}[label=(\arabic*)]
	\item 在整数环 $\mathbb Z$ 中，理想 $(a)$ 是根式理想，当且仅当 $a$ 是无平方因子的或 $a=0$。更一般地，若正整数 $a$ 的素因数分解为
	\[
	a=p_1^{\alpha_1}p_2^{\alpha_2}\cdots p_r^{\alpha_r},\qquad \alpha_i\geq1,
	\]
	则
	\[
	\operatorname{rad}(a)=(p_1p_2\cdots p_r).
	\]
	例如，$\operatorname{rad}(180)=(30)$。注意 $(p_1),(p_2),\ldots,(p_r)$ 正是包含理想 $(a)$ 的素理想，而且它们的交是理想 $(p_1p_2\cdots p_r)$。更一般地，在任意 U.F.D. $R$ 中，若 $a=p_1^{\alpha_1}p_2^{\alpha_2}\cdots p_r^{\alpha_r}$ 是 $a$ 分解为互不相伴不可约元的唯一分解，则
	\[
	\operatorname{rad}(a)=(p_1p_2\cdots p_r).
	\]
	\item $k[x,y]$ 中的理想 $(x^3-y^2)$ 是素理想（第 9.1 节练习 14），因而是根式理想。
	\item 若 $l_1,\ldots,l_m$ 是 $k[x_1,x_2,\ldots,x_n]$ 中的线性多项式，则 $I=(l_1,\ldots,l_m)$ 或者等于 $k[x_1,x_2,\ldots,x_n]$，或者是素理想；因此 $I$ 是根式理想。
\end{enumerate}

\noindent\textbf{命题 14.} 若 $R$ 是 Noether 环，则对任意理想 $I$，$\operatorname{rad}I$ 的某个正幂包含于 $I$。特别地，Noether 环的幂零根基 $N$ 是幂零理想：存在 $k\geq1$ 使 $N^k=0$。

\begin{proof}
	对任意理想 $I$，由于 $R$ 是 Noether 环，理想 $\operatorname{rad}I$ 有有限生成元。设 $a_1,\ldots,a_m$ 生成 $\operatorname{rad}I$。按根式的定义，对每个 $i$，存在正整数 $k_i$ 使得 $a_i^{k_i}\in I$。令 $k=\max\{k_1,\ldots,k_m\}$。理想 $(\operatorname{rad}I)^{km}$ 由形如
	\[
	a_1^{d_1}a_2^{d_2}\cdots a_m^{d_m},\qquad d_1+\cdots+d_m=km
	\]
	的元素生成；这样的每个元素中至少有一个因子满足 $d_i\geq k$。于是 $a_i^{d_i}\in I$，所以 $(\operatorname{rad}I)^{km}$ 的每个生成元都属于 $I$，从而 $(\operatorname{rad}I)^{km}\subseteq I$。
\end{proof}

\subsection*{Zariski 拓扑}

在上一节中我们看到，如果只考虑 $k[\mathbb A^n]$ 中形如 $I(V)$ 的理想，也就是与某个代数集 $V$ 相联系的理想，则映射 $Z$（从这样的理想到代数集）与 $I$（从代数集到理想）互为逆映射：$Z(I(V))=V$ 且 $I(Z(I))=I$。环 $k[\mathbb A^n]/I(V)$ 的元素给出 $V$ 上的 $k$-值函数；由于 $k$ 没有幂零元，非零函数的幂仍为非零函数。换言之，环 $k[\mathbb A^n]/I(V)$ 没有幂零元素，所以由命题 11，理想 $I(V)$ 总是根式理想。

对任意域 $k$，并不总是每个根式理想都等于某个代数集的理想，即不总是形如 $I(V)$。例如，$\mathbb R[x]$ 中的理想 $(x^2+1)$ 是极大理想，因而由推论 13 是根式理想；但它不是任何代数集的理想。如果它是，那么 $x^2+1$ 必须在这个代数集上恒为零，而 $x^2+1$ 在 $\mathbb R$ 中没有零点。任意非代数闭域 $k$ 上都有类似构造：存在 $k[x]$ 中次数至少为 $2$ 的不可约多项式 $p(x)$，它生成 $k[x]$ 中没有 $k$-有理零点的极大（因而根式）理想 $(p(x))$。也许令人惊讶的是，一元多项式没有零点这一现象正是根式理想（即使在任意多个变量中）不成为某个 $I(V)$ 的唯一障碍。下一个定理表明了这一点；它给出“几何”与“代数”之间的基本联系，并说明在代数闭域（如 $\mathbb C$）上，每个根式理想都形如 $I(V)$。在这样的域上，“几何定义的”理想 $I=I(V)$ 与根式理想相同，而后者是理想 $I$ 的一个“纯代数”性质，即 $I=\operatorname{rad}I$。

\noindent\textbf{定理（Hilbert 零点定理）.} 设 $E$ 是代数闭域。则对 $E[x_1,x_2,\ldots,x_n]$ 的每个理想 $I$，都有
\[
I(Z(I))=\operatorname{rad}I.
\]
此外，在对应
\[
\{\text{仿射代数集}\}\ \underset{Z}{\stackrel{I}{\rightleftarrows}}\ \{\text{根式理想}\}
\]
中，映射 $Z$ 与 $I$ 是互为逆的双射。

\begin{proof}
	这将在下一节证明（参见定理 32）。
\end{proof}

\paragraph{例}
零点定理中的映射 $I$ 与 $Z$ 在任意域 $k$ 上都有定义；如上所述，如果 $k$ 不是代数闭域，它们不一定是双射。不过，对任意域 $k$，映射 $Z$ 总是满射，而映射 $I$ 总是单射（参见练习 9）。

零点定理的一个特殊后果是：对任意真理想 $I$，都有 $Z(I)\ne\varnothing$，因为 $\operatorname{rad}I\ne k[\mathbb A^n]$。因此在代数闭域上，真理想中所有多项式总至少有一个公共零点（德语中的 ``Nullstellen''）。

现在说明仿射代数集在仿射 $n$-空间上定义一个拓扑。回忆拓扑空间是一个集合 $X$ 连同 $X$ 的若干子集族 $\mathcal T$，这些子集称为 $X$ 中的闭集，并满足如下公理：
\begin{enumerate}[label=(\roman*)]
	\item 任意交闭集仍为闭集；
	\item 有限并闭集仍为闭集；
	\item 空集与整个空间是闭集。
\end{enumerate}
集合 $U\subseteq X$ 称为开集，若它的补集 $X-U$ 是闭集。拓扑空间的公理也常等价地用开集族来表述。

拓扑空间有许多例子，也有大量关于拓扑的书籍。一个固定集合 $X$ 上可以有许多不同的拓扑；在不同拓扑中闭集族之间不必相关。在任意集合 $X$ 上总至少有两个拓扑：一个是离散拓扑，其中每个子集都是闭集；另一个是平凡拓扑，其中唯一的闭集就是公理所要求的 $\varnothing$ 与 $X$。

令 $X=\mathbb A^n$ 为任意域 $k$ 上的仿射 $n$-空间。由上一节代数集的性质 (3)、(4)、(5)，$\mathbb A^n$ 中所有仿射代数集组成的集合满足上述三个拓扑公理。因此可把这些集合取作 $\mathbb A^n$ 上某个拓扑的闭集。

\noindent\textbf{定义.} 任意域 $k$ 上仿射 $n$-空间的 \emph{Zariski 拓扑}，是以 $\mathbb A^n$ 中的仿射代数集为闭集的拓扑。

Zariski 拓扑相当“粗”，意思是闭集（或开集）相对较少。例如，$\mathbb A^1$ 的 Zariski 拓扑中，闭集只有 $\varnothing$、$k$ 以及有限集（参见第 15.1 节练习 14），所以非空开集就是有限集的补。若 $k$ 是无限域，则在 $\mathbb A^1$ 的 Zariski 拓扑中任意两个非空开集都有非空交。用点集拓扑的语言说，Zariski 拓扑总是 $T_1$ 的（点是闭集），但当 $k$ 是无限域时，Zariski 拓扑绝不是 $T_2$（Hausdorff）的，也就是说两个不同点不能分别落在两个不交开集中（参见练习）。例如当 $k=\mathbb R$ 时，非空 Zariski 开集只是实直线删去有限多个点，任意两个这样的集合都有无穷多个公共点。还要注意，$\mathbb R$ 中 Zariski 开集（闭集）在通常 Euclidean 拓扑下也是开集（闭集）。反之不成立；例如区间 $[0,1]$ 在 Euclidean 拓扑中闭，但在 Zariski 拓扑中不闭。从这个意义上说，$\mathbb R$ 上的 Euclidean 拓扑要细得多；Euclidean 开集（闭集）族真包含 Zariski 开集（闭集）族。

$\mathbb A^n$ 上的 Zariski 拓扑被定义成使 $\mathbb A^n$ 的仿射代数子集成为闭集。换言之，这个拓扑由 $\mathbb A^n$ 的坐标环中的理想的零点集定义。类似的定义也可用于定义 $\mathbb A^n$ 中任一代数集 $V$ 上的 Zariski 拓扑。若 $k[V]$ 是 $V$ 的坐标环，则 $k[V]$ 的不同元素定义 $V$ 上不同的 $k$-值函数，并且可自然地定义
\[
Z:\{k[V]\text{ 的理想}\}\longrightarrow \{V\text{ 的代数子集}\},
\qquad
I:\{V\text{ 的子集}\}\longrightarrow \{k[V]\text{ 的理想}\},
\]
就像 $V=\mathbb A^n$ 的情形一样。例如，若 $J$ 是 $k[V]$ 的理想，则 $Z(J)$ 是 $V$ 中所有使 $J$ 中函数同时为零的点的集合。由此得到的 $V$ 中的零点集很容易验证满足拓扑空间的三个公理，从而定义 $V$ 上的 Zariski 拓扑，其中闭集是代数子集 $Z(J)$。

由格同构定理，$k[V]$ 的理想就是 $k[x_1,\ldots,x_n]$ 中含有 $I(V)$ 的理想模去 $I(V)$ 所得。若 $J$ 在 $k[x_1,\ldots,x_n]$ 中的完全逆像为 $\widetilde J$，则 $\widetilde J$ 在 $\mathbb A^n$ 中的轨迹与 $J$ 在 $V$ 中的轨迹相同。因此这个在 $V$ 上定义 Zariski 拓扑的方式，正是包含 $V\subseteq\mathbb A^n$ 所诱导的子空间拓扑。（回忆在拓扑空间 $X$ 中，子空间 $Y$ 的闭集定义为 $C\cap Y$，其中 $C$ 是 $X$ 的闭集。）上述定义的优点在于它只用 $V$ 的坐标环 $k[V]$ 内在地定义；而 $k[V]$ 的同构类型不依赖于包含 $V$ 的仿射空间，所以 $V$ 上的 Zariski 拓扑也只依赖于 $V$，不依赖于 $V$ 嵌入的环境仿射空间。

若 $V,W$ 是两个仿射代数集，则由于态射 $\varphi:V\to W$ 由多项式函数定义，很容易看出 $\varphi$ 关于 $V$ 与 $W$ 上的 Zariski 拓扑是连续的（参见第 15.1 节练习 27，它说明态射下 Zariski 闭集的逆像是 Zariski 闭集）。事实上，Zariski 拓扑是使点闭且使多项式映射连续的最粗拓扑。不过，也存在关于 Zariski 拓扑连续但不是态射的映射（参见练习 17）。

\noindent\textbf{定义.} 对任意子集 $A\subseteq\mathbb A^n$，$A$ 的 \emph{Zariski 闭包}是包含 $A$ 的最小代数集。若 $A\subseteq V$ 且 $V$ 是代数集，则称 $A$ 在 $V$ 中 \emph{Zariski 稠密}，若 $A$ 的 Zariski 闭包为 $V$。

例如，当 $k=\mathbb R$ 时，由第 15.1 节练习 14，$\mathbb A^1$ 中的代数集是 $\varnothing$、$\mathbb R$ 以及 $\mathbb R$ 的有限子集。因此任意无限实数集的 Zariski 闭包都是整个 $\mathbb A^1$，并且它在 $\mathbb A^1$ 中 Zariski 稠密。

\noindent\textbf{命题 15.} 子集 $A\subseteq\mathbb A^n$ 的 Zariski 闭包是 $Z(I(A))$。

\begin{proof}
	显然 $A\subseteq Z(I(A))$。设 $V$ 是任意包含 $A$ 的代数集：$A\subseteq V$。于是 $I(V)\subseteq I(A)$，从而
	\[
	Z(I(A))\subseteq Z(I(V))=V.
	\]
	因此 $Z(I(A))$ 是包含 $A$ 的最小代数集。
\end{proof}

若 $\varphi:V\to W$ 是代数集的态射，像 $\varphi(V)$ 不必是 $W$ 中的代数子集，也就是说不必在 $W$ 中 Zariski 闭。例如，双曲线 $V=Z(xy-1)\subseteq\mathbb R^2$ 到 $x$-轴的投影像是 $\mathbb R-\{0\}$，如上所述它不是仿射代数集。

下一个结果说明，态射像的 Zariski 闭包由相应的 $k$-代数同态的核决定。

\noindent\textbf{命题 16.} 设 $\varphi:V\to W$ 是代数集的态射，$\overline\varphi:k[W]\to k[V]$ 是相应的坐标环 $k$-代数同态。则：
\begin{enumerate}[label=(\arabic*)]
	\item $\ker\overline\varphi=I(\varphi(V))$；
	\item $\varphi(V)$ 的 Zariski 闭包是 $W$ 中 $\ker\overline\varphi$ 的零点集。特别地，$\overline\varphi$ 单射当且仅当 $\varphi(V)$ 在 $W$ 中 Zariski 稠密。
\end{enumerate}

\begin{proof}
	由于 $\overline\varphi(f)=f\circ\varphi$，所以 $\overline\varphi(f)=0$ 当且仅当 $(f\circ\varphi)(P)=0$ 对所有 $P\in V$ 成立，也就是 $f(Q)=0$ 对所有 $Q=\varphi(P)\in\varphi(V)$ 成立。这正是 $f\in I(\varphi(V))$，故 (1) 成立。由上一命题，$\varphi(V)$ 的 Zariski 闭包是 $I(\varphi(V))$ 的零点集，故 (2) 的第一句话成立。
	
	若 $\overline\varphi$ 单射，则 $\varphi(V)$ 的 Zariski 闭包为 $Z(0)=W$，所以 $\varphi(V)$ Zariski 稠密。反过来，若 $\varphi(V)$ 在 $W$ 中 Zariski 稠密，即 $Z(I(\varphi(V)))=W$，则
	\[
	I(\varphi(V))=I(Z(I(\varphi(V))))=I(W)=0,
	\]
	所以 $\ker\overline\varphi=0$。
\end{proof}

命题 16 表明，定义态射 $\varphi$ 的像的 Zariski 闭包的多项式理想，正是定理 6 中相应的 $k$-代数同态 $\overline\varphi$ 的核。命题 8(1) 允许我们用 Gröbner 基计算这个核。

\paragraph{例（隐式化）}
态射 $\varphi:\mathbb A^n\to\mathbb A^m$ 就是映射
\[
\varphi((a_1,a_2,\ldots,a_n))=(\varphi_1(a_1,a_2,\ldots,a_n),\ldots,\varphi_m(a_1,a_2,\ldots,a_n)),
\]
其中 $\varphi_i$ 是多项式。若 $k$ 是无限域，则 $I(\mathbb A^m)=I(\mathbb A^n)=0$，故可写
\[
k[\mathbb A^m]=k[y_1,\ldots,y_m],\qquad k[\mathbb A^n]=k[x_1,\ldots,x_n].
\]
对应于 $\varphi$ 的 $k$-代数同态 $\overline\varphi:k[\mathbb A^m]\to k[\mathbb A^n]$ 由
\[
y_i\longmapsto \varphi_i(x_1,\ldots,x_n)
\]
定义。像 $\varphi(\mathbb A^n)$ 由所有点 $(b_1,\ldots,b_m)$ 组成，其中
\[
b_i=\varphi_i(a_1,a_2,\ldots,a_n),\qquad a_i\in k.
\]
这是由函数 $\varphi_1,\ldots,\varphi_m$（以 $a_i$ 为参数）参数化的点集。一般这样的参数化点集不是代数集。寻找包含这些点的最小代数集的方程称为\emph{隐式化}，因为这相当于寻找 $b_i$ 所满足的（“最小”的）方程集合，即“隐式”的代数关系。

由命题 16，这个代数集就是 $\varphi(\mathbb A^n)$ 的 Zariski 闭包，是 $\ker\overline\varphi$ 的零点集。由命题 8，该核由 $A\cap k[y_1,\ldots,y_m]$ 给出，其中 $A$ 是 $k[x_1,\ldots,x_n,y_1,\ldots,y_m]$ 中由
\[
y_1-\varphi_1,\ldots,y_m-\varphi_m
\]
生成的理想。如果按字典序
\[
x_1>\cdots>x_n>y_1>\cdots>y_m
\]
计算 $A$ 的约化 Gröbner 基 $G$，则 $G$ 中落在 $k[y_1,\ldots,y_m]$ 里的多项式生成 $\ker\overline\varphi$。这些多项式的零点集定义 $\varphi(\mathbb A^n)$ 的 Zariski 闭包，从而给出隐式化。

作为一个具体例子，考虑 $\mathbb R^2$ 中的点集
\[
A=\{(a^2,a^3)\mid a\in\mathbb R\}.
\]
用 $x,y$ 作 $\mathbb R^2$ 的坐标，用 $t$ 作 $\mathbb R^1$ 的坐标，则 $\mathbb R[x,y,t]$ 中相应的理想为 $(x-t^2,y-t^3)$。对序 $t>x>y$，这个理想的约化 Gröbner 基中唯一落在 $\mathbb R[x,y]$ 中的元素是 $x^3-y^2$，所以 $Z(x^3-y^2)$ 是 $\mathbb R^2$ 中包含 $A$ 的最小代数集。

\paragraph{例（代数集的投影）}
设 $V\subseteq\mathbb A^n$ 是代数集，且 $m<n$。令 $\pi:V\to\mathbb A^m$ 为投影到前 $m$ 个坐标的态射：
\[
\pi((a_1,a_2,\ldots,a_n))=(a_1,a_2,\ldots,a_m).
\]
若 $k[V]$ 中用坐标 $x_1,
\ldots,x_n$，$k[\mathbb A^m]$ 中用坐标 $y_1,
\ldots,y_m$，则与 $\pi$ 对应的 $k$-代数同态为
\[
\overline\pi:k[y_1,
\ldots,y_m]\longrightarrow k[x_1,
\ldots,x_n]/I(V),\qquad y_i\longmapsto x_i.
\]
设 $V=Z(I)$ 且 $I=(f_1,
\ldots,f_s)$。$\pi(V)$ 的 Zariski 闭包是 $\ker\overline\pi=A\cap k[y_1,
\ldots,y_m]$ 的零点集，其中 $A$ 是 $k[x_1,
\ldots,x_n,y_1,
\ldots,y_m]$ 中由 $y_1-x_1,
\ldots,y_m-x_m$ 以及 $I(V)$ 的一组生成元生成的理想。按字典序
\[
x_1>\cdots>x_n>y_1>\cdots>y_m
\]
计算 $A$ 的约化 Gröbner 基，则其中只含 $y_1,
\ldots,y_m$ 的多项式生成 $\pi(V)$ 的 Zariski 闭包的理想。

若 $k$ 是代数闭域，由零点定理可进一步把 $I(V)=\operatorname{rad}I$ 用在这里。此时可直接用 $I$ 的生成元 $f_1,
\ldots,f_s$ 代替 $I(V)$ 的生成元，而得到相同的零点集（参见练习 46）。

作为具体例子，考虑 $\mathbb C^3$ 中
\[
V=Z(xy-z^2,
xz-y,
x^2-z)
\]
投影到前两个坐标。用 $u,v$ 作 $\mathbb C^2$ 中坐标，对序 $x>y>z>u>v$，理想
\[
(u-x,
v-y,
xy-z^2,
xz-y,
x^2-z)
\]
的约化 Gröbner 基中唯一落在 $\mathbb C[u,v]$ 中的多项式是 $u^3-v$。因此包含 $\pi(V)$ 的最小代数集是三次曲线 $v=u^3$。

\subsection*{仿射簇}

现在考虑一个代数集是否可以分解为较小代数集的问题，以及坐标环语言中的相应代数表述。

\noindent\textbf{定义.} 非空仿射代数集 $V$ 称为\emph{不可约的}，若它不能写成
\[
V=V_1\cup V_2
\]
其中 $V_1,V_2$ 是 $V$ 中的真代数集。不可约仿射代数集称为\emph{仿射簇}。

等价地，一个代数集（也就是 Zariski 拓扑中的闭集）不可约，当且仅当它不能写成两个真闭子集的并。

\noindent\textbf{命题 17.}
\begin{enumerate}[label=(\arabic*)]
	\item 仿射代数集 $V$ 不可约，当且仅当 $I(V)$ 是素理想。
	\item 每个非空仿射代数集 $V$ 都可以唯一地写成
	\[
	V=V_1\cup V_2\cup\cdots\cup V_q,
	\]
	其中每个 $V_i$ 不可约，并且 $V_i\not\subseteq V_j$ 对所有 $j\ne i$ 成立（也就是说，这个分解是“极小的”或“不可冗余的”）。
\end{enumerate}

\begin{proof}
	令 $I=I(V)$。先设 $V$ 可约，$V=V_1\cup V_2$，其中 $V_1,V_2$ 是真闭子集。由于 $V_1\ne V$，存在函数 $f_1$ 在 $V_1$ 上为零但不在 $V$ 上恒为零，即 $f_1\in I(V_1)-I$。同理，存在 $f_2\in I(V_2)-I$。于是 $f_1f_2$ 在 $V_1\cup V_2=V$ 上为零，所以 $f_1f_2\in I$，这说明 $I$ 不是素理想。
	
	反过来，若 $I$ 不是素理想，则存在 $f_1,f_2\in k[\mathbb A^n]$ 使 $f_1f_2\in I$，但 $f_1\notin I$ 且 $f_2\notin I$。令
	\[
	V_1=Z(f_1)\cap V,
	\qquad
	V_2=Z(f_2)\cap V.
	\]
	闭集的交仍为闭集，所以 $V_1,V_2$ 是代数集。由于 $f_1,f_2$ 都不在 $V$ 上恒为零，$V_1,V_2$ 都是 $V$ 的真子集。又由于 $f_1f_2\in I$，有
	\[
	V\subseteq Z(f_1f_2)=Z(f_1)\cup Z(f_2),
	\]
	所以 $V=V_1\cup V_2$，即 $V$ 可约。这证明了 (1)。
	
	为证明 (2)，令 $\mathcal S$ 为不能写成有限个不可约代数集并的非空代数集族，反设 $\mathcal S\ne\varnothing$。考虑相应理想集合 $\{I(V)\mid V\in\mathcal S\}$。由于 $k[\mathbb A^n]$ 是 Noether 环，由定理 2 它有极大元 $I_0$。于是 $V_0=Z(I_0)$ 是 $\mathcal S$ 中的极小元。由于 $V_0\in\mathcal S$，按定义它不可能不可约。另一方面，若 $V_0=V_1\cup V_2$，其中 $V_1,V_2$ 是 $V_0$ 的真闭子集，则由 $V_0$ 的极小性，$V_1,V_2$ 都可写成有限个不可约代数集的并，从而 $V_0$ 也可写成有限个不可约代数集的并，矛盾。因此 $\mathcal S=\varnothing$，每个仿射代数集都有仿射簇分解。
	
	为证唯一性，设 $V$ 有两个仿射簇分解（均已去掉冗余项）：
	\[
	V=V_1\cup\cdots\cup V_r=U_1\cup\cdots\cup U_s.
	\]
	则 $V_1$ 包含于 $U_i$ 的并中。由于每个 $V_1\cap U_i$ 是代数集，得到
	\[
	V_1=(V_1\cap U_1)\cup\cdots\cup(V_1\cap U_s).
	\]
	因为 $V_1$ 不可约，必有 $V_1=V_1\cap U_j$ 对某个 $j$ 成立，即 $V_1\subseteq U_j$。对称地，有 $U_j\subseteq V_{j'}$ 对某个 $j'$ 成立。因此 $V_1\subseteq V_{j'}$，由不可冗余性知 $j'=1$，于是 $V_1=U_j$。对每个 $V_i$ 作同样论证，得到 $r=s$ 且
	\[
	\{V_1,
	\ldots,V_r\}=\{U_1,
	\ldots,U_s\}.
	\]
	证明完成。
\end{proof}

\noindent\textbf{推论 18.} 仿射代数集 $V$ 是簇，当且仅当其坐标环 $k[V]$ 是整环。

\begin{proof}
	这立即来自如下事实：$I(V)$ 是素理想，当且仅当商环
	\[
	k[V]=k[\mathbb A^n]/I(V)
	\]
	是整环（第 7 章命题 13）。
\end{proof}

\noindent\textbf{定义.} 若 $V$ 是簇，则整环 $k[V]$ 的分式域称为 $V$ 上的\emph{有理函数域}，记为 $k(V)$。簇 $V$ 的\emph{维数}，记为 $\dim V$，定义为 $k(V)$ 在 $k$ 上的超越次数。

\paragraph{例}
\begin{enumerate}[label=(\arabic*)]
	\item $\mathbb A^n$ 中的单点是仿射簇，因为对应的理想是 $k[\mathbb A^n]$ 中的极大理想。一个点的坐标环同构于 $k$，这也就是它的有理函数域。单点的维数为 $0$。任意有限集是其单点子集的并，这是它分解为仿射子簇的唯一分解。
	\item $\mathbb R^2$ 中的 $x$-轴不可约，因为它的坐标环是
	\[
	\mathbb R[x,y]/(y)\cong\mathbb R[x],
	\]
	这是整环。同样地，$y$-轴以及更一般的 $\mathbb R^2$ 中的直线也不可约（参见第 15.1 节练习 23）。$\mathbb A^n$ 中的线性集合是仿射簇。$x$-轴的有理函数域是 $\mathbb R[x]$ 的分式域 $\mathbb R(x)$，这正是 $\mathbb R(x)$ 被称为有理函数域的原因。$x$-轴（更一般地任意直线）的维数为 $1$。
	\item $\mathbb R^2$ 中 $x$-轴与 $y$-轴的并，即 $Z(xy)$，不是簇：
	\[
	Z(xy)=Z(x)\cup Z(y)
	\]
	是它唯一的子簇分解。相应的坐标环 $\mathbb R[x,y]/(xy)$ 含有零因子。
	\item $\mathbb R^2$ 中双曲线 $xy=1$ 是簇，因为第 15.1 节中看到它的坐标环是整环 $\mathbb R[x,1/x]$。注意，双曲线的两个不相交分支（由 $x>0$ 与 $x<0$ 定义）并不是子簇（也参见练习 12--13）。
	\item 若 $V=Z(l_1,l_2,
	\ldots,l_m)$ 是 $k[x_1,
	\ldots,x_n]$ 中线性多项式 $l_1,
	\ldots,l_m$ 的零点集，且 $V\ne\varnothing$，则 $V$ 是仿射簇（称为线性簇）。注意，判断 $V\ne\varnothing$ 是一个线性代数问题。
\end{enumerate}

本节最后给出若干一般的环论结果；这些结果最初正是受几何中的分解问题所启发。

\subsection*{Noether 环中理想的 primary 分解}

命题 17 的第二部分说明，$k[\mathbb A^n]$ 中形如 $I(V)$ 的任意理想都可以唯一地写成有限个素理想的交；由 Hilbert 零点定理，当 $k$ 是代数闭域时，这尤其适用于所有根式理想。在一大类交换环（包括所有 Noether 环）中，每个理想都有 primary 分解。这是一种类似的分解，但允许出现类似“素数幂”的理想（不过见下面的例子）。这个分解可看作整数 $n\in\mathbb Z$ 分解为素数幂乘积的推广。我们主要考虑 Noether 环的情形。

\noindent\textbf{定义.} 交换环 $R$ 中的真理想 $Q$ 称为\emph{primary 理想}，若每当 $ab\in Q$ 且 $a\notin Q$ 时，就有 $b^n\in Q$ 对某个正整数 $n$ 成立。等价地，若 $ab\in Q$ 且 $a\notin Q$，则 $b\in\operatorname{rad}Q$。

primary 理想的一些基本性质如下。

\noindent\textbf{命题 19.} 设 $R$ 是含 $1$ 的交换环。
\begin{enumerate}[label=(\arabic*)]
	\item 素理想是 primary 理想。
	\item 理想 $Q$ 是 primary，当且仅当 $R/Q$ 中的每个零因子都是幂零元。
	\item 若 $Q$ 是 primary，则 $\operatorname{rad}Q$ 是素理想，并且是包含 $Q$ 的唯一最小素理想。
	\item 若理想 $Q$ 的根式是极大理想，则 $Q$ 是 primary 理想。
	\item 设 $M$ 是极大理想，且 $Q$ 是满足 $M^n\subseteq Q\subseteq M$ 对某个 $n\geq1$ 成立的理想。则 $Q$ 是 primary 理想，且 $\operatorname{rad}Q=M$。
\end{enumerate}

\begin{proof}
	前两个断言直接来自 primary 理想的定义。对 (3)，设 $ab\in\operatorname{rad}Q$。则 $(ab)^m=a^mb^m\in Q$ 对某个 $m$ 成立。由于 $Q$ 是 primary，要么 $a^m\in Q$，从而 $a\in\operatorname{rad}Q$；要么 $(b^m)^n\in Q$ 对某个正整数 $n$ 成立，从而 $b\in\operatorname{rad}Q$。这说明 $\operatorname{rad}Q$ 是素理想；由命题 12，它也就是包含 $Q$ 的最小素理想。
	
	证明 (4) 时传到商环 $R/Q$；由 (2)，只需证明该商环中的每个零因子都是幂零元。于是可化到 $Q=(0)$ 且 $M=\operatorname{rad}(0)$ 是幂零根基同时又是极大理想的情形。由于幂零根基包含于每个素理想中（命题 12），$M$ 是唯一的素理想，从而也是唯一的极大理想。若 $d$ 是零因子，则理想 $(d)$ 是真理想，因而包含在某个极大理想中；这迫使 $d\in M$，所以每个零因子确为幂零元。
	
	最后，若 $M^n\subseteq Q\subseteq M$ 且 $M$ 是极大理想，则 $Q\subseteq M$ 给出 $\operatorname{rad}Q\subseteq\operatorname{rad}M=M$；另一方面 $M^n\subseteq Q$ 给出 $M\subseteq\operatorname{rad}Q$。故 $\operatorname{rad}Q=M$ 是极大理想，$Q$ 由 (4) 是 primary 理想。
\end{proof}

\noindent\textbf{定义.} 若 $Q$ 是 primary 理想，则素理想 $P=\operatorname{rad}Q$ 称为与 $Q$ 相关的素理想；也说 $Q$ 属于 $P$，或称 $Q$ 是 $P$-primary 理想。

容易验证，有限个 $P$-primary 理想的交仍是 $P$-primary 理想（参见练习）。

\paragraph{例}
\begin{enumerate}[label=(\arabic*)]
	\item $\mathbb Z$ 中的 primary 理想是 $0$ 以及 $(p^m)$，其中 $p$ 是素数且 $m\geq1$。
	\item 对任意域 $k$，$k[x,y]$ 中的理想 $(x)$ 是 primary，因为它是素理想。对任意 $n\geq1$，理想 $(x,y)^n$ 是 primary，因为它是极大理想 $(x,y)$ 的幂。
	\item 多项式环 $k[x,y]$ 中的理想 $Q=(x^2,y)$ 是 primary，因为
	\[
	(x,y)^2\subseteq (x^2,y)\subseteq(x,y).
	\]
	类似地，$\mathbb Z[x]$ 中 $Q'=(4,x)$ 是 $(2,x)$-primary 理想。
	\item primary 理想不必是素理想的幂。例如，前一例中的 primary 理想 $Q$ 不是某个素理想的幂。若 $(x^2,y)=P^k$ 对某个素理想 $P$ 与某个 $k\geq1$ 成立，则 $x^2,y\in P^k\subseteq P$，所以 $x,y\in P$。于是 $P=(x,y)$；又因为 $y\notin(x,y)^2$，必须有 $k=1$，从而 $Q=(x,y)$。但 $x\notin(x^2,y)$，矛盾。
	\item 若 $R$ 是 Noether 环，$Q$ 是属于素理想 $P$ 的 primary 理想，则由命题 14，存在 $m\geq1$ 使
	\[
	P^m\subseteq Q\subseteq P.
	\]
	若 $P$ 是极大理想，则命题 19 的最后一句说明反过来也成立。若 $P$ 是非极大的素理想，这不一定成立。例如，考虑 $k[x,y]$ 中的理想 $I=(x^2,xy)$。则
	\[
	(x)^2\subset I\subset (x),
	\]
	且 $(x)$ 是素理想，但 $I$ 不是 primary：$xy\in I$ 且 $x\notin I$，但 $y$ 的任何正幂都不属于 $I$。这个例子也说明，根式为素理想（但不是命题 (4) 中的极大理想）的理想不一定是 primary。
	\item 素理想的幂不必是 primary。例如，考虑商环
	\[
	R=\mathbb R[x,y,z]/(xy-z^2),
	\]
	也就是 $\mathbb R^3$ 中锥面 $z^2=xy$ 的坐标环，并令 $P=(\overline x,\overline z)$ 为 $R$ 中由 $\overline x$ 与 $\overline z$ 生成的理想。这是 $R$ 中的素理想，因为商环
	\[
	R/(\overline x,\overline z)\cong \mathbb R[x,y,z]/(x,z)\cong\mathbb R[y]
	\]
	是整环。然而
	\[
	P^2=(x^2,xz,z^2)=(x^2,xz,xy)=x(x,y,z)
	\]
	不是 primary：$xy=z^2\in P^2$，但 $x\notin P^2$，并且 $y$ 的任何幂都不在 $P^2$ 中。注意 $P^2$ 又给出一个根式为素理想但本身不是 primary 的例子。
	\item 设 $R$ 是 U.F.D.。若 $\pi$ 是 $R$ 中不可约元，则容易看出 $(\pi^n)$ 对 $n=1,2,\ldots$ 是 $(\pi)$-primary 理想。反过来，设 $Q$ 是 $(\pi)$-primary 理想，并令 $n$ 是使 $Q\subseteq(\pi^n)$ 成立的最大整数（这样的整数存在，因为例如某个 $\pi^k\in Q$，所以 $n\leq k$）。若 $q\in Q$ 但 $q\notin(\pi^{n+1})$，则 $q=r\pi^n$ 对某个 $r\in R$ 成立且 $r\notin(\pi)$。由于 $r\notin(\pi)$ 且 $Q$ 是 $(\pi)$-primary，推出 $\pi^n\in Q$。于是 $Q=(\pi^n)$。
\end{enumerate}

在上面的例子中，$k[x,y]$ 中的理想 $(x^2,xy)$ 不是 primary 理想，但可写成 primary 理想的交：
\[
(x^2,xy)=(x)\cap(x,y)^2.
\]

\noindent\textbf{定义.}
\begin{enumerate}[label=(\arabic*)]
	\item 若 $R$ 中理想 $I$ 可以写成有限个 primary 理想的交
	\[
	I=\bigcap_{i=1}^m Q_i,
	\qquad Q_i\text{ 为 primary 理想},
	\]
	则称 $I$ 有一个 \emph{primary 分解}。
	\item 上述 primary 分解称为\emph{极小的}，并且 $Q_i$ 称为 $I$ 的 \emph{primary 分量}，若
	\begin{enumerate}[label=(\alph*)]
		\item 没有一个 primary 理想包含其余 primary 理想的交，即对所有 $i$，$Q_i\not\supseteq\bigcap_{j\ne i}Q_j$；
		\item 相关素理想互不相同：若 $i\ne j$，则 $\operatorname{rad}Q_i\ne\operatorname{rad}Q_j$。
	\end{enumerate}
\end{enumerate}

现在证明在 Noether 环中每个真理想都有极小 primary 分解。这个结果通常称为 Lasker--Noether 分解定理，因为它最初由国际象棋大师 Emanuel Lasker 对多项式环证明，后来由 Emmy Noether 大为简化并推广。

\noindent\textbf{定义.} 交换环 $R$ 中的真理想 $I$ 称为\emph{不可约的}，若它不能非平凡地写成两个其他理想的交；也就是说，若 $I=J\cap K$，其中 $J,K$ 是理想，则必有 $I=J$ 或 $I=K$。

容易看出素理想是不可约的（见第 7.4 节练习 11）。前面例 2 中 $k[x,y]$ 里的理想 $(x,y)^2$ 表明 primary 理想不必不可约，因为它是
\[
(x)+(x,y)^2=(x,y^2)
\]
与
\[
(y)+(x,y)^2=(y,x^2)
\]
的交。不过在 Noether 环中，不可约理想必定是 primary 理想。

\noindent\textbf{命题 20.} 设 $R$ 是 Noether 环。则
\begin{enumerate}[label=(\arabic*)]
	\item 每个不可约理想都是 primary；
	\item $R$ 中每个真理想都是有限个不可约理想的交。
\end{enumerate}

\begin{proof}
	先证 (1)。设 $Q$ 是不可约理想，并设 $ab\in Q$ 且 $b\notin Q$。对固定的 $n$，集合
	\[
	A_n=\{x\in R\mid a^n x\in Q\}
	\]
	容易验证是 $R$ 的理想。显然 $A_1\subseteq A_2\subseteq\cdots$；由于 $R$ 是 Noether 环，这条上升链稳定，即某个 $n>0$ 使 $A_n=A_{n+1}=\cdots$。考虑 $R$ 中两个都包含 $Q$ 的理想
	\[
	I=(a^n)+Q,
	\qquad
	J=(b)+Q.
	\]
	若 $y\in I\cap J$，则 $y=a^nz+q$ 对某个 $z\in R$ 与 $q\in Q$ 成立。由于 $ab\in Q$，有 $aJ\subseteq Q$，特别地 $ay\in Q$。于是
	\[
	a^{n+1}z=ay-aq\in Q,
	\]
	所以 $z\in A_{n+1}=A_n$。这意味着 $a^nz\in Q$，从而 $y\in Q$。因此 $I\cap J=Q$。由于 $Q$ 不可约，且 $(b)+Q\ne Q$（因为 $b\notin Q$），必须有 $a^n\in Q$。这说明 $Q$ 是 primary。
	
	(2) 的证明与命题 17 第二部分的证明相同。令 $\mathcal S$ 为 $R$ 中不能写成有限个不可约理想交的理想集合。若 $\mathcal S$ 非空，则因 $R$ 是 Noether 环，$\mathcal S$ 有极大元 $I$。则 $I$ 本身不可约性不成立，所以 $I=J\cap K$，其中 $J,K$ 是不同于 $I$ 的理想。于是 $I\subset J$ 且 $I\subset K$，由 $I$ 的极大性，$J,K$ 都不在 $\mathcal S$ 中。因此 $J,K$ 都可写成有限个不可约理想的交，从而 $I$ 也可以，矛盾。因此 $\mathcal S=\varnothing$，命题得证。
\end{proof}

由上一命题立即得到：Noether 环中每个真理想都有 primary 分解。如果某个 primary 分量包含其余 primary 分量的交，就可以把它删去，因为这不改变交。因此可假设分解满足极小分解定义中的 (a)。又由于有限个 $P$-primary 理想的交仍为 $P$-primary（练习 31），把分解中属于同一个素理想的所有 primary 理想替换为它们的交，就可进一步假设分解满足极小分解定义中的 (b)。这证明了下面定理的第一句话。

\noindent\textbf{定理 21（Primary 分解定理）.} 设 $R$ 是 Noether 环。则 $R$ 中每个真理想 $I$ 都有极小 primary 分解。若
\[
I=\bigcap_{i=1}^m Q_i=\bigcap_{i=1}^n Q'_i
\]
是 $I$ 的两个极小 primary 分解，则两个分解中的 associated prime 集合相同：
\[
\{\operatorname{rad}Q_1,\operatorname{rad}Q_2,
\ldots,\operatorname{rad}Q_m\}
=
\{\operatorname{rad}Q'_1,\operatorname{rad}Q'_2,
\ldots,\operatorname{rad}Q'_n\}.
\]
此外，属于这个 associated prime 集合中极小元素的 primary 分量 $Q_i$ 由 $I$ 唯一决定。

\begin{proof}
	associated prime 集合唯一性的证明在练习中概述；属于极小素理想的 primary 分量唯一性的证明将在第 4 节给出。
\end{proof}

\noindent\textbf{定义.} 若 $I$ 是 Noether 环 $R$ 中的理想，则 $I$ 的任一 primary 分解中的 associated prime 称为 $I$ 的 \emph{associated prime ideals}。若 $I$ 的一个 associated prime $P$ 不包含 $I$ 的任何其他 associated prime，则称 $P$ 为 \emph{isolated prime ideal}；其余 associated primes 称为 \emph{embedded prime ideals}。

与一个理想 $I$ 相关的素理想给出了关于 $I$ 的大量信息（例如参见练习 41 与 43）。

\noindent\textbf{推论 22.} 设 $I$ 是 Noether 环 $R$ 中的真理想。
\begin{enumerate}[label=(\arabic*)]
	\item 素理想 $P$ 包含 $I$，当且仅当 $P$ 包含 $I$ 的某个 associated prime；也当且仅当 $P$ 包含 $I$ 的某个 isolated prime。也就是说，$I$ 的 isolated primes 正是所有包含 $I$ 的素理想集合中的极小元。特别地，在包含 $I$ 的素理想中，极小元只有有限多个。
	\item $I$ 的根式是 $I$ 的 associated primes 的交，因而也是 $I$ 的 isolated primes 的交。
	\item 存在包含 $I$ 的素理想 $P_1,
	\ldots,P_n$（不一定互异），使得
	\[
	P_1P_2\cdots P_n\subseteq I.
	\]
\end{enumerate}

\begin{proof}
	(1) 的第一句话是练习（参见练习 37），其余部分随之推出。然后 (2) 由 (1) 与命题 12 推出，(3) 由 (2) 与命题 14 推出。
\end{proof}

定理 21 的最后一句说明：不仅 isolated primes，而且属于 isolated primes 的 primary 分量，也由 $I$ 唯一决定。一般而言，一个理想 $I$ 的 primary 分解本身并不唯一。

\paragraph{例}
\begin{enumerate}[label=(\arabic*)]
	\item 令 $I=(x^2,xy)\subseteq\mathbb R[x,y]$。则
	\[
	(x^2,xy)=(x)\cap(x,y)^2=(x)\cap(x^2,y)
	\]
	是 $I$ 的两个极小 primary 分解。$I$ 的 associated primes 是 $(x)$ 以及
	\[
	\operatorname{rad}((x,y)^2)=\operatorname{rad}(x^2,y)=(x,y).
	\]
	由于 $(x)\subset(x,y)$，素理想 $(x)$ 是唯一的 isolated prime；$(x,y)$ 是 embedded prime。素理想 $P$ 包含 $I$，当且仅当 $P$ 包含 $(x)$。与这个 isolated prime 对应的 $(x)$-primary 分量就是 $(x)$，并且出现在两个 primary 分解中；而与 embedded prime $(x,y)$ 对应的 $(x,y)$-primary 分量并不唯一——在第一个分解中是 $(x,y)^2$，在第二个分解中是 $(x^2,y)$。$I$ 的根式是 isolated prime $(x)$。
	
	这个例子说明了术语的来源：一般地，由理想 $I$ 定义的代数空间 $Z(I)$ 的不可约分量，是 $I$ 的 isolated primes 的零点集；embedded primes 的零点集则是这些不可约分量中的不可约子空间（因此“嵌入”在不可约分量中）。在这个例子中，$Z(I)$ 是满足 $x^2=xy=0$ 的点集，即 $\mathbb R^2$ 中的 $y$-轴。这个代数空间只有一个不可约分量，即 $y$-轴；它是 isolated prime $(x)$ 的轨迹。embedded prime $(x,y)$ 的轨迹是原点 $(0,0)$，这是嵌入在 $y$-轴中的不可约子空间。
	\item 设 $R$ 是 U.F.D.。若
	\[
	a=p_1^{e_1}\cdots p_m^{e_m}
	\]
	是元素 $a\in R$ 分解为互异素幂的唯一分解，则
	\[
	(a)=(p_1)^{e_1}\cap\cdots\cap(p_m)^{e_m}
	\]
	是主理想 $(a)$ 的极小 primary 分解。与 $(a)$ 相关的素理想为 $(p_1),
	\ldots,(p_m)$，它们全都是 isolated。理想的 primary 分解是元素分解为素数幂的推广。另见练习 44，其中用极小 primary 分解刻画 U.F.D.。
\end{enumerate}

对任意 Noether 环，理想 $I$ 是根式理想，当且仅当 $I$ 的一个极小 primary 分解中的 primary 分量全是素理想（在这种情况下该 primary 分解唯一），参见练习 43。这推广了前面由命题 17 与 Hilbert 零点定理得到的观察：当 $k$ 是代数闭域时，$k[\mathbb A^n]$ 中任意根式理想都可唯一写成有限个素理想的交；这正是“代数集可以唯一分解成不可约代数集的并”的代数表述。


\subsection*{练习}

\begin{enumerate}[label=\arabic*.]
	\item 直接证明推论 22 的 (3)：考虑不含有限个素理想乘积的理想集合 $\mathcal S$。[若 $I$ 是 $\mathcal S$ 的极大元，说明 $I$ 不是素理想时存在严格包含 $I$ 的理想 $J,K$，使得 $JK\subseteq I$。]
	
	\item 设 $I,J$ 是环 $R$ 的理想。证明：
	\begin{enumerate}[label=(\alph*)]
		\item 若 $J^k\subseteq I$ 对某个 $k\ge1$ 成立，则 $\operatorname{rad}J\subseteq\operatorname{rad}I$。
		\item 若 $I^k\subseteq J\subseteq I$ 对某个 $k\ge1$ 成立，则 $\operatorname{rad}I=\operatorname{rad}J$。
		\item $\operatorname{rad}(IJ)=\operatorname{rad}(I\cap J)=\operatorname{rad}I\cap\operatorname{rad}J$。
		\item $\operatorname{rad}(\operatorname{rad}I)=\operatorname{rad}I$。
		\item $\operatorname{rad}I+\operatorname{rad}J\subseteq\operatorname{rad}(I+J)$，且
		\[
		\operatorname{rad}(I+J)=\operatorname{rad}(\operatorname{rad}I+\operatorname{rad}J).
		\]
	\end{enumerate}
	
	\item 证明两个根式理想的交仍为根式理想。
	
	\item 令 $I=\mathfrak m_1\mathfrak m_2$，其中 $\mathfrak m_1=(x,y)$ 与 $\mathfrak m_2=(x-1,y-1)$ 是 $\mathbb F_2[x,y]$ 中的理想。证明 $I$ 是根式理想。证明 $\mathbb F_2[x,y]$ 中的理想 $(x^3-y^2)$ 是根式理想。
	
	\item 若 $I=(xy,(x-y)z)\subseteq k[x,y,z]$，证明
	\[
	\operatorname{rad}I=(xy,xz,yz).
	\]
	对该理想直接证明 $Z(I)=Z(\operatorname{rad}I)$，证明 $Z(I)$ 不是不可约的，且 $\operatorname{rad}I$ 不是素理想。
	
	\item 举例说明：在不是代数闭的域 $k$ 上，即便 $I$ 是根式理想，包含 $I\subseteq I(Z(I))$ 也可能是真包含。
	
	\item 设 $R,S$ 是环，$\varphi:R\to S$ 是环同态。若 $I$ 是 $R$ 的理想，证明
	\[
	\varphi(\operatorname{rad}I)\subseteq\operatorname{rad}(\varphi(I)).
	\]
	若另外 $\varphi$ 满且 $I$ 含有 $\ker\varphi$，证明等号成立。
	
	\item 设素理想 $P$ 包含理想 $I$。证明 $P$ 包含 $I$ 的根式。
	
	\item 证明对任意域 $k$，零点定理中的映射 $Z$ 总是满射，而映射 $I$ 总是单射。[使用第 15.1 节中 $Z$ 与 $I$ 的性质 (10)。] 给出例子说明在不是代数闭的域上 $Z$ 可以不是单射，$I$ 可以不是满射。
	
	\item 证明若 $k$ 是有限域，则 Zariski 拓扑就是离散拓扑：每个子集都是闭的（也是开的）。
	
	\item 设 $V$ 是 $\mathbb A^n$ 中的簇，$U_1,U_2$ 是 $\mathbb A^n$ 中 Zariski 开集。证明若 $V\cap U_1\ne\varnothing$ 且 $V\cap U_2\ne\varnothing$，则
	\[
	V\cap U_1\cap U_2\ne\varnothing.
	\]
	推出簇的任意非空开子集在 Zariski 拓扑中处处稠密（即其闭包为整个 $V$）。
	
	\item 利用仿射簇的非空开集处处稠密这一事实，证明仿射簇在 Zariski 拓扑中连通。（拓扑空间称为连通，若它不能表示为两个互不相交的真开子集的并。）
	
	\item 证明仿射代数集 $V$ 在 Zariski 拓扑中连通，当且仅当 $k[V]$ 不是两个非零理想的直和。由此推出簇在 Zariski 拓扑中连通。
	
	\item 若 $k$ 是无限域，证明 $\mathbb A^1$ 中的簇是空集、整个空间以及单点集。若 $k$ 是有限域，$\mathbb A^1$ 中的簇是什么？
	
	\item 设 $V$ 是 $\mathbb A^n$ 中的超曲面，且 $I(V)=(f)$，其中 $f\in k[x_1,x_2,
	\ldots,x_n]$ 是非常数多项式。证明 $V$ 是簇当且仅当 $f$ 不可约。
	
	\item 设 $V\subseteq\mathbb A^n$ 是仿射簇，且 $f\in k[V]$。证明 $f$ 的图像（参见第 15.1 节练习 25）是仿射簇。
	
	\item 证明域 $k$ 的元素的任意置换都是从 $\mathbb A^1$ 到自身的 Zariski 连续映射。推出若 $k$ 是无限域，则存在从 $\mathbb A^1$ 到自身的 Zariski 连续映射不是多项式映射。
	
	\item 令 $V$ 是 $k=\mathbb C$ 上 $\mathbb A^n$ 中的仿射代数集。
	\begin{enumerate}[label=(\alph*)]
		\item 证明 $\mathbb C$ 上代数集的态射在欧氏拓扑中连续（把 $\mathbb C^n$ 识别为 $\mathbb R^{2n}$）。
		\item 证明 $V$ 是 $\mathbb C^n$ 中欧氏拓扑闭集（因此 $\mathbb A^n$ 的 Zariski 闭集也是欧氏闭集）。
		\item 给出一个欧氏拓扑闭但不是 Zariski 闭的集合，即不是仿射代数集（所以欧氏拓扑比 Zariski 拓扑更细）。
	\end{enumerate}
	
	\item 给出一个单射 $k$-代数同态 $\varphi:k[W]\to k[V]$，使得其对应态射 $V\to W$ 不是满射。
	
	\item 设 $\varphi:V\to W$ 是仿射代数集的满态射。证明若 $V$ 是簇，则 $W$ 是簇。
	
	\item 设 $V$ 是 $\mathbb A^n$ 中的代数集，$f\in k[V]$。定义
	\[
	V_f=\{v\in V\mid f(v)\ne0\}.
	\]
	\begin{enumerate}[label=(\alph*)]
		\item 证明 $V_f$ 是 $V$ 中的 Zariski 开集（称为 $V$ 中的主开集）。
		\item 令 $J$ 为 $k[x_1,
		\ldots,x_n,x_{n+1}]$ 中由 $I(V)$ 与 $x_{n+1}f-1$ 生成的理想，并令 $W=Z(J)\subseteq\mathbb A^{n+1}$。证明 $J=I(W)$，且投影 $\pi:\mathbb A^{n+1}\to\mathbb A^n$ 到前 $n$ 个坐标，是从 $W$ 到 $V_f$ 的 Zariski 连续双射（所以 $V$ 中的主开集 $V_f$ 可以嵌入某个更大仿射空间中作为闭集）。
		\item 若 $U$ 是 $V$ 的任意开集，证明
		\[
		U=V_{f_1}\cup\cdots\cup V_{f_m}
		\]
		对某些 $f_1,
		\ldots,f_m\in k[V]$ 成立。（这说明主开集构成 Zariski 拓扑的一组基。）
	\end{enumerate}
	
	\item 证明 $GL_n(k)$ 是 $\mathbb A^{n^2}$ 中的开仿射代数集，并且可嵌入 $\mathbb A^{n^2+1}$ 中作为闭仿射代数集。特别地，推出 $\mathbb A^1$ 中非零元素集合 $k^\times$ 可嵌入 $\mathbb A^2$ 中作为双曲线 $xy=1$。[使用前一题。]
	
	\item 若 $k$ 是无限域，证明
	\[
	\{(a,a^2,a^3)\mid a\in k\}\subseteq\mathbb A^3
	\]
	是仿射代数簇。若 $k$ 是有限域，证明该集合总是可约的。
	
	\item 令
	\[
	V=Z(xz-y^2,
	yz-x^3,
	z^2-x^2y)\subseteq\mathbb A^3.
	\]
	若 $k$ 是无限域，证明 $V$ 是仿射簇。[使用第 15.1 节练习 26 与本节练习 20。]
	
	\item 设 $f(x)=x^3+ax^2+bx+c$ 是 $\mathbb Q[x]$ 中不可约三次多项式，判别式为 $D$。令
	\[
	I=(x+y+z+a,
	xy+xz+yz-b,
	xyz+c)\subseteq\mathbb Q[x,y,z].
	\]
	\begin{enumerate}[label=(\alph*)]
		\item 证明 $I$ 是素理想当且仅当 $D$ 不是 $\mathbb Q$ 中的平方；此时 $I$ 是极大理想，且 $\mathbb Q[x,y,z]/I$ 是 $f(x)$ 在 $\mathbb Q$ 上的分裂域。
		\item 若 $D=r^2$，证明 $I=Q_+\cap Q_-$ 是 primary 分解，其中
		\[
		Q_\pm=(I,
		(x-y)(x-z)(y-z)\pm r).
		\]
		证明 $Q_+$ 与 $Q_-$ 是极大理想，且 $\mathbb Q[x,y,z]$ 模掉 $Q_+$ 或 $Q_-$ 都是 $f(x)$ 在 $\mathbb Q$ 上的分裂域。
	\end{enumerate}
	
	\item 拓扑空间 $X$ 称为拟紧，若任意闭集族 $V_i$ 满足交为空，则其中有限多个的交已经为空，即
	\[
	\bigcap_i V_i=\varnothing
	\quad\Rightarrow\quad
	V_{i_1}\cap\cdots\cap V_{i_N}=\varnothing.
	\]
	证明每个仿射代数集都是拟紧的。[把该定义翻译成 $k[x_1,
	\ldots,x_n]$ 中理想的性质。]（拟紧且 Hausdorff 的空间称为紧空间。）
	
	\item 当 $k$ 是无限域时，证明 $k^2$ 上的 Zariski 拓扑不同于先在 $k$ 上取 Zariski 拓扑再形成积拓扑所得的拓扑。[由第 15.1 节练习 14，在积拓扑中，$k\times k$ 中的 Zariski 闭集是有限个形如 $\{a\}\times\{b\}$、$\{a\}\times k$ 与 $k\times\{b\}$ 的集合的并。]
	
	\item 证明下列每个环都有无限多个极小素理想，并且 $(0)$ 不是其中有限多个的交（所以 $(0)$ 在这些环中没有 primary 分解）：
	\begin{enumerate}[label=(\alph*)]
		\item 无限直积环
		\[
		\mathbb Z/2\mathbb Z\times\mathbb Z/2\mathbb Z\times\cdots
		\]
		（这是 Boolean 环，参见第 7.4 节练习 23）。
		\item
		\[
		k[x_1,x_2,
		\ldots]/(x_1x_2,x_3x_4,
		\ldots,x_{2i-1}x_{2i},
		\ldots),
		\]
		其中 $x_1,x_2,
		\ldots$ 是域 $k$ 上的独立变量。
	\end{enumerate}
	
	\item 设 $A,B$ 是理想，且 $AB\subseteq Q$，其中 $Q$ 是 primary 理想。证明若 $A\not\subseteq Q$，则 $B\subseteq\operatorname{rad}Q$。
	
	\item 设 $Q$ 是 $P$-primary 理想，并设 $A$ 是不包含于 $Q$ 的理想。定义
	\[
	A'=\{r\in R\mid rA\subseteq Q\}.
	\]
	证明 $A'$ 是 $P$-primary 理想。
	
	\item 证明若 $Q_1,Q_2$ 是同属素理想 $P$ 的 primary 理想，则 $Q_1\cap Q_2$ 仍是属于 $P$ 的 primary 理想。推出有限个 $P$-primary 理想的交仍是 $P$-primary。
	
	\item 证明若 $Q_1,Q_2$ 是同属极大理想 $M$ 的 primary 理想，则 $Q_1+Q_2$ 与 $Q_1Q_2$ 都是属于 $M$ 的 primary 理想。推出有限个 $M$-primary 理想的和与乘积仍是 $M$-primary。
	
	\item 令
	\[
	I=(x^2,xy,xz,yz)\subseteq k[x,y,z].
	\]
	证明
	\[
	I=(x,y)\cap(x,z)\cap(x,y,z)^2
	\]
	是 $I$ 的 primary 分解，求出 $I$ 的 isolated primes 与 embedded primes，并求 $\operatorname{rad}I$。
	
	\item 设 $\varphi:R\to S$ 是满环同态。证明：含有 $\ker\varphi$ 的理想 $Q\subseteq R$ 是 primary，当且仅当 $\varphi(Q)$ 是 $S$ 中 primary 理想；此时 $\varphi(Q)$ 对应的素理想是 $Q$ 对应素理想 $P$ 的像 $\varphi(P)$。
	
	\item 设 $\varphi:R\to S$ 是环同态。
	\begin{enumerate}[label=(\alph*)]
		\item 设 $I$ 是含有 $\ker\varphi$ 的 $R$ 的理想，且有最小 primary 分解
		\[
		I=Q_1\cap\cdots\cap Q_m,
		\qquad \operatorname{rad}Q_i=P_i.
		\]
		若 $\varphi$ 满，证明
		\[
		\varphi(I)=\varphi(Q_1)\cap\cdots\cap\varphi(Q_m)
		\]
		是 $\varphi(I)$ 的最小 primary 分解，且 $\operatorname{rad}\varphi(Q_i)=\varphi(P_i)$。[使用前一练习。]
		\item 设 $I$ 是 $S$ 的理想，且有最小 primary 分解
		\[
		I=Q_1\cap\cdots\cap Q_m,
		\qquad \operatorname{rad}Q_i=P_i.
		\]
		证明
		\[
		\varphi^{-1}(I)=\varphi^{-1}(Q_1)\cap\cdots\cap\varphi^{-1}(Q_m)
		\]
		是 $\varphi^{-1}(I)$ 的 primary 分解，且 $\operatorname{rad}\varphi^{-1}(Q_i)=\varphi^{-1}(P_i)$；若 $\varphi$ 满，则该分解最小。
	\end{enumerate}
	
	\item 令 $I=(xy,x-yz)\subseteq k[x,y,z]$。证明
	\[
	I=(x,z)\cap(y,x-yz)
	\]
	是最小 primary 分解。[考虑环同态 $\varphi:k[x,y,z]\to k[y,z]$，$x\mapsto yz$，$y\mapsto y$，$z\mapsto z$，并使用前一练习。]
	
	\item 证明：素理想 $P$ 包含理想 $I$，当且仅当 $P$ 包含 $I$ 的某个最小 primary 分解的 associated prime。[使用练习 3 与第 7.4 节练习 11。]
	
	\item 证明根式理想的每个 associated prime 都是 isolated prime。[设根式理想 $I$ 在定理 21 的分解中有 $P_2=\operatorname{rad}Q_2\subseteq P_1=\operatorname{rad}Q_1$。若 $a\in Q_2\cap\cdots\cap Q_m$ 但 $a\notin P_2$，说明某个幂 $a^n$ 属于 $I$，从而 $a\in Q_1$，推出与分解最小性矛盾。]
	
	\item 固定环 $R$ 中元素 $a$。对任意理想 $I$，令
	\[
	I_a=\{r\in R\mid ar\in I\}.
	\]
	\begin{enumerate}[label=(\alph*)]
		\item 证明 $I_a$ 是理想，并且 $I_a=R$ 当且仅当 $a\in I$。
		\item 对理想 $I,J$，证明 $(I\cap J)_a=I_a\cap J_a$。
		\item 若 $Q$ 是 $P$-primary 理想且 $a\notin Q$，证明 $Q_a$ 是 $P$-primary 理想，并且若 $a\notin P$，则 $Q_a=Q$。
	\end{enumerate}
	
	\item 沿用前一题记号。设
	\[
	I=Q_1\cap\cdots\cap Q_m
	\]
	是理想 $I$ 的最小 primary 分解，$P_i$ 为 $Q_i$ 对应的素理想。
	\begin{enumerate}[label=(\alph*)]
		\item 证明
		\[
		I_a=(Q_1)_a\cap\cdots\cap(Q_m)_a
		\]
		且
		\[
		\operatorname{rad}(I_a)=\operatorname{rad}((Q_1)_a)\cap\cdots\cap\operatorname{rad}((Q_m)_a).
		\]
		\item 证明 $\operatorname{rad}(I_a)$ 是那些满足 $a\notin Q_i$ 的素理想 $P_i$ 的交。[使用前一练习。]
		\item 证明若 $\operatorname{rad}(I_a)$ 是素理想，则 $\operatorname{rad}(I_a)=P_j$ 对某个 $j$ 成立。[使用素理想不可约这一事实。]
		\item 对每个 $i=1,
		\ldots,m$，证明存在 $a\in R$ 使得 $\operatorname{rad}(I_a)=P_i$。[说明存在 $a\notin Q_i$ 但 $a\in Q_j$ 对所有 $j\ne i$ 成立。]
		\item 由 (c) 与 (d) 说明，最小 primary 分解的 associated primes 正是所有 $a\in R$ 中使 $\operatorname{rad}(I_a)$ 为素理想者构成的集合；从而这些素理想由 $I$ 唯一决定，与最小 primary 分解的选取无关。
	\end{enumerate}
	
	\item 设 $P_1,
	\ldots,P_m$ 是 Noether 环 $R$ 中理想 $(0)$ 的 associated prime。
	\begin{enumerate}[label=(\alph*)]
		\item 证明 $P_1\cap\cdots\cap P_m$ 是 $R$ 中所有幂零元素的集合。[对 $I=(0)$ 应用推论 22。]
		\item 证明 $P_1\cup\cdots\cup P_m$ 是 $R$ 中所有零因子的集合。[在前一练习中令 $I=(0)$，说明零因子集合由 $\bigcup_{a\in R-(0)}(0)_a=\bigcup_{a\in R-(0)}\operatorname{rad}((0)_a)$ 给出。]
	\end{enumerate}
	
	\item 设 $R$ 是 Noether 环。证明 $R$ 要么是整环，要么含有非零幂零元素，要么至少有两个极小素理想。[使用前一练习。]
	
	\item 证明 Noether 环 $R$ 中的理想 $I$ 是根式理想，当且仅当其最小 primary 分解中的 primary 分量全是素理想；并推出此时最小 primary 分解唯一。[若 $I=Q_1\cap\cdots\cap Q_m$ 是根式理想的最小分解，$Q_i$ 为 $P_i$-primary，说明若 $a\in P_1\cap\cdots\cap P_m$，则某个幂 $a^n\in I$，从而 $a\in I$。推出 $I=P_1\cap\cdots\cap P_m$，并说明这是最小 primary 分解；再证明 $Q_i=P_i$。]
	
	\item 证明 Noether 整环 $R$ 是 U.F.D.，当且仅当对每个 $a\in R$，主理想 $(a)$ 的 isolated associated primes 都是主理想。[参见推论 22 后例 2。为了证明 $R$ 是 U.F.D.，说明不可约元 $a\in R$ 是素元，然后仿照第 8.3 节定理 14 的证明。]
	
	\item 令 $R$ 为开区间 $(-1,1)$ 上所有具有任意阶导数的实值函数构成的环（$C^\infty$ 函数环）。令
	\[
	F(x)=
	\begin{cases}
		e^{-1/x^4},&x\ne0,\\
		0,&x=0,
	\end{cases}
	\]
	（可假设 $F\in R$ 且对所有 $n\ge0$，$F^{(n)}(0)=0$）。令 $(F)$ 为由 $F$ 生成的主理想，$A=\operatorname{rad}((F))$。令 $M$ 为所有在 $x=0$ 处为零的函数构成的极大理想，并令 $P=\bigcap_{n\ge1}M^n$。
	\begin{enumerate}[label=(\alph*)]
		\item 证明 $M=(x)$ 是由函数 $x$ 生成的理想，并且 $M^n=(x^n)$ 由在原点前 $n-1$ 阶导数都为零的函数组成。
		\item 证明 $R$ 不是 Noether 环（比较第 7.4 节练习 33）。[一种方法：令 $G(x)$ 在 $x<0$ 时为 $0$，在 $x\ge0$ 时等于 $F(x)$。令 $I_n$ 为所有在 $x\le1/n$ 时为零的函数构成的理想。使用 $G(x)$ 的平移说明 $I_1\subset I_2\subset I_3\subset\cdots$ 是无限上升链。]
		\item 证明 $P$ 由所有在 $x=0$ 处各阶导数均为零的函数组成（即在 $x=0$ 处 Taylor 级数恒为零的函数），并且 $P$ 是素理想。
		\item 证明 $F\in P$，并推出 $A\subseteq P$。
		\item 证明 $A\ne P$。[令 $G(x)=e^{-1/x^2}$ 当 $x\ne0$，$G(0)=0$。证明 $G\in P$ 但 $G\notin A$。]
		\item 说明存在含有 $(F)$ 的素理想 $Q$，使得 $Q\ne P,M$。证明 $Q\subset P$；即存在非零素理想真包含于 $P$。
	\end{enumerate}
	
	\item 设 $A$ 是 $R=k[x_1,
	\ldots,x_n,y_1,
	\ldots,y_m]$ 中任意理想。
	\begin{enumerate}[label=(\alph*)]
		\item 证明
		\[
		\operatorname{rad}(A\cap k[y_1,
		\ldots,y_m])=\operatorname{rad}A\cap k[y_1,
		\ldots,y_m].
		\]
		\item 设 $(f_1,
		\ldots,f_r)$ 是 $k[x_1,
		\ldots,x_n]$ 中的理想。令 $F_1,
		\ldots,F_l$ 为该理想根式在 $k[x_1,
		\ldots,x_n]$ 中的生成元。设 $J$ 是 $R$ 中理想，令
		\[
		A=J+(f_1,
		\ldots,f_r),
		\qquad
		B=J+(F_1,
		\ldots,F_l).
		\]
		证明 $\operatorname{rad}A=\operatorname{rad}B$。
		\item 由 (a) 与 (b) 推出，在代数闭域 $k$ 上，
		\[
		(y_1-x_1,
		\ldots,y_m-x_m,f_1,
		\ldots,f_r)\cap k[y_1,
		\ldots,y_m]
		\]
		与
		\[
		(y_1-x_1,
		\ldots,y_m-x_m,F_1,
		\ldots,F_l)\cap k[y_1,
		\ldots,y_m]
		\]
		具有相同零点集。[使用 Hilbert 零点定理。]
	\end{enumerate}
	
	\item 求曲线
	\[
	\{(a^2,a^3,a^4)\mid a\in\mathbb C\}
	\]
	在 $\mathbb C^3$ 中的 Zariski 闭包。
	
	\item 证明 $Z(x^3-xyz+z^2)$ 是 $\mathbb R^3$ 中包含点集
	\[
	\{(st,
	s+t,
	s^2t)\mid s,t\in\mathbb R\}
	\]
	的最小代数集。
	
	\item 证明 $Z(x^3z^2-3xy^2z^2-y^6-z^4)$ 是 $\mathbb R^3$ 中包含点集
	\[
	\{(s^2+t^2,
	st,
	s^3)\mid s,t\in\mathbb R\}
	\]
	的最小代数集。
	
	\item 求点集
	\[
	\{(s^4,
	s^3t,
	st^3,
	t^4)\mid s,t\in\mathbb R\}
	\]
	的 Zariski 闭包的定义方程。
	
	\item 证明 Whitney 伞面
	\[
	V=Z(x^2-y^2z)
	\]
	是 $\mathbb R^3$ 中包含
	\[
	S=\{(st,
	s,
	t^2)\mid s,t\in\mathbb R\}
	\]
	的最小代数集。证明 $S$ 在 $V$ 中不是 Zariski 闭的（缺失的点解释了该曲面名称的由来）。在 $\mathbb C$ 上也作同样讨论，但证明此时 $S=V$，因而是闭的。
	
	\item 令
	\[
	V=Z(xz^2-w^3,
	xw^2-y^4,
	y^4z^2-w^5)\subseteq\mathbb C^4.
	\]
	求 $V$ 在投影
	\[
	\pi((x,y,z,w))=(x,y,z)
	\]
	下的像的 Zariski 闭包。
	
	\item 令 $V=Z(xy-1)\subseteq\mathbb A^2$，令 $S$ 为 $V$ 到 $x$ 轴 $\mathbb A^1$ 的投影。
	\begin{enumerate}[label=(\alph*)]
		\item 若 $k=\mathbb R$，证明 $I(V)=(xy-1)\subseteq\mathbb R[x,y]$，且
		\[
		(u-x,xy-1)\cap\mathbb R[u]=0
		\]
		在 $\mathbb R[x,y,u]$ 中成立。使用命题 8 与 16 推出 $S$ 的 Zariski 闭包是 $\mathbb A^1$，并说明 $S$ 本身不是闭集。
		\item 若 $k=\mathbb F_3$，证明
		\[
		I(V)=(xy-1,
		x^3-x,
		y^3-y)\subseteq\mathbb F_3[x,y]
		\]
		且
		\[
		(u-x,
		xy-1,
		x^3-x,
		y^3-y)\cap\mathbb F_3[u]=(u^2-1)
		\]
		在 $\mathbb F_3[x,y,u]$ 中成立。使用命题 8 与 16 推出 $S$ 在 $\mathbb A^1$ 中是 Zariski 闭的。
	\end{enumerate}
	
	\item 回忆两个理想 $I,J$ 的理想商
	\[
	(I:J)=\{r\in R\mid rJ\subseteq I\}
	\]
	（参见第 9.6 节练习 34 以后）。显然 $I\subseteq(I:J)$。
	\begin{enumerate}[label=(\alph*)]
		\item 证明 $Z(I)-Z(J)$ 中的元素都属于 $Z((I:J))$，并推出 $Z(I)-Z(J)$ 的 Zariski 闭包包含于 $Z((I:J))$。
		\item 若 $k$ 是代数闭域且 $I$ 是根式理想，证明 $Z((I:J))$ 正好是 $Z(I)-Z(J)$ 的 Zariski 闭包。
		\item 若 $V,W$ 是仿射代数集，证明
		\[
		(I(V):I(W))=I(V-W).
		\]
	\end{enumerate}
\end{enumerate}
\section{整扩张与 Hilbert 零点定理}

本节中，我们考虑环的整扩张这一重要概念，它是域的代数扩张向环的推广。这将导向 $\mathbb Q$ 的有限扩张中“整数”的定义（这是称为代数数论这一数学分支的基本主题），并且也与代数曲线的切线的存在性有关。

\noindent\textbf{定义.} 设 $R$ 是含于交换环 $S$ 中的子环，并且 $1=1_S\in R$。
\begin{enumerate}[label=(\arabic*)]
	\item 元素 $s\in S$ 称为在 $R$ 上\emph{整}，若 $s$ 是 $R[x]$ 中某个首一多项式的根。
	\item 环 $S$ 称为 $R$ 的\emph{整扩张}，或简称 $S$ 在 $R$ 上整，若每个 $s\in S$ 都在 $R$ 上整。
	\item $R$ 在 $S$ 中的\emph{整闭包}是 $S$ 中所有在 $R$ 上整的元素的集合。
	\item 若 $R$ 等于它在 $S$ 中的整闭包，则称 $R$ 在 $S$ 中\emph{整闭}。整环 $R$ 在其分式域中的整闭包称为 $R$ 的\emph{正规化}。若整环在其分式域中整闭，则称它是\emph{整闭的}或\emph{正规的}。
\end{enumerate}

在给出整扩张的一些例子之前，我们先证明整元素的一些基本性质，它们类似于域上代数元素的性质。

\noindent\textbf{命题 23.} 设 $R$ 是含于交换环 $S$ 中的子环，且 $1\in R$，并设 $s\in S$。则下列条件等价：
\begin{enumerate}[label=(\arabic*)]
	\item $s$ 在 $R$ 上整；
	\item $R[s]$ 是有限生成 $R$-模，其中 $R[s]$ 是 $s$ 的幂的所有 $R$-线性组合所成的环；
	\item 对某个子环 $T$，$s\in T$，$R\subseteq T\subseteq S$，且 $T$ 是有限生成 $R$-模。
\end{enumerate}

\begin{proof}
	先假设 (1) 成立，并设 $s$ 是首一多项式
	\[
	x^n+a_{n-1}x^{n-1}+\cdots+a_0\in R[x]
	\]
	的根。于是
	\[
	s^n=-(a_{n-1}s^{n-1}+a_{n-2}s^{n-2}+\cdots+a_0),
	\]
	所以 $s^n$，进而 $s$ 的所有更高次幂，都可以表示为 $s^{n-1},\ldots,s,1$ 的 $R$-线性组合。因此
	\[
	R[s]=R\cdot 1+Rs+\cdots+Rs^{n-1}
	\]
	作为 $R$-模是有限生成的，这给出 (2)。
	
	若 (2) 成立，则取 $T=R[s]$ 即得 (3)。
	
	现在假设 (3) 成立，并令 $v_1,v_2,\ldots,v_n$ 是 $T$ 的一个有限生成集。对 $i=1,2,\ldots,n$，元素 $sv_i$ 属于 $T$，因为 $T$ 是环，所以它可写成 $v_1,\ldots,v_n$ 的 $R$-线性组合：
	\[
	sv_i=\sum_{j=1}^n a_{ij}v_j,
	\]
	即
	\[
	0=\sum_{j=1}^n(\delta_{ij}s-a_{ij})v_j,\qquad i=1,2,\ldots,n,
	\]
	其中 $\delta_{ij}$ 是 Kronecker delta。若 $B$ 是其第 $i,j$ 项为 $\delta_{ij}s-a_{ij}$ 的 $n\times n$ 矩阵，而 $v$ 是以 $v_1,\ldots,v_n$ 为项的 $n\times1$ 列向量，则这些方程正是 $Bv=0$。由 Cramer 法则可得 $(\det B)v_i=0$ 对所有 $i$ 成立（参见第 11.4 节练习 3）。由于 $1\in T$ 是 $v_1,\ldots,v_n$ 的 $R$-线性组合，遂有 $\det B=0$。但 $B=sI-A$，其中 $A$ 是矩阵 $(a_{ij})$。因此 $s$ 是 $R[x]$ 中首一多项式 $\det(xI-A)$（即 $A$ 的特征多项式）的根，于是 $s$ 是某个系数在 $R$ 中的首一多项式的根，这给出 (1)，证明完毕。
\end{proof}

\noindent\textbf{推论 24.} 设 $R\subseteq S$ 如命题 23，且 $s,t\in S$。
\begin{enumerate}[label=(\arabic*)]
	\item 若 $s$ 与 $t$ 都在 $R$ 上整，则 $s\pm t$ 与 $st$ 也在 $R$ 上整。
	\item $R$ 在 $S$ 中的整闭包是 $S$ 的一个含有 $R$ 的子环。
	\item 整性具有传递性：设 $S$ 是 $T$ 的一个子环；若 $T$ 在 $S$ 上整且 $S$ 在 $R$ 上整，则 $T$ 在 $R$ 上整。
\end{enumerate}

\begin{proof}
	设 $s,t$ 在 $R$ 上整。由命题 23，$R[s]$ 与 $R[t]$ 都是有限生成 $R$-模，设
	\[
	R[s]=Rs_1+Rs_2+\cdots+Rs_n,
	\qquad
	R[t]=Rt_1+Rt_2+\cdots+Rt_m.
	\]
	则
	\[
	R[s,t]=Rs_1t_1+\cdots+Rs_it_j+\cdots+Rs_nt_m
	\]
	是一个含有 $s\pm t$ 与 $st$ 的环，并且它也是有限生成 $R$-模。因此 $s\pm t$ 与 $st$ 也在 $R$ 上整，这证明 (1)，同时也证明 (2)。
	
	为证明 (3)，取 $t\in T$。因为 $t$ 在 $S$ 上整，所以 $t$ 是某个首一多项式
	\[
	p(x)=x^n+a_{n-1}x^{n-1}+\cdots+a_0\in S[x]
	\]
	的根。由于 $a_i\in S$ 在 $R$ 上整，每个环 $R[a_i]$ 都是有限生成 $R$-模，因此环
	\[
	R_1=R[a_0,a_1,\ldots,a_{n-1}]
	\]
	也是有限生成 $R$-模。由于首一多项式 $p(x)$ 的系数在 $R_1$ 中，$t$ 在 $R_1$ 上整，于是环
	\[
	R_1[t]=R[a_0,a_1,\ldots,a_{n-1},t]
	\]
	是有限生成 $R$-模。由命题 23，这意味着 $t$ 在 $R$ 上整，从而得到 (3)。
\end{proof}

推论 24 的第二个断言说明，在 $S$ 中取所有在 $R$ 上整的元素会得到 $S$ 的一个（可能更大的）子环；而该推论最后一个断言说明，取整闭包的过程一步之后便停止：

\noindent\textbf{推论 25.} 设 $R$ 是含于交换环 $S$ 中的子环，且 $1\in R$。则 $R$ 在 $S$ 中的整闭包在 $S$ 中整闭。

\paragraph{例}
\begin{enumerate}[label=(\arabic*)]
	\item 若 $R$ 与 $S$ 是域，则 $S$ 在 $R$ 上整当且仅当 $S$ 在 $R$ 上代数：若 $s\in S$ 是一个系数在 $R$ 中的多项式 $p(x)$ 的根，则它也是把 $p(x)$ 除以其非零首项系数后所得首一多项式的根。
	\item 设 $S$ 是 $R$ 的整扩张，$I$ 是 $S$ 的理想。则 $S/I$ 是 $R/(R\cap I)$ 的整环扩张（把 $s\in S$ 所满足的 $R$ 上首一多项式模掉 $I$，得到 $S/I$ 中 $\bar s$ 在 $R/(R\cap I)$ 上满足的首一多项式）。
	\item 若 $R$ 是 U.F.D.，则 $R$ 是整闭的，理由如下。设 $a/b$ 是 $R$ 的分式域中的元素，其中 $b\ne0$ 且 $a$ 与 $b$ 无公共因子，并设
	\[
	(a/b)^n+r_{n-1}(a/b)^{n-1}+\cdots+r_1(a/b)+r_0=0,
	\]
	其中 $r_0,\ldots,r_{n-1}\in R$。则
	\[
	a^n=b(-r_{n-1}a^{n-1}-\cdots-r_1ab^{n-2}-r_0b^{n-1})
	\]
	表明任何整除 $b$ 的不可约元素都整除 $a^n$，从而整除 $a$。由于 $a/b$ 已约为最低项，这说明 $b$ 必为单位，即 $a/b\in R$。
	\item 域 $k$ 上的多项式环 $k[x,y]$ 在其分式域 $k(x,y)$ 中是整闭的，这是由上例 (3) 得出的。理想 $(x^2-y^3)$ 是素理想（参见第 9.1 节练习 14），所以商环
	\[
	R=k[x,y]/(x^2-y^3)=k[\bar x,\bar y]
	\]
	是整环。不过这个整环不是整闭的，因为 $\bar x/\bar y$ 是 $R$ 的分式域中的元素，并且在 $R$ 上整（因为 $(\bar x/\bar y)^3-\bar x=0$），但它不是 $R$ 的元素。特别地，由前一例知 $R$ 不是 U.F.D.。
\end{enumerate}

下面考虑整环扩张中理想的行为。

\noindent\textbf{定义.} 设 $\varphi:R\to S$ 是交换环同态。
\begin{enumerate}[label=(\alph*)]
	\item 若 $I$ 是 $R$ 的理想，则 $I$ 到 $S$ 的\emph{扩张}是 $S$ 中由 $I$ 的像生成的理想 $\varphi(I)S$。
	\item 若 $J$ 是 $S$ 的理想，则 $J$ 在 $R$ 中的\emph{收缩}是理想 $\varphi^{-1}(J)$。
\end{enumerate}

在特殊情形 $R$ 是 $S$ 的子环且 $\varphi$ 是自然嵌入时，$I\subseteq R$ 的扩张是 $S$ 中的理想 $IS$，而 $J\subseteq S$ 的收缩是 $R$ 的理想 $J\cap R$。

由定义立刻可得：
\begin{enumerate}[label=(\arabic*)]
	\item $I\subseteq IS\cap R$；更一般地，$I$ 包含于其到 $S$ 的扩张再收缩回来的理想中；
	\item $(J\cap R)S\subseteq J$；更一般地，$J$ 包含其在 $R$ 中的收缩再扩张到 $S$ 的理想。
\end{enumerate}
一般而言，这两处都不必取等号（参见练习）。

若 $Q$ 是 $S$ 中的素理想，则其收缩是 $R$ 中的素理想（不过极大理想的收缩不必是极大理想）。另一方面，若 $P$ 是 $R$ 中的素理想，则其扩张在 $S$ 中不必是素理想（甚至不必是真理想）；此外，$P$ 一般也不一定是 $S$ 中某个素理想的收缩（参见练习）。然而，对于整环扩张，情形要受控得多：

\noindent\textbf{定理 26.} 设 $R$ 是含于交换环 $S$ 中的子环，且 $1\in R$，并假设 $S$ 在 $R$ 上整。
\begin{enumerate}[label=(\arabic*)]
	\item 假设 $S$ 是整环。则 $R$ 是域当且仅当 $S$ 是域。
	\item 设 $P$ 是 $R$ 中的素理想。则存在 $S$ 中的素理想 $Q$，使得 $P=Q\cap R$。此外，$P$ 是极大理想当且仅当 $Q$ 是极大理想。
	\item （Going-up 定理）设
	\[
	P_1\subseteq P_2\subseteq\cdots\subseteq P_n
	\]
	是 $R$ 中的素理想链，并假设 $S$ 中存在素理想
	\[
	Q_1\subseteq Q_2\subseteq\cdots\subseteq Q_m
	\]
	满足 $P_i=Q_i\cap R$，$1\le i\le m$，且 $m<n$。则这个上升理想链可以补全：存在 $S$ 中的素理想
	\[
	Q_{m+1}\subseteq\cdots\subseteq Q_n
	\]
	使得对所有 $i$ 都有 $P_i=Q_i\cap R$。
	\item （Going-down 定理）假设 $S$ 是整环且 $R$ 在 $S$ 中整闭。设
	\[
	P_1\supseteq P_2\supseteq\cdots\supseteq P_n
	\]
	是 $R$ 中的素理想链，并假设 $S$ 中存在素理想
	\[
	Q_1\supseteq Q_2\supseteq\cdots\supseteq Q_m
	\]
	满足 $P_i=Q_i\cap R$，$1\le i\le m$，且 $m<n$。则这个下降理想链可以补全：存在 $S$ 中的素理想
	\[
	Q_{m+1}\supseteq\cdots\supseteq Q_n
	\]
	使得对所有 $i$ 都有 $P_i=Q_i\cap R$。
\end{enumerate}

\begin{proof}
	为证明 (1)，先假设 $R$ 是域，并令 $s$ 是 $S$ 的非零元素。由于 $s$ 在 $R$ 上整，故
	\[
	s^n+a_{n-1}s^{n-1}+\cdots+a_1s+a_0=0
	\]
	对某些 $a_0,a_1,\ldots,a_{n-1}\in R$ 成立。因为 $S$ 是整环，我们可假设 $a_0\ne0$（否则消去 $s$ 的因子）。于是
	\[
	s(s^{n-1}+a_{n-1}s^{n-2}+\cdots+a_1)=-a_0.
	\]
	由于 $(-1/a_0)\in R$，这说明
	\[
	(-1/a_0)(s^{n-1}+a_{n-1}s^{n-2}+\cdots+a_1)
	\]
	是 $s$ 在 $S$ 中的逆元，所以 $S$ 是域。反过来，假设 $S$ 是域而 $r$ 是 $R$ 的非零元素。由于 $r^{-1}\in S$ 在 $R$ 上整，有
	\[
	r^{-m}+a_{m-1}r^{-m+1}+\cdots+a_1r^{-1}+a_0=0
	\]
	对某些 $a_0,\ldots,a_{m-1}\in R$ 成立。于是
	\[
	r^{-1}=-(a_{m-1}+\cdots+a_1r^{m-2}+a_0r^{m-1})\in R,
	\]
	所以 $R$ 是域。
	
	(2) 的第一个断言将在推论 50 中给出证明。对于第二个断言，注意整环 $S/Q$ 是 $R/P$ 的整扩张（推论 25 后的例 (2)）。由 (1)，$S/Q$ 是域当且仅当 $R/P$ 是域，也即 $Q$ 是极大理想当且仅当 $P$ 是极大理想。
	
	为证明 (3)，由归纳只需证明：若 $P_1\subseteq P_2$ 且 $Q_1$ 是 $S$ 的素理想，满足 $Q_1\cap R=P_1$，则存在 $S$ 的素理想 $Q_2$，使得 $Q_1\subseteq Q_2$ 且 $Q_2\cap R=P_2$。由于 $\overline S=S/Q_1$ 是 $\overline R=R/P_1$ 的整扩张，(2) 的第一部分说明，存在 $\overline S$ 中的素理想 $\overline Q_2$，满足
	\[
	\overline Q_2\cap\overline R=P_2/P_1.
	\]
	于是 $\overline Q_2$ 在 $S$ 中的原像 $Q_2$ 是一个包含 $Q_1$ 的素理想，并且 $Q_2\cap R=P_2$。
	
	(4) 的证明见第 15.4 节练习 24 的提示。
\end{proof}

\noindent\textbf{推论 27.} 设 $R$ 是环 $S$ 的子环，且 $1\in R$，并假设 $S$ 在 $R$ 上整且作为环在 $R$ 上有限生成。若 $P$ 是 $R$ 的极大理想，则 $S$ 中满足 $Q\cap R=P$ 的极大理想 $Q$ 的个数非零且有限。

\begin{proof}
	由定理 (2)，至少存在一个位于 $P$ 上方的极大理想 $Q$，所以只需说明 $S$ 中这样的极大理想只有有限多个。若 $Q$ 是 $S$ 的极大理想且 $Q\cap R=P$，则 $S/Q$ 是含有域 $R/P$ 的域。为了证明可能的 $Q$ 只有有限多个，只需证明从 $S$ 到含有 $R/P$ 的域且扩张了从 $R$ 到 $R/P$ 的同态的同态只有有限多个。设 $S=R[s_1,
	\ldots,s_n]$，其中每个 $s_i$ 按假设在 $R$ 上整，并令 $p_i(x)$ 是系数在 $R$ 中且以 $s_i$ 为根的首一多项式。若 $Q$ 是 $S$ 的极大理想，则
	\[
	S/Q=(R/P)[\bar s_1,\ldots,\bar s_n]
	\]
	是域 $R/P$ 的扩域，生成元为 $\bar s_1,
	\ldots,
	\bar s_n$。元素 $\bar s_i$ 是把 $p_i(x)$ 的系数模 $P$ 后得到的 $R/P$ 上首一多项式的根。这个首一多项式（在 $R/P$ 的一个固定代数闭包中）只有有限多个可能的根，因此形如 $(R/P)[\bar s_1,
	\ldots,
	\bar s_n]$ 的域扩张只有有限多个可能性，推论得证。
\end{proof}

\subsection*{代数整数}

我们可以用整环扩张的概念定义有理数域 $\mathbb Q$ 的扩域中的“整数”：

\noindent\textbf{定义.} 设 $K$ 是 $\mathbb Q$ 的扩域。
\begin{enumerate}[label=(\arabic*)]
	\item 元素 $\alpha\in K$ 称为\emph{代数整数}，若 $\alpha$ 在 $\mathbb Z$ 上整，即 $\alpha$ 是某个系数在 $\mathbb Z$ 中的首一多项式的根。
	\item $\mathbb Z$ 在 $K$ 中的整闭包称为 $K$ 的\emph{整数环}，记为 $\mathcal O_K$。
\end{enumerate}

代数整数显然在 $\mathbb Q$ 上代数，所以所有代数整数的环就是 $\overline{\mathbb Q}$（$\mathbb Q$ 的一个代数闭包）中的整数环。代数整数的例子包括 $\sqrt2,\sqrt{-1},\sqrt[3]5$ 等，因为这些元素当然都是系数在 $\mathbb Z$ 中的首一多项式的根。代数整数 $\alpha$ 的定义是 $\alpha$ 是 $\mathbb Z[x]$ 中某个首一多项式的根，这一条件似乎不易检验。下一个命题给出一个简单判别法：只需看 $\alpha$ 的最小多项式即可。

\noindent\textbf{命题 28.} 某个 $\mathbb Q$ 的域扩张中的元素 $\alpha$ 是代数整数，当且仅当 $\alpha$ 在 $\mathbb Q$ 上代数，且它的最小多项式 $m_{\alpha,\mathbb Q}(x)$ 具有整数系数。特别地，$\mathbb Q$ 中的代数整数就是整数 $\mathbb Z$，即 $\mathcal O_{\mathbb Q}=\mathbb Z$。

\begin{proof}
	若 $\alpha$ 在 $\mathbb Q$ 上代数且 $m_{\alpha,\mathbb Q}(x)\in\mathbb Z[x]$，则按定义 $\alpha$ 在 $\mathbb Z$ 上整。反过来，假设 $\alpha$ 在 $\mathbb Z$ 上整，并令 $f(x)$ 是 $\mathbb Z[x]$ 中以 $\alpha$ 为根且次数最小的首一多项式。若 $f$ 在 $\mathbb Q[x]$ 中可约，则由 Gauss 引理，$f(x)=g(x)h(x)$，其中 $g(x),h(x)$ 是 $\mathbb Z[x]$ 中次数小于 $f$ 的首一多项式。但这样 $\alpha$ 必为 $g$ 或 $h$ 的根，与 $f$ 的次数最小性矛盾。因此 $f$ 在 $\mathbb Q[x]$ 中不可约，所以 $f(x)=m_{\alpha,\mathbb Q}(x)$，从而 $\alpha$ 的最小多项式具有整数系数。最后，$\alpha=a/b\in\mathbb Q$（$a/b$ 约成最低项且 $b>0$）的最小多项式是 $bx-a$，它是首一的当且仅当 $b=1$。所以 $\alpha\in\mathbb Q$ 是代数整数当且仅当 $\alpha\in\mathbb Z$。
\end{proof}

由于整数 $\mathbb Z$ 是 $\mathbb Q$ 中的代数整数，为了强调（也为了清楚），$\mathbb Z$ 的元素有时称为“有理整数”，以区别于 $\mathbb Q$ 的有限次数扩张（称为数域）中的“整数”。下一个结果给出一般数域中整数环的一些基本结构。

\noindent\textbf{定理 29.} 设 $K$ 是 $\mathbb Q$ 上次数为 $n$ 的数域。
\begin{enumerate}[label=(\arabic*)]
	\item $K$ 的整数环 $\mathcal O_K$ 是 Noether 环，并且是秩为 $n$ 的自由 $\mathbb Z$-模。
	\item 对每个 $\beta\in K$，存在某个非零 $d\in\mathbb Z$，使得 $d\beta$ 是代数整数。特别地，$K$ 是 $\mathcal O_K$ 的分式域。
	\item 若 $\beta_1,\beta_2,
	\ldots,
	\beta_n$ 是 $K$ 的任一 $\mathbb Q$-基，则存在整数 $d$，使得 $d\beta_1,d\beta_2,
	\ldots,d\beta_n$ 是 $\mathcal O_K$ 中一个秩为 $n$ 的自由 $\mathbb Z$-子模的基。$\mathbb Z$-模 $\mathcal O_K$ 的任一基也是 $K$ 作为 $\mathbb Q$ 上向量空间的一组基。
\end{enumerate}

\begin{proof}
	先注意，$\mathcal O_K$ 中元素之间的任何 $\mathbb Z$-线性相关关系都是 $K$ 中的 $\mathbb Q$-线性相关关系；反过来，把 $\mathcal O_K$ 中元素在 $K$ 中的 $\mathbb Q$-线性相关关系乘以系数的公分母，会得到 $\mathcal O_K$ 中的 $\mathbb Z$-线性相关关系。令 $\beta$ 是 $K$ 的任一元素，并设
	\[
	x^k+a_{k-1}x^{k-1}+\cdots+a_0
	\]
	是 $\beta$ 在 $\mathbb Q$ 上的最小多项式。若 $d$ 是这些系数的公分母，则把方程乘以 $d^k$ 得
	\[
	(d\beta)^k+da_{k-1}(d\beta)^{k-1}+\cdots+d^{k-1}a_1(d\beta)+d^ka_0=0,
	\]
	且 $d^ka_0,d^{k-1}a_1,
	\ldots,da_{k-1}\in\mathbb Z$。因此 $d\beta$ 是代数整数，这证明 (2) 的第一部分；(2) 的第二个断言随即得出。
	
	若 $\beta_1,
	\ldots,
	\beta_n$ 是 $K$ 在 $\mathbb Q$ 上的一组基，则存在非零整数 $d$，使得 $d\beta_1,
	\ldots,d\beta_n$ 都属于 $\mathcal O_K$。这些元素在 $\mathbb Q$ 上仍线性无关，因此在 $\mathbb Z$ 上也线性无关，故它们生成 $\mathcal O_K$ 中一个秩为 $n$ 的自由子模，这证明 (3) 的第一部分。
	
	由于 $\mathcal O_K$ 是域 $K$ 的子环，它是无挠 $\mathbb Z$-模。若 $\mathcal O_K$ 包含在某个有限生成 $\mathbb Z$-模中，则 $\mathcal O_K$ 也在 $\mathbb Z$ 上有限生成，从而是自由 $\mathbb Z$-模。令 $L$ 为 $K$ 的 Galois 闭包，则 $\mathcal O_K\subseteq\mathcal O_L$，故只需说明 $\mathcal O_L$ 包含在一个有限生成 $\mathbb Z$-模中。令 $\alpha_1,
	\ldots,
	\alpha_m$ 是 $L$ 的一组 $\mathbb Q$-基。必要时乘以一个整数 $d\in\mathbb Z$，可假设每个 $\alpha_i$ 都是代数整数，即 $\alpha_1,
	\ldots,
	\alpha_m\in\mathcal O_L$。对 $L$ 中每个固定的 $\theta\ne0$，映射
	\[
	T_\theta:L\longrightarrow\mathbb Q,
	\qquad
	T_\theta(a)=\operatorname{Tr}_{L/\mathbb Q}(\theta a)
	\]
	是从 $L$ 到 $\mathbb Q$ 的 $\mathbb Q$-线性变换，其中 $\operatorname{Tr}_{L/\mathbb Q}$ 表示从 $L$ 到 $\mathbb Q$ 的迹映射（参见第 14.2 节练习 18）。这个线性变换非零，因为
	\[
	T_\theta(\theta^{-1})=\operatorname{Tr}_{L/\mathbb Q}(1)=m.
	\]
	于是，从 $L$ 到 $\operatorname{Hom}_{\mathbb Q}(L,\mathbb Q)$ 的映射 $\theta\mapsto T_\theta$ 是 $\mathbb Q$-向量空间的单同态。由于两边在 $\mathbb Q$ 上维数相同，该映射是同构。换言之，$L$ 上的每个线性泛函都形如 $T_\theta$。特别地，存在 $L$ 中元素 $\alpha_1',
	\ldots,
	\alpha_m'$，其相应线性变换 $T_{\alpha_i'}$ 给出 $\alpha_1,
	\ldots,
	\alpha_m$ 的对偶基，即
	\[
	\operatorname{Tr}_{L/\mathbb Q}(\alpha_i'\alpha_j)=
	\begin{cases}
		1,& i=j,\\
		0,& \text{否则}.
	\end{cases}
	\]
	由于 $\alpha_1',
	\ldots,
	\alpha_m'$ 线性无关，它们给出 $L$ 在 $\mathbb Q$ 上的一组基。因此每个元素 $\beta\in\mathcal O_L$ 可写成
	\[
	\beta=a_1\alpha_1'+a_2\alpha_2'+\cdots+a_m\alpha_m'
	\]
	其中 $a_1,
	\ldots,
	a_m\in\mathbb Q$。乘以 $\alpha_i$ 并取迹，得到
	\[
	\operatorname{Tr}_{L/\mathbb Q}(\beta\alpha_i)
	=a_1\operatorname{Tr}_{L/\mathbb Q}(\alpha_1'\alpha_i)+\cdots+a_i\operatorname{Tr}_{L/\mathbb Q}(\alpha_i'\alpha_i)+\cdots+a_m\operatorname{Tr}_{L/\mathbb Q}(\alpha_m'\alpha_i)=a_i.
	\]
	但 $\beta$ 与 $\alpha_i$ 都属于 $\mathcal O_L$，所以 $\beta\alpha_i$ 也属于 $\mathcal O_L$，这意味着
	\[
	a_i=\operatorname{Tr}_{L/\mathbb Q}(\beta\alpha_i)
	\]
	是 $\mathbb Z$ 中的元素（参见第 14.2 节练习 18(d)）。故
	\[
	\beta\in\mathbb Z\alpha_1'+\cdots+\mathbb Z\alpha_m',
	\]
	所以 $\mathcal O_L$ 包含在一个有限生成 $\mathbb Z$-模中，从而证明 $\mathcal O_K$（以及 $\mathcal O_L$）是自由 $\mathbb Z$-模。
	
	由于 $K$ 作为 $\mathbb Q$-向量空间的维数是 $n$，由第 12.1 节定理 5，$\mathcal O_K$ 是秩至多为 $n$ 的自由 $\mathbb Z$-模。又因为 $\mathcal O_K$ 含有秩为 $n$ 的自由 $\mathbb Z$-子模，所以 $\mathcal O_K$ 的 $\mathbb Z$-秩恰为 $n$，这证明 (1)，随后 (3) 的第二个断言也由关于 $\mathbb Z$-线性与 $\mathbb Q$-线性相关关系的前述说明得出。
	
	最后，$\mathcal O_K$ 中的任意理想 $I$ 都是秩为 $n$ 的自由 $\mathbb Z$-模的 $\mathbb Z$-子模，因此是秩至多为 $n$ 的自由 $\mathbb Z$-模；而 $I$ 的一组 $\mathbb Z$-模生成元同时也是 $\mathcal O_K$-生成元。因此 $\mathcal O_K$ 的每个理想都可由至多 $n$ 个元素生成，这说明 $\mathcal O_K$ 是 Noether 环，证明完毕。
\end{proof}

\noindent\textbf{定义.} 数域 $K$ 的\emph{整基}是 $K$ 的整数环作为秩为 $[K:\mathbb Q]$ 的自由 $\mathbb Z$-模的一组基。

若 $P$ 是数域 $K$ 的整数环 $\mathcal O_K$ 中的非零素理想，则 $P\cap\mathbb Z$ 是 $\mathbb Z$ 中的素理想。若 $\alpha\in P$，则 $\alpha$ 在 $\mathbb Q$ 上的最小多项式的常数项属于 $P\cap\mathbb Z$，这说明 $P\cap\mathbb Z=p\mathbb Z$ 也是 $\mathbb Z$ 中的非零素理想。由定理 26，$\mathbb Z$ 中的每个素理想 $(p)$ 都以这种方式出现。由于 $p\mathbb Z$ 是极大理想，由定理 26(2) 还可得，$\mathcal O_K$ 中的非零素理想都是极大理想；又由推论 27，对每个 $p\mathbb Z$，满足 $P\cap\mathbb Z=p\mathbb Z$ 的 $\mathcal O_K$ 中素理想 $P$ 只有有限多个。我们稍后将看到（第 16.3 节推论 16），数域整数环中的每个非零理想都能唯一地写成素理想的乘积；在理想 $p\mathcal O_K$ 的情形中，其不同素因子恰好就是满足 $P\cap\mathbb Z=p\mathbb Z$ 的有限多个 $\mathcal O_K$ 中理想 $P$。这一性质取代了 $\mathcal O_K$ 中元素分解为素元的唯一分解（后者不必成立，因为 $\mathcal O_K$ 不必是 U.F.D.）。我们还将看到，$\mathcal O_K$ 中的 primary 理想是素理想的幂（事实上这与 $\mathcal O_K$ 的理想唯一分解为素理想乘积等价，参见练习）。

\paragraph{例（二次扩张的整数环）}
若 $K$ 是 $\mathbb Q$ 的二次扩张，则 $K=\mathbb Q(\sqrt D)$，其中 $D$ 是某个无平方因子的整数。于是
\[
\mathcal O_{\mathbb Q(\sqrt D)}=\mathbb Z[\omega]=\mathbb Z\cdot1+\mathbb Z\cdot\omega,
\]
其整基为 $1,\omega$，其中
\[
\omega=
\begin{cases}
	\sqrt D,& D\equiv2,3\pmod4,\\[2mm]
	\dfrac{1+\sqrt D}{2},& D\equiv1\pmod4.
\end{cases}
\]
这正是第 7.1 节中引入的二次整数环。由于 $\omega$ 满足 $\omega^2-D=0$（当 $D\equiv2,3\pmod4$ 时），或满足 $\omega^2-\omega+(1-D)/4=0$（当 $D\equiv1\pmod4$ 时），可知 $\omega$ 是 $K$ 中的代数整数，所以 $\mathbb Z[\omega]\subseteq\mathcal O_K$。为证明这就是完整的整数环，令 $\alpha=a+b\sqrt D$，其中 $a,b\in\mathbb Q$，并假设 $\alpha$ 是代数整数。若 $b=0$，则 $\alpha\in\mathbb Q$，所以 $\alpha\in\mathbb Z$。若 $b\ne0$，则 $\alpha$ 的最小多项式是
\[
x^2-2ax+(a^2-b^2D).
\]
由命题 28，$2a$ 与 $a^2-b^2D$ 都是 $\mathbb Z$ 的元素。于是
\[
4(a^2-b^2D)=(2a)^2-(2b)^2D\in\mathbb Z,
\]
故 $4b^2D\in\mathbb Z$。由于 $D$ 是无平方因子的，推出 $2b$ 是整数。写 $a=x/2$，$b=y/2$，其中 $x,y$ 为整数。由于 $a^2-b^2D$ 是整数，有
\[
x^2-y^2D\equiv0\pmod4.
\]
因为模 $4$ 的平方只有 $0$ 与 $1$，并且 $D$ 不被 $4$ 整除，很容易检查只有下列可能：
\begin{enumerate}[label=(\roman*)]
	\item $D\equiv2$ 或 $3\pmod4$，并且 $x,y$ 都为偶数；或
	\item $D\equiv1\pmod4$，并且 $x,y$ 都为偶数或都为奇数。
\end{enumerate}
在情形 (i) 中，$a,b\in\mathbb Z$，故 $\alpha\in\mathbb Z[\omega]$。在情形 (ii) 中，
\[
\alpha=a+b\sqrt D=r+s\omega,
\]
其中 $r=(x-y)/2$，$s=y$ 都是整数，所以同样有 $\alpha\in\mathbb Z[\omega]$。

\paragraph{例（分圆域的整数环）}
$n$ 次单位根的分圆域 $\mathbb Q(\zeta_n)$ 的整数环是 $\mathbb Z[\zeta_n]$，其中 $\zeta_n$ 是任意一个本原 $n$ 次单位根。元素
\[
1,\zeta_n,\ldots,\zeta_n^{\varphi(n)-1}
\]
构成一个整基。显然 $\zeta_n$ 是代数整数，因为它是 $x^n-1$ 的根，所以环 $\mathbb Z[\zeta_n]$ 包含于整数环中。证明它就是 $\mathbb Q(\zeta_n)$ 中的全体代数整数，需要用到超出本书此处范围的代数数论技术。

\subsection*{Noether 正规化引理与 Hilbert 零点定理}

现在我们把整环扩张的代数理论中的一些技巧应用到仿射几何。

\noindent\textbf{定义.} 若 $k$ 是域，某个 $k$-代数中的元素 $y_1,y_2,
\ldots,y_q$ 称为在 $k$ 上\emph{代数无关}，若不存在 $q$ 个变量的非零多项式 $p$（系数在 $k$ 中）使得
\[
p(y_1,y_2,
\ldots,y_q)=0.
\]

因此，$y_1,y_2,
\ldots,y_q$ 代数无关，当且仅当从多项式环 $k[x_1,
\ldots,x_q]$ 到 $k[y_1,
\ldots,y_q]$、由 $x_i\mapsto y_i$ 定义的 $k$-代数同态是同构。$k$ 的一个域扩张中的元素代数无关，当且仅当它们是在 $k$ 上相互独立的超越元。

\noindent\textbf{定理 30（Noether 正规化引理）.} 设 $k$ 是域，并设
\[
A=k[r_1,r_2,
\ldots,r_m]
\]
是有限生成 $k$-代数。则对某个 $q$，$0\le q\le m$，存在 $A$ 中在 $k$ 上代数无关的元素 $y_1,y_2,
\ldots,y_q$，使得 $A$ 在 $k[y_1,y_2,
\ldots,y_q]$ 上整。

\begin{proof}
	对 $m$ 归纳。若 $r_1,
	\ldots,r_m$ 在 $k$ 上代数无关，则取 $y_i=r_i$，$i=1,
	\ldots,m$。否则，存在 $f(x_1,
	\ldots,x_m)\in k[x_1,
	\ldots,x_m]$，使得 $f(r_1,
	\ldots,r_m)=0$。多项式 $f$ 是若干形如 $ax_1^{e_1}x_2^{e_2}\cdots x_m^{e_m}$ 的单项式之和，该单项式的次数为 $e_1+\cdots+e_m$，而 $f$ 的次数 $d$ 是其单项式次数中的最大值。必要时重排变量，可假设 $f$ 是关于 $x_m$ 的非常数多项式，其系数在环 $k[x_1,x_2,
	\ldots,x_{m-1}]$ 中。现在我们作一次变量变换，把 $f$ 变为关于 $x_m$ 的首一多项式，其系数来自 $A$ 的一个由 $m-1$ 个元素在 $k$ 上生成的子环；于是便可应用归纳。
	
	定义整数 $c_i=(1+d)^i$，并对 $1\le i\le m-1$ 定义新变量
	\[
	X_i=x_i-x_m^{c_i}.
	\]
	令
	\[
	g(X_1,X_2,
	\ldots,X_{m-1},x_m)=f(X_1+x_m^{c_1},X_2+x_m^{c_2},
	\ldots,X_{m-1}+x_m^{c_{m-1}},x_m),
	\]
	所以 $g\in k[X_1,
	\ldots,X_{m-1},x_m]$。$f$ 的每个单项式项都给 $g$ 贡献一个常数倍的 $x_m^e$ 项。也容易检查，$c_i$ 的选择保证 $f$ 中不同单项式给出不同的 $e$ 值（例如把新变量中的次数看作以 $b=d+1$ 为底表示的整数）。若 $N$ 是出现的 $x_m$ 的最高次幂，则有
	\[
	g=cx_m^N+\sum_{i=0}^{N-1}h_i(X_1,
	\ldots,X_{m-1})x_m^i
	\]
	其中 $0\ne c\in k$。若现在 $s_i=r_i-r_m^{c_i}$，则
	\[
	\frac1c g(s_1,s_2,
	\ldots,s_{m-1},r_m)=\frac1c f(r_1,r_2,
	\ldots,r_{m-1},r_m)=0,
	\]
	这说明 $r_m$ 在 $B=k[s_1,
	\ldots,s_{m-1}]$ 上整。每个 $r_i$，$1\le i\le m-1$，都在 $B[r_m]$ 上整，因为 $r_i$ 是首一多项式 $x-s_i-r_m^{c_i}$ 的根。因此 $A$ 在 $B[r_m]$ 上整。由整性的传递性，$A$ 在 $B$ 上整。由于 $B$ 是由 $m-1$ 个元素生成的 $k$-代数，归纳完成证明。
\end{proof}

Noether 正规化引理的一个更“几何”的解释见练习 15。下面用正规化引理证明：若 $k$ 是代数闭域，则多项式环 $k[x_1,x_2,
\ldots,x_n]$ 的极大理想都形如
\[
(x_1-a_1,
\ldots,x_n-a_n)
\]
其中 $a_1,
\ldots,a_n\in k$。把 $k[x_1,x_2,
\ldots,x_n]$ 看作 $\mathbb A^n$ 上的多项式函数环时，这就是说极大理想对应于 $\mathbb A^n$ 中各点处的取值映射的核，这类似于紧集上连续函数环的相应结果（参见第 7.4 节练习 33、34）。

\noindent\textbf{定理 31（Hilbert 零点定理——弱形式）.} 设 $k$ 是代数闭域。则 $M$ 是多项式环 $k[x_1,x_2,
\ldots,x_n]$ 中的极大理想，当且仅当
\[
M=(x_1-a_1,
\ldots,x_n-a_n)
\]
对某些 $a_1,
\ldots,a_n\in k$ 成立。等价地，映射 $Z$ 与 $I$ 给出一一对应
\[
\{\mathbb A^n\text{ 中的点}\}
\;
ightleftarrows
\{k[\mathbb A^n]\text{ 中的极大理想}\}.
\]
此外，若 $I$ 是 $k[x_1,x_2,
\ldots,x_n]$ 的任意真理想，则 $Z(I)\ne\varnothing$。

\begin{proof}
	显然 $(x_1-a_1,
	\ldots,x_n-a_n)$ 是 $k[x_1,x_2,
	\ldots,x_n]$ 中的极大理想。反过来，对任意极大理想 $M\subseteq k[x_1,x_2,
	\ldots,x_n]$，令
	\[
	E=k[x_1,x_2,
	\ldots,x_n]/M.
	\]
	则 $E$ 是含有 $k$ 的域，且在 $k$ 上有限生成（由 $\bar x_1,
	\ldots,
	\bar x_n$ 生成）。由 Noether 正规化引理，$E$ 在某个多项式环 $k[y_1,
	\ldots,y_q]$ 上整。于是 $k[y_1,
	\ldots,y_q]$ 由定理 26(1) 是域；而一个或多个变量的多项式环永远不是域，故 $q=0$。因此 $E$ 在 $k$ 上整，所以 $E$ 在 $k$ 上代数。由于 $k$ 代数闭，$E=k$，即 $\bar x_i\in k$，$1\le i\le n$。于是对每个 $i=1,
	\ldots,n$，存在某个 $a_i\in k$，使得 $x_i-a_i\in M$。这意味着极大理想 $(x_1-a_1,
	\ldots,x_n-a_n)$ 包含于 $M$，故二者相等。最后，若 $I$ 是 $k[x_1,x_2,
	\ldots,x_n]$ 的任一真理想，则 $I$ 包含于某个极大理想
	\[
	M=(x_1-a_1,
	\ldots,x_n-a_n),
	\]
	于是 $(a_1,
	\ldots,a_n)\in Z(I)$。
\end{proof}

\noindent\textbf{定理 32（Hilbert 零点定理）.} 设 $k$ 是代数闭域。则对 $k[x_1,x_2,
\ldots,x_n]$ 的每个理想 $I$，有
\[
I(Z(I))=\operatorname{rad}I.
\]
此外，映射 $Z$ 与 $I$ 给出互逆的一一对应
\[
\{\text{仿射代数集}\}
\;
ightleftarrows
\{\text{根式理想}\}.
\]

\begin{proof}
	由于 $\operatorname{rad}I\subseteq I(Z(I))$，只需证明反向包含。由 Hilbert 基定理，$I=(f_1,f_2,
	\ldots,f_m)$。取 $g\in I(Z(I))$。引入一个新变量 $x_{n+1}$，并在 $k[x_1,
	\ldots,x_n,x_{n+1}]$ 中考虑由 $f_1,
	\ldots,f_m$ 与 $x_{n+1}g-1$ 生成的理想 $I'$。在 $\mathbb A^{n+1}$ 中任一点处，若 $f_1,
	\ldots,f_m$ 都为零，则 $g$ 也为零，因为 $g\in I(Z(I))$；因而 $x_{n+1}g-1$ 不为零。所以 $Z(I')=\varnothing$。由弱形式零点定理，$I'$ 不可能是真理想，即 $1\in I'$。于是
	\[
	1=a_1f_1+\cdots+a_mf_m+a_{m+1}(x_{n+1}g-1)
	\]
	其中 $a_i\in k[x_1,
	\ldots,x_{n+1}]$。令 $y=1/x_{n+1}$，并把此等式乘以足够高的 $y$ 的幂，可得
	\[
	y^N=c_1f_1+\cdots+c_mf_m+c_{m+1}(g-y)
	\]
	其中 $c_i\in k[x_1,
	\ldots,x_n,y]$。在这个多项式方程中把 $g$ 代入 $y$，得到 $g^N\in I$（在 $k[x_1,
	\ldots,x_n]$ 中），即 $g\in\operatorname{rad}I$。故 $I(Z(I))\subseteq\operatorname{rad}I$，从而 $I(Z(I))=\operatorname{rad}I$，证明完毕。
\end{proof}

由命题 12 与定理 26(2) 直接可得：若 $S$ 是 $R$ 的整扩张且 $1\in R$，若 $I$ 是 $R$ 的理想，则
\[
(\operatorname{rad}_S IS)\cap R=\operatorname{rad}_R I,
\]
其中 $IS$ 是 $S$ 中由 $I$ 生成的理想，下标表示在哪个环中计算根式。其几何解释如下。

\noindent\textbf{推论 33（Hilbert 零点定理的变形）.} 若 $k$ 是任意域，$\bar k$ 是其代数闭包，$I$ 是 $k[x_1,x_2,
\ldots,x_n]$ 中的理想，则
\[
I_k(Z_{\bar k}(I))=\operatorname{rad}I,
\]
其中 $Z_{\bar k}(I)$ 是 $I$ 中多项式在 $\bar k^n$ 中的零点集，$I_k(Z_{\bar k}(I))$ 是 $k[x_1,x_2,
\ldots,x_n]$ 中所有在 $Z_{\bar k}(I)$ 的点上为零的多项式所成的理想。特别地，$I=(1)$ 当且仅当 $I$ 中多项式在 $\bar k^n$ 中没有公共零点。

\begin{proof}
	因为 $\bar k[x_1,x_2,
	\ldots,x_n]$ 是 $k[x_1,x_2,
	\ldots,x_n]$ 的整扩张（由整元素 $\bar k$ 生成），该推论由定理 32 以及上面关于根式的说明立刻得到。
\end{proof}

由零点定理可得如下代数闭域 $k$ 上几何对象与环论对象之间的字典：
\[
\begin{array}{c|c}
	\text{几何} & \text{代数}\\ \hline
	\text{仿射代数集 }V & \text{坐标环 }k[V]\\
	V\text{ 的点} & k[V]\text{ 的极大理想}\\
	V\text{ 中的仿射代数子集} & k[V]\text{ 的根式理想}\\
	V\text{ 的子簇} & k[V]\text{ 的素理想}\\
	\text{态射 }\varphi:V\to W & k\text{-代数同态 }\bar\varphi:k[W]\to k[V]
\end{array}
\]

\subsection*{计算根式}

存在使用 Gröbner 基来计算多项式环中根式与 primary 分解的算法。虽然这些算法相对初等，但它们有些技术性，因此这里只讨论一些预备结果。

对于由单个多项式 $f\in k[x_1,
\ldots,x_n]$ 定义的超曲面 $V=Z(f)$，确定 $I(V)=\operatorname{rad}(f)$ 是直接的。由于 $k[x_1,
\ldots,x_n]$ 是 U.F.D.，$f$ 唯一地分解为两两非伴随不可约多项式幂的乘积：
\[
f=p_1^{a_1}\cdots p_s^{a_s},
\]
于是 $\operatorname{rad}(f)$ 由 $p_1\cdots p_s$（即 $f$ 的“无平方部分”）生成。

\paragraph{例}
设 $W=Z(J)$，其中 $J=(u^3-uv^2+v^3)\subseteq\mathbb Q[u,v]$。多项式 $x^3-x+1$ 在 $\mathbb Q$ 上不可约，所以 $f=u^3-uv^2+v^3$ 在 $\mathbb Q[u,v]$ 中不可约。因此 $\operatorname{rad}J=J$，且 $I(W)=J$。

对于非主理想 $I$，确定 $\operatorname{rad}I$ 更复杂。下面的命题（基于 Hilbert 零点定理）给出判定某元素是否属于 $\operatorname{rad}I$ 的准则。

\noindent\textbf{命题 34.} 设 $k$ 是任意域。若 $I=(f_1,
\ldots,f_s)$ 是 $k[x_1,
\ldots,x_n]$ 中的真理想，则
\[
f\in\operatorname{rad}I
\]
当且仅当
\[
(f_1,
\ldots,f_s,1-yf)=k[x_1,
\ldots,x_n,y].
\]

\begin{proof}
	由推论 33，$(f_1,
	\ldots,f_s,1-yf)=k[x_1,
	\ldots,x_n,y]$ 当且仅当方程组
	\[
	1-yf(x_1,
	\ldots,x_n)=0,
	\qquad
	f_1(x_1,
	\ldots,x_n)=0,
	\ldots,
	f_s(x_1,
	\ldots,x_n)=0
	\]
	在 $k$ 的代数闭包 $\bar k$ 上没有公共零点。给定 $(a_1,
	\ldots,a_n)\in\bar k^n$，方程 $1-yf(a_1,
	\ldots,a_n)=0$ 有解 $y$，除非 $f(a_1,
	\ldots,a_n)=0$。因此，该方程组没有公共零点，当且仅当对每个满足
	\[
	f_1(a_1,
	\ldots,a_n)=\cdots=f_s(a_1,
	\ldots,a_n)=0
	\]
	的 $(a_1,
	\ldots,a_n)\in\bar k^n$，也有 $f(a_1,
	\ldots,a_n)=0$。等价地，若 $(a_1,
	\ldots,a_n)\in Z_{\bar k}(I)$，则 $f(a_1,
	\ldots,a_n)=0$，即
	\[
	f\in I_k(Z_{\bar k}(I))=\operatorname{rad}I,
	\]
	其中最后等号由推论 33 得出。
\end{proof}

由于固定单项式序下理想的约化 Gröbner 基唯一，立即得到如下算法性方法来判定多项式是否属于某理想的根式。

\noindent\textbf{推论 35.} 设 $I=(f_1,
\ldots,f_s)\subseteq k[x_1,
\ldots,x_n]$。则 $f\in\operatorname{rad}I$ 当且仅当理想
\[
(f_1,
\ldots,f_s,1-yf)\subseteq k[x_1,
\ldots,x_n,y]
\]
相对于任意单项式序的约化 Gröbner 基为 $\{1\}$。

\paragraph{例}
考虑 $I=(x^2-y^2,xy)\subseteq k[x,y]$。在 $k[x,y,t]$ 中，对序 $x>y>t$，理想
\[
(x^2-y^2,xy,1-tx)
\]
的约化 Gröbner 基为 $\{1\}$，说明 $x\in\operatorname{rad}(I)$。为了确定 $I$ 中包含的 $x$ 的最小幂，我们发现 $k[x,y]$ 中理想
\[
(x^2-y^2,xy,x^3)
\]
与 $I$ 有相同的约化 Gröbner 基（即 $\{x^2-y^2,xy,y^3\}$），但 $(x^2-y^2,x^2,xy)$ 的基为 $\{x^2,xy,y^2\}$。由此可知 $x^3\in I$ 而 $x^2\notin I$（等价地，$x^2$ 除以 $\{x^2-y^2,xy,y^3\}$ 作一般多项式除法后有非零余式，而 $x^3$ 的余式为 $0$）。由类似计算（或由对称性），$y\in\operatorname{rad}I$，且 $y^3\in I$ 但 $y^2\notin I$。由于 $(x,y)\subseteq\operatorname{rad}I$，可得 $\operatorname{rad}I=(x,y)$。

关于计算根式的更多结果见练习。

\subsection*{练习}

设 $R$ 是交换环 $S$ 的子环，且 $1\in R$。

\begin{enumerate}[label=\arabic*.]
	\item 利用 U.F.D. 是整闭的这一事实，证明 Gaussian 整数环 $\mathbb Z[i]$ 是 $\mathbb Q(i)$ 中的整数环。
	\item 设 $k$ 是域，令 $t=x/y$ 为整环 $R=k[x,y]/(x^2-y^3)$ 的分式域中的元素。证明 $K=k(t)$ 是 $R$ 的分式域，并且 $k[t]$ 是 $R$ 在 $K$ 中的整闭包。
	\item 设 $k$ 是域，$i$ 与 $j$ 是互素正整数。求整环
	\[
	R=k[x,y]/(x^i-y^j)
	\]
	的正规化（参见第 9.1 节练习 14）。
	\item 设 $k$ 是域，令 $P$ 为多项式环 $k[x,y]$ 中的理想 $(y^2-x^3-x^2)$。证明 $P$ 是素理想，并求整环 $R=k[x,y]/P$ 的正规化。[为证明 $P$ 是素理想，说明 $y^2-x^3-x^2$ 在 U.F.D. $k[x,y]$ 中不可约。然后考虑 $t=y/\bar x\in R$。]
	\item 若 $R$ 是整环，分式域为 $F$，证明 $F$ 是有限生成 $R$-模，当且仅当 $R=F$。
	\item 对下列每一种指定关系，给出具体环 $R\subseteq S$ 以及这些环中的显式理想来展示：
	\begin{enumerate}[label=(\alph*)]
		\item $R$ 的一个理想 $I$，使得 $I\ne IS\cap R$（所以理想的扩张的收缩不必等于 $I$）；
		\item $R$ 的一个素理想 $P$，使得不存在 $S$ 的素理想 $Q$ 满足 $P=Q\cap R$；
		\item $S$ 的一个极大理想 $M$，使得 $M\cap R$ 在 $R$ 中不是极大的；
		\item $R$ 的一个素理想 $P$，其到 $S$ 的扩张 $PS$ 在 $S$ 中不是素理想；
		\item $S$ 的一个理想 $J$，使得 $J\ne (J\cap R)S$（所以理想的收缩的扩张不必等于 $J$）。
	\end{enumerate}
	\item 令 $\mathcal O_K$ 为数域 $K$ 中的整数环。
	\begin{enumerate}[label=(\alph*)]
		\item 假设 $\mathcal O_K$ 的每个非零理想 $I$ 都能写成素理想幂的乘积。证明 $\mathcal O_K$ 的一个理想 $Q$ 是 $P$-primary，当且仅当 $Q=P^m$ 对某个 $m\ge1$ 成立。[先证明：由于 $\mathcal O_K$ 中非零素理想都是极大的，若 $P_1$ 与 $P_2$ 是不同非零素理想，则 $P_1^{m_1}\subseteq P_2^{m_2}$ 蕴含 $P_1=P_2$。]
		\item 假设 $\mathcal O_K$ 的理想 $Q$ 是 $P$-primary，当且仅当 $Q=P^m$ 对某个 $m\ge1$ 成立。假定定理 21 的全部内容，证明 $\mathcal O_K$ 的每个非零理想 $I$ 都能唯一地写成素理想幂的乘积。[证明若 $P_1$ 与 $P_2$ 是不同非零素理想，则 $P_1^{m_1}$ 与 $P_2^{m_2}$ 是互素理想，并使用中国剩余定理。]
	\end{enumerate}
	\item 证明若 $s_1,
	\ldots,s_n\in S$ 都在 $R$ 上整，则环 $R[s_1,
	\ldots,s_n]$ 是有限生成 $R$-模。
	\item 假设 $S$ 在 $R$ 上整，且 $P$ 是 $R$ 中的素理想。证明 $S$ 中由 $P$ 生成的理想 $PS$ 的每个元素 $s$ 都满足方程
	\[
	s^n+a_{n-1}s^{n-1}+\cdots+a_1s+a_0=0,
	\]
	其中系数 $a_0,a_1,
	\ldots,a_{n-1}$ 都属于 $P$。[若 $s=p_1s_1+\cdots+p_ms_m\in PS$，说明 $T=R[s_1,
	\ldots,s_m]$ 满足命题 23(3) 中的假设。沿用命题 23 中证明 $s$ 整的证明，并注意 $s\in PT$，从而 $a_{ij}$ 都是 $P$ 的元素。]
	\item 证明命题 28 的下列推广：设 $R$ 是整闭整环，分式域为 $k$，$\alpha$ 是 $k$ 的某个扩域 $K$ 中的元素。证明 $\alpha$ 在 $R$ 上整，当且仅当 $\alpha$ 在 $k$ 上代数，并且 $\alpha$ 在 $k$ 上的最小多项式 $m_{\alpha,k}(x)$ 的系数属于 $R$。[若 $\alpha$ 整，证明 $\alpha$ 的共轭元，即 $m_{\alpha,k}(x)$ 的根，也都整；于是共轭元的基本对称函数是 $k$ 中在 $R$ 上整的元素。]
	\item 设 $R$ 是整闭整环，分式域为 $k$，并设 $p(x)\in R[x]$ 是首一多项式。证明若
	\[
	p(x)=a(x)b(x)
	\]
	其中 $a(x),b(x)\in k[x]$ 是首一多项式，则 $a(x),b(x)\in R[x]$（与第 9.3 节命题 5，即 Gauss 引理比较）。[参见前一练习。]
	\item 设 $S$ 是整环，并且在前一练习中的环 $R$ 上整。若 $P$ 是 $R$ 中的素理想，令 $s$ 为 $S$ 中由 $P$ 生成的理想 $PS$ 中的任意元素。证明除首项之外，$s$ 在 $k$ 上的最小多项式 $m_{s,k}(x)$ 的系数都属于 $P$。[由练习 10，$m_{s,k}(x)\in R[x]$。练习 9 说明 $s$ 是首一多项式
	\[
	p(x)=x^n+a_{n-1}x^{n-1}+\cdots+a_0
	\]
	的根，其中 $a_0,
	\ldots,a_{n-1}\in P$。使用前一练习说明 $p(x)=m_{s,k}(x)b(x)$，其中 $b(x)\in R[x]$，并在整环 $(R/P)[x]$ 中考虑这个等式。]
	
	接下来的两个练习推广第 7.5 节练习 6，刻画哪些域不是它们任何真子环的分式域。
	
	\item 设 $K$ 是特征 $0$ 的域，$A$ 是 $K$ 的一个子环，并且在满足 $1/2\notin A$ 这一条件下极大。（这样的 $A$ 由 Zorn 引理存在。）令 $F$ 为 $A$ 在 $K$ 中的分式域。
	\begin{enumerate}[label=(\alph*)]
		\item 证明 $K$ 在 $F$ 上代数。[若 $t$ 在 $F$ 上超越，证明 $1/2\notin A[t]$。]
		\item 证明 $A$ 在 $K$ 中整闭。[证明 $1/2$ 不在 $A$ 在 $K$ 中的整闭包中。]
		\item 由 (a) 与 (b) 推出 $K=F$。
	\end{enumerate}
	\item 证明：域 $K$ 是 $K$ 的某个真子环的分式域，当且仅当 $K$ 不是某个有限域的代数闭包的子域。[若 $K$ 含有在 $\mathbb F_p$ 上超越的 $t$，仿照前一练习，用 $1/t$ 代替 $1/2$，证明 $K$ 是某个真子环的商域。]
	
	下一个练习给出 Noether 正规化引理的“几何”解释，说明每个仿射代数集都是某个仿射 $n$-空间的有限覆盖。
	
	\item 设 $V$ 是代数闭域 $k$ 上的仿射代数集。证明对某个 $n$，存在从 $V$ 到 $\mathbb A^n$ 的满态射，其纤维有限；若 $V$ 是簇，则 $n$ 可取为 $V$ 的维数。[由 Noether 正规化引理，有限生成 $k$-代数 $S=k[V]$ 含有一个多项式子代数 $R=k[x_1,x_2,
	\ldots,x_n]$，使得 $S$ 在 $R$ 上整。把定理 6 应用于 $R$ 到 $S$ 的包含，得到从 $V$ 到 $\mathbb A^n$ 的态射 $\varphi$。为说明 $\varphi$ 满且纤维有限，把推论 27 应用于 $R$ 中对应点 $(a_1,
	\ldots,a_n)\in\mathbb A^n$ 的极大理想 $(x_1-a_1,
	\ldots,x_n-a_n)$。]
	\item 设 $V$ 是 $\mathbb C^n$ 中的仿射代数集。证明 $V$ 在 Euclidean 拓扑中紧（即闭且有界），当且仅当它有限。[使用第 15.2 节练习 18、前一练习以及紧集在连续函数下的行为。]
	\item 设 $R$ 是交换环 $S$ 的子环，且 $1_S\in R$，并假设 $S$ 在 $R$ 上整。本练习证明 $R$ 与 $S$ 有相同的 Krull 维数，参见第 16.1 节。
	\begin{enumerate}[label=(\alph*)]
		\item 若
		\[
		P_1\subset P_2\subset\cdots\subset P_n
		\]
		是 $R$ 中不同素理想的链，证明存在 $S$ 中不同素理想的链
		\[
		Q_1\subset Q_2\subset\cdots\subset Q_n
		\]
		使得 $Q_i\cap R=P_i$。
		\item 反过来，若
		\[
		Q_1\subset Q_2\subset\cdots\subset Q_n
		\]
		是 $S$ 中不同素理想的链，且 $P_i=Q_i\cap R$，证明
		\[
		P_1\subset P_2\subset\cdots\subset P_n
		\]
		是 $R$ 中不同素理想的链。[为证明 $P_i$ 不同，转到商环并把问题化为证明：若 $Q$ 是整环 $S$ 中的非零素理想，则 $Q\cap R$ 是 $R$ 中非零素理想。在此情形中，若 $0\ne s\in Q$，说明 $s$ 所满足的 $R[x]$ 中次数最小的多项式的常数项是 $Q\cap R$ 中的一个非零元素。]
	\end{enumerate}
	\item 令 $V=Z(I)$ 与 $W=Z(J)$，其中
	\[
	I=(uv+v)\subseteq\mathbb C[u,v],
	\qquad
	J=(-2y-y^2+2z+z^2,
	2x-yz-z^2)\subseteq\mathbb C[x,y,z].
	\]
	\begin{enumerate}[label=(\alph*)]
		\item 证明 $I$ 与 $J$ 是素理想。推出 $I=I(V)$、$J=I(W)$，并且 $V$ 与 $W$ 是簇。
		\item 证明映射 $\varphi:V\to W$，由
		\[
		\varphi((a_1,a_2))=(a_1^2+a_2,
		a_1+a_2,
		a_1-a_2)
		\]
		定义，是同构。
	\end{enumerate}
	\item 设
	\[
	I=(x^3+y^3+z^3,
	x^2+y^2+z^2,
	(x+y+z)^3)\subseteq k[x,y,z].
	\]
	使用 Gröbner 基证明：若 $\operatorname{ch}(k)\ne2,3$，则 $x,y,z\in\operatorname{rad}I$。
	\item 设
	\[
	I=(x^3+y^3+z^3,
	xy+xz+yz,
	xyz)\subseteq k[x,y,z].
	\]
	使用 Gröbner 基证明 $x,y,z\in\operatorname{rad}I$。
	\item 设
	\[
	I=(x^4+y^4+z^4,
	x+y+z)\subseteq k[x,y,z].
	\]
	\begin{enumerate}[label=(\alph*)]
		\item 使用 Gröbner 基证明：若 $\operatorname{ch}(k)\ne2$，则 $xy+xz+yz\in\operatorname{rad}I$，并确定 $I$ 中包含的 $xy+xz+yz$ 的最小幂。证明 $x,y,z$ 均不属于 $\operatorname{rad}I$。
		\item 若 $J=(x^4+y^4+z^4,
		x+y+z,
		xy+xz+yz)$，证明相对于字典序 $x>y>z$，$J$ 的约化 Gröbner 基为 $\{x+y+z,
		y^2+yz+z^2\}$。推出
		\[
		k[x,y,z]/J\simeq k[y,z]/(y^2+yz+z^2),
		\]
		并且若 $\operatorname{ch}(k)\ne3$，则 $J$ 是根式理想。
		\item 若 $\operatorname{ch}(k)\ne2,3$，证明 $\operatorname{rad}I=J$。
		\item 若 $\operatorname{ch}(k)=3$，证明 $\operatorname{rad}I=(x-y,y-z)$。
		\item 若 $\operatorname{ch}(k)=2$，证明 $I=(x+y+z)$ 是素理想，因而是根式理想。
	\end{enumerate}
	\item 设
	\[
	I=(x^2y+z^3,
	x+y^3-z,
	2y^4z-yz^2-z^3)\subseteq k[x,y,z].
	\]
	使用 Gröbner 基证明 $x,y,z\in\operatorname{rad}I$，并推出 $\operatorname{rad}I=(x,y,z)$。证明 $x^9,y^7,z^9$ 分别是 $x,y,z$ 落在 $I$ 中的最小幂。
	\item 令
	\[
	V=Z(x^3-x^2z-y^2z),
	\qquad
	W=Z(x^2+y^2-z^2)
	\]
	为 $\mathbb C^3$ 中的代数集。证明
	\[
	I(V)=(x^3-x^2z-y^2z),
	\qquad
	I(W)=(x^2+y^2-z^2)
	\]
	在 $\mathbb C[x,y,z]$ 中成立。
	\item 令 $V=Z(x^3+y^3+7z^3)\subseteq\mathbb C^3$。证明
	\[
	I(V)=(x^3+y^3+7z^3)
	\]
	在 $\mathbb C[x,y,z]$ 中成立。
	\item 设
	\[
	I=(xz+y^2+z^2,
	xy-xz+yz-2z^2)
	\]
	且
	\[
	K=I+(x^2-3y^2+yz)\subseteq\mathbb C[x,y,z].
	\]
	\begin{enumerate}[label=(\alph*)]
		\item 由第 15.1 节练习 46，存在从 $\mathbb C[x,y,z]/K$ 到 $\mathbb C[u,v]/(u^3-uv^2+v^3)$ 的单的 $\mathbb C$-代数同态。利用这一点以及命题 34 前的例子，证明 $K$ 是根式理想，并推出 $\operatorname{rad}I\subseteq K$。
		\item 证明 $\operatorname{rad}I\subseteq (y,z)$。
		\item 证明 $K\cap(y,z)=I$，并推出 $I$ 是根式理想，所以若 $V=Z(I)$，则 $I(V)=I$。
		\item 证明 $y(x^2-3y^2+yz)$ 与 $z(x^2-3y^2+yz)$ 属于 $I$，但 $y,z$ 或 $x^2-3y^2+yz$ 中没有一个属于 $I$。
	\end{enumerate}
	\item 设 $I$ 是 $k[x_1,
	\ldots,x_n]$ 中的理想。证明下列条件等价（满足任一条件的理想称为\emph{零维理想}，原因见 (d)）：
	\begin{enumerate}[label=(\alph*)]
		\item 商 $k[x_1,
		\ldots,x_n]/I$ 作为 $k$-向量空间维数有限。
		\item 对每个 $i=1,2,
		\ldots,n$，$I\cap k[x_i]\ne0$。
		\item 若 $G$ 是 $I$ 的任意约化 Gröbner 基，则对每个 $i=1,
		\ldots,n$，存在 $g_i\in G$，其首项为 $x_i^{n_i}$，其中 $n_i\ge1$。
		\item $I$ 中多项式在 $k$ 的代数闭包 $\bar k$ 中的公共零点集 $Z_{\bar k}(I)$ 是有限集。
	\end{enumerate}
	[为证明 (a) 蕴含 (b)，使用嵌入 $k[x_i]/(I\cap k[x_i])\hookrightarrow k[x_1,
	\ldots,x_n]/I$。为证明 (b) 蕴含 (c)，注意某个 $\operatorname{LT}(g_i)$ 整除 $I\cap k[x_i]$ 的一个生成元的首项。为证明 (c) 蕴含 (a)，使用第 9.6 节练习 37。证明 (b) 蕴含 (d)。为证明 (d) 蕴含 (b)，说明点 $Z_{\bar k}(I)$ 的第 $i$ 个坐标 $a_1,
	\ldots,a_N$ 的最小多项式的乘积
	\[
	m_{a_1,k}(x_i)
	\cdots m_{a_N,k}(x_i)
	\]
	是 $I(Z_{\bar k}(I))$ 中的非零多项式，并应用推论 33。]
	\item 设 $I$ 是 $k[x_1,
	\ldots,x_n]$ 中的零维理想，令 $I'$ 为 $I$ 在 $\bar k[x_1,
	\ldots,x_n]$ 中生成的理想，其中 $\bar k$ 是 $k$ 的代数闭包。令 $Z(I)$ 为 $I$ 在 $k^n$ 中的零点集，$Z_{\bar k}(I)$ 为 $I$（等价地 $I'$）在 $\bar k^n$ 中的零点集。
	\begin{enumerate}[label=(\alph*)]
		\item 证明
		\[
		|Z_{\bar k}(I)|=\dim_{\bar k}\bar k[x_1,
		\ldots,x_n]/\operatorname{rad}I'.
		\]
		[说明 $\operatorname{rad}I'$ 是与 $V_{\bar k}$ 中点对应的极大理想的乘积，并使用中国剩余定理。]
		\item 证明
		\[
		|Z(I)|\le \dim_k k[x_1,
		\ldots,x_n]/I.
		\]
		[一种方法：使用第 15.1 节练习 43，并观察
		\[
		\dim_{\bar k}\bar k[x_1,
		\ldots,x_n]/\operatorname{rad}I'\le
		\dim_{\bar k}\bar k[x_1,
		\ldots,x_n]/I'.
		\]]
	\end{enumerate}
	\item 设 $I$ 是 $k[x_1,
	\ldots,x_n]$ 中的零维理想，并假设 $I\cap k[x_i]$ 由非零多项式 $h_i$ 生成（参见练习 26）。令 $r_i$ 为 $h_i$ 的不可约因子的乘积（即 $h_i$ 的“无平方部分”）。
	\begin{enumerate}[label=(\alph*)]
		\item 证明
		\[
		I+(r_1,
		\ldots,r_n)\subseteq\operatorname{rad}I.
		\]
		\item （完美域上零维理想的根式）若 $k$ 是完美域，证明
		\[
		\operatorname{rad}I=I+(r_1,
		\ldots,r_n).
		\]
		[对 $n$ 归纳。写 $r_1=p_1\cdots p_t$，其中 $p_i$ 是 $k[x_1]$ 中互异不可约多项式。若 $J=I+(r_1,
		\ldots,r_n)$，证明 $J=J_1\cap\cdots\cap J_t$，其中 $J_i=J+(p_i)$。证明对每个 $i$，模 $p_i$ 的约化诱导同构
		\[
		k[x_1,
		\ldots,x_n]/J_i\simeq K[x_2,
		\ldots,x_n]/J_i',
		\]
		其中 $K$ 是扩域 $k[x]/(p_i)$，$J_i'\subseteq K[x_2,
		\ldots,x_n]$ 是 $J_i$ 模 $(p_i)$ 的约化。使用第 13.5 节练习 11 说明，由于 $k$ 完美，对每个 $j=2,
		\ldots,n$，$r_j$ 在 $J_i'\cap K[x_j]$ 中的像仍为非零无平方多项式。由归纳推出 $J_i'$ 是根式理想。推出 $J_i$ 是根式理想，最后推出 $J$ 是根式理想。]
		\item 求 $\mathbb Q[x,y]$ 中理想
		\[
		(x^7+x+y^3,
		x^4+y^3+y),
		\quad
		(x^3-xy^2+x,
		x^2y+y^3),
		\quad
		(x^4+y^3,
		x^3-xy+y^2)
		\]
		的根式，以及 $\mathbb Q[x,y,z]$ 中理想
		\[
		(x^2+y^2z,
		x^2y^2+z^3,
		y^2+z^2)
		\]
		的根式。
		\item 令 $k=\mathbb F_p(t)$。说明
		\[
		I=(x^p+t,
		y^p-t)
		\]
		是 $k[x,y]$ 中的零维理想，并且 $I\cap k[x]$ 与 $I\cap k[y]$ 都含有非零无平方多项式，但 $I$ 不是根式理想（所以若 $k$ 不完美，(b) 中的结果不必成立）。[证明 $x+y\in\operatorname{rad}I$ 但 $x+y\notin I$。]
	\end{enumerate}
\end{enumerate}
\section{局部化}

环中“在素理想处局部化”的思想，是代数中用来分离环中理想行为的一个极其有力而且无处不在的工具。它是熟知的“在一点附近局部化”思想的代数类似物，例如在考虑实直线上函数 $f(x)$ 的可微性问题时便会用到这种思想。事实上，这一技术的重要应用之一（也是其发展的原始动机之一）是把仿射代数空间几何中的这类“局部”性质翻译成其坐标环的相应性质。

我们先考虑“分式环”的一个很一般的构造。设 $D$ 是 $R$ 的含 $1$ 的乘法闭子集（即 $1\in D$，且若 $a,b\in D$ 则 $ab\in D$）。下一个结果构造出一个新环 $D^{-1}R$，它是在其中 $D$ 的元素变成单位的“最小”环。这推广了第 7.5 节中分式环的构造，因为现在允许 $D$ 含有零或零因子，所以在这种情形下 $R$ 不必作为子环嵌入 $D^{-1}R$。

\noindent\textbf{定理 36.} 设 $R$ 是含 $1$ 的交换环，$D$ 是 $R$ 的含 $1$ 的乘法闭子集。则存在交换环 $D^{-1}R$ 及环同态
\[
\pi:R\longrightarrow D^{-1}R,
\]
满足如下泛性质：对任意把 $1$ 送到 $1$ 的交换环同态 $\psi:R\to S$，若对每个 $d\in D$，$\psi(d)$ 都是 $S$ 中的单位，则存在唯一同态 $\Psi:D^{-1}R\to S$，使得
\[
\Psi\circ\pi=\psi.
\]

\begin{proof}
	证明与第 7.5 节定理 15 的证明非常类似。在这里，在 $R\times D$ 上定义关系
	\[
	(r,d)\sim(s,e)
	\quad\Longleftrightarrow\quad
	x(er-ds)=0\quad\text{对某个 }x\in D\text{ 成立}.
	\]
	这个关系显然是自反和对称的。若 $(r,d)\sim(s,e)$ 且 $(s,e)\sim(t,f)$，则对某些 $x,y\in D$ 有
	\[
	x(er-ds)=0,
	\qquad
	y(fs-et)=0.
	\]
	将第一个方程乘以 $fy$，第二个方程乘以 $dx$，再相加，得到
	\[
	exy(fr-dt)=0.
	\]
	由于 $D$ 对乘法封闭，故 $(r,d)\sim(t,f)$，所以 $\sim$ 是传递的。
	
	记 $(r,d)$ 关于 $\sim$ 的等价类为 $r/d$，并令 $D^{-1}R$ 为这些等价类的集合。在 $D^{-1}R$ 中定义加法与乘法为
	\[
	\frac{a}{b}+\frac{c}{d}=\frac{ad+bc}{bd},
	\qquad
	\frac{a}{b}\cdot\frac{c}{d}=\frac{ac}{bd}.
	\]
	检查这些运算良定义并使 $D^{-1}R$ 成为交换环，且 $1=1/1$，这是一个练习。对每个 $d\in D$，$d/1$ 是 $D^{-1}R$ 中的单位（即便在 $D^{-1}R$ 是零环的退化情形下也是如此）。
	
	最后定义 $\pi:R\to D^{-1}R$ 为 $\pi(r)=r/1$。容易看出 $\pi$ 是环同态。设 $\psi:R\to S$ 是一个把 $1$ 送到 $1$ 的交换环同态，并且对每个 $d\in D$，$\psi(d)$ 都是 $S$ 中的单位。定义
	\[
	\Psi:D^{-1}R\longrightarrow S,
	\qquad
	\Psi\left(\frac{r}{d}\right)=\psi(r)\psi(d)^{-1}.
	\]
	这个映射良定义：若 $r/d=s/e$，则对某个 $x\in D$ 有 $x(er-ds)=0$。于是
	\[
	\psi(x)\bigl(\psi(er)-\psi(ds)\bigr)=0
	\]
	在 $S$ 中成立；由于 $\psi(x)$ 是 $S$ 中的单位，得到 $\psi(er)-\psi(ds)=0$，从而
	\[
	\psi(r)\psi(d)^{-1}=\psi(s)\psi(e)^{-1}.
	\]
	显然 $\Psi$ 是环同态，且 $\Psi\circ\pi=\psi$。
	
	最后，$\Psi$ 是唯一的，因为 $D^{-1}R$ 的每个元素都可以写成积
	\[
	(r/1)(d/1)^{-1}.
	\]
	$\Psi$ 在每个形如 $x/1$ 的元素上的值由 $\psi$ 唯一决定，即
	\[
	\Psi(x/1)=\Psi(\pi(x))=\psi(x).
	\]
	由于 $\Psi$ 是环同态，它在任意单位 $u$ 的逆 $u^{-1}$ 上的值由 $\Psi(u)$ 唯一决定。因此 $\Psi$ 在 $D^{-1}R$ 的每个元素上的值都唯一确定，证明完毕。
\end{proof}

\noindent\textbf{推论 37.} 在定理 36 的记号下：
\begin{enumerate}[label=(\arabic*)]
	\item
	\[
	\ker\pi=\{r\in R\mid xr=0\text{ 对某个 }x\in D\text{ 成立}\};
	\]
	特别地，$\pi:R\to D^{-1}R$ 是单射当且仅当 $D$ 不含 $R$ 的零因子；
	\item $D^{-1}R=0$ 当且仅当 $0\in D$，从而当且仅当 $D$ 含有幂零元。
\end{enumerate}

\begin{proof}
	由定义，$\pi(r)=0$ 当且仅当 $(r,1)\sim(0,1)$，也即当且仅当对某个 $x\in D$ 有 $xr=0$，这给出 (1)。为证明 (2)，注意 $D^{-1}R=0$ 当且仅当此环的 $1$ 为零，也即 $(1,1)\sim(0,1)$。这发生当且仅当对某个 $x\in D$ 有 $x\cdot1=0$，也即当且仅当 $0\in D$。
\end{proof}

\noindent\textbf{定义.} 环 $D^{-1}R$ 称为 $R$ 关于 $D$ 的\emph{分式环}，或称为 $R$ 在 $D$ 处的\emph{局部化}。

\paragraph{例}
\begin{enumerate}[label=(\arabic*)]
	\item 设 $R$ 是整环，$D=R-\{0\}$。则 $D^{-1}R$ 是第 7.5 节中描述的 $R$ 的分式域 $Q$。更一般地，若 $D$ 是 $R-\{0\}$ 的任意乘法闭子集，则 $D^{-1}R$ 是 $Q$ 的子环，由所有形如 $r/d$ 的元素组成，其中 $r\in R$ 且 $d\in D$。
	
	\item 设 $R$ 是任意含 $1$ 的交换环，$f$ 是 $R$ 的任意元素。令 $D$ 为 $R$ 中 $f$ 的非负次幂组成的乘法集
	\[
	\{f^n\mid n\ge0\}.
	\]
	定义 $R_f=D^{-1}R$。注意 $R_f=0$ 当且仅当 $f$ 是幂零元。若 $f$ 不是幂零元，则 $f$ 在 $R_f$ 中成为单位。不难看出
	\[
	R_f\simeq R[x]/(xf-1),
	\]
	其中 $R[x]$ 是变量 $x$ 上的多项式环（参见练习）。又注意，对任意 $n\ge1$，$R_f$ 与 $R_{f^n}$ 自然同构，因为 $f$ 与 $f^n$ 在这两个环中都是单位。若 $f$ 是零因子，则 $\pi:R\to R_f$ 不把 $R$ 嵌入 $R_f$。例如，令 $R=k[x,y]/(xy)$，并取 $f=x$。则 $x$ 在 $R_x$ 中是单位，而由推论的第一部分，$y$ 被映到 $0$（显式地，$y=xy/x=0$ 在 $R_x$ 中）。在这种情形下，$\pi(R)=k[x]\subset R_f=k[x,x^{-1}]$。
	
	\item （在素理想处局部化）令 $P$ 为任意环 $R$ 中的素理想，并令 $D=R-P$。由素理想的定义，$D$ 是乘法闭的。在这种情形下过渡到环 $D^{-1}R$ 称为把 $R$ 在 $P$ 处局部化，并且 $D^{-1}R$ 记作 $R_P$。$R$ 中不属于 $P$ 的每个元素都在 $R_P$ 中成为单位。例如，若 $R=\mathbb Z$ 且 $P=(p)$ 是素理想，则
	\[
	\mathbb Z_{(p)}=\left\{\frac{a}{b}\in\mathbb Q\mid p\nmid b\right\}\subseteq\mathbb Q,
	\]
	并且每个不能被 $p$ 整除的整数 $b$ 都是单位。
	
	\item 若 $V$ 是任意非空集合，$k$ 是域，令 $R$ 为 $V$ 上任意一个含常值函数的 $k$-值函数环（例如闭区间 $[0,1]$ 上所有连续实值函数组成的环）。对任意 $a\in V$，令 $M_a$ 为 $R$ 中在 $a$ 处为零的函数构成的理想。则 $M_a$ 是由在 $a$ 处取值给出的从 $R$ 到域 $k$ 的环同态的核。由于 $R$ 含有常值函数，取值映射是满射，因此 $M_a$ 是极大理想（从而也是素理想）。在这个素理想处对 $R$ 局部化得到
	\[
	R_{M_a}=\left\{\frac{f}{g}\mid f,g\in R,\ g(a)\ne0\right\}.
	\]
	$R_{M_a}$ 中每个函数都可在 $a$ 处取值，定义为 $(f/g)(a)=f(a)/g(a)$，并且这个值不依赖于代表元 $f/g$ 的选择，因此 $R_{M_a}$ 成为在 $a$ 处有定义的 $k$-值“有理函数”环。
\end{enumerate}

接下来考虑关于定理 36 中映射 $\pi:R\to D^{-1}R$ 的理想扩张与收缩。为简化记号，若 $I$ 是 $R$ 的理想，则用 $eI$ 表示 $I$ 到 $D^{-1}R$ 的扩张（而不用更繁琐的 $D^{-1}R\pi(I)$）；若 $J$ 是 $D^{-1}R$ 的理想，则用 $cJ$ 表示 $J$ 到 $R$ 的收缩。

若 $I$ 是 $R$ 的理想，则容易看出 $eI$ 的每个元素都可以写成 $a/d$ 的形式，其中 $a\in I$ 且 $d\in D$，所以 $I$ 到 $D^{-1}R$ 的扩张也常记作 $D^{-1}I$。

\noindent\textbf{命题 38.} 在上述记号下，有：
\begin{enumerate}[label=(\arabic*)]
	\item 对任意 $D^{-1}R$ 的理想 $J$，有 $J=e(cJ)$。特别地，$D^{-1}R$ 的每个理想都是 $R$ 中某个理想的扩张，并且 $D^{-1}R$ 的不同理想在 $R$ 中有不同收缩。
	\item 对任意 $R$ 的理想 $I$，有
	\[
	c(eI)=\{r\in R\mid dr\in I\text{ 对某个 }d\in D\text{ 成立}\}.
	\]
	并且 $eI=D^{-1}R$ 当且仅当 $I\cap D\ne\varnothing$。
	\item 扩张与收缩给出双射对应
	\[
	\left\{\begin{array}{c}
		R\text{ 的素理想 }P\\
		\text{且 }P\cap D=\varnothing
	\end{array}\right\}
	\overset{e}{\underset{c}{\rightleftarrows}}
	\{D^{-1}R\text{ 的素理想}\}.
	\]
	\item 若 $R$ 是 Noether 环（或 Artin 环），则 $D^{-1}R$ 也是 Noether 环（相应地，Artin 环）。
\end{enumerate}

\begin{proof}
	我们总有 $e(cJ)\subseteq J$。为证明反包含，令 $a/d\in J$。则
	\[
	a/1=d(a/d)\in J,
	\]
	所以 $a\in\pi^{-1}(J)=cJ$。于是 $a/1\in e(cJ)$，从而
	\[
	(a/1)(1/d)=a/d\in e(cJ),
	\]
	故 $J=e(cJ)$。这证明了 (1) 的第一个断言，第二个断言立即得出。
	
	令
	\[
	I'=\{r\in R\mid dr\in I\text{ 对某个 }d\in D\text{ 成立}\}.
	\]
	先证明 $I'\subseteq c(eI)$。若 $r\in I'$，则存在 $d\in D$ 使 $dr=a\in I$。于是 $r/1=a/d\in eI$，所以 $r\in c(eI)$。为证明反包含 $c(eI)\subseteq I'$，设 $r\in c(eI)$，则 $r/1=a/d$，其中 $a\in I$，$d\in D$。于是对某个 $x\in D$ 有 $x(dr-a)=0$，从而 $xdr=xa\in I$；由于 $xd\in D$，得到 $r\in I'$。这证明了 (2) 的第一个断言。现在 $eI=D^{-1}R$ 当且仅当 $1/1\in eI$，当且仅当 $1\in c(eI)=I'$。由 $I'$ 的定义即得 (2) 的第二个断言。
	
	为证明 (3)，先注意若 $Q$ 是 $D^{-1}R$ 中的素理想，则它在任意把 $1$ 送到 $1$ 的同态下的原像是素理想（参见第 7.4 节练习 13），所以 $c$ 把 $D^{-1}R$ 的素理想映到与 $D$ 不相交的 $R$ 的素理想。反过来，令 $P$ 是 $R$ 的与 $D$ 不相交的素理想，并令 $Q=eP$。假设
	\[
	(a/d_1)(b/d_2)\in Q.
	\]
	则 $(ab)/(d_1d_2)\in Q$，所以
	\[
	\frac{ab}{d_1d_2}=\frac{c}{d}
	\]
	对某个 $c\in P$ 和 $d\in D$ 成立。于是对某个 $x\in D$ 有
	\[
	x(dab-d_1d_2c)=0.
	\]
	由于 $c\in P$，有 $xdab\in P$；又因 $P$ 是与 $D$ 不相交的素理想，得 $ab\in P$。由于 $P$ 是素理想，或者 $a\in P$，或者 $b\in P$，因此 $a/d_1\in Q$ 或 $b/d_2\in Q$。这证明 $Q$ 是素理想，并说明 $e$ 把 $R$ 中与 $D$ 不相交的素理想映到 $D^{-1}R$ 的素理想。最后，由 (2) 立即可知，对每个与 $D$ 不相交的 $R$ 的素理想 $P$，都有 $P=c(eP)$。因此 $c$ 与 $e$ 是这些素理想集合之间的互逆对应，从而是双射。这证明 (3)。
	
	由 (1)，$D^{-1}R$ 中每条严格上升（相应地，严格下降）理想链收缩到 $R$ 中一条严格上升（相应地，严格下降）理想链，由此得到 (4)，证明完毕。
\end{proof}

由于 $1\in D$，先局部化理想 $I$，再如 (2) 中那样收缩这个局部化，所得 $R$ 中的理想含有 $I$：
\[
I\subseteq c(eI).
\]

\noindent\textbf{定义.} 设 $R$ 是含 $1$ 的交换环，$D$ 是含 $1$ 的乘法闭子集。$R$ 的理想 $I$ 关于 $D$ 的\emph{饱和}是 $R$ 中的理想 $c(eI)$，其中收缩与扩张都是相对于 $\pi:R\to D^{-1}R$ 来计算的。若 $I=c(eI)$，则称 $I$ 关于 $D$ 是\emph{饱和的}。

粗略地说，命题 38 的 (2) 表明，$I$ 的饱和由 $R$ 中那些若允许使用 $D$ 中元素作分母便会落入 $I$ 的元素组成。若即使允许使用 $D$ 中元素作分母也得不到额外元素，则称这个理想关于 $D$ 是饱和的。

我们可以把关于局部化的结果应用于给出一个算法，用来判定系数在域 $k$ 中的多项式环 $k[x_1,\ldots,x_n]$ 中的理想 $P$ 是否为素理想。基本思想是利用
\[
k[x_1,\ldots,x_i]=k[x_1,\ldots,x_{i-1}][x_i]
\]
这一事实，归纳地考虑理想
\[
P_i=P\cap k[x_1,\ldots,x_i]
\]
是否为素理想。

一般地，设 $R$ 是交换环。若 $P$ 是 $R[x]$ 中的素理想，则 $P\cap R$ 是 $R$ 中的素理想，因此
\[
S=R/(P\cap R)
\]
是整环。令 $F$ 表示它的分式域。我们有两个自然的环同态：
\[
R[x]\longrightarrow (R/(P\cap R))[x]=S[x]\longrightarrow F[x],
\]
其中第一个是自然投影同态，第二个是由 $S\hookrightarrow F$ 诱导的自然包含。注意 $F[x]$ 是 $S[x]$ 关于乘法闭集 $D=S-\{0\}$ 的局部化。下一个命题说明，$P$ 在第一个同态下的像是 $S[x]$ 中关于 $D$ 饱和并且扩张到 $F[x]$ 中成为素理想的素理想；反过来，我们也能通过这些性质判定一个理想是否为 $R[x]$ 中的素理想。

\noindent\textbf{命题 39.} 设 $R$ 是含 $1$ 的交换环，$I$ 是 $R[x]$ 中的理想。则 $I$ 是 $R[x]$ 中的素理想当且仅当：
\begin{enumerate}[label=\roman*.]
	\item $J=I\cap R$ 是 $R$ 中的素理想，即 $S=R/J$ 是整环；并且
	\item 若 $\bar I$ 是 $I$ 在 $S[x]$ 中的像，则 $\bar I F[x]$ 是 $F[x]$ 中的素理想，且满足
	\[
	\bar I F[x]\cap S[x]=\bar I.
	\]
\end{enumerate}

\begin{proof}
	假设 $I$ 是 $R[x]$ 中的素理想，则 $J=I\cap R$ 是 $R$ 中的素理想，且 $S=R/J$ 是整环。由第 9 章命题 2，约化同态
	\[
	R[x]\longrightarrow S[x]=(R/J)[x]
	\]
	的核为 $J[x]$，而 $J[x]\subseteq I[x]$，所以有环同构
	\[
	R[x]/I\simeq S[x]/\bar I.
	\]
	由于 $R[x]/I$ 是整环，故 $\bar I$ 是整环 $S[x]$ 中的素理想。$\bar I\cap S$ 中的元素正是 $R\cap I$ 中元素的像，所以 $\bar I\cap S=0$。由于环 $F[x]$ 是 $S[x]$ 关于乘法闭集 $S-\{0\}$ 的局部化，条件 (ii) 由命题 38(3) 得出。
	
	反过来，若 $I$ 不是素理想，则或者 $J$ 不是 $R$ 中的素理想，或者 $J$ 是 $R$ 中的素理想但 $\bar I$ 不是 $S[x]$ 中的素理想。在后一种情形下，或者 $\bar I F[x]$ 不是 $F[x]$ 中的素理想，或者仍由命题 38(3)，$\bar I$ 不是饱和的。因此若 $I$ 不是素理想，则 (i) 或 (ii) 失败，证明完毕。
\end{proof}

由于 $F[x]$ 是 Euclidean 域，命题 39 中的理想 $\bar I F[x]=(h(x))$ 是主理想，并且它是素理想当且仅当 $h(x)$ 或者为 $0$，或者是 $F[x]$ 中的不可约多项式。设 $h(x)$ 是 $I$ 中的一个元素，其在 $S[x]$ 中的像的首项系数为 $a\in S$。下一个命题说明，$a$ 给出了计算饱和 $\bar I F[x]\cap S[x]$ 所需分母的一个界，并且可用于计算这个饱和。

\noindent\textbf{命题 40.} 设 $S$ 是分式域为 $F$ 的整环，$A$ 是 $S[x]$ 中的非零理想。假设
\[
AF[x]=(h(x)),
\]
其中 $h(x)$ 是 $S[x]$ 中首项系数为 $a\in S$ 的多项式。令 $S_a$ 为 $S$ 关于 $a$ 的幂的局部化。则
\begin{enumerate}[label=(\arabic*)]
	\item
	\[
	AF[x]\cap S[x]=AS_a[x]\cap S[x],
	\]
	\item 若 $\mathcal A$ 表示在多项式环 $S[x,t]$ 中由 $A$ 与 $1-at$ 生成的理想，则
	\[
	AS_a[x]\cap S[x]=\mathcal A\cap S[x].
	\]
\end{enumerate}

\begin{proof}
	先证明
	\[
	AF[x]\cap S_a[x]=AS_a[x].
	\]
	由于 $S_a\subseteq F$，包含关系 $AS_a[x]\subseteq AF[x]\cap S_a[x]$ 是显然的。现在设 $f(x)\in AF[x]\cap S_a[x]$。若 $f(x)$ 的首项为 $sx^N$，$h(x)$ 的首项为 $ax^m$，则由于 $AF[x]=(h(x))$，有 $N\ge m$。于是多项式
	\[
	f(x)-(s/a)x^{N-m}h(x)
	\]
	仍属于 $AF[x]\cap S_a[x]$，并且次数低于 $f(x)$。迭代此过程可见，$f(x)$ 可以写成 $S_a[x]$ 中的一个多项式乘以 $h(x)$，因此 $f(x)\in AS_a[x]$。将等式 $AF[x]\cap S_a[x]=AS_a[x]$ 的两边都与 $S[x]$ 相交，得到命题的第一个断言。
	
	为证明第二个断言，先设 $f(x)\in\mathcal A\cap S[x]$。则可写
	\[
	f(x)=f_1(x,t)b(x)+f_2(x,t)(1-at),
	\]
	其中 $b(x)\in A$，而 $f_1,f_2\in S[x,t]$。代入 $t=1/a$，得到
	\[
	f(x)=f_1(x,1/a)b(x).
	\]
	由于 $f_1(x,1/a)\in S_a[x]$，得 $f(x)\in AS_a[x]\cap S[x]$。反过来，设
	\[
	f(x)=b(x)g(x)\in S[x],
	\]
	其中 $g(x)\in S_a[x]$ 且 $b(x)\in A$。若 $a^N$ 是 $g(x)$ 各系数分母中出现的最高次 $a$ 的幂，则 $a^Ng(x)\in S[x]$。写作
	\[
	f(x)=(at)^Nf(x)+(1-(at)^N)f(x)
	=b(x)t^N(a^Ng(x))+(1-(at)^N)f(x),
	\]
	可见 $f(x)\in\mathcal A\cap S[x]$，得到反包含，证明完毕。
\end{proof}

现在设 $P$ 是 $k[x_1,\ldots,x_n]$ 中的理想。对 $i=1,\ldots,n$，令 $P_i$ 为 $P$ 与 $k[x_1,\ldots,x_i]$ 的交。我们利用命题 39 与命题 40 归纳地判定
\[
P_1,P_2,\ldots,P_n=P
\]
是否分别为其相应多项式环中的素理想。

理想 $P_1$ 在 Euclidean 域 $k[x_1]$ 中是素理想，当且仅当它为 $0$ 或由一个不可约多项式生成。现在假设 $i\ge2$，并且已经证明 $P_{i-1}$ 是 $k[x_1,\ldots,x_{i-1}]$ 中的素理想，于是商环
\[
S=k[x_1,\ldots,x_{i-1}]/P_{i-1}
\]
是整环。若 $F$ 表示 $S$ 的分式域，则由命题 39，$P_i$ 是 $k[x_1,\ldots,x_i]$ 中的素理想，当且仅当它在
\[
(k[x_1,\ldots,x_{i-1}]/P_{i-1})[x_i]=S[x_i]
\]
中的像是一个饱和理想，并且其到 Euclidean 域 $F[x_i]$ 的扩张是素理想。设 $h(x_i)\in S[x_i]$ 是这个理想的一个生成元，且 $a$ 是 $h(x_i)$ 的首项系数。则 $(h(x_i))$ 是 $F[x_i]$ 中的素理想当且仅当 $h(x_i)=0$ 或 $h(x_i)$ 是不可约多项式。由命题 40，$P_i$ 在 $S[x_i]$ 中的像是饱和的，当且仅当它等于 $\mathcal A\cap S[x_i]$，其中 $\mathcal A$ 是 $S[x_i,t]$ 中由 $P_i$ 与 $1-at$ 生成的理想。后一条件可以在 $k[x_1,\ldots,x_i,t]$ 中检验：它等价于检查 $k[x_1,\ldots,x_i,t]$ 中由 $P_i$ 与 $1-at$ 生成的理想与 $k[x_1,\ldots,x_i]$ 的交正好为 $P_i$（参见练习 3）。

把这些观察与第 9 章关于 Gröbner 基的结果结合起来，我们得到下列算法，用来判定 $k[x_1,\ldots,x_n]$ 中的理想 $P$ 是否为素理想（或者等价地，判定相应的仿射代数集是否为簇）。

\paragraph{判定 $k[x_1,\ldots,x_n]$ 中理想何时为素理想的算法}

\begin{enumerate}[label=(\arabic*)]
	\item 对 $P$，关于字典序单项式序
	\[
	x_n>\cdots>x_1
	\]
	计算约化 Gröbner 基
	\[
	G=\{g_1,\ldots,g_m\}.
	\]
\end{enumerate}
由第 9.6 节命题 29，$G$ 中属于 $k[x_1,\ldots,x_i]$ 的元素将是
\[
P_i=P\cap k[x_1,\ldots,x_i]
\]
的约化 Gröbner 基 $\{g_1,\ldots,g_{m_i}\}$。

\begin{enumerate}[label=(\arabic*),resume]
	\item 通过检查 $P_1=0$ 或 $P_1$ 的非零生成元在 $k[x_1]$ 中不可约，判定 $P_1$ 是否为 $k[x_1]$ 中的素理想。
\end{enumerate}
对每个 $i\ge2$，假设 $P_{i-1}$ 已被判定为 $k[x_1,\ldots,x_{i-1}]$ 中的素理想（否则 $P$ 不是 $k[x_1,\ldots,x_n]$ 中的素理想）。令
\[
S=k[x_1,\ldots,x_{i-1}]/P_{i-1},
\]
并令 $F$ 为 $S$ 的分式域。应用步骤 (3) 与 (4) 来判定 $P_i$ 是否为 $k[x_1,\ldots,x_i]$ 中的素理想。

\begin{enumerate}[label=(\arabic*),resume]
	\item 若 $m_i=m_{i-1}$，则 $P_i$ 映到 $S[x_i]$ 中的零理想，从而是素理想。否则 $P_i$ 在 $S[x_i]$ 和 $F[x_i]$ 中的像为非零理想，并由 $g_{m_{i-1}+1},\ldots,g_{m_i}$ 的像生成。对 $F[x_i]$ 中这些生成元应用 Euclidean 算法，找出 $P_i$ 中的一个元素 $h(x_i)$，使其在 $F[x_i]$ 中的像生成 $P_i$ 在 $F[x_i]$ 中的像。判定 $h(x_i)$ 在 $F[x_i]$ 中是否不可约；若不是，则 $P_i$ 与 $P$ 都不是素理想。
\end{enumerate}
（注意，在对 $P_i$ 在 $F[x_i]$ 中的像的生成元应用 Euclidean 算法之后，我们可以在每个等式中乘以 $S$ 的一个单一元素以“清除分母”，使所有余式，特别是最后一个非零余式 $h(x_i)$，都是 $P_i$ 的像中的元素。）

\begin{enumerate}[label=(\arabic*),resume]
	\item 令 $a\in k[x_1,\ldots,x_{i-1}]$ 为 $h(x_i)$ 的首项系数（作为 $x_i$ 的多项式）。在 $k[x_1,\ldots,x_i,t]$ 中，对由 $P_i$ 与 $1-at$ 生成的理想，关于字典序单项式序
	\[
	t>x_i>\cdots>x_1
	\]
	计算约化 Gröbner 基。判定这个约化基中属于 $k[x_1,\ldots,x_i]$ 的元素是否为
	\[
	\{g_1,\ldots,g_{m_i}\};
	\]
	若是，则 $P_i$ 是 $k[x_1,\ldots,x_i]$ 中的素理想；若不是，则 $P_i$ 与 $P$ 不是素理想。
\end{enumerate}

最后我们指出，类似思想（连同一些小的修改，以把关于 Gröbner 基的结果推广到系数在整环 $R$ 中的多项式环 $R[x_1,\ldots,x_n]$）可用于给出算法，以判定例如 $\mathbb Z[x_1,\ldots,x_n]$ 中的理想何时为素理想。

\paragraph{例}
\begin{enumerate}[label=(\arabic*)]
	\item 对任意无限域 $k$，考虑 $k[x,y,z]$ 中的理想
	\[
	P=(xz-y^2,\ yz-x^3,\ z^2-x^2y).
	\]
	由第 15.1 节练习 26 可知 $P$ 是素理想，因为存在从 $k[x,y,z]/P$ 到整环 $k[A^1]$ 的单射（参见第 15.2 节练习 24）。这里我们用本节的思想证明 $P\subseteq\mathbb Q[x,y,z]$ 是素理想。$P$ 关于字典序单项式序 $x>y>z$ 的约化 Gröbner 基为
	\[
	\{x^3-yz,\ x^2y-z^2,\ xy^3-z^3,\ xz-y^2,\ y^5-z^4\}.
	\]
	因此 $P_1=P\cap\mathbb Q[z]=(0)$，且
	\[
	P_2=P\cap\mathbb Q[y,z]=(y^5-z^4).
	\]
	由于 $P_1=0$，理想 $P_1$ 是 $\mathbb Q[z]$ 中的素理想。
	
	接下来检查 $P_2$ 是 $\mathbb Q[y,z]$ 中的素理想，这也可以直接完成（参见练习 4 或第 9.1 节练习 14）。在这种情形下 $S=\mathbb Q[z]$，$F=\mathbb Q(z)$。$P_2$ 在 $F[y]$ 中的像由
	\[
	h(y)=y^5-z^4
	\]
	生成，此多项式在 $\mathbb Q(z)[y]$ 中不可约。$h(y)$ 的首项系数为 $1$，而 $(y^5-z^4,1-t)$ 在 $\mathbb Q[y,z,t]$ 中关于字典序单项式序 $t>y>z$ 的约化 Gröbner 基为
	\[
	\{y^5-z^4,\ 1-t\}.
	\]
	$P_2$ 的约化 Gröbner 基中的元素是这个基中唯一属于 $\mathbb Q[y,z]$ 的元素，所以 $P_2$ 是 $\mathbb Q[y,z]$ 中的素理想。
	
	现在用 $P_2$ 是素理想这一事实来证明 $P$ 是素理想。在这种情形下，$S$ 是整环
	\[
	\mathbb Q[y,z]/P_2=\mathbb Q[y,z]/(y^5-z^4),
	\]
	其分式域 $F$ 由
	\[
	S=\mathbb Q[z]+\mathbb Q[z]y+\mathbb Q[z]y^2+\mathbb Q[z]y^3+\mathbb Q[z]y^4,
	\]
	\[
	F=\mathbb Q(z)+\mathbb Q(z)y+\mathbb Q(z)y^2+\mathbb Q(z)y^3+\mathbb Q(z)y^4
	\]
	给出，其中 $y^5=z^4$。$P$ 在 $S[x]$ 中的像是由元素
	\[
	g_1=x^3-yz,
	\quad
	g_2=yx^2-z^2,
	\quad
	g_3=y^3x-z^3,
	\quad
	g_4=zx-y^2,
	\quad
	y^5-z^4=0
	\]
	生成的理想 $\bar P$。
	
	$F[x]$ 中 $g_1,g_2,g_3,g_4$ 的最大公因子，也即生成 $\bar P$ 在 $F[x]$ 中的像的元素，是不可约多项式
	\[
	x-y^2/z.
	\]
	多项式 $h(x)=zx-y^2$ 属于 $\bar P$，并且其像生成 $F[x]$ 中同一理想，所以可在算法的 (4) 中取 $a=z$。理想
	\[
	(xz-y^2,\ yz-x^3,\ z^2-x^2y,\ 1-zt)
	\]
	在字典序单项式序 $t>x>y>z$ 下的约化 Gröbner 基由 $P$ 的约化 Gröbner 基以及含 $t$ 的元素 $tx-1$ 与 $tz-1$ 组成，因此 $P$ 是 $\mathbb Q[x,y,z]$ 中的素理想。
	
	\item 考虑 $\mathbb Q[x,y,z]$ 中的理想
	\[
	P=(xz-y^3,\ xy-z^2),
	\]
	它关于字典序单项式序 $x>y>z$ 的约化 Gröbner 基为
	\[
	\{xy-z^2,\ xz-y^3,\ y^4-z^3\}.
	\]
	这里 $P_1=0$，且
	\[
	P_2=P\cap\mathbb Q[y,z]=(y^4-z^3)
	\]
	与例 (1) 一样是素理想。在这种情形下
	\[
	S=\mathbb Q[y,z]/P_2
	\]
	由
	\[
	S=\mathbb Q[z]+\mathbb Q[z]y+\mathbb Q[z]y^2+\mathbb Q[z]y^3
	\]
	给出，其中 $y^4=z^3$；其分式域 $F$ 与前例类似。并且 $\bar P=(g_1,g_2)$ 在 $S[x]$ 中，其中
	\[
	g_1=yx-z^2,
	\qquad
	g_2=zx-y^3.
	\]
	$\bar P$ 到 $F[x]$ 的扩张由不可约多项式 $yx-z^2$ 生成；$h(x)=yx-z^2$ 是 $\bar P$ 中的一个元素，它在 $F[x]$ 中具有同一像，首项系数为 $a=y$。对 $\mathbb Q[x,y,z,t]$ 中的理想
	\[
	(xz-y^3,\ xy-z^2,\ 1-yt)
	\]
	使用字典序 $t>x>y>z$，其约化 Gröbner 基为
	\[
	\{x^2-y^2z,\ xy-z^2,\ xz-y^3,\ y^4-z^3,\ ty-1,\ tz^2-x\}.
	\]
	这个基含有元素 $x^2-y^2z$，而它不在 $P$ 的约化 Gröbner 基中，所以 $P$ 不是 $\mathbb Q[x,y,z]$ 中的素理想。这个计算不仅说明 $P$ 不是素理想，而且是通过局部化 $S_a$ 显式地说明 $\bar P$ 在 $S[x]$ 中的像并非饱和。计算出 $a=y$ 使我们可以找出两个不在 $P$ 中而乘积在 $P$ 中的显式元素：$f=x^2-y^2z\notin P$ 且 $y\notin P$，但 $y$ 的某个幂乘以 $f$ 属于 $P$。在此情形下，直接计算可验证 $yf\in P$。
\end{enumerate}

\subsection*{模的局部化}

现在设 $M$ 是 $R$-模，$D$ 如前是 $R$ 中含 $1$ 的乘法闭子集。则构造 $D^{-1}R$ 所用思想可以类似地用于从 $M$ 构造一个 $D^{-1}R$-模 $D^{-1}M$，如下所述。在 $D\times M$ 上定义关系
\[
(d,m)\sim(e,n)
\quad\Longleftrightarrow\quad
x(dn-em)=0\quad\text{对某个 }x\in D\text{ 成立},
\]
容易检查这是一个等价关系。记 $(d,m)$ 的等价类为 $m/d$，并令 $D^{-1}M$ 为等价类的集合。于是直接验证可知，运算
\[
\frac{m}{d}+\frac{n}{e}=\frac{em+dn}{de},
\qquad
\left(\frac{r}{d}\right)\left(\frac{m}{e}\right)=\frac{rm}{de}
\]
是良定义的，并给 $D^{-1}M$ 赋予 $D^{-1}R$-模结构。

\noindent\textbf{定义.} $D^{-1}R$-模 $D^{-1}M$ 称为 $M$ 关于 $D$ 的\emph{分式模}，或称为 $M$ 在 $D$ 处的\emph{局部化}。

注意，局部化 $D^{-1}M$ 也是一个 $R$-模（因为每个 $r\in R$ 通过 $r/1$ 作用在 $D^{-1}M$ 上），并且有 $R$-模同态
\[
\pi:M\longrightarrow D^{-1}M,
\qquad
\pi(m)=\frac{m}{1}.
\]
由等价关系的定义直接得到
\[
\ker\pi=\{m\in M\mid dm=0\text{ 对某个 }d\in D\text{ 成立}\}.
\]
同态 $\pi$ 具有与定理 36 中相类似的泛性质。设 $N$ 是一个 $R$-模，且对每个 $d\in D$，$d$ 在 $N$ 上的左乘都是 $N$ 的双射。若 $\psi:M\to N$ 是任意 $R$-模同态，则存在唯一 $R$-模同态
\[
\Psi:D^{-1}M\longrightarrow N
\]
使得 $\Psi\circ\pi=\psi$。

若 $M,N$ 是 $R$-模，$\varphi:M\to N$ 是 $R$-模同态，则对 $R$ 中任意乘法集 $D$，容易检查有一个诱导的 $D^{-1}R$-模同态
\[
D^{-1}M\longrightarrow D^{-1}N,
\qquad
m/d\longmapsto \varphi(m)/d.
\]

下一个结果说明 $M$ 在 $D$ 处的局部化与张量积有关。

\noindent\textbf{命题 41.} 设 $D$ 是 $R$ 的含 $1$ 的乘法闭子集，$M$ 是 $R$-模。则有 $D^{-1}R$-模同构
\[
D^{-1}M\simeq D^{-1}R\otimes_R M,
\]
也即，$D^{-1}M$ 是由 $R$-模 $M$ 通过标量扩张得到的 $D^{-1}R$-模。

\begin{proof}
	从 $D^{-1}R\times M$ 到 $D^{-1}M$ 的映射
	\[
	(r/d,m)\longmapsto rm/d
	\]
	良定义且 $R$-平衡，所以诱导出从 $D^{-1}R\otimes_R M$ 到 $D^{-1}M$ 的同态。映射
	\[
	m/d\longmapsto (1/d)\otimes m
	\]
	给出一个良定义的逆同态：若 $m/d=m'/d'$ 在 $D^{-1}M$ 中成立，则对某个 $x\in D$ 有 $x(d'm-dm')=0$，于是
	\[
	(1/d)\otimes m=(1/xd'd)\otimes (xd'm)=(1/xd'd)\otimes (xdm')=(1/d')\otimes m'.
	\]
	因此 $D^{-1}M$ 与 $D^{-1}R\otimes_R M$ 作为 $R$-模同构；由于这些互逆同构也是 $D^{-1}R$-模同态，命题得证。
\end{proof}

对环 $R$ 或 $R$-模 $M$ 在 $D$ 处局部化，与环和模上的代数运算配合得非常好，如下命题所示。

\noindent\textbf{命题 42.} 设 $R$ 是含 $1$ 的交换环，$D^{-1}R$ 是 $R$ 关于含 $1$ 的乘法闭子集 $D$ 的局部化。
\begin{enumerate}[label=(\arabic*)]
	\item 局部化与理想的有限和与有限交相容：若 $I,J$ 是 $R$ 的理想，则
	\[
	D^{-1}(I+J)=D^{-1}I+D^{-1}J,
	\qquad
	D^{-1}(I\cap J)=D^{-1}I\cap D^{-1}J.
	\]
	局部化与商相容：
	\[
	D^{-1}R/D^{-1}I\simeq D^{-1}(R/I),
	\]
	其中右边的局部化是相对于 $D$ 在商 $R/I$ 中的像来取的。
	
	\item 局部化与取根式相容：若 $N$ 是 $R$ 的幂零根，则 $D^{-1}N$ 是 $D^{-1}R$ 的幂零根。若 $I$ 是 $R$ 的理想，则
	\[
	\operatorname{rad}(D^{-1}I)=D^{-1}(\operatorname{rad}I).
	\]
	
	\item 命题 38(3) 的对应把 primary 理想对应到 primary 理想。更精确地，设 $Q$ 是 $R$ 中的 $P$-primary 理想。若 $D\cap P\ne\varnothing$，则 $D^{-1}Q=D^{-1}R$。若 $D\cap P=\varnothing$，则 $D^{-1}P$ 是素理想，$Q$ 的扩张 $D^{-1}Q$ 是 $D^{-1}R$ 中的 $D^{-1}P$-primary 理想，并且 $D^{-1}Q$ 收缩回 $R$ 仍为 $Q$。
	
	\item 局部化与模的有限和、有限交以及商相容：若 $L,N$ 是 $R$-模 $M$ 的子模，则
	\begin{enumerate}[label=(\alph*)]
		\item
		\[
		D^{-1}(L+N)=D^{-1}L+D^{-1}N,
		\qquad
		D^{-1}(L\cap N)=D^{-1}L\cap D^{-1}N;
		\]
		\item $D^{-1}N$ 是 $D^{-1}M$ 的子模，且
		\[
		D^{-1}M/D^{-1}N\simeq D^{-1}(M/N).
		\]
	\end{enumerate}
	
	\item 局部化与模的有限直和相容：若 $M,N$ 是 $R$-模，则
	\[
	D^{-1}(M\oplus N)\simeq D^{-1}M\oplus D^{-1}N.
	\]
	
	\item 局部化是正合的（即 $D^{-1}R$ 是平坦 $R$-模）：若
	\[
	0\longrightarrow L\overset{\psi}{\longrightarrow} M\overset{\varphi}{\longrightarrow} N\longrightarrow0
	\]
	是 $R$-模短正合列，则诱导出的 $D^{-1}R$-模列
	\[
	0\longrightarrow D^{-1}L\overset{\psi'}{\longrightarrow} D^{-1}M\overset{\varphi'}{\longrightarrow} D^{-1}N\longrightarrow0
	\]
	也正合。
\end{enumerate}

\begin{proof}
	先证明 (6)。设
	\[
	0\longrightarrow L\overset{\psi}{\longrightarrow} M\overset{\varphi}{\longrightarrow} N\longrightarrow0
	\]
	是 $R$-模短正合列。$D^{-1}N$ 的每个元素都形如 $n/d$，其中 $n\in N$，$d\in D$。由于 $\varphi$ 满射，有 $n=\varphi(m)$ 对某个 $m\in M$ 成立，所以
	\[
	\varphi'(m/d)=\varphi(m)/d=n/d,
	\]
	故 $\varphi':D^{-1}M\to D^{-1}N$ 满射。若 $m/d$ 在 $\varphi'$ 的核中，则对某个 $d_1\in D$ 有 $d_1\varphi(m)=0$。于是 $\varphi(d_1m)=0$，由原列在 $M$ 处正合，得到 $d_1m=\psi(l)$ 对某个 $l\in L$ 成立，所以
	\[
	m/d=d_1m/(d_1d)=\psi(l)/(d_1d)=\psi'(l/(d_1d)).
	\]
	因此 $\ker(\varphi')\subseteq\operatorname{im}(\psi')$。若 $\psi(l)/d\in\operatorname{im}(\psi')$，则
	\[
	\varphi'(\psi(l)/d)=\varphi(\psi(l))/d=0,
	\]
	说明反包含 $\operatorname{im}(\psi')\subseteq\ker(\varphi')$，从而诱导列在 $D^{-1}M$ 处正合。最后，若 $\psi'(l/d)=0$，则对某个 $d_2\in D$ 有 $d_2\psi(l)=0$，也即 $\psi(d_2l)=0$；由于 $\psi$ 是单射，得 $d_2l=0$。于是
	\[
	l/d=d_2l/(d_2d)=0,
	\]
	所以 $\psi'$ 是单射。这证明
	\[
	0\longrightarrow D^{-1}L\overset{\psi'}{\longrightarrow}D^{-1}M\overset{\varphi'}{\longrightarrow}D^{-1}N\longrightarrow0
	\]
	正合。
	
	为证明 (1) 的第一个断言，注意若 $i\in I$、$j\in J$、$d\in D$，则
	\[
	(i+j)/d=i/d+j/d
	\]
	说明 $D^{-1}(I+J)\subseteq D^{-1}I+D^{-1}J$；而若 $i\in I$、$j\in J$、$d_1,d_2\in D$，则
	\[
	i/d_1+j/d_2=(d_2i+d_1j)/(d_1d_2)
	\]
	说明 $D^{-1}I+D^{-1}J\subseteq D^{-1}(I+J)$。对第二个断言，包含
	\[
	D^{-1}(I\cap J)\subseteq D^{-1}I\cap D^{-1}J
	\]
	显然。若
	\[
	a/d\in D^{-1}I\cap D^{-1}J,
	\]
	则对某些 $d_1,d_2\in D$，有 $d_1a\in I$ 且 $d_2a\in J$。于是 $d_1d_2a\in I\cap J$，并且
	\[
	a/d=(d_1d_2a)/(d_1d_2d),
	\]
	这给出包含 $D^{-1}I\cap D^{-1}J\subseteq D^{-1}(I\cap J)$。 (1) 的最后一个断言由把 (6) 应用于正合列
	\[
	0\longrightarrow I\longrightarrow R\longrightarrow R/I\longrightarrow0
	\]
	得到。
	
	为证明 (2)，先设 $a\in\operatorname{rad}I$，所以 $a^n\in I$ 对某个 $n\ge1$ 成立。于是
	\[
	(a/d)^n=a^n/d^n\in D^{-1}I,
	\]
	故
	\[
	D^{-1}(\operatorname{rad}I)\subseteq\operatorname{rad}(D^{-1}I).
	\]
	反过来，若 $a/d\in\operatorname{rad}(D^{-1}I)$，则 $(a/d)^n\in D^{-1}I$ 对某个 $n\ge1$ 成立，也即 $d_1a^n\in I$ 对某个 $d_1\in D$ 成立。于是
	\[
	(d_1a)^n=d_1^{n-1}(d_1a^n)\in I,
	\]
	所以 $d_1a\in\operatorname{rad}I$，从而
	\[
	a/d=d_1a/(d_1d)\in D^{-1}(\operatorname{rad}I),
	\]
	说明
	\[
	\operatorname{rad}(D^{-1}I)\subseteq D^{-1}(\operatorname{rad}I).
	\]
	这证明 (2) 的第二个断言，第一个断言由将其应用于理想 $I=(0)$ 得出。
	
	对于 (3)，先注意 $D\cap P=\varnothing$ 当且仅当 $D\cap Q=\varnothing$（一个包含关系显然，另一个由 $d\in D\cap P$ 推出 $d^n\in D\cap Q$ 对某个 $n$ 成立）。$D\cap P\ne\varnothing$ 的情形以及 $D\cap P=\varnothing$ 时 $D^{-1}P$ 是素理想，已在命题 38 中证明。为说明 $D^{-1}Q$ 是 $D^{-1}R$ 中的 primary 理想，设
	\[
	(a/d_1)(b/d_2)\in D^{-1}Q,
	\qquad
	a/d_1\notin D^{-1}Q.
	\]
	则存在某个 $d\in D$ 使得 $dab\in Q$；由于 $a\notin Q$ 且 $Q$ 是 primary 理想，得 $(db)^n\in Q$ 对某个 $n\ge1$ 成立。于是
	\[
	(b/d_2)^n=d^nb^n/(d^nd_2^n)\in D^{-1}Q,
	\]
	故 $D^{-1}Q$ 是 primary 理想。由 (2)，$D^{-1}Q$ 的根式为 $D^{-1}P$。最后，由命题 38(2)，$D^{-1}Q$ 的收缩是 $R$ 中一个含 $Q$ 的理想，并且它恰由所有满足 $dr\in Q$ 对某个 $d\in D$ 成立的 $r\in R$ 组成。由于 $Q$ 是 $P$-primary，primary 的定义说明：若 $dr\in Q$ 且 $d\notin P$，则 $r\in Q$，因此 $D^{-1}Q$ 的收缩就是 $Q$。
	
	(4) 的证明与 (1) 的证明本质相同，留作练习。
	
	容易看出，若 $R$-模正合列
	\[
	0\longrightarrow L\longrightarrow M\longrightarrow N\longrightarrow0
	\]
	分裂，则 $D^{-1}R$-模正合列
	\[
	0\longrightarrow D^{-1}L\longrightarrow D^{-1}M\longrightarrow D^{-1}N\longrightarrow0
	\]
	也分裂，这给出 (5)。
\end{proof}

命题 38 表明，在乘法闭集 $D$ 处局部化会突出 $R$ 中不含 $D$ 中元素的理想，因为 $R$ 的其他理想扩张到 $D^{-1}R$ 后会变成平凡的。下一个命题用局部化对理想的 primary 分解的作用给出一个更精确的说法。

\noindent\textbf{命题 43.} 设 $R$ 是 Noether 环，且
\[
I=Q_1\cap\cdots\cap Q_m
\]
是 proper 理想 $I$ 的极小 primary 分解，其中 $Q_i$ 是 $P_i$-primary 理想。设 $D$ 是 $R$ 中含 $1$ 的乘法闭集，并把 primary 理想 $Q_1,\ldots,Q_m$ 编号为使得
\[
D\cap P_i=\varnothing\quad(1\le i\le t),
\qquad
D\cap P_i\ne\varnothing\quad(t+1\le i\le m).
\]
则
\[
D^{-1}I=D^{-1}Q_1\cap\cdots\cap D^{-1}Q_t
\]
是 $D^{-1}R$ 中 $D^{-1}I$ 的极小 primary 分解，并且 $D^{-1}Q_i$ 是 $D^{-1}P_i$-primary 理想。此外，对 $1\le i\le t$，$D^{-1}Q_i$ 收缩回 $R$ 为 $Q_i$，并且
\[
c(D^{-1}I)=Q_1\cap\cdots\cap Q_t
\]
是 $D^{-1}I$ 收缩回 $R$ 的极小 primary 分解。

\begin{proof}
	由命题 42(3)，当 $t+1\le i\le m$ 时 $D^{-1}Q_i=D^{-1}R$；当 $1\le i\le t$ 时，$D^{-1}Q_i$ 是 $D^{-1}P_i$-primary 理想，并且拉回为 $Q_i$。由同一命题的 (1)，有
	\[
	D^{-1}I=D^{-1}Q_1\cap\cdots\cap D^{-1}Q_t,
	\]
	而 (3) 说明这是 primary 分解。收缩到 $R$ 得到
	\[
	c(D^{-1}I)=Q_1\cap\cdots\cap Q_t,
	\]
	这也蕴含这些分解都是极小的。
\end{proof}

特别地，我们可以完成定理 21 的证明：

\noindent\textbf{推论 44.} 在 $I$ 的极小 primary 分解中，属于孤立素理想的 primary 理想由 $I$ 唯一决定。

\begin{proof}
	令 $P$ 是属于 $I$ 的素理想集合 $\{P_1,\ldots,P_m\}$ 中的极小元素，并在命题 43 中取 $D=R-P$。则 $D\cap P_i=\varnothing$ 只在 $P=P_i$ 时发生，所以 $I$ 在 $D$ 处局部化后的收缩恰好是属于极小素理想 $P$ 的 primary 理想 $Q$。由于属于 $I$ 的素理想 $\{P_1,\ldots,P_m\}$ 由 $I$ 唯一决定，故属于 $I$ 的孤立素理想的 primary 理想 $Q$ 也由 $I$ 唯一决定。
\end{proof}

局部化在某些素理想上的隔离作用，在素理想 $P$ 处局部化的情形中特别精确（见推论 37 后的例 (3)）。我们先回顾一种重要环的定义（参见第 7.4 节练习 37--39）。

\noindent\textbf{定义.} 含 $1$ 且只有一个极大理想的交换环称为\emph{局部环}。

\noindent\textbf{命题 45.} 设 $R$ 是含 $1$ 的交换环。则下列条件等价：
\begin{enumerate}[label=(\arabic*)]
	\item $R$ 是具有唯一极大理想 $M$ 的局部环；
	\item 若 $M$ 是 $R$ 中所有非单位元素的集合，则 $M$ 是理想；
	\item 存在 $R$ 的极大理想 $M$，使得每个 $1+m$（$m\in M$）都是 $R$ 中的单位。
\end{enumerate}

\begin{proof}
	若 $a\in R$，则理想 $(a)$ 或者等于 $R$，此时 $a$ 是单位；或者是 proper 理想，此时 $(a)$ 包含于某个极大理想中（第 7.4 节命题 11）。因此，若 $R$ 是局部环且 $M$ 是其唯一极大理想，则每个 $a\notin M$ 都是单位，所以 $M$ 恰好由 $R$ 中的非单位组成，这说明 (1) 蕴含 (2)。又若非单位的集合 $M$ 是理想，则这个理想必为 $R$ 中唯一极大理想，因此 (2) 蕴含 (1)。
	
	现在假设 (3) 成立。若 $a$ 是 $R$ 中不属于极大理想 $M$ 的元素，则 $(a)+M=R$，所以对某些 $b\in R$ 与 $m\in M$，有 $ab+m=1$。于是 $ab=1-m$ 由假设为单位，所以 $a$ 也是单位。这说明 $M$ 是 $R$ 中唯一极大理想，故 (3) 蕴含 (1)。反过来，若 $R$ 是局部环，则对任意 $m\in M$，有 $1+m\notin M$，所以 $1+m$ 是单位，因此 (1) 蕴含 (3)。
\end{proof}

\noindent\textbf{命题 46.} 对任意含 $1$ 的交换环 $R$，令 $R_P$ 为 $R$ 在素理想 $P$ 处的局部化，令 $eP$ 为 $P$ 到 $R_P$ 的扩张。
\begin{enumerate}[label=(\arabic*)]
	\item 环 $R_P$ 是局部环，唯一极大理想为 $eP$。$eP$ 到 $R$ 的收缩为 $P$，即 $c(eP)=P$；并且从 $R$ 到 $R_P$ 的映射诱导出整环 $R/P$ 到 $R_P/eP$ 的单射。商 $R_P/eP$ 是域，并与整环 $R/P$ 的分式域同构。
	\item 若 $R$ 是整环，则 $R_P$ 是整环。环 $R$ 嵌入局部环 $R_P$；将 $R$ 与其像等同后，$R_P$ 的唯一极大理想是 $PR_P$。
	\item $R_P$ 中的素理想与 $R$ 中包含于 $P$ 的素理想之间双射对应。
	\item 若 $P$ 是 $R$ 的极小非零素理想，则 $R_P$ 有唯一非零素理想。
	\item 若 $P=M$ 是极大理想，$I$ 是 $R$ 的任意 $M$-primary 理想，则
	\[
	R_M/eI\simeq R/I.
	\]
	特别地，$R_M/eM\simeq R/M$，且对所有 $n\ge1$，
	\[
	(eM)/(eM)^n\simeq M/M^n.
	\]
\end{enumerate}

\begin{proof}
	若 $P'$ 是 $R$ 的素理想，则
	\[
	P'\cap(R-P)=\varnothing
	\]
	当且仅当 $P'\subseteq P$，所以 (3) 立即由命题 38(3) 得出，并且 (4) 随之而来。由于由命题 38(2) 知 $eP\ne R_P$，再由 (3) 得 $R_P$ 是具有唯一极大理想 $eP$ 的局部环，这证明 (1) 的第一个断言。
	
	由命题 38(2)，收缩 $c(eP)$ 是集合
	\[
	\{r\in R\mid dr\in P\text{ 对某个 }d\in R-P\text{ 成立}\},
	\]
	而由于 $P$ 是素理想，$dr\in P$ 且 $d\notin P$ 蕴含 $r\in P$。这说明 $c(eP)=P$，即 (1) 的第二个断言。
	
	从 $R$ 到 $R_P/eP$ 的映射的核为 $c(eP)=P$，所以诱导映射 $R/P\to R_P/eP$ 是单射。由 (1) 的第一部分，商 $R_P/eP$ 是域，所以存在从整环 $R/P$ 的分式域到 $R_P/eP$ 的诱导同态。局部化 $R_P$ 的泛性质说明，存在从 $R_P/eP$ 到 $R/P$ 的分式域的逆同态（因为 $R$ 中不属于 $P$ 的每个元素都映为 $R/P$ 中的单位）。于是 $R_P/eP$ 同构于 $R/P$ 的分式域。
	
	若 $R$ 是整环，则 $R-P$ 不含零因子，所以由推论 37，$R$ 嵌入 $R_P$；若将 $R$ 与其在 $R_P$ 中的像等同，则 $eP=PR_P$，所以 (2) 成立。
	
	为证明 (5)，由命题 42(1)，可过渡到商 $R/I$，从而归约到 $I=0$ 的情形。在这种情形下，$R$ 中的极大理想 $P=M$ 是 $R$ 的幂零根，因此是 $R$ 的唯一极大理想。由命题 45，$R-M$ 的每个元素都是单位，所以 $R_P=R$；(5) 中的每个断言立即成立，命题得证。
\end{proof}

\paragraph{例}
命题 (5) 中的结论在 $P$ 是非极大素理想的情形下一般不成立。例如，取 $R=\mathbb Z$ 中的 $P=(0)$，则 $R/P=\mathbb Z$ 且 $R_P/PR_P=\mathbb Q$；在此情形下，对所有 $n\ge1$，
\[
(PR_P)/(PR_P)^n\simeq P/P^n=0
\]
（参见练习）。

\noindent\textbf{定义.} 设 $M$ 是 $R$-模，$P$ 是 $R$ 的素理想，并令 $D=R-P$。$R_P$-模 $D^{-1}M$ 称为 $M$ 在 $P$ 处的\emph{局部化}，记作 $M_P$。

由命题 41，$M_P$ 也可等同于张量积 $R_P\otimes_R M$。当 $R$ 是整环且 $P=(0)$ 时，$M_{(0)}$ 是 $R$ 的分式域 $F$ 上的模，也即 $F$-向量空间。

元素 $m/1$ 在 $M_P$ 中为零当且仅当 $rm=0$ 对某个 $r\in R-P$ 成立，所以在 $P$ 处局部化会消去 $M$ 中相对于不包含于 $P$ 的素理想 $P'$ 的 $P'$-torsion 元素。特别地，在整环上于 $(0)$ 处局部化会消去 $M$ 的 torsion 子群。

\noindent\textbf{定义.} 若 $R$ 是整环，则 $R$-模 $M$ 的\emph{秩}定义为局部化 $M_{(0)}$ 作为 $R$ 的分式域上的向量空间的维数。

容易看出，这个秩的定义与第 12 章引入的秩的概念一致。

\paragraph{例}
令 $R=\mathbb Z$，并令 $\mathbb Z_{(p)}$ 为 $\mathbb Z$ 在非零素理想 $(p)$ 处的局部化。任意 Abel 群 $M$ 都是 $\mathbb Z$-模，因此可以通过构造 $M_{(p)}$ 将 $M$ 在 $(p)$ 处局部化。这个 Abel 群等于 $M$ 关于那些阶有限且不能被 $p$ 整除的元素所成子群的商。若 $M$ 是有限 Abel 群（或者更一般地，是 torsion Abel 群），则 $M_{(p)}$ 是一个 $p$-群，并且是 $M$ 的 Sylow $p$-子群或 $p$-primary 分量。$M$ 在 $(0)$ 处的局部化 $M_{(0)}$ 是平凡群。举一个具体例子，令 $M=\mathbb Z/6\mathbb Z$ 为 6 阶循环群，视为 $\mathbb Z$-模。则 $M$ 在 $p=2$ 处的局部化是 $\mathbb Z/2\mathbb Z$，在 $p=3$ 处是 $\mathbb Z/3\mathbb Z$，而在 $\mathbb Z$ 的所有其他素理想处化为 $0$。

一般地，在素理想 $P$ 处对模 $M$ 局部化，会得到一个更简单、性质更易判定的模 $M_P$。于是，将 $M_P$ 的这些“局部”性质翻译回关于模 $M$ 本身的“整体”信息，就是一个有意义的问题。例如，最基本的问题，即模 $M$ 是否为 $0$，可以局部地回答：

\noindent\textbf{命题 47.} 设 $M$ 是 $R$-模。则下列条件等价：
\begin{enumerate}[label=(\arabic*)]
	\item $M=0$；
	\item 对 $R$ 的所有素理想 $P$，有 $M_P=0$；
	\item 对 $R$ 的所有极大理想 $\mathfrak m$，有 $M_{\mathfrak m}=0$。
\end{enumerate}

\begin{proof}
	(1) 蕴含 (2) 蕴含 (3) 是显然的，所以只需证明 (3) 蕴含 (1)。设 $m$ 是 $M$ 中的非零元素，并考虑 $m$ 在 $R$ 中的零化子 $I$，即满足 $rm=0$ 的元素 $r\in R$ 组成的理想。由于 $m$ 非零，$I$ 是 $R$ 的 proper 理想。令 $\mathfrak m$ 是 $R$ 中含 $I$ 的极大理想，并考虑 $M$ 在相应局部化 $M_{\mathfrak m}$ 中的元素 $m/1$。若这个元素为 $0$，则对某个 $r\in R-\mathfrak m$ 有 $rm=0$。但这意味着 $r$ 是 $I$ 的元素而不属于 $\mathfrak m$，矛盾。因此 $M_{\mathfrak m}\ne0$，这证明 (3) 蕴含 (1)。
\end{proof}

一般而言，并非所有 $M$ 的局部化都满足的性质也必为 $M$ 所满足。例如，一个环 $R$ 的所有局部化都可以是整环，但 $R$ 本身不必是整环（例如上面的 $\mathbb Z/6\mathbb Z$）。不过，研究各种可能的局部化仍能获得大量信息，这正是这一技术如此有用的原因。例如，若 $R$ 是整环，则每个局部化 $R_P$ 都可视为 $R$ 的分式域 $F$ 中含有 $R$ 的子环；下一个命题说明，$F$ 中同时属于每个局部化的元素只有 $R$ 中的元素。

\noindent\textbf{命题 48.} 设 $R$ 是整环。则 $R$ 是它的局部化的交：
\[
R=\bigcap_P R_P.
\]
事实上，$R$ 是其在极大理想 $\mathfrak m$ 处局部化的交：
\[
R=\bigcap_{\mathfrak m} R_{\mathfrak m}.
\]

\begin{proof}
	如前所述，$R\subseteq\bigcap_{\mathfrak m}R_{\mathfrak m}$。现在设 $a$ 是 $R$ 的分式域 $F$ 中的元素，并且对每个极大理想 $\mathfrak m$ 都有 $a\in R_{\mathfrak m}$。考虑
	\[
	I_a=\{d\in R\mid da\in R\}.
	\]
	容易检查 $I_a$ 是 $R$ 的理想，并且 $a\in R$ 当且仅当 $1\in I_a$，也即 $I_a=R$。假设 $I_a\ne R$。则存在极大理想 $\mathfrak m$ 含 $I_a$。由于 $a\in R_{\mathfrak m}$，有 $a=r/d$，其中 $r\in R$ 且 $d\in R-\mathfrak m$。但这意味着 $d\in I_a$ 且 $d\notin\mathfrak m$，矛盾。因此 $a\in R$，所以 $\bigcap_{\mathfrak m}R_{\mathfrak m}\subseteq R$，我们已证明命题的第二个断言。第一个断言随即显然成立。
\end{proof}

环 $R$ 的另一个可以局部检测的重要性质是正规性。

\noindent\textbf{命题 49.} 设 $R$ 是整环。则下列条件等价：
\begin{enumerate}[label=(\arabic*)]
	\item $R$ 是正规环，即 $R$（在其分式域中）整闭；
	\item 对 $R$ 的所有素理想 $P$，$R_P$ 是正规的；
	\item 对 $R$ 的所有极大理想 $\mathfrak m$，$R_{\mathfrak m}$ 是正规的。
\end{enumerate}

\begin{proof}
	令 $F$ 为 $R$ 的分式域，因此 $R$ 的各个局部化都可视为 $F$ 的子环。
	
	先假设 $R$ 是整闭的，并设 $y\in F$ 在 $R_P$ 上整。则 $y$ 是某个首一 $n$ 次多项式的根，其系数形如 $a_i/d_i$，其中 $d_i\notin P$。元素
	\[
	y'=y(d_0d_1\cdots d_{n-1})^n
	\]
	于是是一个系数在 $R$ 中的首一 $n$ 次多项式的根，也即 $y'$ 在 $R$ 上整。由于 $R$ 假设为正规，得 $y'\in R$，从而
	\[
	y=y'/(d_0d_1\cdots d_{n-1})^n\in R_P,
	\]
	这证明 (1) 蕴含 (2)。蕴含 (2) 蕴含 (3) 是平凡的。现在假设对 $R$ 的所有极大理想 $\mathfrak m$，$R_{\mathfrak m}$ 都正规，并令 $y$ 是 $F$ 中在 $R$ 上整的元素。由于 $R\subseteq R_{\mathfrak m}$，$y$ 特别地也在 $R_{\mathfrak m}$ 上整，所以由假设，对每个极大理想 $\mathfrak m$ 都有 $y\in R_{\mathfrak m}$。于是由前一命题，$y\in R$，这证明 (3) 蕴含 (1)。
\end{proof}

我们现在可以容易地证明 Going-up 定理（见第 3 节）的第一部分，即在推论 27 的证明中用到的部分。

\noindent\textbf{推论 50.} 设 $R$ 是交换环 $S$ 的子环，且 $1\in R$，并假设 $S$ 在 $R$ 上整。若 $P$ 是 $R$ 中的素理想，则存在 $S$ 中的素理想 $Q$，使得
\[
P=Q\cap R.
\]

\begin{proof}
	令 $D=R-P$，则 $D$ 是 $R$ 与 $S$ 的乘法闭子集。于是下图交换：
	\[
	\begin{array}{ccc}
		R & \longrightarrow & D^{-1}R=R_P \\
		\big\downarrow & & \big\downarrow \\
		S & \longrightarrow & D^{-1}S
	\end{array}
	\]
	其中竖直映射为包含。容易看出 $D^{-1}S$ 在 $R_P$ 上整（练习 20）。令 $\mathfrak m$ 为 $D^{-1}S$ 的任一极大理想。则 $\mathfrak m\cap R_P$ 是 $R_P$ 中的极大理想，这是由定理 26(2) 的第二个断言得出的（注意，定理 26(2) 的第一部分在第二个断言的证明中并未使用）。由命题 38(1)，$\mathfrak m\cap R_P$ 是 $P$ 到局部环 $R_P$ 的扩张，并且这个理想收缩到 $R$ 正是 $P$。换言之，沿上边与右边的映射取 $\mathfrak m$ 的原像得到 $P$。若 $Q\subseteq S$ 表示沿下边映射取 $\mathfrak m$ 的原像，则由命题 38(3)，$Q$ 是素理想。由于 $Q\cap R$ 是沿上图左边映射取 $Q$ 的拉回，由图的交换性可知
	\[
	Q\cap R=P.
	\]
\end{proof}

\subsection*{仿射代数簇的局部环}

在本节余下部分，令 $k$ 为代数闭域，$V$ 为 $k$ 上的仿射簇，其坐标环为 $k[V]$。则 $k[V]$ 是整环，所以可以构造它的分式域：
\[
k(V)=\{f/g\mid f,g\in k[V],\ g\ne0\}.
\]
$k(V)$ 的元素称为 $V$ 上的\emph{有理函数}，而 $k(V)$ 称为 $V$ 上的\emph{有理函数域}。当 $k[V]$ 是唯一分解整环时，$f/g$ 有一个本质唯一的“最简”代表；但一般而言，$k(V)$ 中每个分式 $f/g$ 都有许多表示为 $k[V]$ 中两个元素之比的方式。由于 $k[V]$ 是整环，
\[
f/g=f_1/g_1
\]
当且仅当 $fg_1=f_1g$。

$k[V]$ 的元素可视为 $V$ 上的 $k$-值函数；若分母不为零，$k(V)$ 的元素也同样可以这样看待（这有助于解释这个域的术语）。由于 $k(V)$ 的同一元素可能以多种方式写成 $f/g$，我们作如下定义：

\noindent\textbf{定义.} 若存在某些 $f_1,g_1\in k[V]$，使得
\[
f/g=f_1/g_1
\]
且 $g_1(v)\ne0$，则称 $f/g$ 在 $v$ 处\emph{正则}，或称在点 $v\in V$ 处\emph{有定义}。

若 $(f_2,g_2)$ 是另一个这样的对且 $g_2(v)\ne0$，则作为 $k$ 的元素有
\[
f_1(v)/g_1(v)=f_2(v)/g_2(v),
\]
所以只要 $f/g$ 在 $v$ 处正则，就有一个良定义的方式来指定它在 $v$ 处的 $k$ 中取值。

\paragraph{例}
簇
\[
V=Z(xz-yw)\subseteq A^4
\]
的坐标环是
\[
k[V]=k[x,y,z,w]/(xz-yw).
\]
考虑 $k[V]$ 的商域 $k(V)$ 中的元素 $\bar x/\bar y$。由于在 $k[V]$ 中有 $\bar x\bar z=\bar y\bar w$，这个元素也可写为 $\bar w/\bar z$。从 $f$ 的第一个表达式可知，$f$ 在 $V$ 上所有满足 $\bar y\ne0$ 的点处正则；从第二个表达式可知，$f$ 在 $V$ 上所有满足 $\bar z\ne0$ 的点处正则。不难证明这些正是 $f$ 正则的所有点。此外，不存在 $f$ 的单个表达式 $f=a/b$，其中 $a,b\in k[V]$，使得对 $f$ 正则的每个点 $v$ 都有 $b(v)\ne0$（参见练习 25）。

若 $f/g\in k(V)$ 在点 $v$ 处正则，例如 $f/g=f_1/g_1$ 且 $g_1(v)\ne0$，则 $f/g$ 在 $v$ 的 Zariski 开邻域 $V_{g_1}$ 中所有满足 $g_1\ne0$ 的点处也正则。作为 $V$ 上的 $k$-值函数，这意味着若 $f/g$ 在 $v$ 处有定义，则它也在 $v$ 的某个（Zariski 开）邻域中有定义。由于仿射簇的任意非空开集都是 Zariski 稠密的（参见第 2 节练习 11），可知 $V$ 上的每个有理函数都在 $V$ 的一个稠密点集上有定义（在适当意义下“几乎处处”有定义）。此外，表示 $f/g$ 的每一对 $f_1/g_1$ 与 $f_2/g_2$，作为函数在 $v$ 的开邻域 $V_{g_1}\cap V_{g_2}$ 上相等，但这个邻域的“大小”依赖于 $g_1$ 与 $g_2$；一般不存在一个共同的 $v$ 的开邻域，使得 $f/g$ 的所有在 $v$ 处分母非零的代表都同时有定义。

若 $v$ 是 $V$ 中固定点，则有理函数 $f/g$ 在 $v$ 处正则，当且仅当 $f/g=f_1/g_1$ 对某些 $f_1,g_1\in k[V]$ 成立，并且 $g_1\notin I(v)$，其中 $I(v)$ 是 $V$ 上在 $v$ 处为零的函数所成理想。这意味着在 $v$ 处有定义的有理函数集合正是 $k[V]$ 在极大理想 $I(v)$ 处的局部化：

\noindent\textbf{定义.} 对每个点 $v\in V$，$V$ 上在 $v$ 处有定义的有理函数的集合
\[
\mathcal O_{v,V}=\{f/g\in k(V)\mid f/g\text{ 在 }v\text{ 处正则}\}
\]
称为 $V$ 在 $v$ 处的\emph{局部环}。等价地，$V$ 在 $v$ 处的局部环是 $k[V]$ 在极大理想 $I(v)$ 处的局部化。

特别地，$\mathcal O_{v,V}$ 是局部环，唯一极大理想为 $\mathfrak m_{v,V}$，其中
\[
\mathfrak m_{v,V}=\{f/g\in\mathcal O_{v,V}\mid (f/g)(v)=0\},
\]
即 $V$ 上在 $v$ 处有定义并等于 $0$ 的有理函数集合。由于 $\mathcal O_{v,V}$ 是 Noether 整环 $k[V]$ 在素理想处的局部化，$\mathcal O_{v,V}$ 也是 Noether 整环。又由命题 46(5)，注意到
\[
\mathcal O_{v,V}/\mathfrak m_{v,V}\simeq k[V]/I(v)\simeq k.
\]

回忆 $V$ 到 $k$ 的多项式映射也称为 $V$ 到 $k$ 的正则映射。这是因为它们恰好是 $V$ 上处处正则的有理函数：

\noindent\textbf{命题 51.} 若 $V$ 是代数闭域 $k$ 上的仿射簇，则 $V$ 上在 $V$ 的所有点处正则的有理函数恰好是多项式函数 $k[V]$。

\begin{proof}
	这由命题 48 得出：$k[V]$ 在所有极大理想处的局部化在 $k(V)$ 中的交恰为 $k[V]$。
\end{proof}

由于 $k[V]$ 的极大理想与 $V$ 的点双射对应，而局部环 $\mathcal O_{v,V}$ 等于 $k[V]$ 在对应于 $v$ 的极大理想处的局部化，因此 $\mathcal O_{v,V}$ 内在地依赖于环 $k[V]$，与 $V$ 在某个特定仿射空间中的嵌入无关。

设 $\varphi:V\to W$ 是仿射簇之间的态射，其对应的 $k$-代数同态为
\[
\bar\varphi:k[W]\longrightarrow k[V].
\]
若 $v\in V$ 被 $\varphi$ 映到 $w\in W$，则直接可证 $\bar\varphi$ 诱导出相应局部环之间的同态（仍记为 $\bar\varphi$）：
\[
\bar\varphi:\mathcal O_{w,W}\longrightarrow\mathcal O_{v,V},
\qquad
\bar\varphi(h/k)=\bar\varphi(h)/\bar\varphi(k),
\]
并且在这个同态下
\[
\bar\varphi^{-1}(\mathfrak m_{v,V})=\mathfrak m_{w,W}
\]
（具有这一性质的局部环同态称为\emph{局部同态}）。注意，$\bar\varphi$ 一般不能扩张为从整个 $k(W)$ 到 $k(V)$ 的域同态，因为 $\bar\varphi$ 的核中的 $k[W]$ 元素不会映到 $k(V)$ 中的可逆元素。还容易检查，若 $\psi\circ\varphi$ 是态射的复合，则在局部环上有
\[
\overline{\psi\circ\varphi}=\bar\varphi\circ\bar\psi.
\]

局部环 $\mathcal O_{v,V}$ 可用于给出 $V$ 在 $v$ 处“光滑性”（就切线存在性而言）的代数定义，下面说明这一点。先假设 $V=Z(f)$ 是 $A^n$ 中由不可约多项式 $f\in k[x_1,\ldots,x_n]$ 的零点定义的超曲面簇。对 $V$ 上任意点 $v=(v_1,\ldots,v_n)$，令 $D_v(f)(x_1,\ldots,x_n)$ 为线性多项式
\[
D_v(f)(x_1,\ldots,x_n)=\sum_{i=1}^n \frac{\partial f}{\partial x_i}(v)x_i,
\]
其中 $f$ 对 $x_i$ 的偏导数由关于 $x_i$ 的多项式的通常形式求导规则给出（把其他所有变量都视为常数）。多项式
\[
D_v(f)(x_1-v_1,\ldots,x_n-v_n)
\]
是函数 $f$ 在 $v$ 处的一阶 Taylor 多项式，所以给出 $f(x_1,
\ldots,x_n)\in k[x_1,
\ldots,x_n]$ 在 $v$ 处的最佳线性近似。由此可见，若 $T$ 是线性簇
\[
Z(D_v(f)(x_1,
\ldots,x_n)),
\]
即 $D_v(f)$ 为零的点集，则平移 $v+T$ 在 $v$ 处与超曲面 $Z(f)$“相切”。

\paragraph{例}
设 $f=x^2-y\in k[x,y]$，所以 $V=Z(f)$ 就是抛物线 $y=x^2$。我们有
\[
\partial f/\partial x=2x,
\qquad
\partial f/\partial y=-1,
\]
它们在 $v=(3,9)$ 处分别等于 $6$ 与 $-1$。于是
\[
D_{(3,9)}(f)(x,y)=6x-y,
\]
对应的线性簇 $T$ 是过原点的直线 $y=6x$。平移 $(3,9)+T$ 是抛物线在 $(3,9)$ 处的通常切线。$x^2-y$ 在 $(3,9)$ 处的 Taylor 展开为
\[
x^2-y=[6(x-3)-(y-9)]+(x-3)^2.
\]
一阶项是 $D_{(3,9)}(f)(x-3,y-9)$，并给出 $x^2-y$ 在 $(3,9)$ 附近的最佳线性近似。

将这些概念推广到 $A^n$ 中任意仿射簇 $V$ 是直接的。

\noindent\textbf{定义.} 定义 $V$ 在 $v$ 处的\emph{切空间}为线性簇
\[
\mathbb T_{v,V}=Z\bigl(\{D_v(f)(x_1,\ldots,x_n)\mid f\in I(V)\}\bigr).
\]

形式偏导数是 $k$-线性的，并且满足通常的乘积求导法则，所以切空间可以由 $I(V)$ 的生成元计算：若
\[
I(V)=(f_1,f_2,\ldots,f_m),
\]
则
\[
\mathbb T_{v,V}=\bigcap_{i=1}^m Z(D_v(f_i)).
\]
注意 $\mathbb T_{v,V}$ 是向量空间的交，因此是 $k^n$ 的向量子空间。

这个切空间 $\mathbb T_{v,V}$ 的定义虽然清楚地显示了它与簇 $V$ 的切线之间的联系，但看起来依赖于 $V$ 在 $A^n$ 中的嵌入。事实上，切空间完全可以用局部环 $\mathcal O_{v,V}$ 来定义，如下命题所证。

\noindent\textbf{命题 52.} 设 $V$ 是代数闭域 $k$ 上的仿射簇，$v$ 是 $V$ 上一点，其局部环为 $\mathcal O_{v,V}$，相应极大理想为 $\mathfrak m_{v,V}$。则有 $k$-向量空间同构
\[
(\mathbb T_{v,V})^*\simeq \mathfrak m_{v,V}/\mathfrak m_{v,V}^2,
\]
其中 $(\mathbb T_{v,V})^*$ 表示切空间 $\mathbb T_{v,V}$ 的向量空间对偶（参见第 11.3 节）。

\begin{proof}
	令 $(k^n)^*$ 表示 $k^n$ 的 $n$ 维向量空间对偶。由于每个 $D_v(f)$ 都是线性函数，$D_v$ 是从 $k[x_1,\ldots,x_n]$ 到 $(k^n)^*$ 的线性变换。
	
	令 $M_v$ 为 $k[x_1,\ldots,x_n]$ 中由集合 $x_i-v_i$（$1\le i\le n$）生成的极大理想。$M_v$ 在 $k[V]$ 中的像 $M_v/I(V)$ 是 $V$ 上在 $v$ 处为零的函数理想 $I(v)$，并且
	\[
	I(v)^2=M_v^2+I(V).
	\]
	于是 $\mathcal O_{v,V}$ 是 $k[V]$ 在 $I(v)$ 处的局部化；将 $I(v)$ 与它在 $\mathcal O_{v,V}$ 中的像等同，由命题 46(2) 得
	\[
	\mathfrak m_{v,V}=I(v)\mathcal O_{v,V}.
	\]
	由 $D_v$ 的定义，有
	\[
	D_v(x_i-v_i)=x_i,
	\]
	而这些线性函数构成 $(k^n)^*$ 的一组基，因此 $D_v$ 把 $M_v$ 满射到 $(k^n)^*$。$D_v$ 的核由那些在 $v$ 处 Taylor 展开从至少二次项开始的 $k[x_1,\ldots,x_n]$ 的元素组成，它们正是 $M_v^2$ 中的元素。因此 $D_v$ 定义同构
	\[
	D_v:M_v/M_v^2\simeq (k^n)^*.
	\]
	
	切空间 $\mathbb T_{v,V}$ 是 $k^n$ 的向量子空间，所以 $k^n$ 上的每个线性函数都限制为 $\mathbb T_{v,V}$ 上的线性函数。把 $D_v$ 与这个限制映射复合，得到线性变换
	\[
	D:M_v\longrightarrow (k^n)^*\longrightarrow (\mathbb T_{v,V})^*.
	\]
	由于各个映射都是满射，$D$ 也是满射。我们已经看到 $I(v)^2=M_v^2+I(V)$，所以
	\[
	I(v)/I(v)^2\simeq M_v/(M_v^2+I(V)).
	\]
	由命题 46(5)，
	\[
	\mathfrak m_{v,V}/\mathfrak m_{v,V}^2\simeq I(v)/I(v)^2.
	\]
	因此，为证明命题，只需说明
	\[
	\ker D=M_v^2+I(V),
	\]
	因为此时
	\[
	\mathfrak m_{v,V}/\mathfrak m_{v,V}^2
	\simeq M_v/(M_v^2+I(V))
	=M_v/\ker D
	\simeq(\mathbb T_{v,V})^*.
	\]
	
	多项式 $f$ 属于 $\ker D$ 当且仅当 $D_v(f)$ 在 $\mathbb T_{v,V}$ 上为零，也即当且仅当 $f$ 在 $v$ 处展开的 Taylor 多项式的一次项属于 $I(\mathbb T_{v,V})$。由于 $I(V)$ 中函数的一次项生成理想 $I(\mathbb T_{v,V})$，可得 $f$ 属于 $\ker D$ 当且仅当对某个 $g\in I(V)$，$f-g$ 的一次项为零。但这等价于
	\[
	f\in I(V)+M_v^2.
	\]
	所以 $\ker D=I(V)+M_v^2$，命题得证。
\end{proof}

回忆簇 $V$ 的维数按定义是域 $k(V)$ 在 $k$ 上的超越次数。由于每个局部环 $\mathcal O_{v,V}$ 都以 $k(V)$ 为其分式域，$V$ 的维数由其任意局部环的分式域在 $k$ 上的超越次数决定。

\noindent\textbf{定义.} 若 $k$-向量空间 $\mathbb T_{v,V}$ 的维数等于 $\dim V$，则称 $V$ 在点 $v\in V$ 处\emph{非奇异}，或称 $v$ 是 $V$ 的\emph{非奇异点}。等价地（由命题 52），$v$ 是 $V$ 的非奇异点当且仅当
\[
\dim_k(\mathfrak m_{v,V}/\mathfrak m_{v,V}^2)=\dim V.
\]
否则称点 $v$ 为\emph{奇异点}。若 $V$ 在每个点处都非奇异，则称 $V$ \emph{非奇异}或\emph{光滑}。

几何图像是，在非奇异点 $v$ 处，有预期数量的独立切线：曲线上的切线，曲面上的切平面，等等。

一个簇 $V$ 在点 $v$ 处是否非奇异，可以由局部环 $\mathcal O_{v,V}$ 的性质来判定，即是否
\[
\dim_k(\mathfrak m_{v,V}/\mathfrak m_{v,V}^2)=\dim\mathcal O_{v,V}.
\]
具有这一性质的局部环称为\emph{正则局部环}。特别地，奇异性的概念不依赖于 $V$ 在某个特定仿射空间中的嵌入。这种代数解释可用于为抽象代数簇定义光滑性；在抽象代数簇中，例如曲面的切平面这一几何直观并不那么明显。

若 $f_1,\ldots,f_m$ 是 $I(V)$ 的生成元，并且在 $A^n$ 中定义 $V$，则 $V$ 的维数可以由 $I(V)$ 的 Gröbner 基来确定（参见练习 29）。把切空间 $\mathbb T_{v,V}$ 作为 $k$-向量空间来确定其维数，是一个线性代数问题：这个向量空间是 $m$ 个线性方程
\[
D_v(f_i)(x_1,\ldots,x_n)=0
\]
的解集。若这个方程组的系数矩阵
\[
\left(\frac{\partial f_i}{\partial x_j}(v)\right)
\]
的秩为 $r$，则 $\mathbb T_{v,V}$ 是维数为 $n-r$ 的向量空间。由此不难建立下列事实：

\begin{enumerate}[label=\arabic*.]
	\item 对每个 $V\subseteq A^n$ 中的点 $v$，都有
	\[
	\dim V\le \dim_k(\mathbb T_{v,V})\le n.
	\]
	\item $V$ 的奇异点集是 $V$ 的一个 proper Zariski 闭子集。$V$ 的非奇异点集是 $V$ 的一个非空开子集；特别地，$V$ 的非奇异点在 $V$ 中稠密（所以 $V$ 中“多数”点是非奇异的）。
\end{enumerate}
我们还不加证明地陈述下列结果，它进一步联系了 $V$ 的局部几何与 $V$ 的局部环的代数性质：
\begin{enumerate}[label=\arabic*.,resume]
	\item 若 $v$ 是非奇异点，则局部环 $\mathcal O_{v,V}$ 是唯一分解整环；特别地，$\mathcal O_{v,V}$ 是整闭的（参见推论 25 后的例 3）。
\end{enumerate}

若对每个点 $v\in V$，$\mathcal O_{v,V}$ 都是 U.F.D.，则称簇 $V$ 为\emph{阶乘簇}；若对每个 $v\in V$，$\mathcal O_{v,V}$ 都是整闭的，则称 $V$ 为\emph{正规簇}（由命题 49，这等价于 $k[V]$ 整闭）。由上面的 (3)，有
\[
\text{光滑簇}\subseteq\text{阶乘簇}\subseteq\text{正规簇}.
\]
一般而言，上述每个包含都是真包含。不过，当 $V$ 的维数为 $1$，也即 $V$ 是仿射曲线时，这三个性质事实上等价：稍后我们将证明，不可约仿射曲线光滑当且仅当它正规或阶乘（参见第 16.2 节推论 13）。因此，在代数闭域 $k$ 上，
\[
\text{不可约仿射曲线 }C\text{ 光滑当且仅当 }k[C]\text{ 整闭}.
\]
对任意不可约仿射曲线 $C$，$k[V]$ 在 $k(V)$ 中的整闭包 $S$ 也是一条不可约仿射曲线 $\widetilde C$ 的坐标环。于是 $S$ 在 $k[V]$ 上整；由定理 30 与推论 27，可得存在从光滑曲线 $\widetilde C$ 到 $C$ 的态射，它是满的并且具有有限纤维。曲线 $\widetilde C$ 称为 $C$ 的\emph{正规化}或\emph{非奇异模型}，并且可以证明它在同构意义下唯一。注意，存在一条光滑曲线有限地映到 $C$ 上这一问题（一个“几何”问题），是如何由环扩张中整闭包的存在性（一个“代数”问题）解决的。

我们将在第 16.2 节末尾给出不可约仿射曲线光滑性的另一个刻画。

\subsection*{习题}

照常，$R$ 是含 $1$ 的交换环，$D$ 是 $R$ 中的乘法闭集。
\begin{enumerate}[label=\arabic*.]
	\item 设 $M$ 是有限生成 $R$-模。证明 $D^{-1}M=0$ 当且仅当对某个 $d\in D$，有 $dM=0$。
	
	\item 设 $I$ 是 $R$ 中的理想，$D$ 是 $R$ 的乘法闭子集，其分式环为 $D^{-1}R$，并令 $c(eI)$ 为 $I$ 关于 $D$ 的饱和。
	\begin{enumerate}[label=(\alph*)]
		\item 证明 $c(eI)=R$ 当且仅当 $eI=D^{-1}R$，当且仅当 $I\cap D\ne\varnothing$。
		\item 证明 $I=c(eI)$ 饱和当且仅当对每个 $d\in D$，若 $da\in I$ 则 $a\in I$。
		\item 证明扩张与收缩在关于 $D$ 饱和的 $R$ 的理想与 $D^{-1}R$ 的理想之间给出互逆双射。
		\item 令 $I=(2x,3y)\subseteq\mathbb Z[x,y]$。说明 $I$ 关于 $\mathbb Z-\{0\}$ 的饱和是 $(x,y)$。
	\end{enumerate}
	
	\item 若 $I$ 是交换环 $R$ 中的理想，令
	\[
	\varphi:R[x_1,\ldots,x_n]\longrightarrow (R/I)[x_1,\ldots,x_n]
	\]
	为由系数模 $I$ 化得到的、核为 $I[x_1,\ldots,x_n]$ 的环同态。若 $A$ 是 $(R/I)[x_1,\ldots,x_n]$ 中的理想，令 $\widetilde A$ 表示 $A$ 在 $\varphi$ 下的逆像。
	\begin{enumerate}[label=(\alph*)]
		\item 对任意 $i\ge1$，说明 $\varphi$ 下子环 $(R/I)[x_1,\ldots,x_i]$ 的逆像是
		\[
		R[x_1,\ldots,x_i]+I[x_1,\ldots,x_n].
		\]
		\item 证明
		\[
		\varphi(\widetilde A\cap R[x_1,\ldots,x_i])=A\cap (R/I)[x_1,\ldots,x_i].
		\]
	\end{enumerate}
	
	\item 令 $f=y^5-z^4$，把它看作 $\mathbb Q[z]$ 系数的关于 $y$ 的多项式。
	\begin{enumerate}[label=(\alph*)]
		\item 证明 $f$ 在 $\mathbb Q[z]$ 中没有根。
		\item 假设
		\[
		f=(y^2+ay+b)(y^3+cy^2+dy+e).
		\]
		说明 $a,b,c,d,e$ 满足方程组
		\[
		a+c=0,
		\qquad
		ac+b+d=0,
		\qquad
		ad+bc+e=0,
		\qquad
		ae+bd=0,
		\qquad
		be-z^4=0.
		\]
		推出 $e^5=z^{12}$，并得出 $f$ 在 $\mathbb Q[y,z]$ 中不可约。[使用消元。]
	\end{enumerate}
	
	\item 设 $R$ 是 U.F.D.，其分式域为 $F$，且 $p\in R[x]$ 是首一多项式。
	\begin{enumerate}[label=(\alph*)]
		\item 说明 $R[x]$ 中由 $p$ 生成的理想 $pR[x]$ 是素理想当且仅当 $F[x]$ 中由 $p$ 生成的理想 $pF[x]$ 是素理想。[使用 Gauss 引理。]
		\item 证明 $pR[x]$ 是饱和的，即
		\[
		pF[x]\cap R[x]=pR[x].
		\]
	\end{enumerate}
	
	\item 说明
	\[
	I=(y^3-xz,\ xy^2-z^2)
	\]
	不是 $\mathbb Q[x,y,z]$ 中的素理想，并找出显式元素 $a,b\in\mathbb Q[x,y,z]$，使得 $ab\in I$ 但 $a\notin I$ 且 $b\notin I$。
	
	\item 说明
	\[
	P=(y^3-xz,\ xy^2-z^2,\ x^2-yz)
	\]
	是 $\mathbb Q[x,y,z]$ 中的素理想。
	
	\item 说明
	\[
	P=(x^2-yz,\ w^2-x^4z)
	\]
	是 $\mathbb Q[x,y,z,w]$ 中的素理想。
	
	\item 说明
	\[
	P=(xz^2-w^3,\ xw^2-y^4,\ y^4z^2-w^5)
	\]
	是 $\mathbb Q[x,y,z,w]$ 中的素理想。
	
	\item 说明
	\[
	I=(xy-w^3,\ y^2-zw)
	\]
	不是 $\mathbb Q[x,y,z,w]$ 中的素理想，并找出 $a,b$，使得 $ab\in I$ 但 $a,b\notin I$。
	
	\item 令 $R_P$ 为 $R$ 在素理想 $P$ 处的局部化。证明若 $Q$ 是 $R$ 的 $P$-primary 理想，则关于 $Q$ 到 $R_P$ 的扩张与收缩，有 $Q=c(eQ)$。若 $Q$ 是某个包含于 $P$ 的素理想 $P'$-primary 理想，也证明同样结论。
	
	\item 令
	\[
	R=\mathbb R[x,y,z]/(xy-z^2),
	\]
	令 $P=(\bar x,\bar z)$ 为由 $x$ 与 $z$ 的像生成的素理想，并令 $R_P$ 为 $R$ 在 $P$ 处的局部化。证明
	\[
	P^2R_P\cap R=(\bar x),
	\]
	并且它严格大于 $P^2$。
	
	\item 证明：若 $N,N'$ 是 $R$-模 $M$ 的两个 $R$-子模，并且对 $R$ 的每个素理想 $P$（或仅对每个极大理想）在局部化 $M_P$ 中都有 $N_P=N'_P$，则 $N=N'$。
	
	\item 设 $\varphi:M\to N$ 是 $R$-模同态。证明 $\varphi$ 单射（相应地，满射）当且仅当对 $R$ 的每个素理想 $P$（或仅对每个极大理想），诱导的 $R_P$-模同态
	\[
	\varphi_P:M_P\to N_P
	\]
	单射（相应地，满射）。
	
	\item 令 $R=\mathbb Z[\sqrt{-5}]$ 为二次数域 $\mathbb Q(\sqrt{-5})$ 的整数环，令
	\[
	I=(2,1+\sqrt{-5})
	\]
	为 $R$ 中由 $2$ 与 $1+\sqrt{-5}$ 生成的素理想（参见第 8.2 节练习 5）。回忆 $R$ 的每个非零素理想 $P$ 都含有某个素数 $p\in\mathbb Z$。
	\begin{enumerate}[label=(\alph*)]
		\item 若 $P$ 是 $R$ 中不含 $2$ 的素理想，证明 $I_P=R_P$。
		\item 若 $P$ 是 $R$ 中含 $2$ 的素理想，证明 $P=I$ 且
		\[
		I_P=(1+\sqrt{-5})R_P.
		\]
		\item 证明对 $R$ 的每个素理想 $P$，作为 $R_P$-模都有 $I_P\simeq R_P$，但 $I$ 与 $R$ 作为 $R$-模不同构。（这个例子说明，在练习 14 中，给定 $R$-模同态 $\varphi$ 是重要的。）[观察 $I\simeq R$ 作为 $R$-模当且仅当 $I$ 是主理想。]
	\end{enumerate}
	
	\item 证明局部化与张量积相容：对任意 $R$-模 $M,N$ 和 $R$ 中乘法闭集 $D$，存在唯一的 $D^{-1}R$-模同构
	\[
	\varphi:(D^{-1}M)\otimes_{D^{-1}R}(D^{-1}N)\simeq D^{-1}(M\otimes_R N),
	\]
	它把
	\[
	(m/d)\otimes(n/d')
	\]
	送到
	\[
	(m\otimes n)/(dd').
	\]
	
	\item 证明 $R$-模 $A$ 是平坦 $R$-模当且仅当对每个素理想 $P$（或仅对每个极大理想），$A_P$ 是平坦 $R_P$-模。[使用命题 41、练习 14 与 16，以及局部化的正合性。]
	
	\item 在推论 37 后例 (2) 的记号中，证明若 $f$ 在 $R$ 中不是幂零元，则
	\[
	R_f\simeq R[x]/(fx-1).
	\]
	[说明映射 $\varphi:R[x]\to R_f$ 由 $\varphi(r)=r/1$ 与 $\varphi(x)=1/f$ 定义，给出满环同态，并且定理 36 的泛性质给出其逆。]
	
	\item 证明：若 $R$ 是整闭整环，$D$ 是 $R$ 的任意含 $1$ 的乘法闭子集，则 $D^{-1}R$ 是整闭的。
	
	\item 设 $R$ 是环 $S$ 的子环，且 $1\in R$，并且 $S$ 在 $R$ 上整。若 $D$ 是 $R$ 的任意乘法闭子集，证明 $D^{-1}S$ 在 $D^{-1}R$ 上整。
	
	\item 设 $\varphi:R\to S$ 是环同态，$D'$ 是 $S$ 的乘法闭子集。令 $D=\varphi^{-1}(D')$。证明 $D$ 是 $R$ 的乘法闭子集，并且映射
	\[
	\varphi':D^{-1}R\longrightarrow {D'}^{-1}S,
	\qquad
	\varphi'(r/d)=\varphi(r)/\varphi(d)
	\]
	是环同态。
	
	\item 设 $P\subseteq Q$ 是 $R$ 中的素理想，令 $R_Q$ 为 $R$ 在 $Q$ 处的局部化。证明局部化 $R_P$ 同构于 $R_Q$ 在素理想 $PR_Q$ 处的局部化（参见前一练习）。
	
	\item 令 $\varphi:A\to B$ 为满足 $\varphi(1_A)=1_B$ 的交换环同态，$P$ 为 $A$ 中的素理想。关于 $\varphi$，用上标 $c$ 与 $e$ 分别表示理想的收缩与扩张。证明 $P$ 是 $B$ 中某个素理想的收缩，当且仅当
	\[
	P=(P^e)^c.
	\]
	[在 $\varphi(A-P)$ 处局部化 $B$。]
	
	\item （Going-down 定理）设 $S$ 是整环，$R$ 是 $S$ 的整闭子环并且含 $1_S$，$k$ 是 $R$ 的分式域。设 $P_2\subseteq P_1$ 是 $R$ 中的素理想，且 $Q_1$ 是 $S$ 中满足 $Q_1\cap R=P_1$ 的素理想。令 $S_{Q_1}$ 为 $S$ 在 $Q_1$ 处的局部化。
	\begin{enumerate}[label=(\alph*)]
		\item 证明
		\[
		P_2\subseteq P_2S_{Q_1}\cap R.
		\]
		\item 设 $a\in P_2S_{Q_1}\cap R$，并写 $a=s/d$，其中 $s\in P_2S$，$d\in S$，且 $d\notin Q_1$。若 $s$ 在 $k$ 上的极小多项式为
		\[
		x^n+a_{n-1}x^{n-1}+\cdots+a_1x+a_0,
		\]
		其中 $a_0,\ldots,a_{n-1}\in P_2$（参见第 3 节练习 12），证明 $d$ 在 $k$ 上的极小多项式为
		\[
		x^n+b_{n-1}x^{n-1}+\cdots+b_1x+b_0,
		\]
		其中
		\[
		b_i=a_i/a^{n-i},
		\]
		并推出 $b_i\in R$。[使用第 3 节练习 10。]
		\item 证明 $a\in P_2$，并推出
		\[
		P_2S_{Q_1}\cap R=P_2.
		\]
		[说明若 $a\notin P_2$，则 $b_i\in P_2$ 对 $i=0,1,\ldots,n-1$ 成立，这将推出 $d^n\in P_2S\subseteq P_1S\subseteq Q_1$，从而 $d\in Q_1$。]
		\item 证明 $P_2S_{Q_1}$ 包含于 $S_{Q_1}$ 的某个素理想 $P$ 中，并且 $P\cap R=P_2$。[对 $\varphi:R\to S_{Q_1}$ 使用 (c) 与前一练习。]
		\item 令 $Q_2=P\cap S$。证明 $Q_2\subseteq Q_1$ 且 $Q_2\cap R=P_2$。
		\item 结合前一结果用归纳证明 Going-down 定理，即定理 26(4)。
	\end{enumerate}
	
	\item 令 $k$ 是代数闭域，并令
	\[
	V=Z(xz-yw)\subseteq A^4.
	\]
	证明 $f=x/y\in k(V)$ 正则的点集恰好是所有满足 $y\ne0$ 或 $z\ne0$ 的点 $(x,y,z,w)$ 所成集合。[若 $f=a/b$，说明 $ay-bx\in(xz-yw)$ 作为 $k[x,y,z,w]$ 中的多项式成立，并推出 $b\in(y,z)$。] 证明不存在函数 $a/b\in k(V)$，使得对每个 $f$ 正则的点 $v$ 都有 $b(v)\ne0$。
	
	\item （态射的微分）设 $\varphi:V\to W$ 是代数闭域 $k$ 上仿射簇之间的态射，并假设 $\varphi(v)=w$。
	\begin{enumerate}[label=(\alph*)]
		\item 说明 $\varphi$ 诱导从 $k$-向量空间 $M_w/M_w^2$ 到 $k$-向量空间 $M_v/M_v^2$ 的线性映射，并用它说明 $\varphi$ 诱导一个线性映射 $d\varphi$（称为 $\varphi$ 的\emph{微分}），从 $k$-向量空间 $\mathbb T_{v,V}$ 到 $k$-向量空间 $\mathbb T_{w,W}$。
		\item 证明若 $V\subseteq A^n$，$W\subseteq A^m$，并且
		\[
		\varphi=(F_1(x_1,\ldots,x_n),\ldots,F_m(x_1,\ldots,x_n)),
		\]
		则
		\[
		d\varphi:\mathbb T_{v,V}\longrightarrow\mathbb T_{w,W}
		\]
		明确由
		\[
		(d\varphi)(a_1,
		\ldots,a_n)=
		\bigl(D_v(F_1)(a_1,
		\ldots,a_n),
		\ldots,
		D_v(F_m)(a_1,
		\ldots,a_n)\bigr)
		\]
		给出。[若 $g=g(y_1,
		\ldots,y_m)$，说明链式法则推出
		\[
		\frac{\partial(g\circ\varphi)}{\partial x_i}(v)
		=\sum_{j=1}^m
		\frac{\partial g}{\partial y_j}(w)
		\frac{\partial F_j}{\partial x_i}(v),
		\]
		因而
		\[
		D_v(g\circ\varphi)(a_1,
		\ldots,a_n)=D_w(g)(b_1,
		\ldots,b_m),
		\]
		其中 $b_j=D_v(F_j)(a_1,
		\ldots,a_n)$。然后使用事实：若 $g\in I(W)$，则 $g\circ\varphi\in I(V)$。]
		\item 若 $\psi:U\to V$ 是另一个态射且 $\psi(u)=v$，证明相应的
		\[
		d(\varphi\circ\psi):\mathbb T_{u,U}\to\mathbb T_{w,W}
		\]
		与 $d\varphi\circ d\psi$ 相同。
		\item 证明若 $\varphi$ 是同构，则对每个 $\varphi(v)=w$，$d\varphi$ 都是从 $\mathbb T_{v,V}$ 到 $\mathbb T_{w,W}$ 的向量空间同构。
	\end{enumerate}
	
	\item 令 $V=A^1$，且
	\[
	W=Z(xz-y^2,\ yz-x^3,\ z^2-x^2y)\subseteq A^3.
	\]
	令 $\varphi:V\to W$ 为满态射
	\[
	\varphi(t)=(t^3,t^4,t^5)
	\]
	（参见第 1 节练习 26）。对每个 $t\in A^1$，明确描述前一练习中的微分
	\[
	d\varphi:\mathbb T_{t,A^1}\to\mathbb T_{(t^3,t^4,t^5),W};
	\]
	特别地，证明对所有 $t\ne0$，$d\varphi$ 是向量空间同构；而当 $t=0$ 时，它是零映射。用此证明 $V$ 与 $W$ 不同构。
	
	\item 若 $k$ 是域，商环 $k[x]/(x^2)$ 称为 $k$ 上的\emph{对偶数环}。若 $V$ 是 $k$ 上的仿射代数集，说明从 $k[V]$ 到 $k[x]/(x^2)$ 的 $k$-代数同态等价于指定 $V$ 中的一个点 $v$，使得
	\[
	\mathcal O_{v,V}/\mathfrak m_{v,V}=k
	\]
	（称为 $V$ 的 $k$-有理点），并指定 $V$ 在 $v$ 处的切空间 $\mathbb T_{v,V}$ 中的一个元素。
	
	\item （计算簇的维数）令 $P$ 是 $k[x_1,\ldots,x_n]$ 中的素理想，设 $P_0=0$，并令
	\[
	P_i=P\cap k[x_1,\ldots,x_i].
	\]
	定义簇 $V_i=Z(P_i)\subseteq A^i$，其中 $V_0$ 是由单个点组成的零维簇，坐标环为 $k$。
	\begin{enumerate}[label=(\alph*)]
		\item 证明
		\[
		\dim V_{i-1}\le \dim V_i\le \dim V_{i-1}+1.
		\]
		[先给出从 $k(V_{i-1})$ 到 $k[V_i]$ 的单射；再说明 $k[V_i]$ 是由 $k[V_{i-1}]$ 与一个额外生成元生成的 $k$-代数。]
		\item 若 $P_{i-1}$ 在 $k[x_1,\ldots,x_i]$ 中生成的理想等于 $P_i$，说明
		\[
		V_i\simeq V_{i-1}\times A^1,
		\]
		并推出 $\dim V_i=\dim V_{i-1}+1$。
		\item 若 $P_{i-1}$ 在 $k[x_1,\ldots,x_i]$ 中生成的理想真包含于 $P_i$，说明
		\[
		\dim V_i=\dim V_{i-1}.
		\]
		\item 说明 $\dim V$ 等于集合 $\{1,2,
		\ldots,n\}$ 中那些使得 $P_{i-1}$ 在 $k[x_1,
		\ldots,x_i]$ 中生成的理想等于 $P_i$ 的 $i$ 的个数。推出若 $G$ 是 $P$ 关于字典序单项式序
		\[
		x_n>\cdots>x_1
		\]
		的约化 Gröbner 基，且 $G_i=G\cap k[x_1,
		\ldots,x_i]$，其中 $G_0=0$，$N$ 是满足 $G_i\ne G_{i-1}$ 的 $i$（$1\le i\le n$）的个数，则
		\[
		\dim V=n-N.
		\]
	\end{enumerate}
\end{enumerate}

下面十一个练习引入 $R$-模 $M$ 的支集概念，以及它与 $M$ 的伴随素理想的关系。也可参见第 1 节练习 29 到 35 以及第 5 节练习 25 到 30。

\noindent\textbf{定义.} 若 $M$ 是 $R$-模，则使得局部化 $M_P$ 非零的 $R$ 的素理想 $P$ 的集合称为 $M$ 的\emph{支集}，记作 $\operatorname{Supp}(M)$。

\begin{enumerate}[label=\arabic*.,start=30]
	\item 证明 $M=0$ 当且仅当 $\operatorname{Supp}(M)=\varnothing$。[使用命题 47。]
	
	\item 若
	\[
	0\longrightarrow L\longrightarrow M\longrightarrow N\longrightarrow0
	\]
	是 $R$-模正合列，证明 $M_P$ 非零当且仅当局部化 $N_P$ 与 $L_P$ 中至少有一个非零，并推出
	\[
	\operatorname{Supp}(M)=\operatorname{Supp}(L)\cup\operatorname{Supp}(N).
	\]
	特别地，若 $M=M_1\oplus\cdots\oplus M_n$，证明
	\[
	\operatorname{Supp}(M)=\operatorname{Supp}(M_1)\cup\cdots\cup\operatorname{Supp}(M_n).
	\]
	
	\item 设 $P\subseteq Q$ 是 $R$ 中的素理想，且 $M$ 是 $R$-模。证明 $R$-模 $M_Q$ 在 $P$ 处的局部化就是 $M_P$，即
	\[
	(M_Q)_P=M_P.
	\]
	[直接论证，或使用命题 41 与张量积的结合律。]
	
	\item 设 $P\subseteq Q$ 是 $R$ 中的素理想，且 $M$ 是 $R$-模。证明若 $P\in\operatorname{Supp}(M)$，则 $Q\in\operatorname{Supp}(M)$。[使用前一练习。]
	
	\item
	\begin{enumerate}[label=(\alph*)]
		\item 设 $M=Rm$ 是循环 $R$-模。证明 $M_P=0$ 当且仅当存在 $r\in R$，$r\notin P$，使得 $rm=0$。推出 $P\in\operatorname{Supp}(M)$ 当且仅当 $P$ 含有 $m$ 在 $R$ 中的零化子（参见第 10.1 节练习 10）。
		\item 若
		\[
		M=Rm_1+\cdots+Rm_n
		\]
		是有限生成 $R$-模，证明 $P\in\operatorname{Supp}(M)$ 当且仅当 $P$ 属于某个 $\operatorname{Supp}(Rm_i)$，其中 $i=1,
		\ldots,n$。[使用命题 42。] 推出 $P\in\operatorname{Supp}(M)$ 当且仅当 $P$ 含有 $M$ 在 $R$ 中的零化子 $\operatorname{Ann}(M)$。[注意
		\[
		\operatorname{Ann}(M)=\bigcap_{i=1}^n\operatorname{Ann}(Rm_i),
		\]
		然后使用 (a) 与第 7.4 节练习 11。]
	\end{enumerate}
	
	\item 设 $P$ 是 $R$ 的素理想且 $P\cap D=\varnothing$。证明若 $P\in\operatorname{Ass}_R(M)$，则
	\[
	D^{-1}P\in\operatorname{Ass}_{D^{-1}R}(D^{-1}M).
	\]
	[使用命题 38(3) 与命题 42。]
	
	\item 设 $D^{-1}P\in\operatorname{Ass}_{D^{-1}R}(D^{-1}M)$，其中 $P=(a_1,
	\ldots,a_n)$ 是 $R$ 中的有限生成素理想，且 $P\cap D=\varnothing$。
	\begin{enumerate}[label=(\alph*)]
		\item 假设 $m/d\in D^{-1}M$ 在 $D^{-1}R$ 中的零化子为 $D^{-1}P$。证明对某些 $d_1,
		\ldots,d_n\in D$，有
		\[
		d_ia_im=0\in R.
		\]
		\item 令 $d'=d_1d_2\cdots d_n$。证明
		\[
		P=\operatorname{Ann}(d'm),
		\]
		并推出 $P\in\operatorname{Ass}_R(M)$。[包含关系 $P\subseteq\operatorname{Ann}(d'm)$ 是直接的。为证明反包含，说明 $b\in\operatorname{Ann}(d'm)$ 蕴含 $b/1$ 零化 $m/d\in D^{-1}M$，所以 $b/1\in D^{-1}P$，再推出 $b\in P$。]
	\end{enumerate}
	
	\item 设 $M$ 是 Noether 环 $R$ 上的模。使用前两个练习证明，在命题 38(3) 的双射下，满足 $P\cap D=\varnothing$ 的 $\operatorname{Ass}_R(M)$ 中素理想 $P$ 与 $\operatorname{Ass}_{D^{-1}R}(D^{-1}M)$ 中素理想一一对应。
	
	\item 设 $M$ 是 Noether 环 $R$ 上的模，$D$ 是 $R$ 的乘法闭子集。令 $\mathcal S$ 为 $\operatorname{Ass}_R(M)$ 中满足 $P\cap D\ne\varnothing$ 的素理想 $P$ 的子集。本练习证明，局部化映射 $M\to D^{-1}M$ 的核 $N$ 是 $M$ 的唯一子模，满足
	\[
	\operatorname{Ass}_R(N)=\mathcal S,
	\qquad
	\operatorname{Ass}_R(M/N)=\operatorname{Ass}_R(M)-\mathcal S.
	\]
	\begin{enumerate}[label=(\alph*)]
		\item 若 $N'$ 是 $M$ 的子模，且如第 1 节练习 35 中那样满足
		\[
		\operatorname{Ass}_R(N')=\mathcal S,
		\qquad
		\operatorname{Ass}_R(M/N')=\operatorname{Ass}_R(M)-\mathcal S,
		\]
		证明图
		\[
		\begin{array}{ccc}
			M & \xrightarrow{\pi} & M/N'\\
			\varphi\downarrow & & \downarrow \varphi'\\
			D^{-1}M & \xrightarrow{\pi'} & D^{-1}(M/N')
		\end{array}
		\]
		交换，其中 $\pi$ 与 $\pi'$ 是自然投影（参见命题 42(6)），而 $\varphi,\varphi'$ 是局部化同态。
		\item 证明
		\[
		\operatorname{Ass}_{D^{-1}R}(D^{-1}N')=\varnothing,
		\]
		并推出 $D^{-1}N'=0$ 且 $\pi'$ 是单射。[使用前一练习、$\mathcal S$ 的定义以及第 1 节练习 34。]
		\item 若 $x$ 在 $\varphi'$ 的核 $K$ 中，说明
		\[
		\operatorname{Ann}(x)\cap D\ne\varnothing
		\]
		且 $\operatorname{Ass}_R(K)\subseteq\mathcal S$。说明 $\operatorname{Ass}_R(K)\subseteq\operatorname{Ass}_R(M/N')$ 蕴含 $\operatorname{Ass}_R(K)=\varnothing$，并推出 $K=0$。
		\item 证明 $\varphi$ 与 $\pi$ 具有同一核，即 $N=N'$，并且 $M$ 的这个子模是唯一的。
	\end{enumerate}
\end{enumerate}

下面两个练习建立与 $R$-模 $M$ 相关的素理想集合 $\operatorname{Ass}_R(M)$ 与 $\operatorname{Supp}(M)$ 之间的基本关系。
\begin{enumerate}[label=\arabic*.,start=39]
	\item 证明
	\[
	\operatorname{Ass}_R(M)\subseteq\operatorname{Supp}(M).
	\]
	[若 $Rm\simeq R/P$，使用命题 42(4) 与命题 46(1) 说明
	\[
	0\ne (Rm)_P\subseteq M_P.
	\]]
	
	\item 设 $R$ 是 Noether 环，$M$ 是 $R$-模。
	\begin{enumerate}[label=(\alph*)]
		\item 若 $P\in\operatorname{Supp}(M)$，证明 $P$ 含有某个素理想 $Q$，其中 $Q\in\operatorname{Ass}_R(M)$。
		\item 若 $P$ 是 $\operatorname{Supp}(M)$ 中的极小素理想，证明 $P\in\operatorname{Ass}_R(M)$。[使用第 1 节练习 33 说明 $\operatorname{Ass}_{R_P}(M_P)\ne\varnothing$，然后使用练习 37。]
		\item 推出
		\[
		\operatorname{Ass}_R(M)\subseteq\operatorname{Supp}(M),
		\]
		并且这两个集合有相同的极小元素。
	\end{enumerate}
\end{enumerate}

\section{环的素谱}

在本节中，“环”一词均指含 $1$ 的交换环，并且所有环同态 $\varphi:R\to S$ 都假定把 $1_R$ 映到 $1_S$。

我们已经看到，域 $k$ 上仿射代数集 $V$ 的大多数几何性质，都可以翻译为相应坐标环 $k[V]$ 的代数性质，其中 $k[V]$ 是 $V$ 上的 $k$-值函数环。例如，从 $V$ 到 $W$ 的态射对应于从 $k[W]$ 到 $k[V]$ 的 $k$-代数环同态。当域 $k$ 是代数闭域时，这种翻译尤其精确：Hilbert 零点定理建立了 $V$ 中的点 $v$ 与 $k[V]$ 的极大理想 $M=I(v)$ 之间的一一对应；并且如果 $\varphi:V\to W$ 是一个态射，则 $\varphi(v)\in W$ 对应于 $k[W]$ 中的极大理想 $\bar\varphi^{-1}(M)$。在这种发展中，我们通常从仿射代数集的几何性质出发，然后看到其相应坐标环共有的许多代数性质可以对任意交换环定义。现在设想我们反过来做，也就是从一个一般交换环作为代数对象出发，试图仿照 $k[V]$ 与 $V$ 来定义一个相应的“几何”对象。

给定一个交换环 $R$，与 $k[V]$ 和 $V$ 的类比也许最自然地会建议我们把 $R$ 的极大理想 $M$ 的集合定义为相应几何对象的“点”。在这个定义下，如果 $\varphi:R'\to R$ 是一个环同态，那么 $\varphi^{-1}(M)$ 应当对应于极大理想 $M$。不幸的是，一个极大理想在环同态下的逆像一般不必是极大理想。由于素理想在把 $1$ 映到 $1$ 的环同态下的逆像仍是素理想，这提示我们，一个更好的定义也许应当包括 $R$ 的素理想。于是得到如下定义。

\noindent\textbf{定义.} 设 $R$ 是含 $1$ 的交换环。$R$ 的\emph{谱}或\emph{素谱}，记作 $\Spec R$，是 $R$ 的所有素理想的集合。$R$ 的所有极大理想的集合，记作 $\operatorname{mSpec}R$，称为 $R$ 的\emph{极大谱}。

\paragraph{例}
\begin{enumerate}[label=(\arabic*)]
	\item 若 $R$ 是域，则
	\[
	\Spec R=\operatorname{mSpec}R=\{(0)\}.
	\]
	
	\item $\Spec\mathbb Z$ 中的点是素理想 $(0)$ 以及素理想 $(p)$，其中 $p>0$ 是素数；而 $\operatorname{mSpec}\mathbb Z$ 由 $\Spec\mathbb Z$ 中除 $(0)$ 以外的所有素理想组成。
	
	\item $\Spec\mathbb Z[x]$ 的元素如下：
	\begin{enumerate}[label=(\alph*)]
		\item $(0)$；
		\item $(p)$，其中 $p$ 是 $\mathbb Z$ 中的素数；
		\item $(f)$，其中 $f\ne1$ 是一个本原多项式（即其系数的最大公因数等于 $1$），并且在 $\mathbb Q[x]$ 中不可约；
		\item $(p,g)$，其中 $p$ 是 $\mathbb Z$ 中的素数，且 $g$ 是一个模 $p$ 不可约的首一多项式。
	\end{enumerate}
	$\operatorname{mSpec}\mathbb Z[x]$ 的元素正是上面 (d) 中的素理想。
\end{enumerate}

在 $k$ 是代数闭域时，与 $k[V]$ 和 $V$ 的类比中，元素 $f\in k[V]$ 是 $V$ 上取值于 $k$ 的函数，通过在 $V$ 的点 $v$ 处求值而得到。注意，“在 $v$ 处求值”定义了从 $k[V]$ 到 $k$ 的同态，其核为 $I(v)$；而 $f$ 在 $v$ 处的值是商环 $k[V]/I(v)\simeq k$ 中代表 $f$ 的元素。换句话说，$f\in k[V]$ 在 $v\in V$ 处的值可以看作元素 $\bar f\in k[V]/I(v)\simeq k$。一般情形下也可以作类似定义。

\noindent\textbf{定义.} 若 $f\in R$，则 $f$ 在点 $P\in\Spec R$ 处的\emph{值}是元素
\[
f(P)=\bar f\in R/P.
\]

注意，$f$ 在不同点 $P$ 处的值一般位于不同的整环中。还要注意，一般而言 $f\in R$ 不能由它的所有值唯一确定，而只能确定到相差一个 $R$ 的 nilradical 中的元素为止（参见练习 3）。

还有映射 $Z$ 与 $I$ 的类似物，并且也有 Zariski 拓扑的类似物。对 $R$ 的任意子集 $A$，定义
\[
Z(A)=\{P\in X\mid A\subseteq P\}\subseteq\Spec R,
\]
即包含 $A$ 的素理想的集合。显然 $Z(A)=Z(I)$，其中 $I=(A)$ 是由 $A$ 生成的理想；因此只考虑 $Z(I)$ 并无损失，其中 $I$ 是 $R$ 的理想。注意，根据定义，$P\in Z(I)$ 当且仅当 $I\subseteq P$，这又当且仅当对每个 $f\in I$ 都有 $f\in P$。如上把 $f\in R$ 看成 $\Spec R$ 上的一个函数，这就是说，$P\in Z(I)$ 当且仅当对所有 $f\in I$，都有
\[
f(P)=f\bmod P=0\in R/P.
\]
在这个意义下，$Z(I)$ 由 $\Spec R$ 中使 $I$ 中所有函数取值为 $0$ 的点组成。

对 $\Spec R$ 的任意子集 $Y$，定义
\[
I(Y)=\bigcap_{P\in Y}P,
\]
即 $Y$ 中素理想的交。

\noindent\textbf{命题 53.} 设 $R$ 是含 $1$ 的交换环。上面定义的 $R$ 与 $\Spec R$ 之间的映射 $Z$ 与 $I$ 满足：
\begin{enumerate}[label=(\arabic*)]
	\item 对 $R$ 的任意理想 $I$，有
	\[
	Z(I)=Z(\operatorname{rad}(I))=Z(I(Z(I))),
	\qquad
	I(Z(I))=\operatorname{rad} I;
	\]
	\item 对 $R$ 的任意理想 $I,J$，有
	\[
	Z(I\cap J)=Z(IJ)=Z(I)\cup Z(J);
	\]
	\item 若 $\{I_j\}$ 是 $R$ 的任意一族理想，则
	\[
	Z\left(\bigcup I_j\right)=\bigcap Z(I_j).
	\]
\end{enumerate}

\begin{proof}
	若素理想 $P$ 包含理想 $I$，则 $P$ 包含 $\operatorname{rad}I$（第 2 节练习 8），这推出 $Z(I)=Z(\operatorname{rad}(I))$。由于 $\operatorname{rad}I$ 是所有包含 $I$ 的素理想的交（命题 12），由 $I(Y)$ 的定义得 $Z(\operatorname{rad}(I))=Z(I(I))$。类似地，
	\[
	I(Z(I))=\bigcap_{P\in Z(I)}P
	=\bigcap_{I\subseteq P}P
	=\operatorname{rad}I,
	\]
	这完成了 (1) 的证明。
	
	显然 $Z(I\cap J)=Z(I)\cup Z(J)$。设素理想 $P$ 包含 $IJ$。若 $P$ 不包含 $I$，则存在某个 $i\in I$ 使得 $i\notin P$。由于 $iJ\subseteq P$，可知 $J\subseteq P$。这证明 $Z(IJ)=Z(I)\cup Z(J)$，从而完成 (2) 的证明。(3) 的证明是直接的。
\end{proof}

命题中的第一条说明，$\Spec R$ 中每个集合 $Z(I)$ 都可由某个根理想 $I$ 得到；并且由于 $I(Z(I))=\operatorname{rad}I$，这个根理想是唯一的。

命题中的后两条说明，集合
\[
\mathcal T=\{Z(I)\mid I\text{ 是 }R\text{ 的理想}\}
\]
满足第 2 节中闭集拓扑的三个公理。

\noindent\textbf{定义.} 以 $R$ 的理想 $I$ 的集合 $Z(I)$ 为闭集定义在 $\Spec R$ 上的拓扑，称为 $\Spec R$ 上的 \emph{Zariski 拓扑}。

根据定义，$\Spec R$ 中单点集 $\{P\}$ 在 Zariski 拓扑下的闭包由 $R$ 中所有包含 $P$ 的素理想组成。特别地，$\Spec R$ 中点 $P$ 在 Zariski 拓扑下是闭点当且仅当素理想 $P$ 不包含于 $R$ 的任何其他素理想中，也即当且仅当 $P$ 是极大理想（所以 $\Spec R$ 上的 Zariski 拓扑一般不是 Hausdorff 的）。这些点有一个名称。

\noindent\textbf{定义.} $R$ 的极大理想称为 $\Spec R$ 中的\emph{闭点}。

用上述术语来说，$\Spec R$ 中在 Zariski 拓扑下闭的点，正是 $\operatorname{mSpec}R$ 中的点。

拓扑空间的闭子集称为\emph{不可约}的，如果它不能表示为两个真闭子集的并；等价地，如果每个非空开集都是稠密的。与证明命题 17 时所用论证类似，可以得到 $\Spec R$ 中的闭子集 $Y=Z(I)$ 不可约当且仅当 $I(Y)=\operatorname{rad}I$ 是素理想（参见练习 16）。

下面的命题概括了这些结果。

\noindent\textbf{命题 54.} 映射 $Z$ 与 $I$ 定义互逆的一一对应
\[
\left\{
\begin{array}{c}
	\text{$\Spec R$ 的 Zariski 闭子集}
\end{array}
\right\}
\xleftrightarrow[\quad Z\quad]{\quad I\quad}
\left\{
\begin{array}{c}
	\text{$R$ 的根理想}
\end{array}
\right\}.
\]
在这个对应下，$\Spec R$ 中的闭点对应于 $R$ 中的极大理想，而 $\Spec R$ 中的不可约子集对应于 $R$ 中的素理想。

\paragraph{例}
\begin{enumerate}[label=(\arabic*)]
	\item 若 $X=\Spec\mathbb Z$，则 $X$ 是不可约的，且非零素数给出 $X$ 中的闭点。点 $(0)$ 不是闭点；事实上，$(0)$ 的闭包是整个 $X$，即 $(0)$ 在 $\Spec\mathbb Z$ 中稠密。由于这个原因，元素 $(0)$ 称为 $\Spec\mathbb Z$ 中的\emph{一般点}。
	
	由于 $\mathbb Z$ 的每个理想都是主理想，$\Spec\mathbb Z$ 中的 Zariski 闭集为 $\varnothing$、$\Spec\mathbb Z$ 以及 $\mathbb Z$ 中非零素理想的任意有限集合。
	
	\item 设 $X=\Spec\mathbb Z[x]$，如前面例 3 所述。对每个整数素数 $p$，元素 $(p)\in X$ 的 Zariski 闭包由 (d) 型的极大理想 $(p,g)$ 组成。同样，对每个 (c) 型的 $\mathbb Q$-不可约多项式 $f$，元素 $(f)$ 的 Zariski 闭包是那些 (d) 型素理想的集合，其中 $g$ 是 $f$ 在 $\mathbb Z/p\mathbb Z[x]$ 中的某个因子。
\end{enumerate}

\paragraph{例（仿射 $k$-代数）}
设 $R=k[V]$ 是代数闭域 $k$ 上某个仿射代数集 $V\subseteq\mathbb A^n$ 的坐标环。则
\[
R=k[x_1,\ldots,x_n]/I(V),
\]
其中 $I(V)$ 是 $k[x_1,\ldots,x_n]$ 中的根理想。特别地，$R$ 是有限生成 $k$-代数；并且由于 $I(V)$ 是根理想，$R$ 不含非零幂零元。

\noindent\textbf{定义.} 代数闭域 $k$ 上一个有限生成且不含非零幂零元的代数，称为\emph{仿射 $k$-代数}。

若 $R$ 是一个仿射 $k$-代数，则由推论 5 存在满的 $k$-代数同态
\[
\pi:k[x_1,\ldots,x_n]\longrightarrow R,
\]
其核 $I=\ker\pi$ 必为根理想，因为 $R$ 不含非零幂零元。令 $V=Z(I)\subseteq\mathbb A^n$。于是
\[
R\simeq k[x_1,\ldots,x_n]/I=k[V],
\]
即 $R$ 是代数闭域 $k$ 上某个仿射代数集的坐标环。因此，仿射 $k$-代数恰好是代数闭域上仿射代数集的函数环。

由零点定理，$\operatorname{mSpec}R$ 的点与 $V$ 一一对应，而 $\Spec R$ 的点与 $V$ 的子簇一一对应。由定理 6，两个仿射代数集之间的态射与仿射 $k$-代数的（$k$-代数）同态一一对应。用范畴的语言来说，这些结果说明，在代数闭域 $k$ 上存在范畴等价
\[
\left\{
\begin{array}{c}
	\text{仿射代数集}\
	\text{代数集的态射}
\end{array}
\right\}
\longleftrightarrow
\left\{
\begin{array}{c}
	\text{仿射 $k$-代数}\
	\text{$k$-代数同态}
\end{array}
\right\}.
\]
从左到右的映射把仿射代数集 $V$ 送到它的坐标环 $k[V]$。从右到左的映射把仿射 $k$-代数 $R$ 送到 $\operatorname{mSpec}R$。有序对 $(\operatorname{mSpec}R,R)$ 有时称为仿射 $k$-代数 $R$ 的\emph{典范模型}。

在代数闭域 $k$ 上，两个仿射 $k$-代数之间的 $k$-代数同态 $\varphi:R\to S$ 如前例所述具有如下性质：由零点定理，一个 $S$ 中极大理想的逆像是 $R$ 中的极大理想。如前所述，对更一般的环考虑 $\Spec R$ 而不仅仅考虑 $\operatorname{mSpec}R$ 的一个原因，是极大理想在环同态下的逆像一般并非极大理想。当 $R$ 是对应于仿射代数集 $V$ 的仿射 $k$-代数时，空间 $\Spec R$ 不仅包含 $V$ 的“几何点”（以 $\Spec R$ 中闭点的形式），也包含对应于 $V$ 的所有子簇的非闭点（以 $\Spec R$ 中非闭点，即 $R$ 中非极大的素理想 $P$ 的形式）。

一般地，若 $\varphi:R\to S$ 是把 $1_R$ 映到 $1_S$ 的环同态，且 $P$ 是 $S$ 中的素理想，则 $\varphi^{-1}(P)$ 是 $R$ 中的素理想。这定义了映射
\[
\varphi^*:\Spec S\longrightarrow\Spec R,
\qquad
\varphi^*(P)=\varphi^{-1}(P).
\]
若 $Z(I)\subseteq\Spec R$ 是 $\Spec R$ 的 Zariski 闭子集，则容易说明 $(\varphi^*)^{-1}(Z(I))$ 是由 $S$ 中理想 $\varphi(I)S$ 定义的 Zariski 闭子集 $Z(\varphi(I)S)$。由于 $\Spec R$ 中闭子集的逆像是 $\Spec S$ 中的闭子集，诱导映射 $\varphi^*$ 在 Zariski 拓扑下连续。这证明了下面的命题。

\noindent\textbf{命题 55.} 每个把 $1_R$ 映到 $1_S$ 的环同态 $\varphi:R\to S$ 都诱导一个映射
\[
\varphi^*:\Spec S\longrightarrow\Spec R,
\]
并且这个映射关于 $\Spec R$ 与 $\Spec S$ 上的 Zariski 拓扑连续。

从仿射代数集推广到一般环 $R$ 的 $\Spec R$，虽然使事情稍微复杂了一些，但这种更一般的设置至少带来了两个很重要的优点。第一，甚至在交换环 $R$ 含有幂零元时，也可以考虑 $\Spec R$；第二，$\Spec R$ 不必来自任何域 $k$ 上的 $k$-代数；即使它是某个域上的 $k$-代数，域 $k$ 也不必代数闭。仿射 $k$-代数情形中的许多性质在这些更一般的情形中仍然成立，这就允许把“几何”思想应用到这些情形（例如当 $R$ 是有限环时的 $\Spec R$）。

\paragraph{例}
\begin{enumerate}[label=(\arabic*)]
	\item 自然包含 $\varphi:\mathbb Z\to\mathbb Z[i]$ 诱导映射 $\varphi^*:\Spec\mathbb Z[i]\to\Spec\mathbb Z$。$\varphi^*$ 在 $\mathbb Z$ 的非零素理想 $P$ 上的纤维由 $\mathbb Z[i]$ 中包含 $P$ 的素理想组成。若 $P=(p)$，其中 $p=2$ 或 $p$ 是一个模 $4$ 同余于 $3$ 的素数，则这个纤维中只有一个元素；若 $p$ 是模 $4$ 同余于 $1$ 的素数，则纤维中有两个元素：素理想 $(\pi)$ 与 $(\pi')$，其中 $p=\pi\pi'$ 在 $\mathbb Z[i]$ 中（参见第 8.3 节命题 18）。这可以用下图表示：
	
	\begin{center}
		\begin{tikzpicture}[scale=0.95, every node/.style={font=\small}]
			\draw[->] (8.4,1.4) -- (8.4,0.25) node[midway,right] {$\varphi^*$};
			\node at (8.95,1.55) {$\Spec\mathbb Z[i]$};
			\node at (8.95,0) {$\Spec\mathbb Z$};
			\draw (0,1.4) -- (8,1.4);
			\draw (0,0) -- (8,0);
			\node at (0,1.4) {$\bullet\ (0)$};
			\node at (1.6,1.4) {$\bullet\ (1+i)$};
			\node at (3.2,1.72) {$\bullet\ (2+i)$};
			\node at (3.2,1.08) {$\bullet\ (2-i)$};
			\node at (4.8,1.4) {$\bullet\ (p)$};
			\node at (6.4,1.72) {$\bullet\ (\pi)$};
			\node at (6.4,1.08) {$\bullet\ (\pi')$};
			\node at (0,0) {$\bullet$}; \node[below=3pt] at (0,0) {$(0)$};
			\node at (1.6,0) {$\bullet$}; \node[below=3pt] at (1.6,0) {$(2)$};
			\node at (3.2,0) {$\bullet$}; \node[below=3pt] at (3.2,0) {$(5)$};
			\node at (4.8,0) {$\bullet$}; \node[below=3pt] at (4.8,0) {\begin{tabular}{c}$(p)$\\$p\equiv3\ (4)$\end{tabular}};
			\node at (6.4,0) {$\bullet$}; \node[below=3pt] at (6.4,0) {\begin{tabular}{c}$(p)$\\$p\equiv1\ (4)$\end{tabular}};
			\draw[dashed] (0,1.25)--(0,0.15);
			\draw[dashed] (1.6,1.25)--(1.6,0.15);
			\draw[dashed] (3.2,1.25)--(3.2,0.15);
			\draw[dashed] (4.8,1.25)--(4.8,0.15);
			\draw[dashed] (6.4,1.25)--(6.4,0.15);
		\end{tikzpicture}
	\end{center}
	
	\item 若 $k$ 是代数闭域，则 $\Spec k[x]$ 由 $(0)$ 以及理想 $(x-a)$（$a\in k$）组成；自然包含 $\varphi:k[x]\to k[x,y]$ 诱导 Zariski 连续映射
	\[
	\varphi^*:\Spec k[x,y]\longrightarrow\Spec k[x].
	\]
	$\Spec k[x,y]$ 的元素为
	\begin{enumerate}[label=(\alph*)]
		\item $(0)$；
		\item $(f)$，其中 $f$ 是 $k[x,y]$ 中的不可约多项式；
		\item $(x-a,y-b)$，其中 $a,b\in k$
	\end{enumerate}
	（参见练习 4）。素理想 $(0)$ 在 $\Spec k[x,y]$ 中 Zariski 稠密；(b) 中的素理想的 Zariski 闭包由 (c) 中满足 $f(a,b)=0$ 的素理想 $(x-a,y-b)$ 组成；闭点，也即 $\operatorname{mSpec}k[x,y]$ 的元素，是 (c) 中的素理想。
	
	由零点定理，$\Spec k[x,y]$ 中的每个素理想 $P$ 都由相应零点集 $Z(P)$ 唯一确定。素理想 $(0)\in k[x,y]$ 对应于 $\mathbb A^2$。素理想 $(f)$ 对应于满足 $f(x,y)=0$ 的点，而 $P=(f)$ 是包含 $P$ 的所有极大理想的交。极大理想 $(x-a,y-b)$ 对应于点 $(a,b)\in\mathbb A^2$。这些素理想在映射 $\varphi^*$ 下分布在 $\Spec k[x]$ 上，可以几何地画为下图：
	
	\begin{center}
		\begin{tikzpicture}[scale=0.95, every node/.style={font=\small}]
			\draw[->] (7.5,2.3) -- (7.5,0.5) node[midway,right] {$\varphi^*$};
			\node at (8.35,2.55) {$\Spec k[x,y]$};
			\node at (8.1,0.25) {$\Spec k[x]$};
			\draw[->] (0,0.5) -- (7,0.5);
			\node at (0.2,0.85) {$(0)$};
			\node at (3.2,0.15) {$0$};
			\node at (5.2,0.15) {$a$};
			\node at (3.2,0.55) {$\bullet$};
			\node at (5.2,0.55) {$\bullet$};
			\node[below] at (3.2,0.15) {$(x)$};
			\node[below] at (5.2,0.15) {$(x-a)$};
			\draw (3.2,1.0) -- (3.2,3.4) node[above] {$(x=0)$};
			\draw (5.2,1.0) -- (5.2,3.4) node[above] {$(x=a)$};
			\draw (1.2,2.7) -- (6.5,2.7) node[right] {$(y=b)$};
			\draw[domain=1.2:6.4,smooth,variable=\x] plot ({\x},{1.45+0.35*sin(90*\x)}) node[right] {$(f)$};
			\node at (5.2,2.7) {$\bullet$};
			\node[above right] at (5.2,2.7) {$(a,b)$};
			\node at (0.9,3.2) {$(0)$};
		\end{tikzpicture}
	\end{center}
	
	在这个图中，$\Spec k[x]$ 中的素理想 $(x-a)$ 被识别为元素 $a\in k$。$\Spec k[x,y]$ 中的素理想 $(x)$ 对应于 $\mathbb A^2$ 中满足 $x=0$ 的点，即对应于 $\mathbb A^2$ 中的 $y$-轴；素理想 $(y)\in\Spec k[x,y]$ 类似地对应于 $x$-轴。素理想 $(f)\in\Spec k[x,y]$ 对应于 $\mathbb A^2$ 中的不可约曲线 $f(x,y)=0$；这条曲线上的点 $(a,b)\in\mathbb A^2$ 对应于包含 $(f)$ 的极大理想 $(x-a,y-b)\in\Spec k[x,y]$。闭点 $(x-a,y-b)\in\Spec k[x,y]$ 对应于“几何点” $(a,b)\in\mathbb A^2$。
	
	注意，$\Spec k[x,y]$ 捕捉了 $\mathbb A^2$ 中代数集的全部几何：$\mathbb A^2$ 中每个代数集都是上图中 $\Spec k[x,y]$ 的元素所对应的某些不可约代数集的有限并。除处处稠密点 $(0)$ 以外，$\Spec k[x,y]$ 的“几何”图像正是仿射平面 $\mathbb A^2$ 的通常几何。若 $k$ 不是代数闭的，情形稍微复杂一些，但图像相似，参见练习 4。
	
	\item 把 $\Spec\mathbb Z[x]$ 看成由自然包含 $\mathbb Z\to\mathbb Z[x]$ 纤维化在 $\Spec\mathbb Z$ 上，其情形与前一例中 $\Spec k[x,y]$ 的情形非常类似。$\Spec\mathbb Z[x]$ 的元素已在命题 54 后的例 2 中讨论过，可画为下图：
	
	\begin{center}
		\begin{tikzpicture}[scale=0.95, every node/.style={font=\small}]
			\draw[->] (7.8,2.0) -- (7.8,0.55) node[midway,right] {$\varphi^*$};
			\node at (8.65,2.25) {$\Spec\mathbb Z[x]$};
			\node at (8.1,0.25) {$\Spec\mathbb Z$};
			\draw (0,0.55) -- (7.2,0.55);
			\foreach \x/\lab in {0/{(0)},1.5/{(2)},3/{(3)},4.5/{(5)},6.4/{(p)}}
			{\node at (\x,0.55) {$\bullet$}; \node[below=3pt] at (\x,0.55) {\lab};}
			\node at (0,2.0) {$\bullet\ (0)$};
			\draw[domain=1.1:6.6,smooth,variable=\x] plot ({\x},{1.25+0.16*(\x-3.8)*(\x-3.8)}) node[right] {$(f)$};
			\node at (5.7,2.4) {$\bullet\ (x-a)$};
			\node at (1.5,1.4) {$\bullet$};
			\node at (3.0,1.35) {$\bullet$};
			\node at (3.2,1.55) {$\bullet$};
			\node at (4.35,1.25) {$\bullet$};
			\node at (4.55,1.45) {$\bullet$};
			\node at (4.75,1.65) {$\bullet$};
			\node at (6.35,1.7) {$\bullet$};
			\draw[dashed] (1.5,1.15)--(1.5,0.72);
			\draw[dashed] (3,1.15)--(3,0.72);
			\draw[dashed] (4.5,1.15)--(4.5,0.72);
			\draw[dashed] (6.4,1.15)--(6.4,0.72);
		\end{tikzpicture}
	\end{center}
	
	元素 $(0)$ 在 $\Spec\mathbb Z[x]$ 中 Zariski 稠密。$(p)$ 的闭包由 $(p)$ 以及所有闭点 $(p,g)$ 组成，其中 $g$ 是 $\mathbb Z[x]$ 中模 $p$ 不可约的首一多项式。$(f)$ 的闭包由 $(f)$ 以及包含 $(f)$ 的极大理想 $(p,g)$ 组成，这等价于说 $f$ 在商环 $\mathbb Z[x]/(p,g)$ 中的像为 $0$，即不可约多项式 $g$ 是 $f$ 模 $p$ 的一个因子。闭点 $\operatorname{mSpec}\mathbb Z[x]$ 是极大理想 $(p,g)$。
	
	注意，包含 $(f)$ 的极大理想 $(p,g)$ 正是上图中 $\operatorname{mSpec}\mathbb Z[x]$ 中使 $\Spec\mathbb Z[x]$ 上的“函数” $f$ 为零的闭点，这里的函数把素理想 $P$ 送到
	\[
	f(P)=f\bmod P\in\mathbb Z[x]/P.
	\]
	例如，多项式
	\[
	f=x^3-4x^2+x-9\in\mathbb Z[x]
	\]
	符合上图：$f$ 在 $\mathbb Z[x]$ 中不可约，并且在 $\mathbb F_p$ 上分解成不可约因子的情形如下：
	\[
	\begin{aligned}
		f&=x^3+x+1\pmod 2,\\
		f&=x(x+1)^2\pmod 3,\\
		f&=(x+1)(x+2)(x+3)\pmod 5.
	\end{aligned}
	\]
	与 $(f)$ 相交且位于 $(2)$ 上方纤维中的点有一个，即闭点 $(2,x^3+x+1)$。位于 $(3)$ 上方纤维中的闭点有两个，分别为 $(3,x)$ 与 $(3,x+1)$（后一个点带有某种“重数”）。位于 $(5)$ 上方的闭点有三个：$(5,x+1)$、$(5,x+2)$ 与 $(5,x+3)$。在上图中，素数 $p$ 可以取 $p=53$，因为这是大于 $5$ 的第一个使此多项式模 $p$ 有三个不可约因子的素数。注意，虽然图中把素理想 $(f)$ 画成一条光滑曲线，以强调它与前一例中 $\Spec k[x,y]$ 结构的几何相似性，但 $\Spec\mathbb Z$ 中素数上方的纤维是离散的，因此应当谨慎。例如，由于
	\[
	f=(x+2)(x^2+x+6)\pmod 7,
	\]
	$(f)$ 与 $(7)$ 上方纤维的交只含有两个点 $(7,x+2)$ 与 $(7,x^2+x+6)$，每个重数为一。
	
	位于 $\Spec\mathbb Z$ 中 $(p)$ 上方纤维内的 $(f)$ 上的闭点可能个数由多项式 $f$ 在 $\mathbb Q$ 上的 Galois 群控制（参见第 14.8 节）。例如，$f=x^4+1$ 在 $(2)$ 上方纤维中有一个闭点，而对奇素数 $p$，在 $(p)$ 上方纤维中有两个或四个闭点（参见练习 8）。
\end{enumerate}

带有 Zariski 拓扑的空间 $\Spec R$，给出了任意交换环情形下对簇 $V$ 中点的几何推广。现在考虑如何推广 $V$ 上有理函数环的问题。

当 $V$ 是代数闭域 $k$ 上的簇时，坐标环 $k[V]$ 的商域 $k(V)$ 中的元素定义了 $V$ 上的有理函数。$k(V)$ 中每个元素 $\alpha$ 一般可以用许多不同方式写为 $a/f$，其中 $a,f\in k[V]$。$\alpha$ 正则的点的集合 $U$ 是 $V$ 的一个开子集；按定义，它由所有满足 $\alpha$ 可以表示成某个商 $a/f$ 且 $f(v)\ne0$ 的点 $v\in V$ 组成，此时代表元 $a/f$ 定义了局部环 $\mathcal O_{v,V}$ 中的元素。还要注意，同一个代表元 $a/f$ 不仅在 $v$ 处定义 $\alpha$，而且在 $f$ 非零的所有点上都定义 $\alpha$，即在 $V$ 的开子集
\[
V_f=\{w\in V\mid f(w)\ne0\}
\]
上定义。这些开集 $V_f$（称为主开集，参见第 2 节练习 21）随着 $\alpha$ 的各种可能代表元 $a/f$ 而给出 $U$ 的一个开覆盖。命题 51 前的例子中，$V=Z(xz-yw)\subseteq\mathbb A^4$ 上函数 $\alpha=x/y$ 表明，一般情况下，单个代表元并不足以确定整个 $U$；在这个例子中，$U=V_y\cup V_z$，而 $U$ 不能由任何单个 $V_f$ 覆盖（参见第 4 节练习 25）。

把有理函数理解为某些开子集上的正则函数，这一解释可以推广到 $\Spec R$。我们首先定义 $X=\Spec R$ 中集合 $V_f$ 的类似物 $X_f$，并建立其基本性质。

\noindent\textbf{定义.} 对任意 $f\in R$，令 $X_f$ 表示 $X=\Spec R$ 中不包含 $f$ 的素理想的集合。等价地，$X_f$ 是 $\Spec R$ 中元素 $f\in R$ 的值非零的点的集合。集合 $X_f$ 称为 $\Spec R$ 中的\emph{主开集}（或\emph{基本开集}）。

由于 $X_f$ 是 Zariski 闭集 $Z(f)$ 的补集，所以它确实是 $\Spec R$ 中的开集，正如名称所示。下面的命题列出主开集的一些基本性质。回忆，拓扑空间之间的映射称为同胚，如果它连续、双射且逆映射也连续。

\noindent\textbf{命题 56.} 设 $f\in R$，并设 $X_f$ 是 $X=\Spec R$ 中相应的主开集。则：
\begin{enumerate}[label=(\arabic*)]
	\item $X_f=X$ 当且仅当 $f$ 是单位，而 $X_f=\varnothing$ 当且仅当 $f$ 是幂零元；
	\item $X_f\cap X_g=X_{fg}$；
	\item $X_f\subseteq X_{g_1}\cup\cdots\cup X_{g_n}$ 当且仅当
	\[
	f\in\operatorname{rad}(g_1,\ldots,g_n);
	\]
	特别地，$X_f=X_g$ 当且仅当 $\operatorname{rad}(f)=\operatorname{rad}(g)$；
	\item 主开集构成 $\Spec R$ 上 Zariski 拓扑的一组基，即 $X$ 的每个 Zariski 开集都是某些主开集 $X_f$ 的并；
	\item 从 $R$ 到 $R_f$ 的自然映射诱导一个从 $\Spec R_f$ 到 $X_f$ 的同胚，其中 $R_f$ 是 $R$ 在 $f$ 处的局部化；
	\item 任意环的谱都是拟紧的（即每个开覆盖都有有限子覆盖）；特别地，$X_f$ 是拟紧的；
	\item 若 $\varphi:R\to S$ 是任意环同态（且 $\varphi(1_R)=1_S$），则在诱导映射
	\[
	\varphi^*:Y=\Spec S\longrightarrow\Spec R
	\]
	下，$X$ 中主开集 $X_f$ 的全逆像是 $Y$ 中的主开集 $Y_{\varphi(f)}$。
\end{enumerate}

\begin{proof}
	(1)、(2) 与 (7) 留作简单练习。为证明 (3)，注意根据定义，$X_{g_1}\cup\cdots\cup X_{g_n}$ 由那些不包含 $g_1,\ldots,g_n$ 中至少一个元素的素理想 $P$ 组成。因此 $X_{g_1}\cup\cdots\cup X_{g_n}$ 是闭集 $Z((g_1,\ldots,g_n))$ 的补集，其中后者由包含理想 $(g_1,\ldots,g_n)$ 的素理想 $P$ 组成。若 $(g_1,\ldots,g_n)=R$，则 $X_{g_1}\cup\cdots\cup X_{g_n}=X$，无须证明。否则，$X_f\subseteq X_{g_1}\cup\cdots\cup X_{g_n}$ 当且仅当每个满足 $f\notin P$ 的素理想 $P$ 也满足 $P\notin Z((g_1,\ldots,g_n))$。后一条件等价于说：如果素理想 $P$ 包含理想 $(g_1,\ldots,g_n)$，则 $P$ 也包含 $f$；即 $f$ 属于所有包含 $(g_1,\ldots,g_n)$ 的素理想 $P$ 的交。由命题 12，这个交为 $\operatorname{rad}(g_1,\ldots,g_n)$，从而证明 (3)。
	
	若 $U=X-Z(I)$ 是 $X$ 的一个 Zariski 开子集，则 $U$ 是所有 $X_f$（$f\in I$）的并，这证明 (4)。
	
	从 $R$ 到局部化 $R_f$ 的自然环同态建立了 $R_f$ 中素理想与 $R$ 中不包含 $(f)$ 的素理想之间的一一对应（命题 38）。因此，从 $\Spec R_f$ 到 $\Spec R$ 的相应 Zariski 连续映射是连续双射。由于 $R_f$ 的每个理想都是 $R$ 的某个理想的扩张（参见命题 38(1)），可知逆映射也连续，这证明 (5)。
	
	对 (6)，由 (4) 每个开集都是主开集的并，因此只需证明：若 $X$ 被主开集 $X_{g_i}$（$i$ 属于某个指标集 $\mathcal J$）覆盖，则 $X$ 是其中有限多个 $X_{g_i}$ 的并。若由所有 $g_i$ 生成的理想 $I$ 是 $R$ 的真理想，则 $I$ 包含于某个极大理想 $P$ 中。但这时 $X=\Spec R$ 中的点 $P$ 不属于任何主开集 $X_{g_i}$，与 $X$ 被这些 $X_{g_i}$ 覆盖相矛盾。因此 $I=R$，所以 $1\in R$ 可写为有限和
	\[
	1=a_1g_{i_1}+\cdots+a_ng_{i_n}
	\]
	其中 $i_1,\ldots,i_n\in\mathcal J$。考虑有限并 $X_{g_{i_1}}\cup\cdots\cup X_{g_{i_n}}$。若 $X$ 中某点 $P$ 不属于这个并，则 $P$ 是 $R$ 中包含 $g_{i_1},\ldots,g_{i_n}$ 的素理想，从而包含 $1$，矛盾。因此
	\[
	X=X_{g_{i_1}}\cup\cdots\cup X_{g_{i_n}},
	\]
	如所需。(6) 的第二部分由 (5) 推出。
\end{proof}

现在我们为 $X=\Spec R$ 定义簇 $V$ 上有理函数的类似物。如前所见，对簇 $V$，有理函数 $\alpha\in k(V)$ 是某个开集 $U$ 上的正则函数。在每个点 $v\in U$ 处，存在 $\alpha$ 的代表元 $a/f$ 且 $f(v)\ne0$；这个代表元是局部化
\[
\mathcal O_{v,V}=k[V]_{I(v)}
\]
中的元素。这样，$U$ 上的正则函数 $\alpha$ 可以看作从 $U$ 到这些局部化的不交并的函数：点 $v\in U$ 被映到代表元 $a/f\in k[V]_{I(v)}$。此外，同一个代表元不仅可以在 $v$ 处使用，也可以在 $v$ 的整个 Zariski 邻域 $V_f$ 上同时使用（所以“在 $v$ 附近局部地”，$\alpha$ 由 $k[V]$ 中元素的单个商给出）。注意，$a/f$ 是局部化 $k[V]_f$ 中的元素，而后者包含于所有 $w\in V_f$ 的局部化 $k[V]_{I(w)}$ 中。

现在我们把这推广到 $\Spec R$。考虑从 $\Spec R$ 的 Zariski 开子集 $U$ 到局部化 $R_P$（$P\in U$）的不交并的函数 $s$ 的集合，要求 $s(P)\in R_P$，并且 $s$ 局部地由 $R$ 中元素的商给出。更精确地说：

\noindent\textbf{定义.} 设 $U$ 是 $\Spec R$ 的 Zariski 开子集。若 $U=\varnothing$，定义 $\mathcal O(U)=0$。否则，定义 $\mathcal O(U)$ 为所有函数
\[
s:U\longrightarrow\coprod_{Q\in U}R_Q
\]
的集合，其中 $s$ 从 $U$ 到 $Q\in U$ 的局部化 $R_Q$ 的不交并，并且满足如下两个性质：
\begin{enumerate}[label=(\arabic*)]
	\item 对每个 $Q\in U$，有 $s(Q)\in R_Q$；
	\item 对每个 $P\in U$，存在 $P$ 在 $U$ 中的一个开邻域 $X_f\subseteq U$ 以及局部化 $R_f$ 中的一个元素 $a/f^n$，在 $X_f$ 上定义 $s$，即对每个 $Q\in X_f$，有
	\[
	s(Q)=a/f^n\in R_Q.
	\]
\end{enumerate}

若 $s,t$ 是 $\mathcal O(U)$ 的元素，则 $s+t$ 与 $st$ 也都是 $\mathcal O(U)$ 的元素（参见练习 18），所以每个 $\mathcal O(U)$ 都是一个环。此外，每个 $a\in R$ 定义了 $\mathcal O(U)$ 中的一个元素，即 $s(Q)=a\in R_Q$；特别地，$1\in R$ 给出了环 $\mathcal O(U)$ 的单位元。若 $U'$ 是 $U$ 的一个开子集，则存在从 $\mathcal O(U)$ 到 $\mathcal O(U')$ 的自然限制映射，它是环同态（参见练习 19）。

\noindent\textbf{定义.} 设 $R$ 是含 $1$ 的交换环，并令 $X=\Spec R$。
\begin{enumerate}[label=(\arabic*)]
	\item 对 $X$ 的 Zariski 开集 $U$ 给出的环 $\mathcal O(U)$ 以及当 $U'\subseteq U$ 时的限制映射 $\mathcal O(U)\to\mathcal O(U')$，称为 $X$ 上的\emph{结构层}，简单记作 $\mathcal O$（或 $\mathcal O_X$）。
	\item $\mathcal O(U)$ 的元素 $s$ 称为 $\mathcal O$ 在 $U$ 上的\emph{截面}。$\mathcal O(X)$ 的元素称为 $\mathcal O$ 的\emph{整体截面}。
\end{enumerate}

下一个命题推广了命题 51 的结果：簇 $V$ 上处处正则的有理函数恰好是坐标环 $k[V]$ 的元素。

\noindent\textbf{命题 57.} 设 $X=\Spec R$，并设 $\mathcal O=\mathcal O_X$ 是其结构层。$\mathcal O$ 的整体截面就是 $R$ 的元素，即
\[
\mathcal O(X)\simeq R.
\]
更一般地，若 $X_f$ 是 $X$ 中由某个 $f\in R$ 给出的主开集，则 $\mathcal O(X_f)$ 同构于局部化 $R_f$。

\begin{proof}
	设 $a/f^r$ 是局部化 $R_f$ 的一个元素。则由
	\[
	s(Q)=a/f^r\in R_Q,
	\qquad Q\in X_f,
	\]
	定义的映射给出 $\mathcal O(X_f)$ 中的一个元素；显然得到的映射
	\[
	\psi:R_f\longrightarrow\mathcal O(X_f)
	\]
	是环同态。假设对每个 $Q\in X_f$，有 $a/f^r=b/f^s$ 在 $R_Q$ 中成立，即存在某个 $g\notin Q$ 使得
	\[
	g(af^s-bf^r)=0
	\]
	在 $R$ 中成立。令 $I$ 为 $R$ 中所有满足
	\[
	x(af^s-bf^r)=0
	\]
	的元素 $x$ 组成的理想。由 $g\in I$ 可知，对任意 $Q\in X_f$，$I$ 不包含于 $Q$。换句话说，$R$ 中每个包含 $I$ 的素理想也包含 $f$。因此 $f$ 属于所有包含 $I$ 的素理想的交，也就是说 $f\in\operatorname{rad}I$。于是对某个整数 $N\ge0$，有 $f^N\in I$，所以
	\[
	f^N(af^s-bf^r)=0
	\]
	在 $R$ 中成立。这说明 $a/f^r=b/f^s$ 在 $R_f$ 中成立，因此映射 $\psi$ 是单射。
	
	现在设 $s\in\mathcal O(X_f)$。根据定义，$X_f$ 可以被主开集 $X_{g_i}$ 覆盖，在每个 $X_{g_i}$ 上，对每个 $Q\in X_{g_i}$，有
	\[
	s(Q)=a_i/g_i^{n_i}\in R_Q.
	\]
	由命题 56(6)，可以取有限多个 $g_i$；再通过改变 $a_i$，可以假设所有 $n_i$ 都相等（因为若 $n$ 是所有 $n_i$ 的最大值，则
	\[
	a_i/g_i^{n_i}=(a_ig_i^{n-n_i})/g_i^n
	\]
	）。由于对所有 $Q\in X_{g_ig_j}=X_{g_i}\cap X_{g_j}$，有
	\[
	s(Q)=a_i/g_i^n=a_j/g_j^n\in R_Q,
	\]
	把单射性应用于 $R_{g_ig_j}$ 可知
	\[
	a_i/g_i^n=a_j/g_j^n\quad\text{在 }R_{g_ig_j}\text{ 中}.
	\]
	这意味着对某个 $N\ge0$，有
	\[
	(g_ig_j)^N(a_ig_j^n-a_jg_i^n)=0,
	\]
	即
	\[
	a_i g_j^{n+N}g_i^N=a_j g_i^{n+N}g_j^N
	\]
	在 $R$ 中成立；可以假设 $N$ 充分大，使得这个等式对每一对 $i,j$ 都成立。
	
	由于 $X_f$ 是 $X_{g_i}=X_{g_i^{n+N}}$ 的并，由命题 56(3)，$f$ 属于由 $g_i^{n+N}$ 生成的理想的根。于是存在 $M\ge1$ 与 $b_i\in R$，使得
	\[
	f^M=\sum_i b_i g_i^{n+N}.
	\]
	定义
	\[
	a=\sum_i b_i a_i g_i^N\in R.
	\]
	于是
	\[
	g_j^N a_j f^M
	=\sum_i b_i(a_j g_i^{n+N}g_j^N)
	=\sum_i b_i(a_i g_i^N g_j^{n+N})
	=g_j^{n+N}a.
	\]
	因此 $a/f^M=a_j/g_j^n$ 在 $R_{g_j}$ 中成立。所以由 $R_f$ 中元素 $a/f^M$ 定义的 $\mathcal O(X_f)$ 中的元素，在每个 $X_{g_j}$ 上都与 $s$ 相同；由于这些开集覆盖 $X_f$，它在整个 $X_f$ 上与 $s$ 相同。因此映射 $\psi$ 给出同构
	\[
	R_f\simeq\mathcal O(X_f).
	\]
	取 $f=1$ 得 $R\simeq\mathcal O(X)$，证明完毕。
\end{proof}

在仿射簇 $V$ 的情形中，点 $v\in V$ 处的局部环 $\mathcal O_{v,V}$ 是 $k(V)$ 中在 $v$ 处有定义的所有有理函数的集合。换句话说，$\mathcal O_{v,V}$ 是所有包含 $v$ 的开集 $U$ 上正则函数环的并，其中这个并是在 $V$ 的函数域 $k(V)$ 中取的。在更一般的情形 $X=\Spec R$ 中，包含 $P\in\Spec R$ 的开集上的环 $\mathcal O(U)$ 并不包含在如此显然的公共环中。在这种情况下，我们考虑所有对 $(s,U)$ 的集合，其中 $U$ 是包含 $P$ 的 $X$ 的开集且 $s\in\mathcal O(U)$。如果存在包含 $P$ 的开集 $U''\subseteq U\cap U'$，使得 $s$ 与 $s'$ 限制到 $\mathcal O(U'')$ 后相同，则把两个对 $(s,U)$ 与 $(s',U')$ 视为等价。在仿射簇的情形中，这就是说，在点 $v$ 的 Zariski 邻域中定义的两个函数，如果在 $v$ 的某个公共邻域中相同，则它们在 $v$ 处定义同一个正则函数。这些等价类组成的集合定义了环 $\mathcal O(U)$ 的正向极限，记作 $\varinjlim_P\mathcal O(U)$（参见第 7.6 节练习 8）。

\noindent\textbf{定义.} 若 $P\in X=\Spec R$，则所有包含 $P$ 的 $X$ 的开集 $U$ 上的环 $\mathcal O(U)$ 的正向极限
\[
\varinjlim_P\mathcal O(U)
\]
称为结构层在 $P$ 处的\emph{茎}，记作 $\mathcal O_P$。

\noindent\textbf{命题 58.} 设 $X=\Spec R$，并设 $\mathcal O=\mathcal O_X$ 是其结构层。$\mathcal O$ 在点 $P\in X$ 处的茎同构于 $R$ 在 $P$ 处的局部化 $R_P$：
\[
\mathcal O_P\simeq R_P.
\]
特别地，茎 $\mathcal O_P$ 是局部环。

\begin{proof}
	若 $(s,U)$ 代表茎 $\mathcal O_P$ 中的一个元素，则 $s(P)$ 是局部化 $R_P$ 的一个元素。由正向极限的定义，这个元素不依赖于代表元 $(s,U)$ 的选择，因此给出从 $\mathcal O_P$ 到 $R_P$ 的良定义环同态 $\phi$。若 $a,f\in R$ 且 $f\notin P$，则映射
	\[
	s(Q)=a/f\in R_Q
	\]
	定义了 $\mathcal O(X_f)$ 中的一个元素。那么 $(s,X_f)$ 在茎 $\mathcal O_P$ 中的等价类在 $\phi$ 下被映到 $R_P$ 中的 $a/f$，所以 $\phi$ 是满射。
	
	为说明 $\phi$ 也是单射，设 $\mathcal O_P$ 中 $(s,U)$ 与 $(s',U')$ 的等价类满足
	\[
	s(P)=s'(P)\quad\text{在 }R_P\text{ 中}.
	\]
	由 $\mathcal O(U)$ 的定义，在某个 $g\notin P$ 的 $X_g$ 上，$s=a/g^n$。类似地，在某个 $g'\notin P$ 的 $X_{g'}$ 上，$s'=b/(g')^m$。由于
	\[
	a/g^n=b/(g')^m\quad\text{在 }R_P\text{ 中},
	\]
	存在某个 $h\notin P$，使得
	\[
	h(a(g')^m-bg^n)=0
	\]
	在 $R$ 中成立。若 $Q\in X_{gg'h}=X_g\cap X_{g'}\cap X_h$，这个等式说明
	\[
	a/g^n=b/(g')^m\quad\text{在 }R_Q\text{ 中}.
	\]
	因此 $s$ 与 $s'$ 限制到 $X_{gg'h}$ 后相同。由正向极限的定义，$(s,U)$ 与 $(s',U')$ 在茎 $\mathcal O_P$ 中定义同一个元素。这证明 $\phi$ 是单射，从而证明命题。
\end{proof}

命题 58 表明，对 $P\in\Spec R$ 代数地定义的局部化 $R_P$ 起到了仿射簇 $V$ 上点 $v$ 处正则函数局部环 $\mathcal O_{v,V}$ 的作用。若 $\mathfrak m_P$ 表示 $R_P$ 中的极大理想 $PR_P$，并且
\[
k(P)=R_P/\mathfrak m_P
\]
表示相应的商域（由命题 46(1)，它也是 $R/P$ 的分式域），则 $P$ 处的切空间定义为 $k(P)$-向量空间 $\mathfrak m_P/\mathfrak m_P^2$ 的对偶。这是一个代数定义，它推广了簇 $V$ 在点 $v$ 处切空间 $T_{v,V}$ 的定义（由命题 52）。现在可以用它来定义 $\Spec R$ 中的点非奇异是什么意思：点 $P\in\Spec R$ 称为\emph{非奇异}或\emph{光滑}的，如果局部环 $R_P$ 是所谓的“正则局部环”（参见第 16.2 节）。

命题 58 还提示了结构层在 $\Spec R$ 上的一个漂亮几何图像。如果把每个点 $P\in\Spec R$ 看成其上方有局部环 $R_P$，那么在 $X=\Spec R$ 的开集 $U$ 上方，就有这些“茎”的一个“层”（这里的“层”可理解为“丛”，而“茎”可理解为“麦秆的茎”），这有助于解释一些术语。结构层 $\mathcal O(U)$ 中的一个截面 $s$ 是从 $U$ 到这些茎所成丛的映射。下图中阴影区域表示这样一个截面 $s$ 在 $U$ 下的像。

\begin{center}
	\begin{tikzpicture}[scale=0.9, every node/.style={font=\small}]
		\draw[thick] (0,0) .. controls (1,-0.3) and (3,-0.3) .. (4,0)
		.. controls (4.7,0.5) and (4.4,1.1) .. (3.6,1.2)
		.. controls (2.2,1.45) and (1.2,1.4) .. (0.3,1.05)
		.. controls (-0.2,0.8) and (-0.2,0.3) .. (0,0);
		\node at (2,-0.25) {$U$};
		\foreach \x in {0.5,1.1,1.7,2.3,2.9,3.5}
		\draw[gray] (\x,0.15) -- (\x,3.0);
		\draw[very thick, fill=gray!25, opacity=0.8] (0.55,1.35) .. controls (1.2,1.7) and (2.1,1.75) .. (2.8,1.55) .. controls (3.2,1.45) and (3.55,1.65) .. (3.7,1.9) -- (3.7,2.15) .. controls (3.2,1.85) and (2.7,1.95) .. (2.2,2.05) .. controls (1.4,2.15) and (0.9,1.95) .. (0.55,1.6) -- cycle;
		\node at (4.35,2.0) {$s$};
		\node at (3.95,3.05) {$R_P$};
	\end{tikzpicture}
\end{center}

\noindent\textbf{定义.} 设 $R$ 是含 $1$ 的交换环。有序对
\[
(\Spec R,\mathcal O_{\Spec R}),
\]
即带有 Zariski 拓扑的空间 $\Spec R$ 连同结构层 $\mathcal O_{\Spec R}$，称为一个\emph{仿射概形}。

仿射概形的概念给出了仿射代数集几何对任意交换环完全代数化的推广，并且是现代代数几何的出发点。

\paragraph{例}
\begin{enumerate}[label=(\arabic*)]
	\item 若 $F$ 是任意域，则 $X=\Spec F=\{(0)\}$。在这种情形下只有两个开集 $X$ 与 $\varnothing$，它们都是主开集：$X=X_1$ 且 $\varnothing=X_0$。整体截面为 $\mathcal O(X)=F$。只有一个茎：$\mathcal O_{(0)}=F_{(0)}=F$。
	
	\item 若 $R=\mathbb Z$，则由于 $R$ 是 P.I.D.，$X=\Spec\mathbb Z$ 的每个开集都是主开集：
	\[
	X_n=\{(p)\mid p\nmid n\}
	\]
	并且
	\[
	\mathcal O(X_n)=\mathbb Z_n=\mathbb Z[1/n]
	=\left\{a/b\in\mathbb Q\mid \text{若素数 }p\mid b\text{，则 }p\mid n\right\}.
	\]
	对非零素数 $p$，$(p)$ 处的茎是局部环 $\mathbb Z_{(p)}$；$(0)$ 处的茎是 $\mathbb Q$。所有限制映射以及从截面到茎的映射都是自然包含。
	
	\item 对一般整环 $R$，设其商域为 $F$，则茎与截面为
	\[
	\mathcal O(U)=\{a/b\in F\mid b\notin P\text{ 对所有 }P\in U\text{ 成立}\},
	\]
	\[
	\mathcal O_P=R_P=\{a/b\in F\mid b\notin P\},
	\]
	其中 $(0)$ 处的茎是 $F$，即 $\mathcal O_{(0)}=F$。同样，限制映射以及到茎的映射都是包含。
	
	\item 对局部环
	\[
	R=\mathbb Z_{(2)}=\{a/b\in\mathbb Q\mid b\text{ 为奇数}\},
	\]
	有 $\Spec R=\{(0),(2)\}$，其中 $(2)$ 是唯一闭点，且 $\{(0)\}=X_2$ 是主开集。截面 $\mathcal O(\{(0)\})$ 为 $R_2=\mathbb Q$，茎为
	\[
	\mathcal O_{(0)}=R_{(0)}=\mathbb Q,
	\qquad
	\mathcal O_{(2)}=R_{(2)}=R.
	\]
\end{enumerate}

接下来考虑由环同态 $R\to S$ 所产生的对应于 $R$ 与 $S$ 的仿射概形之间的关系。

设 $\varphi:R\to S$ 是一个环同态。我们已在命题 56(7) 中看到，存在诱导连续映射
\[
\varphi^*:Y=\Spec S\longrightarrow X=\Spec R,
\]
并且在此映射下，对 $g\in R$，主开集 $X_g$ 的全逆像是主开集 $Y_{\varphi(g)}$。由此可知，$\varphi$ 也诱导相应截面上的映射，如下所述。设 $Q'\in Y$ 是 $\Spec S$ 的任意元素，并设
\[
Q=\varphi^*(Q')=\varphi^{-1}(Q')\in X
\]
是 $\Spec R$ 中对应的元素。若 $U$ 是 $X$ 中包含 $Q$ 的 Zariski 开集，则
\[
U'=(\varphi^*)^{-1}(U)
\]
是 $Y$ 中包含 $Q'$ 的 Zariski 开集。注意，$\varphi$ 诱导一个自然环同态（记为 $\varphi_Q$）从局部化 $R_Q$ 到局部化 $S_{Q'}$，定义为
\[
\varphi_Q(a/f)=\varphi(a)/\varphi(f)\in S_{Q'}
\]
其中 $f\notin Q$。设 $s\in\mathcal O_X(U)$ 是 $X$ 的结构层的一个截面，并且在 $P\in X$ 的邻域 $X_g$ 中局部由 $a/g^n$ 给出。容易检查复合
\[
s':U'\longrightarrow U\longrightarrow\coprod_{Q\in U}R_Q
\longrightarrow\coprod_{Q'\in U'}S_{Q'}
\]
定义了一个映射，它在邻域 $Y_{\varphi(g)}$ 中局部由元素 $\varphi(a)/\varphi(g)^n$ 给出；因此 $s'\in\mathcal O_Y(U')$ 是 $Y$ 的结构层的一个截面。于是可以直接检查，所得映射
\[
\varphi^\#:\mathcal O_X(U)\longrightarrow\mathcal O_Y(U')
\]
是环同态（把 $\mathcal O_X(U)$ 中的 $1$ 映到 $\mathcal O_Y(U')$ 中的 $1$），并且与 $\mathcal O_X$ 与 $\mathcal O_Y$ 上的限制映射相容（参见练习 20）。同时还得到茎上的诱导环同态：对任意点 $P'\in\Spec S$ 以及对应点 $P=\varphi^*(P')\in\Spec R$，有
\[
\varphi^\#: \mathcal O_{X,P}\longrightarrow\mathcal O_{Y,P'}.
\]
在命题 58 的同构下，从 $R_P\simeq\mathcal O_{X,P}$ 到 $S_{P'}\simeq\mathcal O_{Y,P'}$ 的同态 $\varphi^\#$ 正是由同态 $\varphi$ 诱导出的局部化上的自然环同态 $\varphi_P$。特别地，在 $\varphi^\#$ 下，局部环 $\mathcal O_{Y,P'}$ 中极大理想的逆像是局部环 $\mathcal O_{X,P}$ 中的极大理想。

\noindent\textbf{定义.} 设 $(\Spec R,\mathcal O_{\Spec R})$ 与 $(\Spec S,\mathcal O_{\Spec S})$ 是两个仿射概形。从 $(\Spec S,\mathcal O_{\Spec S})$ 到 $(\Spec R,\mathcal O_{\Spec R})$ 的\emph{仿射概形态射}是一个有序对 $(\varphi^*,\varphi^\#)$，满足：
\begin{enumerate}[label=(\arabic*)]
	\item $\varphi^*:\Spec S\to\Spec R$ 是 Zariski 连续的；
	\item 对 $\Spec R$ 的每个 Zariski 开子集 $U$，都有环同态
	\[
	\varphi^\#:\mathcal O(U)\longrightarrow\mathcal O((\varphi^*)^{-1}(U)),
	\]
	并且这些同态与限制映射交换；
	\item 若 $P'\in\Spec S$ 且对应点 $P=\varphi^*(P')\in\Spec R$，则在茎上的诱导同态
	\[
	\varphi^\#: \mathcal O_{\Spec R,P}\longrightarrow \mathcal O_{\Spec S,P'}
	\]
	下，$\mathcal O_{\Spec S,P'}$ 的极大理想的逆像是 $\mathcal O_{\Spec R,P}$ 的极大理想。
\end{enumerate}

从局部环 $A$ 到局部环 $B$ 的同态 $\psi:A\to B$，如果具有 $B$ 的极大理想的逆像为 $A$ 的极大理想这一性质，则称为局部环的\emph{局部同态}。上述定义中的第三个条件就是要求茎上的诱导同态为局部同态。

用这个术语，定义前的讨论说明，一个环同态 $\varphi:R\to S$ 诱导一个从 $(\Spec S,\mathcal O_{\Spec S})$ 到 $(\Spec R,\mathcal O_{\Spec R})$ 的仿射概形态射。

反过来，设 $(\varphi^*,\varphi^\#)$ 是从 $(\Spec S,\mathcal O_{\Spec S})$ 到 $(\Spec R,\mathcal O_{\Spec R})$ 的一个仿射概形态射。特别地，对 $U=\Spec R$，有 $(\varphi^*)^{-1}(U)=\Spec S$；因此按假设，在整体截面上存在环同态
\[
\varphi^\#: \mathcal O_{\Spec R}(\Spec R)
\longrightarrow
\mathcal O_{\Spec S}(\Spec S).
\]
由命题 57，有环同构
\[
\mathcal O_{\Spec R}(\Spec R)\simeq R,
\qquad
\mathcal O_{\Spec S}(\Spec S)\simeq S.
\]
与这些同构复合说明，$\varphi^\#$ 给出一个环同态 $\varphi:R\to S$。由命题 58，我们有一个局部同态 $\varphi^\#:R_P\to S_{P'}$；并且由它与限制同态的相容性可知，下图交换：
\[
\begin{array}{ccc}
	R & \xrightarrow{\varphi} & S\\
	\downarrow & & \downarrow\\
	R_P & \xrightarrow{\varphi^\#} & S_{P'}
\end{array}
\]
其中两个竖直映射是自然局部化同态。由于假设 $\varphi^\#$ 是局部同态，
\[
(\varphi^\#)^{-1}(P'S_{P'})=PR_P,
\]
由此推出 $\varphi^{-1}(P')=P$。因此，由 $\varphi$ 诱导的从 $\Spec S$ 到 $\Spec R$ 的连续映射与 $\varphi^*$ 相同；并且容易推出 $\varphi$ 也如前所述诱导同态 $\varphi^\#$。这说明存在环同态 $\varphi:R\to S$，同时诱导出 $\varphi^*$ 与 $\varphi^\#$。

我们把这概括为如下定理。

\noindent\textbf{定理 59.} 每个环同态 $\varphi:R\to S$ 都诱导一个仿射概形态射
\[
(\varphi^*,\varphi^\#):(\Spec S,\mathcal O_{\Spec S})
\longrightarrow
(\Spec R,\mathcal O_{\Spec R}).
\]
反过来，每个仿射概形态射都来自这样的一个环同态 $\varphi$。

定理 59 是 $\Spec R$ 情形中定理 6 的类似物；定理 6 把仿射代数集之间的几何问题转化为其坐标环的代数问题。

在仿射概形态射的定义中，要求茎上的同态是局部同态这一条件是必要的：谱上的连续映射连同一组截面上的相容环同态（因而也诱导茎上的同态），并不足以迫使这些映射来自某个环同态。

\paragraph{例}
令 $R=\mathbb Z_{(2)}$ 且 $S=\mathbb Q$，如前面的一组例子所述。定义
\[
\varphi^*:\Spec\mathbb Q\longrightarrow\Spec\mathbb Z_{(2)}
\]
为 $\varphi^*((0))=(2)$（它是 Zariski 连续的）。定义
\[
\varphi^\#:\mathcal O(\Spec R)\longrightarrow\mathcal O(\Spec S)
\]
为包含映射 $\mathbb Z_{(2)}\to\mathbb Q$，并且对所有其他 $U\subseteq\Spec R$ 简单地定义 $\varphi^\#$ 为零映射。直接检查可知，这些同态与限制映射交换。然而，这一族映射并不来自一个环同态，因为对 $(0)\in\Spec S$ 以及 $\varphi^*((0))=(2)\in\Spec R$ 的茎而言，诱导同态
\[
\varphi^\#:\mathcal O_{\Spec R,(2)}\longrightarrow\mathcal O_{\Spec S,(0)}
\]
是包含映射 $\mathbb Z_{(2)}\to\mathbb Q$，它不是局部同态（$(0)$ 的逆像是 $(0)$，而不是极大理想 $2\mathbb Z_{(2)}$）。

定理 59 的证明说明，一个仿射概形态射 $(\varphi^*,\varphi^\#)$ 必然来自 $\varphi^\#$ 在整体截面上定义的环同态。在本例中，整体截面上的同态是 $R$ 到 $S$ 的包含映射。从 $R$ 到 $S$ 的包含映射定义了从 $\Spec S$ 到 $\Spec R$ 的映射，它把 $\Spec S$ 中的 $(0)$ 映到 $\Spec R$ 中的 $(0)$，而不是映到 $\Spec R$ 中的 $(2)$，所以这个映射不同于原来的映射 $\varphi^*$。

前一例说明，如果没有仿射概形态射定义中的第三个（局部同态）条件，定理 59 的逆命题就不成立。由此，定理 59 表明，观察仿射概形的适当背景是局部环空间的范畴。粗略地说，一个\emph{局部环空间}是一个拓扑空间 $X$，连同对 $X$ 的每个开子集 $U$ 给出的一族环 $\mathcal O(U)$（当 $U'\subseteq U$ 时有从 $\mathcal O(U)$ 到 $\mathcal O(U')$ 的相容同态，并且截面满足某些局部条件），使得对 $P\in U$，茎
\[
\mathcal O_P=\varinjlim_P\mathcal O(U)
\]
是局部环。这个范畴中的态射是拓扑空间之间的连续映射，连同相应 $\mathcal O(U)$ 之间的环同态，并满足与仿射概形态射定义中所施加条件完全相同的条件。

一个\emph{概形}是一个局部环空间，其中每个点都位于一个与仿射概形同构的邻域中（并且这些邻域之间还满足某些相容条件），它是现代代数几何中的基本研究对象。这里考虑的仿射概形是构造一般概形时“粘合”在一起的构件，正如通常的 Euclidean 空间是在分析中定义流形时被“粘合”在一起的构件一样。

\subsection*{练习}

所有环均假定为含幺交换环，所有环同态均假定保持恒等元。

\begin{enumerate}[label=\arabic*.]
	\item 若 $N$ 是 $R$ 的 nilradical，证明 $\Spec R$ 与 $\Spec R/N$ 同胚。[说明从 $R$ 到 $R/N$ 的自然同态诱导从 $\Spec R/N$ 到 $\Spec R$ 的 Zariski 连续同构。]
	
	\item 设 $I$ 是环 $R$ 中的理想。证明由典范投影同态 $R\to R/I$ 诱导的连续映射从 $\Spec R/I$ 到 $\Spec R$，把 $\Spec R/I$ 同胚地映到 $\Spec R$ 中的闭集 $Z(I)$ 上。
	
	\item 证明两个元素 $f,g\in R$ 在 $\Spec R$ 中所有元素 $P$ 处的值相同，当且仅当 $f-g$ 属于 $R$ 的 nilradical。特别地，证明仿射 $k$-代数中的元素由其取值唯一确定。
	
	\item 设 $k$ 是任意域，不必代数闭。证明 $k[x,y]$ 中的素理想（即 $\Spec k[x,y]$ 的元素）为
	\begin{enumerate}[label=(\roman*)]
		\item $(0)$；
		\item $(f)$，其中 $f$ 是 $k[x,y]$ 中的不可约多项式；
		\item $(p(x),g(x,y))$，其中 $p(x)$ 是 $k[x]$ 中的不可约多项式，而 $g(x,y)$ 是 $k[x,y]$ 中的不可约多项式，并且模 $p(x)$ 后仍不可约，即 $g(x,y)$ 在商环 $k[x,y]/(p(x))$ 中仍不可约。
	\end{enumerate}
	证明 $\operatorname{mSpec}k[x,y]$ 由 (iii) 中的素理想组成。[使用第 1 节练习 20。]
	
	\item 令 $\mathfrak m=(p(x),g(x,y))$ 是前一练习中的 $k[x,y]$ 的极大理想。说明 $K=k[x,y]/\mathfrak m$ 是 $k$ 的代数域扩张，从而 $k[x,y]$ 也可看成 $K[x,y]$ 的子环。若在典范同态 $k[x,y]\to k[x,y]/\mathfrak m$ 下，$x,y$ 分别映到 $\alpha,\beta\in K$，证明
	\[
	\mathfrak m=k[x,y]\cap (x-\alpha,y-\beta)
	\]
	其中 $(x-\alpha,y-\beta)\subseteq K[x,y]$。
	
	\item 描述 $\Spec\mathbb R[x]$ 与 $\Spec\mathbb C[x]$ 中的元素。描述 $\Spec\mathbb Z_{(2)}[x]$ 中的元素，其中
	\[
	\mathbb Z_{(2)}=\{a/b\in\mathbb Q\mid b\text{ 为奇数}\}
	\]
	是 $\mathbb Z$ 在素理想 $(2)$ 处的局部化。
	
	\item 令 $(f)=(x^5+x+1)$ 为 $\Spec\mathbb Z[x]$ 中的点，并把它看成如命题 55 后例 3 那样纤维化在 $\Spec\mathbb Z$ 上。说明在 $(2)$ 上方纤维中有两个闭点，在 $(5)$ 上方纤维中有三个闭点，在 $(19)$ 上方纤维中有四个闭点，在 $(211)$ 上方纤维中有五个闭点。
	
	\item 令 $(f)=(x^4+1)$ 为 $\Spec\mathbb Z[x]$ 中的点，并把它看成如命题 55 后例 3 那样纤维化在 $\Spec\mathbb Z$ 上。证明在 $(2)$ 上方纤维中有一个闭点；当 $p$ 为奇素数且 $p\equiv1\pmod 8$ 时，在 $p$ 上方纤维中有四个闭点；而对所有其他奇素数 $p$，在 $p$ 上方纤维中有两个闭点（参见第 14 章第 3 节推论 16）。
	
	\item 证明从 $\Spec\mathbb Z[x]$ 到 $\Spec\mathbb Z$ 的 Zariski 连续映射在 $(p)$ 上方的纤维中的元素，与
	\[
	\Spec(\mathbb Z[x]\otimes_{\mathbb Z}\mathbb F_p)
	\]
	中的元素同胚。
	
	\item 设 $X=\Spec R$，并设 $X_f$ 是对应于 $f\in R$ 的主开集。证明
	\[
	X_f\cap X_g=X_{fg}.
	\]
	证明 $X_f=X$ 当且仅当 $f$ 是 $R$ 中的单位，并且 $X_f=\varnothing$ 当且仅当 $f$ 是幂零元。
	
	\item 若 $X_f$ 与 $X_g$ 是 $X=\Spec R$ 中的主开集，证明开集 $X_f\cup X_g$ 是闭集 $Z(I)$ 的补集，其中 $I=(f,g)$ 是 $R$ 中由 $f$ 与 $g$ 生成的理想。
	
	\item 证明 $X=\Spec R$ 的 Zariski 开子集 $U$ 是拟紧的，当且仅当 $U$ 是有限多个主开子集的并。给出一个环 $R$、$\Spec R$ 的一个 Zariski 开子集 $U$，以及 $U$ 的一个不能缩减为有限子覆盖的 Zariski 开覆盖的例子。
	
	\item 设 $\varphi:R\to S$ 是环同态。证明在诱导映射 $\varphi^*$ 从 $Y=\Spec S$ 到 $X=\Spec R$ 下，$X$ 中主开集 $X_f$ 的全逆像是 $Y$ 中的主开集 $Y_{\varphi(f)}$。
	
	\item 设 $R=R_1\times R_2$ 是环 $R_1$ 与 $R_2$ 的直积。证明 $X=\Spec R$ 是两个开子空间 $X_1,X_2$ 的不交并（因此它们也都是闭的），其中 $X_1$ 同胚于 $\Spec R_1$，$X_2$ 同胚于 $\Spec R_2$。
	
	\item 证明 $X=\Spec R$ 不连通，当且仅当 $R$ 是两个非零环的直积，当且仅当 $R$ 含有一个满足 $e\ne0,1$ 的幂等元 $e$（参见前一练习）。
	
	\item 证明 $X=\Spec R$ 是不可约的（即任意两个非空开子集都有非平凡交），当且仅当任意两个非空主开集 $X_f$ 与 $X_g$ 都满足 $X_f\cap X_g\ne\varnothing$。推出 $X=\Spec R$ 不可约当且仅当 $R$ 的 nilradical 是素理想。[使用练习 10。]
	
	\item 令 $G=\langle a\rangle$ 是 $2$ 阶群，令
	\[
	R=\mathbb Z[G]=\{m+na\mid m,n\in\mathbb Z\}
	\]
	是相应的群环，并令 $X=\Spec R$。
	\begin{enumerate}[label=(\alph*)]
		\item 证明 $R$ 的 nilradical 是 $(0)$，但它不是素理想。证明
		\[
		X=X^+\cup X^- ,
		\]
		其中 $X^+=Z(1-a)$ 且 $X^-=Z(1+a)$。[使用 $(1+a)(1-a)=0$。]
		
		\item 证明同态 $\mathbb Z[G]\to\mathbb Z$，由把 $a$ 映到 $1$ 定义，诱导 $X^+$ 与 $\Spec\mathbb Z$ 的一个同胚；而由把 $a$ 映到 $-1$ 定义的同态诱导 $X^-$ 与 $\Spec\mathbb Z$ 的一个同胚。
		
		\item 证明 $X^+\cap X^-$ 由单个元素
		\[
		\mathfrak m=(1+a,1-a)=(2,1-a)
		\]
		组成，并且这是 $X$ 中的闭点。
		
		\item 说明 $(1-a)$ 与 $(1+a)$ 是 $X$ 中唯一的非闭点，其闭包分别为 $X^+$ 与 $X^-$。描述 $X$ 中的闭点 $\operatorname{mSpec}R$，并证明 $\Spec\mathbb Z[\langle a\rangle]$ 可以画为下图：
		
		\begin{center}
			\begin{tikzpicture}[scale=0.95, every node/.style={font=\small}]
				\draw[->] (8.0,1.45) -- (8.0,0.3) node[midway,right] {$\varphi^*$};
				\node at (8.9,1.65) {$\Spec\mathbb Z[\langle a\rangle]$};
				\node at (8.5,0) {$\Spec\mathbb Z$};
				\draw (0,1.8) -- (7.4,1.8);
				\draw (0,1.1) -- (7.4,1.1);
				\node[left] at (0,1.8) {$X^-$};
				\node[left] at (0,1.1) {$X^+$};
				\node at (1.0,1.8) {$\bullet\ (1+a)$};
				\node at (1.0,1.1) {$\bullet\ (1-a)$};
				\node at (3.5,1.8) {$\bullet\ (3,1-a)$};
				\node at (3.5,1.1) {$\bullet\ (3,1+a)$};
				\node at (5.4,1.8) {$\bullet\ (5,1-a)$};
				\node at (5.4,1.1) {$\bullet\ (5,1+a)$};
				\node at (2.2,1.45) {$\bullet$};
				\node[above right=-2pt] at (2.2,1.45) {$\mathfrak m$};
				\draw (0,0) -- (7.4,0);
				\foreach \x/\lab in {1/{(0)},2.2/{(2)},3.5/{(3)},5.4/{(5)}}
				{\node at (\x,0) {$\bullet$}; \node[below=3pt] at (\x,0) {\lab};}
				\draw[dashed] (1.0,1.0)--(1.0,0.15);
				\draw[dashed] (2.2,1.3)--(2.2,0.15);
				\draw[dashed] (3.5,1.0)--(3.5,0.15);
				\draw[dashed] (5.4,1.0)--(5.4,0.15);
			\end{tikzpicture}
		\end{center}
	\end{enumerate}
	
	\item 设 $\mathcal O$ 是 $X=\Spec R$ 上的结构层，设 $U$ 是 $X$ 中开集，并设 $s,t\in\mathcal O(U)$。若在 $X_{f_1}$ 上有 $s=a/f_1^n$，在 $X_{f_2}$ 上有 $t=b/f_2^m$，说明在 $X_{f_1f_2}$ 上有
	\[
	st=\frac{ab}{(f_1f_2)^{n+m}}\,f_1^m f_2^n
	\]
	按通常写法即
	\[
	st=\frac{abf_1^m f_2^n}{(f_1f_2)^{n+m}},
	\qquad
	s+t=\frac{af_2^{m+n}+bf_1^{m+n}}{(f_1f_2)^{n+m}}.
	\]
	推出 $\mathcal O(U)$ 是含单位元的交换环。
	
	\item 设 $\mathcal O$ 是 $X=\Spec R$ 上的结构层，设 $V\subseteq U$ 是 $X$ 中开集，并设 $s\in\mathcal O(U)$。假设 $P\in V$，且在 $X_f\subseteq U$ 上有 $s=a/f^n$。
	\begin{enumerate}[label=(\alph*)]
		\item 说明存在包含 $P$ 的主开集 $X_{f'}\subseteq V\cap X_f$。
		\item 说明对某个 $b\in R$，有 $(f')^m=bf$。
		\item 说明在 $X_{f'}$ 上有
		\[
		s=\frac{ab^n}{(f')^{mn}},
		\]
		并推出把 $s$ 限制到 $V$ 给出从 $\mathcal O(U)$ 到 $\mathcal O(V)$ 的良定义环同态。
	\end{enumerate}
	
	\item 设 $\varphi:R\to S$ 是环同态，设 $X=\Spec R$，$Y=\Spec S$，并设 $V\subseteq U$ 是 $X$ 的 Zariski 开子集。令
	\[
	V'=(\varphi^*)^{-1}(V),
	\qquad
	U'=(\varphi^*)^{-1}(U),
	\]
	它们是关于由 $\varphi$ 诱导的连续映射 $\varphi^*:Y\to X$ 的 $Y$ 中对应 Zariski 开子集。证明截面上的诱导映射
	\[
	\varphi^\#:\mathcal O_X(U)\longrightarrow\mathcal O_Y(U')
	\]
	是环同态。证明 $V'\subseteq U'$，并且 $\varphi^\#$ 与限制映射相容，即图
	\[
	\begin{array}{ccc}
		\mathcal O_X(U)&\xrightarrow{\varphi^\#}&\mathcal O_Y(U')\\
		\downarrow&&\downarrow\\
		\mathcal O_X(V)&\xrightarrow{\varphi^\#}&\mathcal O_Y(V')
	\end{array}
	\]
	交换，其中竖直映射为限制同态。
	
	\item 设 $D$ 是 $R$ 的乘法闭子集。说明局部化同态 $R\to D^{-1}R$ 诱导一个从 $\Spec(D^{-1}R)$ 到所有满足 $P\cap D=\varnothing$ 的 $R$ 中素理想 $P$ 所组成集合的同胚。
	
	\item 说明 $\Spec k[x,y]/(xy)$ 是连通的，但它是两个真闭子集的并，且每个都同胚于 $\Spec k[x]$；因此它不是不可约的（参见练习 16）。
	
	\item 对下列每个环 $R$，列出 $\Spec R$ 的元素、$\Spec R$ 中的开集 $U$、对每个开集 $U$ 结构层在 $U$ 上的截面 $\mathcal O(U)$，以及每个点 $P\in\Spec R$ 处的茎 $\mathcal O_P$：
	\[
	\text{(a) }\mathbb Z/4\mathbb Z,
	\quad
	\text{(b) }\mathbb Z/6\mathbb Z,
	\quad
	\text{(c) }\mathbb Z/2\mathbb Z\times\mathbb Z/3\mathbb Z,
	\quad
	\text{(d) }\mathbb Z/2\mathbb Z\times\mathbb Z/2\mathbb Z\times\mathbb Z/2\mathbb Z.
	\]
	
	\item \begin{enumerate}[label=(\alph*)]
		\item 若 $R$ 的每个理想都是主理想，说明 $\Spec R$ 中每个开集都是主开集。
		\item 说明若
		\[
		R=\mathbb Z[x]/(4,x^2),
		\]
		则 $R$ 含有非主理想，但 $\Spec R$ 中每个开集仍是主开集。
	\end{enumerate}
	
	\item \begin{enumerate}[label=(\alph*)]
		\item 若 $M$ 是 $R$-模，证明 $\operatorname{Supp}(M)$ 是 $\Spec R$ 的 Zariski 闭子集。[使用第 4 节练习 33。]
		\item 若 $M$ 是有限生成 $R$-模，证明
		\[
		\operatorname{Supp}(M)=Z(\operatorname{Ann}(M))\subseteq\Spec R.
		\]
		[使用第 4 节练习 34。]
	\end{enumerate}
	
	\item 设 $M$ 是 Noether 环 $R$ 上的有限生成模。
	\begin{enumerate}[label=(\alph*)]
		\item 证明包含 $\operatorname{Ann}(M)$ 的极小素理想只有有限多个，设为 $P_1,\ldots,P_n$。[使用推论 22。]
		\item 证明 $\{P_1,\ldots,P_n\}$ 也是 $\operatorname{Ass}_R(M)$ 中极小素理想的集合，并且 $\operatorname{Supp}(M)$ 是 $\Spec R$ 中 Zariski 闭集
		\[
		Z(P_1),\ldots,Z(P_n)
		\]
		的并。[使用前一练习与第 4 节练习 40。]
	\end{enumerate}
	
	前一练习给出了 Noether 环 $R$ 上有限生成模 $M$ 的几何图像：在 $\Spec R$ 中每个点 $P$ 上方是局部化 $M_P$（$P$ 上方的茎）。这个茎恰好在 Zariski 闭子集 $Z(P_1),\ldots,Z(P_n)$ 上方的点处非零，其中 $P_i$ 是 $\operatorname{Ass}_R(M)$ 中的极小素理想。这些思想导向了 $\Spec R$ 上与 $M$ 相关的（凝聚）模层的概念（其图像类似于命题 58 后结构层的图像），这是现代代数几何中的一个有力工具。
	
	\item 令 $R=k[x,y]$，并令 $M$ 为 $R$ 中的理想 $(x,y)$。证明
	\[
	\operatorname{Supp}(M)=\Spec R
	\qquad\text{且}\qquad
	\operatorname{Ass}_R(M)=\{0\}.
	\]
	
	下面两个练习说明，在 Noether 环 $R$ 中，对理想 $I$ 而言，primary 分解意义下的相关素理想正是 $\operatorname{Ass}_R(R/I)$ 意义下的相关素理想。
	
	\item 本练习证明 Noether 环 $R$ 中理想 $Q$ 是 $P$-primary 的，当且仅当
	\[
	\operatorname{Ass}_R(R/Q)=\{P\}.
	\]
	\begin{enumerate}[label=(\alph*)]
		\item 设 $Q$ 是 $P$-primary 理想，并令 $M$ 为 $R$-模 $R/Q$。若 $0\ne m\in M$，说明
		\[
		Q\subseteq\operatorname{Ann}(m)\subseteq P,
		\qquad
		\operatorname{rad}\operatorname{Ann}(m)=P.
		\]
		推出若 $\operatorname{Ann}(m)$ 是素理想，则它等于 $P$，从而
		\[
		\operatorname{Ass}_R(R/Q)=\{P\}.
		\]
		[使用第 1 节练习 33。]
		
		\item 对 $R$ 的任意理想 $Q$，设 $0\ne M\subseteq R/Q$。证明 $\operatorname{Ann}(M)$ 的根是 $\operatorname{Supp}(M)$ 中素理想的交。[使用命题 12 与练习 25。]
		
		\item 对 (b) 中的 $M$，证明 $\operatorname{Ann}M$ 的根也等于 $\operatorname{Ass}_R(M)$ 中素理想的交。[使用练习 26(b)。]
		
		\item 若 $Q$ 是 $R$ 的理想且 $\operatorname{Ass}_R(R/Q)=\{P\}$，证明 $\operatorname{rad}Q=P$。[使用 $Q=\operatorname{Ann}(R/Q)$ 以及 (c)。]
		
		\item 若 $Q$ 是 $R$ 的理想且 $\operatorname{Ass}_R(R/Q)=\{P\}$，证明 $Q$ 是 $P$-primary。[若 $ab\in Q$ 且 $a\notin Q$，考虑
		\[
		0\ne M=(Ra+Q)/Q\subseteq R/Q,
		\]
		并说明 $b$ 属于 $\operatorname{Ann}M\subseteq\operatorname{rad}\operatorname{Ann}(M)$。使用第 1 节练习 33--34 说明 $\operatorname{Ass}_R(M)=\{P\}$，再用 (c) 说明 $\operatorname{rad}\operatorname{Ann}(M)=P$，最后推出 $b\in P$。]
	\end{enumerate}
	
	\item 设
	\[
	I=Q_1\cap\cdots\cap Q_n
	\]
	是 Noether 环 $R$ 中理想 $I$ 的一个极小 primary 分解，并设
	\[
	P_i=\operatorname{rad}Q_i,
	\qquad i=1,\ldots,n.
	\]
	本练习证明
	\[
	\operatorname{Ass}_R(R/I)=\{P_1,\ldots,P_n\}.
	\]
	\begin{enumerate}[label=(\alph*)]
		\item 证明自然投影同态诱导一个从 $R/I$ 到
		\[
		R/Q_1\oplus\cdots\oplus R/Q_n
		\]
		的单射，并推出
		\[
		\operatorname{Ass}_R(R/I)\subseteq\{P_1,\ldots,P_n\}.
		\]
		[使用第 1 节练习 34 以及前一练习。]
		
		\item 令
		\[
		Q_i'=\bigcap_{j\ne i}Q_j.
		\]
		说明分解的极小性蕴含
		\[
		0\ne Q_i'/I=(Q_i'+Q_i)/Q_i\subseteq R/Q_i.
		\]
		推出
		\[
		\operatorname{Ass}_R(Q_i'/I)=\{P_i\}.
		\]
		[使用第 1 节练习 33--34 以及前一练习。] 推出 $\{P_i\}\in\operatorname{Ass}_R(R/I)$，从而
		\[
		\operatorname{Ass}_R(R/I)=\{P_1,\ldots,P_n\}.
		\]
		[使用 $Q_i'/I\subseteq R/I$ 以及第 1 节练习 34。]
	\end{enumerate}
	
	\item 令 $I$ 是 $R=k[x,y,z]$ 中的理想
	\[
	I=(x^2,xy,xz,yz).
	\]
	证明 $\operatorname{Ass}_R(R/I)$ 由素理想
	\[
	\{(x,y),(x,z),(x,y,z)\}
	\]
	组成。
	
	\item （二次整数环的素谱）令 $R$ 为二次域 $K=\mathbb Q(\sqrt D)$ 的整数环，其中 $D$ 是无平方因子的整数；令 $P$ 是 $R$ 的非零素理想。本练习说明如何由 $\mathbb Z$ 中的素理想 $(p)$ 明确确定 $R$ 中的素理想，从而特别给出 $\Spec R$ 在 $\Spec\mathbb Z$ 上的纤维描述。
	
	如定理 29 后的讨论与例子所述，有 $R=\mathbb Z[\omega]$，其中若 $D\equiv2,3\pmod 4$，则 $\omega=\sqrt D$；而若 $D\equiv1\pmod 4$，则 $\omega=(1+\sqrt D)/2$。其极小多项式分别为
	\[
	m_\omega(x)=x^2-D,
	\]
	或
	\[
	m_\omega(x)=x^2-x+\frac{1-D}{4},
	\]
	并且 $P\cap\mathbb Z=p\mathbb Z$ 是 $\mathbb Z$ 的非零素理想。
	\begin{enumerate}[label=(\alph*)]
		\item 对 $\mathbb Z$ 中任意素数 $p$，证明环同构
		\[
		R/pR\simeq \mathbb Z[x]/(p,m_\omega(x))
		\simeq\mathbb F_p[x]/(\overline{m_\omega}(x)),
		\]
		其中 $\overline{m_\omega}(x)$ 是 $m_\omega(x)$ 模 $p$ 的化简。推出存在 $R$ 中素理想 $P$，满足 $P\cap\mathbb Z=(p)$（这在本情形下给出定理 26(2) 的另一个证明）。
		
		\item 使用 (a) 中的同构，证明 $P$ 可由 $m_\omega(x)$ 模 $p$ 的分解显式确定：
		\begin{enumerate}[label=(\roman*)]
			\item 若
			\[
			m_\omega(x)=(x-a)^2\pmod p,
			\]
			其中 $a\in\mathbb Z$，则
			\[
			P=(p,\omega-a),
			\qquad
			pR=P^2.
			\]
			说明这种情形只会发生在有限多个素数 $p$ 上，即那些整除 $m_\omega(x)$ 判别式的素数。
			
			\item 若
			\[
			m_\omega(x)=(x-a)(x-b)\pmod p,
			\]
			其中整数 $a,b\in\mathbb Z$ 模 $p$ 不同，则 $P$ 或者是
			\[
			P_1=(p,\omega-a),
			\]
			或者是
			\[
			P_2=(p,\omega-b),
			\]
			且 $P_1,P_2$ 是 $R$ 中不同的素理想，并满足
			\[
			pR=P_1P_2.
			\]
			
			\item 若 $m_\omega(x)$ 模 $p$ 不可约，则
			\[
			P=pR.
			\]
		\end{enumerate}
		
		\item 说明对任意 $D$，$\Spec R$ 在 $\Spec\mathbb Z$ 上方的图像都类似于 $D=-1$ 时 $R=\mathbb Z[i]$ 的情形：在 $\Spec\mathbb Z$ 中 $(0)$ 上方，$\Spec R$ 中恰有一个非闭点 $(0)$；对 (i) 中的素数 $(p)$（称为 ramified primes）以及 (iii) 中的素数（称为 inert primes），其上方在 $\Spec R$ 中恰有一个闭点 $P$；对 (ii) 中的素数（称为 split primes），其上方恰有两个闭点。
	\end{enumerate}
\end{enumerate}


% % ============================================
%  chapter16.tex — 第16章：Artin 环、离散赋值环与 Dedekind 整环
%  根据 764-789.pdf 翻译整理，保持原文结构、证明与练习
% ============================================

\chapter{Artin 环、离散赋值环与 Dedekind 整环}

在本章中，$R$ 始终表示一个带有 $1\neq0$ 的交换环。

\section{Artin 环}

在本节中，我们将研究满足理想降链条件（D.C.C.）的交换环的基本理论，即 \emph{Artin 环}（以 E. Artin 命名）。人们一开始也许会期望这些环具有类似于满足升链条件的交换环（即 Noether 环）的性质，但实际上并非如此。Artin 环的结构受到很强的限制；例如，Artin 环必定也是 Noether 环（定理 3）。非交换 Artin 环在表示论中起着中心作用（参见第 18 章和第 19 章）。

\noindent\textbf{定义.} 对任意交换环 $R$，$R$ 的 \emph{Krull 维数}（或简称\emph{维数}）是不同素理想链
\[
  P_0\subset P_1\subset P_2\subset\cdots\subset P_n
\]
的最大可能长度。若 $R$ 有任意长的不同素理想链，则称 $R$ 的维数是无限的。

有限维的环必须同时满足素理想的升链条件和降链条件（尽管不一定满足所有理想上的相应条件）。域的维数为 $0$，而不是域的主理想整环的维数为 $1$。

我们很快会看到，满足理想降链条件的环总是维数为 $0$（即素理想都是极大理想）。若 $R$ 是整环，并且也是域 $k$ 上的有限生成 $k$-代数，则 $R$ 的维数等于 $R$ 的分式域在 $k$ 上的超越次数（参见练习 11）。特别地，Krull 维数与前面为仿射簇的维数所引入的定义一致。上述定义的优点在于它不涉及任何 $k$-代数结构，并适用于任意交换环 $R$。

\noindent\textbf{定义.} $R$ 的 \emph{Jacobson 根}是 $R$ 的所有极大理想的交，记为 $\Jac R$。

Jacobson 根类似于群的 Frattini 子群，并具有若干相应的性质（参见第 6.1 节练习 24）：

\medskip
\noindent\textbf{命题 1.} 设 $\mathcal J$ 为交换环 $R$ 的 Jacobson 根。
\begin{enumerate}[label=(\arabic*)]
  \item 若 $I$ 是 $R$ 的真理想，则 $(I,\mathcal J)$，即由 $I$ 与 $\mathcal J$ 生成的理想，也是真理想。
  \item Jacobson 根包含 $R$ 的 nilradical：$\rad 0\subseteq\Jac R$。
  \item 元素 $x$ 属于 $\mathcal J$，当且仅当对所有 $r\in R$，$1-rx$ 都是单位。
  \item （Nakayama 引理）若 $M$ 是任意有限生成 $R$-模，且 $\mathcal JM=M$，则 $M=0$。
\end{enumerate}

\begin{proof}
若 $I$ 是 $R$ 中的真理想，则 $I\subseteq M$ 对某个极大理想 $M$ 成立。由于 $\mathcal J\subseteq M$，也有 $(I,\mathcal J)\subseteq M$，这证明了 (1)。

(2) 由两个根的定义以及第 15.2 节命题 12 推出，因为极大理想都是素理想。

假设 $1-rx$ 不是单位，并令 $M$ 为包含 $1-rx$ 的极大理想。由于 $1\notin M$，有 $rx\notin M$，所以 $x$ 不能属于 $\mathcal J$，因为 $\mathcal J\subseteq M$。反过来，假设 $x\notin\mathcal J$，即存在极大理想 $M$ 使 $x\notin M$。于是 $R=(x,M)$，故对某个 $y\in M$ 有 $1=rx+y$。因此 $1-rx=y\in M$，所以 $1-rx$ 不是单位，这证明了 (3)。

为证明 (4)，假设 $M\neq0$，并令 $n$ 为使 $M$ 可由 $n$ 个元素生成的最小整数，设这些生成元为 $m_1,\ldots,m_n$。由于 $M=\mathcal JM$，我们有
\[
  m_n=r_1m_1+r_2m_2+\cdots+r_nm_n
\]
其中 $r_1,r_2,\ldots,r_n\in\mathcal J$。于是
\[
  (1-r_n)m_n=r_1m_1+\cdots+r_{n-1}m_{n-1}.
\]
由 (3)，$1-r_n$ 是单位，所以 $m_n$ 属于由 $m_1,\ldots,m_{n-1}$ 生成的模，这与 $n$ 的最小性矛盾。因此 $M=0$，证明完成。
\end{proof}

\noindent\textbf{定义.} 交换环 $R$ 称为 \emph{Artin 环}，或称为满足\emph{理想降链条件}（或理想上的 D.C.C.），如果 $R$ 中不存在无限下降的理想链；也就是说，只要
\[
  I_1\supseteq I_2\supseteq I_3\supseteq\cdots
\]
是 $R$ 的理想下降链，则存在正整数 $m$，使得对所有 $k\geq m$ 都有 $I_k=I_m$。类似地，$R$-模 $M$ 称为 \emph{Artin 模}，若它满足子模上的 D.C.C.。

由格同构定理立刻可知，Artin 环 $R$ 对任意理想 $I$ 的商环 $R/I$ 仍是 Artin 环。

下面关于 Artin 环的结果与定理 15.2 中的结果平行。证明完全类似，因此留作练习。

\medskip
\noindent\textbf{命题 2.} 下列条件等价：
\begin{enumerate}[label=(\arabic*)]
  \item $R$ 是 Artin 环。
  \item $R$ 的每个非空理想集合在包含关系下含有极小元。
\end{enumerate}

下面的结果给出 Artin 环的主要结构定理。

\medskip
\noindent\textbf{定理 3.} 设 $R$ 是 Artin 环。
\begin{enumerate}[label=(\arabic*)]
  \item $R$ 中只有有限多个极大理想。
  \item 商环 $R/(\Jac R)$ 是有限多个域的直积。更准确地说，若 $M_1,\ldots,M_n$ 是 $R$ 中有限多个极大理想，则
  \[
    R/(\Jac R)\cong k_1\times\cdots\times k_n,
  \]
  其中 $k_i$ 是域 $R/M_i$，$1\leq i\leq n$。
  \item $R$ 的每个素理想都是极大理想，即 $R$ 的 Krull 维数为 $0$。$R$ 的 Jacobson 根等于 $R$ 的 nilradical，并且是一个幂零理想：对某个 $m\geq1$，$(\Jac R)^m=0$。
  \item 环 $R$ 同构于有限多个 Artin 局部环的直积。
  \item 每个 Artin 环都是 Noether 环。
\end{enumerate}

\begin{proof}
为证明 (1)，令 $S$ 为 $R$ 的所有可表示为有限多个极大理想之交的理想组成的集合。由命题 2，$S$ 有一个极小元，设为 $M_1\cap M_2\cap\cdots\cap M_n$。于是对任意极大理想 $M$，我们有
\[
  M\cap M_1\cap M_2\cap\cdots\cap M_n=M_1\cap M_2\cap\cdots\cap M_n,
\]
所以 $M\supseteq M_1\cap M_2\cap\cdots\cap M_n$。由第 7.4 节练习 11，$M\supseteq M_i$ 对某个 $i$ 成立。因此 $M=M_i$，所以 $M_1,\ldots,M_n$ 就是 $R$ 的全部极大理想。

(2) 的证明由中国剩余定理（第 7.6 节）立刻得到，该定理应用于 $M_1,\ldots,M_n$，因为这些极大理想显然两两互素，并且它们的交为 $\Jac R$。

对 (3)，我们先证明 $\mathcal J=\Jac R$ 是幂零的。由 D.C.C.，存在某个 $m>0$，使得对所有正整数 $i$ 都有 $\mathcal J^m=\mathcal J^{m+i}$。反证，假设 $\mathcal J^m\neq0$。令 $S$ 为所有满足 $I\mathcal J^m\neq0$ 的真理想 $I$ 组成的集合，所以 $\mathcal J\in S$。令 $I_0$ 为 $S$ 的一个极小元。存在某个 $x\in I_0$ 使得 $x\mathcal J^m\neq0$，故由极小性必有 $I_0=(x)$。但是现在
\[
  ((x)\mathcal J)\mathcal J^m=x\mathcal J^{m+1}=x\mathcal J^m,
\]
所以由 $(x)$ 的极小性得到 $(x)=(x)\mathcal J$。由上面的 Nakayama 引理，$(x)=0$，矛盾。这证明了 $\Jac R$ 是幂零的。

由于 $\Jac R$ 是幂零的，特别地 $\Jac R\subseteq\rad0$，故由命题 1 的第二个断言，这两个理想相等。

$R$ 中每个素理想 $P$ 都包含 $R$ 的 nilradical，因此由已证结果也包含 $\Jac R$。$P$ 的像是商环
\[
  R/(\Jac R)=k_1\times\cdots\times k_n
\]
中的一个素理想。但是在环的直积 $R_1\times R_2$ 中（其中每个 $R_i$ 都有 $1$），每个理想都形如 $I_1\times I_2$，其中 $I_j$ 是 $R_j$ 的理想，$j=1,2$（参见第 7.6 节练习 3）。于是 $k_1\times\cdots\times k_n$ 中的素理想由某一个分量为 $0$ 的元素组成。特别地，这样的素理想也是 $k_1\times\cdots\times k_n$ 中的极大理想，从而 $P$ 原本就是 $R$ 中的极大理想。这完成了 (3) 的证明。

令 $M_1,\ldots,M_n$ 为 $R$ 的全部不同极大理想，并令 $(\Jac R)^m=0$ 如 (3) 中所述。则
\[
  \prod_{i=1}^n M_i^m\subseteq\left(\prod_{i=1}^n M_i\right)^m\subseteq (\Jac R)^m=0.
\]
由中国剩余定理得到
\[
  R\cong (R/M_1^m)\times (R/M_2^m)\times\cdots\times (R/M_n^m),
\]
且每个 $R/M_i^m$ 是 Artin 环，其唯一极大理想为 $M_i/M_i^m$，这证明了 (4)。

为证明 (5)，由 (4) 只需证明 Artin 局部环是 Noether 环，因此假设 $R$ 是唯一极大理想为 $M$ 的 Artin 环。在这种情形下 $M=\Jac R$，所以对某个正整数 $m$ 有 $M^m=(\Jac R)^m=0$。于是 $R\cong R/M^m$，而在这种情形下容易验证 $R/M^m$ 是 Noether 环当且仅当它是 Artin 环（参见练习 8）。
\end{proof}

\medskip
\noindent\textbf{推论 4.} 环 $R$ 是 Artin 环，当且仅当 $R$ 是 Noether 环且 Krull 维数为 $0$。

\begin{proof}
正向蕴含已在定理 3 中证明。现在假设 $R$ 是 Noether 环，并且 $R$ 的 Krull 维数为 $0$，即 $R$ 的素理想都是极大理想。由于 $R$ 是 Noether 环，由第 15.2 节推论 22(3)，理想 $(0)$ 是有限多个（不一定互异的）素理想的乘积，$(0)=P_1\cdots P_n$，而这些素理想由于 $R$ 维数为 $0$ 都是极大理想。由中国剩余定理，$R$ 同构于有限多个形如 $R/M^m$ 的 Noether 环的直积，其中 $M$ 是 $R$ 的极大理想。如定理 (5) 的证明中一样，$R/M^m$ 是 Artin 环，因此 $R$ 是 Artin 环。
\end{proof}

\paragraph{例}
\begin{enumerate}[label=(\arabic*)]
  \item 设 $n>1$ 为整数。由于环 $R=\ZZ/n\ZZ$ 是有限的，它是 Artin 环。若
  \[
    n=p_1^{a_1}p_2^{a_2}\cdots p_s^{a_s}
  \]
  是 $n$ 分解为不同素数幂的唯一分解，则
  \[
    \ZZ/n\ZZ\cong (\ZZ/p_1^{a_1}\ZZ)\times(\ZZ/p_2^{a_2}\ZZ)\times\cdots\times(\ZZ/p_s^{a_s}\ZZ).
  \]
  每个 $\ZZ/p_i^{a_i}\ZZ$ 是 Artin 局部环，其唯一极大理想为 $(p_i)/(p_i^{a_i})$，所以这就是定理 3(4) 给出的 $\ZZ/n\ZZ$ 的分解。$R$ 的 Jacobson 根是由 $p_1p_2\cdots p_s$（即 $n$ 的无平方因子部分）生成的理想，并且
  \[
    R/(\Jac R)\cong (\ZZ/p_1\ZZ)\times\cdots\times(\ZZ/p_s\ZZ)
  \]
  是域的直积。对 $i=1,\ldots,s$，由 $p_i$ 生成的理想就是 $R$ 的极大理想。

  \item 对任意域 $k$，若 $k$-代数 $R$ 作为 $k$ 上的向量空间是有限维的，则 $R$ 是 Artin 环，因为 $R$ 中的理想特别地都是 $R$ 的 $k$-子空间，故任意理想链的长度都被 $\dim_k R$ 所限制。

  \item 设 $f$ 是域 $k$ 上的多项式环 $k[x]$ 中的非零多项式。则商环 $R=k[x]/(f(x))$ 由前一个例子是 Artin 环。$R$ 作为 Artin 局部环直积的分解为
  \[
    k[x]/(f(x))\cong k[x]/(f_1(x)^{a_1})\times\cdots\times k[x]/(f_s(x)^{a_s}),
  \]
  其中
  \[
    f(x)=f_1(x)^{a_1}\cdots f_s(x)^{a_s}
  \]
  是 $f(x)$ 分解为 $k[x]$ 中不同不可约多项式幂的分解（参见第 9.5 节命题 16）。$R$ 的 Jacobson 根是由 $f(x)$ 的无平方因子部分生成的理想，而 $R$ 的极大理想是由不可约因子 $f_i(x)$ 生成的理想，$i=1,\ldots,s$，这与例 1 类似。
\end{enumerate}

\subsection*{练习}

设 $R$ 是带有 $1$ 的交换环，并设 $\mathcal J$ 为其 Jacobson 根。
\begin{enumerate}[label=\arabic*.]
  \item 假设 $R$ 是 Artin 环且 $I$ 是 $R$ 的理想。证明 $R/I$ 也是 Artin 环。
  \item 证明每个带有 $1$ 的有限交换环都是 Artin 环。
  \item 证明 Krull 维数为 $0$ 的整环是域。
  \item 证明 Artin 整环是域。
  \item 假设 $I$ 是 $R$ 中的幂零理想，并且对某个 $R$-模 $M$ 有 $M=IM$。证明 $M=0$。
  \item 假设 $0\to M'\to M\to M''\to0$ 是 $R$-模的正合列。证明 $M$ 是 Artin $R$-模，当且仅当 $M'$ 与 $M''$ 都是 Artin $R$-模。
  \item 假设 $R=F$ 是域。证明 $R$-模 $M$ 是 Artin 的，当且仅当它是 Noether 的，当且仅当 $M$ 是 $F$ 上的有限维向量空间。
  \item 设 $M$ 是环 $R$ 的极大理想，并假设对某个 $n\geq1$ 有 $M^n=0$。证明 $R$ 是 Noether 环当且仅当 $R$ 是 Artin 环。［观察滤过
  \[
    R\supseteq M\supseteq\cdots\supseteq M^{n-1}\supseteq M^n=0
  \]
  中的每个相继商 $M^i/M^{i+1}$，$i=0,\ldots,n-1$，都是域 $F=R/M$ 上的模。然后使用前两个练习以及第 15.1 节练习 6。］
  \item 设 $M$ 是有限生成 $R$-模。证明若 $x_1,\ldots,x_n$ 是 $M$ 中的元素，且它们在 $M/\mathcal JM$ 中的像生成 $M/\mathcal JM$，则它们生成 $M$。推出若 $R$ 是 Noether 环，且 $a_1,\ldots,a_n$ 在 $\mathcal J/\mathcal J^2$ 中的像生成 $\mathcal J/\mathcal J^2$，则 $\mathcal J=(a_1,
  \ldots,a_n)$。［令 $N$ 为由 $x_1,\ldots,x_n$ 生成的子模，并对模 $A=M/N$ 应用 Nakayama 引理。］
  \item 设 $R=\ZZ_{(2)}$ 是 $\ZZ$ 在素理想 $(2)$ 处的局部化。证明 $\Jac R=(2)$，即由 $2$ 生成的理想。若 $M=\QQ$，证明 $M/2M$ 是有限生成 $R$-模，但 $M$ 在 $R$ 上不是有限生成的。为什么这不与前一练习矛盾？［注意 Nakayama 引理中的假设。］
  \item 设 $V$ 是域 $k$ 上的仿射簇，且 $R=k[V]$ 为其坐标环。令 $d_t(R)$ 表示分式域 $k(V)$ 在 $k$ 上的超越次数，令 $d_p(R)$ 表示以素理想链定义的 $R$ 的 Krull 维数。本练习说明 $d_t(R)=d_p(R)$。由 Noether 正规化引理，$R$ 中存在多项式子环 $R_1=k[y_1,
  \ldots,y_m]$，使得 $R$ 在 $R_1$ 上整。
  \begin{enumerate}[label=(\alph*)]
    \item 证明 $d_t(R_1)=d_t(R)=m$ 且 $d_p(R_1)=d_p(R)$。推出我们可以假设 $R=R_1$。［用上升定理与下降定理（参见第 15.3 节定理 26）证明第二个等式。］
    \item 当 $R=R_1$ 时，通过给出一条长度为 $m$ 的显式素理想链，证明 $d_p(R)\geq d_t(R)$。
    \item 当 $R=R_1$ 时，证明 $R$ 的任意非零素理想都包含某个元素 $f$，使得 $R/(f)$ 在 $k$ 上的超越次数为 $d_t(R)-1$。用归纳法证明 $d_p(R)\leq d_t(R)$，并推出 $d_p(R)=d_t(R)$。
  \end{enumerate}
  \item 设 $R$ 是极大理想为 $M$ 的 Noether 局部环。
  \begin{enumerate}[label=(\alph*)]
    \item 商 $M/M^2$ 是域 $R/M$ 上的模（即向量空间）。证明 $d=\dim_{R/M}(M/M^2)$ 有限。
    \item 证明 $M$ 作为 $R$ 中理想可由 $d$ 个元素生成，且不能由更少的元素生成。［用练习 9。］
    \item 令
    \[
      R=k[x_1,\ldots,x_n]_{(x_1,\ldots,x_n)}
    \]
    为域 $k$ 上多项式环 $k[x_1,
    \ldots,x_n]$ 在极大理想 $(x_1,
    \ldots,x_n)$ 处的局部化，并令 $M$ 为 $R$ 中的极大理想。证明
    \[
      \dim_{R/M}(M/M^2)=n=\dim R.
    \]
    ［参见前一练习。］
  \end{enumerate}
\end{enumerate}

可以证明，对任意极大理想为 $M$ 的 Noether 局部环 $R$，都有 $\dim_{R/M}(M/M^2)\geq\dim R$。若 Noether 局部环 $R$ 满足 $\dim_{R/M}(M/M^2)=\dim R$，则称 $R$ 为\emph{正则局部环}。事实上，正则局部环必然是整环，并且还是整闭的。

\begin{enumerate}[label=\arabic*.,resume]
  \item 若 $R$ 是 Noether 环，证明 $\Spec R$ 上的 Zariski 拓扑是离散的（即每个子集都是 Zariski 开的并且也是 Zariski 闭的），当且仅当 $R$ 是 Artin 环。
  \item 设 $I$ 是多项式环 $k[x_1,x_2,x_3,
  \ldots]$ 中的理想 $(x_1,x_2^2,x_3^3,
  \ldots)$，其中 $k$ 是域，并令
  \[
    R=k[x_1,x_2,x_3,
    \ldots]/I.
  \]
  证明理想 $(x_1,x_2,x_3,
  \ldots)$ 在 $R$ 中的像是 $R$ 中唯一的素理想，但它不是有限生成的。推出 $R$ 是 Krull 维数为 $0$ 的局部环，但不是 Artin 环。
\end{enumerate}

\section{离散赋值环}

在上一节中，我们证明了 Artin 环就是 Krull 维数为 $0$ 的 Noether 环。现在我们考虑最简单的维数为 $1$ 的 Noether 环，即第 8.1 节首次引入的\emph{离散赋值环}：

\noindent\textbf{定义.}
\begin{enumerate}[label=(\arabic*)]
  \item 域 $K$ 上的\emph{离散赋值}是一个函数 $v:K^\times\to\ZZ$，满足
  \begin{enumerate}[label=(\roman*)]
    \item $v$ 是满射；
    \item 对所有 $x,y\in K^\times$，$v(xy)=v(x)+v(y)$；
    \item 对所有满足 $x+y\neq0$ 的 $x,y\in K^\times$，有
    \[
      v(x+y)\geq\min\{v(x),v(y)\}.
    \]
  \end{enumerate}
  子环 $\{x\in K\mid v(x)\geq0\}\cup\{0\}$ 称为 $v$ 的\emph{赋值环}。
  \item 若整环 $R$ 是其分式域上某个离散赋值 $v$ 的赋值环，则称 $R$ 为\emph{离散赋值环}（D.V.R.）。
\end{enumerate}

赋值 $v$ 常被扩张到整个 $K$ 上，定义 $v(0)=+\infty$；在这种情形下，(ii) 和 (iii) 对所有 $a,b\in K$ 成立。

\paragraph{例}
\begin{enumerate}[label=(\arabic*)]
  \item $\ZZ$ 在任意非零素理想 $(p)$ 处的局部化 $\ZZ_{(p)}$ 是一个 D.V.R.，其对应于 $\QQ$ 上如下定义的离散赋值 $v_p$（参见第 7.1 节练习 27）。每个元素 $a/b\in\QQ^\times$ 可唯一写成
  \[
    p^n(a_1/b_1)
  \]
  的形式，其中 $n\in\ZZ$，$a_1/b_1\in\QQ^\times$，并且 $a_1$ 与 $b_1$ 都与 $p$ 互素。定义
  \[
    v_p\left(\frac ab\right)=v_p\left(p^n\frac{a_1}{b_1}\right)=n.
  \]
  容易检查 D.V.R. 的公理成立。我们称 $v_p$ 为 $\QQ$ 上的 \emph{$p$-进赋值}。相应的赋值环是由 $n\geq0$ 的有理数以及 $0$ 组成的集合，也就是分母不被 $p$ 整除的有理数 $a/b$ 的集合，这正是 $\ZZ_{(p)}$。

  \item 对任意域 $F$，设 $f$ 是 $F[x]$ 中的不可约多项式。域 $F(x)$ 中每个非零元素都可唯一写成
  \[
    f^n(a/b)
  \]
  的形式，其中 $n\in\ZZ$，$a/b\in F[x]^\times$，并且 $a$ 与 $b$ 都与 $f$ 互素。于是
  \[
    v_f\left(f^n\frac ab\right)=n
  \]
  定义了 $F(x)$ 上的一个赋值，相应的赋值环是 $F[x]$ 在 $f$ 处的局部化 $F[x]_f$，它由 $F(x)$ 中分母不被 $f$ 整除的有理函数组成。当 $f=x-a$ 是 $F[x]$ 中的一次多项式时，赋值 $v_f$ 给出 $F(x)$ 中元素在 $x=a$ 处的零点阶数（若 $n\geq0$）或极点阶数（若 $n<0$）。

  \item 系数在域 $F$ 中的形式 Laurent 级数环 $F((x))$ 有一个离散赋值 $v$，定义为
  \[
    v\left(\sum_{i\geq n}a_ix^i\right)=n
  \]
  （参见第 7.2 节练习 5）。相应的 D.V.R. 是系数在 $F$ 中的 $x$ 的幂级数环 $F[[x]]$。
\end{enumerate}

注意，$v(1)=v(1)+v(1)$ 蕴含 $v(1)=0$，所以每个离散赋值环 $R$ 都是带有恒等元 $1\neq0$ 的环。由于按定义 $R$ 是域的子环，所以 $R$ 特别是整环。容易看出 D.V.R. 是 Euclidean 整环（参见第 8.1 节例 4），因此特别地也是 P.I.D. 和 U.F.D.。事实上，D.V.R. 的分解和理想结构非常简单，如下一个命题所示。

\medskip
\noindent\textbf{命题 5.} 设 $R$ 是相对于赋值 $v$ 的离散赋值环，并令 $t$ 为 $R$ 中任意满足 $v(t)=1$ 的元素。则：
\begin{enumerate}[label=(\arabic*)]
  \item 非零元素 $u\in R$ 是单位，当且仅当 $v(u)=0$。
  \item 每个非零元素 $r\in R$ 都可写成 $r=ut^n$，其中 $u\in R$ 是单位，$n\geq0$。$R$ 的分式域中每个非零元素 $x$ 都可写成 $x=ut^n$，其中 $u\in R$ 是单位，$n\in\ZZ$。
  \item $R$ 的每个非零理想都是形如 $(t^n)$ 的主理想，其中 $n\geq0$。特别地，$R$ 是 Noether 环。
\end{enumerate}

\begin{proof}
若 $u$ 是单位，则对某个 $v\in R$ 有 $uv=1$，从而 $v(u)+v(v)=v(uv)=0$，并且 $v(u)\geq0$、$v(v)\geq0$，所以 $v(u)=0$。反过来，若 $u$ 非零且 $v(u)=0$，则 $u^{-1}\in K$ 满足 $v(u^{-1})+v(u)=v(1)=0$。因此 $v(u^{-1})=0$，故 $u^{-1}\in R$，所以 $u$ 是单位。这证明了 (1)。

对 (2)，注意若 $v(x)=n$，则 $v(xt^{-n})=0$，所以由 (1) $xt^{-n}=u$ 是 $R$ 中的单位。因此 $x=ut^n$，其中 $x\in R$ 当且仅当 $n=v(x)\geq0$。

若 $I$ 是 $R$ 中的非零理想，取 $r\in I$ 使 $v(r)$ 最小。若 $v(r)=n$，则由 (2)，$r$ 与 $t^n$ 只差一个单位，所以 $t^n\in I$，且 $(t^n)\subseteq I$。若 $a$ 是 $I$ 的任意非零元素，则由 $n$ 的选择有 $v(a)\geq n$。于是 $v(at^{-n})\geq0$，故 $at^{-n}\in R$，这说明 $a\in(t^n)$。因此 $I=(t^n)$，证明了 (3) 的第一个断言。于是 $R$ 中的理想升链显然都是有限的，故 $R$ 是 Noether 环，证明完成。
\end{proof}

\noindent\textbf{定义.} 若 $R$ 是赋值为 $v$ 的 D.V.R.，则 $R$ 中满足 $v(t)=1$ 的元素 $t$ 称为 $R$ 的\emph{一致化参数}（或\emph{局部参数}）。

\medskip
\noindent\textbf{推论 6.} 设 $R$ 是离散赋值环。
\begin{enumerate}[label=(\arabic*)]
  \item 环 $R$ 是整闭局部环，其唯一极大理想由严格正赋值的元素给出：
  \[
    M=\{r\in R\mid v(r)>0\}.
  \]
  $R$ 中每个非零理想都形如 $M^n$，其中 $n\geq0$。
  \item $R$ 的唯一素理想是 $M$ 和 $0$，即 $\Spec R=\{0,M\}$。特别地，D.V.R. 的 Krull 维数为 $1$。
\end{enumerate}

\begin{proof}
任意 U.F.D. 在其分式域中都是整闭的（第 15.3 节例 3），所以 $R$ 是整闭的。其余断言由命题 5 中对 $R$ 的理想的描述立刻得到。
\end{proof}

离散赋值环的定义是以分式域上的赋值给出的，极为显式；因此，看起来似乎很难仅从给定环 $R$ 的“内部”代数性质判断 $R$ 是否为 D.V.R.。事实上，命题 5 与推论 6 中的环论性质刻画了离散赋值环。下面的定理给出了离散赋值环的若干等价代数描述，其中完全没有显式提到赋值。

\medskip
\noindent\textbf{定理 7.} 环 $R$ 的下列性质等价：
\begin{enumerate}[label=(\arabic*)]
  \item $R$ 是离散赋值环；
  \item $R$ 是 P.I.D.，且有唯一非零极大理想 $P$；
  \item $R$ 是 U.F.D.，且只有一个不可约元素 $t$（在相伴意义下唯一）；
  \item $R$ 是 Noether 整环，也是局部环，且其唯一极大理想非零并且是主理想；
  \item $R$ 是 Noether、整闭的整环，同时也是 Krull 维数为 $1$ 的局部环，即 $R$ 有唯一非零素理想：$\Spec R=\{0,M\}$。
\end{enumerate}

\begin{proof}
(1) 蕴含其他每个性质已在上面证明。

若 (2) 成立，则 (3) 立刻成立，因为在 U.F.D. 中不可约元素生成素理想（第 8.3 节命题 12）。

若 (3) 成立，则 $R$ 中每个非零元素都可唯一写成 $ut^n$，其中 $u$ 是单位，$n\geq0$。于是 $R$ 的分式域中每个非零元素都可唯一写成 $ut^n$，其中 $u$ 是单位，$n\in\ZZ$。现在直接检查映射 $v(ut^n)=n$ 是 $R$ 的分式域上的一个离散赋值，且 $R$ 是 $v$ 的赋值环，因此 (1) 成立。

假设 (4) 成立，令 $M=(t)$ 为 $R$ 的唯一极大理想，并令 $M_0=\bigcap_{i\geq1}M^i$。则 $M_0=MM_0$，并且由于 $R$ 是 Noether 环，$M_0$ 是有限生成的。由假设 $M=\Jac R$，所以由 Nakayama 引理 $M_0=0$。若 $I$ 是 $R$ 的任意真非零理想，则存在某个 $n\geq0$，使得 $I\subseteq M^n$ 但 $I\nsubseteq M^{n+1}$。取 $a\in I-M^{n+1}$，并写 $a=t^nu$，其中 $u\in R$。于是 $u\notin M$，所以 $u$ 是局部环 $R$ 中的单位。因此对每个 $a\in I-M^{n+1}$，$(a)=(t^n)=M^n$。这说明 $I=(t^n)$，所以 $R$ 的每个理想都是主理想，从而 (2) 成立。

我们已经证明了 (1)、(2)、(3)、(4) 等价，并且每个都蕴含 (5)。为完成证明，我们证明 (5) 蕴含 (4)，也就是证明 (5) 中的理想 $M$ 是主理想。由于 $0\neq M=\Jac R$ 且 $M$ 因 $R$ 为 Noether 环而有限生成，由 Nakayama 引理（命题 1(4)），$M\neq M^2$。取 $t\in M-M^2$。我们断言 $M=(t)$。由第 15.2 节命题 12，$M$ 是 $R$ 中唯一非零素理想的假设蕴含 $M=\rad(t)$，于是第 15.2 节命题 14 蕴含 $M$ 的某个幂包含于 $(t)$ 中。反证，假设 $(t)\neq M$，于是对某个 $n\geq2$ 有 $M^n\subseteq(t)$ 但 $M^{n-1}\nsubseteq(t)$。则存在元素 $x\in M^{n-1}-(t)$，使得 $xM\subseteq(t)$。注意 $t\neq0$，所以 $y=x/t$ 属于 $R$ 的分式域。而且 $y\notin R$，因为 $x=ty\notin(t)$。然而，由 $x$ 的选择可知 $yM\subseteq R$，并且容易检查 $yM$ 是 $R$ 中的理想。若 $yM=R$，则对某个 $m\in M$ 有 $1=ym$。这导致矛盾，因为这时 $t=xm\in M^2$，与 $t$ 的选择相反。因此 $yM$ 是真理想，故包含在 $R$ 的唯一极大理想中，即 $yM\subseteq M$。现在 $M$ 是有限生成 $R$-模，而 $y$ 通过左乘作用于 $M$，给出一个 $R$-模同态。与第 15.3 节命题 23 的证明中相同的（行列式）方法表明，存在一个系数在 $R$ 中的首一多项式 $p$，使得对所有 $m\in M$ 都有 $p(y)m=0$。由于 $p(y)$ 是含有 $R$ 与 $M$ 的域中的元素，必有 $p(y)=0$。因此 $y$ 在 $R$ 上整。由于按假设 $R$ 是整闭的，推出 $y\in R$，矛盾。因此 $M=(t)$ 是主理想，所以 (5) 蕴含 (4)，定理证毕。
\end{proof}

\medskip
\noindent\textbf{推论 8.} 若 $R$ 是任意 Noether、整闭的整环，且 $P$ 是 $R$ 的极小非零素理想，则 $R$ 在 $P$ 处的局部化 $R_P$ 是离散赋值环。

\begin{proof}
由第 15.4 节的结果，局部化 $R_P$ 是 Noether 的（命题 38(4)）、整闭的（命题 49）、整环（命题 46(2)），并且是具有唯一非零素理想的局部环（命题 46(4)），所以 $R_P$ 满足定理中的 (5)。
\end{proof}

\paragraph{例}
\begin{enumerate}[label=(\arabic*)]
  \item 若 $R$ 是任意主理想整环，则 $R$ 在任意非零素理想 $P=(p)$ 处的局部化 $R_P$ 都是离散赋值环。这由推论 8 立刻推出，因为 $R$ 是整闭的（作为 U.F.D.，参见第 15.3 节例 3），而 P.I.D. 中非零素理想都是极大理想（命题 8.7）。注意 $R_P$ 的分式域 $K$ 与 $R$ 的分式域相同，因此 $R$ 中每个非零素元 $p$ 都在 $K$ 上产生一个赋值 $v_p$，由公式
  \[
    v_p\left(p^n\frac ab\right)=n
  \]
  给出，其中 $a$ 和 $b$ 是 $R$ 中不被 $p$ 整除的元素。这同时推广了上面的例 1 和例 2。

  \item $p$-进整数环 $\ZZ_p$ 是离散赋值环，因为它是 P.I.D.，且有唯一极大理想 $p\ZZ_p$（参见第 7.6 节练习 11）。$\ZZ_p$ 的分式域称为 \emph{$p$-进数域}，记为 $\QQ_p$。元素 $p$ 是 $\ZZ_p$ 的一致化参数，所以 $\QQ_p$ 中每个非零元素都可唯一写成 $p^nu$ 的形式，其中 $n\in\ZZ$，$u\in\ZZ_p^\times$ 是单位（其中 $u=a_0+a_1p+a_2p^2+
  \cdots$，且 $0<a_0<p$，如第 7.6 节练习 11(c) 所述）。相应的 $\QQ_p$ 上的 $p$-进赋值由 $v_p(p^nu)=n$ 给出。
\end{enumerate}

域 $K$ 上的离散赋值 $v$ 定义了 $K$ 上一个相应的度量（或“距离函数”）$d_v$，如下：固定任意实数 $\beta>1$（为验证度量公理，$\beta$ 的实际取值无关紧要），并对所有 $a,b\in K$ 定义
\[
  d_v(a,b)=\lVert a-b\rVert_v,
  \qquad \text{其中}\qquad
  \lVert a\rVert_v=\beta^{-v(a)},
\]
并且规定 $d_v(a,a)=0$。容易检查 $d_v$ 满足度量的三个公理：
\begin{enumerate}[label=(\roman*)]
  \item $d_v(a,b)\geq0$，且等号成立当且仅当 $a=b$；
  \item $d_v(a,b)=d_v(b,a)$，即 $d_v$ 是对称的；
  \item 对所有 $a,b,c\in K$，$d_v(a,b)\leq d_v(a,c)+d_v(c,b)$，即 $d_v$ 满足“三角不等式”。
\end{enumerate}
三角不等式是离散赋值公理 (iii) 的结果。事实上，更强的三角不等式成立：
\begin{enumerate}[label=(\roman*)]
  \item[(iii)$'$] 对所有 $a,b,c\in K$，
  \[
    d_v(a,b)\leq\max\{d_v(a,c),d_v(c,b)\}.
  \]
\end{enumerate}
因此 $d_v$ 有时称为\emph{超度量}。现在可以用 Cauchy 列来形成 $K$ 关于 $d_v$ 的完备化，记为 $K_v$，其方式与由有理数 $\QQ$ 构造实数 $\RR$ 相同。不难证明 $K_v$ 也是一个域，并带有一个离散赋值，该赋值在 $K_v$ 的稠密子集 $K$ 上与 $v$ 一致。

\paragraph{例}
\begin{enumerate}[label=(\arabic*)]
  \item 考虑 $\QQ$ 上的 $p$-进赋值 $v_p$，并取 $\beta=p$。把 $\lVert a\rVert_{v_p}$ 写作 $\lVert a\rVert_p$，于是对与 $p$ 互素的 $a,b$，
  \[
    \left\lVert p^n\frac ab\right\rVert_p=p^{-n}.
  \]
  注意，若整数（或有理数）可被 $p$ 的高次幂整除，则它们的 $p$-进绝对值很小。例如，序列 $1,p,p^2,p^3,
  \ldots$ 在 $p$-进度量下收敛到零。

  不太难看出，$\QQ$ 关于 $p$-进度量的完备化就是 $p$-进数域 $\QQ_p$，而 $\ZZ$ 的完备化就是 $p$-进整数环 $\ZZ_p$。一种看法是检查完备化中的每个元素 $a$ 都可以表示为一个 $p$-进 Laurent 级数：
  \[
    a=\sum_{i=n_0}^{\infty}a_ip^i,
    \qquad n_0\in\ZZ,
    \quad a_i\in\{0,1,
    \ldots,p-1\}\text{ 对所有 }i,
  \]
  然后使用前面的例 2。用这种展开表示时，$p$-进赋值由 $v_p(a)=n_0$ 给出（当 $a_{n_0}\neq0$ 时）。

  \item 类似地，$F(x)$ 关于本节开头例 2 中赋值 $v_x$ 的完备化给出域 $F((x))$，其相应的赋值环是同一组例子中例 3 的 $F[[x]]$。
\end{enumerate}

域 $K$ 关于离散赋值 $v$ 的完备化是一个域 $K_v$，其中的元素可用一致化参数容易地描述。此外，$K_v$ 是一个拓扑空间，其拓扑由度量 $d_v$ 定义。再者，$K_v$ 中元素的 Cauchy 列收敛到 $K_v$ 中的元素（即 $K_v$ 在 $v$-进拓扑下是完备的）。这类似于 $\QQ$ 关于通常 Euclidean 度量的完备化 $\RR$ 的情形。这使得可以把分析中的思想应用于这类环的研究，并且是研究代数数域和代数几何中的重要工具。

\subsection*{分式理想}

我们通过给出 D.V.R. 关于“分式理想”的另一个刻画，来完成对离散赋值环的讨论；分式理想可以对任意整环定义：

\noindent\textbf{定义.} 对任意分式域为 $K$ 的整环 $R$，$R$ 的\emph{分式理想}是 $K$ 的 $R$-子模 $A$，使得对某个非零 $d\in R$ 有 $dA\subseteq R$（等价地，它是形如 $d^{-1}I$ 的子模，其中 $d\in R$ 非零，$I$ 是 $R$ 的理想）。

这两个定义等价，是因为 $dA$ 是 $R$ 的 $R$-子模（即理想）。

$K$ 中分式理想的概念依赖于环 $R$。粗略地说，分式理想就是允许有固定“分母” $d$ 的 $R$ 的理想。$R$ 的理想也是 $R$ 的分式理想（分母取 $d=1$），并且正是那些包含在 $R$ 中的分式理想。为清楚起见，它们有时称为 $R$ 的\emph{整理想}。当 $R$ 是 Noether 整环时，$R$ 的分式理想等同于 $K$ 的有限生成 $R$-子模（参见练习 6）。

对任意 $x\in K$，$R$-模 $Rx=\{rx\mid r\in R\}$ 称为由 $x$ 生成的（循环）\emph{主分式理想}。

若 $A$ 与 $B$ 是分式理想，则它们的乘积 $AB$ 定义为所有形如 $ab$ 的元素的有限和组成的集合，其中 $a\in A$ 且 $b\in B$。若 $A=d^{-1}I$ 且 $B=(d')^{-1}J$，其中 $I,J$ 是 $R$ 的理想，$d,d'\in R$ 非零，则 $AB=(dd')^{-1}IJ$，其中 $IJ$ 是通常的乘积理想。特别地，这说明两个分式理想的乘积仍是分式理想。

\noindent\textbf{定义.} 分式理想 $A$ 称为\emph{可逆的}，若存在分式理想 $B$ 使得 $AB=R$；此时 $B$ 称为 $A$ 的逆，记为 $A^{-1}$。

若 $A$ 是可逆分式理想，则满足 $AB=R$ 的分式理想 $B$ 是唯一的：$AB=AC=R$ 蕴含 $B=B(AC)=(BA)C=C$。

\medskip
\noindent\textbf{命题 9.} 设 $R$ 是整环，且 $A$ 是 $R$ 的分式理想。
\begin{enumerate}[label=(\arabic*)]
  \item 若 $A$ 是非零主分式理想，则 $A$ 可逆。
  \item 若 $A$ 非零，则集合
  \[
    A'=\{x\in K\mid xA\subseteq R\}
  \]
  是 $R$ 的分式理想。一般地，有 $AA'\subseteq R$；且 $AA'=R$ 当且仅当 $A$ 可逆，在这种情形下 $A^{-1}=A'$。
  \item 若 $A$ 是 $R$ 的可逆分式理想，则 $A$ 是有限生成的。
  \item 可逆分式理想全体在乘法下构成一个 Abel 群，恒等元为 $R$。非零主分式理想全体是可逆分式理想群的子群。
\end{enumerate}

\begin{proof}
若 $A=xR$ 是非零主分式理想，则取 $B=x^{-1}R$ 说明 $A$ 可逆，证明了 (1)。

容易看出 $A'$ 是 $K$ 的 $R$-子模。若 $A$ 是非零分式理想，则存在某个非零 $d\in R$ 使 $dA\subseteq R$，所以 $A$ 含有 $R$ 中的非零元素。令 $a$ 是 $A$ 中任意一个也在 $R$ 中的非零元素。由 $A'$ 的定义，有 $aA'\subseteq R$，所以 $A'$ 是分式理想。同样由定义，$AA'\subseteq R$。若 $AA'=R$，则 $A$ 可逆，逆为 $A^{-1}=A'$。反过来，若 $AB=R$，则由 $A'$ 的定义 $B\subseteq A'$。于是 $R=AB\subseteq AA'\subseteq R$，说明 $AA'=R$，证明了 (2)。

若 $A$ 可逆，则由 (2) 有 $AA'=R$，所以对某些 $a_1,
\ldots,a_n\in A$ 和 $a_1',
\ldots,a_n'\in A'$，
\[
  1=a_1a_1'+\cdots+a_na_n'.
\]
若 $a\in A$，则
\[
  a=(aa_1')a_1+\cdots+(aa_n')a_n,
\]
其中每个 $aa_i'\in R$，这是由 $A'$ 的定义得到的。因此 $A$ 由 $a_1,
\ldots,a_n$ 在 $R$ 上生成，所以 $A$ 是有限生成的，证明了 (3)。

最后，显然两个可逆分式理想的乘积仍是可逆的。该乘法交换、结合，并且对任意分式理想 $A$ 有 $RA=A$。可逆分式理想的逆按定义仍是可逆分式理想，这证明了 (4) 的第一个断言。(4) 的第二个断言是直接的，因为 $xR$ 与 $yR$ 的乘积是 $(xy)R$，而 $xR$ 的逆是 $x^{-1}R$。
\end{proof}

\noindent\textbf{定义.} 若 $R$ 是整环，则 $R$ 的可逆分式理想群关于非零主分式理想子群的商群称为 $R$ 的\emph{类群}。$R$ 的类群的阶称为 $R$ 的\emph{类数}。

$R$ 的类群是平凡群且 $R$ 的类数为 $1$，当且仅当 $R$ 是 P.I.D.。$R$ 的类群衡量 $R$ 的理想距离成为主理想有多接近。

分式理想 $A$ 是否可逆，也与 $A$ 作为 $R$-模是否投射有关。回忆，$R$-模 $M$ 在 $R$ 上是投射的，当且仅当 $M$ 是自由模的直和项（第 10.5 节命题 30）。等价地，$M$ 是投射的，当且仅当存在自由 $R$-模 $F$ 以及 $R$-模同态 $f:F\to M$、$g:M\to F$，使得 $f\circ g=1$（第 10.5 节命题 25）。

\medskip
\noindent\textbf{命题 10.} 设 $R$ 是分式域为 $K$ 的整环，且 $A$ 是 $R$ 的非零分式理想。则 $A$ 可逆，当且仅当 $A$ 是投射 $R$-模。

\begin{proof}
先假设 $A$ 可逆，于是如命题 9(2) 中一样，对某些 $a_i\in A$ 与 $a_i'\in A'$ 有
\[
  \sum_{i=1}^n a_ia_i'=1.
\]
令 $F$ 为由 $y_1,
\ldots,y_n$ 生成的自由 $R$-模。定义 $f:F\to A$ 为
\[
  f\left(\sum_{i=1}^n r_iy_i\right)=\sum_{i=1}^n r_ia_i,
\]
并定义 $g:A\to F$ 为
\[
  g(c)=\sum_{i=1}^n(ca_i')y_i.
\]
二者显然都是 $R$-模同态（注意由 $A'$ 的定义 $ca_i'\in R$）。由于
\[
  (f\circ g)(c)=f\left(\sum_{i=1}^n(ca_i')y_i\right)
  =\sum_{i=1}^n(ca_i')a_i
  =c\left(\sum_{i=1}^n a_ia_i'\right)=c,
\]
所以 $f\circ g=1$，且 $A$ 是 $F$ 的直和项，因此是投射的。

反过来，假设 $A$ 非零且投射，所以存在自由 $R$-模 $F$ 以及 $R$-同态 $f:F\to A$、$g:A\to F$，满足 $f\circ g=1$。固定任意 $0\neq a\in A$，并设
\[
  g(a)=\sum_{i=1}^n \widetilde a_i y_i,
\]
其中 $\widetilde a_i\in R$，而 $y_1,
\ldots,y_n$ 是 $F$ 的一组自由生成元集合的一部分。定义 $a_i=f(y_i)$ 且 $a_i'=\widetilde a_i/a\in K$，$i=1,
\ldots,n$。对任意 $b\in A$，由于 $g$ 是 $R$-模同态，我们有 $bg(a)=ag(b)=g(ab)$。写
\[
  g(b)=\sum_{i=1}^n b_i y_i+\sum_{j\in\mathcal I}b_jy_j,
\]
其中 $\{y_j\}$，$j\in\mathcal I$，是 $F$ 的自由生成元集合中其余的元素。于是
\[
  \sum_{i=1}^n (b\widetilde a_i)y_i
  =\sum_{i=1}^n (ab_i)y_i+
  \sum_{j\in\mathcal I}(ab_j)y_j.
\]
在这个等式中比较自由 $R$-模 $F$ 的基元素的系数，得到 $g(b)=\sum_{i=1}^n b_iy_i$，其中 $b_i\in R$，并且对 $i=1,
\ldots,n$ 有 $b\widetilde a_i=ab_i$。特别地，由 $a_i'$ 的定义可知
\[
  ba_i'=b(\widetilde a_i/a)=b_i
\]
对 $A$ 的每个元素 $b$ 都属于 $R$。这说明 $a_i'\in A'$，$i=1,
\ldots,n$。由于 $f\circ g=1$，我们有
\[
  a=f\circ g(a)=f\left(\sum_{i=1}^n\widetilde a_i y_i\right)
  =\sum_{i=1}^n\widetilde a_i a_i
  =\sum_{i=1}^n(aa_i')a_i
  =a\left(\sum_{i=1}^n a_ia_i'\right),
\]
所以 $\sum_{i=1}^n a_ia_i'=1$。于是 $AA'=R$，所以由命题 9，$A$ 可逆，证明完成。
\end{proof}

下一个结果说明，如果整环 $R$ 还是局部环，则分式理想是否可逆决定了 $R$ 是否为 D.V.R.。

\medskip
\noindent\textbf{命题 11.} 假设整环 $R$ 是不是域的局部环。则 $R$ 是离散赋值环，当且仅当 $R$ 的每个非零分式理想都是可逆的。

\begin{proof}
若 $R$ 是一致化参数为 $t$ 的 D.V.R.，则由命题 5，$R$ 的每个非零理想都形如 $(t^n)$，其中 $n\geq0$，且 $R$ 的每个元素 $d$ 都可写成 $ut^m$，其中 $u\in R$ 是单位，$m\geq0$。于是 $R$ 的每个非零分式理想都形如 $t^N R$，其中 $N\in\ZZ$，因而是主分式理想，并由前一命题可逆。

反过来，假设 $R$ 的每个非零分式理想都是可逆的。则由命题 9(3)，$R$ 的每个非零理想都是有限生成的，所以 $R$ 是 Noether 环。令 $M$ 为 $R$ 的唯一极大理想。若 $M=M^2$，则由 Nakayama 引理 $M=0$，这将使 $R$ 成为域，与假设相反。因此存在元素 $t\in M-M^2$。由假设 $M$ 可逆，并且由于 $t\in M$，分式理想 $tM^{-1}$ 是 $R$ 中的非零理想。若 $tM^{-1}\subseteq M$，则 $t\in M^2$，与 $t$ 的选择矛盾。因此 $tM^{-1}=R$，所以 $(t)=M$，而 $M$ 是非零主理想。于是由定理 7 的等价条件 4，$R$ 是 D.V.R.，证明完成。
\end{proof}

我们用代数几何中的一个应用结束本节。

\subsection*{仿射平面曲线的非奇异性和局部环}

设 $k$ 是代数闭域，并设 $C$ 是 $k$ 上的不可约仿射曲线。换言之，$C$ 是一个仿射代数集，其坐标环 $k[C]$ 是整环，且其有理函数域 $k(C)$ 在 $k$ 上的超越次数为 $1$（参见第 15.4 节）。

回忆，按照定义，$C$ 上的点 $v$ 是非奇异的，如果 $\mathfrak m_{v,C}/\mathfrak m_{v,C}^2$ 是 $k$ 上的一维向量空间，其中 $\mathfrak m_{v,C}$ 是在 $v$ 处有定义的 $C$ 上有理函数的局部环 $\mathcal O_{v,C}$ 中的唯一极大理想。

\medskip
\noindent\textbf{命题 12.} 设 $v$ 是 $k$ 上不可约仿射曲线 $C$ 上的点。则 $C$ 在 $v$ 处非奇异，当且仅当局部环 $\mathcal O_{v,C}$ 是离散赋值环。

\begin{proof}
先假设 $v$ 非奇异。则 $\dim_k(\mathfrak m_{v,C}/\mathfrak m_{v,C}^2)=1$，并且由于 $\mathcal O_{v,C}$ 是 Noether 环，由第 1 节练习 12 可知 $\mathfrak m_{v,C}$ 是主理想。因此由定理 7(4)，$\mathcal O_{v,C}$ 是 D.V.R.。反过来，假设 $\mathcal O_{v,C}$ 是 D.V.R.，且 $t$ 是 $\mathcal O_{v,C}$ 的一致化元素。于是 $\mathfrak m_{v,C}$ 中每个元素都可唯一写成 $at$，其中 $a\in\mathcal O_{v,C}$。从 $\mathfrak m_{v,C}$ 到 $\mathcal O_{v,C}/\mathfrak m_{v,C}$ 的映射，定义为把 $at$ 映到 $a\bmod\mathfrak m_{v,C}$，容易检查它是满的 $\mathcal O_{v,C}$-模同态，核为 $\mathfrak m_{v,C}^2$。因此 $\mathfrak m_{v,C}/\mathfrak m_{v,C}^2$ 作为 $\mathcal O_{v,C}/\mathfrak m_{v,C}$-模同构于 $\mathcal O_{v,C}/\mathfrak m_{v,C}$。由于 $\mathcal O_{v,C}/\mathfrak m_{v,C}\cong k$（第 15.4 节命题 46(5)），得到 $\dim_k(\mathfrak m_{v,C}/\mathfrak m_{v,C}^2)=1$，因此 $v$ 是 $C$ 上的非奇异点。
\end{proof}

\noindent\textbf{定义.} 若 $v$ 是 $C$ 上的非奇异点，且对应的离散赋值 $v_v$ 定义在 $k(C)$ 上，则对于 $f\in k(C)$，若 $v_v(f)=n$，则称其为 $f$ 在 $v$ 处的零点阶数（若 $n\geq0$）或极点阶数（若 $n<0$）。

利用命题 12 中曲线上点的非奇异性判别准则，可以证明第 15.4 节首次提到的一个结果：

\medskip
\noindent\textbf{推论 13.} 代数闭域 $k$ 上的不可约仿射曲线 $C$ 是光滑的，当且仅当其坐标环 $k[C]$ 是整闭的。

\begin{proof}
曲线 $C$ 光滑，当且仅当每个局部化 $\mathcal O_{v,C}$ 都是 D.V.R.。由于 $k[C]$ 的 Krull 维数为 $1$（第 1 节练习 11），每个 $\mathcal O_{v,C}$ 亦然。因此由定理 7(5)，每个局部化 $\mathcal O_{v,C}$ 是 D.V.R.，当且仅当 $\mathcal O_{v,C}$ 是整闭的。由第 15.4 节命题 49，这又等价于 $k[C]$ 是整闭的，这证明了该推论。
\end{proof}

\subsection*{练习}

\begin{enumerate}[label=\arabic*.]
  \item 设 $R$ 是相对于其分式域 $K$ 上赋值 $v$ 的离散赋值环。若 $x,y\in K$ 且 $v(x)<v(y)$，证明 $v(x+y)=\min(v(x),v(y))$。［注意 $x+y=x(1+y/x)$。］
  \item 设 $R$ 是唯一极大理想为 $M$、商域为 $F=R/M$ 的离散赋值环。对任意 $n\geq0$，证明 $M^n/M^{n+1}$ 是 $F$ 上的向量空间，且
  \[
    \dim_F(M^n/M^{n+1})=1.
  \]
  \item 假设 $R$ 是整环，也是局部环，其唯一极大理想 $M=(t)$ 非零且为主理想，并假设
  \[
    \bigcap_{n\geq1}(t^n)=0.
  \]
  证明 $R$ 是离散赋值环。［证明 $R$ 中每个非零元素都可写成 $ut^n$ 的形式，其中 $u\in R$ 是单位，$n\geq0$。］
  \item 设 $R$ 是 Noether 局部环，其唯一极大理想 $M=(t)$ 是主理想。证明要么 $R$ 是离散赋值环，要么对某个 $n\geq0$ 有 $t^n=0$。在后一种情形下，证明 $R$ 是 Artin 环。
  \item 假设 $R$ 是 Noether 整环，并且也是 Krull 维数为 $1$ 的局部环。令 $M$ 为 $R$ 的唯一极大理想，并令 $F=R/M$，于是 $M/M^2$ 是 $F$ 上的向量空间。
  \begin{enumerate}[label=(\alph*)]
    \item 证明若 $\dim_F(M/M^2)=1$，则 $R$ 是离散赋值环。
    \item 若 $R$ 的每个非零理想都是 $M$ 的幂，证明 $R$ 是离散赋值环。
  \end{enumerate}
  \item 设 $R$ 是分式域为 $K$ 的整环。证明 $K$ 的每个有限生成 $R$-子模都是 $R$ 的分式理想。若 $R$ 是 Noether 环，证明 $A$ 是 $R$ 的分式理想，当且仅当 $A$ 是 $K$ 的有限生成 $R$-子模。
  \item 若 $R$ 是整环，且 $A$ 是 $R$ 的分式理想，证明若 $A$ 是投射的，则 $A$ 是有限生成的。推出每个非 Noether 的整环都含有一个不是投射的理想。
  \item 设 $R$ 是 Noether 整环，也是非零极大理想为 $M$ 的局部环。证明 $R$ 是 D.V.R.，当且仅当 $R$ 中唯一的 $M$-primary 理想是 $M$ 的幂。
  \item 设
  \[
    C=Z(xz-y^2,
    yz-x^3,
    z^2-x^2y)\subseteq\mathbb A^3
  \]
  是代数闭域 $k$ 上的曲线。若 $v=(0,0,0)\in C$，证明 $\dim_k(\mathfrak m_{v,C}/\mathfrak m_{v,C}^2)=3$，因此 $v$ 是 $C$ 上的奇异点。推出 $k[C]$ 在 $k(C)$ 中不是整闭的，并确定它的整闭包。［参见第 15.4 节练习 27。］
\end{enumerate}

\section{Dedekind 整环}

在上一节中，我们证明了离散赋值环就是整闭的、Krull 维数为 $1$ 的 Noether 整环中的局部环。本节考虑放松局部环条件之后的影响：

\noindent\textbf{定义.} \emph{Dedekind 整环}是 Krull 维数为 $1$ 的 Noether、整闭的整环。

等价地，$R$ 是 Dedekind 整环，当且仅当 $R$ 是 Noether、整闭且不是域的整环，并且其中每个非零素理想都是极大理想。

第一个结果说明，Dedekind 整环是主理想整环类的推广。稍后我们将看到（定理 22），有限生成 Dedekind 整环上模有一个结构定理，它推广了第 12.1 节中为 P.I.D. 证明的相应结果。

\medskip
\noindent\textbf{命题 14.}
\begin{enumerate}[label=(\arabic*)]
  \item 每个主理想整环都是 Dedekind 整环。
  \item 代数数域中的整数环是 Dedekind 整环。
\end{enumerate}

\begin{proof}
P.I.D. 显然是 Noether 的；它是 U.F.D.，所以是整闭的（第 15.3 节例 3）；并且非零素理想都是极大理想（第 8.2 节命题 7），这证明了 (1)。令 $\mathcal O_K$ 为数域 $K$ 中的整数环，即 $\ZZ$ 在 $K$ 中的整闭包。第 15.3 节推论 25 说明 $\mathcal O_K$ 是整闭的；第 15.3 节定理 29 说明 $\mathcal O_K$ 是 Noether 的；而 $\mathcal O_K$ 中非零素理想都是极大理想这一事实在同一定理之后的讨论中已被证明。这证明了 (2)。
\end{proof}

下面的定理给出 Dedekind 整环的一些重要等价刻画。回忆，分式理想的基本性质已在上一节中建立。

\medskip
\noindent\textbf{定理 15.} 设 $R$ 是分式域为 $K\neq R$ 的整环。以下条件对于 $R$ 为 Dedekind 整环是等价的：
\begin{enumerate}[label=(\arabic*)]
  \item 环 $R$ 是 Noether 的、整闭的，并且每个非零素理想都是极大理想。
  \item 环 $R$ 是 Noether 的，并且对 $R$ 的每个非零素理想 $P$，局部化 $R_P$ 是离散赋值环。
  \item $K$ 中 $R$ 的每个非零分式理想都是可逆的。
  \item $K$ 中 $R$ 的每个非零分式理想都是投射 $R$-模。
  \item $R$ 的每个非零真理想 $I$ 都可写成有限多个素理想的乘积：
  \[
    I=P_1P_2\cdots P_n
  \]
  （不一定互异）。当 (5) 中的条件成立时，素理想集合 $\{P_1,
  \ldots,P_n\}$ 是唯一确定的，因此 $R$ 的每个非零真理想 $I$ 都可唯一地（差一个顺序）写成素理想幂的乘积。
\end{enumerate}

\begin{proof}
若 $R$ 满足 (1)，则由推论 8，$R_P$ 是 D.V.R.，所以 (1) 蕴含 (2)。反过来，假设每个 $R_P$ 都是 D.V.R.。则由第 15.4 节命题 49，$R$ 是整闭的；并且由第 15.4 节命题 46(3)，每个非零素理想都是极大理想，所以 (2) 蕴含 (1)。

现在假设 (1) 成立，且 $A$ 是 $R$ 的非零分式理想。令
\[
  A'=\{x\in K\mid xA\subseteq R\}
\]
如命题 9 中所定义。对 $R$ 的任意素理想 $P$，$R$-模在局部化下的行为说明
\[
  (AA')_P=A_P(A')_P=A_P(A_P)'
\]
（参见练习 4）。由于按已证结果 $R_P$ 是 D.V.R.，由命题 11 有 $A_P(A_P)'=R_P$。因此对 $R$ 的所有非零素理想 $P$，$(AA')_P=R_P$，所以 $AA'=R$（第 15.4 节练习 13），于是 $A$ 可逆，说明 (1) 蕴含 (3)。反过来，假设 $R$ 的每个非零分式理想都可逆。则由命题 9(3)，$R$ 的每个理想都是有限生成的，所以 $R$ 是 Noether 环。$R$ 在任意非零素理想 $P$ 处的局部化 $R_P$ 是局部环，其中的非零分式理想都是可逆的（参见练习 4），因此由命题 11 是 D.V.R.。故 (3) 蕴含 (2)，从而 (1)、(2)、(3) 等价。它们与 (4) 的等价性由命题 10 给出。

现在假设 (1) 成立，令 $I$ 是 $R$ 的任意非零真理想。由于 $R$ 是 Noether 环，$I$ 有第 15.2 节定理 21 中的极小 primary 分解
\[
  I=Q_1\cap\cdots\cap Q_n.
\]
伴随素理想 $P_i=\rad Q_i$，$i=1,
\ldots,n$，彼此互异；由于按假设 $R$ 中素理想都是极大理想，这些伴随素理想两两互素，并容易推出 $Q_i$ 也两两互素（练习 5）。于是
\[
  Q_1\cap\cdots\cap Q_n=Q_1\cdots Q_n
\]
（第 7.6 节定理 17），所以 $I$ 是 primary 理想的乘积。$R$ 中的 $P$-primary 理想与局部化 $R_P$ 中的 $PR_P$-primary 理想双射对应（第 15.4 节命题 42(3)）；由于 $R_P$ 是 D.V.R.（因为 (1) 蕴含 (2)），由推论 6 可知，若 $Q$ 是 $R$ 中的 $P$-primary 理想，则 $Q=P^m$ 对某个整数 $m\geq1$ 成立。把这应用于 $Q_i$，$i=1,
\ldots,n$，可知 $I$ 是素理想幂的乘积，这给出 (5) 中的第一个蕴含。

反过来，假设 $R$ 的所有非零真理想都可写成素理想的乘积。我们先说明，对任意整环，把理想分解为可逆素理想的乘积是唯一的；也就是说，若
\[
  P_1\cdots P_n=P'_1\cdots P'_m
\]
是 $I$ 分解为可逆素理想乘积的两个分解，则 $n=m$，且两个素理想集合 $\{P_1,
\ldots,P_n\}$ 与 $\{P'_1,
\ldots,P'_m\}$ 相同。设 $P'_1$ 是集合 $\{P'_1,
\ldots,P'_m\}$ 中的极小元。由于 $P_1\cdots P_n\subseteq P'_1$，素理想 $P'_1$ 包含 $P_1,
\ldots,P_n$ 中的某个素理想，设 $P_1\subseteq P'_1$。类似地，$P_1$ 包含某个 $P'_i$，其中 $i=1,
\ldots,m$。于是 $P'_i\subseteq P_1\subseteq P'_1$；由 $P'_1$ 的极小性，得到 $P'_i=P_1=P'_1$，所以分解变为
\[
  P_1P_2\cdots P_n=P_1P'_2\cdots P'_m.
\]
由于 $P_1$ 可逆，乘以其逆理想得到 $P_2\cdots P_n=P'_2\cdots P'_m$，一个容易的归纳完成证明。特别地，(5) 中的唯一性断言现在由 (5) 的第一个断言推出，因为在 Dedekind 整环中每个分式理想，特别地 $R$ 的每个素理想，都是可逆的。

接下来证明 $R$ 中的可逆素理想都是极大理想。设 $P$ 是 $R$ 中可逆素理想，取 $a\in R$，$a\notin P$。我们要证明 $P+aR=R$。由假设，两个理想 $P+aR$ 和 $P+a^2R$ 可写成素理想乘积，设
\[
  P+aR=P_1\cdots P_n,
  \qquad
  P+a^2R=P'_1\cdots P'_m.
\]
注意，对 $i=1,
\ldots,n$ 有 $P\subseteq P_i$，并且对 $j=1,
\ldots,m$ 有 $P\subseteq P'_j$。在商环 $R/P$（一个整环）中，我们有分解
\[
  (a)=(P_1/P)\cdots(P_n/P),
\]
且每个 $P_i/P$ 是 $R/P$ 中的素理想。由于该乘积是主理想，每个 $P_i/P$ 也是可逆的 $R/P$-理想（参见练习 2）。类似地，
\[
  (a^2)=(P'_1/P)\cdots(P'_m/P)
\]
是分解为可逆素理想的乘积。于是
\[
  (a)^2=(P_1/P)^2\cdots(P_n/P)^2=(P'_1/P)\cdots(P'_m/P)
\]
给出了整环 $R/P$ 中分解为可逆素理想乘积的两个分解，所以由前一段的唯一性结果，$m=2n$，并且
\[
  \{P_1/P,P_1/P,
  \ldots,P_n/P,P_n/P\}=\{P'_1/P,
  \ldots,P'_m/P\}.
\]
因此 $R$ 中素理想 $P'_1,
\ldots,P'_m$ 的集合由 $P_1,
\ldots,P_n$ 每个重复两次组成。这说明
\[
  P+a^2R=(P+aR)^2.
\]
由于 $P\subseteq P+a^2R$ 且 $(P+aR)^2\subseteq P^2+aR$，我们有 $P\subseteq P^2+aR$，因此 $P$ 中每个元素 $x$ 都可写成 $x=y+az$，其中 $y\in P^2$，$z\in R$。于是 $az=x-y\in P$，且由于 $a\notin P$，有 $z\in P$。这说明 $P\subseteq P^2+aP$。显然 $P^2+aP\subseteq P$，所以
\[
  P=P^2+aP=P(P+aR).
\]
由于假设 $P$ 可逆，推出对任意 $a\in R-P$ 都有 $R=P+aR$，这证明 $P$ 是极大理想。

现在证明每个非零素理想都是可逆的。若 $P$ 是非零素理想，取任意非零元素 $a\in P$。由假设，主理想
\[
  Ra=P_1\cdots P_n
\]
可写成素理想的乘积，并且 $P_1,
\ldots,P_n$ 可逆，因为它们的乘积是主理想（再由练习 2）。由于 $P_1\cdots P_n=Ra\subseteq P$，素理想 $P$ 包含某个 $P_i$，$1\leq i\leq n$。由前一段，$P_i$ 是极大理想，所以 $P=P_i$ 可逆。

最后，由于 $R$ 的每个非零真理想都是素理想乘积，推出 $R$ 的每个非零理想都是可逆的；并且由于 $R$ 的每个分式理想都形如 $(d^{-1})I$，其中 $I$ 是 $R$ 中的理想，所以 $R$ 的每个分式理想也可逆。这证明了 (5) 蕴含 (3)，从而完成定理的证明。
\end{proof}

下面的推论立即由命题 14 得到：

\medskip
\noindent\textbf{推论 16.} 若 $\mathcal O_K$ 是代数数域 $K$ 中的整数环，则 $\mathcal O_K$ 中每个非零理想 $I$ 都可唯一写成不同素理想幂的乘积：
\[
  I=P_1^{e_1}P_2^{e_2}\cdots P_n^{e_n},
\]
其中 $P_1,
\ldots,P_n$ 是不同的素理想，并且对 $i=1,
\ldots,n$ 有 $e_i\geq1$。

\paragraph{注记}
这里给出的 Dedekind 整环的发展顺序与历史发展相反。如第 9.3 节中所提到的，在数域整数环中，非零理想唯一分解为素理想乘积替代了非零元素唯一分解为素元素乘积的失败。推论 16 中整数环的这一性质最初导致了理想的定义，而 Dedekind 最初正是用定理 15 中性质 5 来定义我们现在称为 Dedekind 整环的对象。Noether 观察到它们也可以由性质 (1) 来刻画，而我们已把性质 (1) 作为 Dedekind 整环的初始定义。

Dedekind 整环中分解为素理想的唯一分解可用来显式定义 $R$ 上的赋值 $v_P$，使得定理 15(2) 中的赋值环就是局部化 $R_P$（参见练习 6）。现在我们说明理想的唯一分解如何用于为理想定义一种类似于 $\ZZ$ 中整数可除性的可除理论。

\noindent\textbf{定义.} 若 $A$ 和 $B$ 是整环 $R$ 中的理想，则称 $B$ \emph{整除} $A$（并称 $A$ 可被 $B$ 整除），如果存在 $R$ 中理想 $C$ 使得 $A=BC$。

若 $B$ 整除 $A$，则当然 $A\subseteq B$。若 $R$ 是 Dedekind 整环，反过来也成立：$A\subseteq B$ 蕴含 $C=AB^{-1}\subseteq BB^{-1}=R$，所以 $C$ 是 $R$ 中理想且 $BC=A$。

我们还可以定义两个理想 $A$ 和 $B$ 的最大公因子 $(A,B)$ 的概念：$(A,B)$ 整除 $A$ 和 $B$，并且任何同时整除 $A$ 和 $B$ 的理想都整除 $(A,B)$。下一个命题中的第二个断言说明，Dedekind 整环中的整理想总存在这种最大公因子，并给出类似于两个整数最大公因子公式的公式。

\medskip
\noindent\textbf{命题 17.} 设 $R$ 是 Dedekind 整环，$A,B$ 是 $R$ 中两个非零理想，其素理想分解为
\[
  A=P_1^{e_1}\cdots P_n^{e_n},
  \qquad
  B=P_1^{f_1}\cdots P_n^{f_n},
\]
其中 $e_i,f_i\geq0$，$i=1,
\ldots,n$。则：
\begin{enumerate}[label=(\arabic*)]
  \item $A\subseteq B$ 当且仅当 $B$ 整除 $A$（即“包含就是整除”），当且仅当对 $i=1,
  \ldots,n$ 有 $f_i\leq e_i$；
  \item
  \[
    A+B=(A,B)=P_1^{\min(e_1,f_1)}\cdots P_n^{\min(e_n,f_n)},
  \]
  因而特别地，$A$ 与 $B$ 互素，即 $A+B=R$，当且仅当它们没有共同的素理想因子。
\end{enumerate}

\begin{proof}
我们已在上面证明了 (1) 中的第一个断言。若每个 $f_i\leq e_i$，则取
\[
  C=P_1^{e_1-f_1}\cdots P_n^{e_n-f_n}\subseteq R
\]
说明 $B$ 整除 $A$。反过来，若 $B$ 整除 $A$，则把 $C$ 写成素理想乘积并代入 $A=BC$，得到对所有 $i$ 都有 $f_i\leq e_i$，这证明了整个 (1)。由于 $A+B$ 是同时包含 $A$ 与 $B$ 的最小理想，(2) 现在由 (1) 推出。
\end{proof}

\medskip
\noindent\textbf{命题 18（中国剩余定理）.} 设 $R$ 是 Dedekind 整环，$P_1,P_2,
\ldots,P_n$ 是 $R$ 中不同的素理想，且 $a_i\geq0$ 是整数，$i=1,
\ldots,n$。则
\[
  R/P_1^{a_1}\cdots P_n^{a_n}
  \cong R/P_1^{a_1}\times R/P_2^{a_2}\times\cdots\times R/P_n^{a_n}.
\]
等价地，对任意元素 $r_1,r_2,
\ldots,r_n\in R$，存在元素 $r\in R$，唯一地模去 $P_1^{a_1}\cdots P_n^{a_n}$ 中的一个元素，使得
\[
  r\equiv r_1\pmod{P_1^{a_1}},
  \qquad
  r\equiv r_2\pmod{P_2^{a_2}},
  \qquad
  \ldots,
  \qquad
  r\equiv r_n\pmod{P_n^{a_n}}.
\]

\begin{proof}
这由第 7.6 节定理 17 立刻得到，因为前一个命题说明理想 $P_i^{a_i}$ 两两互素。
\end{proof}

\medskip
\noindent\textbf{推论 19.} 设 $I$ 是 Dedekind 整环 $R$ 中的理想。则：
\begin{enumerate}[label=(\arabic*)]
  \item 存在 $R$ 中理想 $J$，它与 $I$ 互素，并且乘积 $IJ=(a)$ 是主理想；
  \item 若 $I$ 非零，则商 $R/I$ 中每个理想都是主理想；等价地，若 $I_1$ 是 $R$ 中包含 $I$ 的理想，则对某个 $b\in R$ 有 $I_1=I+Rb$；
  \item $R$ 中每个理想都可由两个元素生成；事实上，若 $I$ 非零且 $0\neq a\in I$，则对某个 $b\in I$ 有 $I=Ra+Rb$。
\end{enumerate}

\begin{proof}
设
\[
  I=P_1^{e_1}\cdots P_n^{e_n}
\]
是 $I$ 在 $R$ 中的素理想分解。对每个 $i=1,
\ldots,n$，令 $r_i$ 是 $P_i^{e_i}-P_i^{e_i+1}$ 中的一个元素。由前一命题，存在 $a\in R$，使得对所有 $i$ 都有
\[
  a\equiv r_i\pmod{P_i^{e_i+1}}.
\]
因此对所有 $i$，$a\in P_i^{e_i}-P_i^{e_i+1}$，所以由命题 17(1)，$(a)$ 的素理想分解中 $P_i$ 的幂次恰好为 $e_i$：
\[
  (a)=P_1^{e_1}\cdots P_n^{e_n}P_{n+1}^{e_{n+1}}\cdots P_m^{e_m}
\]
对某些与 $P_1,
\ldots,P_n$ 不同的素理想 $P_{n+1},
\ldots,P_m$ 成立。令
\[
  J=P_{n+1}^{e_{n+1}}\cdots P_m^{e_m}
\]
得到 (1)。对 (2)，由中国剩余定理，只需证明在素理想 $P$ 的幂情形下，$R/P^m$ 中每个理想都是主理想；而这立即成立，因为
\[
  R/P^m\cong R_P/P^mR_P
\]
且局部化 $R_P$ 是 P.I.D.。最后，(3) 由 (2) 取 $I=Ra$ 得到。
\end{proof}

推论 19 的第一个断言说明，存在与 $I$ 互素的整理想 $J$，它在 $R$ 的类群中位于 $I$ 的逆类中。甚至可以对 $J$ 施加额外条件，参见练习 11。

\medskip
\noindent\textbf{推论 20.} 若 $R$ 是 Dedekind 整环，则 $R$ 是 P.I.D.（即 $R$ 的类数为 $1$），当且仅当 $R$ 是 U.F.D.。

\begin{proof}
每个 P.I.D. 都是 U.F.D.。反过来，假设 $R$ 是 U.F.D.，令 $P$ 为 $R$ 中任意素理想。由推论 19，$P=Ra+Rb$ 对某些 $0\neq a\in R$ 和 $b\in R$ 成立。由于 $a$ 的不可约因子 $a'$ 中有一个满足 $(a')\subseteq P$（因为它们的乘积属于素理想 $P$），于是由命题 17(1)，$P$ 在 $R$ 中整除 $(a')$。因此 $P=(a')$ 是主理想，因为 $(a')$ 是素理想（第 8.3 节命题 12）。由于 $R$ 中每个理想都是素理想乘积，$R$ 的每个理想都是主理想，即 $R$ 是 P.I.D.。
\end{proof}

推论 20 说明，Dedekind 整环 $R$ 的类数衡量了元素唯一分解失败的程度。代数数论中的一个基本结果是，代数数域的整数环的类数是有限的。然而，对一般 Dedekind 整环，类数不必有限。事实上，对任意 Abel 群 $A$（有限或无限），存在一个 Dedekind 整环，其类群同构于 $A$。

\subsection*{Dedekind 整环上的模与有限生成模基本定理}

接下来转向 Dedekind 整环 $R$ 上模的研究。$R$ 的每个分式理想都是 $R$-模，而下面命题的第一个断言说明，$R$ 的两个分式理想作为 $R$-模同构，当且仅当它们代表 $R$ 的类群中的同一个元素。

\medskip
\noindent\textbf{命题 21.} 设 $R$ 是分式域为 $K$ 的 Dedekind 整环。
\begin{enumerate}[label=(\arabic*)]
  \item 假设 $I$ 与 $J$ 是 $R$ 的两个分式理想。则 $I\cong J$ 作为 $R$-模，当且仅当 $I$ 与 $J$ 相差一个非零主理想：对某个 $0\neq a\in K$ 有 $I=(a)J$。
  \item 更一般地，假设 $I_1,I_2,
  \ldots,I_n$ 与 $J_1,J_2,
  \ldots,J_m$ 是 Dedekind 整环 $R$ 的分式域 $K$ 中的非零分式理想。则
  \[
    I_1\oplus I_2\oplus\cdots\oplus I_n
    \cong
    J_1\oplus J_2\oplus\cdots\oplus J_m
  \]
  作为 $R$-模，当且仅当 $n=m$ 且乘积理想 $I_1I_2\cdots I_n$ 与 $J_1J_2\cdots J_n$ 相差一个主理想：对某个 $0\neq a\in K$，
  \[
    I_1I_2\cdots I_n=(a)J_1J_2\cdots J_n.
  \]
  \item 特别地，
  \[
    I_1\oplus I_2\oplus\cdots\oplus I_n
    \cong
    \underbrace{R\oplus\cdots\oplus R}_{n-1\text{ 个因子}}
    \oplus (I_1I_2\cdots I_n),
  \]
  并且 $R^n\oplus I\cong R^n\oplus J$ 当且仅当 $I$ 与 $J$ 相差一个主理想：$I=(a)J$，$a\in K$。
\end{enumerate}

\begin{proof}
由 $0\neq a\in K$ 乘法给出从 $J$ 到 $(a)J$ 的 $R$-模同构，所以若 $I=(a)J$，则 $I\cong J$ 作为 $R$-模。反过来，注意我们可以假设 $J\neq0$，于是 $I\cong J$ 蕴含 $R\cong J^{-1}I$。这意味着 $J^{-1}I=aR$ 是主理想（生成元 $a$ 由 $1\in R$ 的像给出），即 $I=(a)J$，证明了 (1)。

接下来证明，对任意非零分式理想 $I$ 与 $J$，有
\[
  I\oplus J\cong R\oplus IJ.
\]
必要时把 $I$ 和 $J$ 分别换成同构的 $R$-模 $aI$ 和 $bJ$，我们可以假设 $I$ 与 $J$ 是互素的整理想（参见练习 12），所以 $I+J=R$ 且 $I\cap J=IJ$。容易看出，从 $I\oplus J$ 到 $I+J=R$ 的映射 $(x,y)\mapsto x+y$ 是满的 $R$-模同态，其核为 $I\cap J=IJ$，因此有 $R$-模正合列
\[
  0\to IJ\to I\oplus J\to R\to0.
\]
这个序列分裂，因为 $R$ 是自由模，所以 $I\oplus J\cong R\oplus IJ$，如所断言。

(3) 的第一个断言现在由归纳得到；把该断言与 (1) 结合可知，若对某个非零 $a\in K$ 有 $I_1\cdots I_n=(a)J_1\cdots J_n$，则 $I_1\oplus\cdots\oplus I_n$ 同构于 $J_1\oplus\cdots\oplus J_n$。这证明了 (2) 中的“若”方向。还需要证明 (2) 中的“仅若”方向，因为 (3) 中相应断言是其特殊情形。因此假设
\[
  I_1\oplus I_2\oplus\cdots\oplus I_n
  \cong
  J_1\oplus J_2\oplus\cdots\oplus J_m
\]
作为 $R$-模。

由于 $I\otimes_R K$ 是理想 $I$ 在 $K$ 中的局部化（参见第 15.4 节命题 41），对任意 $K$ 的非零分式理想 $I$ 都有 $I\otimes_R K\cong K$。由于张量积与直和可交换，$(I_1\oplus\cdots\oplus I_n)\otimes_R K\cong K^n$ 是 $K$ 上 $n$ 维向量空间。类似地，$J_1\oplus\cdots\oplus J_m\otimes_R K\cong K^m$，由此推出 $n=m$。

注意，对任意非零元素 $a_1\in I_1$，把 $I_1$ 换成同构的分式理想 $a_1^{-1}I_1$，不会影响 (2) 中断言的成立。因此我们可以假设 $I_1$ 包含 $R$，类似地也可以假设 (2) 中每个分式理想都包含 $R$。令 $\varphi$ 表示从 $I_1\oplus\cdots\oplus I_n$ 到 $J_1\oplus\cdots\oplus J_n$ 的 $R$-模同构。对 $i=1,2,
\ldots,n$，定义
\[
  \varphi((0,
  \ldots,0,1,0,
  \ldots,0))=(a_{1,i},a_{2,i},
  \ldots,a_{n,i})\in J_1\oplus J_2\oplus\cdots\oplus J_n,
\]
其中左边的 $1\in I_i$ 出现在第 $i$ 个位置。由于 $\varphi$ 是 $R$-模同态，得到对每个 $j=1,2,
\ldots,n$，
\[
  J_j=a_{j,1}I_1+a_{j,2}I_2+\cdots+a_{j,i}I_i+\cdots+a_{j,n}I_n.
\]
对 $j=1,2,
\ldots,n$ 取这些理想的乘积，得到对 $\{j_1,j_2,
\ldots,j_n\}$ 的任一排列，
\[
  (a_{j_1,1}a_{j_2,2}\cdots a_{j_n,n})I_1I_2\cdots I_n
  \subseteq
  J_1J_2\cdots J_n.
\]
因此
\[
  dI_1I_2\cdots I_n\subseteq J_1J_2\cdots J_n,
\]
其中 $d$ 是矩阵 $(a_{i,j})$ 的行列式，因为该行列式是若干项
\[
  \epsilon(\sigma)a_{1,\sigma(1)}\cdots a_{n,\sigma(n)}
\]
的和，$\epsilon(\sigma)$ 是 $\{1,2,
\ldots,n\}$ 的置换 $\sigma$ 的符号。类似地，对 $j=1,
\ldots,n$，定义
\[
  \varphi^{-1}((0,
  \ldots,0,1,0,
  \ldots,0))=(b_{1,j},b_{2,j},
  \ldots,b_{n,j})\in I_1\oplus I_2\oplus\cdots\oplus I_n,
\]
其中左边的 $1\in J_j$ 出现在第 $j$ 个位置。两个矩阵 $(a_{i,j})$ 和 $(b_{i,j})$ 的乘积就是恒等矩阵，所以 $d\neq0$，且矩阵 $(b_{i,j})$ 的行列式为 $d^{-1}$。如上同理可得
\[
  d^{-1}J_1J_2\cdots J_n\subseteq I_1I_2\cdots I_n,
\]
这说明
\[
  I_1I_2\cdots I_n=(a)J_1J_2\cdots J_n,
\]
其中 $0\neq a=d^{-1}\in K$，从而完成命题的证明。
\end{proof}

现在我们考虑 Dedekind 整环上的有限生成模，并证明这类模的一个结构定理，它推广了第 12 章中关于 P.I.D. 上有限生成模的结果。

回忆，$M$ 的秩是 $M$ 中 $R$-线性无关元素的最大个数；等价地，是 $M\otimes_R K$ 作为 $K$-向量空间的维数，其中 $K$ 是 $R$ 的分式域（参见第 12.1 节练习 1--4、20）。

\medskip
\noindent\textbf{定理 22.} 设 $M$ 是 Dedekind 整环 $R$ 上的有限生成模。令 $n\geq0$ 表示 $M$ 的秩，并令 $\Tor(M)$ 为 $M$ 的挠子模。则
\[
  M\cong \underbrace{R\oplus R\oplus\cdots\oplus R}_{n-1\text{ 个因子}}\oplus I\oplus\Tor(M)
\]
对 $R$ 的某个理想 $I$ 成立，并且
\[
  \Tor(M)\cong R/P_1^{e_1}\times R/P_2^{e_2}\times\cdots\times R/P_s^{e_s}
\]
对某个 $s\geq0$ 以及（不一定互异的）素理想幂 $P_i^{e_i}$（其中 $e_i\geq1$）成立。理想 $P_i^{e_i}$，$i=1,
\ldots,s$，是唯一的；理想 $I$ 在乘以主理想意义下唯一。

\begin{proof}
先假设 $M$ 是 $R$ 上有限生成无挠模，即 $\Tor(M)=0$。则从 $M$ 到 $M\otimes_R K$ 的自然 $R$-模同态是单射，所以我们可以把 $M$ 看作向量空间 $M\otimes_R K$ 的一个 $R$-子模。若 $M$ 在 $R$ 上秩为 $n$，则 $M\otimes_R K$ 是 $K$ 上 $n$ 维向量空间。令 $x_1,
\ldots,x_n$ 为 $M\otimes_R K$ 在 $K$ 上的一组基，令 $m_1,
\ldots,m_s$ 为 $M$ 的 $R$-模生成元。每个 $m_i$，$i=1,
\ldots,s$，都可写成 $x_1,
\ldots,x_n$ 的 $K$-线性组合。令 $0\neq d\in R$ 为所有这些线性组合中 $K$ 系数的公共分母，并置 $y_i=x_i/d$，$i=1,
\ldots,n$。于是
\[
  M\subseteq Ry_1+\cdots+Ry_n\subset Kx_1+\cdots+Kx_n,
\]
这说明 $M$ 包含在一个秩为 $n$ 的自由 $R$-子模中，并且 $M$ 中每个元素 $m$ 都可唯一写成
\[
  m=a_1y_1+\cdots+a_ny_n
\]
的形式，其中 $a_1,
\ldots,a_n\in R$。定义映射 $\varphi:M\to R$ 为
\[
  \varphi(a_1y_1+\cdots+a_ny_n)=a_n,
\]
这是 $R$-模同态，所以我们有正合列
\[
  0\to\ker\varphi\to M\xrightarrow{\varphi} I_1\to0,
\]
其中 $I_1$ 是 $\varphi$ 在 $R$ 中的像，因此是 $R$ 中的理想。子模 $\ker\varphi$ 也是无挠 $R$-模，其秩至多为 $n-1$（因为它包含在 $Ry_1+\cdots+Ry_{n-1}$ 中）；通过比较秩可知 $I_1$ 非零，并且 $\ker\varphi$ 的秩恰为 $n-1$。由定理 15(4)，$I_1$ 是投射 $R$-模，所以该序列分裂：
\[
  M\cong I_1\oplus(\ker\varphi).
\]
对秩作归纳可知，有限生成无挠 $R$-模同构于 $R$ 的 $n$ 个非零理想的直和：
\[
  M\cong I_1\oplus I_2\oplus\cdots\oplus I_n.
\]
由于 $I_1,
\ldots,I_n$ 每个都是投射 $R$-模，推出任意有限生成无挠 $R$-模都是投射的。

现在若 $M$ 是任意有限生成 $R$-模，则商 $M/\Tor(M)$ 是有限生成且无挠的，因此由刚刚证明的结果是投射的。正合列
\[
  0\to\Tor(M)\to M\to M/\Tor(M)\to0
\]
因此分裂，所以
\[
  M\cong\Tor(M)\oplus (M/\Tor(M)).
\]
由前一段的结果，$M/\Tor(M)$ 同构于 $R$ 的 $n$ 个非零理想的直和；再由命题 21 得到
\[
  M\cong \underbrace{R\oplus R\oplus\cdots\oplus R}_{n-1\text{ 个因子}}\oplus I\oplus\Tor(M)
\]
对某个理想 $I$ 成立。关于理想 $I$ 的唯一性断言也由命题 21(3) 中的唯一性断言立刻得到。

还需证明关于挠子模 $\Tor(M)$ 的断言。设 $N$ 是有限生成挠 $R$-模。令 $I=\Ann(N)$ 为 $N$ 在 $R$ 中的零化子，并设
\[
  I=P_1^{e_1}\cdots P_t^{e_t}
\]
是 $I$ 在 $R$ 中的素理想分解，其中 $P_1,
\ldots,P_t$ 是不同的素理想。则 $N$ 是 $R/I$ 上的模，并且
\[
  R/I\cong R/P_1^{e_1}\times R/P_2^{e_2}\times\cdots\times R/P_t^{e_t}.
\]
于是
\[
  N\cong (N/P_1^{e_1}N)\times (N/P_2^{e_2}N)\times\cdots\times (N/P_t^{e_t}N)
\]
作为 $R$-模。每个 $N/P^eN$ 是
\[
  R/P^e\cong R_P/P^eR_P
\]
上的有限生成模，其中 $R_P$ 是 $R$ 在素理想 $P$ 处的局部化；也就是说，它是 $R_P$ 上被 $P^eR_P$ 零化的有限生成模。由于 $R$ 是 Dedekind 整环，每个 $R_P$ 都是 P.I.D.（甚至是 D.V.R.），所以可以应用 P.I.D. 上有限生成模基本定理，得到每个 $N/P^eN$ 作为 $R_P$-模同构于有限多个形如 $R_P/P^fR_P$ 的模的直和，其中 $f\leq e$。于是每个 $N/P^eN$ 作为 $R$-模同构于有限多个形如 $R/P^fR$ 的模的直和，其中 $f\leq e$。这证明 $N$ 同构于有限多个形如 $R/P_i^{e_i}$ 的模的直和，其中 $P_i$ 为各种素理想。因此 $\Tor(M)$ 可分解为定理中所述的直和。

最后，还需证明 $\Tor(M)$ 分解中的理想 $P_i^{e_i}$，$i=1,
\ldots,s$，是唯一的。这类似于第 12.1 节定理 10 证明中的唯一性论证（也参见第 12.1 节练习 11--12）：对 $R$ 的任意素理想 $P$，商
\[
  P^{i-1}M/P^iM
\]
是域 $R/P$ 上的向量空间，并且差
\[
  \dim_{R/P} P^{i-1}M/P^iM-
  \dim_{R/P} P^iM/P^{i+1}M
\]
就是 $M$ 中同构于 $R/P^i$ 的直和项个数，因此由 $M$ 唯一决定。这完成了定理的证明。
\end{proof}

若 $M$ 是如定理 22 中的 Dedekind 整环 $R$ 上有限生成模，则 $M$ 作为 $R$-模的同构类型由秩 $n$、素理想幂 $P_i^{e_i}$（$i=1,
\ldots,s$，称为 $M$ 的初等因子）以及理想 $I$ 在 $R$ 的类群中的类（称为 $M$ 的 Steinitz 类）决定。注意，P.I.D. 就是类数为 $1$ 的 Dedekind 整环；在这种情形下，$R$ 的每个非零理想 $I$ 作为 $R$-模都同构于 $R$。于是定理 22 化为第 12 章中 P.I.D. 上有限生成模结构定理的初等因子形式。定理 22 中挠 $R$-模的描述也有不变因子形式（参见练习 14）。

下一个结果把 P.I.D. 上有限生成投射模的刻画（第 12.1 节练习 21）推广到 Dedekind 整环。

\medskip
\noindent\textbf{推论 23.} Dedekind 整环上的有限生成模是投射的，当且仅当它是无挠的。

\begin{proof}
我们在定理 22 的证明中已经证明，有限生成无挠 $R$-模是投射的；因此由定理 22 中 $M$ 的分解，$M$ 是投射的，当且仅当 $\Tor(M)$ 是投射的（参见第 10.5 节练习 3）。为完成证明，只需说明没有非零挠 $R$-模是投射的，这留作练习（参见练习 15）。
\end{proof}

\subsection*{练习}

\begin{enumerate}[label=\arabic*.]
  \item 若 $R$ 是整环，证明 $R$ 的每个分式理想可逆，当且仅当 $R$ 的每个整理想可逆。
  \item 设 $R$ 是分式域为 $K$ 的整环，且 $A_1,A_2,
  \ldots,A_n$ 是 $R$ 的分式理想，其乘积是非零主分式理想：对某个 $0\neq x\in K$，有 $A_1A_2\cdots A_n=Rx$。证明对每个 $i=1,
  \ldots,n$，$A_i$ 是可逆分式理想，其逆为
  \[
    (x^{-1})A_1\cdots A_{i-1}A_{i+1}\cdots A_n.
  \]
  \item 设 $R$ 是分式域为 $K$ 的整环，且 $P$ 是 $R$ 中非零素理想。证明 $R_P$ 在 $K$ 中的分式理想正是形如 $AR_P$ 的 $R_P$-模，其中 $A$ 是 $R$ 的一个分式理想。
  \item 设 $R$ 是分式域为 $K$ 的整环，且 $A$ 是 $K$ 中 $R$ 的分式理想。令 $A'=\{x\in K\mid xA\subseteq R\}$ 如命题 9 中所定义。
  \begin{enumerate}[label=(\alph*)]
    \item 对 $R$ 中任意素理想 $P$，证明 $A'$ 在 $P$ 处的局部化 $(A')_P$ 是 $K$ 中 $R_P$ 的分式理想。
    \item 若 $A$ 是有限生成 $R$-模，证明 $(A')_P=(A_P)'$，其中 $(A_P)'$ 是对应于局部化 $A_P$ 的分式 $R_P$-理想 $\{x\in K\mid xA_P\subseteq R_P\}$。
  \end{enumerate}
  \item 若 $Q_1$ 是 $P_1$-primary 理想，$Q_2$ 是 $P_2$-primary 理想，其中 $P_1$ 与 $P_2$ 是 Noether 环 $R$ 中互素的理想，证明 $Q_1$ 与 $Q_2$ 也互素。［用第 15.2 节命题 14。］
  \item 设 $R$ 是分式域为 $K$ 的 Dedekind 整环。
  \begin{enumerate}[label=(\alph*)]
    \item 证明 $K$ 中 $R$ 的每个非零分式理想都可唯一写成不同素理想幂的乘积 $P_1^{a_1}\cdots P_n^{a_n}$，其中 $a_i$ 是非零整数，也可能为负。
    \item 若 $0\neq x\in K$，令 $P^{v_P(x)}$ 为主理想 $(x)$ 在 (a) 中分解时素理想 $P$ 的幂次（若 $P$ 不是出现的素理想之一，则 $v_P(x)=0$）。证明 $v_P$ 是 $K$ 上的一个赋值，其赋值环为 $R_P$，即 $R$ 在 $P$ 处的局部化。
  \end{enumerate}
  \item 设 $R$ 是不是域的 Noether 整环。证明 $R$ 是 Dedekind 整环，当且仅当对 $R$ 的每个极大理想 $M$，不存在 $R$ 的理想 $I$ 使得
  \[
    M^2\subset I\subset M.
  \]
  ［用第 1 节练习 12 以及定理 7 和定理 15。］
  \item 设 $R$ 是 Krull 维数为 $1$ 的 Noether 整环。证明 $R$ 的每个非零理想 $I$ 都可唯一写成 primary 理想的乘积，且这些 primary 理想的根互不相同。［参见定理 15 的证明。使用极小 primary 分解中属于孤立素理想的 primary 分量的唯一性（第 15.2 节定理 21）。］
  \item 假设 $R$ 是整环。证明对每个非零素理想 $P$，$R_P$ 都是 D.V.R.，当且仅当对每个非零极大理想 $M$，$R_M$ 都是 D.V.R.。
  \item 设 $R$ 是不是域的 Noether 整环。证明 $R$ 是 Dedekind 整环，当且仅当非零素理想 $M$ 都是极大理想，并且每个 $M$-primary 理想都是 $M$ 的幂。
  \item 若 $I$ 与 $J$ 是 Dedekind 整环 $R$ 中的非零理想，证明存在 $R$ 中的整理想 $I_1$，它同时与 $I$ 和 $J$ 互素，并且 $I_1I$ 是 $R$ 中的主理想。
  \item 若 $I$ 与 $J$ 是 Dedekind 整环 $R$ 的非零分式理想，证明存在元素 $\alpha,\beta\in K$，使得 $\alpha I$ 和 $\beta J$ 是 $R$ 中互素的非零整理想。
  \item 设 $I$ 与 $J$ 是 Dedekind 整环 $R$ 中的非零理想。证明存在理想 $I_1\subseteq I$，且 $I_1$ 与 $J$ 互素。［用推论 19 找到理想 $I_2$，使得 $I_2I=(a)$ 且 $(I_2,J)=R$。若 $I_2=P_1^{e_1}\cdots P_n^{e_n}$，选取 $b\in R$，使得 $b\in P_i^{e_i}-P_i^{e_i+1}$，并且对每个整除 $J$ 的素理想 $P$，有 $b\equiv1\pmod P$。证明 $(b)=I_2I_1$ 对某个理想 $I_1$ 成立，并考虑 $(a)I_1$ 来证明 $I_1\subseteq I$。］
  \item 证明 Dedekind 整环 $R$ 上每个有限生成挠模都同构于直和
  \[
    R/I_1\oplus R/I_2\oplus\cdots\oplus R/I_n,
  \]
  其中 $I_1,
  \ldots,I_n$ 是 $R$ 的唯一的非零理想，满足
  \[
    I_1\supseteq I_2\supseteq\cdots\supseteq I_n
  \]
  （称为 $M$ 的不变因子）。［参见第 12.1 节。］
  \item 若 $P$ 是 Dedekind 整环 $R$ 中的非零素理想，证明对任意 $n\geq1$，$R/P^n$ 都不是投射 $R$-模。［考虑正合列
  \[
    0\to P^n/P^{n+1}\to R/P^{n+1}\to R/P^n\to0.
  \]
  ］推出若 $M\neq0$ 是有限生成挠 $R$-模，则 $M$ 不是投射的。［参见第 10.5 节练习 3。］
  \item 证明 Dedekind 整环 $R$ 的类数为 $1$，当且仅当每个有限生成投射 $R$-模都是自由的。
  \item 设 $R$ 是 Dedekind 整环。
  \begin{enumerate}[label=(\alph*)]
    \item 证明 $I\sim J$ 当且仅当 $I\cong J$ 作为 $R$-模，在 $R$ 的非零分式理想集合上定义了一个等价关系。令 $C(R)$ 为对应的 $R$-模同构类集合，并令 $[I]\in C(R)$ 表示包含分式理想 $I$ 的等价类。
    \item 证明乘法 $[I][J]=[I\oplus J]$ 给出一个良定的二元运算，使得 $C(R)$ 是 Abel 群，恒等元为 $1=[R]$。
    \item 证明 (b) 中的 Abel 群 $C(R)$ 同构于 $R$ 的类群。
  \end{enumerate}
  \item 若 $R$ 是 Dedekind 整环且 $I$ 是任意非零理想，证明 $R/I$ 只含有限多个理想。特别地，证明 $R/I$ 是 Artin 环。
  \item 假设 $I$ 是 Dedekind 整环 $R$ 中的非零分式理想。显式地把 $I$ 表示为自由 $R$-模的直和项，从而说明 $I$ 是投射的。［考虑 $I\oplus I^{-1}$ 并使用命题 21。］
  \item 设 $I$ 与 $J$ 是 Dedekind 整环 $R$ 中的两个非零分式理想，并且对某个 $n\neq0$ 有 $I^n=J^n$。证明 $I=J$。
  \item 设 $K$ 是代数数域，$\mathcal O_K$ 是 $K$ 中的整数环。若 $P$ 是 $\mathcal O_K$ 中非零素理想，证明对某个素数 $p\in\ZZ$ 与代数整数 $\pi\in\mathcal O_K$，有
  \[
    P=(p,\pi).
  \]
  \item 设 $K=\QQ(\sqrt D)$ 是 $\QQ$ 的二次扩张，其中 $D$ 是无平方因子整数，且 $\mathcal O_K$ 是 $K$ 中的整数环。
  \begin{enumerate}[label=(\alph*)]
    \item 证明 $|\mathcal O_K/(p)|=p^2$。［观察 $\mathcal O_K\cong\ZZ^2$ 作为 Abel 群。］
    \item 用推论 16 说明，$(p)$ 在 $\mathcal O_K$ 中的素理想分解有三种可能：
    \begin{enumerate}[label=(\roman*)]
      \item $(p)=P$ 是素理想，且 $|\mathcal O_K/P|=p^2$；
      \item $(p)=P_1P_2$，其中 $P_1,P_2$ 是不同素理想，且 $|\mathcal O_K/P_1|=|\mathcal O_K/P_2|=p$；
      \item $(p)=P^2$，其中 $P$ 是某个素理想，且 $|\mathcal O_K/P|=p$。
    \end{enumerate}
    （在情形 (i)、(ii)、(iii) 中，素数 $p$ 分别称为在 $\mathcal O_K$ 中惰性、分裂或分歧。分歧素数集合是有限的：若 $D\equiv1,2\pmod4$，则为整除 $D$ 的素数；若 $D\equiv3\pmod4$，则为 $p=2$ 以及整除 $D$ 的素数。参见第 15.5 节练习 31。）
    \item 在 $K=\QQ(\sqrt{-5})$ 的整数环 $\mathcal O_K=\ZZ[\sqrt{-5}]$ 中，确定素数 $p=2,3,5,7,11$ 的素理想分解。
  \end{enumerate}
  \item 令 $\mathcal O$ 为 $\QQ$ 的代数闭包 $\overline\QQ$ 中的整数环。
  \begin{enumerate}[label=(\alph*)]
    \item 证明 $\mathcal O$ 中无限理想列
    \[
      (2)\subset (\sqrt2)\subset (\sqrt[4]2)\subset (\sqrt[8]2)\subset\cdots
    \]
    是严格递增的，因此 $\mathcal O$ 不是 Noether 的。
    \item 证明 $\mathcal O$ 的 Krull 维数为 $1$。［用第 15.3 节定理 26。］
    \item 设 $K$ 是数域，$I$ 是 $\mathcal O_K$ 中任意理想。证明存在 $K$ 的某个有限扩张 $L$，使得 $I$ 延拓到 $\mathcal O_L$ 后成为主理想，即理想 $I\mathcal O_L$ 是主理想（其中 $L$ 依赖于 $I$）——你可以使用 $K$ 的类群是有限群这一事实。［参见练习 20。］
    \item 证明 $\mathcal O$ 是 Bezout 整环（参见第 8.1 节）。
  \end{enumerate}
  \item 设 $F$ 和 $K$ 是代数数域，且 $\QQ\subseteq F\subseteq K$，其整数环分别为 $\mathcal O_F$ 与 $\mathcal O_K$。由于 $\mathcal O_F\subseteq\mathcal O_K$，环 $\mathcal O_K$ 自然是 $\mathcal O_F$ 上的模。
  \begin{enumerate}[label=(\alph*)]
    \item 证明 $\mathcal O_K$ 是秩为 $n=[K:F]$ 的无挠 $\mathcal O_F$-模。［计算在 $\ZZ$ 上的秩。］若 $\mathcal O_K$ 在 $\mathcal O_F$ 上自由，则称 $\mathcal O_K$ 在 $\mathcal O_F$ 上有\emph{相对整基}。
    \item 证明若 $F$ 的类数为 $1$，则 $\mathcal O_K$ 在 $\mathcal O_F$ 上有相对整基。
  \end{enumerate}
  若 $K=\QQ(\sqrt{-5},\sqrt2)$，则 $K$ 的整数环 $\mathcal O_K$ 为
  \[
    \mathcal O_K=\ZZ+\ZZ\sqrt2+\ZZ\sqrt{-10}+\ZZ\omega,
    \qquad
    \omega=(\sqrt{-10}+\sqrt2)/2.
  \]
  \begin{enumerate}[label=(\alph*),resume]
    \item 若 $F_1=\QQ(\sqrt2)$，证明 $\mathcal O_K$ 在 $\mathcal O_{F_1}$ 上有相对整基，并找出一个显式基 $\{\alpha,\beta\}$：
    \[
      \mathcal O_K=\mathcal O_{F_1}\cdot\alpha+\mathcal O_{F_1}\cdot\beta.
    \]
    \item 若 $F_2=\QQ(\sqrt{-5})$，证明
    \[
      P_3=(3,1+\sqrt{-5})=(3,5-\sqrt{-5})
    \]
    是 $\mathcal O_{F_2}$ 中的非主素理想，并且
    \[
      \mathcal O_K=\mathcal O_{F_2}\cdot1+(1/3)P_3\cdot\omega.
    \]
    ［检查 $\sqrt{-10}=-(5-\sqrt{-5})\omega/3$。］推出 $\mathcal O_K$ 作为 $\mathcal O_{F_2}$-模的 Steinitz 类是 $\mathcal O_{F_2}$ 的类群中 $P_3$ 的非平凡类，因此 $\mathcal O_K$ 在 $\mathcal O_{F_2}$ 上没有相对整基。
    \item 判断 $\mathcal O_K$ 在 $K$ 的剩余二次子域 $F_3=\QQ(\sqrt{-10})$ 的整数环上是否有相对整基。
  \end{enumerate}
  \item 设 $C$ 是代数闭域 $k$ 上的非奇异不可约仿射曲线。证明坐标环 $k[C]$ 是 Dedekind 整环。
\end{enumerate}

% % ============================================
%  chapter17.tex — 第17章：同调代数与群上同调导论
%  根据 790-852.pdf 翻译整理，保持原文结构、证明与练习
% ============================================

\chapter{同调代数与群上同调导论}

设 $R$ 是带有 $1$ 的环。在第 10.5 节中我们看到，一个 $R$-模短正合列
\[
  0\to L\xrightarrow{\varphi} M\xrightarrow{\psi} N\to 0
  \tag{17.1}
\]
给出了任意 $R$-模 $D$ 上的一个 Abel 群正合列
\[
  0\to \Hom_R(N,D)\xrightarrow{\psi'}\Hom_R(M,D)\xrightarrow{\varphi'}\Hom_R(L,D),
  \tag{17.2}
\]
而同态 $\varphi'$ 一般不是满射，所以这个序列并不总能延伸为一个短正合列。等价地，从 $L$ 到 $D$ 的同态一般不能提升为从 $M$ 到 $D$ 的同态。在本章中，我们引入若干“同调代数”的技巧，它们给出了以自然方式延长某些正合列的方法。对于上面的情形，可以得到一个涉及“上同调群” $\Ext_R^n(\_,D)$ 的无限正合列（参见定理 8），而这些群度量从 $L$ 到 $D$ 的同态中不能延拓到 $M$ 的那部分。然后我们考虑第 10.5 节中出现的另外两个函子的类似问题，即从 $D$ 到序列 (17.1) 各项取同态，以及将序列 (17.1) 与 $D$ 张量。

在随后的各节中，我们集中讨论这类同调构造的一个重要特例——“有限群的上同调”。我们在这个情形中给出显式计算，并指出这些技巧在建立群论中的若干新结果时的应用。从这个意义上说，第 2--4 节可以看作说明第 1 节中一般理论用途的一个显式“例子”。

上同调群和同调群出现在数学的许多领域中。同调群、上同调群以及同调代数的一般形式概念，大约在 20 世纪中叶从代数拓扑中产生，源于研究拓扑空间的高阶同伦群与基本群之间的关系；不过某些具体上同调群的研究，例如 Schur 关于群扩张的工作（第 4 节中会描述），要早半个世纪。正如代数中常见的情形一样，许多不同领域中共有的思想被抽象为一般理论。同调与上同调的大量语言反映了它的拓扑起源：同调群、链、循环、边界等。

\section{同调代数导论——Ext 与 Tor}

在本节中，我们描述同调代数中的若干一般术语与结果，最终得到所谓的上同调长正合列。随后我们定义与序列 (17.2) 相关的若干（上同调）群，并把一般同调结果应用于得到一个在右端延长该序列的长正合列。然后我们指出，对由从 $D$ 到 (17.1) 中各项取同态所得到的序列，或由与 $D$ 作张量所得到的序列，也有相应的发展。

我们从正合列概念的一个推广开始，即一列 Abel 群同态，使得相邻两个映射的复合为零（也就是说，一个映射的像包含在下一个映射的核中）：

\noindent\textbf{定义.} 设 $C$ 是 Abel 群同态序列
\[
  0\to C^0\xrightarrow{d_1} C^1\to\cdots\to C^{n-1}\xrightarrow{d_n} C^n\xrightarrow{d_{n+1}}\cdots .
  \tag{17.3}
\]
\begin{enumerate}[label=(\arabic*)]
  \item 如果任意两个相邻映射的复合均为零，即对所有 $n$ 都有 $d_{n+1}\circ d_n=0$，则称序列 $C$ 为一个\emph{上链复形}。
  \item 若 $C$ 是上链复形，则其第 $n$ 个\emph{上同调群}是商群 $\ker d_{n+1}/\operatorname{image}d_n$，记为 $H^n(C)$。
\end{enumerate}

有一个完全类似的“对偶”版本，其中同态出现在指标递减的群之间；这时与 (17.3) 对应的序列写作 $\cdots\to C_n\xrightarrow{d_n}\cdots\to C_0\to0$。若任意两个相邻同态的复合为零，则该复形称为\emph{链复形}，其同调群定义为 $H_n(C)=\ker d_n/\operatorname{image}d_{n+1}$。对链复形，记号通常选为下标且递减；对上链复形，指标通常为上标且递增。我们将对二者使用统一的映射记号，因为从上下文可以清楚看出讨论的是链复形还是上链复形。

链复形首先出现在拓扑背景中，随后出现上链复形。考虑到第 2 节中的应用，我们将主要讨论上链和上同调；不过本节中的所有一般结果都有关于链和同调的类似表述。我们还将关注这样的情形：每个 $C^n$ 都是 $R$-模，而 $d_n$ 是 $R$-模同态（简称为 $R$-模复形），于是群 $H^n(C)$ 也都是 $R$-模。

注意，若 $C$ 是上链（分别为链）复形，则 $C$ 是正合列，当且仅当其所有上同调（分别为同调）群都为零。因此，第 $n$ 个上同调（分别为同调）群度量复形在第 $n$ 阶的非正合性。

\noindent\textbf{定义.} 设 $A=\{A^n\}$ 与 $B=\{B^n\}$ 是上链复形。复形同态 $\alpha:A\to B$ 是一族同态 $\alpha_n:A^n\to B^n$，使得对每个 $n$，图
\[
\begin{array}{ccc}
A^n & \longrightarrow & A^{n+1}\\
\downarrow\alpha_n && \downarrow\alpha_{n+1}\\
B^n & \longrightarrow & B^{n+1}
\end{array}
\tag{17.4}
\]
交换。

\medskip
\noindent\textbf{命题 1.} 上链复形同态 $\alpha:A\to B$ 在相应上同调群上诱导群同态 $H^n(A)\to H^n(B)$，$n\geq0$。

\begin{proof}
容易验证，图 (17.4) 的交换性蕴含第一行各阶段映射的像与核都被映到第二行相应映射的像与核中，因此在相应的商群（上同调群）上给出了良定义的映射。
\end{proof}

\noindent\textbf{定义.} 设 $A=\{A^n\}$、$B=\{B^n\}$ 与 $C=\{C^n\}$ 是上链复形。复形的短正合列
\[
  0\to A\xrightarrow{\alpha}B\xrightarrow{\beta}C\to0
\]
是指一列复形同态，使得对每个 $n$，
\[
  0\to A^n\xrightarrow{\alpha_n}B^n\xrightarrow{\beta_n}C^n\to0
\]
都是短正合的。

上链复形的一个主要特征是它们给出上同调长正合列，这就是我们的第一个主要结果：

\medskip
\noindent\textbf{定理 2.} （上同调长正合列）设
\[
  0\to A\xrightarrow{\alpha}B\xrightarrow{\beta}C\to0
\]
是上链复形的短正合列。则存在上同调群的长正合列
\[
  0\to H^0(A)\to H^0(B)\to H^0(C)\xrightarrow{\delta_0}H^1(A)
\]
\[
  \to H^1(B)\to H^1(C)\xrightarrow{\delta_1}H^2(A)\to\cdots .
  \tag{17.5}
\]
其中同一级上同调群之间的映射是命题 1 中的映射。映射 $\delta_n$ 称为\emph{连接同态}。

\begin{proof}
这个证明的细节相当长。对每个 $n$，验证序列
\[
  H^n(A)\to H^n(B)\to H^n(C)
\]
的正合性，是对各映射正合性定义的直接检查，类似于第 10.5 节定理 33 的证明。练习 2 概述了连接同态 $\delta_n$ 的构造。随后还需要一些工作来证明 $\delta_n$ 是同态，并证明序列在 $\delta_n$ 处正合。
\end{proof}

定理 2 中长正合列存在性的一个直接推论是：若上链复形 $A,B,C$ 中任意两个是正合的，则第三个也是正合的（参见练习 6）。

\subsection*{同态与群 $\Ext_R^n(A,B)$}

为了应用定理 2 来分析序列 (17.2)，我们试图构造一个上链复形，使其长正合列 (17.5) 的最初几个上同调群与 (17.2) 中的项一致。为此我们引入 $R$-模的“分解”概念：

\noindent\textbf{定义.} 设 $A$ 是任意 $R$-模。$A$ 的一个\emph{投射分解}是一个正合列
\[
  \cdots\to P_n\xrightarrow{d_n}P_{n-1}\to\cdots\to P_0\xrightarrow{\varepsilon}A\to0,
  \tag{17.6}
\]
其中每个 $P_i$ 都是投射 $R$-模。

每个 $R$-模都有投射分解：令 $P_0$ 是 $A$ 的一组生成元上的任意自由（因而投射）$R$-模，并由第 10 章定理 6 定义一个满射 $R$-模同态 $\varepsilon:P_0\to A$。这开始了分解 $P_0\to A\to0$；$\varepsilon$ 的满射性保证该序列正合。接着令 $K_0=\ker\varepsilon$，并取一个映到 $P_0$ 的子模 $K_0$ 上的自由模 $P_1$；这给出第二阶段 $P_1\to P_0\to A$，按构造也是正合的。如此继续，在第 $n$ 阶取一个自由 $R$-模 $P_{n+1}$ 满射到 $P_n$ 的子模 $\ker d_n$ 上，实际上就得到 $A$ 的一个自由分解。

在 $A$ 的分解中使用投射模的原因之一，是这使得各种映射可以提升（例如参见后面命题 4 的证明）。

一般而言，投射分解是无限长的；但若 $A$ 本身是投射的，则它有一个非常简单的有限长度投射分解，即由 $A$ 到自身的恒等映射给出的 $0\to A\xrightarrow{1}A\to0$。

给定投射分解 (17.6)，我们可以通过把每一项取到 $D$ 中的同态来形成一个相关序列，同时注意这样会反转同态方向。这得到序列
\[
  0\to\Hom_R(A,D)\xrightarrow{\varepsilon'}\Hom_R(P_0,D)\xrightarrow{d_1}\Hom_R(P_1,D)
  \xrightarrow{d_2}\cdots\xrightarrow{d_n}\Hom_R(P_n,D)\xrightarrow{d_{n+1}}\cdots .
  \tag{17.7}
\]
为简化记号，我们仍把由 $\Hom_R(P_{n-1},D)$ 到 $\Hom_R(P_n,D)$ 的诱导映射记作 $d_n$，$n\geq1$；由 $\varepsilon$ 诱导的映射也类似（参见第 10.5 节）。这个序列不一定正合，但它是一个上链复形（这是第 10.5 节定理 33 证明的一部分）。相应的上同调群有一个专门名称。

\noindent\textbf{定义.} 设 $A$ 与 $D$ 是 $R$-模。对 $A$ 的任意投射分解 (17.6)，令 $d_n:\Hom_R(P_{n-1},D)\to\Hom_R(P_n,D)$，$n\geq1$，如 (17.7) 中所示。定义
\[
  \Ext_R^n(A,D)=\ker d_{n+1}/\operatorname{image}d_n,
\]
其中 $\Ext_R^0(A,D)=\ker d_1$。群 $\Ext_R^n(A,D)$ 称为从函子 $\Hom_R(\_,D)$ 导出的第 $n$ 个上同调群。当 $R=\ZZ$ 时，群 $\Ext_{\ZZ}^n(A,D)$ 也简单记作 $\Ext^n(A,D)$。

注意，群 $\Ext_R^n(A,D)$ 也是从 (17.7) 中把项 $\Hom_R(A,D)$ 替换为 $0$ 后得到的上链复形的上同调群（这并不影响上链性质）；也就是说，它们是上链复形
\[
  0\to\Hom_R(P_0,D)\to\cdots
\]
的上同调群。下面将证明这些上同调群不依赖于 $A$ 的投射分解的选择。在此之前，我们先识别第 $0$ 个上同调群并给出若干例子。

\medskip
\noindent\textbf{命题 3.} 对任意 $R$-模 $A$，有
\[
  \Ext_R^0(A,D)\cong\Hom_R(A,D).
\]

\begin{proof}
由于序列 $P_1\to P_0\to A\to0$ 正合，由第 10.5 节定理 33（注意其证明中的第一个注释）可知，相应序列
\[
  0\to\Hom_R(A,D)\to\Hom_R(P_0,D)\to\Hom_R(P_1,D)
\]
也是正合的。因此 $\Ext_R^0(A,D)=\ker d_1=\operatorname{image}\varepsilon'\cong\Hom_R(A,D)$，如所断言。
\end{proof}

\paragraph{例}
\begin{enumerate}[label=(\arabic*)]
  \item 令 $R=\ZZ$，并令 $A=\ZZ/m\ZZ$，其中 $m\geq2$。由命题，
  \[
    \Ext_{\ZZ}^0(\ZZ/m\ZZ,D)\cong\Hom_{\ZZ}(\ZZ/m\ZZ,D),
  \]
  因此 $\Ext_{\ZZ}^0(\ZZ/m\ZZ,D)\cong {}_mD$，其中
  \[
    {}_mD=\{d\in D\mid md=0\}
  \]
  是 $D$ 中阶整除 $m$ 的元素。对于高阶上同调群，我们使用由 $\ZZ$ 上乘以 $m$ 给出的简单投射分解
  \[
    0\to\ZZ\xrightarrow{m}\ZZ\to\ZZ/m\ZZ\to0.
  \]
  对固定的 $\ZZ$-模 $D$ 取同态，得到上链复形
  \[
    0\to\Hom_{\ZZ}(\ZZ/m\ZZ,D)\to\Hom_{\ZZ}(\ZZ,D)\xrightarrow{m}\Hom_{\ZZ}(\ZZ,D)\to0\to\cdots .
  \]
  我们有 $D\cong\Hom_{\ZZ}(\ZZ,D)$（参见第 10.5 节推论 32 后例 4），在此同构下，对任意 Abel 群 $D$，
  \[
    \Ext_{\ZZ}^1(\ZZ/m\ZZ,D)\cong D/mD.
  \]
  由定义和上述上链复形立刻得到，对所有 $n\geq2$，
  \[
    \Ext_{\ZZ}^n(\ZZ/m\ZZ,D)=0.
  \]
  总结为
  \[
  \Ext_{\ZZ}^0(\ZZ/m\ZZ,D)\cong {}_mD,
  \qquad
  \Ext_{\ZZ}^1(\ZZ/m\ZZ,D)\cong D/mD,
  \]
  \[
  \Ext_{\ZZ}^n(\ZZ/m\ZZ,D)=0,
  \qquad n\geq2.
  \]

  \item 相同的 Abel 群可以是若干不同环 $R$ 上的模，而 $\Ext_R$ 上同调群依赖于 $R$。例如，设 $R=\ZZ/m\ZZ$，其中 $m\geq1$。$R$-模 $D$ 等同于指数整除 $m$ 的 Abel 群 $D$，即 $mD=0$。特别地，对 $m$ 的任意因子 $d$，群 $\ZZ/d\ZZ$ 是 $R$-模，并且
  \[
    \cdots\xrightarrow{m/d}\ZZ/m\ZZ\xrightarrow{d}\ZZ/m\ZZ
    \xrightarrow{m/d}\ZZ/m\ZZ\xrightarrow{d}\ZZ/m\ZZ\to\ZZ/d\ZZ\to0
  \]
  是 $\ZZ/d\ZZ$ 作为 $\ZZ/m\ZZ$-模的投射（事实上自由）分解，最后的映射是自然投影 $x\bmod m\mapsto x\bmod d$。对 $\ZZ/m\ZZ$-模 $D$ 取同态，利用同构 $\Hom_{\ZZ/m\ZZ}(\ZZ/m\ZZ,D)\cong D$，去掉第一项，得到上链复形
  \[
    0\to D\xrightarrow{d}D\xrightarrow{m/d}D\xrightarrow{d}D\xrightarrow{m/d}\cdots .
  \]
  因此
  \[
  \Ext_{\ZZ/m\ZZ}^0(\ZZ/d\ZZ,D)\cong {}_dD,
  \]
  \[
  \Ext_{\ZZ/m\ZZ}^n(\ZZ/d\ZZ,D)\cong {}_{m/d}D/dD,
  \qquad n\text{ 为奇数},\ n\geq1,
  \]
  \[
  \Ext_{\ZZ/m\ZZ}^n(\ZZ/d\ZZ,D)\cong {}_dD/(m/d)D,
  \qquad n\text{ 为偶数},\ n\geq2,
  \]
  其中 ${}_kD=\{d\in D\mid kd=0\}$ 表示被 $k$ 杀死的元素集合。特别地，对所有 $n\geq0$，
  \[
    \Ext_{\ZZ/p^2\ZZ}^n(\ZZ/p\ZZ,\ZZ/p\ZZ)\cong\ZZ/p\ZZ,
  \]
  而例如对所有 $n\geq2$，
  \[
    \Ext_{\ZZ}^n(\ZZ/p\ZZ,\ZZ/p\ZZ)=0.
  \]
\end{enumerate}

为了证明上同调群 $\Ext_R^n(A,D)$ 不依赖于 $A$ 的投射分解选择，我们需要能够“比较”分解。下一个命题说明，从 $A$ 到 $B$ 的 $R$-模同态可以提升为从 $A$ 的投射分解到 $B$ 的投射分解的同态；这种提升性质正是分解中各模投射性重要的一个体现。

\medskip
\noindent\textbf{命题 4.} 设 $f:A\to A'$ 是任意 $R$-模同态，并分别取 $A$ 与 $A'$ 的投射分解。则对每个 $n\geq0$，存在 $f$ 的提升 $f_n$，使得如下图交换：
\[
\begin{array}{ccccccccc}
\cdots&\to&P_1&\to&P_0&\to&A&\to&0\\
&&\downarrow f_1&&\downarrow f_0&&\downarrow f&&\\
\cdots&\to&P'_1&\to&P'_0&\to&A'&\to&0 .
\end{array}
\tag{17.8}
\]

\begin{proof}
给定 (17.8) 中的两行和映射 $f$。由于 $P_0$ 是投射的，可将映射 $f\varepsilon:P_0\to A'$ 提升为映射 $f_0:P_0\to P'_0$，使得 $\varepsilon'f_0=f\varepsilon$（第 10.5 节命题 30(2)）。这给出 $f$ 的第一个提升。归纳地继续，假设 $f_n$ 已定义并使图到该处交换。于是 $\operatorname{image} f_nd_{n+1}\subseteq\ker d'_n$。$P_{n+1}$ 的投射性蕴含我们可以把映射 $f_nd_{n+1}:P_{n+1}\to P'_n$ 提升为映射 $f_{n+1}:P_{n+1}\to P'_{n+1}$，使图在下一阶交换。这完成证明。
\end{proof}

命题 4 中的交换图蕴含诱导图
\[
\begin{array}{cccccc}
0\to&\Hom_R(A,D)&\to&\Hom_R(P_0,D)&\to&\Hom_R(P_1,D)\to\cdots\\
&\uparrow f^*&&\uparrow f_0^*&&\uparrow f_1^*\\
0\to&\Hom_R(A',D)&\to&\Hom_R(P'_0,D)&\to&\Hom_R(P'_1,D)\to\cdots
\end{array}
\tag{17.9}
\]
也交换。这个图的两行是上链复形，而交换图表示这些上链复形之间的一个同态。因此由命题 1 得到它们上同调群上的诱导映射。

\medskip
\noindent\textbf{命题 5.} 设 $f:A\to A'$ 是 $R$-模同态，并取命题 4 中的投射分解。则对每个 $n$，存在一个在这些分解所给上同调群上的诱导群同态
\[
  \varphi_n:\Ext_R^n(A',D)\to\Ext_R^n(A,D),
\]
并且映射 $\varphi_n$ 只依赖于 $f$，不依赖于命题 4 中提升 $f_n$ 的选择。

\begin{proof}
上同调群 $\Ext_R^n$ 上映射的存在性来自命题 1 对上链复形同态 (17.9) 的应用。较困难的部分是证明这些映射不依赖于命题 4 中提升 $f_n$ 的选择。容易看出，这等价于证明若 $f$ 是零映射，则上同调群上的诱导映射也全为零。于是设 $f=0$。由模 $P_i$ 的投射性，可以归纳定义 $R$-模同态 $s_n:P_n\to P'_{n+1}$，使得对所有 $n$，
\[
  f_n=d'_{n+1}s_n+s_{n-1}d_n .
  \tag{17.10}
\]
因此映射 $s_n$ 给出 (17.8) 中方块的反向向下对角箭头。（这组映射 $\{s_n\}$ 称为由 $f_n$ 给出的链同态与零链同态之间的\emph{链同伦}，参见练习 4。）取到 $D$ 中的同态后，图 (17.9) 具有由 $s_n$ 诱导的额外向上对角箭头，并且这些诱导同态满足 (17.10) 中的关系（即它们形成上链复形同态之间的同伦）。现在用添加到 (17.9) 中的对角映射，很容易验证 $\Hom_R(P'_n,D)$ 中代表 $\Ext_R^n(A',D)$ 中一个陪集的任意元素，都映到 $\Ext_R^n(A,D)$ 中的零陪集（参见练习 4）。这完成论证。
\end{proof}

还可以验证，命题 5 中同态 $\varphi_0:\Ext_R^0(A',D)\to\Ext_R^0(A,D)$ 在通过命题 3 中同构识别相应群后，正是第 10.5 节中定义的映射 $f^*:\Hom_R(A',D)\to\Hom_R(A,D)$。

\medskip
\noindent\textbf{定理 6.} 群 $\Ext_R^n(A,D)$ 只依赖于 $A$ 与 $D$，即它们与 $A$ 的投射分解选择无关。

\begin{proof}
在命题 4 的记号中，令 $A'=A$，令 $f:A\to A'$ 为恒等映射，并令 (17.8) 的两行是 $A$ 的两个投射分解。对恒等映射的任意提升，所得上同调群同态
\[
  \varphi_n:\Ext_R^n(A',D)\to\Ext_R^n(A,D)
\]
都是同构，理由如下。把 (17.8) 的顶行复制到第二行下面，作为第三行。令 $g$ 是从 $A'$ 到 $A$ 的恒等映射，并由命题 4 将 $g$ 提升为映射 $g_n:P'_n\to P_n$。令
\[
  \psi_n:\Ext_R^n(A,D)\to\Ext_R^n(A',D)
\]
为所得上同调群上的映射。现在 $g_n\circ f_n:P_n\to P_n$ 是恒等映射 $g\circ f$ 的一个提升，并且它们诱导上同调群上的同态 $\varphi_n\circ\psi_n$。然而，由于第一行与第三行相同，对所有 $n$ 取 $P_n$ 到自身的恒等映射，是 $g\circ f$ 的一个特殊提升，而这种选择显然诱导上同调群上的恒等映射。命题 5 的最后断言于是说明 $\varphi_n\circ\psi_n$ 也是 $\Ext_R^n(A,D)$ 上的恒等映射。对称地，$\psi_n\circ\varphi_n$ 是 $\Ext_R^n(A',D)$ 上的恒等映射。故 $\varphi_n$ 与 $\psi_n$ 是同构，证明完成。
\end{proof}

对固定的 $R$-模 $D$ 与固定整数 $n\geq0$，命题 5 与定理 6 说明 $\Ext_R^n(\_,D)$ 定义了从 $R$-模范畴到 Abel 群范畴的一个（反变）函子。

下一个结果说明，$R$-模 $M$ 的子模及相应商模的投射分解可以拼接起来，得到 $M$ 的一个投射分解。

\medskip
\noindent\textbf{命题 7.} （同时分解）设
\[
  0\to L\to M\to N\to0
\]
是 $R$-模短正合列。令 $L=A$ 有如 (17.6) 的投射分解，并令 $N$ 有类似的投射分解，其投射模记作 $\overline P_n$。则存在一个由投射模 $P_n\oplus\overline P_n$ 给出的 $M$ 的分解，使得如下图交换：
\[
\begin{array}{ccccccccc}
&&\vdots&&\vdots&&\vdots&&\\
&&0\to P_1&\to&P_1\oplus\overline P_1&\to&\overline P_1\to0&&\\
&&\downarrow&&\downarrow&&\downarrow&&\\
&&0\to P_0&\to&P_0\oplus\overline P_0&\to&\overline P_0\to0&&\\
&&\downarrow&&\downarrow&&\downarrow&&\\
&&0\to L&\to&M&\to&N\to0&&\\
&&\downarrow&&\downarrow&&\downarrow&&\\
&&0&&0&&0&&
\end{array}
\tag{17.11}
\]
此外，此图的行与列都是正合的，且各行是分裂的。

\begin{proof}
(17.11) 中左、右非零列按假设正合。中间列中的模是投射的（参见第 10.5 节练习 3），行映射取显然的映射，使每一行都是分裂正合列。剩下的只是定义中间列的竖直映射，使整个图交换。这可以从底部向上归纳地直接完成；这个过程的第一步在练习 5 中概述。
\end{proof}

定理 2 与命题 7 现在给出 $\Ext_R$ 的长正合列，它延长了正合列 (17.2)。

\medskip
\noindent\textbf{定理 8.} 设
\[
  0\to L\to M\to N\to0
\]
是 $R$-模短正合列。则存在 Abel 群长正合列
\[
0\to\Hom_R(N,D)\to\Hom_R(M,D)\to\Hom_R(L,D)\xrightarrow{\delta_0}\Ext_R^1(N,D)
\]
\[
\to\Ext_R^1(M,D)\to\Ext_R^1(L,D)\xrightarrow{\delta_1}\Ext_R^2(N,D)\to\cdots ,
\tag{17.12}
\]
其中同一级 $n$ 的群之间的映射如命题 5，连接同态 $\delta_n$ 由定理 2 给出。

\begin{proof}
取命题 7 中该短正合列的同时投射分解，并取到 $D$ 中的同态。为了从所得图中得到上同调群 $\Ext_R^n$，如命题 3 前讨论所述，我们把变换后图中的最低非零行替换为一行零，得到交换图
\[
\begin{array}{ccccccccc}
&&\vdots&&\vdots&&\vdots&&\\
0\to&\Hom_R(\overline P_1,D)&\to&\Hom_R(P_1\oplus\overline P_1,D)&\to&\Hom_R(P_1,D)&\to0\\
&&\uparrow&&\uparrow&&\uparrow&&\\
0\to&\Hom_R(\overline P_0,D)&\to&\Hom_R(P_0\oplus\overline P_0,D)&\to&\Hom_R(P_0,D)&\to0\\
&&\uparrow&&\uparrow&&\uparrow&&\\
&&0&&0&&0&&
\end{array}
\tag{17.13}
\]
(17.13) 的列是上链复形，而行由第 10.5 节命题 29(2) 及其后讨论知是分裂的。因此 (17.13) 是上链复形的短正合列。定理 2 给出上同调群的长正合列，其各项按定义就是群 $\Ext_R^n(\_,D)$，$n\geq0$。第 $0$ 阶项由命题 3 识别，证明完成。
\end{proof}

定理 8 说明了正合列 (17.2) 如何以自然方式延长，并说明群 $\Ext_R^1(N,D)$ 是 (17.2) 在右端不正合的第一个度量——事实上，(17.2) 可以在右端延长为短正合列，当且仅当 (17.12) 中的连接同态 $\delta_0$ 是零同态。特别地，如果对所有 $R$-模 $N$ 都有 $\Ext_R^1(N,D)=0$，则对每个正合列 (17.1)，(17.2) 都在右端正合。我们已经看到（第 10.5 节推论 35），这蕴含 $R$-模 $D$ 是内射的。下一个结果的一部分说明反之亦然，并以 $\Ext_R$ 群刻画内射模。

\medskip
\noindent\textbf{命题 9.} 对 $R$-模 $Q$，下列条件等价：
\begin{enumerate}[label=(\arabic*)]
  \item $Q$ 是内射的；
  \item 对所有 $R$-模 $A$，$\Ext_R^1(A,Q)=0$；
  \item 对所有 $R$-模 $A$ 以及所有 $n\geq1$，$\Ext_R^n(A,Q)=0$。
\end{enumerate}

\begin{proof}
上面已经证明 (2) 蕴含 (1)，而 (3) 蕴含 (2) 显然。因此只需证明若 $Q$ 是内射的，则对所有 $R$-模 $A$ 与所有 $n\geq1$，$\Ext_R^n(A,Q)=0$。取 $A$ 的投射分解
\[
  \cdots\to P_n\to P_{n-1}\to\cdots\to P_0\to A\to0.
\]
由于 $Q$ 是内射的，序列
\[
  0\to\Hom_R(A,Q)\to\Hom_R(P_0,Q)\to\Hom_R(P_1,Q)\to\cdots\to\Hom_R(P_n,Q)\to\cdots
\]
仍是正合的（第 10.5 节推论 35），所以这个上链复形的所有上同调群都是 $0$。特别地，对 $n\geq1$，群 $\Ext_R^n(A,Q)$ 全都平凡，即 (3)。
\end{proof}

对固定的 $R$-模 $D$，定理 8 可视为说明在施加左正合函子 $\Hom_R(\_,D)$ 后，短正合列 $0\to L\to M\to N\to0$ 的右端发生了什么。这就是为什么（反变）函子 $\Ext_R^n(\_,D)$ 称为函子 $\Hom_R(\_,D)$ 的右导出函子。

也可以考虑施加左正合函子 $\Hom_R(D,\_)$ 的效果，即从 $D$ 出发取同态而不是取到 $D$ 中的同态。下一个定理说明，同样的 $\Ext_R$ 群事实上也定义了 $\Hom_R(D,\_)$ 的（协变）右导出函子。

\medskip
\noindent\textbf{定理 10.} 设
\[
  0\to L\to M\to N\to0
\]
是 $R$-模短正合列。则存在 Abel 群长正合列
\[
0\to\Hom_R(D,L)\to\Hom_R(D,M)\to\Hom_R(D,N)\xrightarrow{\gamma_0}\Ext_R^1(D,L)
\]
\[
\to\Ext_R^1(D,M)\to\Ext_R^1(D,N)\xrightarrow{\gamma_1}\Ext_R^2(D,L)\to\cdots .
\tag{17.14}
\]

\begin{proof}
设 $0\to L\to M\to N\to0$ 是 $R$-模短正合列。取 $D$ 的一个投射分解，然后对该分解施加 $\Hom_R(\_,L)$、$\Hom_R(\_,M)$ 与 $\Hom_R(\_,N)$，即可得到一个类似于 (17.13) 的交换图，只是 $L,M,N$ 出现在第二个位置而不是第一个位置。对这个阵列应用长正合列定理即得到 (17.14)。
\end{proof}

定理 10 说明，群 $\Ext_R^1(D,L)$ 度量正合列
\[
  0\to\Hom_R(D,L)\to\Hom_R(D,M)\to\Hom_R(D,N)
\]
能否延长为短正合列——当且仅当 $\gamma_0$ 是零同态时可以。特别地，如果模 $D$ 具有对所有 $R$-模 $B$ 都有 $\Ext_R^1(D,B)=0$ 的性质，则这种延长总是成立；在这种情形下，由第 10.5 节推论 32 可知 $D$ 是投射 $R$-模。与命题 9 中内射 $R$-模的情形一样，这些上同调群的消失事实上刻画投射 $R$-模：

\medskip
\noindent\textbf{命题 11.} 对 $R$-模 $P$，下列条件等价：
\begin{enumerate}[label=(\arabic*)]
  \item $P$ 是投射的；
  \item 对所有 $R$-模 $B$，$\Ext_R^1(P,B)=0$；
  \item 对所有 $R$-模 $B$ 以及所有 $n\geq1$，$\Ext_R^n(P,B)=0$。
\end{enumerate}

\begin{proof}
上面已经证明 (2) 蕴含 (1)，而 (3) 蕴含 (2) 显然。因此只需证明 (1) 蕴含 (3)。若 $P$ 是投射 $R$-模，则由 $P$ 上恒等映射给出的简单正合列
\[
  0\to P\xrightarrow{1}P\to0
\]
是 $P$ 的一个投射分解。取到 $B$ 中的同态给出简单上链复形
\[
  0\to\Hom_R(P,B)\xrightarrow{1}\Hom_R(P,B)\to0\to\cdots\to0\to\cdots,
\]
由定义可知，对所有 $n\geq1$，$\Ext_R^n(P,B)=0$，得到 (3)。
\end{proof}

\paragraph{例}
\begin{enumerate}[label=(\arabic*)]
  \item 由于 $\ZZ^m$ 是自由的，因而是投射的 $\ZZ$-模，由命题 11 得到
  \[
    \Ext_{\ZZ}^n(\ZZ^m,B)=0
  \]
  对所有 Abel 群 $B$、所有 $m\geq1$ 与所有 $n\geq1$ 都成立。
  \item 不难证明，对所有 $n\geq0$，
  \[
    \Ext_R^n(A_1\oplus A_2,B)\cong\Ext_R^n(A_1,B)\oplus\Ext_R^n(A_2,B)
  \]
  （参见练习 10），所以前一个例子结合命题 3 后的例子，就确定了对所有有限生成 Abel 群 $A$ 的 $\Ext_{\ZZ}^n(A,B)$。特别地，对所有有限生成群 $A$、所有 Abel 群 $B$ 以及所有 $n\geq2$，有 $\Ext_{\ZZ}^n(A,B)=0$。
\end{enumerate}

我们选择用 $A$ 的投射分解来定义上同调群 $\Ext_R^n(A,B)$。还有一个平行的发展，使用 $B$ 的内射分解：
\[
  0\to B\to Q_0\to Q_1\to\cdots,
\]
其中每个 $Q_i$ 都是内射的。在这种情形下，$\Ext_R^n(A,B)$ 定义为对 $B$ 的分解施加 $\Hom_R(A,\_)$ 后所得上链序列的第 $n$ 个上同调群。理论以类似于本节发展的方式进行。最终可以证明，用两种方法构造的群 $\Ext_R^n(A,B)$ 之间有自然同构。

\paragraph{例}
\begin{enumerate}[label=(\arabic*)]
  \item 设 $R=\ZZ$，且 $A,B$ 是 $\ZZ$-模，即 Abel 群。回忆 $\ZZ$-模是内射的，当且仅当它可除（第 10.5 节命题 36）。群 $B$ 可嵌入某个内射 $\ZZ$-模 $Q_0$（第 10.5 节推论 37），而 $Q_0$ 模去 $B$ 的像所得的商 $Q_1$ 仍是内射的。因此有 $B$ 的一个内射分解
  \[
    0\to B\to Q_0\to Q_1\to0.
  \]
  对此序列施加 $\Hom_{\ZZ}(A,\_)$，得到上链复形
  \[
    0\to\Hom_{\ZZ}(A,B)\to\Hom_{\ZZ}(A,Q_0)\to\Hom_{\ZZ}(A,Q_1)\to0\to\cdots,
  \]
  由此立刻得到
  \[
    \Ext_{\ZZ}^n(A,B)=0
  \]
  对所有 Abel 群 $A,B$ 与所有 $n\geq2$ 成立。这说明前一个例子的结果在 $A$ 不有限生成时也成立。

  \item 假设 $A$ 是扭 Abel 群。则 $\Ext_{\ZZ}^0(A,\ZZ)\cong\Hom(A,\ZZ)=0$，因为 $\ZZ$ 无扭。序列 $0\to\ZZ\to\QQ\to\QQ/\ZZ\to0$ 给出 $\ZZ$ 的一个内射分解。施加 $\Hom(A,\_)$ 得到上链复形
  \[
    0\to\Hom(A,\ZZ)\to\Hom(A,\QQ)\to\Hom(A,\QQ/\ZZ)\to0\to\cdots,
  \]
  而由于 $\QQ$ 也无扭，这说明
  \[
    \Ext^1(A,\ZZ)\cong\Hom_{\ZZ}(A,\QQ/\ZZ).
  \]
  群 $\Hom(A,\QQ/\ZZ)$ 称为 $A$ 的 Pontriagin 对偶群。若 $A$ 是有限 Abel 群，则 $A$ 的 Pontriagin 对偶同构于 $A$（参见第 5.2 节练习 14）。特别地，对所有非零有限 Abel 群 $A$，$\Ext^1(A,\ZZ)\cong A$ 非零。由前一个例子，对所有 $n\geq2$，$\Ext^n(A,\ZZ)=0$。
\end{enumerate}

我们记录 $\Ext_R^1$ 的一个重要性质，它有助于解释这些上同调群名称的来源。回忆等价扩张在第 10.5 节开头已经定义。

\medskip
\noindent\textbf{定理 12.} 对任意 $R$-模 $N$ 与 $L$，$\Ext_R^1(N,L)$ 与 $N$ 由 $L$ 扩张的等价类集合之间存在一个双射。

虽然我们不证明此结果，但在第 4 节中我们将建立一个类似的双射：群 $G$ 被任意 $\ZZ G$-模 $A$ 扩张的等价类，与某个上同调群的元素之间的双射。

\paragraph{例}
设 $R=\ZZ$ 且 $A=B=\ZZ/p\ZZ$。我们在上面证明了 $\Ext_{\ZZ}^1(\ZZ/p\ZZ,\ZZ/p\ZZ)\cong\ZZ/p\ZZ$，因此由定理 12，$\ZZ/p\ZZ$ 由 $\ZZ/p\ZZ$ 扩张的等价类恰好有 $p$ 个。它们由直和 $\ZZ/p\ZZ\oplus\ZZ/p\ZZ$（对应于 $\Ext_{\ZZ}^1(\ZZ/p\ZZ,\ZZ/p\ZZ)$ 中的平凡类）以及 $p-1$ 个扩张
\[
  0\to\ZZ/p\ZZ\to\ZZ/p^2\ZZ\xrightarrow{i}\ZZ/p\ZZ\to0
\]
给出，其中映射 $i(x)=ix\bmod p$，$i=1,2,\ldots,p-1$。注意，尽管这些作为扩张是不等价的，它们都确定同一个群 $\ZZ/p^2\ZZ$。

\subsection*{张量积与群 $\Tor_n^R(A,B)$}

上同调群 $\Ext_R^n(A,B)$ 确定在施加左正合函子 $\Hom_R(D,\_)$ 与 $\Hom_R(\_,D)$ 后，短正合列在右端的行为。类似地，可以询问在施加右正合函子 $D\otimes_R\_$ 或右正合函子 $\_\otimes_R D$ 后，短正合列在左端的行为。这导致 $\Tor$（同调）群（其名称源于它与扭子模的关系），下面我们简要概述这些左导出函子的发展。在某些方面，该理论与 $\Ext_R$ 理论“对偶”。我们集中讨论 $D\otimes_R\_$ 的情形，其中 $D$ 是右 $R$-模。当 $D$ 是左 $R$-模时，对 $\_\otimes_R D$ 有完全对称的理论；当 $R$ 交换且所有 $R$-模具有相同的左右 $R$-作用时，两种发展所得的同调群同构。

于是设 $D$ 是右 $R$-模。对每个左 $R$-模 $B$，张量积 $D\otimes_R B$ 是 Abel 群，函子 $D\otimes_R\_$ 是协变且右正合的；即对任意左 $R$-模短正合列 (17.1)，
\[
  D\otimes L\to D\otimes M\to D\otimes N\to0
\]
是 Abel 群正合列。这个序列可以在左端延长为长正合列如下。令
\[
  \cdots\to P_n\xrightarrow{d_n}P_{n-1}\to\cdots\to P_0\xrightarrow{\varepsilon}B\to0
\]
是 $B$ 的一个投射分解，并与 $D$ 作张量得到
\[
  \cdots\to D\otimes P_n\xrightarrow{1\otimes d_n}D\otimes P_{n-1}\to\cdots\to D\otimes P_0\xrightarrow{1\otimes\varepsilon}D\otimes B\to0.
  \tag{17.15}
\]
由第 10.5 节定理 39 中的论证可知，(17.15) 是一个链复形——任意两个相邻映射的复合为零——因此可以形成其同调群。

\noindent\textbf{定义.} 设 $D$ 是右 $R$-模，$B$ 是左 $R$-模。对上面任意由左 $R$-模给出的 $B$ 的投射分解，令 $1\otimes d_n:D\otimes P_n\to D\otimes P_{n-1}$，$n\geq1$，如 (17.15) 所示。定义
\[
  \Tor_n^R(D,B)=\ker(1\otimes d_n)/\operatorname{image}(1\otimes d_{n+1}),
\]
其中
\[
  \Tor_0^R(D,B)=(D\otimes P_0)/\operatorname{image}(1\otimes d_1).
\]
群 $\Tor_n^R(D,B)$ 称为从函子 $D\otimes_R\_$ 导出的第 $n$ 个同调群。当 $R=\ZZ$ 时，群 $\Tor_n^{\ZZ}(D,B)$ 也简单记作 $\Tor_n(D,B)$。

因此，$\Tor_n^R(D,B)$ 是从 (17.15) 中去掉项 $D\otimes B$ 后所得链复形的第 $n$ 个同调群。

与命题 3 完全类似的证明（但依赖第 10.5 节定理 39）蕴含如下结果：

\medskip
\noindent\textbf{命题 13.} 对任意左 $R$-模 $B$，有
\[
  \Tor_0^R(D,B)\cong D\otimes_R B.
\]

\paragraph{例}
令 $R=\ZZ$，并令 $B=\ZZ/m\ZZ$，其中 $m\geq2$。由命题，$\Tor_0(D,\ZZ/m\ZZ)$ 同构于 $D\otimes\ZZ/m\ZZ$，故有 $\Tor_0(D,\ZZ/m\ZZ)\cong D/mD$（第 10.4 节推论 12 后例 8）。对高阶群，我们对 $B$ 的投射分解
\[
  0\to\ZZ\xrightarrow{m}\ZZ\to\ZZ/m\ZZ\to0
\]
施加 $D\otimes\_$，并使用同构 $D\otimes\ZZ\cong D$ 与 $D\otimes\ZZ/m\ZZ\cong D/mD$。这给出链复形
\[
  \cdots\to0\to D\xrightarrow{m}D\to D/mD\to0.
\]
于是 $\Tor_1(D,\ZZ/m\ZZ)\cong {}_mD$，即 $D$ 中被 $m$ 零化的子群，并且对所有 $n\geq2$，$\Tor_n(D,\ZZ/m\ZZ)=0$。总结为
\[
\Tor_0(D,\ZZ/m\ZZ)\cong D/mD,
\qquad
\Tor_1(D,\ZZ/m\ZZ)\cong {}_mD,
\]
\[
\Tor_n(D,\ZZ/m\ZZ)=0,
\qquad n\geq2.
\]
与 $\Ext$ 一样，$\Tor$ 群依赖于环 $R$（参见练习 20）。

沿着与 $\Ext_R$ 类似的发展，可以证明：

\medskip
\noindent\textbf{命题 14.}
\begin{enumerate}[label=(\arabic*)]
  \item 同调群 $\Tor_n^R(D,B)$ 与 $B$ 的投射分解的选择无关；
  \item 对每个 $R$-模同态 $f:B\to B'$，存在同调群上的诱导映射
  \[
    \psi_n:\Tor_n^R(D,B)\to\Tor_n^R(D,B'),
  \]
  并且只依赖于 $f$。
\end{enumerate}

有一个与定理 2 类似的同调长正合列，只是所有箭头反向；其证明与上同调情形的论证完全类似。它与同时分解一起给出：

\medskip
\noindent\textbf{定理 15.} 设
\[
  0\to L\to M\to N\to0
\]
是左 $R$-模短正合列。则存在 Abel 群长正合列
\[
  \cdots\to\Tor_2^R(D,N)\xrightarrow{\delta_2}\Tor_1^R(D,L)\to\Tor_1^R(D,M)\to
\]
\[
  \Tor_1^R(D,N)\xrightarrow{\delta_1}D\otimes L\to D\otimes M\to D\otimes N\to0,
\]
其中同一级 $n$ 的群之间的映射如命题 14，映射 $\delta_n$ 称为连接同态。

平坦模也有与命题 9 和 11 对应的刻画，其证明非常类似，留作练习。

\medskip
\noindent\textbf{命题 16.} 对右 $R$-模 $D$，下列条件等价：
\begin{enumerate}[label=(\arabic*)]
  \item $D$ 是平坦 $R$-模；
  \item 对所有左 $R$-模 $B$，$\Tor_1^R(D,B)=0$；
  \item 对所有左 $R$-模 $B$ 与所有 $n\geq1$，$\Tor_n^R(D,B)=0$。
\end{enumerate}

我们把 $\Tor_n^R(A,B)$ 定义为将 $B$ 的投射分解在左侧与 $A$ 张量后所得链复形的同调。也可以通过把 $A$ 的投射分解在右侧与 $B$ 张量后所得链复形来得到同样的群。换言之，$\Tor_n^R(A,B)$ 群同时定义了两个右正合函子 $A\otimes_R\_$ 与 $\_\otimes_R B$ 的（协变）左导出函子：若 $D$ 是左 $R$-模，则右 $R$-模短正合列 $0\to L\to M\to N\to0$ 给出 Abel 群长正合列
\[
  \cdots\to\Tor_2^R(N,D)\xrightarrow{\delta_2}\Tor_1^R(L,D)\to\Tor_1^R(M,D)\to
\]
\[
  \Tor_1^R(N,D)\xrightarrow{\delta_1}L\otimes_R D\to M\otimes_R D\to N\otimes_R D\to0.
\]
特别地，左 $R$-模 $D$ 是平坦的，当且仅当对所有右 $R$-模 $A$，$\Tor_1^R(A,D)=0$。

当 $R$ 交换时，对任意两个具有标准 $R$-模结构的 $R$-模 $A$ 与 $B$，有 $A\otimes_RB\cong B\otimes_RA$（第 10.4 节命题 20），于是 $\Tor_n^R(A,B)\cong\Tor_n^R(B,A)$ 作为 $R$-模。当 $R$ 交换时，$\Tor$ 长正合列是 $R$-模正合列。

\paragraph{例}
\begin{enumerate}[label=(\arabic*)]
  \item 若 $R=\ZZ$，则由于 $\ZZ^m$ 是自由的，因而是平坦的（第 10.5 节推论 42），我们有
  \[
    \Tor_n(A,\ZZ^m)=0
  \]
  对所有 $n\geq1$ 与所有 Abel 群 $A$ 成立。
  \item 由于 $\Tor_n^R(A,B_1\oplus B_2)\cong\Tor_n^R(A,B_1)\oplus\Tor_n^R(A,B_2)$（参见练习 10），前面两个例子一起确定了所有 Abel 群 $A$ 与所有有限生成 Abel 群 $B$ 的 $\Tor_n(A,B)$。
  \item 作为前一个例子的特例，$\Tor_n(A,B)$ 是扭群，并且对每个 Abel 群 $A$、每个有限生成 Abel 群 $B$ 与所有 $n\geq2$，$\Tor_n(A,B)=0$。事实上这些结果在不要求 $B$ 有限生成时也成立。
  \item 正合列 $0\to\ZZ\to\QQ\to\QQ/\ZZ\to0$ 给出长正合列
  \[
    \cdots\to\Tor_1(D,\QQ)\to\Tor_1(D,\QQ/\ZZ)\to D\otimes\ZZ\to D\otimes\QQ\to D\otimes\QQ/\ZZ\to0.
  \]
  由于 $\QQ$ 是平坦 $\ZZ$-模（第 10.5 节推论 42 后例 2），命题说明我们有正合列
  \[
    0\to\Tor_1(D,\QQ/\ZZ)\to D\to D\otimes\QQ,
  \]
  所以 $\Tor_1(D,\QQ/\ZZ)$ 同构于从 $D$ 到 $D\otimes\QQ$ 的自然映射的核，即 $D$ 的扭子群（参见第 10.4 节练习 9）。
\end{enumerate}

下面的结果说明，当 $R=\ZZ$ 时，$\Tor$ 群与扭子群密切相关。$\Tor$ 群最早出现在拓扑背景中扭 Abel 群的应用中，这有助于解释这个术语。

\medskip
\noindent\textbf{命题 17.} 设 $A$ 与 $B$ 是 $\ZZ$-模，并令 $t(A)$ 与 $t(B)$ 分别表示它们的扭子模。则
\[
  \Tor_1(A,B)\cong\Tor_1(t(A),t(B)).
\]

\begin{proof}
当 $A$ 与 $B$ 是有限生成 Abel 群时，这由上面例 3 与例 4 推出。一般情形参见练习 16。
\end{proof}

\medskip
\noindent\textbf{推论 18.} 若 $A$ 是 Abel 群，则 $A$ 无扭，当且仅当对每个 Abel 群 $B$，$\Tor_1(A,B)=0$（等价地，只需对所有有限循环群 $B$ 成立）。


\begin{proof}
由命题，若 $A$ 没有有限阶元素，则对每个 Abel 群 $B$，
\[
  \Tor_1(A,B)=\Tor_1(t(A),B)=\Tor_1(0,B)=0.
\]
反过来，若对所有 $B$ 都有 $\Tor_1(A,B)=0$，则特别地 $\Tor_1(A,\QQ/\ZZ)=0$；而由上面的例子，这个群同构于 $A$ 的扭子群。
\end{proof}

命题 17 与推论 18 的结果在用任意 P.I.D. $R$ 代替 $\ZZ$ 时也成立（参见第 10.5 节练习 26 以及练习 16）。

最后，我们提到，我们所描述的上同调和同调理论，可以通过公理化 $R$-模以及函子 $\Hom_R$ 和张量积的本质性质，在更广泛得多的背景中发展。这导致 Abel 范畴和加性函子的一般概念。在 $R$-模的 Abel 范畴中，任意到 Abel 群范畴的加性函子 $\mathcal F$ 都给出一族导出函子 $\mathcal F_n$，它们也从 $R$-模到 Abel 群，$n\geq0$。于是对每个 $R$-模短正合列 $0\to L\to M\to N\to0$，都有一个（上同调或同调）群的长正合列，其项为 $\mathcal F_n(L)$、$\mathcal F_n(M)$ 与 $\mathcal F_n(N)$；这些长正合列反映函子 $\mathcal F$ 的正合性。若 $\mathcal F$ 左正合或右正合，则第 $0$ 个导出函子 $\mathcal F_0$ 自然等价于 $\mathcal F$（因此 $0$ 次群 $\mathcal F_0(X)$ 同构于 $\mathcal F(X)$）；若 $\mathcal F$ 是正合函子，则对所有 $n\geq1$ 和所有 $R$-模 $X$，$\mathcal F_n(X)=0$。

\subsection*{练习}

\begin{enumerate}[label=\arabic*.]
  \item 给出命题 1 证明的细节。

  \item 本练习定义定理 2 中长正合列的连接映射 $\delta_n$，并证明它是同态。在定理 2 的记号中，令
  \[
    0\to A\xrightarrow{\alpha}B\xrightarrow{\beta}C\to0
  \]
  为上链复形的短正合列；为简单起见，把 $A,B,C$ 的上链映射都记作 $d$。
  \begin{enumerate}[label=(\alph*)]
    \item 若 $c\in C^n$ 代表类 $x\in H^n(C)$，证明存在某个 $b\in B^n$ 使 $\beta_n(b)=c$。
    \item 证明 $d_{n+1}(b)\in\ker\beta_{n+1}$，并推出存在唯一的 $a\in A^{n+1}$ 使得 $\alpha_{n+1}(a)=d_{n+1}(b)$。［用 $c\in\ker d_{n+1}$ 以及图的交换性。］
    \item 证明 $d_{n+2}(a)=0$，并推出 $a$ 在商群 $H^{n+1}(A)$ 中定义了一个类 $\bar a$。［用 $\alpha_{n+2}$ 是单射。］
    \item 证明 $\bar a$ 与 $b$ 的选择无关；即若 $b'$ 是另一个选择，且 $a'$ 是它在 $A^{n+1}$ 中的唯一原像，则 $\bar a=\bar a'$；并证明 $\bar a$ 也与代表类 $x$ 的 $c$ 的选择无关。
    \item 定义 $\delta_n(x)=\bar a$，并证明 $\delta_n$ 是从 $H^n(C)$ 到 $H^{n+1}(A)$ 的群同态。［利用 $\delta_n(x)$ 与 $c,b$ 选择无关来计算 $\delta_n(x_1+x_2)$。］
  \end{enumerate}

  \item 假设有一个 $R$-模交换图，其行正合：
  \[
  \begin{array}{ccccccccc}
   0&\to&A&\xrightarrow{\alpha}&B&\xrightarrow{\beta}&C&\to&0\\
   &&\downarrow f&&\downarrow g&&\downarrow h&&\\
   0&\to&A'&\xrightarrow{\alpha'}&B'&\xrightarrow{\beta'}&C' .&&
  \end{array}
  \]
  \begin{enumerate}[label=(\alph*)]
    \item 若 $c\in\ker h$ 且 $\beta(b)=c$，证明 $g(b)\in\ker\beta'$，并推出对某个 $a'\in A'$ 有 $g(b)=\alpha'(a')$。［用图的交换性。］
    \item 证明 $\delta(c)=a'\bmod\operatorname{image}f$ 是从 $\ker h$ 到商 $A'/\operatorname{image}f$ 的良定义 $R$-模同态。
    \item （蛇引理）证明存在正合列
    \[
      \ker f\to\ker g\to\ker h\xrightarrow{\delta}\operatorname{coker} f\to\operatorname{coker} g\to\operatorname{coker} h,
    \]
    其中 $\operatorname{coker} f$ 是 $A'/\operatorname{image}f$，$\operatorname{coker} g,\operatorname{coker} h$ 类似。
    \item 证明若 $\alpha$ 是单射且 $\beta'$ 是满射（即上面的两行都可以延伸为短正合列），则 (c) 中的正合列可以延伸为
    \[
      0\to\ker f\to\ker g\to\ker h\xrightarrow{\delta}\operatorname{coker} f\to\operatorname{coker} g\to\operatorname{coker} h\to0.
    \]
  \end{enumerate}

  \item 设 $A=\{A^n\}$ 与 $B=\{B^n\}$ 是上链复形，其中两个复形中 $A^n\to A^{n+1}$ 与 $B^n\to B^{n+1}$ 的映射均记为 $d_{n+1}$。若对所有 $n$ 存在模同态 $s_n:A^{n+1}\to B^n$，使得从 $A^n$ 到 $B^n$ 的映射 $\alpha_n-\beta_n$ 满足
  \[
    \alpha_n-\beta_n=d_ns_{n-1}+s_nd_{n+1},
  \]
  则称上链复形同态 $\alpha$ 与 $\beta$ 同伦。这族映射 $\{s_n\}$ 称为从 $\alpha$ 到 $\beta$ 的上链同伦。链复形之间也可以类似定义链同伦。
  \begin{enumerate}[label=(\alph*)]
    \item 证明同伦的上链复形映射诱导相同的上同调映射；也就是说，若 $\alpha$ 与 $\beta$ 是同伦上链复形同态，则对每个 $n\geq0$，从 $H^n(A)$ 到 $H^n(B)$ 的诱导群同态相等。（因此，“同伦”给出两个复形映射诱导相同上同调或同调映射的一个充分条件；该条件一般不是必要的。）［用同伦定义证明，对每个 $z\in\ker d_{n+1}$，$(\alpha_n-\beta_n)(z)\in\operatorname{image}d_n$。］
    \item 证明关系 $\alpha\sim\beta$（若 $\alpha$ 与 $\beta$ 同伦）是在任意上链复形同态集合上的等价关系。
  \end{enumerate}

  \item 如下建立命题 7“同时分解”结果的第一步：假设图 (17.11) 中前两条非零行已经给定，只缺少从 $P_0\oplus\overline P_0$ 到 $M$ 的映射（其中投射模行上的映射是此分裂正合列的显然嵌入与投影）。令 $\mu:\overline P_0\to M$ 是映射 $\overline P_0\to N$ 到 $\overline P_0$ 的提升（因 $\overline P_0$ 投射而存在）。令 $\lambda$ 是图中复合 $P_0\to L\to M$。定义
  \[
    \eta:P_0\oplus\overline P_0\to M,\qquad \eta(x,y)=\lambda(x)+\mu(y).
  \]
  证明在此定义下，(17.11) 的前两条非零行形成交换图。

  \item 设 $0\to A\to B\to C\to0$ 是上链复形的短正合列。证明若 $A,B,C$ 中任意两个正合，则第三个也正合。［用定理 2。］

  \item 证明有限生成 Abel 群 $A$ 是自由的，当且仅当 $\Ext^1(A,\ZZ)=0$。

  \item 证明若 $0\to L\to M\to N\to0$ 是 $R$-模分裂短正合列，则对每个 $n\geq0$，序列
  \[
    0\to\Ext_R^n(N,D)\to\Ext_R^n(M,D)\to\Ext_R^n(L,D)\to0
  \]
  也是短正合且分裂。［用一个分裂同态和命题 5。］

  \item 证明
  \[
    0\to\ZZ/d\ZZ\to\ZZ/m\ZZ\xrightarrow{d}\ZZ/m\ZZ\xrightarrow{m/d}\ZZ/m\ZZ\xrightarrow{d}\ZZ/m\ZZ\to\cdots
  \]
  是 $\ZZ/d\ZZ$ 作为 $\ZZ/m\ZZ$-模的内射分解。［用第 10.5 节命题 36。］
  利用它按对偶群 $\Hom_{\ZZ/m\ZZ}(A,\ZZ/m\ZZ)$ 计算群 $\Ext_{\ZZ/m\ZZ}^n(A,\ZZ/d\ZZ)$。特别地，若 $m=p^2$ 且 $d=p$，给出结果 $\Ext_{\ZZ/p^2\ZZ}^n(\ZZ/p\ZZ,\ZZ/p\ZZ)\cong\ZZ/p\ZZ$ 的另一种推导。

  \item \begin{enumerate}[label=(\alph*)]
    \item 证明投射模 $P_i$ 的任意直和 $\bigoplus_{i\in I}P_i$ 仍是投射的，并证明内射模 $Q_j$ 的任意直积 $\prod_{j\in J}Q_j$ 仍是内射的。
    \item 证明投射分解的任意直和仍是投射分解，并用此证明对任意 $R$-模族 $A_i$（$i\in I$），有
    \[
      \Ext_R^n\left(\bigoplus_{i\in I}A_i,B\right)\cong\prod_{i\in I}\Ext_R^n(A_i,B).
    \]
    ［参见第 10.5 节练习 12。］
    \item 证明内射分解的任意直积仍是内射分解，并用此证明对任意 $R$-模族 $B_j$（$j\in J$），有
    \[
      \Ext_R^n\left(A,\prod_{j\in J}B_j\right)\cong\prod_{j\in J}\Ext_R^n(A,B_j).
    \]
    ［参见第 10.5 节练习 12。］
    \item 证明对任意 $R$-模族 $B_j$（$j\in J$），有
    \[
      \Tor_n^R\left(A,\bigoplus_{j\in J}B_j\right)\cong\bigoplus_{j\in J}\Tor_n^R(A,B_j).
    \]
  \end{enumerate}

  \item （Bass 对 Noether 环的刻画）设 $R$ 是交换环。
  \begin{enumerate}[label=(\alph*)]
    \item 若 $R$ 是 Noether 环，且 $I$ 是 $R$ 的任意非零理想，证明从 $I$ 到内射 $R$-模 $Q_j$（$j\in\mathcal J$）的直和 $\bigoplus_{j\in\mathcal J}Q_j$ 的任意 $R$-模同态 $f:I\to\bigoplus_{j\in\mathcal J}Q_j$ 的像，都包含在某个有限直和的 $Q_j$ 中。
    \item 若 $R$ 是 Noether 环，证明内射 $R$-模的任意直和 $\bigoplus_{j\in\mathcal J}Q_j$ 仍是内射的。［用 Baer 判别法（命题 36）以及第 10.5 节练习 4 与 (a)。］
    \item 设 $I_1\subseteq I_2\subseteq\cdots$ 是 $R$ 的理想升链，其并为 $I$；并令 $I/I_i\to Q_i$（$i=1,2,\ldots$）是商 $I/I_i$ 到内射 $R$-模 $Q_i$ 中的嵌入（由第 10.5 节定理 38）。证明这些嵌入与典范投影 $I\to I/I_i$ 的乘积复合，给出 $R$-模同态
    \[
      f:I\to\bigoplus_{i=1,2,\ldots}Q_i.
    \]
    \item 证明 (b) 的反命题：若内射 $R$-模的任意直和 $\bigoplus_{j\in\mathcal J}Q_j$ 仍是内射的，则 $R$ 是 Noether 环。［若 (c) 中的直和是内射的，用 Baer 判别法把 $f$ 提升为同态 $F:R\to\bigoplus_{i=1,2,\ldots}Q_i$。若 $F(1)$ 在 $Q_i$ 中的分量对所有 $i\geq n$ 为 $0$，证明 $I=I_n$，从而理想升链有限。］
  \end{enumerate}

  \item 证明命题 13：$\Tor_0^R(D,A)\cong D\otimes_R A$。［仿照命题 3 的证明。］

  \item 证明命题 16 中平坦模的刻画。

  \item 假设 $0\to A\to B\to C\to0$ 是 $R$-模短正合列。证明若 $C$ 是平坦 $R$-模，则 $A$ 平坦当且仅当 $B$ 也平坦。［用 $\Tor$ 长正合列。］给出一个例子说明，若 $A$ 和 $B$ 平坦，则 $C$ 未必平坦。

  \item \begin{enumerate}[label=(\alph*)]
    \item 若 $I$ 是 $R$ 的理想，$M$ 是 $R$-模，证明 $\Tor_1^R(M,R/I)$ 同构于映射 $M\otimes_R I\to M$ 的核，该映射把 $m\otimes i$ 送到 $mi$，其中 $i\in I$ 且 $m\in M$。［用与 $0\to I\to R\to R/I\to0$ 相关的 $\Tor$ 长正合列，并注意 $R$ 是平坦的。］
    \item （用 $\Tor$ 的平坦性判别法）证明 $R$-模 $M$ 是平坦的，当且仅当对每个有限生成理想 $I$，$\Tor_1^R(M,R/I)=0$。［用第 10.5 节练习 25。］
  \end{enumerate}

  \item 设 $R$ 是 P.I.D.，$A,B$ 是 $R$-模。若 $t(B)$ 表示 $B$ 的扭子模，证明 $\Tor_1^R(A,t(B))\cong\Tor_1^R(A,B)$，并推出 $\Tor_1^R(A,B)$ 同构于 $\Tor_1^R(t(A),t(B))$。［用第 10.5 节练习 26 证明 $B/t(B)$ 在 $R$ 上平坦，然后把以 $D=A$ 的 $\Tor$ 长正合列应用于短正合列 $0\to t(B)\to B\to B/t(B)\to0$，并用命题 16 后面的评论。］

  \item 令
  \[
    A=\ZZ/2\ZZ\oplus\ZZ/3\ZZ\oplus\ZZ/4\ZZ\oplus\cdots .
  \]
  对任意 Abel 群 $B$，证明
  \[
    \Ext^1(A,B)\cong (B/2B)\times(B/3B)\times(B/4B)\times\cdots .
  \]
  ［用练习 10。］证明 $\Ext^1(A,B)=0$ 当且仅当 $B$ 可除。

  \item 证明 $\ZZ/2\ZZ$ 是投射 $\ZZ/6\ZZ$-模，并推出
  \[
    \Tor_n^{\ZZ/6\ZZ}(\ZZ/2\ZZ,\ZZ/2\ZZ)=0.
  \]

  \item 设 $r\neq0$ 不是交换环 $R$ 中的零因子。
  \begin{enumerate}[label=(\alph*)]
    \item 证明乘以 $r$ 给出商 $R/rR$ 的自由分解
    \[
      0\to R\xrightarrow{r}R\to R/rR\to0.
    \]
    \item 证明 $\Ext_R^0(R/rR,B)={}_rB$，其中 ${}_rB$ 是满足 $rb=0$ 的元素 $b\in B$ 的集合；证明 $\Ext_R^1(R/rR,B)\cong B/rB$，并证明对每个 $R$-模 $B$ 与 $n\geq2$，$\Ext_R^n(R/rR,B)=0$。
    \item 证明 $\Tor_0^R(A,R/rR)=A/rA$，$\Tor_1^R(A,R/rR)={}_rA$，其中 ${}_rA$ 是满足 $ra=0$ 的元素 $a\in A$ 的集合；并证明对每个 $R$-模 $A$ 与 $n\geq2$，$\Tor_n^R(A,R/rR)=0$。
  \end{enumerate}

  \item 证明
  \[
  \Tor_0^{\ZZ/m\ZZ}(A,\ZZ/d\ZZ)\cong A/dA,
  \]
  当 $n\geq1$ 为奇数时，
  \[
  \Tor_n^{\ZZ/m\ZZ}(A,\ZZ/d\ZZ)\cong {}_dA/(m/d)A,
  \]
  当 $n\geq2$ 为偶数时，
  \[
  \Tor_n^{\ZZ/m\ZZ}(A,\ZZ/d\ZZ)\cong {}_{m/d}A/dA.
  \]
  ［用命题 3 后例 2 中的投射分解。］

  \item 令 $R=k[x,y]$，其中 $k$ 是域，并令 $I=(x,y)$。
  \begin{enumerate}[label=(\alph*)]
    \item 令 $\alpha:R\to R^2$ 为 $\alpha(r)=(yr,-xr)$，并令 $\beta:R^2\to R$ 为 $\beta((r_1,r_2))=r_1x+r_2y$。证明
    \[
      0\to R\xrightarrow{\alpha}R^2\xrightarrow{\beta}R\to k\to0
    \]
    （其中 $R\to R/I=k$ 是典范投影）给出 $k$ 作为 $R$-模的自由分解。
    \item 用 (a) 中的分解证明 $\Tor_1^R(k,k)\cong k$。
    \item 证明 $\Tor_1^R(k,I)\cong k$。［用对应于短正合列 $0\to I\to R\to k\to0$ 的长正合列以及 (b)。］
    \item 由 (c) 推出无扭 $R$-模 $I$ 不是平坦的（与第 10.5 节练习 26 比较）。
  \end{enumerate}

  \item （$\Tor$ 的平坦基变换）设 $R,S$ 是交换环，$f:R\to S$ 是一个环同态，使 $S$ 成为 $R$-模，如第 10.4 节推论 12 后例 6。证明若 $S$ 作为 $R$-模是平坦的，则对所有 $R$-模 $A$ 和所有 $S$-模 $B$，有
  \[
    \Tor_n^R(A,B)\cong\Tor_n^S(S\otimes_R A,B).
  \]
  ［证明由于 $S$ 平坦，将 $A$ 的一个 $R$-模投射分解与 $S$ 张量，给出 $S\otimes_R A$ 的一个 $S$-模投射分解。］

  \item （局部化与 $\Tor$）设 $D^{-1}R$ 是交换环 $R$ 关于乘法子集 $D$ 的局部化。证明局部化与 $\Tor$ 可交换；即对所有 $R$-模 $A,B$ 与所有 $n\geq0$，有
  \[
    D^{-1}\Tor_n^R(A,B)\cong\Tor_n^{D^{-1}R}(D^{-1}A,D^{-1}B).
  \]
  ［用前一练习以及 $D^{-1}R$ 在 $R$ 上平坦这一事实，参见第 15.4 节命题 42(6)。］

  \item （平坦性是局部性质）设 $R$ 是交换环。证明 $R$-模 $M$ 平坦，当且仅当对 $R$ 的每个极大（因而每个素）理想处的局部化 $M_P$ 都是平坦 $R_P$-模。［把前一练习与用 $\Tor$ 刻画平坦性的结果结合。］

  \item 若 $R$ 是分式域为 $F$ 的整环，证明对任意 $R$-模 $B$，
  \[
    \Tor_1^R(F/R,B)\cong t(B),
  \]
  其中 $t(B)$ 表示 $B$ 的 $R$-扭子模。
\end{enumerate}

$R$-模 $M$ 称为\emph{有限表示的}，若存在 $R$-模正合列
\[
  R^s\to R^t\to M\to0
\]
对某些整数 $s,t$ 成立。等价地，$M$ 可由 $t$ 个元素生成，且相应 $R$-模同态 $R^t\to M$ 的核可由 $s$ 个元素生成。

\begin{enumerate}[label=\arabic*.,resume]
  \item \begin{enumerate}[label=(\alph*)]
    \item 证明 Noether 环 $R$ 上的每个有限生成模都是有限表示的。［用第 15.1 节练习 8。］
    \item 证明 $R$-模 $M$ 是有限表示且投射的，当且仅当对某个整数 $n\geq1$，$M$ 是 $R^n$ 的直和项。
  \end{enumerate}

  \item 假设 $M$ 是有限表示 $R$-模，且 $0\to A\to B\to M\to0$ 是 $R$-模正合列。本练习证明若 $B$ 是有限生成 $R$-模，则 $A$ 也是有限生成 $R$-模。
  \begin{enumerate}[label=(\alph*)]
    \item 设 $R^s\xrightarrow{\psi}R^t\xrightarrow{\varphi}M\to0$，且 $e_1,\ldots,e_t$ 是 $R^t$ 的一个 $R$-模基。证明存在 $b_1,\ldots,b_t\in B$，使得 $\theta(b_i)=\varphi(e_i)$，$i=1,\ldots,t$。
    \item 若 $f:R^t\to B$ 是由 $f(e_i)=b_i$ 定义的 $R$-模同态，证明 $f(\psi(R^s))\subseteq\ker\theta$。［用 $\varphi\circ\psi=0$。］推出存在一个行正合的交换图
    \[
    \begin{array}{ccccccccc}
    R^s&\xrightarrow{\psi}&R^t&\xrightarrow{\varphi}&M&\to&0\\
    \downarrow g&&\downarrow f&&\Vert&&\\
    0\to A&\xrightarrow{\alpha}&B&\xrightarrow{\theta}&M&\to&0 .
    \end{array}
    \]
    \item 证明 $A/\operatorname{image}g\cong B/\operatorname{image}f$，并用此证明 $A$ 是有限生成的。［对这个同构用练习 3 中的蛇引理。然后证明 $\operatorname{image}g$ 与 $A/\operatorname{image}g$ 都有限生成，并应用第 10.3 节练习 7。］
    \item 若 $I$ 是 $R$ 的理想，推出 $R/I$ 是有限表示 $R$-模，当且仅当 $I$ 是有限生成理想。
  \end{enumerate}

  \item 设 $R$ 是唯一极大理想为 $\mathfrak m$ 的局部环，且 $M$ 是有限表示 $R$-模。设 $m_1,\ldots,m_s$ 是 $M$ 中的元素，其在 $M/\mathfrak mM$ 中的像形成域 $R/\mathfrak m$ 上向量空间 $M/\mathfrak mM$ 的一个基。
  \begin{enumerate}[label=(\alph*)]
    \item 证明 $m_1,\ldots,m_s$ 生成 $M$ 作为 $R$-模。［用 Nakayama 引理。］
    \item 由 (a) 推出存在正合列
    \[
      0\to\ker\varphi\to R^s\xrightarrow{\varphi}M\to0
    \]
    它把 $R^s$ 的一组自由生成元映到元素 $m_1,\ldots,m_s$。推出存在正合列
    \[
      \Tor_1^R(M,R/\mathfrak m)\to(\ker\varphi)/\mathfrak m(\ker\varphi)\to0.
    \]
    ［对与 $R/\mathfrak m$ 张量使用 $\Tor$ 长正合列，并使用任意 $R$-模 $N$ 都有 $N\otimes_R R/\mathfrak m\cong N/\mathfrak mN$（第 10.4 节推论 12 后例 8），以及由 $m_1,
    \ldots,m_s$ 的选择，$\varphi:(R/\mathfrak m)^s\to M/\mathfrak mM$ 是同构这一事实。］
    \item 证明若 $\Tor_1^R(M,R/\mathfrak m)=0$，则 $m_1,
    \ldots,m_s$ 是 $M$ 的一组自由 $R$-模生成元。［用前一练习和 Nakayama 引理证明 $\ker\varphi=0$。］
  \end{enumerate}

  \item 设 $R$ 是唯一极大理想为 $\mathfrak m$ 的局部环。本练习证明有限生成 $R$-模平坦，当且仅当它自由。
  \begin{enumerate}[label=(\alph*)]
    \item 证明 $M=F/K$ 是某个有限生成自由模 $F$ 关于子模 $K\subseteq\mathfrak mF$ 的商。［取自由模 $F$ 使 $F/\mathfrak mF\cong M/\mathfrak mM$。］
    \item 假设 $x\in K$，写作 $x=a_1e_1+\cdots+a_ne_n$，其中 $e_1,
    \ldots,e_n$ 是 $F$ 的一个 $R$-基。令 $I=(a_1,
    \ldots,a_n)$ 为由 $a_1,
    \ldots,a_n$ 生成的 $R$ 的理想。证明若 $M$ 平坦，则 $I=\mathfrak mI$，并推出 $K=0$，所以 $M$ 自由。［用第 10.5 节练习 25(d) 看出 $x\in IK\subseteq\mathfrak mIF$，并推出 $I\subseteq\mathfrak mI$。然后对有限生成理想 $I$ 应用 Nakayama 引理。］
  \end{enumerate}

  \item 设 $R$ 是唯一极大理想为 $\mathfrak m$ 的局部环，$M$ 是 $R$-模，并考虑下列陈述：
  \begin{enumerate}[label=(\roman*)]
    \item $M$ 是自由 $R$-模；
    \item $M$ 是投射 $R$-模；
    \item $M$ 是平坦 $R$-模；
    \item $\Tor_1^R(M,R/\mathfrak m)=0$。
  \end{enumerate}
  \begin{enumerate}[label=(\alph*)]
    \item 证明 (i) 蕴含 (ii)，(ii) 蕴含 (iii)，(iii) 蕴含 (iv)。
    \item 若 $M$ 有限生成，证明 (i)、(ii)、(iii) 等价。（下面练习 34 说明若 $M$ 有限生成但 $R$ 不是局部的，则 (iii) 未必蕴含 (i) 或 (ii)。）［用前一练习。］
    \item 若 $M$ 有限表示，证明 (i)、(ii)、(iii)、(iv) 等价。（下面练习 35 说明若 $M$ 有限生成但不是有限表示，则 (iv) 未必蕴含 (i)、(ii) 或 (iii)。）［用练习 28。］
  \end{enumerate}
  注：Kaplansky 的一个定理（参见 Projective Modules, Annals of Mathematics, 68 (1958), pp. 372--377）说明，在不要求 $M$ 有限生成的条件下，(i) 与 (ii) 等价。

  \item （有限表示模的局部化与 Hom）设 $D^{-1}R$ 是交换环 $R$ 关于乘法子集 $D$ 的局部化，并令 $M$ 是有限表示 $R$-模。
  \begin{enumerate}[label=(\alph*)]
    \item 对任意 $R$-模 $A,B$，证明存在唯一的 $D^{-1}R$-模同态
    \[
      D^{-1}\Hom_R(A,B)\to\Hom_{D^{-1}R}(D^{-1}A,D^{-1}B)
    \]
    它把 $\varphi\in\Hom_R(A,B)$ 映到由 $\varphi$ 诱导的从 $D^{-1}A$ 到 $D^{-1}B$ 的同态。
    \item 对任意 $R$-模 $N$ 与任意 $m\geq1$，证明 $\Hom_R(R^m,N)\cong N^m$ 作为 $R$-模；并推出 $D^{-1}\Hom_R(R^m,N)\cong(D^{-1}N)^m$ 作为 $D^{-1}R$-模。
    \item 假设 $R^s\to R^t\to M\to0$ 正合。证明存在一个行正合的 $D^{-1}R$-模交换图
    {\setlength{\arraycolsep}{2pt}
    \[
    \scalebox{0.82}{$
    \begin{array}{cccccc}
    0\to D^{-1}\Hom_R(M,N)&\to&D^{-1}\Hom_R(R^t,N)&\to&D^{-1}\Hom_R(R^s,N)\\
    \downarrow&&\downarrow&&\downarrow\\
    0\to\Hom_{D^{-1}R}(D^{-1}M,D^{-1}N)&\to&\Hom_{D^{-1}R}((D^{-1}R)^t,D^{-1}N)&\to&\Hom_{D^{-1}R}((D^{-1}R)^s,D^{-1}N).
    \end{array}
    $}
    \]}
    ［第一行先从 $M$ 的表示中各项取到 $N$ 中的 $R$-模同态，使用第 10.5 节定理 33（注意证明中的第一个注释），然后与平坦 $R$-模 $D^{-1}R$ 张量，参见第 15.4 节命题 41 与 42(6)。第二行先将表示与 $D^{-1}R$ 张量，然后取到 $D^{-1}N$ 中的 $D^{-1}R$-模同态。］
    \item 用 (b) 证明，当 $M$ 有限表示时，局部化与取同态可交换；即作为 $D^{-1}R$-模，
    \[
      D^{-1}\Hom_R(M,N)\cong\Hom_{D^{-1}R}(D^{-1}M,D^{-1}N).
    \]
    ［证明图中后两个竖直映射是同构，并推出左边竖直映射也是同构。］（若 $M$ 不是有限表示的，此结果一般不成立。）
  \end{enumerate}

  \item （有限表示模的局部化与 Ext）设 $D^{-1}R$ 是交换环 $R$ 关于乘法子集 $D$ 的局部化。证明若 $M$ 是有限表示 $R$-模，则对每个 $R$-模 $N$ 与每个 $n\geq0$，作为 $D^{-1}R$-模有
  \[
    D^{-1}\Ext_R^n(M,N)\cong\Ext_{D^{-1}R}^n(D^{-1}M,D^{-1}N).
  \]
  ［用 $N$ 的投射分解和前一练习，并注意将分解与 $D^{-1}R$ 张量给出 $D^{-1}R$-模 $D^{-1}N$ 的一个投射分解。］

  \item 设 $R$ 是交换环，$M$ 是有限表示 $R$-模（例如 Noether 环上的有限生成模，或由有限生成理想 $I$ 给出的商 $R/I$，参见练习 26 与 27）。证明下列条件等价：
  \begin{enumerate}[label=(\alph*)]
    \item $M$ 是投射 $R$-模；
    \item $M$ 是平坦 $R$-模；
    \item $M$ 局部自由，即对 $R$ 的每个极大（因而每个素）理想 $P$，局部化 $M_P$ 是自由 $R_P$-模。
  \end{enumerate}
  特别地，证明有限生成投射模与有限表示平坦模相同。［练习 24 与 30 说明 (b) 等价于 (c)。用投射性的 $\Ext$ 判别法以及练习 30 与 32，看出 (a) 等价于 (c)。］

  \item \begin{enumerate}[label=(\alph*)]
    \item 证明交换环 $R$ 的每个 $R$-模都是平坦的，当且仅当 $R$ 的每个有限生成理想 $I$ 都是 $R$ 的直和项；在这种情形下，$R$ 的每个有限生成理想都是主理想且投射（这样的环称为\emph{绝对平坦}）。［用练习 15、前一练习应用于有限表示 $R$-模 $R/I$，以及命题 16 后的评论。］
    \item 证明每个 Boolean 环都是绝对平坦的。［用第 7.4 节练习 24，注意若 $I=Rx$，则 $x$ 是幂等元，所以 $R=Rx\oplus R(1-x)$。］
    \item 令 $R$ 是可数多个 $\ZZ/2\ZZ$ 的直积，$I$ 是可数多个 $\ZZ/2\ZZ$ 的直和。证明 $I$ 是 Boolean 环 $R$ 的一个理想但不是有限生成的，并且循环 $R$-模 $M=R/I$ 是平坦但不是投射的（所以有限生成平坦模未必投射）。
  \end{enumerate}

  \item 令 $R$ 为开区间 $(-1,1)$ 上 $C^\infty$ 函数环在由 $x=0$ 处为 $0$ 的函数构成的极大理想处的局部化所得的局部环（参见第 15.2 节练习 45），令 $\mathfrak m=(x)$ 是 $R$ 的唯一极大理想，并令 $P=\bigcap_{n\geq1}\mathfrak m^n$ 是素理想。置 $M=R/P$。
  \begin{enumerate}[label=(\alph*)]
    \item 证明 $\Tor_1^R(M,R/\mathfrak m)=0$。［用练习 19，取 $r=x$，并注意 $R/P$ 是整环。］
    \item 证明 $M$ 不是平坦的（因而也不是投射的）。［令 $F$ 如第 15.2 节练习 45 中所述。证明由乘以 $F$ 诱导的序列 $0\to R\to R\to R/(F)\to0$ 正合，但与 $M$ 张量后不再正合。］
  \end{enumerate}
\end{enumerate}

\section{群的上同调}

本节中，我们考虑上一节一般技巧在一个重要特例中的应用。

令 $G$ 为群。

\noindent\textbf{定义.} 若 Abel 群 $A$ 上有 $G$ 作为自同构的（左）作用，则称 $A$ 为一个 $G$-模。

注意，$G$-模等同于一个 Abel 群 $A$ 以及一个从 $G$ 到 $A$ 的自同构群的同态 $\varphi:G\to\operatorname{Aut}(A)$。由于 Abel 群等同于 $\ZZ$ 上的模，也容易看出 $G$-模 $A$ 等同于整数群环 $\ZZ G$ 上的模。当 $G$ 是无限群时，环 $\ZZ G$ 由 $G$ 的元素的有限形式和组成，系数在 $\ZZ$ 中。

通常我们常用乘法记号，并把元素 $g\in G$ 对 $a\in A$ 的作用 $g\cdot a$ 简写为 $ga$。

\noindent\textbf{定义.} 若 $A$ 是 $G$-模，令
\[
  A^G=\{a\in A\mid ga=a\text{ 对所有 }g\in G\}
\]
为 $A$ 中被 $G$ 的所有元素固定的元素。

\paragraph{例}
\begin{enumerate}[label=(\arabic*)]
  \item 若对所有 $a\in A$ 与 $g\in G$ 都有 $ga=a$，则称 $G$ 在 $A$ 上平凡作用。在这种情形下 $A^G=A$。除非另有说明，Abel 群 $\ZZ$ 总是假定对任意群 $G$ 都带有平凡 $G$-作用。
  \item 对任意 $G$-模 $A$，$A$ 在 $G$ 作用下的不动点 $A^G$ 显然是 $A$ 的 $\ZZ G$-子模，并且 $G$ 在其上平凡作用。
  \item 若 $V$ 是域 $F$ 上维数为 $n$ 的向量空间，且 $G=\operatorname{GL}_n(F)$，则 $V$ 自然是 $G$-模。在此情形下 $V^G=\{0\}$，因为 $V$ 的任意非零元素都可以被某个线性变换送到 $V$ 的任意另一个非零元素。
  \item 如第 5.5 节中的半直积 $E=A\rtimes G$，在 $A$ 是 Abel 正规子群的情形下给出一个 $G$-模 $A$，其中 $G$ 的作用由同态 $\varphi:G\to\operatorname{Aut}(A)$ 给出。子群 $A^G$ 由 $A$ 中位于 $E$ 的中心里的元素组成。更一般地，若 $A$ 是群 $E$ 的任意 Abel 正规子群，则 $E$ 通过共轭作用于 $A$，这使 $A$ 成为 $E$-模，也成为 $E/A$-模。在此情形下 $A^E=A^{E/A}$ 也由 $A$ 中位于 $E$ 的中心里的元素组成。
  \item 若 $K/F$ 是 Galois 扩张，其 Galois 群为 $G$，则加法群 $K$ 自然是 $G$-模，且 $K^G=F$。类似地，$K$ 中非零元素的乘法群 $K^\times$ 是 $G$-模，其不动点为 $(K^\times)^G=F^\times$。
\end{enumerate}

最后一个例子中的不动点子群在第 14 章的 Galois 理论中起着中心作用。一般地，容易看出 $G$-模短正合列
\[
  0\to A\to B\to C\to0
\]
诱导正合列
\[
  0\to A^G\to B^G\to C^G,
  \tag{17.15}
\]
但它一般不能延伸为短正合列（一般地，商 $C$ 中被 $G$ 固定的一个陪集未必能由 $B$ 中被 $G$ 固定的元素表示）。看出 (17.15) 正合的一种方法，是把 $A^G$ 与一个 $\Hom$ 群联系起来：

\medskip
\noindent\textbf{引理 19.} 假设 $A$ 是 $G$-模，且 $\Hom_{\ZZ G}(\ZZ,A)$ 是从带平凡 $G$-作用的 $\ZZ$ 到 $A$ 的所有 $\ZZ G$-模同态组成的群。则
\[
  A^G\cong\Hom_{\ZZ G}(\ZZ,A).
\]

\begin{proof}
任意 $G$-模同态 $\alpha$ 从 $\ZZ$ 到 $A$ 都由其在 $1$ 上的取值唯一确定。令 $\alpha_a$ 表示满足 $\alpha_a(1)=a$ 的 $G$-模同态。由于 $\alpha_a$ 是 $G$-模同态，
\[
  a=\alpha_a(1)=\alpha_a(g\cdot1)=g\alpha_a(1)=g\cdot a
\]
对所有 $g\in G$ 成立，所以 $a$ 必须属于 $A^G$。反过来，对任意 $a\in A^G$，容易验证映射 $\alpha_a\mapsto a$ 给出从 $\Hom_{\ZZ G}(\ZZ,A)$ 到 $A^G$ 的同构。
\end{proof}

结合上一节的结果，该引理不仅说明序列 (17.15) 正合；它还说明，把 $\ZZ$ 看作 $\ZZ G$-模时，$\ZZ$ 的任意投射分解都会给出延长 (17.15) 的长正合列。这样的一个投射分解是 $\ZZ$ 的\emph{标准分解}或\emph{bar 分解}：
\[
  \cdots\to F_n\xrightarrow{d_n}F_{n-1}\to\cdots\to F_0\xrightarrow{\operatorname{aug}}\ZZ\to0.
  \tag{17.16}
\]
这里对 $n\geq0$，
\[
  F_n=\ZZ G\otimes_{\ZZ}\ZZ G\otimes_{\ZZ}\cdots\otimes_{\ZZ}\ZZ G
\]
（共有 $n+1$ 个因子），它是 $G$-模，作用在纯张量上定义为
\[
  g\cdot(g_0\otimes g_1\otimes\cdots\otimes g_n)
  =(gg_0)\otimes g_1\otimes\cdots\otimes g_n.
\]
不难看出 $F_n$ 是秩为 $|G|^n$ 的自由 $\ZZ G$-模，其 $\ZZ G$-基由元素
\[
  1\otimes g_1\otimes g_2\otimes\cdots\otimes g_n,
  \qquad g_i\in G,
\]
给出。映射 $\operatorname{aug}:F_0\to\ZZ$ 是增广映射
\[
  \operatorname{aug}\left(\sum_{g\in G}a_gg\right)=\sum_{g\in G}a_g,
\]
而映射 $d_1$ 由 $d_1(1\otimes g)=g-1$ 给出。$n\geq2$ 时映射 $d_n$ 更复杂，它们的定义以及 (17.16) 是投射（事实上自由）分解的证明见练习 1--3。

如上一节那样，从 (17.16) 的各项取到 $G$-模 $A$ 中的 $\ZZ G$-模同态（并把第一项替换为 $0$），得到上链复形
\[
  0\to\Hom_{\ZZ G}(F_0,A)\xrightarrow{d_1}\Hom_{\ZZ G}(F_1,A)\xrightarrow{d_2}\Hom_{\ZZ G}(F_2,A)\to\cdots,
  \tag{17.17}
\]
其上同调群按定义就是群 $\Ext_{\ZZ G}^n(\ZZ,A)$。于是，如定理 8 所示，$G$-模短正合列 $0\to A\to B\to C\to0$ 给出一个长正合列，其最初的项由 (17.15) 给出，而高阶项是上同调群 $\Ext_{\ZZ G}^n(\ZZ,A)$。

为了更显式地描述它，可以不直接引用 $\ZZ$ 的标准分解，而重新解释这个上链复形中的各项。$\Hom_{\ZZ G}(F_n,A)$ 的元素由它们在 $F_n$ 的 $\ZZ G$-基元素上的取值唯一确定，而这些基元素可与 $G$ 中元素的 $n$ 元组 $(g_1,g_2,
\ldots,g_n)$ 识别。因此对 $n\geq1$，群 $\Hom_{\ZZ G}(F_n,A)$ 可与从 $G\times\cdots\times G$（$n$ 个因子）到 $A$ 的函数集合识别。对 $n=0$，我们把 $\Hom_{\ZZ G}(\ZZ G,A)$ 与 $A$ 识别。

\noindent\textbf{定义.} 若 $G$ 是有限群且 $A$ 是 $G$-模，定义 $C^0(G,A)=A$；对 $n\geq1$，定义 $C^n(G,A)$ 为从 $G^n=G\times\cdots\times G$（$n$ 个因子）到 $A$ 的所有映射的集合。$C^n(G,A)$ 的元素称为（$G$ 中取值于 $A$ 的）$n$-上链。

每个 $C^n(G,A)$ 都是加法 Abel 群：对 $C^0(G,A)=A$，其群结构由 $A$ 给出；对 $n\geq1$，由函数逐点相加给出：
\[
  (f_1+f_2)(g_1,g_2,
  \ldots,g_n)=f_1(g_1,g_2,
  \ldots,g_n)+f_2(g_1,g_2,
  \ldots,g_n).
\]
在 $\Hom_{\ZZ G}(F_n,A)$ 与 $C^n(G,A)$ 的识别下，(17.17) 中的上链映射 $d_n$ 可以非常显式地给出（也参见练习 3 及其后的评论）：

\noindent\textbf{定义.} 对 $n\geq0$，定义第 $n$ 个\emph{上边界同态} $d_n:C^n(G,A)\to C^{n+1}(G,A)$ 为
\[
\begin{aligned}
(d_nf)(g_1,
\ldots,g_{n+1})={}&g_1\cdot f(g_2,
\ldots,g_{n+1})\\
&+\sum_{i=1}^{n}(-1)^i f(g_1,
\ldots,g_ig_{i+1},
\ldots,g_{n+1})\\
&+(-1)^{n+1}f(g_1,
\ldots,g_n),
\end{aligned}
\tag{17.18}
\]
其中 $f$ 的第 $i$ 个位置中的积 $g_ig_{i+1}$ 在群 $G$ 中相乘。

由定义立刻可知映射 $d_n$ 是群同态。由 (17.17) 是投射分解这一事实，可知 $d_n\circ d_{n-1}=0$，$n\geq1$（也可以只从上面对 $d_n$ 的定义给出一个自含的直接证明，但相当繁琐）。

\noindent\textbf{定义.}
\begin{enumerate}[label=(\arabic*)]
  \item 对 $n\geq0$，令 $Z^n(G,A)=\ker d_n$。$Z^n(G,A)$ 的元素称为 $n$-上循环。
  \item 对 $n\geq1$，令 $B^n(G,A)=\operatorname{image}d_{n-1}$，并令 $B^0(G,A)=1$。$B^n(G,A)$ 的元素称为 $n$-上边界。
\end{enumerate}

由于 $n\geq1$ 时 $d_n\circ d_{n-1}=0$，有 $\operatorname{image}d_{n-1}\subseteq\ker d_n$，所以 $B^n(G,A)$ 总是 $Z^n(G,A)$ 的子群。

\noindent\textbf{定义.} 对任意 $G$-模 $A$，商群
\[
  Z^n(G,A)/B^n(G,A)
\]
称为 $G$ 以 $A$ 为系数的第 $n$ 个上同调群，记作 $H^n(G,A)$，$n\geq0$。

用上链定义 $H^n(G,A)$ 在下面两节中研究低维群 $H^1(G,A)$、$H^2(G,A)$ 及其各种应用时尤其有用。不过应当记住，对所有 $n\geq0$，有 $H^n(G,A)\cong\Ext_{
\ZZ G}^n(\ZZ,A)$。特别地，这些群可以用 $\ZZ$ 的任意投射分解来计算。

\paragraph{例}
\begin{enumerate}[label=(\arabic*)]
  \item 对 $f=a\in C^0(G,A)$，有 $d_0(f)(g)=g\cdot a-a$，所以 $\ker d_0$ 是集合 $\{a\in A\mid g\cdot a=a\text{ 对所有 }g\in G\}$，即 $Z^0(G,A)=A^G$。因此对任意群 $G$ 与 $G$-模 $A$，
  \[
    H^0(G,A)=A^G.
  \]

  \item 设 $G=1$ 是平凡群。则 $G^n=\{(1,1,
  \ldots,1)\}$ 也是平凡的，所以 $f\in C^n(G,A)$ 完全由 $f(1,1,
  \ldots,1)=a\in A$ 决定。把 $f$ 与 $a$ 识别，得到 $C^n(G,A)=A$，对所有 $n\geq0$。于是若 $f=a\in A$，则
  \[
    d_n(f)(1,1,
    \ldots,1)=a+\sum_{i=1}^{n}(-1)^ia+(-1)^{n+1}a
    =\begin{cases}
      0,& n\text{ 为偶数},\\
      a,& n\text{ 为奇数}.
    \end{cases}
  \]
  所以当 $n$ 为偶数时 $d_n=0$，当 $n$ 为奇数时 $d_n=1$ 是恒等映射。因此
  \[
    H^0(1,A)=A^0=A,
    \qquad
    H^n(1,A)=0\quad(n\geq1).
  \]
\end{enumerate}

\paragraph{例：有限循环群的上同调}
设 $G$ 是由 $\sigma$ 生成的 $m$ 阶循环群。令
\[
  N=1+\sigma+\sigma^2+\cdots+\sigma^{m-1}\in\ZZ G.
\]
则 $N(\sigma-1)=(\sigma-1)N=\sigma^m-1=0$，于是有一个特别简单的自由分解
\[
  \cdots\xrightarrow{\sigma-1}\ZZ G\xrightarrow{N}\ZZ G\xrightarrow{\sigma-1}\ZZ G\xrightarrow{N}\ZZ G\xrightarrow{\operatorname{aug}}\ZZ\to0
\]
其中 $\operatorname{aug}$ 表示增广映射（参见练习 8）。从这个分解的各项取到 $A$ 的 $\ZZ G$-模同态（并把第一项替换为 $0$），再利用同构 $\Hom_{
\ZZ G}(\ZZ G,A)=A$，得到上链复形
\[
  0\to A\xrightarrow{\sigma-1}A\xrightarrow{N}A\xrightarrow{\sigma-1}A\xrightarrow{N}\cdots .
\]
它的上同调计算群 $H^n(G,A)$：
\[
  H^0(G,A)=A^G,
\]
且
\[
H^n(G,A)=
\begin{cases}
A^G/NA,& n\geq2\text{ 为偶数},\\
{}_NA/(\sigma-1)A,& n\geq1\text{ 为奇数},
\end{cases}
\]
其中 ${}_NA=\{a\in A\mid Na=0\}$ 是被 $N$ 零化的 $A$ 的子群，因为乘以 $\sigma-1$ 的核是 $A^G$。

特别地，若 $G=(\sigma)$ 在 $A$ 上平凡作用，则 $N\cdot a=ma$，所以在这种情形下
\[
  H^0(G,A)=A,
\]
并且对偶数 $n\geq2$，$H^n(G,A)=A/mA$；对奇数 $n\geq1$，$H^n(G,A)={}_mA$，即 $A$ 中阶整除 $m$ 的元素。再特化为 $m=1$，就得到前面的例 2。

\medskip
\noindent\textbf{命题 20.} 假设对某个整数 $m\geq1$ 有 $mA=0$（即 $G$-模 $A$ 作为 Abel 群的指数整除 $m$）。则对所有 $n\geq0$，
\[
  mZ^n(G,A)=mB^n(G,A)=mH^n(G,A)=0.
\]
特别地，若 $A$ 的指数为素数 $p$，则 Abel 群 $Z^n(G,A)$、$B^n(G,A)$ 与 $H^n(G,A)$ 的指数都整除 $p$，从而这些群都是有限域 $\mathbb F_p=\ZZ/p\ZZ$ 上的向量空间。

\begin{proof}
若 $f\in C^n(G,A)$ 是一个 $n$-上链，则 $f\in A$（若 $n=0$），此时 $mf=0$；或 $f$ 是从 $G^n$ 到 $A$ 的函数（若 $n\geq1$），此时 $mf$ 是从 $G^n$ 到 $mA=0$ 的函数，所以仍有 $mf=0$。因此 $mZ^n(G,A)=mB^n(G,A)=0$，因为它们是 $C^n(G,A)$ 的子群。于是 $mH^n(G,A)=0$，因为 $mZ^n(G,A)=0$；命题中其余断言立即得到。
\end{proof}

由例 1，定理 10 中的长正合列用上同调群 $H^n(G,A)$ 表示为：

\medskip
\noindent\textbf{定理 21.} （群上同调长正合列）假设
\[
  0\to A\to B\to C\to0
\]
是 $G$-模短正合列。则存在长正合列
\[
0\to A^G\to B^G\to C^G\xrightarrow{\delta_0}H^1(G,A)\to H^1(G,B)\to H^1(G,C)\xrightarrow{\delta_1}\cdots
\]
\[
\cdots\xrightarrow{\delta_{n-1}}H^n(G,A)\to H^n(G,B)\to H^n(G,C)\xrightarrow{\delta_n}H^{n+1}(G,A)\to\cdots
\]
的 Abel 群长正合列。

定理 21 中长正合列的许多用途之一，是一种称为\emph{维数平移}的技巧：它可以把 $A$ 的维数 $n+1$ 的上同调群 $H^{n+1}(G,A)$，转化为另一个 $G$-模的维数 $n$ 的上同调群。这个技巧基于寻找一个几乎所有上同调群都为零的 $G$-模。这样的模有一个名称：

\noindent\textbf{定义.} 若 $G$-模 $M$ 满足对所有 $n\geq1$ 都有 $H^n(G,M)=0$，则称 $M$ 对 $G$ 是\emph{上同调平凡的}。

\medskip
\noindent\textbf{推论 22.} （维数平移）假设 $0\to A\to M\to C\to0$ 是 $G$-模短正合列，且 $M$ 对 $G$ 是上同调平凡的。则存在正合列
\[
  0\to A^G\to M^G\to C^G\to H^1(G,A)\to0,
\]
并且
\[
  H^n(G,C)\cong H^{n+1}(G,A)
  \qquad(n\geq1).
\]

\begin{proof}
由于 $M$ 对 $G$ 上同调平凡，定理 21 中长正合列的一段
\[
  H^n(G,M)\to H^n(G,C)\to H^{n+1}(G,A)\to H^{n+1}(G,M)
\]
化为
\[
  0\to H^n(G,C)\to H^{n+1}(G,A)\to0,
\]
这说明 $n\geq1$ 时 $H^n(G,C)\cong H^{n+1}(G,A)$。类似地，定理 21 中长正合列最初的部分给出推论中的第一个断言。
\end{proof}

现在我们指出一个自然构造：由 $G$ 的子群 $H$ 上的一个模产生一个 $G$-模。当 $H=1$ 是平凡群时，这个构造对任意 $G$-模 $A$ 产生一个上同调平凡模 $M$ 以及推论 22 中的一个正合列。

\noindent\textbf{定义.} 若 $H$ 是 $G$ 的子群，且 $A$ 是 $H$-模，定义诱导 $G$-模
\[
  M_H^G(A)=\Hom_{\ZZ H}(\ZZ G,A).
\]
换言之，$M_H^G(A)$ 是所有满足 $f(hx)=hf(x)$（对每个 $x\in G$ 与 $h\in H$）的从 $G$ 到 $A$ 的映射 $f$ 的集合。

元素 $g\in G$ 对 $f\in M_H^G(A)$ 的作用由
\[
  (g\cdot f)(x)=f(xg),
  \qquad x\in G
\]
给出（参见第 10.5 节练习 10）。

回忆若 $H$ 是 $G$ 的子群且 $A$ 是 $H$-模，则由从 $\ZZ H$ 到 $\ZZ G$ 的标量扩张所得的模 $\ZZ G\otimes_{\ZZ H}A$ 是 $G$-模。对有限群 $G$，或更一般地当 $H$ 在 $G$ 中指数有限时，有
\[
  M_H^G(A)\cong\ZZ G\otimes_{\ZZ H}A
\]
（参见练习 10）。当 $G$ 无限时，这未必成立（参见练习 11）。模 $\ZZ G\otimes_{\ZZ H}A$ 有时称为诱导 $G$-模，而模 $M_H^G(A)$ 有时称为余诱导 $G$-模。对于有限群，张量积结合性说明，若 $K\leq H\leq G$，则 $M_H^G(M_K^H(A))=M_K^G(A)$；一般情形下同样结果也成立（由定义结合练习 7 得到）。

\paragraph{例}
\begin{enumerate}[label=(\arabic*)]
  \item 若 $H$ 是 $G$ 的子群，且 $0\to A\to B\to C\to0$ 是 $H$-模短正合列，则
  \[
    0\to M_H^G(A)\to M_H^G(B)\to M_H^G(C)\to0
  \]
  是 $G$-模短正合列，因为 $M_H^G(A)\cong\ZZ G\otimes_{\ZZ H}A$，且 $\ZZ G$ 在 $\ZZ H$ 上自由，因而平坦。
  \item 当 $G$ 有限且 $A$ 是平凡 $H$-模 $\ZZ$ 时，模 $M_H^G(\ZZ)$ 是秩为 $m=|G:H|$ 的自由 $\ZZ$-模。存在一个基 $b_1,
  \ldots,b_m$，使得 $G$ 置换这些基元素的方式与它左乘置换 $G$ 中 $H$ 的左陪集相同；即若令 $b_i\leftrightarrow g_iH$，则 $gb_i=b_j$ 当且仅当 $gg_iH=g_jH$。模 $M_H^G(\ZZ)$ 是 $G$ 在 $\ZZ$ 上的置换模，其稳定子为 $H$。一个有意义的特例是 $G=S_m$ 且 $H=S_{m-1}$，其中 $S_m$ 如通常那样置换 $\{1,2,
  \ldots,m\}$。域上的置换模与诱导模将在第 VI 部分中研究。
  \item 当 $H=1$ 是平凡群时，任意 Abel 群 $A$ 都是 $H$-模。相应的诱导 $G$-模 $M_1^G(A)$ 就是从 $G$ 到 $A$ 的所有映射 $f$ 的集合。对 $g\in G$，映射 $g\cdot f\in M_1^G(A)$ 满足 $(g\cdot f)(x)=f(xg)$，$x\in G$。
  \item 假设 $A$ 是 $G$-模。则有自然映射
  \[
    \varphi:A\to M_1^G(A)
  \]
  从 $A$ 到前一例中的诱导 $G$-模，它把 $a\in A$ 映到函数 $f_a$，其中 $f_a(x)=xa$，$x\in G$。显然 $\varphi$ 是群同态，并且
  \[
    f_{ga}(x)=x(ga)=(xg)a=f_a(xg)=(g\cdot f_a)(x)
  \]
  说明 $\varphi$ 也是 $G$-模同态。由于 $f_a(1)=a$，$f_a$ 是 $G$ 上零函数当且仅当 $a=0$，所以 $\varphi$ 是单射。因此可以把 $A$ 识别为诱导模 $M_1^G(A)$ 的一个 $G$-子模。
  \item 更一般地，若 $A$ 是 $G$-模而 $H$ 是 $G$ 的任意子群，则前一例中的函数 $f_a(x)$ 属于子群 $M_H^G(A)$，因为对所有 $h\in H$ 有
  \[
    f_a(hx)=(hx)(a)=h(xa)=hf_a(x).
  \]
  因此从 $A$ 到 $M_H^G(A)$ 的相应映射是单射 $G$-模同态。
  \item 不动点 $(M_H^G(A))^G$ 是满足对所有 $g\in G$ 都有 $gf=f$ 的从 $G$ 到 $A$ 的映射 $f$，即对所有 $g,x\in G$ 有 $(gf)(x)=f(x)$。由 $G$ 在 $M_H^G(A)$ 上作用的定义，这是方程 $f(xg)=f(x)$ 对所有 $g,x\in G$ 成立。取 $x=1$，说明 $f$ 在整个 $G$ 上为常值：$f(g)=f(1)=a\in A$。常值函数 $f=a$ 是 $M_H^G(A)$ 的元素，当且仅当对所有 $h\in H$ 有 $a=f(hx)=hf(x)=ha$，所以 $(M_H^G(A))^G\cong A^H$。前一例中的元素 $f_a(x)$ 属于子群 $(M_H^G(A))^G$，当且仅当 $xa$ 对 $x\in G$ 是常值，也就是当且仅当 $a\in A^G$。
\end{enumerate}

诱导自 $H$-模 $A$ 的 $G$-模 $M_H^G(A)$ 的重要性质之一是：它相对于 $G$ 的上同调与 $A$ 相对于 $H$ 的上同调相同。

\medskip
\noindent\textbf{命题 23.} （Shapiro 引理）对 $G$ 的任意子群 $H$ 与任意 $H$-模 $A$，有
\[
  H^n(G,M_H^G(A))\cong H^n(H,A),
  \qquad n\geq0.
\]

\begin{proof}
令 $\cdots\to P_n\to\cdots\to P_0\to\ZZ\to0$ 为 $\ZZ$ 的一个由投射 $G$-模组成的分解（例如标准分解）。群 $H^n(G,M_H^G(A))$ 通过取从此分解到 $M_H^G(A)=\Hom_{\ZZ H}(\ZZ G,A)$ 的同态计算。由于 $\ZZ G$ 是自由 $\ZZ H$-模，这个 $G$-模分解也是 $\ZZ$ 的一个由投射 $H$-模组成的分解，因此对其取到 $A$ 中的同态也可用同一分解计算群 $H^n(H,A)$。为看出这两族上同调群同构，使用 Abel 群的自然同构
\[
  \Phi:\Hom_{\ZZ G}(P_n,\Hom_{\ZZ H}(\ZZ G,A))\cong\Hom_{\ZZ H}(P_n,A),
\]
其中
\[
  \Phi(f)(p)=f(p)(1)
\]
对所有 $f\in\Hom_{\ZZ G}(P_n,\Hom_{\ZZ H}(\ZZ G,A))$ 与 $p\in P_n$ 成立。逆同构定义为：对 $f'\in\Hom_{\ZZ H}(P_n,A)$ 与 $p\in P_n$，令 $\Psi(f')(p)$ 是从 $\ZZ G$ 到 $A$ 的映射，它把 $g\in G$ 送到 $f'(gp)$；即
\[
  (\Psi(f')(p))(g)=f'(gp).
\]
注意这是良定义的，因为 $P_n$ 是 $G$-模。（这些映射是伴随结合定理的一个特例，参见练习 7。）由于这些同构与上链映射交换，它们在相应上同调群上诱导同构，即 $H^n(G,M_H^G(A))\cong H^n(H,A)$，如所需。
\end{proof}

\medskip
\noindent\textbf{推论 24.} 对任意 $G$-模 $A$，模 $M_1^G(A)$ 对 $G$ 是上同调平凡的；即对所有 $n\geq1$，
\[
  H^n(G,M_1^G(A))=0.
\]

\begin{proof}
这由命题取 $H=1$ 以及命题 20 前例 2 中对平凡群上同调的计算立刻得到。
\end{proof}

由该推论，上面的例 4 给出 $G$-模短正合列
\[
  0\to A\to M\to C\to0,
\]
其中 $M=M_1^G(A)$ 对 $G$ 是上同调平凡的，而 $C$ 是 $M_1^G(A)$ 模去 $A$ 的像所得的商。推论 22 中的维数平移结果于是成为：

\medskip
\noindent\textbf{推论 25.} 对任意 $G$-模 $A$，
\[
  H^{n+1}(G,A)\cong H^n(G,M_1^G(A)/A)
\]
对所有 $n\geq1$ 成立。

下面考虑若干重要的同态，它们联系各种上同调群。这些同态的一些应用将在随后两节中出现。

一般地，假设有两个群 $G$ 与 $G'$，并且 $A$ 是 $G$-模，$A'$ 是 $G'$-模。若 $\varphi:G'\to G$ 是群同态，则通过定义 $g'\cdot a=\varphi(g')a$（$g'\in G'$，$a\in A$），$A$ 成为 $G'$-模。若现在 $\psi:A\to A'$ 是 Abel 群同态，则我们考虑 $\psi$ 是否为 $G'$-模同态。

\noindent\textbf{定义.} 假设 $A$ 是 $G$-模，$A'$ 是 $G'$-模。群同态 $\varphi:G'\to G$ 与 $\psi:A\to A'$ 称为\emph{相容的}，若当 $A$ 通过 $\varphi$ 成为 $G'$-模时，$\psi$ 是 $G'$-模同态；即对所有 $g'\in G'$ 与 $a\in A$，
\[
  \psi(\varphi(g')a)=g'\psi(a).
\]

相容同态的意义在于它们诱导相关上同调群上的群同态。若 $\varphi:G'\to G$ 与 $\psi:A\to A'$ 是同态，则 $\varphi$ 诱导同态 $\varphi_n:(G')^n\to G^n$，从而给出从 $C^n(G,A)$ 到 $C^n(G',A)$ 的同态，把 $f$ 送到 $f\circ\varphi_n$。映射 $\psi$ 诱导从 $C^n(G',A)$ 到 $C^n(G',A')$ 的同态，把 $f$ 送到 $\psi\circ f$。合起来得到诱导同态
\[
  \lambda_n:C^n(G,A)\to C^n(G',A'),
  \qquad f\mapsto \psi\circ f\circ\varphi_n.
\]
若此外 $\varphi$ 与 $\psi$ 相容，则容易验证诱导映射 $\lambda_n$ 与上边界算子交换：
\[
  \lambda_{n+1}\circ d_n=d_n\circ\lambda_n
\]
对所有 $n\geq0$ 成立。因此 $\lambda_n$ 将上循环映为上循环，将上边界映为上边界，从而在上同调上诱导群同态
\[
  \lambda_n:H^n(G,A)\to H^n(G',A'),
  \qquad n\geq0.
\]

我们考虑这种映射的若干情形。

\paragraph{例}
\begin{enumerate}[label=(\arabic*)]
  \item 假设 $G=G'$ 且 $\varphi$ 是恒等映射。那么群同态 $\psi:A\to A'$ 与 $\varphi$ 相容，正是说 $\psi$ 是 $G$-模同态。因此任意 $G$-模同态 $A\to A'$ 都诱导群同态
  \[
    H^n(G,A)\to H^n(G,A'),
    \qquad n\geq0.
  \]
  特别地，若 $0\to A\to B\to C\to0$ 是 $G$-模短正合列，则对 $n\geq0$，得到从 $H^n(G,A)$ 到 $H^n(G,B)$ 以及从 $H^n(G,B)$ 到 $H^n(G,C)$ 的诱导同态。这些正是定理 21 长正合列中的同态。

  \item （限制同态）若 $A$ 是 $G$-模，则对 $G$ 的任意子群 $H$，$A$ 也是 $H$-模。包含映射 $\varphi:H\to G$ 与恒等映射 $\psi:A\to A$ 是相容同态。上同调上的相应诱导群同态称为\emph{限制同态}：
  \[
    \operatorname{Res}:H^n(G,A)	o H^n(H,A),
    \qquad n\geq0.
  \]
  这个术语来自这样一个事实：从 $C^n(G,A)$ 到 $C^n(H,A)$ 的上链映射，只是把从 $G^n$ 到 $A$ 的映射 $f$ 限制到 $G^n$ 的子群 $H^n$。

  \item （膨胀同态）假设 $H$ 是 $G$ 的正规子群，且 $A$ 是 $G$-模。被 $H$ 固定的元素 $A^H$ 在作用 $(gH)\cdot a=g\cdot a$ 下自然成为商群 $G/H$ 的模。于是投影 $\varphi:G\to G/H$ 与包含 $\psi:A^H\to A$ 显然是相容同态。上同调上的相应诱导群同态称为\emph{膨胀同态}：
  \[
    \operatorname{Inf}:H^n(G/H,A^H)	o H^n(G,A),
    \qquad n\geq0.
  \]

  \item （余限制同态）假设 $H$ 是 $G$ 的指数为 $m$ 的子群，且 $A$ 是 $G$-模。令 $g_1,
  \ldots,g_m$ 是 $H$ 在 $G$ 中的左陪集代表。定义映射
  \[
    \psi:M_H^G(A)	o A,
    \qquad
    f\mapsto\sum_{i=1}^m g_i f(g_i^{-1}).
  \]
  注意，若把某个陪集代表 $g_i$ 改为 $g_ih$，则
  \[
    (g_ih)f((g_ih)^{-1})=g_ihf(h^{-1}g_i^{-1})=g_ihh^{-1}f(g_i^{-1})=g_if(g_i^{-1}),
  \]
  所以 $\psi$ 与陪集代表的选择无关。容易看出 $\psi$ 是 $G$-模同态（甚至是满射的），因此对所有 $n\geq0$ 得到从 $H^n(G,M_H^G(A))$ 到 $H^n(G,A)$ 的群同态。由于 $A$ 也是 $H$-模，由 Shapiro 引理有同构
  \[
    H^n(G,M_H^G(A))\cong H^n(H,A).
  \]
  这两个同态的复合称为\emph{余限制同态}：
  \[
    \operatorname{Cor}:H^n(H,A)	o H^n(G,A),
    \qquad n\geq0.
  \]
  这个同态可以如下显式计算：把 Shapiro 引理证明中的同构 $\Psi$ 与映射 $\psi$ 复合。若 $f\in\Hom_{
\ZZ H}(P_n,A)$ 是代表 $c\in H^n(H,A)$ 的上循环，则代表 $\operatorname{Cor}(c)\in H^n(G,A)$ 的上循环 $\operatorname{Cor}(f)\in\Hom_{\ZZ G}(P_n,A)$ 由
  \[
    \operatorname{Cor}(f)(p)=\sum_{i=1}^m g_i f(g_i^{-1}p)
  \]
  给出。$n=0$ 时尤其简单：可取 $P_0=\ZZ G$。此时 $f\in\Hom_{\ZZ H}(\ZZ G,A)=M_H^G(A)$ 是上循环，当且仅当 $f=a$ 对某个 $a\in A^H$ 为常值；于是 $\operatorname{Cor}(f)$ 是取值为 $\sum_{i=1}^m g_i\cdot a\in A^G$ 的常值函数：
  \[
    \operatorname{Cor}:H^0(H,A)=A^H\to A^G=H^0(G,A),
    \qquad a\mapsto\sum_{i=1}^m g_i\cdot a.
  \]
\end{enumerate}

下一个结果建立限制与余限制同态之间的基本关系。

\medskip
\noindent\textbf{命题 26.} 假设 $H$ 是 $G$ 的指数为 $m$ 的子群。则 $\operatorname{Cor}\circ\operatorname{Res}=m$，即若 $c$ 是某个 $G$-模 $A$ 的 $H^n(G,A)$ 中的上同调类，则
\[
  \operatorname{Cor}(\operatorname{Res}(c))=mc\in H^n(G,A)
\]
对所有 $n\geq0$ 成立。

\begin{proof}
这由上面例 4 中余限制的显式公式推出如下。若 $f\in\Hom_{\ZZ H}(P_n,A)$ 实际上属于 $\Hom_{
\ZZ G}(P_n,A)$，即 $f$ 也是 $G$-模同态，则对 $1\leq i\leq m$，有
\[
  g_if(g_i^{-1}p)=f(p).
\]
由于限制是在上同调上由自然包含 $\Hom_{\ZZ G}(P_n,A)\to\Hom_{\ZZ H}(P_n,A)$ 诱导的映射，对这样的 $f$，我们得到
\[
  \Hom_{\ZZ G}(P_n,A)\xrightarrow{\operatorname{Res}}\Hom_{\ZZ H}(P_n,A)\xrightarrow{\operatorname{Cor}}\Hom_{\ZZ G}(P_n,A),
  \qquad f\mapsto f\mapsto mf.
\]
因此在上同调群上 $\operatorname{Cor}\circ\operatorname{Res}$ 也是乘以 $m$。
\end{proof}

\medskip
\noindent\textbf{推论 27.} 假设有限群 $G$ 的阶为 $m$。则对所有 $n\geq1$ 和任意 $G$-模 $A$，有
\[
  mH^n(G,A)=0.
\]

\begin{proof}
在命题 26 中令 $H=1$，于是 $[G:H]=m$。对任意类 $c\in H^n(G,A)$，有 $mc=\operatorname{Cor}(\operatorname{Res}(c))$。由于 $\operatorname{Res}(c)\in H^n(H,A)=H^n(1,A)$，而由命题 20 前的例 2，当 $n\geq1$ 时 $\operatorname{Res}(c)=0$。故对所有 $n\geq1$，$mc=0$，得到推论。
\end{proof}

\medskip
\noindent\textbf{推论 28.} 若 $G$ 是有限群，则对所有 $n\geq1$ 与所有 $G$-模 $A$，$H^n(G,A)$ 是扭 Abel 群。

\begin{proof}
这由前一推论立即得到。
\end{proof}

\medskip
\noindent\textbf{推论 29.} 假设 $G$ 是有限群，并且 $|G|$ 与 $G$-模 $A$ 的指数互素。则对所有 $n\geq1$，
\[
  H^n(G,A)=0.
\]
特别地，若 $A$ 是有限 Abel 群且 $(|G|,|A|)=1$，则对所有 $n\geq1$，$H^n(G,A)=0$。

\begin{proof}
这因为 Abel 群 $H^n(G,A)$ 由前一推论被 $|G|$ 零化，又由命题 20 被 $A$ 的指数零化。
\end{proof}

注意，前面各推论的陈述对 $n=0$ 一般不成立，因为这时 $H^0(G,A)=A^G$，它甚至未必是扭群。

我们不加证明地提及下列结果。假设 $H$ 是 $G$ 的正规子群，$A$ 是 $G$-模。上同调群 $H^n(H,A)$ 可以赋予 $G/H$-模结构（参见练习 17）。可以证明存在正合列
\[
0\to H^1(G/H,A^H)\to H^1(G,A)\to H^1(H,A)^{G/H}
\xrightarrow{\operatorname{Tra}}H^2(G/H,A^H)\to H^2(G,A),
\]
其中 $H^1(H,A)^{G/H}$ 表示 $H^1(H,A)$ 在 $G/H$ 作用下的不动点，$\operatorname{Tra}$ 是所谓的传递同态。这个正合列把 $G$ 的上同调群与正规子群 $H$ 以及商群 $G/H$ 的上同调群联系起来。换言之，$G$ 的上同调与 $G$ 的滤过 $1\leq H\leq G$ 中因子的上同调有关。更一般地，可以试图把 $G$ 的上同调与 $G$ 的更长滤过中各因子的上同调联系起来。这就是谱序列理论，它是同调代数中的一个重要工具。

\subsection*{Galois 上同调与 profinite 群}

群上同调的一个重要应用出现在群 $G$ 是域扩张 $K/F$ 的 Galois 群时。在这种情形下，有许多感兴趣的群带有 $G$ 的作用，例如 $K$ 的加法群、乘法群 $K^\times$ 等。Galois 群 $G=\Gal(K/F)$ 是包含在 $K$ 中的 $F$ 的有限扩张 $L$ 的 Galois 群的逆极限 $\varprojlim \Gal(L/F)$，并且关于其 Krull 拓扑是紧拓扑群（即 $G$ 的群运算相对于由 $G$ 中有限指数子群 $\Gal(K/L)$ 定义的拓扑是连续的），参见第 14.9 节。在这种情形下，利用 $G$ 的附加拓扑结构是有用的（并且常常是必要的）。例如，$K$ 中含 $F$ 的子域与 $G=\Gal(K/F)$ 的闭子群一一对应；第 14.9 节中讨论的 $\QQ$ 的二次扩张复合的例子说明，一般而言 $G$ 中有许多子群不是闭的。幸运的是，在此背景下定义上同调群所需的修改相对很小，并适用于有限群的任意逆极限（profinite 群）。若 $G$ 是 profinite 群，则
\[
  G=\varprojlim G/N,
\]
其中逆极限取遍 $G$ 的开正规子群 $N$（参见练习 23）。

\noindent\textbf{定义.} 若 $G$ 是 profinite 群，则\emph{离散 $G$-模} $A$ 是带离散拓扑的 $G$-模 $A$，并且 $G$ 在 $A$ 上的作用连续；即映射 $G\times A\to A$，$(g,a)\mapsto g\cdot a$ 连续。

由于 $A$ 赋予离散拓扑，$A$ 的每个子集都是开的，特别地每个元素 $a\in A$ 是开点。$G$ 在 $A$ 上作用的连续性因此等价于：$a$ 在 $G$ 中的稳定子 $G_a$ 是 $G$ 的开子群；由于 $G$ 紧，它因而具有有限指数（参见练习 22）。这又等价于：
\[
  A=\bigcup_H A^H,
\]
其中并取遍 $G$ 的开子群 $H$。

在定义作用于离散 $G$-模 $A$ 的 profinite 群 $G$ 的上同调群 $H^n(G,A)$ 时必须稍加小心，因为这个范畴中没有足够的投射对象。例如，当 $G$ 无限时，自由 $G$-模 $\ZZ G$ 不是离散 $G$-模（$G$ 的作用不连续，参见练习 25）。不过，本节给出的 $H^n(G,A)$ 的显式描述（有时称为离散上同调群）很容易修改——只需要要求上链 $C^n(G,A)$ 是从 $G^n$ 到 $A$ 的连续映射。方程 (17.18) 中上边界映射 $d_n$ 的定义完全相同，上循环、上边界以及相应上同调群的定义也相同。习惯上不为这些上同调群引入单独记号，而是在术语中说明所指的是哪一种上同调。

\noindent\textbf{定义.} 若 $G$ 是 profinite 群，$A$ 是离散 $G$-模，则用连续上链计算得到的上同调群 $H^n(G,A)$ 称为\emph{profinite 上同调群}或\emph{连续上同调群}。当 $G=\Gal(K/F)$ 是域扩张 $K/F$ 的 Galois 群时，Galois 上同调群 $H^n(G,A)$ 总是指用连续上链计算的上同调群。

当 $G$ 是有限群时，每个 $G$-模都是离散 $G$-模，所以 $G$ 的离散上同调群与连续上同调群相同。当 $G$ 无限时，情况未必如此，正如前面提到的无限 profinite 群 $G$ 的自由 $G$-模 $\ZZ G$ 的例子所示。本节的所有主要结果在连续上同调群中仍然有效，只需把“$G$-模”替换为“离散 $G$-模”，把“子群”替换为“闭子群”。例如，群上同调长正合列仍按原样成立；限制同态要求 $G$ 的子群 $H$ 是闭子群（这样 $G^n$ 上的连续映射限制到 $H^n$ 仍连续）；命题 26 要求 $H$ 闭，等等。

可以写作
\[
  G=\varprojlim (G/N),
  \qquad
  A=\bigcup_N A^N,
\]
其中 $N$ 取遍 $G$ 的开正规子群（由于 $G$ 紧，这些子群必为有限指数）。于是 $A^N$ 是离散 $G/N$-模，并且不难证明
\[
  H^n(G,A)=\varinjlim_N H^n(G/N,A^N),
  \tag{17.19}
\]
其中上同调群为连续上同调，直极限取遍 $G$ 的所有开正规子群 $N$（参见练习 24）。由于 $G/N$ 是有限群，这个直极限中的连续上同调群 $H^n(G/N,A^N)$ 正是本节前面考虑的（离散）上同调群。因此，profinite 群 $G$ 的连续上同调的计算总是可以化为有限群上同调的考虑，在那里连续理论与离散理论没有区别。

\subsection*{练习}

\begin{enumerate}[label=\arabic*.]
  \item 令 $F_n=\ZZ G\otimes_{\ZZ}\ZZ G\otimes_{\ZZ}\cdots\otimes_{\ZZ}\ZZ G$（$n+1$ 个因子），$n\geq0$，其 $G$-作用在纯张量上定义为
  \[
    g\cdot(g_0\otimes g_1\otimes\cdots\otimes g_n)=(gg_0)\otimes g_1\otimes\cdots\otimes g_n.
  \]
  \begin{enumerate}[label=(\alph*)]
    \item 证明 $F_n$ 是秩为 $|G|^n$ 的自由 $\ZZ G$-模，其 $\ZZ G$-基为 $1\otimes g_1\otimes g_2\otimes\cdots\otimes g_n$，其中 $g_i\in G$。
  \end{enumerate}
  把 (a) 中的基元素 $1\otimes g_1\otimes g_2\otimes\cdots\otimes g_n$ 记为 $(g_1,g_2,
  \ldots,g_n)$，并对 $n\geq1$ 在这些基元素上定义 $G$-模同态 $d_n$：$d_1(g_1)=g_1-1$，而对 $n\geq2$，
  \[
  \begin{aligned}
  d_n(g_1,\ldots,g_n)&=g_1\cdot(g_2,\ldots,g_n)\\
  &\quad+\sum_{i=1}^{n-1}(-1)^i(g_1,\ldots,g_{i-1},g_ig_{i+1},g_{i+2},\ldots,g_n)\\
  &\quad+(-1)^n(g_1,\ldots,g_{n-1}).
  \end{aligned}
  \]
  定义 $\ZZ$-模收缩同态
  \[
    \ZZ\xrightarrow{s_{-1}}F_0\xrightarrow{s_0}F_1\xrightarrow{s_1}F_2\xrightarrow{s_2}\cdots
  \]
  在一个 $\ZZ$-基上由 $s_{-1}(1)=1$，以及
  \[
    s_n(g_0\otimes\cdots\otimes g_n)=1\otimes g_0\otimes\cdots\otimes g_n
  \]
  给出。
  \begin{enumerate}[label=(\alph*),resume]
    \item 证明
    \[
      \varepsilon s_{-1}=1,
      \qquad d_1s_0+s_{-1}\varepsilon=1,
      \qquad d_{n+1}s_n+s_{n-1}d_n=1\quad(n\geq1),
    \]
    其中 $\varepsilon:F_0\to\ZZ$ 是增广映射 $\varepsilon(\sum_{g\in G}\alpha_gg)=\sum_{g\in G}\alpha_g$。
    \item 证明映射 $s_n$ 是链
    \[
      \cdots\to F_n\xrightarrow{d_n}F_{n-1}\to\cdots\to F_0\xrightarrow{\varepsilon}\ZZ\to0
    \]
    从其自身到其自身的恒等链映射与零链映射之间的链同伦（参见第 1 节练习 4）。
    \item 从 (c) 推出 $(*)$ 的所有 $\ZZ$-模同调群都为零，即 $(*)$ 是 $\ZZ$-模正合列。推出 $(*)$ 是 $\ZZ$ 的一个投射 $G$-模分解。
  \end{enumerate}

  \item 令 $P_n$ 为自由 $\ZZ$-模，其基为 $(g_0,g_1,g_2,
  \ldots,g_n)$，其中 $g_i\in G$；定义 $G$ 在 $P_n$ 上的作用为
  \[
    g\cdot(g_0,g_1,
    \ldots,g_n)=(gg_0,gg_1,
    \ldots,gg_n).
  \]
  对 $n\geq1$ 定义
  \[
    d_n(g_0,g_1,g_2,
    \ldots,g_n)=\sum_{i=0}^n(-1)^i(g_0,
    \ldots,\widehat{g_i},
    \ldots,g_n),
  \]
  其中带帽表示删去该项。
  \begin{enumerate}[label=(\alph*)]
    \item 证明 $P_n$ 是自由 $\ZZ G$-模，其基为 $(1,g_1,g_2,
    \ldots,g_n)$，其中 $g_i\in G$。
    \item 证明 $d_{n-1}\circ d_n=0$，$n\geq1$。［证明缺少 $g_j$ 与 $g_k$ 的项 $(g_0,
    \ldots,
    \widehat g_j,
    \ldots,
    \widehat g_k,
    \ldots,g_n)$ 在 $d_{n-1}d_n(g_0,g_1,
    \ldots,g_n)$ 中出现两次，符号相反。］
    \item 证明 $\varphi:P_n\to F_n$ 定义为
    \[
      \varphi((g_0,g_1,g_2,
      \ldots,g_n))=g_0\otimes(g_0^{-1}g_1)\otimes(g_1^{-1}g_2)\otimes\cdots\otimes(g_{n-1}^{-1}g_n)
    \]
    是 $G$-模同构，其逆 $\psi:F_n\to P_n$ 由
    \[
      \psi(g_0\otimes g_1\otimes\cdots\otimes g_n)=(g_0,g_0g_1,g_0g_1g_2,
      \ldots,g_0g_1\cdots g_n)
    \]
    给出。
    \item 证明若 $\varepsilon(g_0)=1$ 对所有 $g_0\in G$ 成立，则
    \[
      \cdots\to P_n\to\cdots\to P_0\to\ZZ\to0
    \]
    是 $\ZZ$ 的一个自由 $G$-模分解。［证明 (c) 中的同构把本练习中的 $G$-模分解与前一练习中的分解相互对应。］
  \end{enumerate}

  \item 令 $F_n$ 与 $P_n$ 如前两个练习，并令 $A$ 是 $G$-模。
  \begin{enumerate}[label=(\alph*)]
    \item 证明 $\Hom_{\ZZ G}(F_n,A)$ 可与从 $G^n$ 到 $A$ 的映射集合 $C^n(G,A)$ 识别，并且在此识别下，相关上边界映射 $C^n(G,A)\to C^{n+1}(G,A)$ 由方程 (17.18) 给出。
    \item 证明 $\Hom_{\ZZ G}(P_n,A)$ 可与从 $G^{n+1}$ 到 $A$ 的所有映射 $f$ 的集合识别，这些映射满足
    \[
      f(gg_0,gg_1,
      \ldots,gg_n)=gf(g_0,g_1,
      \ldots,g_n).
    \]
  \end{enumerate}
  群 $C^n(G,A)$ 有时称为 $G$ 在 $A$ 中的\emph{非齐次 $n$-上链}群，而前一练习 (b) 中的群称为 $G$ 在 $A$ 中的\emph{齐次 $n$-上链}群。非齐次上链较易描述，因为从 $G^n$ 到 $A$ 的映射没有限制；但齐次上链上的上边界映射 $d_n$ 比非齐次上链上的上边界映射简单（并且在拓扑背景中更自然）。前面练习的结果说明，使用齐次或非齐次上链定义的上同调群 $H^n(G,A)$ 是相同的，并指出了正文中所用上边界映射 $d_n$ 的来源。历史上，$H^n(G,A)$ 最初是用齐次上链定义的。

  \item 假设 $H$ 是群 $G$ 的正规子群，$A$ 是 $G$-模。对每个 $g\in G$，证明映射 $f(a)=ga$（$a\in A^H$）定义了子群 $A^H$ 的一个自同构。

  \item 假设 $G$-模 $A$ 分解为 $G$-子模的直和 $A=A_1\oplus A_2$。证明对所有 $n\geq0$，
  \[
    H^n(G,A)\cong H^n(G,A_1)\oplus H^n(G,A_2).
  \]

  \item 假设
  \[
    0\to A\to M_1\to M_2\to\cdots\to M_k\to C\to0
  \]
  是 $G$-模正合列，其中 $M_1,M_2,
  \ldots,M_k$ 都是上同调平凡的。证明对所有 $n\geq1$，
  \[
    H^{n+k}(G,A)\cong H^n(G,C).
  \]
  ［把正合列分解为一串短正合列并用推论 22。例如，若 $0\to A\xrightarrow{\alpha}M_1\xrightarrow{\beta}M_2\xrightarrow{\gamma}C\to0$ 正合，证明 $0\to A\to M_1\to B\to0$ 与 $0\to B\to M_2\to C\to0$ 都正合，其中 $B=M_1/\operatorname{image}\alpha=M_1/\ker\beta\cong\operatorname{image}\beta=\ker\gamma$。］

  \item （伴随结合性）设 $R,S,T$ 是带 $1$ 的环，$P$ 是左 $S$-模，$N$ 是 $(T,S)$-双模，$A$ 是左 $T$-模。证明
  \[
    \Phi:\Hom_S(P,\Hom_T(N,A))\cong\Hom_T(N\otimes_SP,A)
  \]
  由 $\Phi(f)(n\otimes p)=f(p)(n)$ 定义，是 Abel 群同构。（也参见第 10.5 节定理 43。）

  \item 设 $G$ 是由 $\sigma$ 生成的 $m$ 阶循环群，令 $N=1+\sigma+\sigma^2+\cdots+\sigma^{m-1}\in\ZZ G$。
  \begin{enumerate}[label=(\alph*)]
    \item 证明增广映射 $\operatorname{aug}(\sum_{i=0}^{m-1}a_i\sigma^i)=\sum_{i=0}^{m-1}a_i$ 是从 $\ZZ G$ 到 $\ZZ$ 的 $G$-模同态。
    \item 证明 $\ZZ G$ 中乘以 $N$ 与乘以 $\sigma-1$ 定义了 $\ZZ$ 的一个自由 $G$-模分解
    \[
      \cdots\to\ZZ G\xrightarrow{N}\ZZ G\xrightarrow{\sigma-1}\cdots\to\ZZ G\xrightarrow{N}\ZZ G\xrightarrow{\operatorname{aug}}\ZZ\to0.
    \]
  \end{enumerate}

  \item 假设 $G$ 是由 $\sigma$ 生成的无限循环群。
  \begin{enumerate}[label=(\alph*)]
    \item 证明 $\ZZ G$ 中乘以 $\sigma-1$ 定义了 $\ZZ$ 的一个自由 $G$-模分解
    \[
      0\to\ZZ G\xrightarrow{\sigma-1}\ZZ G\to\ZZ\to0.
    \]
    \item 证明 $H^0(G,A)\cong A^G$，$H^1(G,A)\cong A/(\sigma-1)A$，并且对所有 $n\geq2$，$H^n(G,A)=0$。推出 $H^1(G,\ZZ G)\cong\ZZ$（所以自由模未必上同调平凡）。
  \end{enumerate}

  \item 假设 $H$ 是群 $G$ 的有限指数 $m$ 的子群，$A$ 是 $H$-模。令 $x_1,
  \ldots,x_m$ 为 $H$ 在 $G$ 中左陪集代表：$G=x_1H\cup\cdots\cup x_mH$。
  \begin{enumerate}[label=(\alph*)]
    \item 证明 $\ZZ G=\bigoplus_{i=1}^m x_i\ZZ H=\bigoplus_{i=1}^m\ZZ Hx_i^{-1}$，并且 $\ZZ G\otimes_{\ZZ H}A=\bigoplus_{i=1}^m(x_i\otimes A)$ 作为 Abel 群。
    \item 令 $f_{i,a}$ 是从 $\ZZ G$ 到 $A$ 的函数，定义为
    \[
      f_{i,a}(x)=
      \begin{cases}
      ha,& x=hx_i^{-1},\ h\in H,\\
      0,& \text{否则}.
      \end{cases}
    \]
    证明 $f_{i,a}\in M_H^G(A)=\Hom_{\ZZ H}(\ZZ G,A)$，即 $f_{i,a}(h'x)=h'f_{i,a}(x)$。
    \item 证明映射
    \[
      \phi(f)=\sum_{i=1}^m x_i\otimes f(x_i^{-1})
    \]
    从 $M_H^G(A)$ 到 $\ZZ G\otimes_{\ZZ H}A$ 是 $G$-模同态。［写 $x_i^{-1}g=h_ix_j^{-1}$，并观察
    \[
      x_i\otimes f(x_i^{-1}g)=x_i\otimes h_if(x_j^{-1})=x_ih_i\otimes f(x_j^{-1})=gx_j\otimes f(x_j^{-1}).
    \]
    ］
    \item 证明 $\phi$ 给出 $G$-模同构 $M_H^G(A)\cong\ZZ G\otimes_{\ZZ H}A$。［单射性：$H$-模同态为 $0$ 当且仅当 $f(x_i^{-1})=0$，$i=1,
    \ldots,m$。满射性：证明 $\phi(f_{i,a})=x_i\otimes a$。］
  \end{enumerate}

  \item 证明若 $H$ 在 $G$ 中不是有限指数，则前一练习 (d) 中的同构 $M_H^G(A)\cong\ZZ G\otimes_{\ZZ H}A$ 未必成立。［若 $G$ 是无限循环群，说明 Shapiro 引理蕴含 $H^1(G,M_1^G(\ZZ))=0$，而由练习 9，$H^1(G,\ZZ G)\cong\ZZ$。］

  \item 若 $H$ 是 $G$ 的子群，$A$ 是 Abel 群，令 $M_{G/H}(A)$ 表示从 $H$ 在 $G$ 中左陪集 $gH$ 到 $A$ 的所有映射组成的 Abel 群。
  \begin{enumerate}[label=(\alph*)]
    \item 证明 $M_1^G(A)\cong M_1^H(M_{G/H}(A))$ 作为 $H$-模。［若 $\{g_i\}_{i\in I}$ 是 $H$ 在 $G$ 中左陪集代表的选择，对 $f\in M_1^G(A)$ 与 $F:H\to M_{G/H}(A)$ 定义对应 $F(h)(g_iH)=f(g_ih)$，并检验这是 $H$-模同构。］
    \item 若 $G$-模 $A$ 满足对 $G$ 的所有子群 $H$ 与所有 $n\geq1$ 都有 $H^n(H,A)=0$，称 $A$ 为上同调平凡的。证明对任意 Abel 群 $A$，$M_1^G(A)$ 是上同调平凡的。
    \item 若 $G$ 有限，证明对所有 Abel 群 $A$，$\ZZ G\otimes_{\ZZ}A$ 是上同调平凡的。
  \end{enumerate}

  \item 假设 $A$ 是 $G$-模，$H$ 是 $G$ 的子群。证明对所有 $n\geq0$，由例 3（推论 22 后）中的 $G$-模同态 $A\to M_H^G(A)$ 诱导的从 $H^n(G,A)$ 到 $H^n(G,M_H^G(A))$ 的群同态，再与 Shapiro 引理的同构 $H^n(G,M_H^G(A))\cong H^n(H,A)$ 复合，所得就是限制同态 $H^n(G,A)\to H^n(H,A)$。

  \item 假设 $\varphi:H\to G$ 是子群 $H$ 到 $G$ 的包含映射。若 $A$ 是 $H$-模，$M_H^G(A)$ 是相应诱导 $G$-模，定义群同态 $\psi:M_H^G(A)\to A$ 为取 $f$ 在 $1$ 处的值：$\psi(f)=f(1)$。
  \begin{enumerate}[label=(\alph*)]
    \item 证明 $\varphi$ 与 $\psi$ 是相容同态。
    \item 证明对 $n\geq0$，从 $H^n(G,M_H^G(A))$ 到 $H^n(H,A)$ 的诱导群同态就是 Shapiro 引理中的同构。
  \end{enumerate}

  \item 假设 $H$ 是 $G$ 的正规子群，$A$ 是 $G$-模。对固定 $g\in G$，令 $\psi(a)=ga$，$\varphi(h)=g^{-1}hg$（$h\in H$）。
  \begin{enumerate}[label=(\alph*)]
    \item 证明 $\varphi$ 与 $\psi$ 是相容同态。
    \item 对每个 $n\geq0$，证明由相容同态 $\varphi,\psi$ 诱导的从 $H^n(H,A)$ 到 $H^n(H,A)$ 的同态 $\theta_g$ 是 $H^n(H,A)$ 的自同构。［观察 $\varphi$ 与 $\psi$ 都有逆。］
    \item 说明 $\theta_g$ 作用在 $H^0(H,A)$ 上就是练习 4 中的自同构。
  \end{enumerate}

  \item 设 $A$ 是 $G$-模，对 $g\in G$ 令 $\theta_g$ 表示前一练习中定义的 $H^n(G,A)$ 的自同构。
  \begin{enumerate}[label=(\alph*)]
    \item 证明 $\theta_g$ 作用在 $H^0(G,A)=A^G$ 上是恒等映射。
    \item 证明对 $n\geq1$，$\theta_g$ 作用在 $H^n(G,A)$ 上是恒等映射。［对 $n$ 归纳并用维数平移。对 $n=1$，用推论 22 中的正合列，并结合 (a) 应用于 $C^G$ 上的 $\theta_g$。对 $n\geq2$，用推论 22 中的同构 $H^{n+1}(G,A)\cong H^n(G,C)$。］
  \end{enumerate}

  \item 假设 $H$ 是 $G$ 的正规子群，$A$ 是 $G$-模。对 $n\geq0$，证明 $H^n(H,A)$ 是 $G/H$-模，其中 $gH$ 通过练习 15 中由 $g$ 对 $H$ 的共轭与 $g$ 在 $A$ 上的自然作用诱导的自同构 $\theta_g$ 作用。［用前一练习说明这个陪集作用是良定义的。］

  \item 假设 $G$ 是 $m$ 阶循环群，$H$ 是 $G$ 的指数为 $d$ 的子群，$\ZZ$ 是平凡 $G$-模。用练习 8 中的投射 $G$-模分解证明：
  \begin{enumerate}[label=(\alph*)]
    \item $\operatorname{Cor}:H^n(H,\ZZ)\to H^n(G,\ZZ)$ 在 $n=0$ 时是从 $\ZZ$ 到 $\ZZ$ 的乘以 $d$，在奇数 $n\geq1$ 时是从 $\ZZ/(m/d)\ZZ$ 到 $\ZZ/m\ZZ$ 的乘以 $d$，在偶数 $n\geq2$ 时是从 $0$ 到 $0$；
    \item $\operatorname{Res}:H^n(G,\ZZ)\to H^n(H,\ZZ)$ 在 $n=0$ 时是从 $\ZZ$ 到 $\ZZ$ 的恒等映射，而当 $n\geq1$ 时，视奇偶性而定，是从 $\ZZ/m\ZZ$ 到 $\ZZ/(m/d)\ZZ$ 的自然投影映射，或从 $0$ 到 $0$ 的映射。
  \end{enumerate}

  \item 令 $p$ 是素数，$P$ 是有限群 $G$ 的 Sylow $p$-子群。证明对任意 $G$-模 $A$ 与所有 $n\geq0$，限制映射 $\operatorname{Res}:H^n(G,A)\to H^n(P,A)$ 在 $H^n(G,A)$ 的 $p$-主分量上是单射。推出若 $|A|=p^a$，则限制映射在 $H^n(G,A)$ 上是单射。［用命题 26。］

  \item 令 $p$ 是素数，$G=(\sigma)$ 是 $p^m$ 阶循环群，$W$ 是 $\mathbb F_p$ 上维数 $d>0$ 的向量空间，$\sigma$ 作为线性变换作用在 $W$ 上。假设 $W$ 有一组基，使得 $\sigma$ 的矩阵是特征值为 $1$ 的 $d\times d$ 初等 Jordan 块。
  \begin{enumerate}[label=(\alph*)]
    \item 证明 $d\leq p^m$。［用初等 Jordan 块的最小多项式的事实。］
    \item 证明 $\dim_{\mathbb F_p}W^G=1$。
    \item 证明 $\dim_{\mathbb F_p}(\sigma-1)W=d-1$。
    \item 若 $N=1+\sigma+\cdots+\sigma^{p^m-1}$ 是通常的范元，证明当 $d=p^m$ 时 $NW$ 的维数为 $1$（当 $d<p^m$ 时维数为 $0$），并且 ${}_NW$ 的维数分别为 $d-1$（当 $d=p^m$）与 $d$（当 $d<p^m$）。［令 $R$ 为群环 $\mathbb F_pG$，并证明 $R$ 的每个非零 $R$-子模都包含 $N$。注意 $W$ 是循环 $R$-模，并令 $\phi:R\to W$ 为满同态。推出若 $\phi$ 不是同构，则 $N\in\ker\phi$。］
    \item 推出若 $d=p^m$，则 $H^n(G,W)=0$；若 $d<p^m$，则对所有 $n\geq1$，$H^n(G,W)$ 的阶为 $p$（即这些上同调群为零当且仅当 $W$ 是自由 $\mathbb F_pG$-模）。
  \end{enumerate}

  \item 令 $p$ 是素数，$G=(\sigma)$ 是 $p^m$ 阶循环群，$V$ 是指数为 $p$ 的 $G$-模。令 $V=V_1\oplus V_2\oplus\cdots\oplus V_k$ 是给出 $\sigma$ 的 Jordan 标准形的分解，其中每个 $V_i$ 都是 $\sigma$-不变的，并且 $\sigma$ 在 $V_i$ 上的矩阵是特征值为 $1$ 的 $d_i\times d_i$ 初等 Jordan 块，$d_i\geq1$（参见第 12.3 节）。证明 $|V^G|=p^k$，并且对所有 $n\geq1$，$|H^n(G,V)|=p^s$，其中 $s$ 是维数小于 $p^m$ 的那些 $V_i$ 的个数。［用前一练习和练习 5。］

  \item 假设 $G$ 是拓扑群，即 $G$ 上有一个拓扑，使映射 $G\times G\to G$，$(g_1,g_2)\mapsto g_1g_2$ 以及 $G\to G$，$g\mapsto g^{-1}$ 连续。
  \begin{enumerate}[label=(\alph*)]
    \item 若 $H$ 是 $G$ 的开子群且 $g\in G$，证明陪集 $gH$、$Hg$ 以及子群 $g^{-1}Hg$ 也都是开的。
    \item 证明任意开子群也是闭的。［补集是 (a) 中陪集的并。］
    \item 证明有限指数的闭子群是开的。
    \item 若 $G$ 紧，证明每个开子群 $H$ 都具有有限指数。
  \end{enumerate}

  \item 假设 $G$ 是紧拓扑群。证明下列条件等价：
  \begin{enumerate}[label=(\roman*)]
    \item $G$ 是 profinite 群，即 $G=\varprojlim G_i$ 是有限群 $G_i$ 的逆极限；
    \item 存在 $G$ 中开正规子群族 $\{N_i\}$（$i\in I$），使得 $\bigcap_iN_i=1$，并且在这种情形下 $G\cong\varprojlim(G/N_i)$；
    \item 存在 $G$ 中开子群族 $\{H_j\}$（$j\in\mathcal J$），使得 $\bigcap_jH_j=1$。
  \end{enumerate}
  ［为证明 (iii) 蕴含 (ii)，令 $H$ 是 $G$ 的开子群，并用前一练习 (d) 说明 $N=\bigcap_{g\in G}g^{-1}Hg$ 实际上是有限交，从而 $N\subseteq H\subseteq G$ 且 $N$ 是 $G$ 中开正规子群。］

  \item 假设 $N$ 与 $N'$ 是 profinite 群 $G$ 的开正规子群，且 $N'\subseteq N$。证明投影同态 $\varphi:G/N'\to G/N$ 与包含 $\psi:A^N\to A^{N'}$ 是相容同态，并推出存在从 $H^n(G/N,A^N)$ 到 $H^n(G/N',A^{N'})$ 的诱导同态。

  \item 若 $G$ 是无限 profinite 群，说明 $G$ 在 $A=\ZZ G$ 上的作用不是连续的。［说明 $a\in A$ 的稳定子并不总在 $G$ 中具有有限指数。］
\end{enumerate}

\section[交叉同态与 H1(G,A)]{交叉同态与 $H^1(G,A)$}

本节中，我们更详细地考察上同调群 $H^1(G,A)$，其中 $G$ 是群，$A$ 是 $G$-模。由方程 (17.18) 中上边界映射 $d_1$ 的定义，若 $f\in C^1(G,A)$，则
\[
  (d_1f)(g_1,g_2)=g_1\cdot f(g_2)-f(g_1g_2)+f(g_1).
\]
因此函数 $f:G\to A$ 是 $1$-上循环，当且仅当它满足恒等式
\[
  f(gh)=f(g)+g f(h),
  \qquad g,h\in G.
  \tag{17.20}
\]
等价地，一个 $1$-上循环由一族元素 $\{a_g\}_{g\in G}$ 决定，满足 $a_{gh}=a_g+ga_h$，其中 $g,h\in G$（这时 $1$-上循环 $f$ 是把 $g$ 送到 $a_g$ 的函数）。注意若 $1$ 表示 $G$ 的恒等元，则
\[
  f(1)=f(1^2)=f(1)+1\cdot f(1)=2f(1),
\]
故 $f(1)=0$ 是 $A$ 中的恒等元。因此 $1$-上循环在恒等元处必定是“规范化的”。进而由上循环条件得出，对所有 $g\in G$，
\[
  f(g^{-1})=-g^{-1}f(g).
\]

若 $G$ 在 $A$ 上平凡作用，则上循环条件 (17.20) 只是 $f(gh)=f(g)+f(h)$，即 $f$ 只是从乘法群 $G$ 到加法群 $A$ 的同态。由于这个原因，满足 (17.20) 的从 $G$ 到 $A$ 的函数称为\emph{交叉同态}。

一个 $1$-上链 $f$ 是 $1$-上边界，当且仅当存在某个 $a\in A$ 使得
\[
  f(g)=g\cdot a-a,
  \qquad g\in G.
  \tag{17.21}
\]
（等价地，在上述记号中 $a_g=ga-a$。）注意，由于 $-a\in A$，上边界条件 (17.21) 也可表述为：对某个固定的 $a\in A$ 和所有 $g\in G$，$f(g)=a-g\cdot a$。$1$-上边界称为\emph{主交叉同态}。用这个术语，$H^1(G,A)$ 是交叉同态群模去主交叉同态子群所得的群。

\paragraph{例：Hilbert 第 90 定理}
设 $G=\Gal(K/F)$ 是有限 Galois 域扩张 $K/F$ 的 Galois 群。则乘法群 $K^\times$ 是 $G$-模，并且 $H^1(G,K^\times)=0$。为看出这一点，令 $\{a_\sigma\}$ 是一个 $1$-上循环 $f$ 的取值 $f(\sigma)$，于是 $a_\sigma\in K^\times$ 且
\[
  a_{\sigma\tau}=a_\sigma\sigma(a_\tau)
\]
（这是对群 $K^\times$ 用乘法记号写出的上循环条件）。由自同构的线性无关性（第 14.2 节推论 8），存在 $y\in K$ 使得
\[
  \beta=\sum_{\tau\in G}a_\tau\tau(y)
\]
非零，即 $\beta\in K^\times$。则对任意 $\sigma\in G$，
\[
  \sigma(\beta)=\sum_{\tau\in G}\sigma(a_\tau)\sigma\tau(y)=a_\sigma^{-1}\sum_{\tau\in G}a_{\sigma\tau}\sigma\tau(y)=a_\sigma^{-1}\beta,
\]
其中第二个等号来自上循环条件。因此 $a_\sigma=\beta/\sigma(\beta)$，这是上边界条件 (17.21) 的乘法形式（对应于元素 $a=\beta^{-1}$）。既然每个 $1$-上循环都是 $1$-上边界，便有 $H^1(G,K^\times)=0$。对无限 Galois 扩张，同样结果由上一节方程 (17.19) 成立，因为 $H^1(G,K^\times)$ 是平凡群的直极限。

作为特例，设 $K/F$ 是 Galois 扩张，其 Galois 群 $G$ 是由 $\sigma$ 生成的循环群。上一节显式计算了 $G$ 的上同调群；特别地，对任意 $G$-模 $A$（加法记号），
\[
  H^1(G,A)={}_NA/(\sigma-1)A.
\]
由于在当前情形中该群平凡，可知 $K$ 中元素 $a$ 属于范映射的核，即 $N_{K/F}(a)=1$，当且仅当 $a=\sigma(\beta)/\beta$ 对某个 $\beta\in K$ 成立。（循环情形的直接证明参见第 14.2 节练习 23。）

这个关于循环扩张的著名结果最初由 Hilbert 证明，并作为“第 90 定理”出现在他 1897 年关于数论的著作（称为 Zahlbericht）中。因此，更一般的结果 $H^1(G,K^\times)=0$ 在文献中称为“Hilbert 第 90 定理”。一般而言，高维上同调群 $H^n(G,K^\times)$ 对 $n\geq2$ 可以非平凡（参见练习 13）。

\paragraph{例}
设 $G=\Gal(K/F)$ 是有限 Galois 域扩张 $K/F$ 的 Galois 群，如前一例。则加法群 $K$ 也是 $G$-模，并且对所有 $n\geq1$，$H^n(G,K)=0$。一般证明使用正规基存在性：存在 $a\in K$，使其 Galois 共轭给出 $K$ 作为 $F$ 上向量空间的一组基；等价地，$K\cong\ZZ G\otimes_{\ZZ}F$ 作为 $G$-模。后一同构说明 $K$ 作为 $G$-模是诱导出来的，因此 $H^n(G,K)=0$ 由第 2 节推论 24 得到。循环情形的直接证明参见第 14.2 节练习 26。

若 $G$ 在 $A$ 上平凡作用，则 $g\cdot a-a=0$，所以 $0$ 是唯一的主交叉同态，即 $B^1(G,A)=0$。这证明了下面的结果。

\medskip
\noindent\textbf{命题 30.} 若 $A$ 是 $G$-模且 $G$ 在 $A$ 上平凡作用，则
\[
  H^1(G,A)\cong\Hom(G,A),
\]
即从 $G$ 到 $A$ 的所有群同态构成的群。

若 $G$ 是 profinite 群，则对连续上同调群 $H^1(G,A)$ 同样成立，但需要取从 $G$ 到 $A$ 的连续同态群。

\paragraph{例}
\begin{enumerate}[label=(\arabic*)]
  \item 若 $G$ 在 $A$ 上平凡作用，则 $H^1(G,A)=H^1(G/[G,G],A)$，因为从 $G$ 到 Abel 群 $A$ 的任意群同态都通过换位子群 $[G,G]$ 分解（参见第 5.4 节命题 7(5)）。因此，计算平凡 $G$-作用下的 $H^1$ 可化为计算某个 Abel 群的 $H^1$。
  \item 若 $G$ 是有限群且在 $\ZZ$ 上平凡作用，则 $H^1(G,\ZZ)=0$，因为 $\ZZ$ 没有非零有限阶元素，所以不存在从 $G$ 到 $\ZZ$ 的非零群同态。
  \item 若 $A$ 是素数阶 $p$ 的循环群，且 $G$ 是 $p$-群，则 $G$ 必定在 $A$ 上平凡作用（因为 $A$ 的自同构群阶为 $p-1$），所以在此情形下总有 $H^1(G,A)=\Hom(G,A)$。
  \item 若 $G$ 是有限群且在 $\QQ/\ZZ$ 上平凡作用，则
  \[
    H^1(G,\QQ/\ZZ)=\Hom(G,\QQ/\ZZ)=\widehat G
  \]
  是 $G$ 的对偶群（参见第 5.2 节练习 14）。由于 $\QQ/\ZZ$ 是 Abel 群，$G$ 到 $\QQ/\ZZ$ 的任意同态都通过 $G$ 的换位子商 $G^{\mathrm{ab}}=G/[G,G]$ 分解，所以
  \[
    \Hom(G,\QQ/\ZZ)=\Hom(G^{\mathrm{ab}},\QQ/\ZZ).
  \]
  因此 $\Hom(G,\QQ/\ZZ)\cong\widehat{G^{\mathrm{ab}}}$（再次由练习 14，它非典范地同构于 $G^{\mathrm{ab}}$）。
\end{enumerate}

若 $0\to A\to B\to C\to0$ 是 $G$-模短正合列，则上一节定理 21 中的群上同调长正合列开始为
\[
  0\to A^G\to B^G\to C^G\xrightarrow{\delta_0}H^1(G,A)\to H^1(G,B)\to H^1(G,C)\to\cdots .
\]
连接同态 $\delta_0$ 可如下显式给出：若 $c\in C^G$，则存在 $b\in B$ 映到 $c$；于是 $\delta_0(c)$ 是 $H^1(G,A)$ 中由 $1$-上循环
\[
  \delta_0(c):G\to A,
  \qquad g\mapsto g\cdot b-b
\]
给出的类。注意，由于 $c\in C^G$，对所有 $g\in G$，元素 $g\cdot b-b$（作为 $B$ 中的元素）确实来自 $A$。直接验证 $f=\delta_0(c)$ 满足上循环条件 (17.20)：
\[
  f(gh)=gh\cdot b-b=(g\cdot b-b)+g\cdot(h\cdot b-b)=f(g)+gf(h).
\]
由显式表达式 $f=g\cdot b-b$ 也清楚看出，$\delta_0(c)\in H^1(G,A)$ 映到长正合列下一项 $H^1(G,B)$ 中为 $0$，因为 $f$ 是元素 $b\in B$ 的上边界。

\paragraph{例：Kummer 理论}
设 $F$ 是特征为 $0$ 的域，并且包含所有 $n$ 次单位根组成的群 $\mu_n$，其中 $n\geq1$。令 $K$ 为 $F$ 的一个代数闭包，并令 $G=\Gal(K/F)$。由假设 $\mu_n\subseteq F$，故 $G$ 在 $\mu_n$ 上平凡作用，即 $\mu_n\cong\ZZ/n\ZZ$ 作为 $G$-模。因此 Galois 上同调群 $H^1(G,\mu_n)$ 是从 $G$ 到 $\ZZ/n\ZZ$ 的连续同态群 $\Hom_c(G,\ZZ/n\ZZ)$。若 $\chi$ 是这样的连续同态，则 $\ker\chi$ 是 $G$ 的闭正规子群，因而由 Galois 理论对应于一个 Galois 扩张 $L_\chi/F$。于是 $\Gal(L_\chi/F)\cong\operatorname{image}\chi$，所以 $L_\chi$ 是 $F$ 的一个次数整除 $n$ 的循环扩张。反过来，$F$ 的每个这样的循环扩张都定义 $\Hom_c(G,\ZZ/n\ZZ)$ 中的一个元素，因此 Galois 上同调群 $H^1(G,\mu_n)$ 的元素与 $F$ 的次数整除 $n$ 的循环扩张之间存在双射。

$n$ 次幂同态在 $K^\times$ 上是满射（因为在 $K$ 中总能开 $n$ 次方），其核为 $\mu_n$。因此序列
\[
  1\to\mu_n\to K^\times\xrightarrow{n}K^\times\to1
\]
是离散 $G$-模正合列。相关的 Galois 上同调长正合列给出
\[
  1\to\mu_n^G\to(K^\times)^G\xrightarrow{n}(K^\times)^G\to H^1(G,\mu_n)\to H^1(G,K^\times)\to\cdots .
\]
由 Galois 理论，$\mu_n^G=\mu_n$ 且 $(K^\times)^G=F^\times$，由 Hilbert 第 90 定理，$H^1(G,K^\times)=0$，所以这个正合列成为
\[
  1\to\mu_n\to F^\times\xrightarrow{n}F^\times\to H^1(G,\mu_n)\to0,
\]
从而等价于同构
\[
  H^1(G,\mu_n)\cong F^\times/(F^\times)^n,
\]
其中 $(F^\times)^n$ 表示 $F^\times$ 中 $n$ 次幂元素组成的群。利用上面给出的连接同态的显式形式可以把这个同构写得很明确：对每个 $a\in F^\times$ 与 $\sigma\in G$，元素 $\sqrt[n]{a}\in K^\times$ 在正合列中映到 $a$，并且
\[
  \chi(\sigma)=\frac{\sigma(\sqrt[n]{a})}{\sqrt[n]{a}}
\]
定义了 $H^1(G,\mu_n)$ 中的一个元素（参见练习 11）。这个同态 $\chi$ 的核对应于域 $F(\sqrt[n]{a})$。由前一段的结果，当 $F$ 包含 $n$ 次单位根时，扩张 $L/F$ 是 Galois 的且 Galois 群为阶整除 $n$ 的循环群，当且仅当 $L=F(\sqrt[n]{a})$ 对某个 $a\in F^\times$ 成立。此外，$a$ 在 $F^\times/(F^\times)^n$ 中的类是唯一的，即 $a$ 在乘上 $F$ 中元素的 $n$ 次幂的意义下唯一。这样的扩张称为 Kummer 扩张，参见第 14.7 节和练习 12。

若 $F$ 的特征为素数 $p$，当 $n$ 不被 $p$ 整除时，同样论证仍成立，只需把 $F$ 的代数闭包替换为 $F$ 的可分闭包（$F$ 的最大可分代数扩张）。

\paragraph{例：转移同态}
设 $G$ 是有限群，$H$ 是子群。余限制给出同态 $H^1(H,\QQ/\ZZ)\to H^1(G,\QQ/\ZZ)$，由上面例 4 它给出从 $\widehat{H^{\mathrm{ab}}}$ 到 $\widehat{G^{\mathrm{ab}}}$ 的同态。这给出一个同态
\[
  \operatorname{Ver}:G^{\mathrm{ab}}	o H^{\mathrm{ab}},
\]
称为\emph{转移}（或 Verlagerung）同态（参见练习 14）。为了显式描述此同态，考虑正合列
\[
  0\to\QQ/\ZZ\to M_1^G(\QQ/\ZZ)\to C\to0,
  \tag{17.22}
\]
其中同态把 $a\in\QQ/\ZZ$ 映到第 2 节命题 23 前例 4 中的 $f_a\in M_1^G(\QQ/\ZZ)$（即 $f_a(g)=g\cdot a$）。这是 $G$-模短正合列，因而也是 $H$-模短正合列。对 $H$ 与 $G$ 取上同调所得长正合列的最初部分给出一个交换图；由于 $H^1(H,M_1^G(\QQ/\ZZ))=H^1(G,M_1^G(\QQ/\ZZ))=0$（参见第 2 节练习 12），连接同态是满的。令 $\chi\in H^1(H,\QQ/\ZZ)$，并设 $c\in C^H$ 映到 $\chi$。由交换性，$\chi'=\operatorname{Cor}(\chi)$ 是 $c'=\operatorname{Cor}(c)\in C^G$ 在第二行连接同态下的像。

若 $F\in M_1^G(\QQ/\ZZ)$ 是映到 $c'$ 的任意元素，则 $g\cdot F-F=f_{a'}$ 对唯一的 $a'\in\QQ/\ZZ$ 成立，并且 $\chi'(g)=a'$。由于 $G$ 在 $\QQ/\ZZ$ 上平凡作用，函数 $g\cdot F-F$ 实际上为常值 $a'$，因此可在任意 $x\in G$ 处求值来确定 $\chi'(g)$。取 $g_1,
\ldots,g_m$ 为 $H$ 在 $G$ 中左陪集代表，则 $c'=\sum_{i=1}^m g_i\cdot c$；相应地可取
\[
  F=\sum_{i=1}^m g_i\cdot f
\]
其中 $f\in M_1^G(\QQ/\ZZ)$ 映到 $c$。若对 $h\in H$，$h\cdot f-f=f_a$，则 $\chi(h)=a$。计算 $g\cdot F-F$ 在 $1\in G$ 处的值，得到
\[
  \chi'(g)=\sum_{i=1}^m f(gg_i)-\sum_{i=1}^m f(g_i).
\]
对 $i=1,
\ldots,m$，写
\[
  gg_i=g_{j_i}h(g,g_i),
  \qquad h(g,g_i)
\in H,
  \tag{17.23}
\]
其中所得 $g_{j_i}$ 是 $g_1,
\ldots,g_m$ 的一个排列。于是
\[
  \chi'(g)=\sum_{i=1}^m \chi(h(g,g_i)),
\]
所以转移同态由公式
\[
  \operatorname{Ver}(g)=\prod_{i=1}^m h(g,g_i)
  \tag{17.24}
\]
给出，元素 $h(g,g_i)\in H$ 由方程 (17.23) 定义。特别地，这证明了 (17.24) 中定义的映射是从 $G^{\mathrm{ab}}$ 到 $H^{\mathrm{ab}}$ 的同态，并且与 (17.23) 中代表元 $g_i$ 的选择无关。直接证明这个映射是同态并不完全平凡。当 $G$ 无限而 $H$ 在 $G$ 中指数有限时，同一公式也定义转移同态。

例如，设 $H=n\ZZ$，$G=\ZZ$，并选 $0,1,2,
\ldots,n-1$ 为 $H$ 在 $G$ 中的陪集代表。若 $g=1$，则对 $i=1,2,
\ldots,n-1$，所有 $h(g,g_i)$ 都为 $0$，而 $h(1,n-1)=n$。因此从 $\ZZ$ 到 $n\ZZ$ 的转移映射把 $1$ 送到 $n$，所以只是乘以指数。类似地，从任意循环群 $G$ 到指数为 $n$ 的子群 $H$ 的转移映射是 $n$ 次幂映射。另见练习 8。

对奇素数 $p$ 的循环群 $\mathbb F_p^\times$ 及其子群 $\{\pm1\}$，可知转移映射是同态 $\operatorname{Ver}:\mathbb F_p^\times\to\{\pm1\}$，由
\[
  \operatorname{Ver}(a)=a^{(p-1)/2}=\left(\frac ap\right)=
  \begin{cases}
  +1,& a\text{ 是平方},\\
  -1,& a\text{ 不是平方},
  \end{cases}
\]
给出（符号 $\left(\frac ap\right)$ 称为 Legendre 符号）。

现在考虑半直积 $E=A\rtimes G$，其中 $G$ 在 Abel 群 $A$ 上作用。和第 5.5 节一样，把 $E$ 的元素写为 $(a,g)$，其中 $a\in A$，$g\in G$，群运算为
\[
  (a_1,g_1)(a_2,g_2)=(a_1+g_1a_2,g_1g_2).
\]
注意，$A$ 用加法记号，而 $G$ 与 $E$ 用乘法记号。

\noindent\textbf{定义.} 设 $X$ 是任意群，$Y$ 是 $X$ 的正规子群。列 $1\trianglelefteq Y\trianglelefteq X$ 的\emph{稳定群}是 $X$ 的所有这样的自同构组成的群：它把 $Y$ 映到自身，并且在两个因子 $Y$ 与 $X/Y$ 上均作为恒等作用；即
\[
\operatorname{Stab}(1\trianglelefteq Y\trianglelefteq X)=\{\alpha\in\operatorname{Aut}(X)\mid \alpha(y)=y\text{ 对所有 }y\in Y,
\text{ 且 }\alpha(x)=x\bmod Y\text{ 对所有 }x\in X\}.
\]

在 $Y$ 是 $X$ 的 Abel 正规子群的特例中，由 $Y$ 的元素共轭诱导 $X$ 的（内）自同构，它们稳定列 $1\trianglelefteq Y\trianglelefteq X$；在此情形下，$Y/C_Y(X)$ 同构于 $\operatorname{Stab}(1\trianglelefteq Y\trianglelefteq X)$ 的一个子群，其中 $C_Y(X)$ 是 $Y$ 中位于 $X$ 的中心的元素。

\medskip
\noindent\textbf{命题 31.} 令 $A$ 是 $G$-模，令 $E=A\rtimes G$。对每个上循环 $f\in Z^1(G,A)$，定义 $\alpha_f:E\to E$ 为
\[
  \alpha_f((a,g))=(a+f(g),g).
\]
则映射 $f\mapsto\alpha_f$ 是从 $Z^1(G,A)$ 到 $\operatorname{Stab}(1\trianglelefteq A\trianglelefteq E)$ 上的群同构。在此同构下，上边界子群 $B^1(G,A)$ 映到稳定群的子群 $A/C_A(E)$ 上。

\begin{proof}
容易验证，上循环条件蕴含 $\alpha_f$ 是 $E$ 的自同构，并稳定链 $1\trianglelefteq A\trianglelefteq E$。同样可直接检查 $\alpha_{f_1+f_2}=\alpha_{f_1}\circ\alpha_{f_2}$，所以映射 $f\mapsto\alpha_f$ 是群同态。由 $\alpha_f$ 的定义，该映射是单射。反过来，令 $\alpha\in\operatorname{Stab}(1\trianglelefteq A\trianglelefteq E)$。由于 $\alpha$ 在 $E/A$ 上作用平凡，半直积中每个元素 $(0,g)$ 在 $\alpha$ 下都映到同一 $A$-陪集中的另一个元素 $(a,g)$；定义 $f_\alpha:G\to A$ 为 $f_\alpha(g)=a$。若把 $A$ 识别为 $E$ 中形如 $(a,1)$ 的元素，则 $E$ 中群运算说明
\[
  f_\alpha(g)=\alpha((0,g))(0,g)^{-1}.
\]
由于 $\alpha$ 是 $E$ 的稳定自同构，容易检查 $f_\alpha$ 满足上循环条件。由定义立刻得到 $f_{\alpha_f}=f$，所以 $f\mapsto\alpha_f$ 是同构。

现在 $f$ 是上边界，当且仅当存在某个 $x\in A$ 使得对所有 $g\in G$ 有 $f(g)=x-gx$。这当且仅当
\[
  \alpha_f((a,g))=(a+x-gx,g).
\]
但 $E$ 中由元素 $(x,1)$ 共轭把 $(a,g)$ 映到同一元素 $(a+x-gx,g)$，所以自同构 $\alpha_f$ 是由 $(x,1)$ 的共轭。这证明了命题剩余断言。
\end{proof}

\medskip
\noindent\textbf{推论 32.} 在命题 31 的记号下，对任意 $a\in A$，令 $\varphi_a$ 表示 $E$ 中由 $a$ 共轭给出的自同构。则上循环 $f_1$ 与 $f_2$ 在 $H^1(G,A)$ 中属于同一上同调类，当且仅当对某个 $a\in A$，
\[
  \alpha_{f_1}=\varphi_a\circ\alpha_{f_2}.
\]

命题和推论说明，可以通过寻找稳定列 $1\trianglelefteq A\trianglelefteq E$ 的自同构来计算 $1$-上循环，反之亦然。第一上同调群则由把这些自同构模去内自同构得到，也就是该列的“外稳定自同构”群。

\paragraph{例}
令 $G=\ZZ/2\ZZ$ 通过取逆作用于 $A=\ZZ/4\ZZ$。对应的半直积 $E=A\rtimes G$ 是 $8$ 阶二面体群，它的自同构群同构于 $D_8$；把 $E$ 看作 $D_{16}$ 的正规（指数 $2$）子群，则后者中的共轭限制到 $E$ 给出 $E$ 的 $8$ 个不同自同构（参见第 4.4 节命题 17）。$E$ 的子群 $A$ 在 $E$ 中是特征的，因此 $E$ 的每个自同构都把 $A$ 映到自身，因而也在 $E/A$ 上作用（由于 $|E/A|=2$，必然平凡）。$E$ 的一半自同构取逆 $A$，一半中心化 $A$；事实上，$D_{16}$ 中的 $8$ 阶循环子群（含 $A$）映到一个中心化 $A$ 的 $4$ 阶循环自同构群。因此
\[
  \operatorname{Stab}(1\trianglelefteq A\trianglelefteq E)\cong\ZZ/4\ZZ\cong Z^1(G,A).
\]
由于 $E$ 的中心是 $A$ 中的 $2$ 阶子群，$|A/Z(E)|=2=|B^1(G,A)|$。这证明 $|H^1(G,A)|=2$。

在半直积 $E$ 中，子群 $G$ 是 $A$ 的一个补，即 $E=AG$ 且 $A\cap G=1$；此外，$G$ 的每个 $E$-共轭也是 $A$ 的补。但 $A$ 在 $E$ 中可能有不与 $G$ 在 $E$ 中共轭的补。我们对 $H^1(G,A)$ 的第二种解释说明，这个上同调群刻画 $A$ 在 $E$ 中的补的 $E$-共轭类。

\medskip
\noindent\textbf{命题 33.} 令 $A$ 是 $G$-模，令 $E=A\rtimes G$。对每个 $1$-上循环 $f$，令
\[
  G_f=\{(f(g),g)\mid g\in G\}.
\]
则 $G_f$ 是 $E$ 中 $A$ 的补子群。映射 $f\mapsto G_f$ 是从 $Z^1(G,A)$ 到 $E$ 中 $A$ 的补的集合的双射。两个补在 $E$ 中共轭，当且仅当它们对应的 $1$-上循环在 $H^1(G,A)$ 中属于同一上同调类；因此 $H^1(G,A)$ 与 $A$ 的补的 $E$-共轭类集合之间存在一个双射。

\begin{proof}
由上循环条件，
\[
  (f(g),g)(f(h),h)=(f(g)+g\cdot f(h),gh)=(f(gh),gh),
\]
因此 $G_f$ 对 $E$ 中群运算封闭。前面已经观察到每个上循环满足 $f(1)=0$，所以 $G_f$ 含 $E$ 的恒等元 $(0,1)$。$(f(g),g)$ 的逆为 $(f(g^{-1}),g^{-1})$，所以 $G_f$ 对逆封闭。这证明 $G_f$ 是 $E$ 的子群。由于 $G_f$ 的不同元素代表 $E$ 中 $A$ 的不同陪集，$G_f$ 是 $A$ 的补。不同上循环给出不同陪集代表，因此给出不同补。

反过来，若 $C$ 是 $G$ 中 $A$ 的任意补，则 $C$ 对每个 $g\in G$ 含有 $Ag$ 的唯一陪集代表 $a_gg$。由于 $C$ 对群运算封闭，元素 $(a_gg)(a_hh)=(a_gga_hg^{-1})gh$ 代表陪集 $Agh$，所以 $a_{gh}=a_g(ga_hg^{-1})=a_g(g\cdot a_h)$（用 $A$ 的加法记号写为 $a_{gh}=a_g+g\cdot a_h$）。这说明映射 $f:G\to A$，$f(g)=a_g$ 是上循环，并且 $C=G_f$。故 $1$-上循环与 $A$ 在 $E$ 中的补之间存在双射。

由于 $\operatorname{Stab}(1\trianglelefteq A\trianglelefteq E)$ 正规化 $A$，它置换 $A$ 在 $E$ 中的补。在命题 31 的记号中，对 $1$-上循环 $f_1$ 与 $f_2$，由定义立刻得到
\[
  \alpha_{f_1}(G_{f_2})=G_{f_1+f_2}.
\]
这说明稳定群在补集合上的置换作用是这个群的左正则表示。进一步，若 $a\in A$ 且 $\varphi_a$ 是由 $a$ 共轭得到的稳定自同构，则
\[
  \varphi_a(G_f)=G_{f+\beta_a},
  \tag{17.25}
\]
其中 $\beta_a$ 是 $1$-上边界 $g\mapsto a-g\cdot a$。由于 $G_f$ 是 $A$ 的补，任意 $e\in E$ 可写为 $ag$，其中 $a\in A$，$g\in G_f$。于是 $eG_fe^{-1}=aG_fa^{-1}$，即 $G_f$ 的 $E$-共轭正是 $A$-共轭。现在，补 $G_{f_1}$ 与 $G_{f_2}$ 在 $E$ 中共轭，当且仅当由 (17.25) 对某个 $a\in A$ 有 $G_{f_2}=aG_{f_1}a^{-1}=G_{f_1+\beta_a}$。这说明两个补共轭，当且仅当对应上循环相差一个上边界，即代表 $H^1(G,A)$ 中同一上同调类。证明完成。
\end{proof}

\medskip
\noindent\textbf{推论 34.} 在命题 33 的记号下，$A$ 的所有补在 $E$ 中共轭，当且仅当 $H^1(G,A)=0$。

\medskip
\noindent\textbf{推论 35.} 若 $A$ 是有限 Abel 群且其阶与 $|G|$ 互素，则在任意半直积 $E=A\rtimes G$ 中，$A$ 的所有补都在 $E$ 中共轭。

\paragraph{例}
\begin{enumerate}[label=(\arabic*)]
  \item 令 $A=(a)$ 与 $G=(g)$ 都是 $2$ 阶循环群。群 $G$ 必须在 $A$ 上平凡作用，因此 $A\rtimes G=A\times G$ 是 Klein 四元群。这里 $A\rtimes G$ 是 Abel 群，所以每个子群只与自身共轭；又由于 $H^1(G,A)=\Hom(\ZZ/2\ZZ,\ZZ/2\ZZ)$ 阶为 $2$，故 $E$ 中 $A$ 恰有两个补，即 $(g)$ 与 $(ag)$。
  \item 若 $A=(a)$ 是 $2$ 阶循环群，而 $G=(x)\times(y)$ 是 Klein 四元群，则仍然 $G$ 必须在 $A$ 上平凡作用，所以 $H^1(G,A)=\Hom(\ZZ/2\ZZ\times\ZZ/2\ZZ,\ZZ/2\ZZ)$ 阶为 $4$。$A\times G$ 中 $A$ 的四个补是 $G$、$(ax,y)$、$(x,ay)$ 与 $(ax,ay)$。
  \item 命题 33 也可用于计算 $H^1(G,A)$。令 $A=(r)$ 是 $4$ 阶循环群，$G=(s)$ 是 $2$ 阶循环群，按取逆作用在 $A$ 上：$srs^{-1}=r^{-1}$，如推论 32 后例中。则 $A\rtimes G$ 是 $8$ 阶二面体群 $D_8$。子群 $A$ 在 $D_8$ 中有四个补，即由 $A$ 外每个 $2$ 阶元素生成的群：$(s)$、$(r^2s)$、$(rs)$ 与 $(r^3s)$。前一对和后一对在 $D_8$ 中共轭（两种情形下均经由 $r$），但 $(s)$ 不与 $(rs)$ 共轭。因此 $A$ 在 $A\rtimes G$ 中有两个补的共轭类，所以 $H^1(\ZZ/2\ZZ,\ZZ/4\ZZ)$ 的阶为 $2$。这也可由第 2 节中循环群上同调的计算得到。
\end{enumerate}

\subsection*{练习}

\begin{enumerate}[label=\arabic*.]
  \item 令 $G$ 是 $2$ 阶循环群，$A$ 是 $G$-模。计算下列各情形中 $Z^1(G,A)$、$B^1(G,A)$ 与 $H^1(G,A)$ 的同构类型：
  \begin{enumerate}[label=(\alph*)]
    \item $A=\ZZ/4\ZZ$（平凡作用）；
    \item $A=\ZZ/2\ZZ\times\ZZ/2\ZZ$（平凡作用）；
    \item $A=\ZZ/2\ZZ\times\ZZ/2\ZZ$（任意非平凡作用）。
  \end{enumerate}

  \item 令 $p$ 是素数，$P$ 是 $p$-群。
  \begin{enumerate}[label=(\alph*)]
    \item 证明 $H^1(P,\mathbb F_p)\cong P/\Phi(P)$，其中 $\Phi(P)$ 是 $P$ 的 Frattini 子群（参见第 6.1 节练习）。
    \item 推出 $H^1(P,\mathbb F_p)$ 作为 $\mathbb F_p$ 上向量空间的维数等于 $P$ 的最小生成元个数。［用第 6.1 节练习 26(c)。］
  \end{enumerate}

  \item 若 $G$ 是 $2$ 阶循环群，并通过取逆作用在 $\ZZ$ 上，证明 $|H^1(G,\ZZ)|=2$。［说明在 $E=\ZZ\rtimes G$ 中，$E-\ZZ$ 的每个元素阶为 $2$，且该陪集中有两个共轭类。］

  \item 令 $A$ 是 Klein 四元群，令 $G=\operatorname{Aut}(A)\cong S_3$ 以自然方式作用在 $A$ 上。证明 $H^1(G,A)=0$。［说明在半直积 $E=A\rtimes G$ 中，$G$ 是 $E$ 的 Sylow $3$-子群的正规化子。应用 Sylow 定理说明 $A$ 在 $E$ 中的所有补都共轭。］

  \item 令 $G$ 是 $2$ 阶循环群，作用在阶为 $2^n$ 的初等 Abel $2$-群 $A$ 上。证明 $H^1(G,A)=0$ 当且仅当 $n=2k$ 且 $|A^G|=2^k$。［在 $E=A\rtimes G$ 中说明 $(a,x)$ 是 $2$ 阶元素当且仅当 $a\in A^G$，其中 $G=(x)$。然后比较 $A$ 的补的个数与 $x$ 的 $E$-共轭的个数。］

  \item （Thompson 转移引理）设 $G$ 是偶数阶有限群，$T$ 是 $G$ 的 Sylow $2$-子群，$M\leq T$ 且 $|T:M|=2$，并令 $x$ 是 $G$ 中一个 $2$ 阶元素。如下证明：若 $G$ 没有指数为 $2$ 的子群，则 $M$ 包含 $x$ 的某个 $G$-共轭。
  \begin{enumerate}[label=(\alph*)]
    \item 令 $\operatorname{Ver}:G/[G,G]\to T/[T,T]$ 为转移同态。证明
    \[
      \operatorname{Ver}(x)=\prod_g g^{-1}xg\pmod{[T,T]},
    \]
    其中乘积取遍在 $x$ 左乘下固定的左陪集 $gT$ 的代表。
    \item 证明在左乘下，$x$ 固定 $G$ 中 $T$ 的奇数个左陪集。
    \item 证明若 $G$ 没有指数为 $2$ 的子群，则 $\operatorname{Ver}(x)\in M/[T,T]$。推出对某个 $g\in G$ 必有 $g^{-1}xg\in M$。［在 $2$ 阶群 $T/M$ 中考虑乘积 $\operatorname{Ver}(x)$。］
  \end{enumerate}

  \item 令 $H$ 是 $G$ 的子群，$x\in G$。转移 $\operatorname{Ver}:G/[G,G]\to H/[H,H]$ 可按如下计算：令 $O_1,O_2,
  \ldots,O_k$ 为 $x$ 通过左乘作用于 $H$ 在 $G$ 中左陪集上的不同轨道，令 $O_i$ 长度为 $n_i$，并令 $g_iH$ 为 $O_i$ 的任意代表。
  \begin{enumerate}[label=(\alph*)]
    \item 证明 $O_i=\{g_iH,xg_iH,x^2g_iH,
    \ldots,x^{n_i-1}g_iH\}$，且 $g_i^{-1}x^{n_i}g_i\in H$。
    \item 证明
    \[
      \operatorname{Ver}(x)=\prod_{i=1}^k g_i^{-1}x^{n_i}g_i\pmod{[H,H]}.
    \]
  \end{enumerate}

  \item 假设 $G$ 的中心 $Z(G)$ 在 $G$ 中指数为 $m$。证明 $\operatorname{Ver}(x)=x^m$ 对所有 $x\in G$ 成立，其中 $\operatorname{Ver}$ 是从 $G/[G,G]$ 到 $Z(G)$ 的转移同态。［用前一练习。］

  \item 令 $p$ 是素数，$n\geq3$，令 $V$ 是 $\mathbb F_p$ 上 $n$ 维向量空间，基为 $v_1,v_2,
  \ldots,v_n$。令 $V$ 成为对称群 $S_n$ 的模，其中每个 $\pi\in S_n$ 以自然方式置换基：$\pi(v_i)=v_{\pi(i)}$。
  \begin{enumerate}[label=(\alph*)]
    \item 证明
    \[
    |H^1(S_n,V)|=\begin{cases}0,&p\neq2,\\2,&p=2.
    \end{cases}
    \]
    ［用 Shapiro 引理。］
    \item 对所有素数 $p$，证明 $H^1(A_n,V)=0$。
  \end{enumerate}

  \item 令 $V$ 是前一练习中所述 $S_n$ 在 $\mathbb F_2$ 上的自然置换模，$n\geq3$，并令
  \[
    W=\{a_1v_1+\cdots+a_nv_n\mid a_1+\cdots+a_n=0\}
  \]
  是 $V$ 的“迹零”子模。证明若 $n$ 为偶数，则 $H^1(A_n,W)\neq0$。［说明在半直积 $V\rtimes A_n$ 中，元素 $v_1$ 在 $E=W\rtimes A_n$ 上诱导一个非平凡外自同构，并稳定列 $1\trianglelefteq W\trianglelefteq E$。］

  \item 设 $F$ 是特征不整除 $n$ 的域，$a$ 是 $F$ 中任意非零元素。令 $K$ 是 $F$ 的一个 Galois 扩张，包含 $x^n-a$ 的分裂域，并令 $\sqrt[n]{a}$ 是 $K$ 中固定的一个 $n$ 次根。
  \begin{enumerate}[label=(\alph*)]
    \item 证明 $\sigma(\sqrt[n]{a})/\sqrt[n]{a}$ 是 $n$ 次单位根。
    \item 证明函数 $f(\sigma)=\sigma(\sqrt[n]{a})/\sqrt[n]{a}$ 是 $G$ 中取值于 $K$ 中 $n$ 次单位根群 $\mu_n$ 的 $1$-上循环（注意不假设 $\mu_n\subseteq F$）。
    \item 证明由 $K$ 中 $a$ 的另一个 $n$ 次根得到的 $1$-上循环，与 (b) 中的 $1$-上循环相差一个 $1$-上边界。
  \end{enumerate}

  \item 设 $F$ 是特征不整除 $n$ 且包含 $n$ 次单位根的域，并假设 $L/F$ 是 Galois 扩张，其 Galois 群为指数整除 $n$ 的 Abel 群。证明 $L$ 是 $F$ 的若干循环扩张的复合，且这些循环扩张的次数都是 $n$ 的因子；并用此证明乘法群 $F^\times/(F^\times)^n$ 的子群与这样的扩张 $L$ 之间存在双射。

  \item 扩张 $\CC/\RR$ 的 Galois 群是由复共轭 $\tau$ 生成的 $2$ 阶循环群 $G=(\tau)$。证明
  \[
    H^2(G,\CC^\times)\cong\RR^\times/\RR^+\cong\ZZ/2\ZZ,
  \]
  其中 $\RR^+$ 表示正实数。

  \item 对任意群 $G$，令 $\widehat G=\Hom(G,\QQ/\ZZ)$ 表示其对偶群。
  \begin{enumerate}[label=(\alph*)]
    \item 若 $\varphi:G_1\to G_2$ 是群同态，证明与 $\varphi$ 复合诱导对偶群上的同态 $\widehat\varphi:\widehat{G_2}\to\widehat{G_1}$。
    \item 对 $G$ 中固定的 $g$，说明在 $g$ 处取值给出从 $\widehat G$ 到 $\QQ/\ZZ$ 的同态 $\varphi_g$。
    \item 证明把 $g\in G$ 送到 (b) 中 $\varphi_g$ 的映射，定义从 $G$ 到其二重对偶 $\widehat{\widehat G}$ 的同态。
    \item 证明若 $G$ 是有限 Abel 群，则 (c) 中的同态是 $G$ 与其二重对偶的同构。（由第 5.2 节练习 14，群 $G$ 与其对偶 $\widehat G$ 非典范地同构。本题说明 $G$ 与其二重对偶典范同构——该同构不依赖于 $G$ 的生成元选择。）
    \item 若 $\psi:\widehat{G_2}\to\widehat{G_1}$ 是同态，其中 $G_1,G_2$ 是有限 Abel 群，则由 (a) 与 (d) 有诱导同态 $\varphi:G_1\to G_2$。证明 $\varphi(g_1)=g_2$ 当且仅当对所有 $\chi\in\widehat{G_2}$，若 $\chi'=\psi(\chi)$，则 $\chi(g_2)=\chi'(g_1)$。
  \end{enumerate}

  \item 在计算从 $\mathbb F_p^\times$ 到 $\{\pm1\}$ 的转移映射中使用 Gauss 引理，证明 $2$ 是奇素数 $p$ 模平方，当且仅当 $p\equiv\pm1\pmod 8$。［数一数 $2,4,
  \ldots,p-1$ 中有多少个元素大于 $(p-1)/2$。］
\end{enumerate}


\section[群扩张、因子集与 H2(G,A)]{群扩张、因子集与 $H^2(G,A)$}

若 $A$ 是一个 $G$-模，则由式 (17.18) 中上边界映射 $d_2$ 的定义，函数 $f:G\times G\to A$ 是一个 $2$-上循环，当且仅当它满足恒等式
\[
  f(g,h)+f(gh,k)=g\cdot f(h,k)+f(g,hk)\qquad\text{对所有 }g,h,k\in G.
  \tag{17.26}
\]
等价地，一个 $2$-上循环由一族 $A$ 中元素 $\{a_{g,h}\}_{g,h\in G}$ 决定，满足
\[
  a_{g,h}+a_{gh,k}=g\cdot a_{h,k}+a_{g,hk}
\]
对所有 $g,h,k\in G$ 成立（于是 $2$-上循环 $f$ 就是把 $(g,h)$ 送到 $a_{g,h}$ 的函数）。

一个 $2$-上链 $f$ 是上边界，当且仅当存在函数 $f_1:G\to A$ 使得
\[
  f(g,h)=g f_1(h)-f_1(gh)+f_1(g),
  \qquad\text{对所有 }g,h\in G,
  \tag{17.27}
\]
也就是说，$f$ 是 $1$-上链 $f_1$ 在 $d_1$ 下的像。

本节的主要结果之一，是把 $2$-上循环 $Z^2(G,A)$ 与 $G$ 被 $A$ 扩张时出现的因子集联系起来；这些因子集是在定义扩张中的乘法时，考虑不同陪集代表元的影响而产生的。特别地，我们将证明，$G$ 被 $A$ 扩张的等价类（其中 $G$ 在 $A$ 上的作用固定）与 $H^2(G,A)$ 的元素之间存在一个双射。

我们先观察扩张的一些基本事实。设 $E$ 是 $G$ 被 $A$ 的一个群扩张
\[
  1\to A\xrightarrow{\iota}E\xrightarrow{\pi}G\to 1.
  \tag{17.28}
\]
扩张 (17.28) 如下决定 $G$ 在 $A$ 上的一个作用。对每个 $g\in G$，令 $e_g$ 是 $E$ 中映到 $g$ 的一个元素（这样一组 $G$ 在 $E$ 中的代表元的选择称为 $\pi$ 的一个集合论截面）。元素 $e_g$ 通过共轭作用于 $E$ 的正规子群 $\iota(A)$，把 $\iota(a)$ 映到 $e_g\iota(a)e_g^{-1}$。$E$ 中映到 $g$ 的任何另一个元素都形如 $e_g\iota(a_1)$，其中 $a_1\in A$；又因为 $\iota(A)$ 是 Abel 的，所以此元素在 $\iota(A)$ 上的共轭作用与 $e_g$ 的共轭作用相同，因而与 $g$ 的代表元选择无关。因此 $G$ 作用在 $\iota(A)$ 上；又因 $\iota$ 是单射，也作用在 $A$ 上。由于共轭是自同构，扩张 (17.28) 把 $A$ 定义成一个 $G$-模。

回忆第 10.5 节，两个扩张
\[
  1\to A\xrightarrow{\iota_1}E_1\xrightarrow{\pi_1}G\to1,
  \qquad
  1\to A\xrightarrow{\iota_2}E_2\xrightarrow{\pi_2}G\to1
\]
称为等价的，是指存在群同构 $\beta:E_1\to E_2$，使得图
\[
\begin{array}{ccccccccc}
1&\longrightarrow&A&\xrightarrow{\iota_1}&E_1&\xrightarrow{\pi_1}&G&\longrightarrow&1\\
 &&\Vert&&\downarrow\beta&&\Vert&&\\
1&\longrightarrow&A&\xrightarrow{\iota_2}&E_2&\xrightarrow{\pi_2}&G&\longrightarrow&1
\end{array}
\tag{17.29}
\]
交换。这时我们简称 $\beta$ 是两个扩张之间的等价。如第 10.5 节中所指出的，扩张的等价是自反、对称和传递的。我们还观察到：
\[
  \text{等价扩张定义 $A$ 上相同的 $G$-模结构。}
\]
为了说明这一点，假设 (17.29) 是一个等价，令 $g$ 是 $G$ 的任意元素，令 $e_g$ 是 $E_1$ 中任意一个由 $\pi_1$ 映到 $g$ 的元素。在 $E_1$ 中由共轭给出的 $g$ 对 $A$ 的作用，把每个 $a$ 映为
\[
  \iota_1^{-1}(e_g\iota_1(a)e_g^{-1}).
\]
令 $e'_g=\beta(e_g)$。由于图交换，$\pi_2(e'_g)=g$，所以第二个扩张中 $g$ 在 $A$ 上的作用由 $e'_g$ 的共轭给出。这个共轭把 $a$ 映为
\[
  \iota_2^{-1}(e'_g\iota_2(a)e_g^{'-1}).
\]
由于 $\iota_1,\iota_2$ 与 $\beta$ 都是单射，$g$ 在 $a$ 上的两个作用相等，当且仅当它们在 $E_2$ 中的像相等，也就是说，当且仅当
\[
  \beta\circ\iota_1\bigl(\iota_1^{-1}(e_g\iota_1(a)e_g^{-1})\bigr)=e'_g\iota_2(a)e_g^{'-1}.
\]
这个等式由 $e'_g$ 的定义和图的交换性立即得到。

下面我们看一个如 (17.28) 的扩张如何定义 $Z^2(G,A)$ 中的一个 $2$-上循环。为简单起见，我们通过 $\iota$ 把 $A$ 视为 $E$ 的一个子群，并通过 $\pi$ 把 $G$ 视为 $E/A$。

\noindent\textbf{定义.} 映射 $\mu:G\to E$ 若满足 $\pi\circ\mu(g)=g$ 且 $\mu(1)=0$，也就是说，对每个 $g\in G$，$\mu(g)$ 是 $E$ 的陪集 $Ag$ 的一个代表元，而 $E$ 的单位元（即 $A$ 的零元）代表单位陪集，则称 $\mu$ 是 $\pi$ 的一个\emph{规范化截面}。

固定 (17.28) 中 $\pi$ 的一个截面 $\mu$。$E$ 的每个元素都可唯一地写成 $a\mu(g)$ 的形式，其中 $a\in A$、$g\in G$。对 $g,h\in G$，$E$ 中乘积 $\mu(g)\mu(h)$ 位于陪集 $Agh$ 中，故存在唯一的 $A$ 中元素 $f(g,h)$，使得
\[
  \mu(g)\mu(h)=f(g,h)\mu(gh)
  \qquad\text{对所有 }g,h\in G.
  \tag{17.30}
\]
若进一步 $\mu$ 在单位元处规范化，则还得到
\[
  f(g,1)=0=f(1,g)
  \qquad\text{对所有 }g\in G.
  \tag{17.31}
\]

\noindent\textbf{定义.} 由式 (17.30) 定义的函数 $f$ 称为与截面 $\mu$ 相关联的扩张 $E$ 的\emph{因子集}。若 $f$ 还满足 (17.31)，则称 $f$ 是\emph{规范化因子集}。

我们将在后面的例子中看到，不同的截面 $\mu$ 可能给出相同的因子集 $f$。

现在验证因子集 $f$ 实际上是一个 $2$-上循环。首先注意到，$E$ 中的群运算可写为
\[
\begin{aligned}
  (a_1\mu(g))(a_2\mu(h))
  &=(a_1+\mu(g)a_2\mu(g)^{-1})\mu(g)\mu(h)\\
  &=(a_1+g\cdot a_2)(\mu(g)\mu(h))\\
  &=(a_1+g\cdot a_2+f(g,h))\mu(gh),
\end{aligned}
\tag{17.32}
\]
其中 $g\cdot a_2$ 表示由 $E$ 中共轭给出的 $G$-模作用。现在用 (17.32) 与 $E$ 中的结合律，以两种方式计算乘积 $\mu(g)\mu(h)\mu(k)$：
\[
\begin{aligned}
  (\mu(g)\mu(h))\mu(k)&=(f(g,h)+f(gh,k))\mu(ghk),\\
  \mu(g)(\mu(h)\mu(k))&=(g f(h,k)+f(g,hk))\mu(ghk).
\end{aligned}
\tag{17.33}
\]
因此 (17.33) 两边右端的 $A$ 中因子对所有 $g,h,k\in G$ 都相等，这正是 $f$ 的 $2$-上循环条件 (17.26)。这说明由扩张 $E$ 和任意截面 $\mu$ 得到的因子集，是 $Z^2(G,A)$ 的一个元素。

接着看因子集 $f$ 如何依赖于截面 $\mu$ 的选择。设 $\mu'$ 是同一扩张 $E$ 在 (17.28) 中的另一个截面，并令 $f'$ 是它的因子集。那么对所有 $g\in G$，$\mu(g)$ 与 $\mu'(g)$ 都位于同一陪集 $Ag$ 中，故存在函数 $f_1:G\to A$，使得 $\mu'(g)=f_1(g)\mu(g)$ 对所有 $g$ 成立。于是
\[
  \mu'(g)\mu'(h)=f'(g,h)\mu'(gh)=(f'(g,h)+f_1(gh))\mu(gh).
\]
另一方面，
\[
\begin{aligned}
  \mu'(g)\mu'(h)
  &=(f_1(g)\mu(g))(f_1(h)\mu(h))\\
  &=(f_1(g)+g\cdot f_1(h))(\mu(g)\mu(h))\\
  &=(f_1(g)+g\cdot f_1(h)+f(g,h))\mu(gh).
\end{aligned}
\]
比较这两个关于 $\mu'(g)\mu'(h)$ 的表达式中的 $A$ 中因子，得到
\[
  f'(g,h)=f(g,h)+\bigl(gf_1(h)-f_1(gh)+f_1(g)\bigr)
  \qquad\text{对所有 }g,h\in G,
\]
也就是说，$f$ 与 $f'$ 相差 $f_1$ 的如 (17.27) 中的 $2$-上边界。

我们已经说明，与扩张 $E$ 对应的不同截面所给出的因子集，是 $Z^2(G,A)$ 中相差 $B^2(G,A)$ 中一个上边界的 $2$-上循环。因此，对扩张 $E$ 关联着 $H^2(G,A)$ 中一个良定义的上同调类，它由任一截面 $\mu$ 的 (17.30) 中因子集决定。

若 $G$ 被 $A$ 的扩张 $E$ 是分裂扩张（也就是说 $E=A\rtimes G$ 是由给定的 $G$ 在 $A$ 上的共轭作用得到的半直积），则存在一个从 $G$ 到 $E$ 的截面 $\mu$，它是群同态。在这种情况下，(17.30) 中的因子集 $f$ 恒为 $0$：$f(g,h)=0$ 对所有 $g,h\in G$。因此由分裂扩张定义的 $H^2(G,A)$ 中的上同调类是平凡类。

现在设 $\beta$ 是 (17.28) 中扩张与另一个扩张 $E'$ 之间的等价：
\[
\begin{array}{ccccccccc}
1&\longrightarrow&A&\xrightarrow{\iota}&E&\xrightarrow{\pi}&G&\longrightarrow&1\\
 &&\Vert&&\downarrow\beta&&\Vert&&\\
1&\longrightarrow&A&\xrightarrow{\iota'}&E'&\xrightarrow{\pi'}&G&\longrightarrow&1.
\end{array}
\]
若 $\mu$ 是 $\pi$ 的截面，则 $\mu'=\beta\circ\mu$ 是 $\pi'$ 的截面，所以可用刚才证明的结果来确定对应于 $E'$ 的 $H^2(G,A)$ 中的上同调类。把同态 $\beta$ 作用在式 (17.30) 上，得到
\[
  \beta(\mu(g))\beta(\mu(h))=\beta(f(g,h))\beta(\mu(gh))
  \qquad\text{对所有 }g,h\in G.
\]
由于 $\beta$ 在 $A$ 上限制为恒等映射，这就是
\[
  \mu'(g)\mu'(h)=f(g,h)\mu'(gh)
  \qquad\text{对所有 }g,h\in G,
\]
这说明与 $\mu'$ 相关联的 $E'$ 的因子集，与与 $\mu$ 相关联的 $E$ 的因子集相同。因此等价扩张定义 $H^2(G,A)$ 中相同的上同调类。

接下来说明如何反过来进行：给定 $H^2(G,A)$ 中的一个类，我们构造一个扩张 $E_f$，使得其对应因子集位于 $H^2(G,A)$ 中给定的类内。这个过程推广了第 5.5 节中的半直积构造（它是 $f$ 为代表平凡类的零上循环时的特例）。

首先注意，若一个 $2$-上循环来自如上扩张的因子集，且截面 $\mu$ 是规范化的，则它满足 (17.31) 中的条件。

\noindent\textbf{定义.} 若 $2$-上循环 $f$ 满足 $f(g,1)=0=f(1,g)$ 对所有 $g\in G$ 成立，则称 $f$ 是\emph{规范化 $2$-上循环}。

当 $f$ 是规范化上循环时，$E_f$ 的构造稍简单；为简单起见，我们只在这种情况下说明构造（$f$ 非规范化时所需的小修改在练习 4 中指出）。

我们先说明任意 $2$-上循环 $f$ 都与某个规范化 $2$-上循环位于同一个上同调类中。令 $d_1f_1$ 是 $G$ 上取常值 $f(1,1)$ 的常值函数 $f_1$ 的 $2$-上边界。则 $f(1,1)=d_1f_1(1,1)$，并且由 $2$-上循环条件容易验证 $f-d_1f_1$ 是规范化的。

因此可以假设 $H^2(G,A)$ 中的上同调类由规范化 $2$-上循环 $f$ 表示。令 $E_f$ 为集合 $A\times G$，并在 $E_f$ 上定义二元运算
\[
  (a_1,g)(a_2,h)=(a_1+g\cdot a_2+f(g,h),gh),
  \tag{17.34}
\]
其中照例 $g\cdot a_2$ 表示 $G$ 在 $A$ 上的模作用。直接检查可知群公理成立：由于 $f$ 是规范化的，单位元是 $(0,1)$，逆元由
\[
  (a,g)^{-1}=(-g^{-1}\cdot a-f(g^{-1},g),g^{-1})
  \tag{17.35}
\]
给出。上循环条件通过与前面 (17.32)、(17.33) 类似的计算蕴含结合律——细节留作练习。

由于 $f$ 是规范化 $2$-上循环，$A^*=\{(a,1)\mid a\in A\}$ 是 $E_f$ 的一个子群，而映射 $\iota^*:a\mapsto(a,1)$ 是从 $A$ 到 $A^*$ 的同构。此外，由 (17.34) 与 (17.35) 可得
\[
  (0,g)(a,1)(0,g)^{-1}=(g\cdot a,1)
  \qquad\text{对所有 }g\in G\text{ 与 }a\in A.
  \tag{17.36}
\]
由于 $E_f$ 由 $A^*$ 以及所有 $(0,g)$（$g\in G$）生成，(17.36) 蕴含 $A^*$ 是 $E_f$ 的正规子群。进一步，由 (17.34) 立即知道映射 $\pi^*:(a,g)\mapsto g$ 是从 $E_f$ 到 $G$ 的满同态，其核为 $A^*$，即 $E_f/A^*\cong G$。因此
\[
  1\to A\xrightarrow{\iota^*}E_f\xrightarrow{\pi^*}G\to1
  \tag{17.37}
\]
是 $G$ 被 $A$ 的一个具体扩张；其中 (17.36) 还保证，在这个扩张中由共轭给出的 $G$ 在 $A$ 上的作用，就是确定 $H^2(G,A)$ 中 $2$-上循环 $f$ 时所指定的模作用。扩张序列 (17.37) 表明，这个扩张具有规范化截面 $\mu(g)=(0,g)$，它对应的规范化因子集就是 $f$。注意，这不仅证明了 $H^2(G,A)$ 中的每个上同调类都来自某个扩张 $E$，而且证明了每个规范化 $2$-上循环都是某个扩张的规范化因子集。

最后，设 $f'$ 是同一上同调类中另一个规范化 $2$-上循环，令 $E_{f'}$ 为对应扩张。若 $f$ 与 $f'$ 相差函数 $f_1:G\to A$ 的上边界，即
\[
  f(g,h)-f'(g,h)=g f_1(h)-f_1(gh)+f_1(g)
\]
对所有 $g,h\in G$ 成立，则令 $g=h=1$ 可得 $f_1(1)=0$。定义
\[
  \beta:E_f\to E_{f'},\qquad \beta((a,g))=(a+f_1(g),g).
\]
显然 $\beta$ 是双射，并且
\[
\begin{aligned}
  \beta((a_1,g)(a_2,h))
  &=\beta((a_1+g\cdot a_2+f(g,h),gh))\\
  &=(a_1+g\cdot a_2+f(g,h)+f_1(gh),gh)\\
  &=(a_1+f_1(g)+g\cdot(a_2+f_1(h))+f'(g,h),gh)\\
  &=(a_1+f_1(g),g)(a_2+f_1(h),h)\\
  &=\beta((a_1,g))\beta((a_2,h)),
\end{aligned}
\]
所以 $\beta$ 是从 $E_f$ 到 $E_{f'}$ 的同构。

$\beta$ 在 $A$ 上的限制为 $\beta((a,1))=(a+f_1(1),1)=(a,1)$，所以 $\beta$ 在 $A$ 上是恒等映射。同样，$\beta$ 在 $(a,g)$ 的第二个分量上也是恒等的，因此 $\beta$ 在商群 $G$ 上诱导恒等映射。于是 $\beta$ 定义扩张 $E_f$ 与 $E_{f'}$ 之间的一个等价。这说明扩张 $E_f$ 的等价类只依赖于 $H^2(G,A)$ 中 $f$ 的上同调类。

我们把以上讨论总结为如下定理。

\medskip
\noindent\textbf{定理 36.} 设 $A$ 是一个 $G$-模。则
\begin{enumerate}[label=(\arabic*)]
  \item 函数 $f:G\times G\to A$ 是某个 $G$ 被 $A$ 的扩张 $E$ 的规范化因子集（其中共轭作用由 $A$ 上的 $G$-模作用给出），当且仅当 $f$ 是 $Z^2(G,A)$ 中的规范化 $2$-上循环。
  \item (1) 中所述扩张 $E$ 的等价类与 $H^2(G,A)$ 中的上同调类之间存在双射。该双射把扩张 $E$ 送到 $E$ 关于 $G$ 到 $E$ 的任一规范化截面 $\mu$ 的规范化因子集 $f$ 所在的类；并把 $H^2(G,A)$ 中的上同调类 $c$ 送到由 $c$ 中任一规范化上循环 $f$ 给出的扩张 (17.37) 所定义的扩张 $E_f$。
  \item 在 (2) 的双射下，分裂扩张对应于平凡上同调类。
\end{enumerate}

\medskip
\noindent\textbf{推论 37.} $G$ 被 Abel 群 $A$ 的每个扩张都分裂，当且仅当 $H^2(G,A)=0$。

\medskip
\noindent\textbf{推论 38.} 若 $A$ 是有限 Abel 群且 $(|A|,|G|)=1$，则 $G$ 被 $A$ 的每个扩张都分裂。

\begin{proof}
这立即由第 2 节推论 29 得出。
\end{proof}

我们可以用推论 38 证明一个不要求 $A$ 为 Abel 群的同一结论。

\medskip
\noindent\textbf{定理 39.} （Schur 定理）若 $E$ 是任意有限群，并包含一个正规子群 $N$，使得 $|N|$ 与 $[E:N]$ 互素，则 $N$ 在 $E$ 中有一个补子群。

\noindent\textbf{注.} 回忆，一个阶与指数互素的子群称为 Hall 子群。因此 Schur 定理说，每个正规 Hall 子群都有一个补子群，使群分裂为半直积。

\begin{proof}
对 $|E|$ 作归纳。可以假设 $N\neq1$。令 $p$ 是整除 $|N|$ 的素数，令 $P$ 是 $N$ 的一个 Sylow $p$-子群。令 $E_0$ 为 $P$ 在 $E$ 中的正规化子，并令 $N_0=N\cap E_0$。由 Frattini 论证（第 6.1 节命题 6），$E=E_0N$。由第二同构定理可知，$N_0$ 是 $E_0$ 的一个（正规）Hall 子群，且 $|E_0:N_0|=|E:N|$（参见第 3.3 节练习 10）。

若 $E_0<E$，则由归纳假设应用于 $E_0$ 中的 $N_0$，得到 $E_0$ 包含 $N_0$ 的一个补子群 $K$。由于 $|K|=|E_0:N_0|$，$K$ 也是 $E$ 中 $N$ 的补子群，正合所需。因此可以假设 $E_0=E$，即 $P$ 在 $E$ 中正规。

由于 $P$ 的中心 $Z(P)$ 在 $P$ 中特征，它在 $E$ 中正规（参见第 4.4 节）。若 $Z(P)=N$，则 $N$ 是 Abel 的，此时定理由推论 38 得到。因此可以假设 $Z(P)\neq N$。用横线表示模去 $Z(P)$ 的商群。则 $\overline N$ 是 $\overline E$ 的正规 Hall 子群。由归纳假设，它在 $\overline E$ 中有补子群 $\overline K$。令 $E_1$ 为 $\overline K$ 在 $E$ 中的完全原像。则
\[
  |E_1|=|\overline K|\,|Z(P)|=|E/N|\,|Z(P)|,
\]
所以 $Z(P)$ 是 $E_1$ 的正规 Hall 子群。再次由归纳假设，$Z(P)$ 在 $E_1$ 中有一个补子群；由阶数考虑可知这个补子群也正是 $E$ 中 $N$ 的补子群。证明完毕。
\end{proof}

\noindent\textbf{例.}
\begin{enumerate}[label=(\arabic*)]
  \item 若 $G=\ZZ_2$ 且 $A=\ZZ/2\ZZ$，则 $G$ 在 $A$ 上平凡作用，所以由第 2 节中循环群上同调的计算，$H^2(G,A)=A^G/NA=\ZZ/2\ZZ$。由定理 36，$G$ 被 $A$ 的扩张恰有两个不等价类。它们是 $4$ 阶循环群和 Klein $4$-群；后者是分裂的，因此对应于 $H^2$ 中的平凡类。

  \item 若 $G=(g)\cong\ZZ_2$，而 $A=(a)\cong\ZZ/4\ZZ$ 是一个 $4$ 阶群，且 $G$ 在 $A$ 上平凡作用，则由循环群上同调计算，$H^2(G,A)=A/2A\cong\ZZ/2\ZZ$。同前一例，有两个不等价的 $G$ 被 $A$ 的扩张；显然它们是群 $\ZZ_8$ 与 $\ZZ_4\times\ZZ_2$，其中后一个分裂扩张对应于平凡上同调类。

  若 $E=(r)\times(s)$ 表示 $G$ 被 $A$ 的分裂扩张，其中 $|r|=4$、$|s|=2$，则 $\mu_i(g)=r^is$，$i=0,\ldots,3$，给出 $G$ 在 $E$ 中的四个规范化截面。截面 $\mu_0,\mu_2$ 都给出零因子集 $f$。截面 $\mu_1,\mu_3$ 都给出因子集 $f'$，其中 $f'(g,g)=a^2\in A$。$f$ 与 $f'$ 都给出位于 $H^2(G,A)$ 的平凡上同调类中的规范化 $2$-上循环。与零 $2$-上循环 $f$ 对应的扩张 $E_f$ 是以元素 $(a,1)$ 与 $(1,g)$ 为 $\ZZ_4\times\ZZ_2$ 的通常生成元（阶分别为 $4$ 与 $2$）的群。然而在 $E_{f'}$ 中，$(a,1)$ 的阶为 $4$，而 $(1,g)$ 的阶也为 $4$，因为
  \[
    (1,g)^2=(f'(g,g),g^2)=(a^2,1).
  \]
  $2$-上循环 $f$ 与 $f'$ 相差 $f_1$ 的上边界，其中 $f_1(1)=1$ 且 $f_1(g)=a$。同构 $\beta(a,g)=(a+f_1(g),g)$ 从 $E_f$ 到 $E_{f'}$，把 $E_f$ 的生成元 $(a,1)$ 与 $(1,g)$ 分别送到 $E_{f'}$ 的生成元 $(a,1)$ 与 $(a,g)$，并给出这两个扩张之间的显式等价。

  $G$ 通过取逆作用在 $A$ 上的情形在练习 3 中处理。

  \item 设 $G=\ZZ_2$，而 $A$ 是 Klein $4$-群。若 $G$ 非平凡地作用于 $A$，则 $G$ 交换 $A$ 的两个非单位元，比如 $a$ 与 $b$，并固定第三个非单位元 $c$。于是 $A^G=NA=\{1,c\}$，所以 $H^2(G,A)=0$，从而 $G$ 被 $A$ 的每个扩张 $E$ 都分裂。也可直接看出这一点：由于作用非平凡，这样的群必须是非 Abel 群，因而必须是 $D_8$。从第 2.5 节中 $D_8$ 的格可以看出，对每个 Klein $4$-群，$D_8$ 中都有一个不包含在该 $4$-群中的 $2$ 阶子群，并且该子群分裂该扩张。

  若 $G$ 平凡地作用于 $A$，则 $H^2(G,A)=A/2A\cong A$，所以在这种情况下有 $4$ 个不等价的 $G$ 被 $A$ 的扩张。这些在练习 1 中讨论。
\end{enumerate}

\noindent\textbf{例.} （$8$ 阶群与 $H^2(\ZZ_2\times\ZZ_2,\ZZ/2\ZZ)$）令 $G=\{1,a,b,c\}$ 为 Klein $4$-群，令 $A=\ZZ/2\ZZ$。$2$-群 $G$ 必平凡地作用于 $A$。$H^2(G,A)$ 的元素分类 $8$ 阶扩张 $E$，其中 $E$ 关于某个 $\ZZ_2$ 子群的商群同构于 Klein $4$-群。虽然这样的群按群同构只有四个，但我们将看到存在八个不等价扩张。

由于 $G\times G$ 有 $16$ 个元素，$|C^2(G,A)|=2^{16}$。这里的上循环条件 (17.26) 化为
\[
  f(g,h)+f(gh,k)=f(h,k)+f(g,hk)
  \qquad\text{对所有 }g,h,k\in G.
  \tag{17.38}
\]
对于上循环子群 $Z^2(G,A)$，有下列关系：
\begin{enumerate}[label=(\arabic*)]
  \item $f(g,1)=f(1,g)=f(1,1)$，对所有 $g\in G$；
  \item $f(g,1)+f(g,a)+f(g,b)+f(g,c)=0$，对所有 $g\in G$；
  \item $f(1,h)+f(a,h)+f(b,h)+f(c,h)=0$，对所有 $h\in G$。
\end{enumerate}
第一条关系由在 (17.38) 中令 $h=k=1$ 以及令 $g=h=1$ 得到。后两条关系分别由在 (17.38) 中令 $g=h$ 与 $h=k$，再使用关系 (1) 与 (2) 得到。由此，每个 $2$-上循环 $f$ 可由 $\mathbb F_2$ 中的向量 $(\alpha,\beta,\gamma,\delta,\epsilon)$ 表示，其中
\[
  \alpha=f(1,g)=f(g,1)\quad\text{对所有 }g\in G,
\]
且
\[
  \beta=f(a,a),\qquad \gamma=f(a,b),\qquad \delta=f(b,a),\qquad \epsilon=f(b,b),
\]
因为上面的关系于是决定 $f$ 的其余值：
\[
\begin{array}{lll}
  f(a,c)=\alpha+\beta+\gamma, & f(b,c)=\alpha+\delta+\epsilon, & f(c,a)=\alpha+\beta+\delta,\\
  f(c,b)=\alpha+\gamma+\epsilon, & f(c,c)=\alpha+\beta+\gamma+\epsilon. &
\end{array}
\]
因此 $|Z^2(G,A)|\leq2^5$。虽然最终可以证明满足这些关系的每个函数都是 $2$-上循环（因而阶正好为 $32$），但这将由下面的其他考虑推出。

上循环 $f$ 是上边界，当且仅当存在函数 $f_1:G\to A$，使得
\[
  f(g,h)=f_1(h)-f_1(gh)+f_1(g),
  \qquad\text{对所有 }g,h\in G.
\]
容易看出这个上边界条件等价于以下条件：
\begin{enumerate}[label=(\roman*)]
  \item $f(g,1)=f(1,g)=f(g,g)$，对所有 $g\in G$；
  \item 只要 $g,h$ 是不同的非单位元且 $g',h'$ 也是不同的非单位元，就有 $f(g,h)=f(g',h')$。
\end{enumerate}
这些关系等价于 $\alpha=\beta=\epsilon$ 且 $\gamma=\delta$。因此 $B^2(G,A)$ 由向量 $(\alpha,\alpha,\gamma,\gamma,\alpha)$ 组成，于是 $H^2(G,A)$ 的维数至多为 $3$（即阶至多为 $2^3=8$）。容易看出，$\{(0,\beta,\gamma,0,\epsilon)\}$（其中 $\beta,\gamma,\epsilon\in\mathbb F_2$）给出 $Z^2(G,A)/B^2(G,A)$ 的一组代表元，并且这些代表上循环都是规范化的。现在我们通过显式给出八个不等价群扩张来证明 $|H^2(G,A)|=8$（并且也证明 $|Z^2(G,A)|=2^5$）。

设 $E$ 是 $G$ 被 $A$ 的一个扩张，为简单起见假设 $A\leq E$。若 $\mu:G\to E$ 是一个截面，则 $E$ 关于 $\mu$ 的因子集满足
\[
  \mu(g)\mu(h)=f(g,h)\mu(gh).
\]
群 $E$ 由 $\mu(a),\mu(b)$ 与 $A$ 生成，而 $A$ 包含在 $E$ 的中心中，因为 $G$ 在 $A$ 上平凡作用。因此 $E$ 是 Abel 群，当且仅当 $\mu(a)\mu(b)=\mu(b)\mu(a)$；由上面的关系，这当且仅当 $f(a,b)=f(b,a)$。若 $g$ 是 $G$ 的一个非单位元，则也可从上面的关系看出，$\mu(g)$ 在 $E$ 中是 $2$ 阶元素，当且仅当 $f(g,g)=0$。由于 $A$ 包含在 $E$ 的中心中，任意非单位陪集 $A\mu(g)$ 中的两个元素有相同的阶（或者为 $2$，或者为 $4$）。

含有一个 $2$ 阶正规子群且商群同构于 Klein $4$-群的 $8$ 阶群有四个：$\ZZ_2\times\ZZ_2\times\ZZ_2$、$\ZZ_4\times\ZZ_2$、$D_8$ 与 $Q_8$。

群 $E\cong\ZZ_2\times\ZZ_2\times\ZZ_2$ 是 $G$ 被 $A$ 的分裂扩张，其因子集为 $f=0$。

当 $E\cong Q_8$ 时，在四元数群的通常记号中 $A=(-1)$。在这个（非 Abel）群中，每个非单位陪集都由 $4$ 阶元素组成；这个性质是 $Q_8$ 所特有的。因此所得因子集 $f$ 满足 $f(g,g)\neq0$ 对 $G$ 的所有非单位元 $g$ 成立。

当 $E\cong\ZZ_4\times\ZZ_2=(x)\times(y)$ 时，必须有 $A=(x^2)$。陪集 $Ax$ 与 $Axy$ 都由 $4$ 阶元素组成，陪集 $Ay$ 由 $2$ 阶元素组成。因此 $\mu(a),\mu(b),\mu(c)$ 中恰有一个元素阶为 $2$，另外两个阶为 $4$。这提示定义三个从 $E$ 到 $G$ 的同态，在生成元上给出为
\[
\begin{array}{ll}
  \pi_1(y)=a, & \pi_1(x)=b,\\
  \pi_2(y)=b, & \pi_2(x)=a,\\
  \pi_3(y)=c, & \pi_3(x)=a.
\end{array}
\]
这些同态都满射到 $G$，核为 $A$，并且分别使 $\mu(a)$、$\mu(b)$、$\mu(c)$ 是 $E$ 中 $2$ 阶元素。任何在 $A$ 上为恒等的 $E$ 到自身的同构，都必须把由 $2$ 阶元素组成的 $A$ 的唯一非单位陪集 $Ay$ 映到它自身。因此任何与由 $\pi_1$ 定义的扩张等价的扩张，也必须把 $y$ 映到 $a$（因为等价在 $G$ 上是恒等的）。于是由 $\pi_1,\pi_2,\pi_3$ 定义的三个扩张不等价。

当 $E\cong D_8=(r,s)$ 时，情形类似。此时 $A=(r^2)$，陪集 $As$ 与 $Asr$ 由 $2$ 阶元素组成，而陪集 $Ar$ 由 $4$ 阶元素组成。在这种情况下，$\mu(a),\mu(b),\mu(c)$ 中恰有一个元素阶为 $4$，另外两个阶为 $2$，于是得到三个同态，在生成元上定义为
\[
\begin{array}{ll}
  \pi_1(r)=a, & \pi_1(s)=b,\\
  \pi_2(r)=b, & \pi_2(s)=a,\\
  \pi_3(r)=c, & \pi_3(s)=a.
\end{array}
\]
同前，相应扩张两两不等价。

$G$ 被 $A$ 的 $8$ 个不等价扩张的存在性证明 $|H^2(G,A)|=8$，因而这些就是所有不等价扩张的完整列表。特别地，由把 $y$ 映到 $a$、把 $x$ 映到 $c$ 的同态 $\pi'_1$ 所定义的扩张 $E'_1\cong\ZZ_4\times\ZZ_2$，必定等价于上面的扩张 $E_1$（其他两个同构于 $\ZZ_4\times\ZZ_2$ 的扩张以及三个 $D_8$ 扩张也类似）。这证明这些群存在某些外自同构；参见练习 9。

\noindent\textbf{注.} 对任意素数 $p$，基本 Abel 群 $E_{p^m}$ 以有限域 $\mathbb F_p$ 为系数的上同调群，可通过把它们与直积各因子的上同调群联系起来而确定，如第 2 节末所述。一般地，$H^2(E_{p^m},\mathbb F_p)$ 是 $\mathbb F_p$ 上维数为 $\frac12m(m+1)$ 的向量空间。当 $p=2$ 且 $m=2$ 时，这就是上面的结果 $H^2(\ZZ_2\times\ZZ_2,\ZZ/2\ZZ)\cong(\ZZ/2\ZZ)^3$。

\subsection*{交叉积代数与 Brauer 群}

设 $F$ 是一个域。回忆，$F$-代数 $B$ 是一个环，它在其中心中含有域 $F$，且 $B$ 的单位元就是 $F$ 的单位元，参见第 10.1 节。

\noindent\textbf{定义.} 若 $F$-代数 $A$ 不含非平凡的真（双边）理想，则称 $A$ 是\emph{单的}。若单 $F$-代数 $A$ 的中心为 $F$，则称 $A$ 是\emph{中心单 $F$-代数}。

最简单的中心单 $F$-代数包括系数在 $F$ 中的 $n\times n$ 矩阵代数 $M_n(F)$。

若 $K/F$ 是有限 Galois 域扩张，其 Galois 群为 $G=\operatorname{Gal}(K/F)$，则可使用 $Z^2(G,K^\times)$ 中的规范化 $2$-上循环来构造某些中心单 $F$-代数。从 $2$-上循环构造这些代数，并以 $H^2(G,K^\times)$ 对它们分类（参见下文定理 42），是同调方法在数论中的重要应用。这一构造在 $G$ 为循环群时的情形，是导致抽象上同调发展的先驱之一。

设 $f=\{a_{\sigma,\tau}\}_{\sigma,\tau\in G}$ 是 $Z^2(G,K^\times)$ 中的规范化 $2$-上循环。令 $B_f$ 为以 $u_\sigma$（$\sigma\in G$）为基的 $K$ 上向量空间：
\[
  B_f=\left\{\sum_{\sigma\in G}\alpha_\sigma u_\sigma\mid \alpha_\sigma\in K\right\}.
  \tag{17.39}
\]
在 $B_f$ 上定义乘法为
\[
  u_\sigma\alpha=\sigma(\alpha)u_\sigma,
  \qquad
  u_\sigma u_\tau=a_{\sigma,\tau}u_{\sigma\tau}
  \tag{17.40}
\]
其中 $\alpha\in K$，$\sigma,\tau\in G$。第二个等式说明 $a_{\sigma,\tau}$ 给出了 $B_f$ 中元素 $u_\sigma$ 的一个“因子集”，这也是使用这一术语的原因之一。用这个乘法可得
\[
  (u_\sigma u_\tau)u_\rho=a_{\sigma,\tau}a_{\sigma\tau,\rho}u_{\sigma\tau\rho},
  \qquad
  u_\sigma(u_\tau u_\rho)=\sigma(a_{\tau,\rho})a_{\sigma,\tau\rho}u_{\sigma\tau\rho}.
\]
由于
\[
  a_{\sigma,\tau}a_{\sigma\tau,\rho}=\sigma(a_{\tau,\rho})a_{\sigma,\tau\rho}
\]
是上循环条件 (17.26) 的乘法形式，所以 (17.40) 中定义的乘法是结合的。

由于上循环是规范化的，$a_{1,\sigma}=a_{\sigma,1}=1$ 对所有 $\sigma\in G$ 成立，并由 (17.40) 得知元素 $u_1$ 是 $B_f$ 的单位元。把 $K$ 与 $B_f$ 中的元素 $\alpha u_1$ 认同，可见 $B_f$ 是一个含有域 $K$ 的 $F$-代数；若 $n=[K:F]=|G|$，则 $B_f$ 在 $F$ 上的维数为 $n^2$。

\medskip
\noindent\textbf{命题 40.} 由 (17.39) 中 $K$-向量空间基 $u_\sigma$ 和 (17.40) 中乘法定义的 $F$-代数 $B_f$ 是一个中心单 $F$-代数。

\begin{proof}
还需证明 $B_f$ 的中心为 $F$，并且 $B_f$ 不含非零真理想。设 $x=\sum_{\sigma\in G}\alpha_\sigma u_\sigma$ 是 $B_f$ 中心中的元素。那么对 $\beta\in K$，等式 $x\beta=\beta x$ 说明若 $\alpha_\sigma\neq0$，则 $\sigma(\beta)=\beta$。对任意 $\sigma\neq1$，存在不被 $\sigma$ 固定的元素 $\beta\in K$，所以 $\alpha_\sigma=0$ 对所有 $\sigma\neq1$ 成立，因而 $x=\alpha_1u_1$。于是 $xu_\tau=u_\tau x$ 当且仅当 $\tau(\alpha_1)=\alpha_1$；若这对所有 $\tau$ 成立，则必须有 $\alpha_1\in F$。因此 $B_f$ 的中心是 $F$。

为证明 $B_f$ 是单的，设 $I$ 是 $B_f$ 的非零理想，令
\[
  x=\sum_{i=1}^m\alpha_{\sigma_i}u_{\sigma_i}
\]
是 $I$ 中具有最少非零项数 $m$ 的非零元素。若 $m>1$，则存在 $\beta\in K^\times$，使得 $\sigma_m(\beta)\neq\sigma_{m-1}(\beta)$。于是元素
\[
  x-\sigma_m(\beta)x\beta^{-1}
\]
属于理想 $I$，其 $u_{\sigma_{m-1}}$ 的系数非零，而 $u_{\sigma_m}$ 的系数为 $0$，故非零项数少于 $x$，矛盾。因此 $m=1$，即 $x=\alpha u_\sigma$，其中 $\alpha\in K$ 且 $\sigma\in G$。此元素是单位，其逆形如 $\sigma^{-1}(\alpha^{-1})u_\sigma^{-1}$ 乘以适当的因子，因此 $I=B_f$。证明完毕。
\end{proof}

\noindent\textbf{定义.} 由 (17.39) 与 (17.40) 定义的中心单 $F$-代数 $B_f$ 称为因子集 $\{a_{\sigma,\tau}\}$ 的\emph{交叉积代数}。

若 $f'=\{a'_{\sigma,\tau}\}$ 是 $H^2(G,K^\times)$ 中与 $\{a_{\sigma,\tau}\}$ 位于同一上同调类的规范化上循环，则存在元素 $b_\sigma\in K^\times$，使得
\[
  a'_{\sigma,\tau}=a_{\sigma,\tau}\bigl(\sigma(b_\tau)b_{\sigma\tau}^{-1}b_\sigma\bigr),
\]
这就是上边界条件 (17.27) 的乘法形式。若 $B_{f'}$ 是由该上循环如 (17.39)、(17.40) 定义的具有 $K$-基 $u'_{\sigma}$ 的 $F$-代数，则由 $u'_{\sigma}\mapsto b_\sigma u_\sigma$ 定义的 $K$-向量空间同态 $\varphi$ 满足
\[
\begin{aligned}
  \varphi(u'_{\sigma}u'_{\tau})
  &=\varphi(a'_{\sigma,\tau}u'_{\sigma\tau})
    =a'_{\sigma,\tau}b_{\sigma\tau}u_{\sigma\tau}\\
  &=b_\sigma\sigma(b_\tau)u_\sigma u_\tau
    =(b_\sigma u_\sigma)(b_\tau u_\tau)
    =\varphi(u'_{\sigma})\varphi(u'_{\tau}).
\end{aligned}
\]
所以 $\varphi$ 是从 $B_{f'}$ 到 $B_f$ 的 $F$-代数同构。

我们已经说明，$H^2(G,K^\times)$ 中的每个上同调类 $c$ 定义了一个中心单 $F$-代数的同构类，即代表 $c$ 的任意规范化上循环 $\{a_{\sigma,\tau}\}$ 的交叉积代数的同构类。下一个结果说明，平凡上同调类对应于包含 $M_n(F)$ 的同构类。

\medskip
\noindent\textbf{命题 41.} $H^2(G,K^\times)$ 中平凡上同调类的交叉积代数同构于矩阵代数 $M_n(F)$，其中 $n=[K:F]$。

\begin{proof}
若 $a\in K$，则乘以 $a$ 定义了 $K$ 作为 $F$ 上 $n$ 维向量空间时的一个线性变换 $T_a$。类似地，每个自同构 $\sigma\in G$ 定义了 $K$ 的一个 $F$-线性变换 $T_\sigma$；选取 $K$ 在 $F$ 上的一个基后，可把 $T_a$ 与 $T_\sigma$ 都视为 $M_n(F)$ 的元素。若 $B_0$ 表示平凡因子集（所有 $\sigma,\tau\in G$ 都有 $a_{\sigma,\tau}=1$）的交叉积代数，考虑加法映射
\[
  \varphi:B_0\to M_n(F),\qquad \varphi(a u_\sigma)=T_aT_\sigma.
\]
由于对 $a\in F$ 有 $T_{\alpha a}=aT_\alpha$，映射 $\varphi$ 是一个 $F$-向量空间同态。若 $x\in K$，则
\[
  T_\sigma T_\alpha(x)=T_\sigma(\alpha x)=\sigma(\alpha x)=\sigma(\alpha)\sigma(x)=T_{\sigma(\alpha)}T_\sigma(x),
\]
所以作为 $K$ 上的线性变换，有 $T_\sigma T_\alpha=T_{\sigma(\alpha)}T_\sigma$。由 $u_\sigma u_\tau=u_{\sigma\tau}$ 推出
\[
\begin{aligned}
  \varphi((\alpha u_\sigma)(\beta u_\tau))
  &=\varphi(\alpha\sigma(\beta)u_{\sigma\tau})
    =T_{\alpha\sigma(\beta)}T_{\sigma\tau}\\
  &=T_\alpha T_{\sigma(\beta)}T_\sigma T_\tau
    =T_\alpha T_\sigma T_\beta T_\tau
    =\varphi(\alpha u_\sigma)\varphi(\beta u_\tau),
\end{aligned}
\]
这说明 $\varphi$ 是从 $B_0$ 到 $M_n(F)$ 的 $F$-代数同态。由于 $\ker\varphi$ 是 $B_0$ 的理想且 $\varphi\neq0$，由命题 40 得 $\ker\varphi=0$，故 $\varphi$ 是单射。由于 $B_0$ 与 $M_n(F)$ 作为 $F$ 上向量空间都有维数 $n^2$，所以 $\varphi$ 是 $F$-代数同构。
\end{proof}

\noindent\textbf{例.} 若 $K=\CC$ 且 $F=\RR$，则 $G=\operatorname{Gal}(\CC/\RR)$ 阶为 $2$，由复共轭 $\tau$ 生成。我们有 $|H^2(G,\CC^\times)|=2$。对应于平凡类的中心单 $\RR$-代数 $B_0$ 是 $\CC u_1\oplus\CC u_\tau$，满足 $u_\tau(a+bi)=(a-bi)u_\tau$ 且 $u_\tau^2=u_1$。它在映射
\[
  \varphi((a+bi)u_1+(c+di)u_\tau)
  =aI+bT_i+cT_\tau+dT_iT_\tau
  =\begin{pmatrix}a+c&-b+d\\ b+d&a-c\end{pmatrix}
\]
下同构于矩阵代数 $M_2(\RR)$。

代表非平凡上同调类的规范化上循环 $f$ 由取值
\[
  a_{1,1}=a_{1,\tau}=a_{\tau,1}=1,
  \qquad a_{\tau,\tau}=-1
\]
定义。相应的中心单 $\RR$-代数 $B_f$ 为 $\CC v_1\oplus\CC v_\tau$。元素 $v_1$ 是 $B_f$ 的单位元，且有关系
\[
  v_\tau(a+bi)=(a-bi)v_\tau,
  \qquad v_\tau^2=-v_1.
\]
令 $v_1=1$ 且 $v_\tau=j$，可见 $B_f$ 作为 $\RR$-代数同构于实 Hamilton 四元数代数 $\RR+\RR i+\RR j+\RR k$。

单代数有丰富的理论。我们不加证明地提及以下结果。令 $A$ 是在 $F$ 上有限维的中心单 $F$-代数。
\begin{enumerate}[label=\Roman*.]
  \item 若 $F\subseteq B\subseteq A$，其中 $B$ 是单 $F$-代数，定义 $B$ 在 $A$ 中的中心化子 $B^c$ 为 $A$ 中与 $B$ 的所有元素交换的元素集合。定义反代数 $B^{\operatorname{opp}}$ 为具有相反乘法的集合 $B$，也就是说，$B^{\operatorname{opp}}$ 中的乘积 $b_1b_2$ 由 $B$ 中的乘积 $b_2b_1$ 给出。$B^c$ 与 $B^{\operatorname{opp}}$ 都是单 $F$-代数，并且
  \begin{enumerate}[label=\alph*.]
    \item $(\dim_F B)(\dim_F B^c)=\dim_F A$；
    \item $A\otimes_F B^{\operatorname{opp}}\cong M_r(B^c)$，作为 $F$-代数，其中 $r=\dim_F B$；
    \item 若 $B$ 是中心单 $F$-代数，则 $B\otimes_F B^c\cong A$。
  \end{enumerate}

  \item 若 $A'$ 是 Artin（满足左理想降链条件）的单 $F$-代数，则 $A\otimes_F A'$ 是 Artin 单 $F$-代数，其中心为 $(A')^c$。

  \item 有 $A\cong M_r(\Delta)$，其中 $\Delta$ 是中心为 $F$ 的某个除环，$r\geq1$。除环 $\Delta$ 与整数 $r$ 由 $A$ 唯一决定。同样的陈述对任意 Artin 单 $F$-代数也成立。
\end{enumerate}
最后一个结果是下一章中更详细描述的 Wedderburn 定理的一部分。

\noindent\textbf{定义.} 若 $A$ 是中心单 $F$-代数，且包含 $F$ 的域 $L$ 满足 $A\otimes_F L\cong M_m(L)$ 对某个 $m\geq1$ 成立，则称 $L$\emph{分裂} $A$。

由 (II) 可知，$\Delta$ 的每个极大交换子代数都是一个域 $E$，满足 $E=E^c=E^{\operatorname{opp}}$；若 $[E:F]=m$，则得到 $\dim_F\Delta=m^2$。把 (II) 应用于 $A=\Delta$ 与 $B=E$，也可见 $\Delta\otimes_F E\cong M_m(E)$。还可以证明，中心单 $F$-代数 $A$ 的一个极大子域 $E$ 也满足 $E=E^c=E^{\operatorname{opp}}$，所以再次由 (II) 得 $A\otimes_F E\cong M_r(E)$（其中 $r^2=\dim_F A$）。

若 $A=M_r(\Delta)$，则域 $L$ 分裂 $A$ 当且仅当 $L$ 分裂 $\Delta$，如下所示。若 $\Delta\otimes_F L\cong M_n(L)$，则
\[
  A\otimes_F L\cong M_r(\Delta)\otimes_F L
  \cong M_r(\Delta\otimes_F L)
  \cong M_r(M_n(L))
  \cong M_{rn}(L).
\]
反过来，若 $A\otimes_F L\cong M_n(L)$，则
\[
  M_n(L)\cong M_r(\Delta)\otimes_F L\cong M_r(\Delta\otimes_F L).
\]
由 (II) 与 (III)，$\Delta\otimes_F L\cong M_s(\Delta')$，其中 $\Delta'$ 是某个除环。结合前一同构，(III) 中的唯一性说明 $\Delta'\cong L$，于是同构 $\Delta\otimes_F L\cong M_s(L)$ 表明 $L$ 分裂 $\Delta$。

由上面的讨论可见，$\Delta$ 的一个极大交换子域分裂 $\Delta$，也分裂任意 $r\geq1$ 的 $A\cong M_r(\Delta)$。由此不难证明，每个在 $F$ 上有限维的中心单 $F$-代数都可由 $F$ 的某个有限 Galois 扩张分裂。

在 (I) 中取 $A$ 为交叉积代数 $B_f$，取 $B=K$，得到 $K=K^c=K^{\operatorname{opp}}$ 且 $B_f\otimes_F K\cong M_n(K)$。特别地，交叉积代数 $B_f$ 总是被 $K$ 分裂。

\noindent\textbf{例.} 在上面的 Hamilton 四元数例子中，有 $B_f\otimes_{\RR}\CC\cong M_2(\CC)$。我们有 $B_f\otimes_{\RR}\CC=\CC+\CC i+\CC j+\CC k$，到 $M_2(\CC)$ 的一个显式同构 $\varphi$ 由
\[
  \varphi(i)=\begin{pmatrix}\sqrt{-1}&0\\0&-\sqrt{-1}\end{pmatrix},
  \qquad
  \varphi(j)=\begin{pmatrix}0&-1\\1&0\end{pmatrix}
\]
给出，并按 $\CC$-线性延拓。

由 (III)，每个中心单 $F$-代数 $A$ 作为 $F$-代数同构于 $M_r(\Delta)$，其中除环 $\Delta$ 在 $F$-同构意义下唯一决定，称为 $A$ 的\emph{除环部分}。

\noindent\textbf{定义.} 两个中心单 $F$-代数 $A$ 与 $B$ 称为\emph{相似}，若 $A\cong M_r(\Delta)$ 且 $B\cong M_s(\Delta)$ 对同一个除环 $\Delta$ 成立，也就是说，若 $A$ 与 $B$ 有相同的除环部分。

令 $[A]$ 表示 $A$ 的相似类。由 (II)，若 $A$ 与 $B$ 是中心单 $F$-代数，则 $A\otimes_F B$ 仍是中心单 $F$-代数，因此可在相似类上定义乘法
\[
  [A][B]=[A\otimes_F B].
\]
类 $[F]$ 是此乘法的单位元，张量积的结合性说明该乘法是结合的。由 (Ib) 在取 $B=A$ 时应用（于是因 $A$ 是中心的，$B^c=F$），得到 $[A][A^{\operatorname{opp}}]=[F]$，所以在这个乘法下逆元存在。

\noindent\textbf{定义.} 由中心单 $F$-代数的相似类组成、乘法为 $[A][B]=[A\otimes_F B]$ 的群，称为 $F$ 的\emph{Brauer 群}，记为 $\operatorname{Br}(F)$。

若 $L$ 是 $F$ 的任意扩域，则由 (II)，代数 $A\otimes_F L$ 是中心单 $L$-代数。容易检查，映射 $[A]\mapsto[A\otimes_F L]$ 是从 $\operatorname{Br}(F)$ 到 $\operatorname{Br}(L)$ 的良定义同态。该同态的核由满足 $A\otimes_F L\cong M_m(L)$（对某个 $m\geq1$）的代数 $A$ 的类组成，也就是由被 $L$ 分裂的代数 $A$ 组成。

\noindent\textbf{定义.} 若 $L/F$ 是域扩张，则相对 Brauer 群 $\operatorname{Br}(L/F)$ 是由被 $L$ 分裂的中心单 $F$-代数的相似类组成的群。等价地，$\operatorname{Br}(L/F)$ 是同态 $[A]\mapsto[A\otimes_F L]$ 从 $\operatorname{Br}(F)$ 到 $\operatorname{Br}(L)$ 的核。

下面的定理概括了这一领域中的若干主要结果，并说明 Brauer 群与上面构造的交叉积代数之间的基本联系。

\medskip
\noindent\textbf{定理 42.} 设 $K/F$ 是次数为 $n$ 的 Galois 扩张，$G=\operatorname{Gal}(K/F)$。
\begin{enumerate}[label=(\arabic*)]
  \item 维数 $\dim_F A=n^2$ 的中心单 $F$-代数 $A$ 被 $K$ 分裂，当且仅当 $A\otimes_F K\cong M_n(K)$，也当且仅当 $A$ 同构于 (17.39)、(17.40) 中的交叉积代数 $B_f$。
  \item $F$-同构类的中心单 $F$-代数 $A$（满足 $A\otimes_F K\cong M_n(K)$）与 $H^2(G,K^\times)$ 的元素之间存在双射。在此双射下，包含规范化上循环 $f$ 的类 $c\in H^2(G,K^\times)$ 对应于由 (17.39)、(17.40) 定义的交叉积代数 $B_f$ 的同构类，而平凡上同调类对应于 $M_n(F)$。
  \item 每个在 $F$ 上有限维并被 $K$ 分裂的中心单 $F$-代数，都相似于某个维数为 $n^2$ 且被 $K$ 分裂的代数。(2) 中的双射还建立了 $\operatorname{Br}(K/F)$ 与 $H^2(G,K^\times)$ 之间的双射，并且这是群同构。
  \item 被 $K$ 分裂的中心单除 $F$-代数的 $F$-同构类集合与 $H^2(G,K^\times)$ 之间存在双射。
\end{enumerate}

如前所述，每个在 $F$ 上有限维的中心单 $F$-代数都可由 $F$ 的某个有限 Galois 扩张分裂，因此
\[
  \operatorname{Br}(F)=\bigcup_K \operatorname{Br}(K/F),
\]
其中并集取遍 $F$ 的所有有限 Galois 扩张。由此可知，$\operatorname{Br}(F)$ 与 $H^2(\operatorname{Gal}(F'/F),(F')^\times)$ 之间存在一个双射，其中 $F'$ 表示 $F$ 的一个可分代数闭包。这里 $\operatorname{Gal}(F'/F)$ 被视为一个 profinite 群，而上同调群指连续 Galois 上同调。

这个结果和定理 42 的一个推论是，$F$ 上中心单除代数 $\Delta$ 的 $F$-同构类的一组完全代表，可以从 $F$ 的有限 Galois 扩张的交叉积代数的除代数部分得到。那些在 $K$ 上分裂的除代数出现在 $K/F$ 的交叉积代数中。

\noindent\textbf{例.} 由于 $H^2(\operatorname{Gal}(\mathbb F_{q^a}/\mathbb F_q),\mathbb F_{q^a}^{\times})=0$（参见练习 10），有 $\operatorname{Br}(\mathbb F_{q^a}/\mathbb F_q)=0$，因此也有 $\operatorname{Br}(\mathbb F_q)=0$。作为结果，每个有限除代数都是一个域（直接证明参见第 13.6 节练习 13），并且每个有限中心单 $\mathbb F_q$-代数都同构于一个全矩阵环 $M_r(\mathbb F_q)$。

\subsection*{习题}

\begin{enumerate}
  \item 令 $A=\{1,a,b,c\}$ 为 Klein $4$-群，令 $G=(g)$ 为 $2$ 阶循环群，平凡地作用在 $A$ 上。
  \begin{enumerate}[label=(\alph*)]
    \item 证明 $|C^2(G,A)|=2^8$。
    \item 证明上边界是常值函数，并推出 $|B^2(G,A)|=4$。
    \item 用上循环条件证明 $|Z^2(G,A)|\leq2^4$。
    \item 若 $E=\ZZ_4\times\ZZ_2=(x)\times(y)$，证明由 $\pi(x)=g$、$\pi(y)=1$ 以及 $\iota_1(a)=x^2$、$\iota_1(b)=y$（相应地，$\iota_2(b)=x^2$、$\iota_2(a)=y$，以及 $\iota_3(c)=x^2$、$\iota_3(a)=y$）定义的扩张 $1\to A\to E\to G\to1$，连同分裂扩张 $\ZZ_2\times\ZZ_2\times\ZZ_2$，给出 $G$ 被 Klein $4$-群的 $4$ 个不等价扩张。通过显式给出相应上循环，推出 $H^2(G,A)$ 的阶为 $4$。
  \end{enumerate}

  \item 令 $A=\ZZ/4\ZZ$，令 $G$ 为 $2$ 阶循环群，平凡地作用在 $A$ 上。
  \begin{enumerate}[label=(\alph*)]
    \item 证明 $|C^2(G,A)|=2^8$。
    \item 用上边界条件证明 $|B^2(G,A)|=2^3$。
    \item 用上循环条件证明 $|Z^2(G,A)|\leq2^4$。
    \item 通过给出两个不等价的 $G$ 被 $A$ 的扩张及其相应上循环，证明 $|H^2(G,A)|=2$。
  \end{enumerate}

  \item 令 $A=\ZZ/4\ZZ$，令 $G$ 为 $2$ 阶循环群，并令 $G$ 通过取逆作用在 $A$ 上。
  \begin{enumerate}[label=(\alph*)]
    \item 证明有四个上边界，并且只有零上边界是规范化的。
    \item 通过直接计算上循环群和上边界群，证明 $|H^2(G,A)|=2$。
    \item 展示两个不同的上同调类及其相应扩张群。
    \item 证明对于一个给定的 $G$ 被 $A$ 的扩张，若扩张群同构于 $D_8$，则有四个规范化截面，它们的因子集全都是零 $2$-上循环。
    \item 证明对于一个给定的 $G$ 被 $A$ 的扩张，若扩张群同构于 $Q_8$，则有十六个截面，其中四个是规范化的，并且这些规范化截面全都有相同的因子集。
  \end{enumerate}

  \item 对一个非规范化的 $2$-上循环 $f$，在集合 $A\times G$ 上用式 (17.34) 中同一二元运算定义扩张群 $E_f$。在这种情况下，通过证明单位元现在为 $(-f(1,1),1)$ 且逆元由
  \[
    (a,x)^{-1}=(-x^{-1}\cdot a-f(x^{-1},x)-f(1,1),x^{-1})
  \]
  给出，验证两个群公理。（结合律的验证本质上与规范化 $2$-上循环的情形相同。）还证明集合
  \[
    A^{**}=\{(a-f(1,1),1)\mid a\in A\}
  \]
  是 $E_f$ 的一个子群，并且映射 $\iota^{**}:a\mapsto(a-f(1,1),1)$ 是从 $A$ 到 $A^{**}$ 的同构。证明这个扩张 $E_f$，连同嵌入 $\iota^{**}$ 与到 $G$ 的通常投影映射 $\pi^*$，等价于由与 $f$ 同一类中的一个规范化上循环得到的扩张。

  \item 证明给定扩张 $1\to A\to E\to G\to1$ 与自身的等价，在复合下构成一个群，并且该群同构于稳定群 $\operatorname{Stab}(1\to\iota(A)\to E)$。（因此命题 31 蕴含 $Z^1(G,A)$ 是该扩张与自身的等价群。）

  \item （Gaschütz 定理）令 $p$ 是素数，令 $A$ 是有限群 $G$ 的 Abel 正规 $p$-子群，令 $P$ 是 $G$ 的一个 Sylow $p$-子群。证明 $G$ 是 $G/A$ 被 $A$ 的分裂扩张，当且仅当 $P$ 是 $P/A$ 被 $A$ 的分裂扩张。（注意由第 4.5 节练习 37，$A\leq P$。）［用 Sylow 定理证明若 $G$ 在 $A$ 上分裂，则 $P$ 也在 $A$ 上分裂。反过来，证明经定理 36 与 $P/A$ 被 $A$ 的扩张相关联的一个规范化 $2$-上循环，是 $H^2(G/A,A)$ 中某个规范化 $2$-上循环在限制同态
  \[
    \operatorname{Res}:H^2(G/A,A)\to H^2(P/A,A)
  \]
  下的像。然后使用命题 26 以及乘以 $|G:P|$ 是 $A$ 的一个自同构这一事实。］

  \item
  \begin{enumerate}[label=(\alph*)]
    \item 通过展示 $A_4$ 被 $2$ 阶循环群的一个非分裂扩张，证明 $H^2(A_4,\ZZ/2\ZZ)\neq0$。［见第 4.5 节练习 11。］
    \item 通过证明 $\operatorname{SL}_2(\mathbb F_5)$ 是 $A_5$ 被 $2$ 阶循环群的一个非分裂扩张，证明 $H^2(A_5,\ZZ/2\ZZ)\neq0$。［用第 4.5 节命题 21 与命题 23。］
  \end{enumerate}

  \item 有限群 $G$ 的 Schur 乘子定义为群 $H^2(G,\CC^\times)$，其中复数乘法群 $\CC^\times$ 是平凡 $G$-模。证明 Schur 乘子是有限群。［证明每个上同调类都含有一个取值在 $n$ 次单位根中的上循环，其中 $n=|G|$，方法如下：若 $f$ 是任意上循环，则由推论 27 有 $f^n\in B^2(G,\CC^\times)$。定义 $k\in C^2(G,\CC^\times)$ 为 $k(g_1,g_2)=f(g_1,g_2)^{1/n}$（任取 $n$ 次根）。证明 $k\in B^2(G,\CC^\times)$ 且 $fk^{-1}$ 取值于 $1$ 的 $n$ 次根群。］

  \item 用定理 39 后面的例子中 Klein $4$-群被 $\ZZ_2$ 的扩张分类，证明如下事实（使用该例中的记号）：
  \begin{enumerate}[label=(\alph*)]
    \item 存在 $\ZZ_4\times\ZZ_2$ 的一个（外）自同构，它交换陪集 $Ax$ 与 $Axy$，并固定陪集 $Ay$。
    \item 存在 $D_8$ 的一个外自同构，它交换陪集 $As$ 与 $Asr$，并固定陪集 $Ar$。
  \end{enumerate}

  \item 设 $\mathbb F_q$ 是有限域，$G=\operatorname{Gal}(\mathbb F_{q^d}/\mathbb F_q)=(\sigma_q)$，其中 $\sigma_q$ 是 Frobenius 自同构，并令 $N$ 是循环群 $G$ 的通常范数元素。
  \begin{enumerate}[label=(\alph*)]
    \item 用 Hilbert 定理 90 证明
    \[
      |{}_N\mathbb F_{q^d}^{\times}|=\frac{q^d-1}{q-1},
    \]
    并推出从 $\mathbb F_{q^d}$ 到 $\mathbb F_q$ 的范数映射是满射。
    \item 证明 $H^n(G,\mathbb F_{q^d}^{\times})=0$ 对所有 $n\geq1$ 成立。
  \end{enumerate}
\end{enumerate}

% \setchapterimage{最后两章是有限群表示论及其若干应用的导论。我们已经在第一部分看到，群在集合上的作用，即置换表示，是揭示群结构的基本工具。Cayley 定理和 Sylow 定理，以及第 6.1 节和第 6.2 节中的许多结果与应用，都是建立在群作用于集合之上的。Galois 理论一章发展了数学中最优美的对应之一：一个群作为域的自同构群而作用时，会给出 Galois 群的子群格与域扩张的中间子域格之间的对应。在最后这两章中，我们研究群作为向量空间上的线性变换而作用。我们主要关心的是利用这些线性作用来提供关于群本身的信息。
% 	
% 	在第三部分中我们看到，模是环的“表示对象”，其意义是：$R$-模公理规定了 $R$ 在某个 Abel 群 $M$ 上的“环作用”，并保持 $M$ 的 Abel 群结构。在 $M$ 是 $F[x]$-模的情形中，$x$ 作为从向量空间 $M$ 到自身的一个线性变换而作用。在第 12 章中，主理想整环上有限生成模的分类为我们提供了关于 $M$ 上这些线性变换的大量信息（例如标准形）。在第 16 章中，我们利用 Dedekind 整环中的理想结构，把第 12 章的结果推广到这类整环上有限生成模的分类。在本部分中，我们将遵循与研究 $F[x]$-模相似的过程，把多项式环替换为 $G$ 的群环 $FG$，并对某些域 $F$ 上所有有限生成 $FG$-模进行分类（Wedderburn 定理）。随后我们利用这个分类推出一些关于有限群的结果，例如第 19 章中 Burnside 关于 $p^aq^b$ 阶群可解性的定理。
% }
%\part{有限群表示论引论}
% % ============================================
%  chapter18.tex — 第18章：表示论与特征标理论
%  根据 853-893.pdf 翻译整理，保持原文结构、证明与练习
% ============================================


\chapter{表示论与特征标理论}

\section{线性作用与群环上的模}

除非明确另有说明，本书剩余部分所考虑的群都将是有限群。在本节中，始终设 $F$ 是一个域，$G$ 是一个有限群。我们首先引入基本术语。回忆，如果 $V$ 是 $F$ 上的向量空间，则 $\GL(V)$ 是从 $V$ 到自身的非奇异线性变换所成的群（运算为复合）；若 $n\in\ZZ^+$，则 $\GL_n(F)$ 是所有以 $F$ 中元素为项的可逆 $n\times n$ 矩阵所成的群（运算为矩阵乘法）。

\noindent\textbf{定义.} 设 $G$ 是有限群，$F$ 是域，$V$ 是 $F$ 上的向量空间。
\begin{enumerate}[label=(\arabic*)]
  \item $G$ 的一个\emph{线性表示}是任意同态 $G\to\GL(V)$。该表示的\emph{次数}是 $V$ 的维数。
  \item 设 $n\in\ZZ^+$。$G$ 的一个\emph{矩阵表示}是任意同态 $G\to\GL_n(F)$。
  \item 一个线性表示或矩阵表示称为\emph{忠实的}，如果它是单射。
  \item $G$ 在 $F$ 上的\emph{群环}是所有形如
  \[
    \sum_{g\in G} a_g g,\qquad a_g\in F
  \]
  的形式和所成的集合；加法按分量定义，乘法由 $(ag)(\beta h)=(a\beta)(gh)$（其中 $a$ 与 $\beta$ 在 $F$ 中相乘，$gh$ 是 $G$ 中的乘积）并借分配律推广到和（参见第 7.2 节）。
\end{enumerate}

除非我们专门讨论置换表示，术语“表示”总是指“线性表示”。当我们希望强调域 $F$ 时，就说 $F$-表示，或说 $G$ 在 $F$ 上的向量空间 $V$ 上的表示。

回忆，若 $V$ 是 $n$ 维有限维向量空间，则选定 $V$ 的一个基便给出同构 $\GL(V)\cong\GL_n(F)$。这样，$G$ 在有限维向量空间上的任意线性表示给出一个矩阵表示，反之亦然。在绝大多数情形中，我们的线性表示都将是有限次数的，并将在不给出特别说明时在线性表示与矩阵表示之间自由切换（只有在需要给出二者之间的显式对应时才指定基）。此外，给定一个有限次数线性表示 $\varphi:G\to\GL(V)$，相应的矩阵表示给出数值不变量（例如对 $g\in G$ 的 $\det\varphi(g)$），这些不变量与选择用来给出 $\GL(V)$ 和 $\GL_n(F)$ 之间同构的基无关。利用这些不变量将是我们后续发展的基础。

在给出表示的例子以前，我们更详细地回忆群环 $FG$（群环是在第 7.2 节中引入的，该节也讨论了若干记号与例子）。设 $G$ 的元素为 $g_1,g_2,\ldots,g_n$。$FG$ 中每个元素都形如
\[
  \sum_{i=1}^n a_i g_i,
  \qquad a_i\in F.
\]
两个形式和相等当且仅当群元素的相应系数全都相等。$FG$ 中的加法和乘法定义如下：
\[
  \sum_{i=1}^n a_i g_i+\sum_{i=1}^n b_i g_i
  =\sum_{i=1}^n (a_i+b_i)g_i,
\]
\[
  \left(\sum_{i=1}^n a_i g_i\right)\left(\sum_{j=1}^n b_j g_j\right)
  =\sum_{k=1}^n\left(\sum_{g_i g_j=g_k}a_i b_j\right)g_k,
\]
其中系数 $a_i$ 与 $b_j$ 的加法和乘法都在 $F$ 中进行。注意，由乘法定义知
\[
  FG\text{ 是交换环当且仅当 }G\text{ 是 Abel 群。}
\]
群 $G$ 出现在 $FG$ 中（把 $g_i$ 认作 $1g_i$），域 $F$ 也出现在 $FG$ 中（把 $\beta$ 认作 $\beta g_1$，其中 $g_1$ 是 $G$ 的恒等元）。在这些认同下，
\[
  \beta\left(\sum_{i=1}^n a_i g_i\right)=\sum_{i=1}^n (\beta a_i)g_i,
  \qquad \beta\in F.
\]
于是
\[
  FG\text{ 是 }F\text{ 上的向量空间，以 }G\text{ 的元素为一组基。}
\]
特别地，$FG$ 是 $F$ 上维数等于 $|G|$ 的向量空间。$F$ 中的元素与 $FG$ 的所有元素交换，即 $F$ 包含在 $FG$ 的中心中。当我们希望强调后两个性质时，我们称 $FG$ 是一个 $F$-代数（一般地，$F$-代数是一个环 $R$，它在中心中包含 $F$，因而 $R$ 同时是一个环和一个 $F$-向量空间）。

注意，$FG$ 中的运算类似于 $F$-代数 $F[x]$ 中的运算（虽然 $F[x]$ 在 $F$ 上是无限维的）。在一些著作中 $FG$ 记作 $F[G]$，但后一种记号目前较少使用。

\noindent\emph{脚注.} 上面显示的形式和可以看成从 $G$ 到 $F$ 的函数的一种写法，该函数在群元素 $g_i$ 上取值 $a_i$。同样的“形式性”也用于构造自由模（见第 10.3 节定理 6）。

\paragraph{例}
\begin{enumerate}[label=(\arabic*)]
  \item 若 $G=\langle g\rangle$ 是 $n\in\ZZ^+$ 阶循环群，则 $FG$ 的元素形如
  \[
    \sum_{i=0}^{n-1} a_i g^i.
  \]
  映射 $F[x]\to F\langle g\rangle$ 把 $x^k$ 送到 $g^k$，$k\geq0$，并按 $F$-线性延拓，得到一个满环同态，其核等于由 $x^n-1$ 生成的理想。因此
  \[
    F\langle g\rangle\cong F[x]/(x^n-1).
  \]
  这是 $F$-代数同构，即一个 $F$-线性的环同构。

  \item 在前例记号下，令 $r=1+g+g^2+\cdots+g^{n-1}$，于是 $r$ 是 $F\langle g\rangle$ 的非零元素。注意 $rg=g+g^2+\cdots+g^{n-1}+1=r$，故 $r(1-g)=0$。因此环 $F\langle g\rangle$ 含有零因子（当 $n>1$ 时）。更一般地，若 $G$ 是任意阶大于 $1$ 的群，则对任意非恒等元 $g\in G$，$F\langle g\rangle$ 是 $FG$ 的子环，因而 $FG$ 也含有零因子。

  \item 令 $G=S_3$ 且 $F=\QQ$。元素
  \[
    r=5(1\ 2)-7(1\ 2\ 3),\qquad
    s=-4(1\ 2\ 3)+12(1\ 3\ 2)
  \]
  是 $\QQ S_3$ 的典型成员。它们的和与积为
  \[
    r+s=5(1\ 2)-11(1\ 2\ 3)+12(1\ 3\ 2),
  \]
  \[
    rs=-20(2\ 3)+28(1\ 3\ 2)+60(1\ 3)-84
  \]
  （回忆，置换的乘积或复合从右到左计算）。第 7.2 节中给出了群环 $\QQ D_8$ 中两个元素之和与乘积的一个显式例子。
\end{enumerate}

在给出具体表示的例子之前，我们讨论 $G$ 的表示与 $FG$-模之间的对应（这样之后便可同时给出二者的例子）。这个讨论与第 10.1 节中对 $F[x]$-模的处理非常相似。

首先设 $\varphi:G\to\GL(V)$ 是 $G$ 在 $F$ 上向量空间 $V$ 上的一个表示。如上，写 $G=\{g_1,\ldots,g_n\}$，于是对每个 $i\in\{1,\ldots,n\}$，$\varphi(g_i)$ 是 $V$ 到自身的线性变换。通过如下定义环元素在 $V$ 元素上的作用，把 $V$ 做成一个 $FG$-模：
\[
  \left(\sum_{i=1}^n a_i g_i\right)\cdot v
  =\sum_{i=1}^n a_i\varphi(g_i)(v),
  \qquad \sum_{i=1}^n a_i g_i\in FG,
  \ v\in V.
\]
我们验证模公理 2(b)（见第 10.1 节）的一个特殊情形，它精确地说明了 $\varphi$ 是群同态这一事实在哪里被用到：
\[
\begin{aligned}
 (g_i g_j)\cdot v
 &=\varphi(g_i g_j)(v) &&\text{（由作用的定义）}\
 &=(\varphi(g_i)\circ\varphi(g_j))(v) &&\text{（因为 }\varphi\text{ 是群同态）}\
 &=\varphi(g_i)(\varphi(g_j)(v)) &&\text{（由线性变换复合的定义）}\
 &=g_i\cdot(g_j\cdot v) &&\text{（由作用的定义）。}
\end{aligned}
\]
这个论证按线性性推广到 $FG$ 的任意元素，证明模公理 2(b) 一般成立。其余模公理容易检查。

注意，$F$ 是 $FG$ 的子环，域元素 $a$ 在向量上的作用与环元素 $a1$ 在向量上的作用相同；也就是说，$FG$-模作用延拓了 $F$ 在 $V$ 上的作用。

反过来，设给定一个 $FG$-模 $V$。我们如下得到相应的 $F$ 上向量空间以及 $G$ 的表示。由于 $V$ 是 $FG$-模，它也是 $F$-模，即 $F$ 上向量空间。此外，对每个 $g\in G$，我们得到从 $V$ 到 $V$ 的一个映射，记为 $\varphi(g)$，由
\[
  \varphi(g)(v)=g\cdot v,
  \qquad v\in V
\]
定义，其中 $g\cdot v$ 是环元素 $g$ 在元素 $v$ 上给定的作用。由于 $F$ 中元素与每个 $g\in G$ 交换，由模公理知对所有 $v,w\in V$ 与所有 $\alpha,\beta\in F$，
\[
\begin{aligned}
  \varphi(g)(\alpha v+\beta w)
  &=g\cdot(\alpha v+\beta w)\\
  &=g\cdot(\alpha v)+g\cdot(\beta w)\\
  &=\alpha(g\cdot v)+\beta(g\cdot w)\\
  &=\alpha\varphi(g)(v)+\beta\varphi(g)(w),
\end{aligned}
\]
即每个 $\varphi(g)$ 是线性变换。进一步，由模公理 2(b) 得
\[
  \varphi(g_i g_j)(v)=(\varphi(g_i)\circ\varphi(g_j))(v),
\]
（这本质上是上面计算的逆序）。这证明 $\varphi$ 是群同态（特别地，$\varphi(g^{-1})=\varphi(g)^{-1}$，故 $G$ 的每个元素都映到一个非奇异线性变换，即 $\varphi:G\to\GL(V)$）。

这个讨论说明，$FG$-模与二元组 $(V,\varphi)$ 之间存在双射：
\[
  \{V\text{ 是 }FG\text{-模}\}\longleftrightarrow
  \left\{\begin{array}{l}
  V\text{ 是 }F\text{ 上的向量空间，}\
  \varphi:G\to\GL(V)\text{ 是一个表示。}
  \end{array}\right.
\]
因此，给出 $F$ 上向量空间 $V$ 上的表示 $\varphi:G\to\GL(V)$，等价于给出一个 $FG$-模 $V$。在这个对应下，我们说模 $V$ \emph{实现}（affords）$G$ 的表示 $\varphi$。

回忆第 10.1 节：若向量空间 $M$ 通过线性变换 $T$ 做成 $F[x]$-模，则 $M$ 的 $F[x]$-子模正是 $M$ 的 $T$-稳定子空间。在当前情形中，若 $V$ 是实现表示 $\varphi$ 的 $FG$-模，则 $V$ 的子空间 $U$ 称为 $G$-不变或 $G$-稳定的，如果对所有 $g\in G$ 与所有 $u\in U$ 都有 $g\cdot u\in U$（即 $\varphi(g)(u)\in U$）。容易推出
\[
  V\text{ 的 }FG\text{-子模正是 }V\text{ 的 }G\text{-稳定子空间。}
\]

\paragraph{例}
\begin{enumerate}[label=(\arabic*)]
  \item 令 $V$ 是 $F$ 上的一维向量空间，并令 $gv=v$ 对所有 $g\in G$、$v\in V$ 成立，从而把 $V$ 做成 $FG$-模。这个模实现表示 $\varphi:G\to\GL(V)$，其中 $\varphi(g)=I$ 是恒等线性变换，对所有 $g\in G$。相应的矩阵表示（相对于 $V$ 的任意基）是把每个群元素送到 $1\times1$ 单位矩阵的同态。以后我们称它为 $G$ 的\emph{平凡表示}。平凡表示次数为 $1$；若 $|G|>1$，它不是忠实的。

  \item 令 $V=FG$，把这个环看成自身上的左模。则 $V$ 实现一个次数等于 $|G|$ 的表示。若取 $G$ 的元素作为 $V$ 的基，则每个 $g\in G$ 在\emph{左正则置换表示}下置换这些基元素：
  \[
    g\cdot g_i=gg_i.
  \]
  相对于 $V$ 的这个基，群元素 $g$ 的矩阵在第 $i$ 行第 $j$ 列为 $1$，当且仅当 $gg_j=g_i$；其余位置为 $0$。这个（线性或矩阵）表示称为 $G$ 的\emph{正则表示}。注意 $G$ 的每个非恒等元都诱导基上的非恒等置换，所以正则表示总是忠实的。

  \item 设 $n\in\ZZ^+$，$G=S_n$，令 $V$ 是 $F$ 上 $n$ 维向量空间，基为 $e_1,e_2,\ldots,e_n$。令 $S_n$ 作用在 $V$ 上：对每个 $\sigma\in S_n$ 定义
  \[
    \sigma\cdot e_i=e_{\sigma(i)},\qquad 1\leq i\leq n,
  \]
  即 $\sigma$ 置换基元素的下标。这给出一个从 $S_n$ 到 $\GL(V)$ 的（单射）同态（即 $S_n$ 的一个次数 $n$ 的忠实表示），因而把 $V$ 做成 $FS_n$-模。与前例一样，相对于基 $e_1,
  \ldots,e_n$，$\sigma$ 的矩阵在第 $i$ 行第 $j$ 列为 $1$，当且仅当 $\sigma\cdot e_j=e_i$（其余项为 $0$）。因此 $\sigma$ 的矩阵在第 $i$ 行第 $j$ 列为 $1$ 当且仅当 $\sigma(j)=i$。

  例如考虑 $FS_3$ 在 $F$ 上三维向量空间（基为 $e_1,e_2,e_3$）上的作用。令 $\alpha=(1\ 2)$，$\tau=(1\ 2\ 3)$，并令 $r=2\alpha-3\tau\in FS_3$。则
  \[
  \begin{aligned}
    r\cdot(\alpha_0e_1+\beta e_2+\gamma e_3)
    &=2(\alpha_0e_{\alpha(1)}+\beta e_{\alpha(2)}+\gamma e_{\alpha(3)})
      -3(\alpha_0e_{\tau(1)}+\beta e_{\tau(2)}+\gamma e_{\tau(3)})\\
    &=2(\alpha_0e_2+\beta e_1+\gamma e_3)-3(\alpha_0e_2+\beta e_3+\gamma e_1)\\
    &=(2\beta-3\gamma)e_1-\alpha_0e_2+(2\gamma-3\beta)e_3.
  \end{aligned}
  \]

  \item 若 $\psi:H\to\GL(V)$ 是 $H$ 的任意表示，$\varphi:G\to H$ 是任意群同态，则复合 $\psi\circ\varphi$ 是 $G$ 的一个表示。例如，令 $V$ 是前例中维数为 $n$ 的 $FS_n$-模。若 $\pi:G\to S_n$ 是 $G$ 的任意置换表示，则 $\pi$ 与上面的表示复合给出 $G$ 的一个线性表示。换言之，$V$ 在作用
  \[
    g\cdot e_i=e_{\pi(g)(i)},\qquad g\in G
  \]
  下成为 $FG$-模。注意，正则表示 (2) 正是这里 $n=|G|$ 且 $\pi$ 为 $G$ 的左正则置换表示的特例。

  \item 任意从 $G$ 到乘法群 $F^\times=\GL_1(F)$ 的同态都是次数 $1$ 的（矩阵）表示。例如，设 $G=\langle g\rangle\cong\ZZ/n\ZZ$ 是 $n$ 阶循环群，$\zeta$ 是 $F$ 中一个固定的 $n$ 次单位根。令 $g^i\mapsto\zeta^i$，$i\in\ZZ$。当且仅当 $\zeta$ 是本原 $n$ 次单位根时，这个 $\langle g\rangle$ 的表示是忠实的。

  \item 在许多情形下，给出群 $G$ 的显式矩阵表示比展示一个 $FG$-模更容易。例如，回忆二面体群 $D_{2n}$ 有表示
  \[
    D_{2n}=\langle r,s\mid r^n=s^2=1,
    \ rs=sr^{-1}\rangle.
  \]
  若 $R$ 和 $S$ 是满足关系 $R^n=S^2=I$ 与 $RS=SR^{-1}$ 的矩阵，则映射 $r\mapsto R$、$s\mapsto S$ 唯一延拓为从 $D_{2n}$ 到由 $R$ 和 $S$ 生成的矩阵群的同态，因而给出 $D_{2n}$ 的一个表示。矩阵 $R,S\in M_2(\RR)$ 的一个显式例子可如下得到：在 $x,y$ 平面中以原点为中心画一个正 $n$ 边形，使直线 $y=x$ 为其对称轴之一。设 $R$ 是平面旋转 $2\pi/n$ 弧度的矩阵，$S$ 是关于直线 $y=x$ 的反射矩阵，则二者都把该正 $n$ 边形映到自身。因此这些矩阵作为该 $n$ 边形的对称而作用，并满足上述关系。这些矩阵容易算出（参见第 1.6 节练习 25），于是映射
  \[
    r\mapsto R=\begin{pmatrix}
    \cos(2\pi/n)&-\sin(2\pi/n)\\
    \sin(2\pi/n)&\cos(2\pi/n)
    \end{pmatrix},\qquad
    s\mapsto S=\begin{pmatrix}0&1\\1&0\end{pmatrix}
  \]
  唯一延拓为 $D_{2n}$ 到 $\GL_2(\RR)$ 的一个（次数 $2$）表示。由于矩阵 $R$ 和 $S$ 的阶分别为 $n$ 和 $2$，它们生成 $\GL_2(\RR)$ 中一个 $2n$ 阶子群，故此表示是忠实的。

  \item 使用四元数群的通常生成元和关系
  \[
    Q_8=\langle i,j\mid i^4=j^4=1,
    \ i^2=j^2,
    \ j^{-1}ij=i^{-1}\rangle,
  \]
  可类似地得到（参见第 1.6 节练习 26）一个从 $Q_8$ 到 $\GL_2(\CC)$ 的表示 $\varphi$，定义为
  \[
    \varphi(i)=\begin{pmatrix}\sqrt{-1}&0\\0&-\sqrt{-1}\end{pmatrix},
    \qquad
    \varphi(j)=\begin{pmatrix}0&-1\\1&0\end{pmatrix}.
  \]
  这个 $Q_8$ 的表示是忠实的。

  \item 四元数群 $Q_8$ 的一个 $4$ 维表示可由实 Hamilton 四元数 $\HH$ 得到（参见第 7.1 节）。群 $Q_8$ 是 $\HH$ 的单位群的一个子群，$Q_8$ 的每个元素都通过左乘作用在 $4$ 维实向量空间 $\HH$ 上。由于实数在 $\HH$ 的中心中（即 $\HH$ 是一个 $\RR$-代数），左乘是 $\RR$-线性的。因此这个线性作用给出从 $Q_8$ 到 $\GL_4(\RR)$ 的同态。相对于 $\HH$ 的基 $1,i,j,k$，容易写出 $Q_8$ 各元素的显式矩阵。例如，左乘 $i$ 作用为 $1\mapsto i$、$i\mapsto-1$、$j\mapsto k$、$k\mapsto-j$；左乘 $j$ 作用为 $1\mapsto j$、$i\mapsto-k$、$j\mapsto-1$、$k\mapsto i$，所以
  \[
  i\mapsto
  \begin{pmatrix}
  0&-1&0&0\\
  1&0&0&0\\
  0&0&0&-1\\
  0&0&1&0
  \end{pmatrix},
  \qquad
  j\mapsto
  \begin{pmatrix}
  0&0&-1&0\\
  0&0&0&1\\
  1&0&0&0\\
  0&-1&0&0
  \end{pmatrix}.
  \]
  这个 $Q_8$ 的表示也是忠实的。

  \item 设 $H$ 是群 $G$ 的正规子群，并且 $H$ 是某个素数 $p$ 的初等 Abel $p$-群。则 $V=H$ 是 $\FF_p$ 上的向量空间，其中标量 $a$ 作用在向量 $v$ 上为 $av=v^a$（见第 10.1 节）。$G$ 的每个元素通过共轭在 $V$ 上作用，这个作用是 $\FF_p$-线性的，因为 $gv^ag^{-1}=(gvg^{-1})^a$；$G$ 在 $V$ 上的这个作用把 $V$ 做成 $\FF_pG$-模（初等 Abel $p$-群的自同构在第 4.4 节和第 10.1 节讨论过）。这个表示的核是 $G$ 中与 $H$ 的每个元素都交换的元素集合 $C_G(H)$（它总是包含 Abel 群 $H$ 本身）。因此，一个群在自身子集上的作用常常给出有限域上的线性表示。有限域上的群表示称为\emph{模表示}，它们对于研究群的内部结构是基本的。

  \item 对 $FG$-子模的一个例子，令 $G=S_n$，令 $V$ 为例 3 中描述的 $FS_n$-模。令 $N$ 是 $V$ 中所有坐标都相等的向量所成的子空间：
  \[
    N=\{a_1e_1+a_2e_2+\cdots+a_ne_n\mid a_1=a_2=\cdots=a_n\}.
  \]
  这是一个 $1$ 维 $S_n$-稳定子空间。每个 $\sigma\in S_n$ 固定 $N$ 中每个向量，因此子模 $N$ 实现 $S_n$ 的平凡表示。作为练习可以证明：若 $n\geq3$，则 $N$ 是 $V$ 中唯一的 $1$ 维 $S_n$-稳定子空间，即唯一的 $1$ 维 $FS_n$-子模（$N$ 称为 $FS_n$ 的迹子模）。

  $V$ 的另一个 $FS_n$-子模是所有坐标和为零的向量所成的子空间 $I$：
  \[
    I=\{a_1e_1+a_2e_2+\cdots+a_ne_n\mid a_1+a_2+\cdots+a_n=0\}.
  \]
  同样，$I$ 是 $S_n$-稳定子空间（因为每个 $\sigma\in S_n$ 只是置换每个向量的坐标，因而保持系数和不变）。由于 $I$ 是从 $V$ 到 $F$ 的线性变换（把一个向量送到其系数之和，称为增广映射——参见第 7.3 节）的核，$I$ 的维数为 $n-1$。

  \item 若 $V=FG$ 是上面例 2 中的 $G$ 的正则表示，则 $V$ 有如前例一样的维数为 $1$ 与 $|G|-1$ 的 $FG$-子模：
  \[
    N=\{a_1g_1+a_2g_2+\cdots+a_ng_n\mid a_1=a_2=\cdots=a_n\},
  \]
  \[
    I=\{a_1g_1+a_2g_2+\cdots+a_ng_n\mid a_1+a_2+\cdots+a_n=0\}.
  \]
  实际上 $N$ 和 $I$ 是 $FG$ 的双边理想（不仅是左理想——注意 $N$ 在 $FG$ 的中心中）。理想 $I$ 称为 $FG$ 的\emph{增广理想}，$N$ 称为 $FG$ 的\emph{迹理想}。
\end{enumerate}

回忆，在研究向量空间 $V$ 到自身的线性变换 $T$ 时，我们把 $V$ 做成 $F[x]$-模（其中 $x$ 在 $V$ 上作为 $T$ 作用）；我们的目标是把 $V$ 分解为循环子模的直和。这样便能找到 $V$ 的一个基，使得 $T$ 关于该基的矩阵处于某种标准形。改变 $V$ 的基并不改变模 $V$，但会通过相似改变 $T$ 的矩阵表示（即改变 $\GL(V)$ 与 $\GL_n(F)$ 之间的同构）。我们引入类似术语来描述两个 $FG$-模在换基意义下何时相同。

\noindent\textbf{定义.} $G$ 的两个表示称为\emph{等价的}（或\emph{相似的}），如果实现它们的 $FG$-模是同构模。不等价的表示称为\emph{不等价的}。

设 $\varphi:G\to\GL(V)$ 和 $\psi:G\to\GL(W)$ 是等价表示（这里 $V$ 与 $W$ 必须是同一域 $F$ 上的向量空间）。设 $T:V\to W$ 是二者之间的 $FG$-模同构。由于 $T$ 特别是 $F$-模同构，$T$ 是向量空间同构，故 $V$ 与 $W$ 维数相同。此外，对所有 $g\in G$、$v\in V$，由于 $T$ 是 $FG$-模同构，有 $T(g\cdot v)=g\cdot T(v)$。由环元素作用的定义，这意味着
\[
  T(\varphi(g)v)=\psi(g)(T(v)),
\]
即
\[
  T\circ\varphi(g)=\psi(g)\circ T,
  \qquad g\in G.
\]
特别地，若我们把 $V$ 和 $W$ 认作同一个向量空间，则 $G$ 在向量空间 $V$ 上的两个表示 $\varphi$ 和 $\psi$ 等价，当且仅当存在某个 $T\in\GL(V)$ 使得
\[
  T\circ\varphi(g)\circ T^{-1}=\psi(g),
  \qquad g\in G.
\]
这个 $T$ 是对所有 $\varphi(g)$（$g\in G$）同时进行的换基。

用矩阵术语说，两个表示 $\varphi$ 和 $\psi$ 等价，当且仅当存在一个固定可逆矩阵 $P$，使得
\[
  P\varphi(g)P^{-1}=\psi(g),
  \qquad g\in G.
\]
上面的线性变换 $T$ 或矩阵 $P$ 称为\emph{交织}表示 $\varphi$ 与 $\psi$（它给出了把 $\varphi$ 变为 $\psi$ 的“规则”）。

为了研究 $FG$-模分解为子模（直和）的情形，我们需要一些术语。由于下一节将更一般地讨论直和分解，我们对任意环陈述这些定义。

\noindent\textbf{定义.} 设 $R$ 是环，$M$ 是非零 $R$-模。
\begin{enumerate}[label=(\arabic*)]
  \item 若 $M$ 的子模只有 $0$ 和 $M$，则称 $M$ 是\emph{不可约的}（或\emph{单的}）；否则称 $M$ 是\emph{可约的}。
  \item 若 $M$ 不能写成 $M_1\oplus M_2$，其中 $M_1,M_2$ 是非零子模，则称 $M$ 是\emph{不可分解的}；否则称 $M$ 是\emph{可分解的}。
  \item 若 $M$ 是不可约子模的直和，则称 $M$ 是\emph{完全可约的}。
  \item 一个表示称为不可约、可约、不可分解、可分解或完全可约，按照实现它的 $FG$-模是否具有相应性质而定。
  \item 若 $M$ 是完全可约 $R$-模，则 $M$ 的任一直和因子称为 $M$ 的一个\emph{组成分量}（即若存在 $M$ 的子模 $N'$ 使得 $M=N\oplus N'$，则 $N$ 是 $M$ 的一个组成分量）。
\end{enumerate}

一个不可约模按定义既是不可分解的，也是完全可约的。稍后我们将给出不可分解但不可约的模的例子。

若 $R=FG$，则不可约 $FG$-模 $V$ 是一个非零 $F$-向量空间，它没有非平凡的真 $G$-不变子空间。例如，若 $\dim_F V=1$，则 $V$ 必然不可约（它的子空间只有 $0$ 和 $V$）。

设 $V$ 是有限维 $FG$-模且 $V$ 可约。令 $U$ 是一个 $G$-不变子空间。取 $U$ 的一个基并扩充成 $V$ 的一个基。于是对每个 $g\in G$，$g$ 在 $V$ 上作用的矩阵 $\varphi(g)$ 相对于这个基形如
\[
  \varphi(g)=\begin{pmatrix}
  \varphi_1(g)&\psi(g)\\
  0&\varphi_2(g)
  \end{pmatrix},
\]
其中 $\varphi_1=\varphi|_U$（相对于 $U$ 的所选基），$\varphi_2$ 是 $G$ 在 $V/U$ 上的表示（$\psi$ 不一定是同态——$\psi(g)$ 不必是方阵）。因此，可约表示正是那些有相应矩阵表示使所有矩阵同为分块上三角形的表示。

进一步假设 $FG$-模 $V$ 可分解，$V=U\oplus U'$。取 $U$ 的一个基和 $U'$ 的一个基的并作为 $V$ 的基。相对于这个基，每个 $g\in G$ 的矩阵形如
\[
  \varphi(g)=\begin{pmatrix}
  \varphi_1(g)&0\\
  0&\varphi_2(g)
  \end{pmatrix},
\]
（即对所有 $g\in G$ 都有 $\psi(g)=0$）。因此，可分解表示正是那些有相应矩阵表示使所有矩阵同为分块对角形的表示。

\paragraph{例}
\begin{enumerate}[label=(\arabic*)]
  \item 如上所述，所有次数 $1$ 的表示都是不可约的、不可分解的且完全可约的。特别地，平凡表示和上面例 5 中的表示都具有这些性质。
  \item 若 $|G|>1$，则 $G$ 的正则表示是可约的（增广理想与迹理想是非零真子模）。稍后我们将判定该表示何时完全可约，以及它如何分解为直和。
  \item 对 $n>1$，上面例 10 中描述的 $FS_n$-模是可约的，因为 $N$ 和 $I$ 是非零真子模。模 $N$ 是不可约的（因为一维），而若域 $F$ 的特征不整除 $n$，则 $I$ 也是不可约的。
  \item 例 6 中二面体群 $D_{2n}=G$ 的次数 $2$ 表示在 $n\geq3$ 时是不可约的。不存在 $G$-不变的一维子空间，因为旋转 $2\pi/n$ 弧度不会把 $\RR^2$ 中任何直线映到自身。类似地，例 7 中 $Q_8$ 的次数 $2$ 复表示是不可约的，因为给定矩阵 $\varphi(i)$ 恰有两个一维特征空间（对应于不同特征值 $\pm\sqrt{-1}$），而这些特征空间并不在矩阵 $\varphi(j)$ 下不变。例 8 中次数 $4$ 表示 $\varphi:Q_8\to\GL_4(\RR)$ 也可证明是不可约的（见练习）。不过我们将看到，若把 $\varphi$ 看成复表示 $\varphi:Q_8\to\GL_4(\CC)$（只是把矩阵的实项看作复项），则存在复矩阵 $P$，使得对所有 $g\in Q_8$，$P^{-1}\varphi(g)P$ 是 $2\times2$ 分块矩阵的直和。因此，一个在域 $F$ 上不可约的表示，在扩域后可能变成可约的。
  \item 令 $G=\langle g\rangle$ 为 $n$ 阶循环群，假设 $F$ 含有所有 $n$ 次单位根。由群代数例 1，$F\langle g\rangle\cong F[x]/(x^n-1)$。因此 $FG$-模正是被 $x^n-1$ 零化的 $F[x]$-模。后者的（有限维）模按等价由 Jordan 标准形定理描述。

  若 $g$ 在 $F\langle g\rangle$-模 $V$ 上作用的极小多项式在 $F$ 中有不同根，则存在 $V$ 的一个基，使得 $g$（因而其所有幂）由对角矩阵表示（参见第 12.3 节推论 25）。在这种情况下，$V$ 是完全可约的 $F\langle g\rangle$-模（是 $1$ 维 $\langle g\rangle$-不变子空间的直和）。一般地，$g$ 在 $V$ 上作用的极小多项式整除 $x^n-1$，所以若 $x^n-1$ 在 $F$ 中有不同根，则 $V$ 是完全可约的 $F\langle g\rangle$-模。多项式 $x^n-1$ 在 $F$ 中有不同根当且仅当 $F$ 的特征不整除 $n$。这给出每个 $F\langle g\rangle$-模都完全可约的一个充分条件。

  若 $g$ 在 $V$ 上作用的极小多项式没有不同根（于是 $F$ 的特征整除 $n$），则 $g$ 的 Jordan 标准形必须有大小大于 $1$ 的初等 Jordan 块。由于每个线性变换都有唯一的 Jordan 标准形，$g$ 不能用对角矩阵表示，即 $V$ 不是完全可约的。由第 12.3 节中关于循环模的结果可知，在任一大小大于 $1$ 的 Jordan 块中，$g$ 的（一维）特征空间没有 $\langle g\rangle$-不变补空间，即 $V$ 可约但不完全可约。

  具体地，设 $p$ 是素数，$F=\FF_p$，令 $g$ 的阶为 $p$。设 $V$ 是 $\FF_p$ 上二维空间，基为 $v,w$，定义 $g$ 在 $V$ 上作用为
  \[
    g\cdot v=v,
    \qquad g\cdot w=v+w.
  \]
  这个 $V$ 的自同态在 $\GL(V)$ 中确有阶 $p$，且 $g$ 相对于该基的矩阵是初等 Jordan 块
  \[
    \varphi(g)=\begin{pmatrix}1&1\\0&1\end{pmatrix}.
  \]
  现在 $V$ 是可约的（$\operatorname{span}\{v\}$ 是 $\langle g\rangle$-不变子空间），但 $V$ 是不可分解的（上述 $2\times2$ 初等 Jordan 矩阵不相似于对角矩阵）。
\end{enumerate}

有限群表示论中的第一个基本结果说明例 5 如何推广到非循环群。

\medskip
\noindent\textbf{定理 1.} （Maschke 定理）设 $G$ 是有限群，$F$ 是特征不整除 $|G|$ 的域。若 $V$ 是任意 $FG$-模，$U$ 是 $V$ 的任意子模，则 $V$ 有一个子模 $W$ 使得 $V=U\oplus W$（即每个子模都是直和因子）。

\noindent\textbf{注记.} 当 $F$ 的特征为 $0$ 时，Maschke 定理的假设适用于任意有限群。

\begin{proof}
Maschke 定理证明的想法是构造一个 $FG$-模同态
\[
  \pi:V\to U
\]
它是到 $U$ 上的投影，即满足以下两个性质：
\[
  \text{(i) }\pi(u)=u\text{ 对所有 }u\in U,
  \qquad
  \text{(ii) }\pi(\pi(v))=\pi(v)\text{ 对所有 }v\in V
\]
（事实上 (ii) 由 (i) 以及 $\pi(V)\subseteq U$ 推出）。

先假设我们能构造出这样的 $FG$-模同态，并令 $W=\ker\pi$。由于 $\pi$ 是模同态，$W$ 是子模。我们如下看到 $W$ 是 $U$ 的直和补。若 $v\in U\cap W$，则由 (i) 有 $v=\pi(v)$，而由 $W$ 的定义有 $\pi(v)=0$。这说明 $U\cap W=0$。为了证明 $V=U+W$，取任意 $v\in V$，写
\[
  v=\pi(v)+(v-
  \pi(v)).
\]
按定义，$\pi(v)\in U$。由 $\pi$ 的性质 (ii)，
\[
  \pi(v-\pi(v))=\pi(v)-\pi(\pi(v))=\pi(v)-\pi(v)=0,
\]
即 $v-\pi(v)\in W$。这说明 $V=U+W$，从而 $V=U\oplus W$。因此，要证明 Maschke 定理，只需找到这样的 $FG$-模投影 $\pi$。

由于 $U$ 是子空间，它在 $V$ 中有向量空间直和补 $W_0$（取 $U$ 的一个基 $B_1$，扩充为 $V$ 的一个基 $B$，令 $W_0$ 为 $B-B_1$ 的张成空间）。于是作为向量空间，$V=U\oplus W_0$，但 $W_0$ 未必 $G$-稳定（即未必是 $FG$-子模）。令 $\pi_0:V\to U$ 是与此直和分解对应的从 $V$ 到 $U$ 的向量空间投影，即
\[
  \pi_0(u+w)=u,
  \qquad u\in U,
  \ w\in W_0.
\]
证明的关键思想是把 $\pi_0$ 在 $G$ 上“平均”以形成一个 $FG$-模投影 $\pi$。对每个 $g\in G$ 定义
\[
  g\pi_0g^{-1}:V\to U,
  \qquad
  (g\pi_0g^{-1})(v)=g\cdot\pi_0(g^{-1}\cdot v),
  \quad v\in V,
\]
其中点号表示环 $FG$ 元素的作用。由于 $\pi_0$ 把 $V$ 映到 $U$ 中，而 $U$ 在 $g$ 作用下稳定，故 $g\pi_0g^{-1}$ 把 $V$ 映到 $U$ 中。$g$ 与 $g^{-1}$ 都作为 $F$-线性变换作用，因此 $g\pi_0g^{-1}$ 是线性变换。此外，若 $u$ 在 $G$-稳定空间 $U$ 中，则 $g^{-1}u$ 也在 $U$ 中，由 $\pi_0$ 的定义有 $\pi_0(g^{-1}u)=g^{-1}u$。由此得到对所有 $g\in G$，
\[
  g\pi_0g^{-1}(u)=u,
  \qquad u\in U,
\]
（即 $g\pi_0g^{-1}$ 也是 $V$ 到 $U$ 的向量空间投影）。

令 $n=|G|$，并把 $n$ 看作 $F$ 中的元素（$n=1+\cdots+1$，共 $n$ 次）。由假设，$n$ 在 $F$ 中不为零，因而在 $F$ 中有逆元。定义
\[
  \pi=\frac1n\sum_{g\in G}g\pi_0g^{-1}.
\]
由于 $\pi$ 是从 $V$ 到 $U$ 的线性变换之和的标量倍，它也是从 $V$ 到 $U$ 的线性变换。进一步，定义 $\pi$ 的和式中的每一项在子空间 $U$ 上都限制为恒等映射，因此 $\pi|_U$ 是 $1/n$ 乘以 $n$ 个恒等映射之和。这些观察证明：
\[
  \pi:V\to U\text{ 是线性变换},
  \qquad \pi(u)=u\ (u\in U),
  \qquad \pi^2(v)=\pi(v)\ (v\in V).
\]
还需证明 $\pi$ 是 $FG$-模同态（即 $FG$-线性的）。只需证明对所有 $h\in G$，$v\in V$，有 $\pi(hv)=h\pi(v)$。事实上
\[
\begin{aligned}
  \pi(hv)
  &=\frac1n\sum_{g\in G}g\pi_0(g^{-1}hv)\\
  &=\frac1n\sum_{k\in G}hk\pi_0(k^{-1}v)\\
  &=h\left(\frac1n\sum_{k\in G}k\pi_0(k^{-1}v)\right)
  =h\pi(v),
\end{aligned}
\]
因为当 $g$ 遍历 $G$ 的所有元素时，$k=h^{-1}g$ 也遍历 $G$ 的所有元素，而模元素 $h$ 可由模的分配律提出求和号。这建立了 $FG$-模投影 $\pi$ 的存在性，从而完成证明。
\end{proof}

Maschke 定理的应用将针对有限生成 $FG$-模。不过，与 $F[x]$-模的情形不同，有限生成 $FG$-模自动是有限维向量空间（差别在于 $FG$ 本身是有限维的，而 $F[x]$ 不是）。设 $V$ 是 $FG$-模。若 $V$ 是 $F$ 上有限维向量空间，则 $V$ 当然作为 $FG$-模有限生成（任意 $F$-基给出 $FG$ 上的生成元集）。反过来，若 $V$ 作为 $FG$-模由 $v_1,
\ldots,v_k$ 有限生成，则容易看出 $V$ 作为向量空间由有限集合 $\{g\cdot v_i\mid g\in G,
1\leq i\leq k\}$ 张成。因此
\[
  FG\text{-模有限生成当且仅当它是有限维的。}
\]

\medskip
\noindent\textbf{推论 2.} 若 $G$ 是有限群，$F$ 是特征不整除 $|G|$ 的域，则每个有限生成 $FG$-模都是完全可约的（等价地，$G$ 的每个有限次数 $F$-表示都是完全可约的）。

\begin{proof}
令 $V$ 是有限生成 $FG$-模。如上所述，$V$ 在 $F$ 上有限维，所以可对其维数归纳。若 $V$ 不可约，则它完全可约，结论成立。于是设 $V$ 有一个非零真 $FG$-子模 $U$。由 Maschke 定理，$U$ 有一个 $FG$-子模补 $W$，即 $V=U\oplus W$。由归纳假设，$U$ 和 $W$ 各自都是不可约子模的直和，因此 $V$ 也是。归纳完成。
\end{proof}

\medskip
\noindent\textbf{推论 3.} 设 $G$ 是有限群，$F$ 是特征不整除 $|G|$ 的域，$\varphi:G\to\GL(V)$ 是 $G$ 的有限次数表示。则存在 $V$ 的一个基，使得对每个 $g\in G$，$\varphi(g)$ 关于该基的矩阵为分块对角形式
\[
  \begin{pmatrix}
  \varphi_1(g)&&&0\\
  &\varphi_2(g)&&\\
  &&\ddots&\\
  0&&&\varphi_m(g)
  \end{pmatrix},
\]
其中 $\varphi_i$ 是 $G$ 的不可约矩阵表示，$1\leq i\leq m$。

\begin{proof}
由推论 2，可写 $V=U_1\oplus U_2\oplus\cdots\oplus U_m$，其中 $U_i$ 是 $V$ 的不可约 $FG$-子模。令 $\mathcal B_i$ 是 $U_i$ 的一个基，令 $\mathcal B$ 是这些 $\mathcal B_i$ 的并。对每个 $g\in G$，$\varphi(g)$ 关于基 $\mathcal B$ 的矩阵具有推论中给出的形式，其中 $\varphi_i(g)$ 是 $\varphi(g)|_{U_i}$ 关于基 $\mathcal B_i$ 的矩阵。
\end{proof}

Maschke 定理的逆命题也成立。即，若 $F$ 的特征整除 $|G|$，则 $G$ 存在（有限生成）$FG$-模不是完全可约的。具体地，正则表示（即模 $FG$ 本身）不是完全可约的。

在第 18.2 节中，我们将讨论 $FG$-模分解为不可约子模直和时，其组成分量唯一性的问题。

\subsection*{练习}

设 $F$ 是域，$G$ 是有限群，$n\in\ZZ^+$。
\begin{enumerate}[label=\arabic*.]
  \item 证明：若 $\varphi:G\to\GL(V)$ 是任意表示，则 $\varphi$ 给出 $G/\ker\varphi$ 的一个忠实表示。
  \item 设 $\varphi:G\to\GL_n(F)$ 是矩阵表示。证明映射 $g\mapsto\det(\varphi(g))$ 是次数 $1$ 表示。
  \item 证明 $G$ 的次数 $1$ 表示与 Abel 群 $G/G'$（其中 $G'$ 是 $G$ 的换位子群）的次数 $1$ 表示一一对应。
  \item 设 $V$ 是一个（可能无限维）$FG$-模（$G$ 是有限群）。证明对每个 $v\in V$，存在一个包含 $v$ 的 $FG$-子模，其维数 $\leq |G|$。
  \item 证明：若 $|G|>1$，则每个不可约 $FG$-模的维数 $<|G|$。
  \item 对第二组例子中描述的下列每个表示，写出每个 $g\in G$ 的矩阵 $\varphi(g)$：
  \begin{enumerate}[label=(\alph*)]
    \item 例 3 中描述的 $S_3$ 表示（令该例中 $n=3$）；
    \item 例 6 中描述的 $D_8$ 表示（即令该例中 $n=4$，并写出所有群元素的所有正弦、余弦值）；
    \item 例 7 中描述的 $Q_8$ 表示；
    \item 例 8 中描述的 $Q_8$ 表示。
  \end{enumerate}
  \item 令 $V$ 是第二组例子中例 3 所描述的 $S_4$ 的 $4$ 维置换模。令 $\pi:D_8\to S_4$ 是由 $D_8$ 通过左乘作用在其子群 $\langle s\rangle$ 的左陪集集合上得到的置换表示。按例 4 所述，通过 $\pi$ 把 $V$ 做成 $FD_8$-模，并相对于基 $e_1,
  \ldots,e_4$ 写出此表示中 $r$ 和 $s$ 的 $4\times4$ 矩阵。
  \item 令 $V$ 是第二组例子中例 3 和例 10 所描述的 $FS_n$-模。
  \begin{enumerate}[label=(\alph*)]
    \item 证明：若 $v\in V$ 且对所有 $\sigma\in S_n$ 都有 $\sigma\cdot v=v$，则 $v$ 是 $e_1+e_2+\cdots+e_n$ 的 $F$-倍数。
    \item 证明：若 $n\geq3$，则 $V$ 有唯一的 $1$ 维子模，即由 $e_1+e_2+\cdots+e_n$ 的所有 $F$-倍数组成的子模 $N$。
  \end{enumerate}
  \item 证明第二组例子中例 8 所述 $Q_8$ 在 $\HH$ 上的 $4$ 维表示是不可约的。[证明任一 $Q_8$-稳定子空间都是左理想。]
  \item 证明 $\GL_2(\RR)$ 没有与 $Q_8$ 同构的子群。[可直接用 $Q_8$ 的生成元和关系计算。可把一个生成元化为有理标准形以简化计算。]
  \item 令 $\varphi:S_n\to\GL_n(F)$ 是由第二组例子中例 3 的置换模给出的矩阵表示，矩阵相对于基 $e_1,
  \ldots,e_n$ 计算。证明对所有 $\sigma\in S_n$，$\det\varphi(\sigma)=\epsilon(\sigma)$，其中 $\epsilon(\sigma)$ 是置换 $\sigma$ 的符号。[在对换上检查。]
  \item 假设 $F$ 的特征不是 $2$。令 $H$ 是 $M_n(F)$ 中所有这样的 $T$ 所成的集合：$T$ 的每行每列恰有一个非零项，其他项为零，且非零项为 $\pm1$。证明 $H$ 是 $\GL_n(F)$ 的子群，且 $H\cong E_{2^n}\rtimes S_n$（半直积），其中 $E_{2^n}$ 是 $2^n$ 阶初等 Abel 群。
\end{enumerate}

接下来的若干练习探讨一个重要结果，称为 Schur 引理，以及它的一些推论。
\begin{enumerate}[label=\arabic*.,resume]
  \item 设 $R$ 是环，$M$ 与 $N$ 是单（即不可约）$R$-模。
  \begin{enumerate}[label=(\alph*)]
    \item 证明从 $M$ 到 $N$ 的每个非零 $R$-模同态都是同构。[考虑其核与像。]
    \item 证明 Schur 引理：若 $M$ 是单 $R$-模，则 $\Hom_R(M,M)$ 是一个除环（回忆 $\Hom_R(M,M)$ 是从 $M$ 到 $M$ 的所有 $R$-模同态所成的环，乘法为复合）。
  \end{enumerate}
  \item 设 $\varphi:G\to\GL(V)$ 是 $G$ 的表示。$\varphi$ 的中心化子定义为所有从 $V$ 到自身的线性变换 $A$ 的集合，满足对所有 $g\in G$ 都有 $A\varphi(g)=\varphi(g)A$（即与所有 $\varphi(g)$ 交换的 $V$ 的线性变换）。
  \begin{enumerate}[label=(\alph*)]
    \item 证明：$V$ 到 $V$ 的线性变换 $A$ 在 $\varphi$ 的中心化子中，当且仅当它是从 $V$ 到自身的 $FG$-模同态（所以 $\varphi$ 的中心化子就是环 $\Hom_{FG}(V,V)$）。
    \item 证明：若 $z$ 在 $G$ 的中心中，则 $\varphi(z)$ 在 $\varphi$ 的中心化子中。
    \item 假设 $\varphi$ 是不可约表示（所以 $V$ 是单 $FG$-模）。证明：若 $H$ 是 $\GL(V)$ 的任意有限 Abel 子群，且对所有 $A\in H$ 都有 $A\varphi(g)=\varphi(g)A$，则 $H$ 是循环群（换言之，环 $\Hom_{FG}(V,V)$ 的单位群中任意有限 Abel 子群都是循环的）。[由前一练习，$\Hom_{FG}(V,V)$ 是除环，所以问题化为证明任意除环的非零元素乘法群中任意有限 Abel 子群都是循环群。证明任意除环中由 Abel 子群生成的除子环是域，并用第 9.5 节命题 18。]
    \item 证明：若 $\varphi$ 是忠实不可约表示，则 $G$ 的中心是循环群。
    \item 由 (d) 推出：若 $G$ 是 Abel 群，$\varphi$ 是任意不可约表示，则 $G/\ker\varphi$ 是循环群。
  \end{enumerate}
  \item 展示有限循环群的所有 $1$ 维复表示；务必判定哪些是不等价的。
  \item 展示有限 Abel 群的所有 $1$ 维复表示。推出有限 Abel 群的不等价次数 $1$ 复表示的个数等于该群的阶。[先把 Abel 群分解为循环群的直积，再使用前一练习。]
  \item 证明复数域上 Abel 群表示的 Schur 引理的如下变形：若 $G$ 是 Abel 群，则 $G$ 的任一不可约复表示 $\varphi$ 都是次数 $1$ 的，且 $G/\ker\varphi$ 是循环群。[这可不借助练习 14，只用观察：对任意 $g\in G$，$\varphi(g)$ 的特征空间都是 $G$-稳定的。你证明 $\varphi$ 次数为 $1$ 的论证也应适用于无限 Abel 群。]
  \item 证明复表示 Schur 引理的一般形式：若 $\varphi:G\to\GL_n(\CC)$ 是不可约矩阵表示，且 $A$ 是一个与所有 $\varphi(g)$ 交换的 $n\times n$ 矩阵，则 $A$ 是标量矩阵。推出若 $\varphi$ 是忠实不可约复表示，则 $G$ 的中心是循环群，且对 $G$ 中心的所有元素 $z$，$\varphi(z)$ 都是标量矩阵。[如前一练习，$A$ 的特征空间是 $G$-稳定的。]
  \item 证明：若 $G$ 是 Abel 群，则 $G$ 的任一有限维复表示等价于一个到对角矩阵的表示（即复数域上任意有限个两两交换矩阵可同时对角化）。[这可不借助 Maschke 定理，而通过考察特征空间完成。]
  \item 证明任意有限群 $G$ 的次数 $1$ 复表示的个数等于 $|G:G'|$，其中 $G'$ 是 $G$ 的换位子群。[用练习 3 和 16。]
  \item 设非循环 Abel 群 $G$ 通过共轭作用在初等 Abel $p$-群 $V$ 上，其中 $p$ 是不整除 $|G|$ 的素数。
  \begin{enumerate}[label=(\alph*)]
    \item 证明：若 $W$ 是不可约 $\FF_pG$-子模，则存在某个非恒等元 $g\in G$，使得 $W\leq C_V(g)$（这里 $C_V(g)$ 是 $V$ 中在 $g$ 共轭作用下固定的元素子群）。
    \item 证明：$V$ 由子群 $C_V(g)$ 生成，其中 $g$ 遍历 $G$ 的所有非恒等元。
  \end{enumerate}
  \item 设 $p$ 是素数，$P$ 是 $p$-群，$F$ 是特征为 $p$ 的域。证明 $P$ 在 $F$ 上唯一的不可约表示是平凡表示。[先对 $p$ 阶群做此证明，使用 $F$ 含有所有 $p$ 次单位根（即只有 $1$）这一事实。若 $P$ 的阶不是 $p$，取 $P$ 的中心中一个 $p$ 阶元素 $z$，证明 $z$ 在不可约表示的核中，并对 $P/\langle z\rangle$ 用归纳。]
  \item 设 $p$ 是素数，$P$ 是非平凡 $p$-群，$F$ 是特征为 $p$ 的域。证明正则表示不是完全可约的。[用前一练习。]
  \item 设 $p$ 是素数，$P$ 是非平凡 $p$-群，$F$ 是特征为 $p$ 的域。证明正则表示是不可分解的。
\end{enumerate}

\section{Wedderburn 定理及若干推论}

在本节中，我们给出一个著名的由 Wedderburn 给出的分类定理。它特别描述了当 $F$ 的特征不整除 $G$ 的阶时，群代数 $FG$ 的结构。由这个分类定理，我们将推出各种推论，包括：对每个有限群 $G$，非同构不可约 $FG$-模只有有限多个。这个结果与 Maschke 定理一起，在某种意义上完成了有限群在这类域上的表示论中的 Holder 纲领。本书余下部分将发展确定与处理不可约表示的技巧，并把这些知识应用于获取群论信息。

\medskip
\noindent\textbf{定理 4.} （Wedderburn 定理）设 $R$ 是带 $1$ 的非零环（不一定交换）。则以下条件等价：
\begin{enumerate}[label=(\arabic*)]
  \item 每个 $R$-模都是投射的；
  \item 每个 $R$-模都是内射的；
  \item 每个 $R$-模都是完全可约的；
  \item 把环 $R$ 看成左 $R$-模时，有直和分解
  \[
    R=L_1\oplus L_2\oplus\cdots\oplus L_n,
  \]
  其中每个 $L_i$ 是单模（即单左理想），并且 $L_i=Re_i$，其中某些 $e_i\in R$ 满足
  \begin{enumerate}[label=(\roman*)]
    \item 若 $i\neq j$，则 $e_i e_j=0$；
    \item 对所有 $i$，$e_i^2=e_i$；
    \item $\sum_{i=1}^n e_i=1$；
  \end{enumerate}
  \item 作为环，$R$ 同构于除环上矩阵环的直积，即
  \[
    R=R_1\times R_2\times\cdots\times R_r,
  \]
  其中 $R_j$ 是 $R$ 的双边理想，且 $R_j$ 同构于某个除环 $\Delta_j$ 上所有 $n_j\times n_j$ 矩阵所成的环，$j=1,2,
  \ldots,r$。整数 $r$、整数 $n_j$ 以及除环 $\Delta_j$（至多相差同构）都由 $R$ 唯一决定。
\end{enumerate}

\begin{proof}
Wedderburn 定理的一个证明在练习 1 到 10 中列出提纲。
\end{proof}

\noindent\textbf{定义.} 满足定理 4 中任一（等价）性质的环 $R$ 称为\emph{具有极小条件的半单环}。

满足定理 4 等价条件的环 $R$ 还满足左理想的极小条件或降链条件（D.C.C.）：
\[
  \text{若 } I_1\supseteq I_2\supseteq\cdots
  \text{ 是 }R\text{ 的左理想降链，}
\]
则存在 $N\in\ZZ^+$ 使得对所有 $k\geq N$ 有 $I_k=I_N$（这解释了上述定义中该术语的使用）。我们处理的环都将有这个极小条件。例如，群代数总有这个性质，因为在任意严格下降的理想链中，理想的向量空间维数（它们是 $FG$ 的 $F$-子空间）严格下降，因此严格下降链长度至多为 $FG$ 的维数（即 $|G|$）。因此我们将用术语“半单”表示“具有极小条件的半单”。Wedderburn 定理结论 (5) 中的环 $R_i$ 称为 $R$ 的\emph{Wedderburn 分量}，$R$ 的直积分解称为其\emph{Wedderburn 分解}。注意，交换环的 Wedderburn 定理是第 16.1 节中 Artin 环分类的推论。具有极小条件的交换半单环是 Jacobson 根为零的 Artin 环，因而是域的直积（这些域就是其 Wedderburn 分量）。

应注意条件 (5) 是一个双边条件，它完全描述了 $R$ 的整体结构（环直积中的运算按分量加法与乘法定义）。特别地，它蕴含半单环也在右理想上满足极小条件。理解结论 (5) 中直积 $R_1\times\cdots\times R_r$ 的元素的一个有用方式，是把它们看成大小 $n\times n$ 的分块对角矩阵（其中 $n=\sum_{i=1}^r n_i$），形如
\[
  \begin{pmatrix}
  A_1&&&0\\
  &A_2&&\\
  &&\ddots&\\
  0&&&A_r
  \end{pmatrix},
\]
其中 $A_i$ 是 $\Delta_i$ 上的任意 $n_i\times n_i$ 矩阵。

回忆第 10.5 节：一个 $R$-模 $Q$ 是内射的，若每当 $Q$ 是某个 $R$-模 $M$ 的子模时，$M$ 都有子模 $N$ 使得 $M=Q\oplus N$。因此 Maschke 定理推出：

\medskip
\noindent\textbf{推论 5.} 若 $G$ 是有限群，$F$ 是特征不整除 $|G|$ 的域，则群代数 $FG$ 是半单环。

在得到关于不变量 $n,r,\Delta_j$ 等如何与某些域 $F$ 上群环 $FG$ 中的不变量相关的更精确信息之前，我们先研究矩阵环的结构（即 Wedderburn 定理结论 (4) 与 (5) 中描述的环）。我们引入环论中广泛使用的一些术语。回忆，环 $R$ 的中心是与 $R$ 中所有元素交换的元素所成的子环，记为 $Z(R)$（若环有 $1$，中心含有 $1$）。

\noindent\textbf{定义.}
\begin{enumerate}[label=(\arabic*)]
  \item 环 $R$ 中的非零元素 $e$ 称为\emph{幂等元}，如果 $e^2=e$。
  \item 幂等元 $e_1$ 和 $e_2$ 称为\emph{正交的}，如果 $e_1e_2=e_2e_1=0$。
  \item 幂等元 $e$ 称为\emph{本原的}，如果它不能写成两个（交换的）正交幂等元之和。
  \item 幂等元 $e$ 称为\emph{本原中心幂等元}，如果 $e\in Z(R)$ 且 $e$ 不能写成环 $Z(R)$ 中两个正交幂等元之和。
\end{enumerate}

命题 6 描述矩阵环的理想结构，命题 8 把这些结果推广到矩阵环的直积。

\medskip
\noindent\textbf{命题 6.} 设 $\Delta$ 是除环，$n\in\ZZ^+$，$R$ 是所有 $\Delta$ 上 $n\times n$ 矩阵所成的环，$I$ 是单位矩阵（即 $R$ 的 $1$）。
\begin{enumerate}[label=(\arabic*)]
  \item $R$ 的双边理想只有 $0$ 和 $R$。
  \item $R$ 的中心由标量矩阵 $aI$ 组成，其中 $a$ 在 $\Delta$ 的中心中：
  \[
    Z(R)=\{aI\mid a\in Z(\Delta)\},
  \]
  并且这是同构于 $Z(\Delta)$ 的域。特别地，若 $\Delta$ 是域，则 $R$ 的中心是所有标量矩阵所成的子环。$R$ 中唯一的中心幂等元是 $I$（特别地，$I$ 是本原的）。
  \item 令 $e_i$ 为在第 $i,i$ 位置为 $1$、其余位置为零的矩阵。则 $e_1,
  \ldots,e_n$ 是正交本原幂等元，并且 $\sum_{i=1}^n e_i=I$。
  \item $L_i=Re_i$ 是由第 $i$ 列任意而其他列全为零的矩阵组成的左理想。$L_i$ 是单左 $R$-模。每个单左 $R$-模都同构于 $L_1$（特别地，所有 $L_i$ 都是同构的 $R$-模），且作为左 $R$-模有
  \[
    R=L_1\oplus\cdots\oplus L_n.
  \]
\end{enumerate}

在证明这个命题之前，下面的结果很有用。

\medskip
\noindent\textbf{引理 7.} 设 $R$ 是任意非零环。
\begin{enumerate}[label=(\arabic*)]
  \item 若 $M$ 和 $N$ 是单 $R$-模，$\varphi:M\to N$ 是非零 $R$-模同态，则 $\varphi$ 是同构。
  \item （Schur 引理）若 $M$ 是单 $R$-模，则 $\Hom_R(M,M)$ 是除环。
\end{enumerate}

\begin{proof}
为证明 (1)，注意由于 $\varphi$ 非零，$\ker\varphi$ 是 $M$ 的真子模。由 $M$ 的单性，$\ker\varphi=0$。同样，$\varphi$ 的像是单模 $N$ 的非零子模，故 $\varphi(M)=N$。这证明 $\varphi$ 双射，所以 (1) 成立。

由 (1)，环 $\Hom_R(M,M)$ 的每个非零元素都是同构，因此有逆。这给出 (2)。
\end{proof}

\begin{proof}[命题 6 的证明]
令 $A$ 是 $R$ 中任意矩阵，其第 $i,j$ 项为 $u_{ij}$。令 $E_{ij}$ 是第 $i,j$ 位置为 $1$、其余位置为零的矩阵。以下直接计算留作练习：
\begin{enumerate}[label=(\roman*)]
  \item $E_{ij}A$ 是这样的矩阵：第 $i$ 行等于 $A$ 的第 $j$ 行，其余各行为零。
  \item $AE_{ij}$ 是这样的矩阵：第 $j$ 列等于 $A$ 的第 $i$ 列，其余各列为零。
  \item $E_{pq}AE_{rs}$ 是第 $p,s$ 项为 $u_{qr}$、其余项为零的矩阵。
\end{enumerate}
为证明 (1)，设 $J$ 是 $R$ 的任意非零双边理想，取 $A\in J$，其第 $q,r$ 项非零。给定任意 $p,s\in\{1,
\ldots,n\}$，由 (iii) 得
\[
  E_{ps}=\frac1{u_{qr}}E_{pq}AE_{rs}\in J.
\]
由于 $\{E_{ps}\mid 1\leq p\leq n,
1\leq s\leq n\}$ 的 $\Delta$-线性组合给出整个 $R$，故 $J=R$。这证明 (1)。

为证明 (2)，设 $A\in Z(R)$。则对所有 $i,j$ 有 $E_{ij}A=AE_{ij}$。由上面的 (i) 与 (ii)，立即推出 $A$ 的所有非对角项为零，且所有对角项相等。因此 $A=aI$，其中 $a\in\Delta$。此外，$A$ 还必须与所有标量矩阵集合 $\{\beta I\mid \beta\in\Delta\}$ 交换，即 $a$ 必须与 $\Delta$ 的所有元素交换。最后，由于 $Z(R)$ 是域，显然它只含有一个幂等元（即 $I$）。这证明 (2) 的各部分。

在 (3) 中，显然 $e_1,
\ldots,e_n$ 是正交幂等元且和为 $I$。它们的本原性将在证明 (4) 后给出。

接着证明 (4)。由上面的 (ii)，$Re_i=RE_{ii}$ 是第 $i$ 列任意、其他列为零的矩阵集合。此外，若 $A$ 是 $Re_i$ 的任意非零元素，则显然 $RA\subseteq Re_i$。反向包含也成立，因为若 $u_{pi}$ 是 $A$ 的一个非零项，则由 (i) 得
\[
  e_i=E_{ii}=\frac1{u_{pi}}E_{ip}A\in RA.
\]
这证明对任意非零元素 $A\in Re_i$ 都有 $Re_i=RA$，所以 $Re_i$ 必是单 $R$-模。

设 $M$ 是任意单 $R$-模。由于 $1m=m$ 对所有 $m\in M$ 成立，且 $I=\sum_{i=1}^n e_i$，存在某个 $i$ 和某个 $m\in M$ 使得 $e_i m\neq0$。对这个 $i$ 与 $m$，映射 $re_i\mapsto re_i m$ 是从单 $R$-模 $Re_i$ 到单模 $M$ 的非零 $R$-模同态。由引理 7(1)，它是同构。由 (ii)，映射 $r\mapsto rE_{i1}$ 给出 $Re_i\cong Re_1$。最后，每个矩阵都是其列的直和，所以 $R=L_1\oplus\cdots\oplus L_n$。这完成 (4) 的证明。

剩下证明 (3) 中的幂等元是本原的。若 $e_i=a+b$，其中 $a$ 与 $b$ 是正交幂等元，则我们将看到
\[
  L_i=Re_i=Ra\oplus Rb.
\]
这将与 $L_i$ 是单 $R$-模矛盾。为证明上述直和，先注意由于 $ab=ba=0$，有 $ae_i=a\in Re_i$ 与 $be_i=b\in Re_i$。对所有 $r\in R$，有 $re_i=ra+rb$，故 $Re_i=Ra+Rb$。而且 $Ra\cap Rb=0$，因为若 $ra=sb$，则 $ra=raa=sba=0$（回忆 $a=a^2$ 且 $ba=0$）。这完成全部证明。
\end{proof}

\medskip
\noindent\textbf{命题 8.} 设 $R=R_1\times R_2\times\cdots\times R_r$，其中 $R_i$ 是除环 $\Delta_i$ 上 $n_i\times n_i$ 矩阵所成的环，$i=1,2,
\ldots,r$。
\begin{enumerate}[label=(\arabic*)]
  \item 把 $R_i$ 与直积的第 $i$ 个分量认同。令 $z_i$ 是这样的 $r$ 元组：第 $i$ 位置为 $R_i$ 的单位元，其他位置为零。则 $R_i=z_iR$；对任意 $a\in R_i$，有 $z_i a=a$ 且对所有 $j\neq i$ 有 $z_j a=0$。元素 $z_1,
  \ldots,z_r$ 是 $R$ 的所有本原中心幂等元。它们两两正交，并且 $\sum_{i=1}^r z_i=1$。
  \item 令 $N$ 是任意左 $R$-模，令 $z_iN=\{z_ix\mid x\in N\}$，$1\leq i\leq r$。则 $z_iN$ 是 $N$ 的左 $R$-子模，每个 $z_iN$ 是 $R_i$-模，且对所有 $j\neq i$，$R_j$ 在 $z_iN$ 上平凡作用，并且
  \[
    N=z_1N\oplus z_2N\oplus\cdots\oplus z_rN.
  \]
  \item 单 $R$-模正是单 $R_i$-模且 $R_j$ 对 $j\neq i$ 平凡作用的模，具体意义如下。令 $M_i$ 是唯一的单 $R_i$-模（参见命题 6）。通过令 $R_j$ 对所有 $j\neq i$ 平凡作用，可把 $M_i$ 看成 $R$-模。则 $M_1,
  \ldots,M_r$ 是两两非同构单 $R$-模，任一单 $R$-模都同构于其中某个 $M_i$。显式地，$R$-模 $M_i$ 同构于单左理想 $(0,
  \ldots,0,L^{(i)},0,
  \ldots,0)$，其中第 $i$ 个分量 $L^{(i)}$ 由第一列任意、其余列为零的矩阵组成。
  \item 对任意 $R$-模 $N$，$R$-子模 $z_iN$ 是单 $R$-模的直和，且每个单模同构于 (3) 中的 $M_i$。特别地，若 $M$ 是单 $R$-模，则存在唯一的 $i$ 使得 $z_iM=M$；对此指标 $i$ 有 $M\cong M_i$，而对所有 $j\neq i$ 有 $z_jM=0$。
  \item 若每个 $\Delta_i$ 都等于域 $F$，则 $R$ 是 $F$ 上维数为 $\sum_{i=1}^r n_i^2$ 的向量空间，并且 $\dim_F Z(R)=r$。
\end{enumerate}

\begin{proof}
在 (1) 中，由于环直积中的乘法按分量进行，显然 $z_i$ 乘以 $R$ 的元素 $(a_1,
\ldots,a_r)$ 得到第 $i$ 位置为 $a_i$、其他位置为零的 $r$ 元组。因此 $R_i=z_iR$，$z_i$ 是 $R_i$ 的单位元，若 $a\in R_j$ 且 $j\neq i$，则 $z_i a=0$。显然 $z_1,
\ldots,z_r$ 是两两正交的中心幂等元，且和为 $R$ 的单位元。$R$ 的中心幂等元按定义是 $Z(R)=F_1\times F_2\times\cdots\times F_r$ 中的幂等元，其中 $F_i$ 是 $R_i$ 的中心。由命题 6，$F_i$ 是域 $Z(\Delta_i)$。若 $w=(w_1,
\ldots,w_r)$ 是任意中心幂等元，则 $w_i\in F_i$，且由 $w^2=w$ 得 $w_i^2=w_i$。由于域中方程 $x^2=x$ 的解只有 $0$ 与 $1$，$R$ 中唯一的中心幂等元是各项为 $0$ 或 $1$ 的 $r$ 元组。因此 $z_1,
\ldots,z_r$ 是本原中心幂等元；又因每个中心幂等元都是这些元素之和，它们构成 $R$ 的本原中心幂等元的全集。这证明 (1)。

为证明 (2)，令 $N$ 是任意左 $R$-模。首先注意，对任意 $z\in Z(R)$，集合 $\{zx\mid x\in N\}$ 是 $N$ 的 $R$-子模。特别地，$z_iN$ 是 $R$-子模。令 $z_ix\in z_iN$，并令 $a\in R_j$，其中 $j\neq i$。由 (1)，$a=az_j$，所以
\[
  a z_i x=(az_j)(z_i x)=az_i z_jx=0,
\]
因为 $z_i z_j=0$。因此 $R_j$ 对 $z_iN$ 平凡作用，$j\neq i$。对每个 $x\in N$，由 (1) 有
\[
  x=1x=z_1x+\cdots+z_rx,
\]
故 $N=z_1N+\cdots+z_rN$。最后，这个和是直和，因为例如若 $x\in z_1N\cap(z_2N+\cdots+z_rN)$，则 $x=z_1x$，而 $z_1$ 乘以 $z_2N+\cdots+z_rN$ 中任意元素都为零。这证明 (2)。

在 (3) 中，先注意若 $M$ 是 $R_i$-模，当定义 $R_j$ 对 $M$ 平凡作用（$j\neq i$）时，$M$ 成为 $R$-模。对这样的模 $M$，$R$-子模与 $R_i$-子模相同。因此每个 $M_i$ 是单 $R$-模，因为它是单 $R_i$-模。

接着令 $M$ 是单 $R$-模。由 (2)，
\[
  M=z_1M\oplus\cdots\oplus z_rM.
\]
由于 $M$ 没有非平凡真 $R$-子模，必有唯一的 $i$ 使得 $M=z_iM$ 且对所有 $j\neq i$ 有 $z_jM=0$。因此单 $R$-模 $M$ 被所有 $R_j$（$j\neq i$）零化。这说明 $M$ 的 $R$-子模与 $R_i$-子模相同，所以 $M$ 是单 $R_i$-模。由命题 6，$M$ 作为 $R_i$-模同构于 $M_i$。由于 $R_j$ 对 $M$ 和 $M_i$ 都平凡作用（$j\neq i$），这个 $R_i$-模同构也可看成 $R$-模同构。

若 $i\neq j$，且假设 $\theta:M_i\to M_j$ 是 $R$-模同构。若 $s_i\in M_i$，则 $s_i=z_i s_i$，所以
\[
  \theta(s_i)=\theta(z_i s_i)=z_i\theta(s_i)=0,
\]
因为 $\theta(s_i)\in M_j$ 且 $z_i$ 在 $M_j$ 上平凡作用。这与 $\theta$ 是同构矛盾，证明 $M_1,
\ldots,M_r$ 两两非同构。

最后，(3) 中描述的左理想被所有 $R_j$（$j\neq i$）平凡作用，并且由命题 6，它作为 $R_i$-模至多同构地唯一单。因此这个左理想是同构于 $M_i$ 的单 $R$-模。这证明 (3)。

对 (4)，我们已经证明若 $M$ 是任意单 $R$-模，则存在唯一的 $i$ 使得 $z_iM=M$，且对所有 $j\neq i$ 有 $z_jM=0$。此外，我们已经证明对这个指标 $i$，单 $R$-模 $M$ 同构于 $M_i$。现在令 $N$ 是任意 $R$-模。则 $z_iN$ 是 $R_i$-模，且被所有 $R_j$（$j\neq i$）平凡作用。由 Wedderburn 定理，$z_iN$ 是单 $R$-模的直和。由于每个这些单直和因子都被所有 $R_j$（$j\neq i$）平凡作用，它们都同构于 $M_i$。这证明 (4)。

在 (5) 中，若每个 $\Delta_i$ 等于域 $F$，则作为 $F$-向量空间，
\[
  R=M_{n_1}(F)\times\cdots\times M_{n_r}(F).
\]
每个矩阵环 $M_{n_i}(F)$ 在 $F$ 上维数为 $n_i^2$，故 $R$ 在 $F$ 上维数为 $\sum_{i=1}^r n_i^2$。此外，每个 $M_{n_i}(F)$ 的中心是一维的（由命题 6(2) 它同构于 $F$），故 $Z(R)$ 在 $F$ 上维数为 $r$。命题证毕。
\end{proof}

现在我们把 Wedderburn 定理（以及上述环论计算）应用到群代数 $FG$。首先，为了应用 Wedderburn 定理，我们需要 $F$ 的特征不整除 $|G|$。事实上，由于接下来几节将处理数值数据，使 $F$ 的特征为 $0$ 会更方便。其次，若能迫使出现在 $FG$ 的 Wedderburn 分解中的所有除环都等于域 $F$，事情就会简化；我们将证明，要求 $F$ 代数闭足以保证这一点。为简化记号，本书余下的大部分内容将取 $F=\CC$。读者容易检查，任意特征为 $0$ 的代数闭域（例如所有代数数的域）都可全程代替 $\CC$。

由推论 5，环 $\CC G$ 是半单的，所以由 Wedderburn 定理，
\[
  \CC G\cong R_1\times R_2\times\cdots\times R_r,
\]
其中 $R_i$ 是某个除环 $\Delta_i$ 上的 $n_i\times n_i$ 矩阵环。把这个直积的元素看成 $n\times n$ 分块矩阵（$n=\sum_{i=1}^r n_i$），第 $i$ 个块的项来自 $\Delta_i$，域 $\CC$ 在这个直积中作为标量矩阵出现，并包含在 $\CC G$ 的中心中。注意每个 $\Delta_i$ 是 $\CC$ 上的有限维向量空间。下一个结果说明这蕴含每个 $\Delta_i=\CC$。

\medskip
\noindent\textbf{命题 9.} 若 $\Delta$ 是一个除环，它是代数闭域 $F$ 上的有限维向量空间，并且 $F\subseteq Z(\Delta)$，则 $\Delta=F$。

\begin{proof}
由于 $F\subseteq Z(\Delta)$，对每个 $a\in\Delta$，由 $a$ 与 $F$ 生成的除环是一个域。此外，由于 $\Delta$ 在 $F$ 上有限维，域 $F(a)$ 是 $F$ 的有限扩张。因为 $F$ 代数闭，它没有非平凡有限扩张，所以对所有 $a\in\Delta$ 都有 $F(a)=F$，即 $\Delta=F$。
\end{proof}

这个命题证明了 $\CC G$ 的 Wedderburn 分解中每个 $R_i$ 都是 $\CC$ 上矩阵环：
\[
  \CC G\cong M_{n_1}(\CC)\times M_{n_2}(\CC)\times\cdots\times M_{n_r}(\CC).
\]
现在命题 8(5) 蕴含
\[
  \sum_{i=1}^r n_i^2=|G|.
\]

本节最后一个应用是证明 $r$（即 $\CC G$ 的 Wedderburn 分量的个数）等于 $G$ 的共轭类数。为此先注意命题 8(5) 断言 $r=\dim_{\CC}Z(\CC G)$。我们用另一种方法计算这个维数。

令 $K_1,
\ldots,K_s$ 是 $G$ 的不同共轭类（回忆它们划分 $G$）。对 $G$ 的每个共轭类 $K_i$，令
\[
  X_i=\sum_{g\in K_i}g\in\CC G.
\]
注意当 $i\neq j$ 时，$X_i$ 与 $X_j$ 没有共同项，故它们是 $\CC G$ 中线性无关的元素。此外，由于任意群元素的共轭作用置换每个共轭类中的元素，
\[
  h^{-1}X_i h=X_i,
\]
即 $X_i$ 与所有群元素交换。这证明 $X_i\in Z(\CC G)$。

我们证明这些 $X_i$ 构成 $Z(\CC G)$ 的一组基，从而证明 $s=\dim_{\CC}Z(\CC G)=r$。既然 $X_i$ 线性无关，只需证明它们张成 $Z(\CC G)$。令
\[
  X=\sum_{g\in G}a_g g
\]
是 $Z(\CC G)$ 中任意元素。由于 $h^{-1}Xh=X$，有
\[
  \sum_{g\in G} a_g h^{-1}gh=\sum_{g\in G}a_g g.
\]
由于 $G$ 的元素构成 $\CC G$ 的一组基，上述两个和中 $g$ 的系数相等。由于 $h$ 任意，固定群元素 $g$ 的同一共轭类中的每个元素在 $X$ 中具有相同系数，因此 $X$ 可写成若干 $X_i$ 的线性组合。

我们把这些结果总结如下。

\medskip
\noindent\textbf{定理 10.} 设 $G$ 是有限群。
\begin{enumerate}[label=(\arabic*)]
  \item
  \[
    \CC G\cong M_{n_1}(\CC)\times M_{n_2}(\CC)\times\cdots\times M_{n_r}(\CC).
  \]
  \item $\CC G$ 恰有 $r$ 个不同的不可约模同构类型，它们的复维数为 $n_1,n_2,
  \ldots,n_r$（因此 $G$ 恰有 $r$ 个对应次数的不等价不可约复表示）。
  \item $\sum_{i=1}^r n_i^2=|G|$。
  \item $r$ 等于 $G$ 的共轭类数。
\end{enumerate}

\medskip
\noindent\textbf{推论 11.}
\begin{enumerate}[label=(\arabic*)]
  \item 设 $A$ 是有限 Abel 群。$A$ 的每个不可约复表示都是 $1$ 维的（即是从 $A$ 到 $\CC^\times$ 的同态），并且 $A$ 有 $|A|$ 个不等价不可约复表示。此外，$A$ 的每个有限维复矩阵表示都等价于一个到对角矩阵群中的表示。
  \item 任意有限群 $G$ 的不等价（不可约）次数 $1$ 复表示的个数等于 $|G/G'|$。
\end{enumerate}

\begin{proof}
若 $A$ 是 Abel 群，则 $\CC A$ 是交换环。由于当 $k>1$ 时，$k\times k$ 矩阵环不是交换的，所以必须有每个 $n_i=1$。因此 $r=|A|$（也就是 $A$ 的共轭类数）。由于每个 $\CC A$-模都是不可约子模的直和，存在一个基，使得矩阵相对于该基都是对角的。这证明了推论的第一部分。

对一般群 $G$，每个次数 $1$ 表示 $\varphi$ 是从 $G$ 到 $\CC^\times$ 的同态。因此 $\varphi$ 通过 $G/G'$ 分解。反过来，$G/G'$ 的每个次数 $1$ 表示与自然投影 $G\to G/G'$ 复合，给出 $G$ 的一个次数 $1$ 表示。因此 $G$ 的次数 $1$ 表示正是 Abel 群 $G/G'$ 的不可约表示。由 (1) 立即得到 (2)。
\end{proof}

\paragraph{例}
\begin{enumerate}[label=(\arabic*)]
  \item 有限 Abel 群 $A$ 的不可约复表示（即从 $A$ 到 $\CC^\times$ 的同态）可显式描述如下：把 $A$ 分解为循环群的直积
  \[
    A\cong C_1\times\cdots\times C_n,
  \]
  其中 $|C_i|=|\langle x_i\rangle|=d_i$。把每个 $x_i$ 映到一个（不一定本原的）$d_i$ 次单位根，并延拓到 $x_i$ 的所有幂，得到同态。由于每个 $x_i$ 的像有 $d_i$ 种选择，由这个过程定义的 $A$ 到 $\CC^\times=\GL_1(\CC)$ 的不同同态个数等于 $|A|$。由推论 11，这些就是 $A$ 的所有不可约表示。注意，为了实现这些表示，域必须含有相应的单位根。下面的一个练习探讨循环群在 $\QQ$ 上的不可约表示。

  \item 令 $G=S_3$。由定理 10，$G$ 的不可约复表示个数为 $3$（即 $S_3$ 的共轭类数）。由于次数平方和为 $6$，次数必为 $1,1,2$。两个次数 $1$ 表示立即可见：平凡表示，以及把置换送到其符号的 $S_3$ 到 $\{\pm1\}$ 的表示（即若 $\sigma$ 为偶置换则 $\sigma\mapsto+1$，若 $\sigma$ 为奇置换则 $\sigma\mapsto-1$）。次数 $2$ 表示可通过如下分解三个基向量上的置换表示得到：令 $S_3$ 通过置换指标作用在向量空间 $V$ 的基向量 $e_1,e_2,e_3$ 上。向量 $t=e_1+e_2+e_3$ 是非零固定向量，所以 $t$ 张成一个 $1$ 维 $G$-不变子空间（它是平凡表示的一份拷贝）。由 Maschke 定理存在一个 $2$ 维 $G$-不变补空间 $I$。注意这个置换表示不是次数 $1$ 表示之和：否则它可由对角矩阵表示，置换在其作用中就会交换——这是不可能的，因为该表示忠实而 $G$ 非 Abel。因此 $I$ 不能进一步分解，所以 $I$ 实现这个不可约二维表示。实际上，$I$ 是第 1 节中描述的“增广”子模：
  \[
    I=\{w\in V\mid w=a_1e_1+a_2e_2+a_3e_3,
    \ a_1+a_2+a_3=0\}.
  \]
  显然 $e_1-e_2$ 和 $e_2-e_3$ 是 $I$ 中的线性无关向量，因此它们构成这个二维空间的一组基。相对于 $I$ 的这个基，我们得到 $S_3$ 的一个矩阵表示；例如这个矩阵表示在 $S_3$ 的两个元素上的值为
  \[
    (1\ 2)\mapsto\begin{pmatrix}-1&1\\0&1\end{pmatrix},
    \qquad
    (1\ 2\ 3)\mapsto\begin{pmatrix}0&-1\\1&-1\end{pmatrix}.
  \]

  \item 我们分解任意有限群在 $\CC$ 上的正则表示。回忆，这是由左 $\CC G$-模 $\CC G$ 本身实现的表示。由定理 10，$\CC G$ 首先是双边理想的直积：
  \[
    \CC G\cong M_{n_1}(\CC)\times M_{n_2}(\CC)\times\cdots\times M_{n_r}(\CC).
  \]
  现在由命题 6(4)，每个 $M_{n_i}(\CC)$ 进一步分解为 $n_i$ 个同构的单左理想的直和。这些左理想给出不可约 $\CC G$-模同构类的一个完整集合。因此 $G$ 在 $\CC$ 上的正则表示分解为 $G$ 的所有不可约表示的直和，其中每个不可约表示出现的重数等于它的次数。
\end{enumerate}

我们记录 $\CC G$ 的另一个性质，它将在第 19.2 节证明。

\medskip
\noindent\textbf{定理 12.} 有限群 $G$ 的每个复不可约表示的次数都整除 $G$ 的阶；也就是说，在定理 10 的记号下，对 $i=1,2,
\ldots,r$，每个 $n_i$ 都整除 $|G|$。

在下一节中，我们将用群元素描述 $\CC G$ 的本原中心幂等元。

\subsection*{练习}

设 $G$ 是有限群，$R$ 是带 $1$ 的环。
\begin{enumerate}[label=\arabic*.]
  \item 证明 Wedderburn 定理中的条件 (1) 与 (2) 等价。
  \item 证明 Wedderburn 定理中 (3) 蕴含 (2)。[设 $Q$ 是 $R$-模 $N$ 的子模。用 Zorn 引理证明存在子模 $M$，在满足 $Q\cap M=0$ 的条件下极大。若 $Q+M=N$，则 (2) 成立；否则令 $M_1$ 是 $N$ 中某个单模的完全原像，该单模位于 $N/M$ 中且不包含在 $(Q+M)/M$ 中，并说明 $M_1$ 与 $M$ 的极大性矛盾。]
  \item 证明 Wedderburn 定理中 (4) 蕴含 (3)。[设 $N$ 是非零 $R$-模。先通过考虑循环子模证明 $N$ 含有单子模。然后把 Zorn 引理应用于单子模直和的集合（适当有序），证明 $N$ 含有极大完全可约子模 $M$。若 $M\neq N$，令 $M_1$ 是 $N$ 中 $N/M$ 的一个单模的完全原像，并与 $M$ 的极大性矛盾。]
  \item 证明 Wedderburn 定理中 (5) 蕴含 (4)。[用命题 6 和命题 8 证明中的方法，把每个 $R_i$ 作为左 $R$-模分解。]
\end{enumerate}

接下来的六个练习建立一些关于环与模的一般结果，它们蕴含 Wedderburn 定理剩余的蕴含：(2) 蕴含 (5)。在这些练习中假设 $R$ 满足 (2)：每个 $R$-模都是内射的。
\begin{enumerate}[label=\arabic*.,resume]
  \item 证明 $R$ 在左理想上满足降链条件（D.C.C.）。推出 $R$ 是有限个左理想的直和。[若不然，则证明作为左 $R$-模，$R$ 是无限多个非零子模的直和。把元素 $1$ 写入这个直和中推出矛盾。]
  \item 证明 $R=R_1\times R_2\times\cdots\times R_r$，其中 $R_j$ 是双边理想且是单环（即没有非零真双边理想）。证明每个 $R_j$ 有单位元并在左理想上满足 D.C.C.。[用前一练习证明 $R$ 有极小双边理想 $R_1$。作为左 $R$-模，$R=R_1\oplus R'$，其中 $R'$ 是左理想。证明 $R'$ 是右理想，并使用 D.C.C. 归纳进行。]
  \item 令 $S$ 是带 $1$ 的单环，它在左理想上满足 D.C.C.，令 $L$ 是 $S$ 中的极小左理想。证明作为左 $S$-模，$S\cong L^n$，其中 $L^n=L\oplus\cdots\oplus L$（$n$ 个因子）。[由单性论证 $LS=S$，所以 $1=l_1s_1+\cdots+l_ns_n$，其中 $l_i\in L$，$s_i\in S$，且 $n$ 极小。证明映射 $(x_1,
  \ldots,x_n)\mapsto x_1s_1+\cdots+x_ns_n$ 是左 $S$-模的满同态；用 $L$ 与 $n$ 的极小性证明它是单射。]
  \item 设 $A$ 是任意带 $1$ 的环，$L$ 是任意左 $A$-模，$P$ 是 $n$ 个 $L$ 的直和。
  \begin{enumerate}[label=(\alph*)]
    \item 证明环同构 $\Hom_A(P,L^n)\cong M_n(D)$，其中 $D=\Hom_A(L,L)$（环 $\Hom_A(X,X)$ 中乘法是函数复合，参见第 10.2 节命题 2(4)）。
    \item 推出若 $L$ 是单 $A$-模，则 $\Hom_A(P,P)$ 同构于除环上的矩阵环。[用 Schur 引理和 (a)。]
    \item 证明环同构 $\Hom_A(A,A)\cong A^{\mathrm{opp}}$，其中 $A^{\mathrm{opp}}$ 是 $A$ 的反环（元素与加法同 $A$，但 $A^{\mathrm{opp}}$ 中乘积 $x\cdot y$ 的值为在 $A$ 中计算的 $yx$），参见第 17.4 节末尾。[任意同态由 $1$ 的像决定。]
  \end{enumerate}
  \item 证明：若 $S$ 是带 $1$ 的单环，并且在左理想上满足 D.C.C.，则对某个除环 $\Delta$，有 $S\cong M_n(\Delta)$。（此结果与练习 6 一起完成了 Wedderburn 定理中 (2) 蕴含 (5) 的存在性部分。）[用练习 7 和 8 证明 $S^{\mathrm{opp}}\cong\Hom_S(P,L)\cong M_n(D)$，其中 $D$ 是除环。然后证明 $S\cong M_n(\Delta)$，其中 $\Delta$ 是除环 $D^{\mathrm{opp}}$。]
  \item 证明上一练习同构 $S\cong M_n(\Delta)$ 中的 $\Delta$ 和 $n$ 由 $S$ 唯一决定（这证明 Wedderburn 定理中的唯一性陈述），如下进行。设 $S=M_n(\Delta)\cong M_{n'}(\Delta')$ 为环，其中 $\Delta$ 和 $\Delta'$ 是除环。
  \begin{enumerate}[label=(\alph*)]
    \item 证明 $\Delta\cong\Hom_S(L,L)$，其中 $L$ 是 $S$ 中的极小左理想。推出 $\Delta\cong\Delta'$。[用命题 6(4)。]
    \item 证明除环 $\Delta$ 上有限生成（左）模有“基”（线性无关生成集），且任意两组基具有相同基数。推出 $n=n'$。[模仿第 11.1 节推论 4(2) 的证明。]
  \end{enumerate}
  \item 证明：若 $R$ 是带 $1$ 的环且每个 $R$-模都是自由模，则 $R$ 是除环。
  \item 设 $F$ 是域，$f(x)\in F[x]$，$R=F[x]/(f(x))$。求出 $f(x)$ 在 $F[x]$ 中分解的充要条件，使得 $R$ 是半单环。当 $R$ 半单时，描述其 Wedderburn 分解。[见第 9.5 节命题 16。]
  \item 令 $G$ 是 $n$ 阶循环群，令 $R=\QQ G$。描述 $R$ 的 Wedderburn 分解，并求 $G$ 在 $\QQ$ 上的不可约表示的个数与次数。特别地，证明若 $n=p$ 是素数，则 $G$ 在 $\QQ$ 上恰有一个非平凡不可约表示，且该表示次数为 $p-1$。[回忆第 1 节第一例中 $\QQ G=\QQ[x]/(x^n-1)$。用第 9.5 节命题 16 与第 13.6 节的结果。]
  \item 令 $p$ 是素数，令 $F=\FF_p$ 是 $p$ 阶域。令 $G$ 是 $3$ 阶循环群，令 $R=FG$。分别对 $p=2$ 和 $p=7$ 描述 $R$ 的 Wedderburn 分解，并求 $G$ 在 $F$ 上的不可约表示的个数与次数。
  \item 证明：若 $P$ 是某个素数 $p$ 的 $p$-群，则 $P$ 有忠实不可约复表示当且仅当 $Z(P)$ 是循环群。[用第 1 节练习 18、定理 6.1(2) 和例 3。]
  \item 证明：若 $V$ 是不可约 $FG$-模，且 $F$ 是代数闭域，则 $\Hom_{FG}(V,V)$ 作为环同构于 $F$。
  \item 令 $F$ 是域，$R=M_n(F)$，令 $M$ 是唯一的不可约 $R$-模。证明 $\Hom_R(M,M)$ 作为环同构于 $F$。
  \item 求出 $M_n(\Delta)$ 的所有双边理想。
\end{enumerate}

\section{特征标理论与正交关系}

一般而言，对阶很大的群，表示很难计算；若不是完全不可能写出，也会非常笨重。例如，一个次数 $100$ 的矩阵表示涉及含有 $10000$ 个项的矩阵，而且即使只在群的一组生成元上描述表示，也可能需要若干个 $100\times100$ 矩阵。然而也有一些引人注目的例子，其中高次数表示已经被计算出来并被有效使用。其中一个例子是 Z. Janko 在 1965 年构造单群 $J_1$（关于单群的存在问题在第 6.2 节末尾讨论过）。Janko 正在研究单群的某些性质，并发现若任何单群具有这些性质，则它必然具有阶 $175560$，且由两个元素生成。进一步，他证明具有这些性质的假想单群必须在域 $\FF_{11}$ 上有一个 $7$ 维表示，并且两个生成元映到两个显式的 $7\times7$ 矩阵。注意，对任意单群 $S$，每个不是把所有群元素都映到单位矩阵的表示 $S\to\GL_n(F)$ 都是忠实表示，因此 $S$ 同构于它在 $\GL_n(F)$ 中的像。特别地，Janko 的计算表明，若满足这些性质的单群存在，则它唯一。M. Ward 能够证明这两个矩阵确实生成 $\GL_7(\FF_{11})$ 中一个阶为 $175560$ 的子群，因此确实存在满足 Janko 性质的单群。

类似地，S. Norton、R. Parker 和 J. Thackray 使用 $\FF_2$ 上的 $112$ 维表示构造了阶为
\[
  86775571046077562880
\]
的单群 $J_4$。该群被证明由两个元素生成，并且在他们的分析过程中计算出了这两个生成元在 $\GL_{112}(\FF_2)$ 中的显式矩阵。

1981 年，R. Griess 构造了最大的散在单群，即所谓的 Monster 群，其阶为
\[
  2^{46}\cdot3^{20}\cdot5^9\cdot7^6\cdot11^2\cdot13^3\cdot17\cdot19\cdot23\cdot29\cdot31\cdot41\cdot47\cdot59\cdot71.
\]
他的证明涉及一个 $\CC$ 上 $196884$ 维代数的自同构计算，并由此通过此次数的表示构造了 Monster 群。

类比地，对大次数 $n$ 的置换表示 $\varphi:G\to S_n$，通常也很难写出相关的显式置换。不过存在一些数值不变量，例如置换 $\varphi(g)$ 的符号与循环型；这些数值不变量可能比置换本身更容易计算（也就是说，可能无需实际写出置换本身，就能确定元素的循环型，就像第 14.8 节中计算 $\QQ$ 上 Galois 群那样）。在给定情形下，仅这些不变量就可能提供足够信息来进行某些分析，例如证明某个给定群不是单群（如第 6.2 节中所示）。此外，刚提到的不变量不依赖于集合 $\{1,2,
\ldots,n\}$ 的标号（即它们独立于 $S_n$ 中的“换基”），并且它们在 $G$ 中共轭的元素上相同。

在本节中，我们说明如何把数值不变量附加到线性表示上。这些不变量只依赖于表示的等价类（同构类型）。换言之，对每个表示 $\varphi:G\to\GL_n(F)$，我们将给每个矩阵 $\varphi(g)$ 附加 $F$ 的一个元素，并且看到在许多情形中，即使不知道矩阵 $\varphi(g)$，也可以计算这个数。此外，我们将看到这些不变量与 $\varphi$ 的相似类无关（即若把表示 $\varphi$ 替换为等价表示，则对固定 $g\in G$ 该不变量相同），并且在某种意义上它们刻画了 $G$ 的表示的相似类。

本节中 $G$ 始终是有限群，暂时 $F$ 是任意域。所有考虑的表示都假设有限维。

\noindent\textbf{定义.}
\begin{enumerate}[label=(\arabic*)]
  \item \emph{类函数}是任意从 $G$ 到 $F$ 的函数，它在 $G$ 的共轭类上为常数；即 $f:G\to F$ 满足
  \[
    f(g^{-1}xg)=f(x),
    \qquad g,x\in G.
  \]
  \item 若 $\varphi$ 是由 $FG$-模 $V$ 实现的 $G$ 的表示，则 $\varphi$ 的\emph{特征标}是函数
  \[
    \chi:G\to F,
    \qquad \chi(g)=\operatorname{tr}\varphi(g),
  \]
  其中 $\operatorname{tr}\varphi(g)$ 是 $\varphi(g)$ 相对于 $V$ 的某个基的矩阵的迹（即该矩阵对角项之和）。特征标称为不可约或可约，按照其表示不可约或可约而定。特征标的次数是实现它的任一表示的次数。
\end{enumerate}

在定义第二部分的记号下，我们也称 $\chi$ 是由 $FG$-模 $V$ 实现的特征标。一般地，特征标不是从群到域的加法群或乘法群的同态。

\paragraph{例}
\begin{enumerate}[label=(\arabic*)]
  \item 平凡表示的特征标是函数 $\chi(g)=1$，对所有 $g\in G$。这个特征标称为 $G$ 的\emph{主特征标}。
  \item 对次数 $1$ 表示，通常把特征标与表示等同（把 $1\times1$ 矩阵与其唯一项等同）。因此对 Abel 群，不可约复表示与其特征标相同（参见推论 11）。
  \item 令 $\pi:G\to S_n$ 是置换表示，令 $\varphi$ 是向量空间 $V$ 的基 $e_1,
  \ldots,e_n$ 上由此得到的线性表示：
  \[
    \varphi(g)(e_i)=e_{\pi(g)(i)}
  \]
  （参见第 1 节例 4）。相对于这个基，$\varphi(g)$ 的矩阵在第 $i,i$ 对角项为 $1$，当且仅当 $\pi(g)$ 固定 $i$；否则该矩阵在第 $i,i$ 位置为 $0$。因此若 $\eta$ 是 $\varphi$ 的特征标，则
  \[
    \eta(g)=g\text{ 在 }\{1,2,
    \ldots,n\}\text{ 上的固定点个数。}
  \]
  特别地，若 $\pi$ 是由左乘作用在 $G$ 的某个子群 $H$ 的左陪集集合上得到的置换表示，则所得特征标称为 $G$ 在 $H$ 上的\emph{置换特征标}。
  \item 例 3 中 $\pi$ 是 $G$ 的正则置换表示的特殊情形值得记录：若 $\varphi$ 是 $G$ 的正则表示（由模 $FG$ 实现），$\rho$ 是其特征标，则
  \[
    \rho(g)=\begin{cases}
      0,&g\neq1,\\
      |G|,&g=1.
    \end{cases}
  \]
  $G$ 的正则表示的特征标称为 $G$ 的\emph{正则特征标}。注意这给出特征标取值为 $0$ 的具体例子，而且它不是从 $G$ 到 $F$ 或 $F^\times$ 的群同态。
  \item 令 $\varphi:D_{2n}\to\GL_2(\RR)$ 是第 1 节第二组例 6 中描述的显式矩阵表示。若 $\chi$ 是 $\varphi$ 的特征标，则取给定 $2\times2$ 矩阵的迹可得 $\chi(r)=2\cos(2\pi/n)$ 且 $\chi(s)=0$。由于 $\varphi$ 把 $D_{2n}$ 的恒等元送到 $2\times2$ 单位矩阵，$\chi(1)=2$。
  \item 令 $\varphi:Q_8\to\GL_2(\CC)$ 是第 1 节第二组例 7 中描述的显式矩阵表示。若 $\chi$ 是 $\varphi$ 的特征标，则取迹可得 $\chi(i)=0$ 与 $\chi(j)=0$。由于元素 $-1\in Q_8$ 映到 $2\times2$ 单位矩阵的负矩阵，$\chi(-1)=-2$。由于 $\varphi$ 把 $Q_8$ 的恒等元送到 $2\times2$ 单位矩阵，$\chi(1)=2$。
  \item 令 $\varphi:Q_8\to\GL_4(\RR)$ 是第 1 节第二组例 8 中描述的矩阵表示。若 $\chi$ 是 $\varphi$ 的特征标，则由所列矩阵可见 $\chi(i)=\chi(j)=0$。由于 $\varphi$ 把 $Q_8$ 的恒等元送到 $4\times4$ 单位矩阵，$\chi(1)=4$。
\end{enumerate}

对 $n\times n$ 矩阵 $A$ 和 $B$，直接计算给出 $\operatorname{tr}AB=\operatorname{tr}BA$。若 $A$ 可逆，则这蕴含
\[
  \operatorname{tr}A^{-1}BA=\operatorname{tr}B.
\]
因此表示的特征标与实现它的向量空间基的选择无关，即
\[
  \text{等价表示具有相同特征标。}\tag{18.1}
\]
令 $\varphi$ 是 $G$ 的次数 $n$ 表示，其特征标为 $\chi$。由于对所有 $g,x\in G$，$\varphi(g^{-1}xg)=\varphi(g)^{-1}\varphi(x)\varphi(g)$，取迹得到
\[
  \text{表示的特征标是类函数。}\tag{18.2}
\]
由于 $n\times n$ 单位矩阵的迹为 $n$，而 $\varphi$ 把 $G$ 的恒等元送到恒等线性变换（或单位矩阵），
\[
  \chi(1)\text{ 是 }\varphi\text{ 的次数。}\tag{18.3}
\]

若 $V$ 是 $FG$-模，其对应表示的特征标为 $\chi$，则群环 $FG$ 的每个元素都作为从 $V$ 到 $V$ 的线性变换而作用。因此每个 $\sum_{g\in G}a_g g\in FG$，当它作为 $V$ 到 $V$ 的线性变换来看时，都有一个迹。$g\in G$ 在 $V$ 上作用的迹按定义是 $\chi(g)$。由于矩阵线性组合的迹等于迹的线性组合，$\sum_{g\in G}a_g g$ 在 $V$ 上作用的迹为 $\sum_{g\in G}a_g\chi(g)$。注意这个 $FG$ 上的迹函数是 $G$ 的特征标 $\chi$ 到从 $FG$ 到 $F$ 的 $F$-线性变换的唯一延拓。由此我们也把 $G$ 的特征标看作定义在群环 $FG$ 上。

注意上面例 3 中，若域 $F$ 有特征 $p>0$，则特征标值模 $p$ 后可能为零，即使固定点个数非零。为了避免这类异常，并使用当 $F$ 代数闭时 Wedderburn 定理的推论，我们再次把域专门取为复数域（或任意特征为 $0$ 的代数闭域）。由前一节结果
\[
  \CC G\cong M_{n_1}(\CC)\times M_{n_2}(\CC)\times\cdots\times M_{n_r}(\CC).
  \tag{18.4}
\]
在本节余下部分固定如下记号：
\[
\begin{array}{l}
  M_1,M_2,
  \ldots,M_r\text{ 是不等价不可约 }\CC G\text{-模，}\
  \chi_i\text{ 是由 }M_i\text{ 实现的特征标， }1\leq i\leq r.
\end{array}
\tag{18.5}
\]
因此 $r$ 是 $G$ 的共轭类数；必要时可重标 $M_1,
\ldots,M_r$，使得 $\chi_i$ 的次数为 $n_i$（这也等于 $M_i$ 在 $\CC$ 上的维数）。

现在每个（有限维）$\CC G$-模 $M$ 都同构（等价）于不可约模的直和：
\[
  M\cong a_1M_1\oplus a_2M_2\oplus\cdots\oplus a_rM_r,
  \tag{18.6}
\]
其中 $a_i$ 是非负整数，表示不可约模 $M_i$ 在该直和分解中的重数；也就是说
\[
  a_iM_i=\underbrace{M_i\oplus\cdots\oplus M_i}_{a_i\text{ 次}}.
\]
注意，若表示 $\varphi$ 由模 $M$ 实现，且 $M=M_1\oplus M_2$，则可取 $M$ 的一个基为 $M_1$ 的一组基与 $M_2$ 的一组基的并。相对于这个基，矩阵表示形如
\[
  \varphi(g)=\begin{pmatrix}
  \varphi_1(g)&0\\
  0&\varphi_2(g)
  \end{pmatrix},
\]
其中 $\varphi_i$ 是由 $M_i$ 实现的表示，$i=1,2$。立即可见，若 $\psi$ 是 $\varphi$ 的特征标，$\psi_i$ 是 $\varphi_i$ 的特征标，则 $\psi(g)=\psi_1(g)+\psi_2(g)$，即 $\psi=\psi_1+\psi_2$。归纳得到：
\[
  \text{一个表示的特征标是其直和分解中组成分量特征标之和。}\tag{18.7}
\]
若 $\psi$ 是上面 (18.6) 中模 $M$ 实现的特征标，则
\[
  \psi=a_1\chi_1+a_2\chi_2+\cdots+a_r\chi_r.
  \tag{18.8}
\]
因此每个（复）特征标都是不可约（复）特征标的非负整数和。反过来，通过取模的直和可知，每个这样的特征标和都是 $G$ 的某个复表示的特征标。

接下来证明特征标与复表示等价类之间的对应是双射。令 $z_1,z_2,
\ldots,z_r$ 是前一节描述的 $\CC G$ 的本原中心幂等元。由于它们正交（或等价地，它们是 $\CC G$ 分解为 $r$ 个子环直积时某一位置为 $1$、其余位置为零的 $r$ 元组），$z_1,
\ldots,z_r$ 是 $\CC G$ 中线性无关的元素。如上，每个不可约特征标 $\chi_i$ 是 $\CC G$ 上的函数。由命题 8(3) 有：
\begin{enumerate}[label=(\alph*)]
  \item 若 $j\neq i$，则 $z_jM_i=0$，即 $z_j$ 在 $M_i$ 上作为零矩阵作用，因此 $\chi_i(z_j)=0$；
  \item $z_i$ 在 $M_i$ 上作为恒等映射作用，因此 $\chi_i(z_i)=n_i$。
\end{enumerate}
因此 $\chi_1,
\ldots,\chi_r$ 是线性无关集 $z_1,
\ldots,z_r$ 的对偶基的非零倍数，故它们是线性无关函数。现在若 (18.6) 中的 $\CC G$-模 $M$ 还可以以另一种方式分解为不可约模，例如
\[
  M\cong b_1M_1\oplus b_2M_2\oplus\cdots\oplus b_rM_r,
\]
则得到关系
\[
  a_1\chi_1+a_2\chi_2+\cdots+a_r\chi_r
  =b_1\chi_1+b_2\chi_2+\cdots+b_r\chi_r.
\]
由不可约特征标的线性无关性，$b_i=a_i$ 对所有 $i$ 成立。因此，在 $M$ 分解为不可约子模直和的任一分解中，不可约模 $M_i$ 的重数都相同，$1\leq i\leq r$。特别地，
\[
  \text{两个表示等价当且仅当它们具有相同特征标。}\tag{18.9}
\]

这种唯一性也可用另一种方式看出。首先用命题 8(2) 唯一地把任意有限维 $\CC G$-模 $M$ 分解为
\[
  M=z_1M\oplus z_2M\oplus\cdots\oplus z_rM.
\]
由同一命题的 (4)，$z_iM$ 是单模直和，每个单模同构于 $M_i$。通过数维数，$M_i$ 在 $z_iM$ 的直和分解中的重数等于
\[
  \frac{\dim z_iM}{\dim M_i}.
\]
这证明 $M$ 的任何分解为单子模直和的分解中，$M_i$ 的重数都是唯一确定的。

注意，与 $F[x]$-模分解为循环子模类似，一个 $\CC G$-模可能有许多分解为不可约直和的方式——只有重数是唯一的（也见练习）。更精确地说，与单个线性变换的 Jordan 标准形比较，直和因子 $a_iM_i=M_i\oplus\cdots\oplus M_i$（$a_i$ 次）即子模 $z_iM$，类似于对应一个特征值的广义特征空间。这个 $M$ 的子模是唯一的（正如广义特征空间唯一），称为 $M$ 的 $\chi_i$-等型分量。在 $\chi_i$-等型分量内部，直和项 $M_i$ 类似于一维特征空间；就像自同态的特征空间没有唯一基一样，那里也没有唯一的分解。若 $G=\langle g\rangle$ 是有限循环群，则 $G$ 的等型分量与 $g$ 的广义特征空间相同。

观察：$G$ 上所有复值类函数的向量空间有一组基，由在某个给定共轭类上为 $1$、在其他所有共轭类上为 $0$ 的函数构成。这样的函数有 $r$ 个，其中 $r$ 是 $G$ 的共轭类数，所以类函数复向量空间的维数为 $r$。由于 $G$ 的（复）不可约特征标个数等于共轭类数，且这些特征标是线性无关的类函数，我们看到
\[
  \text{不可约特征标构成所有复类函数空间的一组基。}\tag{18.10}
\]

特征标理论的下一步，是在类函数空间上放置 Hermitian 内积结构，并证明不可约特征标关于此内积构成一组标准正交基。对类函数 $\theta$ 与 $\psi$，定义
\[
  (\theta,\psi)=\frac1{|G|}\sum_{g\in G}\theta(g)\overline{\psi(g)},
\]
其中横线表示复共轭。容易检查 $(\ ,\ )$ 是 Hermitian 的：对 $\alpha,\beta\in\CC$，
\begin{enumerate}[label=(\alph*)]
  \item $(\alpha\theta_1+\beta\theta_2,
  \psi)=\alpha(\theta_1,
  \psi)+\beta(\theta_2,
  \psi)$；
  \item $(\theta,
  \alpha\psi_1+\beta\psi_2)=\overline{\alpha}(\theta,
  \psi_1)+\overline{\beta}(\theta,
  \psi_2)$；
  \item $(\theta,
  \psi)=\overline{(\psi,
  \theta)}$。
\end{enumerate}

我们的主要目标是证明不可约特征标关于这个 Hermitian 形式构成复类函数空间的一组标准正交基（我们已经知道它们是一组基）。一旦在下一个命题中显式确定本原中心幂等元，这个事实将由它们的正交性推出。

\medskip
\noindent\textbf{命题 13.} 设 $z_1,
\ldots,z_r$ 是 $\CC G$ 中的正交本原中心幂等元，标号方式使得 $z_i$ 在不可约 $\CC G$-模 $M_i$ 上作为恒等映射作用；令 $\chi_i$ 是由 $M_i$ 实现的特征标。则
\[
  z_i=\frac{\chi_i(1)}{|G|}\sum_{g\in G}\overline{\chi_i(g)}\,g.
\]

\begin{proof}
令 $z=z_i$，写
\[
  z=\sum_{g\in G}a_g g.
\]
回忆本节例 4：若 $\rho$ 是 $G$ 的正则特征标，则
\[
  \rho(g)=\begin{cases}
  0,&g\neq1,\\
  |G|,&g=1,
  \end{cases}
  \tag{18.11}
\]
并回忆第 2 节最后一个例子：
\[
  \rho=\sum_{j=1}^r\chi_j(1)\chi_j.
  \tag{18.12}
\]
为求系数 $a_g$，把 $\rho$ 作用于 $zg^{-1}$，并用 $\rho$ 的线性性和 (18.11)，得到
\[
  \rho(zg^{-1})=a_g|G|.
\]
另一方面，用 (18.12) 计算 $\rho(zg^{-1})$ 得
\[
  \sum_{j=1}^r\chi_j(1)\chi_j(zg^{-1})=a_g|G|.
  \tag{18.13}
\]
令 $\varphi_j$ 是由 $M_j$ 实现的不可约表示，$1\leq j\leq r$。由于可把 $\varphi_j$ 看成从 $\CC G$ 到 $\End M_j$ 的代数同态，有
\[
  \varphi_j(zg^{-1})=\varphi_j(z)\varphi_j(g^{-1}).
\]
我们已经观察到：若 $j\neq i$，则 $\varphi_j(z)=0$；若 $j=i$，则 $\varphi_i(z)$ 是 $M_i$ 上的恒等自同态。因此
\[
  \varphi_j(zg^{-1})=\begin{cases}
  0,&j\neq i,\\
  \varphi_i(g^{-1}),&j=i.
  \end{cases}
\]
这证明 $\chi_j(zg^{-1})=\chi_i(g^{-1})\delta_{ij}$，其中 $\delta_{ij}$ 为 Kronecker delta。代入 (18.13) 得
\[
  a_g=\frac1{|G|}\chi_i(1)\chi_i(g^{-1}).
\]
由下面命题 14 可知 $\chi_i(g^{-1})=\overline{\chi_i(g)}$，这正是命题中 $g$ 的系数。证明完成。
\end{proof}

不可约特征标的标准正交性将通过下列计算直接由中心本原幂等元的正交性推出：
\[
\begin{aligned}
  z_i\delta_{ij}=z_iz_j
  &=\frac{\chi_i(1)\chi_j(1)}{|G|^2}
    \sum_{g,h\in G}\overline{\chi_i(g)}\,\overline{\chi_j(h)}\,gh\\
  &=\frac{\chi_i(1)\chi_j(1)}{|G|^2}
    \sum_{y\in G}\left(\sum_{x\in G}\overline{\chi_i(xy^{-1})}\,\overline{\chi_j(x^{-1})}\right)y,
\end{aligned}
\]
（从前一个和式得到后一个和式，作代换 $x=g$ 与 $y=gh$）。由于 $G$ 的元素是 $\CC G$ 的基，我们可把它的系数与命题 13 中 $z_i$ 的系数比较（比较 $g$ 的系数），得到
\[
  \delta_{ij}\frac{\chi_i(1)}{|G|}\overline{\chi_i(g)}
  =\frac{\chi_i(1)\chi_j(1)}{|G|^2}
  \sum_{x\in G}\overline{\chi_i(xg^{-1})}\,\overline{\chi_j(x^{-1})}.
\]
化简并把 $g$ 换为 $g^{-1}$，得到
\[
  \delta_{ij}\frac{\chi_i(g)}{\chi_j(1)}
  =\frac1{|G|}\sum_{x\in G}\chi_i(xg)\chi_j(x^{-1}),
  \qquad g\in G.
  \tag{18.14}
\]
令 $g=1$，得到
\[
  \delta_{ij}=\frac1{|G|}\sum_{x\in G}\chi_i(x)\chi_j(x^{-1}).
  \tag{18.15}
\]
如果 $\chi_j(x^{-1})=\overline{\chi_j(x)}$，则右边正是内积 $(\chi_i,\chi_j)$；这就是下一个命题的内容。

\medskip
\noindent\textbf{命题 14.} 若 $\psi$ 是 $G$ 的任意特征标，则对所有 $x\in G$，$\psi(x)$ 是若干单位根之和，并且
\[
  \psi(x^{-1})=\overline{\psi(x)}.
\]

\begin{proof}
设 $\varphi$ 是特征标为 $\psi$ 的表示，固定 $x\in G$，令 $|x|=k$。由于 $\varphi(x)$ 的极小多项式整除 $X^k-1$（因而有不同根），存在底层向量空间的一组基，使得 $\varphi(x)$ 关于这个基的矩阵是对角矩阵，且对角项为 $k$ 次单位根。由于 $\psi(x)$ 是对角项之和（且不依赖于基的选择），$\psi(x)$ 是若干单位根之和。此外，若 $\epsilon$ 是单位根，则 $\epsilon^{-1}=\overline{\epsilon}$。因此，一个对角项为单位根的对角矩阵的逆矩阵，是对角项为那些单位根的复共轭的对角矩阵。由于一个和的复共轭等于各项复共轭之和，
\[
  \psi(x^{-1})=\operatorname{tr}\varphi(x^{-1})
  =\overline{\operatorname{tr}\varphi(x)}=
  \overline{\psi(x)}.
\]
\end{proof}

注意，在命题 14 的证明中，我们首先固定一个群元素 $x$，然后选择表示空间的一组基，使得 $\varphi(x)$ 是对角矩阵。总是可以对角化单个元素，但当且仅当 $\varphi$ 相似于若干次数 $1$ 表示之和时，才可以同时对角化所有 $\varphi(x)$。

把上述命题与式 (18.15) 结合，得到：

\medskip
\noindent\textbf{定理 15.} （群特征标的第一正交关系）设 $G$ 是有限群，$\chi_1,
\ldots,\chi_r$ 是 $G$ 在 $\CC$ 上的不可约特征标。则关于上面的内积 $(\ ,\ )$，有
\[
  (\chi_i,\chi_j)=\delta_{ij},
\]
并且不可约特征标构成类函数空间的一组标准正交基。特别地，若 $\theta$ 是任意类函数，则
\[
  \theta=\sum_{i=1}^r (\theta,\chi_i)\chi_i.
\]

\begin{proof}
我们已经证明不可约特征标构成类函数空间的一组标准正交基。若 $\theta$ 是任意类函数，写 $\theta=\sum_{i=1}^r a_i\chi_i$，其中 $a_i\in\CC$。由 Hermitian 内积的线性性得 $a_i=(\theta,\chi_i)$，如所断言。
\end{proof}

我们列出第二正交关系而不证明；本书应用中不会用到它。

\medskip
\noindent\textbf{定理 16.} （群特征标的第二正交关系）在上述记号下，对任意 $x,y\in G$，
\[
  \sum_{i=1}^r \chi_i(x)\overline{\chi_i(y)}
  =\begin{cases}
  |C_G(x)|,&x\text{ 与 }y\text{ 在 }G\text{ 中共轭},\\
  0,&\text{否则。}
  \end{cases}
\]

\noindent\textbf{定义.} 对 $G$ 上任意类函数 $\theta$，$\theta$ 的\emph{范数}是 $(\theta,\theta)^{1/2}$，记为 $\|\theta\|$。

当类函数用不可约特征标表示为 $\theta=\sum a_i\chi_i$ 时，其范数容易计算为
\[
  \|\theta\|=\left(\sum |a_i|^2\right)^{1/2}.
\]
因此
\[
  \text{一个特征标范数为 }1\text{ 当且仅当它不可约。}
\]

最后，观察到计算特征标 $\theta$ 与 $\psi$ 的内积可如下简化。若 $K_1,
\ldots,K_r$ 是 $G$ 的共轭类，大小分别为 $d_1,
\ldots,d_r$，代表元分别为 $g_1,
\ldots,g_r$，则和式 $(\theta,\psi)$ 中值 $\theta(g_i)\overline{\psi(g_i)}$ 出现 $d_i$ 次，对 $K_i$ 中每个元素出现一次。合并这些项得到
\[
  (\theta,\psi)=\frac1{|G|}\sum_{i=1}^r d_i\theta(g_i)\overline{\psi(g_i)},
\]
这是只在 $G$ 的共轭类代表元上求和。特别地，$\theta$ 的范数满足
\[
  \|\theta\|^2=(\theta,
  \theta)=\frac1{|G|}\sum_{i=1}^r d_i|\theta(g_i)|^2.
\]

\paragraph{例}
\begin{enumerate}[label=(\arabic*)]
  \item 令 $G=S_3$，令 $\pi$ 是本节开头例子中描述的次数 $3$ 置换特征标。回忆 $\pi(\sigma)$ 等于 $\sigma$ 固定集合 $\{1,2,3\}$ 中元素的个数。$S_3$ 的共轭类可由 $1$、$(1\ 2)$ 和 $(1\ 2\ 3)$ 代表，大小分别为 $1,3,2$，且 $\pi(1)=3$、$\pi((1\ 2))=1$、$\pi((1\ 2\ 3))=0$。因此
  \[
    \|\pi\|^2=\frac16\left[1\pi(1)^2+3\pi((1\ 2))^2+2\pi((1\ 2\ 3))^2\right]
    =\frac16(9+3+0)=2.
  \]
  这说明 $\pi$ 是两个不同不可约特征标之和，且每个出现重数 $1$。令 $\chi_1$ 是 $S_3$ 的主特征标，所以对所有 $\sigma\in S_3$，$\chi_1(\sigma)=\overline{\chi_1(\sigma)}=1$。则
  \[
    (\pi,\chi_1)=\frac16\left[\pi(1)\chi_1(1)+3\pi((1\ 2))\chi_1((1\ 2))+2\pi((1\ 2\ 3))\chi_1((1\ 2\ 3))\right]
    =\frac16(3+3+0)=1.
  \]
  所以主特征标作为 $\pi$ 的组成分量出现重数 $1$。这证明 $\pi=\chi_1+\chi_2$，其中 $\chi_2$ 是 $S_3$ 的某个次数 $2$ 不可约特征标（与前面对此表示的分解一致）。这也说明 $\chi_2$ 在 $\sigma\in S_3$ 上的值等于 $\sigma$ 的固定点个数减 $1$。

  \item 令 $G=S_4$，$\pi$ 是次数 $4$ 的自然置换特征标（所以仍有 $\pi(\sigma)$ 等于 $\sigma$ 的固定点个数）。$S_4$ 的共轭类可由 $1$、$(1\ 2)$、$(1\ 2\ 3)$、$(1\ 2\ 3\ 4)$ 和 $(1\ 2)(3\ 4)$ 代表，大小分别为 $1,6,8,6,3$。同样计算：
  \[
  \begin{aligned}
    \|\pi\|^2
    &=\frac1{24}[1\pi(1)^2+6\pi((1\ 2))^2+8\pi((1\ 2\ 3))^2
      +6\pi((1\ 2\ 3\ 4))^2+3\pi((1\ 2)(3\ 4))^2]\\
    &=\frac1{24}(16+24+8+0+0)=2.
  \end{aligned}
  \]
  因此 $\pi$ 有两个不同不可约组成分量。若 $\chi_1$ 是 $S_4$ 的主特征标，则
  \[
  \begin{aligned}
    (\pi,\chi_1)
    &=\frac1{24}[1\pi(1)+6\pi((1\ 2))+8\pi((1\ 2\ 3))
      +6\pi((1\ 2\ 3\ 4))+3\pi((1\ 2)(3\ 4))]\\
    &=\frac1{24}(4+12+8+0+0)=1.
  \end{aligned}
  \]
  这证明次数 $4$ 的置换特征标是主特征标与一个次数 $3$ 不可约特征标之和。

  \item 令 $G=D_8$，其中
  \[
    D_8=\langle r,s\mid s^2=r^4=1,
    \ rs=sr^{-1}\rangle.
  \]
  $D_8$ 的共轭类由 $1,s,r,r^2,sr$ 代表，大小分别为 $1,2,2,1,2$。令 $\varphi$ 是第 1 节例 6 中由把正方形嵌入 $\RR^2$ 得到的次数 $2$ 矩阵表示：
  \[
    \varphi(s)=\begin{pmatrix}0&1\\1&0\end{pmatrix},
    \quad
    \varphi(r)=\begin{pmatrix}0&-1\\1&0\end{pmatrix},
    \quad
    \varphi(r^2)=\begin{pmatrix}-1&0\\0&-1\end{pmatrix},
    \quad
    \varphi(sr)=\begin{pmatrix}0&-1\\-1&0\end{pmatrix}.
  \]
  令 $\psi$ 是这个表示的特征标（把实矩阵看成复矩阵）。由于 $\psi$ 实值，计算得
  \[
  \begin{aligned}
    \|\psi\|^2
    &=\frac18[1\psi(1)^2+2\psi(s)^2+2\psi(r)^2+1\psi(r^2)^2+2\psi(sr)^2]\\
    &=\frac18(4+0+0+4+0)=1.
  \end{aligned}
  \]
  因此表示 $\varphi$ 是不可约的（即使允许用复矩阵作相似变换）。
\end{enumerate}

我们已经看到两个特征标之和仍是特征标。具体地，若 $\psi_1$ 与 $\psi_2$ 是表示 $\varphi_1$ 与 $\varphi_2$ 的特征标，则 $\psi_1+\psi_2$ 是 $\varphi_1+\varphi_2$ 的特征标。

\medskip
\noindent\textbf{命题 17.} 若 $\psi_1$ 与 $\psi_2$ 是特征标，则它们的乘积 $\psi_1\psi_2$ 也是特征标。

\begin{proof}
令 $V_1$ 与 $V_2$ 是实现特征标 $\psi_1$ 与 $\psi_2$ 的 $\CC G$-模，定义 $W=V_1\otimes_{\CC}V_2$。由于每个 $g\in G$ 都作为线性变换作用在 $V_1$ 与 $V_2$ 上，作用
\[
  g(v_1\otimes v_2)=(gv_1)\otimes(gv_2)
\]
在线性延拓后，由第 11.2 节命题 17 给出 $W$ 上的良定义线性变换。容易检查这个作用也把 $W$ 做成 $\CC G$-模。由第 11.2 节练习 38，$W$ 实现的特征标为 $\psi_1\psi_2$。
\end{proof}

下一章将包含更多显式特征标计算，以及群特征标在证明某些类群的定理中的应用。

\subsection*{关于 Fourier 分析与群特征标的若干说明}

这段简短讨论旨在指出上述结果与数学其他领域之间的一些联系。

到目前为止所描述的群表示理论，是称为\emph{调和分析}的数学领域的一个特殊分支。读者可能已经熟悉 Fourier 级数的基本理论，它也属于这一范畴。我们作若干观察，说明有限群表示论如何对应于某些无限群的“Fourier 级数”（特别是圆周上的 Fourier 级数）。为了数学上精确，需要 Lebesgue 积分来保证某些（Hilbert）空间的完备性；不过读者可把“Lebesgue”替换为“Riemann”来体会其思想。

令 $G$ 是复平面中单位圆上的点所成的乘法群：
\[
  G=\{z\in\CC\mid |z|=1\}.
\]
我们通常把 $G$ 看成 $\RR$ 中区间 $[0,2\pi]$ 且两端点识别，即加法群 $\RR/2\pi\ZZ$（同构为：实数 $x$ 对应复数 $e^{ix}$）。注意 $G$ 有平移不变测度，即 Lebesgue 测度，圆的测度为 $2\pi$。对有限群，计数测度是平移不变测度（所以子集 $H$ 的测度是其中元素个数 $|H|$），而关于这个计数测度的有限群上的积分只是有限和。

空间
\[
  L^2(G)=\{f:G\to\CC\mid f\text{ 可测且 }|f|^2\text{ 在 }G\text{ 上可积}\}
\]
扮演无限群 $G$ 的群代数的角色。这个空间在函数卷积下成为一个带 $1$ 的交换环：对 $f,g\in L^2(G)$，乘积 $f*g:G\to\CC$ 定义为
\[
  (f*g)(x)=\frac1{2\pi}\int_0^{2\pi} f(x-y)g(y)\,dy,
  \qquad x\in G.
\]
回忆，对有限群 $H$，群代数在形式上也是 $H$ 上复值函数所成的环，乘法为卷积；这些函数写作形式和——元素 $\sum a_g g\in\CC G$ 表示把 $g$ 送到 $a_g\in\CC$ 的函数。

$G$ 到 $\GL_1(\CC)$ 的连续同态的完整集合为
\[
  e_n(x)=e^{inx},
  \qquad x\in[0,2\pi],
  \ n\in\ZZ.
\]
回忆，对有限 Abel 群，所有不可约表示都是 $1$ 维的；对 $1$ 维表示，特征标与表示可等同。

环 $L^2(G)$ 带有 Hermitian 内积：对 $f,g\in L^2(G)$，
\[
  (f,g)=\frac1{2\pi}\int_0^{2\pi} f(t)\overline{g(t)}\,dt.
\]
在这个内积下，$\{e_n\mid n\in\ZZ\}$ 是标准正交基（这里“基”在分析意义下使用：它们线性无关，且唯一与全部这些函数正交的函数是 $0$）。此外，
\[
  L^2(G)=\widehat{\bigoplus_{n\in\ZZ}} E_n,
\]
其中 $E_n$ 是由 $e_n$ 张成的一维子空间，直和号上的帽表示在 $L^2$-拓扑中取直和的闭包，等号表示 $L^2$ 意义下的相等。回忆，有限 Abel 群的群代数是不可约一维子模的直和，每个出现重数一。这些事实蕴含 Fourier 分析中的著名结果：$[0,2\pi]$ 上每个平方可积函数 $f(x)$ 都有 Fourier 级数
\[
  f(x)=\sum_{n=-\infty}^{\infty} c_n e^{inx},
\]
其中 Fourier 系数 $c_n$ 由
\[
  c_n=(f,e_n)=\frac1{2\pi}\int_0^{2\pi} f(t)e^{-int}\,dt
\]
给出。

这段简短描述说明有限群表示论如何推广到某些无限群，而我们已经证明的结果在后一背景中可能早已为人熟悉。事实上，对任意（不一定 Abel 的）紧 Lie 群也有完全类似的理论——这里不可约（复）表示不必是一维的，但它们全都是有限维的，且 $L^2(G)$ 分解为它们的直和，每个不可约表示出现的重数等于其次数。这个理论中（至少在入门层面）的重点常常在于能够把函数表示成（Fourier）级数，然后利用这些级数解决其他问题（例如求解微分方程）。底层群提供了建立这种“调和分析”的“对称性”，而不是自身成为主要研究对象。

\subsection*{练习}

设 $G$ 是有限群。除非另有说明，所有表示与特征标都在 $\CC$ 上。
\begin{enumerate}[label=\arabic*.]
  \item 证明对任意交换环上的 $n\times n$ 矩阵 $A$ 与 $B$，有 $\operatorname{tr}AB=\operatorname{tr}BA$。
  \item 在 (a) 到 (c) 中，令 $\psi$ 是指定表示 $\varphi$ 实现的特征标。
  \begin{enumerate}[label=(\alph*)]
    \item 令 $\varphi$ 是第 1 节第二组例子中例 6 描述的 $D_{10}$ 的次数 $2$ 表示（这里 $n=5$），并证明 $\|\psi\|^2=1$（因此 $\varphi$ 不可约）。
    \item 令 $\varphi$ 是第 1 节第二组例子中例 7 描述的 $Q_8$ 的次数 $2$ 表示，并证明 $\|\psi\|^2=1$（因此 $\varphi$ 不可约）。
    \item 令 $\varphi$ 是第 1 节第二组例子中例 8 描述的 $Q_8$ 的次数 $4$ 表示，并证明 $\|\psi\|^2=4$（因此虽然 $\varphi$ 在 $\RR$ 上不可约，但在 $\CC$ 上 $\varphi$ 分解为某个次数 $2$ 不可约表示的两份）。
  \end{enumerate}
  \item 若 $\chi$ 是 $G$ 的不可约特征标，证明 $\CC G$-模的 $\chi$-等型子空间是唯一的。
  \item 证明：若 $N$ 是任意不可约 $\CC G$-模，且 $M=N\oplus N$，则 $M$ 有无限多种分解为两个 $N$ 的直和的方式。
  \item 证明类函数是特征标，当且仅当它是不可约特征标的正整数线性组合。
  \item 设 $\varphi:G\to\GL(V)$ 是特征标为 $\psi$ 的表示。令 $W=\{v\in V\mid \varphi(g)(v)=v\text{ 对所有 }g\in G\}$，即 $V$ 中被 $G$ 的所有元素逐点固定的子空间。证明 $\dim W=(\psi,\chi_1)$，其中 $\chi_1$ 是 $G$ 的主特征标。
  \item 假设 $V$ 是 $\CC G$-模，$G$ 通过置换基 $\mathcal B=\{e_1,
  \ldots,e_n\}$ 作用在 $V$ 上。把 $\mathcal B$ 写成 $G$ 在 $\mathcal B$ 上的轨道的不交并 $\mathcal B_1,
  \ldots,\mathcal B_t$。
  \begin{enumerate}[label=(\alph*)]
    \item 证明作为 $\CC G$-模，$V$ 分解为 $V_1\oplus\cdots\oplus V_t$，其中 $V_i$ 是 $\mathcal B_i$ 的张成空间。
    \item 证明若 $v_i$ 是 $\mathcal B_i$ 中向量之和，则由 $v_i$ 张成的一维子空间是 $V_i$ 中唯一实现平凡表示的 $\CC G$-子模（换言之，$V_i$ 中任意在 $G$ 作用下固定的向量都是 $v_i$ 的倍数）。[用 $G$ 在 $\mathcal B_i$ 上传递这一事实。也见第 1 节练习 8。]
    \item 令 $W=\{v\in V\mid \varphi(g)(v)=v\text{ 对所有 }g\in G\}$ 为 $V$ 中被 $G$ 所有元素逐点固定的子空间。推出 $\dim W=t$，即 $G$ 在 $\mathcal B$ 上的轨道个数。
  \end{enumerate}
  \item 证明如下结果（有时称为 Burnside 引理，虽然其起源归于 Frobenius）：设 $G$ 是 $S_n$ 的子群，对每个 $\sigma\in G$，令 $\operatorname{Fix}(\sigma)$ 表示 $\sigma$ 在 $\{1,
  \ldots,n\}$ 上的固定点个数。令 $t$ 是 $G$ 在 $\{1,
  \ldots,n\}$ 上的轨道个数。则
  \[
    t|G|=\sum_{g\in G}\operatorname{Fix}(g).
  \]
  [用前两个练习。]
  \item 令 $G$ 是有限集合 $\Omega$ 上一个非平凡传递置换群，令 $\psi$ 是由 $\Omega$ 得到的复线性表示所实现的特征标（参见第 1 节例 4），所以 $\psi(\sigma)$ 是 $\sigma$ 在 $\Omega$ 上的固定点个数。现在令 $G$ 作用在 $\Omega\times\Omega$ 上：
  \[
    g(\omega_1,
    \omega_2)=(g\cdot\omega_1,
    g\cdot\omega_2),
  \]
  并令 $\pi$ 是由这个作用得到的线性表示所实现的特征标。
  \begin{enumerate}[label=(\alph*)]
    \item 证明 $\pi=\psi^2$。
    \item 证明 $G$ 在 $\Omega\times\Omega$ 上的轨道个数由内积 $(\psi,
    \psi)$ 给出。[由前面练习，$\Omega\times\Omega$ 上的轨道个数等于 $(\pi,
    \chi_1)$，其中 $\chi_1$ 是主特征标。]
    \item 回忆，若 $G$ 在 $\Omega\times\Omega$ 上的作用恰有 $2$ 个轨道，则称 $G$ 在 $\Omega$ 上是二重传递的（它总是至少有 $2$ 个轨道，因为对角线 $\{(\omega,
    \omega)\mid \omega\in\Omega\}$ 是一个轨道）。证明若 $G$ 在 $\Omega$ 上二重传递，则 $\psi=\chi_1+\chi_2$，其中 $\chi_1$ 是主特征标，$\chi_2$ 是 $G$ 的一个非主不可约特征标。
    \item 令 $\Omega=\{1,2,
    \ldots,n\}$，令 $G=S_n$ 以自然方式作用在 $\Omega$ 上。说明相应线性表示的特征标分解为主特征标加上一个次数 $n-1$ 的不可约特征标。
  \end{enumerate}
  \item 令 $\psi$ 是群 $G$ 的任意二维表示的特征标，令 $x$ 是 $G$ 中阶为 $2$ 的元素。证明 $\psi(x)=2,0$ 或 $-2$。把这个结论推广到 $n$ 维表示。
  \item 令 $\chi$ 是 $G$ 的一个不可约特征标。证明对 $G$ 中心中的每个元素 $z$，都有 $\chi(z)=\epsilon\chi(1)$，其中 $\epsilon$ 是某个单位根。[用 Schur 引理。]
  \item 令 $\psi$ 是 $G$ 的某个表示 $\varphi$ 的特征标。证明对 $g\in G$，以下成立：
  \begin{enumerate}[label=(\alph*)]
    \item 若 $\psi(g)=\psi(1)$，则 $g\in\ker\varphi$；
    \item 若 $|\psi(g)|=\psi(1)$ 且 $\varphi$ 忠实，则 $g\in Z(G)$（这里 $|\psi(g)|$ 是复绝对值）。[用命题 14 证明的方法。]
  \end{enumerate}
  \item 设 $\varphi:G\to\GL(V)$ 是表示，$\chi:G\to\CC^\times$ 是次数 $1$ 表示。证明 $\chi\varphi:G\to\GL(V)$ 定义为 $(\chi\varphi)(g)=\chi(g)\varphi(g)$ 是表示（注意用复数 $\chi(g)$ 乘线性变换 $\varphi(g)$ 是良定义的）。证明 $\chi\varphi$ 不可约当且仅当 $\varphi$ 不可约。若 $\psi$ 是由 $\varphi$ 实现的特征标，证明 $\chi\psi$ 是由 $\chi\varphi$ 实现的特征标。推出任一不可约特征标与次数 $1$ 特征标的乘积仍是不可约特征标。
\end{enumerate}

接下来的若干练习研究代数共轭特征标的概念。这些练习可看作命题 14 的推广及其若干推论。特别地，我们得到一个群论刻画，说明何时群的所有不可约特征标都取值于 $\QQ$。

令 $F$ 为 $\CC$ 中所有在 $\QQ$ 上代数的元素所成的子域（代数数域）。因此 $F$ 是包含于 $\CC$ 的 $\QQ$ 的代数闭包，所有在 $\CC$ 上建立的结果在 $F$ 上不加改变地成立。
\begin{enumerate}[label=\arabic*.,resume]
  \item 注意由于 $F\subseteq\CC$，每个表示 $\varphi:G\to\GL_m(F)$ 也可看作复表示。证明若 $\varphi$ 是 $F$ 上不可约的表示，则把它看作较大域 $\CC$ 上表示时仍不可约（注意若 $F$ 不是代数闭的这不成立——参见上面练习 2(c)）。说明 $G$ 在 $F$ 上的不可约特征标集合与在 $\CC$ 上的不可约特征标集合相同（即它们作为 $G$ 上类函数完全是同一集合）。推出每个复表示都等价于一个在 $F$ 上的表示。[由于 $F$ 和 $\CC$ 都是特征为 $0$ 的代数闭域，其不可约特征标都由第一正交关系刻画。]
\end{enumerate}

令 $\varphi:G\to\GL_m(F)$ 是任意表示，特征标为 $\psi$。令 $\QQ(\varphi)$ 表示由所有矩阵 $\varphi(g)$（$g\in G$）的所有项生成的 $F$ 的子域。
\begin{enumerate}[label=\arabic*.,resume]
  \item 证明 $\QQ(\varphi)$ 是 $\QQ$ 的有限扩张。
\end{enumerate}

现在令 $K$ 是包含 $\QQ(\varphi)$ 的 $\QQ$ 的任意 Galois 扩张，令 $\sigma\in\Gal(K/\QQ)$。事实上，由于 $K$ 的每个自同构都延拓为 $F$ 的自同构，我们可假设 $\sigma$ 是 $F$ 的任意自同构。定义映射 $\varphi^\sigma:G\to\GL_n(F)$，令 $\varphi^\sigma(g)$ 是对矩阵 $\varphi(g)$ 的每个项施加域自同构 $\sigma$ 得到的 $n\times n$ 矩阵。
\begin{enumerate}[label=\arabic*.,resume]
  \item 证明 $\varphi^\sigma$ 是表示。还证明 $\varphi^\sigma$ 的特征标为 $\psi^\sigma$，其中 $\psi^\sigma(g)=\sigma(\psi(g))$。
  \item 证明 $\varphi$ 不可约当且仅当 $\varphi^\sigma$ 不可约。
\end{enumerate}
表示 $\varphi^\sigma$（或特征标 $\psi^\sigma$）称为 $\varphi$（或 $\psi$）在 $\sigma$ 下的\emph{代数共轭}；若存在 $F$ 的某个自同构 $\sigma$ 使得 $\varphi_1^\sigma=\varphi_2$（或 $\psi_1^\sigma=\psi_2$），则称两个表示 $\varphi_1,
\varphi_2$（或特征标 $\psi_1,
\psi_2$）代数共轭。这里需要注意这个标准记号，因为指数记号通常表示右作用，而 $F$ 的自同构在表示上是左作用：$\varphi^{\sigma\tau}=(\varphi^\tau)^\sigma$。

令 $\QQ(\psi)$ 是由所有 $\psi(g)$（$g\in G$）生成的 $F$ 的子域。令 $|G|=n$，令 $\epsilon$ 是 $F$ 中一个本原 $n$ 次单位根。
\begin{enumerate}[label=\arabic*.,resume]
  \item 证明 $\QQ(\psi)\subseteq\QQ(\epsilon)$。推出 $\QQ(\psi)$ 是 $\QQ$ 的 Galois 扩张，且其 Galois 群是 Abel 群。[见命题 14。]
\end{enumerate}

回忆第 14.5 节，$\Gal(\QQ(\epsilon)/\QQ)\cong(\ZZ/n\ZZ)^\times$，其 Galois 自同构由生成元 $\epsilon$ 上的作用给出：$\sigma_a:\epsilon\mapsto\epsilon^a$，其中 $a$ 是与 $n$ 互素的整数。
\begin{enumerate}[label=\arabic*.,resume]
  \item 证明：若 $\sigma_a\in\Gal(\QQ(\epsilon)/\QQ)$ 是上面定义的域自同构，则对所有 $g\in G$，有 $\psi^{\sigma_a}(g)=\psi(g^a)$。[用命题 14 的方法。]
  \item 证明：若 $g\in G$ 与 $g^a$ 对所有与 $n$ 互素的整数 $a$ 都共轭，则对 $G$ 的每个特征标 $\psi$，都有 $\psi(g)\in\QQ$。[用前一练习与 $\QQ$ 是所有 $\sigma_a$ 的不动域这一事实。]
  \item 证明：元素 $g\in G$ 与 $g^a$ 对所有与 $|G|$ 互素的整数 $a$ 共轭，当且仅当 $g$ 与 $g^{a'}$ 对所有与 $|g|$ 互素的整数 $a'$ 共轭。
  \item 说明对任意正整数 $n$，对称群 $S_n$ 的每个特征标都是有理值的（即对所有 $g\in S_n$ 与 $S_n$ 的所有特征标 $\psi$，都有 $\psi(g)\in\QQ$）。
\end{enumerate}

接下来两个练习建立练习 20 的逆命题。
\begin{enumerate}[label=\arabic*.,resume]
  \item 证明：群 $G$ 中元素 $x$ 与 $y$ 共轭，当且仅当对 $G$ 的所有不可约特征标 $\chi$，都有 $\chi(x)=\chi(y)$。
  \item 设 $g\in G$，并假设 $G$ 的每个不可约特征标在 $g$ 上都取有理值。证明对每个与 $|G|$ 互素的整数 $a$，$g$ 与 $g^a$ 共轭。[若对某个与 $|G|$ 互素的 $a$，$g$ 与 $g^a$ 不共轭，则由前一练习存在不可约特征标 $\chi$ 使得 $\chi(g)\neq\chi(g^a)$。由假设 $\chi(g)\in\QQ$ 推出矛盾。]
  \item 描述 $n$ 阶循环群的哪些不可约特征标彼此代数共轭。
  \item 证明 $Q_8$ 和 $D_8$ 的每个不可约特征标都是有理值的。证明 $D_{10}$ 有一个不可约特征标不是有理值的。
  \item 令 $G=H\times K$，令 $\varphi:H\to\GL(V)$ 是 $H$ 的不可约表示，特征标为 $\chi$。则 $G\to H\to\GL(V)$ 给出 $G$ 的不可约表示，其中第一个映射是自然投影 $\pi_H$；该表示的特征标 $\widetilde\chi$ 满足
  \[
    \widetilde\chi((h,k))=\chi(h).
  \]
  同样，$K$ 的任一不可约特征标 $\psi$ 给出 $G$ 的不可约特征标 $\widetilde\psi$，其中
  \[
    \widetilde\psi((h,k))=\psi(k).
  \]
  \begin{enumerate}[label=(\alph*)]
    \item 证明乘积 $\widetilde\chi\widetilde\psi$ 是 $G$ 的不可约特征标。[证明其范数为 $1$。]
    \item 证明 $G$ 的每个不可约特征标都由直因子的不可约特征标的这种乘积得到。[用定理 10 的 (3) 或 (4)。]
  \end{enumerate}
  \item （$\GL_2(\QQ)$ 的有限子群）设 $G$ 是 $\GL_2(\QQ)$ 的有限子群。
  \begin{enumerate}[label=(\alph*)]
    \item 说明 $\GL_2(\QQ)$ 不含阶为 $n$ 的元素，其中 $n=5,7$ 或 $n\geq9$。推出 $|G|=2^a3^b$。[用有理标准形。]
    \item 说明 Klein 四元群是 $\GL_2(\QQ)$ 中唯一的非循环 Abel 子群。由此和 (a) 推出 $|G|\mid24$。
    \item 说明 $\GL_2(\QQ)$ 的有限子群只有阶为 $1,2,3,4,6$ 的循环群、Klein 四元群，以及阶为 $6,8,12$ 的二面体群。[用第 4.5 节中小阶群分类与第 1 节练习 10 把 $G$ 限定在此列表中。反过来说明列表中每个群都有一个二维忠实有理表示。]
  \end{enumerate}
\end{enumerate}


% % ============================================
%  chapter19.tex — 第19章：特征标理论的例子与应用
%  根据 894-918.pdf 翻译整理，保持原文结构、证明与练习
% ============================================

\providecommand{\Ind}{\operatorname{Ind}}
\chapter{特征标理论的例子与应用}

\section{小阶群的特征标}

有限群的特征标表是按如下方式排成的特征标值的表：在最上面一行列出 $r$ 个共轭类的代表元，在第一列列出不可约特征标。表中第 $\chi_i$ 行、第 $g_j$ 列的项为 $\chi_i(g_j)$。有限群的特征标表在行和列的置换意义下唯一。通常把主特征标放在第一行，把恒等元放在第一列，并按次数递增的顺序列出特征标。在我们的例子中，我们将在每个类的下面列出该共轭类的大小，于是整张表将有 $r+2$ 行和 $r+1$ 列（虽然严格地说，特征标表是特征标值组成的 $r\times r$ 矩阵）。这样可以方便地用第一正交关系检查“行的正交性”：若这些类由 $g_1,\ldots,g_r$ 代表，大小分别为 $d_1,\ldots,d_r$，则
\[
  (\chi_i,\chi_j)=\frac{1}{|G|}\sum_{k=1}^r d_k\chi_i(g_k)\overline{\chi_j(g_k)}.
\]
第二正交关系说，特征标表任意两个不同列的 Hermite 内积为零（即它给出了“列的正交性”）。

许多特征标表可见 Conway、Curtis、Norton、Parker 与 Wilson 的 \emph{Atlas of Finite Groups}，Clarendon Press, 1985。这些表包括 Monster 单群 $M$ 的特征标表。群 $M$ 有 $194$ 个不可约特征标。$M$ 的一个非主不可约特征标的最小次数为 $196883$，最大次数约为 $2\times 10^{26}$。尽管如此，仍然可以计算这些特征标在 $M$ 的所有共轭类上的取值。

作为特征标表的第一个例子，令 $G=\langle x\rangle$ 为 $2$ 阶循环群。于是 $G$ 有 $2$ 个共轭类和两个不可约特征标：
\[
\begin{array}{c|cc}
\text{类：} & 1 & x\\
\text{大小：} & 1 & 1\\\hline
\chi_1 & 1 & 1\\
\chi_2 & 1 & -1
\end{array}
\]
\[
\text{表 }Z_2\text{ 的特征标表}
\]
这个 Abel 群的特征标与表示相同，而任意 Abel 群的不可约表示已在第 18.2 节末尾的例 1 中描述。

类似地，若 $G=\langle x\rangle$ 是 $3$ 阶循环群，且 $\zeta$ 是一个固定的本原三次单位根（于是 $\zeta^2=\overline{\zeta}$），则 $G$ 的特征标表如下：
\[
\begin{array}{c|ccc}
\text{类：} & 1 & x & x^2\\
\text{大小：} & 1 & 1 & 1\\\hline
\chi_1 & 1 & 1 & 1\\
\chi_2 & 1 & \zeta & \zeta^2\\
\chi_3 & 1 & \zeta^2 & \zeta
\end{array}
\]
\[
\text{表 }Z_3\text{ 的特征标表}
\]

接着我们构造 $S_3$ 的特征标表。回忆第 18.2 节例 2，$S_3$ 有 $3$ 个不可约特征标，其值已在该例和第 18.3 节末尾的例 1 中描述。
\[
\begin{array}{c|ccc}
\text{类：} & 1 & (1\ 2) & (1\ 2\ 3)\\
\text{大小：} & 1 & 3 & 2\\\hline
\chi_1 & 1 & 1 & 1\\
\chi_2 & 1 & -1 & 1\\
\chi_3 & 2 & 0 & -1
\end{array}
\]
\[
\text{表 }S_3\text{ 的特征标表}
\]

接下来考虑 $D_8$，采用第 18.3 节例 3 的记号。由推论 11，$D_8$ 有四个次数为 $1$ 的特征标。此外，在例 3 中我们构造了一个不可约的次数 $2$ 表示。由于这些表示次数平方和为 $8$，这便列出了全部不可约表示（或者，由于有 $5$ 个共轭类，也知道有 $5$ 个不可约表示）。若用上横线表示到交换子商群（它是 Klein 四元群）的像，则 $\bar r=\bar r^2$。次数 $1$ 的表示（也就是它们的特征标）通过把生成元 $s$ 与 $r$ 送到 $\pm1$ 来计算（而乘积类被映到相应取值的乘积）。次数 $2$ 的不可约表示的矩阵已在第 18.3 节例 3 中算出，这个表示的特征标可直接从这些矩阵读出。因此 $D_8$ 的特征标表如下：
\[
\begin{array}{c|ccccc}
\text{类：} & 1 & r^2 & s & r & sr\\
\text{大小：} & 1 & 1 & 2 & 2 & 2\\\hline
\chi_1 & 1 & 1 & 1 & 1 & 1\\
\chi_2 & 1 & 1 & -1 & 1 & -1\\
\chi_3 & 1 & 1 & 1 & -1 & -1\\
\chi_4 & 1 & 1 & -1 & -1 & 1\\
\chi_5 & 2 & -2 & 0 & 0 & 0
\end{array}
\]
\[
\text{表 }D_8\text{ 的特征标表}
\]

现在计算 $8$ 阶四元数群的特征标表。我们使用通常的表示
\[
  Q_8=\langle i,j\mid i^4=1,\, i^2=j^2,\, i^{-1}ji=j^{-1}\rangle,
\]
并令 $k=ij$ 且 $i^2=-1$。$Q_8$ 的共轭类由 $1,-1,i,j,k$ 代表，大小分别为 $1,1,2,2,2$。由于 $Q_8$ 的交换子商群是 Klein 四元群，所以有四个次数为 $1$ 的特征标。剩下的一个不可约特征标必须次数为 $2$，以使次数平方和为 $8$。设 $\chi_5$ 是 $Q_8$ 的次数 $2$ 不可约特征标。可以检查第 18.1 节第二组例子中例 7 显式描述的表示 $\varphi:Q_8\to\GL_2(\CC)$ 实现 $\chi_5$；但我们说明如何仅借助正交关系得到 $\chi_5$ 的值，而不需要知道这些显式矩阵。若 $\varphi$ 是次数 $2$ 的不可约表示，则由 Schur 引理（参见第 18.1 节练习 18），$\varphi(-1)$ 是一个 $2\times2$ 标量矩阵；又由于 $-1$ 在 $Q_8$ 中有阶 $2$，所以它是 $\pm$ 单位矩阵。因此 $\chi_5(-1)=\pm2$。令 $\chi_5(i)=a$，$\chi_5(j)=b$，$\chi_5(k)=c$。正交关系给出
\[
  1=(\chi_5,\chi_5)=\frac18\bigl(2^2+(\pm2)^2+2a\bar a+2b\bar b+2c\bar c\bigr).
\]
由于 $a\bar a,b\bar b,c\bar c$ 都是非负实数，它们必全为零。此外，因为 $\chi_5$ 与主特征标正交，得到
\[
  0=(\chi_1,\chi_5)=\frac18(2+(\pm2)+0+0+0),
\]
故 $\chi_5(-1)=-2$。$Q_8$ 的完整特征标表如下：
\[
\begin{array}{c|ccccc}
\text{类：} & 1 & -1 & i & j & k\\
\text{大小：} & 1 & 1 & 2 & 2 & 2\\\hline
\chi_1 & 1 & 1 & 1 & 1 & 1\\
\chi_2 & 1 & 1 & -1 & 1 & -1\\
\chi_3 & 1 & 1 & 1 & -1 & -1\\
\chi_4 & 1 & 1 & -1 & -1 & 1\\
\chi_5 & 2 & -2 & 0 & 0 & 0
\end{array}
\]
\[
\text{表 }Q_8\text{ 的特征标表}
\]
注意 $D_8$ 与 $Q_8$ 有相同的特征标表，因此
\[
  \text{非同构的群可能有相同的特征标表。}
\]
还要注意，$Q_8$ 的次数 $2$ 表示的值也可通过对特征标表的每一列应用第二正交关系而容易算出。我们把这个检验留作练习。还要注意，尽管 $D_8$ 与 $Q_8$ 的次数 $2$ 不可约特征标有相同的（实数）值，$D_8$ 的次数 $2$ 表示可由实矩阵实现，而可以证明 $Q_8$ 在 $\RR$ 上没有忠实的二维表示（参见第 18.1 节练习 10）。

作为下一个例子，我们构造 $S_4$ 的特征标表。$S_4$ 的共轭类由 $1,(1\ 2),(1\ 2\ 3),(1\ 2\ 3\ 4),(1\ 2)(3\ 4)$ 代表，大小分别为 $1,6,8,6,3$。由于 $S_4'=A_4$，有两个次数为 $1$ 的特征标：主特征标和取值为置换符号的特征标。

为了得到次数 $2$ 的不可约特征标，令 $V$ 为由 $(1\ 2)(3\ 4)$ 与 $(1\ 3)(2\ 4)$ 生成的 $4$ 阶正规子群。$S_4/V\cong S_3$ 的任意表示 $\varphi$，与自然投影 $S_4\to S_4/V$ 复合，给出 $S_4$ 的一个表示；若前者不可约，则后者也不可约。令 $\varphi$ 为该投影与 $S_3$ 的不可约二维表示的复合，并令 $\chi_3$ 为其特征标。$1$ 和 $(1\ 2)(3\ 4)$ 的类映到 $S_3$ 商中的恒等元，$(1\ 2)$ 和 $(1\ 2\ 3\ 4)$ 映到对换，$(1\ 2\ 3)$ 映到 $3$-循环。因此 $\chi_3$ 的值可直接从 $S_3$ 表中次数 $2$ 特征标的值读出。

由于 $S_4$ 有 $5$ 个不可约特征标，而次数平方和为 $24$，剩余的两个不可约特征标都必须次数为 $3$。在第 18.3 节例 2 中已计算出其中一个，称为 $\chi_4$。回忆
\[
  \chi_4(\sigma)=(\sigma\text{ 的不动点个数})-1.
\]
剩余的不可约特征标 $\chi_5$ 是 $\chi_4\chi_2$。可利用第 18.3 节命题 17 或练习 13 看出这个乘积确实是一个特征标。第一正交关系验证它不可约。
\[
\begin{array}{c|ccccc}
\text{类：} & 1 & (1\ 2) & (1\ 2\ 3) & (1\ 2\ 3\ 4) & (1\ 2)(3\ 4)\\
\text{大小：} & 1 & 6 & 8 & 6 & 3\\\hline
\chi_1 & 1 & 1 & 1 & 1 & 1\\
\chi_2 & 1 & -1 & 1 & -1 & 1\\
\chi_3 & 2 & 0 & -1 & 0 & 2\\
\chi_4 & 3 & 1 & 0 & -1 & -1\\
\chi_5 & 3 & -1 & 0 & 1 & -1
\end{array}
\]
\[
\text{表 }S_4\text{ 的特征标表}
\]

由 $S_4$ 的特征标表可以很容易地计算 $A_4$ 的特征标表。注意 $A_4$ 有 $4$ 个共轭类。此外，$|A_4:A_4'|=3$，所以 $A_4$ 有三个次数为 $1$ 的特征标，每个次数 $1$ 表示的核都含有 $V=A_4'$。剩下的不可约特征标必须次数为 $3$。直接在 $A_4$ 中应用正交关系，可以检验 $S_4$ 的特征标 $\chi_4$ 限制到 $A_4$（等于 $\chi_5|_{A_4}$）是不可约的。这个不可约性检验确实是必要的，因为一个群的不可约表示限制到子群时未必仍然不可约（例如，$S_3$ 的不可约次数 $2$ 表示限制到任意真子群时必变为可约，因为这些真子群都是 Abel 群）。$A_4$ 的特征标表如下：
\[
\begin{array}{c|cccc}
\text{类：} & 1 & (1\ 2)(3\ 4) & (1\ 2\ 3) & (1\ 3\ 2)\\
\text{大小：} & 1 & 3 & 4 & 4\\\hline
\chi_1 & 1 & 1 & 1 & 1\\
\chi_2 & 1 & 1 & \zeta & \zeta^2\\
\chi_3 & 1 & 1 & \zeta^2 & \zeta\\
\chi_4 & 3 & -1 & 0 & 0
\end{array}
\]
\[
\text{表 }A_4\text{ 的特征标表}
\]
其中 $\zeta$ 是 $\CC$ 中的一个本原三次单位根。

作为最后一个例子，我们构造如下 $S_5$ 的特征标表：
{\scriptsize
\[
\begin{array}{c|ccccccc}
\text{类：} & 1 & (1\ 2) & (1\ 2\ 3) & (1\ 2\ 3\ 4) & (1\ 2\ 3\ 4\ 5) & (1\ 2)(3\ 4) & (1\ 2)(3\ 4\ 5)\\
\text{大小：} & 1 & 10 & 20 & 30 & 24 & 15 & 20\\\hline
\chi_1 & 1 & 1 & 1 & 1 & 1 & 1 & 1\\
\chi_2 & 1 & -1 & 1 & -1 & 1 & 1 & -1\\
\chi_3 & 4 & 2 & 1 & 0 & -1 & 0 & -1\\
\chi_4 & 4 & -2 & 1 & 0 & -1 & 0 & 1\\
\chi_5 & 5 & -1 & -1 & 1 & 0 & 1 & -1\\
\chi_6 & 5 & 1 & -1 & -1 & 0 & 1 & 1\\
\chi_7 & 6 & 0 & 0 & 0 & 1 & -2 & 0
\end{array}
\]
\[
\text{表 }S_5\text{ 的特征标表}
\]
}
共轭类及其大小已在第 4.3 节中计算。由于 $|S_5:S_5'|=2$，有两个次数为 $1$ 的特征标：主特征标和“符号”特征标。

$S_5$ 在 $5$ 个点上的自然置换给出一个次数为 $5$ 的置换特征标。与 $S_4$ 和 $S_3$ 的情形一样，正交关系表明它的范数平方为 $2$，且包含主特征标。因此 $\chi_3$ 是置换特征标减去主特征标（并且像较小的对称群一样，$\chi_3(\sigma)$ 等于 $\sigma$ 的不动点个数减 $1$）。按与 $S_4$ 中相同的论证，得到 $\chi_4=\chi_3\chi_2$ 也是一个不可约特征标。

为了得到 $\chi_5$，回忆 $S_5$ 有六个 Sylow $5$-子群。它在这些 Sylow $5$-子群上按共轭作用给出一个忠实的次数 $6$ 置换表示。若 $\psi$ 是相关线性表示的特征标，则由于 $\sigma\in S_5$ 固定一个 Sylow $5$-子群当且仅当它正规化该子群，我们有
\[
  \psi(\sigma)=\sigma\text{ 所正规化的 Sylow }5\text{-子群的个数。}
\]
$S_5$ 中 Sylow $5$-子群 $\langle(1\ 2\ 3\ 4\ 5)\rangle$ 的正规化子是
\[
  \langle (1\ 2\ 3\ 4\ 5),(2\ 3\ 5\ 4)\rangle,
\]
而 Sylow $5$-子群的所有正规化子都在 $S_5$ 中共轭于这个群。这个正规化子只含有恒等元、$5$-循环、$4$-循环以及两个不交对换的乘积。没有其他循环型正规化任何 Sylow $5$-子群，因此在任何其他类上，$\psi$ 为零。为了计算 $\psi$ 在剩余三个非恒等类上的值，注意（在 $S_6$ 中观察）任意一个在 $6$ 个点上的忠实作用都有如下事实：$5$ 阶元必须是一个 $5$-循环（因此固定 $1$ 个点）；固定一个点的 $4$ 阶元必须是一个 $4$-循环（因此固定 $2$ 个点）；作为 $4$ 阶元平方的 $2$ 阶元也恰好固定 $2$ 个点。这给出了 $\psi$ 的所有取值。现在直接计算表明
\[
  \|\psi\|^2=2\qquad\text{且}\qquad (\chi_1,\psi)=1.
\]
因此 $\chi_5=\psi-\chi_1$ 是次数 $5$ 的不可约特征标。由与 $\chi_4$ 相同的理论，得到 $\chi_6=\chi_5\chi_2$ 是另一个不可约特征标。

由于有 $7$ 个共轭类，还有一个剩余的不可约特征标，其次数为 $6$。它的值可从正则特征标 $\rho$ 的分解中立即得到（参见第 18.2 节例 3 和第 18.3 节例 4）：
\[
  \chi_7=\frac{\rho-\chi_1-\chi_2-4\chi_3-4\chi_4-5\chi_5-5\chi_6}{6}.
\]
由正交关系直接计算可检验 $\chi_7$ 不可约。注意，$\chi_7$ 的值是在没有显式给出具有这个特征标的表示的情况下计算出来的。

\subsection*{练习}
\begin{enumerate}[label=\arabic*.]
  \item 计算 $Z_2\times Z_2$、$Z_2\times Z_3$ 和 $Z_2\times Z_2\times Z_2$ 的特征标表。说明为什么 $Z_2\times Z_3$ 的表中含有本原 $6$ 次单位根。
  \item 计算 $D_{16}$ 的不可约特征标的次数。
  \item 计算 $A_5$ 的不可约特征标的次数。推出 $S_5$ 的次数 $6$ 不可约特征标在限制到 $A_5$ 时不是不可约的。［$A_5$ 的共轭类已在第 4.3 节中求出。］
  \item 利用本节中的特征标表，对 (a) 到 (d) 的每一部分，用第一正交关系把指定的置换特征标（参见第 18.3 节例 3）写成不可约特征标之和：
  \begin{enumerate}[label=(\alph*)]
    \item $S_3$ 的子群 $A_3$ 的置换特征标；
    \item $S_4$ 的子群 $\langle(1\ 2\ 3\ 4)\rangle$ 的置换特征标；
    \item $S_4$ 的子群 $V_4$ 的置换特征标；
    \item $S_5$ 的子群 $\langle(1\ 2\ 3),(1\ 2),(4\ 5)\rangle$ 的置换特征标（这个子群是 $S_5$ 的一个 Sylow $3$-子群的正规化子）。
  \end{enumerate}
  \item 假设对一个群的任意特征标 $\psi$，$\psi^2$ 也总是特征标（其中 $\psi^2(g)=(\psi(g))^2$）——这是第 18.3 节命题 17 的一个特殊情形。利用本节中的特征标表，对 (a) 到 (e) 的每一部分，写出指定特征标 $\chi$ 的平方 $\chi^2$ 的值，并用第一正交关系把 $\chi^2$ 写成不可约特征标之和：
  \begin{enumerate}[label=(\alph*)]
    \item $\chi=\chi_3$，$S_3$ 表中的次数 $2$ 特征标；
    \item $\chi=\chi_5$，$Q_8$ 表中的次数 $2$ 特征标；
    \item $\chi=\chi_5$，$S_4$ 表中的最后一个特征标；
    \item $\chi=\chi_4$，$S_5$ 表中的第二个次数 $4$ 特征标；
    \item $\chi=\chi_7$，$S_5$ 表中的最后一个特征标。
  \end{enumerate}
  \item 计算 $A_5$ 的特征标表。
  \item 证明 $S_6$ 有一个次数 $5$ 的不可约特征标。
  \item 计算 $D_{10}$ 的特征标表。（这个表含有非实项。）
  \item 计算 $D_{12}$ 的特征标表。
  \item 计算 $S_3\times S_3$ 的特征标表。
  \item 计算 $Z_3\times S_3$ 的特征标表。
  \item 计算 $Z_2\times S_4$ 的特征标表。
  \item 计算 $S_3\times S_4$ 的特征标表。
  \item 设 $n$ 是整数且 $n\geq3$。证明 $D_{2n}$ 的每个不可约特征标次数为 $1$ 或 $2$，并求出每种次数的不可约特征标数。［$D_{2n}$ 的共轭类已在第 4.3 节练习 31 和 32 中求出，它的交换子子群已在第 5.4 节中计算。］
  \item 证明特征标表是可逆矩阵。［使用正交关系。］
  \item 对 $A_5$ 和 $D_{10}$ 的每一个，描述哪些不可约特征标是代数共轭的（参见第 18.3 节的练习）。
  \item 设 $p$ 是任意素数，$P$ 是一个 $p^3$ 阶非 Abel 群（在同构意义下 $P$ 有两种选择；当 $p$ 为奇素数时，这些群已在第 5.5 节分类 $p^3$ 阶群时构造）。本练习确定 $P$ 的特征标表，并说明两个同构类型具有相同的特征标表（该论证包括本节已经处理的 $p=2$ 情形）。
  \begin{enumerate}[label=(\alph*)]
    \item 证明 $P$ 有 $p^2$ 个次数 $1$ 特征标。
    \item 证明 $P$ 有 $p-1$ 个次数 $p$ 不可约特征标，并且这些特征标与 $p^2$ 个次数 $1$ 特征标一起给出 $P$ 的全部不可约特征标。［使用第 18.2 节定理 10(3) 和定理 12。］
    \item 推出（不论同构类型如何）群 $P$ 有 $p^2+p-1$ 个共轭类，其中 $p$ 个大小为 $1$（即中心类），而 $p^2-1$ 个各自大小为 $p$。还推出大小为 $p$ 的类正是 $P$ 中心的非恒等陪集（即若 $x\in P-Z(P)$，则 $x$ 的共轭类是陪集 $xZ(P)$ 中的 $p$ 个元素）。
    \item 证明若 $\chi$ 是次数 $p$ 的不可约特征标，则实现 $\chi$ 的表示是忠实的。
    \item 固定 $P$ 的中心的一个生成元 $z$，并令 $\varepsilon$ 是 $\CC$ 中一个固定的本原 $p$ 次单位根。证明若 $\chi$ 是次数 $p$ 的不可约特征标，则 $\chi(z)=p\varepsilon^i$，其中某个 $i\in\{1,2,\ldots,p-1\}$。进一步证明对所有 $x\in P-Z(P)$ 都有 $\chi(x)=0$。（于是注意，次数 $p$ 特征标全都是代数共轭的。）［使用与构造 $Q_8$ 的特征标表时相同的推理。］
    \item 证明对每个 $i\in\{1,2,\ldots,p-1\}$，存在唯一的次数 $p$ 不可约特征标 $\chi_i$ 满足 $\chi_i(z)=p\varepsilon^i$。推出 $P$ 的特征标表被唯一决定，并描述它。［回忆第 6.1 节：不论同构类型如何，都有 $P'=Z(P)$ 且 $P/P'\cong Z_p\times Z_p$。由此可以写出次数 $1$ 的特征标。第 (e) 部分描述了次数 $p$ 特征标。］
  \end{enumerate}
\end{enumerate}

\section{Burnside 定理与 Hall 定理}

在本节中，我们给出特征标理论的一个“理论性”应用：Burnside 的 $p^a q^b$ 定理。我们还证明 Philip Hall 对有限可解群的刻画，其证明是群论证明，但在归纳的第一步依赖 Burnside 定理。

\subsection*{Burnside 定理}

下面的结果由 Burnside 于 1904 年证明。虽然关于它的纯群论证明近来已经被发现（见 B. Huppert 与 N. Blackburn 的 \emph{Finite Groups III}, Springer-Verlag, 1982, 定理 2.8），但这里给出的 Burnside 原始证明，在具备当前表示论知识之后，非常易懂、优美且相当简短。

\noindent\textbf{定理 1.} \textbf{(Burnside)} 对素数 $p$ 与 $q$，每个 $p^a q^b$ 阶群都是可解的。

在开始证明 Burnside 定理本身之前，我们先建立一些一般性的结果。这些预备命题的一个简单推论是：任意有限群的不可约特征标的次数整除群的阶。直接导向 Burnside 定理证明的具体结果出现在引理 6 和引理 7 中。很容易推出，Burnside 定理的一个最小阶反例必是非 Abel 单群，而正是这两个特征标论引理通过证明正规子群的存在而给出矛盾。

我们首先回忆第 15.3 节中代数整数的定义。

\noindent\textbf{定义.} 元素 $\alpha\in\CC$ 称为一个\emph{代数整数}，如果它是某个系数在 $\ZZ$ 中的首一多项式的根。

证明 Burnside 定理所需的基本结果是：

\noindent\textbf{命题 2.} 设 $\alpha\in\CC$。
\begin{enumerate}[label=(\arabic*)]
  \item 下列条件等价：
  \begin{enumerate}[label=(\roman*)]
    \item $\alpha$ 是代数整数；
    \item $\alpha$ 在 $\QQ$ 上代数，且 $\alpha$ 在 $\QQ$ 上的极小多项式有整数系数；
    \item $\ZZ[\alpha]$ 是一个有限生成 $\ZZ$-模（其中 $\ZZ[\alpha]$ 是 $\CC$ 中由 $\ZZ$ 和 $\alpha$ 生成的子环，即所有 $\alpha$ 的非负幂的 $\ZZ$-线性组合所成的环）。
  \end{enumerate}
  \item $\CC$ 中的代数整数构成一个环，而 $\QQ$ 中的代数整数就是 $\ZZ$ 的元素。
\end{enumerate}

\noindent\textbf{证明.} 这些已在第 15.3 节中建立。（第 15.3 节中关于整扩张和代数整数性质的部分可以独立于第 15 章其余部分阅读。）

\noindent\textbf{推论 3.} 对有限群 $G$ 的每个特征标 $\psi$，对所有 $x\in G$，$\psi(x)$ 都是代数整数。

\noindent\textbf{证明.} 由第 18.3 节命题 14，$\psi(x)$ 是若干单位根之和。每个单位根都是代数整数，所以结果立即由命题 2(2) 得到。

在开始主证明之前，我们还需要一些特征标论的预备引理。对任意有限群 $G$ 采用如下记号：$\chi_1,\ldots,\chi_r$ 是 $G$ 的不同不可约（复）特征标，$\mathcal K_1,\ldots,\mathcal K_r$ 是 $G$ 的共轭类，$\varphi_i$ 是特征标为 $\chi_i$ 的不可约矩阵表示。

\noindent\textbf{命题 4.} 对每个 $i$，在 $\{\mathcal K_1,\ldots,\mathcal K_r\}$ 上定义复值函数 $\omega_i$ 为
\[
  \omega_i(\mathcal K_j)=\frac{|\mathcal K_j|\chi_i(g)}{\chi_i(1)},
\]
其中 $g$ 是 $\mathcal K_j$ 的任意元素。则对所有 $i,j$，$\omega_i(\mathcal K_j)$ 都是代数整数。

\noindent\textbf{证明.} 我们首先证明，若 $I$ 是单位矩阵，则
\[
  \sum_{g\in\mathcal K_j}\varphi_i(g)=\omega_i(\mathcal K_j)I. \tag{19.1}
\]
为此，令 $X$ 为 (1) 左端。如第 18.2 节中所见，每个 $x\in G$ 通过共轭作用置换 $\mathcal K_j$ 的元素，因此 $X$ 与所有 $\varphi_i(g)$ 交换。由 Schur 引理（第 18.1 节练习 18），$X$ 是一个标量矩阵：
\[
  X=\alpha I\qquad\text{对某个 }\alpha\in\CC.
\]
只需证明 $\alpha=\omega_i(\mathcal K_j)$。但
\[
  \operatorname{tr}X=\sum_{g\in\mathcal K_j}\chi_i(g)=|\mathcal K_j|\chi_i(g).
\]
因此 $\alpha\chi_i(1)=\operatorname{tr}X=|\mathcal K_j|\chi_i(g)$，这正是建立 (1) 所需的等式。

现在固定 $g\in\mathcal K_s$，定义 $a_{ijs}$ 为有序对 $g_i,g_j$ 的个数，其中 $g_i\in\mathcal K_i$，$g_j\in\mathcal K_j$ 且 $g_i g_j=g$。注意 $a_{ijs}$ 是整数。它与 $g\in\mathcal K_s$ 的选择无关，因为若 $x^{-1}gx$ 是 $g$ 的共轭元，则每个乘积为 $g$ 的有序对 $g_i,g_j$ 对应到乘积为 $x^{-1}gx$ 的有序对 $x^{-1}g_i x,x^{-1}g_jx$（反之亦然）。

接着我们证明对所有 $i,j,t\in\{1,\ldots,r\}$，
\[
  \omega_t(\mathcal K_i)\omega_t(\mathcal K_j)=\sum_{s=1}^r a_{ijs}\omega_t(\mathcal K_s). \tag{19.2}
\]
为此注意，由 (1)，(2) 左端是下列等式左端标量矩阵的对角项：
\[
\begin{aligned}
  \left(\sum_{g\in\mathcal K_i}\varphi_t(g)\right)
  \left(\sum_{h\in\mathcal K_j}\varphi_t(h)\right)
  &=\sum_{g\in\mathcal K_i}\sum_{h\in\mathcal K_j}\varphi_t(gh)\\
  &=\sum_{s=1}^r\sum_{g\in\mathcal K_s} a_{ijs}\varphi_t(g)\\
  &=\sum_{s=1}^r a_{ijs}\sum_{g\in\mathcal K_s}\varphi_t(g)\\
  &=\sum_{s=1}^r a_{ijs}\omega_t(\mathcal K_s)I.
\end{aligned}
\]
比较这些标量矩阵的项得到 (2)。

现在 (2) 蕴含：对每个 $t\in\{1,
\ldots,r\}$，由 $\ZZ$ 和 $\omega_t(\mathcal K_1),\ldots,\omega_t(\mathcal K_r)$ 生成的 $\CC$ 的子环是一个有限生成 $\ZZ$-模（作为 $\ZZ$-模由 $1,\omega_t(\mathcal K_1),\ldots,\omega_t(\mathcal K_r)$ 生成）。由于 $\ZZ$ 是主理想整环，子模 $\ZZ[\omega_t(\mathcal K_i)]$ 也是有限生成 $\ZZ$-模，因此由命题 2，$\omega_t(\mathcal K_i)$ 是代数整数。证明完毕。

\noindent\textbf{推论 5.} 有限群 $G$ 的每个复不可约表示的次数都整除 $G$ 的阶，即对 $i=1,2,\ldots,r$，有 $\chi_i(1)\mid |G|$。

\noindent\textbf{证明.} 在命题 4 的记号下，取 $g_j\in\mathcal K_j$，则
\[
\begin{aligned}
  \frac{|G|}{\chi_i(1)}
  &=\frac{|G|}{\chi_i(1)}(\chi_i,\chi_i)\\
  &=\sum_{j=1}^r\frac{|\mathcal K_j|\chi_i(g_j)\overline{\chi_i(g_j)}}{\chi_i(1)}\\
  &=\sum_{j=1}^r \omega_i(\mathcal K_j)\overline{\chi_i(g_j)}.
\end{aligned}
\]
右端是代数整数，左端是有理数，因而是整数。这证明了推论。

下面两个引理直接导向 Burnside 定理。

\noindent\textbf{引理 6.} 若群 $G$ 有一个共轭类 $\mathcal K$ 和一个不可约矩阵表示 $\varphi$，其特征标为 $\chi$，满足 $(|\mathcal K|,\chi(1))=1$，则对 $g\in\mathcal K$，要么 $\chi(g)=0$，要么 $\varphi(g)$ 是标量矩阵。

\noindent\textbf{证明.} 由假设，存在 $s,t\in\ZZ$ 使得 $s|\mathcal K|+t\chi(1)=1$。因此
\[
  s|\mathcal K|\chi(g)+t\chi(1)\chi(g)=\chi(g).
\]
两边同除以 $\chi(1)$，并注意由推论 3 与命题 4，$\chi(g)$ 和 $|\mathcal K|\chi(g)/\chi(1)$ 都是代数整数，因此 $\chi(g)/\chi(1)$ 也是代数整数。令
\[
  a_1=\frac{\chi(g)}{\chi(1)},
\]
令 $a_1,a_2,\ldots,a_n$ 为它在 $\QQ$ 上的所有代数共轭（即 $a_1$ 在 $\QQ$ 上极小多项式的根）。由于 $a_1$ 是 $\chi(1)$ 个单位根之和除以整数 $\chi(1)$，每个 $a_i$ 也都是 $\chi(1)$ 个单位根之和除以 $\chi(1)$。因此对所有 $i$，$a_i$ 的复绝对值都不超过 $1$。现在 $b=\prod_{i=1}^n a_i\in\QQ$，且 $b$ 是代数整数（$\pm b$ 是 $a_1$ 的不可约多项式的常数项），故 $b\in\ZZ$。但
\[
  |b|=\prod_{i=1}^n |a_i|\leq 1,
\]
所以 $b=0,\pm1$。由于所有 $a_i$ 彼此共轭，$b=0\Longleftrightarrow a_1=0\Longleftrightarrow \chi(g)=0$。又 $b=\pm1\Longleftrightarrow |a_i|=1$ 对所有 $i$ 成立。因此，要么 $\chi(g)=0$，要么 $|\chi(g)|=\chi(1)$。前一种情形引理已证，故假设 $|\chi(g)|=\chi(1)$。

取一个与 $\varphi$ 等价的矩阵表示 $\varphi_1$，使 $\varphi_1(g)$ 为对角矩阵：
\[
  \varphi_1(g)=\begin{pmatrix}
  \varepsilon_1&&0\\
  &\ddots&\\
  0&&\varepsilon_n
  \end{pmatrix}.
\]
于是 $\chi(g)=\varepsilon_1+\cdots+\varepsilon_n$。由三角不等式，若某些 $\varepsilon_i\neq\varepsilon_j$，则 $|\varepsilon_1+\cdots+\varepsilon_n|<n=\chi(1)$。这与当前情形不符，所以必须有 $\varphi_1(g)=\varepsilon I$（其中 $\varepsilon=\varepsilon_i$ 对所有 $i$）。由于标量矩阵只相似于自身，$\varphi(g)=\varepsilon I$。证明完毕。

\noindent\textbf{引理 7.} 若 $G$ 的某个非恒等共轭类 $\mathcal K$ 的大小 $|\mathcal K|$ 是某个素数的幂，则 $G$ 不是非 Abel 单群。

\noindent\textbf{证明.} 反设 $G$ 是非 Abel 单群，并令 $|\mathcal K|=p^c$。取 $g\in\mathcal K$。若 $c=0$，则 $g\in Z(G)$，这与非 Abel 单群中心平凡矛盾。如前，令 $\chi_1,\ldots,\chi_r$ 为 $G$ 的所有不可约特征标，其中 $\chi_1$ 为主特征标，令 $\rho$ 为 $G$ 的正则特征标。把 $\rho$ 分解为不可约特征标，得到
\[
  0=\rho(g)=1+\sum_{i=2}^r \chi_i(1)\chi_i(g). \tag{19.3}
\]
若对每个满足 $\chi_j(g)\neq0$ 的 $j>1$ 都有 $p\mid\chi_j(1)$，则写 $\chi_j(1)=pd_j$。此时 (3) 变为
\[
  0=1+p\sum_j d_j\chi_j(g).
\]
于是 $\sum_j d_j\chi_j(g)=-1/p$ 是代数整数，矛盾。这证明存在某个 $j$，使得 $p\nmid\chi_j(1)$ 且 $\chi_j(g)\neq0$。若 $\varphi$ 是特征标为 $\chi_j$ 的表示，则 $\varphi$ 是忠实的（因为假设 $G$ 是单群）；并且由引理 6，$\varphi(g)$ 是标量矩阵。由于 $\varphi(g)$ 与所有矩阵交换，$\varphi(g)\in Z(\varphi(G))$。这迫使 $g\in Z(G)$，与 $G$ 是非 Abel 单群矛盾。引理证明完毕。

现在证明 Burnside 定理。令 $G$ 是某些素数 $p,q$ 的 $p^a q^b$ 阶群。若 $p=q$ 或者任一指数为 $0$，则 $G$ 是幂零的，因而可解。因此可假设并非如此。按归纳法，令 $G$ 是一个最小阶反例。若 $G$ 有真非平凡正规子群 $N$，则由归纳假设，$N$ 和 $G/N$ 都可解，因此 $G$ 也可解（参见第 3.4 节或命题 6.10）。故可假设 $G$ 是非 Abel 单群。令 $P\in\Syl_p(G)$。由第 4 章定理 8，存在 $g\in Z(P)$ 且 $g\neq1$。由于 $P\leq C_G(g)$，$g$ 的共轭类大小（等于 $|G:C_G(g)|$）与 $p$ 互素，即是 $q$ 的幂。这违反引理 7，从而完成 Burnside 定理的证明。

\subsection*{Philip Hall 定理}

回忆，有限群的一个子群称为 Hall 子群，如果它的阶与指数互素。对群 $G$ 的任意子群 $H$，若子群 $K$ 满足 $G=HK$ 且 $H\cap K=1$，则称 $K$ 是 $H$ 在 $G$ 中的一个补子群。

\noindent\textbf{定理 8.} \textbf{(P. Hall)} 设 $G$ 的阶为 $p_1^{a_1}p_2^{a_2}\cdots p_t^{a_t}$，其中 $p_1,\ldots,p_t$ 是不同素数。若对每个 $i\in\{1,\ldots,t\}$，都存在 $G$ 的子群 $H_i$ 使 $|G:H_i|=p_i^{a_i}$，则 $G$ 可解。

Hall 定理也可表述为：若对每个 $i\in\{1,\ldots,t\}$，$G$ 的一个 Sylow $p_i$-子群有补子群，则 $G$ 可解。Hall 定理的逆命题也成立——这正是第 6.1 节练习 33。

我们首先需要一些初等引理。

\noindent\textbf{引理 9.} 若 $G$ 是阶大于 $1$ 的可解群，则存在 $P\normal G$，其中 $P$ 是某个素数 $p$ 的非平凡 $p$-群。

\noindent\textbf{证明.} 这是第 6.1 节末尾关于可解群极小正规子群练习的特殊情形。也可直接看出：令 $P$ 为 $G$ 的导列中最后一个非平凡项 $G^{(n-1)}$ 的非平凡 Sylow 子群（其中 $G$ 的可解长度为 $n$）。此时 $G^{(n-1)}$ 是 Abel 群，所以 $P$ 是 $G^{(n-1)}$ 的特征子群，因而正规于 $G$。

\noindent\textbf{引理 10.} 设 $G$ 的阶为 $p_1^{a_1}p_2^{a_2}\cdots p_t^{a_t}$，其中 $p_1,\ldots,p_t$ 是不同素数。设 $G$ 有子群 $H,K$，使得对每个 $i\in\{1,\ldots,t\}$，或者 $p_i^{a_i}$ 整除 $|H|$，或者 $p_i^{a_i}$ 整除 $|K|$。则 $G=HK$ 且 $|H\cap K|=(|H|,|K|)$。

\noindent\textbf{证明.} 固定某个 $i\in\{1,\ldots,t\}$，先设 $p_i^{a_i}$ 整除 $H$ 的阶。由于 $HK$ 是若干右陪集 $H$ 的不交并，而每个这样的右陪集大小都等于 $|H|$，可知 $p_i^{a_i}$ 整除 $|HK|$。类似地，若 $p_i^{a_i}$ 整除 $|K|$，由于 $HK$ 是若干左陪集 $K$ 的不交并，同样有 $p_i^{a_i}$ 整除 $|HK|$。因此 $|G|\mid |HK|$，从而 $G=HK$。由于
\[
  |HK|=\frac{|H||K|}{|H\cap K|},
\]
得到 $|H\cap K|=(|H|,|K|)$。

现在开始证明 Hall 定理，对 $|G|$ 归纳。注意，若 $t=1$，则对任意群假设都平凡满足（$H_1=1$）；若 $t=2$，则由 Sylow 定理，对任意群假设也满足（$H_1$ 是 $G$ 的 Sylow $p_2$-子群，$H_2$ 是 $G$ 的 Sylow $p_1$-子群）。若 $t=1$，则 $G$ 幂零，从而可解；若 $t=2$，则由 Burnside 定理，$G$ 可解。因此假设 $t\geq3$。

固定 $i$。由前一引理，对所有 $j\in\{1,\ldots,t\}-\{i\}$，
\[
  |H_i:H_i\cap H_j|=p_j^{a_j}.
\]
因此 $H_i$ 的每个 Sylow $p_j$-子群在 $H_i$ 中都有补子群 $H_j\cap H_i$。由归纳假设，$H_i$ 可解。

由引理 9，可选择 $P\normal H_1$，其中 $|P|=p_i^b>1$，对某个 $i>1$。由于 $t\geq3$，存在指标 $j\in\{1,\ldots,t\}-\{1,i\}$。由引理 10，
\[
  |H_1\cap H_j|=p_2^{a_2}\cdots p_{j-1}^{a_{j-1}}p_{j+1}^{a_{j+1}}\cdots p_t^{a_t}.
\]
因此 $H_1\cap H_j$ 含有 $H_1$ 的一个 Sylow $p_i$-子群。由于 $P$ 是 $H_1$ 的正规 $p_i$-子群，$P$ 包含在 $H_1$ 的每个 Sylow $p_i$-子群中，所以 $P\leq H_1\cap H_j$。由引理 10，$G=H_1H_j$，故每个 $g\in G$ 可写为 $g=h_1h_j$，其中 $h_1\in H_1$，$h_j\in H_j$。于是
\[
  gH_jg^{-1}=h_1h_jH_j(h_1h_j)^{-1}=h_1H_jh_1^{-1},
\]
从而
\[
  \bigcap_{g\in G}gH_jg^{-1}=\bigcap_{h_1\in H_1}h_1H_jh_1^{-1}.
\]
现在 $P\leq H_j$，且对所有 $h_1\in H_1$ 都有 $h_1Ph_1^{-1}=P$。因此
\[
  1\neq P\leq \bigcap_{h_1\in H_1}h_1H_jh_1^{-1}.
\]
于是 $N=\bigcap_{g\in G}gH_jg^{-1}$ 是 $G$ 的一个非平凡真正规子群。由此 $N$ 和 $G/N$ 都满足定理的假设（参见第 3.3 节中的练习）。由归纳假设，$N$ 与 $G/N$ 都可解，所以 $G$ 可解。这完成 Hall 定理的证明。

\subsection*{练习}
\begin{enumerate}[label=\arabic*.]
  \item 证明对所有 $n$，对称群 $S_n$ 的每个特征标都是整数值的（即对所有 $g\in S_n$ 以及 $S_n$ 的所有特征标 $\psi$，都有 $\psi(g)\in\ZZ$）。［见第 18.3 节练习 22。］
  \item 设 $G$ 是有限群，且每个极大子群的指数或者是素数，或者是素数的平方。证明 $G$ 可解。（单群 $\GL_3(\FF_2)$ 具有如下性质：每个极大子群的指数或者是 $7$，或者是 $8$，即或者是素数，或者是素数的三次幂——参见第 6.2 节。）［令 $p$ 为整除 $|G|$ 的最大素数，令 $P$ 为 $G$ 的 Sylow $p$-子群。若 $P\normal G$，对 $G/P$ 应用归纳。否则令 $M$ 为包含 $N_G(P)$ 的极大子群。使用第 4.5 节练习 51 证明 $p=3$，并推出 $|G|=2^a3^b$。］
  \item 假设 $G$ 是一个拥有 Abel 子群 $H$ 的有限群，且 $H$ 的指数是某个素数的幂。证明 $G$ 可解。
  \item 把前一练习中的“Abel”替换为“幂零”，重复该练习。
  \item 使用 Philip Hall 定理证明中的思想，在所有 Sylow 子群都是 Abel 群的特殊情形下证明 Burnside 的 $p^a q^b$ 定理（不使用特征标理论）。
\end{enumerate}

\section{诱导特征标理论导论}

设 $G$ 是有限群，$H$ 是 $G$ 的子群，$\varphi$ 是子群 $H$ 在任意域 $F$ 上的表示。在本节中，我们说明如何由其子群的表示 $\varphi$ 得到 $G$ 的一个表示，称为\emph{诱导表示}。我们还确定这个诱导表示的特征标，即\emph{诱导特征标}，如何用 $\varphi$ 的特征标表示，并通过在具体群中计算若干诱导特征标来说明这个公式。最后，我们应用诱导特征标理论证明不存在 $3^3\cdot7\cdot13\cdot409$ 阶单群；在第 6.2 节末尾关于单群存在性问题的讨论中曾提到这个群阶。诱导表示和诱导特征标理论标志着更高级表示论的开端。本节旨在作为导论，而不是全面处理；我们列入的结果也是为这一目的而选择的。

首先注意，子群 $H$ 的表示 $\varphi$ 不一定能延拓为 $G$ 的一个表示 $\Phi$，并满足 $\Phi|_H=\varphi$。例如，$A_3\normal S_3$，而 $A_3$ 有一个忠实的次数 $1$ 表示（参见第 1 节）。由于 $S_3$ 的每个次数 $1$ 表示都以 $A_3=S_3'$ 为核的一部分，所以 $A_3$ 的这个表示不能延拓为 $S_3$ 的表示。作为另一个不能延拓到整个群的子群表示的例子，取 $G$ 为任意单群，令 $\varphi$ 为 $H$ 的某个表示，使得 $\ker\varphi$ 是 $H$ 的一个真非平凡正规子群。如果 $\varphi$ 延拓为 $G$ 的表示 $\Phi$，则 $\Phi$ 的核将是 $G$ 的一个真非平凡子群，这与 $G$ 是单群矛盾。我们将看到，诱导特征标的方法由子群 $H$ 的给定表示 $\varphi$ 产生 $G$ 的一个表示 $\Phi$，但一般 $\Phi|_H\neq\varphi$（事实上，除非 $H=G$，否则 $\Phi$ 的次数将大于 $\varphi$ 的次数）。

我们在第 10.4 节推论 9 后的例 5 中看到，由于 $FH$ 是 $FG$ 的子环，环 $FG$ 是一个 $(FG,FH)$-双模；因此对任意左 $FH$-模 $V$，Abel 群 $FG\otimes_{FH}V$ 是左 $FG$-模（称为 $V$ 从 $FH$ 到 $FG$ 的标量扩张）。在有限群表示论中，这个扩张有一个专门名称。

\noindent\textbf{定义.} 设 $H$ 是有限群 $G$ 的子群，$V$ 是一个实现 $H$ 的表示 $\varphi$ 的 $FH$-模。$FG$-模 $FG\otimes_{FH}V$ 称为 $V$ 的\emph{诱导模}，它所实现的 $G$ 的表示称为 $\varphi$ 的\emph{诱导表示}。若 $\psi$ 是 $\varphi$ 的特征标，则诱导表示的特征标称为\emph{诱导特征标}，记作 $\Ind_H^G(\psi)$。

\noindent\textbf{定理 11.} 设 $H$ 是有限群 $G$ 的子群，$g_1,\ldots,g_m$ 是 $G$ 中 $H$ 的不同左陪集的代表。设 $V$ 是实现 $H$ 的次数 $n$ 矩阵表示 $\varphi$ 的 $FH$-模。$FG$-模
\[
  W=FG\otimes_{FH}V
\]
在 $F$ 上维数为 $nm$，且 $W$ 有一组基，使它实现如下对每个 $g\in G$ 定义的矩阵表示 $\Phi$：
\[
\Phi(g)=
\begin{pmatrix}
\varphi(g_1^{-1}gg_1)&\cdots&\varphi(g_1^{-1}gg_m)\\
\vdots&&\vdots\\
\varphi(g_m^{-1}gg_1)&\cdots&\varphi(g_m^{-1}gg_m)
\end{pmatrix},
\]
其中每个 $\varphi(g_i^{-1}gg_j)$ 是出现在 $\Phi(g)$ 的第 $i,j$ 个块位置的 $n\times n$ 块；当 $g_i^{-1}gg_j\notin H$ 时，规定 $\varphi(g_i^{-1}gg_j)$ 为零块。

\noindent\textbf{证明.} 首先注意，$FG$ 是一个自由右 $FH$-模：
\[
  FG=g_1FH\oplus g_2FH\oplus\cdots\oplus g_mFH.
\]
由于张量积与直和交换（第 10.4 节定理 17），作为 Abel 群有
\[
  W=FG\otimes_{FH}V\cong (g_1\otimes V)\oplus(g_2\otimes V)\oplus\cdots\oplus(g_m\otimes V).
\]
由于 $F$ 位于 $FG$ 的中心中，这也是 $F$-向量空间同构。因此，若 $v_1,v_2,\ldots,v_n$ 是 $V$ 的一个实现矩阵表示 $\varphi$ 的基，则
\[
  \{g_i\otimes v_j\mid 1\leq i\leq m,
  1\leq j\leq n\}
\]
是 $W$ 的一个基。这说明 $W$ 的维数为 $mn$。把基按每组 $n$ 个元素排序为
\[
  g_1\otimes v_1,\ldots,g_1\otimes v_n,
  g_2\otimes v_1,
  \ldots,g_2\otimes v_n,
  \ldots,
  g_m\otimes v_1,\
  \ldots,g_m\otimes v_n.
\]
我们计算每个 $g$ 作用在 $W$ 上相对于这个基的矩阵表示 $\Phi(g)$。固定 $j$ 与 $g$，并令 $gg_j=g_i h$，其中 $i$ 是某个指标，$h\in H$。于是对每个 $k$，
\[
\begin{aligned}
  g(g_j\otimes v_k)
  &=(gg_j)\otimes v_k\\
  &=g_i h\otimes v_k\\
  &=g_i\otimes hv_k\\
  &=\sum_{\ell=1}^n a_{\ell k}(h)(g_i\otimes v_\ell),
\end{aligned}
\]
其中 $a_{\ell k}$ 是 $h$ 相对于基 $\{v_1,\ldots,v_n\}$ 作用在 $V$ 上的矩阵的 $\ell,k$ 项。换言之，$g$ 在 $W$ 上的作用把第 $j$ 块的 $n$ 个基向量送到第 $i$ 块的基向量，并在该块上具有矩阵 $\varphi(h)$。由于 $h=g_i^{-1}gg_j$，这就描述了定理中的块矩阵 $\Phi(g)$，证毕。

\noindent\textbf{推论 12.} 在定理 11 的记号下：
\begin{enumerate}[label=(\arabic*)]
  \item 若 $\psi$ 是 $V$ 实现的特征标，则诱导特征标由
  \[
    \Ind_H^G(\psi)(g)=\sum_{i=1}^m \psi(g_i^{-1}gg_i)
  \]
  给出，其中当 $g_i^{-1}gg_i\notin H$ 时，规定 $\psi(g_i^{-1}gg_i)=0$；
  \item 若 $g$ 在 $G$ 中不与 $H$ 的任何元素共轭，则 $\Ind_H^G(\psi)(g)=0$。特别地，若 $H$ 是 $G$ 的正规子群，则 $\Ind_H^G(\psi)$ 在 $G-H$ 的所有元素上为零。
\end{enumerate}

\noindent\textbf{注记.} 由于 $H$ 的特征标 $\psi$ 在 $H$ 的共轭类上为常数，对所有 $h\in H$ 都有 $\psi(g)=\psi(h^{-1}gh)$。当 $h$ 遍历 $H$ 的所有元素时，$xh$ 遍历陪集 $xH$ 的所有元素。因此诱导特征标的公式也可写为
\[
  \Ind_H^G(\psi)(g)=\frac1{|H|}\sum_{x\in G}\psi(x^{-1}gx),
\]
其中每个固定陪集中的元素给出相同的特征标值 $|H|$ 次（这解释了因子 $1/|H|$），并且当 $x^{-1}gx\notin H$ 时仍规定 $\psi(x^{-1}gx)=0$。

\noindent\textbf{证明.} 由上面计算出的 $g$ 的矩阵，$\Phi(g)$ 对角线上各块 $\varphi(g_i^{-1}gg_i)$ 除了 $g_i^{-1}gg_i\in H$ 时以外都是零。因此块矩阵 $\Phi(g)$ 的迹是所有满足 $g_i^{-1}gg_i\in H$ 的矩阵 $\varphi(g_i^{-1}gg_i)$ 的迹之和。由于 $\varphi(g_i^{-1}gg_i)$ 的迹是 $\psi(g_i^{-1}gg_i)$，第 (1) 部分成立。

若对所有陪集代表 $g_i$ 都有 $g_i^{-1}gg_i\notin H$，则 $\Ind_H^G(\psi)(g)$ 的和中每一项都是零。特别地，若 $g$ 不在正规子群 $H$ 中，则 $g$ 的任意共轭元也不在 $H$ 中，所以 $\Ind_H^G(\psi)$ 在 $g$ 上为零。

\paragraph{例}
\begin{enumerate}[label=(\arabic*)]
  \item 令
  \[
    G=D_{12}=\langle r,s\mid r^6=s^2=1,\,rs=sr^{-1}\rangle
  \]
  为 $12$ 阶二面体群，令 $H=\{1,s,r^3,sr^3\}$，于是 $H$ 同构于 Klein 四元群，且 $|G:H|=3$。按定理 11 的记号，我们展示 $H$ 的某个具体表示 $\varphi$ 的诱导表示中 $r$ 与 $s$ 的矩阵。设 $H$ 在 $\QQ$ 上二维向量空间上的表示相对于某个基 $v_1,v_2$ 为
  \[
    \varphi(s)=\begin{pmatrix}-1&0\\0&1\end{pmatrix}=A,
    \qquad
    \varphi(r^3)=\begin{pmatrix}1&0\\0&-1\end{pmatrix}=B,
    \qquad
    \varphi(sr^3)=\begin{pmatrix}-1&0\\0&-1\end{pmatrix}=C.
  \]
  因此 $n=2$，$m=3$，诱导表示 $\Phi$ 次数为 $nm=6$。取 $g_1=1,g_2=r,g_3=r^2$ 为 $G$ 中 $H$ 的左陪集代表，于是 $g_k=r^{k-1}$。则
  \[
    g_i^{-1}rg_j=r^{-(i-1)+1+(j-1)}=r^{j-i+1},
    \qquad
    g_i^{-1}sg_j=sr^{(i-1)+(j-1)}=sr^{i+j-2}.
  \]
  因此诱导表示的 $6\times6$ 矩阵为
  \[
    \Phi(r)=\begin{pmatrix}
    0&0&B\\ I&0&0\\0&I&0
    \end{pmatrix},
    \qquad
    \Phi(s)=\begin{pmatrix}
    A&0&0\\0&0&C\\0&C&0
    \end{pmatrix},
  \]
  其中 $2\times2$ 矩阵 $A,B,C$ 如上，$I$ 是 $2\times2$ 单位矩阵，$0$ 表示 $2\times2$ 零矩阵。

  \item 若 $H$ 是 $G$ 的任意子群，$\psi_1$ 是 $H$ 的主特征标，则 $\Ind_H^G(\psi_1)(g)$ 对每个满足 $g_i^{-1}gg_i\in H$ 的陪集代表 $g_i$ 计数 $1$。由于 $g_i^{-1}gg_i\in H$ 当且仅当 $g$ 在左乘作用下固定左陪集 $g_iH$，$\Ind_H^G(\psi_1)(g)$ 是 $g$ 在 $G$ 关于 $H$ 左陪集的置换表示中固定的点数。因此由第 18.3 节例 3 可知：若 $\psi_1$ 是 $H$ 的主特征标，则 $\Ind_H^G(\psi_1)$ 是 $G$ 在 $H$ 的左陪集上的置换特征标。在特殊情形 $H=1$ 中，这意味着若 $\chi_1$ 是平凡子群 $H=1$ 的主特征标，则 $\Ind_1^G(\chi_1)$ 是 $G$ 的正则特征标。这也说明，诱导特征标一般不是不可约的，即使被诱导的特征标本身不可约。

  \item 令 $G=S_3$，令 $\psi$ 是 $A_3=\langle x\rangle$ 的一个非主线性特征标，使得 $\psi(x)=\zeta$，其中 $\zeta$ 是一个本原三次单位根（$A_3=Z_3$ 和 $S_3$ 的特征标表出现在第 1 节中）。令 $\Psi=\Ind_{A_3}^{S_3}(\psi)$。于是 $\Psi$ 的次数为 $1\cdot |S_3:A_3|=2$，并且由推论，$\Psi$ 在所有对换上为零。若 $y$ 是任意对换，则 $1,y$ 是 $S_3$ 中 $A_3$ 的左陪集代表集，且 $y^{-1}xy=x^2$。因此 $\Psi(x)=\psi(x)+\psi(x^2)$ 等于 $\zeta+\zeta^2=-1$。这说明若 $\psi$ 是 $A_3$ 的两个非主不可约特征标之一，则 $\psi$ 的诱导特征标就是 $S_3$ 的（唯一）次数 $2$ 不可约特征标。特别地，一个子群的不同特征标可能诱导出整个群的同一特征标。

  \item 令 $G=D_8$ 有通常的生成元与关系，并令 $H=\langle s\rangle$。令 $\psi$ 是 $H$ 的非主不可约特征标，令 $\Psi=\Ind_H^G(\psi)$。取 $1,r,r^2,r^3$ 为 $H$ 的左陪集代表。由定理 11，$\Psi(1)=4$。由于 $\psi(s)=-1$，可直接计算出 $\Psi(s)=-2$。由推论 12(2)，得到 $\Psi(r)=\Psi(r^2)=\Psi(sr)=0$。在第 1 节 $D_8$ 特征标表的记号中，由正交关系得到 $\Psi=\chi_2+\chi_4+\chi_5$（也可通过观察检验）。
\end{enumerate}

本节剩余部分取域 $F$ 为复数域：$F=\CC$。

在最后给出诱导特征标在单群上的一个应用之前，我们先计算一类重要群的特征标。

\noindent\textbf{定义.} 有限群 $G$ 称为具有 Frobenius 核 $Q$ 的 \emph{Frobenius 群}，如果 $Q$ 是 $G$ 的真非平凡正规子群，并且对 $Q$ 的所有非恒等元 $x$ 都有 $C_G(x)\leq Q$。

鉴于本节开头提到的对单群的应用，我们将限制在如下 Frobenius 群 $G$：其阶为 $q^a p$，其中 $p,q$ 是不同素数，Frobenius 核 $Q$ 是 $q^a$ 阶初等 Abel $q$-群，并且循环群 $G/Q$ 通过共轭作用不可约地作用在 $Q$ 上。换言之，我们假设 $Q$ 是若干 $q$ 阶循环群的直积，且 $G$ 中包含在 $Q$ 中的正规子群只有 $1$ 和 $Q$，即 $Q$ 是 $G$ 的极小正规子群。例如，$A_4$ 是这种类型的 Frobenius 群，其 Frobenius 核为它的 Sylow $2$-子群 $V_4$。此外，若 $p,q$ 是不同素数，$p<q$，且 $G$ 是 $pq$ 阶非 Abel 群（当 $p\mid q-1$ 时总存在），则 $G$ 是 Frobenius 群，其 Frobenius 核为它的 Sylow $q$-子群（由 Sylow 定理它是正规子群）。我们基本确定这些 Frobenius 群的特征标表。更一般 Frobenius 群的类似结果见练习。

\noindent\textbf{命题 13.} 设 $G$ 是阶为 $q^a p$ 的 Frobenius 群，其中 $p,q$ 是不同素数，Frobenius 核 $Q$ 是阶为 $q^a$ 的初等 Abel $q$-群，并且循环群 $G/Q$ 通过共轭作用不可约地作用在 $Q$ 上。则下列结论成立：
\begin{enumerate}[label=(\arabic*)]
  \item $G=QP$，其中 $P$ 是 $G$ 的 Sylow $p$-子群。$G$ 的每个非恒等元的阶为 $p$ 或 $q$。每个 $p$ 阶元都与 $P$ 的某个元素共轭，每个 $q$ 阶元都属于 $Q$。$P$ 的非恒等元代表 $p-1$ 个不同的 $p$ 阶元共轭类，每个这样的类大小为 $q^a$。$q$ 阶元素有 $(q^a-1)/p$ 个不同共轭类，每个这样的类大小为 $p$。
  \item $G'=Q$，所以 $G$ 的次数 $1$ 特征标数为 $p$，且每个次数 $1$ 特征标的核都含有 $Q$。
  \item 若 $\psi$ 是 $Q$ 的任意非主不可约特征标，则 $\Ind_Q^G(\psi)$ 是 $G$ 的一个不可约特征标。此外，$G$ 的每个次数 $>1$ 的不可约特征标都等于某个 $Q$ 的非主不可约特征标 $\psi$ 的 $\Ind_Q^G(\psi)$。$G$ 的每个不可约特征标次数为 $1$ 或 $p$，而次数 $p$ 的不可约特征标个数为 $(q^a-1)/p$。
\end{enumerate}

\noindent\textbf{证明.} 注意，由阶的考虑，$QP=G$。由 Frobenius 群的定义并由于 $Q$ 是 Abel 群，对 $Q$ 的每个非恒等元 $h$，有 $C_G(h)=Q$。若 $x$ 是一个 $pq$ 阶元，则 $x^p$ 是 $q$ 阶元，因而属于 $G$ 的唯一 Sylow $q$-子群 $Q$。但于是 $x$ 与 $x^p$ 交换，所以 $x\in C_G(x^p)=Q$，矛盾。因此 $G$ 没有 $pq$ 阶元素。由 Sylow 定理，每个 $p$ 阶元都与 $P$ 的某个元素共轭，每个 $q$ 阶元都属于 $Q$。$P$ 的两个不同元素在 $G$ 中不共轭，因为若对某些 $x,y\in P$ 有 $g^{-1}xg=y$，则在 Abel 群 $\overline G=G/Q$ 中有 $\bar g^{-1}\bar x\bar g=\bar y$，所以 $\bar x=\bar y$。于是 $x=y$，因为 $P\cap Q=1$。因此 $p$ 阶元素恰有 $p-1$ 个共轭类，它们由 $P$ 的非恒等元代表。若 $x$ 是 $P$ 的非恒等元，则 $C_G(x)=P$，所以 $x$ 的共轭类含有 $|G:P|=q^a$ 个元素。最后，若 $h$ 是 $Q$ 的非恒等元，则 $C_G(h)=Q$，且若 $P=\langle x\rangle$，$h$ 的共轭类是 $\{h,h^x,\ldots,h^{x^{p-1}}\}$。这证明了 (1) 的所有部分。

由于 $G/Q$ 是 Abel 群，$G'\leq Q$。由于 $G$ 非 Abel 且 $Q$ 按假设是 $G$ 的极小正规子群，必有 $G'=Q$。现在 (2) 由第 18.2 节推论 11 得出。

令 $\psi$ 是 $Q$ 的非主不可约特征标，并令 $\Psi=\Ind_Q^G(\psi)$。我们用正交关系说明 $\Psi$ 不可约。令 $1,x,\ldots,x^{p-1}$ 为 $Q$ 在 $G$ 中的陪集代表。由推论 12，$\Psi$ 在 $G-Q$ 上为零，所以
\[
\begin{aligned}
  \|\Psi\|^2
  &=\frac1{|G|}\sum_{h\in Q}\Psi(h)\overline{\Psi(h)}\\
  &=\frac1{|G|}\sum_{h\in Q}\left|\sum_{i=0}^{p-1}\psi(x^i h x^{-i})\right|^2\\
  &=\frac{p}{|G|}\sum_{h\in Q}\psi(h)\overline{\psi(h)}\\
  &=\frac{p|Q|}{|G|}=1,
\end{aligned}
\]
其中第二行来自诱导特征标 $\Psi$ 的定义，第三行因为第二行求和中 $Q$ 的每个元素恰出现 $p$ 次，最后一行来自 $Q$ 中第一正交关系，因为 $\psi$ 是 $Q$ 的不可约特征标。这证明 $\Psi$ 是 $G$ 的不可约特征标。

我们通过计数这样得到的不同不可约特征标的数目，证明 $G$ 的每个次数 $>1$ 的不可约特征标都是 $Q$ 的某个非主次数 $1$ 特征标的诱导特征标。由 (1) 和 (2)，$G$ 的不可约特征标数（即共轭类数）为 $p+(q^a-1)/p$，次数 $1$ 的特征标数为 $p$。因此 $G$ 的次数 $>1$ 的不可约特征标数为 $(q^a-1)/p$。群 $P$ 在 $Q$ 的非主不可约特征标集合 $\mathcal C$ 上作用如下：对每个 $\psi\in\mathcal C$ 和每个 $x\in P$，定义 $\psi^x$ 为
\[
  \psi^x(h)=\psi(xhx^{-1})\qquad\text{对所有 }h\in Q.
\]
由于 $\psi$ 是从 $Q$ 到 $\CC^\times$ 的非平凡同态（回忆，Abel 群 $Q$ 的所有不可约特征标次数都为 $1$），容易推出 $\psi^x$ 也是同态。因此 $\psi^x\in\mathcal C$，所以 $P$ 置换 $\mathcal C$ 的元素。现在令 $x$ 是循环群 $P$ 的生成元。于是 $1,x,\ldots,x^{p-1}$ 是 $Q$ 在 $G$ 中的左陪集代表。由推论 12，对这组陪集代表，若 $\psi\in\mathcal C$，则 $\Ind_Q^G(\psi)$ 在任意 $h\in Q$ 上的值是和
\[
  \psi(h)+\psi^x(h)+\cdots+\psi^{x^{p-1}}(h).
\]
因此诱导特征标 $\Ind_Q^G(\psi)$ 限制到 $Q$ 后，分解为 $Q$ 的不可约特征标之和：
\[
  \Ind_Q^G(\psi)|_Q=\psi+\psi^x+\cdots+\psi^{x^{p-1}}.
\]
若 $\psi_1,\psi_2$ 属于 $P$ 在 $\mathcal C$ 上作用的不同轨道，则诱导特征标 $\Ind_Q^G(\psi_1)$ 和 $\Ind_Q^G(\psi_2)$ 限制到 $Q$ 后是不同的特征标（它们没有共同的不可约组成因子）。因此从 $P$ 在 $\mathcal C$ 上的不同轨道中的元素诱导出的特征标，是 $G$ 的不同不可约特征标。Abel 群 $Q$ 有 $q^a-1$ 个非主不可约特征标（即 $|\mathcal C|=q^a-1$），且 $|P|=p$，所以 $P$ 在 $\mathcal C$ 上至少有 $(q^a-1)/p$ 个轨道，从而至少得到这么多不同的次数 $p$ 不可约特征标。由于 $G$ 恰有 $(q^a-1)/p$ 个次数 $>1$ 的不可约特征标，$G$ 的每个次数 $>1$ 不可约特征标必有次数 $p$，且必是某个 $\mathcal C$ 中元素的诱导特征标。证明完毕。

对于最后的例子，我们需要诱导特征标的两个性质。这些性质列于下一个命题中，其证明是直接练习，很容易由诱导特征标公式或由诱导模的定义以及张量积性质得出。

\noindent\textbf{命题 14.} 设 $G$ 是群，$H$ 是 $G$ 的子群，$\psi$ 与 $\psi'$ 是 $H$ 的特征标。
\begin{enumerate}[label=(\arabic*)]
  \item （特征标诱导的可加性）$\Ind_H^G(\psi+\psi')=\Ind_H^G(\psi)+\Ind_H^G(\psi')$。
  \item （特征标诱导的传递性）若 $H\leq K\leq G$，则
  \[
    \Ind_K^G(\Ind_H^K(\psi))=\Ind_H^G(\psi).
  \]
\end{enumerate}

由命题 14 的第 (1) 部分可知，若 $\sum_{i=1}^t n_i\psi_i$ 是 $H$ 的特征标的任意整线性组合，且对所有 $i$ 有 $n_i\geq0$，则
\[
  \Ind_H^G\left(\sum_{i=1}^t n_i\psi_i\right)
  =\sum_{i=1}^t n_i\Ind_H^G(\psi_i). \tag{*}
\]
$H$ 的形如 $\sum_{i=1}^t n_i\psi_i$ 的类函数（其中系数是任意整数，不必非负）称为 $H$ 的\emph{广义特征标}或\emph{虚特征标}。对 $H$ 的广义特征标，我们用方程 (*) 定义其 $G$ 的诱导广义特征标，现在也允许负系数 $n_i$。这样，函数 $\Ind_H^G$ 成为从 $H$ 的广义特征标加法群到 $G$ 的广义特征标加法群的群同态（它把特征标映到特征标）。这意味着推论 12 中的诱导特征标公式在 $\psi$ 是 $H$ 的广义特征标时同样成立。

\subsection*{关于 $3^3\cdot7\cdot13\cdot409$ 阶群的应用}

现在我们以证明下列结果结束：
\[
  \text{不存在 }3^3\cdot7\cdot13\cdot409\text{ 阶单群。}
\]
如本节开头所说，在第 6.2 节末尾关于单群存在性问题的讨论中提到过这个阶的单群。可以通过涉及次数 $819$ 的置换表示的论证证明不存在这个阶的单群（参见第 6.2 节练习）。我们在这里给出一个特征标论证明，因为这些方法说明了有限群理论中的一些重要思想。这个方法基于 M. Suzuki 的奠基性论文 \emph{The nonexistence of a certain type of simple group of odd order}, Proc. Amer. Math. Soc., 8 (1957), pp. 686--695，该文处理更一般的群。由于我们处理的是一个具体的群阶，论证更简单，数值上也更显式，但仍保留了 Suzuki 工作的一些关键思想。此外，Suzuki 的论文及其后继论文 W. Feit、M. Hall 与 J. Thompson 的 \emph{Finite groups in which the centralizer of any non-identity element is nilpotent}, Math. Zeit., 74 (1960), pp. 1--17，是冗长而困难的 Feit--Thompson 定理（参见第 3.4 节）的原型。我们的讨论也传达了这些基本论文的一些风味。特别地，这些论文都遵循一种基本发展路线：首先确定 Sylow 子群的结构与嵌入方式，然后应用特征标理论（极大依赖诱导特征标）。

在本节剩余部分，假设 $G$ 是一个 $3^3\cdot7\cdot13\cdot409$ 阶单群。我们列出 $G$ 的一些性质，它们可用第 6.2 节中源于 Sylow 定理的方法验证。细节留作练习。
\begin{enumerate}[label=(\arabic*)]
  \item 令 $q_1=3$，令 $Q_1$ 为 $G$ 的 Sylow $3$-子群，令 $N_1=N_G(Q_1)$。则 $Q_1$ 是一个 $3^3$ 阶初等 Abel $3$-群，且 $N_1$ 是一个阶为 $3^3\cdot13$ 的 Frobenius 群，Frobenius 核为 $Q_1$，并且 $N_1/Q_1$ 通过共轭不可约地作用在 $Q_1$ 上。
  \item 令 $q_2=7$，令 $Q_2$ 为 $G$ 的 Sylow $7$-子群，令 $N_2=N_G(Q_2)$。则 $Q_2$ 是 $7$ 阶循环群，$N_2$ 是 $7\cdot3$ 阶非 Abel 群（所以 $N_2$ 是以 $Q_2$ 为 Frobenius 核的 Frobenius 群）。
  \item 令 $q_3=13$，令 $Q_3$ 为 $G$ 的 Sylow $13$-子群，令 $N_3=N_G(Q_3)$。则 $Q_3$ 是 $13$ 阶循环群，$N_3$ 是 $13\cdot3$ 阶非 Abel 群（所以 $N_3$ 是以 $Q_3$ 为 Frobenius 核的 Frobenius 群）。
  \item 令 $q_4=409$，令 $Q_4$ 为 $G$ 的 Sylow $409$-子群，令 $N_4=N_G(Q_4)$。则 $Q_4$ 是 $409$ 阶循环群，$N_4$ 是 $409\cdot3$ 阶非 Abel 群（所以 $N_4$ 是以 $Q_4$ 为 Frobenius 核的 Frobenius 群）。
  \item $G$ 的每个非恒等元都有素数阶，并且对每个 $i=1,2,3,4$ 以及每个 $g\in G-N_G(Q_i)$ 都有 $Q_i\cap Q_i^g=1$。$G$ 的非恒等共轭类为：
  \begin{enumerate}[label=(\alph*)]
    \item $2$ 个 $3$ 阶元素类（每个类大小为 $7\cdot13\cdot409$）；
    \item $2$ 个 $7$ 阶元素类（每个类大小为 $3^3\cdot13\cdot409$）；
    \item $4$ 个 $13$ 阶元素类（每个类大小为 $3^3\cdot7\cdot409$）；
    \item $136$ 个 $409$ 阶元素类（每个类大小为 $3^3\cdot7\cdot13$），
  \end{enumerate}
  因而 $G$ 中有 $145$ 个共轭类。
\end{enumerate}
由于每个群 $N_i$ 都是满足命题 13 假设的 Frobenius 群，$N_i$ 的次数 $>1$ 特征标个数可从该命题读出：
\begin{enumerate}[label=(\roman*)]
  \item $N_1$ 有 $2$ 个次数 $13$ 的不可约特征标；
  \item $N_2$ 有 $2$ 个次数 $3$ 的不可约特征标；
  \item $N_3$ 有 $4$ 个次数 $3$ 的不可约特征标；
  \item $N_4$ 有 $136$ 个次数 $3$ 的不可约特征标。
\end{enumerate}
从现在起，为简化记号，对 $G$ 的任意子群 $H$ 和 $H$ 的任意广义特征标 $\mu$，令
\[
  \mu^*=\Ind_H^G(\mu),
\]
所以星号总表示从某个子群诱导到整个群 $G$，而该子群将由上下文明确。

下面的引理是证明中的关键点。它说明推论 12 中诱导特征标的消失性质（连同 Sylow 子群 $Q_i$ 的平凡交性质，即对所有 $g\in G-N_G(Q_i)$ 有 $Q_i\cap Q_i^g=1$）如何把某些广义特征标的内积与它们的诱导广义特征标的内积联系起来。对于这些计算，重要的是广义特征标在恒等元上为零（这解释了为什么我们考虑相同次数特征标的差）。

\noindent\textbf{引理 15.} 对任意 $i\in\{1,2,3,4\}$，令 $q=q_i$，$Q=Q_i$，$N=N_i$，并令 $p=|N:Q|$。令 $\psi_1,\ldots,\psi_4$ 是 $N$ 的任意次数 $p$ 不可约特征标（不必互异），并令 $\alpha=\psi_1-\psi_2$，$\beta=\psi_3-\psi_4$。则 $\alpha$ 与 $\beta$ 是 $N$ 的广义特征标，并且在 $N$ 中每个阶不等于 $q$ 的元素上为零。此外，$\alpha^*$ 与 $\beta^*$ 是 $G$ 的广义特征标，并且在 $G$ 中每个阶不等于 $q$ 的元素上为零，且
\[
  (\alpha^*,\beta^*)_G=(\alpha,\beta)_N
\]
（其中 $(\ ,\ )_H$ 表示在群 $H$ 中计算的类函数通常 Hermite 内积）。换言之，从 $N$ 到 $G$ 的诱导在这样的 $N$ 的广义特征标 $\alpha,\beta$ 上保持内积。

\noindent\textbf{证明.} 由命题 13，存在 $Q$ 的非主次数 $1$ 特征标 $\lambda_1,\ldots,\lambda_4$，使得对 $j=1,\ldots,4$ 有 $\psi_j=\Ind_Q^N(\lambda_j)$。因此由推论 12，每个 $\psi_j$ 都在 $N-Q$ 上消失，故 $\alpha$ 与 $\beta$ 也如此。注意，由于对所有 $j$ 都有 $\psi_j(1)=p$，有 $\alpha(1)=\beta(1)=0$。由诱导的传递性，对所有 $j$，
\[
  \psi_j^*=\Ind_N^G(\psi_j)=\Ind_Q^G(\lambda_j).
\]
再次把推论 12 应用于后一诱导特征标，看到 $\psi_j^*$ 在所有不与 $Q$ 的某个元素共轭的元素上消失，因此 $\alpha^*$ 与 $\beta^*$ 也如此。由于诱导特征标 $\psi_j^*$ 全都次数为 $|G:Q|$，广义特征标 $\alpha^*$ 与 $\beta^*$ 在恒等元上为零。因此 $\alpha^*$ 与 $\beta^*$ 在 $G$ 中所有非 $q$ 阶元素上消失。最后，若 $g_1,\ldots,g_m$ 是 $N$ 在 $G$ 中的左陪集代表且 $g_1=1$，则由于对 $k>1$ 有 $Q\cap Q^{g_k}=1$（由上面的 (5)），由诱导（广义）特征标的公式立即得到，对所有非恒等元 $x\in Q$（即对所有 $N$ 中的 $q$ 阶元素 $x$），有 $\alpha^*(x)=\alpha(x)$ 且 $\beta^*(x)=\beta(x)$。此外，由 Sylow 定理，$G$ 中每个 $q$ 阶元素都属于 $Q$ 的某个共轭，因此集合 $Q-\{1\}$ 的 $G$-共轭把 $G$ 中 $q$ 阶元素划分为 $|G:N|$ 个不交子集。由于 $\alpha^*$ 与 $\beta^*$ 是 $G$ 上类函数，当 $x$ 遍历这些子集中的任意一个时，$\sum \alpha^*(x)\overline{\beta^*(x)}$ 相同。这些事实蕴含
\[
\begin{aligned}
  (\alpha^*,\beta^*)_G
  &=\frac1{|G|}\sum_{x\in G}\alpha^*(x)\overline{\beta^*(x)}\\
  &=\frac1{|G|}\sum_{\substack{x\in G\\ |x|=q}}\alpha^*(x)\overline{\beta^*(x)}\\
  &=\frac1{|G|}\sum_{\substack{x\in N\\ |x|=q}} |G:N|\alpha^*(x)\overline{\beta^*(x)}\\
  &=\frac1{|N|}\sum_{x\in N}\alpha(x)\overline{\beta(x)}=(\alpha,\beta)_N.
\end{aligned}
\]
证明完毕。

下一个引理在 $N_i$ 的次数 $>1$ 不可约特征标与 $G$ 的某些非主不可约特征标之间建立对应。

\noindent\textbf{引理 16.} 对任意 $i\in\{1,2,3,4\}$，令 $q=q_i$，$Q=Q_i$，$N=N_i$，并令 $p=|N:Q|$。令 $\psi_1,\ldots,\psi_k$ 是 $N$ 的不同次数 $p$ 不可约特征标。则存在 $G$ 的两两不同不可约特征标 $\chi_1,\ldots,\chi_k$，它们都有相同次数，并且存在一个固定符号 $\varepsilon=\pm1$，使得对所有 $j=2,3,\ldots,k$，有
\[
  \psi_1^*-\psi_j^*=\varepsilon(\chi_1-\chi_j).
\]

\noindent\textbf{证明.} 令 $\alpha_j=\psi_1-\psi_j$，$j=2,3,\ldots,k$，则 $\alpha_j$ 满足引理 15 的假设。由于 $\psi_1\neq\psi_j$，由引理 15，
\[
  2=\|\alpha_j\|^2=(\alpha_j,\alpha_j)_N=(\alpha_j^*,\alpha_j^*)_G=\|\alpha_j^*\|^2
\]
对所有 $j$ 成立。因此 $\alpha_j^*$ 必有 $G$ 的两个不同不可约特征标作为其不可约组成因子。由于 $\alpha_j^*(1)=0$，它必是两个具有相同次数的不同不可约特征标之差。特别地，若 $k=2$，引理成立（这正是 $q=3$ 与 $q=7$ 的情形）。所以假设 $k>2$，并写
\[
  \alpha_2^*=\psi_1^*-\psi_2^*=\varepsilon(\chi-\chi'),
  \qquad
  \alpha_3^*=\psi_1^*-\psi_3^*=\varepsilon_1(\theta-\theta')
\]
其中 $\chi,\chi',\theta,\theta'$ 是 $G$ 的不可约特征标，$\varepsilon,\varepsilon_1$ 是符号。如上所证，$\chi\neq\chi'$ 且 $\theta\neq\theta'$。必要时交换 $\theta$ 与 $\theta'$，可假设 $\varepsilon=\varepsilon_1$。于是
\[
  \alpha_2^*-\alpha_3^*=\psi_3^*-\psi_2^*=\varepsilon(\theta-\theta'-\chi+\chi').
\]
由引理 15，$\psi_2^*-\psi_3^*=(\psi_2-\psi_3)^*$ 也恰有两个不同不可约组成因子，因此或者 $\theta=\chi$，或者 $\theta'=\chi'$。必要时把 $\varepsilon$ 替换为 $-\varepsilon$，可假设 $\theta=\chi$，从而现在有
\[
  \alpha_2^*=\psi_1^*-\psi_2^*=\varepsilon(\chi-\chi'),
  \qquad
  \alpha_3^*=\psi_1^*-\psi_3^*=\varepsilon(\chi-\theta')
\]
其中 $\chi,\chi',\theta'$ 是 $G$ 的不同不可约特征标，符号 $\varepsilon$ 已确定。标记 $\chi=\chi_1$，$\chi'=\chi_2$，$\theta'=\chi_3$。现在类似地检查可知，对每个 $j\geq3$，存在 $G$ 的不可约特征标 $\chi_j$，使得
\[
  \alpha_j^*=\psi_1^*-\psi_j^*=\varepsilon(\chi_1-\chi_j),
\]
并且 $\chi_1,
\ldots,\chi_k$ 两两不同。由于所有 $\chi_j$ 与 $\chi_1$ 有相同次数，证明完成。

我们说明，未必有 $\chi_j=\psi_j^*$ 对任何 $j$ 成立；所成立的只是 $N$ 的不可约特征标之差诱导为 $G$ 的不可约特征标之差。

由引理 16 得到的 $G$ 的不可约特征标 $\chi_j$ 称为与 $Q$ 相联系的\emph{例外特征标}。

\noindent\textbf{引理 17.} 若 $i,j$ 是 $\{1,2,3,4\}$ 中不同元素，则与 $Q_i$ 相联系的例外特征标全都不同于与 $Q_j$ 相联系的例外特征标。

\noindent\textbf{证明.} 令 $\chi$ 是与 $Q_i$ 相联系的例外特征标，$\theta$ 是与 $Q_j$ 相联系的例外特征标。由构造，存在 $Q_i$ 的不同不可约特征标 $\psi,\psi'$，使得 $\psi^*-\psi'^*=\chi-\chi'$；并且存在 $Q_j$ 的不同不可约特征标 $\lambda,\lambda'$，使得 $\lambda^*-\lambda'^*=\theta-\theta'$。令 $\alpha=\psi-\psi'$，$\beta=\lambda-\lambda'$。由引理 15，$\alpha^*$ 在 $G$ 中所有阶不等于 $q_i$ 的元素上为零（包括恒等元），而 $\beta^*$ 在 $G$ 中所有阶不等于 $q_j$ 的元素上为零。因此显然 $(\alpha^*,\beta^*)=0$。容易推出，$\alpha^*$ 的两个不可约组成因子也分别与 $\beta^*$ 的两个不可约组成因子正交。这证明了引理。

现在很容易说明这样的单群 $G$ 不存在。由引理 16 和 $G$ 的性质 (i) 到 (iv)，我们可以计数例外特征标的个数：
\begin{enumerate}[label=(\roman*)]
  \item 与 $Q_1$ 相联系的例外特征标有 $2$ 个；
  \item 与 $Q_2$ 相联系的例外特征标有 $2$ 个；
  \item 与 $Q_3$ 相联系的例外特征标有 $4$ 个；
  \item 与 $Q_4$ 相联系的例外特征标有 $136$ 个。
\end{enumerate}
记与 $Q_i$ 相联系的例外特征标的共同次数为 $d_i$，$i=1,\ldots,4$。由引理 17，例外特征标给出了 $G$ 的 $144$ 个非主不可约特征标；因此这些特征标连同主特征标，就是 $G$ 的全部不可约特征标（$G$ 的共轭类数为 $145$）。不可约特征标次数平方和等于 $G$ 的阶：
\[
  1+2d_1^2+2d_2^2+4d_3^2+136d_4^2=1004913.
\]
化简得到
\[
  d_1^2+d_2^2+2d_3^2+68d_4^2=502456. \tag{19.4}
\]
最后，由于单群 $G$ 的每个非主不可约表示都是忠实的，并且 $N_1$ 的忠实表示的最小次数为 $13$，每个 $d_i\geq13$。由于 $d_4<\sqrt{502456/68}<86$ 且 $d_4$ 整除 $|G|$，可知
\[
  d_4\in\{13,21,27,39,63\}.
\]
此外，由推论 5，每个 $d_i$ 都整除 $|G|$，因此每个 $d_i$ 的可能性很少。现在检查方程 (4) 没有解（用计算机做这件事尤其容易）。这个矛盾完成了证明。

\subsection*{练习}

在所有练习中，所有表示都取在复数域上。
\begin{enumerate}[label=\arabic*.]
  \item 令 $G=S_3$，令 $H=A_3$，令 $V$ 是实现 $A_3$ 的自然置换表示的三维 $\CC H$-模。更明确地说，令 $V$ 有基 $e_1,e_2,e_3$，并令 $\sigma\in A_3$ 在 $V$ 上按 $\sigma e_i=e_{\sigma(i)}$ 作用。取 $1$ 和 $(1\ 2)$ 为 $S_3$ 中 $A_3$ 的左陪集代表，写出定理 11 中描述的 $S_3$ 在诱导模 $W$ 上作用的显式矩阵，对 $S_3$ 的每个元素都写出。
  \item 在 (a) 到 (f) 的每一部分中，都给出了某个群 $G$ 的子群 $H$ 的一个特征标 $\psi$。计算诱导特征标 $\Ind_H^G(\psi)$ 在 $G$ 的所有共轭类上的值，并利用第 1 节中的特征标表把 $\Ind_H^G(\psi)$ 写成不可约特征标之和：
  \begin{enumerate}[label=(\alph*)]
    \item $\psi$ 是 $S_3$ 的子群 $\langle(1\ 2)\rangle$ 的唯一非主次数 $1$ 特征标；
    \item $\psi$ 是 $D_8$ 的子群 $\langle r\rangle$ 的次数 $1$ 特征标，定义为 $\psi(r)=i$，其中 $i\in\CC$ 是 $-1$ 的平方根；
    \item $\psi$ 是 $D_8$ 的子群 $\langle r\rangle$ 的次数 $1$ 特征标，定义为 $\psi(r)=-1$；
    \item $\psi$ 是 $S_4$ 的子群 $V_4=\langle(1\ 2),(3\ 4)\rangle$ 的任意非主次数 $1$ 特征标；
    \item $\psi=\chi_4$ 是第 1 节中 $H=S_4$ 的特征标表中两个次数 $3$ 特征标的第一个，且 $H$ 是 $G=S_5$ 的子群；
    \item $\psi$ 是 $S_5$ 的子群 $V_4=\langle(1\ 2),(3\ 4)\rangle$ 的任意非主次数 $1$ 特征标。
  \end{enumerate}
  \item 使用命题 13，显式写出下列各群的特征标表：
  \begin{enumerate}[label=(\alph*)]
    \item $10$ 阶二面体群；
    \item $57$ 阶非 Abel 群；
    \item 一个 $56$ 阶非 Abel 群，它有正规初等 Abel Sylow $2$-子群。
  \end{enumerate}
  \item 令 $H$ 是 $G$ 的子群，令 $\varphi$ 是 $H$ 的一个表示，并假设 $N$ 是 $G$ 的正规子群，满足 $N\leq H$ 且 $N$ 包含在 $\varphi$ 的核中。证明 $N$ 也包含在 $\varphi$ 的诱导表示的核中。
  \item 令 $N$ 是 $G$ 的正规子群，令 $\psi_1$ 是 $N$ 的主特征标。令 $\Psi$ 为诱导特征标 $\Ind_N^G(\psi_1)$，于是由前一练习，我们可以把 $\Psi$ 看作 $G/N$ 的一个表示的特征标。证明 $\Psi$ 是 $G/N$ 的正则表示的特征标。
  \item 令 $Z$ 是 $G$ 的中心的任意子群，令 $|G:Z|=m$，令 $\psi$ 是 $Z$ 的特征标。证明
  \[
    \Ind_Z^G(\psi)(g)=
    \begin{cases}
      m\psi(g),& g\in Z,\\
      0,& g\notin Z.
    \end{cases}
  \]
  \item 令 $\varphi$ 是 $G$ 的子群 $H$ 的一个矩阵表示，并按定理 11 陈述中的显示公式对每个 $g\in G$ 定义矩阵 $\Phi(g)$。直接证明 $\Phi$ 是一个表示，即证明对所有 $x,y\in G$，都有 $\Phi(xy)=\Phi(x)\Phi(y)$。
  \item 令 $G$ 是一个具有 Frobenius 核 $Q$ 的 Frobenius 群。假设 $Q$ 和 $G/Q$ 都是 Abel 群，但 $G$ 不是 Abel 群（即 $G\neq Q$）。令 $|Q|=n$，$|G:Q|=m$。
  \begin{enumerate}[label=(\alph*)]
    \item 证明 $G/Q$ 是循环群，并证明 $G=QC$，其中 $C$ 是 $G$ 的某个循环子群，满足 $C\cap Q=1$（即 $G$ 是 $Q$ 与 $C$ 的半直积，且 $|C|=m$）。［令 $q$ 为 $n$ 的一个素因子，并令 $G/Q$ 通过共轭作用在初等 Abel $q$-群 $\{h\in Q\mid h^q=1\}$ 上。对这个 $\FF_qG/Q$-模的一个不可约组成因子应用第 18.1 节练习 14(e) 和 Frobenius 群的定义。］
    \item 证明 $n$ 与 $m$ 互素。［若素数 $p$ 同时整除 $Q$ 的阶与指数，令 $P$ 为 $G$ 的 Sylow $p$-子群。则 $P\cap Q\normal P$，且 $P\cap Q$ 是 $Q$ 的 Sylow $p$-子群。考虑子群 $Z(P)\cap Q$ 在 $G$ 中的中心化子（由第 6.1 节定理 1，这个交非平凡）。］
    \item 证明 $G$ 没有 $qp$ 阶元素，其中 $q$ 是 $n$ 的任意非平凡因子，$p$ 是 $m$ 的任意非平凡因子。［像命题 13 中那样论证。］
    \item 证明 $G$ 中包含在 $Q$ 中的非恒等共轭类个数为 $(n-1)/m$，每个这样的类大小为 $m$。［像命题 13 中那样论证。］
    \item 证明 $C$ 的两个不同元素在 $G$ 中不共轭。推出 $C$ 的非恒等元代表 $m-1$ 个不同的 $G$ 共轭类，且每个这样的类大小为 $n$。进一步推出 $G-Q$ 的每个元素都共轭于 $C$ 的某个元素，且 $G$ 有 $m+(n-1)/m$ 个共轭类。
    \item 证明 $G'=Q$，并推出 $G$ 有 $m$ 个不同的次数 $1$ 特征标。［为证明 $Q\leq G'$，令 $C=\langle x\rangle$，并论证映射 $h\mapsto[h,x]=x^{-1}h^{-1}xh$ 是从 $Q$ 到 $Q$ 的同态，其核平凡，因而该映射满。］
    \item 证明若 $\psi$ 是 $Q$ 的任意非主不可约特征标，则 $\Ind_Q^G(\psi)$ 是 $G$ 的一个不可约特征标。证明 $G$ 的每个次数 $>1$ 的不可约特征标都等于某个 $Q$ 的非主不可约特征标 $\psi$ 的 $\Ind_Q^G(\psi)$。推出 $G$ 的每个不可约特征标次数或者为 $1$，或者为 $m$，且次数 $m$ 的不可约特征标数为 $(n-1)/m$。［检查命题 13(3) 的证明在对相应数值作适当改变后，给出了这个更一般的结果。］
  \end{enumerate}
  \item 使用前一练习，显式写出
  \[
    \langle(1\ 2\ 3\ 4\ 5),(2\ 3\ 5\ 4)\rangle
  \]
  的特征标表，它是 $S_5$ 中 Sylow $5$-子群的正规化子（这个群是 $20$ 阶 Frobenius 群）。
  \item 令 $N$ 是 $G$ 的正规子群，令 $\psi$ 是 $N$ 的特征标，令 $g\in G$。定义 $\psi^g$ 为
  \[
    \psi^g(h)=\psi(ghg^{-1})\qquad\text{对所有 }h\in N.
  \]
  \begin{enumerate}[label=(\alph*)]
    \item 证明 $\psi^g$ 是 $N$ 的特征标（$\psi$ 与 $\psi^g$ 称为 $N$ 的 $G$-共轭特征标）。证明 $\psi^g$ 不可约当且仅当 $\psi$ 不可约。
    \item 证明映射 $\psi\mapsto\psi^g$ 是 $G$ 在 $N$ 的特征标集合上的一个右群作用，并且 $N$ 包含在这个作用的核中。
    \item 证明若 $\psi_1$ 与 $\psi_2$ 是 $N$ 的 $G$-共轭特征标，则 $\Ind_N^G(\psi_1)=\Ind_N^G(\psi_2)$。还证明，若 $\psi_1$ 与 $\psi_2$ 是 $N$ 的不是 $G$-共轭的特征标，则 $\Ind_N^G(\psi_1)\neq\Ind_N^G(\psi_2)$。［使用命题 13(3) 证明中的论证。］
  \end{enumerate}
  \item 证明若 $G=A_4$ 且 $N=V_4$ 是它的 Sylow $2$-子群，则 $N$ 的任意两个非主不可约特征标都是 $G$-共轭的（参见前一练习）。
  \item 令 $G=D_{2n}$ 由其通常的生成元与关系给出。证明若 $\psi$ 是 $H=\langle r\rangle$ 的任意次数 $1$ 特征标，且 $\psi\neq\psi^s$，则 $\Ind_H^G(\psi)$ 是 $D_{2n}$ 的一个不可约特征标。证明 $D_{2n}$ 的每个不可约特征标都是 $\langle r\rangle$ 的某个次数 $1$ 特征标的诱导特征标。
  \item 证明命题 14 的两个部分。
  \item 证明下列结果，称为 Frobenius 互反律：令 $H<G$，令 $\psi$ 是 $H$ 的任意特征标，令 $\chi$ 是 $G$ 的任意特征标。则
  \[
    (\psi,\chi|_H)_H=(\Ind_H^G(\psi),\chi)_G.
  \]
  ［用诱导特征标 $\Ind_H^G(\psi)$ 的公式展开右端，或仿照第 17.2 节中 Shapiro 引理的证明。］
  \item 假设 $G$ 是一个 $3^3\cdot7\cdot13\cdot409$ 阶单群，其 Sylow 子群及其正规化子由本节性质 (1) 到 (5) 描述。证明由 $G$ 在子群 $N_4$ 的左陪集上的作用得到的次数 $819$ 置换特征标分解为 $\chi_0+\gamma+\gamma'$，其中 $\chi_0$ 是 $G$ 的主特征标，$\gamma$ 与 $\gamma'$ 是 $G$ 的不同次数 $409$ 不可约特征标。［使用第 18.3 节练习 9 说明这个置换特征标 $\pi$ 满足 $\|\pi\|^2=3$。］
\end{enumerate}

% % ============================================
%  appendices_I_II_faithful.tex — 附录 I 与附录 II
%  根据 919-932.pdf 翻译整理，保持原文结构、证明与练习
% ============================================

\chapter*{附录 I\\笛卡尔积与 Zorn 引理}
\addcontentsline{toc}{chapter}{附录 I：笛卡尔积与 Zorn 引理}

本附录第 1 节包含任意集合族的笛卡尔积的定义。在正文中，我们主要关注有限多个（或偶尔可数多个）集合的积。我们说明一般定义在这些情形中如何与熟悉的笛卡尔积的“有序 $n$-元组”概念相一致。第 2 节包含对 Zorn 引理及相关主题的讨论。

\section*{1. 笛卡尔积}
\addcontentsline{toc}{section}{1. 笛卡尔积}

若集合 $I$ 的元素被用来标记某个集合族，则称 $I$ 为一个\emph{指标集}。特别地，若 $A$ 与 $I$ 是集合，我们可以通过规定对所有 $i\in I$ 有 $A_i=A$ 来形成集合族 $\{A_i\mid i\in I\}$。因此任何集合都可以作为指标集；我们使用这个术语是为了强调其元素被用作指标。

\noindent\textbf{定义.}
\begin{enumerate}[label=(\arabic*)]
  \item 设 $I$ 是指标集，并设 $\{A_i\mid i\in I\}$ 是一族集合。\emph{选择函数}是任意函数
  \[
    f:I\longrightarrow \bigcup_{i\in I}A_i
  \]
  使得对所有 $i\in I$，都有 $f(i)\in A_i$。

  \item 设 $I$ 是指标集，并且对所有 $i\in I$，令 $A_i$ 为一个集合。集合族 $\{A_i\mid i\in I\}$ 的\emph{笛卡尔积}是从 $I$ 到 $\bigcup_{i\in I}A_i$ 的所有选择函数的集合，记为 $\prod_{i\in I}A_i$（其中若 $I$ 或任何一个集合 $A_i$ 为空，则笛卡尔积为空集）。这个笛卡尔积的元素写作 $\prod_{i\in I}a_i$，这里它表示满足对每个 $i\in I$ 都有 $f(i)=a_i$ 的选择函数 $f$。

  \item 对每个 $j\in I$，集合 $A_j$ 称为笛卡尔积 $\prod_{i\in I}A_i$ 的第 $j$ 个\emph{分量}，而 $a_j$ 是元素 $\prod_{i\in I}a_i$ 的第 $j$ 个\emph{坐标}。

  \item 对 $j\in I$，$\prod_{i\in I}A_i$ 到第 $j$ 个坐标 $A_j$ 上的\emph{投影映射}定义为
  \[
    \prod_{i\in I}a_i\longmapsto a_j.
  \]
\end{enumerate}

笛卡尔积 $\prod_{i\in I}A_i$ 中的每一个选择函数 $f$ 都可以被看作一种从每个集合 $A_i$ 中“选择”一个元素 $f(i)$ 的方式。

若 $I=\{1,2,\ldots,n\}$，其中 $n\in\ZZ^+$，且 $f$ 是从 $I$ 到 $A_1\cup\cdots\cup A_n$ 的选择函数，其中每个 $A_i$ 都非空，则我们可以把 $f$ 对应到唯一的有序 $n$-元组：
\[
  f\longmapsto (f(1),f(2),\ldots,f(n)).
\]
注意，由选择函数的定义，对所有 $i$ 都有 $f(i)\in A_i$，所以以上 $n$-元组在第 $i$ 个位置有一个 $A_i$ 中的元素。

反过来，给定一个 $n$-元组 $(a_1,a_2,\ldots,a_n)$，其中对所有 $i\in I$ 有 $a_i\in A_i$，存在唯一的从 $I$ 到 $\bigcup_{i\in I}A_i$ 的选择函数 $f$ 与之相联系，即
\[
  f(i)=a_i,\qquad \text{对所有 } i\in I.
\]
显然，从 $n$-元组到选择函数的这个映射，是前一段所描述的映射的逆。因此，有序 $n$-元组与 $\prod_{i\in I}A_i$ 的元素之间存在一个双射。此后，当 $I=\{1,2,\ldots,n\}$ 时，我们将写
\[
  \prod_{i=1}^n A_i\quad\text{或}\quad A_1\times A_2\times\cdots\times A_n
\]
来表示笛卡尔积，并把其元素描述为有序 $n$-元组。

若 $I=\ZZ^+$，我们类似地写 $\prod_{i=1}^{\infty}A_i$ 或 $A_1\times A_2\times\cdots$ 来表示这些 $A_i$ 的笛卡尔积。我们将把元素写成有序元组 $(a_1,a_2,\ldots)$，即第 $i$ 项属于 $A_i$ 的无限序列。

注意，当 $I=\{1,2,\ldots,n\}$ 或 $I=\ZZ^+$ 时，我们使用了 $I$ 上的自然顺序来把笛卡尔积的元素排列成 $n$-元组。$I$ 的任何其他顺序（或有限、可数指标集上的任何顺序）都会给出同一个笛卡尔积的元素的另一种不同表示。

\noindent\textbf{例.}
\begin{enumerate}[label=(\arabic*)]
  \item $A\times B=\{(a,b)\mid a\in A,\ b\in B\}$。

  \item $\RR^n=\RR\times\RR\times\cdots\times\RR$（$n$ 个因子）是通常的实数项 $n$-元组的集合，即 Euclid $n$-空间。

  \item 假设 $I=\ZZ^+$，且对所有 $i\in I$，$A_i$ 都是同一个集合 $A$。笛卡尔积 $\prod_{i\in\ZZ^+}A$ 是 $A$ 中元素的所有无限序列 $a_1,a_2,a_3,\ldots$ 的集合。特别地，若 $A=\RR$，则笛卡尔积 $\prod_{i\in\ZZ^+}\RR$ 是所有实数序列的集合。

  \item 假设 $I$ 是任意指标集，且对所有 $i\in I$，$A_i$ 都是同一个集合 $A$。笛卡尔积 $\prod_{i\in I}A$ 正是从 $I$ 到 $A$ 的所有函数的集合，其中函数 $f:I\to A$ 对应于笛卡尔积中的元素 $\prod_{i\in I}f(i)$。这个笛卡尔积常常（特别是在拓扑学书中）记为 $A^I$。注意，对每个固定的 $j\in I$，到第 $j$ 个坐标的投影映射把函数 $f$ 送到 $f(j)$，即是在 $j$ 处取值。

  \item 设 $R$ 是一个环，$x$ 是 $R$ 上的一个不定元。系数在 $R$ 中、关于 $x$ 的多项式环 $R[x]$ 的定义，可以用笛卡尔积来给出，而不必用更直观且熟悉的“形式和”来表述（在第 7 章和第 9 章中，我们以后一种形式引入它们，因为这是我们想象和使用它们的方式）。令 $I$ 为指标集 $\ZZ^+\cup\{0\}$，并令 $R[x]$ 为笛卡尔积 $\prod_{i=0}^{\infty}R$ 的子集，它由元素 $(a_0,a_1,a_2,\ldots)$ 组成，其中只有有限多个 $a_i$ 非零。若 $(a_0,a_1,a_2,\ldots,a_n,0,0,\ldots)$ 是这样的序列，我们用更熟悉的“形式和” $\sum_{i=0}^n a_ix^i$ 来表示它。定义这些序列的加法与乘法，使多项式加法和乘法的通常规则成立。
\end{enumerate}

\medskip
\noindent\textbf{命题 1.} 设 $I$ 是一个非空可数集合，并且对每个 $i\in I$，令 $A_i$ 为一个集合。笛卡尔积的基数等于这些集合基数的乘积：
\[
  \left|\prod_{i\in I}A_i\right|=\prod_{i\in I}\lvert A_i\rvert,
\]
其中若某个 $A_i$ 是无限集，或者若 $I$ 是无限集且无穷多个 $A_i$ 的基数 $\geq 2$，则等式两边都是无穷。特别地，
\[
  \lvert A_1\times A_2\times\cdots\times A_n\rvert
  =\lvert A_1\rvert\times\lvert A_2\rvert\times\cdots\times\lvert A_n\rvert.
\]

\begin{proof}
为了计算选择函数的个数，注意每个 $i\in I$ 都可以被映到 $A_i$ 的 $\lvert A_i\rvert$ 个元素中的任意一个，并且当 $i\neq j$ 时，选择函数在 $i$ 与 $j$ 处的取值可以完全独立地选择。因此，选择函数的个数就是这些 $A_i$ 的基数的乘积，如所断言。

对于有限多个集合的笛卡尔积 $A_1\times A_2\times\cdots\times A_n$，这可以从 $n$-元组表示中容易地看出：$A_1\times A_2\times\cdots\times A_n$ 的元素是 $n$-元组 $(a_1,a_2,\ldots,a_n)$，而每个 $a_i$ 可以被选为 $A_i$ 的 $\lvert A_i\rvert$ 个元素中的任意一个。由于这些选择在 $i\neq j$ 时彼此独立，笛卡尔积中有 $\lvert A_1\rvert\cdot\lvert A_2\rvert\cdots\lvert A_n\rvert$ 个元素。
\end{proof}

\begin{center}
\textbf{练习}
\end{center}

\begin{enumerate}[label=\arabic*.]
  \item 设 $I$ 与 $J$ 是任意两个指标集，并令 $A$ 为任意集合。对任意函数 $\varphi:J\to I$，定义
  \[
    \varphi^*:\prod_{i\in I}A\longrightarrow\prod_{j\in J}A
    \qquad\text{为}\qquad
    \varphi^*(f)=f\circ\varphi
  \]
  对所有选择函数 $f\in\prod_{i\in I}A$ 成立。
  \begin{enumerate}[label=(\alph*)]
    \item 令 $I=\{1,2\}$，令 $J=\{1,2,3\}$，并令 $\varphi:J\to I$ 由 $\varphi(1)=2$、$\varphi(2)=2$、$\varphi(3)=1$ 定义。明确描述在这个 $\varphi^*$ 下，$A\times A\times A$ 中的一个 3-元组如何映到 $A\times A$ 中的一个有序对。
    \item 令 $I=J=\{1,2,\ldots,n\}$，并假设 $\varphi$ 是 $I$ 的一个置换。用 $A\times A\times\cdots\times A$ 中的 $n$-元组描述函数 $\varphi^*$。
  \end{enumerate}
\end{enumerate}

\section*{2. 偏序集与 Zorn 引理}
\addcontentsline{toc}{section}{2. 偏序集与 Zorn 引理}

在正文的少数地方，我们将会把 Zorn 引理作为一种“无限归纳”形式使用，在这些地方需要知道存在某个关于特定性质极大的集合。例如，Zorn 引理被用来证明每个向量空间都有一组基。在这种情形下，向量空间 $V$ 的基是 $V$ 的一个子集，它在由线性无关向量组成的集合中是极大的（这种极大性保证这些向量张成 $V$）。对于有限维空间，这可以通过归纳证明；然而，对于任意维数的空间，需要 Zorn 引理来建立这个结论。通过得到完全一般性成立的结果，理论在若干地方常常会变得稍微整洁一些，尽管正文中的主要结果并不需要使用它。

正文中一个具体例子是：某个有助于简化事情的极大对象由 Zorn 引理构造出来，这就是一个域的代数闭包。域 $F$ 的代数闭包是 $F$ 的一个扩张，它在所有代数扩张的集合中是极大的。这样的域包含（在同构意义下）所有在 $F$ 上代数的元素，因此所有涉及这些代数元素的操作都可以在这一个较大的域中进行。在任何具体情形中，可以通过把所涉及的代数元素添加到基域 $F$ 中来避免使用代数闭包；然而，在复杂证明中，这会变得繁琐（而且常常会遮蔽问题本身）。对于本书例子中出现的具体域，可以避免使用 Zorn 引理来构造代数闭包（例如，任意复数子域或任意有限域的代数闭包的构造都不需要它）。

Zorn 引理的第一次使用出现在第 7.4 节命题 11 的证明中。

为了陈述 Zorn 引理，我们需要一些术语。

\noindent\textbf{定义.} 非空集合 $A$ 上的\emph{偏序}是 $A$ 上的一个关系 $\leq$，满足
\begin{enumerate}[label=(\arabic*)]
  \item 对所有 $x\in A$，$x\leq x$（自反性），
  \item 若 $x\leq y$ 且 $y\leq x$，则对所有 $x,y\in A$ 有 $x=y$（反对称性），
  \item 若 $x\leq y$ 且 $y\leq z$，则对所有 $x,y,z\in A$ 有 $x\leq z$（传递性）。
\end{enumerate}

我们通常说 $A$ 在顺序 $\leq$ 下是一个偏序集，或说 $A$ 被 $\leq$ 偏序。

\noindent\textbf{定义.} 设非空集合 $A$ 被 $\leq$ 偏序。
\begin{enumerate}[label=(\arabic*)]
  \item $A$ 的子集 $B$ 称为一条\emph{链}，如果对所有 $x,y\in B$，或者 $x\leq y$，或者 $y\leq x$。
  \item $A$ 的子集 $B$ 的一个\emph{上界}是元素 $u\in A$，使得对所有 $b\in B$，都有 $b\leq u$。
  \item $A$ 的一个\emph{极大元}是元素 $m\in A$，使得若对任意 $x\in A$ 有 $m\leq x$，则 $m=x$。
\end{enumerate}

在文献中，一条链也称为 tower，或称为全序的、线性序的，或简称有序的子集。

下面的一些例子强调上界与极大元之间的区别。还要注意，如果 $m$ 是 $A$ 的极大元，并不一定有对所有 $x\in A$ 都满足 $x\leq m$（即 $m$ 不一定是最大元）。

\noindent\textbf{例.}
\begin{enumerate}[label=(\arabic*)]
  \item 设 $A$ 是某个集合 $X$ 的幂集（即所有子集的集合），并令 $\leq$ 为集合包含关系：$\subseteq$。注意这只是偏序，因为 $X$ 的某些子集可能不可比较，例如单点集：若 $x\neq y$，则 $\{x\}\nsubseteq\{y\}$ 且 $\{y\}\nsubseteq\{x\}$。在这种情形中，链的一个例子是 $X$ 的一族子集，例如
  \[
    X_1\subseteq X_2\subseteq X_3\subseteq\cdots.
  \]
  $A$ 的任意子集 $B$ 都有一个上界 $b$，即
  \[
    b=\bigcup_{x\in B}x.
  \]
  这个偏序集 $A$ 有一个（唯一的）极大元，即 $X$。

  在许多情形中，集合 $A$ 由某个集合 $X$ 的某些（但不一定是全部）子集组成（即 $A$ 是 $X$ 的幂集的一个子集），且 $A$ 上的顺序仍然是包含关系。上界和极大元的存在性取决于 $A$ 的性质。

  \item 令 $A$ 为 $\ZZ^+$ 的所有真子集的集合，在 $\subseteq$ 下排序。在这种情形中，链不一定有极大元，例如链
  \[
    \{1\}\subseteq\{1,2\}\subseteq\{1,2,3\}\subseteq\cdots
  \]
  没有上界。集合 $A$ 确实有极大元：例如，对任意 $n\in\ZZ^+$，$\ZZ^+-\{n\}$ 是 $A$ 的一个极大元。

  \item 令 $A=\RR$，带通常的 $\leq$ 关系。在这个例子中，$A$ 的每个子集都是一条链（包括 $A$ 本身）。$A$ 的一个子集有上界的概念，与 $\RR$ 的一个子集被某个实数从上方界定的通常概念相同（因此某些集合，如有限长度区间，有上界；而另一些集合，如正实数集合，没有上界）。集合 $A$ 没有极大元。
\end{enumerate}

\medskip
\noindent\textbf{Zorn 引理.} 若 $A$ 是一个非空偏序集，且其中每条链都有一个上界，则 $A$ 有一个极大元。

Zorn 引理独立于通常的（Zermelo--Fraenkel）集合论公理，是一个非平凡的结果；这里所谓独立，是指如果集合论公理是一致的，$^{1}$ 那么这些公理连同 Zorn 引理也是一致的；并且如果集合论公理是一致的，那么这些公理连同 Zorn 引理的否定也是一致的。Zorn 引理中使用“引理”一词是历史性的。

为了完整性（并将 Zorn 引理与其他课程中出现的表述联系起来），我们列出 Zorn 引理的另外两个等价表述。

\medskip
\noindent\textbf{选择公理.} 任意非空集合族的非空集合的笛卡尔积非空。换言之，若 $I$ 是任意非空（指标）集合，并且对所有 $i\in I$，$A_i$ 是非空集合，则存在一个从 $I$ 到 $\bigcup_{i\in I}A_i$ 的选择函数。

\noindent\textbf{定义.} 设 $A$ 是非空集合。$A$ 上的一个\emph{良序}是 $A$ 上的一个全序，使得 $A$ 的每个非空子集都有最小元，即对每个非空 $B\subseteq A$，存在某个 $s\in B$，使得对所有 $b\in B$ 有 $s\leq b$。

\medskip
\noindent\textbf{良序原理.} 每个非空集合 $A$ 都有一个良序。

\medskip
\noindent\textbf{定理 2.} 在通常的（Zermelo--Fraenkel）集合论公理下，下列命题等价：
\begin{enumerate}[label=(\arabic*)]
  \item Zorn 引理；
  \item 选择公理；
  \item 良序原理。
\end{enumerate}

\begin{proof}
这由初等集合论推出。关于这些等价性以及其他若干等价性，可参见 Hewitt 与 Stromberg 的 \emph{Real and Abstract Analysis}，Springer-Verlag，1965，第 3 节。
\end{proof}

\noindent$^{1}$ 参见 P. J. Cohen 在 \emph{Proc. Nat. Acad. Sci.}，50(1963) 与 51(1964) 中的论文。\\
$^{2}$ 这并不知道是真的！

\begin{center}
\textbf{练习}
\end{center}

\begin{enumerate}[label=\arabic*.]
  \item 令 $A$ 为 $\RR$ 的所有有限子集的集合，按包含关系排序。讨论上界、极小元与极大元的存在性（或不存在性）（其中极小元的定义与极大元类似）。解释为什么这不是一个良序。

  \item 令 $A$ 为 $\RR$ 的所有无限子集的集合，按包含关系排序。讨论上界、极小元与极大元的存在性（或不存在性）。解释为什么这不是一个良序。

  \item 说明给定集合上的下列偏序不是良序：
  \begin{enumerate}[label=(\alph*)]
    \item $\RR$，带通常的关系 $\leq$。
    \item $\RR^+$，带通常的关系 $\leq$。
    \item $\RR^+\cup\{0\}$，带通常的关系 $\leq$。
    \item $\ZZ$，带通常的关系 $\leq$。
  \end{enumerate}

  \item 说明 $\ZZ^+$ 在通常的关系 $\leq$ 下是良序的。
\end{enumerate}

\chapter*{附录 II\\范畴论}
\addcontentsline{toc}{chapter}{附录 II：范畴论}

范畴论为讨论大类数学对象的性质提供语言和数学基础，例如“所有集合”的类或“所有群”的类，同时避开 Russell 悖论等问题。在这个框架中，可以探讨各类对象研究中使用的概念和方法之间的共性：同态、同构等；还可以引入研究类与类之间关系的工具：函子、范畴等价等。于是可以给出“自然”变换和“自然”同构的精确定义，无论是在一个给定类内部，还是在两个类之间。（在正文中，我们把“自然”描述为“与坐标无关”。）一个类内部自然同构的典型例子，是第 11.3 节中任意有限维向量空间与其二重对偶之间的同构。事实上，S. Eilenberg 与 S. MacLane 在 1945 年引入范畴与函子的主要动机之一，就是要给诸如此类情形中的“自然”概念以精确含义。范畴论也在形式化新概念方面起了基础性作用，例如概形（参见第 15.5 节），它们是当代研究的主要领域（例如代数几何）的基本概念。这类工作的先驱包括 A. Grothendieck、K. Morita 及其他人。

我们对范畴论的处理应被看作是对一些基本语言的介绍。由于我们没有讨论集合论的 Zermelo--Fraenkel 公理或类的 Godel--Bernays 公理，我们不涉及范畴论的基础。不过，为了与集合论公理保持一致，我们隐含地假设存在一个宇宙集合 $U$，它包含人们在“普通”数学中会遇到的所有集合、群、环等（因此“所有集合”的范畴隐含地意为“$U$ 中所有集合”的范畴）。关于进一步学习，可参阅集合论、逻辑或范畴论的书籍，例如 S. MacLane 的 \emph{Categories for the Working Mathematician}，Springer-Verlag，1971。

我们组织本附录的方式是：在可能的地方，每个新概念的例子都按照这些术语和结构在正文中出现的顺序来使用。例如，第一个函子的例子涉及集合与群，第二个例子使用环，等等。这样，本附录可以在学习早期阅读；随着读者逐渐熟悉更多种类的数学结构，重新阅读这些例子也会得到更深入的理解。

\section*{1. 范畴与函子}
\addcontentsline{toc}{section}{1. 范畴与函子}

我们从本附录的基本概念开始。

\noindent\textbf{定义.} 一个\emph{范畴} $\mathcal C$ 由一类对象以及对象之间的态射集合组成。对每个有序对象对 $A,B$，有一个从 $A$ 到 $B$ 的态射集合 $\operatorname{Hom}_{\mathcal C}(A,B)$；对每个有序对象三元组 $A,B,C$，有一个态射的复合律，即映射
\[
  \operatorname{Hom}_{\mathcal C}(A,B)\times\operatorname{Hom}_{\mathcal C}(B,C)
  \longrightarrow \operatorname{Hom}_{\mathcal C}(A,C),
\]
其中 $(f,g)\mapsto gf$，并且 $gf$ 称为 $g$ 与 $f$ 的复合。对象与态射满足以下公理：对对象 $A,B,C,D$，
\begin{enumerate}[label=(\roman*)]
  \item 若 $A\neq B$ 或 $C\neq D$，则 $\operatorname{Hom}_{\mathcal C}(A,B)$ 与 $\operatorname{Hom}_{\mathcal C}(C,D)$ 是不交集合；
  \item 态射复合满足结合律，即对每个 $f\in\operatorname{Hom}_{\mathcal C}(A,B)$、$g\in\operatorname{Hom}_{\mathcal C}(B,C)$ 与 $h\in\operatorname{Hom}_{\mathcal C}(C,D)$，有
  \[
    h(gf)=(hg)f;
  \]
  \item 每个对象都有一个恒等态射，即对每个对象 $A$，存在态射 $1_A\in\operatorname{Hom}_{\mathcal C}(A,A)$，使得对每个 $f\in\operatorname{Hom}_{\mathcal C}(A,B)$ 有 $f1_A=f$，并且对每个 $g\in\operatorname{Hom}_{\mathcal C}(B,A)$ 有 $1_Ag=g$。
\end{enumerate}

态射也称为箭头。恒等态射对每个对象是唯一的，这一点作为练习可证（论证与群的恒等元唯一性的论证相同）。当范畴由上下文清楚给出时，我们将把 $\operatorname{Hom}_{\mathcal C}(A,B)$ 写作 $\operatorname{Hom}(A,B)$。

我们在全文中使用的术语对所有范畴都是共同的：从 $A$ 到 $B$ 的态射记为 $f:A\to B$ 或 $A\xrightarrow{f}B$。对象 $A$ 是 $f$ 的\emph{定义域}，$B$ 是 $f$ 的\emph{陪域}。从 $A$ 到 $A$ 的态射称为 $A$ 的\emph{自同态}。态射 $f:A\to B$ 称为\emph{同构}，如果存在态射 $g:B\to A$，使得 $gf=1_A$ 且 $fg=1_B$。

有子范畴 $\mathcal C$ 属于 $\mathcal D$ 的自然概念，也就是说，$\mathcal C$ 的每个对象也是 $\mathcal D$ 中的对象，并且对 $\mathcal C$ 中的对象 $A,B$，有包含关系
\[
  \operatorname{Hom}_{\mathcal C}(A,B)\subseteq \operatorname{Hom}_{\mathcal D}(A,B).
\]

\noindent\textbf{例.}

在下面每个例子中，我们都把验证范畴公理的细节留作练习。
\begin{enumerate}[label=(\arabic*)]
  \item $\mathbf{Set}$ 是所有集合的范畴。对任意两个集合 $A$ 与 $B$，$\operatorname{Hom}(A,B)$ 是从 $A$ 到 $B$ 的所有函数的集合。态射复合就是熟悉的函数复合：$gf=g\circ f$。$\operatorname{Hom}(A,A)$ 中的恒等元是映射 $1_A(a)=a$，对所有 $a\in A$ 成立。这个范畴含有所有有限集合的范畴作为子范畴。

  \item $\mathbf{Grp}$ 是所有群的范畴，其中态射是群同态。注意，群同态的复合仍然是群同态。$\mathbf{Grp}$ 的一个子范畴是 $\mathbf{Ab}$，即所有 Abel 群的范畴。类似地，$\mathbf{Ring}$ 是所有带 $1$ 的非零环的范畴，其中态射是把 $1$ 送到 $1$ 的环同态。所有带 $1$ 的交换环的范畴 $\mathbf{CRing}$ 是 $\mathbf{Ring}$ 的子范畴。

  \item 对固定的环 $R$，范畴 $R\text{-}\mathbf{mod}$ 由所有左 $R$-模组成，态射为 $R$-模同态。

  \item $\mathbf{Top}$ 是这样的范畴：对象为拓扑空间，态射为拓扑空间之间的连续映射（参见第 15.2 节）。注意，从一个空间到自身的恒等（集合）映射在每种拓扑下都是连续的，所以 $\operatorname{Hom}(A,A)$ 总有一个恒等元。

  \item 令 $\varnothing$ 为\emph{空范畴}，它没有对象也没有态射。令 $1$ 表示只有一个对象 $A$ 和一个态射的范畴：$\operatorname{Hom}(A,A)=\{1_A\}$。令 $2$ 为有两个对象 $A_1$ 与 $A_2$，且只有一个非恒等态射的范畴：$\operatorname{Hom}(A_1,A_2)=\{f\}$ 且 $\operatorname{Hom}(A_2,A_1)=\varnothing$。注意，对象 $A_1$ 与 $A_2$ 以及态射 $f$ 是“原始”的，意思是 $A_1$ 与 $A_2$ 不被定义为集合，而 $f$ 只是从 $A_1$ 到 $A_2$ 的一个箭头（字面意义上的箭头）；它没有被定义为某个集合的元素上的集合映射。可以以这种方式继续，定义 $N$ 为有 $N$ 个对象 $A_1,A_2,\ldots,A_N$ 的范畴，唯一的非恒等态射是对每个 $j>i$ 从 $A_i$ 到 $A_j$ 的唯一箭头（因此箭头复合被唯一决定）。

  \item 若 $G$ 是一个群，按如下方式形成范畴 $G$。唯一的对象是 $G$，并且 $\operatorname{Hom}(G,G)=G$；两个函数 $f$ 与 $g$ 的复合是群 $G$ 中的乘积 $gf$。注意，$\operatorname{Hom}(G,G)$ 有一个恒等态射：群 $G$ 的恒等元。
\end{enumerate}

\noindent\textbf{定义.} 设 $\mathcal C$ 与 $\mathcal D$ 是范畴。
\begin{enumerate}[label=(\arabic*)]
  \item 我们说 $\mathcal F$ 是从 $\mathcal C$ 到 $\mathcal D$ 的\emph{协变函子}，如果
  \begin{enumerate}[label=(\alph*)]
    \item 对 $\mathcal C$ 中每个对象 $A$，$\mathcal F A$ 是 $\mathcal D$ 中的一个对象；并且
    \item 对每个 $f\in\operatorname{Hom}_{\mathcal C}(A,B)$，有 $\mathcal F(f)\in\operatorname{Hom}_{\mathcal D}(\mathcal F A,\mathcal F B)$，
  \end{enumerate}
  并满足以下公理：
  \begin{enumerate}[label=(\roman*)]
    \item 若 $gf$ 是 $\mathcal C$ 中态射的复合，则在 $\mathcal D$ 中有 $\mathcal F(gf)=\mathcal F(g)\mathcal F(f)$；
    \item $\mathcal F(1_A)=1_{\mathcal F A}$。
  \end{enumerate}

  \item 我们说 $\mathcal F$ 是从 $\mathcal C$ 到 $\mathcal D$ 的\emph{反变函子}，如果 (1) 中的条件成立，但性质 (b) 与公理 (i) 分别替换为：
  \begin{enumerate}
    \item[(b')] 对每个 $f\in\operatorname{Hom}_{\mathcal C}(A,B)$，$\mathcal F(f)\in\operatorname{Hom}_{\mathcal D}(\mathcal F B,\mathcal F A)$；
  \end{enumerate}
  \begin{enumerate}
    \item[(i')] 若 $gf$ 是 $\mathcal C$ 中态射的复合，则在 $\mathcal D$ 中有 $\mathcal F(gf)=\mathcal F(f)\mathcal F(g)$，
  \end{enumerate}
  即反变函子反转箭头。
\end{enumerate}

\noindent\textbf{例.}

在每个例子中，验证函子公理的工作留作练习。更多函子的例子见本节末尾的练习。
\begin{enumerate}[label=(\arabic*)]
  \item 恒等函子 $I_{\mathcal C}$ 把任意范畴 $\mathcal C$ 映到自身，它把对象和态射都送到自身。更一般地，如果 $\mathcal C$ 是 $\mathcal D$ 的子范畴，则包含函子把 $\mathcal C$ 映入 $\mathcal D$，它把对象和态射都送到自身。

  \item 设 $\mathcal F$ 是从 $\mathbf{Grp}$ 到 $\mathbf{Set}$ 的函子，它把任意群 $G$ 映到同一个集合 $G$，并把任意群同态 $\varphi$ 映到同一个集合映射 $\varphi$。这个函子称为\emph{遗忘函子}，因为它“移除”或“遗忘”群以及它们之间同态的结构。类似地，也有从 $\mathbf{Ab}$、$R\text{-}\mathbf{mod}$、$\mathbf{Top}$ 等范畴到 $\mathbf{Set}$ 的遗忘函子。

  \item \emph{Abel 化函子}把 $\mathbf{Grp}$ 映到 $\mathbf{Ab}$，它把每个群 $G$ 送到 Abel 群 $G^{\mathrm{ab}}=G/G'$，其中 $G'$ 是 $G$ 的换位子群（参见第 5.4 节）。每个群同态 $\varphi:G\to H$ 被映到商群上的诱导同态：
  \[
    \overline\varphi:G^{\mathrm{ab}}\longrightarrow H^{\mathrm{ab}}
    \qquad\text{由}\qquad
    \overline\varphi(xG')=\varphi(x)H'
  \]
  给出。换位子群的定义保证 $\overline\varphi$ 良定义，并且函子公理成立。

  \item 设 $R$ 是一个环，$D$ 是一个左 $R$-模。对每个左 $R$-模 $N$，集合 $\operatorname{Hom}_R(D,N)$ 是一个 Abel 群；若 $R$ 是交换环，则它也是一个 $R$-模（参见第 10.2 节命题 2）。若 $\varphi:N_1\to N_2$ 是一个 $R$-模同态，则对每个 $f\in\operatorname{Hom}_R(D,N_1)$，有 $\varphi\circ f\in\operatorname{Hom}_R(D,N_2)$。因此
  \[
    \varphi':\operatorname{Hom}_R(D,N_1)\longrightarrow\operatorname{Hom}_R(D,N_2),
    \qquad \varphi'(f)=\varphi\circ f.
  \]
  这说明映射
  \[
    \operatorname{Hom}(D,\_):N\longrightarrow\operatorname{Hom}_R(D,N),
    \qquad
    \operatorname{Hom}(D,\_):\varphi\longrightarrow\varphi'
  \]
  是从 $R\text{-}\mathbf{Mod}$ 到 $\mathbf{Grp}$ 的协变函子。若 $R$ 是交换的，它把 $R\text{-}\mathbf{Mod}$ 映到自身。

  \item 在前一个例子的记号下，我们观察到，若 $\varphi:N_1\to N_2$，则对每个 $g\in\operatorname{Hom}_R(N_2,D)$，有 $g\circ\varphi\in\operatorname{Hom}_R(N_1,D)$。因此
  \[
    \varphi':\operatorname{Hom}_R(N_2,D)\longrightarrow\operatorname{Hom}_R(N_1,D),
    \qquad \varphi'(g)=g\circ\varphi.
  \]
  在这个情形中，映射
  \[
    \operatorname{Hom}(\_,D):N\longrightarrow\operatorname{Hom}_R(N,D),
    \qquad
    \operatorname{Hom}(\_,D):\varphi\longrightarrow\varphi'
  \]
  定义了一个反变函子。

  \item 当 $D$ 是一个右 $R$-模时，映射 $D\otimes_R\_:N\mapsto D\otimes_R N$ 定义了从 $R\text{-}\mathbf{Mod}$ 到 $\mathbf{Ab}$ 的协变函子（若 $R$ 是交换的，则到 $R\text{-}\mathbf{Mod}$）。这里态射 $\varphi:N_1\to N_2$ 映到态射 $1\otimes\varphi$。

  同样，当 $D$ 是一个左 $R$-模时，$\_\otimes_R D:N\mapsto N\otimes_R D$ 定义了从右 $R$-模范畴到 $\mathbf{Ab}$ 的协变函子（若 $R$ 是交换的，则到 $R\text{-}\mathbf{Mod}$），其中态射 $\varphi$ 映到态射 $\varphi\otimes1$。

  \item 设 $K$ 是一个域，并令 $K\text{-}\mathbf{fdVec}$ 为所有有限维 $K$-向量空间的范畴，其中该范畴中的态射是 $K$-线性变换。我们定义从 $K\text{-}\mathbf{fdVec}$ 到自身的二重对偶函子 $\mathcal D^2$。回忆第 11.3 节，$V$ 的对偶空间 $V^*$ 定义为 $V^*=\operatorname{Hom}_K(V,K)$；$V$ 的二重对偶是 $V^{**}=\operatorname{Hom}_K(V^*,K)$。于是 $\mathcal D^2$ 在对象上定义为把向量空间 $V$ 映到其二重对偶 $V^{**}$。若 $\varphi:V\to W$ 是有限维空间的线性变换，则
  \[
    \mathcal D^2(\varphi):V^{**}\longrightarrow W^{**}
    \qquad\text{由}\qquad
    \mathcal D^2(\varphi)(E_v)=E_{\varphi(v)}
  \]
  定义，其中对每个 $v\in V$，$E_v$ 表示“在 $v$ 处求值”。由第 11.3 节定理 19，$E_v\in V^{**}$，并且 $V^{**}$ 的每个元素都可唯一地写成某个 $v\in V$ 的 $E_v$。由于 $\varphi(v)\in W$，有 $E_{\varphi(v)}\in W^{**}$，所以 $\mathcal D^2(\varphi)$ 是良定义的。
\end{enumerate}

从 $\mathcal C$ 到 $\mathcal D$ 的函子 $\mathcal F$ 称为\emph{忠实的}（或称为\emph{满的}），如果对 $\mathcal C$ 中每一对对象 $A$ 与 $B$，映射
\[
  \mathcal F:\operatorname{Hom}(A,B)\longrightarrow\operatorname{Hom}(\mathcal F A,\mathcal F B)
\]
是单射（或满射）。因此，例如，遗忘函子是忠实的但不是满的。

\begin{center}
\textbf{练习}
\end{center}

\begin{enumerate}[label=\arabic*.]
  \item 设 $N$ 是一个群，并令 $\mathbf{Nor}\text{-}N$ 为所有含有 $N$ 作为正规子群的群的集合。对象 $A$ 与 $B$ 之间的态射是任意把 $N$ 映入 $N$ 的群同态。
  \begin{enumerate}[label=(\alph*)]
    \item 证明 $\mathbf{Nor}\text{-}N$ 是一个范畴。
    \item 说明投影同态 $G\mapsto G/N$ 如何用来定义从 $\mathbf{Nor}\text{-}N$ 到 $\mathbf{Grp}$ 的一个函子。
  \end{enumerate}

  \item 设 $H$ 是一个群。按如下方式在对象和态射上定义从 $\mathbf{Grp}$ 到自身的映射 $\mathcal H_{\times}$：
  \[
    \mathcal H_{\times}:G\longmapsto H\times G,
  \]
  并且若 $\varphi:G_1\to G_2$，则
  \[
    \mathcal H_{\times}(\varphi):H\times G_1\longrightarrow H\times G_2
    \qquad\text{由}\qquad
    (h,g)\longmapsto (h,\varphi(g))
  \]
  定义。证明 $\mathcal H_{\times}$ 是一个函子。

  \item 说明把环映到其单位群的映射（即 $R\mapsto R^{\times}$）定义了一个从 $\mathbf{Ring}$ 到 $\mathbf{Grp}$ 的函子。用明确的例子说明这个函子既非忠实也非满。

  \item 说明对每个 $n\geq1$，映射 $\mathcal{GL}_n:R\mapsto GL_n(R)$ 定义了一个从 $\mathbf{CRing}$ 到 $\mathbf{Grp}$ 的函子。［在态射上定义 $\mathcal{GL}_n$ 时，把每个环同态作用到矩阵的每个项上。］

  \item 补全说明例 7 中所描述的二重对偶映射满足函子公理的细节。
\end{enumerate}

\section*{2. 自然变换与泛性质}
\addcontentsline{toc}{section}{2. 自然变换与泛性质}

如本附录导言中所提到的，范畴论产生的动机之一，是要给“自然”同构这一概念一个精确定义。我们现在给出这个定义，并看到正文中提到的一些自然映射如何是范畴论概念的实例。我们同样给出“泛箭头”和“泛对象”的范畴论定义，并从这个角度看正文中出现的一些泛性质。

\noindent\textbf{定义.} 设 $\mathcal C$ 与 $\mathcal D$ 是范畴，并且 $\mathcal F,\mathcal G$ 是从 $\mathcal C$ 到 $\mathcal D$ 的协变函子。从 $\mathcal F$ 到 $\mathcal G$ 的一个\emph{自然变换}或\emph{函子态射}是一个映射 $\eta$，它给 $\mathcal C$ 中每个对象 $A$ 指派一个 $\operatorname{Hom}_{\mathcal D}(\mathcal F A,\mathcal G A)$ 中的态射 $\eta_A$，且具有如下性质：对 $\mathcal C$ 中每一对对象 $A$ 与 $B$，以及每个 $f\in\operatorname{Hom}_{\mathcal C}(A,B)$，都有
\[
  \mathcal G(f)\eta_A=\eta_B\mathcal F(f),
\]
即下图交换：
\[
\begin{array}{ccc}
\mathcal F A & \xrightarrow{\eta_A} & \mathcal G A\\
\mathcal F(f)\downarrow && \downarrow\mathcal G(f)\\
\mathcal F B & \xrightarrow{\eta_B} & \mathcal G B.
\end{array}
\]
若每个 $\eta_A$ 都是同构，则称 $\eta$ 为函子的\emph{自然同构}。

考虑特例 $\mathcal C=\mathcal D$，并且 $\mathcal C$ 是 $\mathbf{Set}$ 的一个子范畴，而 $\mathcal F$ 是恒等函子。如果每当 $\mathcal G$ 把对象 $A$ 映到对象 $\mathcal G A$ 时，都存在一个从 $A$ 到 $\mathcal G A$ 的态射 $\eta_A$，并且每当存在从 $A$ 到 $B$ 的态射 $f$ 时，态射 $\mathcal G(f)$ 作为从 $\mathcal G A$ 到 $\mathcal G B$ 的映射与 $f$ 相容，则存在一个从恒等函子到 $\mathcal G$ 的自然变换 $\eta$。事实上，$\mathcal G(f)$ 在 $\mathcal G A$ 的子集 $\eta_A(A)$ 到 $\mathcal G B$ 的子集 $\eta_B(B)$ 上由 $f$ 唯一决定。若 $\eta$ 是一个自然同构，则 $\mathcal G$ 在每个态射上的值完全由 $\eta$ 决定，即
\[
  \mathcal G(f)=\eta_Bf\eta_A^{-1}.
\]
在这种情形下，函子 $\mathcal G$ 完全由 $\eta$ 指定。我们将看到，前一节中的一些函子例子正是这样产生的。

\noindent\textbf{例.}
\begin{enumerate}[label=(\arabic*)]
  \item 对任意范畴 $\mathcal C$ 与 $\mathcal D$，以及从 $\mathcal C$ 到 $\mathcal D$ 的任意函子 $\mathcal F$，恒等映射是从 $\mathcal F$ 到自身的自然同构：对 $\mathcal C$ 中每个对象 $A$，$\eta_A=1_{\mathcal F A}$。

  \item 设 $R$ 是一个环，$\mathcal F$ 是从 $R\text{-}\mathbf{Mod}$ 到自身的任意函子。零映射是从 $\mathcal F$ 到自身的自然变换：对每个 $R$-模 $A$，$\eta_A=0_A$，其中 $0_A$ 是从 $A$ 到自身的零映射。这不是自然同构。

  \item 设 $\mathcal F$ 是从 $\mathbf{Grp}$ 到自身的恒等函子，设 $\mathcal G$ 是 Abel 化函子（例 3），这里把它看作从 $\mathbf{Grp}$ 到自身的映射。对每个群 $G$，令 $\eta_G:G\to G/G'$ 为到商群的通常投影映射。于是 $\eta$ 是关于这两个函子的一个自然变换（但不是同构）。（我们称映射 $\eta_G$ 为自然投影映射。）

  \item 令 $\mathcal G=\mathcal D^2$ 为从域 $K$ 上有限维向量空间范畴到自身的二重对偶函子（例 7）。则存在一个从恒等函子到 $\mathcal G$ 的自然同构 $\eta$，由
  \[
    \eta_V:V\longrightarrow V^{**}\qquad\text{为}\qquad \eta_V(v)=E_v
  \]
  给出，其中对每个 $v\in V$，$E_v$ 是“在 $v$ 处求值”。

  \item 令 $\mathcal{GL}_n$ 为从 $\mathbf{CRing}$ 到 $\mathbf{Grp}$ 的函子，定义如下。每个对象（交换环）$R$ 被 $\mathcal{GL}_n$ 映到由 $R$ 中元素组成的 $n\times n$ 可逆矩阵群 $GL_n(R)$。对每个环同态 $f:R\to S$，令 $\mathcal{GL}_n(f)$ 为把 $f$ 作用到矩阵每个项上的矩阵映射。由于 $f$ 把 $1$ 送到 $1$，可知 $\mathcal{GL}_n(f)$ 把可逆矩阵送到可逆矩阵（参见第 1 节练习 4）。令 $\mathcal G$ 为从 $\mathbf{CRing}$ 到 $\mathbf{Grp}$ 的函子，它把每个环 $R$ 映到其单位群 $R^\times$，并把每个环同态 $f$ 映到其在单位群上的限制（仍记为 $f$）。行列式是从 $\mathcal{GL}_n$ 到 $\mathcal G$ 的自然变换，因为行列式对所有环都由同一个多项式定义，因此下图交换：
  \[
  \begin{array}{ccc}
  GL_n(R) & \xrightarrow{\det} & R^\times\\
  \mathcal{GL}_n(f)\downarrow && \downarrow f\\
  GL_n(S) & \xrightarrow{\det} & S^\times.
  \end{array}
  \]
\end{enumerate}

设 $\mathcal C,\mathcal D$ 与 $\mathcal E$ 是范畴，令 $\mathcal F$ 为从 $\mathcal C$ 到 $\mathcal D$ 的函子，令 $\mathcal G$ 为从 $\mathcal D$ 到 $\mathcal E$ 的函子。函子复合 $\mathcal G\mathcal F$ 从 $\mathcal C$ 到 $\mathcal E$ 有显然的定义。当 $\mathcal E=\mathcal C$ 时，复合 $\mathcal G\mathcal F$ 把 $\mathcal C$ 映到自身，而 $\mathcal F\mathcal G$ 把 $\mathcal D$ 映到自身。我们说 $\mathcal C$ 与 $\mathcal D$ 是\emph{同构的}，如果存在某些 $\mathcal F$ 与 $\mathcal G$，使得 $\mathcal G\mathcal F$ 是恒等函子 $I_{\mathcal C}$，且 $\mathcal F\mathcal G=I_{\mathcal D}$。由第 10.1 节中的讨论，范畴 $\ZZ\text{-}\mathbf{Mod}$ 与 $\mathbf{Ab}$ 是同构的。第 10.1 节中的观察还推出，初等 Abel $p$-群的范畴与 $\FF_p$ 上向量空间的范畴是同构的。在实际中，我们倾向于把这样的同构范畴等同起来。下面关于范畴同构的推广，给出了两个范畴“相似”的更宽泛且更有用的概念。

\noindent\textbf{定义.} 范畴 $\mathcal C$ 与 $\mathcal D$ 称为\emph{等价的}，如果存在从 $\mathcal C$ 到 $\mathcal D$ 的函子 $\mathcal F$ 和从 $\mathcal D$ 到 $\mathcal C$ 的函子 $\mathcal G$，使得函子 $\mathcal G\mathcal F$ 自然同构于 $I_{\mathcal C}$（$\mathcal C$ 的恒等函子），并且 $\mathcal F\mathcal G$ 自然同构于恒等函子 $I_{\mathcal D}$。

范畴等价是自反的、对称的和传递的，这是一个练习。第 15.5 节中的仿射 $k$-代数例子给出了一个范畴等价（在该例中，需要修改自然变换定义中的箭头方向，以适应反变函子）。另一个例子（需要一些证明）是：对带 $1$ 的交换环 $R$，左模范畴 $R\text{-}\mathbf{Mod}$ 与 $M_{n\times n}(R)\text{-}\mathbf{Mod}$ 是等价的。

最后，我们引入泛箭头与泛对象的概念。

\noindent\textbf{定义.}
\begin{enumerate}[label=(\arabic*)]
  \item 设 $\mathcal C$ 与 $\mathcal D$ 是范畴，$\mathcal F$ 是从 $\mathcal C$ 到 $\mathcal D$ 的函子，$X$ 是 $\mathcal D$ 中的对象。从 $X$ 到 $\mathcal F$ 的一个\emph{泛箭头}是一对 $(U(X),t)$，其中 $U(X)$ 是 $\mathcal C$ 中的一个对象，而 $t:X\to\mathcal F U(X)$ 是 $\mathcal D$ 中的一个态射，满足如下性质：对 $\mathcal C$ 中任意对象 $A$，若 $\varphi$ 是 $\mathcal D$ 中从 $X$ 到 $\mathcal F A$ 的任意态射，则存在唯一的 $\mathcal C$ 中的态射 $\Phi:U(X)\to A$，使得 $\mathcal F(\Phi)t=\varphi$，即下图交换：
  \[
  \begin{array}{ccc}
     & X & \\
     t\swarrow && \searrow\varphi\\
     \mathcal F U(X) & \xrightarrow{\mathcal F(\Phi)} & \mathcal F A.
  \end{array}
  \]

  \item 设 $\mathcal C$ 是一个范畴，$\mathcal F$ 是从 $\mathcal C$ 到所有集合的范畴 $\mathbf{Set}$ 的函子。函子 $\mathcal F$ 的一个\emph{泛元素}是一对 $(U,t)$，其中 $U$ 是 $\mathcal C$ 中的对象，而 $t$ 是集合 $\mathcal F U$ 的一个元素，满足如下性质：对 $\mathcal C$ 中任意对象 $A$ 和集合 $\mathcal F A$ 中任意元素 $g$，存在唯一的 $\mathcal C$ 中的态射 $\varphi:U\to A$，使得 $\mathcal F(\varphi)(t)=g$。
\end{enumerate}

\noindent\textbf{例.}
\begin{enumerate}[label=(\arabic*)]
  \item （泛箭头：自由对象）设 $R$ 是带 $1$ 的环。我们把如下陈述翻译成泛箭头的语言：若 $U(X)$ 是集合 $X$ 上的自由 $R$-模，则任意从 $X$ 到 $R$-模 $A$ 的集合映射都唯一地通过 $R$-线性延拓为一个从 $U(X)$ 到 $A$ 的 $R$-模同态（参见第 10.3 节定理 6）：令 $\mathcal F$ 为从 $R\text{-}\mathbf{Mod}$ 到 $\mathbf{Set}$ 的遗忘函子，于是 $\mathcal F$ 把 $R$-模 $A$ 映到集合 $A$，即作为集合 $A=\mathcal F A$。令 $X$ 为任意集合（即 $\mathbf{Set}$ 中的对象），令 $U(X)$ 为以 $X$ 为基的自由 $R$-模，并令 $t:X\to\mathcal F U(X)$ 为把每个 $b\in X$ 送到 $U(X)$ 中基元素 $b$ 的集合映射。那么自由 $R$-模的泛性质正是 $(U(X),t)$ 是从 $X$ 到遗忘函子 $\mathcal F$ 的泛箭头这一结果。

  类似地，自由群、向量空间（即域上的自由模）、多项式代数（即自由 $R$-代数）等，都是泛箭头的实例。

  \item （泛箭头：分式域）令 $\mathcal F$ 为从域的范畴到整环的范畴的遗忘函子，其中两个范畴中的态射都是单的环同态。对任意整环 $X$，令 $U(X)$ 为其分式域，并令 $t$ 为 $X$ 到 $U(X)$ 中的包含映射。那么 $(U(X),t)$ 是从 $X$ 到函子 $\mathcal F$ 的泛箭头（参见第 7.5 节定理 15(2)）。

  \item （泛对象：张量积）这个例子涉及第 10.4 节中两个模的张量积的构造。令 $\mathcal C=R\text{-}\mathbf{Mod}$ 为交换环 $R$ 上 $R$-模的范畴，并令 $M$ 与 $N$ 是 $R$-模。对每个 $R$-模 $A$，令 $\operatorname{Bilin}(M,N;A)$ 表示从 $M\times N$ 到 $A$ 的所有 $R$-双线性函数的集合。在对象上定义从 $R\text{-}\mathbf{Mod}$ 到 $\mathbf{Set}$ 的函子
  \[
    \mathcal F:A\longrightarrow \operatorname{Bilin}(M,N;A),
  \]
  并且若 $\varphi:A\to B$ 是一个 $R$-模同态，则
  \[
    \mathcal F(\varphi)(h)=\varphi\circ h
    \qquad\text{对每个 }h\in\operatorname{Bilin}(M,N;A).
  \]
  令 $U=M\otimes_R N$，并令 $t$ 为双线性函数
  \[
    t:M\times N\longrightarrow M\otimes_R M
    \qquad\text{由}\qquad
    t(m,n)=m\otimes n
  \]
  给出，于是 $t$ 是集合 $\operatorname{Bilin}(M,N;M\otimes_R N)=\mathcal F U$ 的一个元素。那么 $(M\otimes_R N,t)$ 是 $\mathcal F$ 的一个泛元素，因为对任意 $R$-模 $A$，以及任意双线性映射 $g:M\times N\to A$（即 $\mathcal F A$ 的任意元素），存在唯一的 $R$-模同态 $\varphi:M\otimes_R N\to A$，使得
  \[
    g=\varphi\circ t=\mathcal F(\varphi)(t).
  \]
\end{enumerate}

\begin{center}
\textbf{练习}
\end{center}

\begin{enumerate}[label=\arabic*.]
  \item 令 $\mathbf{Nor}\text{-}N$ 为第 1 节练习 1 中描述的范畴，并令 $\mathcal F$ 为从 $\mathbf{Nor}\text{-}N$ 到 $\mathbf{Grp}$ 的包含函子。描述一个从 $\mathbf{Nor}\text{-}N$ 到 $\mathbf{Grp}$ 的函子 $\mathcal G$，使得由 $\eta_G:G\to G/N$ 定义的变换 $\eta$ 是从 $\mathcal F$ 到 $\mathcal G$ 的自然变换。

  \item 设 $H$ 与 $K$ 是群，且 $\mathcal H_{\times}$ 与 $\mathcal K_{\times}$ 是第 1 节练习 2 中描述的从 $\mathbf{Grp}$ 到自身的函子。令 $\varphi:H\to K$ 为群同态。
  \begin{enumerate}[label=(\alph*)]
    \item 说明映射 $\eta_A:H\times A\to K\times A$，由 $\eta_A(h,a)=(\varphi(h),a)$ 给出，确定了一个从 $\mathcal H_{\times}$ 到 $\mathcal K_{\times}$ 的自然变换 $\eta$。
    \item 说明变换 $\eta$ 是自然同构，当且仅当 $\varphi$ 是群同构。
  \end{enumerate}

  \item 把第 5.4 节命题 7(5) 中所描述的换位子商群的泛性质，表述为某个函子 $\mathcal F$ 的泛箭头。
\end{enumerate}

\end{document}
